% 本文件由 LaTeX 排版 Agent 根据有效 translation.json 自动生成。
% author=Augustine of Hippo; work_id=1701; title=论若望福音的讲道集（讲座）

\chapter{\WorkTitle{论\BibleBook{若望}{福音}的讲道集（讲座）}}

% structure: p0001 source=h1 target=omitted reason=page-title-already-rendered-by-parent

\Emph{讲道一} \BibleRef{若 1:1-5}\newline{} \Emph{讲道二} \BibleRef{若 1:6-14}\newline{} \Emph{讲道三} \BibleRef{若 1:15-18}\newline{} \Emph{讲道四} \BibleRef{若 1:19-33}\newline{} \Emph{讲道五} \BibleRef{若 1:33}\newline{} \Emph{讲道六} \BibleRef{若 1:32-33}\newline{} \Emph{讲道七} \BibleRef{若 1:34-51}\newline{} \Emph{讲道八} \BibleRef{若 2:1-4}\newline{} \Emph{讲道九} \BibleRef{若 2:1-2}\newline{} \Emph{讲道十} \BibleRef{若 2:12-21}\newline{} \Emph{讲道十一} \BibleRef{若 2:23-3:5}\newline{} \Emph{讲道十二} \BibleRef{若 3:6-21}\newline{} \Emph{讲道十三} \BibleRef{若 3:22-29}\newline{} \Emph{讲道十四} \BibleRef{若 3:29-36}\newline{} \Emph{讲道十五} \BibleRef{若 4:1-42}\newline{} \Emph{讲道十六} \BibleRef{若 4:43-54}\newline{} \Emph{讲道十七} \BibleRef{若 5:1-18}\newline{} \Emph{讲道十八} \BibleRef{若 5:19}\newline{} \Emph{讲道十九} \BibleRef{若 5:19-30}\newline{} \Emph{讲道二十} \BibleRef{若 5:19}\newline{} \Emph{讲道二十一} \BibleRef{若 5:20-23}\newline{} \Emph{讲道二十二} \BibleRef{若 5:24-30}\newline{} \Emph{讲道二十三} \BibleRef{若 5:19-40}\newline{} \Emph{讲道二十四} \BibleRef{若 6:1-14}\newline{} \Emph{讲道二十五} \BibleRef{若 6:15-44}\newline{} \Emph{讲道二十六} \BibleRef{若 6:41-59}\newline{} \Emph{讲道二十七} \BibleRef{若 6:60-72}\newline{} \Emph{讲道二十八} \BibleRef{若 7:1-13}\newline{} \Emph{讲道二十九} \BibleRef{若 7:14-18}\newline{} \Emph{讲道三十} \BibleRef{若 7:19-24}\newline{} \Emph{讲道三十一} \BibleRef{若 7:25-36}\newline{} \Emph{讲道三十二} \BibleRef{若 7:37-39}\newline{} \Emph{讲道三十三} \BibleRef{若 7:40-8:11}\newline{} \Emph{讲道三十四} \BibleRef{若 8:12}\newline{} \Emph{讲道三十五} \BibleRef{若 8:13-14}\newline{} \Emph{讲道三十六} \BibleRef{若 8:15-18}\newline{} \Emph{讲道三十七} \BibleRef{若 8:19-20}\newline{} \Emph{讲道三十八} \BibleRef{若 8:21-25}\newline{} \Emph{讲道三十九} \BibleRef{若 8:26-27}\newline{} \Emph{讲道四十} \BibleRef{若 8:28-32}\newline{} \Emph{讲道四十一} \BibleRef{若 8:31-36}\newline{} \Emph{讲道四十二} \BibleRef{若 8:37-47}\newline{} \Emph{讲道四十三} \BibleRef{若 8:48-59}\newline{} \Emph{讲道四十四} \BibleRef{若 9}\newline{} \Emph{讲道四十五} \BibleRef{若 10:1-10}\newline{} \Emph{讲道四十六} \BibleRef{若 10:11-13}\newline{} \Emph{讲道四十七} \BibleRef{若 10:14-21}\newline{} \Emph{讲道四十八} \BibleRef{若 10:22-42}\newline{} \Emph{讲道四十九} \BibleRef{若 11:1-54}\newline{} \Emph{讲道五十} \BibleRef{若 11:55-12}\newline{} \Emph{讲道五十一} \BibleRef{若 12:12-26}\newline{} \Emph{讲道五十二} \BibleRef{若 12:27-36}\newline{} \Emph{讲道五十三} \BibleRef{若 12:37-43}\newline{} \Emph{讲道五十四} \BibleRef{若 12:44-50}\newline{} \Emph{讲道五十五} \BibleRef{若 13:1-5}\newline{} \Emph{讲道五十六} \BibleRef{若 13:6-10}\newline{} \Emph{讲道五十七} \BibleRef{若 13:6-10}, \BibleRef{歌 5:2-3}\newline{} \Emph{讲道五十八} \BibleRef{若 13:10-15}\newline{} \Emph{讲道五十九} \BibleRef{若 13:16-20}\newline{} \Emph{讲道六十} \BibleRef{若 13:21}\newline{} \Emph{讲道六十一} \BibleRef{若 13:21-26}\newline{} \Emph{讲道六十二} \BibleRef{若 13:26-31}\newline{} \Emph{讲道六十三} \BibleRef{若 13:31-32}\newline{} \Emph{讲道六十四} \BibleRef{若 13:33}\newline{} \Emph{讲道六十五} \BibleRef{若 13:34-35}\newline{} \Emph{讲道六十六} \BibleRef{若 13:36-38}\newline{} \Emph{讲道六十七} \BibleRef{若 14:1-3}\newline{} \Emph{讲道六十八} \BibleRef{若 14:1-3}\newline{} \Emph{讲道六十九} \BibleRef{若 14:4-6}\newline{} \Emph{讲道七十} \BibleRef{若 14:7-10}\newline{} \Emph{讲道七十一} \BibleRef{若 14:10-14}\newline{} \Emph{讲道七十二} \BibleRef{若 14:10-14}\newline{} \Emph{讲道七十三} \BibleRef{若 14:10-14}\newline{} \Emph{讲道七十四} \BibleRef{若 14:15-17}\newline{} \Emph{讲道七十五} \BibleRef{若 14:18-21}\newline{} \Emph{讲道七十六} \BibleRef{若 14:22-24}\newline{} \Emph{讲道七十七} \BibleRef{若 14:25-27}\newline{} \Emph{讲道七十八} \BibleRef{若 14:27-28}\newline{} \Emph{讲道七十九} \BibleRef{若 14:29-31}\newline{} \Emph{讲道八十} \BibleRef{若 15:1-3}\newline{} \Emph{讲道八十一} \BibleRef{若 15:4-7}\newline{} \Emph{讲道八十二} \BibleRef{若 15:8-10}\newline{} \Emph{讲道八十三} \BibleRef{若 15:2-12}\newline{} \Emph{讲道八十四} \BibleRef{若 15:13}\newline{} \Emph{讲道八十五} \BibleRef{若 15:14-15}\newline{} \Emph{讲道八十六} \BibleRef{若 15:15-16}\newline{} \Emph{讲道八十七} \BibleRef{若 15:17-19}\newline{} \Emph{讲道八十八} \BibleRef{若 15:20-21}\newline{} \Emph{讲道八十九} \BibleRef{若 15:22-23}\newline{} \Emph{讲道九十} \BibleRef{若 15:23}\newline{} \Emph{讲道九十一} \BibleRef{若 15:24-25}\newline{} \Emph{讲道九十二} \BibleRef{若 15:26-27}\newline{} \Emph{讲道九十三} \BibleRef{若 16:1-4}\newline{} \Emph{讲道九十四} \BibleRef{若 16:44-47}\newline{} \Emph{讲道九十五} \BibleRef{若 16:8-11}\newline{} \Emph{讲道九十六} \BibleRef{若 16:12-13}\newline{} \Emph{讲道九十七} \BibleRef{若 16:12-13}\newline{} \Emph{讲道九十八} \BibleRef{若 16:12-33}\newline{} \Emph{讲道九十九} \BibleRef{若 16:13}\newline{} \Emph{讲道一百} \BibleRef{若 16:13-15}\newline{} \Emph{讲道一百零一} \BibleRef{若 16:16-23}\newline{} \Emph{讲道一百零二} \BibleRef{若 16:23-28}\newline{} \Emph{讲道一百零三} \BibleRef{若 16:29-33}\newline{} \Emph{讲道一百零四} \BibleRef{若 17:1}\newline{} \Emph{讲道一百零五} \BibleRef{若 17:1-5}\newline{} \Emph{讲道一百零六} \BibleRef{若 17:6-8}\newline{} \Emph{讲道一百零七} \BibleRef{若 17:9-13}\newline{} \Emph{讲道一百零八} \BibleRef{若 17:14-19}\newline{} \Emph{讲道一百零九} \BibleRef{若 17:20}\newline{} \Emph{讲道一百一十} \BibleRef{若 17:21-23}\newline{} \Emph{讲道一百一十一} \BibleRef{若 17:24-26}\newline{} \Emph{讲道一百一十二} \BibleRef{若 18:1-12}\newline{} \Emph{讲道一百一十三} \BibleRef{若 18:13-27}\newline{} \Emph{讲道一百一十四} \BibleRef{若 18:28-32}\newline{} \Emph{讲道一百一十五} \BibleRef{若 18:33-40}\newline{} \Emph{讲道一百一十六} \BibleRef{若 19:1-16}\newline{} \Emph{讲道一百一十七} \BibleRef{若 19:17-22}\newline{} \Emph{讲道一百一十八} \BibleRef{若 19:23-24}\newline{} \Emph{讲道一百一十九} \BibleRef{若 19:24-30}\newline{} \Emph{讲道一百二十} \BibleRef{若 19:31-20:9}\newline{} \Emph{讲道一百二十一} \BibleRef{若 20:10-29}\newline{} \Emph{讲道一百二十二} \BibleRef{若 20:30-21:11}\newline{} \Emph{讲道一百二十三} \BibleRef{若 21:12-19}\newline{} \Emph{讲道一百二十四} \BibleRef{若 21:19-25}\par

\SourceRecord{Translated by John Gibb.；Nicene and Post-Nicene Fathers, First Series；Vol. 7.；Edited by Philip Schaff.；Buffalo, NY: Christian Literature Publishing Co.,；1888.；<http://www.newadvent.org/fathers/1701.htm>.}

% structure: p0001 source=h1 target=section reason=page-title

\section{\WorkTitle{讲道集1（若望福音1:1-5）}}

1. 当我留意我们刚才从宗徒读经中所读到的：\Quote{属血气的人，不能领受天主圣神的事}\BibleRef{格前 2:14}，并考虑到在当前的集会中，我亲爱的诸位，你们中间必然有许多属血气的人，他们只按肉体认识，还不能提升自己到属神的理解，我深感困难，不知在主所赐的恩佑下，我如何能够表达，或按我渺小的能力，解释从福音所读到的：\Quote{在起初已有圣言，圣言与天主同在，圣言就是天主}；因为属血气的人不能领悟这事。那么，弟兄们，我们为此就闭口不言吗？那么，为什么还要诵读它，如果我们对此缄默不语？或者，为什么还要听它，如果它不被解释？又为什么解释它，如果它不被理解？因此，另一方面，既然我不怀疑你们中间有些人不只能在被解释时领受它，甚至在解释之前就能理解它，我将不辜负那些能够领受的人，以免因害怕我的话浪费在那些不能领受的人耳中。最后，天主的慈悲将与我们同在，以至于或许足够供应众人，各人按自己能领受的领受，而说话的人按自己能说的说。因为按事情本身来说，谁能呢？我敢说，我的弟兄们，也许连若望本人也没有按事情本身来说，而只是按他所能说的来说；因为是人谈论天主，固然受天主默感，但仍然是个人。因为他受了默感，所以他说了一些；如果他没有受默感，他就什么也不会说；但因是一个人受了默感，他并没有说全部，而只是说了人所能说的。\par

2. 因为这位若望，亲爱的弟兄们，是那些山中的一座，关于这些山曾记载说：\Quote{愿群山为祢的百姓收纳和平，丘陵收纳正义。}山是高贵的灵魂，丘陵是渺小的灵魂。但山收纳和平，是为了让丘陵能够收纳正义。什么是丘陵所收纳的正义？信德，因为\Quote{义人因信德而生活}。然而，较小的灵魂本不会收纳信德，除非那更大的、被称为山的灵魂，被智慧本身所光照，使他们能够将小灵魂所能领受的传给小灵魂；而丘陵因信德生活，因为山收纳了和平。正是由这些山本身，曾对教会说：\Quote{愿你们平安}；而山本身在向教会宣告和平时，并没有使自己与那赐和平者分离\BibleRef{若 20:19}，好使他们能真诚地、而非虚伪地宣告和平。\par

3. 因为另有其他的山，它们造成沉船，若有人将船驶向它们，船便撞得粉碎。因为在危难中的人看见陆地时，很容易冒险去接近它；但有时陆地出现在山上，而岩石隐伏在山下；当有人驶向山时，他撞在岩石上，在那里找不到安息，而是船毁。因此曾有一些山，在人间显得伟大，他们制造了异端和分裂，分裂了天主的教会；但那些分裂了天主子民的人，并非那被说\Quote{愿群山为祢的百姓收纳和平}的山。因为那些分裂了合一的人，是如何收纳了和平呢？\par

4. 但那些收纳和平以向百姓宣告和平的人，使智慧本身成为默观的对象，正如人心所能把握那\Quote{眼所未见，耳所未闻，人心所未曾想到过的}\BibleRef{格前 2:9}。如果它没有上升到人的心中，如何上升到若望的心中呢？难道若望不是一个人吗？或者，也许它并没有上升到若望的心中，而是若望的心上升到它那里？因为那上升到人心的东西，是从下面到人；但人心所上升到的那里，是上面，从人出发。同样，弟兄们，能否说，如果它上升到若望的心中（如果可以说的话），那么它上升到他的心中，只是在他不再是人的程度上。说\Quote{不再是人}是什么意思？在他开始成为天使的程度上。因为所有圣人都是天使，因为他们是天主的使者。因此，对于属血肉和属血气的人，他们不能领受属于天主的事，宗徒是怎样说的呢？\Quote{因为你们说：我是属保禄的，我是属阿颇罗的，你们岂不是人吗？}\BibleRef{格前 3:4} 他所责备的这些人，因为他们是人，他愿意将他们变成什么？你们愿知道他要将他们变成什么吗？在圣咏中听：\Quote{我曾说：你们是神，都是至高者的儿子。}因此，天主召唤我们，要我们不再是人。但那时将更好，我们不再是人，如果我们首先承认我们是人这一事实，即为了从谦卑上升到那高处；免得当我们自以为是什么，而其实什么也不是时，我们不仅没有领受我们所不是的，反而连我们所是的也失去了。\par

5. 是故，弟兄们，在这些山中的，也有若望，他曾说：\BibleQuote{在起初已有圣言，圣言与天主同在，圣言就是天主。}这座山已领受了和平；他正默观圣言的天主性。这座山是怎样的呢？何其崇高？他超越了大地上所有的高峰，超越了天空所有的平原，超越了星辰所有的高度，超越了所有天使的唱诗班和军团。因为除非他超越了一切受造之物，否则他无法到达那藉以创造万物的那一位。你无法想象他超越了何等境界，除非你看见他到达了何处。你询问天地吗？它们被造。你询问天上和地上的事物吗？当然更是被造。你询问神灵、天使、总领天使、上座者、宰制者、异能者、率领者吗？这些也是被造的。因为圣咏在列举这一切时，这样结束：\BibleQuote{他一发言，万有造成；他一出命，各物生成。}如果\BibleQuote{他一发言，万有造成，}那么就是藉着圣言造成的；但如果藉着圣言造成，那么若望的心怎能达到他所说的\BibleQuote{在起初已有圣言，圣言与天主同在，圣言就是天主}，除非他超越了所有藉圣言造成的事物？这是何等的一座山！何其神圣！何其崇高，在那些为天主的百姓领受了和平、使小山领受公义的山中！\par

6. 那么，弟兄们，请想一想，若望是否正是那些山中的一座，我们刚才所歌唱的关于它们：\BibleQuote{我举目向山，我的救助从何而来。}所以，我的弟兄们，如果你们愿意明白，就举目向这座山，也就是说，提升你们自己到这位福音作者那里，提升到他的意义里。但是，因为这些山虽领受了和平，但将自己的希望寄托在人身上的人不能享有和平，所以不要如此举目向山，以致认为你的希望应寄托在人身上；而要说：\BibleQuote{我举目向山，我的救助从何而来，}这样你立刻加上：\BibleQuote{我的救助来自上主，他创造了天地。}因此，让我们举目向山，我们的救助从何而来；然而，我们的希望并不应寄托在山上，因为山领受了他们要为我们施予的；所以，应当将我们的希望寄托在那使山也领受的地方。当我们举目向圣经时，因为圣经是通过人施予的，我们就是举目向山，我们的救助从何而来；但尽管如此，因为写圣经的人是，他们并非自己发光，而是\BibleQuote{祂是真光，照亮每一个来到世上的人。}\BibleRef{若 1:9} 洗者若翰也是一座山，他说：\BibleQuote{我不是基督，}\BibleRef{若 1:30} 以免任何人将希望寄托在那座山上，而从光照该山的那里坠落。他也承认说：\BibleQuote{从祂的满盈中，我们都领受了。}\BibleRef{若 1:16} 所以你们应该说：\BibleQuote{我举目向山，我的救助从何而来，} 但不要将临到你们的救助归给山；而要接着说：\BibleQuote{我的救助来自上主，他创造了天地。}\par

7. 因此，弟兄们，愿我劝告的结果是，你们明白在举心向圣经时（当福音高声宣报：\BibleQuote{在起初已有圣言，圣言与天主同在，圣言就是天主}，以及所读的其余部分），你们就是举目向山。因为若非这些山说了这些话，你们根本无法想出该如何思考它们。所以，从山来了你们的救助，使你们甚至听到了这些事；但你们还不能理解所听到的。你们当呼求来自上主、创造天地的救助；因为山只能这样说话，不能凭自己光照，因为它们自己也是藉着听闻而被光照。从那里，说这些话的若望领受了——他曾躺在主的怀里，并从主的怀里汲取了他所能给我们饮用的。但他赐给我们饮用的是言语。那么，你们应该从赐给你们饮用的那一位所汲取的源头去领受理解；这样你们可以举目向山，从那里你们的援助将到来，以便从那里你们能领受那杯，就是那话语，赐给你们饮用；然而，既然你们的救助来自上主，创造天地的，你们可以从他充满胸怀的那源头充满你们的胸怀；因此你们说：\BibleQuote{我的救助来自上主，他创造了天地：}那么，让能充满的来充满吧。弟兄们，这就是我所说的：让每个人以合适的方式举心，并领受所讲的话。但也许你们会说，我比天主更临在于你们。愿你们绝无此念！祂更临在于你们；因为我显现于你们的眼前，祂统御你们的良心。那么，将你们的耳朵给我，将你们的心给祂，好使你们两者都得充满。看，你们的眼睛和那些肉体的感官，你们举向我们；然而不是向我们，因为我们不是那些山，而是向福音本身，向那位福音作者；但你们的心，当向主以被充满。此外，让每个人如此举心，以致看见他举起了什么，以及向哪里。我说\Quote{他举起了什么，以及向哪里}是什么意思？让他省察他举起的是何种心，因为他是向主举起它，免得它被肉欲的重担所累，还未举起就坠落。但每个人难道不看见他带着肉体的重担吗？他当以节制努力洁净那可以向上主举起的东西。因为\BibleQuote{心里洁净的人是有福的，因为他们要看见天主。}\BibleRef{玛 5:8}\par

8. 但我们且看，这些话\Quote{在起初已有圣言，圣言与天主同在，圣言就是天主。}（\BibleRef{若 1:1-2}）有何益处。我们说话时也发出言语。难道这样的言语是与天主同在的吗？我们发出的那些言语不是响过就消失了吗？那么，天主的圣言也响过就终止了吗？若是如此，万物怎是藉祂造成的？没有祂，什么也未曾造成？如果祂响过就消失了，祂所创造的万物又怎受祂管理？那么，那既是发出又不消失的言语是怎样的呢？我亲爱的，请听，这是大事。因着日常谈话，言语在此处对我们变得卑贱，因为通过它们的响过与消失，它们被轻视，似乎只是言语而已。但在人自身内有一个存留于内的言语；因为声音从口中发出。有一种真正属神的方式所说的言语，那是你从声音中所理解的，而非声音本身。注意，我说\Quote{天主}这个词时，我发出的言语何其短——四个字母两个音节！难道天主就只是这四个字母两个音节吗？或者，那被意指的对象，与这言语的微不足道同样贵重吗？当你听见\Quote{天主}时，你心中发生了什么事？当我说\Quote{天主}时，我心中发生了什么事？在我们的思想中有一个伟大而完美的实体，超越一切可变的受造物，无论肉体或灵魂。如果我向你们说：\Quote{天主是可变的还是不变的？}你们会立刻回答：\Quote{我断不相信或想象天主是可变的：天主是不变的。}你们的灵魂，虽然微小，虽然或许还是属血肉的，却只能回答我说天主是不变的：但每一个受造物都是可变的；那么，你们是如何藉着精神的一瞥，进入那超越受造物的领域，从而自信地回答说\Quote{天主是不变的}呢？那么，当你们思想一个实体，活的、永恒的、全能的、无限的、无所不在、处处完整、无处受限制时，你们心中是什么呢？当你们思想这些性质时，这就是你们心中关于天主的言语。但这难道就是那由四个字母两个音节组成的声音吗？因此，凡发出并消失的一切，都是声音、字母、音节。祂的言语是声音，发出就消失；但那被声音所意指的，在说话者思想时存在于他内，在听者理解时存在于他内，当声音消失时，它却存留。\par

9. 请留意那言语。你们可以在心中有一个言语，好像是心灵中所生的构思，以至心灵生出构思；而构思，可以说是心灵的产物，内心的孩子。因为先是你的心产生构思，要建造某座建筑，在地上立起某件大事；构思已形成，工程尚未完成：你看见你要做什么；但别人不会赞叹，直到你建造起那堆东西，使那建筑成形并完成；于是人们观看那令人赞叹的建筑，赞叹建筑师的构思；他们因所见而惊讶，并为未见而喜悦：谁能看见构思呢？如果因为某座大建筑，一个人的构思就受到赞美，你们愿意看到天主的构思，即主耶稣基督，也就是天主的圣言是怎样的吗？请看这世界的建筑。观看那由圣言所造成的，然后你们就会明白圣言的本质。看这世界的两个物体：天和地。谁能用言语述说天的美丽？谁能用言语述说地的丰饶？谁能恰如其分地颂扬季节的变化？谁能恰如其分地颂扬种子的力量？你们看到哪些我没有提及，以免我列举太多，或许比你们自己心中能想起的还少。那么，从这建筑，来估量那造成它的圣言的本质：而且不只是它；因为所有这些事物都是可见的，因为它们与肉体感官有关。连天使也是由那圣言造成的；连总领天使也是由那圣言造成的，诸异能、座者、宰制者、率领者，都是藉那圣言造成的；万物是藉祂造成的。因此，请估量这是怎样的圣言。\par

10. 或许现在有人回答我说：\Quote{谁能领会这圣言？}那么，当你们听见\Quote{圣言}时，不要想象某种微不足道的东西，也不要以爲那是你们每天听见的言语——\Quote{他说了这样的话}、\Quote{他说了那样的话}、\Quote{你对我说了那样的话}；因为由于不断重复，\Quote{言语}这个词，可说是变得毫无价值。而当你们听见\Quote{在起初已有圣言}时，免得你们想象某种无价值的东西，如同你们惯常想到听人说话时那样，请听你当如何想：\Quote{圣言就是天主。}\par

11. 现在或许有不相信的亚略派（\OriginalTerm{亚略派}{Arian}）出来说\Quote{天主的圣言是受造的}。天主的圣言怎么可能是受造的呢？既然天主藉圣言造成万物。如果天主的圣言本身也是受造的，那么它是由另一个什么圣言造成的呢？但若你说有圣言的圣言，我说，那造成它的就是天主的独生子。但若你不说有圣言的圣言，就承认那造成万物的不是受造的。因为那造成万物的不能由自己造成。那么，请相信福音作者。因为他本可以说：\Quote{起初天主造了圣言}，正如梅瑟说：\Quote{起初天主造了天和地}（\BibleRef{创 1:1}），并这样列举一切：\Quote{天主说：\InnerQuote{让它造成}，就造成了。}\EditorialNote{创世纪第一章}如果\Quote{说}，谁说的？天主。造成了什么？某个受造物。在天主说话与受造物造成之间，那造成它的若不是圣言，是什么呢？因为天主说：\Quote{让它造成，就造成了。}这圣言是不变的；虽然可变的事物由它造成，圣言本身却是不变的。\par

12. 那么，不要相信那藉以造成万物的本身是被造的，免得你们不被那使万物更新的圣言所更新。因为你们已经藉圣言而被造，但你们还必须藉圣言被更新。然而，如果你们对圣言的信仰是错误的，你们就不能藉圣言被更新。虽然圣言的创造已经临到你们，使你们由祂所造，但你们因自己而被败坏；如果你们因自己而被败坏，就让那创造你们的使你们更新；如果你们因自己而变得更坏，就让那创造你们的使你们再造。但是，如果你们对圣言持有错误的看法，祂怎能藉圣言再造你们呢？福音作者说：\Quote{在起初已有圣言}；而你们说：\Quote{在起初圣言被造成}。他说：\Quote{万物是藉着祂造成的}；而你们说圣言本身是被造的。福音作者本可以说：\Quote{在起初圣言被造成}；但他说的是什么呢？\Quote{在起初已有圣言}。如果祂是，祂就不是被造的；这样万物才能藉祂而被造，没有祂，什么也没有造成。那么，如果\Quote{在起初已有圣言，圣言与天主同在，圣言就是天主}；如果你们无法想象这是什么，就等待你们长大。那是干粮；先接受奶水，使你们得到滋养，才能接受干粮。\par

13. 弟兄们，要留意接下来的话：\Quote{万物是藉着祂造成的，没有祂，什么也没有造成}，这样你们就不会以为\Quote{\OriginalTerm{虚无}{nothing}}是某种东西。因为许多人错误地理解\Quote{没有祂，什么也没有造成}，往往幻想\Quote{虚无}是某种东西。的确，罪不是由祂造成的；显然罪是虚无，人犯罪时也变成虚无。偶像也不是由圣言造成的；——它确实具有某种人的形状，但人本身是由圣言造成的；——因为偶像中人的形状不是由圣言造成的，正如经上所载：\Quote{我们知道偶像算不得什么。}\BibleRef{格前 8:4}因此，这些事物不是由圣言造成的；但是，凡按本性造成的，凡属于受造物的，一切固定在天上的，从高处照耀的，在天之下飞行的，在宇宙中运动的，所有的受造物：弟兄们，我要说得更明白，使你们能理解我；我要说，从天使到一只虫子。在受造物中，有什么比天使更尊贵？有什么比虫子更低微？那造天使的也造了虫子；但天使适合天上，虫子适合地上。那创造的也安排了。如果祂把虫子放在天上，你们可能会指责；如果祂愿意天使从腐朽的肉中产生，你们也可能会指责：然而天主几乎就是这样做的，祂是不应被指责的。因为所有从肉身生的人，不是虫吗？天主从这些虫子中造了天使。因为如果主亲自说：\Quote{但我是一条虫，不是一个人，}谁还会迟疑，不说\WorkTitle{约伯传}中所写的：\Quote{何况人是什么不过是腐臭，人子是什么不过是条虫子？}\BibleRef{约 25:6}首先他说：\Quote{人是腐臭，}然后：\Quote{人子是条虫子：}因为虫子从腐臭而生，因此\Quote{人是腐臭，}\Quote{人子是条虫子。}请看，那\Quote{在起初已有圣言，圣言与天主同在，圣言就是天主！}的，为了你们竟甘愿成为什么？祂为什么为了你们成为这样？为使你们这些不能咀嚼的可以吮吸。因此，弟兄们，完全按这意义理解\Quote{万物是藉着祂造成的，没有祂，什么也没有造成。}因为每一个受造物，无论大小，都由祂造成：天上和地下的，有形和无形，都由祂造成。因为任何形状、结构、部分的配合、任何有重量、数目、度量的实体，都只藉着那圣言，藉着那创造者的圣言而存在，关于这圣言，经上说：\Quote{祢以度量、数目和重量安排了万物。}\BibleRef{智 11:21}\par

14. 因此，不要让任何人欺骗你们，如果你们偶尔受苍蝇困扰。因为有些人曾被魔鬼戏弄，并被苍蝇所捕获。正如捕鸟者习惯在陷阱中放苍蝇欺骗饥饿的鸟，这些人也被魔鬼用苍蝇欺骗了。某一个人正受苍蝇困扰；一个\OriginalTerm{摩尼教徒}{Manichæan}发现他处于烦恼中，当他说他不能忍受苍蝇，并且极其憎恨它们时，摩尼教徒立刻说：\Quote{谁造了它们？}由于他正受困扰，憎恨它们，他不敢说：\Quote{天主造了它们，}虽然他是一名天主教徒。那另一个立刻加上：\Quote{如果天主没有造它们，谁造了它们？}\Quote{诚然，}天主教徒回答说：\Quote{我相信魔鬼造了它们。}那另一个立刻说：\Quote{如果魔鬼造了苍蝇，如我所见你承认，因为你明白这事很好，那么谁造了蜜蜂，它比苍蝇稍大一点？}天主教徒不敢说天主造了蜜蜂而没有造苍蝇，因为情况非常相似。从蜜蜂他把他引向蝗虫；从蝗虫到蜥蜴；从蜥蜴到鸟；从鸟到羊；从羊到牛；从牛到大象，最后到人；并且说服那人，人不是天主造的。于是这可怜的人，因被苍蝇困扰，自己变成了一只苍蝇，成了魔鬼的财产。事实上，他们说，\OriginalTerm{贝耳则步}{Beelzebub}的意思是\Quote{苍蝇之王}；关于这些，经上记载：\Quote{死苍蝇使香膏失去芬芳。}\BibleRef{训 10:1}\par

15. 那么，弟兄们，我为什么说这些呢？要关闭你们内心的耳朵，不受仇敌诡计的欺骗。要明白天主创造了万物，并按其秩序安排了它们。那么，为什么我们从一个天主所造的受造物那里遭受许多恶事呢？因为我们得罪了天主吗？天使们会遭受这些吗？也许我们，在天使那样的生命中，也不会怕这类事。为了惩罚你，该控告你的罪，而不是审判者。因为，因着我们的骄傲，天主指定了那微小而可鄙的受造物来折磨我们；如此，既然人变得骄傲，并自夸反对天主，且虽是凡人，却压迫凡人，虽是人也未承认自己的同类——既然他抬举自己，就藉着蚊虫被降卑。你为什么因人的骄傲而膨胀？有人责备了你，你就充满了怒气。赶走蚊虫，好使你安睡：要明白你是谁。因为，弟兄们，为使你们知道，正是为了驯服我们的骄傲，这些事物才被造来烦扰我们，天主本可用熊、狮子、蛇来降卑法郎骄傲的子民；祂却使苍蝇和青蛙降在他们身上，\EditorialNote{出谷纪第八章}，好让最卑贱的受造物制伏他们的骄傲。\par

16. \BibleQuote{万事万物，}然后，弟兄们，\BibleQuote{万物是藉着祂造的，凡受造的没有一样不是藉着祂造的。} 但万物是怎样藉着祂造的呢？\BibleQuote{那受造的，在祂内有生命。} 也可以这样读：\BibleQuote{那在祂内受造的，是生命；} 如果我们这样读，那么一切都是生命。因为有什么不是在祂内受造的呢？因为祂是天主的智慧，圣咏上说：\BibleQuote{在智慧中，祢造了万物。} 那么，如果基督是天主的智慧，圣咏又说：\BibleQuote{在智慧中，祢造了万物：} 既然万物是藉着祂造的，那么万物也是在祂内造的。那么，如果万物都是在祂内造的，亲爱的弟兄们，而那在祂内受造的，是生命，那么大地是生命，树木也是生命。我们的确说树木是生命，但指的是十字架的树木，我们从那里获得了生命。那么，石头是生命。这样理解这段经文是不恰当的，如同那最卑劣的摩尼教派再次悄然潜入我们中间，说石头有生命，墙壁有灵魂，绳有灵魂，羊毛和衣服也有灵魂。因为他们惯于在癫狂中如此谈论；当他们被反驳和驳倒时，他们就会以某种方式引用圣经说：\Quote{为什么说\BibleQuote{那在祂内受造的，是生命}？} 因为如果万物都是在祂内造的，那么万物都是生命。不要被他们迷惑；要这样读：\BibleQuote{那受造的；}在这里稍作停顿，然后继续：\BibleQuote{在祂内有生命。} 这是什么意思呢？大地是受造的，但那受造的大地本身并不是生命；但在智慧本身之中，存在一种精神性的理念，大地就是根据这理念受造的：这就是生命。\par

17. 尽我所能，我要向你们解释我的意思，亲爱的。一个木匠制作一个箱子。首先他在构思中有一个箱子；因为若不在构思中有它，他怎能通过工艺产生它呢？但理念中的箱子并不是眼睛所见的那个箱子。它在构思中无形地存在，将在作品中成为可见。看，它在作品中造出来了；它是否在构思中就不存在了呢？一个在作品中造出，另一个仍存在于构思中；因为那个箱子可能会腐烂，而另一个会依照存在于构思中的那个被造出来。那么，要注意构思中的箱子和事实上的箱子。实际的箱子不是生命，构思中的箱子是生命；因为工匠的灵魂，在所有这些被造出来之前，它存在于其中，是活的。所以，亲爱的弟兄们，因为天主的智慧，万物都是藉着它造的，在造物之前，它按构思包含了一切，因此那些藉着这个构思本身所造的东西，并不立刻就是生命，但无论什么造物，在祂内就是生命。你看见大地，有一个在构思中的大地；你看见天空，有一个在构思中的天空；你看见太阳和月亮，这些也在构思中存在：但在外表上它们是形体，在构思中它们是生命。要明白，若你们能以某种方式理解，因为一件大事被说出来了。如果我不伟大，我说话的人，或藉着我说话的人，但仍然是从一个伟大的权威而来。因为这些事不是由我这个小人物说的；我提及的祂不是小人物。各人按他所能的、按他所能的程度领受；凡不能领受任何一部分的，让他滋养他的心，好使他能成为有能力领受的。他如何滋养呢？让他用奶滋养，好使他能到达坚固的食物。让他不要离开藉着肉身降生的基督，直到他到达惟独由圣父所生的基督，即天主圣言与天主同在，万物是藉着祂造的；因为那就是生命，在祂内是人的光。\par

18. 因为接着是：\BibleQuote{生命是人的光；} 并且正是从这个生命，人被光照。牲畜不被光照，因为牲畜没有理性心智以看见智慧。但人是按照天主的肖像造的，并且有理性心智，藉此他能领悟智慧。那么，那生命，万物是藉着它造的，本身就是光；但不是每一种动物的光，而是人的光。因此稍后祂说：\BibleQuote{那才是真光，照亮每一个来到世界的人。} 若翰洗者被那光照亮；若望宗徒自己也被同样的光照亮。那说\BibleQuote{我不是基督；但有一位在我以后来的，我连解祂鞋带也不配}的人，被那光充满。\BibleRef{若 1:26} 那说\BibleQuote{在起初已有圣言，圣言与天主同在，圣言就是天主}的人，也被那光照亮。因此那生命是人的光。\par

但也许你们中有些人迟钝的心还不能接受那光明，因为被自己的罪所重压，以致不能看见。不要因此就认为光明在任何方面是缺席的，因为他们不能看见它；他们自己因罪恶而成了黑暗。\BibleQuote{光明在黑暗中照耀，黑暗未能理解它。}因此，弟兄们，如同一个盲人置身于阳光之下，太阳临在于他，但他却缺席于太阳。同样，每一个愚昧的人，每一个不义的人，每一个不虔敬的人，都是心中盲目。智慧临在；但它是临在于一个盲人，却从他的眼睛缺席；不是因为它对他缺席，而是因为他缺席于它。那么他该做什么呢？让他变得纯洁，使他能看见天主。就好像一个人因眼睛被尘土、眼屎或烟气弄得污秽刺痛而不能看见，医生会对他说：\Quote{从你眼中清除一切不好的东西，使你能看见你眼睛的光。}尘土、眼屎和烟气就是罪恶与不义：那么除去这一切，你就能看见那临在的智慧；因为天主就是那智慧，并且经上记载：\BibleQuote{心里纯洁的人是有福的，因为他们要看见天主。}\BibleRef{玛 5:8}\par

\SourceRecord{Translated by John Gibb.；Nicene and Post-Nicene Fathers, First Series；Vol. 7.；Edited by Philip Schaff.；Buffalo, NY: Christian Literature Publishing Co.,；1888.；<http://www.newadvent.org/fathers/1701001.htm>.}

% structure: p0001 source=h1 target=section reason=page-title

\section{讲道集2（\BibleBook{若望}{福音}1:6-14）}

弟兄们，我们应当尽可能讲解圣经经文，尤其是圣福音，不遗漏任何部分，以便我们自己能按各自的能力汲取滋养，并从那滋养我们的源泉中为你们服务。上个主日，我们记得，我们讲解了第一部分，即：\BibleQuote{在起初已有圣言，圣言与天主同在，圣言就是天主。圣言在起初就与天主同在。万有是藉着他而造成的；凡受造的，没有一样不是由他而造成的。在他内有生命，这生命是人的光。光在黑暗中照耀，黑暗不能胜过它。}我相信，到目前为止，我在经文的讲解上已经进展到这一步：让所有在场的人回忆当时所说的话；你们那些不在场的人，相信我以及那些选择在场的人。因此如今——因为我们不能总是重复一切，出于对那些希望听到后续内容的人的公道，并且因为重复之前的思想对他们是一个负担，剥夺了他们后续的内容——让那些上次缺席的人不要要求重述，而要同那些在场的人一起聆听现在的讲解。\par

2. 接着是：\BibleQuote{曾有一人，由天主派遣，名字叫若翰。}确实，亲爱的弟兄们，之前所讲的那些话，是关于基督不可言宣的天主性，而且几乎是不可言宣地讲的。因为谁能领悟\BibleQuote{在起初已有圣言，圣言与天主同在}呢？不要因日常言语的习惯而让\Quote{圣言}这个名称看起来平凡，因为紧接着说：\BibleQuote{圣言就是天主。}这个圣言就是昨天我们详细谈论的那一位；我相信天主临在，并且即使只有如此多的谈论，也触及了你们的心。\BibleQuote{在起初已有圣言。}他是相同的，并以同样的方式存在；正如他是，他永远是；他不能改变；也就是说，他\Emph{是}。这个他的名，他对他的仆人梅瑟说：\BibleQuote{我是自有者；那\Emph{自有者}派了我来。}\BibleRef{出 3:14}那么，当你们看到所有有死的事物都是可变的，当你们看到不仅身体因其性质而改变，通过出生、增长、衰微、死亡，而且连灵魂本身也因各种意志的作用而伸张和分裂，当你们看到人们若能专心致志于智慧的光和热就能获得智慧，若因某种邪恶影响而远离它就会失去智慧，谁能领悟呢？因此，当你们看到所有这些都是可变的，那\OriginalTerm{是者}{quod est}是什么，除非是那超越一切可变之物的东西？那么谁能领会呢？或者，无论以何种方式运用心灵的力量去触摸那\OriginalTerm{是者}{quod est}的人，能到达他心中以任何方式触摸到的那位吗？这就像一个人从远处看到他的故乡，中间隔着大海；他看到要去的地方，但没有办法去。同样，我们渴望到达我们那稳固的所在，那里\OriginalTerm{是者}{quod est}是，因为只有这一位永远如其所是；这个世界的海洋阻断了我们的道路，虽然我们已经看到我们要去的地方；因为许多人甚至没有看到他们要去哪里。为了使有一条路我们可以走，祂从我们想去的那一位那里来了。祂做了什么？祂设立了一棵树，使我们能穿过大海。因为没有人能渡过这个世界的海洋，除非被基督的十字架承载。就连视力弱的人有时也会拥抱这个十字架；那从远处看不到他要往哪里去的人，不要离开它，它会将他载过去。\par

3. 因此，我的弟兄们，我希望把这个印在你们心中：如果你们愿意过虔诚和基督徒的生活，要紧紧抓住基督，照着祂为我们成为的那一位，以便你们能到达祂按其本性与本来面目所在的那一位。祂接近我们，为使我们能成为这个；因为祂为我们成了那一位，软弱的人可以依靠祂，渡过这个世界的海洋，到达他们的故乡；在那里不需要船，因为没有海要渡。那么，与其用心灵看见那\OriginalTerm{是者}{quod est}，却离开基督的十字架，倒不如用心灵看不见它，却不离开基督的十字架。超越这个，最好的，如果可能的话，是我们既看到我们应当去的地方，又紧紧抓住那在我们前行时承载我们的。那些高山——他们被称为高山，是神圣正义之光照耀的杰出心灵——他们能够做到这一点，并且看见了那\OriginalTerm{是者}{quod est}。因为若翰看见时就说：\BibleQuote{在起初已有圣言，圣言与天主同在，圣言就是天主。}他们看见了，为了到达他们从远处看见的那一位，他们没有离开基督的十字架，也没有轻视基督的卑微。但那些不能理解这事的弱小者，若不离开基督的十字架、苦难和复活，就会被同一艘船引导到他们看不见的地方，而他们也将与那些看见的人一同到达那里。\par

4. 但确实有一些这世上的哲学家，曾藉着受造物寻求造物主；因为藉着受造物可以寻见祂，正如宗徒清楚所说的：\BibleQuote{其实，自从天主创世以来，祂那看不见的美善，即祂永远的大能和神性，都可凭祂所造的万物，辨认洞察出来，以致他们无可推诿。}接着又说：\BibleQuote{因为他们虽然认识天主，}他没有说因为他们不认识，而是说\BibleQuote{因为他们虽然认识天主，却没有以祂为天主而光荣祂，也不感谢祂，反而在思念中成了虚妄，他们那愚昧的心就昏暗了。}怎样昏暗的呢？接着他更明白地说：\BibleQuote{他们自称为智者，反而成了愚昧。}\BibleRef{罗 1:20}他们看见了他们必须到达的地方；但对赐予他们所看见的那一位忘恩负义，他们想把所见的一切归于自己；变得骄傲后，他们便丧失了所见的一切，转而去拜偶像和图像，去崇拜魔鬼，敬拜受造物而轻视造物主。但这些人因蒙蔽而行了那些事，并且变得骄傲，以致被蒙蔽：当他们骄傲时，就自称是智者。因此，关于他们，他所说的\BibleQuote{他们虽然认识天主}，就是指他们看见了若望所说的：万物是藉着天主的圣言造成的。因为这些事在哲学家的著作中也找得到：天主有一位独生子，万物都是藉着祂造成的。他们能看见那存在的，但只是远远地看见；他们不愿抓住基督的谦卑——这艘船可安全地将他们载到那他们从远处看见的地方，而基督的十字架在他们看来是可鄙的。海必须渡过，而你竟轻视那木头？啊，骄傲的智慧！你嘲笑被钉的基督；祂就是你从远处看见的那一位：\BibleQuote{在起初已有圣言，圣言与天主同在。}但祂为什么被钉呢？因为你需要祂屈尊的木头。你因骄傲而膨胀，被远远地逐出那父家；这世界的波涛阻断了道路，除了藉着木头，再也无法回到父家。忘恩负义的人！你嘲笑那来到你面前使你得以返回的那一位：祂成了道路，而且是穿越大海的道路；\BibleRef{玛 14:25}祂在海面上行走，以显示在海中有条道路。但你自己无法在海中行走，就乘船吧，被木头载着吧：相信被钉的那一位，你就能到达那里。祂为你被钉，为教导你谦卑；因为若祂以天主的身份来临，就不会被认出。因为若祂以天主的身份来临，祂就不会来到那些不能看见天主的人那里。因为按祂的天主性，祂既不前来也不离去；祂无所不在，不被任何地方容纳。那么，祂是怎样来临的呢？祂显现为人。\par

5. 因此，因为祂如此是人，以致天主隐藏于祂内，有一位伟大的人被派在祂前面，藉着他的见证，人们能发现祂不仅是一个人。这人是谁？\BibleQuote{有一个人。}那人怎能谈论天主的真理呢？\BibleQuote{是由天主派遣来的。}他叫什么名字？\BibleQuote{名字叫若翰。}他为什么来？\BibleQuote{他为作证而来，为给光作证，为使众人藉着他而信。}这位要为光作证的人是怎样的人呢？那位若翰是伟大的，拥有极大的功德，极大的恩宠，极大的崇高！你尽可钦佩，钦佩吧；但把他当作一座山。但山若是未被光覆盖，仍处于黑暗中。因此，你只钦佩若翰，好能听见接下来所说的：\BibleQuote{他不是那光。}免得你若将山当作光，就在山上触礁，而得不到安慰。但你该钦佩什么呢？钦佩山本身。但你要举目仰望那光照山的那一位，山之所以被高举，正是为了首先接受光线，并将光线传给你的眼睛。因此，\BibleQuote{他不是那光。}\par

6. 那么他为什么来呢？\BibleQuote{而是为给光作证。}为什么？\BibleQuote{为使众人藉着他而信。}他要为怎样的光作证呢？\BibleQuote{那光是真的。}为什么加上\BibleQuote{真的}呢？因为一个被光照的人也被称为光；但真光是那光照的光。连我们的眼睛也称为光；然而，若非夜里点灯，或白日太阳出现，这些光就空开着。同样，若翰是光，但不是真光；因为若未被光照，他就是黑暗；但因被光照，他成了光。因为他若未被光照，就会是黑暗，正如所有那些曾是不敬神的人，宗徒对信仰的他们说过：\BibleQuote{你们从前原是黑暗。}但现在，因为他们相信了，怎么样？——\BibleQuote{如今你们在主内却是光。}\BibleRef{弗 5:8}除非他加了\BibleQuote{在主内}，我们就不能明白。他说：\BibleQuote{在主内是光：}你们从前不是在主内的黑暗。\BibleQuote{因为你们从前原是黑暗，}那里他没有加\BibleQuote{在主内}。因此，黑暗在你们内，光在主内。这样，\BibleQuote{他不是那光，而是被派遣来为光作证。}\par

7. 但那光在哪里？\Quote{他才是真光，照亮那进入世界的每一个人。}如果每一个人，那么也包括\Person{若翰}。因此，真光照亮了他，使他成为指向那光的见证人。亲爱的诸位，请理解：他来到软弱的心智、受伤的心灵、目光昏暗的灵魂面前。他正是为此而来。灵魂怎能看见那圆满的存有呢？正如常见的那样，通过某个被光照的物体，我们才知道肉眼无法直视的太阳已经升起。因为即使眼睛受伤的人也能看到被太阳照亮的墙壁、山峦、树木或类似的东西；通过另一个被照亮的物体，那升起便向那些尚不能直视它的人显明。同样，所有那些\Person{基督}来到他们面前的人，并不适合看见他：他将他的光芒投射在\Person{若翰}身上；藉着\Person{若翰}承认自己被照射和照亮，而不是声称自己是照射和照亮者，那照亮者、那光照者、那充满者便被知晓了。那是谁呢？\Quote{他照亮每一个人，}他说，\Quote{那进入世界的。}因为人若没有远离那光，就不需要被照亮；但正因如此，他需要在此被照亮，因为他离开了那原本可以一直照亮他的光。\par

8. 那么，如果他来到这里，他在哪里？\Quote{他在世界之中。}他既在这里，又来到这里；他按天主性在这里，按肉体来到这里；因为当他按天主性在这里时，不能被愚昧、盲目和邪恶的人看见。这些邪恶的人就是黑暗，关于这黑暗曾说过：\Quote{光在黑暗中照耀，黑暗却没有接受它。}\BibleRef{若 1:5}看，他现在就在这里，以前也在这里，永远在这里；他从不离开，不到别处去。你需要有一些方法，藉以看见那从不离开你的那一位；你需要不离开那从不离开任何地方的那一位；你需要不抛弃，你就不会被抛弃。不要跌倒，他的太阳就不会向你落下。如果你跌倒，他的太阳就在你身上落下；但如果你站立，他就与你同在。但你并没有站立：记住你是如何跌倒的，那在你之前跌倒的是如何把你拉下去的。因为他把你拉下去，不是凭暴力，不是凭攻击，而是凭你自己的意志。因为如果你没有同意罪恶，你就会站立，就会保持光明。但现在，因为你已经跌倒，心中受了伤——那能看见那光的器官——他便以你所能看见的方式来到你面前；他如此显现为人，以致他从人寻求见证。天主从人寻求见证，而天主以人为见证——天主以人为见证，但这是为了人：我们如此软弱。我们藉着灯寻求白昼；因为\Person{若翰}自己被称为灯，主说：\Quote{他是一盏燃烧发光的灯；你们曾经情愿在他的光中暂时欢乐；但我有比若翰更大的见证。}\BibleRef{若 5:35}\par

9. 因此，他表明，为了人的缘故，他愿意藉着一盏灯向相信的人显明自己，使他的仇敌因同一盏灯而蒙羞。有些仇敌试探他，说：\Quote{告诉我们，你凭什么权柄做这些事？}\Quote{我也，}他说，\Quote{要问你们一个问题；回答我。\Person{若翰}的圣洗是从哪里来的？是从天上，还是从人来的？}他们感到困扰，彼此说：\Quote{如果我们说\InnerQuote{从天上来}，他会对我们说：\InnerQuote{你们为什么不信他？}}（因为他已为\Person{基督}作证，并说：\Quote{我不是\Person{基督}，而是他。}）\Quote{但如果我们说\InnerQuote{从人来的}，我们怕群众，因为他们会拿石头打我们：因为他们认为\Person{若翰}是一位先知。}他们害怕被石头砸，但更害怕承认真理，就对真理说了谎言；\Quote{邪恶将谎言加诸自身。}因为他们说：\Quote{我们不知道。}主见他们自封了门，便没有给他们开门，因为他们没有叩门。因为经上说：\Quote{叩门，就给你们开门。}\BibleRef{玛 7:7}这些仇敌不仅没有叩门让门为他们打开，反而因否认他们知道而自封了门。主对他们说：\Quote{我也不告诉你们我凭什么权柄做这些事。}他们便因\Person{若翰}而蒙羞；在他们身上应验了这话：\Quote{我为我的受傅者预备了一盏灯。我要使他的仇敌披上羞耻。}\par

10. \Quote{他在世界之中，世界也是藉着他造的。}不要以为他在世界中就像大地在世界中、天空在世界中、太阳在世界中、月亮和星辰、树木、牲畜和人在世界中那样。他不是那样在世界中。而是以怎样的方式呢？如同工匠治理他所造的工。因为他造世界不像木匠造箱子。木匠造的箱子在木匠之外，所以它在被造时放在另一个地方；虽然工匠在旁边，他坐在另一个地方，并且外在于他所造的物件。但天主临在于世界而造它；他无所不在，造它，并不退到别处，也不像从外面处理他所造的材料。凭着他的尊威的临在，他造他所造的；他的临在治理他所造的。因此，他在世界中作为世界的造物主；因为\Quote{世界是藉着他造的，世界却不认识他。}\par

11. \Quote{世界是藉着祂造成的}是什么意思？天、地、海和其中所有的一切，都称为世界。此外，另有意义：那些爱世界的人也称为世界。\Quote{世界是藉着祂造成的，世界却不认识祂。}难道诸天不认识自己的造主吗？难道天使不认识自己的造主吗？就连魔鬼也承认祂，星辰不也承认祂吗？万物从各方面都作见证。但不认识祂的是谁呢？是因爱世界而被称为世界的人。我们以心居住；由于他们爱世界，便配以他们所居住的那名字来称呼。好比我们说：这房子不好，或这房子好；我们并非在说不好或好的时候指责或赞美墙壁；而是说坏的房子是指有坏居民的房子，好的房子是指有好居民的房子。同样，我们称那些因爱世界而居住在世界中的人为世界。他们是谁？就是爱世界的人；因为他们以心居住在世界。至于那些不爱世界的人，的确，他们在肉身中寄居世界，但他们的心却居住在天上，正如宗徒所说：\Quote{我们的户籍是在天上。}因此\Quote{世界是藉着祂造成的，世界却不认识祂。}\par

12. \Quote{祂来到自己的地方}——因为这一切都是藉着祂造成的——\Quote{自己的人却没有接纳祂。}他们是谁？是祂所造的人。犹太人，就是祂起初使之超越万民的人。因为其他民族崇拜偶像、服事魔鬼，但那民族是出于亚巴郎的后裔，并且由于祂屈尊所取的那肉身，在卓越的意义上属于祂自己。\Quote{祂来到自己的地方，自己的人却没有接纳祂。}他们全然没有接纳祂吗？难道没有人接纳祂？难道没有人得救？因为除非有人接纳来临的基督，否则无人得救。\par

13. 但若望接着说：\BibleQuote{凡接纳祂的。}祂赐给了他们什么？极大的仁慈！极大的怜悯！祂是天主的独生子，却不愿独自一人。许多人没有儿子时，在晚年收养一个儿子，借此用意志获得本性所拒绝的：这是人所行的。但若有人有独生子，他更以他为乐；因为那独生子将拥有一切，不会有人与他分遗产而使他贫穷。天主却不是这样：祂所生的独生子，并藉着祂创造万有的那一位，祂派遣到这个世界，为叫祂不独自一人，而能有收养的弟兄。因为我们并非如独生子那样由天主所生，而是因祂的恩宠被收养。因为祂，那独生子，来解除我们被缠绕的罪，这些罪的负担阻碍了我们的收养：祂亲自解开了那些祂愿意使之成为自己弟兄的人，并使他们成为共同继承人。因为宗徒这样说：\BibleQuote{但若是儿子，便是藉着天主成为继承人。}又说：\BibleQuote{天主的继承人，与基督同为继承人。}祂不怕有共同继承人，因为祂的遗产不会因许多人占有而变窄。正是那些人，当祂占有时，成了祂的产业，而祂也成了他们的产业。请听他们如何成为祂的产业：\BibleQuote{上主对我说：祢是我的儿子，我今日生祢。祢向我请求，我必将万民赐祢作产业。}请听祂如何成为他们的产业：祂在圣咏中说：\BibleQuote{上主是我的产业，是我杯中的分。}让我们占有祂，也让祂占有我们：祂占有我们如主；我们占有祂如救恩，占有祂如光明。那么，祂赐给了那些接纳祂的人什么呢？\BibleQuote{祂赐给他们权能，成为天主的子女，就是那些信祂名字的人；}好使他们抓住木杖而渡过大海。\par

14. 他们是如何出生的呢？因为他们成为天主的子女和基督的弟兄，他们当然是出生的。因为若不出生，怎能成为儿子呢？但人的儿子是由血肉、人欲和婚姻的拥抱而生的。然而，他们是如何出生的呢？\Quote{他们不是由血气，}好比男女的混合而生。\Quote{血气}一词并非拉丁文；但因在希腊文中是复数，译者宁可用这样的表达，并且说文法上不通的拉丁文，以便使听众中的软弱者明白。因为若他用单数的\Quote{血}，就不能说明他想要表达的意思；因为人是因男女的混合之血而生的。那么，我们就这样说，不要害怕文法家的戒尺，只要我们达到确定而稳固的真理。那理解而又责备的人，对于他所理解的，是忘恩负义的。\Quote{不是由血气，也不是由肉欲，也不是由人欲。}宗徒以\Quote{肉}指女人；因为当女人由亚当的肋骨被造时，亚当说：\BibleQuote{这才是我的骨中之骨，肉中之肉。}\BibleRef{创 2:23} 宗徒又说：\BibleQuote{爱自己妻子的，就是爱自己；因为从来没有人恨过自己的肉身。}\BibleRef{弗 5:28} 因此，\Quote{肉}被用来指女人，正如\Quote{灵}有时被用来指丈夫。为什么？因为一个治理，一个被治理；一个应当吩咐，一个应当服事。因为当肉命令而灵服事时，家就倒过来了。还有什么比女人管辖男人的家更糟呢？但男人命令、女人服从的家，才是有秩序的。同样，当灵命令、肉服从时，那人才是有秩序的。\par

15. 这些人\BibleQuote{不是由血气，也不是由肉欲，也不是由男欲，而是由天主生的。}但为使人生于天主，天主首先生于人。因为基督是天主，基督亦生于人。的确，祂在世上只寻求一位母亲，因为祂在天上已有了一位父亲：那创造我们的祂生于天主，那再造我们的祂生于女人。那么，人啊，不要惊奇你因恩宠成为儿子，你依照祂的圣言生于天主。圣言自己首先选择生于人，好使你为得救恩而生于天主，并对你自己说：天主愿意生于人并非无缘故的，而是因为祂视我为重要，好使我成为不朽，并为我而生为一个有死的人。因此，当祂说了\Quote{生于天主}后，免得我们对此恩宠感到惊愕和战栗，这恩宠如此伟大以致难以置信人竟能生于天主，祂仿佛向你保证，说：\BibleQuote{圣言成了血肉，寄居在我们中间。}那么，你为什么惊奇人竟能生于天主呢？请看天主自己生于人：\BibleQuote{圣言成了血肉，寄居在我们中间。}\par

16. 但因\BibleQuote{圣言成了血肉，寄居在我们中间}，祂藉着自己的诞生制作了眼药，来洗净我们心灵的眼睛，并使我们能藉着祂的谦卑看见祂的威严。因此\BibleQuote{圣言成了血肉，寄居在我们中间}：祂治愈了我们的眼睛；接着呢？\BibleQuote{我们见了祂的光荣。}非经祂血肉的谦卑治愈，无人能见祂的光荣。那么，我们为何不能看见呢？请思考，亲爱的诸位，并领会我所说的。好像尘土、泥土溅入人的眼睛，伤了眼睛，使它不能见光；那受伤的眼睛被涂上药膏；被泥土所伤，也用泥土医治。因为所有眼药和药物都取自泥土。你被尘土弄瞎，也被尘土治愈：那么，血肉伤了你，血肉也治愈你。灵魂因顺从肉体的私欲而成为属血肉的；因此心灵的眼睛被弄瞎了。\BibleQuote{圣言成了血肉}：那位医生为你制作了眼药。正如祂藉血肉来临以熄灭血肉的恶习，藉死亡以消灭死亡；因此在你内发生了这样的事，以致当\BibleQuote{圣言成了血肉}时，你能够说：\BibleQuote{我们见了祂的光荣。}什么样的光荣？就如祂成为人子那样？那是祂的谦卑，而非祂的光荣。然而，当人的眼目因血肉而被治愈后，它得以看见什么呢？\BibleQuote{我们见了祂的光荣，正如父独生者的光荣，满溢恩宠和真理。}关于恩宠和真理，我们将在同一福音的另一处更详尽地谈论，若主赐予我们机会。目前让这些事足够，愿你们在基督内得到建树：在信德中得到安慰，在善行中保持警醒，并注意不要离开那可使你渡海的木头。\par

\SourceRecord{Translated by John Gibb.；Nicene and Post-Nicene Fathers, First Series；Vol. 7.；Edited by Philip Schaff.；Buffalo, NY: Christian Literature Publishing Co.,；1888.；<http://www.newadvent.org/fathers/1701002.htm>.}

% structure: p0001 source=h1 target=section reason=page-title

\section{\WorkTitle{论道第三篇（\BibleRef{若 1:15-18}）}}

我们曾奉主的名起意，并向你们，亲爱的诸位，应许要论述天主的恩宠和真理，天主的独生子，我们的主救主耶稣基督，就是充满了这恩宠和真理而显现给圣人们的，并要说明，这作为新约的一部分，如何与旧约相区别。那么，请你们注意，使我按我所领受的，从天主那里得到什么，你们也按你们所能领受的，听到同样的话。因为只有当种子撒在你们心里时，鸟儿不把它啄去，荆棘不把它窒息，炎阳不把它晒焦，并且有每日劝勉的雨水和你们自己的善念降下，如同在田里用耙子所做的那样，使土块破碎，种子被覆盖并得以发芽，你们才能结出农夫喜悦的果子。但如果你们以好种子和好雨水回报，不结果子却结荆棘，那么种子不受责备，雨水也不担过失；但荆棘所应得的是火。\BibleRef{玛 13:3}\par

我认为无需花费很多时间来努力说服你们，我们是基督徒；如果真是基督徒，就因这名分属于基督。我们在额上带着祂的记号；如果我们也在心里带着它，就不以此为耻。祂的记号就是祂的谦卑。贤士们藉着一颗星认识了他；\BibleRef{玛 2:2}这记号是由主赐下的，是属天的、美丽的。祂不愿星成为信友额上的记号，而是祂的十字架。藉着它，祂被贬抑，也藉着它，祂受光荣；藉着它，祂举扬了谦卑的人，甚至藉着祂在受贬抑时所降下的那个十字架。那么，我们属于福音，我们属于新约。\Quote{法律是藉梅瑟传授的，恩宠和真理却是由耶稣基督而来的。}我们询问宗徒，他对我们说，因为我们不在法律之下，而在恩宠之下。\BibleRef{罗 6:14}\Quote{因此，天主派遣了自己的儿子，生于女人，生于法律之下，为要把在法律之下的人赎出来，使我们获得义子的地位。}\BibleRef{迦 4:4}请看，基督为此而来，为要把在法律之下的人赎出来，使我们如今不在法律之下，而在恩宠之下。那么，谁赐下了法律？那赐下法律者，也赐下了恩宠；但法律是由一个仆人传送的，而恩宠是祂亲自降下的。人们是怎样成为在法律之下的呢？由于不遵守法律。因为遵守法律的人，不是法律之下，而是与法律同在；但在法律之下的人，不被高举，反被法律压下。因此，所有的人都处于法律之下，由法律而定罪；为此，法律在他们头上，为显示罪恶，而非除去罪恶。法律命令，法律的赐予者就在法律所命令的事上显示怜悯。人们试图凭自己的力量履行法律的命令，却因自己的鲁莽和狂妄而跌倒；不与法律同在，反在法律之下，成为有罪的人；并且因他们自己的力量无法履行法律，既已在法律之下成为有罪，他们就呼求解救者的帮助；法律所带来的罪责使骄傲者生病。骄傲者的病痛成了谦卑者的忏悔。如今病人承认自己生病；让医生来医治病人。\par

谁是医生？我们的主耶稣基督。谁是我们的主耶稣基督？就是那位被钉死祂的人所看见的那一位。祂被捉拿、掌掴、鞭打、唾弃、戴荆棘冠、悬在十字架上、死亡、被长矛刺穿、从十字架上卸下、安放在坟墓里。那位耶稣基督我们的主，正是那位，完全是我们的伤口的医生。那位被凌辱的、钉在十字架上的，当祂悬在十字架上时，迫害者摇头说：\Quote{如果他是天主子，让他从十字架上下来吧！}\BibleRef{玛 27:39}——祂，别无他人，是我们的全备医生。那么，为什么祂没有向嘲弄祂的人显示祂是天主子呢？如果祂允许自己被举在十字架上，至少当他们说：\Quote{如果他是天主子，让他从十字架上下来吧！}的时候，祂就该下来，向他们表明祂就是他们所敢嘲弄的天主子？祂不愿意。为什么祂不愿意？是因为祂不能吗？显然祂能。因为哪一个更大，是从十字架上下来，还是从坟墓中复活？但祂容忍了那些凌辱祂的人；因为十字架不是作为能力的证明，而是作为忍耐的榜样被取用的。祂在那里医治了你的创伤，就在祂长久背负自己的创伤的地方；祂在那里治愈你于永死，就在祂俯允接受暂时死亡的地方。祂是死了呢，还是在祂内死亡死了？那是怎样的一种死亡，竟致死亡死了！\par

4. 然而，是否就是我们的主耶稣基督自己——祂的整个自我——被看见、被抓住并被钉在十字架上呢？那同一的自我吗？犹太人看见的是同一个，但不是整个；这不是整个基督。那么什么是呢？\BibleQuote{在起初已有圣言。}在什么起初？\BibleQuote{圣言与天主同在。}什么圣言？\BibleQuote{圣言就是天主。}那么，也许这圣言是由天主所造的吗？不。因为\BibleQuote{祂在起初就与天主同在。}那么，其他由天主所造的东西不像圣言吗？不：因为\BibleQuote{万物是藉着祂造的；凡受造的，没有一样不是藉着祂造的。}万物是怎样藉着祂造的？因为\BibleQuote{在祂内有生命，这生命就是人的光；}在受造之前就有生命。受造的不是生命；但在技艺中，也就是说，在天主的智慧中，在受造之前它就是生命。受造的会消逝；在智慧中的不会消逝。因此，在受造物之中有生命。是什么样的生命，既然灵魂也是身体的生命？我们的身体有自己的生命；当它失去生命时，身体就死亡。那么，生命就是这样吗？不；而是\BibleQuote{这生命就是人的光。}它是牲畜的光吗？因为光是人类和牲畜的光。有一种人的光：让我们看看人如何不同于牲畜，然后我们就明白什么是人的光。你除了理智外与牲畜没有区别；不要在其他方面夸耀。你以自己的力量自夸吗？你被野兽超越。你以自己的速度自夸吗？你被苍蝇超越。你以自己的美貌自夸吗？孔雀的羽毛是多么美丽！那么你在哪方面更好？在天主的肖像上。天主的肖像在哪里？在心灵中，在理智中。那么，如果你在这方面比牲畜好，因为你有理智可以理解牲畜无法理解的事；人因此比牲畜好，人的光就是心灵的光。心灵的光超越心灵，超越所有心灵。这就是那生命，万物藉着它而被造。\par

5. 它在哪里？在这里吗？它与父在一起，而不在这里吗？或者，更真实地说，它既与父在一起，也在这里吗？那么，如果它在这里，为什么看不见？因为\BibleQuote{光在黑暗中照耀，黑暗却没有接纳它。}噢，人们，不要成为黑暗，不要成为不信、不义、不正直、贪婪、贪爱这个世界的人：因为这些就是黑暗。光并不缺席，而是你们远离了光。一个盲人在阳光下，太阳就在他面前，但他自己却远离太阳。那么，不要成为黑暗。因为这或许就是我们将要谈论的恩宠，使得我们不再成为黑暗，并且宗徒可以对我们说：\BibleQuote{我们从前是黑暗，但现在在主内是光明。}\BibleRef{弗 5:8}因为那时人的光，即心灵的光，没有被看见，所以必须有一个人为光作证，他不在黑暗中，而是已经受到光照；然而，因为他受了光照，却不是光本身，\BibleQuote{而是为光作证。}因为\BibleQuote{他不是那光。}那么光是什么？\BibleQuote{那才是真光，照亮每个人来到世上。}那么那光在哪里？\BibleQuote{它在这世上。}它怎样\BibleQuote{在这世上}？像太阳、月亮、灯的光一样，那光这样在世上吗？不。因为\BibleQuote{世界是藉着祂造的，世界却不认识祂。}也就是说，\BibleQuote{光在黑暗中照耀，黑暗却没有接纳它。}因为世界是黑暗；因为爱世界的人就是世界。难道受造物不认识它的创造主吗？诸天藉着一颗星作证；\BibleRef{玛 2:2}海洋作证，当主行走其上时承载了祂；\BibleRef{玛 14:26}风作证，听从祂的命令而平静；\BibleRef{玛 23:27}大地作证，当祂被钉在十字架上时震动。\BibleRef{玛 27:51}如果所有这些都作了证，那么世界在什么意义上不认识祂？除非世界指那些爱世界的人，那些心中住在世界里的人。而且世界是邪恶的，因为世界的居民是邪恶的；正如一所房子是邪恶的，不是因为墙壁，而是因为它的居民。\par

6. \BibleQuote{祂来到自己的地方，}也就是说，祂来到属于祂自己的地方；\BibleQuote{自己的人却不接受祂。}那么，希望是什么？除非\BibleQuote{凡接受祂的，祂给他们权能成为天主的子女。}如果他们成为子女，他们就生了；如果生了，是怎么生的？不是从血肉，\BibleQuote{也不是从血，也不是从肉欲，也不是从人的意愿，而是从天主生的。}因此，让他们欢欣，他们是从天主生的；让他们相信他们是从天主生的；让他们接受他们是从天主生的证据：\BibleQuote{圣言成了血肉，寄居在我们中间。}如果圣言不以从人生为耻，人怎以从天主生为耻呢？因为祂做了这事，祂医治了我们；因为祂医治了我们，我们得以看见。因为这一点，\BibleQuote{圣言成了血肉，寄居在我们中间，}成了我们的良药，好使我们因泥土而瞎眼，也因泥土而被医治；被医治后，我们能看见什么？\BibleQuote{我们看见了祂的光荣，}他说，\BibleQuote{正如父独生者的光荣，充满恩宠和真理。}\par

7. \Quote{若翰为祂作证，呼喊说：\InnerQuote{这就是我曾说的：那在我以后来的，成了在我以前的。}}祂在我以后来，却在我以前。\Quote{祂成了在我以前的}是什么意思？祂在我以前。不是在我被造以前被造，而是被置于我之上，这就是\Quote{祂成了在我以前的}。为什么祂成了在你以前，却在你以后来？\Quote{因为祂在我以前。}在你以前，若翰哪！在你以前有何了不起！你为祂作证倒是很好；但我们还是听祂亲自说的：\Quote{在亚巴郎以前，我就存在。}\BibleRef{若 8:58}但亚巴郎也是人类中间出生的：在他以前有许多人，以后也有许多人。请听圣父对圣子说：\Quote{在路济弗尔以前，我已生了祢。}那在路济弗尔以前受生的，祂亲自光照一切。有一个名叫路济弗尔的，他堕落了；因为他原是天使，成了魔鬼；关于他，圣经说：\BibleQuote{路济弗尔，清晨升起的，堕落了。}\BibleRef{依 14:27}为什么他叫\Emph{路济弗尔}？因为他是受光照而发光的。但为何他又成了黑暗？因为他没有存于真理中。\BibleRef{若 8:44}因此，祂在路济弗尔以前，在每一个受光照者以前；因为每一个能受光照者，都必须由那光照一切者所光照。\par

8. 因此接着就说：\Quote{从祂的圆满中，我们都领受了。}你们领受了什么？\Quote{以及恩宠上加恩宠。}因为按福音原文如此，我们对照希腊文抄本可以发现。祂不是说：从祂的圆满中我们都领受了恩宠上加恩宠；而是这样说：\Quote{从祂的圆满中我们都领受了，以及恩宠上加恩宠，}——也就是说，我们领受了；这表示祂要我们明白，我们从祂的圆满中领受了某种未明言的东西，此外还有恩宠上加恩宠。因为我们从祂的圆满中首先领受了恩宠；然后又领受了恩宠，即恩宠上加恩宠。我们首先领受的是什么恩宠？信德：行走在信德中，我们行走在恩宠中。我们如何获得这恩宠？凭我们什么先前的功德？不要有人自夸，而要返回自己的良心，探询自己思想的隐秘处，回忆自己行为的系列；他不要考虑自己现在是什么（如果他现在是某种人），而要考虑他过去是什么，以便成为某种人：他会发现自己只配受惩罚。因此，如果你只配受惩罚，而祂来不是惩罚罪，而是赦罪，那么你所领受的是恩宠，而不是报酬。为什么称为恩宠？因为是白白赐予的。因为你并没有凭先前的功德购得你所领受的。那么，罪人领受了这最初的恩宠，即他的罪得到赦免。他凭什么配得？让他质问正义，他找到惩罚；让他质问仁慈，他找到恩宠。但天主也曾通过先知应许了这事；因此，当祂来赐下所应许的时，祂不仅赐下恩宠，也赐下真理。真理如何显明？因为所应许的成就了。\par

9. 那么，\Quote{恩宠上加恩宠}是什么？藉着信德，我们使天主对我们和好；既然我们本不配得赦罪，而我们这些不配的人领受了如此大的恩惠，这便称为恩宠。什么是恩宠？白白赐予的。什么是\Quote{白白赐予的}？赐予，而非付还。如果是应得的，那就是付给工资，而非赐予恩宠；但如果是应得的，你就是善的；但如果（正如事实）你是恶的，却信了那使不虔敬的人成义者\BibleRef{罗 4:5}（什么是\Quote{使不虔敬的人成义}？从不虔敬的造出虔诚的），那么你当思量，律法按理当使你受什么惩罚，而你藉恩宠得到了什么。但既获得了信德的恩宠，你便因信德而成义（因为义人因信德而生活）；并因凭信德生活而得天主欢心。因凭信德生活而得天主欢心后，你将获得不朽作为赏报，即永生。而那也是恩宠。因为你凭什么功德得永生？因恩宠。因为如果信德是恩宠，永生便好像是信德的工资：天主似乎是将永生作为应得的赏赐（应得给谁？给信者，因为他凭信德赚得了它）；但因为信德本身就是恩宠，所以永生也是恩宠上加恩宠。\par

10. 请听宗徒保禄承认恩宠，后来又要求还债。保禄如何承认恩宠？他说：\Quote{我从前是亵渎者、迫害者和施暴者，但我蒙受了仁慈。}\BibleRef{弟前 1:13}他说那蒙受仁慈的人是不配的，但他之所以蒙受，并非由于自己的功绩，而是由于天主的仁慈。现在请听他，这位先领受了无功的恩宠，后来要求还债的人说：\Quote{因为我已被奠祭，我离世的时刻已经临近。这场好仗，我已打完；这场赛跑，我已跑到终点；这信德，我已保持了。从此以后，正义的冠冕已为我预备下了。}\BibleRef{弟后 4:6}现在他要求还债，索取应得的。请看随后的话：\Quote{在那日，主，正义的审判者，要赏给我。}他先前领受恩宠时，需要仁慈的圣父；而为了恩宠的报酬，则需要正义的审判者。那没有定不敬神者罪的那一位，岂会定信者之罪呢？然而，你若正确思考，正是祂首先赐给你信德，使你蒙恩；因为你并非出于自己而蒙恩，以至于有任何东西应归于你。因此，后来赐予永生的赏报时，祂是加冕祂自己的恩赐，而不是你的功绩。所以，弟兄们，我们都从祂的圆满中领受了；从祂仁慈的圆满、良善的丰盈中我们领受了。领受了什么？罪的赦免，使我们因信德而成义。此外还有什么？恩宠上加恩宠；也就是说，为了这使我们因信德而生活的恩宠，我们将领受另一种恩宠。那么，这难道不是恩宠吗？因为若我说这也是应得的，我就把某些东西归于自己，仿佛它应归我。但天主在我们身上加冕的是祂自己仁慈的恩赐；条件是我们要恒心持守那起初所领受的恩宠。\par

11. \Quote{因为律法是藉梅瑟颁赐的；}这律法定了罪人的罪。因为宗徒说什么？\Quote{律法闯入，为使过犯增多。}这对骄傲的人是有益的，因为过犯增多了，他们便多多归功于自己，仿佛把自己的力量看得很高；而他们若没有那命令者的帮助，便无法完成正义。天主为压制他们的骄傲，赐下了律法，仿佛说：看，你们去履行吧，不要以为缺少一位命令者。命令者并不缺少，但缺少一位履行者。\par

12. 那么，若缺少一位履行者，他是为何不能履行呢？因为生来就带有罪与死亡的遗产。由亚当所生，他带出了那在亚当内所成孕的。第一个人堕落了，所有从他而生的人，都从他那里继承了肉体的私欲。需要另一个人出生，不继承任何私欲。一个人和一个人：一个至于死亡的人，一个至于生命的人。宗徒这样说：\Quote{既然死亡藉着一个人，死者的复活也藉着一个人。}哪一个人带来死亡，哪一个人带来死者的复活？不要急于分辨；他接着说：\Quote{因为正如在亚当内众人都死了，照样在基督内众人都要复活。}\BibleRef{格前 15:21}谁属于亚当？所有由亚当所生的人。谁属于基督？所有藉着基督而重生的人。为什么众人都在罪中？因为除了藉着亚当，无人出生。然而，他们由亚当所生是出于定罪的必然；藉着基督而重生则是出于意愿和恩宠。人并非被迫藉着基督而重生：他们并非因自己的意愿而由亚当所生。然而，所有属于亚当的人，都是带着罪的罪人；所有属于基督的人，都被成义，并且正义不在于他们自己，而在于祂。因为在他们自己身上，你若问，他们属于亚当；在祂内，你若问，他们属于基督。为什么？因为祂，元首，我们的主耶稣基督，并非带着罪的遗产而来；但祂仍然带着可朽的血肉而来。\par

13. 死亡是罪的刑罚；在主内，死亡是仁慈的恩赐，而非罪的刑罚。因为主没有任何可使他公正地死去的理由。祂自己说：\Quote{看，这世界的首领来了，他在我内一无所获。}那么你为何要死？\Quote{但为叫众人知道，我实行我父的旨意，起来，我们离开这里罢。}\BibleRef{若 14:30}祂自己没有任何必须死去的理由，却死了；你既有这样的理由，难道拒绝死吗？不要拒绝以平等的心承担你所应得的，当祂未曾拒绝受苦，为将你从永死中拯救出来时。一个人和一个人；但前者只是人，后者却是天主而人。前者是罪的人，后者是正义的人。你在亚当内死了，在基督内复活；因为两者都归于你。现在你已相信了基督，仍须偿还你因亚当所欠的债。但罪的锁链不会永远束缚你，因为你的主的暂时死亡杀死了你的永死。这就是恩宠，弟兄们，这就是真理，因为曾被预许，且已显明。\par

14. 这恩宠并不存在于旧约中，因为律法威胁，却不带来援助；命令，却不医治；显明软弱，却不除去软弱；但它为那将要带着恩宠和真理而来的医生预备了道路；就像一位医生，即将来治疗某人，先派他的仆人，使他发现病人被捆绑着。病人并不健康，他不愿被治好，并且为了不被治好，竟夸耀自己健康。律法被派遣了，它捆绑了他；他发现自己被控告，现在他对那绷带喊叫。主来了，用有些苦涩和锐利的药物医治：因为祂对病人说：忍受；祂说：忍耐；祂说：不要爱世界，要有耐心，让节制的火医治你，让你的伤口忍受迫害的剑。你虽被捆绑，岂不极其惊恐吗？祂，自由而无捆绑的，饮下了祂所赐给你的；祂先受苦，为安慰你，仿佛说：你害怕为自己所受的，我先为你受了。这就是恩宠，伟大的恩宠。谁能以相称的方式赞美它呢？\par

15. 我讲论基督的谦卑，各位弟兄。谁能讲论基督的尊威和基督的天主性呢？在解释和讲论基督的谦卑时，我们发现无论如何都不足够，实际上是全然不足的：我们将祂完全交托给你们的思考，我们并不试图在你们的听闻中充满祂。请思考基督的谦卑。但你们说，除非你们宣告，谁能向我们解释呢？让祂在内心宣告。那住在内心的比在外面呼喊的更善于宣告。愿祂亲自向你们显示祂谦卑的恩宠，祂已开始居住在你们心中。但现在，如果在解释和陈述祂的谦卑时我们有所欠缺，谁能讲论祂的尊威呢？如果\Quote{圣言成了血肉}扰乱了我们，谁能解释\Quote{在起初已有圣言}呢？因此，各位弟兄，请持守基督的完整。\par

16. \Quote{法律是藉梅瑟颁布的，恩宠和真理是由耶稣基督而来的。}法律是由一个仆人颁布的，并使人类有罪；恩赦是由一位皇帝颁布的，并释放了有罪者。\Quote{法律是藉梅瑟颁布的。}不要让仆人将超过他所做的事归于自己。他被拣选担任伟大的职务，如同他家中忠信的仆人，但仍是仆人，他能依照法律行事，却不能解除法律的罪责。\Quote{法律}于是\Quote{是藉梅瑟颁布的：恩宠和真理是由耶稣基督而来的。}\par

17. 并且恐怕有人会说：难道恩宠和真理不是通过看见天主的梅瑟而来的吗？他立刻加上：\Quote{从来没有人见过天主。}那么，天主如何为梅瑟所知呢？因为主向祂的仆人启示了自己。什么主？就是那同一位基督，祂预先藉祂的仆人颁布法律，为使自己带着恩宠和真理而来。\Quote{因为从来没有人见过天主。}那么，祂如何向那位仆人显现，正如他所能接受的呢？但他说：\Quote{那独生子，}他说，\Quote{在圣父怀中的，祂将祂启示出来。}\Quote{在圣父怀中}是什么意思？是在圣父的隐秘处。因为天主并没有像我们衣服上的怀，也不应当被想成是坐着的，像我们那样，也没有束上带子以有一个怀；但因为我们的怀是在内部的，所以圣父的隐秘处被称为圣父的怀。那位认识圣父的，既然在圣父的隐秘处，就将祂启示出来。\Quote{因为从来没有人见过天主。}祂于是来了，并叙述了祂所见的一切。梅瑟看见了什么？梅瑟看见了一朵云，看见了一位天使，看见了一火。所有这些都是受造物：它承载了其主的形象，但并未显示主本身的临在。因为你们在法律中明确记载：\Quote{梅瑟曾与主面对面谈话，如同朋友与朋友谈话。}按照同样的圣经，你们发现梅瑟说：\Quote{如果我在祢眼中蒙恩，求祢清楚地将祢自己显示给我，使我看见祢。}他说这话还嫌不够：他得到的答复是：\Quote{你不能看见我的面。}因此，各位弟兄，是一位天使与梅瑟交谈，承载着主的形象；天使所做的一切事情都预许了未来的恩宠和真理。那些仔细研究法律的人知道这一点；当我们有机会也论及此事时，我们不会忽略向你们讲论，亲爱的各位弟兄，正如主向我们启示的那样。\par

18. 但你们要知道，所有那些以形体所见的事物都不是天主的本体。因为我们是藉肉体的眼睛看见那些事物的：天主的本体如何被看见呢？请问福音：\BibleQuote{心里洁净的人是有福的，因为他们要看见天主。}\BibleRef{玛 5:8}有人被自己内心的虚妄所欺骗，曾说：圣父是不可见的，但圣子是可见的。如何可见？如果是因为祂的肉身，因为祂取了肉身，事情是明显的。因为在那些看见基督肉身的人中，有些相信了，有些把祂钉在十字架上；那些相信的人在祂被钉时也曾怀疑；除非他们在复活后触摸了那肉身，他们的信德不会恢复。那么，如果因为祂的肉身，圣子是可见的，这一点我们也承认，这是天主教的信仰；但如果在他们所说的祂取肉身之前，即祂降生成人之前，他们就大大受骗，严重错误了。因为那些可见的、有形的显现是通过受造物发生的，其中可能显示了一个形象：但绝不是显示了本体本身。请注意这个简单的证明，亲爱的各位弟兄。天主的智慧不能被眼睛看见。各位弟兄，如果基督是天主的智慧和天主的德能\BibleRef{格前 1:24}；如果基督是天主的圣言，并且人的言语不能被眼睛看见，那么天主的圣言能被这样看见吗？\par

19. 所以，要从你们心中驱除肉性的思想，好使你们真正在恩宠之下，属于\WorkTitle{新约}。因此，在\WorkTitle{新约}中应许了永生。阅读\WorkTitle{旧约}，就会看到同样的诫命曾加给一个仍然肉性的民族，如同加给我们一样。因为崇拜唯一天主也是加给我们的诫命。\BibleQuote{不可妄呼上主你的天主的名}也是加给我们的，这是第二条诫命。\BibleQuote{遵守安息日}更是加给我们的，因为命令要神性地遵守。因为犹太人奴役性地遵守安息日，将其用于奢侈和醉酒。她们的女人在那天在阳台上跳舞，还不如纺毛线好呢？愿天主禁止，弟兄们，我们称那为遵守安息日。基督徒神性地遵守安息日，戒避奴役性的工作。因为戒避奴役性的工作是什么意思？戒避罪恶。我们如何证明？请问主。\BibleQuote{凡犯罪的，就是罪恶的奴隶。}\BibleRef{若 8:34}因此，神性地遵守安息日是加给我们的。现在，所有这些诫命更是加给我们，并要遵守：\BibleQuote{不可杀人。不可奸淫。不可偷盗。不可作假见证。孝敬你的父亲和你的母亲。不可贪恋你邻人的财物。不可贪恋你邻人的妻子。}\BibleRef{出 20:3}难道所有这些不也是加给我们的吗？但问报酬是什么，你会发现那里说：\BibleQuote{使你的仇敌从你面前被赶出，使你领受天主向你祖先所应许的土地。}\BibleRef{肋 26:1}因为他们不能理解看不见的事，就被可见的事所束缚。为什么被束缚？免得他们完全灭亡，陷入偶像崇拜。因为，我的弟兄们，正如我们所读的，他们忘记了天主在他们眼前所行的伟大奇迹。海被分开；波涛中开了一条路；追赶的敌人被他们经过的同一波涛淹没：\BibleRef{出 14:21}然而当\Person{梅瑟}，天主的人，从他们眼前离开后，他们就要求一个偶像，说：\BibleQuote{给我们造神像，在我们前面引路，因为这个人抛弃了我们。}他们全部的希望在人身上，而不是在天主身上。看，那人死了；难道那位从\Place{埃及}地拯救他们的天主死了吗？当他们为自己造了一头牛犊的像，向它朝拜说：\BibleQuote{以色列，这就是你的\OriginalTerm{神}{gods}，领你出\Place{埃及}地的神。}\BibleRef{出 32:1}这么快就忘记了如此明显的恩宠！除了肉性的应许，还能用什么方式束缚这样的人民呢？\par

20. 十诫中所命令的同样的事情也是我们被命令要遵守的；但同样的应许并未加给我们。什么应许给我们？永生。\BibleQuote{永生就是：认识你，唯一的真天主，和你所派遣的耶稣基督。}\BibleRef{若 17:3}认识天主是应许的：即恩宠上加恩宠。弟兄们，我们现在相信，我们尚未看见；因为信德的报酬将是看见我们所相信的。先知们知道这事，但在他来临之前是隐藏的。因为有一位叹息的恋人，在圣咏中说：\BibleQuote{我向上主求得一事，我仍要追求这事。}你问他寻求什么？或许他肉性地寻求流奶流蜜之地，虽然这是应当神性地寻求和渴望的；或者寻求仇敌的屈服，敌人的死亡，或此世的权势和财富。因为他爱火燃烧，深深叹息，灼热而渴慕。让我们看看他渴望什么：\BibleQuote{我向上主求得一事，我仍要追求这事。}他追求的是什么呢？他说：\BibleQuote{使我一生的岁月，常住在上主的殿里。}假设你住在上主的殿里，你的喜乐从何而来？他说：\BibleQuote{使我得以瞻仰上主的荣美。}\par

21. 我的弟兄们，你们为什么呼喊？为什么欢乐？为什么爱？除非这爱的一点火花在那里？我问你们，你们渴望什么？难道能用眼睛看见？能用手触摸吗？难道是某种悦目的美丽吗？殉道者们不是被强烈地爱着吗？当我们纪念他们时，我们不是因爱而燃烧吗？弟兄们，我们在他们身上爱什么呢？是被野兽撕裂的肢体吗？如果问肉体的眼睛，还有什么更令人反感？如果问心灵的眼睛，还有什么更美丽？在你们眼中，一个非常俊美的小偷看起来如何？你们的眼睛是多么震惊！肉体的眼睛震惊吗？如果你们质问它们，没有任何东西比那个身体更匀称、更优美；四肢的对称和色彩的美吸引眼睛；然而，当你们听说他是个小偷时，你们的心灵就远离他。另一方面，你们看见一个弯腰的老人，拄着拐杖，几乎不能移动，满脸皱纹。你们听说他是正直的，你们就爱他，拥抱他。弟兄们，这就是应许给我们的赏报：爱这样的赏报，渴慕这样的王国，向往这样的国度，如果你们希望到达我们的主所带来的，就是恩宠和真理。但是如果你们贪求来自天主的身体赏报，你们仍处于法律之下，因此你们将不能满全法律。因为当你们看见那些暂时的事物赐给得罪天主的人时，你们的脚步就动摇，你们对自己说：看，我敬拜天主，每天跑向教会，我的膝盖因祈祷而磨损，却常常生病；有些人犯谋杀、抢劫之罪，却欢乐富足，他们很好。难道你们向天主寻求的就是这些东西吗？当然你们属于恩宠。所以，如果天主赐给你们恩宠，因为祂白白地赐予，你们也该白白地爱。不要为了赏报而爱天主；让祂自己作你们的赏报。让你们的灵魂说：\BibleQuote{我向上主祈求一事，我仍恳切追求：愿我一生住在上主的殿里，欣赏上主的美善。}不要害怕你们的享受会因饱足而消逝：那美的享受将如此：它永远呈现在你们面前，而你们永远不会满足；事实上，你们将永远满足，却又永不满足。因为如果我说你们将不会满足，那就意味着饥饿；如果我说你们将满足，我怕饱足：在既无饱足又无饥饿之处，我不知道该说什么；但是天主有祂所能显示给那些不知如何表达、却相信将要领受的人的东西。\par

\SourceRecord{Translated by John Gibb.；Nicene and Post-Nicene Fathers, First Series；Vol. 7.；Edited by Philip Schaff.；Buffalo, NY: Christian Literature Publishing Co.,；1888.；<http://www.newadvent.org/fathers/1701003.htm>.}

% structure: p0001 source=h1 target=section reason=page-title

\section{讲道集4（\BibleRef{若 1:19-33}）}

圣善的弟兄们，你们已屡次听见，也深知洗者若翰，他比妇女所生的更伟大，又在对主的承认中更为谦卑，因而得了作新郎之友的恩宠；他热切于新郎，并非为自己，也不是寻求自己的光荣，而是那审判者的光荣，他作为先驱走在他前面。因此，先前的先知们得以预言基督，但此人得以用手指指出祂。因为基督在降来之前，未被相信先知的人所认识，即使祂已降临，他们仍不认识祂。因为祂初次是谦卑地、隐蔽地降临；越是谦卑，越是隐蔽；但人民，由于骄傲而轻视天主的谦卑，竟将救主钉死，使自己成为祂的定罪者。\par

2. 但那位初次因谦卑而隐蔽降临的，难道不会因被高举而显现地再来吗？你们刚听了圣咏：\BibleQuote{天主必显现而来，我们的天主绝不缄默。}祂曾缄默，为受审判；祂将不缄默，当祂开始审判时。经上不会说\BibleQuote{祂必显现而来}，除非祂初次曾隐蔽而来；也不会说\BibleQuote{祂绝不缄默}，除非祂先前曾缄默。祂如何缄默呢？请问依撒意亚：\BibleQuote{祂被牵去如羊待宰，又如羔羊在剪毛者前无声，祂不开口。}\BibleRef{依 53:7}\BibleQuote{但祂必显现而来，绝不缄默。}如何\Quote{显现}？\BibleQuote{火要在祂面前行，周围有强烈的暴风。}那暴风将刮走禾场上所有糠秕——禾场如今正在打谷；那火将焚烧暴风所刮走的一切。但现在祂缄默：在审判中缄默，但在诫命中不缄默。因为如果基督缄默，这些福音有何用？宗徒之声有何用？圣咏之歌有何用？先知之言有何用？在这所有之中，基督并不缄默。但现在祂在报复上缄默，但在警告上不缄默。然而祂将光荣地降临以施行报复，并使祂自己向所有不信祂的人显现。但现在，因为祂在临在时隐藏，故应受轻视。因为若不受轻视，祂就不会被钉死；若不被钉死，祂就不会流祂的血——那赎买我们的代价。但为给我们付出代价，祂被钉死；为被钉死，祂被轻视；为被轻视，祂以谦卑显现。\par

3. 但因祂仿佛在黑夜中、在可朽的肉体内显现，祂为自己点了一盏灯，好使人看见祂。那盏灯就是若翰，\BibleRef{若 5:35}关于他你们近来听了许多事；而这段福音经文包含了若翰的话：首先是，且是最主要的，他承认自己不是基督。但若翰的卓越如此伟大，以致人们可能相信他是基督；在此他显出了谦卑的凭据，即他说他不是，尽管他可能被相信是基督。因此，\BibleQuote{这是若翰的见证：当犹太人从耶路撒冷派遣司祭和肋未人到他那里问他：\InnerQuote{你是谁？}}但是他们若不被他那敢于施洗的权威之卓越所感动，就不会派遣人来。\BibleQuote{他承认了，并不否认。}他承认了什么？\BibleQuote{他承认说：\InnerQuote{我不是基督。}}\par

4. \BibleQuote{他们又问他说：那么，你是谁？你是厄里亚吗？}因为他们知道厄里亚要在基督之前到来。因为基督的名号没有犹太人是不知道的。他们不认为他就是基督，但他们并非认为基督不会来临。当他们盼望祂降临时，他们却因祂的临在而绊倒，如同跌在一块矮石上。因为祂当时还是一块小石头，但已非人手从山上凿出，正如达尼尔先知所说，他看见一块石头非人手从山上凿出。但随后呢？\BibleQuote{那石头，}他说\BibleQuote{渐渐长大，成了一座大山，充满了整个大地。}\BibleRef{达 2:34}所以，我亲爱的弟兄们，请注意我所说的：基督在犹太人面前，已从山上凿出。先知希望藉着山理解犹太王国。但犹太王国并没有充满整个大地。石头是从那里凿出的，因为主是从那里为降临人间而生。为何非人手？因为童贞女产生基督没有人的合作。那时，那块非人手凿出的石头已在犹太人眼前，但它是谦卑的。并非无理由：因为那石头尚未长大而充满全地——祂将自己的王国，即教会，充满了整个大地。那时因为它尚未长大，他们便绊倒在它上面，如绊倒在石头上；在他们身上发生了经上所记的：\BibleQuote{凡跌在这石头上的，必被摔碎；这石头落在谁身上，必要把他压碎。}\BibleRef{路 20:18}起初他们跌倒在谦卑的祂身上；作崇高者的祂将降临在他们身上；但为在祂高举来临时把他们压碎，祂先在他们身上将他们摔碎。他们绊倒在祂身上，被摔碎了；不是被压碎，而是被摔碎；祂将高举降临，并要压碎他们。但犹太人应得宽恕，因为他们绊倒在一块尚未长大的石头上。那么，那些绊倒在大山本身的人，又是何等样人呢？你们已知道我说的是谁。那否认分散在全世界的教会的人，不是绊倒在谦卑的石头上，而是绊倒在大山本身；因为这石头长大成了山。瞎眼的犹太人没有看见谦卑的石头，但看不见山，是何等大的瞎眼！\par

5. 当时他们看见\Person{基督}是卑微的，并不认识祂。祂被一盏灯指给他们。因为首先，那在妇人所生中没有比他更大的，说：\Quote{我不是\Person{基督}。}有人问他说：\Quote{你是\Person{厄里亚}吗？}他回答说：\Quote{我不是。}因为\Person{基督}在祂之前派遣\Person{厄里亚}：而他说：\Quote{我不是，}给我们留下了一个问题。因为恐怕人们由于理解不足，以为\Person{若翰}与\Person{基督}所说的相矛盾。因为在某处，主\Person{耶稣基督}在福音中论及自己说了某些话，祂的门徒回答祂说：\Quote{那么，经师们（即精通法律的人）怎么说\Person{厄里亚}必须先来呢？}主回答说：\Quote{\Person{厄里亚}已经来了，他们任意待了他；并且，如果你们愿意知道，\Person{洗者若翰}就是那人。}主\Person{耶稣基督}说：\Quote{\Person{厄里亚}已经来了，\Person{洗者若翰}就是那人。}但\Person{若翰}被问时，承认自己不是\Person{厄里亚}，正如他承认自己不是\Person{基督}一样。他承认自己不是\Person{基督}是真的，同样他承认自己不是\Person{厄里亚}也是真的。那么，我们如何比较报信者的话与法官的话呢？绝不要以为报信者说假话；因为他所讲的乃是从法官那里听到的。那么，他为什么说\Quote{我不是\Person{厄里亚}，}而主却说\Quote{他是\Person{厄里亚}}呢？因为主\Person{耶稣基督}愿意在他身上预示自己的来临，并说\Person{若翰}具有\Person{厄里亚}的精神。\Person{若翰}之于第一次来临，正如\Person{厄里亚}之于第二次来临。正如法官有两次来临，同样有两个报信者。法官确实是同一位，但报信者有两位，却没有两位法官。第一次，法官必须来受审判。祂在祂之前派遣了第一位报信者；祂称他为\Person{厄里亚}，因为\Person{厄里亚}在第二次来临中将是\Person{若翰}在第一次来临中所是的。\par

6. 请留意，亲爱的弟兄们，我所说的是多么真实。当\Person{若翰}受孕，或更确切地说，当他出生时，\Person{圣神}预言这事将应验在他身上：祂说：\Quote{他将成为至高者的前驱，具有\Person{厄里亚}的精神和能力。}见：\BibleRef{路 1:17} \Quote{具有\Person{厄里亚}的精神和能力}是什么意思？是指具有那与\Person{厄里亚}同样的\Person{圣神}。那么为何说是代替\Person{厄里亚}？因为\Person{厄里亚}之于第二次来临，正是\Person{若翰}之于第一次来临。因此，\Person{若翰}按字面意思正确地回答了。因为主用比喻的话说：\Quote{\Person{厄里亚}，他就是\Person{若翰}；}但他，如我所说，按字面意思说：\Quote{我不是\Person{厄里亚}。}\Person{若翰}没有说谎，主也没有说谎；报信者的话和法官的话都没有虚假，只要你们理解。但谁能理解呢？那效法报信者的谦卑，并承认法官的崇高的人。因为没有什么比报信者更卑微。我的弟兄们，\Person{若翰}的功德无过于这谦卑，因为他虽能欺骗人，被人认为是\Person{基督}，被当作\Person{基督}接受（他的恩宠和卓越如此之大），然而他公开承认说：\Quote{我不是\Person{基督}。}\Quote{你是\Person{厄里亚}吗？}如果他回答说我是\Person{厄里亚}，那就好像\Person{基督}已经第二次来临审判，而不是第一次来受审判。仿佛是说：\Person{厄里亚}还要来，他说：\Quote{我不是，}\Quote{\Person{厄里亚}。}但你们要留心意\Person{若翰}所前来的那卑微者，免得你们感受到\Person{厄里亚}所前来的那崇高者。因为主也是如此完成了这句话：\Quote{\Person{洗者若翰}就是那将要来的。}他作为\Person{厄里亚}将来临的预像而来。那么，\Person{厄里亚}将亲自作为\Person{厄里亚}而来，现在在预像中他是\Person{若翰}。现在\Person{若翰}亲自是\Person{若翰}，在预像中是\Person{厄里亚}。两位报信者彼此交换了他们的预像，而保留了他们各自的本人；但法官是同一位主，无论由这位报信者还是那位报信者作前驱。\par

7. \Quote{他们便问他说：那么，你是谁？你是\Person{厄里亚}吗？他说：不是。他们又对他说：你是先知吗？他回答说：不是！于是他们对他说：你是谁？好叫我们给那派遣我们来的人一个答复。你说你自己是什么？他说：我是旷野中呼喊者的声音。}见：\BibleRef{依 40:3} 这是\Person{依撒意亚}所说的。这预言在\Person{若翰}身上应验了，\Quote{我是旷野中呼喊者的声音。}呼喊什么？\Quote{你们要预备主的道路，修直我们天主的途径。}你们难道不觉得一个报信者会呼喊：\Quote{走开，让路。}报信者喊的是\Quote{走开，}\Person{若翰}却说\Quote{来吧。}报信者使人远离法官；\Person{若翰}却召唤人到法官面前。是的，确实，\Person{若翰}召唤人到那卑微者面前，免得他们经历祂作为被高举的法官时将是什么样子。\Quote{我是旷野中呼喊者的声音：预备主的道路，正如先知\Person{依撒意亚}所说的。}他没有说：我是\Person{若翰}，我是\Person{厄里亚}，我是先知。他说的什么？他被称为：\Quote{旷野中呼喊者的声音，预备主的道路：我是这个预言本身。}\par

8. \BibleQuote{而他们之中被派来的，是法利塞人，}就是犹太人的首要人物；\BibleQuote{他们问祂说：\InnerQuote{你既不是基督，又不是厄里亚，也不是先知，那么你为什么施洗呢？}}仿佛他们觉得施洗是僭越，好像要探究：你以什么身份施洗？我们问你是否是基督；你说不是。我们问你是否是祂的先驱，因为我们知道在基督来临之前，厄里亚会来；你回答说不是。我们问你，是否你是许久之前来的某个传报者，即先知，并且接受了那权柄，你说你不是先知。若翰并不是先知；他比先知更大。主曾为他作证说：\BibleQuote{你们到旷野里去看什么？被风吹动的芦苇吗？}当然是指他并非被风吹动的；因为若翰不是那样被风摇动的人；凡是被风吹动的，是受每一阵诱惑的风吹动。\BibleQuote{但你们出去看什么？穿细软衣服的人吗？}因为若翰穿着粗糙的衣服；就是说，他的外衣是骆驼毛的。\BibleQuote{看，穿细软衣服的人是在王宫里。}那么你们出去不是要看穿细软衣服的人。\BibleQuote{但你们出去看什么？一位先知吗？是的，我告诉你们，一位比先知更大的就在这里；} \BibleRef{玛 11:7}因为先知们早就预言了基督，而若翰指出祂就在眼前。\par

9. \BibleQuote{你既不是基督，又不是厄里亚，也不是先知，那么你为什么施洗呢？若翰回答他们说：我用水施洗，但有一位站在你们中间，是你们不认识的。}因为，很真实地，祂没有显现，因为祂谦卑，因此灯被点亮了。请看若翰如何退让，他本可能被当作别的人。\BibleQuote{祂就是在我以后来的，成了在我以前的，}（即，如我们已经说过的，\Quote{在我以上}），\BibleQuote{我连解祂的鞋带也不配。}他多么谦卑自下！因此他被高举；\BibleQuote{因为凡自谦的，必被高举。}\BibleRef{路 14:11}因此，圣洁的弟兄们，你们应当注意，如果若翰如此谦卑，以至于说：\BibleQuote{我不配解祂的鞋带，}那么那些说：\Quote{我们施洗；我们所给的是我们的，我们的是圣洁的}人，需要如何谦卑呢？他说：不是我，而是祂；他们却说：我们。若翰不配解祂的鞋带；如果他说他配，那仍然是多么谦卑！如果他说他配，并且这样说：\Quote{祂在我以后来，成了在我以前的，我只配解祂的鞋带，}他就已经大大谦卑了。但当他承认连这事也不配时，他真是充满圣神，他以仆人之态认主，并得以由仆人成为朋友。\par

10. \BibleQuote{这些事发生在约旦河对岸的伯达尼，那里若翰施洗。次日，若翰看见耶稣向他走来，就说：看，天主的羔羊；看，除免世界罪过者！}谁也不要妄自尊大，说是自己除免世界的罪过。现在请看若翰所指的骄傲之人。那时异端尚未出生，却已被指出；那时他就在河中向他们呼喊，现在他由福音向他们呼喊。耶稣来了，他说什么？\BibleQuote{看，天主的羔羊！}如果纯洁就是羔羊，那么若翰也是羔羊，难道他不是纯洁的吗？但谁是纯洁的呢？纯洁到什么程度？所有人都从那枝芽而来，达味歌唱时甚至哀叹：\BibleQuote{看，我在罪孽中成形，我母亲在罪恶中怀了我。}唯独那来临的羔羊，不是这样。因为祂不是在罪孽中受孕，因为祂不是由可朽坏者受孕；祂的母亲也没有在罪恶中怀祂，因为童贞女因信德怀了祂，因信德生了祂。因此，\BibleQuote{看，天主的羔羊。}祂不是亚当的枝条：祂只从亚当取了肉躯，却没有承当亚当的罪。祂没有从我们的团块中取罪，祂除免了我们的罪。\BibleQuote{看，天主的羔羊，除免世界罪过者！}\par

11. 你们知道有些人有时说：\Quote{我们，圣洁的人，除去人的罪；因为如果施洗者自己不圣洁，他自己满身是罪，怎能除去别人的罪呢？}为反驳这些议论，我们不要讲自己的话，让我们读若翰说的：\BibleQuote{看，天主的羔羊；看，除免世界罪过者！}不要让人对人的狂妄信赖；不要让麻雀飞往山上，而要信靠主；即使它举目向山，看救助从哪里来，也当明白它的救助来自创造天地的主。若翰是如此卓越，以至于人对他说：\BibleQuote{你是基督吗？}他说：不是。\BibleQuote{你是厄里亚吗？}他说：不是。\BibleQuote{你是先知吗？}他说：不是。那么你为什么施洗？\BibleQuote{看，天主的羔羊；看，除免世界罪过者！这就是我对你们说过的那位：在我以后来的一位，成了在我以前的，因为祂在我以前。}\BibleQuote{在我以后来，}因为祂诞生较晚；\BibleQuote{成了在我以前的，}因为祂在我以上；\BibleQuote{祂在我以前，}因为\BibleQuote{在起初已有圣言，圣言与天主同在，圣言就是天主。}\par

12. \Quote{我原不认识祂，}他说；\Quote{但为使祂显示给以色列，为此我来以水施洗。} \BibleQuote{若翰又作证说：\InnerQuote{我看见圣神彷佛鸽子从天降下，停在祂身上。我原本不认识祂，但那位派遣我来以水施洗的告诉我说：\InnerQuote{你看见圣神降下并停在谁身上，谁就是那以圣神施洗的。}我看见了，便作证：这一位就是天主子。}}亲爱的诸位，请留意片刻。若翰何时认识了基督？因为他被派遣以水施洗。他们问：为什么？他说：为使祂显示给以色列。若翰的洗礼有什么益处？我的弟兄们，如果它在任何方面有益处，它现在就会保留，人们就会受若翰的洗礼，然后再来到基督的洗礼。但他说了什么？\Quote{为使祂显示给以色列}——就是显示给以色列本身，给以色列人民，使基督显示给他们——因此他以水施洗。若翰接受了洗礼的职事，好藉着悔改的水为预备主的道路，他自己并不是主；但哪里主已被认识，为祂预备道路就是多余的，因为对那些认识祂的人，祂自己成了道路；因此若翰的洗礼不长久。但主是如何被指出的呢？谦卑地，使若翰能领受一个洗礼，而主自己也应受那个洗礼。\par

13. 主需要受洗吗？我立刻回答任何问这问题的人：主需要诞生吗？主需要被钉十字架吗？主需要死亡吗？主需要被埋葬吗？如果祂为我们承担了如此大的屈辱，难道祂不能也接受洗礼吗？祂接受仆人的洗礼有什么益处呢？是为了使你不轻看接受主的洗礼。亲爱的弟兄们，请注意。在教会中将要出现一些更高恩宠的望教者。有时你会看到一个望教者持守贞洁，告别世界，放弃所有财产，分施给穷人；虽然他只是一个望教者，但在救恩的教导上，或许比许多信友更好。我们应当担心这样的人会对自己说，关于那使罪得赦的圣洗圣事：\Quote{我还要领受什么？看，我比这信友更好，比那个（他心中想着那些已婚的，或可能无知的，或保留财产的信友，而他却将财产给了穷人）更好，}他认为自己比那些已受洗的人更好，因此不屑于领受洗礼，说：\Quote{难道我要领受这个人所有的，而心里却想着那些我所轻视的人？}并且，他视之为一种屈辱去领受那些比他低劣的人所领受的，因为他似乎已经比他们更好了；然而，他所有的罪仍然在他身上，若不来到那赦免所有罪的救恩的洗礼，他不能带着他所有的优点进入天国。但是主为了邀请这样的优秀者到祂的洗礼，使罪得以赦免，祂自己来到了祂仆人的洗礼；虽然祂没有罪需要赦免，也没有任何需要洗净的东西，祂却接受了仆人的洗礼；藉此，祂对那骄傲自大、轻视、或许不屑于同无知的人一起领受那带来救恩的洗礼的儿子说话，对他说：\Quote{你如何自大？你如何自夸？你的优秀有多大？你的恩宠有多大？难道能比我的更大吗？若我来到仆人那里，你竟不屑于来到主这里吗？若我接受了仆人的洗礼，你竟不屑于受主施洗吗？}\par

14. 但是，我的弟兄们，为使你们知道，主来到若翰这里，并非出于任何罪的锁链的必然性——如同其他福音传述者所说，当主来到他那里受洗时，若翰自己说：\Quote{祢到我这里来吗？我当受祢的洗。}\BibleRef{玛 3:14} 主如何回答他？\Quote{你暂且容许吧：让我们完成全部义德？}这话是什么意思，\Quote{让我们完成全部义德}？我来是为人类死，难道我不该为人类受洗吗？\Quote{让我们完成全部义德}是什么意思？就是让全部谦逊得以完成。那么怎样？难道祂不愿意从一个好仆人手中接受圣洗，却接受了恶仆人的苦难吗？请留心。主受了圣洗——如果若翰为此目的而施洗，好藉他的圣洗使主彰显祂的谦逊，那么是否不该有任何人受若翰的圣洗呢？但有许多人受了若翰的圣洗。当主受了若翰的圣洗时，若翰的圣洗就停止了。若翰立即被投入狱中。此后我们再也找不到任何人受那圣洗。那么，如果若翰来施洗是为了这个目的，即让我们看见主的谦逊，好叫我们不屑于从主那里接受那主曾从仆人接受过的圣洗，若翰是否应当只为主一人施洗呢？但是，如果若翰只为主一人施洗，有些人就会认为若翰的圣洗比基督的圣洗更圣洁，好像只有基督才配受若翰的圣洗，而人类却只配受基督的圣洗。请注意，亲爱的弟兄们。我们受的是基督的圣洗，不仅是我们，还有全世界，而且这要持续到终局。我们当中有谁能在任何方面与基督相比？连祂的鞋带，若翰都自称不配解。那么，如果那位如此卓越的人、那位又是天主的基督，是唯一受若翰圣洗的人，人们会说什么呢？\Quote{若翰的圣洗是何等伟大！他的圣洗是伟大的，是不可言喻的圣事；看，唯独基督配受若翰的圣洗。}这样，仆人的圣洗就显得比主的圣洗更大了。其他人也受了若翰的圣洗，好使若翰的圣洗不致显得比基督的圣洗更好；但是主也受了圣洗，好叫其他仆人因主接受了仆人的圣洗，就不致不屑于接受主的圣洗。为此目的，若翰被派遣了。\par

15. 但是，若翰认识基督吗，还是不认识祂？如果他不认识祂，那么当基督来到约旦河时，他为什么说：\Quote{我当受祢的洗}？这就是说，我知道你是谁。那么，如果他早已认识祂，他肯定是看见鸽子降下时认识祂的。显然，鸽子是在主从圣洗的水中上来之后才降在主身上的。\Quote{主受了圣洗，从水里上来，天开了，他看见鸽子降在祂身上。}那么，如果鸽子是在圣洗之后降下，而主受圣洗之前若翰就对祂说：\Quote{祢到我这里来吗？我当受祢的洗}；也就是说，在他认识祂之前，他就对祂说了\Quote{祢到我这里来吗？我当受祢的洗}。那么，他怎么说：\Quote{我也不认识祂；但是那差我来用水施洗的，对我说：你看见圣神降下，如同鸽子，停在祂身上，那一位就是要以圣神施洗的}？这不是一个小问题，我的弟兄们。如果你们看到了这个问题，你们就看到了不小的事；剩下的就是主要赐予我们解答。不过，我要说，如果你们看到了这个问题，那就不是小事。看，若翰摆在你们眼前，站在约旦河边。看，洗者若翰。看，主来了，尚未受圣洗，还没有受圣洗。听若翰的声音：\Quote{祢到我这里来吗？我当受祢的洗。}看，他已经认识主，并愿意受祂的圣洗。主受了圣洗，从水里上来；天开了，圣神降下；那时若翰才认识祂。那么，如果他是第一次认识祂，为什么他先前说：\Quote{我当受祢的洗}？但如果他并不是那时才第一次认识祂，因为他已经认识祂，那么他所说的\Quote{我不认识祂；但那差我来用水施洗的，对我说：你看见圣神降下，如同鸽子，停在祂身上，那一位就是要以圣神施洗的}又是什么意思？\par

16. 我的弟兄们，这个问题若在今天解决，会压迫你们，我不怀疑，因为我已经讲了许多话。但你们要知道，这个问题的性质，仅凭它就能扑灭多纳徒派。我说了这许多，我的诸位爱者，是为引起你们的注意，如同我惯常所做的；也是为使你们为我们祈祷，使主恩赐我们能宣讲适宜的话，并使你们堪当领受适宜的话。今天暂且请你们推迟这个问题。但同时，我简短地说这一点，直到我给出更完整的解答：要和平地探求，不争吵，不争辩，不口角，不怀敌意；你们自己寻求，也询问别人，并且说：\Quote{今天我们主教向我们提出了这个问题，他将在日后解决它，如果主愿意的话。}但无论它是否解决，都当认为我已经提出了在我看来重要的东西；因为它确实看起来相当重要。若翰说：\Quote{我需要受祢的洗}，似乎他认识基督。因为如果他不认识他愿意受其洗的那一位，他说话就草率了，他说：\Quote{我需要受祢的洗。}因此他认识他。若他认识他，那么这句话是什么意思：\Quote{我不认识他，但派遣我来以水施洗的那位对我说：你看见圣神降下并停留在谁身上，就像鸽子一样，那一位就是将以圣神施洗的}？我们该说什么呢？说我们不知道鸽子何时来的吗？恐怕他们以此托辞，就让我们阅读其他论述此事更为清晰的福音作者，我们最明确地发现，鸽子正是在主从水中上来的时候降下的。他在受洗之时天开了，他看见圣神降下。若是在他已受洗之后若翰才认识他，那么他对前来受洗的那一位说\Quote{我需要受祢的洗}是什么意思呢？期间你们自己思考，彼此商议，互相讨论。愿上主我们的天主恩赐，在你们从我这里听到之前，先向你们中的某些人启示这解释。尽管如此，弟兄们，你们要知道，通过解决这个问题，多纳徒派的指控，如果他们还有羞耻之心，就会哑口无言，他们关于圣洗的恩宠的嘴将被封住，关于这件事他们散布迷雾迷惑无知者，并为飞鸟设下网罗。\par

\SourceRecord{Translated by John Gibb.；Nicene and Post-Nicene Fathers, First Series；Vol. 7.；Edited by Philip Schaff.；Buffalo, NY: Christian Literature Publishing Co.,；1888.；<http://www.newadvent.org/fathers/1701004.htm>.}

% structure: p0001 source=h1 target=section reason=page-title

\section{讲道集 5（若望福音 1:33）}

我們已按主的意願，抵達了我們所應許的日子。祂也將賜予我們，使我們能抵達應許的實現。因為那時，我們所說的話，若對我們和你們有益，便是出於祂；但出於人的話是虛假的，正如我們的主耶穌基督親自所說的：\Quote{說謊言的，是出於自己。}\BibleRef{若 8:44}沒有人擁有任何出於自己的東西，除了虛假和罪。但若人擁有任何真理和正義，那是出於那源泉，我們在這曠野中應當渴望那源泉，以便如同被它的幾滴露珠沾濕，在此旅途中得著安慰，以免我們在路上疲乏，而能到達祂的安息和滿足的豐盈。所以，如果\Quote{說謊言的，是出於自己}，那麼說真理的，是出於天主。若翰是真實的，基督就是真理；若翰是真實的，但每個真實的人都是從真理而來成為真實的。所以，若若翰是真實的，而人除了從真理而來不能成為真實，那麼他是從誰成為真實的呢？除非是從那位說\Quote{我是真理}的那一位？\BibleRef{若 14:6}因此，真理不能反對真實的人，真實的人也不能反對真理。真理派遣了真實的人，他是真實的，因為被真理所派遣。如果那是派遣若翰的真理，那麼就是基督派遣了他。但基督與聖父所做的事，聖父也做；聖父與基督所做的事，基督也做。聖父不做任何與聖子分離的事，聖子也不做任何與聖父分離的事：不可分離的愛，不可分離的合一；不可分離的尊威，不可分離的能力，按照祂親口所說的話：\Quote{我與聖父原是一體。}\BibleRef{若 10:30}那麼，是誰派遣了若翰？如果我們說是聖父，我們說得對；如果我們說是聖子，我們也說得對；但更清楚地說，我們說是聖父和聖子。但聖父和聖子所派遣的，是一位天主所派遣的；因為聖子說：\Quote{我與聖父原是一體。}那麼，祂如何不認識派遣祂的那一位呢？因為他說：\Quote{我不認識祂；但那位派遣我用水施洗的，對我說：}我問若翰：\Quote{誰派遣你用水施洗？祂對你說了什麼？}\Quote{你看見聖神如同鴿子降下，停在祂身上的那一位，就是那以聖神施洗的。}若翰啊，派遣你的那一位對你說的，就是這個嗎？顯然就是這個；那麼，是誰派遣了你？也許是聖父。真實的天主是聖父，真理是天主聖子：如果聖父沒有聖子派遣了你，那麼天主沒有真理派遣了你；但你如果是真實的，因為你說真理，並出於真理說話，那麼聖父並非沒有聖子而派遣了你，而是聖父和聖子一同派遣了你。所以，如果聖子與聖父一同派遣了你，你如何不認識派遣你的那一位呢？你在真理中所看見的那一位，親自派遣了你，好使祂能在肉身中被認出，並說：\Quote{你看見聖神如同鴿子降下，停在祂身上的那一位，就是那以聖神施洗的。}\par

2. 若翰聽見這話，是為認識他所不認識的那一位，還是更充分地認識他已認識的那一位？因為如果他完全不認識祂，當祂來到河邊受洗時，他就不會對祂說：\Quote{我需要受你的洗，你反倒到我這裡來嗎？}\BibleRef{玛 3:14}所以他知道祂。但鴿子是在何時降下的？是在主受洗後，從水中上來的時候。但如果派遣祂的那一位說：\Quote{你看見聖神如同鴿子降下，停在祂身上的那一位，就是那以聖神施洗的，}而他不認識祂，但當鴿子降下時，他學會認識祂，而鴿子降下的時間是主從水中上來的時候；但若翰在主來到他那裡到水邊時已經認識主：這就向我們顯明，若翰在某種程度上認識，在某種程度上起初不認識主。除非我們如此理解，他就是說謊者。他承認主並說：\Quote{你到我這裡來受洗嗎？}以及\Quote{我需要受你的洗}，這怎能是真實的呢？他這樣說時是真實的嗎？而當他說：\Quote{我不認識祂；但那位派遣我用水施洗的，對我說：你看見聖神如同鴿子降下，停在祂身上的那一位，就是那以聖神施洗的}時，他又怎能是真實的呢？主藉著鴿子被顯明，不是給那不知道祂的人，而是給那在某種程度上認識祂、在某種程度上不認識祂的人。我們要發現的是，若翰關於祂所不知道的、並藉著鴿子所學到的是什麼。\par

3. 為什麼若翰被派遣來施洗？我記得，親愛的，我已經根據我的能力向你們解釋過了。因為如果若翰的洗禮對於我們的救恩是必要的，那麼它甚至現在也應當被使用。因為我們不能認為現在人們不獲救，或現在獲救的人更少，或那時有一種救恩，現在有另一種。如果基督被改變了，救恩也被改變了；如果救恩在基督內，而基督自身是同一的，那麼對我們來說就是同樣的救恩。但為什麼若翰被派遣來施洗？因為基督必須受洗。為什麼基督必須受洗？為什麼基督必須誕生？為什麼基督必須被釘十字架？因為如果祂來是為了指出謙卑的道路，並使祂自己成為謙卑的道路；那麼在一切事上，謙卑都必須由祂來完成。祂藉此屈尊賦予祂自己的洗禮權威，好使祂的僕人知道他們應以何種熱忱奔赴主的洗禮，因為祂自己沒有拒絕接受僕人的洗禮。這個恩惠賜給了若翰，使那洗禮被稱為他的洗禮。\par

4. 请留意这事，运用你的辨别力，并认识它，亲爱的诸位。若翰所领受的圣洗被称为若翰的圣洗：唯独他领受了这样的恩赐。在他之前没有一位义人，在他之后也没有一个人，如此领受了一个圣洗以被称为他自己的圣洗。他确实领受了它，因为凭自己他什么也不能做：因为若有人凭自己说话，他就是凭自己说谎言。他除了从主耶稣基督那里领受，还能从哪里领受呢？从他那里他领受了施洗的权能，之后他用这权能为人施洗。不要惊奇；因为基督对待若翰，如同对待他自己的母亲一样。因为论到基督，经上说：\BibleQuote{万物是借着他造的。}\BibleRef{若 1:3} 如果万物是借着他造的，那么玛利亚也是由他造的，基督后来由玛利亚出生。请留意，亲爱的诸位；正如他创造了玛利亚，并由玛利亚受造，同样，他赐予了若翰的圣洗，并受若翰的洗。\par

5. 因此，他领受若翰的圣洗正是为了这个目的：他领受了来自卑微者的卑微之物，从而劝勉卑微者领受那崇高之物。但是，为什么他不独自受若翰的洗呢？如果若翰——基督由他受洗——被差来正是为了预备主的道路，也就是为基督自己预备道路，那么为什么他不独自受洗呢？这一点我们已经解释过，但我们再次提及，因为它对当前的问题很有必要。如果我们的主耶稣基督独自受了若翰的圣洗——请牢牢记住我们所说的；不要让世界有力量从你们心中抹去天主圣神在那里写下的；不要让忧虑的荆棘有力量窒息正在你们心中播下的种子；因为我们为什么被迫重复同样的事？不正是因为我们不能确信你们心中的记忆吗？——如果主独自受了若翰的圣洗，就会有人这样认为：若翰的圣洗大于基督的圣洗。因为他们会说，那个圣洗是如此伟大，以致唯独基督配受它。因此，为使主给我们立下谦逊的榜样，为使我们可以获得圣洗的救恩，基督接受了那为他并非必要、但为了我们却是必要的事。再者，唯恐基督从若翰领受的受尊于基督的圣洗，其他人也被允许受若翰的洗。但是，对于那些受若翰之洗的人，那圣洗并不足够：因为他们还受了基督的圣洗；因为若翰的圣洗不是基督的圣洗。那些领受基督圣洗的人不再寻求若翰的圣洗；那些领受若翰圣洗的人却寻求基督的圣洗。因此，若翰的圣洗对基督是足够的。它怎能不充足呢，连它本身都不是必要的？对他而言，没有圣洗是必要的；但为劝勉我们领受他的圣洗，他领受了他仆人的圣洗。并且，唯恐仆人的圣洗被尊于主的圣洗，其他的同仆也受了这个仆人的圣洗。但是，那些受了那圣洗的同仆，理应也接受主的圣洗；而那些受了主圣洗的人，并不需要同仆的圣洗。\par

6. 因此，若翰领受了一个可正当地称为若翰圣洗的圣洗，但主耶稣基督不愿将自己的圣洗赐给任何人——不是说没有人应受主的圣洗，而是说主自己应始终施洗：这样做的目的是，主借仆人施洗；也就是说，那些被主的仆人施洗的人，是主在施洗，不是那些仆人。因为以仆人的身份施洗是一回事，以权能施洗是另一回事。因为圣洗的特性源自那藉其权能而赐予者，而非藉其服务而赐予者。正如若翰是怎样的，他的圣洗也是怎样的：一位义人的义洗；但这个人从主那里领受了恩宠，且如此大的恩宠，以致他堪当作为审判者的前驱，用手指指出他，并应验那预言的话：\BibleQuote{在旷野里有呼喊者的声音：预备主的道路。}\BibleRef{依 40:3} 正如主是怎样的，他的圣洗也是怎样的：因此，主的圣洗是神圣的，因为主就是天主。\par

7. 但是主耶稣基督，如果他愿意，本可以将权能赐予他的一位仆役，让那人仿佛代替他施行他自己的圣洗，并且从他自身转移施洗的权能，将其赋予他的一位仆役，并给予转移给仆役的圣洗与主施予时同样的权能。他之所以不愿这样做，是为了使受洗者的希望寄托于他们承认自己因之受洗的那一位。因此，他不愿仆役将自己的希望寄托于仆役。所以，当宗徒看见人们希望将希望寄托于他自身时，他喊道：\BibleQuote{保禄曾为你们被钉死在十字架上吗？或者你们曾因保禄之名受洗吗？} \BibleRef{格前 1:13} 保禄是以仆役的身份施洗，而不是以权能本身；但主是以权能施洗。请留意。他既能赐予仆役这权能，却不愿。因为如果他将这权能赐予仆役——即是说，那属于主的东西成为他们的——那么将有多少仆役就有多少圣洗；以至于我们称若翰的圣洗，也会称伯多禄的圣洗、保禄的圣洗、雅各伯的圣洗、多默的圣洗、玛窦的圣洗、巴尔多禄茂的圣洗；因为我们称那圣洗为若翰的圣洗。但或许有人反对说：向我们证明那圣洗被称为若翰的圣洗。我要从真理自身的言词来证明，当他问犹太人时：\BibleQuote{若翰的洗礼是从哪里来的？是从天上来的，还是从人来的？} \BibleRef{玛 21:25} 因此，为了避免如同从主领受施洗权能的仆役那样称呼许多圣洗，主将施洗的权能保留给自身，并将职务赐予他的仆役。仆役说他施洗；他说得对，正如宗徒所说：\BibleQuote{我也给斯特法纳一家施了洗} \BibleRef{格前 1:16} ，但他是作为仆役。因此，即使他是恶的，而他恰好有施洗的职务，人们不认识他，但天主认识他；天主将权能保留给自身，仍允许藉着他施行圣洗。\par

8. 但是若翰在主身上所不知道的。他知道他是主，他知道他应该受他的洗；他明认他是真理，而他自己是真的人，是由真理派遣的：这些他知道。但他所不知道的，是什么在他内呢？就是主将把他圣洗的权能保留给自身，不传或不转移给任何仆役；但无论善仆役以职务方式施洗，或恶仆役施洗，受洗者不应知道自己受洗，除非藉着那位将施洗权能保留给自身的主。为使你们知道，弟兄们，若翰在主身上所不知道的，他藉着鸽子学会了：因为他知道主；但主将施洗的权能保留给自身，不赐给任何仆役，他还不知道。关于此事他说：\Quote{我不认识他。} 为使你们知道他在那里学会了这事，请留意下文：\Quote{但那派遣我来用水施洗的，他告诉我说：你看见圣神降下，如同鸽子，并停在他身上，那一位就是。} 那一位是谁？是主吗？但他已经知道主。那么，假设若翰只说了这些：\Quote{我不认识他；但那派遣我来用水施洗的，他告诉我说——} 我们问：他说了什么？接着是：\Quote{你看见圣神降下，如同鸽子，并停在他身上。} 我不说接着的话。同时请留意：\Quote{你看见圣神降下，如同鸽子，并停在他身上，那一位就是。} 但那一句\Quote{他}是谁？那派遣我来的，藉着鸽子要教导我什么？就是他本人是主。我已经知道谁派遣了我；我已经知道我对他说过的那一位：\Quote{你到我这里来受洗吗？我更需要受你的洗。} 那么，我如此知道主了，以至于我愿由他受洗，而不愿他由我受洗；随后他对我说：\BibleQuote{你暂且容许吧，因为我们应当这样履行全部义德。} \BibleRef{玛 3:15} 我来是为受苦；难道我不是来受洗吗？\Quote{让全部义德得以履行，} 我的天主对我说。让全部义德得以履行；让我教导完全的谦卑。我知道在我未来子民中将有骄傲的人；我知道有些人那时将在某种恩宠上杰出，以致当他们看到普通人受洗时，因为他们认为自己更好——或是在节欲上，或是在施舍上，或是在教义上——或许不屑于领受那些比他们低下的人所领受的。我需要医治他们，使他们不致藐视来到主的圣洗，因为我曾来到仆役的圣洗。\par

9. 若翰已知道此事，他已认识主。那么，鸽子教导了什么？祂借着鸽子——即借这样来临的圣神——要教导什么？祂对若翰说：\BibleQuote{你看见圣神降下，如同鸽子，停在祂身上，那一位就是祂}，但祂是谁？是主？我知道。但你是否已知道这件事：那位有权施洗的主，不将权能赐给任何仆人，而是为自己保留，使所有由仆人施洗的人，不将圣洗归于仆人，而归于主？你是否已知道这件事？我并不知道；那么祂向我说了什么？\BibleQuote{你看见圣神降下，如同鸽子，停在祂身上，那一位就是用圣神施洗的}。祂没有说：\Quote{祂是主}；没有说：\Quote{祂是基督}；没有说：\Quote{祂是天主}；没有说：\Quote{祂是耶稣}；没有说：\Quote{祂是由童贞玛利亚所生，在你之后、在你之前的那个}。祂没有说这些，因为若翰已知道这些。但他不知道什么？不知道圣洗的这大权柄，主自己将拥有并为自己保留，无论是身在尘世、还是身体在天堂、而尊威亲临，免得保禄说：我的圣洗；免得伯多禄说：我的圣洗。因此请看，留意宗徒们的话。没有一位宗徒说过：我的圣洗。虽然宗徒们有一共同的福音，你却看见他们说：我的福音；你未见他们说：我的圣洗。\par

10. 因此，我的弟兄们，若翰学到了。让我们也学习若翰借着鸽子所学到的。因为鸽子教导若翰，并没有不教导教会，教会曾被告知：\BibleQuote{我的鸽子，只有一个}。\BibleRef{歌 6:8} 愿鸽子教导鸽子；愿鸽子知道若翰从鸽子所学到的。圣神以鸽子的形像降下。但若翰在鸽子中所学的这一点，为何他是在鸽子中学到的？因为他应当学习，而或许他不那么应当学习，却应当借着鸽子学习。我的弟兄们，关于鸽子，我该说什么？或口舌或心灵的能力何时足以照我所愿地讲论？也许我的愿望不及我讲论的本分；即使我能照我所愿地讲论，我岂能照我所应当地讲论？我宁愿听一位比我高明的人讲论这事，而不是由我向你们讲论。\par

11. 若翰学习认识他所认识的那一位；但他是在祂身上学习他所不知道的；至于他所知道的，他并不学习。他知道什么？主。他不知道什么？主的圣洗的权能不将从主归于任何人，而圣洗的职事显然会如此；权能不从主归于任何人，职事则归于善人和恶人。愿鸽子不因恶人的职事而退缩，却要注视主的权能。在主是良善的情况下，恶仆能伤害你什么？若法官心怀善意，心怀恶意的传令官能给你设置什么障碍？若翰借着鸽子学到了这一点。他学到的是什么？让他自己重复。他说：\BibleQuote{那一位向我说：你看见圣神降下，如同鸽子，停在祂身上，那一位就是用圣神施洗的}。愿那些诱惑者不要欺骗你，鸽子，他们说：我们施洗。鸽子，承认鸽子所教导的：\BibleQuote{那一位就是用圣神施洗的}。借着鸽子，我们被教导那就是祂；而你以为你是由祂的权能受洗，其实是借着谁的服务而你受洗？若你这样想，你尚不在鸽子的身体内；若你不在鸽子的身体内，难怪你没有淳朴；因为借着鸽子，主要象征了\OriginalTerm{淳朴}{simplicitas}。\par

12. 因此，我的弟兄们，若翰藉着鸽子的纯朴学到了\Quote{这就是那以圣神施洗的}，除非是为了表明那些分裂教会的人不是鸽子？他们是鹰和鸢。鸽子不撕裂。而你们看到，他们因着所谓的迫害而使我们遭受仇恨。的确，他们遭受了身体上的迫害（如果可以这样称呼），因为那是主的鞭笞，明显是暂时的管教，免得祂若他们不承认并改正，最终永远定他们的罪。真正迫害教会的是那些以诡计迫害的人；他们用舌头的剑击打更重的心；他们尽可能地杀害基督在人的内里，流更苦的血。他们似乎害怕当权者的审判。如果你品行好，当权者能对你怎样？如果你作恶，就该害怕当权者；\Quote{因为他不是无故地佩剑}\BibleRef{罗 13:4}，宗徒说。不要拔出你用来打击基督的剑。基督徒，你在基督徒身上迫害什么？皇帝在你身上迫害了什么？他迫害了肉身；你在基督徒身上迫害了圣神。你没有杀害肉身。然而，他们也不饶恕肉身；凡他们能杀的，都用刀杀了；他们既不饶恕自己的，也不饶恕外人。这是众所周知的。当权者因合法而遭人怨恨；依法行事者行事招人怨恨；违法行事者却不招人怨恨。我的弟兄们，你们每个人要留意基督徒拥有什么。他的人性与许多人共有，他的基督徒身份使他与许多人区别开来，而他的基督徒身份比他的本性更严格地属于他。因为，作为基督徒，他按天主的肖像被更新，而人本是按天主的肖像被造的（\BibleRef{哥 3:10}）；但作为人，他可能是坏的，可能是外邦人，可能是偶像崇拜者。你在基督徒身上迫害的正是他更好的部分；你想从他那里夺走他赖以生存的东西。因为他按生命之神暂时生活，这神使他的身体有活力，但他按从他主那里领受的圣洗永远生活；你想从他那里夺走他从主领受的，你想夺走他赖以生存的东西。强盗对于他们要掠夺的人，目的在使自己致富，剥夺受害者所有的一切；但你从他那里夺取，而你自己却不会因此更多，因为夺取他的并不会增加你的。但是，说实话，他们正像那些夺取自然生命的人：他们夺走他人的生命，而自己却没有两条命。\par

13. 那么，你想夺走什么？在你想要再付洗的人身上，什么使你不悦？你不能给他他已经有的，反而使他否认他所拥有的。外邦人的迫害者曾犯下更大的残暴吗？刀剑指向殉道者，放出野兽，点上火：这些是为了什么？为了使受苦者被迫说：我不是基督徒。你教导你所想要再付洗的人什么呢？除非他先说：我不是基督徒。迫害者用火焰，你用舌头；你用引诱做了他通过杀害没有做到的事。你给予什么，又要给谁？如果他告诉你真相，不撒谎，被你引诱，他会说：我有。你问：你有圣洗吗？他说：我有。只要他说\Quote{我有}，你就说\Quote{我不给}。那就不要给，因为你想给的不能附着于我；因为我领受的不能从我这里夺走。但是，等着，让我看看你要教我什么。他说：首先，说\Quote{我没有}。但这是我已经有的；若我说\Quote{我没有}，我撒谎；因为有的我确实有。他说：你没有。教我我没有它。一个坏人给了你。如果基督是坏的，那么一个坏人给了我。他说：基督不是坏的，但基督没有给你。那么谁给了？回答：我知道我是从基督领受的。他说：给你的人不是基督，而是一个\OriginalTerm{叛教者}{traditor}。我会看到谁是天主的仆人；我会看到谁是天主的传报者。关于职务本身，我不争辩；我留意审判者：也许在你对天主的仆人的异议中，你说了谎。但我拒绝讨论；让两者的主处理祂自己仆人的案件。如果我要你举证，你什么也不能提供；事实上，你在说谎，已经证明你无法举证。但我不把我的案件建立在这上面，免得因我热心地为无辜者辩护，你就推断我把希望寄托在无辜者身上。那些人是怎样的人，随他们去吧；我是从基督领受的，我是由基督付洗的。他说：不，不是基督，而是那位主教给你付了洗，而且那位主教与他们共融。我是由基督付洗的，我知道。你怎么知道？鸽子教导了我，就是若翰所见的。邪恶的鸢，你不能把我从鸽子的怀中撕裂。我被列于鸽子的肢体中，因为鸽子所教导的，我就知道。你对我说：这个人或那个人给你付了洗；藉着鸽子，对我和对你说：\Quote{这就是那施洗的}。我应该相信谁，是鸢还是鸽子？\par

14. 确实地告诉我，好使你们被那盏灯所挫败，那灯也曾挫败过你们的前辈仇敌，他们与你们相似，即法利塞人。当他们质问主，祂凭借什么权柄做那些事时，祂说：\Quote{我也要问你们这个问题：告诉我，若翰的圣洗圣事，是从哪里来的？是从天上，还是从人来的？} 他们正预备施展诡计，却被这问题缠住，便开始彼此议论说：\Quote{如果我们回答说，是从天上来的，祂会对我们说：那么你们为什么不信他呢？} 因为若翰曾论及主说：\Quote{请看天主的羔羊，除免世罪者！}\BibleRef{若 1:29} 那么，你们为什么询问我凭什么权柄行事呢？哦，豺狼，我所行的，是凭羔羊的权柄。但为使你们认识羔羊，你们为什么不信若翰呢？他曾说：\Quote{请看天主的羔羊，除免世罪者！}\BibleRef{若 1:29} 他们知道若翰曾论及主所说的话，便彼此说：\Quote{如果我们说若翰的圣洗圣事是从天上来的，祂会对我们说：那么你们为什么不信他呢？如果我们说是从人来的，百姓会拿石头砸我们，因为他们都认为若翰是先知。} 因此，他们害怕人；因此，他们羞于承认真理。黑暗以黑暗回应；但他们却被光明战胜。他们如何回答呢？\Quote{我们不知道；} 关于他们知道的事，他们说：\Quote{我们不知道。} 主说：\Quote{我也不告诉你们，我凭什么权柄做这些事。}\BibleRef{玛 21:23} 第一批敌人就这样被挫败了。怎么挫败的？借着那盏灯。那灯是谁？若翰。我们能证明他就是那盏灯吗？我们能证明；因为主说：\Quote{他是燃烧而照耀的灯。}\BibleRef{若 5:35} 我们也能证明敌人是因他而被挫败的吗？请听圣咏：\Quote{我为我预备了一盏灯，是为我的基督。我要使他的敌人蒙羞。}\par

15. 趁我们还在这生命的黑暗中行走，我们借着信德的灯行走；让我们也持守若翰这盏灯，并借着他挫败基督的敌人；事实上，让基督亲自借着祂自己的灯挫败祂的敌人。让我们提出主曾向犹太人提出的问题，让我们问说：\Quote{若翰的圣洗圣事，是从哪里来的？是从天上，还是从人来的？} 他们会说什么呢？注意，如果他们不像被灯挫败的敌人。他们会说什么呢？如果他们说，是从人来的，连他们自己的人也会用石头砸他们；但如果他们说，是从天上来的，让我们对他们说：那么，你们为什么不信他呢？他们也许说，我们信他。那么，你们为什么说你们施洗呢？而若翰说：\Quote{这就是那施洗者。} 但是，他们说，侍奉如此伟大审判官的施行者应当是正义的。而我也说，大家都说，侍奉如此伟大审判官的施行者应当是正义的；让施行者无论如何，如果他们愿意，成为正义的；但如果那些坐在梅瑟座位上的不愿成为正义的，我的主人使我安全，祂的圣神曾说：\Quote{这就是那施洗者。} 祂怎样使我安全呢？\Quote{经师和法利塞人坐在梅瑟的座位上：凡他们所说的，你们要实行；但不要效法他们的行为，因为他们只说而不做。}\BibleRef{玛 23:2} 如果施行者是正义的，我就把他与保禄同列，与伯多禄同列；我与那些正义的施行者同列：因为事实上，正义的施行者不寻求自己的光荣；他们是施行者，不愿被人视为审判官，他们厌恶人将希望寄托在他们身上；因此，我把正义的施行者与保禄同列。因为保禄说什么？\Quote{我栽种了，阿颇罗浇灌了，但使之生长的是天主。所以，栽种的不算什么，浇灌的也不算什么，只有那使之生长的天主才算什么。}\BibleRef{格前 3:6} 但骄傲的施行者被与魔鬼同列；然而基督的恩赐并不会被污染，它纯净地流过他，液体般地穿过他，到达肥沃的土地。假设他是石头，不能从水养育果实；即使通过石头的渠道，水仍然流过，水流到园地；在石头的渠道里它不产生任何生长，但尽管如此，它给园子带来许多果实。因为圣事的精神德能就像光：被光照的人纯净地接受它，即使它穿过不洁的人，也不会被玷污。让施行者无论如何成为正义的，不寻求自己的光荣，而寻求祂的光荣，他们是祂的施行者；让他们不要说：\Quote{这圣洗是我的；} 因为它不是他们的。让他们听从若翰。看，若翰充满了圣神；他的圣洗圣事是从天上来的，不是从人来的；但他拥有它多久呢？他亲自说：\Quote{预备主的道路。}\BibleRef{若 1:23}\par

16. 然而，他们习惯对我们说什么呢？\Quote{看，在若翰之后，圣洗圣事已被授予了。} 因为在那问题在天主教会内得到恰当处理之前，许多人，包括伟大而良善的人，都曾在此事上犯错；但因为他们都是鸽子的肢体，他们没有自绝于教会，在他们身上发生了宗徒所说的：\Quote{如果你们在某些事上存有异议，天主也会将此启示给你们。}\BibleRef{斐 3:15} 因此，那些自绝于教会的人成了不可教导的。那么他们习惯说什么呢？看，在若翰之后圣洗圣事已被授予；在异端者的圣洗之后，难道就不该再授予吗？因为某些领受了若翰圣洗的人，被保禄命令受洗，\BibleRef{宗 19:3} 因为他们没有领受基督的圣洗。那么，他们说，你们为什么夸耀若翰的功劳，而似乎贬低异端者的可怜呢？我也承认异端者是邪恶的；但异端者所授予的是基督的圣洗，那是若翰所未曾授予的。\par

17. 我回到若翰身上，说：\Quote{这就是那施洗者。} 因为若翰比异端者更好，正如若翰比醉汉更好，若翰比凶手更好。如果我们应当由更坏的人施洗，因为宗徒们是由更好的人施洗——他们中无论谁由醉汉施洗，我不说由凶手，我不说由某个恶人的帮手，我不说由抢夺他人财物者，我不说由欺压孤儿者，或拆散婚姻者；我不提这些；我讲的是每年发生、每天发生的事；我讲的是所有人都被召唤去做的，甚至在这城中，当有人对他们说：让我们充当无理性者，让我们享乐，像在正月初一的朔日这样的日子我们不应禁食：这些就是我所讲的，这些琐碎的日常行为——当一个人由醉汉施洗时，谁更好？若翰还是醉汉？请回答，如果你能，说醉汉比若翰更好！你绝不敢这样做。那么你，作为一个清醒的人，由你的醉汉施洗吧。因为如果宗徒们是由若翰施洗，那么清醒的人岂不更应当由醉汉施洗？或者你说：醉汉与我合一？难道若翰，新郎的朋友，不与新郎合一吗？\par

18. 但我对你自己说，无论你是谁，你比若翰更好吗？你绝不敢说：我比若翰更好。那么，如果你自己的人更好，就让他们由你施洗吧。因为如果洗礼是由若翰施行的，你就羞愧，因为洗礼不是由你施行的。你会说：但我拥有并教导基督的洗礼。那么，现在承认审判者，不要做骄傲的传报者。你施行基督的洗礼，所以洗礼不是由你施行；由若翰施行了，因为他没有施行基督的洗礼，而是他自己的；因为他如此接受了它，以致它是他自己的。所以，你不比若翰更好；但通过你所施行的洗礼比若翰的更好；因为一个是基督的，另一个是若翰的。由保禄施行的和由伯多禄施行的，是基督的；即便由犹达斯施行的，也是基督的。犹达斯施行了洗礼，犹达斯之后洗礼没有重复；若翰施行了洗礼，若翰之后洗礼重复了：因为如果由犹达斯施行的，是基督的洗礼；但由若翰施行的，是若翰的洗礼。我们不把犹达斯置于若翰之上；但基督的洗礼，甚至由犹达斯的手施行，我们也置于若翰的洗礼之上，即使那是若翰的手正确施行的。因为关于主在受苦之前，曾说过他施洗比若翰更多；随后又补充说：\Quote{其实，耶稣自己并没有施洗，而是他的门徒施洗。} \BibleRef{若 4:1} 他，又不是他：他凭权能，他们凭服务；他们履行施洗的服务，施洗的权能留在基督内。他的门徒们就施洗了，犹达斯还在门徒中；那么，那些由犹达斯施洗的人，不是再次受洗；而那些由若翰施洗的人，他们又再次受洗了吗？显然有重复，但不是同一洗礼的重复。因为若翰施洗的人，若翰施洗；犹达斯施洗的人，基督施洗。同样，那么，由醉汉施洗的人，由凶手施洗的人，由奸夫施洗的人，如果是基督的洗礼，就是由基督施洗的。我不怕奸夫、醉汉或凶手，因为我注视着鸽子，透过它向我说：\Quote{这就是那施洗者。}\par

19. 但是，我的弟兄们，说——我不说犹达斯——但任何人比那一位更好，乃是疯狂的，关于那一位曾说过：\Quote{在妇女所生者中，没有兴起比施洗若翰更大的。} \BibleRef{玛 11:11} 所以，没有仆人比他更受偏爱；但主的洗礼，甚至通过一个邪恶的仆人施行，也比一个本是朋友的仆人的洗礼更受偏爱。请听宗徒保禄提到的那类人，假弟兄，出于嫉妒宣讲天主的圣言，以及他论及他们的话：\Quote{我就因此欢喜，而且还要欢喜。} \BibleRef{斐 1:15} 他们宣讲基督，确是出于嫉妒，但仍然宣讲基督。不要考虑\Quote{为什么}，而要考虑\Quote{谁}：出于嫉妒，基督被宣讲给你们。看基督，避免嫉妒。不要模仿邪恶的宣讲者，而要模仿被宣讲的那位善者。因此，基督曾被一些人出于嫉妒而宣讲。嫉妒是什么？一种可怕的恶。魔鬼因这恶被摔下；这个恶毒的瘟疫使他摔下；某些基督的宣讲者被它占据，然而宗徒允许他们宣讲。为什么？因为他们宣讲基督。但那嫉妒的人恨；那恨的人，关于他说了什么？请听宗徒若望：\Quote{凡恨自己弟兄的，便是凶手。} \BibleRef{若一 3:15} 看，若翰之后洗礼被施行，凶手之后洗礼没有被施行；因为若翰给了自己的洗礼，凶手给了基督的洗礼。那圣事如此神圣，连凶手的服务也不能玷污它。\par

20. 我不拒绝\Person{若翰}，反而我相信\Person{若翰}。我相信\Person{若翰}什么？是他通过鸽子所学到的？他通过鸽子学到了什么？\BibleQuote{这就是以圣神施洗的那一位。}因此，弟兄们，现在要紧紧抓住这一点，并铭刻在你们心中；因为如果我想今天更详尽地解释为何是通过鸽子？时间就不够了。因为我认为，我已在某种程度上向你们，圣洁的弟兄们，说明了：一件必须学习的事是通过鸽子灌输给\Person{若翰}的，一件关于基督的事，是\Person{若翰}所不知道的，尽管他已经认识基督；但是，为何这件事必须通过鸽子来指明，如果可能简短说明，我是想说的；但因为说来话长，我不愿加重你们的负担，既然你们的祈祷帮助我履行了我的诺言；借着你们虔诚的关注和善愿的再次帮助，同样也会对你们变得清楚，为何\Person{若翰}在那件他关于主所学习的事上，即\BibleQuote{这就是以圣神施洗的那一位。}并且他没有将施洗的权力转让给他的任何一个仆人——为何这件事他必须仅通过鸽子来学习。\par

\SourceRecord{Translated by John Gibb.；Nicene and Post-Nicene Fathers, First Series；Vol. 7.；Edited by Philip Schaff.；Buffalo, NY: Christian Literature Publishing Co.,；1888.；<http://www.newadvent.org/fathers/1701005.htm>.}

% structure: p0001 source=h1 target=section reason=page-title

\section{讲道集第六篇（\BibleRef{若 1:32-33}）}

1. 圣善的弟兄们，我向你们承认，我曾担心寒冷会使你们在聚会时变得冷淡；但你们以这次拥挤的集会证明了你们在精神上是热切的，我便不怀疑你们也在为我在祈祷，好使我偿还我所亏欠的。因为我曾以基督的名向你们许诺，由于之前时间短促，未能讲解，我今天要讨论为何天主愿意以鸽子的形象显示圣神。为了解释这事，今日曙光已照临我们；我看出，出于听闻的渴望和虔诚的热忱，你们聚集的人数比平时更多。愿天主藉我们的口满全你们的期待。因为你们的聚集是出于你们的爱；但爱什么呢？若是爱我们，那也好；因为我们愿意被你们所爱，但不是为着我们自己。因为我们爱你们于基督内，你们也当报以爱我们于基督内，愿我们的爱共同叹息朝向天主；因为鸽子的叫声是叹息或呻吟。\par

2. 如今，若鸽子的叫声是呻吟，如同我们大家所知道的，鸽子是在爱中呻吟，那么请听宗徒所说的，并不要惊奇圣神愿意以鸽子的形象显示：\Quote{因为按我们应该怎样祈祷，}他说，\Quote{我们不知道；但圣神亲自以不可言喻的叹息代我们转求。}\BibleRef{罗 8:26}那么，我的弟兄们，我们该说什么呢？难道说圣神在自己内，与圣父和圣子同享完美永恒的真福时呻吟吗？因为圣神是天主，如同天主子是天主，圣父也是天主。我说了\Quote{天主}三次，但不是三个天主；因为确实，是天主\Emph{三次}而非三个天主；因为圣父、圣子、圣神是一个天主：这你们知道得很清楚。那么，并不是在祂自己内，在祂与那圣三、那真福、祂永恒的本体中，圣神呻吟；而是在我们内祂呻吟，因为祂使我们呻吟。圣神教导我们呻吟，这并不是小事，因为祂使我们知道我们在异地作客，祂教导我们渴慕祖国；借着那渴望，我们呻吟。那在这世上安好的人，或者更好说那自以为安好的人，他在肉性之事的喜乐中欢腾，在暂时之物的丰裕中欢腾，在空虚的幸福中欢腾，他有乌鸦的叫声；因为乌鸦的叫声充满喧哗，并非呻吟。但那知道自己处于这个必死生命的压力下，是一个\Quote{离开主}的朝圣者\BibleRef{格后 5:6}，他尚未拥有那应许给我们的永恒真福，但他在望德中拥有它，并将要在现实中拥有它，那时主要公开在光荣中来临，祂先前曾隐晦地谦卑而来；我说，那知道这事的人便叹息。并且只要他为此叹息，他就叹息得好；是圣神教他叹息，他从鸽子学来的。的确，许多人因着现世的苦难而叹息。他们或许被损失击碎，或被身体的病患压倒，或被关在监牢里，或被锁链捆绑，或被海中的波浪抛掷，或被仇敌的阴谋围困。因此他们叹息，但不是以鸽子的呻吟，不是出于对天主的爱，不是在于圣神。因此，当这些人从同样的苦难中被拯救出来时，他们大声欢呼，由此显明他们是乌鸦，不是鸽子。从方舟中放出一只乌鸦，它不再回来；放出一只鸽子，它回来了。诺厄放出了这两只鸟\BibleRef{创 8:6}。他那里有乌鸦，也有鸽子。那方舟包含了两种；若方舟是教会的预象，你们确实看到，在现今世界的洪水中，教会必然包含两种，既是乌鸦，也是鸽子。谁是乌鸦？那些寻求自己利益的人。谁是鸽子？那些寻求基督之事的人。\BibleRef{斐 2:21}\par

3. 因此，当祂派遣圣神时，祂以两种方式有形地显示了祂——藉着鸽子与火：藉着鸽子在祂受洗时降在主身上，藉着火在门徒们聚集时。因为主复活升天后，与门徒们共度了四十天，到了五旬节日，祂按照所应许的派遣了圣神给他们。相应地，那时圣神降临充满了那地方，先有从天上来的响声，好像一阵大风吹过，正如我们在《宗徒大事录》中所读到的，它说：\BibleQuote{有像舌头似的火焰显现给他们，}它说，\BibleQuote{分开停在他们每人头上；他们就充满圣神，按圣神所赐给他们的话说起别的语言来。}这里我们看到鸽子降在主身上；那里，分开的舌头降在聚集的门徒们身上：前者显示纯朴，后者显示热忱。有些人被称为纯朴，却只是懒惰；他们被称为纯朴，实为迟钝。但充满圣神的斯德望并非如此：他纯朴，因为他没有伤害任何人；他热忱，因为他责备不虔敬的人。他在犹太人面前没有保持缄默。他的言语是燃烧的：\BibleQuote{你们这些执拗、心与耳未受割损的人，时常反抗圣神。}强大的冲击力，却是无胆汁的鸽子发怒。要知道他虽激烈却无胆汁，请注意，听到这些话，那些乌鸦立刻拿起石头，冲向这只鸽子。他们开始用石头砸斯德望；而那位稍早前激昂、灵里热情如火的人——他曾如发起进攻般击打敌人，像一个充满暴力的人用你们所听到的那样如火如荼的话攻击他们：\BibleQuote{你们这些执拗、心与耳未受割损的人，}——谁听到那些话都可能以为斯德望若被允许，会立刻将他们消灭——但当扔向他的石头到达时，他屈膝跪下说：\BibleQuote{主，不要将这罪归于他们。}\BibleRef{宗 7:51}他持守了鸽子的合一。因为他的师傅，鸽子曾降在祂身上，已在他之前行了同样的事：祂悬在十字架上时说：\BibleQuote{父，宽恕他们，因为他们不知道自己在做什么。}\BibleRef{路 23:34}因此，鸽子表明：被圣神圣化的人应无诡诈；而火则指示我们，他们的纯朴不应保持冷淡。也不要因舌头是分开的而困扰；因为舌头是多样的，因此显现为分开的舌头。它说：\BibleQuote{分开的舌头，像火焰一样，且停在每人身上。}舌头虽有分别，但舌头的分别并不意味着分裂。不要因分开的舌头而害怕分离；要在鸽子中认出合一。\par

4. 因此，圣神降在主身上时应当以这种方式显现，好让每个人明白：若他拥有圣神，就当像鸽子一样纯朴，与弟兄们有真正的和平，就是鸽子之吻所象征的和平。乌鸦也有它们的吻，但在乌鸦那里是虚假的和平，在鸽子那里才是真正的和平。所以，并非每个说\BibleQuote{平安与你们同在}的人都该被当作鸽子听从。那么，如何区分乌鸦之吻与鸽子之吻呢？乌鸦亲吻，却撕裂；鸽子的本性是绝不撕裂的。因此，哪里有撕裂，亲吻中就没有真正的和平。那些没有撕裂教会的人才有真正的和平。乌鸦以腐肉为生，鸽子却不如此；它靠地上的果实生活，它的食物是无辜的。弟兄们，鸽子的这一点确实值得赞叹。麻雀是很小的鸟，但至少它们也杀苍蝇。鸽子却不做这类事，因为它不以死物为食。那些撕裂教会的人以死物为食。天主是强有力的；让我们祈祷，那些被他们吞食却未察觉的人能重新活过来。许多人承认他们确实重生了，因为在他们归来时，我们每天都在基督的名下与他们一同喜乐。要纯朴，但要以这样的方式保持热忱，让你们的热情在你们的舌头上。不要缄默，用炽热的舌头说话，点燃那些冷淡的人。\par

5. 为什么，我的弟兄们？谁看不见他们所看不见的？这不足为奇，因为那些不愿从那里回转的人，就像从方舟放出的乌鸦一样。因为谁看不见他们所看不见的？他们甚至对圣神本人忘恩负义。看，鸽子降在主身上，降在受洗的主身上：于是，那神圣而真实的天主圣三显现出来，对我们而言是唯一的天主。因为主从水中上来，正如我们在福音中所读到的：\BibleQuote{天忽然为祂开了，祂看见天主圣神好像鸽子降下，落在祂身上；立刻有声音说：这是我的爱子，我所喜悦的。}\BibleRef{玛 3:16}天主圣三极其明显地出现：圣父在声音中，圣子在人身中，圣神在鸽子中。在这圣三中，让我们看见，如同我们确实看见的，宗徒们被派遣到哪里，以及那些人不看的是多么奇妙。并非他们真的看不见，而是他们真的对迎面而来的事实闭上眼睛：就是门徒们因父、及子、及圣神之名被派遣出去，而派遣他们的是那一位，关于祂曾说：\BibleQuote{这人才是施洗者。}事实上，这话是对祂的仆人们说的，是祂自己保留了这权柄。\par

6. 若翰在祂身上所看见并认识的，正是他先前所不知道的。并非他不知道祂是天主子，或不知道祂是主，或不知道祂是基督；也并非他不知道这一点，即祂要用水和圣神施洗。这些他是知道的；但是，祂这样做时，要把权柄保留给自己，而不转移给祂的任何施行者，这正是他在鸽子身上所学到的。因为基督只将这项权柄保留给自己，而没有赐给祂的任何施行者，尽管祂屈尊借施行者施洗；我所说的这项权柄，正是教会的合一所赖以建立的，这合一预表在鸽子身上，正如经上所记载的：\Quote{我的鸽子是独一的，是她母亲的唯一者。} \BibleRef{歌 6:8} 因为，正如我曾说过的，诸位弟兄，倘若主将权柄转移给祂的施行者，那么有多少施行者就有多少圣洗，而圣洗的合一也就不复存在了。\par

7. 诸位弟兄，请注意：在我们的主耶稣基督来到祂受洗时（因为在受洗之后，鸽子降下，若翰借此认识到祂的独特之处，因为曾告诉他：\Quote{你看见圣神降下，如同鸽子，停在谁身上，谁就是要用圣神施洗的那一位。}），若翰知道祂就是用圣神施洗的那一位；但是，这独特之处在于，权柄不会从祂转移给另一个人，尽管祂赐予权柄，这正是他在那里学到的。我们如何证明若翰已经知道主要用圣神施洗呢？这样，他必须理解为从鸽子学到的就是：主将要用圣神施洗，并且权柄不会从祂转移给任何其他人？我们如何证明这一点？鸽子是在主受洗后降下的；但在主来到约旦河受若翰的洗之前，我们已经说过若翰认识祂，从那些话中可知，他说：\Quote{祢到我这里来受洗吗？我本当受祢的洗。} 是的，他认识祂是主，认识祂是天主子；但我们如何证明他已经知道同一位就是要用圣神施洗的呢？在祂来到河边以前，当许多人涌到若翰那里受洗时，他对他们说：\Quote{我固然以水洗你们，但在我以后来的那一位，比我更强，我连给祂解鞋带也不配，祂要以圣神和火洗你们。} \BibleRef{玛 3:14} 他早已知道这些。那么，他从鸽子那里学到了什么，以致他后来不成为撒谎者（愿天主不许我们这样想），如果这不是：在基督身上有某种独特之处，尽管许多施行者，无论是义人还是不义的人，都施洗，但圣洗的功效只归于那鸽子降在祂身上的那一位，也就是经上所说的\Quote{这是那用圣神施洗的一位。} 伯多禄可以施洗，但施洗的是祂；保禄可以施洗，但施洗的仍是祂；犹达斯可以施洗，但施洗的仍是祂。\par

8. 因为如果圣洗的神圣性是根据施行者功德的差异，那么功德不同，圣洗也会不同；领受者就会认为，他越是觉得给他施洗的人更好，他所领受的也越好。圣人本身——请明白，诸位弟兄，就是那些属于鸽子的人，他们在耶路撒冷圣城中有份，就是教会中的善人，关于他们，宗徒说：\Quote{主认识那些属于祂的人。} \BibleRef{弟后 2:19} ——他们被赋予不同的恩宠，并非都有相同的功德。有些人比另一些人更圣洁，有些人比另一些人更好。因此，如果一个人从一位圣善的义人那里领受圣洗，另一个人从另一位在天主面前功德较劣、地位较低、节制较差、生活较差的人那里领受圣洗，那么他们所领受的如何会是同一、同等、相似的，如果不是因为\Quote{施洗的是祂？} 正如当善者和更善者施行圣洗时，一个人不是领受好的，另一个人领受更好的；而是尽管施行者一个是善者，另一个是更善者，他们领受的却是同一、平等的，不是在一个情况下更好，在另一个情况下更差；同样，当恶人施行圣洗时，由于教会的无知或容忍（因为恶人要么不为人知，要么被容忍；糠秕被容忍，直到最后禾场被完全扬净），所赐予的是同一的，并不因施行者不同而不同，而是相同和平等的，因为\Quote{施洗的是祂。}\par

9. 因此，亲爱的诸位，让我们看看那些人不想看什么；不是他们可能看不到什么，而是他们悲伤地看到什么，仿佛它对他们关闭了。门徒们作为仆役，奉父、及子、及圣神之名被派往哪里去付洗？他们被派往哪里？祂说：\BibleQuote{去，给万民授洗。}弟兄们，你们已经听到那嗣业是如何来的：\BibleQuote{你向我请求，我必将万民赐你作嗣业，将地极赐你为产业。}你们已经听到\BibleQuote{法律出自熙雍，上主的话来自耶路撒冷。}\BibleRef{依 2:3} 因为在那里门徒们被告知：\BibleQuote{去，奉父、及子、及圣神之名给万民授洗。}\BibleRef{玛 28:19} 当我们听到\BibleQuote{去，给万民授洗}时，我们变得专注。因谁的名？\BibleQuote{因父、及子、及圣神之名。}这是唯一的天主；因为经上不是说要奉父、及子、及圣神之\InnerQuote{名}（复数），而是\BibleQuote{因父、及子、及圣神之名}。哪里你听到一个名，那里就有唯一的天主；正如论到亚巴郎的后裔所说的，宗徒保禄解释说：\BibleQuote{地上万民都要因你的后裔蒙受祝福；他没有说\InnerQuote{后裔们}，如同许多，而是如同一个，\InnerQuote{因你的后裔}，就是基督。}因此，正如宗徒愿意向你们指出，因为那处没有说\Quote{后裔们}，基督是唯一的；同样，这里当说\Quote{因名}，不说\Quote{因诸名}，正如\Quote{因后裔}，不说\Quote{因后裔们}，便证明父、及子、及圣神是唯一的天主。\par

10. 但是，看，门徒们对主说：我们被告知要因什么名付洗；祢使我们成为仆役，并对我们说：\BibleQuote{去，奉父、及子、及圣神之名付洗。}我们要往哪里去？往哪里？祢没有听到吗？往我的嗣业去。祢问：我们要往哪里去？往我用我的血所买来的地方。那么往哪里？祂说：往万民那里。我以为祂说：\BibleQuote{去，奉父、及子、及圣神之名给非洲人付洗。}感谢天主，主解决了这个问题，鸽子教导了我们。感谢天主，宗徒们被派往万民；如果往万民，那么就是往所有语言。圣神表明了这一点，在语言中被分开，在鸽子中被结合。这里语言被分开，那里鸽子结合它们。万民的语言是一致的，也许只有非洲的语言不一致。有什么比这更明显的，我的弟兄们？在鸽子中是合一，在语言中是万民的共通。因为从前语言因骄傲而变得不和谐，然后从一种语言变成了许多种语言。因为洪水之后，某些骄傲的人，好像试图对抗天主，仿佛有什么对于天主是高的，或者有什么能给骄傲以安全，他们建造了一座塔，显然是为了不被洪水毁灭，如果之后再有洪水的话。因为他们听说并认为一切不义都被洪水冲走了；他们不肯戒绝不义；他们寻求塔的高度作为对抗洪水的防御；他们建造了一座高塔。\BibleQuote{天主看到他们的骄傲，并挫败了他们的计谋，使他们的言语彼此无法理解，因此语言因骄傲而多样化。}\BibleRef{创 11:1} 如果骄傲造成了语言的差异，基督的谦卑将这些差异合而为一。教会如今正在把那座塔所分裂的重新联合起来。从一种语言，变成了许多种；不要惊奇：这是骄傲的作为。从许多种语言，变成了一种；不要惊奇：这是爱德的作为。因为虽然语言的声音各不相同，但内心所呼求的天主是唯一的，所保存的和平是唯一的。那么，圣神在表明合一时，除了藉着鸽子，还能怎样彰显呢？好使已进入和平状态的教会能被称为\BibleQuote{我的鸽子是唯一的}？谦卑应当怎样被代表呢，若不是藉着一只无辜、忧伤的鸟，而不是藉着一只骄傲、欢跃的鸟，如同乌鸦？\par

11. 但是也许他们会说：好吧，既然它是鸽子，而且鸽子是唯一的，那么就不可能在那唯一的鸽子之外有洗礼。因此，如果鸽子与你们同在，或者你们自己就是鸽子，那么当我来到你们这里时，你们就把我所没有的给我。你们知道这就是他们所说的；但是你们很快就会看到这不是鸽子的声音，而是乌鸦的喧嚷。因为请稍加留意，亲爱的诸位，要害怕他们的诡计；不，要提防他们，并且倾听反对者的话，只是为了拒绝它们，而不是吞下它们并纳入腹中。对此，你们当效法主在别人向祂递上苦胆时所行的：\BibleQuote{祂尝了，便吐了出来。}\BibleRef{玛 27:34} 你们也要这样，听到后便抛弃。他们到底说什么？让我们看看。看，他说：\Quote{你是鸽子。}哦，天主教会，正是对你说：\BibleQuote{我的鸽子是唯一的，是她母亲的独一者}；这肯定是对你说的。住口，不要问我；先证明是否对我说过；如果是对我说的，我就立刻听到了。\Quote{是对你说的，}他说。我以天主教会的口吻回答：\Quote{是对我说的。}而这个回答，弟兄们，从我一个人的口中发出，我相信也发自你们的心，我们一同肯定：是的，是对天主教会说的：\BibleQuote{我的鸽子是唯一的，是她母亲的独一者}。他说，在这鸽子之外，没有洗礼：我在这鸽子之外受了洗，所以我没有洗礼；如果我没有洗礼，为什么当我来到你们这里时，你们不给我洗礼？\par

12. 我也要提出问题；让我们暂且搁置询问这话是对谁说的：\Quote{我的鸽子是唯一的，是她母亲唯一的}；我们仍在询问——这话要么是对我说的，要么是对你（指听者）说的；让我们推迟询问这话是对谁说的。我所问的是：如果鸽子是纯朴的、无罪的、无苦胆的、亲吻时温和的、爪不凶猛的，那么贪婪者、掠夺者、奸诈者、昏醉者、无耻者，是否属于这鸽子的肢体？他们是这鸽子的肢体吗？他说：断乎不可。而谁又会真的这样说呢，弟兄们？且不提别的，若我只提掠夺者，他们可能是鹰的肢体，而非鸽子的肢体。鸢鹰掠夺和抢夺，鹰也如此，乌鸦也如此；鸽子不掠夺也不撕裂，因此，那些抢夺和抢劫的人不是鸽子的肢体。难道你们中间没有一个掠夺者吗？那么为什么圣洗依然有效，而在这情况下是鹰（而非鸽子）施予了它？为什么你们不自己在抢劫者、奸淫者、醉汉之后重新施洗？为什么不在你们中间的贪婪者之后重新施洗？难道这些人都是鸽子的肢体？你们如此羞辱你们的鸽子，以至于使那些具有秃鹫本性的人成为她的肢体。那么，弟兄们，我们说什么呢？在天主教会内有好人和坏人，但只有坏人与他们同在。但或许我这样说带有敌意：这一点也之后再加审视。他们确实说，在他们中间有好人和坏人；因为，若他们断言他们只有好人，那就让他们自己的人相信吧，我同意。让他们说：我们中间只有圣洁、正义、贞洁、清醒的人；没有奸淫者、没有高利贷者、没有欺骗者、没有假发誓者、没有酗酒者；让他们这样说吧，因为我不在意他们的舌头，我触摸他们的心。但由于他们为我们所熟知，也为你们和他们自己所熟知，正如你们在天主教会内为你们自己所知，也为他们所知，我们既不要指责他们，也不要让他们自夸。我们承认，在教会内有好人也有坏人，就像麦子和糠粃。有时，被麦子施洗的人是糠粃，而被糠粃施洗的人是麦子。否则，如果被麦子施洗的人所受的圣洗有效，而被糠粃施洗的人所受的圣洗无效，那么\Quote{这就是施洗者}这话就不真实了。但如果\Quote{这就是施洗者}是真实的，那么由糠粃所给的圣洗有效，并且他施洗的方式与鸽子相同。因为坏人（施行圣洗者）不是鸽子，也不属于鸽子的肢体，绝不可能被断言为如此，无论在我们天主教会内还是在他们那里，如果他们断言他们的教会就是鸽子。那么我们该如何理解呢，弟兄们？既然显而易见且众所周知，而且他们必须承认，尽管违反他们的意愿，即当在他们中间有坏人施行圣洗时，圣洗不是根据那些坏人给的；在我们这里也是如此，当坏人施行圣洗时，圣洗也不是根据他们给的。鸽子不会根据乌鸦施洗；那么，乌鸦为什么要根据鸽子施洗呢？\par

13. 亲爱的，思考一下，为什么通过鸽子也有某种象征被指出，即鸽子——也就是以鸽子形状的圣神——在受圣洗时来到主面前并停留在祂身上，而若翰因着鸽子的降临得知，主内里有一种专属于祂自己的施洗权能？因为正是藉着这专属于祂自己的权能，如我所说，教会的平安得到了保证。然而，也可能有人拥有鸽子之外的圣洗；但圣洗在鸽子之外而能使他得益，是不可能的。亲爱的，思考并理解我所说的，因为藉着这种欺骗，他们误导那些迟钝冷淡的弟兄们。让我们更纯朴、更热忱。看，他们说，我领受了，还是没有？我回答：你领受了。那么，如果我领受了，就没有你能给我的东西；我安全了，即使按你自己的见证。因为我肯定我已领受，你也承认我已领受：藉着双方的承认，我安全了：那么你对我允诺什么呢？既然你承认我已领受了你自称拥有的东西，你为什么还要使我成为天主教徒？但当我说：\Quote{到我这里来}时，我是在说，你并不拥有，而你却承认我拥有。你为什么说：\Quote{到我这里来}？\par

14. 鸽子教导我们。她从主的头回答，并说：\Quote{你有圣洗，但没有我为之叹息的爱德。}这是怎么回事，他说，我有圣洗，却没有爱德？我有圣事，却没有爱德？不要喊叫：告诉我，分裂合一的人怎么可能有爱德？他说，我有圣洗。你有；但那圣洗，没有爱德，对你毫无益处；因为没有爱德，你什么都不是。圣洗本身，即使在一个什么都不是的人身上，也不是什么都没有。圣洗的确是某种东西，而且是伟大的东西，为了那位之故，经上说：\Quote{这就是施洗者}。但恐怕你幻想那伟大的东西能对你有什么益处，如果你不在合一中，那么在他受圣洗之后，鸽子降临，仿佛在暗示：\Quote{如果你有圣洗，要留在鸽子内，免得你所拥有的对你没有益处。}那么，我们说，到鸽子这里来吧；不是要你开始拥有以前没有的，而是使你确实拥有的开始有益于你。因为你在外面拥有圣洗而导致毁灭；如果你在里面拥有它，它就开始有益于你而得救。\par

15. 圣洗不仅对你们无益，反而有害。连圣物也可能有害。善人得圣物，是为得救；恶人得之，是为受定案。因为我们确实知道，弟兄们，我们领受的是什么；我们所领受的无论如何是圣的，无人说它不圣。宗徒说什么？\BibleQuote{但那吃喝不相称的，就是吃喝自己的定案。}（\BibleRef{格前11:29}）他并没有说那件事物本身是坏的，而是说恶人因领受不当，把他所领受的善事领受为定案。主递给犹大的那一小块饼是坏的吗？决不是！医生不会给人毒药；医生给的是健康；但因人无尊严地领受，那领受的人因不处于平安中，便将健康领受为毁灭。同样，那领受圣洗的人也是如此。他说：\Quote{我（圣洗）是为我自己领的。}你领了圣洗，我承认。好好留意你所拥有的；正是因着你所拥有的那同一件事，你将受定案。为什么？因为你拥有属于鸽子的，却不在鸽子内。如果你拥有鸽子的事物而在鸽子内，你就安全。设想你是一个士兵：如果你有将军的印记在营内，你安全服役；但如果你有印记出界，那印记不仅无助于你的服役，你反而要作为逃兵受罚。所以，来吧，来吧，不要说：\Quote{我已经领了，我已经够了。}来吧，鸽子正在呼唤你，以她的叹息呼唤你。我的弟兄们，我对你们说：要藉着呻吟而非争吵来呼唤；要藉着祈祷、邀请、禁食来呼唤；让他们因你们的爱德明白你们怜恤他们。我不怀疑，我的弟兄们，如果他们看见你们的悲伤，他们将惊愕，并重新活过来。所以，来吧，来吧；不要害怕；你若不来，就该害怕；不，不要害怕，宁可哀叹你自己。来吧，你若来，必要喜乐；你固然要在旅途的困苦中呻吟，但要在希望中喜乐。来吧，到鸽子那里去；曾有人对她说过：\Quote{我的鸽子是唯一的，是她母亲的独生女。}你们难道没有看见那唯一只鸽子在基督的头上，你们难道没有看见在整个世界中的那些舌头？藉着鸽子与藉着舌头是同一位圣神；如果藉着鸽子是同一位圣神，藉着舌头也是同一位圣神，那么圣神是赐给了整个世界。你自绝于那圣神，好使你与乌鸦一同鼓噪，而不是与鸽子一同叹息。那么，来吧。\par

16. 但你可能忧虑，说：\Quote{我在外面领了圣洗；我因此害怕我有罪，因为我在外面领了圣洗。}你已经知道该哀叹什么了。你说得对，你是有罪的，不是因为你的领受，而是因为你在外面领受。那么，保存你所领受的；改正你在外面领受的事实。你领受了属于鸽子的事物，却不在鸽子内。有两件事是对你说的：你领受了，以及你在鸽子之外领受了。就你领受了而言，我赞成；就你在外面领受了而言，我不赞成。那么，保存你所领受的；它没有被改变，而是被承认：它是我君王的印记，我不会亵渎它。我要纠正逃兵，而不是改变印记。\par

17. 不要因你的圣洗而夸耀，因为我称它为真实的圣洗。看，我说它是如此；整个天主教会都说它是如此；鸽子注视着它，承认它，并且呻吟，因为你是在外面拥有它；她在其中看到她可以承认的东西，也看到她可以纠正的东西。它是真实的圣洗，来吧。你夸耀它是真实的，然而你却不肯来吗？那么，那些不属于鸽子的恶人又怎么样呢？鸽子对你说：\Quote{连那些恶人——我在他们中间呻吟，他们不属于我的肢体，而我必须在他们中间呻吟——他们不也拥有你所夸耀的东西吗？许多醉汉不也领了圣洗吗？许多贪婪的人不也领了圣洗吗？许多拜偶像的，更坏的，是暗中拜偶像的人，不也领了圣洗吗？外教人不是公开到偶像那里去，或者至少曾经去吗？而现在基督徒暗中寻求占卜者，咨询占星家。然而这些人都有圣洗；但鸽子在乌鸦中间呻吟。}那么，你为什么在拥有它这件事上夸耀呢？你所拥有的，恶人也拥有。要有谦卑、爱德、和平；要拥有那你还未拥有的善事，好使你已拥有的善事对你有益处。\par

18. 因为你们所拥有的，连\OriginalTerm{西满·马古斯}{Simon Magus}也拥有过——\WorkTitle{宗徒大事录}可以作为见证，这部正典书籍每年都应在教会中诵读。你们知道，每年在主受难之后的季节，都要诵读那本书，其中记载了宗徒如何悔改，从迫害者变成了宣讲者；\EditorialNote{宗徒大事录第九章}同样也记载了在五旬节那天，圣神以火舌的形状降下。\EditorialNote{宗徒大事录第二章}我们读到，在撒玛黎雅，许多人因斐理伯的宣讲而信了主；他被认为是宗徒之一，或是执事之一，因为我们读到那里任命了七位执事，其中包括斐理伯的名字。那么，通过斐理伯的宣讲，撒玛黎雅人信了；撒玛黎雅开始充满信徒。西满·马古斯也在那里。他用巫术迷惑了民众，以致人们把他当作天主的大能。然而，他被斐理伯所行的奇迹所打动，也信了；但他是如何信的，随后发生的事件证实了。西满也受了圣洗。在耶路撒冷的宗徒们听说了这事，就派伯多禄和若望到撒玛黎雅去；他们发现许多人已受了圣洗，但还没有人领受圣神——如同当时圣神降临的方式，使领受者能说各种语言，作为外邦人将要相信的明显标记——于是他们为这些人祈祷并覆手，他们就领受了圣神。这位西满——他不是鸽子，而是教会中的乌鸦，因为他寻求自己的利益，而非耶稣基督的利益；因此他更爱基督徒所拥有的权能，而非正义——西满，我说，看见圣神是通过宗徒的覆手而赐予的（不是宗徒赐予的，而是因他们的祈祷而赐予的），便对他们说：\Quote{你们愿意给我多少钱，好使我也能覆手给人，赐予圣神？伯多禄对他说：\InnerQuote{愿你的钱与你一同丧亡，因为你以为天主的恩赐可以用钱买到。}}他对谁说的\Quote{\InnerQuote{愿你的钱与你一同丧亡}}？无疑是对一个已受圣洗的人。他已经领受了圣洗，但他没有依附于鸽子的怀抱。你们要明白他没有依附；请听宗徒伯多禄的话，因为他接着说道：\Quote{你在这种信仰中没有份也没有权，因为我看出你正怀着苦胆的苦楚。}\BibleRef{宗 8:5}鸽子没有苦胆；西满却有，因此他从鸽子的怀抱中被分离出来。圣洗对他有什么益处呢？所以，不要夸耀你们的圣洗，好像仅凭圣洗就能使你们得救。不要发怒，除掉你们的苦胆，来到鸽子那里。在这里，那东西才会对你们有益，否则不仅无益，反而有害。\par

19. 你也不要说：\Quote{我不会来，因为我在外面受了圣洗。}那么，开始拥有爱德，开始结出果实，让果实能在你身上被找到，鸽子就会把你带进去。我们在圣经中读到这一点。方舟是用不朽坏的木头建造的。不朽坏的木材就是圣人们，即属于基督的信徒。正如圣殿中被用来建造的活石被认为是有信仰的人，同样，不朽坏的木材就是那些在信仰中坚持到底的人。那么，在这同一艘方舟中，木材是不朽坏的。方舟就是教会，鸽子就在那里施圣洗；因为方舟漂浮在水上，不朽坏的木材在方舟内受了圣洗。我们发现有些木材是在外面受圣洗的，如同世上所有的树木一样。然而，水是相同的，并非另一种；所有的水都来自天上，或来自深渊的泉源。正是同样的水，既使方舟内不朽坏的木材受了圣洗，也使外面的木材受了圣洗。鸽子被放出，起初找不到落脚之处；它就返回方舟，因为到处是水，它宁愿返回，也不愿重受圣洗。但乌鸦在水尚未干时就放出，它重受了圣洗，却不愿返回，就死在那水中。愿天主使我们免于乌鸦的死亡。因为乌鸦为何不返回，难道不是因为被水冲走了吗？但另一方面，鸽子找不到落脚之处，水在四周向它呼唤：\Quote{来吧，来吧，在这里浸一下；}正如这些异端者呼喊道：\Quote{来吧，来吧，这里你有它；}鸽子找不到落脚之处，就返回了方舟。诺厄第二次又放出它，正如方舟差遣你们出去对他们说话；后来鸽子做了什么？因为外面有受过圣洗的木材，它带回了橄榄枝。那枝子既有叶子也有果实。愿你们不仅拥有言语，也不仅拥有叶子；愿有果实，你们就返回方舟，不是靠自己，乃是鸽子召你们回来。你们在外边叹息，好把她们召回里面。\par

20. 此外，关于这橄榄的果实，如果仔细考察，你们会明白它是什么。橄榄的果实象征爱德。我们如何证明这一点？正如所有液体都无法压住油，它反而冲破一切，向上浮起并超越它们；同样，爱德不能被压到低处，而必定要显现在高处。因此宗徒论到它说：\Quote{我还要指示你们一条更高超的道路。}既然我们说过油超越其他液体，为免有人以为这不是指爱德，宗徒说：\Quote{我指示你们一条更高超的道路，}让我们听听随后的话：\Quote{我若能说人间的语言和天使的语言，却没有爱德，我就成了发声的锣，或发响的钹。}现在去吧，多纳图，呼喊：\Quote{我善于辞令；}现在去吧，呼喊：\Quote{我有学问。}能有多大善辞令？能有多大学问？你说了天使的语言吗？即使你能说天使的语言，却没有爱德，我听见的不过是发声的锣和发响的钹。我要的是实在；让我在叶子中找到果实；但愿他们不仅有言语，而且有橄榄，让他们返回方舟。\par

21. 但你会说：我拥有圣事。你说得对；圣事是天主的；你领受了圣洗圣事，我承认这点。但宗徒说了什么？\Quote{纵使我通晓一切奥秘，拥有先知话和一切信德，甚至能够移山，}你可能会说这话，\Quote{我相信，这对我足够了。}但雅各伯说了什么？\Quote{魔鬼也相信，且战栗。}\BibleRef{雅 2:19}信德是强大的，但若没有爱德，便一无所益。魔鬼宣认了基督。因此，他们由于相信，而非由于爱，说出了：\Quote{我们与祢有什么相干？}\BibleRef{谷 1:24}他们有信德，却没有爱德，所以他们是魔鬼。不要以信德自夸；你在这方面与魔鬼是同一层次的。不要对基督说：我与祢有什么相干？因为基督的合一在向你说话。来，学习平安，回到鸽子的怀抱中。你在外边领受了圣洗圣事；拥有果实，你就会回到方舟。\par

22. 但你会说：\Quote{你们为什么寻找我们，如果我们是坏人？}为要使你们成为好人。我们寻找你们的原因，正是因为你们是坏的；因为如果你们不是坏的，我们就已经找到你们，也就不再寻找你们了。那好人已被找到；那坏人仍在被寻找。因此，我们正在寻找你们；你们回到方舟吧。\Quote{但我已经领受了圣洗圣事。}\Quote{纵使我通晓一切奥秘，拥有先知话和一切信德，甚至能移山，却没有爱德，我什么也不是。}让我在那里看到果实；让我在那里看到橄榄，你们就被召回到方舟。\par

23. 但你说什么？\Quote{看，我们遭受了许多痛苦。}但愿你们是为基督而受这些痛苦，而不是为了你们自己的荣誉！听随后的话：他们有时确实夸耀，因为他们做了许多施舍，周济穷人；因为他们遭受磨难：但这是为多纳图，不是为基督。想想你们如何受苦，因为如果你们为多纳图受苦，那是为一个骄傲的人；你们若为多纳图受苦，就不在鸽子内。多纳图不是新郎的朋友；因为他若是，他就会寻求新郎的荣耀，而不是自己的荣耀。看新郎的朋友说：\Quote{这就是那施洗者。}你们为之受苦的那个人不是新郎的朋友。你们没有婚宴礼服；如果你们来赴宴，会被赶出去；不，你们已经被赶出去了，因此你们是可怜的：终于回来吧，不要自夸。听宗徒说什么：\Quote{纵使我把我所有的财产分施给穷人，并舍身被火焚烧，却没有爱德。}看看你们所没有的。\Quote{纵使我舍身被火焚烧，}他说，而且也是为了基督的名；但由于有许多人骄傲地这样做，而非出于爱德，因此，\Quote{纵使我舍身被火焚烧，却没有爱德，于我毫无益处。}\BibleRef{格前 13:2}那些在迫害时期受苦的殉道者是因为爱德而行此事；但这些人出于他们的虚荣和骄傲而行；因为在没有迫害者的情况下，他们自投身于毁灭。那么，来吧，使你们能拥有爱德。\Quote{但我们有自己的殉道者。}什么殉道者？他们不是鸽子；因此他们试图飞翔，却跌落在岩石上。\par

24. 那么，我的弟兄们，你们看，一切都在反对他们：所有神圣的书页，所有先知话，整部\WorkTitle{福音}，所有\WorkTitle{宗徒书信}，鸽子的每一声叹息，然而他们不醒悟，还没有从睡眠中醒来。但如果我们就是鸽子，让我们叹息，让我们坚持，让我们盼望；天主的怜悯必与你们同在，使圣神的火在你们的纯朴中燃烧；他们将会来。不可绝望；祈祷，宣讲，去爱；主有能力成就一切。他们已经开始感到羞愧；许多人已感到羞愧，并且脸红；基督会帮助，使其余的人也感到羞愧。然而，我的弟兄们，至少让糠秕留在那里；让所有的麦子聚集在一起；让他们当中结出果实的，借着鸽子回到方舟。\par

25. 处处失败，他们如今又拿什么来指控我们，找不到话可说？他们夺去了我们的房屋，夺去了我们的田产。他们拿出遗嘱。\Quote{看，\Person{该乌斯·塞维乌斯}将一份产业捐赠给了\Person{福斯蒂诺斯}主持的教会。}  \Person{福斯蒂诺斯}是哪个教会的主教？这教会是什么？他说，\Quote{捐赠给了\Person{福斯蒂诺斯}主持的教会。} 但\Person{福斯蒂诺斯}主持的不是教会，而是教派。然而，鸽子就是教会。为何呼喊？我们没有吞占房屋；让鸽子拥有它们。应当查问谁是鸽子，让她拥有它们。因为你们知道，我的弟兄们，那些房屋不属于\Person{奥古斯定}；如果你们不知道，且以为我喜欢拥有它们，天主知道，是的，祂知道我对那些田产的心思，甚至我在那事上所受的苦；祂知道我的呻吟，因为祂屈尊将鸽子的一点分赐给我。看，那些田产在那里；你们凭什么权利主张拥有它们？凭神权，还是人权？让他们回答：我们在圣经中有神权，在君主的法律中有人权。每个人拥有他所拥有的，凭什么权利？岂不是凭人权？因为凭神权，\BibleQuote{大地属于主，和其中的一切。} 穷人和富人都是天主用同一泥土造的；同一土地同样供养穷人和富人。然而，凭人权，某人说，这块田产是我的，这栋房屋是我的，这个仆人是我的。因此，凭人权，就是凭皇帝们的权利。为什么这样？因为天主正是通过这世界的皇帝和君王将这些人权分给了人类。你希望我们宣读皇帝的法律，并依照这些法律来处置田产吗？如果你们要凭人权拥有产业，就让我们宣读皇帝的法律；让我们看看他们是否允许异端者拥有什么了。但你说，皇帝与我何干？你拥有土地正是凭从皇帝来的权利。或者取消由皇帝创制的权利，然后谁敢说，那块田产是我的，或那个奴隶是我的，或这栋房子是我的？然而，为了使他们拥有这些东西，人们接受了来自君王的权利，你难道要我们宣读法律，好让你因即使拥有一个花园而高兴，并将其归因于鸽子的仁慈，使你甚至能在那里继续拥有产业吗？因为有一些众所周知的法律可以宣读，其中皇帝们指示：那些在天主教会共融之外、擅自冒用基督徒之名、且不愿在平安中敬拜平安之作者的人，不得以教会名义拥有任何东西。\par

26. 但我们与皇帝有何相干？不过我已经说过，我们正在讨论人权。然而宗徒要我们服从君王，要我们尊敬君王，并说：\BibleQuote{尊敬君王。}\BibleRef{伯前 2:17} 不要说，我与君王何干？就像那样，你与产业何干？产业是通过从君王而来的权利所享有的。你说，我与君王何干？那就不要说产业是你的；因为你已将它们归之于那些同一的人权，人凭它们享有产业。但他说，我所关心的是神权。好，让我们读福音；让我们看看鸽子降临其上的基督的天主教会扩展到多远，这鸽子教导说：\BibleQuote{这一位就是那施洗的。} 那么，一个人怎能凭神权拥有产业，若他说：\Quote{我施洗。} 而鸽子说：\BibleQuote{这一位就是那施洗的。} 而圣经说：\BibleQuote{我的鸽子是唯一的，她母亲独有的。} 你为什么撕裂鸽子？——不，而是撕裂了你自己的肠子，因为当你们自己被撕碎时，鸽子保持完整。所以，我的弟兄们，如果他们在每一点上都被击退，无话可说，我将告诉他们该做什么；让他们来到天主教会，并与我们一同拥有大地，以及那位造天和造地者。\par

\SourceRecord{Translated by John Gibb.；Nicene and Post-Nicene Fathers, First Series；Vol. 7.；Edited by Philip Schaff.；Buffalo, NY: Christian Literature Publishing Co.,；1888.；<http://www.newadvent.org/fathers/1701006.htm>.}

% structure: p0001 source=h1 target=section reason=page-title

\section{讲道集第七篇（若 1:34-51）}

1. 我们因你们的人数而欢欣，因为你们怀着热忱而来，人数之多超出我们的期望。这正是一切劳苦与危险中的喜乐与安慰：你们对天主的爱、虔诚的热忱、确切的望德和心灵的热火。你们在诵读圣咏时听见说：\Quote{穷苦和匮乏的人在世上向天主呼号。}这是那声音——正如你们常听到并当记住的——不是一个人的，却也是一个人的；不是一个人的，因为信友们众多——如同遍及全世界的糠秕中许多谷粒呻吟——却也是一个人的，因为他们都是基督的肢体，同为一个身体。这个贫穷匮乏的子民，不知与世俗同乐；他们的哀伤在内里，喜乐也在内里，无人看见，惟有倾听那呻吟者、加冕那仰望者的一位。世俗的欢乐是虚妄的。它被满怀期待地盼望，但到来时却不能持守。例如，今天在这个城市是失丧者欢乐的日子，明天自然就过去了；他们自己明天也不再是今天的样子。一切事物都过去、飞逝、如烟消散；爱慕这些的人有祸了！因为每个灵魂都追随它所爱的。\BibleQuote{凡有血肉的都如草，一切荣华都如田野的花：草会枯萎，花会凋谢；但主的话永远常存。}\BibleRef{依 40:1}看哪，你们若渴望永远存留，就该爱什么。但你们可能回答说：我如何能把握天主的话？\BibleQuote{圣言成了血肉，居住在我们中间。}\BibleRef{若 1:14}\par

2. 因此，亲爱的诸位，让我们的匮乏与贫穷成为对那些自以为富足的人的悲伤。因为他们的欢乐如同疯子的欢乐。正如疯子多半在疯狂中欢乐、发笑，却为清醒的人悲伤，同样，我们——若已接受了来自天上的良药，因为我们都是疯子，如今仿佛痊愈了，因为我们不再爱那些所爱过的——让我们向天主为那些仍在疯狂中的人呻吟，因为祂也能拯救他们。他们需要观看自己，并对自己不满；他们渴望观看，却不知道观看自己。只要他们稍微转目注视自己，就会看见自己的混乱。但在那之前，让我们的追求不同，让灵魂的娱乐不同；我们的悲伤比他们的欢乐更有价值。就弟兄的人数而言，难以想象有人会被那庆典所迷惑；但就姊妹的人数而言，我们感到悲伤，这更是悲伤的原因，因为她们不更倾向于教会，至少羞耻本应阻止她们出现在公共场合。愿那看见的祂看顾这事；愿祂的仁慈临在医治一切。让我们这些聚在一起的人饱飨天主的宴席，让祂的话成为我们的欢乐。因为祂已邀请我们到祂的福音中，祂是我们的食粮，没有比祂更甘甜的，只要人心中有健康的味觉。\par

3. 然而，亲爱的弟兄们，我想你们记得这部福音是按适当的段落依次诵读的；我认为你们没有忘记最近所讲的，尤其是关于若翰和鸽子的事。即关于若翰：他藉着鸽子学到了关于主的新知识，尽管他已认识主。藉着天主圣神的感动，我们得知，若翰确实已经认识主，但主自己要亲自授洗，祂不会将施洗的能力从自己转给任何人，这是他藉着鸽子学到的，因为曾对他说：\BibleQuote{你看见圣神降下，如同鸽子，停留在祂身上，这就是那要以圣神施洗的那位。}\BibleRef{若 1:33}\Quote{这就是那位}是什么意思？不是另一位，虽然是藉着另一位。但为什么要用鸽子呢？已经说了许多，我不能也不必重述一切——但主要是为了表示和平，因为那些在外受洗的树木，鸽子在它们身上找到果子，便带回方舟，正如你们记得诺厄从方舟放出的鸽子，它漂浮在洪水之上，被圣洗洗净，却没有沉没。当它被放出时，衔回一根橄榄枝；不只是叶子，还有果子。\BibleRef{创 8:8}因此，我们应当为那些在外面受洗的弟兄盼望这事：愿他们有果子；鸽子不会让他们留在外面，而会带他们回方舟。因为一切果子的总和就是爱德，没有爱德，一个人无论有什么都算不得什么。这是宗徒详尽所说的，我们已提及并重述。因为他说：\BibleQuote{我若说人间的语言和天使的语言，却没有爱德，我就成了鸣的锣、响的钹；我若有全知、明了一切奥秘和一切知识，有全备的信德——但他说全备的信德是什么意思？——甚至能移山，却没有爱德，我就算不得什么。我若把一切财产分给穷人，并舍身被火焚烧，却没有爱德，为我毫无益处。}\BibleRef{格前 13:1}但他们绝不能说他们拥有爱，因为他们分裂了合一。这些事已经说了，让我们看下文。\par

4. 若翰因看见了而作证。他作了什么证？\Quote{这就是天主子。}因此应当由祂授洗，祂是天主的独生子，而非义子。义子是独生子的仆役；独生子有权柄，义子只有服务。如果一个不属于子女行列的仆役——因为他生活邪恶、行为败坏——施行了圣洗，我们的安慰是什么？\Quote{这就是那位施洗的。}\par

5. \BibleQuote{第二天，若翰和他的两个门徒站在那里；他注视着行走的耶稣，说：看，天主的羔羊！} 的确，以一种特殊的意义，羔羊；因为门徒也被称为羔羊：\BibleQuote{看，我派遣你们如同羔羊进入狼群。}\BibleRef{玛 10:16} 他们也被称为光：\BibleQuote{你们是世界的光；}\BibleRef{玛 5:14} 但祂以另一种意义被称为光，关于祂曾说过：\BibleQuote{那才是照亮每人的真光，光照每一个来到世界的人。}\BibleRef{若 1:9} 同样，祂以特殊意义被称为鸽子，独无玷污，无罪恶；不是一位已洗去罪恶的，而是一位从未有玷污的。为什么？因为若翰曾论及主说：\BibleQuote{看，天主的羔羊。} 难道若翰自己不也是一只羔羊吗？他不是一位圣人吗？他不是新郎的朋友吗？因此，若翰以特殊意义论祂说：\BibleQuote{这是天主的羔羊。} 因为唯独藉这羔羊的血，人才能被救赎。\par

6. 我的弟兄们，如果我们承认我们的价值，就是羔羊的血，那么，如今那些庆祝不知名女人的血之节日的人，是多么忘恩负义啊！他们说，金子从女人的耳朵上被夺走，血流了出来，金子被放在秤或天平上，血这边占了很大优势。如果女人的血足以重过金子，那么那藉祂而世界被造的羔羊的血所具有的胜过世界的力量又有多大呢？的确，那不知名的神祇曾被血安抚，以压下重量。不洁的神祇知道耶稣基督要来，他们从天使那里听到祂的来临，从先知那里听到，并且期待祂。因为如果他们不期待，为什么他们喊道：\BibleQuote{我们与祢有什么关系？祢是来在时期以前毁灭我们吗？我们知道祢是谁，祢是天主的圣者。}\BibleRef{谷 1:24} 他们期待祂会来，但不知道时期。但你从圣咏中关于耶路撒冷听到了什么？\BibleQuote{因为祢的仆人们喜爱她的石头，并怜悯她的尘土。}他说，\BibleQuote{祢要起来，并怜悯熙雍，因为怜悯她的时期已来到。} 当天主怜悯的时期来到时，羔羊来了。什么样的羔羊连狼都惧怕？什么样的羔羊被宰杀时却杀死了狮子？因为魔鬼被称为狮子，四处巡游，咆哮，寻找可吞噬的人。\BibleRef{伯前 5:8} 藉着羔羊的血，狮子被征服了。看，基督徒的戏剧。而且还有：他们以肉身的眼睛注视虚妄，我们以心灵的眼睛注视真理。弟兄们，不要以为我们的主天主没有给我们留下戏剧；因为如果没有戏剧，你们今天为什么聚集在一起？看，我们所说的话你们看见了，你们就欢呼；如果你们没有看见，你们就不会欢呼。这是一件在全世界都值得看的大事：狮子被羔羊的血征服；基督的肢体从狮子的牙齿下被救出，并被结合于基督的身体。因此，某个神祇伪造了这赝品，要他的像用血来买，因为他知道人类终将要用宝贵的血来救赎。因为恶神为自己伪造了一些荣誉的影子，以欺骗那些跟随基督的人。我的弟兄们，甚至那些用护身符、咒语、敌人的诡计来诱惑人的人，将基督的名字混入他们的咒语中；因为他们现在不能诱惑基督徒，给他们毒药时会加入一些蜂蜜，使甜掩盖苦，喝下去导致毁灭。以至于我知道那\Person{皮勒阿图斯}的司祭有时习惯说，\Person{皮勒阿图斯}本人也是基督徒。为什么这样呢，弟兄们？除非他们不能以其他方式诱惑基督徒。\par

7. 那么，不要到别处去寻找基督，只应在基督愿意向你们宣讲祂自己的地方去寻找；并要按照祂愿意向你们宣讲的方式，紧抓住祂，将祂铭刻在你们心中。这是抵御敌人一切攻击和一切陷阱的壁垒。不要害怕，除非得到许可，他不会试探；可以肯定，除非得到许可或被派遣，他什么也不做。他像一个恶天使，被控制他的力量派遣；当他请求某事时，他就被允许。弟兄们，这事发生，只是为了义人受考验，恶人受惩罚。那么，你为何害怕呢？行走在上主你的天主内；要确信，祂不愿你受的苦，你决不会受；祂允许你受的苦，是惩戒者的鞭策，而非定罪者的惩罚。我们正被教育以承受永恒的产业，难道我们还拒绝被鞭策吗？我的弟兄们，如果一个孩子拒绝接受父亲责打的掌击或鞭打，难道他不会被称为骄傲、不可救药、对父亲的管教忘恩负义吗？地上的父亲教育他的儿子是为了什么？为了使他不失去为他所获得、所积攒的暂时之物，这些是他不愿儿子失去的，而离开的人不能永远保留它们。他教导一个儿子，不是要与他一同拥有，而是为了让他以后拥有。我的弟兄们，如果一个父亲教导一个将要继承他的儿子，并且教导他必须经历这一切，就像那告诫他的人注定要经历一样，那么，你希望我们的天父如何教育我们呢？我们不是要继承祂，而是要到祂那里去，并与祂永远住在一份不腐朽、不死、不被风暴摧毁的产业中。祂自己既是产业，又是父亲。我们将拥有祂，难道不该接受训练吗？让我们聆听父亲的训诲。当我们头痛时，让我们不要求助于迷信的调解人、占卜者和虚妄的疗法。我的弟兄们，我岂不为你们哀哭吗？我天天发现这些事；我能做什么呢？我尚未说服基督徒们，他们的希望应放在天主身上。看，如果有人死了，他曾经接受了这些疗法中的一种（有多少人带着疗法死去，又有多少人没有疗法而活！），那么灵魂以何等的信心走向天主呢？他失去了基督的标记，接受了魔鬼的标记。也许他可以说他没有失去基督的标记。那么，你可以同时拥有基督的标记和魔鬼的标记。基督不愿共有所有权，祂愿意单独拥有祂所买来的。祂以如此高昂的代价买来，为要独自拥有；而你却使祂与你因罪将自己卖给的那个魔鬼成为合伙人。\Quote{祸哉，心怀二意的人！}\BibleRef{德 2:12}即那些在心中将一部分给天主，一部分给魔鬼的人。天主因魔鬼在那里占有一份而发怒，便离去，而魔鬼将占据全部。因此，宗徒所说并非徒然：\Quote{也不可给魔鬼留地步。}\BibleRef{弗 4:27}所以，弟兄们，让我们认识羔羊吧；让我们认识我们的赎价。\par

8. \Quote{若翰站着，他的两个门徒也在那里。}看，若翰的两个门徒：既然若翰，新郎的朋友，是那样的人，他不寻求自己的光荣，而是为真理作证。难道他愿意他的门徒留在他身边而不跟随主吗？相反，他亲自指示他的门徒应当跟随谁。因为他们把他当作是羔羊；而他说：\Quote{你们为什么注意我呢？我不是羔羊；看，天主的羔羊！}关于这羔羊，他也早已说过：\Quote{看，天主的羔羊。}天主的羔羊赐给我们什么益处呢？他说：\Quote{请看，除免世罪者！}那两位与若翰在一起的门徒听了这话，便跟随了祂。\par

9. 让我们看看接下来：\Quote{看，天主的羔羊！}若翰说了这话，那两个门徒听见他说话，便跟随了耶稣。耶稣转过身来，看见他们跟着，就问他们说：\Quote{你们找什么？}他们回答说：\Quote{辣彼！——意即师傅——祢住在哪里？}他们并非那样跟随祂，以至于立即依附祂；因为很明显，当他们依附祂时，是祂从船上召唤了他们。因为两人中的一位是安德肋，正如你们刚才所听到的，安德肋是伯多禄的兄弟；我们从福音中知道，主从船上召唤了伯多禄和安德肋，说：\Quote{来跟从我，我要使你们成为渔人的渔夫。}\BibleRef{玛 4:19}从那时起，他们依附了祂，不再离开。这一次，这两个人跟随了祂，并非不再离去，而是为了看看祂住在哪里，并应验圣经：\Quote{使你的脚磨损祂的门槛，起来不断到祂那里去，受教于祂的诫命。}\BibleRef{德 6:36}祂指示了他们祂所住的地方；他们去了，并与祂同住。多么有福的一天，多么有福的一夜！谁能告诉我们从主那里听到的事呢？让我们也在心中建造一座房屋，让祂进来教导我们，并与我们交谈。\par

10. \BibleQuote{你们找什么？}他们对祂说：\BibleQuote{辣彼（译出来就是师傅），祢住在哪里？祂对他们说：你们来看。他们去了，看见了祂住的地方，并且那一天就与祂同住；那时大约是第十时辰。}我们是否认为这与福音作者告诉我们是什么时辰毫无关系？他是否要我们对此不予以任何注意，不探求任何东西？那是第十时辰。这个数字表示法律，因为法律是在十条诫命中赐下的。但时候已到，法律应藉着爱来满全，因为犹太人不能因恐惧而满全法律。因此主说：\BibleQuote{我来不是为废除法律，而是为成全。}\BibleRef{玛 5:17} 因此，适当地，在第十时辰，这两个人藉着新郎的朋友的证言跟随了祂，并且主在第十时辰听到了\BibleQuote{辣彼（译出来就是师傅）。}如果主在第十时辰听到了辣彼，而第十的数字与法律相关，那么法律的师傅无非就是法律的颁布者。谁也不能说一个人颁布了法律，另一个人教导法律：因为那颁布法律的同样教导它；祂是自己法律的师傅，并教导它。仁慈在祂的舌头上，因此祂仁慈地教导法律，正如论及智慧时所说的：\BibleQuote{法律和仁慈在她的舌上。}\BibleRef{箴 31:26} 不要害怕你不能满全法律，快投奔仁慈。如果你不能满全法律，就利用那盟约，利用那担保，利用那位天上的、精通法律者为你所规定和拟定的祈祷。\par

11. 凡有案件并希望向皇帝祈求的人，都会找一位精通法律并在学校训练过的人，为他们撰写祈愿书；以免万一他们以不当的方式请求，不仅得不到所寻求的，反而得到惩罚而非恩惠。因此，当宗徒们想要恳求，却不知如何接近皇帝天主时，他们对基督说：\BibleQuote{主，请教给我们祈祷；}也就是说：\Quote{祢是我们精通法律的师傅，我们的陪审者，是的，天主的共座者，请为我们编制祈祷词。}于是主从天界法律的书卷教导他们，教他们如何祈祷；在祂所教导的当中，祂设定了一个条件：\BibleQuote{求祢宽恕我们的罪过，如同我们宽恕别人一样。}\BibleRef{路 11:1} 如果你不按着法律寻求，你就成了有罪的。你成了有罪的，难道不在皇帝面前颤抖吗？献上谦卑的祭献，献上仁慈的祭献；祈祷说：\Quote{宽恕我吧，因为我也宽恕。}但如果你说，就要做。因为你要做什么？如果你在自己的祈祷中说了谎，你要往哪里去？不像在法庭中所说的，你将要失去诏书的利益；但你将不会获得诏书。因为法庭的法律是：凡在他的祈愿中说谎的，将不能从他所获得的中得到任何利益。但这只是在人间，因为一个人可能被欺骗：皇帝可能被你欺骗，当你向他递交祈愿时；因为你说了你想要的，而你所对说话的人不知道它是真是假；他把你送到你的对手那里，如果可能的话去驳斥你，这样如果你在法官面前被证实说了假话（因为他无法不批准诏书，不知道你是否说谎），你将在你提交诏书的地方失去它的利益。但是天主，祂知道你是说谎还是说实话，祂不会让你在审判中失去利益，但祂不允许你获得它，因为你竟敢向真理说谎。\par

12. 那么，你将做什么？告诉我。在每一点上成全法律，以致在何事上也不得罪，是困难的；因此，罪恶的状况是确定的；你难道要拒绝使用补救吗？看，我的弟兄们，主为灵魂的病患预备了怎样的补救！那么，当你头痛时，我们称赞你，如果你把福音放在你的头上，而不是去找护身符。因为人的软弱已经发展到如此地步，那些求助于护身符的人的状况如此可悲，以致当我们看见一个人躺在床上，被热病和疼痛折磨，除了福音放在他的头旁之外，别无指望时，我们就欢喜；这并不是因为这样做是出于这个目的，而是因为福音被优先于护身符。如果，那么，将福音放在头旁以缓解头痛，难道不应当将它放在心里以医治它脱离罪吗？那么，就这样做吧。做什么？让它放在心里，让心得到医治。好——好，你不再关心身体的安全，只向天主祈求。如果祂知道这对你有利，祂会赐给你；如果祂不赐给你，你拥有了也无益。有多少人卧病在床，因此成为无辜！因为如果他们康复，就会出去行恶。健康对多少人是伤害！那个出去到窄路杀人的强盗，他生病该多好！那个夜间起来挖邻居墙的人，他发高烧该多好！如果他病了，就相对无辜；健康了，却成了罪恶的。天主知道什么对我们有益，让我们只努力做到这一点：使我们的心纯全无过犯；当身体受鞭打时，让我们向祂祈求缓解。宗徒保禄恳求主取走他肉体上的刺，主没有应允。他烦乱了吗？他充满了忧伤，并说自己被遗弃了吗？他反而说自己没有被遗弃，因为那他所希望被取走的没有被取走，好使软弱得以医治。因为他从医师的话中找到了这个：\BibleQuote{我的恩宠为你足够了，因为我的能力在软弱中才成全。}\BibleRef{格后 12:8} 那么，你怎么知道天主不愿意医治你呢？现在鞭打你对你还有益。你怎么知道医生用手术刀切割病患部位时，病情有多严重？他不知道尺度，该做什么，该做多少吗？他切割时，患者的喊叫能阻止医生按照其技艺切割的手吗？一个喊叫，另一个切割。那个不听患者喊叫的医生是残忍的吗？他岂不是仁慈地追踪伤口，好医治病人？我的弟兄们，我说这些话，是为了当我们在天主的管教下时，没有人寻求任何其他帮助，只寻求天主的帮助。小心不要灭亡；小心不要离开羔羊，而被狮子吞噬。\par

13. 那么，我们已经宣告了，为什么是在第十时辰。让我们看接下来是什么：\BibleQuote{听了若翰的话而跟随了他的那两个人中的一个，是西满伯多禄的兄弟安得肋。他先找到了自己的兄弟西满，并向他说：\InnerQuote{我们找到了默西亚。}——默西亚意即基督。} 默西亚，在希伯来语；基督，在希腊语；在拉丁语，是受傅者。\OriginalTerm{傅油}{\GreekText{Χρῖσμα}} 在希腊语是傅油；因此，基督就是受傅者。祂被特别地傅油，超越地傅油；所有基督徒都受傅，但祂是超越地受傅。听祂在圣咏中如何说：\BibleQuote{为此，天主，你的天主，以喜乐油傅了你，胜过你的同伴。} 因为所有圣者都是祂的同伴，但祂在特殊的意义上是至圣者，特别受傅，特别为基督。\par

14. \BibleQuote{他遂领他到耶稣跟前，耶稣注视着他，说：你是若望的儿子西满，你要叫刻法——意即伯多禄。} 主说伯多禄是谁的儿子，这并非什么大事。对主来说，什么是大事呢？祂知道祂所有圣人的名字，是祂在创世以前就预定了的。你们奇怪祂对一个人说：\Quote{你是这个人的儿子，你要被称为这个或那个}？祂把西满的名字改为伯多禄，这有什么了不起呢？伯多禄来自\OriginalTerm{盘石}{petra}，而盘石就是教会；所以，在伯多禄的名字中，教会就被预表了。除非那在盘石上建造的，谁能安全呢？主亲自说的什么？祂说：\BibleQuote{凡听了我的话而实行的，我要把他比作一个聪明人，他把自己的房屋建在盘石上}（他就不屈于诱惑）。\BibleQuote{雨淋，水冲，风吹，袭击那座房屋，它并不坍塌，因为根基是建在盘石上。可是，凡听了我的话而不实行的}（现在让我们各自恐惧警惕），\BibleQuote{我要把他比作一个愚昧人，他把房屋建在沙土上：雨淋，水冲，风吹，袭击那座房屋，它就坍塌了，且坍塌得很惨。}\BibleRef{玛 7:24} 那在沙土上建造的人，进入教会有什么益处呢？因为听而不行，他确实建造了，但建在沙土上。如果他不听，他什么也不建；但如果他听，他就建造。但我们问：建造在哪里？因为如果他听了又实行，他就建在盘石上；如果他听了却不实行，他就建在沙土上。有两种建造者：建在盘石上的，和建在沙土上的。那么，那些不听的人呢？他们安全吗？祂说他们安全，是因为他们不建造吗？他们在雨、风、洪水之下是赤裸的；当这些来到时，在房屋倒塌之前，他们自己就先被冲走了。唯一的保障，就是既建造，又建在盘石上。如果你愿意听而不行，你建造的是一座废墟；当诱惑来到时，它推倒房屋，连你一起冲走。但如果你不听，你就是赤裸的；你自己也被这些诱惑拖走。所以，听而实行，这是唯一的补救。今天，也许有多少人因听而不行，被这节日的洪流冲走了！因为听而不行，洪水来到，这每年的节日；急流满了，它将过去而干涸，但被它冲走的人是有祸的！亲爱的，你们要知道：除非人听而实行，他不是建在盘石上，他不属于主所如此称赞的那个伟大的名字。因为祂已引起了你们的注意。如果西满先前就已被称为伯多禄，你们就不会如此清楚地看到这盘石的奥迹，就会以为这称呼是偶然的，而非出于天主的眷顾；因此，天主愿意他先被称为别的名字，好让这名字的改变，向我们推荐圣事的真义。\par

15. \Quote{次日，祂愿意往加里肋亚去，遇见斐理伯，就对他说：\InnerQuote{你跟从我吧！}那斐理伯是安德肋和伯多禄城里的人。斐理伯找到纳塔乃耳}（斐理伯已被主召叫了）；\Quote{并给他说：\InnerQuote{我们找到了梅瑟在法律上所写，先知们所预报的那位，就是耶稣，若瑟的儿子。}} 他被称为祂母亲所许配的那人的儿子。因为祂是在圣母仍是童贞女时受孕并出生的，所有基督徒都从福音清楚地知道。斐理伯对纳塔乃耳说了这些，又添上了地点：\Quote{从纳匝肋来的。} 纳塔乃耳对他说：\Quote{从纳匝肋能出什么好事吗？} 这是什么意思，弟兄们？并非像有些人所读的，因为同样也常被读作：\Quote{从纳匝肋能出好事吗？} 因为斐理伯接下来说的话：\Quote{你来看！} 但斐理伯的话可以适合两种读法：无论你读作肯定的，\Quote{从纳匝肋能出什么好事吗？}对此斐理伯回答\Quote{你来看；}或者读作疑问的，把整句当作问句：\Quote{从纳匝肋能出好事吗？你来看。} 既然这样，无论以这种方式或那种方式读，后面的话并非不相融，我们就需要探究我们该采用哪种解释。\par

16. 纳塔乃耳是怎样的人，我们从下面的话可以证明。请听他是怎样的人；主亲自为他作证。主是伟大的，由若翰的见证为众人所知；有福的纳塔乃耳，由真理的见证为众人所知。因为主虽然没有得到若翰的见证，祂自己为自己作证，因为真理足以为自己作证。但因为人们不能接受真理，他们藉着一盏灯寻求真理，因此若翰被派遣来给他们指出主。请听主为纳塔乃耳作的见证：\Quote{纳塔乃耳对他说：\InnerQuote{从纳匝肋能出好事吗？}斐理伯对他说：\InnerQuote{你来看！}耶稣看见纳塔乃耳向祂走来，就指着他说：\InnerQuote{看，这确是一个以色列人，在他内毫无诡诈。}} 伟大的见证！关于安德肋、伯多禄、斐理伯，都没有说过关于纳塔乃耳所说的这句话：\Quote{看，这确是一个以色列人，在他内毫无诡诈。}\par

17. 那么，弟兄们，我们该怎么办？此人（指\Person{纳塔乃耳}）应居宗徒之首吗？不仅\Person{纳塔乃耳}未被列为宗徒之首，他甚至在十二人中既非中间也非末位，尽管天主圣子曾为他作了如是的证言，说：\Quote{看，这确是一个以色列人，在他内毫无诡诈。} 需要问原因吗？按主所暗示的，我们找到一个合理的理由。因为我们应当明白，\Person{纳塔乃耳}是一位博学且精通法律的人，正因如此，主不愿将他列于门徒之中，因为祂拣选了无学问的人，为要藉他们使世人羞愧。请听宗徒所说的话：\Quote{弟兄们！你们看看你们是怎样蒙召的：按肉眼来看，你们中有智慧的不多，有权势的不多，出身高贵的不多；但天主却拣选了世上愚拙的，为羞辱那有智慧的；拣选了世上软弱的，为羞辱那强壮的；拣选了世上卑贱的和受人轻视的，以及那些一无所有的，为推翻那有的。} \BibleRef{格前 1:20} 如果拣选了一位有学问的人，他或许会说，他之所以被拣选，是因为他的学识使他堪当被选。我们的主\Person{耶稣基督}，为折断骄傲者的颈项，并不藉渔夫寻求演说家，而是藉渔夫赢得了皇帝。\Person{西彼廉}作为演说家固然伟大，但在其之前，渔夫\Person{伯多禄}藉着他不仅使演说家，也使皇帝相信。起初，没有高贵的人被选，没有有学问的人被选，因为天主拣选了世上软弱的，为羞辱那强壮的。因此，此人心地伟大、毫无诡诈，唯独因此未被拣选，免得主被任何人视为拣选了有学问的人。正是由于他精通法律，当他听到\Quote{从纳匝肋}这话时——因他查考了圣经，知道救主应从那里期待，这是其他经师和法利塞人难以知道的——这位精通法律的人，听到\Person{斐理伯}说：\Quote{我们找到了那一位，梅瑟在法律上所记载、先知们所记载的，就是纳匝肋的耶稣，若瑟的儿子。}——这位极其精通圣经的人，一听到\Quote{纳匝肋}这名，便满怀希望，说：\Quote{从纳匝肋能出好事。}\par

18. 现在我们来看看关于此人的其余部分。\Quote{看，这确是一个以色列人，在他内毫无诡诈。} 什么叫\Quote{在他内毫无诡诈}？也许他没有罪？也许他没有病？也许他不需要医师？决不可能。没有人在此世出生到如此地步，以致不需要那医师。那么，\Quote{在他内毫无诡诈}这句话是什么意思？让我们在主的名下更仔细地探究一下——这马上就会显明。主说\OriginalTerm{诡诈}{Dolus}；每一位懂拉丁文的人都知道，诡诈是指一件事做出来而另一件事伪装的。请注意，亲爱的。诡诈并非\OriginalTerm{痛苦}{dolor}。我这样说是因为许多不熟练拉丁文的弟兄这样说话，他们说\Quote{Dolus}折磨他，把\OriginalTerm{诡诈}{Dolus}当作\OriginalTerm{痛苦}{dolor}使用。\Quote{诡诈}是欺骗，是虚伪。当一个人心中隐藏一件事，口里却说另一件事，这就是诡诈，他仿佛有两颗心；他仿佛有一个心的密室，他在那里看见真理，另有一个密室，他在那里构思谎言。为使你们知道这就是诡诈，在\BibleBook{}{圣咏}中说：\BibleQuote{诡诈的嘴唇。} 什么是\BibleQuote{诡诈的嘴唇}？接着又说：\BibleQuote{他们存着二心说邪恶的话。} 什么是\Quote{存着二心}，不就是双心吗？所以，如果诡诈不在\Person{纳塔乃耳}内，那么医师判断他是可治愈的，而非痊愈的。一个痊愈的人是一回事，一个可治愈的人是另一回事，一个无可救药的人又是另一回事：患病但非无可救药的人称为可治愈；患病且无可救药的人称为不可治愈；但已经痊愈的人不需要医师。因此，来治愈的医师看见他是可治愈的，因为他内里没有诡诈。如果他是个罪人，他内里怎会没有诡诈呢？他承认自己是个罪人。因为如果他是个罪人，却说自己是个义人，他口里就有了诡诈。所以，在\Person{纳塔乃耳}身上，祂称赞了对罪的承认，而非判断他无罪。\par

19. 因此，当那些自以为义的法利塞人责怪主——因为医生与病人同在一处——并且说：\Quote{看，他与谁一同吃饭？与税吏和罪人。}时，医生回答这些狂人：\Quote{健康的人不需要医生，而是有病的人；我来不是召义人，而是召罪人。}\BibleRef{玛 11:11}这就是说，因为你们自以为是义人，而你们是罪人；你们自以为健康，而你们正在衰败，你们便远离药物，不持守健康。因此，那位曾请主吃饭的法利塞人，在自己眼中是健康的；但那位患病的妇人冲进她未被邀请的房屋，被渴望健康的勇气所推动，没有接近主的头，也没有接近手，而是接近脚；用眼泪洗脚，用头发擦干，亲吻它们，用香膏涂抹——她作为罪人，与主的足迹和解了。坐在那里吃饭的法利塞人，自以为健康，却责备医生，心里说：\Quote{这人若是先知，必知道摸他的是谁。}他怀疑主不知道，因为主没有推开她以防止被不洁的手触摸；但主知道，他允许自己被触摸，好让触摸本身得以医治。主看见法利塞人的心，便讲了一个比喻：\Quote{有一个债主，有两个欠债人；一个欠五百银币，另一个欠五十；因为他们无力偿还，他就宽免了二者。哪一个更爱他？}他回答说：\Quote{我想，是那得宽免更多的。}主转向那妇人，对西满说：\Quote{你看见这妇人吗？我进了你的家，你没有给我水洗脚；但她用眼泪洗了我的脚，用头发擦干；你没有给我亲吻；她不住地亲吻我的脚；你没有给我抹油；她用香膏抹了我的脚。所以，我告诉你：她许多的罪都赦免了，因为她爱得多；但那得赦免少的，爱得也少。}\BibleRef{路 7:36}这就是说，你病得更重，却自以为健康；你以为赦免得少，其实你欠得更多。她做对了，因为她里面没有\OriginalTerm{诡诈}{guile}，配得药物。\Quote{里面没有诡诈}是什么意思？她承认自己的罪。主也在纳塔乃耳身上称赞这一点，说他里面没有诡诈；因为许多法利塞人虽罪孽深重，却自称义人，带着诡诈，这使他们无法得到医治。\par

20. 耶稣于是看见这个人，在他里面没有\OriginalTerm{诡诈}{guile}，便说：\Quote{看，这确是一个以色列人，在他内毫无诡诈。}纳塔乃耳对他说：\Quote{你从哪里认识我？}耶稣回答他说：\Quote{斐理伯叫你以前，当你还在无花果树下时，我就看见了你。}纳塔乃耳回答他说：\Quote{辣彼，你是天主子，你是以色列的君王。}纳塔乃耳可能从\Quote{当你还在无花果树下时，我就看见了你，在斐理伯叫你以前}这句话中领悟了某种伟大的事；因为他所说的话\Quote{你是天主子，你是以色列的君王}与后来伯多禄的话并无不同，那时主对他说：\Quote{约纳的儿子西满，你是有福的，因为不是血和肉启示了你，而是我在天之父。}在那里，他称他为磐石，并赞美教会在此信德上的坚强支持。这里，纳塔乃耳已经说：\Quote{你是天主子，你是以色列的君王。}为什么呢？因为有人对他说：\Quote{斐理伯叫你以前，当你还在无花果树下时，我就看见了你。}\par

21. 我们必须探究这无花果树是否有什么含义。请听，我的弟兄们。我们发现无花果树被诅咒，因为它只有叶子，没有果实。\BibleRef{玛 20:19}在人类之初，当亚当和厄娃犯罪后，他们用无花果树叶编了裙子。\BibleRef{创 3:7}所以无花果树叶象征罪恶。纳塔乃耳当时在无花果树下，如同在死亡的阴影之下。主看见了他，正是关于他，经上说：\Quote{坐在死亡阴影下的人，有光照耀他们。}\BibleRef{依 9:2}那么，对纳塔乃耳说了什么呢？你对我说：\Quote{纳塔乃耳，你从哪里认识我？}现在你对我说话，因为斐理伯叫了你。主已经知道那被宗徒所召的人属于他的教会。啊，教会！啊，以色列，在你内毫无诡诈的！如果你是毫无诡诈的以色列子民，你已经藉着宗徒认识了基督，正如纳塔乃耳藉着斐理伯认识了基督。但他的慈悲在你认识他以前就已看见你，那时你正卧在罪中。难道是我们先寻求基督，而不是他先寻求我们吗？难道是病人来找医生，而不是医生来找病人吗？那只迷失的羊，难道不是牧人撇下九十九只在荒野，去寻找它，找到后欢欢喜喜地扛在肩上吗？那块丢失的银币，难道不是那妇人点上灯，打扫整个屋子，直到找到它吗？她找到后，对邻居说：\Quote{与我同乐吧，因为我找到了那枚丢失的银币。}\BibleRef{路 15:4}同样，我们如同羊迷失，如同银币丢失；我们的牧人找到了羊，但他是去寻找羊；那妇人找到了银币，但她是去寻找银币。那妇人是谁？基督的肉身。那灯是什么？\Quote{我为我的基督预备了一盏灯。}因此，我们被寻找，好使我们被找到；既被找到，我们就说话。我们不要骄傲，因为在被找到之前，我们是迷失的，若不被寻找，我们早已失丧。那么，不要让我们所爱的人，并渴望将他们争取到公教会的和平中来的人对我们说：\Quote{你对我们有何意图？为什么你寻找我们，如果我们是有罪的人？}我们寻找你，正是为叫你不致丧亡；我们寻找你，因为我们曾被寻找；我们渴望找到你，因为我们已被找到。\par

22. 当纳塔乃耳说\Quote{祢从哪里认识我呢？}主就对他说：\Quote{在斐理伯叫你以前，当你还在无花果树下时，我已看见了你。}噢，你这毫无诡诈的以色列，不论你是谁；噢，你这因信德而生活的子民，在我藉我的宗徒召叫你以前，当你还在死亡的阴影下时，你并未看见我，我却看见了你。主于是对他说：\Quote{因为我告诉你说：我看见你在无花果树下，你就信了；你将要看见比这些更大的事。}这\Quote{比这些更大的事}是什么呢？祂又对他说：\Quote{我实实在在告诉你们：你们要看见天开了，天主的使者在人子身上，上去下来。}弟兄们，这比\Quote{在无花果树下我看见了你}更为重大。因为主在我们被召时使我们成义，比祂看见我们躺在死亡的阴影下更为重大。若我们仍留在祂看见我们的地方，对我们有什么益处呢？我们岂不仍躺在那里？这更大的事是什么？我们何时看见天使在人子身上上去下来呢？\par

23. 我之前已经谈过这些上去下来的天使；但恐怕你们已经忘记，我将简单重提，以唤醒你们的记忆。若我是引入而非回忆这个主题，就需要更多言辞。雅各伯在梦中看见一架梯子，并在梯子上看见天使上去下来；他把他所枕的石头立作柱子，浇油在上面。\BibleRef{创 28:12}你们已经听说默西亚就是基督；你们已经听说基督就是受傅者。因为雅各伯立那石头，那浇过油的石头，并非为了来朝拜它；否则那就是偶像崇拜，而不是指出基督。所行的是对基督的预指，按照预指所应许的方式，而基督被指出来了。一块石头被傅油，但不是为偶像。一块石头被傅油；为何是石头？\Quote{看，我在熙雍安放一块石头，一块精选的、宝贵的基石；凡信赖祂的，决不会蒙羞。}为何傅油？因为\OriginalTerm{基督}{Christus}来自\OriginalTerm{傅油}{chrisma}。但他当时在梯子上看见了什么？天使上去下来。所以这就是教会，弟兄们：天主的使者就是善讲者，宣讲基督；这就是\Quote{在人子身上上去下来}的意义。他们如何上去，又如何下来？我们有一个例子；请听宗徒保禄。我们在他身上所见到的，让我们相信其他真理宣讲者也是如此。看保禄上去：\BibleQuote{我知道一个在基督内的人，十四年前被提到第三层天上去（或在身内，或在身外，我不知道：天主知道），并且他听到了不可言传的话，是人不能说出来的。}\BibleRef{格后 12:2}你们听见他上去，也听见他下来：\BibleQuote{我从前不能对你们像对属神的人说话，只能对你们像对属血肉的人，像对基督内的婴孩一样；我给你们喝奶，而不是吃肉。}\BibleRef{格前 3:1}看，那上去的人下来了。请问他是否上到第三层天？请问他是否下来给婴孩喂奶？听他说他下来了：\BibleQuote{我在你们中间成了小孩，好像乳母抚育自己的孩子。}\BibleRef{得前 2:7}因为我们看见乳母和母亲都俯就婴孩，虽然她们能说拉丁语，却缩短话语，以某种方式摇晃舌头，以便从规范的语言中编出幼稚的昵语；因为若他们按规则说话，婴儿听不懂，也不会受益。若有一位很会讲话的父亲，且是那样一位演说家以致法庭回响着他的雄辩、审判席为之震动，若他有一个小儿子，回家时他便放下那曾上去的法庭雄辩，而以童言俯就他的小孩。请在同一个地方听宗徒在同一句话中上去和下来：\BibleQuote{因为我们或疯狂，那是为了天主；或清醒，那是为了你们。}\BibleRef{格后 5:13}\Quote{我们疯狂}是什么意思？即我们看见了那些人不该说出来的事。\Quote{我们为了你们而清醒}是什么意思？\Quote{我曾决定在你们中间什么也不知道，只知道耶稣基督，并祂被钉在十字架上。}若主自己上去下来，显然祂的宣讲者以效法上去，以宣讲下来。\par

24. 若我们比平时更长久地留住你们，是为了让危险的时刻过去：我们猜想那些人现在已经结束了他们的虚妄。但是，弟兄们，我们既享用了救恩的盛宴，就当完成剩下的事，好使我们能以属灵的喜乐虔敬地充满主日，并将真理的喜乐与虚妄的喜乐相比较；我们若感到惊骇，就让我们哀恸；若哀恸，就让我们祈祷；若祈祷，就使我们蒙垂听；若蒙垂听，我们也赢得他们。\par

\SourceRecord{Translated by John Gibb.；Nicene and Post-Nicene Fathers, First Series；Vol. 7.；Edited by Philip Schaff.；Buffalo, NY: Christian Literature Publishing Co.,；1888.；<http://www.newadvent.org/fathers/1701007.htm>.}

% structure: p0001 source=h1 target=section reason=page-title

\section{论题八（若 2:1-4）}

1. 我们的主耶稣基督所行的使水变成酒的奇迹，对于那些知道这是天主作为的人而言，并不希奇。因为那位在婚宴上、在那六口石缸里（祂命令将水注满）造酒的主，每年在葡萄树上所做的也是同样的事。正如仆人倒进水缸的水因主的作为变成了酒，同样，云层倾注的雨水也因同一位主的作为变成了酒。但我们并不惊异于后者，因为它每年发生；因着一再重复，它失去了神奇性。然而，它比那水缸里所行的事更引人深思。因为谁若细察天主管辖和治理这整个世界的事工，岂能不因这些奇迹而惊愕和折服？倘若他思考任何一颗种籽那强韧的力量，那便是伟大之事，令他肃然起敬。但由于人们专注于其他的事，丧失了以日常赞美天主造物主的思考，天主便将某些超凡的行动保留给自己，好借这些奇迹震惊他们，使他们仿佛从沉睡中醒来，敬拜祂。死人复活了，人们惊叹；每天有那么多人生下来，却无人惊奇。如果我们更加审慎地反省，使一个未曾存在的人存在，比使一个已存在的人复活更值得惊奇。然而，同一天主，我们的主耶稣基督的圣父，借着祂的圣言成就这一切；祂既创造，也治理。先前的奇迹，祂借着与祂同在的圣言——天主——而行；后来的奇迹，祂借着同一圣言降生成人，并为我们成为人而行。正如我们因那人耶稣所行的事而惊叹，也让我们因耶稣天主所行的事而惊叹。诸天、大地、海洋、天上的一切装饰、地中丰盈的宝藏、海中的丰饶——所有我们眼前可及之物——都是耶稣天主所造。我们观看这些事物，若圣神在我们内，它们就如此使我们喜悦，以致我们赞美那设计者；但并非如此——我们转向受造物而远离造物主，仿佛面朝被造物，背朝造物主。\par

2. 这些事我们确实看见，就在我们眼前。但那些我们看不见的，如天使、异能、德能、宰制者，以及这超乎诸天、超出我们视野的整个宇宙中的一切居民，又当如何？然而，天使也常常在必要时向人显现。难道天主不也是借着祂的圣言，即祂的独生子、我们的主耶稣基督，创造了这一切吗？人的灵魂本身，虽不可见，却借着它在肉身中的作为，使那些恰当反思的人大为惊叹——它若不是由天主所造，又能由谁所造？若不是通过天主子，又能通过谁呢？暂且不提人的灵魂：任何禽兽的灵魂，看它如何调节那庞大的躯体，发挥感官——眼睛看，耳朵听，鼻子嗅，味觉分辨滋味 — 简言之，各肢体执行各自的职能！难道做这些的是身体，而不是灵魂——身内的居住者吗？灵魂不为肉眼所见，却借着这些行动激发惊叹。现在请你思量人的灵魂，天主赐予它理智，以认识造物主，辨别好坏，即是非：看它借着身体行多少事！观察这整个按照同一人类社会秩序安排的世界——有怎样的行政、怎样有序的权威等级、怎样的公民身份、怎样的法律、习俗、技艺！这一切都是由灵魂所成就，而这灵魂的能力却不可见。当灵魂离开身体时，身体只是一具尸体：首先，它仿佛凭借灵魂像一种防腐剂，身体得以暂时免于腐烂。因为所有的肉体都是可朽的，除非被灵魂像调味品一样保存，就会腐败。但人的灵魂与禽兽的灵魂共有此性质；更值得赞叹的是我已提及的那些属于心灵和理智的性质，在其中它也按造物主的肖像得以更新，人是按天主的肖像被造的。（\BibleRef{哥 3:10}）当这身体穿上不朽，这必死的穿上不死的时，灵魂的这份能力将是什么？（\BibleRef{格前 15:54}）如果它在可朽的肉体中运作就有这样的能力，那么在复活后借着属灵的身体又会有怎样的能力？然而这灵魂，如我所说，有着可赞叹的本性和实体，是不可见的、理智的；这灵魂也是由天主耶稣所造，因为祂是天主的圣言。\Quote{万物是借着祂造的；凡受造的，没有一样不是借着祂造的。}\par

3. 因此，当我们看见耶稣天主行了如此的事迹，为何要惊叹于人耶稣将水变成酒呢？因为祂成为人，并没有失去祂是天主。人被加于祂，天主并未失落于祂。这奇迹是由同一位行了一切的那位而行。因此，我们不要因天主行了这事而惊奇，而要爱祂，因为祂在我们的中间行了这事，并为了使我们恢复。因为祂通过这事例本身给了我们某些暗示。我想，祂来到婚宴不是没有原因的。除了奇迹本身，这事本身还含有某种奥秘和圣事性的意义。让我们敲门，使祂为我们开启，并以无形的酒充满我们：因为我们曾是水，祂使我们成为酒，使我们成为有智慧的；因为祂赐给我们祂信德的智慧，而我们之前是愚昧的。这智慧或许在于：连同对天主的尊敬、对祂威严的赞美、对祂最具权能的仁慈的爱德，我们得以理解在这奇迹中所行的事。\par

4. 主受邀请，前来参加婚宴。他来到这屋子参加婚宴，有何希奇呢？因为他来到这世界，正是来参加婚宴。的确，他若不来参加婚宴，他在这里就没有新娘。但宗徒说什么？\Quote{我曾把你们许配给一个丈夫，把你们当作贞洁的童女献给基督。}为什么他害怕基督新娘的童贞会被魔鬼的诡计败坏？他说：\Quote{我害怕，}他说，\Quote{免得你们的心被败坏，失去对基督的单纯和贞洁，就如蛇用诡计诱惑了厄娃一样。}\BibleRef{格后 11:3} 因此，他在这里有一位新娘，是他用自己的血赎回的，并赐予圣神作为抵押。他将她从魔鬼的奴役中释放：他为她的罪而死，并为了她的成义而复活。\BibleRef{罗 4:25} 谁能为自己的新娘做出这样的奉献？男人可以献给新娘各种地上的装饰——金银、宝石、房屋、奴隶、地产、农场——但谁会献出自己的血？因为如果一个人把自己的血献给新娘，他就不能活着娶她为妻。但主无惧死亡，为她献出自己的血，他复活后将要娶她为妻，他早已在童贞女的子宫中与她结合。因为圣言是新郎，人性是新妇；两者合而为一，天主子，同时也是人子。童贞玛利亚的子宫，他在其中成为教会的头，就是他的洞房：他从那里出来，如新郎从洞房出来，正如圣经预言：\Quote{又如巨人，欢然奔路。}他从洞房出来，如同新郎；受邀请前来参加婚宴。\par

5. 他之所以似乎不承认他的母亲（她从她怀中作为新郎出来），是因为一个确凿无疑的奥秘，他对她说：\Quote{女人，这于我和你有什么关系？我的时辰还没有到。}这是什么意思？他前来参加婚宴，难道是为了教导人轻慢自己的母亲吗？确实，他所出席的那场婚宴的主人，是为了生子而娶妻，他也希望自己生养的儿女尊敬他：那么，耶稣前来参加婚宴，难道是为了羞辱自己的母亲吗？而婚宴的目的正是为了生子，而天主命令子女要孝敬父母？毫无疑问，弟兄们，这里隐藏着某种奥秘。此事如此重要，以至于有些人——宗徒正如我们之前所提醒的，曾告诫我们要提防他们说：\Quote{我害怕，免得你们的心被败坏，失去对基督的单纯和贞洁，就如蛇用诡计诱惑了厄娃一样，}——他们削弱福音的可信度，断言耶稣并非由童贞玛利亚所生，并试图从这段话中找出支持他们错误的论据，以至于说：他对她说\Quote{女人，这于我和你有什么关系？}她怎么可能是他的母亲？为此，我们必须回答他们，并向他们说明主为何说这话，免得他们在疯狂中自以为发现了有损于健康信仰的东西，从而使童贞新娘的贞洁受损，也就是说，使教会的信仰受损。因为实际上，弟兄们，那些宁愿相信谎言而不相信真理的人，他们的信仰已经败坏。这些人表面上尊敬基督，却否认他具有肉身，这无异于宣称他是说谎者。那些在人心建立谎言的人，难道不是把真理从他们心中赶出去吗？他们让魔鬼进来，把基督赶出去；他们引进奸夫，把新郎关在门外，显然是蛇的伴郎，甚至可以说是蛇的拉皮条者。因为他们说话的目的，正是为了让蛇占据，让基督被排除。蛇如何占据？当谎言占据时。当虚谎占据时，蛇就占据；当真理占据时，基督就占据。因为他亲自说过：\Quote{我是真理；}\BibleRef{若 14:6} 至于那另一个，他说：\Quote{他不站在真理中，因为在他内没有真理。}\BibleRef{若 8:44} 基督就是真理，你们应当承认在他内一切都是真实的：真实的圣言，与圣父平等的天主，真实的灵魂，真实的肉身，真实的人，真实的天主，真实的诞生，真实的受难，真实的死亡，真实的复活。如果你说其中任何一样是虚假的，腐烂就进入，谎言的蛆虫就由蛇的毒液滋生，什么都不会存留。\par

6. 那么，有人问，主所说的\BibleQuote{女人，这于我和你有什么关系？}是什么意思？主可能在后续经文中解释了祂为什么说这话：祂说：\BibleQuote{我的时辰，}祂说，\BibleQuote{还没有到。}因为祂这样说：\BibleQuote{女人，这于我和你有什么关系？我的时辰还没有到。}我们必须寻求知道为什么说这话。但首先让我们藉此抵制异端者。古蛇说什么？那从古时以嘶嘶声注入毒液者说什么？他说什么？他说耶稣没有女人作祂的母亲。你从哪里证明这点？他说，从耶稣说了\BibleQuote{女人，这于我和你有什么关系？}而证明。谁记述了这事，使我们该相信耶稣说了这话？谁记述的？除了福音作者\Person{若望}，没有别人。但同一位福音作者\Person{若望}说：\BibleQuote{耶稣的母亲也在那里。}因为他这样告诉我们：\BibleQuote{第三天，在加里肋亚的加纳有婚宴，耶稣的母亲也在那里。耶稣和祂的门徒也被请去赴席。}我们有福音作者说出的两句话。福音作者说：\BibleQuote{耶稣的母亲在那里}；也正是同一位福音作者告诉了我们耶稣对祂母亲说了什么。弟兄们，请看，他是如何告诉我们耶稣回答祂的\Emph{母亲}的，他先说了\BibleQuote{祂的母亲对祂说}，好使你们能保持内心的童贞，不受蛇舌的伤害。这里在同一位福音、同一位福音作者的记录中，我们被告知：\BibleQuote{耶稣的母亲在那里}，以及\BibleQuote{祂的母亲对祂说}。谁记述了这些？福音作者\Person{若望}。耶稣回答祂的母亲说了什么？\BibleQuote{女人，这于我和你有什么关系？}谁记述了这事？正是同一位福音作者\Person{若望}。啊，最忠实与讲真话的福音作者，你告诉我耶稣说\BibleQuote{女人，这于我和你有什么关系？}为什么你加上了祂的母亲，而祂却并不承认她？因为你曾说过\BibleQuote{耶稣的母亲在那里}，以及\BibleQuote{祂的母亲对祂说}；为什么你不如说，玛利亚在那里，玛利亚对祂说。你叙述了这两件事：\BibleQuote{祂的母亲对祂说}，以及\BibleQuote{耶稣回答她说：女人，这于我和你有什么关系？}你为什么这样做？岂不因为两者都是真的吗？现在，那些人愿意相信福音作者在一处，即他告诉我们耶稣对祂母亲说\BibleQuote{女人，这于我和你有什么关系？}的同时，却不愿相信他在另一处，即他说\BibleQuote{耶稣的母亲在那里}，以及\BibleQuote{祂的母亲对祂说}。但谁能抵抗那蛇并持守真理，使心灵的童贞不被魔鬼的诡计所败坏？是那相信两者皆真的人，即相信耶稣的母亲在那里，并且耶稣那样回答祂的母亲。但如果他尚不明白耶稣是在何种意义上说\BibleQuote{女人，这于我和你有什么关系？}那就让他暂且相信祂说了这话，而且是对祂的母亲说的。让他先有虔诚的信德，然后他才会在理解上结出果实。\par

7. 我问你们，忠信的基督徒们，耶稣的母亲在那里吗？你们答：她在。你们从哪里知道？答：福音如此说。耶稣回答祂的母亲说了什么？你们答：\BibleQuote{女人，这于我和你有什么关系？我的时辰还没有到。}你们从哪里知道这事？答：福音如此说。如果你们渴望为新郎保持贞洁的童贞，就不容任何人败坏你们这信德。但如果有人问你们，祂为什么这样回答祂的母亲，让明白的人说明；但尚未明白的人，要最坚定地相信耶稣作了这个回答，而且是向祂的母亲作的。藉此虔诚，他必能藉祈祷叩击真理之门，不以争辩接近，从而也学会明白耶稣为何如此回答。只是这一点：当他自以为知道，或因不知道而羞愧的时候，要谨慎，免得被迫相信福音作者在说\BibleQuote{耶稣的母亲在那里}时说了谎，或者耶稣自己为我们的罪受了虚假的死，并为我们的成义显示了虚假的伤痕；以及祂说\BibleQuote{你们如果固守我的话，就确是我的门徒，也会认识真理，而真理必会使你们获得自由}时说了谎。\BibleRef{若 8:31}因为如果祂有虚假的母亲、虚假的肉体、虚假的死、死时虚假的伤、复活时虚假的痕，那么使相信祂的人得自由的就不是真理，而是谎谬了。相反，让谎谬向真理屈服，让那些想要被认为讲真话的人蒙羞，因为他们试图证明基督是个欺骗者，并且当别人说\Quote{我们不信你，因为你撒谎}时，他们却断言真理本身说了谎。然而，如果我们问他们：你们从哪里知道基督说了\BibleQuote{女人，这于我和你有什么关系？}他们回答说他们相信福音。那么，当福音说\BibleQuote{耶稣的母亲在那里}和\BibleQuote{祂的母亲对祂说}时，他们为什么不信？如果福音在此处说谎，我们怎能相信它在那里，即耶稣说了\BibleQuote{女人，这于我和你有什么关系？}这些可怜的人为什么不忠信地相信主确实这样回答，不是对陌生人，而是对祂的母亲？并且虔诚地寻求知道祂为什么这样回答？相差甚远：一个人说：\Quote{我想知道为什么基督这样回答祂的母亲}，另一个人说：\Quote{我知道基督不是对祂的母亲这样回答}。愿意理解封闭的事是一回事，不愿意相信敞开的事是另一回事。说\Quote{我想知道为什么基督这样回答祂的母亲}的人，希望他所信的福音为他开启；但说\Quote{我知道基督不是对祂的母亲这样回答}的人，正是指责他所信其中记载基督这样回答的福音为虚假。\par

8. 如今，各位弟兄，若你们认为合适，既已驳斥了那些人——他们始终在自己盲目中徘徊，除非因谦卑而得痊愈——就让我们探究主为何如此回答祂的母亲。祂以非凡方式由圣父而生，无母；由母亲而生，无父；作为天主，祂无母；作为人，祂无父；在万世之前，祂无母；在时期的终结，祂无父。祂所说的话是对祂母亲的回答，因为\BibleQuote{耶稣的母亲在那里}，且\BibleQuote{祂的母亲对祂说}。这一切福音都记载了。我们正是在那里得知\BibleQuote{耶稣的母亲在那里}，也是在那里得知祂对她说：\BibleQuote{女人，这于我和你有什么关系？我的时辰尚未来到。}让我们相信全部；对于尚不明白的，让我们探求。首先要注意，免得如同摩尼教徒因主说\BibleQuote{女人，这于我和你有什么关系？}而找到他们谎言的借口，占星家们也因祂说\BibleQuote{我的时辰尚未来到}而找到他们欺骗的借口。如果主说这话是出于占星家的意思，那么我们就犯了亵渎罪，因我们焚烧了他们的书籍。但如果我们做得正确——正如宗徒时代所行的——（\BibleRef{宗 19:19}）那么主说\BibleQuote{我的时辰尚未来到}就不是按照他们的观念。因为那些空谈者和受骗的诱惑者说：你看，基督处于命运之下，正如祂说\BibleQuote{我的时辰尚未来到}。那么，我们应先回答谁——是异端者还是占星家？因为两者都出于蛇，企图败坏教会心中无玷信德所持守的童贞。让我们先回答我们所提出的那些人，其实在很大程度上我们已经回答了他们。但恐怕他们以为我们对主回答母亲的话无话可说，我们就进一步准备反驳他们；因为我认为已经说的足以驳斥他们。\par

9. 那么，为何圣子对母亲说：\BibleQuote{女人，这于我和你有什么关系？我的时辰尚未来到？}我们的主耶稣基督既是天主又是人。按祂是天主，祂没有母亲；按祂是人，祂有母亲。所以，她是祂的肉身的母亲，祂的人性的母亲，祂为我们而取的软弱的母亲。但祂将要行的奇迹，是按祂的天主性而行，而非按祂的软弱；是按祂为天主的身份，而非按祂出生为软弱的身份。然而，天主的软弱比人更强。（\BibleRef{格前 1:25}）于是祂的母亲向祂要求一个奇迹；而祂，将要行天主性的事工，如此不承认一个人类的子宫；实际上是在说：\Quote{在我内行奇迹的那一位并非由你所生，你没有生我的天主性；但因为我的软弱由你而生，当那同一软弱挂在十字架上时，我将承认你。}的确，这就是\BibleQuote{我的时辰尚未来到}的含义。因为那时祂才承认了，事实上祂一直知道。祂在预定中认识祂的母亲，甚至在祂由她出生之前；甚至作为天主，祂创造了那一位——祂作为人将由她所造——之前，祂已认识她是祂的母亲；但在某个时辰，出于奥秘，祂不承认她；而在某个尚未到来的时辰，同样出于奥秘，祂又承认她。因为那时祂承认她，正是她所生的那一位将要死去之时。那使玛利亚存在的，并未死去，而是那由玛利亚所取的死了；不是天主性的永恒，而是肉身的软弱正在死去。所以祂如此回答，在信徒的信德中区分了祂来的\Emph{谁}和\Emph{如何}。因为祂本是天主、天地的主，却经由一个女人作母亲而来。按祂是世界的主、天地的主，当然祂也是玛利亚的主；但按那所说\BibleQuote{生于女人，生于法律之下}，祂是玛利亚的儿子。同一位既是玛利亚的主，又是玛利亚的儿子；同一位既是玛利亚的创造者，又是从玛利亚所造的。不要惊奇祂既是儿子又是主。因为正如祂被称为玛利亚的儿子，同样祂也被称为达味的儿子；祂是达味的儿子，因为祂是玛利亚的儿子。请听宗徒明言：\BibleQuote{按肉身来说，祂是从达味的后裔生的。}（\BibleRef{罗 1:3}）也请听祂被称为达味的主；让达味自己宣告：\BibleQuote{上主对我主说：你坐在我右边。}耶稣亲自引用这段经文给犹太人，并用它驳斥了他们。（\BibleRef{玛 22:45}）那么，祂如何既是达味的儿子，又是达味的主？按肉身是达味的儿子，按天主性是达味的主；同样，按肉身是玛利亚的儿子，按威严是玛利亚的主。如今，因为她不是祂天主性的母亲，而祂要行的奇迹正是藉着天主性，所以祂回答她：\BibleQuote{女人，这于我和你有什么关系？}但不要以为我否认你是我的母亲：\BibleQuote{我的时辰尚未来到}；因为在那时辰我将承认你，当你作为母亲所生的软弱被挂在十字架上时。让我们证实这一点。当主受苦时，同一位福音作者告诉我们——他认识主的母亲，并在这婚宴上让我们认识她——同一位，我说，告诉我们：\BibleQuote{那里，靠近十字架，有耶稣的母亲；耶稣对祂的母亲说：女人，看，你的儿子！又对门徒说：看，你的母亲！}祂将母亲托付给门徒照顾；将母亲托付给他，仿佛祂将比她先死，且要在她死前复活。那人将她作为人托付给人照顾。玛利亚所生的正是这人性。那时辰业已来到，就是祂当时所说的\BibleQuote{我的时辰尚未来到}的那个时辰。\par

10. 依我之见，弟兄们，我们已经回答了异端者。现在让我们回答占星术士。他们如何试图证明耶稣受命运支配？因为他们说，祂亲自说过：\Quote{我的时辰还没有到。}因此我们相信祂；如果祂曾说\Quote{我没有时辰，}祂就会排除占星术士；但看，他们说，祂说：\Quote{我的时辰还没有到。}那么，如果祂曾说\Quote{我没有时辰，}占星术士就会被排除，没有任何毁谤的根据；但现在祂说了\Quote{我的时辰还没有到，}我们怎能反驳祂自己的话？真是奇妙，占星术士因相信基督的话，竟试图说服基督徒相信基督生活在命运的时辰之下。好吧，让他们相信基督所说：\Quote{我有权舍弃我的生命，也有权再取回来；没有人夺去我的生命，而是我甘愿舍弃，并且我有权再取回来。}\BibleRef{若 10:18}那么这权能受命运支配吗？让他们指给我们一个人，他有权力决定何时死、活多久：他们永远做不到。因此，让他们相信天主所说：\Quote{我有权舍弃我的生命，也有权再取回来；}并让他们探究为何说\Quote{我的时辰还没有到；}且不要因这些话，将命运强加于诸天的创造者、众星的创造者和统治者。因为即使命运来自众星，众星的创造者也不能受它们的宿命支配。此外，不仅基督没有你们所谓的命运，就连你、我、他，或任何一个人，也都没有。\par

11. 然而，他们自己被欺骗，又欺骗别人，向人提出谬论。他们设下陷阱捉人，而且是在大街上。那些张网捕捉野兽的人，尚且只在树林和旷野中做这事；为了捉人而张网在广场上，这是何等可怜虚妄！当人把自己卖给另一个人时，他们得到金钱；但这些占星术士却付出金钱，为把自己卖给虚妄。因为他们去找占星术士，为自己买来主人，就是占星术士乐意给他们的：无论是土星、木星、水星，或其他任何不敬神的名号。那人本来是自由地走进去，为的是付出金钱后出来时却成了奴隶。不，更确切地说，如果他本是自由的，他就不会进去；但他进去，是因为他的主人\Quote{谬误}和他的女主人\Quote{贪婪}把他拖了进去。因此，真理也说：\Quote{凡犯罪的，就是罪的奴隶。}\BibleRef{若 8:34}\par

12. 那么，祂为什么说\Quote{我的时辰还没有到？}而是因为，祂有权决定何时去世，却尚未视这权能的使用为合适。正如我们，弟兄们，例如说：\Quote{现在我们出去举行圣事的指定时辰到了。}如果我们提前出去，难道不是荒谬无理吗？而我们只在适当的时候行事，因此当我们这样表达时，我们在这行动中是否顾念命运？那么，\Quote{我的时辰还没有到}是什么意思？当我晓得那是我受苦的合适时刻，当我的受苦将有益处时，我就甘愿受苦。那个时辰尚未到来：为要你保留两者，即\Quote{我的时辰还没有到；}和\Quote{我有权舍弃我的生命，也有权再取回来。}那么，祂来了，有权决定何时去世。确实，如果祂在拣选门徒之前就去世，那是不合宜的。假如一个人自己无权决定他的时辰，他可能在拣选门徒之前就去世；如果他恰好在门徒已被拣选和教导之后去世，那也是加给他的，而非他自己的作为。但与此相反，祂来了，有权决定何时去、何时回、前进多远，而且坟墓的区域为祂是敞开的，不仅是在死亡时，也是在复活时；我说，祂为了向我们显示祂教会不朽的希望，在元首身上显示了肢体所当期待的事。因为那在元首内复活了的，也将在祂所有的肢体内复活。那么，那时辰尚未到来，适当的时候尚未到。必须召叫门徒，宣讲天国，在奇迹中彰显主的性体，在祂对必死之人的深切同情中彰显祂的人性。因为那因是人而饥饿的祂，因是天主而以五个饼喂饱数千人；那因是人而睡眠的祂，因是天主而命令风和浪。所有这一切都必须先行展示，为叫圣史们有可写的内容，好使教会有所宣讲。但当祂做了祂认为足够的事后，祂的时辰就来了，不是出于必须，而是出于意愿——不是出于状况，而是出于权能。\par

13. 那么，弟兄们？因为我们已回答了这些和那些，我们是否对水缸的象征意义闭口不言？水变成酒象征什么？管宴席的象征什么？新郎象征什么？耶稣的母亲在奥秘中象征什么？婚宴本身象征什么？我们必须谈论这一切，但不可加重你们的负担。我原打算昨天也因基督的名向你们讲道，那时例行的讲道应当向你们宣讲，我所爱的，但我被某些必要的事务所阻。那么，如果你们愿意，圣善的弟兄们，让我们把有关这段寓意解释的隐意的事推迟到明天，不加重你们和我们自身的软弱。你们中许多人或许今天是因为节日的隆重而聚集，并非来听讲道。让那些明天来的人来听，这样我们既不辜负那些好学的人，也不使那些挑剔的人感到负担。\par

\SourceRecord{Translated by John Gibb.；Nicene and Post-Nicene Fathers, First Series；Vol. 7.；Edited by Philip Schaff.；Buffalo, NY: Christian Literature Publishing Co.,；1888.；<http://www.newadvent.org/fathers/1701008.htm>.}

% structure: p0001 source=h1 target=section reason=page-title

\section{讲道集第9篇（\BibleRef{若 2:1-2}）}

愿主我们的天主亲临，使我们得以向你们实践所作的许诺。因为昨日，圣洁的弟兄们，若你们记得，时间短促，我们未能完成已开始的讲道，便将那些奥秘地隐藏在这段福音事迹中的深义，靠天主的助佑，留待今日阐释。因此，我们现在无须再停留于称赞天主的奇迹。因为祂是同一的天主，在整个受造界中每日行奇迹，这些奇迹因着它们的频繁发生，在人眼中变得轻忽；而非因其施行之易。然而，那些由同一的主——即为我们降生成人的圣言——所行的非凡事工，却更令人们惊异，不是因为它们比祂每日在受造界中所行的更大，而是因为这些每日发生的事仿佛循着自然之序；而那些事则显示出一种仿佛直接临在的能力之功效，展现于人眼前。我们说过，你们记得，一个死人复活了，人们惊叹；而每日从无有中出生的人，却无人惊奇。同样，水变成酒，谁不惊奇？尽管天主每年在葡萄树中行此事。然而，主耶稣所行的一切事工，不仅以其奇迹性唤起我们的心，也以信德之道建树我们的心，因此我们应当仔细探究这些事工的涵义与象征。因为所有这些事的含义，我们记得，留待今日思考。\par

主来到他所受邀的婚宴，除了奥秘的象征意义外，还愿向我们保证婚姻是祂自己的建立。因为将会有一些人，如宗徒所说，\Quote{禁止嫁娶}（\BibleRef{弟前 4:3}），并断言婚姻是邪恶的，是出于魔鬼的设立；然而同一的主在福音中，当被问及人是否可以因任何缘故休妻时，祂回答：除非因奸淫的缘故，否则不可。在祂的回答中，你们若记得，祂说：\BibleQuote{天主所结合的，人不可拆散}（\BibleRef{玛 19:6}）。凡正确受教于公教信仰的人都知道，天主建立了婚姻；正如夫妇的结合来自天主，离婚则来自魔鬼。但若发生奸淫，人可休妻，因为她先选择了不再为妻，未向丈夫保守婚姻的忠贞。那些向天主发愿守贞的女子，虽然在教会中享有更高的尊荣与圣洁地位，但也并非没有婚姻。因为她们与整个教会一同归于一个婚姻，一个以基督为新郎的婚姻。为此，主受邀请来到婚宴，是为确认婚姻的贞洁，并显示婚姻圣事的奥秘。因为那婚宴中的新郎，有人对他说：\Quote{你却把好酒留到现在}，这新郎代表了主的位格。因为好酒——即福音——基督一直留到现在。\par

现在我们开始揭示这些奥迹的隐藏含义，以祂——我们曾因祂的名向你们许诺的那位——所能赐予我们的程度。在古代，有预言，且没有时代被剥夺了预言的安排。但预言，由于基督在其中未被领悟，便是水。因为在水里，酒以某种方式隐含其中。宗徒告诉我们应当如何理解这水：他说：\Quote{直到今日，}又说：\Quote{诵读《梅瑟》时，那同一的面纱仍搁在他们的心上；它未揭开，因为只在基督内才废除。但当你们转向主时，}又说：\Quote{这面纱就被除去}（\BibleRef{格后 3:14}）。他所说的面纱，是指预言的遮蔽，使其不被领悟。你们何时皈依主，面纱就被除去；同样，淡而无味也被除去，因为你们皈依了主；那时水对你们变成了酒。阅读所有先知书；如果基督在其中未被领悟，你们能找到什么如此乏味愚蠢？在其中领悟基督，那么你们所读的不仅有滋味，甚至使你们沉醉；将心灵提升，超脱身体，以致忘记过去的，而向前的直跑（\BibleRef{斐 3:13}）。\par

4. 因此，从古时起，甚至从人类生育系列开始运行的时候，关于基督的预言就没有缄默；但是预言的含意隐藏在其中，因为当时它还是水。我们如何证明在往昔所有时代，直到主来临的时代，关于他的预言从未中断？从主自己的话。因为他从死者中复活后，发现他所跟随的门徒对他自己心存怀疑。因为他们看见他死了，并且不盼望他复活；他们的一切希望都已消失。那窃贼在受到称赞后，凭什么理由被判定为当天就与主同在乐园？因为当他被钉在十字架上时，他宣认了基督，而门徒却怀疑他。主发现他们摇摆不定，并且在某种程度上自责，因为他们曾期待从他获得救赎。然而他们为他哀伤，认为他是无辜被除掉的，因为他们知道他是无罪的。这正是门徒自己在主复活后所说的话，当时主在路上遇见几个门徒，他们忧愁，说：\BibleQuote{你独在耶路撒冷作客，还不知道这几天在那里所发生的事吗？耶稣对他们说：什么事？他们回答说：关于纳匝肋人耶稣，他是一位先知，在天主和众百姓面前，言行都有大能；我们的司祭和首领竟把他交出去，判了死罪，把他钉在十字架上。但我们原希望他就是那要拯救以色列的；可是，这些事发生到今天已是第三天了。}那两个人中的一个，正往附近的村庄去，主在路上遇见了他，他说了这些话和其他话之后，耶稣回答说：\BibleQuote{唉，无知的人哪，心迟钝，不相信先知们所说的一切话。基督岂不该受这一切苦难，而进入他的光荣吗？于是从梅瑟和众先知起，凡经上所指着自己的话，都给他们讲解明白了。}同样，在另一处，当主要求门徒用手触摸他，好相信他确实身体复活了时，他说：\BibleQuote{这就是我还与你们同在时所告诉你们的话：凡梅瑟法律、先知和圣咏上论及我所记载的，都必须应验。于是开启他们的明悟，使他们理解圣经，又对他们说：经上曾这样记载：基督必须受苦，第三天从死者中复活；并且必须因他的名宣讲补赎及罪过的赦免，从耶路撒冷开始，直到万邦。}\par

5. 当福音的这些话被理解时——它们当然是清楚的——所有隐藏在主这个奇迹中的奥秘都将被揭示。请注意他说的话：凡论及他所记载的，都必须应验在他身上。这些记载在哪里？他说：\BibleQuote{在法律书、先知书和圣咏上。}他没有省略旧约圣经的任何部分。这些是水；因此门徒被主称为\Emph{无知}的人，因为他们当时尝到的仍是水，不是酒。他如何将水变成酒？当他开启他们的明悟，并从梅瑟起，通过所有先知，给他们讲解圣经时，他们如今沉醉其中，说：\BibleQuote{当他在路上给我们讲解圣经时，我们的心不是火热的吗？}因为他们在这以前不认识基督的这些书中认识了他。这样，我们的主耶稣基督将水变成了酒，如今有了滋味，而以前没有；如今令人沉醉，而以前没有。因为如果他命令将水从水缸里倒出，然后亲自从受造物的隐秘处取出酒来——就像他使饼增多，满足数千人那样（因为五个饼本身不足以满足五千人，甚至不能装满十二筐，而是主的全能如同饼的源泉）——同样地，他可以在水倒出后倒入酒。但若他这样做，他似乎就拒绝了旧约圣经。然而，当他亲自将水变成酒时，他向我们表明旧约圣经也出自他，因为水缸是按照他的命令装满的。的确，旧约圣经也来自主；但除非在其中理解基督，它就没有滋味。\par

6. 但是请注意他自己说的话：\BibleQuote{凡梅瑟法律、先知和圣咏上论及我所记载的。}我们知道，法律从我们有记载的时代开始，也就是从世界的开始：\BibleQuote{在起初，天主创造了天地。}\BibleRef{创 1:1}从那起直到我们现在所处的时代，共有六个时代，这是第六个，正如你们常听到和知道的。第一个时代由亚当到诺厄；第二个时代由诺厄到亚巴郎；正如玛窦福音所遵循和区分的，第三个时代由亚巴郎到达味；第四个时代由达味到迁至巴比伦；第五个时代由迁至巴比伦到若翰洗者；\BibleRef{玛 1:17}第六个时代由若翰洗者到世界终结。此外，天主在\Emph{第六天}按自己的肖像造了人，因为在这第六个时代里，通过福音、按创造我们的那位的肖像，我们的心志得以更新；\BibleRef{哥 3:10}水变成了酒，使我们在法律和先知中品味到已彰显的基督。因此，\BibleQuote{那里有六口石缸，}他命令装满水。这六口石缸象征着六个时代，这些时代并非没有预言。那些六个时期，好像用关节分开和隔开，如果没有基督充满它们，就会像空器皿一样。我为什么说这些时期会徒然地流逝，除非在其中宣讲主耶稣？预言得以应验，水缸充满了；但是，为使水变成酒，基督必须在整个预言中被理解。\par

7. 这是什么意思：\Quote{它们各可装二至三梅特瑞塔}？这句话无疑给我们传达了一个奥秘的意义。因为他藉\Quote{\OriginalTerm{梅特瑞塔}{metretæ}}指的是一些度量，就像说罐子、瓶子之类的东西。\OriginalTerm{梅特瑞塔}{Metreta}是一种度量的名称，其名称来自\Quote{度量}一词。因为\OriginalTerm{度量}{\GreekText{μέτρον}}就是希腊文的度量，\Quote{梅特瑞塔}一词由此而来。然后，\Quote{它们各可装}，\Quote{二至三梅特瑞塔}。弟兄们，我们该说什么呢？如果他只说\Quote{各装三个}，我们的心立刻就会想到天主圣三的奥秘。或许我们不应立即拒绝这一含义的应用，因为他曾说\Quote{各装二或三个}；因为当圣父和圣子被称呼时，圣神必然被理解。因为圣神不单是圣父的，也不单是圣子的，而是圣父和圣子的神。因为经上记载：\BibleQuote{凡爱世界的，在他内就没有圣父之神。}\BibleRef{若一 2:15}又说：\BibleQuote{谁若没有基督的神，就不属于祂。}\BibleRef{罗 8:9}所以，圣神就是圣父和圣子的神。因此，当圣父和圣子被提及时，圣神也被理解，因为祂是圣父和圣子的神。当提及圣父和圣子时，就好像提到了\Quote{两个梅特瑞塔}；但由于圣神就在其中，就成了\Quote{三个梅特瑞塔}。这就是为什么经上说\Quote{这六个水缸各装二或三梅特瑞塔}，而不是\Quote{一些装两个，另一些装三个}。这好像他在说：当我说各装两个时，我意思是圣父和圣子之神也要连同他们被理解；当我说各装三个时，我就更明确地宣告同一的天主圣三。\par

8. 因此，谁称呼圣父和圣子，就应当由此理解圣父和圣子之间的相互之爱，那就是圣神。或许圣经在查考之下（我不是说今天就能向你们证明此事，或者好像找不到另一个证据）——但圣经在搜查之下，或许确实向我们显示圣神就是爱德。你们不要以为爱德是无价值的东西。实际上，它怎么能是无价值的呢？当一切被认为有价值的东西都称为\OriginalTerm{宝贵}{chara}？因此，如果非廉价的东西是宝贵的，那么还有什么比\OriginalTerm{宝贵本身}{charitas}更宝贵呢？宗徒如此推崇爱德，以致他说：\BibleQuote{我给你们指示一条更高超的道路。我若能说人间的语言和天使的语言，但我若没有爱德，我就成了发声的锣或发响的钹。我若有先知之恩，又明白一切奥秘和各种知识；我若有全备的信德，甚至能移山，但我若没有爱德，我什么也不算。我若把我所有的财产全施舍了，又舍身投火被烧，但我若没有爱德，为我毫无益处。}\BibleRef{格前 13:1}那么，爱德何等伟大，若缺少它，其他一切都徒然；若拥有它，则一切皆正！然而，保禄宗徒用丰富而饱满的言辞赞美爱德，说的还不如撰写此福音的若望宗徒简短的一句话。因为他不迟疑地说：\Quote{天主是爱}。经上也记载：\BibleQuote{天主的爱藉着所赐给我们的圣神，已倾注在我们心中了。}\BibleRef{罗 5:5}那么，谁称呼圣父和圣子，能不由此理解圣父和圣子之爱呢？当人开始拥有这爱时，他就将拥有圣神；若没有这爱，他就不会拥有圣神。正如你的身体若没有灵魂，即你的神，就是死的；同样，你的灵魂若没有圣神，即没有爱德，也被视为死的。因此，\Quote{水缸各装两个梅特瑞塔}，因为在所有时期的预言中都宣告了圣父和圣子；但圣神也在其中，所以又加上\Quote{或三个}。\BibleQuote{我与父原是一体。}\BibleRef{若 10:30}但千万别以为当我们听到\BibleQuote{我与父原是一体}时，圣神不在其中。然而，既然他提到了圣父和圣子，就让水缸装\Quote{两个梅特瑞塔}；但请注意\Quote{或三个}。\Quote{你们要去，使万民成为门徒，因父、子、圣神之名给他们授洗。}因此，当说\Quote{两个}时，并未明说但已隐含天主圣三；而当说\Quote{或三个}时，也明说了天主圣三。\par

9. 但还有一个含义不可忽略，我将予以阐述：让人各选其喜爱者。我们不隐瞒所领受的。因为这是主的宴席，仆人不应欺骗客人，尤其当你们像现在这样饥饿时，你们的渴望是明显的。自古以来所分施的预言，其目的是万民的救恩。诚然，梅瑟只被派往以色列民，法律也由他赐给那一民族；先知也出自那一民族，就连时代的分期也按那一民族而定；因此，水缸也被称为\Quote{按犹太人取洁礼}；然而，预言也向所有其他民族宣告，这是显而易见的，因为基督隐藏在他内，万民都要因他而蒙福，正如主曾向亚巴郎应许说：\BibleQuote{地上万民都要因你的后裔而蒙福。}\BibleRef{创 22:18}但这当时还未被理解，因为水尚未变成酒。因此，预言被分施给万民。但为更愉快地显明这一点，在时间允许的范围内，让我们提及几个关于不同时代的事实，这些事实分别由水缸所代表。\par

10. 在最起初，亚当和厄娃是全人类、而非仅是犹太人的祖先；凡在亚当身上关于基督所预表的，无疑关系到万民，因为万民的救恩是在基督内。关于第一缸水，我怎么可能比宗徒论及亚当和厄娃时说得更好呢？因为当我引述的并非自己的话，而是宗徒的话时，没有人会说我对这意义理解有误。那么，宗徒所提及的那事，当他说道：\Quote{二人要成为一体：这是一个伟大的奥迹！}\BibleRef{厄 3:31} 时，其中包含了何等伟大的关于基督的奥迹啊！惟恐有人以为这奥迹的伟大存在于各自娶妻的人身上，他又说：\Quote{但我是指基督和教会说的。}这伟大的奥迹是什么，即\Quote{二人要成为一体}？圣经在\BibleBook{}{创世纪}中论到亚当和厄娃时，说了这些话：\Quote{为此，人要离开父亲和母亲，依附自己的妻子，二人成为一体。}\BibleRef{创 2:24} 如今，如果基督依附于教会，以致二人成为一体，那么祂是如何离开祂的父亲和母亲呢？祂离开祂的圣父，其意义是：当祂以天主的形体存在时，祂不以自己与天主同等为僭夺，反而空虚了自己，取了奴仆的形体。\BibleRef{斐 2:6} 在这个意义上，祂离开了圣父，并非祂遗弃或远离了圣父，而是祂没有以与圣父同等的形体显现于人前。祂又是如何离开了祂的母亲呢？藉着离开犹太人的会堂——按肉身祂是从那里出生的——并依附于祂从万民中所聚集的教会。因此，第一缸水当时就包含着关于基督的预言；但是，只要我所说的这些事情尚未在万民中宣讲，那预言就只是水，还没有变成酒。既然主藉着宗徒光照了我们，使我们通过这一句话\Quote{二人要成为一体；一个伟大的奥迹，关于基督和教会}找到我们寻求的，我们现在就被允许到处寻求基督，并从所有的水缸中喝到酒。亚当沉睡，为使厄娃被造；基督死了，为使教会得以形成。当亚当沉睡时，厄娃从他肋旁被造；当基督死了时，长矛刺透祂的肋旁，为使圣事涌流而出，教会便由此形成。对每个人来说，难道不明显吗：当时所行的事预表了将来的事，因为宗徒说亚当本人就是那未来者的预像？他说：\Quote{他是那未来者的预像。}\BibleRef{罗 5:14} 一切都是奥秘地预表的。因为事实上，天主本可以在亚当醒着时取用他的肋骨造出女人。或者，莫非需要他沉睡，免得在取肋骨时感到肋旁的疼痛？有谁睡得如此深沉，以致骨头被抽出也不会醒来？或者是由于是天主抽出，那人便没有感觉？其实，那位能在人沉睡时无痛取出肋骨的主，同样能在人醒着时这样做。但无疑，第一缸水正在被装满，那是关于将来之事的预言在当时的安排。\par

11. 基督也在诺厄以及那象征全世界的方舟中被预表。为何方舟里关着各种动物，岂不是为预表万民？因为天主本可以重新创造每一种动物。当它们尚未存在时，祂不是说过：\Quote{地要生出}，地就生出了吗？祂本可从同样的本源重新造出，正如祂当初所造的；祂用一句话造了，也能用一句话再造：除非祂是为我们揭示一个奥迹，并充满预言之安排的第二个水缸，使世界藉着木头以预像的方式得救，因为世界的生命将被钉在木头上。\par

12. 如今，在第三缸水上，正如我之前所提过的，对亚巴郎说过：\Quote{地上万民都要因你的后裔获得祝福。}有谁看不出亚巴郎独生子是谁的预像？他背负着用于祭献自己的木头，前往他要被献祭的地方。因为主背负了自己的十字架，正如\WorkTitle{福音}告诉我们的。关于第三缸水，这样说就足够了。\par

13. 至于达味，我为何要说他的预言延伸至万民呢？因为我们刚才还听到了圣咏——实在很难找出一篇未宣报这事的圣咏。但显然，正如我所说，我们刚才唱道：\Quote{上主，求祢起来审判大地，因为祢要在万民中承受产业。}这也正是多纳徒派被逐出婚宴的原因：就像那个没有婚宴礼服的人，他虽被邀请而来，却因没有礼服以荣耀新郎而被逐出宾客的行列；因为那寻求自己荣耀、而非基督荣耀的人，没有婚宴礼服：他们拒绝与那位新郎的朋友——他说：\Quote{那一位就是施洗者}——保持一致。于是，对那没有婚宴礼服的人，责难地反对他说：\Quote{朋友，你怎么到这里来了？}\BibleRef{玛 22:13} 正如那人无言以对，这些人也是如此。因为当内心静默时，口舌的喧嚷有何用呢？因为他们内心知道，并且在自己里面，他们无话可说。内里是静默的，外面却制造喧嚷。但是，无论他们愿意与否，他们甚至在自身中间也听到这歌唱：\Quote{上主，求祢起来审判大地，因为祢要在万民中承受产业}；而他们不与万民共融，岂不是承认自己是无继承权的吗？\par

14. 现在，弟兄们，我所说的预言延伸到所有民族（因为我想向你们展示这个表达的另一含义：\BibleQuote{每口可盛两或三梅特累塔}），——我说，这预言延伸到所有民族，正如我们刚才提醒你们的，在\Person{亚当}身上显示出来，\BibleQuote{他是那要来者的预象。}谁不知道从他的后代生出了所有民族呢？在他名字的四个字母中，地球的四角通过它们的希腊名称被指示出来吗？因为如果用希腊语表达东方、西方、北方和南方，正如圣经在不同地方提到它们，你会发现这些单词的首字母构成了\Person{亚当}这个词：因为希腊语中世界的四个方位被称为\OriginalTerm{东方}{Anatole}、\OriginalTerm{西方}{Dysis}、\OriginalTerm{北方}{Arktos}、\OriginalTerm{南方}{Mesembria}。如果你把这四个单词像四行诗一样一个接一个写下来，大写字母就形成了\Person{亚当}这个词。同样的事情也体现在\Person{诺厄}身上，由于方舟，其中包含所有动物，象征所有民族；同样在\Person{亚巴郎}身上，他更清楚地被告知：\BibleQuote{地上万民都要因你的后裔蒙受祝福；}同样在\Person{达味}身上，从他的圣咏中，忽略其他表达，我们刚才歌唱：\BibleQuote{天主，求祢起来审判大地，因为祢要继承万民。}现在，对哪一位天主说了\BibleQuote{起来，}除了那位睡了的那位？\BibleQuote{天主，求祢起来审判大地。}好像是在说，祢曾睡了，被大地审判了；起来审判大地。而这预言延伸到何处：\BibleQuote{因为祢要继承万民？}\par

15. 此外，在第五时代，就好像在第五个水缸里，\Person{达尼尔}看见一块石头从山上被凿出，非人手所凿，打碎了地上所有王国；他看见那石头长大，变成一座大山，充满了整个地面。\BibleRef{达 2:34} 还有比这更明显的吗，弟兄们？石头从山上被凿出：这块石头就是工匠弃而不用的石头，反而成了屋角的基石。它是从哪座山上被凿出的呢？如果不是从犹太王国，也就是我们的主耶稣基督按肉身所生的那一个？它是非人手所凿的，没有人的劳作；因为基督是由一位童贞女所生，没有丈夫的拥抱。被凿出的山没有充满整个地面；因为犹太王国没有拥有所有民族。但另一方面，基督的王国我们看见正占据着整个世界。\par

16. 第六时代属于\Person{若翰洗者}，在妇人所生的人中，没有兴起一个比他更大的；论到他，曾说\BibleQuote{他比先知还大。}\BibleRef{玛 11:11} 若翰如何表明基督被派到所有民族呢？当犹太人来到他那里受洗，以免他们因亚巴郎的名号而自夸时，他对他们说：\BibleQuote{毒蛇的种类，谁指示你们逃避那将要来的忿怒呢？你们要结出与悔改相称的果实来；}也就是说，要谦卑；因为他在对骄傲的人说话。他们为什么骄傲呢？因为按肉体的后裔，而不是效法他们祖先\Person{亚巴郎}的果实。他对他们说了什么？\BibleQuote{不要心里说：我们有亚巴郎为父。我告诉你们，天主能从这些石头中给亚巴郎兴起子孙来。}\BibleRef{玛 3:9} 他用石头意指所有民族，不是因为他们持久的坚强，如那被工匠弃用的石头那样，而是因为他们愚钝和愚蠢的麻木，因为他们变得像他们所习惯崇拜的东西：因为他们崇拜无感觉的偶像，他们自己也同样无感觉。\BibleQuote{愿制造它们的，以及一切信赖它们的，都变得像它们一样。} 因此，当人们开始崇拜天主时，他们听到对他们说什么？\BibleQuote{好使你们成为你们在天之父的子女，因为他使太阳升起，光照善人，也光照恶人；降雨给义人，也给不义的人。}\BibleRef{玛 5:45} 因此，如果一个人变得像他所崇拜的东西，那么\BibleQuote{天主能从这些石头中给亚巴郎兴起子孙来}是什么意思？让我们自问，我们将看到这是事实。因为我们是来自那些民族，但如果天主没有从石头中给亚巴郎兴起子孙，我们本不会从他们而来。我们因效法亚巴郎的信德而成为他的子女，不是因肉体的出生。因为正如他们因堕落而被剥夺继承权，我们因效法而被收养。所以，弟兄们，这第六个水缸的预言也延伸到所有民族；因此论到一切，曾说\BibleQuote{每口可盛两或三梅特累塔。}\par

17. 然而，我们如何证明万民都属于\Quote{每缸两或三米特瑞特}的呢？在某种程度上，这是一个计数的问题：他先说同样的水缸\Quote{每缸两个}，后又曾说\Quote{每缸三个}；显然是为了向我们暗示其中的奥秘。\Quote{每缸两米特瑞特}如何呢？就是割损与未受割损。圣经提到这两类人，不遗漏任何种类的人，它说：\BibleQuote{割损与未受割损；}\BibleRef{哥 3:11}在这两个名称中，你拥有万民：它们就是每缸两米特瑞特。在这两堵来自不同方向的墙壁上，\BibleQuote{基督成了角石，为在自己内缔造和平。}\BibleRef{弗 2:14}让我们也在这万民的情况下展示\Quote{每缸三米特瑞特}。诺厄有三个儿子，人类藉他们得以复兴。因此主说：\BibleQuote{天国好比酵母，妇人取来藏在三斗面中，直到全部发酵。}\BibleRef{路 13:21}这妇人若不是主的肉身，是什么呢？酵母若不是福音，是什么呢？那三斗若不是万民，又是什么呢？这是由于诺厄的三个儿子。因此，\Quote{六口缸，每口盛两或三米特瑞特}就是六个时期，包含着关乎万民的预言，无论是以两类人代表，即犹太人和希腊人（如宗徒常提到的），还是以三类人代表，由于诺厄的三个儿子。因为这预言被表明为延伸到万民。正因这种延伸，它被称为\Quote{尺度}，正如宗徒所说：\BibleQuote{我们领受了尺度，为伸展到你们那里。}\BibleRef{格后 10:13}因为他在向外邦人宣讲福音时说：\BibleQuote{一个尺度，为伸展到你们那里。}\par

\SourceRecord{Translated by John Gibb.；Nicene and Post-Nicene Fathers, First Series；Vol. 7.；Edited by Philip Schaff.；Buffalo, NY: Christian Literature Publishing Co.,；1888.；<http://www.newadvent.org/fathers/1701009.htm>.}

% structure: p0001 source=h1 target=section reason=page-title

\section{讲道集第十篇（\BibleBook{若望}{福音} 2:12-21）}

1. 在圣咏中，你们听见了穷苦人的叹息，他们的肢体在全地忍受磨难，直到世界末日。弟兄们，你们首要的事是成为这些肢体中的一员，并与他们同在：因为一切磨难都要过去。\Quote{那些欢笑的人有祸了！}\BibleRef{路 6:25}\Quote{有福的，}真理说，\Quote{是那些哀恸的人，因为他们要受安慰。}天主成了人：人将成为什么，天主竟为了人成了人？愿这望德在此生的一切磨难和诱惑中安慰我们。因为仇敌不停地迫害；当他没有公开咆哮时，就暗中图谋。他怎样图谋呢？\Quote{他们因愤怒而行事诡诈。}因此他被称为狮子和龙。但基督得到了什么话？\Quote{你要践踏狮子和龙。}狮子，指公开的咆哮；龙，指隐藏的诡计。龙把亚当逐出乐园；同一魔鬼如狮子般迫害教会，正如伯多禄所说：\Quote{你们的仇敌魔鬼，如同咆哮的狮子巡遊，寻找可吞食的人。}\BibleRef{伯前 5:8}不要以为魔鬼失去了它的凶残。当他甜言蜜语地谄媚时，更应提防他。但在这一切诡计和诱惑中，我们该做什么？唯有我们在圣咏中所听见的：\Quote{至于我，当他们烦扰我时，我披上苦衣，以斋戒谦卑我的灵魂。}有一位听祈祷的，不要犹豫去祈祷；但听的那位住在内心。你不必把目光投向某座山；不必仰面朝星、日、月；也不可认为你在海边祈祷就会被垂听：反而要厌恶这样的祈祷。只要洁净你心灵的内室，无论你在何处，无论你在哪里祈祷。听你的那位就在里面，在隐秘处，圣咏作者称之为自己的怀中，当他这样说时：\Quote{我的祈祷要转回我的怀中。}听你的那位不在你之外，你不必远行，也不必举身，好像用手去够祂。反而，你若高举自己，必要跌倒；你若谦卑自己，祂就亲近你。我们的主天主就在这里，天主的圣言，成了血肉的圣言，圣父之子，天主子，人子；崇高者创造我们，卑微者更新我们，行走在人中，带着人性，隐藏神性。\par

2. \Quote{祂下去，}正如福音作者所说，\Quote{到葛法翁，祂和祂的母亲、祂的弟兄、祂的门徒；他们在那里住了不多几日。}看，祂有母亲、弟兄和门徒：祂有母亲，因此也有弟兄。因为我们的圣经惯常称那些为弟兄，不仅是出于同一男人和女人，或同一母亲，或同一父亲（尽管母亲不同）；或者，实际上，是姑表或姨表亲：我们的圣经并非只称这些人为弟兄。圣经必须按照它说话的方式来理解。它有它自己的语言；一个不懂这种语言的人会困惑，说：主从何而有弟兄？因为玛利亚难道又生了第二次？决不然！在她身上开始了童贞女的尊荣。她能成为母亲，但不能成为与男人同房的女人。她被称作\OriginalTerm{女人}{mulier}（通常指\Emph{妻子}）。因为圣经本身的语言也是如此。因为厄娃，她刚被从丈夫的肋旁造出，尚未与丈夫同房时，如你们所知，也被称为女人：\Quote{他造了她成为一个\OriginalTerm{女人}{mulier}}。那么，弟兄从何而来？玛利亚的亲属，不论何种关系，都是主的弟兄。我们如何证明这一点？从圣经本身。罗特被称为\Quote{亚巴郎的弟兄}；他是他兄弟的儿子。读圣经，你们会发现亚巴郎是罗特的叔父（父系），然而他们被称为弟兄。为什么？只因他们是亲属。叙利亚人拉班是雅各伯的舅父（母系），因为他是黎贝加（依撒格的妻子、雅各伯的母亲）的兄弟。\BibleRef{创 28:5}读圣经，你们会发现舅父和外甥也被称为弟兄。\BibleRef{创 29:12}当你们知道这个规则，就会发现玛利亚的所有血亲都是基督的弟兄。\par

3. 但那些门徒更是弟兄；因为即使那些亲属，若不是门徒，也不是弟兄：若他们不承认他们的弟兄为他们的师傅，这弟兄关系也毫无益处。因为在某处，当有人告诉祂说，祂的母亲和祂的弟兄站在外面，那时祂正对门徒说话，祂说道：\Quote{谁是我的母亲？谁是我的弟兄？祂伸手指着门徒说，这些人就是我的弟兄；}又说道：\Quote{凡遵行我天父旨意的，他就是我的母亲、兄弟、姊妹。}\BibleRef{玛 12:46}因此，玛利亚也不例外，因为她遵行了天父的旨意。主在她身上所尊崇的，是她遵行了天父的旨意，而非血肉生了血肉。亲爱的，要留意。此外，当主因行奇迹异能而受到群众景仰，并显露出隐藏在血肉之下的事时，某些赞叹的灵魂说：\Quote{怀过你的胎是有福的！}祂却说：\Quote{不如说，那些听天主的话而遵行的人是有福的。}\BibleRef{路 11:27}这就是说，即使我的母亲，你们所称有福的，她之所以有福，是因为她持守天主的话：并非因在她内圣言成了血肉并居住在我们中间；而是因她持守了那创造她的同一圣言，这圣言在她内成了血肉。愿人不因肉身子孙而欢欣，却要因在精神上与天主结合而喜乐。我们说这些话，是因为福音作者说，祂在葛法翁住了几日，与祂的母亲、祂的弟兄和祂的门徒同住。\par

4. 随后如何？\BibleQuote{犹太人的逾越节近了，耶稣便上耶路撒冷去。}叙述者就他所忆，讲述了另一件事。\BibleQuote{祂在圣殿里看见卖牛、羊、鸽子的，并兑换银钱的人坐着；祂用绳索做成鞭子，把众人连牛带羊都赶出圣殿，倒出兑换银钱者的钱，推翻他们的桌子，并对卖鸽子的人说：\InnerQuote{把这些东西拿出去，不要使我父的圣殿成为商场。}}弟兄们，我们听到了什么？请看，那圣殿当时仍是一个预象，而主却将一切寻求自己利益、前来做买卖的人从其中驱逐出去。他们在那里卖什么呢？是那个时代祭献中所需要的东西。因为你们知道，亲爱的，祭献曾赐予那百姓，因为他们当时尚存肉性的心和石心，以免他们堕落去拜偶像；他们在那里献牛、羊、鸽子为祭：你们知道这事，因为你们已读过。那么，在圣殿里出售那些为在圣殿中献祭而购买之物，本不是大罪；然而，祂仍将他们驱逐出去。如果他们买卖的是合法且不违背正义的事（因为买得体面之物，卖之并非不合法），祂尚且将那些人赶走，决不肯使祈祷的圣殿成为买卖之所；那么，倘若祂在那里发现醉汉，主会做什么呢？既然天主的圣殿不应成为贸易之所，难道应成为酗酒之所吗？但当我们说这话时，他们向我们咬牙切齿；你们所听到的圣咏却安慰我们：\BibleQuote{他们向我咬牙切齿。}然而，我们知道我们如何可得痊愈，尽管鞭笞的打击加在基督身上，因为祂的圣言正承受鞭笞：\BibleQuote{鞭笞，}祂说，\BibleQuote{聚集起来攻击我，他们却不自知。}祂被犹太人的鞭子鞭笞；如今祂被假基督徒的亵渎所鞭笞：他们为他们的主加增鞭笞，却不自知。让我们，因祂帮助，如同圣咏作者那样行：\BibleQuote{至于我，当他们烦扰我时，我便穿上苦衣，以斋戒卑微我的灵魂。}\par

5. 然而，弟兄们，我们说（因为祂没有饶恕那些人：那位将要被他们鞭笞的，首先鞭笞了他们），祂给了我们一个标记，即祂用绳索做成鞭子，用它击打那些妄图将天主的圣殿当作买卖的无序之徒。因为的确，每个人都用自己的罪孽为自己绞出一条绳索：\BibleQuote{祸哉，那些以绳索牵引罪孽的人？}谁做长绳索？就是罪上加罪的人。罪如何加在罪上？当已犯的罪用其他的罪掩盖时。一个人犯了偷盗：为了不被发现，他去寻求占星者。犯偷盗已足够，你为何还要罪上加罪？看，两桩罪已犯。当你被禁止去占星者那里时，你辱骂主教：看，三桩罪。当你听说有人论你说\Quote{将他从教会中逐出}，你说\Quote{我要投靠多纳图派}：看，你加了第四桩罪。绳索在增长；要畏惧那绳索。你在此世受这绳索抽打而被纠正，这是好的；免得最后有人对你说：\BibleQuote{捆起他的手和脚，把他丢在外面的黑暗中。} \BibleRef{玛 22:3} 因为，\BibleQuote{各人被自己的罪孽绳索缠住。} \BibleRef{箴 5:22} 前者是主的话，后者是另一段圣经的话；但二者都是主的话。人受自己的罪孽所捆，被丢在外面的黑暗中。\par

6. 然而，若要在预象中寻求此事的奥秘，卖牛的是谁？卖绵羊和鸽子的又是谁？他们是在教会中谋求自己利益，而非基督之事的人。他们把一切视为买卖，自己却不愿被救赎：他们不愿被买，却想出售。的确，他们若能因基督的血被救赎，从而进入基督的平安，那真是好事。现今，在这世上获取任何短暂易逝之物，无论是金钱、口腹之乐，还是由人的赞美而来的荣誉，有何益处？它们不都是风和烟吗？它们不都经过并消失吗？它们不都像一条直奔大海的河流吗？那落入其中的有祸了，因为他将被卷入大海。因此，我们应该抑制自己对这类欲望的一切情感。我的弟兄们，寻求这些东西的人，就是出卖者。因为那\Person{西满}也曾想购买圣神，正因为他打算出售圣神；他认为宗徒们正是那种被主用鞭子赶出圣殿的商人。他本人就是这样的人，想要购买他所可能出售的：他属于卖鸽子的人。原来圣神是以鸽子的形象出现的。\BibleRef{玛 3:16}那么，弟兄们，谁是卖鸽子的人？岂不就是那些说\Quote{我们赐予圣神}的人吗？但他们为什么这样说？以什么价格出售？以对自己的荣誉为代价。他们收取暂时的荣誉席位作为代价，好让人看到他们是卖鸽子的人。让他们当心那用细绳做的鞭子。鸽子不是用来卖的：它是白白赐予的；因为它被称为恩宠或恩惠。所以，我的弟兄们，正如你们看到那些卖主，普通小贩，每个人都夸耀他所卖的：他们设了多少摊位！\Person{普里米亚诺}在\Place{迦太基}有一个摊位，\Person{马克西米亚诺}有另一个，\Person{罗加图斯}在\Place{毛里塔尼亚}有另一个，他们在\Place{努米底亚}也有，这一派那一派，我们如今无法一一列举。因此，有人四处寻找购买鸽子，而每个人都在自己的摊位夸耀他所卖的。让这样的人的心远离每一个卖主，让他到能白白领受的地方来。是的，弟兄们，他们并不羞愧，因他们这些苦毒恶毒的分裂，使自己成了如此多的派别，自以为是他们不是的东西，自高自大，以为自己算什么，其实什么都不是。\BibleRef{迦 6:3}但在他们身上应验了什么？既然他们不肯改正，却应验了你们在圣咏中所听到的：\Quote{他们被撕裂，却毫不悔悟}？\par

7. 那么，谁卖牛呢？那些曾将圣经分施给我们的人，被视为牛。宗徒们是牛，先知们是牛。因此宗徒说：\Quote{牛在场上踹谷的时候，不可笼住它的嘴。天主岂是关心牛吗？祂岂不是为我们说的吗？的确，是为我们说的，因为耕田的应该怀着希望去耕，打场的也应该怀着分享的希望。}\BibleRef{格前 9:9}这些牛，便为我们留下了圣经的叙述。因为他们分施的并非出于自己，而是寻求主的荣耀。如今，你们在那篇圣咏中听见了什么？\Quote{愿那些希望祂仆人平安的人，不断地说：愿主被尊崇。}天主的仆人，天主的子民，天主的教会。愿希望那教会平安的人尊崇主，而不是仆人：\Quote{愿他们不断地说：愿主被尊崇。}谁应当说？\Quote{那些希望祂仆人平安的人。}那个子民、那个仆人的声音，明显就是你们所听见的圣咏中哀歌的声音，你们听见时深受感动，因为你们属于那个子民。一人所唱的，从众人心中回响。那在那些声音中如同在镜子中认出自己的人有福了。那么，谁希望祂仆人、祂子民、祂称为祂的\Quote{独一者}的平安，并希望祂从狮子口中被救出：\Quote{把我的独一者从犬类的权势中救出？}就是那些常说\Quote{愿主被尊崇}的人。这些牛，便尊崇了主，而不是自己。请看这头牛尊崇它的主人，因为\Quote{牛认识它的主人}\BibleRef{依 1:3}；看这头牛是如何害怕人离开牛的主人而依赖牛：它多么畏惧那些愿意信赖它的人：\Quote{保禄曾为你们被钉死吗？或者你们是因保禄的名受洗的吗？}\BibleRef{格前 1:13}我所给的，我不是赐予者：你们白白地领受了；鸽子从天降下。\Quote{我栽种了，}他说，\Quote{\Person{阿颇罗}浇灌了，但使之生长的是天主；所以栽种的不算什么，浇灌的也不算什么，但使之生长的是天主。}\BibleRef{格前 3:6}\Quote{愿那些希望祂仆人平安的人，不断地说：愿主被尊崇。}\par

8. 然而这些人，正是借着圣经欺骗民众，好从他们手中获得荣誉和赞美，且使人不归向真理。他们在欺骗那些他们向其寻求荣誉的民众时，正是借着圣经，实际上是在卖牛：他们也卖羊，即普通民众本身。他们又把它们卖给谁，除非是卖给魔鬼？因为如果教会是基督的唯一独有者，那么从教会中被割裂的任何部分，是谁带走？岂不是那咆哮巡游的狮子，\Quote{寻找可吞食的人}？\BibleRef{伯前 5:8}那从教会中被割裂的人有祸了！至于她，她将保持完整。\Quote{因为主认识属于祂的人。}\BibleRef{弟后 2:19}然而这些人，尽其所能地卖牛和羊，也卖鸽子：让他们提防自己罪过的鞭子。但当他们因这些不义而遭受类似事情时，让他们承认主已用细绳做了鞭子，并告诫他们要改变自己，不再作贩卖者：因为他们若不改变，终将听见说：\Quote{捆起这人的手脚，把他丢在外面的黑暗中。}\par

9. \Quote{那时，门徒们就想起经上记载的话：\BibleQuote{我为祢殿宇所发的热忱把我耗尽。}} 因着对天主殿宇的这种热忱，主把这些人都赶出了圣殿。弟兄们，愿每一位基督徒，作为基督的肢体，都被对天主殿宇的热忱所吞噬。谁是为天主殿宇所吞噬的人？就是那努力纠正他所见到的一切错误，愿意使之改善，不袖手旁观的人；他若不能改善，便忍受它，为它哀叹。麦粒在打谷场上被扬出，是为了在糠秕分离后得以进入谷仓。如果你是一粒麦子，不要在入仓之前就从场上被扬出，免得你在被收进仓里之前就被鸟儿叼走。因为天上的飞鸟，即空中的权势，正等着从打谷场上叼走些什么；它们只能叼走那些被扬出的东西。因此，愿对天主殿宇的热忱吞噬你：愿对天主殿宇的热忱吞噬每一位基督徒，就是这殿宇中作为肢体的热忱。因为你自己的家并不比那个你获得永息的家更重要。你进入自己的家是为了暂时的休息，你进入天主的殿宇是为了永久的安息。那么，如果你忙于照看自己家里不要有错，难道就容忍你在天主的殿宇中——那里为你摆设了救恩和无尽的安息——看到任何错处而不尽力去阻止吗？例如，你看见一个兄弟冲向剧场？阻止他，警告他，使他难过，如果对天主殿宇的热忱吞噬了你。你看见别人奔跑着想要喝醉，甚至是在圣地，那在任何地方都不合宜的事情？阻止你能阻止的人，约束你能约束的人，恐吓你能恐吓的人，温柔地吸引你能吸引的人，但不可保持沉默。是朋友吗？要温柔地劝诫。是妻子吗？要严厉地管束。是女仆吗？即使用责打也要遏制。尽你所能，为你所承担的部分，如此你便完成了。\BibleQuote{我为祢殿宇所发的热忱把我耗尽。} 但如果你冷漠、懈怠，只顾自己，仿佛自己就足够了，心里说：\Quote{我何必管别人的罪？我只要照管好自己的灵魂就够了，这灵魂我要为天主保持纯洁。} 那么，你难道不记得那个藏起塔冷通而不去投资的家仆吗？他受责罚是因为他丢掉了它吗？不，是因为他没有用它获利。\BibleRef{玛 25:25} 因此，我的弟兄们，你们要听，免得你们无所事事。我要给你们劝告：愿那内在的那位赐予它；因为虽是通过我，却是祂在赐予。你们都知道自己在自己家里、与朋友、与佃户、与客户、与尊贵者、与卑微者该做什么：按天主给予的入口，按祂为祂的话语开门，不要停止为基督赢得人；因为你们是被基督赢得的。\par

10. \BibleQuote{犹太人对祂说：祢做这些事，给我们显什么标记呢？} 主回答说：\BibleQuote{你们拆毁这座圣殿，三天之内我要把它重建起来。} 犹太人便说：\BibleQuote{这圣殿建造了四十六年，祢说三天之内要把它重建起来吗？} 他们是属血肉的，思想属血肉的事；但祂却是在讲论属神的事。然而谁能明白祂所说的是哪座圣殿呢？但我们不必远求；祂已透过圣史使我们得知，圣史告诉了我们祂所说的是哪座圣殿。圣史说：\BibleQuote{但祂所说的是祂身体的圣殿。} 显而易见，主被杀后，第三天复活了。这事如今我们所有人都知道：若向犹太人隐藏，那是因为他们在外面；但对我们却是敞开的，因为我们知道我们所信的是谁。那圣殿的拆毁与重建，我们即将在每年的庆节中庆祝；为此我们劝勉你们准备自己，你们中那些慕道者，好能领受恩宠；现在就是时候，现在就当立志，使那时得以诞生。这件事我们已经知道了。\par

11. 但或许有人会问我们：圣殿建造了四十六年，这一事实是否含有某种奥义？的确，关于此事可说很多；但我们可以简要地说一些容易明白的话。弟兄们，如果我没有记错的话，我们昨天说过，亚当是一个人，却也是全人类。我们是这样说的，如果你们记得。他仿佛被粉碎了，分散了，如今正在被聚集，并通过一种灵性的共融与和谐，仿佛重新被结合为一。而\Quote{那叹息的穷人}，作为一人，就是同一个亚当，但他在基督内正在被更新：因为那一位无罪的亚当来了，要在祂自己的肉身上毁灭亚当的罪，并使亚当得以重新恢复天主的肖像。基督的肉体也是来自亚当：那座被犹太人拆毁、主在三天内重建的圣殿也是来自亚当。因为祂复活了自己的肉身：看，祂就是这样与圣父同等的天主。我的弟兄们，宗徒说：\BibleQuote{谁使祂从死人中复活？} 他说的是谁？是圣父。\BibleQuote{祂服从至死，}他说，\BibleQuote{甚至死在十字架上；因此，天主也高高举扬了祂，赐给祂一个名字，超越一切名字。} \BibleRef{斐 2:8} 那位被复活并举扬的就是主。是谁复活了祂？是圣父，就是祂在圣咏中所说的：\BibleQuote{求祢举起我，我要报复他们。} 因此，是圣父复活了祂。难道祂没有复活自己吗？难道圣父能离开圣言做什么吗？圣父离开祂的独生子能做什么？因为请听，祂也是天主。\BibleQuote{你们拆毁这座圣殿，三天之内我要把它重建起来。} 难道祂是说：拆毁圣殿，三天内圣父要把它重建起来吗？但正如圣父复活时，圣子也复活；同样，圣子复活时，圣父也复活：因为圣子曾说过：\BibleQuote{我与圣父原是一体。} \BibleRef{若 10:30}\par

12. 现在，数字四十六是什么意思？同时，亚当如何遍布全球，你们昨天已经听过解释，通过四个希腊词的首字母。因为如果你们写下这四个词，上下排列，即世界四方的名称：\OriginalTerm{东方}{\GreekText{ἀνατολή}}、\OriginalTerm{西方}{\GreekText{δύσις}}、\OriginalTerm{北方}{\GreekText{ἄρχτος}}、\OriginalTerm{南方}{\GreekText{μεσημβρία}}，这就是整个地球——因此主说，当祂来审判时，要从四方聚集祂的选民（\BibleRef{谷 13:27}）——我说，如果你们取这四个希腊词——\OriginalTerm{东方}{\GreekText{ἀνατολή}}、\OriginalTerm{西方}{\GreekText{δύσις}}、\OriginalTerm{北方}{\GreekText{ἄρχτος}}、\OriginalTerm{南方}{\GreekText{μεσημβρία}}，即Anatole、Dysis、Arctos、Mesembria——这些词的首字母组成Adam。那么，我们如何在其中也找到数字四十六呢？因为基督的肉身来自亚当。希腊人用字母计算数字。我们称为A的字母，他们在他们的语言中称为Alpha（\OriginalTerm{Alpha}{\GreekText{α}}），而Alpha（\GreekText{α}）称为一。而在数字中他们写下Beta（\OriginalTerm{Beta}{\GreekText{β}}），即他们的\Emph{\OriginalTerm{b}{\GreekText{β}}}，在数字中称为二。他们写下Gamma（\OriginalTerm{Gamma}{\GreekText{γ}}），在数字中称为三。他们写下Delta（\OriginalTerm{Delta}{\GreekText{δ}}），在数字中称为四；并且如此，借助所有字母，他们有数字。我们称为M的字母，他们称为My（\OriginalTerm{My}{\GreekText{μ}}），表示四十；因为他们说My（\GreekText{μ}）\OriginalTerm{四十}{\GreekText{τεσσαράχοντα}}。现在看这些字母组成的数字，你就会发现圣殿建造了四十六年。因为单词Adam有Alpha（\OriginalTerm{Alpha}{\GreekText{α}}），即一；它有Delta（\OriginalTerm{Delta}{\GreekText{δ}}），即四；加起来是五；它又有Alpha（\OriginalTerm{Alpha}{\GreekText{α}}），即一；加起来是六；它还有My（\OriginalTerm{My}{\GreekText{μ}}），即四十；这就有了四十六。我的弟兄们，这些事是由我们之前的教父告诉我们的，而数字四十六是由他们从字母中发现的。而且，因为我们的主耶稣基督从亚当取了身体，并非从亚当取罪；从祂取了属肉体的圣殿，并非那必须从圣殿中赶出的罪孽；并且犹太人把祂从亚当所取的肉身钉在十字架上（因为玛利亚来自亚当，主的肉身来自玛利亚）；而且，祂要在三天内复活那同一肉身，就是他们即将在十字架上杀害的：他们拆毁了那用了四十六年建成的圣殿，而祂在三天内把它重建起来。\par

13. 我们赞美主、我们的天主，祂召集我们共赴属灵的喜乐。让我们常怀谦卑之心，愿我们的喜乐与祂同在。我们不要因世间任何顺遂而得意，要知道我们的幸福要等到这些事都过去之后才会来临。如今，我的弟兄们，让我们的喜乐在于望德：不要有人为眼前之事而喜，免得他滞留途中。愿喜乐全然来自将来的望德，愿渴望全然属于永生。愿一切叹息都向往基督。愿惟有那至美者——祂爱污秽者使其成为美丽——成为我们一切的渴望；让我们只追随祂，只为祂渴望叹息；\Quote{愿他们常说：愿主被尊崇，那些愿意祂仆人平安的人。}\par

\SourceRecord{Translated by John Gibb.；Nicene and Post-Nicene Fathers, First Series；Vol. 7.；Edited by Philip Schaff.；Buffalo, NY: Christian Literature Publishing Co.,；1888.；<http://www.newadvent.org/fathers/1701010.htm>.}

% structure: p0001 source=h1 target=section reason=page-title

\section{讲道集第11篇（\BibleRef{若 2:23-3:5}）}

1. 主适当地为我们安排，使这段经文今日按顺序出现：因为我想你们已经注意到，亲爱的诸位，我们已着手按时序思考并讲解\BibleBook{若望}{福音}。今日适逢其会，你们当从福音中听到：\BibleQuote{人除非由水和圣神重生，不能见到天主的国。}（\BibleRef{若 3:5}）此刻正是我们劝勉你们这些仍为慕道者的时候，你们虽然相信了基督，却仍负着罪债。没有人能带着罪债见到天主的国；因为只有罪获得赦免的人，才能与基督一同为王；而罪除非借着水和圣神重生，是无法赦免的。让我们观察这些话的含义，使懒惰的人能明白他们必须以何等迫切的心情卸下重担。假如他们背负着沉重的石头、木料，甚至某种收益，或是搬运谷物、酒、金钱，他们会奔跑着卸下负担；他们背负着罪的负担，却懒得奔跑。你们必须奔跑以卸下这重担；它压着你们，它要淹死你们。\par

2. 看，你们已经听见，当我们的主耶稣基督\BibleQuote{在耶路撒冷过逾越节庆日时，许多人看见祂所行的神迹，便信了祂的名。}（\BibleRef{若 2:23}）\BibleQuote{许多人信了祂的名}；接下来呢？\BibleQuote{耶稣却不信任他们。}（\BibleRef{若 2:24}）这是什么意思呢？\BibleQuote{他们信了祂的名}，或信靠，而\BibleQuote{耶稣却不信任他们}？或许是因为他们并没有真正相信祂，只是假装相信，所以耶稣不信任他们？但福音的作者若没有作真实见证，就不会说\BibleQuote{许多人信了祂的名}。这真是一件伟大而奇妙的事：人相信基督，基督却不信任人。尤其奇妙的是，祂既是天主子，自然是自愿受苦。若祂不愿意，就绝不会受苦，因为若祂不愿意，就不会诞生；若祂只愿意诞生而不愿意死，祂也能按意愿行事，因为祂是全能圣父的全能圣子。让我们用事实来证明。当他们想要捉拿祂时，祂离开了他们。福音说：\BibleQuote{他们企图把祂从山顶推下去，祂却从他们中间穿过，安然无恙地离开了。}（\BibleRef{路 4:30}）当叛徒犹达斯出卖祂后，他们来捉拿祂时——犹达斯以为他能交付他的师傅和主——在那里主也显明祂是自愿受苦，而非出于不得已。因为当犹太人想要捉拿祂时，祂对他们说：\BibleQuote{你们找谁？他们回答说：纳匝肋人耶稣。祂说：我就是。他们一听这话，就倒退跌倒在地。}（\BibleRef{若 18:4}）在此，祂以回答将他们击倒在地，显示了祂的能力；而祂被他们捉拿，则显示了祂的意愿。因此，祂是出于怜悯而受苦。因为\BibleQuote{祂被交付，是为了我们的过犯；祂复活，是为了使我们成义。}（\BibleRef{罗 4:25}）请听祂自己的话：\BibleQuote{我有权舍去我的生命，也有权再取回来；没有人能夺去我的生命，而是我自愿舍去，为能再取回来。}（\BibleRef{若 10:18}）既然祂有这样的能力，以言语宣告、以行为证明，那么，耶稣不信任他们是什么意思？好像他们会违背祂的意愿加害于祂，或做任何违背祂意愿的事？尤其是他们已经信了祂的名。况且，论到这些人，福音作者说\BibleQuote{他们信了祂的名}，却又说\BibleQuote{耶稣却不信任他们}。为什么呢？\BibleQuote{因为祂认识众人，不需要别人为祂作证，因为祂自己知道人心里有什么。}（\BibleRef{若 2:25}）工匠比祂的作品更知道作品内部。人的创造主知道人心里有什么，而受造的人自己却不知道。我们不是从伯多禄身上证明了吗？他不知道自己心里有什么，当他说：\BibleQuote{同祢一起，甚至到死。}请听主知道人心里有什么：\BibleQuote{你愿意同我一起到死吗？我实实在在告诉你：鸡叫以前，你要三次不认我。}（\BibleRef{若 13:37-38}）因此，人不知道他自己心里有什么，但人的创造主知道人心里有什么。然而，许多人信了祂的名，耶稣却不信任他们。我们能说什么呢，弟兄们？也许接下来的内容会向我们揭示这些话语的奥秘。他们信了祂，这是显而易见的，真实的，无人怀疑；福音这样说，诚实的福音作者作证。再者，耶稣不信任他们，也是显而易见的，没有一个基督徒怀疑；因为福音也这样说，同一位诚实的福音作者作证。那么，为什么他们信了祂的名，耶稣却不信任他们呢？让我们看看接下来是什么。\par

3. \Quote{有一个法利塞人，名叫尼苛德摩，是犹太人的首领。这人夜里来见耶稣，对祂说：辣彼——你知道师傅被称为辣彼——我们知道你是由天主而来的师傅，因为没有人能行你所行的这些神迹，除非天主与他同在。} 这尼苛德摩原是属于那些因看见祂所行的神迹奇事而信了祂名的人。因为这正是上文所说的：\Quote{那时，耶稣在耶路撒冷过逾越节庆日，有许多人信了祂的名。} 他们为什么信？他接着说：\Quote{看见祂所行的神迹。} 关于尼苛德摩，他又说什么？\Quote{有一个犹太人的首领，名叫尼苛德摩，他夜里来见耶稣，对祂说：辣彼，我们知道你是由天主而来的师傅。} 因此这人也在祂的名内有了信仰。他为什么信？他接着说：\Quote{因为没有人能行你所行的这些神迹，除非天主与他同在。} 所以，如果尼苛德摩是属于那些信了祂名的人，那么现在让我们就此尼苛德摩的情况，考虑一下为什么耶稣不肯把自己交托给他们。\Quote{耶稣回答他说：我实实在在告诉你：人若不重生，就不能见到天主的国。} 因此，耶稣把自己交托给那些重生的人。请看，那些人已经信了祂，但耶稣却不把自己交托给他们。所有\OriginalTerm{慕道者}{catechumens}都是如此：他们已经信了基督的名，但耶稣却不把自己交托给他们。亲爱的诸位，请留心听，并要明白。如果我们问一位慕道者：\Quote{你信基督吗？} 他回答说：\Quote{我信}，并划十字圣号；他已经把基督的十字架印在额上，不以主的十字架为耻。请看，他已在祂的名内相信了。我们问他：\Quote{你吃人子的肉，喝人子的血吗？} 他不知道我们在说什么，因为耶稣尚未把自己交托给他。\par

4. 因此，既然尼苛德摩属于那一类人，他来见主，却是在夜里来的；这或许与主题有关。他来见主，却是在夜里来；他来见光明，却是在黑暗中来的。但那由水和圣神重生的人，从宗徒那里听到什么呢？\BibleQuote{你们从前是黑暗，如今在主内却是光明；当像光明之子一样行走} \BibleRef{弗 5:8} 又说：\BibleQuote{但我们属于白昼的，应当清醒。} \BibleRef{得前 5:8} 因此，重生的人本属于黑夜，如今却属于白昼；本是黑暗，如今却是光明。现在耶稣把自己交托给他们，他们来见耶稣，不像尼苛德摩在夜里，不在黑暗中寻求白昼。因为这样的人现在也承认：耶稣已临近他们，在他们身上完成了救恩；因为祂说过：\BibleQuote{谁不吃我的肉，不喝我的血，在他内就没有生命。} \BibleRef{若 6:54} 慕道者虽然已经在额上印有十字架的记号，他们已属于大户人家；但要从仆人成为儿子。因为他们已属于大户人家，但尚有所不足。但以色列百姓是在什么时候吃\OriginalTerm{玛纳}{manna}的呢？是在他们渡过\OriginalTerm{红海}{Red Sea}之后。至于红海所代表的意义，请听宗徒的话：\BibleQuote{弟兄们，我不愿意你们不知道，我们的祖先都曾在云柱下，都从海中经过，} 他们经过海是为了什么？如果你问他，他接着会说：\BibleQuote{都曾在云中和海中受洗归于梅瑟。} \BibleRef{格前 10:1} 如果海的\OriginalTerm{预像}{figure}有如此大的效力，那么圣洗的真实形态该有多么大的效力！如果那在预像中成就的事，使百姓在渡海后来到玛纳，那么基督藉着祂圣洗的真理，将祂自己的子民通过祂自己领过海后，要赐给他们什么呢？藉着祂的圣洗，祂领那些信从的人过海；他们的一切罪过，如同追逐他们的仇敌，都被消灭，正如所有埃及人都在海中灭亡一样。祂领他们到哪里去，我的弟兄们？耶稣藉着圣洗领他们到哪里去，就是梅瑟当年领他们经过海时所预像的那里？到哪里？到玛纳那里。玛纳是什么？祂说：\BibleQuote{我是从天上降下的生活的食粮。} \BibleRef{若 6:51} 那些信徒已经渡过红海，他们领受玛纳吗？为什么是红海？除了是海，为什么还称为\Quote{红}？这\Quote{红海}预表基督的圣洗。基督的圣洗如何是红的？岂不是因基督的血而成圣吗？那么，祂领那些信从并受洗的人到哪里？到玛纳那里。请看，我说\Quote{玛纳}：犹太人，即以色列子民，所领受的是什么，是众所周知的，是他们从天上降下的玛纳，没有人不知道；然而慕道者不知道基督徒所领受的是什么。那么，让他们为自己的无知而羞愧吧；让他们渡过红海，吃玛纳吧，好使他们既然信了耶稣的名，同样耶稣也把自己交托给他们。\par

5. 所以，我的弟兄们，请注意这位夜间来到耶稣面前的人是如何回答的。虽然他来到了耶稣面前，但由于他是在夜间来的，所以他仍从自己肉身的黑暗中发言。他不明白他从主那里听到的，不明白他从那\BibleQuote{光照每一個人來到這世界的真光}那里听到的。\BibleRef{若 1:9} 主已经对他说：\Quote{人若不重生，不能见到天主的国。尼苛德摩对祂说：人老了，如何能重生呢？} 圣神对他说话，他却思念肉身。他思念自己的肉身，因为他还未思念基督的肉身。因为主耶稣曾说：\Quote{除非人吃我的肉，喝我的血，在他内便没有生命。} 有些跟随祂的人因而反感，彼此说：\Quote{这话生硬，谁能听得下去？} 因为他们幻想，耶稣说这话的意思是他们将能像宰杀羔羊一样把祂切碎煮来吃。他们被祂的话吓住，就退去了，不再跟随祂。圣史这样记载：\Quote{主自己与那十二人留下；他们对祂说：看，那些人离开了你。祂说：你们也要离开吗？}——祂愿意向他们表明，祂是他们所必需的，而不是他们为基督所必需。愿无人幻想，当他告诉基督他是基督徒时，他能吓倒基督；仿佛你若成为基督徒，基督会更蒙福。你成为基督徒，这为你是一件好事；但若你不成为基督徒，基督也不会受损。请听圣咏之声：\Quote{我曾对主说：你是我的天主，因为你不需要我的财物。} 正因如此，\Quote{你是我的天主，因为你不需要我的财物。} 若你没有天主，你将变小；若你与天主同在，天主并不会变得更大。祂不会因你而变得更大，但你离了祂就会变小。所以，要在祂内成长；不要退去，使祂仿佛有所减损。你若来到祂跟前，你将得到更新；你若离开祂，你将遭受损失。你来到祂跟前时，祂保持完整，你跌倒时，祂也保持完整。因此，当祂对门徒们说：\Quote{你们也要离开吗？} 彼得，那磐石，以众人的声音回答说：\Quote{主，我们往哪里去？你有永生的话。} 主的肉在他的口中甘甜。然而，主给他们解释并说：\Quote{使人生活的乃是圣神。} 在祂说了\Quote{除非人吃我的肉，喝我的血，在他内便没有生命}之后，为避免他们作血肉的理解，祂说：\Quote{使人生活的乃是圣神，肉身一无所益；我对你们所说的话，就是神，就是生命。} \BibleRef{若 6:54}\par

6. 这位夜间来到耶稣面前的尼苛德摩，并未品味这神和这生命。耶稣对他说：\Quote{人若不重生，不能见到天主的国。} 而他，品味他自己的肉身，尚未在口中品味基督的肉，就说：\Quote{人老了，如何能重生呢？难道他能再入母腹而重生吗？} 这个人只知道一种出生，即出自亚当和厄娃的出生；他还不知道那出自天主和教会的出生：他只认识那些产生死亡的双亲，还不认识产生生命的双亲；他只认识产生后继者的双亲，还不认识那产生永存之人的永生的双亲。\par

既然有两种出生，那么他只理解了一种。一种是出于地的，另一种是出于天的；一种是出于肉身的，另一种是出于圣神的；一种是出于可朽的，另一种是出于永恒的；一种是出于男女的，另一种是出于天主和教会的。然而，这两种出生各自单一，不可重复其中之一或另一。尼苛德摩正确地理解了肉身的出生；你们也要这样理解圣神的出生，正如尼苛德摩理解肉身的出生一样。尼苛德摩理解了什么呢？\Quote{难道人能再入母腹而重生吗？} 因此，若有人告诉你要在圣神内重生第二次，你要用尼苛德摩的话回答：\Quote{难道人能再入母腹而重生吗？} 我已由亚当所生，亚当不能再生我一次。我已由基督所生，基督不能再生我。正如从母腹出生不可重复，同样，从圣洗也不可重复。\par

7. 那由天主教会所生的人，就好像是由撒辣所生，由那自由妇所生；那由异端所生的人，就好像是由婢女所生，但仍是亚巴郎的后裔。亲爱的，请思考这奥秘是多么伟大。天主作证说：\BibleQuote{我是亚巴郎的天主，依撒格的天主，雅各伯的天主。} 难道没有其他圣祖吗？在这些圣祖之前，不是有圣诺厄吗？他在全人类中，独与他全家，堪当从洪水中得救——在他内，以及在他儿子们内，教会已经预表了。他们乘着方舟，逃脱了洪水。后来还有我们所知的伟人，圣经所称许的梅瑟，\BibleQuote{他在祂的全家中是忠信的。} \BibleRef{户 12:7} 然而，那三位圣祖却被特别提名，仿佛唯独他们配得天主的欢心：\BibleQuote{我是亚巴郎的天主，依撒格的天主，雅各伯的天主：这是我的名直到永远。} \BibleRef{出 3:6} 崇高的奥秘！是主能开启我们口和你们的心，使我们能照祂所愿启示的讲论，也使你们能按对你们有益的去领受。\par

8. 圣祖就是这三：亚巴郎、依撒格和雅各伯。你们知道雅各伯的儿子共有十二个，从此产生了以色列民族；因为雅各伯本人就是以色列，以色列民族由十二支派组成，对应于以色列的十二个儿子。亚巴郎、依撒格、雅各伯三位圣祖，一个民族。三位圣祖，可以说是民族的开端；民族在他们身上被预表；而先前的民族就是现在的民族。因为在犹太民族中预表了基督徒民族。那里是预像，这里是真理；那里是影子，这里是实体：正如宗徒所说：\BibleQuote{这些事发生在他们身上，乃是预像。} 这是宗徒的声音：\BibleQuote{它们被写下来，} 他说，\BibleQuote{是为了我们这些世末已临到的人。} \BibleRef{格前 10:11} 现在请回想亚巴郎、依撒格和雅各伯。关于这三位的例子，我们发现自由妇人生子，婢女也生子；我们看到自由妇人的后裔，也看到婢女的后裔。婢女所象征的并非善事：\Quote{赶出这婢女和她的儿子，} 他说，\Quote{因为婢女的儿子不能与自由妇人的儿子一同继承产业。} 宗徒叙述了这事，并说亚巴郎的两个儿子预表了旧约和新约两个盟约。属于旧约的是贪恋暂时事物、贪恋世界的人；属于新约的是贪恋永生的人。因此，地上的耶路撒冷是天上的耶路撒冷——我们众人的母亲——的影子，她是天上的；这些都是宗徒的话。关于那座我们旅居时远离的城，你们知道很多，现在也听了很多。但在这些生育中，在这些胎孕的果实中，在这些自由妇人和婢女的世代中，我们发现一件奇妙的事：即四种人；在这四种人中，未来基督徒民族的预像得以完成，所以针对那三位圣祖所说的话并不奇怪：\BibleQuote{我是亚巴郎的天主，依撒格的天主，雅各伯的天主。} 因为在所有基督徒中，弟兄们，请注意，要么善人生于恶人，要么恶人生于善人；要么善人生于善人，要么恶人生于恶人；你们找不到超过这四种类型。我将再次重复：要注意，要记住，要激发你们的心，不要迟钝；要接受，免得你们被掳，关于所有基督徒有四种类型。要么善人生于善人，要么恶人生于恶人；要么善人生于恶人，要么恶人生于善人。我想这是清楚的。善人生于善人：如果施洗者是善的，而受洗者也正确地相信，并正确地被列在基督的肢体中。恶人生于恶人：如果施洗者是恶的，而受洗者以心怀二意接近天主，不遵守他们在教会中听到的劝诫的道德，以至于不是糠秕，而是麦粒。亲爱的，你们知道这样的人有多少。恶人生于善人：有时一个通奸者给人付洗，而受洗者却成了义人。善人生于恶人：有时施洗者是圣的，受洗者却不愿意遵守天主的道。\par

9. 弟兄们，我想这在教会中是众所周知的，我们所说的也因日常的例子而显明；但让我们看看我们之前的圣祖们的情况，看他们也有这四种类型。善人生于善人：阿纳尼雅给保禄付洗。恶人生于恶人：宗徒声明，有些福音宣讲者，他说，不以纯洁的动机宣讲福音，然而他容忍他们在基督徒团体中，说：\BibleQuote{那又如何呢？无论如何，无论是出于假意，还是出于真诚，基督被宣讲了，为此我欢喜。} \BibleRef{斐 1:18} 难道他心怀恶意，为别人的恶而欢喜吗？不，而是因为藉着恶人，真理被宣讲了，藉着恶人的口，基督被宣讲了。如果这些人给像他们一样的人付洗，那就是恶人给恶人付洗；如果他们给主所告诫的那样的人付洗，当主说：\BibleQuote{凡他们所吩咐你们的，你们要遵守实行；但不要效法他们的行为，} \BibleRef{玛 23:3} 那么他们是恶人给善人付洗。善人给恶人付洗，就如西满术士被\Person{斐理伯}——一位圣洁的人——付洗一样。\BibleRef{宗 8:13} 所以，我的弟兄们，这四种类型是众所周知的。看，我再次重复它们，要记住它们，数算它们，默想它们；提防恶的，持守善的。善人生于善人：当圣洁的人由圣洁的人付洗时；恶人生于恶人：当施洗者和受洗者都生活不义和不敬虔时；善人生于恶人：当施洗者是恶的，而受洗者是善的时；恶人生于善人：当施洗者是善的，而受洗者是恶的时。\par

10. 我们在这三个名字中如何发现这一点，\BibleQuote{我是亚巴郎的天主、依撒格的天主、雅各伯的天主}？我们将婢女视为恶者，将自由之女视为善者。自由之女生善者；撒辣生了依撒格；婢女生恶者；哈加尔生了依市玛耳。仅在亚巴郎的事例中，我们就看到两种情形：既有善者出于善者，也有恶者出于恶者。但是我们在哪里看到恶者出于善者的预表？黎贝加，依撒格的妻子，是一位自由之女：请读，她生了双生子；一个是善的，另一个是恶的。天主借圣经明确宣告：\BibleQuote{我爱了雅各伯，却恨了厄撒乌。}黎贝加生了这两个，雅各伯和厄撒乌：一个被拣选，另一个被弃绝；一个承受产业，另一个被剥夺继承权。天主没有从厄撒乌中形成祂的子民，而是从雅各伯中形成。种子是同一位，但所怀的胎却不同；子宫是同一个，但所生的子嗣却各异。那位生了雅各伯的自由之女，不也是生了厄撒乌的同一个自由之女吗？他们在母亲的腹中争斗；当他们在那里争斗时，有话说给黎贝加：\BibleQuote{两个民族在你腹中。}两个人，两个民族；一个好民族，一个坏民族：但他们仍在同一个腹中争斗。教会中有多少恶人！同一个母腹怀抱着他们，直到最后他们被分开：善人向恶人呼喊，恶人也向善人呼喊，两者在同一个母亲的肠腹中一同争斗。他们会永远在一起吗？最终要进入光明；这里以奥秘预表的诞生将要显明；那时将会显明\BibleQuote{我爱了雅各伯，却恨了厄撒乌。}\par

11. 所以，弟兄们，我们现在已经发现：善者出于善者——自由之女所生的依撒格；恶者出于恶者——婢女所生的依市玛耳；善者出于恶者——黎贝加所生的厄撒乌；那么我们在哪里找到恶者出于善者呢？还剩下雅各伯，这样这四种类型可以在三位圣祖身上完成。雅各伯有自由之女为妻，也有婢女为妻：自由之女生育，婢女也生育，这样就有了以色列的十二个儿子。如果你数算他们所有人，看他们由谁所生，他们并不全都是自由之女所生，也不全都是婢女所生；但他们都是同一粒种子所出。那么，我的弟兄们，那些由婢女所生的儿子难道没有与他们的弟兄一同承受应许之地吗？我们在那里发现雅各伯的好儿子中有婢女所生的，也有自由之女所生的。他们从婢女腹中所生对他们毫无妨碍，因为他们认识自己在父亲里的种子，因此他们与弟兄一同拥有国度。所以，正如雅各伯的儿子们，他们由婢女所生并不妨碍他们拥有国度，并与他们的弟兄平等地承受应许之地；他们由婢女所生没有妨碍他们，而是父亲的种子得胜了：同样，凡是由恶人施洗的，看似如同由婢女所生；然而，因为他们属于\OriginalTerm{圣言}{Word of God}的种子——这在雅各伯身上预表——他们不要气馁，他们将与弟兄们一同承受产业。所以，那由好种子所生的不要惧怕；只要他不要模仿那婢女，如果他是由婢女所生的话。不要模仿那邪恶、骄傲的婢女。因为雅各伯由婢女所生的儿子们为何能与他们的弟兄一同承受应许之地，而由婢女所生的依市玛耳却被赶出继承之外呢？为何？岂不是因为他骄傲，而他们谦卑吗？他骄傲地昂起脖子，想要在与他的弟兄玩耍时引诱他。\par

12. 那里有一个伟大的奥秘。他们在一起玩耍，依市玛耳和依撒格；撒辣看见他们玩耍，就对亚巴郎说：\BibleQuote{赶出这婢女和她的儿子，因为婢女的儿子不能与我的儿子依撒格一同承受产业。}当亚巴郎为此忧愁时，主向他确认了他妻子的话。现在这里显然是一个奥秘，这件事以某种方式孕育着未来的事。她看见他们玩耍，就说：\BibleQuote{赶出这婢女和她的儿子。}这是什么意思，弟兄们？因为依市玛耳在与孩子依撒格玩耍时，对他做了什么恶事呢？那玩耍是一种嘲弄；那玩耍意味着欺骗。现在请留意，亲爱的，这个伟大的奥秘。宗徒称之为迫害；那玩耍，那游戏，他称之为迫害：因为他说道：\BibleQuote{然而，从前那按血肉生的，迫害了那按圣神生的，现在也是这样。}也就是说，按血肉生的迫害那按圣神生的。谁是按血肉生的？爱世界的人，爱今世生活的人。谁是按圣神生的？爱天国的人，爱基督的人，渴望永生、自由敬拜天主的人。他们玩耍，宗徒却称之为迫害。因为在他说了这些话之后：\BibleQuote{然而，从前那按血肉生的，迫害了那按圣神生的，现在也是这样。}宗徒接着说，并指出了他所说的是何种迫害：\BibleQuote{但经上说了什么？赶出婢女和她的儿子，因为婢女的儿子不能与我的儿子依撒格一同承受产业。}我们查找经上何处说了这话，以查看是否依市玛耳对依撒格有任何迫害发生在这之前；我们发现这话是撒辣看见孩子们一起玩耍时说的。撒辣所看见的玩耍，宗徒称之为迫害。因此，那些用玩耍引诱你们的人，反而是更加迫害你们。他们说：\Quote{来，来这里受洗，这里有真正的洗礼。}不要玩耍，只有一个真正的洗礼；那另一个是游戏：你们将被引诱，那将是对你们的一种严重迫害。你们最好将王国当作礼物送给依市玛耳；但依市玛耳不会接受，因为他想要玩耍。守住你父亲的产业，听这话：\BibleQuote{赶出这婢女和她的儿子，因为婢女的儿子不能与我的儿子依撒格一同承受产业。}\par

13. 这些人竟敢说，他们常遭受天主教君王或天主教王侯的迫害。他们受的是什么迫害？身体的痛苦。即使他们有时受过苦，并且如何受苦的，让他们自己知道，并与自己的良心解决；他们毕竟只受了身体的苦。他们所施加的迫害更为严重。当依市玛耳想与依撒格玩耍，当他向你谄媚，当他提供另一个圣洗时，你要当心。回答他说：我已经有了圣洗。因为如果这个圣洗是真的，那么要给你另一个圣洗的人就是在嘲弄你。要当心灵魂的迫害。虽然多纳徒党人在天主教王侯手中有时也受过一些苦，但那只是身体的痛苦，而不是精神欺骗的痛苦。请听并看在旧约历史的事实中，所有未来事物的预兆和迹象。我们看到撒辣使她的婢女哈加尔受苦：撒辣是自由的。当她的婢女开始骄傲时，撒辣向亚巴郎抱怨说：\Quote{赶走这婢女；}她竟在我面前昂头。他的妻子向亚巴郎抱怨，好像是他做的。但亚巴郎，他并非因情欲而与婢女结合，而是出于生养子女的责任，因为撒辣将婢女给了他，要她为他生子，就对撒辣说：\Quote{看，她是你的婢女；你随意待她吧。}撒辣就苦苦折磨她，她便从她面前逃走了。看，自由的女人折磨了婢女，宗徒却不称之为迫害；奴隶与主人玩耍，他却称之为迫害。这种折磨不被称作迫害，那种玩耍却被称作迫害。弟兄们，你们怎么看？难道你们不明白这预示什么吗？因此，当天主愿意激起权柄反对异端者、分裂者、扰乱教会者、像厌恶基督那样攻击祂者、亵渎圣洗者时，让他们不要惊讶；因为天主激起他们，是为了让哈加尔被撒辣责打。让哈加尔认识自己，并屈服她的颈项；因为当她受辱离开主母时，一位天使遇见她，对她说：\Quote{哈加尔，撒辣的婢女，你怎么了？}当她抱怨主母时，她听到天使说了什么？\BibleQuote{回到你主母那里去。}\BibleRef{创 16:9}她因此受折磨，是为叫她回去；但愿她回去，因为她的后裔，正如雅各伯的儿子们，将同他们的弟兄一起获得产业。\par

14. 但他们惊讶于基督徒的权柄被激起反对那可憎的扰乱教会者。难道他们不该被激动吗？不然他们该如何向天主交账？亲爱的，请注意我说的话：这世上的基督徒君王关心的是，他们的母亲教会（他们由她而灵性重生）能在他们的时代享平安。我们读了达内尔的异象和先知史。三个青年在火中赞美主：拿步高王惊讶于青年赞美天主，而火焰对他们无害；当他惊讶时，拿步高王（他既不是犹太人，也未受割礼，曾立了自己的像并强迫所有人朝拜它）因见到三个青年的赞美而深受感动，看到天主的威严临在火中，他说了什么？\BibleQuote{我要向全地的一切民族、语言颁布一道谕令。}怎样的谕令？\BibleQuote{凡说话亵渎沙得客、默沙客、阿贝德乃哥的天主的，必被处死，他们的房屋要成为废墟。}\EditorialNote{达内尔第三章}看，一个外邦君王怀着愤怒的义怒行事，以使以色列的天主不被亵渎，因为祂能将三个青年从火中救出；而他们却不愿让基督徒君王严厉对待基督被藐视拒绝的事——那基督不仅救了三个青年，更是救了整个世界，连同这些君王自己，脱离地狱之火！因为那三个青年，我的弟兄，是从暂时的火中被救的。难道祂不是同一位天主，即玛加伯的天主和三个青年的天主吗？后者祂从火中救出；前者身体在火刑中灭亡，但心灵坚守法律的诫命。后者公开被救，前者秘密地得了殉道者的冠冕。\EditorialNote{玛加伯下第七章}从地狱的火焰中被救，比从人间权势的火窑中被救更为伟大。那么，如果拿步高因天主从火中救出三个青年而赞美、颂扬、光荣天主，并且如此光荣以致在他的王国中颁布谕令：\BibleQuote{凡说话亵渎沙得客、默沙客、阿贝德乃哥的天主的，必被处死，他们的房屋要成为废墟，}那么，那些看到自己不是三个青年从火焰中被救，而是自己从地狱中被救的君王，当他们看到基督（他们藉祂而得救）在基督徒身上被藐视凌辱，当听到有人对基督徒说\Quote{你说你不是基督徒}时，他们怎能不被激动呢？人们愿意行这样的事，却不愿承受这样的惩罚。\par

15. 请看他们做些什么，受些什么。他们杀害灵魂，在肉身受苦；他们造成永死，却抱怨自己遭受暂时的死亡。然而他们遭受的是怎样的死亡呢？他们向我们提出他们的一些被迫害的殉道者。看，\Person{玛尔库鲁斯}被从岩石上抛下；看，\Place{巴加亚}的\Person{多纳图}被投入井中。罗马官府何时曾规定过把人从岩石上抛下这样的刑罚？但我们这一方的人又怎样回答呢？实际发生了什么我不知道，但我们的人说了什么？他们说那些人自己跳下去，并把恶名归咎于官府。让我们回想罗马官府的惯例，并看看我们该相信谁。我们的人断言那些人自己跳下去。如果他们不是那些人的门徒——那些人如今在没有受迫害时也自己跳下悬崖——那么我们就不要相信我们的人的指控：那些人做出这些人之常做的事，又有什么奇怪呢？罗马官府从未使用过这样的刑罚：难道他们没有权力公开处死他们吗？但那些人，他们希望死后受尊崇，却找不到一种更能使他们出名的死法。总之，事实如何，我不知道。即使你们在肉身受苦，哦，多纳图派，是在天主教会手中，正如\Person{哈加尔}在\Person{撒辣}手中受苦一样；\BibleQuote{回到你的主母那里去}。一个确实有必要讨论的问题使我们耽搁得太久，以致完全无法讲解全部福音课文。在此期间，亲爱的弟兄们，愿这一点足够你们领受，免得因谈论其他事，已说的话被排除在你们心外。持守这些事，宣讲这些事；你们自己既已炽热，就前往那里，点燃那些冷淡的人。\par

\SourceRecord{Translated by John Gibb.；Nicene and Post-Nicene Fathers, First Series；Vol. 7.；Edited by Philip Schaff.；Buffalo, NY: Christian Literature Publishing Co.,；1888.；<http://www.newadvent.org/fathers/1701011.htm>.}

% structure: p0001 source=h1 target=section reason=page-title

\section{第十二篇讲道（\BibleBook{若望}{福音} 3:6-21）}

亲爱的各位，我们观察到，昨日我们以预告激发你们的注意，使你们比往常更踊跃、更众多地聚集。但此刻，若你们愿意，让我们先履行按序而来的福音读经的讲道义务。然后，亲爱的各位，你们将听到我们已为教会和平所成就的事，以及我们期望进一步完成的事。因此，现在，让你们全部的心思都专注在福音上；不要有人思考别的事。因为若全神贯注的人尚且难以领悟，那么因杂念而分心的人岂不更会失去已领受的？此外，亲爱的各位，你们记得，在上主日，因主的恩助，我们讲述了关于属神的重生。我们已吩咐将那段读经再次读给你们听，以便我们现今能借你们在基督之名下的祈祷，完成当时未讲完的部分。\par

属神的重生是唯一的，正如肉身的生育也是唯一的。尼苛德摩向主说真话：人老了，不能重新进入母胎而出生。他确实说人老了不能这样做，仿佛他在婴儿时就能做似的。但无论是刚出母腹，还是年事已高，人都不可能再回到母腹中出生。然而，正如肉身的生育，女人的子宫只能一次分娩孩子，同样，为属神的出生，教会的子宫只能一次使人受洗。因此，若有人说：\Quote{但这个人是在异端中出生的，那个人是在分裂中出生的；}那么，一切都被切除了，如果你们还记得我们曾向你们讲解过的关于我们三位祖先的事，天主愿意被称为他们的天主，并非因为他们独自如此，而是因为在它们身上完整地预表了未来子民的形象。因为我们发现：婢女所生者被剥夺继承权，自由妇人所生者成为继承人；同样，自由妇人所生者被剥夺继承权，婢女所生者成为继承人。依市玛耳，婢女所生，被剥夺继承权；依撒格，自由妇人所生，成为继承人；厄撒乌，自由妇人所生，被剥夺继承权；雅各伯的儿子们，婢女所生，成为继承人。因此，在这三位祖先身上，预见了整个未来子民的形象；天主说：\InnerQuote{我是亚巴郎的天主，依撒格的天主，雅各伯的天主；这是我的名，直到永远。}（\BibleRef{出 3:6}）不如我们纪念向亚巴郎本人所应许的：因为这同样应许给了依撒格和雅各伯。我们读到什么？\InnerQuote{地上万民都要因你的后裔蒙受祝福。}（\BibleRef{创 22:18}）那时，这一个人相信了他尚看不见的事；如今人们看见了，却被弄瞎了。向一个人所应许的，在万民中实现了；而那些不愿看见已实现之事的人，正从万民的共融中分离自己。但他们不愿看见又有什么用呢？他们无论愿意与否，都看见了；公开的真理撞击着他们紧闭的眼睛。\par

是对尼苛德摩的回答——他属于那些信了耶稣的人——说：耶稣却不信任他们。确实，祂不信任某些人，尽管他们已经信了祂。经上记着：\BibleQuote{许多人看见祂所行的神迹，就信了祂的名。但耶稣却不信任他们。因为祂不需要谁对人有见证；祂自己知道人心里有什么。}看，他们已经信了耶稣，但耶稣却不信任他们。为什么？因为他们还没有由水和圣神重生。由此我们曾劝勉，并现今劝勉我们的弟兄——慕道者。因为若问他们，他们已经信了耶稣；但因尚未领受祂的体和血，耶稣尚未信任他们。他们必须做什么才能使耶稣信任他们？他们必须由水和圣神重生；那正在为他们分娩的教会必须产生他们。他们已被孕育；他们必须被带到光中；他们有乳房可受哺育；但愿他们不怕出生后被窒息；但愿他们不离开母亲的乳房。\par

4. 没有人能回到母亲的腹中再出生。但是，有人是从婢女所生的吗？好吧，从前由婢女所生的人，难道他们曾回到自由的腹中重新出生吗？\Person{亚巴郎}的后裔也存在于\Person{依市玛耳}身上；但为使\Person{亚巴郎}从婢女得一个儿子，这是出于他妻子的建议。这孩子是丈夫的后裔，不是出于子宫，而是完全出于妻子的意愿。他是由婢女所生的这一事实，难道是他被剥夺继承权的原因吗？那么，如果他是因婢女之子而被剥夺继承权，就没有婢女之子能被接纳进入继承。\Person{雅各伯}的儿子们被接纳进入继承，但\Person{依市玛耳}却被排除在外，不是因为由婢女所生，而是因为他向他的母亲、向他母亲的儿子骄傲；因为他的母亲其实是\Person{撒辣}，而不是\Person{哈加尔}。一个提供了她的子宫，另一个的意愿被附加了：\Person{亚巴郎}不会做\Person{撒辣}不愿意的事，因此他更是\Person{撒辣}的儿子。但是，因为他向他的兄弟骄傲，在玩耍时骄傲，即是在嘲笑他；\Person{撒辣}说了什么？\BibleQuote{把婢女和她的儿子赶出去，因为婢女的儿子不能和我的儿子\Person{依撒格}一同继承产业。} \BibleRef{创 21:9} 所以，使他被拒绝的不是婢女的子宫，而是奴隶的颈项。因为生来自由者如果是骄傲的，就是奴隶，而且更糟的是，他是一个坏主母的奴隶，即骄傲本身的奴隶。因此，我的弟兄们，回答那个人说，人不能第二次出生；要无畏地回答，人不能第二次出生。任何第二次做的事就是嘲笑，任何第二次做的事就是戏弄。那是\Person{依市玛耳}在戏弄，让他被赶出去。因为\Person{撒辣}观察到他们在戏弄，圣经说，就对\Person{亚巴郎}说：\BibleQuote{把婢女和她的儿子赶出去。} 男孩们的戏弄使\Person{撒辣}不悦。她在他们的戏弄中看到了某种奇怪的东西。难道有儿子的人不喜欢看到他们玩耍吗？她看到了并反对。她在他们的戏弄中看到了某种东西；她看到了其中的嘲笑，觉察到那奴隶的骄傲；她不悦，就赶走了他。婢女的儿女，当邪恶时，被赶出去；自由的妇人的儿女，如果像\Person{厄撒乌}，也被赶出去。所以，任何人都不要因出身于好父母而自恃；任何人都不要因由圣人施洗而自恃。让那由圣人施洗的人仍要谨慎，免得他是\Person{厄撒乌}，而不是\Person{雅各伯}。那么，弟兄们，我要说，宁可让那些寻求自己利益、贪爱世界的人施洗——这就是婢女之名所表示的意思——而自己却属灵地寻求\Person{基督}的继承，好比\Person{雅各伯}由婢女所生的儿子，也比让圣人施洗后变得骄傲，以致成为将被赶出的\Person{厄撒乌}，即使是由自由的妇人所生的，要好得多。你们要紧紧持守这一点，弟兄们。我们不是在奉承你们，不要让你们的任何希望寄托在我们身上；我们既不奉承自己，也不奉承你们；每人要担负自己的重担。我们有责任说话，免得我们不幸被定罪；你们有责任去听，并且用心去听，免得我们所给予的从你们手中被索回；不，当它被索回时，愿它被发现是赚得，而非亏损。\par

5. 主对\Person{尼苛德摩}说，并向他解释：\BibleQuote{我实实在在告诉你：人除非由水和圣神重生，不能进入天主的国。} 祂说，你明白的是肉性的出生，当你说：人岂能再进入母亲的腹中呢？为进入天主的国，出生必须是由水和圣神。如果一个人是为获得一个人类父亲的暂时继承而出生，就让他由肉性的母亲的腹中出生；如果一个人是为获得天主为父的永恒继承而出生，就让他由教会的腹中出生。一个父亲，作为一个将要死的人，通过他的妻子生一个儿子来接替他；但天主从教会生儿子，不是要接替祂，而是要同祂永远在一起。祂接着说：\BibleQuote{那由肉身生的就是肉身；那由圣神生的就是神。} 那么我们是属灵地出生的，并且我们是由圣言和圣事生为神。圣神临在好使我们出生；圣神无形地临在，你由祂所生，因为你也必须无形地出生。因为祂接着说：\BibleQuote{你不要惊奇，因为我对你说：你们必须重生。风随意向哪里吹，你听到它的声音，却不知道它从哪里来，往哪里去。} 没有人看见圣神；那么我们是怎样听见圣神的声音呢？圣咏响起，这是圣神的声音；福音响起，这是圣神的声音；圣言响起，这是圣神的声音。\BibleQuote{你听到它的声音，却不知道它从哪里来，往哪里去。} 但如果你是由圣神所生的，你也要这样，那不由圣神所生的人不知道你从哪里来，往哪里去。因为祂接着说：\BibleQuote{凡由圣神生的，也是这样。}\par

6. \Quote{尼苛德摩回答祂说：\InnerQuote{这事怎能成就呢？}} 事实上，按血肉之见，他并不知道如何成就。在他身上发生主所说的事：他听见圣神的声音，却不知道从何处来，往何处去。\Quote{耶稣回答他说：\InnerQuote{你是以色列的师傅，连这事都不知道吗？}} 啊，弟兄们！难道我们以为主是故意嘲弄这位犹太人的师傅吗？主知道祂在做什么；祂愿意这人生于圣神。没有人能生于圣神，除非他谦卑，因为谦卑本身使我们生于圣神；\Quote{因为主亲近心灵破碎的人。} 这人因他的师傅身份而自高自大，自以为身为犹太人的教师是件了不起的事。耶稣挫败他的骄傲，好使他能生于圣神：祂嘲笑他是个无知的人；并非主要显得比他优越。天主与人相比，真理与谬误相比，有何可比性？基督大于尼苛德摩！这能说吗，可以说吗，可以想吗？若说\Quote{基督大于天使，} 那真是荒谬：因为万有是由祂而造，祂远超一切受造物。但祂仍斥责那人的骄傲：\Quote{你是以色列的师傅，连这事都不知道吗？} 仿佛祂说：看，你一无所知，你是骄傲的首领；生于圣神吧！因为你若生于圣神，就会遵守天主的道，跟随基督的谦逊。的确，祂远高于众天使，以致\Quote{祂虽具有天主的形体，却没有以自己与天主同等为应当把持的，反而空虚自己，取了奴仆的形体，成为人形，形态上完全像人：祂谦卑自己，听命至死}（且免得你认为任何一种死都可以），\Quote{甚至死在十字架上。} \BibleRef{斐 2:6} 祂悬在十字架上，他们讥笑祂。祂本可以从十字架上下来，但祂忍耐，为能从坟墓中复活。主，忍受了骄傲的奴隶；\BibleRef{玛 11:30} 医生忍受了病人。若祂这样做，那些应当生于圣神的人该当如何行呢！——若祂这样做，祂是天上真正的师傅，不单是人的师傅，也是天使的师傅。因为若天使有学问，那是因天主的圣言而有的。若他们因天主的圣言而有学问，请问他们因何而有学问；你便会发现：\Quote{在起初已有圣言，圣言与天主同在，圣言就是天主。} 人的颈项被除去，只有又硬又僵的颈项，好使它温顺地负起基督的轭，关于这轭曾说过：\Quote{我的轭是柔和的，我的担子是轻松的。} \BibleRef{玛 11:30}\par

7. 祂接着说：\Quote{我若对你们说地上的事，你们尚且不信；我若对你们说天上的事，你们怎能信呢？} 弟兄们，祂说了什么地上的事呢？\Quote{人若不重生}——这是地上的事吗？\Quote{圣神随意吹拂，你听见它的声音，却不知道从何处来，往何处去}——这是地上的事吗？因为若祂说这话是指风，某些人就是这样理解的，当被问及主所说的\Quote{地上的事}是什么时，他们便困惑地说，主所说的\Quote{圣神随意吹拂}和\Quote{你听见它的声音，却不知道从何处来，往何处去}，是指风而言。那么祂称什么为地上的呢？祂在谈论属神的诞生，并接着说：\Quote{凡由圣神而生的，就是这样。} 那么，弟兄们，我们中间有谁看不见，例如南风从南往北刮，或另一阵风从东往西吹？那么，我们怎么不知道它从何处来，往何处去呢？那么，祂说了什么地上的事，使人不相信呢？难道是关于重建圣殿的事吗？的确，祂的身体取自大地，祂准备将那取自尘世之躯的大地复活起来。他们不相信祂将要复活大地。\Quote{我若对你们说地上的事，你们尚且不信，}祂说，\Quote{我若对你们说天上的事，你们怎能信呢？} 也就是说，你们若不信我能重建你们所拆毁的圣殿，怎会信人能由圣神重生呢？\par

8. 祂继续说：\Quote{没有人升过天，除了那从天降下、仍在天上的人子。}看哪，祂在地上，也同时在天上；在肉身中在地上，藉祂的\OriginalTerm{天主性}{divinity}在天上；是的，藉祂的天主性无所不在。由母亲所生，却不离开圣父。我们理解基督有两种\OriginalTerm{诞生}{nativities}：一种是神性的，另一种是人性的：一种是使我们得以存在的，另一种是使我们得以再造的：两者都奇妙；前者无母，后者无父。但因为祂从\Person{亚当}取了身体——因为\Person{玛利亚}出自亚当——并要复活那同一个身体，所以祂说\Quote{拆毁这座圣殿，三天内我要把它重建起来}时，所说的是属世的事。但祂说\Quote{人除非由水和圣神重生，不能看见天主的国}时，所说的却是属天的事。那么，兄弟们！天主愿意成为人子，也愿意人成为天主子。祂为我们而降下，让我们为祂而上升。因为只有祂降下和上升，就是那位说\Quote{没有人升过天，除了那从天降下的一位}的。那么，难道祂所使之成为天主子的人不会升天吗？当然会：这是赐给我们的应许：\Quote{他们要像天主的使者一样。}\BibleRef{玛 22:30}那么，为什么除了那降下的一位，没有人升天呢？因为只有一位降下，只有一位上升。其余的人呢？我们该如何理解？岂不是他们要成为祂的肢体，使那一位得以上升？因此，接着说\Quote{没有人升过天，除了那从天降下、仍在天上的人子。}你们惊奇祂既在地上也在天上吗？祂也使祂的门徒如此。听宗徒\Person{保禄}说：\Quote{但我们的家乡是在天上。}\BibleRef{斐 3:20}如果宗徒保禄，一个人，在肉身中行走在地上，却拥有在天上的家乡，难道那天与地的主宰不能同时在天上和地上吗？\par

9. 因此，如果除了祂没有人降下和上升，那么其他人还有什么希望呢？其他人的希望在于：祂降下，是为了使那些通过祂上升的人在祂内并与祂一起成为一体。\Quote{祂并不是说\InnerQuote{及给众子孙}，}宗徒说，\Quote{如同指许多人一样；而是指一个，\InnerQuote{及给你的子孙}，就是基督。}祂又对信徒说：\Quote{你们属于基督；如果属于基督，就是\Person{亚巴郎}的子孙。}祂所说的一体，也就是我们众人。因此，在圣咏中，有时许多人歌唱，表明从许多人中形成一人；有时一人歌唱，表明从许多人中形成的一人。因此，在那池塘中，只有一个人被治愈；其他下去的人都没有被治愈。这一个人表明了教会的合一。那些憎恨合一、在人间自结派别的人有祸了！让他们听那愿意使他们在一人内、借着一人、为一人而成为一体者的话：让他们听那说\Quote{你们不要制造许多}者的话：\Quote{我栽种，\Person{阿颇罗}浇灌，但使之生长的是天主。所以，栽种的不算什么，浇灌的也不算什么，只在那使之生长的天主。}\BibleRef{格前 3:6}他们曾说：\Quote{我属于保禄，我属于阿颇罗，我属于\Person{刻法}。}祂却说：\Quote{基督被分裂了吗？}要成为一体，成为一件事，成为一个人：\Quote{没有人升过天，除了那从天降下的一位。}看！他们曾对保禄说：我们愿属于你。保禄对他们说：我不愿你们属于保禄，而要属于那位连同你们也属于祂的一位。\par

10. 因为祂降下并死了，藉着那死亡将我们从死亡中解救出来：被死亡所杀，祂杀死了死亡。兄弟们，你们知道，这死亡是因魔鬼的嫉妒而进入世界的。\Quote{天主并没有创造死亡，}经上说，\Quote{也不喜欢活人灭亡；但祂创造了万物，使它们存在。}但这经文说些什么呢？\Quote{但因魔鬼的嫉妒，死亡进入了全世界。}\BibleRef{智 1:2}对于魔鬼为诱惑我们而提供的死亡，人并非被迫接受；因为魔鬼没有强迫的力量，只有诱惑的狡猾。如果你没有同意，魔鬼就一无所获：你自己的同意，人啊，引你走向死亡。从必死的人生出必死的人；从永恒者我们变成了必死者。从\Person{亚当}而来，众人都是必死的；但\Person{耶稣}、天主子、天主的\OriginalTerm{圣言}{Word}、万物藉祂而受造、与圣父同等的独生子，却成为了必死的：\Quote{因为圣言成了血肉，居住在我們中間。}\par

11. 祂忍受了死亡；但死亡被祂挂在十字架上，而必死的人从死亡中被解救出来。主唤起一件大事，这大事曾以预像在古时的子民中成就：祂说：\Quote{正如梅瑟在旷野里举起了蛇，人子也应照样被举起来，使凡信祂的人不致丧亡，反而获得永生。}这是一个伟大的奥迹，凡读过的人都知道。但愿那些没有读过的人，以及那些已经忘记也许曾听到或读过的人，也一同倾听。以色列子民在旷野中被蛇咬伤，无助地倒下；他们遭受了巨大的灾祸，许多人死亡：因为这是天主打发来责罚和鞭打他们，为要教训他们。在这事中显示了一个伟大的奥迹：未来之事的预像；主亲自在这段经文中作证，因此没有人能给出与真理关于自身的启示不同的解释。梅瑟奉主的命令制造了一条铜蛇，把它挂在旷野的杆子上，并告诫以色列子民：凡被蛇咬的，一望那挂在杆子上的铜蛇，就必得痊愈。这事就这样行了：人受了伤；他们一望，就痊愈了。\BibleRef{户 21:6} 咬人的蛇是什么？是罪，来自肉身的必死性。被高举的蛇是什么？是主在十字架上的死亡。因为死亡由蛇而来，所以用蛇的形象来预表。蛇的咬伤是致死的，主的死亡是赋予生命的。人凝视蛇，好使蛇不再有能力。这是什么意思？人凝视死亡，好使死亡不再有能力。但这是谁的死亡？生命的死亡；若可以这样说，生命的死亡；是的，因为可以这样说，但这样说令人惊奇。然而，既然这是要做成的事，难道不该说出来吗？难道我该迟疑于说出主屈尊为我所行的事吗？基督不就是生命吗？然而基督却被挂在十字架上。基督不就是生命吗？然而基督却死了。但在基督的死亡中，死亡死了。生命死了，使死亡死亡；生命的圆满吞没了死亡；死亡被吸收在基督的身体里。同样，在复活时，当我们凯旋，我们将歌唱：\Quote{死亡啊，你的胜利在哪里？死亡啊，你的刺在哪里？}\BibleRef{格前 15:54} 同时，弟兄们，为使我们从罪中得到治愈，让我们现在凝视被钉的基督；因为祂说：\Quote{正如梅瑟在旷野里举起了蛇，人子也应照样被举起来，使凡信祂的人不致丧亡，反而获得永生。}正如那些仰望那铜蛇的人没有被蛇咬伤致死，同样，那些以信德仰望基督死亡的人，就从罪的咬伤中得到治愈。但那些人是从死亡得到治愈，获得暂时生命；而这里祂说：\Quote{使他们获得永生。}至于预像和实体之间的区别：预像带来暂时生命；实体（预像所预表的）带来永生。\par

12. \Quote{因为天主没有派遣子到世界上来审判世界，而是为叫世界借着祂而得救。}那么，就医生而言，祂来是为医治病人。凡不遵守医生嘱咐的人，是自己毁灭自己。祂来是为拯救世界：祂被称为世界的救主，岂不正是因为祂来是为拯救世界，而不是为审判世界吗？你不愿由祂而得救；你就必被自己审判。为什么我说\Quote{必被审判}呢？请看祂所说的话：\Quote{信祂的人，不被审判；但不信的人。}你以为祂接下去会说\Quote{已被审判}吗？祂说：\Quote{已经受了审判。}审判尚未显现，但已经发生了。因为主认识那些属于祂的人：祂知道谁是坚持到底为得冠冕，谁是为得火焰；祂认识祂打谷场上的麦子，也认识糠秕；认识好谷粒，也认识莠子。不信的人已经受了审判。为什么受了审判？\Quote{因为他没有信从天主独生子的名字。}\par

13. \BibleQuote{审判就在于此：光明来到了世界，世人却爱黑暗甚于光明，因为他们的行为是邪恶的。} 我的弟兄们，主要找到谁的行为是善的呢？没有人的行为是善的：祂发现所有人的行为都是邪恶的。那么，怎么有些人行了真理，并来到了光中呢？因为接下来祂说：\BibleQuote{履行真理的，却来就光明，为显示出他的行为是在天主内完成的。} 怎样有些人行了善工而来到光中，即来到基督面前？又怎样有些人喜爱了黑暗？因为如果祂发现所有人都是罪人，并治愈所有人的罪，那条主之死被预表的铜蛇治愈了被咬的人，并且因蛇的咬伤而竖起了蛇，即为了必死之人的缘故主之死，祂发现他们是不义的；那么，我们如何理解\BibleQuote{审判就在于此：光明来到了世界，世人却爱黑暗甚于光明，因为他们的行为是邪恶的。}？这是怎么回事？事实上，谁的行为是善的呢？祢不是来使不敬虔的人成义的吗？\Quote{但是他们喜爱了} 祂说，\Quote{黑暗甚于光明。} 祂在那里强调了重点：因为许多人喜爱自己的罪；许多人承认自己的罪；而承认自己罪、控告自己罪的人，现在是与天主一同工作。天主控告你的罪：如果你也控告自己，你就与天主联合。仿佛有两样事物：人和罪人。你被称为人，是天主的作为；你被称为罪人，是人自己的作为。抹去你所做的，好让天主拯救祂所做的。你应该痛恨在你内你自己的作为，而爱在你内天主的作为。当你的行为开始令你不悦时，从那时起你的善工就开始了，因为你指责自己的邪恶行为。承认邪恶行为是善工的开端。你行真理，并来到光中。你怎样行真理呢？你不抚慰、不哄骗、不奉承自己，也不说\Quote{我是义人}，而你是恶人：这样，你开始行真理。你来到光中，好让你的行为显明是在天主内完成的；因为你的罪，正是那令你不悦的事物，如果不是天主光照你，祂的真理向你显示，就不会令你不悦。但是那喜爱自己罪的人，即使在受劝诫之后，仍憎恨那劝诫他的光，并逃避光，以免他所喜爱的行为被证明是恶的。但行真理的人在自身内控告自己的邪恶行为，不宽容自己，不放纵自己，好让天主宽恕他：因为他希望天主宽恕的那个罪，他自己承认了，并且他来到光中；他感谢光向他显示了他在自身内应当憎恨什么。他向天主说：\BibleQuote{求祢转过祢的脸，不要看我的罪过。} 然而他怎能这样说，除非他加上\BibleQuote{因为我承认我的过犯，我的罪恶常在我的眼前？} 愿那你不愿在天主面前的事物，常在你面前。但若你将罪置于身后，天主就会将它重新放到你眼前；祂将在不再有悔改果实的时候这样做。\par

14. 奔跑吧，我的弟兄们，免得黑暗抓住你们。醒来吧，为你们的救恩，趁还有时间醒来吧；不要让任何人被拦在天主圣殿之外，不要让任何人被拦在主的工作之外，不要让任何人被从不断的祈祷中叫走，不要让任何人被夺去惯常的虔敬。醒来吧，趁还是白昼：白昼照耀，基督就是白昼。祂乐意赦免罪，但只赦免那些承认罪的人；祂准备惩罚那些自我辩护、自夸为义人、自以为是什么其实是虚无的人。但那行在祂的爱与慈悲中的人，即使远离那些重大和致命的罪行，如谋杀、偷盗、奸淫，但仍因那些似乎小罪——言语上、心思上、对许可之事的不节制——他在承认罪中行真理，并在善工中来到光中：因为许多小罪，若被忽视，足以致命。微小的是汇成江河的水滴；微小的是沙粒；但若许多沙粒聚在一起，堆积起来就会压碎人。舱底被忽视的积水与奔腾的巨浪造成同样的后果。它从舱底慢慢渗入；长期渗入而不抽出去，就会使船沉没。那么，这抽出去是什么呢？就是通过善工、叹息、禁食、施舍、宽恕，以便罪不致淹没我们。然而，此生的道路是艰难的，充满了诱惑：在顺境中，不要让它抬高我们；在逆境中，不要让它压垮我们。那赐予今世幸福的天主，赐予它是为安慰你，而非为毁灭你。同样，那在此生责打你的天主，是为你的改善，而非为你的定罪。忍受那为训练而纠正你的父亲吧，免得你感到那在惩罚中审判你的法官。这些事我们每天都告诉你们，而且必须常常说，因为它们是美好而有益的。\par

\SourceRecord{Translated by John Gibb.；Nicene and Post-Nicene Fathers, First Series；Vol. 7.；Edited by Philip Schaff.；Buffalo, NY: Christian Literature Publishing Co.,；1888.；<http://www.newadvent.org/fathers/1701012.htm>.}

% structure: p0001 source=h1 target=section reason=page-title

\section{第13篇（\WorkTitle{若望福音} 3:22-29）}

1. 按\WorkTitle{若望福音}的诵读次序——你们关心自身进步的诸位或可记得——是依循定序进行的，因此今天我们所读的这段经文，正待我们讲解。你们应记得，在前几篇讲道中，自\WorkTitle{若望福音}的开端直至今日的课目，我们均已讲解。虽然你们或许已遗忘许多，但至少应记得我们已完成了分内的工作。你们从中所听到关于\Person{若翰}洗礼的内容，即使未能全部记住，但我相信你们已听到足以存留于心的话。此外，也曾论及为何圣神以鸽子的形状显现，以及那个极难解的疑问是如何得解的，即主身上有何事是\Person{若翰}所不知晓，而藉鸽子才晓得的——当时\Person{若翰}已认识主，因为\Person{耶稣}前来受洗时，他对祂说：\BibleQuote{我当受祢的洗，祢反而来就我吗？}主回答他说：\BibleQuote{你暂且允许，以成全一切义德。}\BibleRef{玛 3:14}\par

2. 如此，我们的诵读次序迫使我们回到这同一位\Person{若翰}。他就是\Person{依撒意亚}所预言的那一位：\BibleQuote{有声音在旷野呼喊：预备主的道路，修直祂的途径。}\BibleRef{依 40:3} 他为他的主（因主认为他配得）和他的朋友作了这样的见证。主，甚至他的朋友，也亲自为\Person{若翰}作证。因为论及\Person{若翰}，祂说：\BibleQuote{妇人所生的，没有兴起一个比洗者若翰更大的。}但祂将自己置于\Person{若翰}之上时，祂之所更伟大之处，在于祂是天主。祂说：\BibleQuote{然而在天国里最小的，也比他大。}\BibleRef{玛 11:11} 在年龄上较小，但在权能、神性、尊威、光耀上更大——正如\BibleQuote{在起初已有圣言，圣言与天主同在，圣言就是天主。}在前文里，\Person{若翰}曾为主作证，他确实称祂为天主子，但未说祂是天主，也未否认；他缄默了祂的天主性，未否认祂是天主；然而他并非完全缄默了祂的天主性，或许我们会在今天的课文中发现这一点。他曾称祂为天主子；但人也曾被称作天主子。他宣称祂如此崇高，以至于他自己不配解祂的鞋带。这种伟大性使我们领悟许多：他不配解祂鞋带的那一位，就是那一位，在妇人所生的当中没有兴起比\Person{若翰}更大的。他实在比所有人及天使都更大。因为我们发现有一位天使禁止一个人俯伏在他脚前。例如，在\WorkTitle{默示录}中，一位天使向此福音的作者\Person{若望}显示某些事时，\Person{若望}因异象的伟大而惊惧，俯伏在天使脚前。但天使说：\BibleQuote{起来，不可如此！你应朝拜天主，因为我是你的同仆，也是你弟兄们的同仆。}\BibleRef{默 22:8} 那么，一位天使尚且禁止人俯伏在他脚前，难道不明显表明，那位连一个在妇人所生中没有比他更大的人都自认不配解其鞋带的那一位，必超越所有天使吗？\par

3. 然而，\Person{若望}可以更明确地说出，我们的主耶稣基督就是天主。我们在本段中可能会发现，我们一直在歌颂的也许就是祂：\BibleQuote{上主统治了大地}；那些想象祂只在\Place{非洲}统治的人对此充耳不闻。但他们不要以为，当说到\BibleQuote{因为天主是普世的君王}时，就不是指着基督说的。因为除了我们的主耶稣基督，谁还是我们的君王呢？祂就是我们的君王。你们在刚才歌唱的那篇圣咏的同一节中听到了什么？\BibleQuote{请歌颂我们的天主，请歌颂；请歌颂我们的君王，请歌颂。} 他称之为天主的，也称之为我们的君王：\BibleQuote{请歌颂我们的天主，请歌颂；请歌颂我们的君王，请歌颂，你们要以明悟歌颂。} 为的是叫你们不要误解你们所歌颂的那位只在一个地方统治，他说：\BibleQuote{因为天主是普世的君王。} 祂既然出现在大地的一隅，在\Place{耶路撒冷}、在\Place{犹太}，行走在人间，诞生、吸吮母乳、成长、饮食、醒来、睡眠、坐在井旁、疲惫；被捉拿、受鞭打、被吐唾沫、戴荆棘冠、挂在木头上、被枪刺破、死去、埋葬，又怎能是普世的君王呢？当时在地方上被看见的是血肉，只有血肉对血肉之眼是可见的；不朽的威严隐藏在可朽的血肉之中。我们若穿透血肉的结构，以何种眼睛才能看见那不朽的威严呢？有另一只眼，有一只内在的眼。例如，\Person{多俾亚}并非没有眼睛，当他的肉眼失明时，他仍给儿子传授生活的规诫。\EditorialNote{多俾亚传第四章} 儿子握住父亲的手，好让父亲用脚行走，而父亲却给儿子教导行义路的忠告。这里我看见眼睛，那里我理解眼睛。传授生活规诫者的眼睛，比那握手的眼睛更好。耶稣向\Person{斐理}说话时，也要求这样的眼睛：\BibleQuote{这么长久我与你们在一起，你们还不认识我吗？} 祂要求这样的眼睛时说：\BibleQuote{斐理，看见我的，就是看见了父。} 这些是明悟之眼，这些是心神之眼。正因如此，圣咏说到\BibleQuote{因为天主是普世的君王}时，立刻加上\BibleQuote{你们要以明悟歌颂。} 因为当我说\BibleQuote{请歌颂我们的天主}时，我是说天主是我们的君王。但是，我们的君王你们曾在人中看见，作为人；你们看见祂受苦、被钉、死亡：在那血肉中隐藏着某种东西，是你们可能已用血肉之眼看见的。那里隐藏着什么？\BibleQuote{你们要以明悟歌颂。} 不要用肉眼寻求那由心神所看见的。用舌头\BibleQuote{歌颂}吧，因为祂在你们中间作为血肉；但是因为\BibleQuote{圣言成了血肉，居住在我们中间}，所以将声音还给血肉，将心神的凝视归于天主。\BibleQuote{你们要以明悟歌颂}，这样你们就会看见\BibleQuote{圣言成了血肉，居住在我们中间}。\par

4. 现在，让\Person{若望}也作证：\BibleQuote{此后，耶稣和门徒到了\Place{犹太}地方，在那里同他们住下，并施洗。} 祂受了圣洗，但祂付洗。祂受洗所用的圣洗，并非祂所用的圣洗。主由仆人受洗，却施予圣洗，显示谦逊之路，并引导人归于主的圣洗，即祂自己的圣洗，借着谦逊的榜样，连祂自己也不拒绝从仆人受洗。而且在仆人的圣洗中，为主预备了道路；主自己也受洗，为那些来到祂面前的人，将自己作为道路。让我们听祂亲自说：\BibleQuote{我是道路、真理、生命。} 如果你寻求真理，就持守道路，因为道路与真理是同一的。你所行走的道路，与你要前往的\Emph{终点}是同一的：你不是沿着一条路走到另一个目标；你不是借着别的东西作为道路而来到基督那里，你乃是借着基督而来到基督那里。如何借着基督到基督？借着人而基督，到天主而基督；借着成了血肉的圣言，到那在起初就与天主同在的圣言；从人所吃的，到天使每日所吃的。因为经上记载：\BibleQuote{他给他们降下天粮，人吃了天使的粮食。} 什么是天使的粮食？\BibleQuote{在起初已有圣言，圣言与天主同在，圣言就是天主。} 人如何吃了天使的粮食？\BibleQuote{圣言成了血肉，居住在我们中间。}\par

5. 但尽管我们说天使进食，弟兄们，不要以为这是用牙齿进行的。因为如果你们这样想，那么天使所食的天主就好像被撕碎了。谁撕碎义德呢？然而，仍有人问我：谁能食义德呢？那么，经上怎么说的？\BibleQuote{饥渴慕义的人是有福的，因为他们要得饱饫。}你们按肉体所吃的食物，为了滋养你们而朽坏；它被消耗以修复你们的损耗。食义德吧：当你们被滋养时，它仍完整无损。正如用这肉眼观看有形光体，我们的眼睛得到滋养，然而这仍是肉眼所见的有形之物。许多人长期处在黑暗中，他们的视力因禁光（如同禁食）而减弱。眼睛被剥夺了它们的食物（因为它们以光为食），因禁食而疲倦、衰弱，以致不能忍受那原本滋养它们的光；如果光长时间缺席，它们就熄灭，视觉本身也仿佛在它们内死去。那么，光会因每天被那么多眼睛所食而减少吗？你们的眼睛得到滋养，光却保持完整。既然天主能在有形光体上向肉眼显示这事，难道祂不能向洁净的心灵显示那另一道永不疲倦、保持完整、毫不减损的光吗？哪道光？\BibleQuote{在起初已有圣言，圣言与天主同在。}我们看看这是否是光。\BibleQuote{因为在祢那里有生命的泉源，在祢的光明中我们才能看到光明。}在地上，泉源是一回事，光是另一回事。当你口渴时，你寻找泉源；为了到达泉源，你寻找光；如果不是白天，你就点灯以到达泉源。那泉源本身就是光：对于口渴者是泉源，对于盲者是光。让眼睛开启以看见光，让心灵的口唇开启以饮用泉源；你所饮的，你所见的，你所听的，天主成为你的一切；因为祂就是你所爱的这一切。如果你着眼于有形之物，则天主不是饼，不是水，不是这光，不是衣服，也不是房屋。因为这一切都是可见的、各自分离的事物。饼是什么，水不是；衣服是什么，房屋不是；这些东西是什么，天主不是，因为它们是可见之物。天主对你是这一切：如果你饥饿，祂是你的饼；如果你口渴，祂是你的水；如果你在黑暗中，祂是你的光；因为祂保持不朽。如果你赤身露体，祂是你的不朽外衣，当这可朽坏的穿上了不可朽坏的，这可死的穿上了不死的时候。关于天主，一切都能说，但没有任何话是相称地说的。没有什么比这种表达的贫乏更广泛的了。你为祂寻找一个合适的名字，你找不到；你想无论如何谈论祂，你却看到祂是一切。羔羊和狮子有何相似？两者都论及基督。\BibleQuote{请看天主的羔羊！}如何是狮子？\BibleQuote{犹大支派中的狮子已得胜了。}\BibleRef{默 5:5}\par

6. 让我们听听若翰：\BibleQuote{耶稣施洗。}我们说耶稣施洗。耶稣如何？主如何？天主子如何？圣言如何？然而，圣言成了血肉。\BibleQuote{若翰也在哀嫩靠近撒林施洗。}有一个湖，叫\Quote{哀嫩}。我们怎么知道那是一个湖？\BibleQuote{因为那里水多，众人就都来受洗。那时若翰还没有被下在监里。}如果你们记得（瞧，我再说一次），我曾告诉过你们为什么若翰施洗：因为主必须受洗。为什么主必须受洗？因为将来有许多人会藐视圣洗，为了显得他们拥有比他们所见的其他信徒更大的恩宠。例如，一个望教者，现在过着节制的生活，可能会轻视一个已婚者，并自认比那个信徒更好。那个望教者可能会在心里说：\Quote{我需要受洗做什么，仅仅得到别人已有的东西，而我已比他更好？}因此，恐怕骄傲的颈项会使某些因自己义德的功劳而十分得意的人坠入毁灭，主愿意由仆人受洗，仿佛在对祂的杰出子民说：\Quote{你们为什么自夸？为什么因有人有智慧、有人有学识、有人有贞洁、有人有忍耐的勇气而高举自己？你们能拥有与我——赐给你们这些的那一位——一样多吗？然而我却由仆人受了洗，你们却不屑于由主受洗。}这就是\BibleQuote{以成全一切义德}的意义。\par

7. 但有人会说：\Quote{那么，若翰只给主施洗就够了；何必还需要若翰给其他人施洗呢？}现在我们也说过这一点：如果若翰只给主施洗，人们不会没有这种想法：若翰的圣洗比主的圣洗更好。事实上，他们会说：\Quote{若翰的圣洗如此伟大，以至于只有基督才配受此洗。}因此，为了表明主将要赐予的圣洗比若翰的圣洗更好——使人明白一个是仆人的圣洗，另一个是主的圣洗——主受洗是为立谦卑的榜样；但祂并非唯一受若翰施洗的人，以免若翰的圣洗看似比主的圣洗更好。为此，我们的主耶稣基督指明了道路，正如你们所听到的，弟兄们，免得有人自诩拥有某方面丰厚的恩宠而藐视领受主的圣洗。因为无论望教者进步多大，他仍背负着自己的罪债；除非他来到圣洗，否则罪债不得赦免。正如以色列子民直到来到红海才摆脱埃及人，同样，无人能摆脱罪恶的重压，直到来到圣洗之池。\par

8. \Quote{那时，若翰的门徒与犹太人关于取洁发生了争论。}若翰施洗，基督也施洗。若翰的门徒被激动了，人们纷纷追随基督，也有人来到若翰那里。来到若翰那里的人，被他打发到耶稣那里受洗；但那些由基督施洗的人，却未被打发到若翰那里。若翰的门徒惊惶起来，便开始与犹太人争辩，如同常发生的那样。要知道犹太人宣称基督是更大的，人们应当投奔祂的洗礼。若翰的门徒尚未明白此事，便为若翰的洗礼辩护。他们来到若翰本人面前，好让他解决这问题。亲爱的各位，要明白。在此我们被教导要看见谦逊的益处，当人们争论有误时，也显明若翰是否愿意荣耀自己。现在他大概说了：\Quote{你们说的是实话，你们争辩得对；我的洗礼是更好的，我亲自给基督施了洗。}在基督受洗之后，若翰本可以这样说。如果他想要高举自己，这是多么好的机会！但他更知道在谁面前谦卑自己：他向那位他知道出生在自己之后的主，甘愿通过承认祂而让位。他明白自己的救恩在于基督。他刚才已经说过：\Quote{我们都从祂的丰满中领受了；}这就是承认祂是天主。因为如果祂不是天主，众人怎能从祂的丰满中领受？因为如果祂是人，而不是天主，那么祂自己也从天的丰满中领受，所以祂就不是天主。但如果所有的人都从祂的丰满中领受，祂就是泉源，他们是饮水者。从泉源饮水的人，既渴又饮。泉源从不干渴，它从不需要自己。人需要泉源。他们以干渴的胃和焦干的嘴唇跑到泉源那里以获得滋润。泉源流动以解渴，主耶稣也是如此。\par

9. 那么，让我们看若翰如何回答：\Quote{他们来到若翰那里，对他说：辣彼，从前同你在约旦河对岸、你所见证的那位，请看，祂也在施洗，众人都到祂那里去了。}这就是说，你说什么？难道不该阻止他们，好让他们更多地到你这里来？\Quote{他回答说：人不能领受什么，除非自天赐予他。}你们认为若翰说这话是指谁？是指他自己。\Quote{作为一个人，我领受了，}他说，\Quote{来自天上。}亲爱的各位，要注意：\Quote{人不能领受什么，除非自天赐予他。你们自己为我作证，我曾说过：我不是基督。}这等于说：\Quote{你们为什么欺骗自己？看你们是怎样把这个问题放在我面前的。你们对我说了什么？\InnerQuote{辣彼，从前同你在约旦河对岸、你所见证的那位。}那么你们知道我为祂作了怎样的见证。难道我现在要说祂不是我所宣告的那位吗？因为我从天上领受了一些，为了成为什么，你们难道希望我空虚，以至于违背真理？\InnerQuote{人不能领受什么，除非自天赐予他。你们自己为我作证，我曾说过：我不是基督。}}不是基督；但如果你比祂大呢，既然你给祂施了洗？\Quote{我是被差遣的：}我是先驱，祂是审判者。\par

10. 但要听一个更强而有力、更有表现力的见证。你们看我们正在讨论什么；你们看，爱任何人代替基督就是犯奸淫。我为什么这样说？让我们听若翰的声音。人们可能误会他，以为他是他所不是的那位。他拒绝虚假的荣誉，以保存完整的真理。看他宣称基督是什么；他自己又自称是什么？\Quote{拥有新娘的是新郎。}要贞洁，爱新郎。但你是谁，竟对我们说：\Quote{拥有新娘的是新郎？但新郎的朋友，站着听祂，因新郎的声音而大大喜乐。}主我们的天主会按着我心中的动荡帮助我，因为充满悲伤，来抒发我所感受的痛苦；但我以基督亲自恳求你们，在思想中想象我无法说出的事情；因为我知道我的悲伤不能以合适的庄严表达出来。现在我看到许多犯奸淫者想要得到那位新娘，她被如此大的代价赎买，在丑陋中被爱，好使她变得美丽，被这样一位赎买、解救并装饰；那些犯奸淫者用他们的言语努力使自己被爱，而不是新郎。关于那一位，经上记载：\Quote{这就是那施洗的。}\BibleRef{若 1:33}谁从我们中间出去说：\Quote{我就是那施洗的}？谁从我们中间出去说：\Quote{我所给的是圣洁的}？谁从这里出去说：\Quote{你们由我而生是有益的}？让我们听新郎的朋友，而不是反对新郎的犯奸淫者；让我们听一个怀着妒意的人，但不是为了自己。\par

11. 弟兄们，请在心内回到自己的家中。我讲论肉性的事，讲论属世的事；我按人的方式说话，是为了你们肉身的软弱。你们中许多人有过妻子，许多人想要有妻子，许多人虽然不想要，却仍有过妻子；许多人根本不想要妻子，却是从你们父亲的妻子所生。这是一种触动每个人心的情感。在人间事务中，没有人对人群如此疏离，以至于感受不到我所说的。比方说，有一个人出门远行，把他的新娘托付给朋友照顾：\Quote{请你注意，你是我的好朋友；要当心，免得我不在时，另有别人或许代替我被爱。} 那么，一个人作为他朋友的新娘或妻子的监护人，确实努力不让别人被爱，但如果他自己希望被爱来代替他的朋友，并渴望享受那个托付给他照顾的人，他在全人类面前会是多么可憎！让他看见她凭窗凝望，或与某人谈话时有些过于轻率，他便像嫉妒的人一样禁止她。我看到他嫉妒，但让我看看他为谁嫉妒：是为他不在的朋友，还是为他现在的自己。要想到我们的主\Person{耶稣基督}就是这样做的。祂将新娘托付给祂的朋友照顾；祂动身往远方去，为要得一个王国，正如祂自己在福音中说的，\BibleRef{路 19:12}，但祂仍以祂的尊威临在。让那远渡海洋的朋友被欺骗吧；如果他受骗，那欺骗他的人有祸了！为什么人企图欺骗天主——天主察看所有人的心，鉴察所有人的秘密？但有些异端者出现，说：\Quote{是我施予，是我圣化，是我使人成义；不要到别的派别那里去。} 他确实在嫉妒，但要看为谁。\Quote{不要到偶像那里去，}他说——他嫉妒得对；\Quote{也不要到占卜者那里去，}——仍然嫉妒得对。让我们看看他为谁嫉妒：\Quote{我所施予的是神圣的，因为是我施予；我施洗的人就受了洗；我不施洗的人就没有受洗。} 听那新郎的朋友，学会为你的朋友嫉妒；听那\Quote{施洗者}的声音。为什么想要攫取不属于你的东西？那留下祂新娘的人难道就真的不在了吗？你们难道不知道，那从死者中复活的人正坐在圣父的右边吗？如果犹太人鄙视悬挂在木上的祂，你们难道鄙视坐在天上的祂吗？亲爱的，请确信，我为这件事忍受极大的悲痛；但是，正如我已经说过的，我把其余的交由你们去想。即使我说一整天，也无法表达。如果我哀哭一整天，也不够。如果我像先知所说的，有\BibleQuote{眼泪的泉源}，即使我变成眼泪，变成完全的眼泪，变成舌头，变成完全的舌头，也不够。\par

12. 让我们回来看看这位若翰说了什么：\BibleQuote{拥有新娘的是新郎；}她不是我的新娘。难道你不为这婚姻而欢喜吗？是，他说，我确实欢喜：\BibleQuote{但是新郎的朋友，站在那里并听到他，因新郎的声音而大大喜乐。} 不是因为我的声音，他说，我欢喜，而是因为新郎的声音。我处在听者的位置；祂是发言者；我是必须被光照的人，祂是光；我如同耳朵，祂是圣言。因此，新郎的朋友站着听祂。为什么站着？因为他不会跌倒。怎么会不跌倒？因为他谦卑。看他站在坚实的地上；\BibleQuote{我不配解祂的鞋带。} 你谦卑得好；理所当然你不跌倒；理所当然你站着，听祂，因新郎的声音而大大喜乐。同样，宗徒也是新郎的朋友；他也嫉妒，但不是为自己，而是为新郎。听他嫉妒时的话：\BibleQuote{我为你们嫉妒，}他说，\BibleQuote{以天主的嫉妒：}不是出于我自己，也不是为我自己，而是以天主的嫉妒。为什么？怎样？你嫉妒谁，为谁？\BibleQuote{因为我曾把你们许配给一个丈夫，把你们作为贞洁的童女献给基督。} 那么你为什么害怕？为什么嫉妒？\BibleQuote{我害怕，}他说，\BibleQuote{免得你们的心，正如蛇用诡计欺骗了厄娃一样，从基督内的贞洁中受到败坏。} \BibleRef{格后 11:2} 整个教会被称为童女。你看见教会的肢体是多样的，他们蒙受并因各样的恩赐而欢喜：有些男人已婚，有些女人已婚；有些是鳏夫，不再想娶妻，有些是寡妇，不再想嫁人；有些男人从小保持贞洁，有些女人向天主誓守童贞：恩赐是多样的，但所有这些是一个童女。这童贞在哪里？因为它不在身体内。它属于少数女人；如果童贞也可以指男人，在教会中少数男人也拥有身体的神圣完整；然而，这样的肢体是更尊贵的。但是，其他肢体不但在身体上，而且都在心灵上保持童贞。心灵的童贞是什么？完整的信德、坚定的望德、真诚的爱德。这就是那位为新郎嫉妒的人所害怕会被蛇败坏的童贞。因为，正如身体的某个部分受损一样，舌头的诱惑玷污了心的童贞。让那有正当理由想要保持身体童贞的人，当心不要在心智上受到败坏。\par

13. 弟兄们，我还能说什么呢？连异端者也有守贞者，异端者中有许多守贞者。让我们看看他们是否爱新郎，以致能保全这贞洁。这贞洁是为谁保守的呢？\Quote{为基督。} 让我们看看它究竟是为基督，还是为多纳图；让我们看看这贞洁是为谁保存的：你能轻易证明。看，我给你们指出新郎，因为祂显明了自己。\Person{若翰}为祂作证：\BibleQuote{这就是那施洗的。} 啊，你这守贞者，如果为了这个新郎你保存你的贞洁，为什么你跑到那说\Quote{我就是施洗的}那里去，而新郎的朋友却告诉你\BibleQuote{这就是那施洗的}呢？再者，你的新郎拥有整个世界；那么，你为何要因与世界的一部分接触而受玷污呢？新郎是谁？\BibleQuote{因为天主是全地的君王。} 你的新郎拥有全部，因为祂赎买了全部。看看祂以何代价赎买了它，好让你明白祂赎买了什么。祂付了什么代价？祂倾流了祂的血。祂在哪里倾流，在哪里流出祂的血？在祂的苦难中。难道你不就是在向你的新郎歌唱，或假装歌唱，当整个世界被赎买时：\BibleQuote{他们刺穿了我的手和我的脚，他们数点了我所有的骨头；他们却注视着我，他们观看着我，他们分了我的外衣，并为我的长衣拈阄}？你是新娘，要认出你新郎的长衣。他们在哪件长衣上拈阄？请教福音；看看你是许配给了谁，看看你从哪里接受信物。请教福音；看看它在主的苦难中告诉了你什么。\BibleQuote{那里有一件长衣}，让我们看看是怎样的；\BibleQuote{从上面织成的一片。} 这从上面织成的长衣象征什么？不就是爱德吗？这件长衣象征什么？不就是合一吗？想一想这件长衣，连基督的迫害者都没有把它分开。因为经上说：\BibleQuote{他们彼此说，我们不要撕破它，让我们为它拈阄。} 看，这就是圣咏所说的！基督的迫害者没有撕裂祂的衣裳；而基督徒却分裂了教会。\par

14. 但是，弟兄们，我该说什么呢？让我们清楚地看看祂赎买了什么。因为祂在那里付了赎价，就在祂付价之处。祂付了多大的代价？如果祂只为非洲付了代价，那么让我们成为多纳图派，而不被称为多纳图派，而是基督徒；因为基督只赎买了非洲：尽管即使在这里也有非多纳图派的人。但祂并没有沉默祂在交易中赎买了什么。祂已经结算了账目：感谢天主，祂没有欺骗我们。那位新娘需要听，然后才明白她向谁许下了贞洁。在那篇圣咏中，它说：\BibleQuote{他们刺穿了我的手和我的脚，他们数点了我所有的骨头；} 其中最清楚地宣告了主的苦难——那篇圣咏每年最后一周，在全体人民听中，在基督苦难临近时宣读；这篇圣咏既在他们中间也在我们中间宣读——在那里，我说，注意，弟兄们，祂赎买了什么：让货物账单被宣读：你们听祂赎买了什么：\BibleQuote{普世各邦都要记起，并归向上主；万民各族都要在祂面前敬拜：因为国度属于祂，祂要治理万民。} 看哪，这就是祂所赎买的！看哪！\BibleQuote{因为天主是全地的君王，} 就是你的新郎。那么，你为什么想要让如此富有的一位沦为乞丐呢？承认祂：祂赎买了全部；而你却说：\Quote{你在这里有一份。} 哦，愿你们蒙新郎喜悦；愿那说话者未受玷污，而且更糟的是，心受玷污，而非身体！你们爱的是人而不是基督；爱那说\Quote{就是我施洗}的，却听不进新郎的朋友说\BibleQuote{这就是那施洗的}，也听不进他说\BibleQuote{有新娘的是新郎。} 他说：我没有新娘；但我是什么呢？\BibleQuote{但新郎的朋友，站着听祂，因新郎的声音而大大欢喜。}\par

15. 显然，那么，我的弟兄们，守贞、节制、施舍对那些人毫无益处。所有在教会中受赞扬的行为对他们都无益处；因为他们撕裂了合一，即那件爱德的\Quote{长衣}。他们做了什么？他们中有许多能言善辩的人；伟大的口才，如流的口舌。他们说天使的话语吗？让他们听新郎的朋友，他为新郎而不是为自己发热心：\BibleQuote{即使我说人天使的话，却没有爱德，我就成了发声的铜，或响亮的钹。} \BibleRef{格前 13:1}\par

16. 但他们说什么呢？\Quote{我们有圣洗。} 你们有，但不是自己的。有是一回事，拥有是另一回事。圣洗你们有，因为你们领受而受洗，领受而受光照，只要你们自己没有成为黑暗；而你们施予时，是以仆役的身份施予，而非主人；是作为传令官宣告，而非作为法官。法官通过传令官说话，然而在案卷上写的不是\Quote{传令官说，}而是\Quote{法官说。} 因此，要看你们所施予的是否凭权柄属于你们。但如果你们领受了，就要同新郎的朋友一起承认：\BibleQuote{除非从天上赐给他，人什么也得不到。} 同新郎的朋友一起承认：\BibleQuote{有新娘的是新郎；新郎的朋友侍立静听，} 啊，但愿你们也侍立静听他，而不是跌倒，去听你们自己！因为听从他，你们就会站立并聆听；因为你们会说话，而你们的头已充满骄傲。我说，教会如果我是新娘，如果我领受了聘礼，如果被那血的代价所赎回，我就听新郎的声音；我也听新郎朋友的声音，只要他荣耀我的新郎，而不是自己。让朋友说：\BibleQuote{有新娘的是新郎；新郎的朋友侍立静听他，因新郎的声音而喜乐万分。} 看，你们有圣事；我承认你们有。你们有外形，但你们是砍下的树枝；你们有外形，我要的是根。没有根，外形就不结果实；但根在哪里？岂不就在爱德中吗？你们听听那砍下的树枝的外形：让保禄说话：\Quote{我若明白一切奥秘，}他说，\Quote{且有一切预言和一切信德}（多么大的信德！），\Quote{以致能移山，却没有爱德，我什么也不算。}\par

17. 那么，不要让人给你们讲无稽之谈。\Quote{彭提乌斯行了奇迹；多纳图斯祈祷，天主从天上应允了他。} 首先，他们要么是被骗，要么是骗人。最后，就算他能移山：\BibleQuote{却没有爱德，}宗徒说，\BibleQuote{我什么也不算。} 让我们看看他是否有爱德。我原本相信他有，如果他没有分裂合一。因为对于那些我可称为行奇迹者的人，我的天主提醒了我，说：\BibleQuote{在末后时期，将有假先知兴起，行异兆和奇事，以致如果可能，连被选者也要被迷惑；看，我已预先告诉了你们。} \BibleRef{谷 13:22} 因此新郎告诫了我们，我们连奇迹也不应受骗。有时，一个逃兵会吓唬一个纯朴的乡下人；但不愿受惊或被诱的人所注意的是，他是否属于军营，是否具有他所标记的身份。那么，我的弟兄们，让我们持守合一：没有合一，即使行奇迹者也算不了什么。以色列人民在合一中，却没有行奇迹：法郎的术士不在合一中，却行了与梅瑟相似的作为。\BibleRef{出 7:12} 以色列人民，如我所说，没有行奇迹。谁同天主一起得救？是那些行奇迹的人，还是那些不行奇迹的人？宗徒伯多禄使死人复活：\Person{西满·玛古斯}行了许多事：那里有些基督徒既不能做伯多禄所做的，也不能做西满所做的；他们因什么而喜乐？因他们的名字被记录在天上。因为这是我们的主耶稣基督在门徒回来时，因外邦人的信德所说的话。门徒们自己确实夸耀说：\BibleQuote{主，因着你的名字，连魔鬼也屈服于我们。} 他们承认得对，他们将光荣归于基督的名字；然而他对他们说什么？\BibleQuote{你们不要因魔鬼屈服于你们而欢喜，却要因你们的名字被记录在天上而欢喜。} \BibleRef{路 10:17} 伯多禄驱逐魔鬼。某个老寡妇，某个平信徒，有爱德而持守信德的完整，确实不做这事。伯多禄是身体上的眼睛，那人是手指，然而他在与伯多禄相同的身体内；如果手指的能力不如眼睛，它却没有从身体上砍下。作一个手指而留在身体内，好过作一个眼睛而被从身体中拔出。\par

18. 因此，我的弟兄们，不要让任何人欺骗你们，不要让任何人引诱你们：要爱基督的平安，他为你们被钉十字架，而他是天主。保禄说：\BibleQuote{栽种的不算什么，浇灌的也不算什么，只算那使之生长的天主。} \BibleRef{格前 3:7} 我们中有人说自己算什么吗？如果我们说自己算什么，而不把光荣归于他，我们就是奸夫；我们渴望自己被爱，而非新郎。你们要爱基督，也要爱我们在基督内的我们，你们也在他内被我们所爱。让肢体彼此相爱，但都要活在元首之下。确实，我的弟兄们，我不得不说了很多，却仍然说得少：我没能完成这段；天主会在适当的时候帮助我们完成。我不愿再加重你们心灵的负担；我愿他们自由地为那些仍然耳聋而不明白的人叹息和祈祷。\par

\SourceRecord{Translated by John Gibb.；Nicene and Post-Nicene Fathers, First Series；Vol. 7.；Edited by Philip Schaff.；Buffalo, NY: Christian Literature Publishing Co.,；1888.；<http://www.newadvent.org/fathers/1701013.htm>.}

% structure: p0001 source=h1 target=section reason=page-title

\section{第十四篇讲道（\BibleRef{若 3:29-36}）}

1. 这段神圣福音的教训向我们显示了我们的主耶稣基督天主性的卓越，以及那获得\Quote{新郎的朋友}称号之人的谦卑；为使我们能区分\Quote{那人}——只是人——与\Quote{那人}——也是天主。因为\Quote{那人}也是天主者，便是我们的主耶稣基督：在万世之前是天主，在此世的世代中成为人：出自圣父的天主，出自童贞玛利亚的人，然而同是一位主、救主耶稣基督，圣子，既是天主又是人。但若翰，一位出类拔萃蒙受恩宠的人，被派遣在祂之前，是由那\Quote{光}所光照的人。因为关于若翰，经上说：\BibleQuote{他不是那光，但他是要为光作证。}他本人确实可以被称作光，而且名正言顺；但那是被光照的光，而非光照人的光。那光照人的光，与被光照的光，是不同的事物：因为连我们的眼睛也被称为光（\Emph{lumina}），但当我们在黑暗中睁开眼睛时，它们却看不见。但那光照人的光，是自有的、为自身的光，不需要另一道光来照耀；但其余的一切都需要它才能发光。\par

2. 因此若翰承认了祂：正如你们所听到的，当耶稣正在收纳许多门徒时，他们向若翰报告，似乎是激起他的嫉妒——因为他们好像是出于嫉妒讲述此事——说：\Quote{看，祂收的门徒比你还多。}若翰承认了自己是什么，并因此堪当属于祂，因为他不敢断言自己就是耶稣之所是。现今若翰说：\BibleQuote{人不能领受什么，除非是上天赐予他的。}因此基督赐予，人领受。\BibleQuote{你们自己可以给我作证，我曾说过：我不是基督，而是被派遣在祂前面的。那有新娘的，是新郎；但新郎的朋友，站着听祂，因新郎的声音而大为欢喜。}他不是凭自己给自己喜乐。凡愿凭自己喜乐的，必将忧愁；但凡愿以天主为喜乐的，必将永远喜乐，因为天主是永恒的。你愿拥有永恒的喜乐吗？紧贴那永恒者。若翰声明自己就是这样的人。\BibleQuote{新郎的朋友因新郎的声音而欢喜，}不是因为自己的声音，而且\BibleQuote{站着听。}因此，他若跌倒，就听不见祂了：因为论及那跌倒的一位，经上说：\BibleQuote{他没有在真理中站住}；\BibleRef{若 8:44}这是指魔鬼说的。所以新郎的朋友应当\BibleQuote{站着听。}什么叫站着？就是存留在他所领受的恩宠中。他听见那使他欢喜的声音。若翰就是这样：他知道自己因何欢喜；他没有窃取自己所不是的；他承认自己是被光照的，而非光照人的。\BibleQuote{但那光是真光，}圣史说，\BibleQuote{照亮每一个来到这世上的人。}如果\BibleQuote{每一个人，}那么也包括若翰自己，因为他也是人。何况，虽然妇人中所生的没有兴起比若翰更大的，但他自己也是妇人所生的一员。他能与那位因祂的旨意而由独一无二、非凡的诞生所生者相比吗？因为主的两种降生都是绝无仅有的，无论是神性的还是人性的：按神性祂没有母亲，按人性祂没有父亲。所以若翰不过是其余人中的一个：然而他得了更大的恩宠，以致妇人所生的没有兴起比他更大的；他给了我们的主耶稣基督如此伟大的见证，称祂为新郎，而自己为新郎的朋友，然而不配解新郎的鞋带。亲爱的诸位，关于这一点你们已经听了很多；让我们来看下面的话，因为有点难懂。但正如若翰自己所说：\BibleQuote{人不能领受什么，除非是上天赐予他的，}凡是我们没有理解的，让我们求那从天赐予者吧：因为我们是人，除非那位不是人的（天主）赐给我们，否则我们不能领受什么。\par

3. 现在以下是紧接着的话：若翰说：\BibleQuote{所以，我的这份喜乐已经满足了。}他的喜乐是什么？就是因新郎的声音而欢喜。它在我内满足了：我有了我的恩宠；我不更多为自己妄取，以免也失去我已领受的。这喜乐是什么？\Quote{欢欢喜喜地因新郎的声音而喜乐。}那么，一个人可以明白，他不应该因自己的智慧而喜乐，而应因他从天主所领受的智慧而喜乐。他不要祈求更多，以免失去他所找到的。因为许多人，因自称为明智，反而成了愚妄的。宗徒谴责他们，论到他们说：\BibleQuote{因为天主的事，他们原本知道，天主也给他们显明了。}你们听他论到某些忘恩负义、不虔敬的人所说的话：\BibleQuote{自从创造世界以来，祂那看不见的美善，都可凭祂所造的万物辨认出来，祂永恒的德能以及神性，以致他们无可推诿。}为什么无可推诿？\BibleQuote{因为他们虽然认识天主}（他没有说：\Quote{因为他们不认识天主}）\BibleQuote{却没有尊崇祂为天主，也没有感谢祂；反而在推想中变为虚妄，他们的愚心就昏暗了；他们自称为明智，反而成了愚妄的。}\BibleRef{罗 1:19}\par

4. \Quote{他必兴盛，我必衰微。}这是什么意思？他必受高举，我必受贬抑。耶稣如何能兴盛？天主如何能兴盛？完美的不会增长。天主既不增长也不减少。因为如果他增长，他就不是完美的；如果他减少，他就不是天主。而耶稣既是天主，又如何能兴盛？如果就人的状况而言，因为他屈尊成为人，曾是婴孩；而且，虽是圣言，却躺在马槽里；虽是母亲的造物主，却吸吮她的乳汁：那么耶稣在肉身上年龄增长，这也许就是为什么说\Quote{他必兴盛，我必衰微。}但为何如此？就肉身而言，若翰与耶稣年龄相同，相差六个月；他们一同长大；如果我们的主耶稣基督愿意在死前在世更久，而若翰也与他同在，那么，正如他们一同长大，他们也会一同衰老：那么，又如何说\Quote{他必兴盛，我必衰微}呢？尤其是，我们的主耶稣基督现在已三十岁，一个三十岁的人还会增长吗？从这个年龄开始，人开始走下坡，趋向壮年，再到老年。再者，即使他们两人都是少年，他也不会说\Quote{他必兴盛，}而会说：我们一起增长。但现在两人都是三十岁。六个月的间隔在年龄上没有差别；这种差别是通过阅读而非容貌发现的。\par

5. 那么，\Quote{他必兴盛，我必衰微}是什么意思？这是一个伟大的奥秘！在主耶稣来临之前，人们自夸；他作为人来，是要减少人的荣耀，增加天主的荣耀。他无原罪而来，却发现众人都在罪中。如果他来是为消除罪，那么天主可以白施恩宠，人应当认罪。因为人的认罪是人的谦卑；天主的怜悯是天主的崇高。因此，既然他来是为赦免人的罪，让人承认自己的卑微，让天主显示他的怜悯。\Quote{他必兴盛，我必衰微：}就是说，他必施予，我必领受；他必受荣耀，我必认罪。让人知道自己的状况，并向天主认罪；且听宗徒对那骄傲自大、一心抬高自己的人说：\Quote{你有什么不是领受的呢？既然是领受的，为什么还夸口好像不是领受的呢？}\BibleRef{格前 4:7}那么，让人明白他所领受的；当他把自己没有的东西当作自己的时候，就让他衰微：因为对他而言，天主在他内受光荣是好的。让他衰微在自己内，好能在天主内兴盛。基督与若翰用他们的死亡预示了这些见证和这真理。因为若翰被斩首而衰微；基督被钉在十字架上而受高举；所以即使在死亡中，也显明了\Quote{他必兴盛，我必衰微}。再者，基督诞生时，白昼开始变长；若翰诞生时，白昼开始变短。因此，他们的诞生和死亡为若翰的话作证，他说：\Quote{他必兴盛，我必衰微。}愿天主的荣耀在我们内增加，我们自己的荣耀减少，好让我们的荣耀也在天主内增加！因为这就是宗徒所说的，这就是圣经所说的：\Quote{凡夸耀的，应当因主而夸耀。}\BibleRef{格前 1:31}你要夸耀自己吗？你会增长，但你的邪恶会增长得更坏。因为谁变得更坏，就该正当地衰微。因此，让那永远完美的天主增长，并在你内增长。因为你越理解天主，越领悟他，他就似乎在你内增长；但他本身不增长，因为是永远完美的。你昨天懂得一点，今天懂得更多，明天将懂得更多：天主的真光在你内增加；仿佛天主在增长，而他是永远完美的。正如一个人的眼睛从过去的盲目得医治，开始看见一点微光，第二天看见更多，第三天更多：对他来说，光似乎在增长；然而光本身是完美的，无论他看见与否。内在的人也是如此：他确实在天主内进步，天主似乎在他内增长；但人本身在衰微，以便从自己的荣耀中降下，升入天主的荣耀。\par

6. 我们刚才所听的，现在显明而清晰。\Quote{从天上来的是在万有之上。}看他对基督说了什么。他自己呢？\Quote{属于地的是出于地，谈论地的事。从天上来的是在万有之上}——这是基督；而\Quote{属于地的是出于地，谈论地的事}——这是若翰。难道就是这样吗：若翰属于地，谈论地的事？他给基督作的全部见证是谈论地的事吗？当若翰为基督作证时，从他口中出来的岂非天主的声音？那么他如何谈论地的事呢？这是指人而言。就人自身而言，他属于地，谈论地的事；而当他谈论一些神圣的事时，他是被天主光照的。因为他若不被光照，他就是谈论地的事的地。天主的恩宠本身是独立的，人的本性本身是独立的。只要考察人的本性：人生下来长大，学会人的习俗。他除了地、地的事，还知道什么？他谈论人的事，知道人的事，思念人的事；属血肉的，他属血肉地判断，属血肉地推测：看！完全是人的样子。愿天主的恩宠来到，光照他的黑暗，正如经上说：\Quote{你要点亮我的烛光，主；我的天主，求你光照我的黑暗。}愿它接纳人的心思，转向自己的光；他立刻开始说，如宗徒所说：\Quote{然而不是我，而是天主的恩宠与我同在；}\BibleRef{格前 15:10}和\Quote{我如今活着，却不是我在生活，而是基督在我内生活。}\BibleRef{迦 2:20}这就是说：\Quote{他必兴盛，我必衰微。}这样，若翰：就若翰而言，他属于地，谈论地的事；凡是你们从若翰听到的神圣之事，都是出自光照者的，而不是出自领受者的。\par

7. \BibleQuote{他从天上来，是在万有之上；他把他所看见所听见的，见证出来，却没有一个人接受他的见证。} 祂从天上来，是在万有之上，就是我们的主\Person{耶稣基督}；关于祂，前面曾说：\BibleQuote{没有人上过天，除了那自天降下而仍在天上的人子。} 祂是在万有之上；\BibleQuote{他把他所看见所听见的，说出来。} 况且，祂有一位圣父，祂自己就是天主子；祂有一位圣父，祂也听自圣父。祂从圣父所听见的，是什么呢？谁能够阐发这个？什么时候我的舌头，什么时候我的心才能胜任，无论是心去理解，还是舌去说出，那圣子从圣父所听到的呢？难道是圣子听见了圣父的圣言吗？不，圣子就是\OriginalTerm{圣言}{Word}。你看，人的一切努力在此何等疲竭；你看，我们内心的一切猜测，我们昏暗理智的一切紧张，在此都失败了。我听见圣经说，圣子所说的是祂从圣父所听见的；我又听见圣经说，圣子自己就是圣父的圣言：\BibleQuote{在起初已有圣言，圣言与天主同在，圣言就是天主。} 我们所说的言语是短暂易逝的：一旦你的话语从口中发出声音，它就过去了；它发出响声，然后归于沉寂。你能追上你的声音，并留住它让它停住吗？然而，你的思想存留下来，而从那个存留的思想，你发出许多转瞬即逝的言语。我们要说什么呢，弟兄们？当天主说话时，祂发出了声音，或响声，或音节吗？如果祂发出了，祂用的是什么语言？是希伯来语，还是希腊语，还是拉丁语？在有民族区分的地方，语言是必要的。但那里没有人能说天主用这种或那种语言说了话。观察你自己的心。当你孕育一个将要说出的话语—因为我要说，如果我能，我们可以在自己身上观察到什么，并非借此理解那个—嗯，当你孕育一个要说出的话语，你的意思是说出一个事物，而那个事物的概念本身已经是你心里的一个话语：它还没有出来，但已经在心里诞生，等待着出来。但你考虑到那将要受话的人，就是你要与他说话的人：如果他是拉丁人，你就寻找拉丁语表达；如果是希腊人，你就想希腊语词；如果是布匿人，你就考虑你是否知道布匿语。由于听者的多样性，你求助于各种语言来说出所孕育的话语；但那个概念本身并不受任何特定语言的束缚。因此，当天主说话时，祂不需要一种语言，也不采取任何形式的言辞，圣子如何听见祂呢？既然天主的说话就是圣子自己？事实上，正如你心中有你所说的话语，它与你同在，并且就是那精神的概念本身（因为正如你的灵魂是精神，你所孕育的话语也是精神；因为它尚未接受声音而被音节分开，而是留在你心里，在心灵的明镜中）；同样，圣父发出了祂的圣言，即生了圣子。而你的确在你心中按时生了话语；圣父无时间地生了圣子，祂藉着圣子创造了所有时间。因此，圣子既是天主的圣言，祂对我们说话时说的不是祂自己的话，而是圣父的话；当祂说圣父的话时，祂愿意把自己说给我们听。这就是\Person{若翰}所说的话，既合适又必要；我们已按我们的能力解释了。凡心尚未获得对如此伟大之事适当感知的人，他有地方可转向，有门可敲，有对象可问，可寻求，可领受。\newline{}\par

8. \BibleQuote{那從天上來的，是在萬有之上；祂所見所聞的，就為此作證，卻沒有人接受祂的見證。}如果沒有人，祂來是為什麼？他的意思是，某一類的人沒有人。有一些人預備受天主的義怒，與魔鬼一同被定罪；這些人中，沒有人接受基督的見證。因為如果根本沒有人，沒有任何一個人接受，那麼這些話\BibleQuote{但接受祂見證的人，就蓋上了印，證明天主是真實的}是什麼意思呢？那麼，肯定不是\Emph{沒有人}，如果你自己說\BibleQuote{接受祂見證的人，就蓋上了印，證明天主是真實的}。也許若翰被問時，會回答說：我知道我所說的，當我說\Emph{沒有人}時。事實上，有些人是生來就屬於天主義怒的，並且預先被知道。因為天主知道誰將信、誰將不信；祂知道誰將在他們所信的道上堅持到底，誰將跌倒；所有將得永生的人，都被天主數算；祂已經知道分別出來的人。如果祂知道這事，並藉著祂的聖神將這事啟示給先知，祂也把這事給了若翰。如今若翰觀察，不是用他的眼睛——因為就他自己而言，他是屬地的，說屬地的事——而是藉著他從天主所領受的聖神恩寵，他看到了一個不敬虔、不信的百姓。默想那百姓的不信，他說：\BibleQuote{那從天上來的，祂的見證，沒有人接受。}沒有人？誰當中？就是那些將在左邊的人，那些將聽到\BibleQuote{你們這些可咒罵的，離開我，到那為魔鬼和他的使者預備的永火裡去！}的人。那些接受的是誰？就是那些將在右邊的人，那些將聽到\BibleQuote{我父所祝福的，你們來承受那從創世以來為你們預備的國度吧。}的人。因此，他在聖神中觀察到一個區分，但在人類中卻是混雜的；那尚未在空間上分開的，他在理解中、在心靈的視野中將之分開；他看到了兩個百姓，一個是信者的，一個是不信者的。他將思想專注於不信者，說：\BibleQuote{那從天上來的，是在萬有之上；祂所見所聞的，祂為此作證，卻沒有人接受祂的見證。}然後他將思想從左邊轉向右邊，繼續說：\BibleQuote{接受祂見證的人，就蓋上了印，證明天主是真實的。}\BibleQuote{蓋上了印，證明天主是真實的}是什麼意思？豈不是說人是虛謊的，天主是真實的嗎？因為沒有任何人能說任何真理，除非他被那不能說謊的所光照。所以，天主是真實的；但基督是天主。你願意證明這一點嗎？接受祂的見證，你就會發現。因為\BibleQuote{那接受祂見證的人，就蓋上了印，證明天主是真實的。}誰是真實的？那從天上來、在萬有之上的，就是天主，是真實的。但如果你還不明白祂是天主，你就還沒有接受祂的見證；接受吧，你就蓋上印；你就能確信地理解，明確地承認，天主是真實的。\par

9. \BibleQuote{因為天主所派遣的，講論天主的話。}祂自己就是真天主，並有天主派遣了祂：天主派遣了天主。兩者結合，一位真天主，由天主所派遣。分別問及他們每一位，祂是天主；問及二者，他們是天主。不是單獨每一位是天主，而兩者是眾神；而是每一位是天主，二者也是天主。因為在那裡聖神的愛德是那麼偉大，合一的平安是那麼深厚，以至於當你問及他們每一位時，回答你是：天主；當你問及天主聖三時，你得到的回答是：天主。因為人的心靈，當它依附於天主時，就成為一個心靈，正如宗徒明言：\BibleQuote{那與主結合的，便是與祂成為一神；}\BibleRef{格前 6:17}那麼，與聖父結合的同等聖子，與祂一同成為一位天主，豈不更加如此！請聽另一個見證。你知道有多少人相信了，他們變賣一切所有，放在宗徒腳前，按各人的需要分給各人；聖經論及那聖徒的聚集說：\BibleQuote{他們一心一意在主內。}\BibleRef{宗 4:32}如果愛德使這麼多靈魂成為一個靈魂，這麼多人心成為一顆心，那麼聖父與聖子之間的愛德該是多麼偉大！肯定比那些成為一心的人之間的愛德更大。如果眾弟兄的心因愛德而成為一心，眾弟兄的靈魂因愛德而成為一靈，你還能說天主聖父和天主聖子是兩位嗎？如果他們是兩位天主，那他們之間就沒有最高的愛德。因為如果在這裡的愛德如此之大，以至於使你的靈魂和你朋友的靈魂成為一個靈魂，那麼聖父和聖子怎能不是一位天主呢？願真誠的信德遠離這種想法。總之，那愛德是多麼卓越，從此可知：許多人的靈魂是許多的，但如果他們彼此相愛，就成為一個靈魂；但對於人而言，仍可稱為許多靈魂，因為合一不那麼緊密。但在那裡，你理當說一位天主；不能說兩位或三位天主。由此，向你顯示了愛德的至上和超越的卓越，以至於沒有比這更大的了。\par

10. \Quote{因为天主所派遣的，讲论天主的话。}当然，他是指着基督说的，为将自己与基督区别开来。那么又怎样呢？难道天主没有派遣若翰本人吗？难道他自己没有说：\Quote{我是祂前面被派遣的}？以及\Quote{那派遣我以水施洗的}？难道关于若翰不是如此说的：\Quote{看，我派遣我的使者在你面前，他必预备你的道路}？\BibleRef{拉 3:1}他（就是那被称为超越先知的人）难道不是自己讲论天主的话吗？那么，如果天主也派遣了他，他也讲论天主的话，我们如何理解他明确地指着基督说\Quote{天主所派遣的讲论天主的话}呢？但请看他又加上：\Quote{因为天主赐圣神是没有限量的。}这是什么意思，\Quote{因为天主赐圣神是没有限量的}？我们发现天主确实限量赐予圣神。请听宗徒所说：\Quote{按基督恩赐的尺度。}\BibleRef{弗 4:7}祂赐给人是有限量的，赐给独生子却是没有限量的。祂如何限量赐给人呢？\Quote{有人由圣神领受了智慧的言语，有人由同一圣神领受了知识的言语，有人由同一圣神领受了信德，有人领受了先知之恩，有人领受了辨别神恩的能力，有人领受了各种异语，有人领受了治病的奇恩。难道都是宗徒？难道都是先知？难道都是教师？难道都是行奇迹的？难道都有治病的恩赐？难道都说异语？难道都解释异语？}\BibleRef{格前 12:8}这人有一种恩赐，那人有另一种；那人所有的，这人没有：这是有限量的，是恩赐的某种分配。因此，赐给人是有限量的，而他们之间的和谐共融形成一体。就如手领受一种恩赐来工作，眼另一种来看，耳另一种来听，脚另一种来行走；然而那完成一切的是同一灵魂，在手里工作，在脚里行走，在耳里听，在眼里看；同样，信徒的恩赐也是多样的，如同分给肢体一样，分给他们，各按自己适当的尺度。但基督，祂给予，却不受限量领受。\par

11. 现在请听接下来的话：因为祂论及圣子曾说：\BibleQuote{因为天主赐圣神是没有限量的：圣父爱圣子，已将一切交在祂手中，}祂又加上：\BibleQuote{已将一切交在祂手中，}为叫你也知道这里是以何等区别说\BibleQuote{圣父爱圣子。}为什么？难道圣父不爱若翰吗？然而祂并未将一切交在他手中。难道圣父不爱保禄吗？然而祂并未将一切交在他手中。\BibleQuote{圣父爱圣子：}但如同父亲爱，不像主人爱仆人；如同爱独生子，不像爱义子。因此\BibleQuote{已将一切交在祂手中。}\Quote{一切}是什么意思？就是圣子应当与圣父一样。祂生祂与自己同等，在祂内，与天主同等并非僭夺。\BibleQuote{圣父爱圣子，已将一切交在祂手中。}所以，祂既屈尊将圣子派遣给我们，我们不要以为那派遣给我们的，是比圣父小的。圣父在派遣圣子时，派遣了祂的另一自我。\par

12. 但门徒们仍然以为圣父比圣子更伟大，因为他们只看见血肉，而不懂得祂的天主性，就对祂说：\BibleQuote{主，请将父显示给我们，我们就满足了。}这无异于说：\Quote{我们已经认识你，且因认识你而称颂你：因为我们感谢你将自己显示给我们。但我们至今还不认识父：因此我们的心火热，被一种神圣的渴望所占据，就是渴望看见那派遣你的父。请将祂显示给我们，我们就再也不向你要什么了：因为当祂被显示时，我们就满足了，没有比祂更大的了。}美好的渴望，美好的愿望，但智慧不足。现在主耶稣自己，视他们为寻求大事的小人物，而祂自己在小人物中是大人物，但在小人物中仍是小人物，就对那说过这话的门徒之一斐理伯说：\BibleQuote{斐理伯，这么长久的时间我与你们在一起，你还不认识我吗？}斐理伯或许可以回答：我们认识了你，但我们何曾对你说：将你自己显示给我们？我们认识了你，但我们寻求认识的是父。祂立刻加上：\BibleQuote{看见了我的，就看见了父。}\BibleRef{若 14:8}那么，既然与父平等的那一位被派遣了，我们不要从肉体的软弱来估计祂，而应想到那穿上肉体却不被肉体压制的威严。因为，祂与父同为天主，却成了人而生活在人中间，为叫你通过这成了人的一位，能够成为接受天主的。因为人不能接受天主。人能看见人；天主却无法领悟。为什么不能领悟天主？因为没有心灵的眼睛来领悟祂。内心有些紊乱，有些不安；人的身体眼睛健全，但心灵的眼睛却病了。祂成了人，呈现在肉体的眼前，使你能因相信那有形可见的，而得到医治，得以看见那原本不能以精神看见的祂。\BibleQuote{斐理伯，这么长久的时间我与你们在一起，你还不认识我吗？看见了我的，就看见了父。}为什么他们没有看见祂？看，他们确实看见了祂，却没有看见父：他们看见了肉体，但威严是隐藏的。那些爱祂的门徒所看见的，那些钉死祂的犹太人也看见了。因此，祂内在的一切，且如此内在于肉体，以至于当祂来到肉体时，仍与父同在。\par

13. 肉性的思想不明白我所说的：让它暂缓理解，从信德开始；让它听接下来所说的：\Quote{信子的人，有永生；不信子的人，不得见生命，天主的义怒却常在他身上。}他没有说，天主的义怒临到他；而是说\Quote{天主的义怒常在他身上。}凡生为有死之人，都带有天主的义怒。什么义怒？就是亚当最初领受的义怒。因为如果第一个人犯了罪，听到判决\Quote{你必定要死，}他就成了有死的，而我们也开始生为有死之人；我们生来就带有天主的义怒。那一位出自这苗裔，没有罪，却穿上了肉身和死亡。如果祂与我们一同分享了天主的义怒，难道我们迟迟不愿与祂分享天主的恩宠吗？因此，那不信子的人，\Quote{天主的义怒常在他身上。}什么义怒？就是宗徒所说的：\Quote{我们从前也都在本性上为义怒之子，和其余的人一样。}\BibleRef{弗 2:3}所以众人都是义怒之子，因为来自死亡的诅咒。要信基督，因为你被造为可死的，好使你接受祂——那位不朽的；当你接受了祂的不朽，你就不再是可死的了。祂活了，你曾死；祂死了，为叫你活。祂给我们带来了天主的恩宠，除去了天主的义怒。天主战胜了死亡，免得死亡战胜人。\par

\SourceRecord{Translated by John Gibb.；Nicene and Post-Nicene Fathers, First Series；Vol. 7.；Edited by Philip Schaff.；Buffalo, NY: Christian Literature Publishing Co.,；1888.；<http://www.newadvent.org/fathers/1701014.htm>.}

% structure: p0001 source=h1 target=section reason=page-title

\section{讲道集第15篇（\BibleRef{若 4:1-42}）}

亲爱的，这并不是新事：圣史若望如同雄鹰，飞得更高，翱翔于尘世的黑暗迷雾之上，以更坚定的目光注视真理之光。从他的福音中，经由我们的职务，在主帮助下，我们已经处理并讨论了许多内容；今天所读的段落依序而来。我即将要说的，在主允许下，你们中许多人将以这样的方式听见，以至于你们是在复习已知的，而非学习未知的。然而尽管如此，你们的注意力不应松懈，因为这不是获取知识，而是复习知识。这段经文已被诵读，我们手中拥有要论述这段经文的权柄——就是主耶稣在雅各伯井旁与撒玛黎雅妇人的对话。那里所说的是伟大的奥迹，是伟大事物的预象；喂养饥饿的人，使疲倦的灵魂得滋润。\par

主知道这事后，\BibleQuote{当祂听说法利塞人获悉祂收的门徒比\Person{若翰}还多，并为他们施洗（虽然\Person{耶稣}不亲自施洗，而是祂的门徒施洗），祂就离开\Place{犹太}，再次往\Place{加里肋亚}去。}我们不应在此过多论述，免得因沉溺于明显之事，而缺乏探究和揭示晦暗之事的时间。当然，如果主看到他们得知祂收的门徒更多、施洗更多这一事实，对法利塞人的救恩如此有益，以至于他们会成为祂的门徒，并渴望受祂的洗，那么祂就不会离开\Place{犹太}，而是会为了他们而留在那里。但正因为祂知道他们知道这一事实，同时也知道他们的嫉妒，以及他们得知此事不是为了跟随，而是为了迫害祂，祂就离开了那里。的确，即使在场，祂也能使他们抓不住祂，如果祂不愿意；祂有能力不被杀害，如果祂不愿意，因为如果祂不愿意，祂也有能力不诞生。但因为在祂作为人所做的一切事上，祂都是为将要相信祂的人树立榜样（这样，天主的任何仆人如果看到迫害者的愤怒，或那些企图将他的灵魂引入邪恶者的愤怒，而退到别处，他并不犯罪；但如果天主的仆人这样做，他可能显得犯下罪过，除非主先做了这事），那位善师这样做是为了教导我们，并非因为祂害怕。\par

也许你们会惊讶为什么说\BibleQuote{\Person{耶稣}施洗比\Person{若翰}还多，}说了这话之后，又补充说\BibleQuote{虽然\Person{耶稣}不亲自施洗，而是祂的门徒施洗。}这是怎么回事？难道前面的陈述是假的，然后被这个补充纠正了？还是两者都真，即\Person{耶稣}既施洗又不施洗？事实上祂施洗，因为是祂洁净；祂不施洗，因为不是祂用手触摸。门徒提供身体的职务；祂提供祂尊威的帮助。既然祂从未停止洁净，何时祂会停止施洗？同一位\Person{若翰}，以洗者的身份，论到祂说：\BibleQuote{这一位就是施洗者。}因此，\Person{耶稣}仍然在施洗；只要我们继续受洗，\Person{耶稣}就在施洗。让人毫不畏惧地来到下面的施行者那里；因为他有一位上面的师傅。\par

但或许有人说，\Person{基督}确实施洗，但只以圣神，不以身体。仿佛任何人有形的、可见的圣洗圣事竟是由别人而非祂的恩赐所灌注。你想知道是祂施洗，不仅以圣神，也以水吗？听宗徒说：\BibleQuote{正如\Person{基督}，}他说，\BibleQuote{爱了教会，并为了她舍弃了自己，用水借圣言洗涤她，使她洁净，好将教会呈献给自己，光荣的教会，没有瑕疵或皱纹，或任何类似的东西。}\BibleRef{弗 5:25} 祂洗涤她。如何？\BibleQuote{用水借圣言的洗涤。}什么是\Person{基督}的洗礼？用水借圣言的洗涤。除去水，就不是圣洗；除去圣言，也不是圣洗。\par

关于初步的情景，主因之与那妇人交谈，我们就说这么多；让我们看看剩下的事情；那些事情充满奥迹，并孕育着圣事。\BibleQuote{祂必须经过\Place{撒玛黎雅}。于是祂来到\Place{撒玛黎雅}的一座城，名叫\Place{息哈尔}，靠近\Person{雅各伯}给他儿子\Person{若瑟}的那块地。\Person{雅各伯}的泉源就在那里。}那是一口井；但每口井都是泉源，但不是每个泉源都是井。因为水从地中流出，自献于取它的人使用时，它被称为泉源；但如果可及且在地面，只称为泉源；但如果深且向下，则称为井，但并不失去泉源之名。\par

6. \BibleQuote{耶稣因行路疲倦，就那样坐在井旁。那时大约是第六时辰。}现在开始奥迹。因为耶稣疲倦并非无意义；天主的德能疲倦并非无意义；祂疲倦并非无意义，疲乏者因祂而复苏；祂疲倦并非无意义，因祂的缺席我们疲倦，因祂的临在我们坚强。然而耶稣疲倦，行路疲倦；祂坐下，且在井旁；是第六时辰，疲倦的祂坐下。这一切都暗示某事，旨在表明某事，它们使我们热切，鼓励我们叩门。愿祂亲自给我们和你们开启；祂竟屈尊劝勉我们说：\BibleQuote{叩门，就给你们开启。}耶稣是为你们行路疲倦。我们找到耶稣是德能，也找到耶稣是软弱：我们找到强壮的耶稣和软弱的耶稣：强壮，因为\BibleQuote{在起初已有圣言，圣言与天主同在，圣言就是天主：这圣言在起初就与天主同在。}你愿看见这天主圣子如何强壮吗？\BibleQuote{万物是藉着祂造成的，没有祂，什么也没有造成：}且无需劳动而造成了。那么什么能比祂更强壮？祂不劳动而造成了万物。你愿认识祂软弱吗？\BibleQuote{圣言成了血肉，居住在我们中间。}基督的德能创造了你，基督的软弱再造了你。基督的德能使不存在的存在；基督的软弱使存在的免于灭亡。祂以德能塑造了我们，以软弱寻找了我们。\par

7. 于是，祂如软弱者哺育软弱者，如同母鸡哺育小鸡；因为祂将自己比作母鸡：\BibleQuote{我多少次，}祂对\Place{耶路撒冷}说，\BibleQuote{愿意聚集你的子女，有如母鸡把小鸡聚集在翅膀底下，但你们不愿意！}\BibleRef{玛 23:37}弟兄们，你们看，母鸡如何因小鸡而变得软弱。没有别的鸟，当它是母亲时，能立刻被认出如此。我们看见各种麻雀在我们眼前筑巢；我们看见燕子、鹳鸟、鸽子每日筑巢；但我们不知道它们是父母，除非看见它们在巢上。但母鸡为它的雏鸡变得如此衰弱，即使小鸡不跟随着它，若你不见雏鸡，你仍立刻认出它是母亲。它双翅下垂，羽毛蓬乱，声音嘶哑，全身如此萎靡卑微，以致如我所说，即使你不见它的幼雏，你仍感到它是母亲。耶稣就是这样软弱，行路疲倦。祂的行路就是为我们而取的人性。因为那位无所不在、无所不在者，怎能行路呢？祂往哪里去，或从哪里来，除非祂取了可见血肉的形状，才能到我们这里来？因此，既然祂屈尊这样来到我们中间，藉所取的血肉以仆人的形状显现，这血肉的取用就是祂的行路。所以，\BibleQuote{行路疲倦}，除了在肉体内疲倦，还能是什么？耶稣在肉身内软弱：但你们不要软弱；要在祂的软弱中坚强，因为\BibleQuote{天主的软弱比人更强壮。}\par

8. 在这种形象之下，\Person{亚当}，那将来之人的预像，给了我们这奥迹一个伟大的指示；确切说，是天主在他内给的指示。因为亚当在沉睡中堪当接受妻子，那妻子是从他的肋骨造成的：因为从在十字架上睡眠的基督，即将诞生教会——即从祂的肋旁，当祂睡眠时；因为当祂挂在十字架上时，从被长枪刺透的肋旁，流出了教会的圣事。但是，弟兄们，我为什么选择这样说呢？因为正是基督的软弱使我们坚强。一个显著的形象先前在\Person{亚当}身上出现了。天主可以从男人取肉以造成女人，而且这似乎更合适。因为被造的是较弱的性别，软弱应是由肉而不是由骨造成；因为骨是肉中较坚强的部分。祂没有取肉来造女人；却取了一根骨头，女人由骨头形成，肉填补了骨头的位置。祂本可以骨还骨；祂本可以取一根肋骨，而是肉，来造女人。那么，这意味着什么？女人，可以这样说，是由骨头造成的，是强的；亚当，可以这样说，是由肉造成的，是弱的。这是基督和教会；祂的软弱是我们的力量。\par

9. 但为什么在第六时辰？因为在世界的第六时代。在福音里，把每一个时代算作一个时辰：第一个时代，从\Person{亚当}到\Person{诺厄}；第二个，从\Person{诺厄}到\Person{亚巴郎}；第三个，从\Person{亚巴郎}到\Person{达味}；第四个，从\Person{达味}到迁徙\Place{巴比伦}；第五个，从迁徙\Place{巴比伦}到\Person{若翰}的圣洗；从此是第六时代正在展开。你为何惊奇？耶稣来了，且藉着谦卑自己，来到井旁。祂疲倦地来了，因为祂带着软弱的肉身。在第六时辰，因为在世界的第六时代。到井旁，因为到了我们这居所的深处。为此圣咏说：\BibleQuote{我从深处呼求祢，上主。}祂坐下，如我所说，因为祂被谦卑了。\par

10. \BibleQuote{有一个妇人来了。}这是教会尚未成义但即将成义的预象，因为这就是本次讲论的主题。她来时不识主，却找到了主，并与之交谈。让我们看看是什么事，以及为何如此。\BibleQuote{有一个撒玛黎雅妇人来打水。}撒玛黎雅人不属于犹太民族：他们是外邦人，虽然住在邻近地区。要详细讲述撒玛黎雅人的起源需要很长时间；为了避免因冗长的论述而耽搁必要的事，只说我们认为撒玛黎雅人是外邦人就够了。并且，免得你以为我说这话是鲁莽而非真实，请听主耶稣自己论及那个撒玛黎雅人时所说的话，祂是十个被洁净的癞病人之一，只有他回来谢恩：\BibleQuote{不是十个人洁净了吗？那九个在哪里？除了这个外邦人，再没有别人回来归光荣于天主。}\BibleRef{路 17:17}这与实体的预象相关，因为这位预表教会的妇人来自外邦人：因为教会要来自外邦人，即犹太民族之外的人。因此，让我们在那妇人身上听见我们自己，在她身上认出我们自己，在她身上为我们自己感谢天主。因为她是一个预象，而非实体；因为她首先显出了预象，然后成为了实体。因为她相信了那藉着她将预象摆在我们面前的主。\BibleQuote{她于是来打水。}她只是来打水，如同人们习惯做的那样，无论男女。\par

11. \BibleQuote{耶稣对她说：\InnerQuote{请给我喝。}因为祂的门徒已进城买食物去了。于是撒玛黎雅妇人对祂说：\InnerQuote{你既是犹太人，怎么向我一个撒玛黎雅妇人要水喝呢？}原来犹太人与撒玛黎雅人素不来往。}可见他们是外邦人：事实上，犹太人不使用他们的器皿。那妇人带着打水的器皿，所以见一个犹太人向她讨水喝，就感到奇怪——这是犹太人不习惯做的事。但讨水喝的那位，渴慕的正是这妇人本身的信德。\par

12. 最后，请听是谁在要水喝：\BibleQuote{耶稣回答说：\InnerQuote{你若知道天主的恩赐，和对你说\InnerQuote{请给我喝}的是谁，你或许会求祂，而祂必赐给你活水。}}祂要水喝，却应许赐水；祂渴望如将得者，祂丰盈如将施者。\BibleQuote{你若知道，}祂说，\BibleQuote{天主的恩赐。}天主的恩赐就是圣神。但此刻祂仍谨慎地对妇人说话，逐步进入她的内心。或许祂正在教导她。有哪句劝勉比这更甘饴、更仁慈呢？\BibleQuote{你若知道天主的恩赐，}等等：至此祂让她悬疑。通常从泉源涌出的水称为活水；从雨水积聚在水池或水槽中的水不称为活水。它可能来自泉源，但若积聚在某处，不与源头相通，且水流中断，与泉源之水道分离，便不称为\BibleQuote{活水}：只有流动而取用的水才称为活水。那泉中就有这样的水。那么，为什么祂要应许赐予祂正在祈求的东西呢？\par

13. 妇人在悬疑中对祂说：\BibleQuote{主，你没有打水的器具，井又深。}请看她是如何理解活水的，仅仅是指那泉中的水。\BibleQuote{你要给我活水，而我带着打水的器具，你却没有。活水就在这里；你怎能给我呢？}她以另一种方式理解，并属肉体地领会，她是在以某种方式敲门，为让师傅开启那封闭的奥秘。她在无知中敲门，并非出于诚挚的目的；她仍是怜悯的对象，尚未受教。\par

14. 主更清楚地谈论那活水。妇人曾说过：\BibleQuote{难道你比我们的祖先雅各伯还大吗？他给了我们这口井，他自己、他的子女和牲畜都喝过这井的水。}你不能给我这井的活水，因为你没有打水的器具：也许你应许了另一泉源？你能比我们的祖先更好吗？他挖了这口井，自己和他的家人都在使用。那么，让主来宣告祂所称的活水是什么。\BibleQuote{耶稣回答说：\InnerQuote{凡喝这水的，还要再渴；但谁若喝了我所赐给他的水，他将永远不渴；我所赐给他的水，将在他内成为涌到永生的水泉。}}主已更公开地讲论：\BibleQuote{它将在他内成为涌到永生的水泉；喝这水的人将永远不渴。}还有比这更明显的吗？祂所应许的不是可见的水，而是无形之水？祂不是按属肉体的意思，而是按属灵的意义说话，这难道还不明显吗？\par

15. 然而，妇人仍然心系肉身：她因想到不再口渴而欣喜，以为主是按属肉体的意思应许了她；这确实会在将来应验，但是在死者的复活之时。她现在就渴望得到。天主确实曾赐予祂的仆人厄里亚，四十天不饥不渴。难道祂不能永远赐予吗？既然祂有能力赐予四十天？然而，她为此叹息，渴望没有匮乏，没有劳苦。要不断地来到那水泉，背负着重担来满足自己的需要，而当所打的水用完后又必须再来——这对她是每日的劳苦；因为她的需要得以缓解，而非消除。耶稣应许的这份恩赐令她喜悦；于是她求祂赐予活水。\par

16. 然而，我们不可忽略主所许诺的是一件属灵的事。所谓\Quote{凡喝这水的，还要再渴}是什么意思？就这水而言，这是真的；就这水所象征的而言，也是真的。因为井中的水是在黑暗深处的世俗享乐：人用情欲的器皿从中打水。他们俯身放下情欲，以触及从井底打上来的享乐，并享受那享乐和先前放下以取它的情欲。因为若没有预先放出情欲，便无法得到享乐。那么，把情欲看作器皿，把享乐看作从井底打上的水：当人得到这世界的享乐时，那就是他的肉、他的饮料、他的沐浴、他的表演、他的恋爱；难道他不会再渴吗？因此，祂说：\Quote{凡喝这水的，还要再渴}；但若他接受我赐的水，\Quote{他将永远不渴}。经上说：\Quote{我们必因祢殿中的美物饱足。}那么，祂要赐的是什么水呢？岂不就是那水，经上说：\Quote{在祢那里有生命的泉源}？因为那些\Quote{因祢殿中的肥甘沉醉}的人，又怎么会渴呢？\par

17. 祂所许诺的是一种喂养和圣神的丰富满盈；但妇人还不明白；她既不明白，又如何回答呢？\Quote{妇人对祂说：先生，请给我这水，使我不渴，也不用来这里打水。}缺乏迫使她劳苦，她的软弱在诉苦。但愿她听到了那邀请：\BibleQuote{凡劳苦和负重担的，你们都到我这里来，我要使你们安息！}\BibleRef{玛 11:28} 事实上，这正是耶稣对她说的，好使她不再劳苦；但她还不明白。\par

18. 最后，为要她明白，\Quote{耶稣对她说：你去，叫你的丈夫，然后到这里来。}这\Quote{叫你的丈夫}是什么意思？祂是想通过她的丈夫来给她那水吗？或者，因为她不明白，祂想通过她的丈夫来教导她？或许正如宗徒论及妇女时所说：\BibleQuote{她们若愿意学什么，可以在家里问自己的丈夫。}但宗徒这话是指没有耶稣在场教导的情况。总之，这是对宗徒禁止在教会中发言的妇女说的。\BibleRef{格前 14:34} 但当主亲自在场、亲自对她说话时，何需通过她的丈夫对她说话呢？祂岂是通过她的丈夫对玛利亚说话？玛利亚坐在祂脚前听祂的话，而玛尔大忙于许多侍奉，埋怨她妹妹的福分。\BibleRef{路 10:40} 因此，我的弟兄们，让我们听和明白主对妇人说的\Quote{叫你的丈夫}是什么意思。因为祂或许也对我们的灵魂说：\Quote{叫你的丈夫。}让我们也探问灵魂的丈夫。怎么，耶稣自己不已经是灵魂真正的丈夫吗？愿理解力在场，因为我们即将说的话只有细心的听众才能领会；因此愿理解力在场以领会，或许那理解力本身将被发现是灵魂的丈夫。\par

19. 耶稣见那妇人没有领悟，愿意她领悟，便对她说：\Quote{请你的丈夫来。}\Quote{因为你不明白我所说的原因，是你的理智不在场：我是按圣神说话，你却按肉身听话。我所说的事与耳朵的愉悦、眼睛、嗅觉、味觉、触觉无关；唯由心灵接收，唯由理智汲取。那理智不在你内，你怎能领悟我所说的话？\InnerQuote{请你的丈夫来}，把你的理智带来。你拥有灵魂算得了什么？不算什么，因为野兽也有灵魂。你比野兽好在哪？在于拥有理智，而野兽没有。}那么，\Quote{请你的丈夫来}是什么意思？\Quote{你不理解我，你不明白我：我对你讲论天主的恩赐，你的想念却是属肉体的；你不是不希望属肉体的口渴，我却是向心神说话：你的理智缺失了。\InnerQuote{请你的丈夫来。}不要像那无理智的骡马。}所以，我的兄弟们，拥有灵魂却没有理智，即是不使用理智，不按理智生活，就是野兽的生命。因为我们与野兽共享某些东西，就是那使我们活在肉身内的东西，但它必须被理智管辖。因为灵魂的动欲，它随从肉身，渴望不受约束地放纵于肉欲的享乐，被理智主宰。哪个应被称为丈夫？是统治的，还是被统治的？无疑，当生命井然有序时，理智统治灵魂，因为理智本身属于灵魂。因为理智并非灵魂以外的某物，而是灵魂的一个部分：如眼睛并非肉身以外的某物，而是肉身的器官。但眼睛虽是肉身的器官，却唯有它能享受光；其他肉身肢体虽可浸在光中，却无法感受光：唯有眼睛既沐浴其中又享受它。同样，在我们灵魂内有一种被称为理智的东西。灵魂的这个东西，称为理智和心智，被更高的光照亮。那更高的光，照亮人心智的，就是天主；因为\BibleQuote{那是那真光，照亮每个来到这世界的人。}这样的光就是基督，这样的光正与那妇人说话：然而她的理智不在场，无法被那光照亮；不仅是光照在她身上，更是要她享受它。因此主说：\Quote{请你的丈夫来}，仿佛在说：我愿意光照，但此处却没有我可光照的人；把理智带来，好使你受教，受其管辖。因此，把没有理智的灵魂当作妇人；拥有理智如同拥有丈夫。但这丈夫不能好好\Emph{管辖}妻子，除非他自己被更高的管辖。\BibleQuote{因为女人的头是男人，而男人的头是基督。}\BibleRef{格前 11:3}男人的头正与那妇人谈话，而男人却不在场。因此主仿佛说：把你的头带来，好让他接受他的头，便说：\Quote{请你的丈夫来，并到这里来；}也就是说，在这里，在此处：因为你如同缺席，既然你不理解在此处之真理的声音；要在此处，但不可独处；带着你的丈夫在此。\par

20. 丈夫尚未被叫来，她仍然不明白，仍然挂念肉身；因为男人缺席：\Quote{我没有，}她说，\Quote{丈夫。}主继续说并说出奥秘。你可以理解那妇人当时确实没有丈夫；她与某个男人同居，并非合法丈夫，与其说是丈夫，不如说是情夫。主对她说：\Quote{你说得对：我没有丈夫。}那么祢怎么说：\Quote{请你的丈夫来}？现在请听主如何深知她没有丈夫。\Quote{他对她说，}等等。免得妇人以为主说\Quote{你说得对：我没有丈夫}只是因为祂从她那里得知了这事，而非因祂自己的神性知道，请听你没有说过的事：\BibleQuote{因为你曾有五个丈夫，而现在你所有的并不是你的丈夫；这话你说得对。}\par

21. 祂再次敦促我们更精确地探究这五位丈夫的问题。实际上，许多人不太荒谬也不太不可能地将这妇人的五位丈夫理解为梅瑟的五卷书，因为撒玛黎雅人使用这些书卷，且遵守同一法律：他们正是从这法律获得割损。但鉴于随后的经文限制了我们的理解——\Quote{你现在所有的，不是你的丈夫}——在我看来，我们更容易将这五种身体感官理解为灵魂的前五位丈夫。因为人出生后，在能运用心智和理智之前，只受肉身的感官支配。在小孩子身上，灵魂寻求或回避所听见、所见、所嗅、所尝和所触摸到的事物。它寻求一切愉悦之物，回避一切冒犯之物，即这五种感官。起初，灵魂依照这五种感官生活，如同五位丈夫；因为它是被它们支配的。但为何称它们为丈夫？因为它们合法且正当：它们确由天主创造，是天主赐予灵魂的恩赐。灵魂在受这五位丈夫支配时仍然是软弱的，生活在这五位丈夫之下；但当灵魂达到运用理性的年龄，若被高尚的纪律和智慧的教导所接管，这五位男人便由一位真正的、合法的、比他们更美好的丈夫取代其统治，这丈夫既治理得更好，又永恒地治理，永恒地培育和教导她。因为这五种感官支配我们，不是为永恒，而是为了那些需要寻求或回避的暂时事物。但当受到智慧浸润的理智开始统治灵魂时，它现在不仅知道如何避免坑洼、在平地上行走——这是眼睛即使在灵魂软弱时也能向灵魂显示的——也不仅是被悦耳的声音所吸引、排斥刺耳的声音；不仅是对宜人的香气喜悦、拒绝难闻的气味；不仅是陶醉于甘甜、不悦于苦涩；不仅是因为柔软而舒适、因粗糙而受伤。因为所有这一切对处在软弱中的灵魂是必要的。那么，那理智使用什么规则呢？不是辨别黑白，而是辨别正义与非正义、善与恶、有益与无益、贞洁与不洁，以便爱前者而避后者；辨别爱德与仇恨，以便常处于前者而不处于后者。\par

22. 这位丈夫尚未在那妇人身上接替前五位丈夫。而在他未接替之处，谬误便当权。因为当灵魂开始能够运用理智时，它要么受智慧心智统治，要么受谬误统治：但谬误不是统治，而是败坏。因此，在那五种感官之后，那妇人仍然游荡，谬误使她摇摆不定。这谬误不是合法的丈夫，而是姘夫：因此主对她说：\Quote{你说得对：我没有丈夫。因为你曾有过五位丈夫。}肉身的五种感官起初支配了你；你到了运用理性的年龄，却还未达到智慧，反而陷入了谬误。因此，在那五位丈夫之后，\Quote{你现在所有的，不是你的丈夫。}如果不是丈夫，那是什么呢？岂不正是姘夫吗？所以，\Quote{叫}——不是叫那姘夫，而是\Quote{你的丈夫}——好叫你用理智接纳我，而不是因谬误对我形成虚假的观念。因为那妇人仍在谬误中，正想着那水，而主此时说的是圣神。她为何出错？岂不是因为她有姘夫，而非丈夫？因此，撇开那败坏你的姘夫，\Quote{去，叫你的丈夫。}叫吧，然后来，好叫你理解我。\par

23. \Quote{那妇人向祂说：\InnerQuote{先生，我看你是位先知。}}丈夫开始到来，但尚未完全到来。她视主为先知，而主确实是先知；因为祂自己说过：\BibleQuote{先知除了在自己的本乡外，是没有不受尊重的。} \BibleRef{路 4:24} 再者，关于祂，梅瑟也曾说：\BibleQuote{我要从他们的兄弟中给他们兴起一位像你一样的先知。} \BibleRef{申 18:18} 像你，即指肉身的形态，而非祂威严的崇高。因此，我们发现主耶稣被称为先知。所以这妇人如今离正确不远了。\Quote{我看，}她说，\Quote{你是位先知。}她开始呼唤丈夫，并摒弃姘夫；她开始询问一件常困扰她的事。因为撒玛黎雅人与犹太人之间有争执，因为犹太人在撒罗满所建的圣殿中敬拜天主；但撒玛黎雅人因地处偏远，不在那里敬拜。因此，犹太人因在圣殿中敬拜天主，便自夸比撒玛黎雅人优越。\Quote{因为犹太人与撒玛黎雅人没有来往：}因为后者对他们说：你们凭什么夸口，自以为比我们优越？只因你们有一座我们没有的圣殿吗？难道我们的祖先，那些悦乐天主的，是在那圣殿中敬拜的吗？他们不是在这座山上，就是我们所在的这里敬拜的吗？那么，他们——撒玛黎雅人——说，我们在这座山上祈祷，就如我们的祖先祈祷过的那样，我们做得更好。两方都因无知而争辩，因为他们没有丈夫；他们因圣殿和山而互相敌视。\par

24. 然而，主现在教导那妇人什么呢？因她的丈夫已经临在了。\Quote{那妇人对祂说：主，我看祢是位先知。我们的祖先在这山上朝拜；而你们却说，人应该朝拜的地方是在耶路撒冷。耶稣对她说：女人，你相信我吧。} 因为教会将要来临，正如\WorkTitle{雅歌}所说：\Quote{她将要来临，并从信德的起始越过。} 她来临是为了穿过；她不能穿过，除非从信德的起始。她的丈夫临在，所以她如今正确地听到：\Quote{女人，你相信我吧。} 因为现在你里面有那能信的东西了，既然你的丈夫临在。当你称我为先知时，你已经开始用理智临在。女人，你相信我吧；因为如果你们不信，你们就不会领悟。因此，\Quote{女人，你相信我吧，因为时候要到，那时你们既不在这山上，也不在耶路撒冷朝拜圣父。你们朝拜你们所不认识的；我们朝拜我们所认识的，因为救恩出自犹太人。但时候将要来到。} 什么时候？\Quote{现在就是了。} 那么，是什么时候？\Quote{当真正朝拜的人以心神和真理朝拜圣父的时候，} 不在这山上，也不在圣殿里，而是以心神和真理。\Quote{因为圣父正寻找这样朝拜祂的人。} 为什么圣父寻找这样朝拜祂的人，不是在山上、不是在圣殿里，而是以心神和真理？\Quote{天主是神。} 如果天主是身体，那么祂应当在山上被朝拜，因为山是物质的；祂应当在圣殿里被朝拜，因为圣殿是物质的。\Quote{天主是神；朝拜祂的人必须以心神和真理朝拜。}\par

25. 我们听到了，并且是明显的；我们曾走出门外，如今被召回到内里。你曾说：但愿我能找到一座又高又孤独的山！因为我想，既然天主在高处，祂从高处会更愿意听我。因为你在一座山上，你就以为自己离天主近，祂会迅速听你，仿佛从最近处呼叫他？祂居住在高处，却垂顾卑微者。\Quote{主是近的。} 对谁？也许对高者？\Quote{对痛悔心碎的人。} 这是奇妙的事：祂居住在高处，却仍然接近卑微者；\Quote{祂垂顾卑微之事，但高傲之事祂从远处知晓；} 祂远远看见骄傲者，他们越是自高，祂就越不接近他们。那么，你曾寻找一座山吗？下来，好使你接近祂。但你愿意上升吗？上升，但不要寻找山。\Quote{登山的路径，}它说，\Quote{是在他心里，在涕泣之谷中。} 谷就是谦卑。因此要一切皆在内里进行。即使你也许寻找某个高处、某个圣地，也要在时间里为天主造一座圣殿。\Quote{因为天主的圣殿是圣的，这圣殿就是你们。}\BibleRef{格前 3:17} 你愿意在圣殿里祈祷吗？在你自己内祈祷。但先成为天主的圣殿，因为祂在祂的圣殿里垂听祈祷者。\par

26. \Quote{时候来到，现在就是了，那时真正朝拜的人将以心神和真理朝拜圣父。我们朝拜我们所认识的；你们朝拜你们所不认识的，因为救恩出自犹太人。} 祂将一件大事归于犹太人；但不要认为祂是指那些虚假的犹太人。要明白那墙，有另一面墙与之连接，它们得以结合，靠在那角石上，就是基督。因为有一面墙来自犹太人，另一面来自外邦人；这些墙远相隔，直到它们在角石上合一。现在那些异乡人是天主盟约的陌生人和外人。\BibleRef{弗 2:11} 据此，经上说：\Quote{我们朝拜我们所认识的。} 确实，这话是就犹太人而言，但并非所有犹太人，不是被摈弃的犹太人，而是像宗徒们、像先知们、像所有那些变卖自己财产、并将财产价银放在宗徒脚前的圣人一样。\Quote{因为天主并没有弃绝祂所预知的人民。}\BibleRef{罗 11:2}\par

27. 那妇人听了这话，就继续说。她已称祂为先知；她观察到与她说话的那位所说的事更加属于先知；她如何回答？请看：\Quote{那妇人对祂说：我知道默西亚将要来到，祂称为基督；当祂来了，祂将告诉我们一切。} 这是什么？刚才她说：犹太人为圣殿争论，我们为这座山争论；当祂来了，祂将轻视这座山，推翻圣殿；祂将教导我们一切，使我们知道如何以心神和真理朝拜。她知道谁能教导她，但她还不认识现在正在教导她的那一位。但现在她配得接受祂的显现。默西亚就是受傅者：受傅者，在希腊文中是基督；在希伯来文中是默西亚；由此，在布匿语中，Messe 意为傅油。因为希伯来语、布匿语和叙利亚语是同源且邻近的语言。\par

28. 于是，\Quote{那妇人对祂说：我知道默西亚将要来到，祂称为基督；当祂来了，祂将告诉我们一切。耶稣对她说：同你讲话的我就是。} 她呼求自己的丈夫；他成为妇人的头，基督成为男人的头。如今这妇人在信德中得以确立，并被治理，好似即将正当地生活。她听了这话：\Quote{同你讲话的我就是。} 她还能说什么呢？当主耶稣愿意向这妇人显示自己时，祂曾对她说：\Quote{你相信我吧？}\par

29. \Quote{门徒们立刻来了，并惊奇祂同那妇人谈话。} 祂是在寻找那失落的她，祂来是为了寻找那失落的：他们对此感到惊奇。他们对一件好事感到惊奇，他们并未猜疑坏事。\Quote{却没有人说：祢寻找什么，或祢为什么同她谈话？}\par

30. \BibleQuote{那妇人遂撇下了她的水罐。} 她听见了\BibleQuote{同你谈话的我就是祂}，将主基督接受到她心中，她除了撇下她的水罐、跑去传布福音之外，还能做什么呢？她抛弃了情欲，急忙去宣告真理。凡愿意传布福音的人，要学习：让她们在井旁丢弃水罐。你们还记得我之前关于水罐所说的话：它是一个打水的器皿，称为\OriginalTerm{水罐}{hydria}，按其希腊名称，因为水在希腊语中是\OriginalTerm{水}{hydor}；正如若从拉丁语称为\OriginalTerm{水器}{aquarium}一样。她当时丢弃了她的水罐，那水罐对她已不再有用，反成了重担，因为她是如此渴求以那水得饱饫。她抛开重担，为使人认识基督，\BibleQuote{她就往城里去，对那些人说：你们来看！有一个人说出了我所做过的一切事。} 她循序渐进，免得那些人发怒愤慨并且迫害她。\BibleQuote{莫非这就是基督？他们就出城，来到祂那里。}\par

31. \BibleQuote{与此同时，祂的门徒恳求祂说：老师，请吃。} 因为他们去买了食物，已经回来了。\BibleQuote{但祂说：我有食物吃，那是你们不知道的。因此门徒彼此说：难道有人给祂拿了什么吃的吗？} 那妇人不明白关于水的事，又何足为奇？看，门徒也还不明白关于食物的事。但祂听见了他们的想法，便像老师一样教导他们，不是绕着圈子，如同祂在寻找那妇人丈夫时那样，而是直截了当：\BibleQuote{我的食物，} 祂说，\BibleQuote{就是承行派遣我者的旨意。} 因此，就那妇人而言，承行派遣祂者的旨意甚至是祂的饮料。正是为此祂曾说：\BibleQuote{我渴了，请给我水喝；} 也就是说，要在她内工作信德，并喝她的信德，将她移植到祂自己的身体内，因为祂的身体就是教会。因此祂说：\BibleQuote{我的食物就是承行派遣我者的旨意。}\par

32. \BibleQuote{你们不是说：还有四个月才到收割吗？} 祂为工作炽热，并筹划派遣工人。你们数着还有四个月到收割；我指示你们另一种庄稼，已经发白，准备就绪。看，我告诉你们：\BibleQuote{举目观看，田里的庄稼已经发白，可以收割了。} 因此祂将要派遣收割者。\BibleQuote{因为这话是真实的：一人收割，另一人播种；好叫播种的和收割的一同欢乐。我派遣你们去收割你们没有劳苦过的；别人劳苦了，而你们进入了他们的劳苦。} 那么怎样呢？祂派遣收割者，难道没有派遣撒种者吗？收割者往哪里去？往别人已经劳苦过的地方。因为哪里已经有了劳苦，那里必定已经有播种；播下的如今已经成熟，需要镰刀和脱粒。那么，收割者将被派往哪里？往先知们先前已经宣讲过的地方，因为他们是撒种者。如果他们不是撒种者，那妇人又从哪里得到这话：\BibleQuote{我知道\OriginalTerm{默西亚}{Messias}将要来}？那妇人如今是成熟的果实，庄稼已经发白，需要镰刀。\BibleQuote{我派遣你们}，然后。往哪里？\BibleQuote{去收割你们没有播种的：别人播种了，而你们进入了他们的劳苦。} 谁劳作了？亚巴郎、依撒格和雅各伯。读他们的劳苦；在他们所有的劳苦中，都有对基督的预言，正因如此，他们是撒种者。梅瑟和所有其他圣祖以及所有先知，在撒种的寒冷季节受了多少苦！因此，在犹太地区庄稼现在已经成熟。恰如那里的谷物成熟，当成千上万的人带来了他们财物的价钱，放在宗徒脚前，卸下了他们世上行囊的负担，开始跟随主基督。确实，庄稼熟了。用它做了什么？从那庄稼中散出了一些谷粒，种遍了整个世界；另一个庄稼正在发芽，将在世界终结时收割。关于那庄稼，经上说：\BibleQuote{流泪撒种的，必欢呼收割。} 但是，对于那庄稼，派遣的不是宗徒，而是天使。祂说：\BibleQuote{收割者，} \BibleQuote{就是天使。} \BibleRef{玛 13:39} 那庄稼于是在莠子中生长，等待在世界终结时被净化。但门徒们最初被派遣去的庄稼（先知们在那里劳作的）已经成熟。然而，弟兄们，请注意所说的话：\BibleQuote{可以一同欢乐，无论是播种的还是收割的。} 他们在时间上劳苦不同，但将要一样平等地享受欢乐；他们要一同领受永生的赏报。\par

33. \BibleQuote{那城里有许多撒玛黎雅人信了祂，因为那妇人作证说：祂说出了我所做的一切事。撒玛黎雅人来到祂那里，求祂在他们那里住下，祂就在那里住了两天。还有更多的人因祂的话而信了，并对那妇人说：现在我们信，不是因为你的话；因为我们亲自听见了，并知道这确实是世界的救主。} 这一点也必须稍加注意，因为这段讲道已近尾声。那妇人首先宣告了祂，撒玛黎雅人相信了她的见证；他们求祂与他们住下，祂就在那里住了两天，更多的人相信了。当他们相信之后，就对那妇人说：\BibleQuote{现在我们信，不是因为你的话；而是我们亲自认识祂，并知道这确实是世界的救主。} 最初是经由传言，然后是由于祂的临在。如今对于那些外教人、尚未成为基督徒的人也是同样。基督由基督徒朋友介绍给他们；正如同因那妇人的报告（即教会的报告），他们来到基督面前，因这报告而相信。祂与他们住了两天，即赐给他们爱德的两条诫命；更多的人相信，而且更坚定地信靠祂，因为祂确实是世界的救主。\par

\SourceRecord{Translated by John Gibb.；Nicene and Post-Nicene Fathers, First Series；Vol. 7.；Edited by Philip Schaff.；Buffalo, NY: Christian Literature Publishing Co.,；1888.；<http://www.newadvent.org/fathers/1701015.htm>.}

% structure: p0001 source=h1 target=section reason=page-title

\section{讲道集第十六篇（\BibleRef{若 4:43-54}）}

1. 今天的福音段落接续昨日的，这正是我们讲道的主题。这段经文的含义虽然不难探究，却值得宣讲，值得赞叹与颂扬。因此，在诵读这段福音时，我们必须请你们留意，而非费力地阐释它。\par

现在，耶稣在\Place{撒玛黎雅}住了两天后，\BibleQuote{往加里肋亚去了，}祂在那里长大。圣史接着说道：\Quote{因为}\BibleQuote{耶稣亲自作证说：先知在自己的本乡不受尊荣。}耶稣之所以在两天后离开撒玛黎雅，并非因为祂在撒玛黎雅不受尊荣；因为撒玛黎雅不是祂的本乡，\Place{加里肋亚}才是。因此，祂那么快就离开撒玛黎雅，来到祂长大的加里肋亚，祂又怎能作证说\BibleQuote{先知在自己的本乡不受尊荣}呢？相反，似乎更可能是，如果祂不屑于去加里肋亚而留在撒玛黎雅，祂或许就作证说先知在自己的本乡不受尊荣了。\par

2. 亲爱的诸位，现在要留意，当主提示并赐予我当说的话时，这里向我们暗示了不小的奥秘。你们知道摆在我们面前的问题；你们去寻找它的解答。但为使解答更令人愿意接受，让我们重复主题。困扰我们的问题是：为何圣史说：\Quote{因为}\BibleQuote{耶稣亲自作证说：先知在自己的本乡不受尊荣。}受此推动，我们回到前文，以发现圣史说这话的意图；我们发现他在叙事的前文中叙述说，两天后耶稣从\Place{撒玛黎雅}往\Place{加里肋亚}去了。那么，圣史啊，难道你正是为此说耶稣作证说先知在自己的本乡不受尊荣，仅仅因为祂两天后离开撒玛黎雅，急忙来到加里肋亚吗？相反，我本会更以为，如果耶稣在自己的本乡不受尊荣，祂就不该急忙前往，而应离开撒玛黎雅。但若我没有弄错，或者更确切地说，因为这是真的，我没有弄错；因为圣史比我更清楚地看见他所说的，比我更看见真理，他是在主的怀中汲取的：因为这位圣史正是\Person{若望}，他在所有门徒中斜靠在主的胸膛上，主虽然爱众人，却更爱他超过其余。那么，难道是他弄错了，而我的意见正确吗？相反，如果我有虔敬之心，让我顺服地聆听他所说的，好使我堪当如同他所想的那样去想。\par

3. 因此，亲爱的诸位，请听我对此事的看法，但这并不妨碍你们自己的判断，若你们已形成更好的见解。因为我们同有一位师傅，在同一所学校里同为同窗。这就是我的意见，请看看我的意见是否真实，或接近真理。他在撒玛黎雅住了两天，撒玛黎雅人便信了他；他在加里肋亚住了许多天，加里肋亚人却仍不信他。请回顾那段落，或回忆昨天的读经和讲道。他来到撒玛黎雅，起初在那里曾被那在雅各伯井旁与他谈论伟大奥秘的妇人宣讲。撒玛黎雅人看见并听见他后，因那妇人的话而信了他，更因他自己的话而更坚定地相信，甚至更多的人都信了：经上就是这样记载的。他在那里过了两天（这天数奥秘地指示了那两条诫命的数目，全部法律和先知都系于其上，正如你们记得我们昨天向你们暗示的），他就往加里肋亚去，来到加里肋亚的加纳城，就是他曾把水变成酒的地方。在那里，当他把水变成酒时，正如若望自己所记载的，他的门徒就信了他；但当然，那屋里挤满了宾客。如此伟大的奇迹行了，然而只有他的门徒信了他。他现在又回到这加里肋亚城。\BibleQuote{看，有一个王臣，他的儿子患病，来到他跟前，恳求他下去}到那城或家中，\BibleQuote{医治他的儿子，因为他快要死了。}那恳求的人难道没有相信吗？你们想听我说什么？请问主对他有何看法。被恳求后，他回答说：\BibleQuote{除非你们看到神迹和奇事，你们总是不信。}他向我们显示了一个不冷不热、或信德冷淡、或毫无信德的人；但他急于通过医治他的儿子来试探基督是何等样人，他是谁，他能做什么。我们确实听见了恳求者的话；我们没有看见那怀疑者的心；但那位既听见言语又看见心的主告诉了我们这件事。简而言之，福音作者本人以其叙述的见证向我们表明，那希望主到他家医治他儿子的人，尚未相信。因为在他被告知他的儿子痊愈了，并发现他的儿子正是在主说\BibleQuote{你回去吧，你的儿子活了}的那个时刻痊愈的，然后经上说\BibleQuote{他自己和他全家就都信了。}那么，如果他和他全家相信的原因，是因为他被告知他的儿子痊愈了，并发现他们告诉他的时刻与基督预告的时刻一致，这就说明当他提出请求时他尚未相信。撒玛黎雅人没有等待任何神迹，他们仅凭主的话就相信了；但他的同乡却理应听到这样的话：\BibleQuote{除非你们看到神迹和奇事，你们总是不信；}而且即便在那里，尽管行了如此伟大的奇迹，也只有\BibleQuote{他自己和他全家}相信了。仅凭他的讲论，许多撒玛黎雅人就相信了；在那奇迹发生的地方，只有那一家相信了。那么，弟兄们，主在这里向我们显示什么呢？当时犹太的加里肋亚是主的家乡，因为他是在那里长大的。但现在这情形预示着什么——因为\Quote{\OriginalTerm{奇兆}{prodigy}}之所以如此称呼，并非没有原因，而是因为它们预示或预兆什么：\Quote{\OriginalTerm{奇兆}{prodigy}}这个词的语源，就好比是\Emph{porrodicium, quod porro dicat}，即宣告未来之事，预示将来之事——现在所有这些情形都预示了某事，预告了某事；我们暂且假定主耶稣基督按血肉的家乡（因为他在世上没有家乡，除非按他在世上所取的血肉而言）；我再说，假定主的家乡是指犹太民族。看，在自己的家乡他没有荣耀。请观察此时此刻犹太人的群众；观察那现今分散在全世、被连根拔起的民族；观察那些被折下、被砍断、散落、枯干的枝条；因这些枝条被折下，野橄榄枝得以接上；看看犹太人的群众：他们甚至现在对我们说什么？\Quote{你们所敬拜和钦崇的那位是我们的同胞。}我们回答说：\BibleQuote{先知在自己的家乡不受尊敬。}简而言之，那些犹太人看见主在地上行走并行了奇迹；他们看见他使瞎子看见、聋子听见、哑巴说话、瘫子康复、在海面上行走、命令风和浪、复活死人：他们看见他行了如此伟大的神迹，然而之后几乎没有人相信。我是对天主的子民说话；我们这么多人相信了，我们见过什么神迹呢？因此，当时所发生的事预兆了现今正在发生的事。犹太人就像加里肋亚人；我们就像那些撒玛黎雅人。我们听了福音，同意了它，借着福音信了基督；我们没有看见任何神迹，也不要求任何神迹。\par

因为，虽是被选和圣洁的十二人之一，但他是个以色列人，属于主的本族，就是那个\Person{多默}，他想要把手指探入伤口的地方。\Person{主}责备他，正如责备这位官长一样。对官长祂说：\BibleQuote{除非你们看见神迹和奇事，总不信。}；对\Person{多默}祂说：\BibleQuote{因为你看见了我，才相信。}祂在\Place{撒玛黎雅人}之后来到\Place{加里肋亚人}那里，\Place{撒玛黎雅人}相信了祂的话，祂在他们面前没有行神迹，很快便放心地离开他们，因为祂的天主性亲临并没有离开他们。现在，当\Person{主}对\Person{多默}说：\BibleQuote{来，伸过你的手，不要无信，但要相信。}，他触摸了伤口的地方，便喊叫说：\BibleQuote{我的主，我的天主。}，他被责备，并对他说：\BibleQuote{因为你看见了我，才相信。}为什么？岂不是\BibleQuote{因为先知在自己的家乡没有荣耀}吗？但既然这位\Person{先知}在陌生人中有荣耀，接下来是什么？\BibleQuote{那些没有看见而相信的，才是有福的。}\BibleRef{若 20:29}我们就是这里预言的那些人；\Person{主}预先赞扬的，祂也屈尊在我们身上实现了。他们看见祂，就是那些钉死祂的人，用手触摸祂，于是少数人相信了；我们既没有看见也没有触摸祂，我们听见并相信了。但愿我们得享祂所应许的福，在此世和来世得以实现：在此世，因为我们被偏爱胜过祂自己的家乡；在来世，因为我们被接上，代替了那些被折断的枝条！\par

因为祂显示祂要折断这些枝条，并接上这\OriginalTerm{野橄榄}{wild olive}，这是因着那\Person{百夫长}的信德而感动，他对祂说：\BibleQuote{主啊，我不配祢到我舍下来，只要祢说一句话，我的小孩就会痊愈。因为我也是一个在权下的人，有兵在我以下；我对这个说：去，他就去；对那个说：来，他就来；对我的仆人说：做这事，他就去做。}\Person{耶稣}转身对跟随祂的人说：\BibleQuote{我实在告诉你们，在以色列我没有见过这么大的信德。}为什么在\Place{以色列}没有见过这么大的信德？\BibleQuote{因为先知在自己的家乡没有荣耀。}\Person{主}难道不能对那位\Person{百夫长}说，正如祂对这位官长所说的：\BibleQuote{去吧，你的儿子活了}吗？请看分别：这位官长渴望\Person{主}下到他家；那位\Person{百夫长}却自称不配。对一人说：\BibleQuote{我去治好他}；对另一人说：\BibleQuote{去吧，你的儿子活了。}对一人祂应许亲临；对另一人祂以一句话治愈。官长强求祂的亲临；\Person{百夫长}自认不配祂的亲临。这里是向高傲让步；那里是向谦卑让步。正如祂对官长说：\BibleQuote{去吧，你的儿子活了}；不要使我疲倦。\BibleQuote{除非你们看见神迹和奇事，总不信。}你渴望我亲临你家中，我能以一言命令；不要因神迹而相信。外邦的\Person{百夫长}相信我能以一言行事，祂还没做就相信了；而你，\BibleQuote{除非你们看见神迹和奇事，总不信。}因此，若是如此，就让它们像骄傲的枝条一样被折断，让谦卑的\OriginalTerm{野橄榄}{wild olive}被接上；然而让根保留，那些被砍下，这些被接上。根保留在哪里？在圣祖们那里。因为\Place{以色列}民族是基督自己的家乡，因为祂按肉身是出自他们；但这棵树的根是\Person{亚巴辣罕}、\Person{依撒格}和\Person{雅各伯}，即圣祖们。他们在哪里？在天主那里安息，享有极大荣耀；以至于那穷人获得提升后，离开身体后，被提升到\Person{亚巴辣罕}的怀中，他被那骄傲的富人在远处看到，是在\Person{亚巴辣罕}的怀中。因此根保留，根被称赞；但骄傲的\OriginalTerm{自然的枝条}{natural branches}理应被折断并枯干；因它们的折断，谦卑的\OriginalTerm{野橄榄}{wild olive}找到了位置。\par

现在听听，\OriginalTerm{自然的枝条}{natural branches}如何被折断，\OriginalTerm{野橄榄}{wild olive}如何被接上，借着这位\Person{百夫长}本人，我认为适宜提及他，以便与这位官长比较。\BibleQuote{我实在告诉你们，在以色列我没有见过这么大的信德。因此我告诉你们：将有许多人从东方和西方来。}\OriginalTerm{野橄榄}{wild olive}多么广泛地占据了大地！这世界曾是一片\OriginalTerm{苦涩的森林}{bitter forest}；但因着谦卑，因着这\InnerQuote{我不配}——将有许多人从东方和西方来。就算他们来了，他们会怎么样？因为他们若来，就被从森林中砍下；他们将被接在哪里，以免枯干？\BibleQuote{他们将坐席，}祂说，\BibleQuote{与亚巴辣罕、依撒格和雅各伯一同。}在什么宴席上？免得你邀请他们不就永生，而是酗酒？\BibleQuote{他们将坐席吗？在天国里。}那些出自\Person{亚巴辣罕}血统的人将会如何？那满树的枝条将会怎样？除了被折断，好让这些被接上，还能怎样？给我们显示他们将被折断：\BibleQuote{但本国的子民将被赶入外方的黑暗。}\BibleRef{玛 8:5}\par

7. 因此，愿那位先知在我们中间受到尊敬，因为祂在自己的家乡没有受到尊敬。祂在祂所形成的地方没有受到尊敬；愿祂在祂所塑造的地方受到尊敬。因为万物的造物主，就是在那地方，以仆人的形象被造成。祂被造成的那座城，熙雍，那个犹太民族，当祂与圣父同在并作为天主的圣言时，正是祂自己所创造的；因为\BibleQuote{万物是藉着祂造成的，凡受造的，没有一样不是藉着祂造成的。}关于那个人，我们今天听到说：\BibleQuote{天主与人之间的唯一中保，就是人基督耶稣。}（\BibleRef{弟前 2:5}）\WorkTitle{圣咏}也预言说：\Quote{人要说：我的母亲是熙雍。}有一个人，就是天主与人之间的中保之人，说：\Quote{我的母亲熙雍。}为什么说\Quote{我的母亲是熙雍}？因为祂从那里取了肉身，从那里有了童贞玛利亚，从她的胎中祂取了仆人的形象；祂屈尊显现得最为卑微。一个人说：\Quote{我的母亲是熙雍}；而这个说\Quote{我的母亲是熙雍}的那个人，是在她内被造成的，在她内成了人。因为祂在她之前是天主，在她内成了人。那在她内成了人的那位，\Quote{亲自奠立了她；至高者在卑微的她内成了人。}因为\BibleQuote{圣言成了血肉，居住在我们中间。}\Quote{祂，至高者，亲自奠立了她。}现在，因为祂奠立了这个地方，愿祂在这里受到尊敬。祂出生的地方拒绝了祂；愿祂所重生的地方接待祂。\par

\SourceRecord{Translated by John Gibb.；Nicene and Post-Nicene Fathers, First Series；Vol. 7.；Edited by Philip Schaff.；Buffalo, NY: Christian Literature Publishing Co.,；1888.；<http://www.newadvent.org/fathers/1701016.htm>.}

% structure: p0001 source=h1 target=section reason=page-title

\section{讲道集第十七篇（\\BibleRef{若 5:1-18}）}

1. 天主显奇迹并不足为奇；若是由人施行，那才是奇迹。我们更应当欢喜，因为我们的主救主耶稣基督降生成人，且在人间施行神圣作为，胜于惊奇。祂为了人而降生成人，这比祂在人中所行的，对我们的救恩更为重要；祂治愈灵魂的缺陷，比治愈必死身体的软弱更为重要。然而，灵魂不认识那要治愈它的那一位，肉身却有眼睛可见有形之事，但内心还没有健全的眼睛去认出隐藏在肉身中的天主，于是祂行了灵魂所能见的事，以治愈那使灵魂不能看见的。\par

祂来到一个地方，那里躺着许多病人——瞎眼的、瘸腿的、枯干的。祂是灵魂和肉体的医师，来治愈一切相信者的灵魂，祂从那些病人中拣选了一个来医治，以此象征合一。若我们以平常心，仅以人的理解和才智来看待这事，就能力而言，祂所行的并非大事；就良善而言，祂也做得太少。那里躺着那么多病人，却只有一个被治愈，而祂本可一句话就叫他们全部起来。那么，我们岂不必须明白：祂的能力和良善正在做那些能让灵魂通过祂的作为而理解永恒救恩的事，胜于让身体获得暂时的健康？因为那真实的身体健康，从主所期待的健康，终将在死人复活时实现。那时活着将不再死；被治愈将不再生病；满足将不再饥渴；更新将不再衰老。然而，当时藉着我们的主救主耶稣基督那些作为而被开的盲眼，后来在死亡中又闭上了；瘫痪者得到力量的肢体，在死亡中又松弛了；凡在必死之肢体上暂时痊愈的，最终都归于无有。但相信的灵魂却进入了永生。因此，对于那将要相信的灵魂——祂来赦免其罪，祂谦卑自己以治愈其疾病——祂藉着治愈这个软弱的病人，给出了一个有意义的证据。关于此事和此证据的深邃奥迹，只要主恩准我们，在你们专心并以祈祷扶助我们的软弱下，我将尽力讲述。但凡我所不能的，那帮助我尽其所能的那一位，必会补足给你们。\par

2. 关于这个周围有五道廊子的池子，里面躺着许多病人，我记得我经常讲论；你们中多数人会和我一起回忆我将要说的事，而非初次听闻。但回顾已知之事决非无益，好使不知者受教，已知者得坚固。因此，作为已知之事，必须简略提及，而非从容灌输。这个池子和这水，在我看来，似乎象征着犹太民族。因为\\BibleBook{若望}{默示录}清楚地向我们指明；当他看见许多水并询问那是什么时，他被告知那是各民族（\\BibleRef{默 17:15}）。那么，那水——即那民族——被梅瑟五书所围，如同被五道廊子所围。然而，那些书卷带出病人，却未治愈他们。因为法律定人的罪，而不赦免罪人。因此，没有恩宠的经文使人成为罪人，而恩宠却释放那认罪的。这正是宗徒所说的话：\\BibleQuote{因为如果有一种法律能够赐予生命，那么义德便确实出于法律。}（\\BibleRef{迦 3:21}）他接着说：\\BibleQuote{但圣经把所有的人都圈在罪中，为使那藉着信德于耶稣基督的恩许，赐予那些相信的人。}（\\BibleRef{迦 3:21}）还有什么比这更明显的？这些话岂不向我们解释了那五道廊子和那许多病人吗？五道廊子就是法律。为什么五道廊子不治愈病人？因为\\BibleQuote{如果有一种法律能够赐予生命，那么义德便确实出于法律。}那么，廊子为何容纳那些它们所不治愈的人呢？因为\\BibleQuote{圣经把所有的人都圈在罪中，为使那藉着信德于耶稣基督的恩许，赐予那些相信的人。}\par

3. 那么，所行的是什么呢？使那些不能在廊子中被治愈的人，在水被搅动之后得以治愈。因为水忽然被搅动，而搅动它的却不可见。你们可以相信，这是由天使的能力所行，但并非没有奥迹的寓意。水被搅动后，那能下去的人就跳进去，而他独自被治愈；此后进去的人，则徒劳无功。那么，这意味着什么？除非是：基督来到犹太民族中间，藉着行大事，教导有益之事，搅动了罪人，祂的临在搅动了水，并激起它走向祂自己的死亡。但搅动的那一位是隐藏的。因为\\BibleQuote{如果他们认识祂，决不会钉死光荣之主}（\\BibleRef{格前 2:8}）。因此，下到那搅动的水里，就是相信主的死亡。那里只有一人被治愈，象征合一；此后来的，不会被治愈，因为凡在合一之外的，都不能被治愈。\par

4. 现在让我们看看祂在那个人身上想要表明什么——正如我先前所说的，祂自己保持合一的奥迹，愿意从众多病患中治愈那人。祂从这人的年数中发现了所谓的软弱年数：\Quote{他患病已有三十八年}。这个数目更多地关联到软弱而非健康，这一点需要更仔细地解释。我希望你们留心；主会帮助我们，使我能够恰当地说话，使你们能够充分地聆听。数字四十被推荐给我们注意，作为一个被某种完美所祝圣的数字。亲爱的，我想这是你们所熟知的。圣经多次见证这一事实。禁食因这个数字而被祝圣，正如你们所知道的。因为梅瑟禁食了四十天，厄里亚也一样；我们的主和救主耶稣基督也亲自完成了禁食的这个数目。梅瑟代表法律，厄里亚代表先知，主代表福音。正是为此，这三人在那座山上显现，祂在那里向门徒们显示自己面容和衣袍的光辉。因为祂出现在梅瑟和厄里亚中间，正如福音从法律和先知那里得到见证。\BibleRef{罗 3:21} 因此，无论是在法律、先知还是福音中，数字四十都在禁食方面被推荐给我们注意。现在，禁食，就其宽泛和一般的意义来说，是戒绝世上的不义和不法的快乐，这是完美的禁食：\BibleQuote{弃绝不虔敬和世俗的情欲，在今世自守、公义、敬虔度日。}宗徒将什么赏报与这禁食相连呢？他接着说：\BibleQuote{等候那有福的盼望，以及我们伟大的天主和救主耶稣基督荣耀的显现。}\BibleRef{铎 2:12} 因此，在今世，我们好像庆祝四十天的斋戒，当我们正当地生活，戒绝不义和不法的快乐时。但由于这斋戒不会没有赏报，我们等候\BibleQuote{那有福的盼望，以及伟大天主和我们的救主耶稣基督荣耀的启示。}在那盼望中，当盼望的实现成为现实时，我们将领受我们的工资，一块\OriginalTerm{德纳}{denarius}。因为那是给在葡萄园中劳作的工人所给的工资，\BibleRef{玛 20:10} 我想你们记得；因为我们不必像对待完全无知和没有经验的人一样重复一切。那么，一块德纳，其名称来自数字\Emph{十}，被给予，这与四十相加成为五十；因此，我们在复活节前以劳苦守\OriginalTerm{四旬期}{Quadragesima}，但在复活节后以喜乐守\OriginalTerm{五旬期}{Quinquagesima}，因为已经领受了我们的工资。现在，对于这个，似乎是对属于数字四十的善工的有益劳苦，加上了安息和福乐的德纳，使它成为数字五十。\par

5. 主耶稣自己更明显地显示了这一点，当祂复活后与门徒在地上同行了四十天；并在第四十天升天，在十天之后派遣了工资，即圣神。这些事是以预象来成就的，并借着某种预象预演了现实本身。我们借着有意义的象征得到滋养，以便能够达到永恒的现实。我们是工人，仍在葡萄园中劳作；当白日结束，工作完成时，工资将被支付。但哪一个工人能坚持到领工资，除非他在劳动时得到滋养？即使你自己，也不会只给工人工资；你岂不也给他一些东西来恢复他劳力中的力量？你确实喂养你将要给工资的人。同样，主也在圣经那些有意义的象征中，在劳动时喂养我们。因为如果我们失去理解神圣奥迹的喜乐，我们就会在劳动中昏厥，并且没有人会得到赏报。\par

6. 那么，工作如何在数字四十中得以成全？原因可能是，因为法律是在十条诫命中赐下的，并要在全世界宣讲：这全世界，我们要注意，由四个部分组成：东、西、南、北；因此数字十乘以四成为四十。或者，可能是因为法律由福音所成全，而福音有四卷书：因为在福音中记载说：\BibleQuote{我不是来废除法律，而是来成全它。}那么，无论是这个原因，还是那个原因，或是其他更可能的原因，虽向我们隐藏，却不向更有学识的人隐藏；然而，确定的是，在数字四十中，表示善工中的某种成全，这些善工最主要的是通过戒绝世上的不法情欲来实行的，即一般意义上的禁食。\par

也听宗徒所说：\BibleQuote{爱是法律的满全。}\BibleRef{罗 10:10} 这爱从何而来？藉天主的恩宠，藉圣神。因为我们无法从自身获得，仿佛靠自己产生。它是天主的恩赐，而且是一份大恩赐：因为他说：\BibleQuote{天主的爱藉着所赐予我们的圣神，已倾注在我们心中。}\BibleRef{罗 5:5} 因此爱成全了法律，最真切地说：\Quote{爱是法律的成全。} 让我们探究这爱，主是如何推荐它让我们思考的。记住我所陈述的：我想解释那无力之人的三十八年，为什么三十八这个数字是软弱的数字，而非健康的数字。现在，如我所说，爱满全了法律。四十这个数字属于法律在所有工作中的成全；但在爱中，两条诫命交给了我们遵守。我恳求你们，记在眼前，铭刻记忆，我所说的话；不要轻视圣言，免得你们的灵魂成为被践踏的路，所撒的种子不能发芽，\Quote{天空的飞鸟会来把它啄去。} 要领悟它，存记在心中。主赐给我们的爱的诫命有两条：\BibleQuote{你当全心、全灵、全意爱上主你的天主；} 以及，\BibleQuote{你当爱近人如你自己。全部法律和先知都系于这两条诫命。}\BibleRef{玛 22:37} 寡妇将\Quote{两个小钱}，她所有的一切，投入天主的奉献中，这是有理由的。主人取\Quote{两个}银钱，为那被强盗打伤的人使他痊愈，这是有理由的。耶稣与撒玛黎雅人共度两天，以坚固他们于爱中，这是有理由的。因此，虽然一般的某种好事由数字二象征，但尤其向我们展示了爱的双重特性。因此，如果四十这个数字拥有法律的成全，而法律只在爱的双重诫命中得以满全，那么你们为什么惊异那差二就满四十的人竟软弱患病呢？\par

7. 因此，现在我们来看这无力之人被主治愈的神圣奥迹。主亲自来了，他是爱的导师，充满爱，正如所预言的，\Quote{缩短}了\Quote{在地上的言语}，并显示了法律和先知都系于爱的两条诫命。\Person{梅瑟}和他的四十，\Person{厄里亚}和他的四十，都系于此；主在此见证中也引用了这个数字。这无力之人被主亲自治愈；但在治愈他之前，主对他说什么？\Quote{你愿意痊愈吗？} 那人回答说没有人把他放入水池。他确实需要一个\Quote{人}来治愈他，但那\Quote{人}也是天主。\BibleQuote{因为天主只有一个，在天主与人之间的中保也只一个，就是那人基督耶稣。}\BibleRef{弟前 2:5} 那么，所需的那人来了：为何治疗要延迟？他说：\Quote{起来}；\Quote{拿起你的床，行走。} 他说了三点：\Quote{起来，拿起你的床，行走。} 但那个\Quote{起来}不是要作工的诫命，而是治愈的行动。那人在痊愈后，接受了两个诫命：\Quote{拿起你的床，行走。} 我问你们，为什么说\Quote{行走？}不够，或者说\Quote{起来？}也就够了？因为那人痊愈起来后，不会留在原地。他起来不就是为了离开吗？我的印象是，那发现这人缺乏两件事的主，给了他这两条诫命：因为藉着命令他做两件事，仿佛弥补了所缺乏的。\par

8. 那么，我们如何在主的这两条命令中找到爱德的两条诫命？祂说：\Quote{拿起你的床，行走吧。}弟兄们，请与我一同忆起这两条诫命是什么。因为你们应当完全熟悉它们，不仅在我们诵念时想起，更应永不从心中抹去。愿这永远是你们至上的思想：你们必须爱天主和你们的近人：\Quote{全心、全灵、全意爱你的天主，并爱近人如你自己。}这些必须时时默想、沉思、持守、实践并完成。在享乐的次序中，爱天主优先；但在行动的次序中，爱近人居先。因为那以两条诫命命令你这爱德的那一位，并没有吩咐你先爱近人再爱天主，而是先爱天主，然后爱你的近人。然而你，既然你还未看见天主，便借着爱近人而挣得看见祂；借着爱近人，你净化了眼睛以看见天主，正如若望明确所说的：\Quote{如果你不爱你所看见的弟兄，怎么能爱你所看不见的天主呢？}\BibleRef{若一 4:20}看，你被告知：\Quote{爱天主。}如果你对我说：\Quote{请把祂指给我看，好让我爱祂；}我将回答什么呢？只能以同一位若望所说的：\Quote{从来没有人见过天主。}并且，为使你不至于认为自己完全与看见天主隔绝，他说：\Quote{天主是爱；那住在爱内的，就住在天主内。}\BibleRef{若一 4:16}因此，爱你的近人；注视你爱近人的泉源；在那里，你将尽你所能看见天主。那么，开始爱你的近人。\Quote{把你食物分给饥饿的人，将无家可归的穷人领到你的屋里；你若看见赤身露体的，给他衣服；不要轻视你血亲中的人。}而这样做之后，其结果你将得到什么？\Quote{那时，你的光明将如晨光般升起。}\BibleRef{依 58:7}你的光明就是你的天主，是为你\Quote{晨光}，因为祂将在世界的黑夜之后来到你这里：因祂不升不落，永远长存。当你回归时，祂将是你的晨光；当你背离祂时，祂曾为你落下。因此，在我看来，祂在这\Quote{拿起你的床}中，似乎说了：爱你的近人。\par

9. 但是，为何以拿起床来表达爱近人，仍然隐藏着，并且依我看，需要加以阐释：除非，或许，以床——一把迟钝无感的物件——来指称我们的近人会冒犯我们。愿我的近人不要因被一个无灵魂无感觉的东西来象征而生气。主自己，我们的救主耶稣基督，被称为\OriginalTerm{角石}{corner-stone}，为在自己内建造两者。祂也被称为\OriginalTerm{磐石}{rock}，从中流出水来：\Quote{那磐石就是基督。}\BibleRef{格前 10:4}那么，如果基督被称为磐石，近人被称为木头，有什么稀奇呢？然而，并非任何木头；正如那磐石也非任何磐石，而是流出水来给口渴者饮的磐石；也非任何石头，而是角石，它本身将来自不同方向的两堵墙结合起来。同样，你也不可把近人当作任何木头，而是一张床。那么，请问，床有什么意义呢？什么意义？不就是软弱的人曾躺在上面，而当他痊愈后，便扛起床吗？宗徒怎么说？\Quote{你们要彼此担待重担，这样就满了全基督的法律。}\BibleRef{迦 6:2}基督的法律就是爱德，而爱德除非我们彼此担待重担，就不能满全。他说：\Quote{彼此以爱德含忍，尽力保持圣神的合一，在和平的联系中。}\BibleRef{弗 4:2}当你软弱时，你的近人担待了你；你痊愈了，你就担待你的近人。这样，人啊，你将补足你所欠缺的。\Quote{那么，拿起你的床。}但当你拿起它时，不要停留在原地；\Quote{行走罢。}借着爱你的近人，关心你的近人，你完成了你的行进。你往哪里去？岂不是往主天主那里去？我们应以全心、全灵、全意爱祂。因为我们尚未到达主那里，但我们的近人与我们同在。那么，当你行走时，担待他，好使你能来到你渴望永居的那一位那里。因此，\Quote{拿起你的床，行走罢。}\par

10. 那人这样做了，犹太人便气愤。因为他们看见一个人在安息日扛着床，却没有责怪主在安息日治好他，以便祂能回答他们：如果你们中有一人有一头牲畜掉进井里，难道不在安息日把它拉上来，救回牲畜吗？因此，如今他们并非反对祂在安息日治好一个人，而是反对那人扛着床。但如果治愈不可推迟，难道也该命令劳作吗？他们说：\Quote{你不可做你所做的事，} \Quote{拿起你的床。}他辩护时，把治愈他的那一位放在谴责者面前，说：\Quote{那使我痊愈的，祂对我说：\InnerQuote{拿起你的床，行走吧。}}难道我不应从那使我得治愈的那一位接受命令吗？他们却说：\Quote{对你说\InnerQuote{拿起你的床，行走吧}的那人是谁？}\par

11. \Quote{但那痊愈的人不知道是谁}对他说了这话。\Quote{因为耶稣}做完这事，给了他这个命令后，\Quote{就从人群中转身离开了。}请看这也是如何应验的。我们背负邻人，走向天主；但我们走向的那一位，我们尚未看见：正因如此，那人还不认识耶稣。其中向我们暗示的奥秘是：我们相信那尚未看见的祂；而为了不被看见，祂在人群中侧身而去。在人群中看见基督是困难的；我们的心灵需要某种独处；正是通过某种静观的独处，才能看见天主。人群有喧闹；这种看见需要隐秘。\Quote{拿起你的床榻}——你被背负，也要背负邻人；\Quote{然后行走}，好能抵达终点。不要在人群众中寻找基督：祂并非群众中的一位；祂超越一切群众。那大鱼首先从海中升起，祂坐在天上，为我们转求；如同大司祭独自进入了帐幔内；群众站在外面。你要行走，背负你的邻人：你若学会了背负，你——原是惯于被背负的。一言以蔽之，即使现在你还不认识耶稣，尚未看见耶稣：此后又如何？由于那人没有停止拿起床榻行走，\Quote{后来耶稣在圣殿中遇见了他。}他在人群中未曾看见耶稣，却在圣殿中看见了祂。主耶稣固然既在人群中也在圣殿中看见了他；但瘫子不认识人群中的耶稣，却在圣殿中认识了祂。于是那人来到主前：在圣殿中看见了祂，在祝圣过的、神圣的地方看见了祂。主又对他说什么？\Quote{看，你已经痊愈了；不要再犯罪，免得遭遇更坏的事。}\par

12. 于是，那人看见耶稣，认出祂是自己痊愈的根源后，便不懒惰地宣讲他所看见的那一位：\Quote{他就离开，告诉犹太人，使他痊愈的是耶稣。}他向他们传报，他们却向他发怒；他宣讲自己的救恩，他们却不寻求自己的救恩。\par

13. 犹太人迫害主耶稣，因为祂在安息日做了这些事。让我们听听主现在对犹太人所作的答复。我已经告诉你们，祂通常如何回答关于在安息日医治人的问题：他们不会在安息日忽视自己的牲畜，无论是解救它们还是喂养它们。关于搬动床榻，祂又怎样回答？一件明显的身体工作在犹太人眼前完成了；不是身体的治愈，而是一项身体的工作，这工作看起来不如治愈那样必要。那么，让主公开宣告：安息日的圣事，即守一日的记号的标记，是暂时赐给犹太人的，但这圣事的应验已在祂自己内来到。\Quote{我父}，祂说，\Quote{一直工作到现在，我也工作。}这话在他们中间引起了巨大骚动：水因主的来临而被搅动，但搅动者却未被看见。然而有一位大病的人要被这搅动的水治愈，整个世界要被主的死亡治愈。\par

14. 那么，让我们看看真理所作的答复：\Quote{我父一直工作到现在，我也工作。}那么，圣经所说\Quote{天主在第七日歇了祂所做的一切工作}是假的吗？主耶稣自己曾对犹太人说：\Quote{如果你们相信梅瑟，你们也就会相信我，因为他写的是关于我。}祂现在说这话，难道是与梅瑟所传的圣经相矛盾吗？那么，请看看，梅瑟是不是有意用\Quote{天主在第七日安息了}来象征什么。因为天主在创造自己的工作时并没有疲惫，需要像人一样休息。祂怎么会疲惫呢？祂凭借话语创造了万有。然而，两者都是真的：\Quote{天主在第七日歇了祂所做的一切工作}，以及耶稣所说的\Quote{我父一直工作到现在}。但谁能用言语向人阐明呢？人是软弱的，向软弱的人讲话；无学问的，向寻求学问的人讲解；即使他偶然理解了一些，他也难以把它表达出来并向人们阐明——人们或许勉强接受，即使所接受的能被阐明。我说，我的弟兄们，谁能用言语阐明天主如何在安息中工作，又在工作中安息？我请求你们，在你们行路时，把这事暂且搁下；因为这种看见需要天主的圣殿，需要圣地。背负你们的邻人，行走吧。你们要在那不需要人言的地方看见祂。\par

15. 或许我们可以更恰当地说，这句话\Quote{天主在第七天休息了}是以一个伟大的奥迹预示了主、我们的救主耶稣基督自己，祂曾说：\Quote{我父至今工作，我也工作。}因为主耶稣当然是天主。祂是天主圣言，你们已经听到\Quote{在起初已有圣言}；且不是任何言语，而是\Quote{圣言就是天主，万物是藉着祂造成的}。祂或许被预示为要在第七天从祂的一切工作中休息。因为，请读福音，看看耶稣行了何等伟大的工程。祂在十字架上成就了我们的救恩，使先知所预言的一切在祂身上应验。祂被戴上荆棘冠冕，挂在木头上，说：\Quote{我渴了}，接收了蘸满醋的海绵，以应验经上的话：\Quote{在我渴时，他们拿醋给我喝。}当祂的一切工程完成后，在一周的第六天，祂低下了头，交付了灵魂，并在安息日那天在坟墓里安息了。因此，祂仿佛对犹太人说：\Quote{你们为什么指望我不在安息日工作？安息日是为你们设立的，作为我的预像。你们观察天主的工程：我曾在它们被造成时在场，万物都是藉着我造成的，我知道它们。\InnerQuote{我父至今工作}。圣父造了光，但祂说话就有了光；如果祂说话，那是藉着祂的圣言造的：我就是祂的圣言，我曾是，我现在仍是；那世界在我之内被造于那些工程中，世界在我之内被统治于这些工程中。我父在创造世界时工作，至今在统治世界时仍工作：因此，祂在创造时藉着我创造，在统治时藉着我统治。}祂说这话，但对谁说的？对聋、瞎、瘸、瘫的人，他们不承认医生，仿佛在疯狂中失去了理智，想要杀祂。\par

16. 此外，福音作者继续说了什么？\Quote{因此犹太人越发想要杀害祂，因为祂不但破坏了安息日，并且又称天主是自己的父}——不是以普通方式，而是如何？\Quote{使自己与天主平等。}因为我们都对天主说：\Quote{我们在天上的父}；我们也读到犹太人说过：\Quote{你实在是我们的父亲。}\BibleRef{依63:16} 所以，他们发怒不是因为祂说天主是祂的父，而是因为祂以与常人截然不同的方式说。看，犹太人明白亚略派所不明白的。事实上，亚略派说圣子不与父平等，因此这异端被逐出教会。看，那些盲目的、杀害基督的人，仍然明白基督的话。他们不明白祂是基督，也不明白祂是天主子；但他们确实明白，在这话中暗示了一位与天主平等的天主子。他们不知道祂是谁，但他们仍然承认这样一位被宣告，因为\Quote{祂说天主是自己的父，使自己与天主平等。}难道祂不与天主平等吗？祂并没有使自己平等，而是父生了祂使之平等。如果祂使自己平等，那就成了强夺。因为那想使自己与天主平等、而其实不是的，就堕落了，从天使变成了魔鬼，\BibleRef{依14:14} 并向人递上那使他自身倾覆的骄傲之杯。这堕落者因嫉妒人站立得住，对人说：\Quote{尝一尝，你们就会像天主一样}\BibleRef{创3:5} ——就是说，你们藉着篡夺来攫取你们所没有变成的，因为我也因强夺而堕落了。他没有提出这点，而是劝诱这点。然而，基督是由父所生而与父平等，不是受造的；是由父的性体所生。因此宗徒如此宣告祂：\Quote{祂虽具有天主的形象，却不以自己与天主同等为强夺的。}\Quote{不以自己与天主同等为强夺}是什么意思？祂没有篡夺与天主的平等，而是处在祂被生的那种平等中。那么我们如何能到达那平等的天主呢？\Quote{祂空虚自己，取了奴仆的形象。}\BibleRef{斐2:6} 但祂空虚自己，不是由于失去祂所是的，而是由于取了祂所不是的。犹太人轻视这奴仆的形象，不能理解主基督与父平等，尽管他们对祂自己所肯定的这一点毫不怀疑，因此他们发怒；然而祂仍然容忍他们，并在他们攻击祂时寻求医治他们。\par

\SourceRecord{Translated by John Gibb.；Nicene and Post-Nicene Fathers, First Series；Vol. 7.；Edited by Philip Schaff.；Buffalo, NY: Christian Literature Publishing Co.,；1888.；<http://www.newadvent.org/fathers/1701017.htm>.}

% structure: p0001 source=h1 target=section reason=page-title

\section{第十八篇（若望福音 5:19）}

圣史\Person{若望}，在他的同伴与其他福音作者中间，从主领受了这份特殊而独特的恩赐——在筵席上他斜倚在主胸前，以此表明他正从主最深的心怀汲取更深奥的奥秘——为要讲论关于天主圣子的事，这些事或许能唤起幼小心灵的注意，却尚不能充满他们，因为他们还不能领受；而对那些成长稍大、达到内在成人一定程度的心灵，他藉这些话赐予他们一些既有操练又有滋养的东西。你们在诵读时已经听到，并且记得这篇讲论是如何引起的。因为昨天读到：\BibleQuote{因此犹太人便图谋杀害耶稣，因为祂不但破坏了安息日，而且还称天主为自己的父，使自己与天主平等。}\BibleRef{若 5:18}这令犹太人不悦，却令圣父喜悦。无疑，这也令那些如同尊敬圣父一样尊敬圣子的人喜悦；因为若它不令他们喜悦，他们便不会蒙喜悦。因为天主不会因你喜悦而变得更伟大，但你若是不喜悦，你便会变得更渺小。如今，针对他们的这种诽谤，或出于无知，或出于恶意，主所说的话完全不是他们所能理解的，而是能令他们激动和困扰，并在困扰时，也许寻求那位医生。祂说出了应当写下来的话，以便日后甚至我们也能读到。如今我们已经看到当犹太人听到这些话时心中发生了什么；当我们自己听到时，让我们更仔细地思考会发生什么。因为异端以及某些乖僻的主张，引诱灵魂并将它们抛入深渊，无非是在圣经不被正确理解，且其中未被正确理解的部分被鲁莽而大胆地断言之时才产生的。所以，亲爱的诸位，我们应当极其谨慎地聆听那些我们尚且幼小而不易理解的事，并且要以虔诚的心和战栗，正如经上记载的，持守这健全的准则：在我们能够理解的事上，按照我们所领受的信德，如同食物般欢欣；而尚未能理解的，我们应搁置疑惑，暂时推迟理解；也就是说，即使我们尚不知道它是什么，我们仍丝毫不怀疑它是善且真的。至于我，弟兄们，你们必须考虑我是谁，竟敢向你们讲话，以及我承担了什么：因为我承担了探讨属神的事，而我是一个属人的人；探讨属灵的事，而我是属血肉的；探讨永恒的事，而我是必朽的。但愿我也能，亲爱的诸位，远离虚妄的骄傲，倘若我的言论在天主家中是健全的，\BibleQuote{这原是天主的教会，真理的柱石和基础。}\BibleRef{弟前 3:15}我按自己的尺度取出我所放置于你们面前的：哪里敞开着，我与你们一同看见；哪里关闭着，我与你们一同叩门。\par

现在，犹太人被激怒并愤慨：的确，是合理的，因为一个人竟敢使自己与天主平等；但在这点上是不合理的，因为在这人身上他们不认识天主。他们看见了血肉，却不知道天主；他们观察了居所，却不认识居住者。那血肉是一座圣殿，天主居住在它内。使血肉与圣父平等的不是耶稣，将仆人的形体与主比较的不是祂；不是祂为我们所成为的，而是祂造了我们时所是的。因为基督是谁（我向天主教徒说），你们知道，因为你们正确相信了；不是圣言独自，也不是血肉独自，而是圣言成了血肉，居住在我们中间。我再提及你们所知道的关于圣言的：\BibleQuote{在起初已有圣言，圣言与天主同在，圣言就是天主。}\BibleRef{若 1:1}这里是与圣父的平等。但\BibleQuote{圣言成了血肉，居住在我们中间。}\BibleRef{若 1:14}比这血肉，圣父是更大的。所以，圣父既平等又更大；与圣言平等，比血肉更大；与藉着祂造我们的那位平等，比那位为我们而被造的那位更大。凭着这健全的公教准则——你们应当特别知道，你们知道的人应持守它，你们的信德绝不可从它滑落，任何人的辩论都不能从你们心中夺走它——让我们衡量我们所能理解的事；而那些我们或许还不理解的事，让我们推迟，日后当我们有能力时再以此准则衡量。那么，我们知道祂，与圣父平等，是天主圣子，因为我们知道祂在起初是作为天主的圣言。那么，犹太人为何要杀害祂？\BibleQuote{因为祂不但破坏了安息日，而且还称天主为自己的父，使自己与天主平等。}\BibleRef{若 5:18}看见了血肉，没有看见圣言。所以，让祂藉着血肉发言反对他们，让内在的居住者藉着祂的居所发言，使凡能看见的，就知道谁住在里面。\par

3. 那么祂对他们说了什么？\BibleQuote{那时耶稣回答他们说：}因祂使自己与天主平等而愤怒，\BibleQuote{我实实在在告诉你们：子凭自己不能做什么，只有看见父做什么，才能做。}犹太人如何回答这些话，没有记载；也许他们什么也没说。然而，有些希望被视为基督徒的人却并不沉默，反而从这些话中生出与我们相悖的某些意见，这些意见不可轻忽，既为他们也为我们。即是说，亚略异端者虽然断言那取了肉身的子，不是因肉身而是取肉身之前就比父小，并且不与父同质，却从这些话中获取曲解的借口，回答我们说：\Quote{你们看，主耶稣见犹太人因他使自己与天主圣父平等而激动愤慨，便附加了这样的话，以表明他并不与天主平等。因为犹太人，}他们说：\Quote{就对基督发怒，因为他使自己与天主平等；而基督愿纠正他们这种印象，并表明子不与父，即不与天主平等，便这样说，仿佛他说：你们为什么生气？为什么愤慨？我并不与天主平等，因为\InnerQuote{子凭自己不能做什么，除非看见父做什么}。如今，}他们说：\Quote{那\InnerQuote{凭自己不能做什么，除非看见父做什么}的，当然更小，而非平等。}\par

4. 在这扭曲和败坏的己心准则中，让异端者听我们说，并非斥责，而是仍如询问一般，让他向我们解释他所想的。因为，我猜想，无论你是谁（因为我们可视为他亲临此），你与我们同持\BibleQuote{在起初已有圣言}这信念。他说：我持守。并且\BibleQuote{圣言与天主同在}？他说：这我也持守。那么继续，并持守那更强的后续断言，即\BibleQuote{圣言就是天主}。他说：连这我也持守；然而，这一个是较大的天主，那一个是较小的天主。这颇有异教之风：我以为我是在与基督徒交谈。若有较大的天主和较小的天主，那么我们所敬拜的就是两位天主，而非一位。他说：为什么？难道你们不也肯定两位天主，彼此平等吗？我并未肯定这一点；因为我将此平等理解为其中蕴含不可分割的爱；若爱不可分割，则有完全的合一。因为若天主置于人心中的爱能使众人之心成为一心，能使众灵魂成为一魂，正如\BibleBook{宗徒}{大事录}所记载的信主并彼此相爱的人那样：\BibleQuote{他们一心一意归向天主：}\BibleRef{宗 4:32}因此，若我的灵魂和你的灵魂因同思同恋而成为一个灵魂，那么天主圣父和天主圣子之爱泉中岂不更是独一天主？\par

5. 但对于这些使你心乱的话，请收回你的思绪，与我一起反省那关于圣言我们所寻求的。我们已持守\BibleQuote{圣言就是天主}；我再加上另一件事，即当福音作者说\BibleQuote{这在起初与天主同在}后，立刻接着说\BibleQuote{万有是藉着祂造成的}。现在我要以诘问来催逼你，我要搬动你反对你自己，并控告你自己：只须记住关于圣言的这一点，即\BibleQuote{圣言就是天主，万有是藉着祂造成的}。现在听那些使你断言子更小的话，诚然，因为祂说\BibleQuote{子凭自己不能做什么，只有看见父做什么，才能做}。他说：正是如此。请向我稍加解释：我想，你就是这样想的：父做某些事，子观察父怎样做，以便祂自己也能做那些祂看见父做的事。你似乎设立了两位工匠：父和子犹如师傅和学徒，正如手艺人的父亲惯常教儿子他们的手艺。看，我降格到你的肉性理解：此刻我像你一样思考，让我们看看这种想法是否与我们刚才一同述说、一同持守的关于圣言的事相协调，即\BibleQuote{圣言就是天主}，以及\BibleQuote{万有是藉着祂造成的}。那么，假设父作为工匠做某些工作，子作为学徒，祂\BibleQuote{凭自己不能做什么，只有看见父做什么，才能做}：祂仿佛急切地注视着父的手，以便看见父塑造任何东西，就同样凭自己的作品塑造类似的东西。但此处父所做的一切，祂都希望子注意祂，并且自己也照样作；父是藉着谁做的呢？来吧！现在是时候让你坚持你起初的见解了，那见解你曾与我一同背诵并持守：\BibleQuote{在起初已有圣言，圣言与天主同在，圣言就是天主，万有是藉着祂造成的}。但你，在与我一起持守万有是藉着圣言造成的之后，却又以你肉性的才智和幼稚的幻想，想象天主在做某事，而圣言在注意，以便当天主做完后，圣言也能做类似的事。那么，天主做什么事能不用圣言？因为若祂做任何事，那么万有就不是藉着圣言造成的；你放弃了你所持守的立场。但若万有是藉着圣言造成的，就纠正你的误解。父行了一切，而且只藉着圣言行：圣言如何注意看父在没有圣言的情况下做事，以便圣言能照样做？凡父所做的，祂都是藉着圣言做的；否则，\BibleQuote{万有是藉着祂造成的}就是假的。但\BibleQuote{万有是藉着祂造成的}是真的。也许这在你看来还不够？好吧，\BibleQuote{而没有祂，什么也没有造成。}\par

6. 那么，从这肉性的智慧中退避吧，让我们探讨这句话是怎会说的：\Quote{圣子什么也不能凭自己作，唯有看见圣父所作的，才能作。}让我们探讨，若我们堪当领悟。因为我承认这是一件大事，且全然困难：看见圣父藉着圣子作事：并非圣父与圣子各自作祂们个别的工作，而是圣父藉着圣子作祂所作的一切工作；以至于没有任何工作是圣父不藉着圣子作的，也没有任何工作是圣子不藉着圣父作的，因为\BibleQuote{万物是藉着祂造成的，凡受造的，没有一样不是由祂造成的。}这些真理极其坚固地建立在信德的根基上，那么，这\Quote{看见}的本性是什么呢？你寻求，我想，认识圣子作事：先寻求认识圣子看见。因为祂究竟说了什么？祂说：\Quote{圣子什么也不能凭自己作，唯有看见圣父所作的，才能作。}注意祂所说的：\Quote{唯有看见圣父所作的。}看见在先，作事在后：祂看见为了作事。至于你，你为何现在寻求知道祂如何作事，而你还未懂得祂如何看见？你为何奔向那后来的，而留下那先来的？祂声明自己看见并作事，而非作事并看见；因为\Quote{祂什么也不能凭自己作，唯有看见圣父所作的。}你要我向你解释祂如何作事吗？向我解释祂如何看见。你若不能解释这个，我也不能解释那个。你若尚未有能力理解这个，我也不能理解那个。因此，让我们各人寻求，各人叩门，好使各人堪当领受。为何你，仿佛你是有学问的，不公平地责怪我，一个无学问的人？我在作事方面，你在看见方面，彼此都是无学问的，让我们询问这位导师，不是在学校里幼稚地争吵。然而，我们已一同学会\BibleQuote{万物是藉着祂造成的。}因此显然，并非圣父作一类不同的工作，圣子看见了才作类似的其他工作；而是圣父藉着圣子作完全相同的工作，因为万物是藉着圣言造成的。至于天主如何作事，谁知道呢？祂如何造了，我不说世界，而是你自己的眼睛——你因肉性的依恋而用这眼睛把有形事物与无形事物相比较？因为你把那些你用这肉眼惯常看见的东西，想成天主的模样。但若能用这肉眼看见天主，祂就不会说：\BibleQuote{心里洁净的人是有福的，因为他们要看见天主。}因此，你有一个身体的眼去看工匠，但你还没有一个心的眼去看天主；于是你惯常在一个工匠身上看见的，你便转移到天主身上。把地上的事物留在地上；把你的心高举起来。\par

7. 那么，亲爱的，我们将解释我们所问的：圣言如何看见，圣言怎样看见圣父，圣言的看见是什么？我不敢，也不冒昧，为我自己或为你承诺解释这个；虽然我估计你的程度，但我仍知道自己的。因此，若你愿意，不要再拖延，让我们浏览这段经文，看看属血肉的心如何因主的话而困扰；困扰，好叫它们不继续固守它们所把持的。让这从它们手中夺去，如同从孩童手中夺去一件玩具，他们用它自娱自乐而伤害自己，使他们如同长大的人，可以有更有益的东西栽种在他们里面，并能进步，而不是在地上爬行。起来，寻求，叹息，渴望，并叩那关闭的门。但若我们还不渴望，还不热切寻求，还不叹息，我们只会把珍珠丢给所有不加分辨的人，或自己找到珍珠却不顾它的种类。因此，亲爱的，我愿在你们心中激起渴望。好的品格导向正确的理解：一种生活引向另一种生活。一种生活是地上的，另一种是天上的；有兽性的生活，有人性的生活，有天使的生活。兽性的生活被地上的享乐所激动，寻求唯独地上的享乐，并以无度的欲望匍匐其后；天使的生活是唯独天上的；人的生活介于天使与野兽之间。若人按照肉性生活，他就与野兽同等；若他按照圣神生活，他就加入天使的共融。当你按照圣神生活时，甚至在天使的生活中检查你是小还是长大。因为若你仍是个小孩，天使们向你说：\Quote{长大吧：我们吃饼；你被奶滋养，以信德的奶，好使你达到见面的干粮。}但若仍有对污秽享乐的渴望，若思想仍充满欺诈，若谎言不被避免，若伪誓加诸谎言，这样污秽的心怎敢说：\Quote{向我解释圣言如何看见}——即使我有能力，即使我现在自己看见？此外，虽然我本人或许不是这样的品格，我仍远远不及这个看见，那么那个被地上的欲望压住的人，还未被来自天上的这种渴望所攫住，岂不更沉重！厌恶与渴望之间有巨大的差别；再者，渴望与享受之间也有差别。若你像野兽那样生活，你厌恶；天使们则完全享受。反之，若你不像野兽那样生活，你不再厌恶：你渴望某物，却没有得到；你因这渴望本身，已开始了天使的生活。愿它在你们内成长，并在你们内得以成全；愿你们领受这个，不是从我，而是从那造了我又造了你们的天主！\par

8. 然而\Person{主}也没有把我们遗弃于机遇，因为祂说：\BibleQuote{子不能由自己作什么，只有看见父作什么，才作什么}，祂要我们明白：\Person{圣父}不是先作一些子可以看见的工程，然后子再作另一些工程；而是\Person{圣父}与\Person{圣子}作完全相同的工程。因为祂接着说：\BibleQuote{凡父所作的，子也照样作}。不是父作了工程以后，子再照样作别的工程；而是\BibleQuote{凡父所作的，子也照样作}。如果子所作的正是父所作的，那么父是藉着子而作；如果父藉着子作祂所作的，那么父不作一些工程、子另作一些工程；而是父与子的工程是相同的工程。儿子如何也作相同的工程？既是\BibleQuote{相同的}，又是\BibleQuote{照样}。免得你以为它们是相同的，但方式不同，祂说\BibleQuote{相同的}，又说\BibleQuote{照样}。它们怎能是相同的却又不是照样的？举一个我以为对你们不算太难理解的例子：当我们书写字母时，首先由我们的心形成，然后由我们的手形成。确实如此：否则你们为何都同意，正是因为你们察觉到了？正如我所说的，这对我们众人都是明摆着的。字母首先由我们的心做成，然后由我们的身体做成；手服务，心命令；心和手都做出相同的字母。你们以为心作一些字母，手作另一些吗？手的确作相同的字母，但不是照样：我们的心理性地形成它们，而我们的手有形可见地形成它们。你看，相同的东西被做出，但不是照样。因此，主说\BibleQuote{凡父所作的，子也作}还不够，还必须加上\BibleQuote{也照样作}。因为如果你理解为：正如你理解你的心作什么，你的手也作什么，只是方式不同？然而在这里，祂加上了\BibleQuote{子也照样作}。如果祂既作这些工程，又照样作，那么就醒寤吧！让犹太人被粉碎，让基督徒相信，让异端者被说服：\Person{圣子}与\Person{圣父}是平等的。\par

9. \BibleQuote{因为父爱子，把自已所作的一切事指示给祂}。这就是那个\BibleQuote{指示}。\BibleQuote{指示}，好像是对谁呢？当然是向那看见的。我们又回到那无法解释的事：\Person{圣言}如何看见。看，人是\Person{圣言}造的；但人有眼睛、耳朵、手、身体的各种器官：他能用眼睛看，用耳朵听，用手工作；器官各不相同，作用各不相同。一个器官不能做另一个器官的事；然而，由于身体的统一，眼睛既为自己也为耳朵看，耳朵既为自己也为眼睛听。我们能否设想，类似的情形也存在于\Person{圣言}中，因为万物都是藉着祂造的？圣经在\BibleBook{}{圣咏}中说：\BibleQuote{你们愚昧的人，省悟吧！你们蠢笨的人，终要明白！那造耳朵的，难道自己听不见吗？那造眼睛的，难道自己看不见吗？}因此，如果\Person{圣言}是造眼睛的那一位，因为万物都是藉着\Person{圣言}造的；如果\Person{圣言}是造耳朵的那一位，因为万物都是藉着\Person{圣言}造的：我们就不能说\Person{圣言}听不见，\Person{圣言}看不见；免得\BibleBook{}{圣咏}责备我们说：\BibleQuote{蠢笨的人，终要明白！}因此，如果\Person{圣言}听见和看见，如果\Person{圣子}听见和看见，我们还需在祂身上分别寻找眼睛和耳朵吗？祂是凭一部分听，凭另一部分看吗？祂的耳朵不能做祂眼睛所做的吗？祂的眼睛不能做祂耳朵能做的吗？难道祂不全是视觉、全是听觉吗？也许是的；不，不是也许，而是确实是的；然而，祂的看见和祂的听见，在方式上远非我们所能想象。在\Person{圣言}中，看见和听见同时存在：在祂内，看见和听见不是不同的东西；而是听见就是看见，看见就是听见。\par

10. 我们以眼见，以耳闻，又怎能知道这事？我们也许回到自己，如果我们不是那些罪人，他们曾被告知：\BibleQuote{罪人们，回到你们心里去吧！}\BibleRef{依 46:8} 回到你们心里：为何离开你们自己，又自取灭亡？为何走上孤寂之路？你们徘徊游荡而迷失；回来吧。到何处？到主那里。这很快就能做到：首先回到你们自己的心；你们流浪在外，成了你们自己的放逐者；你们不认识你们自己，却还在问是谁创造了你们！回来，回到你们心里，将你们自己从身体中提升出来：你们的身体是你们的居所；你们的心甚至通过你们的身体来感知。但你们的身体不是你们的心；甚至离开你们的身体，回到你们心里。在你们的身体中，你们发现眼睛在一个地方，耳朵在另一个地方：你们在心里找到这个吗？或者你们心里没有耳朵吗？否则主为何说：\BibleQuote{有耳听的，就听吧！}\BibleRef{路 8:8} 或者你们心里没有眼睛吗？否则宗徒为何说：\BibleQuote{你们心灵的眼目被照亮？}\BibleRef{弗 1:18} 回到你们心里；在那里看看，或许你们能感知到一点关于天主的事，因为天主的肖像就在其中。基督居于内在的人中，在内在的人中，你们按天主的肖像得到更新，在祂自己的肖像中认出它的创造者。看身体的全部感官如何将知觉带给内在的心，报告它们在外界所感知的；看内在的统帅有多少仆役，以及即使没有这些仆役，它自身能做什么。眼睛向心报告黑白之物；耳朵向同一颗心报告悦耳与刺耳的声音；鼻孔向同一颗心宣告香气与臭气；味觉向同一颗心宣告苦与甜；触觉向同一颗心宣告光滑与粗糙；而心向自己宣告正义与不义之事。你们的心看见、听见并判断感官感知到的一切其他事物；并且，对于感官所不及的，它辨别正义与不义之事、恶与善之事。向我指出你们心的眼睛、耳朵、鼻孔。被传递到你们心的事物各不相同，但在那里却没有不同的肢体。在你们的肉体中，你们在一个地方听，在另一个地方看；在你们心里，你们在哪里看，就在那里听。如果这是肖像，那拥有肖像的祂何其更有能力！因此，圣子既听又看；圣子本身就是听和看：对于祂，听与存在是同一回事；\Quote{存在；}看与存在也是同一回事。\Quote{存在。} 对于你们，看与存在不是同一回事；因为如果你们失去视力，你们仍能存在；如果你们失去听力，你们仍能存在。\par

11. 我们以为我们已经敲了门吗？在我们内是否升起什么，使我们甚至能略微猜测光芒可能从哪里来？弟兄们，我认为，当我们谈论这些事并默想它们时，我们是在操练自己。当我们操练自己，却仿佛被自己的重量弹回到习惯的思想中时，我们就如同视力弱的人，当他们被带到光前来看时，若他们先前完全没有视力，而通过医师的殷勤照料开始以某种方式恢复视力。当医师要测试恢复的进展时，他试图向他们展示他们渴望看见、但在瞎眼时不能看见的东西；当视力稍许恢复后，他们被带到光前；当他们看到光时，却仿佛被那强光击退；当医师指出对象时，他们回答：\Quote{刚才我看见了，但现在不能了。} 那么医师做什么？他将他们带回惯常的方式，并施用药膏以滋养他们对那曾瞬间看见之物的渴望，以便藉着这渴望能更彻底地治愈他们；若为了恢复健康而使用任何刺痛的药膏，让病人勇敢地忍受，并因着对光的爱而振奋，对自己说：\Quote{何时我才能用强健的眼睛看见那有病弱眼睛时不能看见的呢？} 他催促医师，恳求他医治自己。因此，弟兄们，如果或许类似的事已在你们心中发生，如果你们曾举起心去看圣言，却被它的光击退，落回了惯常的方式；那么，祈求医师施用猛烈的药膏，即正义的规诫。有你们可以看见的东西，但没有你们借以看见的能力。你们先前不信我，以为没有你们可以看见的东西：现在，你们如同被理性的引导带到它面前：你们靠近了，睁大眼睛去看它，感到悸动，却又退缩回去了。你们确知有你们可以看见的东西，但你们尚未准备好看见它。因此，接受治愈吧。什么是眼药？不说谎，不发虚誓，不奸淫，不偷盗，不诈骗。但你们习惯了这些，且带着某种痛苦才从旧习惯中被拉出来：这正是刺痛之处，然而它却治愈。因为我坦率地告诉你们，出于对我和对你们的恐惧，如果你们放弃治愈，并藐视成为可享受这光的准备，因你们眼目的软弱，你们将爱黑暗；因爱黑暗，你们将留在黑暗中；因留在黑暗中，你们将被投到外面的黑暗里：那里将有哀号和切齿。如果对光的爱在你们身上没有产生效果，就让对痛苦的恐惧产生一些效果吧。\par

12. 我以为已经说得足够多了，却还没有讲完福音的章节；如果我继续说明余下的部分，会加重你们的负担，而且我担心连已经灌输的也可能丢失；因此，亲爱的诸位，现在这对你们就足够了。我们是负有义务的人，不是现在，而是终生如此，因为我们为你们而活。然而，你们要以善生安慰我们在这世上如此软弱、劳苦且充满危险的今生；不要以恶行折磨和消耗我们。因为，如果因你们的恶行而受冒犯，我们逃离你们、与你们分离、不再到你们那里去，你们岂不抱怨说：即使我们病了，你们或许会照顾我们；即使我们软弱，你们或许会来探望我们？请看，我们确实照顾你们；请看，我们确实探望你们；但不要让我们遇到你们从宗徒那里所听到的情况：\Quote{我怕我对你们是白白劳碌了。} \BibleRef{迦 4:11}\par

\SourceRecord{Translated by John Gibb.；Nicene and Post-Nicene Fathers, First Series；Vol. 7.；Edited by Philip Schaff.；Buffalo, NY: Christian Literature Publishing Co.,；1888.；<http://www.newadvent.org/fathers/1701018.htm>.}

% structure: p0001 source=h1 target=section reason=page-title

\section{第19篇讲道（\BibleRef{若 5:19-30}）}

在前一次讲道中，就所论及的主题以及我们贫乏的理解所能达到的程度，我们曾因福音上的话而谈论过：\Quote{子不能由自己做什么，除非他看见父做什么}——对于子——即\OriginalTerm{圣言}{Word}，因为子就是圣言——来说，\Quote{看见}是什么意思；以及万物既然都是藉着圣言造成的，那么如何理解子先看见父行事，然后自己也做所看见的事，因为父除非藉着子，什么都不做？因为万物都是藉着他造的，没有他，什么也没有受造。然而，我们并没有向你们交付任何完全解释清楚的东西，那是因为我们自己也没有任何清楚阐明的理解。因为，的确，甚至在理解得以进展的地方，言语有时也会匮乏；在理解毫无完美之处的地方，言语更是多么不足！因此，现在，如同主所赐予我们的，让我们简要地过一遍这段经文，并在今天完成应尽的职分。如果或许还有一点时间或力量剩余，我们将（尽我们和你们所能）重新思考圣言\Quote{看见}和\Quote{被显示}是什么意思；因为，事实上，这里所说的一切，如果按照人的感觉、肉性地理解，充满虚妄幻想的灵魂只会为我们制造出圣父和圣子的某些形象，就好像两个男人，一个在显示，另一个在观看；一个在说话，另一个在听——所有这些都只是心中的偶像。如果现在偶像终于从它们自己的庙宇中被推倒了，那么它们更该从基督徒的心中推倒！\par

他说：\Quote{子不能由自己做什么，除非他看见父做什么。}这是真的：要持守这一点，同时不要忘记你们在福音开头所领受的，即\Quote{在起初已有圣言，圣言与天主同在，圣言就是天主}，尤其是\Quote{万物都是藉着他造的。}要联合你们刚才所听到的和那时所听到的，并使两者在你们心中和谐一致。因此，\Quote{子不能由自己做什么，除非他看见父做什么}，这样的说法仍然意味着：父所做的，他只藉着子去做，因为子是他的圣言：并且，\Quote{在起初已有圣言，圣言与天主同在，圣言就是天主}；还有，\Quote{万物都是藉着他造的。}因为他所做的，无论什么，子也同样做；不是别的事，而是这些事；并且不是以不同的方式，而是以同样的方式。\par

\Quote{因为父爱子，凡自己所做的都显示给他。}似乎与上面所说的\Quote{除非他看见父做什么}相关，也与这个\Quote{凡自己所做的都显示给他}相关。但如果父显示他所做的，而子不能做除非父已经显示，并且父不能显示除非他已经做了，那么就会得出结论：父并非通过子做一切；此外，如果我们坚定地持守父通过子做一切，那么父在做事之前就先显示给子。因为如果父在做完之后才显示给子，好让子去做那些已经显示（已经完成）的事，那么毫无疑问，有些事是父不藉着子而做的。但父不藉着子做什么，因为天主子是天主的圣言，万物都是藉着他造的。因此，剩下的可能是：父将要做的事，他显示为将要做的事，好由子去完成。因为如果子做那些父已经显示为已完成的事，那么父显示的那些事一定不是藉着子做的。因为这些事除非先完成，就不能显示给子，而子除非先看到显示，也不能做它们；因此它们是在没有子的情况下完成的。然而，\Quote{万物都是藉着他造的；}这句话是真实的；因此，它们是在被造\Emph{之前}就被显示了。但我们说过，这要先搁置，待简要过一遍经文之后再回来看，如果我们，如我们所说，还有一些时间和力量剩余，可以重新考虑那些被延迟的问题。\par

现在请注意一个更广泛、更困难的问题。他说：\Quote{还要把比这些更大的工程显示给他，叫你们惊奇。}\Quote{比这些更大的。}比哪些更大？回答很快出现：比你们刚才听到的身体疾病的治愈更大。因为整个这段讲论的起因是关于那个患病三十八年的人，他被基督的话治愈；关于这个治愈，主可以说：\Quote{还要把比这些更大的工程显示给他，叫你们惊奇。}因为有更大的，父将把它们显示给子。这不是\Quote{已经显示}（表示过去），而是\Quote{将要显示}（表示未来），或即将显示。又出现一个困难问题：那么，为什么还有某些东西在父那里，尚未显示给子呢？当他讲这些话时，是否还有某些东西在父那里是子所不知道的？因为，如果它是\Quote{将要显示}，也就是说\Quote{即将显示}，那么他还没有显示；而且他将在同一时间显示给子和这些人，因为下文是\Quote{叫你们惊奇。}这是一件很难理解的事：永恒的父如何在时间中显示某些东西给与他同永恒的子，而子本来就知道父那里的一切。\par

5. 但什么是更大的工作呢？或许这是容易理解的。\Quote{因为就如\Person{圣父}，}\Person{主}说，\Quote{复活死人并赐予他们生命，同样，\Person{圣子}也随祂的意愿赐予人生命。}那么，复活死人比医治病人是更大的工作。但\Quote{就如\Person{圣父}复活死人并赐予他们生命，同样，\Person{圣子}也随祂的意愿赐予人生命。}因此，\Person{圣父}复活一些人，\Person{圣子}复活另一些人？但万有都是藉祂而成：因此\Person{圣子}所作的就是\Person{圣父}所作的众人；因为\Person{圣子}不作其他不同方式的事，而是\Quote{这些}并以\Quote{相同的方式。}因此必须清楚理解并持守。但要记住：\Quote{\Person{圣子}随祂的意愿赐予人生命。}在这里，也要知道\Person{圣子}的能力，还有祂的意愿。\Person{圣子}赐予生命给祂所愿意的人，\Person{圣父}也赐予生命给祂所愿意的人——\Person{圣子}与\Person{圣父}所赐予的是同样的人；因此\Person{圣父}和\Person{圣子}的能力是相同的，意愿也是相同的。那么接着呢？\Quote{因为\Person{圣父}不审判任何人，但将全部审判交给了\Person{圣子}，为叫众人尊敬\Person{圣子}，如同尊敬\Person{圣父}一样：}\Person{主}加上了这话，作为前面句子的理由。一个重大的问题摆在我们面前；请专心注意。\Person{圣子}赐予生命给祂所愿意的人，\Person{圣父}赐予生命给祂所愿意的人；\Person{圣子}复活死人，正如\Person{圣父}复活死人。并且还有，\Quote{\Person{圣父}不审判任何人。}如果在审判中死人必须复活，怎能说\Person{圣父}复活死人，如果祂不审判任何人，因为\Quote{祂将全部审判交给了\Person{圣子}}呢？但在那审判中死人复活；有的复活得生命，有的复活受惩罚。如果\Person{圣子}做这一切，而\Person{圣父}不做，因为\Quote{祂不审判任何人，但将全部审判交给了\Person{圣子}，}这似乎与先前所说的相矛盾，即\Quote{就如\Person{圣父}复活死人并赐予他们生命，同样，\Person{圣子}也随祂的意愿赐予人生命。}因此，\Person{圣父}和\Person{圣子}一起复活；如果他们一起复活，他们就一起赐予生命：所以他们也一起审判。那么，怎能说\Quote{因为\Person{圣父}不审判任何人，但将全部审判交给了\Person{圣子}}这句话是真的呢？同时，让现在提出的这些问题吸引你们的心思；\Person{主}会使它们得到解答时令你们喜悦。因为事实如此，弟兄们：每一个问题，除非它激发心思思考，否则在解释时就不会带来喜悦。愿\Person{主}亲自与我们同行，或许祂会在祂所隐藏的那些事中稍微启示自己。因为祂将祂的光隐藏在云彩中；要像鹰一样飞越覆盖整个大地的每一片朦胧雾气，并在\Person{主}的话语中见到最清澈的光，是困难的。因此，或许祂会用祂光线的热量驱散我们的黑暗，并屈尊在后续中稍微启示自己，让我们推迟这些问题，看看后面的话。\par

6. \Quote{凡不尊敬\Person{圣子}的，也不尊敬那派遣祂的\Person{圣父}。}这是真理，且是清楚的。那么，既然\Quote{祂将全部审判交给了\Person{圣子}，}如上文祂所说的，\Quote{为叫众人尊敬\Person{圣子}，如同尊敬\Person{圣父}一样，}如果有些人对\Person{圣父}尊敬而对\Person{圣子}不尊敬，会怎样呢？这不可能，祂说：\Quote{凡不尊敬\Person{圣子}的，也不尊敬那派遣祂的\Person{圣父}。}因此，人不能说，我尊敬了\Person{圣父}，因为我不知道\Person{圣子}。如果你尚未尊敬\Person{圣子}，你也没有尊敬\Person{圣父}。因为尊敬\Emph{\Person{圣父}}是什么意思呢？岂不正是因为祂有一个儿子吗？当你被教导要尊敬\Person{天主}，因为祂是\Person{天主}时，这是一回事；但当你被教导要尊敬祂，因为祂是\Emph{\Person{圣父}}时，则是另一回事。当你被教导要尊敬祂，因为祂是\Person{天主}时，你被引导视祂为造物主、全能者、至高、永恒、不可见、不变的神灵；但当你被教导要尊敬祂，因为祂是\Emph{\Person{圣父}}时，这等同于尊敬\Person{圣子}；因为如果没有\Person{圣子}，就不能被称为\Emph{\Person{圣父}}，同样，如果没有\Person{圣父}，也不能被称为\Emph{\Person{圣子}}。但恐怕你也许确实尊敬\Person{圣父}为较大者，而尊敬\Person{圣子}为较小者——正如你可能对我说：\Quote{我确实尊敬\Person{圣父}，因为我知道祂有一个儿子；我在\Emph{\Person{圣父}}这个称呼上也没有错误，因为我不将\Person{圣父}理解为没有\Person{圣子}，然而我也尊敬\Person{圣子}，但作为较小者，}——\Person{圣子}亲自纠正你，并召叫你，说：\Quote{为叫众人尊敬\Person{圣子}，}不是较低程度地，而是\Quote{如同他们尊敬\Person{圣父}一样。}因此，\Quote{凡不尊敬\Person{圣子}的，也不尊敬那派遣祂的\Person{圣父}。}\Quote{我，}你说，\Quote{希望给\Person{圣父}更大的尊敬，给\Person{圣子}较小的尊敬。}你既给了\Person{圣子}较小的尊敬，便是夺去了\Person{圣父}的尊敬。因为存这种心思，你必然认为\Person{圣父}不愿意或不能够生一个与自己相等的\Person{圣子}：如果祂不愿意，祂就缺乏意志；如果祂不能够，祂就缺乏能力。因此，难道你看不出，存这种心思，你在哪里想给\Person{圣父}更大的尊敬，就在哪里侮辱了\Person{圣父}吗？因此，要尊敬\Person{圣子}如同尊敬\Person{圣父}一样，如果你愿意同时尊敬\Person{圣父}和\Person{圣子}。\par

7. \Quote{我实实在在告诉你们：凡听我的话，而信派遣我来者的，就有永生，不受审判，而是已通过，}不是\Emph{正在通过}，而是已经通过了，\Quote{从死亡进入生命。}且注意这一点：\Quote{凡听我的话，而}——他不是说信我，而是——\Quote{信派遣我来者。}让他听圣子的话，好使他能信圣父。为什么听见你的话，却信了别人？当我们听见任何人的话，难道不是信那说话的人吗？难道不是把自己的信德借给那说话的人吗？那么，他说：\Quote{凡听我的话，而信派遣我来者，}如果不是因为\Quote{他的话在我内}，又是什么意思呢？\Quote{听我的话}不就是\Quote{听我}吗？同样，\Quote{信派遣我来者}，因为信他，就是信他的话；但信他的话，也就是信我，因为我是圣父的圣言。因此圣经是和谐的，一切安排妥当，毫不冲突。所以，抛弃你心中的争论，了解圣经的一致吧。你以为真理会说自相矛盾的话吗？\par

8. \Quote{凡听我的话，而信派遣我来者的，就有永生，不受审判，而是已出死入生。}你们记得我们前面所确立的：\Quote{就如圣父唤醒死者，使他们活，同样圣子也随意使人活。}他已经开始启示自己；看，甚至现在，死者正在复活。因为\Quote{凡听我的话，而信派遣我来者的，就有永生，不会进入审判。}证明他已经复活了。\Quote{而已通过}，他说，\Quote{从死亡进入生命。}那从死亡进入生命的，毫无疑问已经复活了。因为他不可能从死亡进入生命，除非他先是在死亡中而不是在生命中；但当他通过后，他将在生命中，而不是在死亡中。他因此曾经死了，现在又活了；他曾迷失，又被找到。\BibleRef{路 15:32}因此，现在确实发生复活，人们从死亡进入生命；从不信之死到信德之生；从虚伪之死到真理之生；从不义之死到正义之生。因此，确实有一种死者的复活。\par

9. 愿他更圆满地开启它，并开始光照我们！\Quote{我实实在在告诉你们：时刻将到，且现在就是。}我们确实期待末日的死者复活，因为我们如此相信；是的，我们不仅期待，而且明确有义务期待：因为我们相信死者会在末日复活，这不是虚假的事。因此，主耶稣愿意向我们启示一种死者的复活在死者的复活之前，那不同于拉匝禄（\BibleRef{若 11:43}）、寡妇的儿子（\BibleRef{路 6:14}）和会堂长女儿（\BibleRef{玛 5:41}）的复活，因为他们复活后又死了（因为在他们身上发生了一种死者的复活在死者的复活之前）；但是，正如他在这里所说，\Quote{有永生，不受审判，而已出死入生。}进入什么生命？进入永生。那么，不像拉匝禄的身体：他确实从坟墓之死进入人的生命，但并非进入永生，因为他还要死去；而将在世界末日复活的死者，将进入永生。因此，我们的主耶稣基督，我们天上的导师，圣父的圣言和真理，愿意向我们显示一种死者的复活通向永生，在死者的复活通向永生之前，他说：\Quote{时刻将到。}毫无疑问，你们，浸润着对肉身复活的信德，期待世界末日的时刻，为了不让你们在这里期待，他补充说：\Quote{且现在就是。}因此，他所说的\Quote{时刻将到}，不是指最后的那时刻，那时\BibleQuote{在总领天使的号令和天主的号角声中，主亲自从天降下，在基督内死去的人将首先复活；然后我们这些活着存留的人将和他们一起被提到云中，在空中迎接基督：这样我们就常与主同在。}\BibleRef{得前 4:15}那时刻将要来到，但并非现在。但考虑这是什么时刻：\Quote{时刻将到，且现在就是。}在那时刻发生什么？除了死者的复活，还有什么？是什么样的复活？是那种他们复活后永生的复活。这也将在最后的时刻发生。\par

10. 那么，我们如何理解这两种复活呢？我们是否可能理解为，那些现在复活的人将来不再复活；有些人的复活是现在，有些人的复活是将来？并非如此。因为如果我们正确相信了，我们就已经在这复活中复活了；而我们自己，既已复活，还在期待末日的另一次复活。此外，我们现在都已复活而进入永生，只要我们恒心持守同一信德；那时，我们也要复活而进入永生，当我们与天使平等的时候。\BibleRef{路 20:36} 但是，愿他自己分辨并开启我们所敢说的话：如何会有在复活之前的复活，不是不同的人，而是相同的人；而且不像拉匝禄的复活，而是进入永生。他将清楚地开启。请听导师，他正像曙光临于我们，如同我们的太阳滑入我们的心中；不是肉身的眼所渴望看到的那种，而是心灵的眼热切渴望得以开启的那位。那么，让我们倾听他说：\Quote{实实在在，我告诉你们：时候要来，现在就是了，那时死人}——你们看到是在宣告复活——\Quote{要听见天主子的声音；听见的人就要活着。} 他为什么加上\Quote{听见的人就要活着}？难道他们若不活着能听见吗？那么，只要说\Quote{时候要来，现在就是了，那时死人要听见天主子的声音}就够了。我们会立刻明白他们是活的，因为他们若不活着就不能听见。不，他说，不是因为他们活着而听见，而是因听见而重获生命：\Quote{要听见，而听见的人就要活着。} 那么，\Quote{要听见}是什么意思，难道不是\Quote{要听从}吗？就肉耳的听觉而言，并非所有听见的人都活着。事实上，许多听见却不相信；因听见而不相信，他们不听从；因不听从，他们就不活着。所以这里，那些\Quote{要听见}的人就是那些\Quote{要听从}的人。那么，听从的人就要活着：让他们确信无疑，\Emph{要活着}。基督，天主的圣言，被传报给我们；天主子，万物藉着他而造成，他为了救恩的计划，确实取了肉身，由童贞女所生，在肉身中是婴孩，在肉身中是青年，在肉身中受苦，在肉身中死亡，在肉身中复活，在肉身中升天，应许肉身的复活，也應許心神的復活——心神先于肉身，肉身在心神之后。凡听见并听从的，必要活着；凡听见却不听从的，即听见而轻视、听见而不信的，必不活着。为什么不活着？因为他没有听见。\Quote{没有听见}是什么意思？就是不听从。因此，\Quote{听见的人就要活着}。\par

11. 现在请将你们的思绪转向我们说过必须暂缓讨论的事，好让现在，如果可能，得以开启。关于这同一复活，他立即接续说：\Quote{因为父怎样在自己内有生命，同样也赐给子在自己内有生命。} 那句\Quote{父在自己内有生命}是什么意思？他没有在其他地方拥有生命，只在自身内。事实上，他的生命就在他内，不是从别处来，也不是从另一个人得来。他并不，好比说，借用生命，也不好比说，分有生命，分有一种不是他自己本身的生命；而是\Quote{在自己内有生命}，以至生命本身对他而言就是他自己。如果我能更进一步，或多或少地谈论这件事，通过提出例子来启迪你们的理解，这将有赖于天主的助佑和你们专注的虔诚。天主活着，灵魂也活着；但天主的生命是不变的，灵魂的生命是可变的。在天主内既无增加也无减少；但他始终在自身内是同样的，永远是他是的那一位：不是现在这样，以后另一样，先前又另一样。但灵魂的生命极其多样：它曾愚蠢地活着，它现在智慧地活着；它曾不义地活着，它现在义地活着；时而记忆，时而遗忘；时而学习，时而不能学习；时而失去所学的，时而领会所失去的。灵魂的生命是可变的。当灵魂在不义中活着时，那就是它的死亡；当它再次成为义的，它就分有了另一种生命，那不是它自己的生命，因为它上升走向天主，并紧紧依恋天主，得以从祂成义。因为经上说：\Quote{那信靠使不虔敬者成义之主的，他的信德就被视为义。} \BibleRef{罗 4:5} 因离弃天主，它变为不义；因归向祂，它成为义。难道你们不觉得它有点像某种冷的东西，靠近火就变暖，离开火就变冷？某种暗的东西，靠近光就变亮，离开光就变暗？灵魂就是这样；天主却不是这样。此外，人可能会说他现在眼睛里有光。那么，让你们的眼睛来说，如果它们能够，用它们自己的声音说：\Quote{我们自己在内有光。} 我回答：你们说你们在自身内有光是不正确的：你们有光，但在天上；你们有光，但在月亮上，在蜡烛里，如果碰巧是黑夜，而不是在你们自己内；因为闭眼时，你们失去了睁眼时所感知的。你们在自身内没有光；如果你们能在太阳落山时还保留光，那是黑夜，享受黑夜的光；如果蜡烛被拿走时还能保留光，但既然蜡烛被拿走时你们仍留在黑暗中，你们在自身内就没有光。因此，在自身内有光就是不需要从别处来的光。看，谁理解他在说\Quote{父怎样在自己内有生命，同样也赐给子在自己内有生命}时所表明的子与父相等，就能看到在父与子之间只有这个区别：父在自己内有生命，无人赐给祂；而子在自己内有生命，是由父赐给的。\par

12. 但这里又出现一片需驱散的云雾。我们不要灰心，要奋力前行。此处是心灵的牧场；我们不可轻看，好能生存。看，你说：你自己承认，圣父赐给了圣子生命，使圣子在自己内拥有生命，如同圣父在自己内拥有生命一样；圣父不欠缺，圣子也不欠缺；圣父是生命，圣子也是生命；两者合一是一个生命，而非两个生命；因为天主是唯一的，而非两位神；而此\Emph{拥有生命}即是成为生命。那么，为何说圣父\Quote{赐给}圣子生命呢？并非圣子先前没有生命，从圣父领受生命才能生存；若是这样，他就不是在自己内拥有生命。看，我是在谈论灵魂。灵魂存在；即使它不智慧、不正义、不虔诚，它仍是灵魂。它存在为一回事，智慧、正义、虔诚为另一回事。因此，在某些方面它尚未智慧、尚未正义、尚未虔诚。然而它并非虚无，也非无生命；因为它以其自身的某些行动显示自己是活的，尽管它并未显示自己智慧、虔诚或正义。因为若它没有生命，就不会推动身体，不会命令脚行走、手工作、眼观看、耳听闻；不会开口说话，也不会移动舌头区分言语。因此，通过这些行动，它显示自己有生命，且是优于身体的某种东西。但通过这些行动，它是否在任何方面显示自己是智慧、虔诚、正义的呢？愚昧、邪恶、不义之人不也行走、工作、观看、听闻、说话吗？但当灵魂上升到超越自身、高于自身、且其存在所由来的那一位时，它便获得智慧、正义、圣德；此前它缺乏这些，就是死的，没有使自身生活的生命，只有使身体生活的生命。因为灵魂中使身体生活的是一回事，使灵魂自己生活的则是另一回事。灵魂当然优于身体，但天主优于灵魂本身。灵魂，即使愚昧、不虔、不义，仍是身体的生命。但由于它的生命是天主，正如它存在于身体内时，赋予身体活力、优美、活动及肢体的功能；同样，当天主——它的生命——存在于灵魂内时，便赋予灵魂智慧、虔敬、正义、爱德。因此，灵魂赋予身体的，与天主赋予灵魂的，属于不同种类：灵魂赋予生命，也被赋予生命。它赋予生命时自己可能是死的，即使它本身未被赋予生命。但当圣言来临，倾注于听众，他们不仅听闻，而且服从时，灵魂便从死亡中复活，进入自己的生命——即从不义、愚昧、不虔，转向它的天主，那就是它的智慧、正义、光明。让它转向祂，并被祂光照。\Quote{你们前来近，}他说，\Quote{到祂那里去。}我们要得到什么？\Quote{你们必要光明。}因此，如果因\Quote{前来近}而被光照，因\Quote{远离}而变昏暗，那么你们的光不在你们自己内，而在你们的天主内。来近祂，好能复活；若远离祂，你们就要死亡。若因来近祂而生活，因远离祂而死亡，那么你们的生命不在你们自己内。因为你们的生命就是你们的光。\BibleQuote{因为在你那里有生命的泉源，因你的光明我们才能看见光明。}\par

13. 那么，灵魂在蒙受光照之前是一回事，蒙受光照之后，因分享更美者而变得更好，但并非如此——天主的圣言、天主的圣子与祂在获得生命之前有所不同，以致祂是因分享而拥有生命；不，祂在自己内拥有生命，因此祂本身就是生命本身。那么，祂说\Quote{赐给圣子在自己内有生命}是什么意思呢？我简短地说：祂\Emph{生了}圣子。因为祂并非先无生命而后\Emph{获得}生命，而是因被生而为生命。圣父不是因被生而是生命；圣子因被生而是生命。圣父没有父亲；圣子出自天主圣父。圣父的存在不由任何而来，但祂之为父，是因为圣子。同样，圣子之为子，是因为圣父：祂的存在来自圣父。因此祂说了这话：\Quote{赐给圣子生命，使祂在自己内有生命}。就好像祂在说：\Quote{圣父，祂在自己内是生命，生了圣子，使祂在自己内是生命}。确实，祂愿意人们将这个赐给(\Emph{\OriginalTerm{赐给}{dedit (has given)}})理解为与生(\Emph{\OriginalTerm{生}{genuit (has begotten)}})相同。这就好比我们对某人说：\Quote{天主赐给你存在}。对谁？如果是对某个已经存在的人，那么祂就没有赐给他存在，因为能够接受者在被赐予之前已经存在。因此，当你听到\Quote{祂赐给你存在}时，你并不存在来接受，而是你接受了存在，由此你才得以存在。建筑者赐给这房屋存在。但祂赐给它什么呢？祂赐它成为房屋。赐给谁？赐给这房屋。赐给什么？成为房屋。祂怎能赐给房屋使之成为房屋呢？因为如果房屋已经存在，那么当房屋已经存在时，祂赐给它成为房屋是什么意思呢？那么\Quote{赐它成为房屋}是什么意思呢？意思是祂成就了它成为房屋。那么，祂赐给圣子什么呢？赐祂成为圣子，生了祂成为生命——也就是说，\Quote{赐给祂在自己内有生命}，使祂成为不需要生命的生命，以免祂被理解为因分享而拥有生命。因为如果祂因分享而拥有生命，祂可能会因失去而丧失生命。关于圣子，不要接受、不要想、也不要相信这是可能的。因此，圣父持续生命，圣子持续生命：圣父，在自己内的生命，不是来自圣子；圣子，在自己内的生命，但来自圣父。由圣父所生，使祂在自己内生活；但圣父不是被生的，在自己内的生命。祂所生的圣子并不小于祂，以致需要靠成长而相等。因为那创造时间的完美者，当然不会借助时间来达到自己的完美。在一切时间之前，祂与圣父同永恒。因为\Emph{圣父}从未没有圣子；但圣父是永恒的，因此圣子也是同永恒的。灵魂，你怎么样？你曾是死的，失去了生命；那么通过圣子听从圣父吧。起来，接受生命，在那位在自己内有生命者内，你可以接受不在你内的生命。那么赐给你生命的是圣父和圣子；第一次复活完成，当你起来分享你自己所不是的生命，并因分享而得以生活。从你的死亡起来，到你的生命，即你的天主，并从死亡过渡到永生。因为圣父在自己内有永生；除非祂生了如此一位在自己内有生命的圣子，就不能如圣父复活死人并使他们活一样，圣子也随祂意愿使人活。\par

14. 然而，身体的复活又如何呢？因为那些听见而生活的人，除非借着听见，从何处得生活？因为\BibleQuote{新郎的朋友站着听他，因新郎的声音而大大欢喜}\BibleRef{若 3:29}，不是因他自己的声音；也就是说，他们因分享而听见并生活，而非因进入存在；凡听见的都生活，因为凡服从的都生活。主啊，也请告诉我们一些关于肉身复活的事；因为曾有人否认它，声称唯有这个因信德而成就的复活才是复活。主刚才提到了这个复活，并燃起了我们的渴望，因为\Quote{死人将听见天主子的声音，并将生活}。并不是\Emph{有些}听见的人将生活，\Emph{另一些}将死亡；而是\Quote{凡听见的都生活}，因为凡服从的都生活。看，我们看见了一个精神的复活；因此我们不要放弃对肉身复活的信仰。主耶稣啊，除非你向我们宣告此事，我们又能以谁反对那些主张相反的人呢？因为确实，所有试图在人身上嫁接某种宗教的派别都允许这种精神的复活；否则，可能会对他们说：如果灵魂不复活，你为什么对我说话？你想在我里面做什么？如果你不将较差的变成较好的，你为什么说话？如果你不将不义者变为义人，你为什么说话？但如果你将不义者变为义人，将不敬神者变为敬神者，将愚昧者变为智者，你就是承认我的灵魂会复活，若我听从你并相信。所以，所有创立任何派别的人，即使是虚假宗教，当他们希望被相信时，不得不承认这种精神的复活：所有人都对此认同；但许多人否认肉身的复活，并断言复活已经在信德中发生了。宗徒抵抗这样的人，说：\BibleQuote{其中有依默纳约和非肋托，他们偏离了真理，说复活已过，并颠覆了一些人的信德}\BibleRef{弟后 2:17}。他们说复活已经发生了，但这样一来，另一个复活就不必期待了；他们责备那些期待肉身复活的人，仿佛所应许的复活已经在相信的行为中，即在精神中完成了。宗徒责备这些人。他为什么责备他们？他们岂非肯定了主刚才所说的：\Quote{时辰来到，且现在就是，死人将听见天主子的声音，凡听见的将生活}？但是，耶稣对你说：至今我是在谈论精神的生命；我尚未谈论身体的生命；但我谈论的是那作为身体之生命的生命，即灵魂的生命，身体的生命存在于其中。因为我知道有身体躺在坟墓里；我也知道你们的身体将躺在坟墓里。我不是在谈论那个复活，而是在谈论这个；在此复活中，你们要复活，免得在那复活中你们复活受罚。但为使你们知道我也谈论那个复活，我添加了什么？\Quote{因为父怎样在自己内有生命，也同样赐给了圣子在自己内有生命}。这生命，父是它，子是它，关乎什么？关乎灵魂还是身体？确实不是身体感受那智慧的生命，而是理性的心智。因为并非每个灵魂都有能力理解智慧。事实上，走兽有灵魂，但走兽的灵魂不能理解智慧。因此，是人灵魂能感知这生命，即父在自己内有，并赐给圣子在自己内有的生命；因为那是\Quote{照亮}不是所有灵魂，而是\Quote{来到这世界的每个人}的\Quote{真光}。所以，当我向心智本身说话时，让它听见，即让它服从并生活。\par

15. 因此，主啊，关于肉身的复活，请不要保持沉默；免得人不相信它，而我们继续成为推理者，而非宣讲者。但是\Quote{因为父怎样在自己内有生命，也同样赐给了圣子在自己内有生命}。让听见的人理解；让他们相信，好使他们理解；让他们服从，好使他们生活。并为了不让他们以为复活在此已经完成，让他们进一步听这话：\Quote{并且赐给了祂施行审判的权柄}。谁赐给了？圣父。赐给了谁？圣子；即，祂赐给了圣子在自己内有生命，就赐给了同一者施行审判的权柄。\Quote{因为祂是人子}。因为这是基督，既是天主子，又是人子。\Quote{在起初已有圣言，圣言与天主同在，圣言就是天主。这圣言在起初与天主同在}。看，祂如何赐给祂在自己内有生命！但因为\Quote{圣言成了血肉，寄居在我们中间}，成为人，由童贞玛利亚所生，祂是人子。那么，祂作为人子接受了什么？施行审判的权柄。什么审判？就是世界末日的审判。那时也将有复活，但那是身体的复活。所以，天主藉着基督、天主子，复活灵魂；祂藉着同一位基督、人子，复活身体。\Quote{赐给了祂权柄}。祂本不应有这权柄，除非祂接受了它；否则祂就是一个没有权柄的人。但那位是天主子的，也是人子。因为藉着位格的合一，人子与天主子成为一位格，天主子就是人子所是的同一性体。但它有什么特性，以及为什么，必须加以区分。人子有灵魂和身体。天主子，即\OriginalTerm{天主圣言}{Word of God}，拥有一个人，如同灵魂拥有身体。正如灵魂拥有身体并不构成两个位格，而是一个人；同样，圣言，拥有一个人，并不构成两个位格，而是一个基督。什么是人？一个理性的灵魂，拥有一个身体。什么是基督？天主圣言，拥有一个人。我知道我在说什么，知道说话的我是什么，以及对谁说话。\par

16. 现在请听关于身体的复活，不是听我讲，而是主要说话，因为有些人已经借着从死亡中复活，借着紧紧持守生命而复活了。是何种生命？是一种不知死亡的生命。为何不知死亡？因为它不知易变性。为何不知易变性？因为它是生命本身。\Quote{并且赐给祂权柄施行审判，因为祂是人子。} 是什么审判，何种审判？\Quote{不要惊奇这事} 我所说的——赐给祂权柄施行审判——\Quote{因为时刻将到。} 祂没有加上\Quote{现在就是了}：因此祂的意思是让我们知道世界终穷的某一时刻。现在是死人复活的时刻，将来在世界终穷也是死人复活的时刻：但现在是心灵上复活，将来是肉体上复活；现在是借着天主的圣言、天主子在心灵上复活，将来是借着成为肉身的圣言、人子在肉体上复活。因为来审判的不是圣父自己，尽管圣父并没有离开圣子。那么，圣父自己如何不来呢？在于祂在审判中不会\Emph{被看见}。\Quote{他们要仰望他们所刺透的那一位。} \BibleRef{若 19:37} 那站在审判者面前的形像将要做审判者：那受过审判的形像将要审判；因为它曾受不义的审判，它将要公义地审判。仆人的形像将要出现，而且那形像将要显现。因为天主的形像如何能向义人和不义的人显现呢？如果审判只在义人之间，那么天主的形像也许可以向义人显现。但由于审判是针对义人和不义的人，而恶人不允许看见天主——因为\Quote{心里纯洁的人是有福的，因为他们要看见天主，} \BibleRef{玛 5:8}——这样的审判者将要显现，可以被祂将要加冕的人和祂将要定罪的人看见。因此，仆人的形像将被看见，天主的形像将被隐藏。天主子将隐藏在仆人内，人子将显现，因为\Quote{赐给祂权柄施行审判，因为祂是人子。} 并且因为只有祂以仆人的形像显现，而圣父不显现，因为圣父没有取仆人的形像；因此祂上面说：\Quote{圣父不审判任何人，但把一切审判都交给了子。} 所以它曾被延迟是恰当的，使提出者自己成为解释者。因为先前是隐藏的；现在，如我所想，已经显明了，即\Quote{祂赐给祂权柄施行审判，} 且\Quote{圣父不审判任何人，但把一切审判都交给了子：} 因为审判将要借着圣父所没有的形像进行。是什么审判？\Quote{不要惊奇这事，因为时刻将到：} 不是现在这为灵魂复活的时候，而是将来为身体复活的时候。\par

17. 让祂更清楚地宣告这事，使否认身体复活的异端者找不到诡辩的托词，尽管含义已经清楚地照耀出来。当上面说到\Quote{时刻将到}时，祂加上了\Quote{现在就是了}；但现在，\Quote{时刻将到}，祂没有加上\Quote{现在就是了}。然而，让祂借着公开的真理，撕裂一切诡辩攻击的把柄、圈套和钉子，以及所有迷惑人的反对网罗。\Quote{不要惊奇这事：因为时刻将到，凡在坟墓里的……} 还有什么更明显的？还有什么更清楚的？身体是在坟墓里的；灵魂却不在坟墓里，无论义人的还是不义人的。义人的灵魂曾在亚巴郎的怀抱里；不义之人的灵魂在地狱里受折磨：两者都不在坟墓里。上面，当祂说\Quote{时刻将到，现在就是了}时，我恳求你们留心听。你们知道，弟兄们，我们获得肚腹的食粮需要辛劳；心灵的食粮需要何等更大的辛劳！你们劳苦地站着听，但我们要更劳苦地站着讲。如果我们为你们的缘故劳苦，你们也该为你们自己的缘故与我们一同劳苦。此外，当祂说\Quote{时刻将到}，又加上\Quote{现在就是了}时，祂接着说了什么？\Quote{死人将要听见天主子的声音，听见的人就必活着。} 祂没有说\Quote{所有的死人都要听见，听见的人就必活着}；因为祂的意思是让不义的人被理解。那么，难道所有不义的人都听从福音吗？宗徒公开说：\Quote{但不是所有人都听从福音。} \BibleRef{罗 10:16} 但听见的人必活着，因为所有听从福音的人借着信德将进入永生：然而并非所有人都听从；这是现今的事。但确实，在终末，\Quote{凡在坟墓里的，} 无论义人还是不义人，\Quote{必听见祂的声音，就出来。} 为什么祂不说\Quote{并且要活着}呢？的确，所有人都会出来，但并非所有人都要活着。因为祂在上面所说\Quote{凡听见的必活着}中，祂的意思是让我们明白，在那种听见和听从本身中就有永生和蒙福的生命，并非所有从坟墓出来的人都会拥有。因此，这里无论在提及坟墓时，还是在表达从坟墓中\Quote{出来}时，我们都清楚地理解了身体的复活。\par

18. \Quote{凡听见祂声音的，都要出来。} 若凡听见的都出来，审判又在哪里？好像一切混乱，我看不到区分。诚然祢领受了审判的权柄，因为祢是人子；看，祢将临于审判；尸体将要复活；但请告诉我们关于审判本身的事，即恶人与善人的分离。那么，再听这个：\Quote{行善的，复活进入生命；作恶的，复活进入审判。} 当祂上面谈到心意和灵魂的复活时，祂作了任何区别吗？没有，因为凡\Quote{听见的都要生活}；因着听见，即因着服从，他们将要生活。但并非所有从坟墓中复活出来的都进入永生——只有那些行善的；而行恶的，进入审判。因为祂在这里将审判用作惩罚。也将有一种分离，不像现在这样。因为现在我们被分离，不是按地方，而是按品性、情感、愿望、信德、望德、爱德。现在我们一起与不义者生活，虽然所有人的生命并不相同：我们在暗中被区分，在暗中被分离；如同禾场上的谷粒，而非谷仓里的谷粒。在禾场上，谷粒既分离又混杂：分离，因为从糠秕中分离出来；混杂，因为还没有簸净。那时将有公开的分离；按品性区分生命，如同在智慧中一样，在身体中也将如此。行善的将去与天主的众天使同住；作恶的，将去与魔鬼及其天使一同受苦。奴仆的形体将要过去。因为祂显现自己，正是为了执行审判。审判之后，祂将离开，将带领那以祂为头的身体，并把天国交给天主。\BibleRef{格前 15:24} 那时将公开显现那恶人不能看见的天主的形体，他们必须看见奴仆的形体。祂在另一处也如此说：\Quote{这些人要进入永火}（论及左边的一些人），\Quote{而义人进入永生}；\BibleRef{玛 25:46} 这生命祂在另一处论及：\Quote{永生就是：认识祢，唯一的真天主，和祢所派遣的耶稣基督。} \BibleRef{若 17:3} 那时祂将显现，\Quote{祂虽具有天主的形体，却没有以自己与天主同等为应当紧持不舍的。} \BibleRef{斐 2:6} 那时祂将显现自己，正如祂应许要显示给爱祂的人。因为\Quote{爱我的人，}祂说，\Quote{必遵守我的命令；爱我的人，我父也必爱他；我也要爱他，并将我自己显示给他。} \BibleRef{若 14:21} 祂亲身与那些祂正在对他们说话的人同在；但他们看见了奴仆的形体，没有看见天主的形体。他们被祂自己的驮兽带到祂的寓所去医治；但现在既已痊愈，他们将看见，因为祂说，\Quote{我要将我自己显示给他。} 祂如何被显示与父相等？当祂对斐理伯说：\Quote{看见我的，就是看见了父。} \BibleRef{若 14:19}\par

19. \Quote{我由自己不能作什么；我怎样听见，就怎样审判；我的审判是正义的。} 否则我们会对祂说：\Quote{祢将审判，而父不审判，因为\InnerQuote{审判权全交给了子。}因此，祢不是按父来审判。} 所以祂补充说：\Quote{我由自己不能作什么；我怎样听见，就怎样审判；我的审判是正义的，因为我不寻求自己的意愿，只寻求那派遣我来者的意愿。} 毫无疑问，子愿意赐生命给谁，就赐给谁。祂不寻求自己的意愿，只寻求派遣祂者的意愿。不是我自己的，我本性的意愿；不是我的，不是人子的；不是我的，以免抵抗天主。因为人若行他们自己的意愿，而非天主的意愿，就是做他们喜欢的事，而非天主所命的事；但当他们做他们喜欢的事，却仍遵循天主的意愿时，他们就不是行自己的意愿，虽然他们做了他们愿意做的事。欣然地做你所受命的事，这样你既做了你所愿意的，又不是行你自己的意愿，而是行那命令者的意愿。\par

20. 那么如何？\Quote{我怎样听见，就怎样审判。} 子\Quote{听见}，而父\Quote{指示}给祂，子看见父所作的事。但我们将这些事延后，以便在时间与精力允许的情况下，在读完这段经文之后，尽可能更清晰、更充分地加以处理。如果我说我还能再多讲，你们也许不能再听下去了。再者，或许由于你们渴望听，你们会说：\Quote{我们能。} 那么，我最好承认我的软弱，既然已经疲倦，不能再多讲，而不在你们已经饱足的时候，继续向你们倾注那些你们无法好好消化的事。至于这个我延到今天所承诺的事，若有机会，靠主的帮助，你们仍是我的债权人，直到明天。\par

\SourceRecord{Translated by John Gibb.；Nicene and Post-Nicene Fathers, First Series；Vol. 7.；Edited by Philip Schaff.；Buffalo, NY: Christian Literature Publishing Co.,；1888.；<http://www.newadvent.org/fathers/1701019.htm>.}

% structure: p0001 source=h1 target=section reason=page-title

\section{《论若望福音》第二十篇（若 5:19）}

1. 我们主耶稣基督的话，尤其是那按福音作者若望所记载的——他并非无端地倚靠在主的胸怀，为能汲取那更高智慧的秘密，并以福音宣讲将他所爱而饮下的再吐出来——是如此奥秘深邃，令心术不正的人困扰，却操练那些心地正直的人。因此，亲爱的诸位，请留意所诵读的这几句话。让我们看看，若我们凭祂的恩赐和帮助——祂愿意我们听到祂的话，当时这些话被听到并被记下，以便现在可以诵读——能理解祂在你们刚才听到祂所说的话中的意思：\Quote{我实实在在告诉你们：子不能凭自己做什么，除非祂看见父做什么；因为父所做的事，子也同样做。}\par

2. 现在需要提醒你们这段言论的起因，根据这段经文之前的内容——当时主治愈了那些躺在撒罗满池五廊中的一个人，并对他说：\Quote{拿起你的床，回家去吧。}但祂是在安息日做的；因此犹太人心烦意乱，错误地控告祂是法律的破坏者和违犯者。于是祂对他们说：\Quote{我父到现在一直工作，我也工作。}\BibleRef{若 5:17}因为他们从肉体意义理解安息日的遵守，幻想天主从创造世界的工作之后一直沉睡至今；所以祂祝圣了那一天，从那天开始祂仿佛从工作中安息。然而，对我们的祖先曾颁赐安息日的圣事（\BibleRef{出 20:8}），我们基督徒按精神遵守：避免一切奴性的工作，即一切罪（因为主说：\Quote{凡犯罪的，就是罪的奴隶}），并在心里安息，即心灵的平安。虽然在此生我们竭力追求这安息，但只有在离开此生之后，我们才能达至完全的安息。但天主之所以被称为安息，是因为祂在完成一切受造物之后，不再创造新的受造物。此外，圣经称之为安息，是为告诫我们，在善工之后我们将得安息。因为在\BibleBook{}{创世纪}中如此记载：\Quote{天主看了他所造的一切，认为非常好，天主在第七天安息了，}这是为使你，人啊，考虑到天主自己在善工之后也被称为安息了，就当在行善之后才期望安息；正如天主在按自己的肖像和模样造了人，并在人身上完成了所有非常美好的工作之后，在第七天安息了，照样，你也只有在恢复你被造时所具有的那肖像——你因犯罪而失去了那肖像——之后，才可期望安息。因为实际上，不能说天主辛劳了，祂\Quote{一说，便成就了。}有谁在如此轻易的工作之后，还像劳苦之后那样渴望安息呢？如果祂命令，有人违抗祂，如果祂命令，却没有成就，祂劳苦使其成就，那么祂才算是在劳苦之后安息。但在同一本\BibleBook{}{创世纪}中，我们读到：天主说：要有光，就有了光；天主说：要有穹苍，穹苍就形成了；其余的一切都因祂的话立刻成就。对此\BibleBook{}{圣咏}也作证说：\Quote{祂一发言，万有造成；祂一下令，诸天创造。}——祂既在命令时从不劳苦，又怎会在创造世界之后需要安息，好像劳苦之后享受闲暇？因此，这些说法是神秘的，这样安排是为使我们期望在此生之后得安息，只要我们行了善工。为此，主制止了犹太人的狂妄，驳斥了他们的错误，并表明他们对天主的思想不正确。当他们因祂在安息日医治人而恼怒时，祂对他们说：\Quote{我父到现在一直工作，我也工作：}所以不要以为我父在安息日如此安息，以致此后不再工作；而是正如祂现在工作，我也照样工作。但如父不必劳苦，同样子也不必劳苦。天主\Quote{一说，便成就了；}基督对那瘫痪的人说：\Quote{拿起你的床，回家去吧，}这事便成就了。\par

3. 然而，公教信仰认为：圣父与圣子的工程是不可分离的。亲爱的诸位，如果可能，我愿向你们讲述这一点；但按照主的话：\Quote{谁能领悟，就领悟吧！}（\BibleRef{玛 19:12}）那不能领悟的人，不要归咎于我，而要归咎于自己的迟钝；并要转向那位开启人心的主，求他倾注他所自由赐予的恩惠。最后，如果有人未能理解，是因为我没有按应当的方式宣讲，那就请原谅人性的软弱，并恳求天主的仁慈吧。因为我们内有一位导师——基督。凡你们不能通过你们的耳朵和我的口舌领受的，要在心中转向他：他既教导我该说什么，也按他愿意的分量分施给你们。他知道该给什么，该给谁；他必帮助寻求者，给叩门者开门。即使他没有给，也不要有人以为自己被遗弃了。因为他或许延迟给一些东西，但绝不使人饥饿。如果他不在那一刻给予，那是在锻炼寻求者，并非轻视恳求者。那么，请你们注视并留意我愿说的，即使我可能无法完全说出。公教信仰，由天主之神在祂的圣人们内所确认，反对一切异端的乖谬，确认圣父与圣子的工程是不可分离的。我所说的是什么呢？正如圣父与圣子是不可分离的，同样，圣父与圣子的工程也是不可分离的。圣父与圣子如何不可分离？因为祂自己说过：\Quote{我与父原是一体。}（\BibleRef{若 10:30}）因为圣父与圣子并不是两位神，而是一位天主：圣言与那圣言之所从出者，是一，是独一者；圣父与圣子由爱德相连，是一位天主；而爱德之神也是一位，因此圣父、圣子、圣神组成了天主圣三。所以，不仅圣父与圣子的工程，连圣神的工程也是如此；正如位格之间的平等与不可分离，工程也是不可分离的。我将更清楚地告诉你们\Quote{工程不可分离}是什么意思。公教信仰并不说天主圣父造了某物，圣子造了另一物；而是说圣父所造的，圣子也造了，圣神也造了。因为万物都是藉着圣言造成的；当\Quote{祂一说，它们就造成了}，是藉着圣言造成的，是藉着基督造成的。因为\Quote{在起初已有圣言，圣言与天主同在，圣言就是天主；万物是藉着祂而造成的。}如果万物是藉着祂造成的，那么天主说\Quote{要有光}，就有了光；祂是在圣言中造成的，是藉着圣言造成的。\par

4. 看，我们现在听到了福音，其中祂回答了那些愤慨的犹太人，因为他们\Quote{不单破坏了安息日，而且称天主为自己的父，使自己与天主平等。}（\BibleRef{若 5:18}）上文正是这样记载的。因此，当天主子、真理，回应他们错误的愤怒时，祂说：\Quote{我实实在在告诉你们：子不能由自己作什么，而只作祂看见父所作的事。}仿佛祂在说：\Quote{你们为什么因我说天主是我的父，并使我与天主平等而恼怒？我是如此平等的：是祂生了我的那种平等；我是如此平等的：祂不是出于我，而是我出于祂。}这正包含在以下的话中：\Quote{子不能由自己作什么，而只作祂看见父所作的事。}也就是说，无论子要作什么，祂作的能力都来自父。为什么祂作的能力来自父？因为祂作儿子的身份来自父。为什么祂作儿子的身份来自父？因为祂的能力来自父，祂的存在来自父。因为对于子来说，\Quote{能}与\Quote{是}是同一回事。人却并非如此。你们一定要从人性的软弱（它远远低于天人）的比较中提升你们的心；如果你们中有人或许达到了奥秘，并在仿佛强光的光芒中感到敬畏，略有所见，而不完全无知，那么，不要以为他领悟了全部，免得他变得骄傲，失去已获得的知识。在人身上，\Quote{存在}与\Quote{能够}是不同的。因为有时人\Quote{存在}却\Quote{不能}行他所愿的；有时人又\Quote{存在}得如此，以致他\Quote{能}行他所愿的；因此，他的\Quote{存在}与\Quote{能够}是不同的。如果人的\OriginalTerm{存在}{esse}与\OriginalTerm{能够}{posse}是同一回事，那么他愿的时候就能了。但在天主那里，并非如此：祂的实体\Quote{存在}是一回事，祂的能力\Quote{能够}是另一回事；凡属于祂的，凡祂所是的，都与祂同体，因为祂是天主；并非祂以某种方式\Quote{存在}，以另一种方式\Quote{能够}；祂同时拥有\Quote{存在}与\Quote{能够}，因为祂同时拥有\Quote{愿意}与\Quote{作成}。既然子的能力出于父，那么子的实体也出于父；既然子的实体出于父，那么子的能力也出于父。在子里，能力与实体并无区别：能力就是实体本身；实体是\Quote{存在}，能力是\Quote{能够}。因此，因为子出于父，祂说：\Quote{子不能由自己作什么。}因为祂不是由自己而成为子，所以祂不能由自己而能。\par

5. 祂似乎使自己变得低微，当祂说：\Quote{子不能由自己做什么，除非祂看见父做什么。} 于是异端者的虚妄昂起了头；他们确实说子小于圣父、权柄较小、尊威较小、能力较小，不明白基督话语的奥迹。但是，亲爱的，请留意，看他们如何因基督的话而在肉性理智中困惑。这就是我稍早所说的话：天主的话扰乱所有乖僻的心，正如它锻炼虔诚的心，尤其是由福音作者若望所说的话。因为祂所说的是深奥的话，不是随意的话，也不是容易明白的话。所以，一个异端者若碰巧听到这些话，立刻起来对我们说：\Quote{看，子小于圣父；听子的话，祂说：\InnerQuote{子不能由自己做什么，除非祂看见父做什么。}} 且慢；正如经上所载：\BibleQuote{你要温顺地听道，好能领悟。} \BibleRef{德 5:13} 好吧，假设因为我主张圣父与子的能力和尊威相等，我在听到这些话时就感到困惑：\Quote{子不能由自己做什么，除非祂看见父做什么。} 好吧，我因这些话而困惑，就要问你，你这自认为立刻领悟了的人一个问题。我们知道在福音中，子曾在海面上行走；\BibleRef{玛 14:25}祂何时看见圣父在海面上行走？现在他困惑了。那么，放下你对这些话的理解，让我们一同来查考。我们该怎么办呢？我们已听见主的话：\Quote{子不能由自己做什么，除非祂看见父做什么。} 子在海面上行走，圣父从未在海面上行走。然而确实\Quote{子不能由自己做什么，除非祂看见父做什么。}\par

6. 那么，请与我一同回到我刚才所说的话，看看是否可以这样理解，使我们双方都能避开这个问题。因为我看到我如何能按照公教信仰，不绊倒也不失足地避开；而你，另一方面，四面被困，正在寻找出路。看看你是从哪条路进来的。也许你还不明白我所说的\Quote{看看你是从哪条路进来的}：请听祂自己说：\BibleQuote{我就是门。} \BibleRef{若 10:7} 所以你寻找出去的方法并非没有原因；你只发现你没有从门进来，而是从墙上跌进来的。因此，你尽力从你的跌倒中起来，从门进来，好使你进去时不绊倒，出去时不迷路。经由基督而来，不要把你心里所想的话拿出来说；而是祂所显示的，你就说。请看公教信仰如何避开这个问题。子在海上行走，把肉身的脚放在波浪上：肉身行走，神性引导。但当肉身行走、神性引导时，圣父缺席了吗？若缺席，子自己怎能说：\BibleQuote{但住在我的父，祂自己做这些事？} \BibleRef{若 14:10} 如果圣父住在子内，自己做祂的工作，那么那次海上行走就是由圣父并通过子所作的。因此，那次行走是圣父与子不可分离的工作。我看见两者在其中工作。圣父没有撇下子，子也没有离开圣父。如此，子所做的一切，没有圣父就不能做；因为圣父所做的一切，没有子也不能做。\par

7. 我们已经避开了这个问题。请注意，我们正当地说：圣父、子和圣神的工作是不可分离的。但照你的理解，看，天主造了光，子看见圣父造光——根据你肉性的理解，你宁愿认为祂较小，因为祂说：\Quote{子不能由自己做什么，除非祂看见父做什么。} 天主父造了光；子造了别的什么光呢？天主父造了穹苍，即水与水之间的天；子看见了祂——根据你这迟钝而朦胧的理解。好吧，既然子看见圣父造穹苍，并且说：\Quote{子不能由自己做什么，除非祂看见父做什么，}那就请你指出子造的另一个穹苍。你失去了根基吗？但那\BibleQuote{被建造在宗徒和先知的基础上的，耶稣基督自己就是屋角基石，他们在基督内被带入和平之境；} \BibleRef{弗 2:14} 他们不争辩，也不在异端中游荡。因此我们明白：光是由天主父（所造），但是通过子（所造）；穹苍是由天主父（所造），但是通过子（所造）。因为\BibleQuote{万物是借着祂而造的，凡受造的，没有一样不是借着祂而造的。} 抛弃你的理解吧，它不该被称为理解，而显然是愚蠢。天主父造了世界；子造了别的什么世界呢？请把子的世界指给我看。我们所在的这个世界是谁的？告诉我们，是谁造的？如果你说：\Quote{是子造的，不是圣父造的}，那么你就背离了圣父；如果你说：\Quote{是圣父造的，不是子造的}，福音就这样回答你：\BibleQuote{世界是借着祂而造的，世界却不认识祂。}\par

8. 为此，圣父和圣子的工程是不可分的。此外，这句话：\Quote{子不能凭自己做什么}\BibleRef{若 5:19}，就等于说\Quote{子不是出于自己}。因为若他是子，他就是被生的；若他被生，他就是出于那生他的。然而，圣父生了与他相等的子。那生的并不缺少什么；他生了同永生的子，并不需要时间来生；他从自己产生了\OriginalTerm{圣言}{Word}，并不需要母亲来生；圣父在生的时候，在年岁上并不先于圣子，以致他生了一个比自己年轻的子。但或许有人说，天主在众多世代之后，于晚年生了子。然而，正如圣父没有年岁，圣子也没有成长；圣父没有变老，圣子也没有长大，而是相等者生了相等者，永生者生了永生者。有人会说，永生者如何生了永生者？如同暂时的火焰生出暂时的光。生光的火焰与它所生的光同时存在：生光的火焰在时间上并不先于所生的光；火焰从何时开始，光就从何时开始。给我看没有光的火焰，我便给你看没有圣子的天主圣父。因此，\Quote{子不能凭自己做什么，除非他看见父做什么}\BibleRef{若 5:19}，意思是，对于子而言，\Emph{看见}和\Emph{由父所生}是同一回事。他的看见和他的本体并无不同；他的能力和他的本体也无不同。凡他所有的，都是出于圣父；凡他能做的，都是出于圣父；因为他的所能和所是乃是一回事，并且都出于圣父。\par

9. 此外，他继续用自己的话，搅扰那些误解此事的人，以便将迷途者召回正确的理解。在他说了\Quote{子不能凭自己做什么，除非他看见父做什么}之后；恐怕有人对此产生属肉体的理解，偷偷潜入，使人思想偏斜，想像如同两个工匠，一个是师傅，一个是学徒，学徒专心观察师傅做……比如说，一个箱子，这样，当师傅做箱子时，学徒便依照他所看见的师傅工作的样子，做另一个箱子；恐怕，我说，属肉体的人在天主统一性的事上，为自己塑造成这种双重的观念，于是他接着说：\Quote{因为凡父所做的，子也同样做}\BibleRef{若 5:19}。并非父做某些事，子做另一些相似的事，而是同样的事，以同样的方式。因为他不是说，\Quote{凡父所做的，子也做另一些相似的事}；而是说：\Quote{凡父所做的，子也同样做}。凡父所做的，子也做：圣父造了世界，圣子造了世界，圣神造了世界。若有三神，便有三个世界；若只有一天主——圣父、圣子、圣神，那么只有一个世界，由圣父藉着圣子在圣神内造成。因此，子所做的事，也是父所做的事，并且不是以不同的方式；他既做这些事，也以同样的方式做。\par

10. 在他说了\Quote{做这些事}之后，为什么又加上\Quote{同样做}呢？以免另一个歪曲的理解或谬误在人的思想中产生。你看见，比方说，人的工作：人有心神和身体；心神管理身体，但心神和身体之间有很大差别：身体可见，心神不可见；心神的能力和德能与任何一种身体（即使是天上的身体）的能力和德能之间有很大差别。然而，心神管理它的身体，身体也行事；心神似乎做什么，身体就做什么。这样，身体似乎做心神所做的同一件事，但不是\Quote{同样地}做。如何做同一件，却不是同样地？心神在自己内形成言语；它命令舌头，舌头便发出心神所形成的话语：心神造了，舌头也造了；身体的主人造了，仆人（身体）也造了；但为了使仆人能造，它从主人那里领受了要造什么，并在主人命令时造出。同一件事由两者做成，但它是同样地吗？如何不是同样地？有人会说。你看，我的心神所形成的言语，留在我内；而我的舌头所造的言语，穿过被击打的空气，就不再存在了。当你在心神中说了一句话，并用舌头说出来之后，回到你的心神中，看看你所造的那句话是否还在。它是否如同留在你心神中那样，留在你的舌头上？舌头所发出的，是舌头通过发声造的，心神则通过思想造；但舌头所发出的已经过去，心神所思想的仍然存留。因此，身体造了心神所造的，但不是同样地。因为心神造了心神所能持守的，而舌头所造的则发出声音，穿过空气刺激耳朵。你追捕音节，能使它们存留吗？然而，圣父和圣子并非如此；而是\Quote{做这些事}，并且\Quote{同样地做}\BibleRef{若 5:19}。若天主造了存留的天，这存留的天是圣子所造。若天主圣父造了会死的人，这同一会死的人是圣子所造。凡圣父所造而持久的事物，这些持久的事物圣子也造了，因为他以同样的方式造了；凡圣父所造而暂时的事物，这些暂时的事物圣子也造了，因为他不但造了相同的事物，而且以同样的方式造了。因为圣父藉着圣子造了万物，因为圣父藉着圣言造了万有。\par

11. 在圣父和圣子中寻找分离，你找不着；不，即使你已升至高处，也找不着；不，即使你已达到超乎你理智之上，也找不着。因为如果你在你游荡的理智自造的事物中打转，你是在与自己的想象交谈，而不是与天主的圣言交谈；你自己的想象欺骗了你。也越过肉身，去理解理智；也越过理智，去理解天主。除非你越过了理智，否则你到不了天主；更何况你若滞留于肉身，又怎能到达天主呢！那些思念肉身的人，离理解天主何其远！——因为即使他们理解了理智，也到不了那里。当人的思想属于肉身时，人就远离天主；肉身与理智之间有天壤之别，而理智与天主之间的差别更大。若你专注于理智，你就在中间：你若将注意力向下，便有肉身；若向上，便有天主。从肉身中提升自己，甚至越过你自己。请看圣咏所说，你便被提醒应如何思想天主：\Quote{我的眼泪，} 它说，\Quote{日夜成了我的食粮，因为有人天天对我说：\InnerQuote{你的天主在哪里？}} 正如外邦人可能说：\Quote{请看我们的神，你的天主在哪里？} 他们确实向我们显示可见的；我们敬拜不可见的。而我们可以向谁显示呢？向一个没有视力去看的人吗？无论如何，如果他们用眼睛看见他们的神，我们也有另一双眼睛来看我们的天主：因为\Quote{心地纯洁的人是有福的，因为他们要看见天主。} \BibleRef{玛 5:8} 因此，当他说他苦恼时，因为有人天天对他说：\Quote{你的天主在哪里？} \Quote{我记起了这些事，} 他说，\Quote{因为天天有人对我说：\InnerQuote{你的天主在哪里？}} 他仿佛想要抓住他的天主，就说：\Quote{这些事，} 他说，\Quote{我记起了，并在自己之上倾注了我的灵魂。} 因此，为能到达我的天主，就是那位关于他有人对我说：\Quote{你的天主在哪里？我倾注了我的灵魂}的那位，不在我的肉身之上，而是\Quote{在我之上}；我超越了自己，为能到达他：因为那造我的他在我之上；没有人能到达他，除非他超越了自己。\par

12. 试想肉身：它是可朽的、属土的、软弱的、可败坏的；抛弃它吧。是的，也许你会说，但肉身是暂时的。那么再思想其他的身体，即天上的；它们更伟大、更美好、更辉煌。再仔细看它们。它们从东向西运行，不停留；它们被眼睛看见，不仅被人，也被田间的野兽看见。也越过它们吧。而你会说，我行走在地上，怎能越过天上的物体呢？你不是在肉身中越过它们，而是在理智中。也抛弃它们吧：尽管它们如此发光，它们是物体；尽管它们从天上闪耀，它们是物体。好了，现在也许你认为，在考虑了所有这些之后，你已无处可去了。而你说，我要越过天上的物体到哪里去？我要用理智越过什么？你考虑过所有这些了吗？你说，我考虑过了。你是用什么来考虑它们的？让那考虑者亲自显现。那考虑这一切、分辨、区分、并以智慧的天平权衡它们的，实在是理智。无疑，那么，你用来观看这一切的理智，比你所观看的这一切更好。这个理智，就是一个精神，而不是一个身体。也越过它吧。并且为让你看见你要越过到哪里，首先，将那个理智本身与肉身相比。但愿你绝不要降格去这样比较！把它与太阳、月亮和星辰的光辉相比；理智的光辉更伟大。首先，观察理智的敏捷；看看这思考的理智的闪烁是否不比照耀的太阳的光芒更猛烈。用理智，你看见太阳升起。它的运行与你的理智相比是多么缓慢！太阳将要做什么，你一瞬间就能想到。它将从东到西；明天从另一边升起。你的思绪已经完成了这些，太阳却还落在后面，而你已走完了全程。因此，理智是伟大的。但我怎么说\Emph{是}呢？也越过它吧。因为理智，尽管它比每一种身体都好，它本身是可变的。它时而知，时而不知；时而忘记，时而记起；时而愿意，时而不愿意；时而错误，时而正确。因此，越过一切可变之物；不仅越过一切可见的，也越过一切变化的。因为你已越过了可见的肉身；越过了可见的天、太阳、月亮和星辰。也越过一切变化的。因为当你处理完那些可见的事物，来到你的理智时，你发现了你理智的可变性。天主难道是可变吗？那么，越过你的理智。倾注你的灵魂\Quote{在你之上}，为能到达天主，就是那一位关于他有人对你说：\Quote{你的天主在哪里？}\par

13. 不要以为你要做超乎人的能力的事。圣史若望本人就这样做了。他超越了肉身，超越了他所踏的大地，超越了他所注视的海洋，超越了飞鸟翱翔的空气，超越了太阳、月亮、星辰，超越了一切不可见的神体，超越了他自己的理智，藉着他理性灵魂的理性本身。超越这一切之后，在他之上倾注他的灵魂，他到了哪里？他看到了什么？\BibleQuote{在起初已有圣言，圣言与天主同在。} 因此，如果你在光中看不见分离，为什么要在工程中寻求分离呢？请看天主，看他的圣言紧附于那发言的圣言，以致发言者不是藉音节发言，而是他的发言是智慧光明的照耀。关于智慧本身有何记载？\BibleQuote{它是永恒光明的辉映。} \BibleRef{智 7:26} 观察太阳的光辉。太阳在天上，将它的光明散布在一切陆地和一切海洋之上，它不过是一个有形的光。\par

如果你真能将光辉与太阳分开，那就能将圣言与圣父分开。我指的是太阳。一盏灯细微的火焰，一口气就能吹灭，却能向四周的一切散播光芒：你看见火焰所产生的光向外扩散，你看见它的放射，却看不见分离。那么，亲爱的弟兄们，请明白：\Person{圣父}、\Person{圣子}、\Person{圣神}在自身内是不可分离的合一；这个天主圣三是唯一的天主；唯一天主的一切作为，都是\Person{圣父}、\Person{圣子}、\Person{圣神}的作为。其余所有的内容，即与我们的\Person{主耶稣基督}的讲论相关的部分，既然明天也应当向你们作一篇讲道，那么请你们出席来听。\par

\SourceRecord{Translated by John Gibb.；Nicene and Post-Nicene Fathers, First Series；Vol. 7.；Edited by Philip Schaff.；Buffalo, NY: Christian Literature Publishing Co.,；1888.；<http://www.newadvent.org/fathers/1701020.htm>.}

% structure: p0001 source=h1 target=section reason=page-title

\section{训道篇第21篇（若望福音5:20-23）}

1. 昨日，主赏赐我们尽量讨论，并按照我们的能力领会，圣父与圣子的工程如何是不可分离的；圣父不作一些工程，圣子作另一些，而是圣父借着圣子，借着祂的圣言，成就一切，如经上记载：\BibleQuote{万物是借着祂造的；凡受造的，没有一样不是借着祂造的。}今日让我们看看后文。让我们向同一的主祈求怜悯，并希望，若祂认为适宜，我们能明白真理；但若我们不能做到这一点，至少不至于陷入谬误。因为不知道比迷途更好；但知道比不知道更好。因此，在一切事上，我们应当努力追求知道。若能知道，愿光荣归于天主；若暂时不能达到真理，也不要走向虚假。因为我们应当深思我们是什么，我们在讨论什么。我们是血肉之人，行走在此生中；虽已由天主圣言的种子重生，在基督内更新，但尚未完全脱离亚当。的确，我们那沉重压着灵魂的、可朽坏的部分（\BibleRef{智 9:15}）显明是，并明明白白是来自亚当的；但我们里面的精神部分，提升灵魂的，是来自天主的恩赐和祂的慈悲，祂派遣了祂的独生子参与我们的死亡，并引领我们到达祂自己的不朽。圣子是我们的导师，使我们可以不犯罪；也是我们的辩护者，若我们犯了罪并告解、悔改；为我们转求，若我们向天主祈求任何美善；并同圣父一起赐予，因为圣父和圣子是同一的天主。但祂作为人向人讲说这些：隐藏的天主，显现的人，为使他们——显现的人——成为神；天主子成为人子，为使人的儿子成为天主的儿子。祂以祂何等智慧之工成就此事，我们在祂自己的话语中感知。因为祂向小孩子如同小孩子说话，但祂自己如此幼小，却又伟大，我们虽幼小，在祂内却伟大。祂说话，实为抚育和滋养吃奶的婴孩，使他们因爱而成长。\par

2. 祂曾说过：\BibleQuote{子不能凭自己做什么，唯有看见父做什么，才能做什么。}我们却理解并非父单独做什么，子看见后才做同样的事；但祂说：\BibleQuote{子不能凭自己做什么，唯有看见父做什么，才能做什么，}我们理解子完全出自父——祂的全部实体和全部能力都出自生祂的父。但刚才，当祂说过祂像父一样做这些事，以免我们理解为父做一些，子做另一些，而是子以同等的能力做父所做的同样的事，而父借着子做，祂接着说出了我们今天所读的：\BibleQuote{因为父爱子，把自己所作的一切指示给祂看。}有死的思想再次被搅动。父将自己所作的事指示给子看；因此，有人说，父单独做事，以便子能看见祂做什么。再次，在人的思想中，出现如两位工匠——比如，木匠教儿子自己的技艺，并给他看自己所做的，好让儿子也能做。祂说：\BibleQuote{指示给祂看自己所作的一切。}难道因此，当祂在做时，子不做，以便能看见父做？然而，当然，\BibleQuote{万物是借着祂造的；凡受造的，没有一样不是借着祂造的。}由此我们看出父如何指示子祂所做的，因为父所做的没有一样不是借着子做的。父做了什么？祂造了世界。祂是否将已造的世界指示给子看，以便子也造类似之物？那么让我们看看子所造的世界。然而，\BibleQuote{万物是借着祂造的；凡受造的，没有一样不是借着祂造的，}以及\BibleQuote{世界也是借着祂造的。}如果世界是借着祂造的，万物都是借着祂造的，而父除了借着子什么也不做，那么父在哪里指示子祂所做的呢？难道不就是在子本身内、借着子而行动的吗？在哪里，父的工程能被指示给子看，好像祂在外面坐着做，而子注意看父的手如何做？那不可分离的圣三在哪里？那被称为\BibleQuote{天主的德能和天主的智慧}的圣言在哪里？\BibleRef{格前 1:24} 关于同一智慧，圣经所说的：\BibleQuote{祂是永恒之光的辉映}？\BibleRef{智 7:26} 关于祂再次所说的：\BibleQuote{祂有力地达至地极，并温柔地治理万有}？\BibleRef{智 8:1} 无论父做什么，祂都借着子做：借着祂的智慧和能力；不是从外面指示子祂可能看见的，而是在子本身内，祂指示祂所做的。\par

3. 圣父看见什么，或者更确切地说，圣子在圣父内看见什么，以便自己也做？也许我能说出来，但请给我找一个能理解的人；或者也许我能思考却说不出来；或者也许我甚至不能思考。因为那神性超越我们，如同天主超越人，如同不朽的超越必死的，如同永恒的超越暂时的。愿祂激励并赐予我们，并从生命之源中屈尊滴下一些露珠落在我们的干渴上，使我们不致在这旷野中干涸！让我们对祂说：主，我们已学会称祂为父。我们敢于这样说，因为祂自己愿意；只要我们如此生活，以致祂不对我们说：\BibleQuote{如果我是父亲，我的荣耀在哪里？如果我是主，我的敬畏在哪里？} 那么，让我们对祂说：\Quote{我们的父。} 我们对谁说\Quote{我们的父}？是对基督的父。那么，向基督的父说\Quote{我们的父}的人，向基督说的是除了\Quote{我们的兄弟}之外还能是什么？然而，祂并不像基督的父那样同样成为我们的父；因为基督从未如此连接我们，以至于不区分祂和我们。因为祂是与父平等的圣子，与父同在的永恒圣子，并与父同永恒；但我们藉着圣子成为儿子，藉着\OriginalTerm{独生子}{Only-begotten}被收养。因此，从未从我们的主耶稣基督的口中听到，当祂对门徒说话时，称至高天主祂的父为\Quote{我们的父}，而是要么说\BibleQuote{我的父}，要么说\BibleQuote{你们的父。} 但祂没有说\Quote{我们的父}；以至于在某处祂用了这两种说法：\BibleQuote{我升到我的天主那里去，} 祂说，\BibleQuote{也升到你们的天主那里去。} 为什么祂不说\Quote{我们的天主}？此外，祂说了\BibleQuote{我的父和你们的父，} 祂没有说\Quote{我们的父。} 祂如此结合以致于区分，如此区分以致于不分离。祂愿意我们在祂内成为一体，但圣父和祂自己是一体。\par

4. 那么，无论我们如何理解，无论我们如何看见，我们都不会像圣子那样看见，即使当我们与天使同等时。因为即使我们看不见，我们也是某种存在；但我们看不见时，除了是不看见的人之外，还是什么？而为了看见，我们转向我们可以看见的那一位，在我们内形成了一种先前没有的看见，尽管我们原先就存在。因为一个人在不看见时\Emph{是}存在的；而同一个人在看见时被称为一个看见的人。那么，对他而言，看见与作为人不是同一回事；因为如果那是同一回事，他不在看见时就不会是人。但既然他在不看见时是人，并且寻求看见他所看不见的，他就是一个寻求者，并转向去看；当他正确转向并看见了，他就成为一个看见的人，而先前是一个不看见的人。因此，看见对他而言是一件来去的事；当他转向时，它来到他那里；当他转离时，它离开他。圣子也是这样吗？我们绝不要这样想。祂从未处于一种状态：是圣子却不看见，然后后来被造成看见；但看见圣父对祂而言与是圣子是同一回事。因为我们因转向罪恶而失去光照，因转向天主而接受光照。因为光照我们的光是一回事；我们被光照的是另一回事。但光照我们的光本身既不转离自身，也不失去它的光明，因为它作为光而存在。圣父就这样向圣子显示祂所做的事，使得圣子在圣父内\Emph{看见}一切，并在圣父内\Emph{是}一切。因为祂藉着看见而被生，并且藉着被生而看见。然而，并非祂曾经没有被生，而后被生；也非祂曾经不看见，而后看见。但祂的看见在于什么，祂的存在也在于同样的，祂的被生也在于同样的，祂的延续也在于同样的，祂的不变也在于同样的，祂的无始无终的永存也在于同样的。所以我们不要从肉性的意义上理解，以为圣父坐着做一件工作，并把它显示给圣子；而圣子看见圣父所做的工作，在另一个地方或用其他材料做另一件工作。因为\BibleQuote{万有是藉着祂造成的，没有祂，什么也没有造成。} 圣子是圣父的\OriginalTerm{圣言}{Word}。圣父所说的没有一句不是在圣子内说的。因为通过在圣子内说出祂将藉着圣子做的事，祂生了圣子，并藉着祂创造了万有。\par

5. \Quote{祂将显示比这些更大的工作，叫你们希奇。} 这里我们又感到困惑。谁能配得考察如此伟大的奥秘呢？但现在，既然祂屈尊向我们说话，祂亲自开启了它。因为祂不会说祂不让我们理解的话；既然祂屈尊说话，无疑祂引起了注意：难道祂会撇下那些被祂唤起而留心听的人吗？我们说过，子的认知不是时间性的——子的认知与子本身并非不同；祂的看见与祂本身也不是不同；而是看见本身即是子，认知以及父的智慧即是子；并且那智慧与看见是永恒的，与那从其而来的那一位同等永恒；它不是随时间变化的东西，也不是本来不存在而后来产生的东西，也不是存在过而又消失的东西。那么，在这种情况下，时间起了什么作用，以致祂说：\Quote{祂将显示比这些更大的工作给祂}？\Quote{祂将显示}，就是说，\Quote{祂将要显示}。\Emph{已经显示}与\Emph{将要显示}不同：我们称\Emph{已经显示}为过去的行为；\Emph{将要显示}为未来的行为。那么，弟兄们，我们在这里该怎么办呢？看哪，我们曾宣告祂与父同等永恒，在祂内没有任何因时间而变化的东西，也没有任何通过时刻或空间的移动，我们曾宣告祂永远与父同在，看见父，并且因看见而存在；我说，祂在这里又提到时间，对我们说：\Quote{祂将显示比这些更大的工作给祂。} 那么，父将要显示给子一些子还不知道的东西吗？那么，我们该怎么办？我们如何理解这一点？看哪，我们的主耶稣基督在上，也在下。祂何时在上？当祂说：\Quote{凡父所做的，子也同样做。} 我们如何知道祂现在在下？从这里：\Quote{祂将显示比这些更大的工作给祂。} 主耶稣基督，我们的救主，天主的圣言，万物藉祂而受造，父将要显示什么给祢，是祢还不知道的呢？父有什么对祢是隐藏的呢？在父内的什么对祢是隐藏的，而父对祢并不隐藏？祂将要显示什么更大的工作给祢？或者说，祂要显示给祢的工作比哪些工作更大？因为当祂说\Quote{比这些更大}时，我们应当首先理解那些比之更大的工作。\par

6. 让我们再次回忆这次讲论是从何开始的。那是当那个患病三十八年的人被治愈，耶稣命令现已痊愈的他拿起床榻回家的时候。确实，为此原因，与祂说话的犹太人被激怒。祂在言语上说话，而在意义上保持沉默；向那些理解的人暗暗地暗示了意义，却向那些发怒的人隐藏了这事。为此原因，我说，犹太人因主在安息日做了这事而发怒，便给了这次讲论的契机。因此，我们不要听到这些事好像忘记了前面所说的话，而是要回顾那个软弱无力三十八年的人突然痊愈，而犹太人惊奇并恼怒。他们从安息日寻求黑暗，甚于从奇迹寻求光明。那么，在他们发怒时，祂对他们说：\Quote{父将显示比这些更大的工作给祂。} \Quote{比这些更大：}比哪些？比你们所看见的，一个患病三十八年的人被治愈；父将要显示比这些更大的工作给子。更大的工作是什么？祂接着说：\BibleQuote{因为父怎样复活死者并赐给他们生命，子也照样随自己的意愿赐予生命。} 显然这些是更大的。死者复活比病人痊愈大得多：这些是更大的。但父何时将要显示这些给子？子不知道它们吗？说话的那一位，难道不知道如何复活死者吗？祂还需要学习如何使死者复活——我说的是祂，万物藉祂而受造的那一位？祂使我们存在时，我们尚未存在，难道祂还需要学习如何使我们复活吗？那么，祂的话是什么意思呢？\par

7. 但现在祂屈尊就卑，那刚才还以天主身份发言的祂，如今开始以人的身份发言。尽管如此，那位是天主的人仍是同一人，因为天主成了人；但祂成了祂所不是的，却没有失去祂所是的。因此，人加于天主，使那天主是谁的人成为人，但并非叫祂从此以后是人而不是天主。那么，让我们也听祂作为我们的弟兄，我们曾听祂作为我们的造物主。我们的造物主，因为祂是起初的圣言；我们的弟兄，因为祂生于童贞玛利亚：造物主，在亚巴郎之前，在亚当之前，在地之前，在天之前，在一切有形与无形的受造物之前；但作为弟兄，祂出自亚巴郎的后裔，出于犹大支派，出于以色列童贞女。因此，如果我们认识这位同时以天主和人的身份向我们发言的祂，就让我们理解天主和人的言语；因为有时祂向我们讲述属于尊威的事，有时讲述属于谦卑的事。因为那同一位是崇高的，却成了卑微的，为使我们这些卑微的成为崇高的。那么，祂说什么？\BibleQuote{父将显示}给我\BibleQuote{比这些更大的事，叫你们惊奇。}因此，祂将显示给我们，不是给祂。既然父要显示给我们，因此祂说：\BibleQuote{叫你们惊奇。}事实上，祂已解释了祂说\BibleQuote{父将显示}给我的意思。祂为何不说：父将显示给你们，而说：祂将显示给子？因为我们也都是子的肢体；并且，正如我们这些肢体所学习的，祂自己在某种意义上也在祂的肢体中学习。祂如何在我们内学习？如同祂在我们内受苦。我们从何证明祂在我们内受苦？从那来自天上的声音：\BibleQuote{扫禄，扫禄，你为什么迫害我？}\BibleRef{宗 9:4}难道不正是祂自己，在世界的终结要坐在审判者的位上，将义人放在右边，恶人放在左边，说：\BibleQuote{你们蒙我父降福的，来承受国度；因为我饿了，你们给了我吃}？当他们回答：\BibleQuote{主，我们几时看见祢饿了？}祂将对他们说：\BibleQuote{你们对我这些最小兄弟中的一个所做的，就是对我做的。}\BibleRef{玛 25:31}此时让我们质问祂，并对祂说：主，祢什么时候会成为一个学习者，既然祢教导万有？的确，祂立刻在我们信德中回答我们：当我的最小兄弟中的一个学习时，我就学习。\par

8. 让我们欢喜，并感恩，因为我们不仅成了基督徒，更成了基督。弟兄们，你们明白并领会天主赐予我们的恩宠吗？惊奇吧，喜乐吧，我们成了基督。因为如果祂是头，我们便是肢体：那完整的人便是祂和我们。这正是宗徒保禄所说的：\BibleQuote{使我们不再作小孩子，被各种教义之风摇动，飘来飘去。}但之前祂说过：\BibleQuote{直到我们众人一同达到信仰的统一，并认识天主子，成为完美的人，达到基督圆满年龄的尺度。}\BibleRef{弗 4:14}那么，基督的圆满便是头和肢体。头和肢体是什么？基督和教会。我们若骄傲地自认这一点，确实是僭越，除非祂自己屈尊许诺了它，祂借着同一位宗徒说：\BibleQuote{你们却是基督的身体，各自作肢体。}\BibleRef{格前 12:27}\par

9. 那么，每当父显示给基督的肢体，就是显示给基督。一个伟大却真实的奇迹发生了。有一个显示给基督的，是基督所知道的，但它经由基督显示给基督。这是一件奇妙而伟大的事，但圣经如此说。我们难道要反驳天主的宣言吗？我们岂不更该理解它们，并因祂白白赐予我们的恩赐而感谢祂？我所说的\BibleQuote{经由基督显示给基督}是什么？是经由头显示给肢体。看，你可以在你自己身上看到这一点。假设你闭着眼睛要拿起某物，你的手不知道往哪里去；但你的手总归是你的肢体，因为它没有与身体分离。睁开你的眼睛，现在手看见往哪里去；头显示，肢体跟随。那么，如果在你自己身上能发现这样的事，即你的身体显示给你的身体，并且通过你的身体有东西显示给你的身体，那么就不要惊奇说：经由基督显示给基督。因为头显示使肢体看见，头教导使肢体学习；然而，一个人，头和肢体。祂不愿分离自己，却屈尊使自己与我们结合。祂离我们很远，非常远。有什么比受造物和造物主更远？有什么比天主和人更远？有什么比正义和不义更远？有什么比永恒和必朽更远？看，起初的圣言，与天主同在的天主，万物藉着祂而造，离我们如此之远。那么，祂如何成为近的，以致祂成为我们所是的，而我们在祂内？\BibleQuote{圣言成了血肉，居住在我们中间。}\BibleRef{若 1:14}\par

10. 因此，祂将要向我们显示这一点；祂曾向那些在肉身中看见祂的门徒显示了这一点。这是什么？\BibleQuote{父怎样使死人复活，赐予他们生命，子也照样赐予生命给祂所愿意的人。}难道父复活一些，子复活另一些吗？事实上，万物都是藉着祂造的。我们要说什么呢？我的弟兄们？基督复活了\Person{拉匝禄}；父复活了哪个死人，好让基督看见如何复活拉匝禄呢？当基督复活拉匝禄时，难道不是父复活了他吗？或者是子独自的作为，没有父？你们自己读那段经文，就能看见祂呼求父使拉匝禄复活。\BibleRef{若 11:41}作为人，祂呼唤父；作为天主，祂与父一同行事。因此，复活的拉匝禄也是由父和子一同复活的，在圣神的恩赐和恩宠中；那奇妙的工作是天主圣三完成的。因此，我们不可这样理解：\BibleQuote{父怎样使死人复活，赐予他们生命，子也照样赐予生命给祂所愿意的人，}以致认为有些人是被父复活赐命的，有些人是被子复活赐命的；而是子所复活赐命的，正是父所复活赐命的同一群人，因为\BibleQuote{万物是藉着祂造的，凡受造的，没有一样不是藉着祂造的。}为表明祂虽然由父所赐，却拥有同等的权能，因此祂说：\BibleQuote{子也照样赐予生命给祂所愿意的人，}好在此事上显示祂的意愿；又免得有人说：\Quote{父藉着子复活死人，但父是强有力的，且拥有权能，而子如同仆人、天使那样，藉着另一个的权能行事，}祂在说\BibleQuote{子也照样赐予生命给祂所愿意的人}时，指出了祂的权能。并非父愿意的与子不同；而是父与子有同一本体，因此也有同一意愿。\par

11. 这些被父和子赐予生命的死人是谁呢？可是我们所谈过的那些人——\Person{拉匝禄}、或那位寡妇的儿子（\BibleRef{路 7:14}）、或会堂长的女儿（\BibleRef{路 8:54}）？因为我们知道，这些人是由主基督复活的。祂要向我们意指的乃是另一件事——即我们众人所期盼的死人复活；不是某些人曾有过的那种复活，为使其余的人相信。因为拉匝禄复活后还要再死；我们却要复活，永远活着。成就这种复活的是父呢，还是子？不，确实是父在子内。因此，子是（在父内），而父在子内。我们从哪里证明祂是谈论这种复活呢？当祂说了：\BibleQuote{父怎样使死人复活，赐予他们生命，子也照样赐予生命给祂所愿意的人，}免得我们在此理解为祂为行奇迹而非为永生所行的复活，祂接着说道：\BibleQuote{因为父不审判任何人，却把审判权完全交给了子。}这是什么意思？祂是在谈论死人复活，说：\BibleQuote{父怎样使死人复活，赐予他们生命，子也照样赐予生命给祂所愿意的人，}然后立即添加了一个理由，关于审判，说：\BibleQuote{因为父不审判任何人，却把审判权完全交给了子。}祂为何这样说了，除非是要指明祂所谈论的是那将在审判时发生的死人复活？\par

12. \BibleQuote{因为，}祂说，\BibleQuote{圣父不审判任何人，但已把一切审判交给了圣子。}（\BibleRef{若 5:22}）稍前，我们曾认为圣父做了圣子不做的事，因为祂说：\BibleQuote{圣父爱圣子，并将自己所作的一切指示给祂。}（\BibleRef{若 5:20}）好像圣父在做，而圣子在观看。这样，一种肉性的概念潜入我们心中，仿佛圣父做了圣子没做，而圣子只是看着圣父展示祂所做的事。然后，正如圣父做了圣子没做的事，我们刚才又看到圣子做了圣父没做的事。祂如何扭转我们，使我们的心忙碌！祂引导我们来回，不允许我们停留在肉体的某一处，好藉着改变来锻炼我们，藉着锻炼来净化我们，藉着净化使我们能够领受，并在我们能够领受时充满我们。这些话与我们有什么关系？祂刚才在说什么？祂现在在说什么？稍前，祂说圣父将祂所做的一切指示给圣子。我仿佛看到圣父在做，圣子在等待看；不久之后，我又看到圣子在做事，圣父闲置着：\BibleQuote{因为圣父不审判任何人，但已把一切审判交给了圣子。}（\BibleRef{若 5:22}）因此，当圣子即将审判时，圣父会闲置而不审判吗？这是什么？我该如何理解？主啊，祢说什么？祢是天主圣言，我是人。祢说\BibleQuote{圣父不审判任何人，但已将一切审判交给了圣子}？我在另一处读到祢说：\BibleQuote{我不审判任何人；有一个寻找并审判的。}（\BibleRef{若 8:15}）祢说\BibleQuote{有一个寻找并审判的}，除非是指圣父？祂查问你们的过犯，并为此审判。那么这里\BibleQuote{圣父不审判任何人，但已将一切审判交给了圣子}该如何理解？让我们问问\Person{伯多禄}；让我们听他书信中的话：\BibleQuote{基督为我们受了苦，}（他说，）\BibleQuote{给我们留下榜样，叫我们追随祂的足迹；祂没有犯过罪，口中也没有诡诈；祂被辱骂时，不还骂；受难时，不恐吓，却将自己交托给那按公义审判的。}（\BibleRef{伯前 2:21}）那么\BibleQuote{圣父不审判任何人，但已将一切审判交给了圣子}这话如何真实？我们在这里感到困惑，因困惑而奋发，好藉奋发而净化。让我们尽己所能，藉祂自己的恩赐，深入这些话语的奥秘。可能我们行事鲁莽，因为我们想讨论和查考天主的话。然而，它们为何被说出，岂非为了被认识？为何发出声响，岂非为了被听见？为何被听见，岂非为了被理解？那么，愿祂大大坚固我们，并按照祂认为值得的程度赐予我们一些恩惠；如果我们尚未深入泉源，就让我们饮用溪流。看，\Person{若望}自己像一条溪流般涌向我们，从高处将圣言传给我们。祂将其降低，并使之平易，使我们不再畏惧那崇高的，却能亲近那谦卑的。\par

13. 一定有一种真实的、强烈的意义，如果我们能领会的话，即\BibleQuote{圣父不审判任何人，但已把一切审判交给了圣子。}（\BibleRef{若 5:22}）这是因为在审判中，除了圣子，没有谁能显现给人类。圣父将隐藏，圣子将显现。圣子将以什么形式显现？以祂升天时的形式。因为祂在天主的形式中，与圣父一同隐藏；在仆人的形式中，显明给人。因此，不是\BibleQuote{圣父不审判任何人，但已把一切审判交给了圣子}？而是指公开的审判，在此公开审判中圣子将审判，因为祂将显现给那些受审判的人。\WorkTitle{圣经}更清楚地告诉我们，是圣子要显现。祂复活后第四十天，在门徒注视下升天；他们听见天使的声音：\BibleQuote{\Place{加里肋亚}人，}它说，\BibleQuote{你们为什么站着望天呢？这位被接升天的耶稣，你们看见祂怎样升天，也要怎样降来。}（\BibleRef{宗 1:3}）他们怎样看见祂去？在肉身中，他们触摸过、亲手检验过，甚至通过触摸证实了祂的伤；在那身体中，祂四十天与他们一同出入，真实地显现给他们，不是虚假；不是幻影、阴影或幽灵，而是如祂自己所说，没有欺骗他们：\BibleQuote{你们摸一摸，看看！因为鬼神没有肉和骨头，如同你们看见我有的一样。}（\BibleRef{路 24:39}）那个身体现在确实配得上天上的居所，不再受死亡支配，也不因岁月而改变。祂不是从婴儿长到那个年龄，再从壮年衰残到老年：祂保持升天时的样子，要来见那些祂曾愿意在祂来之前向他们传布圣言的人。因此，祂将带着人性形体而来，恶人将看见这形体；右边的和左边分开的都要看见：正如经上记载：\BibleQuote{他们要瞻望他们所刺透的。}（\BibleRef{匝 12:13}）如果他们瞻望他们所刺透的，他们就要看见被长枪刺穿的那个身体；因为长枪刺不透圣言。因此，恶人将能看见这个他们能伤害的身体。隐藏在身体内的天主，他们看不见：审判之后，那些在右边的人将看见祂。所以，这就是祂说\BibleQuote{圣父不审判任何人，但已把一切审判交给了圣子}的意思——即圣子将显现，在人的身体内显现为公开的审判；对右边的人说：\BibleQuote{你们蒙我圣父降福的，来承受国度吧！}和对左边的人说：\BibleQuote{你们到永火里去吧！那是为魔鬼和他的使者预备的。}（\BibleRef{玛 25:34}）\par

14. 请看，那人的形象将为虔敬者与恶人、义人与不义者、信徒与不信者、喜乐者与哀恸者、信赖者与蒙羞者所看见；看，它将被看见。当那形象在审判中出现，审判完结之后，经上说：圣父不审判任何人，但把全部审判交给了圣子，其原因正是：圣子在审判中将以那从我们取来的形象显现。此后将如何？何时才能看见那一切信徒渴望看见的天主的形象？何时才能看见那在起初、与天主同在的圣言、万物藉以造成的圣言？何时才能看见那宗徒所说的\BibleQuote{虽具有天主的形象，却不以自己与天主同等为抢夺}的天主的形象？\BibleRef{斐 2:6} 因为那形象是伟大的，在其中圣父与圣子的本质得以辨认；不可言喻，不可理解，尤其对小孩子而言。这形象何时才能被看见？看，右边是义人，左边是不义者；他们都同样看见那人，看见人子，看见那被刺透的，看见那被钉十字架的；他们看见那自卑的，那由童贞女所生的，犹大支派的羔羊。但他们何时才能看见那与天主同在的圣言？那同一位仍在那里，但仆人的形象将显现。仆人的形象将向仆人显示；天主的形象是为儿子们保留的。因此让仆人成为儿子吧；让右边的人进入旧约所应许的永恒基业，那基业是殉道者虽未看见却相信的，为了其应许他们毫不迟疑地倾流了鲜血；让他们去那里并在那里看见。他们何时去那里？让主亲自说：\BibleQuote{这些人要进入永火，但义人却要进入永生。}\BibleRef{玛 25:46}\par

15. 注意，祂提到了永生。祂是否告诉我们将在那里看见并认识圣父与圣子？如果我们永远活着，却看不见那圣父与圣子呢？听，在另一处，祂提到了永生，并说明了什么是永生：\BibleQuote{不要害怕；我不欺骗你们；我向那爱我的人所应许的不是无缘无故的，说：\InnerQuote{谁爱我的诫命并遵守它，他就是爱我；爱我的人必蒙我父所爱，我也要爱他，并将我自己显现给他。}}\BibleRef{若 14:21} 让我们回答主说：主，我们的天主，这是何等大事？何等大事？祢将祢自己显示给我们吗？那么，祢没有显示给犹太人吗？钉死祢的人不是看见祢了吗？但在审判时，当祢将我们放在祢右边时，祢将显现自己；那些站在左边的人不也看见祢吗？祢说\InnerQuote{将我自己显现给我们}是什么意思？我们现在不是看见祢在说话吗？祂回答说：我将以天主的形象显现自己；现在你们只看见仆人的形象。忠实的人，我不欺骗你们；相信你们将看见。你们爱，却尚未看见；难道爱本身不会引导你们看见吗？爱吧，坚持去爱；祂说，我不会辜负你们的爱，我净化了你们的心。因为我为什么净化了你们的心？不就是要你们看见天主吗？因为\BibleQuote{心里纯洁的人是有福的，因为他们要看见天主。}\BibleRef{玛 5:8} 那仆人仿佛与主争辩说：\Quote{但祢并没有说明，当祢说\InnerQuote{义人进入永生}时，祢没有说\InnerQuote{他们将看见我在天主的形象中，并看见与我同等的父。}} 请注意祂在另一处所说的：\BibleQuote{永生就是：认识祢，唯一的真天主，和祢所派遣的耶稣基督。}\BibleRef{若 17:3}\par

16. 接着，在所述审判之后，圣父不审判任何人，而将一切交给了圣子，此后将如何？后面是什么？\BibleQuote{为使众人尊敬圣子，如同尊敬圣父一样。}\BibleRef{若 5:23} 犹太人尊敬圣父，却蔑视圣子。因为圣子被视为仆人，圣父被尊为天主。但圣子将显现与圣父同等，使众人尊敬圣子如同尊敬圣父一样。因此，我们现在就拥有这信德。不要让犹太人说：\Quote{我尊敬圣父；与圣子何干？} 要回答他：\Quote{不尊敬圣子的，也不尊敬圣父。你处处说谎；你亵渎圣子，也亏负了圣父。因为圣父派遣了圣子，你却蔑视圣父所派遣的。你怎能尊敬派遣者，而亵渎被派遣者呢？}\par

17. 看，有人说，圣子被派遣了，而圣父更大，因为祂派遣了。远离血肉吧；旧人暗示时间上的陈旧。让那古老、永恒、无始、对你来说常新的，把你的理解从时间转向它。圣子因为被说成是被派遣的就更小吗？我听说是派遣，不是分离。但是，他说，在人中，我们看到派遣者比被派者大。就算是这样；但人间事物欺骗人，神圣事物洁净他。不要看人间事物，那里面派遣者显得更大，被派者更小；尽管如此，人间事物自身也反对你。例如，如果一个人想向一位女子求婚，由于不能亲自去做，就派一位朋友去替他求。还有许多情况，较小的被选去派较大的。那么，你为什么现在要因为这一个派遣、另一个被派遣，而提出吹毛求疵的反对呢？太阳发出光芒，但并不分离它；月亮发出光辉，但并不分离它；灯发出光，但并不分离它：我在那里看到的是发出，不是分离。因为如果你从人间事物中寻找例子，哦，异端的虚妄，尽管如我所说，连人间事物在某些情况下也驳斥你，并证明错误；然而，想想在人间事物中（你希望从中为神圣事物提取例子）是多么不同。一个人派遣别人时，他自己留在后面，只有被派的人前进。那个派遣的人与被他派的人一起走吗？然而，派遣圣子的圣父并没有离开圣子。听主亲自说：\Quote{看，时辰要来，那时人人都要回到自己那里去，你们要撇下我独自一人；但我不是独自一人，因为圣父与我同在。}\BibleRef{若 16:32} 祂怎么派遣了与祂同来的祂？祂怎么派遣了祂从未离开的祂？在另一处祂说：\Quote{住在我内的圣父做这些工作。}\BibleRef{若 14:10} 看，圣父在祂内，在祂内工作。派遣的圣父没有离开被派的圣子，因为被派者和派遣者是一体。\par

\SourceRecord{Translated by John Gibb.；Nicene and Post-Nicene Fathers, First Series；Vol. 7.；Edited by Philip Schaff.；Buffalo, NY: Christian Literature Publishing Co.,；1888.；<http://www.newadvent.org/fathers/1701021.htm>.}

% structure: p0001 source=h1 target=section reason=page-title

\section{\WorkTitle{讲道集 22（\BibleBook{若望}{福音} 5:24-30）}}

承昨日与前日的讲道，今日的福音选段接续而来。我们必须尽力依序阐释，并非依据其重要性，而是按我们能力所及；因为你们领受之时，并非依照那涌流之泉的丰沛，而是按你们有限的容量；而我们也在你们耳中宣讲，并非如泉水所倾注的那样多，而是将我们所能领受的，传达到你们心中——那事理本身在你们心里所结的果实，比我们在你们耳中所播的更为丰硕。所探讨的乃是伟大的事理，并非由大师，实由极为渺小之人；然而，那为我们的缘故、本为伟大却成为渺小的祂，赐予我们希望与信心。因为倘若我们未受祂鼓励，未蒙祂邀请去理解祂；倘若祂因我们卑微而遗弃我们，因为我们若未分享祂的死亡性，便无法分享祂的天主性，而祂若未分享我们的死亡性、来到我们中间向我们宣讲祂的福音；倘若祂不愿与我们分享我们身上卑微至极的部分——那么我们或许会想，那位承受了我们渺小的祂，并不愿意将祂自身的伟大赐予我们。我说这话，以免有人指责我们处理这些事过于冒昧，或对自己绝望，以为不能凭天主的恩赐理解天主子屈尊向我们所说的话。因此，祂既屈尊向我们说了，我们就应相信祂愿意我们理解。但我们若不明白，那未被求问便赐下祂圣言的祂，一旦被求问，必赐予领悟。\par

2. 看，祂话语中的这些奥秘是何等深奥。\Quote{我实实在在告诉你们：听我的话，并相信派遣我来的那位，便有永生。}我们岂不都在追求永生？祂说：\Quote{听我的话，并相信派遣我来的那位，便有永生。}那么，祂要我们听祂的话，却不要我们理解么？既然听并相信便有永生，那么理解岂不更是如此？然而，虔诚的行为是信德，信德的果实是理解，好使我们达至永生，那时将不再有福音为我们诵读；待一切读经的篇章和读经者、宣讲者的声音都消逝之后，那此时为我们分施福音的祂，将亲自显现给所有属于祂的人——他们此时以洁净的心和不朽的身体与祂同在，再不死去——祂洗净他们、光照他们，使他们现在活着并看见\BibleQuote{在起初已有圣言，圣言与天主同在}。因此，让我们此时省察我们是谁，并思忖我们听的是谁。基督是天主，祂正与人交谈。祂愿意他们领悟祂，愿祂使他们有能力领悟；祂愿意他们看见祂，愿祂开启他们的眼睛。然而，祂对我们说话并非无缘无故，而是因为祂所应许我们的乃是真实的。\par

3. \Quote{听我的话，}祂说，\Quote{并相信派遣我来的那位，便有永生，且不受审判，而是由死亡进入了生命。}我们何时由死亡进入生命，以致不受审判？在此生中，有从死亡到生命的过渡；在此生中——这还不是生命——有从此处由死亡到生命的过渡。这过渡是什么？祂说：\Quote{听我的话，并相信派遣我来的那位。}你遵循这些，相信并过渡。一个人站着时能过渡么？显然，身体站着，心思却过渡。他在哪里，从何处过渡，又过渡到哪里？他从死亡过渡到生命。看一个站着的人，这里所说的一切都可能在他身上发生。他站着，他听；也许他先前不信，因听而相信；片刻前他不信，此刻他信了；他仿佛从不信的区域，经过心灵的移动（而非身体的移动），进入了信德的区域——这是一种进入更佳之境的移动；因为那些再次背弃信德的人，是移入更糟之境。看，在此生中——正如我刚才所说，这还不是生命——有从死亡到生命的过渡，以致不受审判。但我为何说这还不是生命？如果这是生命，主就不会对某人说：\Quote{你若愿意进入生命，就当遵守诫命。}\BibleRef{玛 19:17}因为祂并非对他说：你若愿意进入永生；祂没有加上\Emph{永}字，仅仅说\Quote{生命}。因此，此生不应被称为生命，因为它不是真正的生命。什么是真正的生命？不就是永生么？请听宗徒对弟茂德所说：\Quote{你要嘱咐那些今世富足的人，不要心高气傲，也不要寄望于无常的财富，而要寄望于生活的天主，祂厚赐百物给我们享用；叫他们行善，在善事上富足，乐于施舍，甘心济贫。}他为何说这话？请听下文：\Quote{好叫他们为自己积存美好的根基，以备将来，能得着真正的生命。}\BibleRef{弟前 6:17}若他们应当为自己积存美好的根基以备将来，以便得着真正的生命，那么他们目前所处的生命无疑是虚假的生命。因为倘若你已经拥有了真正的生命，为何还渴望得着真正的？难道真正的生命还需要去得着？那么，必须离开虚假的。离开的路是什么？往哪里去？听、信，你便由死亡进入生命，且不受审判。\par

4. 这是什么意思，\BibleQuote{你不受审判}？\BibleRef{格后 5:10} 谁比宗徒保禄更好，他说：\BibleQuote{我们众人都必须出现在基督的审判台前，好使每个人按他在肉身所行的，或善或恶，领受报应。}保禄说：\BibleQuote{我们众人都必须出现在基督的审判台前}；而你竟敢向自己保证你不会受审判？愿这事远离我，你说，我竟敢向自己保证这事。但我相信那应许者。救主说话，真理应许，祂亲自对我说：\BibleQuote{凡听我的话，并相信派遣我来者的，就有永生，且出死入生，并不得进入审判。}于是，我听了我的主的话，并信了；所以如今，当我曾是\Emph{不信者}时，我成了信者；正如祂告诫我的，我出死入生，我不受审判；不是因我的狂妄，而是因祂的应许。然而，保禄说话与基督相反吗？仆人反对他的主，门徒反对他的师傅，人反对天主；以至于主说\BibleQuote{凡听见而相信的，出死入生}，宗徒却说\BibleQuote{我们众人都必须出现在基督的审判台前}？否则，如果出现审判台前的人不受审判，我不知道如何理解。\par

5. 上主我们的天主随后启示了这一点，并借着\WorkTitle{圣经}提醒我们，当提到审判时，应如何理解。我劝你们，因此，要留心。有时审判指的是惩罚，有时指的是区分。按那种用审判指区分的话法，\BibleQuote{我们众人都必须出现在基督的审判台前}，好使人\BibleQuote{能在那里按他在肉身所行的，或善或恶领受报应。}因为这本身是一种区分，将善物分给善人，将恶物分给恶人。因为如果审判总是被取为坏的意义，\BibleBook{}{圣咏}就不会说：\BibleQuote{审判我，天主。}或许有人听见某一句话说\BibleQuote{\Emph{审判我}，天主}时感到惊讶。因为人惯于说：\Quote{宽恕我，天主；}\Quote{怜悯我，天主。}谁说\BibleQuote{审判我，天主}？有时在\BibleBook{}{圣咏}中，这节经文甚至被放在停顿处，由读经者读出，由会众回应。难道它不会如此打动某些人的心，以致他害怕歌唱并对天主说\BibleQuote{审判我，天主}吗？然而会众满怀信心地歌唱，并不认为自己从神圣言语中学到的东西是在祈求坏事；即使他们不完全明白，也相信他们所唱的是好事。然而，\BibleBook{}{圣咏}本身并没有让人对它的意义一无所知。因为接着，它在接下来的话中显示了它所说的是哪种审判；不是定罪的审判，而是区分的审判。因为\BibleBook{}{圣咏}说：\BibleQuote{审判我，天主。}\BibleQuote{审判我，天主，并分辨我的案件与不圣洁的民族}是什么意思？那么根据这种分辨的审判，\BibleQuote{我们众人都必须出现在基督的审判台前。}然而，按照定罪的审判，祂说：\BibleQuote{凡听我的话，并相信派遣我来者的，就有永生，不得进入审判，而由死入生。}\BibleQuote{不得进入审判}是什么意思？就是不进入定罪。让我们从\WorkTitle{圣经}证明，\Emph{审判}这个词被用于指\Emph{惩罚}；尽管就在这段经文里，稍后你还会听到同一术语\Emph{审判}被用于仅指定罪和惩罚。然而，宗徒在某处写信给那些滥用身体的人，你们中的信友知道；因为他们滥用身体，受到了主鞭打的责罚。他对他们说：\BibleQuote{你们中间有许多人软弱、患病，并且长眠。}因此许多人甚至死了。他接着说：\BibleQuote{因为如果我们自己审判自己，就不会被主审判；}也就是说，如果我们自己责备自己，就不会被主责备。\BibleQuote{但当受审判时，我们受主的责罚，以免我们与世人一同被定罪。}因此，有些人在这里按惩罚受审判，以便在那里被宽免；有些人在这里被宽免，以便在那里更受折磨；还有一些人，他们承受惩罚本身而无惩罚的鞭打，如果他们没有因天主的鞭打而改正；这样，因为他们在这里藐视那责打人的父，他们在那里要经验到那惩罚的审判官。因此，有一种审判，即天主，也就是天主子，最终要将魔鬼及其使者，以及所有不信者和不虔敬者与他一同投入其中。这种审判，那如今相信并由死入生的人，将不会进入。\par

6. 因为，免得你以为因着信仰，你不会依照肉身死亡，或者，免得你按肉身理解，对自己说：\Quote{我的主曾对我说：\InnerQuote{谁听我的话，并相信那派遣我的，就已出死入生。}那么，我已相信，我就不必死了。}——要知道，你必会付清那债务，即死亡，这是因亚当的惩罚你所欠下的。因为，我们众人当时都在他内，他领受了这个判决：\BibleQuote{你必定要死。}（\BibleRef{创 2:17}）天主的判决不能失效。然而，在你偿付了旧人的死亡之后，你将被接纳进入新人的永生，并从死亡进入生命。同时，现在就开始生命的转化。你的生命是什么？信德：\Quote{义人因信德而生活。}不信的人呢？他们是死的。主所说的那个躯体内的死者中，就有这样的人：\BibleQuote{任凭死人埋葬他们的死人。}（\BibleRef{玛 8:22}）因此，即使在今生，也有死的和活的；在某种意义上，所有人都活着。谁是死的？那些没有相信的人。谁是活的？那些已经相信的人。宗徒对死者说了什么？\Quote{你这睡眠的，醒来罢。}但是，有反对者说，他说的是睡眠，不是死亡。请听下文：\Quote{你这睡眠的，醒来罢，并从死者中起来。}好像睡眠者会说：我要往哪里去？\BibleQuote{基督必要光照你。}（\BibleRef{弗 1:14}）基督光照了你，你现在相信了，立刻就从死亡进入生命：要持守你所进入的，你就不会进入审判。\par

7. 祂亲自解释了这一点，并接着说：\Quote{我实实在在告诉你们。}恐怕因祂说过\Quote{已出死入生}，我们便把这理解为将来复活的事，而祂愿意表明，相信的人已经过去了，那从死亡进入生命，就是从不信进入信德，从不义进入正义，从骄傲进入谦卑，从仇恨进入爱德，于是祂现在说：\Quote{我实实在在告诉你们，时候将到，且现在就是。}还有什么更明显的呢？\Quote{现在就是，死人将要听见天主子的声音，那听见的人就要活着。}我们已经谈过这些死者。我们想什么，我的弟兄们？在这人群中，难道没有死人听我说话吗？他们相信并按真信德行事，就活着，不是死的。但那些要么不信，要么像魔鬼那样信而战栗（\BibleQuote{魔鬼也信，且战栗。}）（\BibleRef{雅 2:19}）且生活邪恶，承认天主子却没有爱德的人，更应被视为死的。然而，这个时刻仍在进行。因为主所说的时刻，并不是一天十二个时辰中的一个。从祂说话时起，直到现在，甚至直到世界终结，同一个时刻正在过去；若望在他的书信中论及这个时刻说：\BibleQuote{孩子们，现在是最末的时刻了。}（\BibleRef{若一 2:18}）因此，现在是这个时候。谁是活的，就让他活；谁是死的，就让他活；让那躺在死中的，听见天主子的声音，起来并活着。主在拉匝禄的坟墓前呼喊，那死了四天的人就起来了。那个在坟墓里发臭的人，出来到了空气里。他被埋葬了，一块石头盖在他上面；救主的声音击碎了石头的坚硬；而你的心如此坚硬，天主的圣言还没有击碎它！在你心里起来，从你的坟墓里出来。因为你如同在坟墓里一样，躺在你心里，被恶习的重压如同石头一样压着。起来，出来。什么是\Quote{起来，出来}？相信并承认。因为相信的人已经起来了；承认的人已经出来了。为什么我们说承认的人已经出来了？因为他在承认之前是隐藏的；但当他承认时，他就从黑暗进入光明。他承认之后，对仆人们说了什么？在拉匝禄的尸体旁说了什么？\Quote{解开他，让他走。}如何？正如对祂的仆人宗徒们所说的：\BibleQuote{你们在地上所释放的，在天上也要被释放。}（\BibleRef{玛 18:18}）\par

8. \Quote{时候将到，且现在就是，死人将要听见天主子的声音，那听见的人就要活着。}他们将从什么泉源活着？从生命。从什么生命？从基督。我们如何证明泉源就是基督生命？祂说：\BibleQuote{我是道路、真理、生命。}（\BibleRef{若 14:6}）你愿意行走吗？\Quote{我是道路。}你愿意不被欺骗吗？\Quote{我是真理。}你愿意不死吗？\Quote{我是生命。}你的救主对你说：除了到我这里来，你没有地方可去；除了借着我，你没有路可走。因此，这个时刻现在正在进行，这件事正在明确发生，且永不停止。那些曾死的人，起来了；他们过渡到生命；在天主子声音中，他们活着；因着祂而活着，同时持守对祂的信德。因为圣子有生命，祂由此使相信的人得以活着。\par

9. 祂如何拥有呢？就如圣父拥有一样。请听祂自己说：\Quote{因为圣父怎样在自己内有生命，同样祂也将生命赐给圣子，叫圣子在自己内有生命。}弟兄们，我将尽我所能地讲。因为这些话正是那些使渺小悟性困惑的言语。祂为什么要加上\OriginalTerm{在自己内}{in Himself}呢？说\Quote{因为圣父怎样有生命，同样祂也将生命赐给圣子，叫圣子有生命}就已经足够了。祂加上了\Quote{在自己内}：因为圣父\Quote{在自己内有生命}，而圣子在自己内有生命。祂要我们在祂所说的\Quote{在自己内}这句话中明白某种含义。在这里，一个奥秘隐藏在这个词里；让我们敲门，好使门能打开。哦，主啊，祢说的是什么呢？祢为何加上了\Quote{在自己内}？难道祢使其复活的保禄宗徒没有生命吗？他有的，主说。至于那些已死的人要复活，以及因祢的话相信而进入生命的人：当他们进入以后，他们不是在祢内拥有生命吗？他们将拥有生命；因为不久之前我也说过：\Quote{谁听我的话，并相信派遣我来的那一位，就有永生。}因此，那些信祢的人拥有生命；而祢并没有说\Quote{在他们自己内}。但当祢论及圣父时：\Quote{就如圣父在自己内有生命}；再者，当祢论及祢自己时，祢说：\Quote{同样祂也将生命赐给圣子，叫圣子在自己内有生命。}就如祂拥有，祂也赐予拥有。祂在哪里拥有？\Quote{在自己内。}祂赐予在哪里拥有？\Quote{在自己内。}保禄在哪里拥有生命？不在他自己内，而在基督内。信友，你在哪里拥有？不在你自己内，而在基督内。让我们看看宗徒是否这样说：\Quote{我如今活着，但不是我，而是基督在我内活着。}\BibleRef{迦 2:20}我们的生命，作为我们自己的，也就是说，出于我们自己意志的生命，只会是邪恶、有罪、不义的；但我们内那良善的生命来自天主，而非我们自己；是由天主赐予我们，而非我们自己。但基督在自己内有生命，就如圣父一样，因为祂是天主的圣言。对祂而言，并非时而邪恶、时而良善地活着；但人却时而邪恶、时而良善地活着。那曾邪恶活着的人，是在他自己的生命里；那良善活着的人，已进入了基督的生命。你成为生命的分享者；你并非你所领受的那一位，而是接受者；但天主子并非如此，好像祂起初没有生命，然后领受了生命。因为如果祂这样领受了生命，祂就不会在自己内有生命了。因为\Quote{在自己内}究竟是什么意思呢？即祂自己就是那生命本身。\par

10. 我也许可以更清楚地说明这件事。有人点亮一支蜡烛：那支蜡烛，就那闪耀的小火焰而言——那个火在自己内有光；但你的眼睛，在蜡烛不在时原本闲置无见，现在也有了光，但不是在自己内。再者，如果它们转离蜡烛，它们就变暗；如果转向它，它们就被照亮。但无疑那火存在多久就照耀多久：如果你要从它取走光，你也同时熄灭了它；因为没有光它就不能存留。但基督是光，不可熄灭，与圣父同永恒，永远明亮，永远照耀，永远燃烧：因为如果祂不燃烧，难道会在圣咏中说\Quote{也没有能逃避祂的热气的}吗？但你因你的罪曾是冷漠的；你转向好得以温暖；如果你转离，你将变得冷漠。在你的罪中你曾是黑暗的；你转向好被光照；如果你转离，你将变成黑暗。因此，因为在你内你曾是黑暗，当你被光照时，你将\Emph{成为}光，却是在光内。因为宗徒说：\Quote{你们从前是黑暗，如今在主内却是光明。}\BibleRef{弗 5:8}当他说\Quote{如今却是光明}时，他加上了\Quote{在主内}。因此，在你内是黑暗，\Quote{在主内却是光明}。如何成了\Quote{光明}？因为藉分享那光，你是光。但如果你要从那光照你的光离开，你便回到你的黑暗。基督并非如此，天主的圣言并非如此。但如何不是呢？\Quote{就如圣父在自己内有生命，同样祂也将生命赐给圣子，叫圣子在自己内有生命；}如此祂活着，不是藉分享，而是永不改变，并且完全是生命本身。\Quote{同样祂也将生命赐给圣子。}就如祂拥有，祂同样赐予。区别何在？因为一位赐予，另一位领受。祂领受时是否已经存在？我们是否该理解基督曾一度没有光？而祂自己就是圣父的智慧，关于这智慧说：\Quote{它是永远光明的辉映。}\BibleRef{智 7:26}因此，所说的\Quote{赐给圣子}，就如说\Quote{生了圣子}；因为藉着生祂赐予了。正如祂赐祂存在，同样祂赐祂是生命，也赐祂在自己内有生命。在自己内有生命是什么意思？无需从别处获得生命，而是祂自己就是生命的满盈，其他相信的人从中藉此在活着时拥有生命。因此，\Quote{赐给祂在自己内有生命}。赐给谁？如同赐给祂自己的圣言，如同赐给那\Quote{在起初已有圣言，圣言与天主同在}的一位。\par

11. \Emph{后来}，因为祂成了人，祂给了祂什么？\Quote{并赐给祂行审判的权柄，因为祂是人子。}就祂是天主子而言，\Quote{就如父在自己内有生命，同样也赐给子在自己内有生命；}就祂是人子而言，\Quote{祂赐给祂行审判的权柄。}这就是我昨天向你们解释的，我亲爱的，即在审判时人将被看见，而天主将不被看见；但在审判之后，天主将被那些在审判中得胜的人看见，而被恶人却看不见。因此，既然在审判中，人将以祂升天时所具有的那一形象被看见，为此祂在上面已说过：\Quote{父不审判任何人，但将一切审判交与了子。}祂在这里也重复了同样的话，当祂说：\Quote{并赐给祂行审判的权柄，因为祂是人子。}正如你要说：\Quote{赐给祂行审判的权柄。}以什么方式？难道祂以前没有行审判的权柄吗？既然\Quote{在起初已有圣言，圣言与天主同在，圣言就是天主；}既然\Quote{万物是藉着祂造成的，}难道祂还没有行审判的权柄吗？是的，但按此，我说：\Quote{祂赐给祂行审判的权柄，因为祂是人子：}按此，祂接受了审判的权柄，\Quote{因为祂是人子。}因为就祂是天主子而言，祂一直拥有这权柄。被钉的那位接受了；曾处于死亡的那位，现在在生命之中；天主的圣言从未处于死亡，而一直在生命之中。\par

12. 如今，关于复活，我们中或许有人说：看，我们已经复活了；那听见基督、相信并已从死亡进入生命的人，也将不会进入审判。时辰来到，且现在就是，那听见天主圣子声音的人将要生活：他原是死的，他听见了；看，他复活了。那么，所说的将来还有复活，是什么意思？请省察自己，不要急于下结论，免得你匆匆跟随其后。的确，有现在发生的这种复活；不信者原是死的，不义者原是死的；义人活着，他们从不信之死进入信德的生命。但不要因此就相信将来不会有身体的复活；要相信也有身体的复活。因为请听在这因信德而有的复活宣告之后接着说的，免得有人认为这是唯一的复活，或者陷入那些败坏他人思想之人的绝望和错误，他们\Quote{说复活已过，}关于这些人，宗徒说：\Quote{他们颠覆了一些人的信德。}\BibleRef{弟后 2:18}因为我相信他们向人说了这样的话：\Quote{看，当主说：\InnerQuote{那信我的人已从死亡进入生命；}}复活已在相信的人——他们从前是不信者——身上发生了；怎么还能有第二次复活？感谢我们的主天主，祂支持动摇者，指引困惑者，坚定怀疑者。请听下面的话，如今你已没有理由为自己制造死亡的黑暗。如果你相信了，就相信全部。你说，我要相信什么全部？请听祂说什么：\Quote{你们不要惊奇这事，}即祂赐给子行审判的权柄。我说，在世界末日，祂说。如何是末日？\Quote{不要惊奇这事，因为时辰要来。}这里祂没有说\Quote{且现在就是，}因为那要在世界末日才来。\par

13. 你说：你从哪里向我证明祂所说的就是那复活本身？如果你耐心听，你很快就会亲自证明。那么让我们继续：\Quote{不要惊奇这事，因为时辰要来，那时所有在坟墓里的。}有什么比这复活更明显的？刚才祂没有说\Quote{在坟墓里的}，而是说：\Quote{死人将要听见天主圣子的声音；那听见的人将要生活。}祂没有说：有些人要生活，有些人要受罚；因为所有相信的人都要生活。但关于坟墓，祂说了什么？\Quote{所有在坟墓里的都要听见祂的声音，并且要出来。}祂没有说\Quote{要听见并生活。}因为如果他们邪恶地生活，躺在坟墓里，他们起来是为死亡，不是为生命。那么，让我们看看谁要出来。虽然，刚才死人因听见并相信而活了，那里没有作区分：没有说：死人要听见天主圣子的声音；当他们听见后，有些要生活，有些要受罚；而是\Quote{所有听见的都要生活：}因为相信的人要生活，有爱德的人要生活，他们中没有一个要死。但关于坟墓，\Quote{他们要听见祂的声音，并且出来：行过善的，到生命的复活；行过恶的，到审判的复活。}这就是审判，即那惩罚，祂刚才曾说过：\Quote{那信我的人已从死亡进入生命，并且不会进入审判。}\par

14. \BibleQuote{我凭自己什么也不能做；我听见了才审判，我的审判是公正的。}如果祢听见了才审判，祢听谁的话呢？如果是听圣父的，但确实 \BibleQuote{圣父不审判任何人，却把审判权全交给了圣子。}祢何时，作为圣父的传报者，宣告祢所听见的呢？我说我所听见的，因为圣父是什么，我就是什么：因为，说话确实是祂的职能；因为我就是圣父的圣言。基督向你们说这话。接着，关于你们自己。\BibleQuote{我听见了才审判}是什么意思，不就是 \Quote{我是}的意思吗？因为基督如何听见呢？弟兄们，我恳求你们，让我们考察一下。基督是听见圣父的话吗？圣父如何对祂说话呢？毫无疑问，如果圣父对祂说话，圣父就对祂使用言语；因为每个人对某人说话时，都是借着一个词说话。圣父如何对圣子说话呢，既然圣子是圣父的圣言？圣父对我们所说的一切，都是借着祂的圣言说的：圣父的圣言就是圣子；那么，圣父用什么别的词对圣言本身说话呢？天主是唯一的，只有一个圣言，在唯一圣言内包含一切。那么，\BibleQuote{我听见了才审判}是什么意思呢？正如我来自圣父，我也如此审判。所以 \BibleQuote{我的审判是公正的。}主耶稣啊，如果祢凭自己什么都不做——像属血肉的人所想的那样；如果祢凭自己什么都不做，那么祢刚才怎么说 \BibleQuote{圣子也赐生命给祂所愿意的人}呢？现在祢却说：我自己什么都不做。但圣子宣告的不就是祂来自圣父吗？那来自圣父的，就不是凭自己。如果圣子是凭自己，祂就不是圣子：祂是来自圣父的。圣父之所有，并非来自圣子；圣子之所有，来自圣父。与圣父平等；但仍是圣父的圣子，而非圣子的圣父。\par

15. \BibleQuote{因为我不寻求我的意愿，而是寻求派遣我者的意愿。}独生子说：\BibleQuote{我不寻求我的意愿，}可是人却渴望做自己的意愿！那与圣父平等的，竟这样谦卑自己；而躺在最底层、除非有人伸手就不能起来的，却如此高抬自己！那么，让我们做圣父的意愿、圣子的意愿、圣神的意愿；因为这位天主圣三有一个意愿、一个能力、一个尊威。然而，圣子说：\BibleQuote{我来不是为执行我的意愿，而是为执行派遣我者的意愿，}正是因为基督不是凭自己，而是来自圣父。但祂为了显现为人而拥有的，是取自祂自己所造的受造物。\par

\SourceRecord{Translated by John Gibb.；Nicene and Post-Nicene Fathers, First Series；Vol. 7.；Edited by Philip Schaff.；Buffalo, NY: Christian Literature Publishing Co.,；1888.；<http://www.newadvent.org/fathers/1701022.htm>.}

% structure: p0001 source=h1 target=section reason=page-title

\section{讲道集 23（若 5:19-40）}

1. 主在福音某处论及，祂话语的明智听众应当像一个人，他想要建造房屋，便深挖直到抵达磐石上稳固的根基，并在那里安全地建立他所建造的，以对抗洪水的冲击；这样，当洪水来临时，它反而被建筑物的力量击退，不会因其冲击力而毁坏那房屋。\BibleRef{玛 7:24} 让我们视天主的圣经为一片田地，我们想在其上建造什么。让我们不要怠惰，也不要以表面为满足；让我们深挖直到抵达那磐石：\BibleQuote{并且那磐石就是基督。} \BibleRef{格前 10:4}\par

2. 今日所读的段落向我们论及主的见证，祂并不需要人的见证，却有比人更大的见证；祂告诉了我们这见证是什么：\BibleQuote{我所做的工，}祂说，\BibleQuote{为我作证。}然后祂补充说：\BibleQuote{派遣我的圣父也为我作证。}祂所做的这些工，祂说也是从圣父领受的。因此，这些工作证，圣父作证。若翰没有作证吗？他确实作了证，但如同灯盏；不是为了满足朋友，而是为了挫败敌人：因为圣父很久以前就预言道：\BibleQuote{我为我的受傅者预备了一盏灯；我要使祂的敌人蒙羞，但我的圣化将在祂身上兴盛。}假如你在夜间被留在黑暗中，你注视着那灯，欣赏那灯，因它的光而欢欣。但那灯却说有一轮太阳，你应当因它而欢欣；尽管它在夜间燃烧，它却吩咐你期盼白昼。因此，并非不需要那人的见证。因为如果不需要他，为何派遣他呢？但是，反之，免得人停留在灯那里，以为灯的光足够他，因此主既没有说这灯是多余的，也没有说你们应当停留在灯那里。天主的圣经发出另一见证：天主无疑在那里为祂的圣子作了证，而犹太人在那圣经中寄托了他们的希望——即在天主的法律中，由祂的仆人梅瑟所传。\BibleQuote{你们查考圣经，}祂说，\BibleQuote{因你们认为其中有永生；正是这圣经为我作证；你们却不肯到我这里来得生命。}为什么你们以为在圣经中有永生？请问圣经为谁作证，并且明白什么是永生。而因为为了梅瑟的缘故，他们愿意拒绝基督，视祂为梅瑟律例和诫命的敌人，祂如同借着另一盏灯来谴责那些人。\par

3. 事实上，所有人都是灯盏，因为他们可以点燃，也可以熄灭。再者，当灯盏明智时，他们因圣神而发光闪耀；然而，如果它们曾燃烧而后熄灭，它们甚至会发臭。天主的仆人因祂仁慈的油而保持为好灯盏，并非凭自己的力量。天主的白白恩宠实在是灯盏的油。\BibleQuote{因为我比他们众人更劳碌，}某一盏灯说；免得他显得凭自己的力量燃烧，他补充说：\BibleQuote{然而不是我，而是天主的恩宠同我在一起。} \BibleRef{格前 15:10} 因此，在主来临之前，一切预言都是一盏灯。关于这灯，宗徒伯多禄说：\BibleQuote{我们有更确实的先知之言，你们在这话上留意，如同在黑暗中发光的灯盏，直到破晓，晨星在你们心中升起。} \BibleRef{伯后 1:19} 因此，先知们是灯盏，一切预言是一盏大灯。宗徒们呢？他们不也是灯盏吗？他们当然是。只有祂不是灯盏。因为祂并非被点燃又熄灭；因为\BibleQuote{父怎样在自己内有生命，也赐给子在自己内有生命。}宗徒们，我说，也是灯盏；他们感恩，因为他们既被真理之光点燃，又因爱德之灵燃烧，并借天主的恩宠之油得到供应。如果他们不是灯盏，主就不会对他们说：\BibleQuote{你们是世界的光。}因为祂说了\BibleQuote{你们是世界的光}之后，祂显明他们不应认为自己是那样的光——即如经上所说：\BibleQuote{那是真光，照亮进入这世界的每一个人。}但这是关于主说的，当时祂被与若翰洗者区分开。的确，关于若翰洗者，经上说过：\BibleQuote{他不是那光，而是要为光作证。} \BibleRef{若 1:9} 免得你说，他不是光，基督却称\BibleQuote{他是一盏灯}？——我回答，与那另一个光相比，他不是光。因为\BibleQuote{那是真光，照亮进入这世界的每一个人。}因此，当祂也对门徒说\BibleQuote{你们是世界的光}时，免得他们以为有什么属于他们是唯独归于基督的，从而灯盏被骄傲之风熄灭，当祂说了\BibleQuote{你们是世界的光}后，祂随即补充说：\BibleQuote{建在山上的城是不能隐藏的；人点灯，也不放在斗底下，而是放在灯台上，就照亮家中所有的人。}但倘若祂未称宗徒为灯，而是点灯的人，他们须把灯放在灯台上呢？听祂称他们自己为灯。\BibleQuote{照样，你们的光也当在人前照耀，}祂说，\BibleQuote{使他们看见你们的善行，光荣你们在天上的父。} \BibleRef{玛 5:14}\par

4. 因此，梅瑟为基督作证，若翰也为基督作证，所有其他先知和宗徒都为基督作证。在这一切见证之前，祂列上了自己工程的见证。因为通过那些人，也是天主而非其他为祂的圣子作证。然而，天主以另一种方式为祂的圣子作证。天主通过圣子自己启示祂的圣子，祂通过圣子启示自己。倘若一个人能够达到祂那里，他就不需要灯盏；并且通过真正地深挖，他会将他的建筑带到磐石上。\par

5. 弟兄们，今日的读经是容易的；但因为有昨日所亏欠的（因为我知道我延迟了，并非撤回了，而且主已恩准我今天仍能向你们说话），请回想你们应当要求什么，或许，在保持虔诚和有益的谦卑的同时，我们可以在某种程度上伸展自己，不是反对天主，而是向着祂，并举起我们的灵魂，将其倾注在我们之上，如同圣咏作者一般，他曾被问说：\Quote{你的天主在哪里？}\Quote{这些事，}他说，\Quote{我默想了，并将我的灵魂倾注在我之上。}因此，让我们将灵魂举向天主，而不是反对天主；因为也有一句话说：\Quote{上主，我向祢举起了我的灵魂。}并且让我们在祂的帮助下举起它，因为它是沉重的。它为何沉重呢？因为那朽坏的身体压重了灵魂，那地上的帐幕压抑了多虑的心思。\BibleRef{智 9:15}那么，让我们试试看，是否能使我们的心思从许多事物中抽离，以便专注于一件事，并提升到那唯一者（确实，正如我所说，除非那愿意我们灵魂被提升到祂那里的那一位帮助我们，我们无法做到）。这样，我们或许能在某种程度上领会：天主的圣言，圣父的独生子，与圣父同永生、同等的那一位，除了看见父所做的，什么也不做，而父自己除了通过子（子看见父所做的）什么也不做。因为主耶稣——在我看来，祂愿意在这里向留心的人启示一些伟大的事，并倾注给那些能够领受的人，同时又激发那些不能领会的人勤勉，以便他们虽然尚未理解，却可通过正直的生活成为能领受的人——已向我们表明：人的灵魂和理性心智（即在人内，不在兽内）除藉由天主的实体本身外，别无其他方式可得力量、光明和幸福；灵魂本身从身体获得一些东西，又使身体服从于它，而身体的感官可由形体事物得到抚慰和喜乐，并且由于灵魂与身体在此生中的这种结合和互相拥抱，当身体感官得到抚慰时，灵魂就喜乐，当它们受到冒犯时，灵魂就忧伤；然而，使灵魂本身幸福的幸福，除非藉分享那永远存在、不变的生命、那永恒的天主实体，无法实现；正如灵魂（低于天主）使身体（低于灵魂）生活，同样，唯有那高于灵魂的才能使那个灵魂幸福生活。因为灵魂高于身体，而天主高于灵魂。灵魂将其些微给予它的下级，同时从上级得到些微给予。让灵魂事奉它的主，免得被自己的仆役践踏。弟兄们，这就是基督宗教，它被传遍全世界，而它的敌人惊惶失措；他们在被打败的地方抱怨，在得胜的地方猛烈攻击。这就是基督宗教：敬拜独一天主，而非众神，因为只有独一天主能使灵魂幸福。灵魂因分享天主而幸福。并非因分享一个神圣的灵魂而使软弱的灵魂幸福，也非因分享一位天使而使神圣的灵魂幸福；但如果软弱的灵魂寻求幸福，让它寻求那使神圣灵魂幸福的那一位。因为你不是因一位天使而幸福，而是天使与你同由同一源头而幸福。\par

6. 既已预先阐明并坚固这些道理——即：理性灵魂唯独因天主而得福乐，身体唯独因灵魂而得生命，灵魂乃是介于天主与身体之间的某种存在——请你们将思想转向并同我一起回忆，不是今天读过的段落（关于此我们已经说得足够），而是昨天的段落；这三天我们一直在反复探讨、尽力挖掘，直至抵达那磐石。圣言基督，天主的圣言与天主同在，圣言与圣言天主，基督、天主与圣言是一位天主。向着这目标前进吧！啊，灵魂，藐视甚至超越万物，向着这目标前进！没有比被称为理性心智的受造物更强大、更崇高的了：超越此者，唯有造物主。但我说基督是圣言，基督是天主的圣言，圣言基督是天主；然而基督不仅仅是圣言，因为\BibleQuote{圣言成了血肉，居住在我们中间}：\BibleRef{若 1:14}；因此基督既是圣言又是血肉。因为当\BibleQuote{祂具有天主的形体时，不以自己与天主同等为僭越}。而我们这些卑微者，软弱且匍匐在地，无法上升到天主，难道就被遗弃了吗？绝对不可。\BibleQuote{祂空虚了自己，取了奴仆的形体}；\BibleRef{斐 2:6}因此并非丧失了天主的形体。祂本是天主却成了人，以接受祂所不是的，而非丧失祂所是的：如此，天主成了人。你们既有为你们软弱的，又有为你们成全的。让基督以祂的人性提升你们，以祂的天主性人领导你们，引导你们到达祂的天主性。这就是基督全部的道理和\OriginalTerm{安排}{dispensatio}，弟兄们，别无其他：为使灵魂复活，也使身体复活。因为两者都死了：身体因软弱而死，灵魂因邪恶而死。既然两者都死了，两者也都可复活。哪两者？灵魂与身体。那么，灵魂能靠什么复活，除非靠天主基督？身体能靠什么，除非靠人基督？因为在基督内也有一个人类的灵魂，整个灵魂；不仅有灵魂的非理性部分，还有那被称为心智的理性部分。因为曾有某些异端者被逐出教会，他们幻想基督的身体内没有理性心智，而只有像野兽那样的动物生命；因为没有理性心智，生命只是动物生命。但因他们被驱逐，并被真理驱逐，请接受完整的基督：圣言、理性心智与血肉。这就是完整的基督。让你的灵魂藉着天主的部分从邪恶中复活，你的身体藉着人的部分从朽坏中复活。至爱的诸位，请听我看来这段经文深奥之处：请看基督在此所说，祂来唯一的目的，就是使灵魂从邪恶中复活，身体从朽坏中复活。我已说过我们的灵魂凭何复活：凭天主的本质本身；我们的身体凭何复活：凭主耶稣基督人性的安排。\par

7. \BibleQuote{我实实在在告诉你们：子不能凭自己做什么，惟有看见父做什么，祂才做；因为父所做的事，子也照样做。}是的，天、地、海；天上、地上、海中的万物；有形无形的，地上的走兽、田间的植物、水中的游鱼、空中的飞鸟、天上的星辰；此外还有天使、德能、宝座、宰制、率领、掌权者：\BibleQuote{万有都是藉着祂造的。}难道天主造了这一切，然后给子看造好的，让子也造一个充满这一切的另一个世界吗？当然不是。相反，祂说什么？\BibleQuote{因为凡父所做的事，}不是别的，\BibleQuote{子也照样做，}不是不同地，\BibleQuote{而是照样。} \BibleQuote{父爱子，将祂自己所作的一切事指给祂看。}父指给子看，为使灵魂复活，因为灵魂是由父和子复活的；灵魂除非以天主为生命，否则不能生活。既然如此，灵魂除非以天主为生命就不能生活，正如灵魂是身体的生命；那么父指给子看，即父所做的，祂是藉着子做的。因为父不是先做然后指给子看，而是藉着指示，藉着子而做。因为子看见父在做任何事之前就已指示；从父的指示和子的看见，父藉着子所做的事就完成了。如此，灵魂若能看到那合一之\OriginalTerm{联结}{coniunctio unitatis}——父指示、子看见，受造物由父的指示和子的看见而成，并且那由父的指示和子的看见而成的东西，既不是父也不是子，而是低于父和子，即凡由父藉着子所造之物——就能复活。谁看见这个呢？\par

8. 看哪，我们再次自谦于血肉的想象，并降卑到你们的层次，如果我们曾从你们那里有所提升的话。你希望向你的儿子展示某事，使他能做你所做的事；你将要去做，并以此展示那事。因此，你所要做的事，为向你的儿子展示它，你当然不是借你的儿子去做；而是你独自做那件事，当做成后，他可以看到，并同样地再做另一件类似的事。但在这里并非如此；你为什么固守你自己的相似，而抹去你内天主的相似呢？在那里，情况完全不同。找一个情形，你向你的儿子展示你在做之前所做的事；这样，在你展示之后，它将由你的儿子完成。或许现在你想到类似的事：\Quote{看，你说，我想造一所房子，并希望它由我儿子建造；在我亲自建造之前，我向我儿子指出我打算做的事；两者他都做，我也借着他——我向他指出我意愿的那位——来做。}你确实从先前的相似中退却了，但你仍陷于极大的不似之中。因为，看，在你造房子之前，你告知你的儿子，并向他指出你打算做的事；这样，在你建造之前展示之后，他就可以制造你所展示的，于是你便借着他来制造；但你将对你的儿子说一些话，这些话语将在你与他之间传递；在展示者与观看者之间，在说话者与听者之间，飞扬着清晰可辨的声音，这声音既不是你所是的，也不是他所是的。那声音，确实，从你口中发出来，通过空气的振动触及你儿子的耳朵，并充满听觉，将你的思想传达给他心中；那声音，我再说，不是你自己，也不是你的儿子。一个记号从你的心灵给出给你儿子的心灵，但那记号既不是你的心灵，也不是你儿子的心灵，而是别的东西。难道我们以为圣父如此对圣子说话吗？在圣父与圣言之间竟有言语吗？那么是如何的呢？或者，无论圣父要对圣子说什么，若祂借一个言语来说，圣子本身就是圣父的圣言，祂会借一个言语对圣言说话吗？或者，既然圣子是伟大的圣言，圣父与圣子之间还需要传递更小的言语吗？难道是这样，某种如同暂存、飘逝受造物的声音，必须从圣父的口中发出，并撞击圣子的耳朵？天主有身体，以致这声音如同从祂嘴唇发出吗？圣言也有肉体的耳朵，声音可以进入吗？放下所有形体的观念，注视单纯，如果你心志专一。但你如何才能心志专一？如果你不纠缠于世界，而是从世界解脱出来。因为通过解脱，你将变得心志专一。并且，如果你能，就看看我所说的；如果不能，就相信你所看不见的。你对你的儿子说话；你借着一个言语说话；你，和你的儿子，都不是那发声的言语。\par

9. 你说，我有另一种展示的方式；因为我的儿子受过如此良好的教导，他无需我说话便能听见，但我通过点头向他示意该做什么。看，你随意用点头向他示意，然而心灵确实在自身内持有它所要展示的。你用什么来点头呢？用身体——即用嘴唇、表情、眉毛、眼睛、手。这一切都不是你的心灵本身；它们也是媒介；通过这些标记，有某种东西被理解，而这东西既不是你的心灵，也不是你儿子的心灵；但你用身体所做的一切都在你的心灵之下，也在你儿子的心灵之下；你的儿子也无法知道你的心灵，除非你通过身体给他标记。那么，我说什么呢？在那里并非如此；在那里一切都是单纯的。圣父向圣子显示祂所做的事，并通过显示而生了圣子。我明白我所说的；但因为我也看出是对谁说的，愿这种领悟迟早在你们内形成，好使你们把握它。如果你们现在不能领悟天主是什么，至少领悟天主不是什么：如果你们认为天主不是与祂自身不同的某物，你们就会大有进步。天主不是身体，不是大地，不是天空，不是月亮，或太阳，或星辰——不是这些有形之物。因为如果不是天上的事物，更何况地上的事物呢！将一切身体置于问题之外。此外，请听另一件事：天主不是可变的灵。因为我承认——且必须承认，因为是福音在说——\BibleQuote{天主是神}。但要超越一切可变的灵，超越一切现在知道、现在不知道的灵；现在记得、现在忘记的灵；愿意以前所不愿的、不愿以前所愿意的灵；或现在经历这些变化或可能经历这些变化的灵：超越这一切。你在天主内找不到任何可变性；也找不到任何东西以前是某种样子，而现在则不同。因为凡发现有转变之处，那里就有某种死亡发生了：因为事物不再是它以前的样子，就是一种死亡。据说灵魂是不朽的；确实如此，因为它永远活着，其内有某种持续的生命，但这生命是可变的。根据这生命的可变性，它可被称为有死的；因为如果它曾明智地生活，然后变得愚拙，它就为更坏而死了；如果它曾愚拙地生活，然后变得明智，它就为更好而死了。因为圣经教导我们，有一种为更坏的死亡，也有一种为更好的死亡。无论如何，那些为更坏而死的人，就是经上所说的：\BibleQuote{让死人埋葬他们的死人} \BibleRef{玛 8:22}，以及\BibleQuote{你这睡着的人，醒来，从死人中起来，基督必光照你} \BibleRef{弗 5:14}，以及出自我们面前的这段经文：\BibleQuote{死人将要听见，听见的人将要活着}。因为他们是为更坏而死的，所以他们又活过来。通过活过来，他们为更好而死，因为通过复活，他们将不再是以前的样子；但成为以前所不是的，那就是死亡。但若为更好，也许就不能称为死亡？宗徒却称那是死亡：\BibleQuote{你们若是与基督同死，脱离了世上的元素，为什么还像仍在世上活着那样判断世事呢？} \BibleRef{哥 2:20}，以及：\BibleQuote{因为你们已经死了，你们的生命与基督一同藏在天主内}。他愿意我们死去，好使我们活着，因为我们活着是为了死去。因此，无论什么死去，无论是从更好到更坏，还是从更坏到更好，都不是天主；因为至善不能变得更好，真正的永恒也不能变得更好。因为真正的永恒，是没有时间的地方。但若现在是这样，现在是那样呢？一旦时间被引入，它就不是永恒的。为使你们知道，天主并非如此，不像灵魂那样——灵魂确实是不朽的——然而天主的宗徒却说：\BibleQuote{唯独祂具有不朽}，除非他公开说，唯独祂具有不变性，因为唯独祂拥有真正的永恒？因此，那里没有可变性。\par

10. 在你自身内认出我愿在内在说的东西；不是像在身体内那样内在，因为某种意义上可以说\Quote{在你自身内}。因为在你内有健康、你的年龄（不论如何），但这是就身体而言。在你内有你的手和你的脚；但有一件事在你内，是内在的；另一件事在你内，如同在你的衣服内。但将你的衣服和你自己放在外面，下到你自己内，去到你的隐秘处，你的心灵，在那里，如果你能，看看我要说什么。因为如果你远离你自己，你怎能接近天主？我曾谈论天主，你以为你会明白。我现在谈论灵魂，谈论你自己：明白这个，我将在那里考验你。因为当我打算从你自己的心灵给你一些与你的天主相似的比喻时，我并不远求实例；因为人并非按身体，而是按那同一个心灵，按天主的肖像受造的。让我们在天主自己的肖像中寻求天主；让我们在祂自己的形象中认出造物主。在那里，内在，如果我们能，让我们找到我们所谈论的——圣父如何向圣子显示，以及圣子如何看见圣父所显示的，在圣父通过圣子创造任何事物之前。但当我已说了，而你们已明白时，你们不可认为所谈论的正如我们的比喻一样，好使你们在其中保持虔敬，这是我愿你们保持的，并恳切劝戒你们保持的：就是说，如果你们不能领悟天主是什么，不要以为知道祂不是什么是一件小事。\par

11. 请看，在你心中，我见到两样东西：你的\OriginalTerm{记忆}{memory}和你的\OriginalTerm{思想}{thought}，它们仿佛是你灵魂的观看能力与\OriginalTerm{视觉}{vision}。你看见某物，并用双眼察觉它，并把它托付给记忆保管。在那里，在内里，是你所托付给记忆的东西，如同藏在仓库里、宝库里、如同某种密室和内室。你想着别的事，你的注意力在别处；你所看见的虽在你的记忆里，却不为你所见，因为你的思想转向了另一件事。我立刻证明这一点。我对你们这些认识的人说话；我提名叫\Place{迦太基}；凡认识它的人立刻就在心中看到了迦太基。难道有同你们的心一样多的迦太基吗？你们都是借着这个名字、借着这些为你们所熟悉的音节（从我口中涌出）而看见了它：你们的耳朵被触动；灵魂的感官经由身体被触动，心神从另一对象转向这个词，就看见了迦太基。迦太基是在那时那地被造出来的吗？它早已在那里，只是隐藏在记忆里。为什么隐藏在那里？因为你的心专注于别的事；但当你思想转回记忆中的东西时，便从中形成了一种心神的视觉。先前没有视觉，但存有记忆；视觉是由思想回转于记忆而形成的。于是，你的记忆就将迦太基显示给你的思想；那先前在你内的东西（当你将心神指向记忆之前），当你思想转向它时，它就将之呈现给你思想的注意力。看，一种展示由记忆完成，一种视觉在思想中产生；其间没有言语，没有从身体发出的信号：你既不点头，也不写字，也不发声；而思想却看见了记忆所展示的东西。但那展示的与它所展示的对象，都属同一实体。然而，为使你的记忆能有迦太基在里面，那影像经由双眼被引入，因为你看见了你所储存于记忆的东西。同样，你看见了你所记忆的树，同样，山、河；同样，朋友、敌人、父亲、母亲、兄弟、姐妹、儿子、邻居的面容；同样，书中所写的字句、书本身；同样，这个教堂：这一切你看见了，在看见后托付给记忆；并且，如同在那里储存了你可以随意借着思想看见的东西，即使它们不再在肉体的眼前。你在迦太基时看见了迦太基；你的灵魂经由眼睛接受了影像；这影像被储存在你的记忆里；而你，这位曾在迦太基的人，在你内保存了某种东西，使你在即使不在那里时也能与自己一同看见它。这一切你都是从外部接收的。圣父所展示给圣子的，圣子并非从外部接收：一切都发生在内里，因为外间根本不会有任何受造物，除非圣父藉圣子造了它。每一受造物都由天主所造；在被造之前不存在。因此，并非在受造并被保存于记忆之后，圣父才把它展示给圣子，如同记忆向思想展示一样；相反，圣父展示它以使它被造，圣子看见它被造成；并且圣父藉展示而创造了它，因为他是藉圣子的看见而造的。因此，我们不应惊讶于经上说：\BibleQuote{但他看见父所作的，}而非\Emph{展示}。因为由此暗示，在父那里，\Emph{作}与\Emph{展示}是同一件事；从而我们可以明白，他是藉圣子的看见而行一切事。那\OriginalTerm{展示}{showing}与那\OriginalTerm{看见}{seeing}都不是时间性的。因为所有时间都由圣子所造，它们自然不可能在时间中的某一点被展示给他去造。但圣父的展示生出圣子的看见，正如圣父生出圣子一样。因为展示产生看见，而非看见产生展示。如果我们能更纯粹、更完美地审视此事，或许我们会发现：圣父不是一物，他的展示是另一物；圣子也不是一物，他的看见是另一物。但如果我们几乎不能领悟这一点——如果我们几乎不能解释记忆如何向思想展示它从外部接收的东西——那么，我们如何能领会或解释圣父如何向圣子展示那不从别处来的、或者说那不外乎他自身的呢！我们只是小子：我告诉你们天主不是什么，并未给你们显示天主是什么。那么，我们该做什么才能领悟他的所是？你们能靠我或藉我做到吗？我对这些小子说话，对你们也对我自己说；有一位我们可以藉着他：我们刚才歌唱过、刚听过：\BibleQuote{将你的挂虑托给上主，他必扶持你。}你之所以不能，人啊，是因为你是个小子；既是小子，就需被喂养；被喂养了，你就长成；当初作为小子所不能的，当你长成时就会看见；但为使你能被喂养，\BibleQuote{把你的挂虑托给上主，他必扶持你。}\par

12. 因此现在让我们简要地过一遍剩下的内容，你们要看主如何使我们明白我在此向你们所推荐的那些事。\Quote{圣父爱圣子，并将自己所作的一切事指示给祂。}祂自己复活灵魂，但藉着圣子，好使复活的灵魂能享见天主的本体，即圣父和圣子的本体。\Quote{祂还要指示祂比这些更大的事业。}比哪些更大？比治愈身体更大。我们已经讨论过这一点，现在不必再停留。身体的复活归于永恒，比这治愈那瘫子的身体（只能持续一时）更大。\Quote{祂还要指示祂比这些更大的事业，使你们惊奇。} \Quote{祂还要指示}，好像是暂时性的行动，因此是对在时间中受造的人说的，因为天主圣言不是受造的，万世万物是由祂造成的。但基督在时间中成为了人。我们知道在哪个执政官年份，童贞玛利亚因圣神受孕而生了基督。因此，祂作为天主创造了时间，却作为人在时间中受造。因此，正如在时间中，\Quote{祂要指示祂更大的事业}，即身体的复活，\Quote{使你们惊奇}于圣子所行的身体复活。\par

13. 然后祂回到灵魂的复活：\Quote{因为父怎样复活死者，并使他们活起来，子也照样任意复活祂所愿意的人；}但这是按圣神说的。父使活起来，子也使活起来；父复活祂所愿意的，子复活祂所愿意的；但父复活的和子复活的是同一些人，因为万有都是藉着祂造的。\Quote{因为父怎样复活死者，并使他们活起来，子也照样任意复活祂所愿意的人。}这是论及灵魂的复活；但身体的复活呢？祂回来又说：\Quote{因为父不审判任何人，却将一切审判都交给了子。}灵魂的复活是藉着父和子永恒不变的本体成就的。但身体的复活是藉着圣子的人性的安排成就的，这安排是暂时的，不是与父同永恒的。因此，当祂提到审判，即在其中有身体的复活时，祂说：\Quote{因为父不审判任何人，却将一切审判都交给了子；}但关于灵魂的复活，祂说：\Quote{父怎样复活死者，并使他们活起来，子也照样任意复活祂所愿意的人。}那么，这是父和子共同的工作。但关于身体的复活：\Quote{父不审判任何人，却将一切审判都交给了子，为叫众人尊敬子，如同尊敬父一样。}这是指灵魂的复活。\Quote{为叫众人尊敬子。}怎样？\Quote{如同尊敬父一样。}因为子以与父同样的方式成就灵魂的复活；子使活起来，如同父一样。因此，在灵魂的复活中，\Quote{叫众人尊敬子如同尊敬父一样。}但关于身体的复活的尊敬呢？\Quote{那不尊敬子的，就是不尊敬派遣祂来的父。}祂没有说\Quote{如同}，而是说\Quote{尊敬}和\Quote{尊敬}。因为那作为人的基督受到尊敬，但不是如同父天主那样。为什么？因为关于这一点，祂说过：\Quote{父比我大。}（\BibleRef{若 14:28}）子何时受到如同父一样的尊敬？当\Quote{圣言在起初就与天主同在，万物是藉着祂造成的}。因此，在第二种尊敬中，祂说什么？\Quote{那不尊敬子的，就是不尊敬派遣祂来的父。}子不是被派遣的，而是因为祂成为了人。\par

14. \Quote{我实实在在告诉你们：}祂又回到灵魂的复活，藉着反复提醒，使我们能领会祂的意思，因为我们跟不上祂那如飞絮般急迫的讲论。看，天主的圣言在我们这里停留；看，祂似乎与我们软弱的人同居。祂再次回到论及灵魂的复活。\Quote{我实实在在告诉你们：那听我的话，并相信派遣我来的那位，就有永生；}但这永生是从父而来的。\Quote{因为那听我的话，并相信派遣我来的那位，就有永生}，因相信派遣子来的父而得的永生。\Quote{并且不受审判，而已出死入生。}但使祂活起来的是父，是祂所相信的父。怎么，你不使活起来吗？看，子也\Quote{任意复活祂所愿意的人。} \Quote{我实实在在告诉你们：时候要到，凡在坟墓里的，都要听见天主子的话，凡听见的就要活起来。}这里祂没有说\Quote{他们相信那派遣我来的，因此就要活；}而是说：藉着听见天主子的话，\Quote{凡听见的}，即凡听从天主子的话的，\Quote{就要活起来。}因此，他们既从父得生命，当他们相信父时；又从子得生命，当他们听见天主子的话时。为什么他们既从父又从子得生命？\Quote{因为父怎样在自己内有生命，照样也赐给子在自己内有生命。}\par

15. 祂已经讲完了灵魂的复活；接下来要更清楚地讲论身体的复活。\BibleQuote{并且赐给祂行审判的权柄}，不只是凭信德和智慧复活灵魂，而是要行审判。但为何如此？\BibleQuote{因为祂是人子。}因此，圣父通过人子做一些事，这些事不是出于祂自己的本性（圣子与此本性相等），例如：祂应被诞生、被钉十字架、死亡并复活；因为这些没有一样是圣父所经历的。同样，身体的复活也是如此。因为灵魂的复活是圣父出于自己的本性，通过圣子的本性（圣子与祂相等）而实现的；因为灵魂分享了那不变的光，但身体没有；但身体的复活，圣父是通过人子实现的。因为\BibleQuote{祂赐给祂行审判的权柄，因为祂是人子}，依照祂上面所说的：\BibleQuote{因为圣父不审判任何人。}为了表明祂所说的是关于身体的复活，祂接着说：\BibleQuote{你们不要惊奇这事，因为时辰来到}——不是\Emph{且现在是}，而是\BibleQuote{时辰来到，凡在坟墓里的（你们昨天已经足够听了解释）都要听见祂的声音，就出来。}到哪里去？受审判：\BibleQuote{行过善的，复活进入生命；行过恶的，复活进入审判。}而祢独自做这事吗？因为圣父把一切审判都交给了圣子，且不审判任何人？祂说：是的，我做的。但祢如何做这事？\BibleQuote{我由自己不能做什么；我怎样听见，就怎样审判；我的审判是正义的。}当祂论及灵魂的复活时，祂没有说\Emph{我听见}，而是说\Emph{我看见}。因为\Emph{我听见}指圣父发命令。因此，如今作为人，正如祂比圣父小；出自仆人的\OriginalTerm{形体}{form}，而非天主的\OriginalTerm{形体}{form}：\BibleQuote{我怎样听见，就怎样审判；我的审判是正义的。}为何人的审判是正义的？我的弟兄们，请留意：\BibleQuote{因为我不寻求自己的旨意，只寻求那派遣我者的旨意。}\par

\SourceRecord{Translated by John Gibb.；Nicene and Post-Nicene Fathers, First Series；Vol. 7.；Edited by Philip Schaff.；Buffalo, NY: Christian Literature Publishing Co.,；1888.；<http://www.newadvent.org/fathers/1701023.htm>.}

% structure: p0001 source=h1 target=section reason=page-title

\section{\WorkTitle{讲道集第24篇}（\BibleRef{若 6:1-14}）}

我们的主耶稣基督所行的奇迹，确实是神圣的工作，激发人的心智从可见的事物提升到对天主的认识。然而，由于祂并非可以用肉眼看见的那种实体，而祂在治理整个世界和管理整个受造界时所行的奇迹，因着它们惯常的恒常性而被轻视，以至于几乎没有人肯思考天主在每一粒种籽中所显示的奇妙惊人的工作；祂按其仁慈，为自己保留了一些超越通常进程和自然秩序的工作，在合宜的时候施行，好叫那些轻视祂日常工作的人，因着观看这些——并非更大、却是罕见的工作——而惊愕。诚然，治理整个世界比用五个饼使五千人吃饱是更大的奇迹；然而无人对前者感到惊奇；但对后者人们感到惊奇，不是因为它更大，而是因为它罕见。因为现在还有谁喂养整个世界，岂不是那从几粒种籽创造庄稼的那位吗？祂因此像天主那样创造。因为，祂从几粒种籽使田地的出产增多的同一来源，祂也在祂手中使五个饼增多。德能确实在基督的手中；但那五个饼犹如种子，并非交托给土地，而是由创造土地的那位增多的。因此，在这奇迹中，有那被带到感官面前、使心智被唤醒而注意的事物，有那呈现在眼前、使理智被操练的事物，好叫我们藉着祂可见的工作赞美不可见的天主；并藉着信德被提升、被信德所净化，我们或许渴望以不可见的方式看见那不可见者，我们藉着可见的事物认识祂。\par

2. 然而，在基督的奇迹中只观察到这些是不够的。让我们质问奇迹本身，它们关于基督告诉我们什么：因为它们有自己的言语，如果它们能被理解的话。因为基督本身就是圣言，圣言的行动对我们而言也是一个言语。因此，关于这个奇迹，既然我们听到了它多么伟大，让我们也探究它多么深奥；让我们不仅以它的表面为乐，也让我们寻求认识它的深处。这个我们在外面赞叹的奇迹，里面有其含义。我们看见了，我们注视了某种伟大、荣耀且全然神圣、只能由天主完成的事：我们因工作赞美了工作者。但是，正如如果我们某处看到优美的书法，仅仅赞美书写者的手——因为他写出的字母均匀、整齐、优雅——是不够的，如果我们不也阅读他通过那些字母传达给我们的信息；同样，仅仅审视这一作为的人，或许因其美丽而喜悦，以赞美工作者；但理解的人，却好像是阅读它。因为一幅画被观看的方式与文字被观看的方式不同。当你看到一幅画，看到并赞美它就是全部；当你看到文字时，这还不是全部，因为你还被提醒要去阅读它。此外，当你看到文字时，如果你碰巧不能阅读，你会说：\Quote{我们想这里写的是什么？}你问它是什么，而你已经看见它是什么。那个你请教你所看见的是什么的人，会向你显示另一件事。他有与你不同的眼睛。难道你不是同样看见字母的形状吗？但你不是同样理解这些符号。你看见并赞美；但他看见、赞美、阅读并理解。因此，既然我们已经看见并赞美，让我们也阅读并理解。\par

3. 主在山上：更让我们理解，主在山上就是圣言在高处。因此，在山上所做的事，并非如同低处的，也不可草率放过，而应向上仰望。祂看见群众，知道他们饥饿，仁慈地喂养了他们：不仅凭着祂的良善，也凭着祂的德能。因为单凭良善有何用？那里没有饼来喂养饥饿的群众。难道德能不伴随良善吗？那样群众就会饿着肚子留下。总之，与主同在的门徒们也饥饿，他们自己也愿意喂养群众，以免他们空着肚子离开，但他们没有东西喂养他们。主问，从哪里可以买饼来喂饱群众。圣经说：\BibleQuote{但他说这话，是为试探他；因为他自己知道要怎样行。} 也就是，试探祂所问的门徒斐理伯。那么，试探他有什么益处呢？除非要显示门徒的无知？而且，也许在显示门徒的无知时，祂预示了更多的东西。那么，当五个饼的圣事本身开始向我们说话并表明其意义时，这一点就会显明：因为我们会在那里看到，为何主在此行动中愿意通过问祂自己所知道的事来显示门徒的无知。因为我们有时会问我们所不知道的事，愿意聆听以学习；有时我们会问我们知道的事，想看看我们所问的人是否也知道。主两者都知道；知道祂所问的事，因为祂知道祂自己要做什么；祂也同样知道斐理伯不知道这事。那么，祂为什么问呢？却是要显示斐理伯的无知。而祂为什么这样做，我们将如我所说，后来才会明白。\par

4. \Person{安德肋}说：\Quote{这里有一个孩童，带着五个大麦饼和两条鱼，但为这么多的人，这些算得什么？}当\Person{斐理伯}被问到时，曾说两百银币的饼也不够叫这许多人各吃一点，那里有一个孩童，带着五个大麦饼和两条鱼。\Quote{\Person{耶稣}说：\InnerQuote{叫众人坐下。}那里有许多青草，他们便坐下，人数约有五千。\Person{主耶稣}拿起饼来，谢了恩，}吩咐将饼掰开，摆在坐下的人面前。不再是五个饼，而是祂所添加的，那创造增多的那位所加的。\Quote{鱼也照样，随他们所要。}众人吃饱了还不够，还剩下碎块；这些碎块被吩咐收拾起来，免得糟蹋了：\Quote{他们便将那五个大麦饼的碎块，就是众人吃了剩下的，装满十二个篮子。}\par

5. 简而言之：五个饼预表\WorkTitle{梅瑟五书}；它们不是小麦饼而是大麦饼，这很恰当，因为它们属于旧约。你知道大麦的结构，我们很难得其精粹；因为精粹被一层外壳包着，且外壳坚韧而紧贴，需费力才能剥去。旧约的文字也是如此，披着属肉体的圣事的外衣；但如果我们得其精粹，它便滋养并满足我们。于是，有一个孩童带来了五个饼和两条鱼。若问这孩童是谁，也许他就是以色列子民，以一种幼稚的方式，只是携带着，却没有吃。因为他们所携带的东西，封闭时是重担，打开后却成了滋养。至于两条鱼，似乎预表旧约中两个崇高的人物，即司铎与首领，他们被受傅以圣化并治理百姓。最后，在奥秘中，那位被这些人所预表的亲自来了；那位由大麦精粹所指向、却为外壳所隐藏的，终于来了。祂来了，以祂一个位格兼具司铎和首领双重身份：作为司铎，祂将自己作为祭品献给天主，为我们；作为首领，我们受祂治理。那些原来封闭携带的，如今已开启。愿感谢归于祂。祂亲自实现了旧约所预许的。祂吩咐将饼掰开；在掰开中它们增多。再没有比这更真实的了。因为当\WorkTitle{梅瑟五书}被诠释时，它们藉着被掰开——即被开启和阐述——产生了多少书卷呢？但正因为在那大麦中，起初子民的无知被遮蔽了，关于这子民经上记载：\BibleQuote{当诵读梅瑟时，帕子还盖在他们心上；}\BibleRef{格后 3:15}因为帕子尚未除去，因为基督尚未到来；当基督悬在十字架上时，圣殿的帐幔尚未撕裂；我之所以说，因为子民的无知系于法律，因此主藉此印证显明了门徒的无知。\par

6. 因此，没有什么是无意义的；一切都有意义，但需要能理解的人：就连这吃饱的人数，也预表了在法律之下的子民。为何是五千人？因为他们是在法律之下，这法律在\WorkTitle{梅瑟五书}中展开。为何病人被放在五个廊子那里，却没有得痊愈？然而，祂在那里治愈了那个患病的人，就是此处用五个饼使群众吃饱的那位。此外，他们坐在草地上；因此是按肉体理解的，并安息于肉体。\BibleQuote{凡有血肉的都如草。}\BibleRef{依 40:6}而那些碎块是什么？岂不是人们无法吃的吗？我们理解它们是一些更隐密的意义，是群众无法领会的。那么，除了将这些群众无法领会的隐密意义托付给那些也适合教导他人的人——正如宗徒们一样——还能怎样呢？为何装满十二个篮子？这既奇妙地发生了，因为是一件大事；又有益地发生了，因为是一件属灵的事。当时看见的人惊奇；但我们听到，并不惊奇。因为这事行出来，是为让他们看见，但写下来，是为让我们听到。眼睛为他们的作用，信德为我们起同样的作用。我们，即是以理智领悟我们眼睛所不能见的；并且我们比他们更有福，因为关于我们经上记载：\BibleQuote{那些没有看见而相信的，才是有福的。}\BibleRef{若 20:29}我还要说，也许我们领悟了那群众没有领悟的。我们确实得到了滋养，因为我们能获得大麦的精粹。\par

7. 最后，看见这奇迹的人想什么呢？他说：\Quote{那些人看见他所行的标记，就说\InnerQuote{这确实是一位先知}}。也许他们因此仍以为基督是一位先知，即因为他们坐在草地上。但他却是先知们的主，是先知们的完成者，是先知们的圣化者，然而也是一位先知：因为曾对梅瑟说过：\BibleQuote{我要为他们兴起一位像你一样的先知。}相似之处是按肉身，而非按尊威。而主的这许诺应理解为关于基督本身的，这在\BibleBook{宗徒}{大事录}中有清楚的阐明和宣读。\BibleRef{宗 7:37}主论及自己说：\BibleQuote{先知除了在自己的本乡以外，没有不受尊敬的。}\BibleRef{若 4:44}主是一位先知，主是天主的圣言；没有先知不以天主的圣言而说预言：天主的圣言与先知们同在，天主的圣言是一先知。以前的时代获得了由天主的圣言所感动和充满的先知；我们却得到了天主的圣言本身作为我们的先知。但基督就如他是天使之主一样，同样是先知，是先知之主。因为他也被称为伟大的谋士的天使。然而，先知在别处说什么？\BibleQuote{不是使者，也不是天使，而是他亲自来拯救他们；}\BibleRef{依 35:4}这就是说，他不会派遣使者来拯救他们，也不会派遣天使，而是他要亲自来。谁来？天使本身吗？他当然不是藉着一位天使来拯救他们，除非他本身就是这样的天使，同时也是天使之主。因为天使意为使者。如果基督没有带来信息，他就不会被称作天使；如果基督没有说预言，他就不会被称作先知。他劝勉我们持守信德并把握永生；他宣告了现在的事，预言了将来的事；因为他宣告了现在，因此他是天使或使者；因为他预言了将来，因此他是先知；而且，作为天主的圣言，他成了肉身，因此他是天使和先知之主。\par

\SourceRecord{Translated by John Gibb.；Nicene and Post-Nicene Fathers, First Series；Vol. 7.；Edited by Philip Schaff.；Buffalo, NY: Christian Literature Publishing Co.,；1888.；<http://www.newadvent.org/fathers/1701024.htm>.}

% structure: p0001 source=h1 target=section reason=page-title

\section{第25讲（若 6:15-44）}

1. 紧接昨日福音课的是今日的，今日的讲论应归于你们。当那奇迹发生，耶稣以五饼喂饱五千人，群众惊奇，说祂是来到世上的大先知时，接着是这样：\Quote{耶稣既知道他们来要强捉祂，立祂为王，就独自又退避到山上。}因此，应当明白，主与门徒坐在山上，看见群众向祂走来时，祂曾从山上下来，在山麓喂饱了群众。因为若祂先前没有从山上下来，又怎能再退避到那里呢？主从高处下来喂饱群众，是有含义的。祂喂饱了他们，就升天了。\par

2. 但当祂知道他们想强捉祂、立祂为王时，祂为什么升天呢？怎么，祂不是王，怕被立为王吗？祂当然不是那种由人立的王，而是将王国赐予人的王。难道耶稣——祂的行为就是言语——在此不也向我们表示什么吗？因此，他们想强捉祂、立祂为王，而祂为此独自退避到山上，这一行为在祂岂是沉默的吗？它不说什么，不意味着什么吗？或许他们想强捉祂，是要提前祂王国的时刻？因为祂现在来，并非立即为王，像我们祈祷的\Quote{愿祢的国来临}那样。祂确实永远与圣父一同为王，因为祂是天主之子、天主圣言、万物由之而造的圣言。但先知们预言了祂的王国，按祂作为人而成为基督、并使祂的信徒成为基督徒的那个方面而言。因此，将有一个基督徒的王国，目前正被聚集、被预备、被基督的血所赎买。祂的王国终将显现，当祂的圣人的光荣被揭示时，在祂所施行的审判之后——祂自己在上文说那是人子要施行的审判。关于这个王国，宗徒也说过：\Quote{当祂将王国交给天主圣父时。}\BibleRef{格前 15:24}与此相关，祂自己也说：\Quote{你们这些受我父祝福的，来承受自创世以来为你们预备的王国吧。}\BibleRef{玛 25:34}但门徒和信祂的群众以为祂就是这样来即刻为王的；因此，他们想强捉祂、立祂为王；他们想提前那个祂藏在祂自己那里的时刻，要在适当的时候显明它，并在适当的时候于世界末日宣告它。\par

3. 你们要知道他们想立祂为王——即提前并立刻显明基督的王国，而基督必须先受审判，然后审判别人——当祂被钉十字架时，那些寄望于祂的人对祂的复活已失去希望，祂从死者中复活后，遇见两个绝望地交谈的门徒，他们呻吟着彼此谈论所发生的事；祂装作陌生人，他们的眼被蒙蔽，认不出祂，祂就加入他们的谈话；他们向祂叙述他们谈论的事，说祂是一位先知，在言行上大有能力，被司祭长所杀；\Quote{而我们，}他们说，\Quote{曾希望祂就是要拯救以色列的那一位。}\BibleRef{路 24:13}你们希望得对；你们所希望的是真实的事；以色列的救赎在祂内。但你们为什么着急呢？你们想抓住它。以下也表明这是他们的感受：当门徒问祂关于结局时，他们对祂说：\Quote{祢要在这时显现吗？以色列国什么时候来到？}因为他们现在渴望它，现在想要它；也就是说，他们想强捉祂、立祂为王。但祂对门徒说（因为祂还得独自升天）：\Quote{父凭自己的权柄所定的时候、日期，不是你们可以知道的；但当圣神降临于你们身上时，你们将获得从上而来的能力，并要在耶路撒冷、全犹太和撒玛黎雅，直到地极，为我作证人。}\BibleRef{宗 1:6}你们希望我现在就显现王国；让我先聚集我所要显现的；你们爱崇高，你们将获得崇高，但要跟随我经过谦卑。关于祂也是这样预言的：\Quote{万民的聚集将环绕祢，为此祢要重返高处}；也就是说，为使万民的聚集环绕祢，使祢聚集许多人，祢要重返高处。祂就这样做了；祂喂饱了人，就升天了。\par

4. 但为什么说\Quote{祂退避}？因为祂若非自愿就不能被捉住，也不能被强捉，非自愿也不能被认出。但你们要知道这事是奥秘地、并非出于必需而是出于特意而为的，你们将在以下立刻看到：祂向寻找祂的同一群众显现，对他们说了许多话，并大谈天上的食粮；当谈论食粮时，祂不是与那些祂曾退避以免被祂捉住的人在一起吗？那么，难道祂不能在那个时候如此行事，使祂不被他们捉住，正如后来祂与他们说话时一样吗？因此，退避是有含义的。\Quote{祂退避}是什么意思？祂的崇高无法被理解。因为对于你不理解的事，你说：它从我这里逃走了。因此，\Quote{祂独自又退避到山上——从死者中首生者，升到诸天之上，并为我们转求。}\par

5. 同时，祂，那唯一的大司祭在上（祂已进入了帐幔内，百姓站在外面；因为旧约法律下那一年一次这样做的司祭，正是预表祂）：祂当时在上，门徒在船上经历着什么？因为那船预表教会，而祂在高天。如果我们首先不理解那船关于教会所经历的事，那么那些事件便不是有意义的，而仅仅是暂时性的；但如果我们看到那些预象在教会中所表达的真实含义，那么基督的行动显然就是一种言语。\Quote{但他说，天已晚了，祂的门徒下到海边；他们上了船，就渡海往葛法翁去了。}祂宣布那后来完成的事为迅速完成——\Quote{他们渡海往葛法翁去了。}祂回到如何来的：他们渡过湖面航行而来。当他们航行到祂已经说过他们来到的地方时，祂用总括的方法解释所遭遇的事。\Quote{那时天已黑了，耶稣还没有来到他们那里。}祂说\Quote{黑了}是对的，因为光还没有来到他们那里。\Quote{那时天已黑了，耶稣还没有来到他们那里。}随着世界末日临近，错谬增加、恐怖倍增、不义增加、不信增加；简言之，那光——由同一福音作者若望充分清楚地表明就是爱德，以致他说，\BibleQuote{凡恨自己弟兄的，就是在黑暗中；}\BibleRef{若壹 2:11}——那光，我说，常常熄灭；弟兄间仇恨的黑暗增加，天天增加，耶稣尚未来到。它如何显为增加？\Quote{因为罪恶要增多，许多人的爱情要冷淡。}黑暗增加，耶稣尚未来到。黑暗增加、爱情冷淡、罪恶增多——这些就是震撼船的波浪；风暴和风就是毁谤者的喧嚷。从此爱情冷淡；从此波浪汹涌，船被颠簸。\par

6. \Quote{忽然狂风大作，海就翻腾起来。}黑暗增加，辨别力减弱，不义增长。\Quote{当他们划了约二十五或三十斯塔狄时。}与此同时，他们奋力前进，继续前行；那些风、风暴、波浪和黑暗既未能使船无法前进，也未能使船破裂沉没；反而在这一切患难中继续前行。因为，尽管罪恶增多，许多人的爱情冷淡，波浪汹涌，黑暗增长，狂风怒吼，船仍在前进；\BibleQuote{因为那坚持到底的，必要得救。}\BibleRef{玛 24:12} 那斯塔狄的数目也不可轻看。因为当说到\Quote{他们划了二十五或三十斯塔狄，耶稣来到他们那里}时，确实不可能毫无意义。说\Quote{二十五}就足够了，同样\Quote{三十}；尤其是这是一个估计，而非叙述者的断言。一个简单的估计，如果他说将近三十斯塔狄或将近二十五斯塔狄，会使真理受损害吗？但从二十五他变成了三十。让我们来审视二十五这个数。它由什么组成？由什么构成？由\OriginalTerm{五数}{quinary}（五）构成。五数属于法律。同样的，梅瑟五书、那容纳病人的五个廊子、喂饱五千人的五个饼，都是五数。因此，二十五数表示法律，因为五乘以五，即五乘五，等于二十五，即五的平方。但这法律在福音到来之前缺乏圆满。此外，圆满包含在六数中。因此，天主在六天内完成了，或者说\Emph{圆满了}世界；同样，五乘以六，使法律因福音而得以圆满，即六乘五等于三十。因此，耶稣来到那些遵守法律的人那里。祂是如何来的？祂在水面上行走，将世上一切翻腾踏在脚下，压制一切高处。如此进行，只要时间延续，只要世纪流转。苦难增加、灾祸增加、忧愁增加，这一切翻腾而升高：耶稣继续行走，踏在波浪之上。\par

7. 然而苦难如此巨大，以致连那些信靠耶稣、努力坚持到底的人也大大害怕，唯恐失败；当基督踏着波浪、践踏世界的野心和高度时，基督徒深感恐惧。这些事难道没有预先告诉他吗？他们看见耶稣在水面上行走，也正当\Quote{他们害怕了}；就像基督徒，尽管对来世抱有希望，却常因人间事务的崩溃而不安，当他们看到这世界的高傲被践踏时。他们打开福音，打开圣经，发现这一切都早已预言；这是主的作为。祂践踏世界的高度，为要由谦卑的人得到荣耀。关于它的高傲，经上预言：\Quote{你要毁灭最坚固的城邑，}和\Quote{敌人的枪矛已到尽头，你已毁灭了城邑。}那么，你们基督徒为什么害怕呢？基督说：\Quote{是我，不要害怕。}为什么你们为这些事惊慌？为什么害怕？我预先说了这些事，我实行它们，它们必然发生。\Quote{是我，不要害怕。因此他们愿意接祂上船。}认出祂并欢喜，他们从恐惧中被释放。\Quote{立刻船就到了他们所要去的地方。}到了岸，就结束了；从水到旱地，从动荡到稳固，从路途到终点。\par

8. \BibleQuote{第二天，站在海那一边的群众，} 就是门徒们所到之处，\BibleQuote{看见那里没有别的船，只有门徒们所上的那一只，而且耶稣没有同门徒们一起上船，只有门徒们自己去了；然而从提庇黎雅来了别的船，靠近他们祝谢主后吃饼的地方：所以当群众看见耶稣不在那里，门徒们也不在时，他们也上了船，往葛法翁去找耶稣。} 然而他们对这样大的奇迹略有所知。因为他们看见门徒独自上了船，并且那里没有别的船。但也有船从他们吃饼的地方附近来，群众就坐这些船跟随祂。那时祂没有同门徒一同上船，那里也没有别的船。那么，耶稣怎么会突然到了海那边呢？除非祂在海面上行走，显示一个奇迹。\par

9. \BibleQuote{当群众找到了祂。} 看，祂向那些祂曾逃往山上去躲避的人群显现，因为祂害怕他们强迫祂为王。祂用各种方式向我们证明并使我们知道，这一切都是奥秘中说的，并在伟大的圣事（或奥迹）中做的，以表示重要的事物。看，那曾是逃往山上的人群中的祂，现在不是正对同一群人说话吗？让他们现在抓住祂，让他们现在立祂为王。\BibleQuote{当他们过了海找到祂时，就对祂说：\InnerQuote{辣彼，你什么时候到了这里？}}\par

10. 在奇迹的圣事后，祂开始讲论，如果能的话，那些吃饱的人可以得到进一步的饱足，祂可以用讲论充满他们的心灵，因为祂用饼充满了他们的肚腹，只要他们能接受。如果他们不接受，就让他们不接受的部分被收集起来，以免碎块失落。因此，让祂说话，让我们听。\BibleQuote{耶稣回答说：\InnerQuote{我实实在在告诉你们：你们找我，并不是因为见了神迹，而是因为吃饼吃饱了。}} 你们找我，是为了肉体，不是为了精神。有多少人寻找耶稣，只是为了祂能给他们现世的利益！一个人有事务在手，他寻求圣职人员的转求；另一个人被比他更有权势的人压迫，他逃到教会。另一个人希望为他与一位他影响力很小的人调停。一个人这样，一个人那样，教会天天充满了这样的人。很少有人为了耶稣本身而寻求耶稣。\BibleQuote{你们找我，并不是因为见了神迹，而是因为吃饼吃饱了。你们不要为那可朽坏的食粮劳碌，而要为那存留到永生的食粮劳碌。} 你们为了别的寻求我，要为了我自身而寻求我。因为祂暗示了真理，即祂自己就是那食粮：这在后续中清楚地显露出来。\BibleQuote{那是人子所要赐给你们的。} 我相信，你们期望再次吃饼，再次坐下，再次饱餐。但祂说过：\BibleQuote{不要为那可朽坏的食粮，而要为那存留到永生的食粮。} 正如祂对那撒玛黎雅妇人所说的：\BibleQuote{如果你知道那向你求喝的是谁，你或许会求祂，祂就会给你活水。} 当她说：\BibleQuote{你从哪里得到它？因为祢没有打水的器具，而且井很深。} 祂回答那撒玛黎雅妇人：\BibleQuote{如果你知道那向你求喝的是谁，你早就求祂，祂就会给你水，谁喝这水将永不再渴；因为谁喝这水还会再渴。} 她就高兴并愿意接受，仿佛不再受身体口渴的苦，因为她打水劳累。就这样，在这样的一番对话中，祂谈到了精神的饮料。完全同样地，这里也是如此。\par

11. 因此，\Quote{这食物，不是那必朽坏的，而是那存留到永生的，就是人子将要赐给你们的；因为天主圣父已印证了他。}不要将这位人子当作其他人子，关于他们曾说过：\Quote{人子们将信赖你翅膀的庇荫。}这位人子因某种圣神的恩宠而与众不同；按肉身说，他是人子，从众人中被拣选出来：他就是人子。这位人子也是圣子；这个人甚至是天主。在另一处，当他询问他的门徒时，他说：\Quote{人们说，我，人子，是谁？他们回答说：有人说是若翰，有人说是厄里亚，有人说是耶肋米亚，或是一位先知。他对他们说：但你们说我是谁？伯多禄回答说：你是基督，永生天主之子。}\BibleRef{玛 16:13}他自己宣称是人子，伯多禄宣称他是永生天主之子。他极合适地提到了他在怜悯中显示自己是什么；门徒也极合适地提到了他在光荣中继续是什么。天主的圣言将他的谦卑交托给我们注意；那人承认了他主的荣耀。的确，弟兄们，我认为这是公正的。他为了我们而谦卑自己，让我们光荣他。因为他成为人子并非为了自己，而是为了我们。因此他成为人子，当\Quote{圣言成了血肉，居住在我们中间。}为此，\Quote{天主圣父印证了他。}什么是印证？就是加上某个特殊的标记。印证就是印上一个无法与其他混淆的标记。印证就是给一件事物加上标记。当你给任何东西加上标记时，你这样做是为了避免它与其他东西混淆，以免你无法认出来。因此，\Quote{圣父} 则 \Quote{印证了他。}什么是\Quote{印证了}？赐给他某种特殊的东西，使他与所有其他人无可比拟。因此论到他说：\Quote{天主，你的天主，用喜乐油膏了你，胜过你的同伴。}那么，什么是印证？不就是使他独一无二吗？这就是\Quote{胜过你的同伴}的含义。所以，他说，不要因为我的人子身份而轻视我，而要从我寻求\Quote{不是那必朽坏的食物，而是那存留到永生的食物。}因为我成为人子，并非像你们一样：我成为人子，是如此的，即天主圣父印证了我。这是什么意思？他\Quote{印证了我}？赐给我某种我独有的东西，使我不与世人混淆，反而世人要因我而得救。\par

12. \Quote{他们于是对祂说：我们要做什么，好叫我们做天主的工作？}因为祂曾对他们说：\Quote{不要为那必朽坏的食物劳碌，而要为那存留到永生的食物劳碌。}\Quote{我们该做什么？}他们问；遵守什么，我们才能完成这条诫命？\Quote{耶稣回答他们说：这是天主的工作，就是信祂所派遣的那一位。}这也就是吃那食物，不是那必朽坏的，而是那存留到永生的。你预备牙齿和肚腹做什么？相信，你就已经吃了。信德确实与行为不同，正如宗徒所说：\Quote{人成义是因着信德，不在乎法律的行为。}\BibleRef{罗 3:28}有些行为看起来是好的，却没有对基督的信德；但那些行为并不好，因为它们没有被指向那个使行为成为善的终点；\Quote{因为基督是法律的终向，使凡信祂的人都获得正义。}\BibleRef{罗 10:4}因此，祂不愿意将信德与行为分开，反而宣称信德本身是工作。因为正是那同一信德，藉着爱德而工作。\BibleRef{迦 5:6}祂并没有说：这是你们的工作；而是说：\Quote{这是天主的工作，就是信祂所派遣的那一位。}这样，凡夸耀的，应当因主而夸耀。因为祂邀请他们相信，而他们却仍然要求他们可以相信的征兆。看犹太人是否不要求征兆。\Quote{他们于是对祂说：你行什么征兆，好叫我们看见而相信你？你做什么工作？}用五个饼吃饱他们，这是小事吗？他们确实知道这事，但他们更喜欢来自天上的玛纳胜过这食物。但主耶稣宣称自己是那样的，即祂比梅瑟更卓越。因为梅瑟不敢说自己给了\Quote{不是那必朽坏的食物，而是那存留到永生的食物。}耶稣应许了比梅瑟所给的更大的东西。的确，藉着梅瑟应许了一个王国，流奶流蜜的地方，暂时的和平，丰富的子女，健康的身体，以及其他一切事物，确实是暂时的好处，但在形象上是属灵的；因为在旧约中，它们是应许给旧人的。因此，他们考虑梅瑟应许的事物，他们也考虑基督应许的事物。前者应许地上饱满的肚腹，但那是必朽坏的食物；后者应许了\Quote{不是那必朽坏的食物，而是那存留到永生的食物。}他们注意到应许更大事物的那一位，但似乎尚未见到祂做更大的事。因此，他们考虑梅瑟做了什么工作，他们希望由应许他们如此伟大事物的那一位做出一些更大的工作。他们说：你做什么，好叫我们相信你？并且，为使你们知道他们将这些以前的奇迹与这个相比，并认为耶稣所做的这些奇迹是较小的，他们说：\Quote{我们的祖先在旷野里吃过玛纳。}但什么是玛纳？也许你们轻视它。\Quote{正如经上所载：祂给了他们玛纳吃。}藉着梅瑟，我们的祖先从天上得到了饼，但梅瑟没有对他们说：\Quote{要为那不衰朽的食物劳碌。}你应许了\Quote{不衰朽、存留到永生的食物}，但你没有做梅瑟所做的那些工作。他给的不是大麦饼，而是来自天上的玛纳。\par

13. \BibleQuote{那时耶稣对他们说：我实实在在告诉你们：并不是梅瑟给了你们那从天上来的饼，而是我的父赐给你们那从天上来的真饼；因为那从天上降下来并赐生命给世界的就是真饼。}那么真饼就是赐生命给世界的的那一位；也就是我稍前所说过的食粮——\BibleQuote{不要为那必朽坏的食物劳作，而要为那留存到永生的食物劳作。}因此，那玛纳所预表的就是这食粮，所有那些标记都是指向我的标记。你们渴望关于我的标记，难道你们轻看那被标记的那一位吗？不是梅瑟给了从天上来的饼：是天主赐予饼。但是什么饼？也许是玛纳？不，而是玛纳所预表的饼，即主耶稣自己。我的父赐给你们真饼。\BibleQuote{因为天主的饼就是那从天上降下来并赐生命给世界的那一位。于是他们对祂说：主啊，求祢常把这饼赐给我们。}如同那位撒玛黎雅妇人，曾对她说：\BibleQuote{凡喝这水的，将永远不再渴。}她立即按肉体理解，为摆脱缺乏，便说：\BibleQuote{主啊，请把这水赐给我；}同样地，这些人也说：\BibleQuote{主啊，求祢把这饼赐给我们；}这饼能使我们饱足，却不会匮乏。\par

14. \BibleQuote{耶稣对他们说：我就是生命的饼；到我这里来的，永远不饿；信我的，永远不渴。} \BibleQuote{到我这里来的}就是与\BibleQuote{信我的}意思相同；\BibleQuote{永远不饿}应理解为与\BibleQuote{永远不渴}相同。因为两者都象征那永不缺乏的永恒满足。你们渴望从天上来的饼；它就在你们面前，你们却不吃。\BibleQuote{但我曾对你们说过：你们也看见了我，却仍然不信。}但我并没有因此失去我的子民。\BibleQuote{难道你们的不信能使天主的信实失效吗？}\BibleRef{罗 3:3}因为，请看下文：\BibleQuote{凡父赐给我的人，必到我这里来；到我这里来的，我决不把他抛弃在外。}那是什么样内在的境界，以至于不再被抛弃在外？高贵的内室，甜蜜的隐退！没有厌倦、没有恶念的苦涩、没有诱惑的侵扰与悲伤之打断的隐秘居所！那不就是那有德的仆人将要进入的隐秘居所吗？主要对他说：\BibleQuote{进入你主人的福乐吧！}\BibleRef{玛 25:23}\par

15. \BibleQuote{凡到我这里来的，我必不抛弃。因为我从天降下，不是为执行我的旨意，而是为执行那派遣我来者的旨意。}难道这就是祢不抛弃那来到祢面前的人的原因吗？因为祢从天降下，不是为执行自己的旨意，而是为执行那派遣祢者的旨意？伟大的奥秘！我恳求你们，让我们一起敲门；也许会有一些东西出来喂养我们，正如那使我们喜悦的。那伟大而甜美的隐秘居所：\BibleQuote{凡到我这里来的。}请留意，请思考：\BibleQuote{凡到我这里来的，我决不抛弃。}为什么？\BibleQuote{因为我从天降下，不是为执行我的旨意，而是为执行那派遣我来者的旨意。}难道这正是祢不抛弃那来到祢面前之人的原因吗？因为祢从天降下，不是为执行自己的旨意，而是为执行那派遣祢者的旨意？正是这个原因。我们为何要问是否同一原因？就是同一原因；祂自己这样说的。因为怀疑祂的意思与祂所说的不同，是不对的。\BibleQuote{凡愿意到我这里来的，我决不抛弃。}仿佛你问：为什么？祂回答说：\BibleQuote{因为我来不是为执行我的旨意，而是为执行那派遣我来者的旨意。}恐怕灵魂离开天主的原因就是骄傲；是的，我毫不怀疑。因为经上记载：\BibleQuote{骄傲是万恶之始；人的骄傲之始是背离天主。}经上这样记载，是坚定而确实的，是真实的。因此，正如圣经过论骄傲的必死之人，穿着破烂的肉身之衣，被可朽坏的身体之重压着，而且自高自大，忘记自己披着什么皮囊——我问，圣经对他说什么？\BibleQuote{为何尘土和灰烬骄傲呢？}为何骄傲！让圣经告诉我们。\BibleQuote{因为在他有生之日，他把他的内心深处抛了出来。}\BibleRef{德 10:14}什么是\BibleQuote{抛了出来}？就是\BibleQuote{远远地扔掉}。这就是\BibleQuote{出去}。因为\BibleQuote{进入内在}，就是渴慕内心深处；\BibleQuote{抛出来}，就是\BibleQuote{出去}。骄傲的人抛出内心深处，谦卑的人切慕内心深处。如果我们因骄傲而被抛弃，就让我们因谦卑而回归。\par

16. 骄傲是万病的根源，因为骄傲是万罪的根源。当医生除去身体上的一种紊乱时，如果他仅仅治好了某种特定原因所产生的病症，却没有除去原因本身，他似乎暂时治好了病人，但原因还存在，疾病就会复发。举例来说，更明确地讲，体内某种体液产生皮屑或疮疡；接着发高烧，并伴有不小的疼痛；用某些治疗药物来抑制皮屑，并减轻疮疡的热度；药物用上了，也见效了；你看见那个浑身疮疡和皮屑的人被治好了；但因为那体液没有被排除，它又复发了溃疡。医生看出这一点，便清除体液，除去原因，就不会再有疮疡了。不义从哪里增多？从骄傲。治好骄傲，就不会再有更多的不义。因此，为使万病的根源，即骄傲，得以治愈，圣子降下，自甘卑微。人啊，你为何骄傲？天主为你成了卑微。你或许羞于效法一个卑微的人；无论如何，效法卑微的天主吧。圣子以人的形象而来，成了卑微。你被教导要变得谦卑，不是从人变成畜生。祂是天主，却成了人；你，人啊，要承认你是人。你全部的谦卑就是认识你自己。因此，因为天主教导谦卑，祂说：\BibleQuote{我来不是为行我的旨意，乃是行那派遣我者的旨意。}这就是对谦卑的称赞。骄傲行自己的旨意，谦卑行天主的旨意。因此，\BibleQuote{凡到我这里来的，我必不把他赶出去。}为什么？\BibleQuote{因为我来不是为行我的旨意，乃是行那派遣我者的旨意。}我谦卑而来，我来教导谦卑，我来作谦卑的老师：到我这里来的，与我成为一体；到我这里来的，成为谦卑的；凡依附我的，必会谦卑，因为他不行自己的旨意，乃行天主的旨意；因此他必不被赶出去，因为他骄傲时曾被赶出去。\par

17. 看看那些在\WorkTitle{圣咏}中为我们所推荐的内里之事：\BibleQuote{但人子们将投靠于祢翅膀的荫庇。}看看什么是进入内里；看看什么是逃到祂的保护之下；看看什么是甚至在圣父的杖下奔跑，因为祂所收纳的每一个儿子，祂都鞭打管教。\BibleQuote{但人子们必投靠于祢翅膀的遮蔽之下。}什么是内里？\BibleQuote{他们必因祢殿里的丰盛而饱足，}当祢差遣他们进入内里，进入他们主的喜乐时；\BibleQuote{他们必因祢殿里的丰盛而饱足；祢必叫他们喝祢喜乐的溪水。因为祢那里有生命的泉源。}不是在祢之外，而是在祢之内，有生命的泉源。\BibleQuote{并在祢的光中我们必得见光。求祢向认识祢的人施慈爱，向心里正直的人施行公义。}那跟随他们主旨意的人，不求自己的事，只求主耶稣基督的事，他们就是心里正直的人，他们的脚必不动摇。因为\BibleQuote{天主善待以色列，善待心里正直的人。至于我，他说，我的脚几乎失闪。}为什么？\BibleQuote{因为我嫉妒恶人，看见恶人享平安。}那么天主善待谁呢？除非心里正直的人？因为当我的心弯曲时，天主使我反感。为什么反感？因为祂赐福给恶人，因此我的脚步踌躇，仿佛我白白事奉了天主。所以，我的脚几乎失闪，因为我心里不正直。那么什么是心里正直？就是遵行天主的旨意。一个人富有，另一个人劳苦；一个生活邪恶却富有，另一个生活正直却受苦。让那生活正直而受苦的人不要生气；他拥有那富足之人所没有的内里之物：所以不要忧愁，不要烦恼，不要丧胆。那个富人在自己的箱子里有金子；这个人在自己的良心中有天主。现在比较金子和天主，箱子和良心。前者拥有会朽坏的东西，并且放在会朽坏的地方；后者拥有不会朽坏的天主，并且放在不能被夺走的地方：只要他是心里正直；因为那时他进入内里，不再出去。因此，他说了什么？\BibleQuote{因为祢那里有生命的泉源：}不在我们这里。因此我们必须进入内里，好使我们活；我们不可，好像满足于灭亡，也不愿满足于自己的东西，以致枯干，而必须把我们的口放在那真正的泉源，那里水永不断绝。因为亚当想凭自己的主意生活，他也因那先前因骄傲而堕落者而跌倒，那者邀请他喝他自己骄傲的杯。因此，因为\BibleQuote{祢那里有生命的泉源，并在祢的光中我们必得见光，}让我们在内里喝，让我们在内里看。为什么有一处从那出去？听为什么：\BibleQuote{不要让骄傲的脚临到我。}因此，那被骄傲的脚所临到的人，就出去了。表明他因此出去了。\BibleQuote{也不要让罪人的手动摇我；}因为骄傲的脚。你为什么说这个？\BibleQuote{所有作恶的人都跌倒了。}他们在哪里跌倒？就在他们的骄傲中。\BibleQuote{他们被赶出，不能站立。}那么，如果骄傲赶出了那些不能站立的人，谦卑就送进那些能永远站立的人。此外，那说\BibleQuote{卑微的骨头必欢乐}的人，之前曾说：\BibleQuote{求祢使我听到欢喜和快乐。}他说的\BibleQuote{使我听到}是什么意思？因听见祢我快乐；因祢的声音我快乐；因在内里喝我快乐。因此我不跌倒；因此\BibleQuote{卑微的骨头必欢乐}；因此\BibleQuote{新郎的朋友站着，听他。}因此他站着，因为他听。他喝内里的泉源，因此他站着。那些不愿意喝内里泉源的人，\BibleQuote{他们跌倒了；他们被赶出，不能站立。}\par

18. 如此，谦逊之师来不是为执行自己的旨意，而是为执行派遣祂者的旨意。让我们到祂那里去，进入祂内，接在祂身上，好使我们不再执行自己的旨意，而是执行天主的旨意：祂必不将我们赶出去，因为我们是祂的肢体，因为祂愿意藉着教导我们谦逊而成为我们的元首。最后，请听祂亲自谈论说：\Quote{凡劳苦和负重担的，你们都到我这里来；你们背起我的轭，向我学习，因为我是良善心谦的：}你们学了这事，\Quote{你们必会找到你们灵魂的安息，}\BibleRef{玛 11:28}从那里你们不会被赶出去；\Quote{因为我从天降下，不是为执行自己的旨意，而是为执行派遣我者的旨意；}我教导谦逊；唯有谦逊的人能到我这里来。唯有骄傲能赶走；保持谦逊而不离弃真理的人怎能出去呢？弟兄们，关于这隐藏的意义，能说的现已说了；这意义确是隐藏的，我不知道是否已用适当的言词为你们阐述清楚：为何祂不把来到祂面前的人赶出去；因为祂来不是为执行自己的旨意，而是为执行派遣祂者的旨意。\par

19. \Quote{这就是，}祂说，\Quote{就是派遣我的父的旨意：凡祂赐给我的，叫我一个也不失落。}保持谦逊的人已被赐给祂；祂接纳这样的人：不保持谦逊的人远离谦逊之师。\Quote{凡祂赐给我的，叫我一个也不失落。}\Quote{所以你们的天父不愿这些小子中的一个丧亡。}骄傲的人可能丧亡；但小子中没有一个丧亡，因为\Quote{你们若不变成如同这个小孩，决不能进入天国。}\Quote{凡父赐给我的，叫我一个也不失落，并且在末日我要使他复活。}请看祂在这里如何描绘那双重复活。\Quote{到我这里来的}立刻复活了，在我的肢体中成了谦逊的；但我在末日还要使他按肉身复活。\Quote{因为这是我父的旨意：凡看见子，并信从祂的，必得永生；并且在末日我要使他复活。}祂上面说：\Quote{凡听我的话，并相信派遣我来的那位；}但现在说：\Quote{凡看见子，并信从祂的。}祂没有说：看见子，并信从父；因为信从子等同于信从父。因为，\Quote{就如父在自己内有生命，同样也赐给子在自己内有生命。凡看见子，并信从祂的，必得永生：}藉着信从并进入生命，正如藉着第一次复活那样。并且，因为那不是唯一的复活，祂说：\Quote{并且在末日我要使他复活。}\par

\SourceRecord{Translated by John Gibb.；Nicene and Post-Nicene Fathers, First Series；Vol. 7.；Edited by Philip Schaff.；Buffalo, NY: Christian Literature Publishing Co.,；1888.；<http://www.newadvent.org/fathers/1701025.htm>.}

% structure: p0001 source=h1 target=section reason=page-title

\section{第26篇释义（若 6:41-59）}

1. 当我们的主\Person{耶稣基督}，正如我们在宣读福音时所听到的，曾说祂就是自天降下的食粮时，犹太人便窃窃私语说：\Quote{这人不是\Person{若瑟}的儿子，他的父母我们不是认识么？怎么他竟说：\InnerQuote{我是从天上降下来的呢？}}这些犹太人远离天上的食粮，不知如何饥渴它。他们心的颚是病弱的；耳朵开着却是聋的，他们看见却站着是盲的。的确，这食粮要求内心人的饥饿：因此祂在另一处说：\BibleQuote{饥渴慕义德的人是有福的，因为他们要得饱饫。}\BibleRef{玛 5:6}但宗徒\Person{保禄}说\Person{基督}为我们成了义德。\BibleRef{格前 1:30}因此，饥渴这食粮的人，就是饥渴义德——那从天降下的义德，天主所赐的义德，而非人靠自己成就的义德。因为如果人不是为自己制造义德，同一位宗徒就不会论及犹太人说：\Quote{因为他们不认识天主的义德，企图建立自己的义德，就不服从天主的义德。}\BibleRef{罗 10:3}这些人就是如此，他们不明白自天降下的食粮，因为满足于自己的义德，就不饥渴天主的义德。这是什么？天主的义德和人的义德？天主的义德在此不是指天主是义的那义德，而是天主赐予人的义德，使人藉着天主成为义。但犹太人的义德又是什么？是靠自己的力量成就的义德，他们依仗它，就宣称自己好像靠自己的德能履行了法律。但除非有恩宠的帮助，就是自天降下的食粮的帮助，没有人能履行法律。\Quote{因为法律的圆满}，如宗徒简而言之说，\Quote{就是爱德。}\BibleRef{罗 13:10}爱德，即爱，不是爱钱财，而是爱天主；不是爱地或天，而是爱那造了天地的祂。人从哪里能有那爱呢？让我们听听同一句：\Quote{天主的爱}，他说，\Quote{藉着赐予我们的圣神，已倾注在我们心中了。}\BibleRef{罗 5:5}因此，主将要赐予圣神时，说祂就是自天降下的食粮，劝勉我们信祂。因为信祂就是吃生活的食粮。信的人吃；他无形地饱足，因为无形地他重生。内在的婴孩，内在的新人。他更新之处，就在那里得到饮食。\par

2. 主对这些抱怨者回答了什么？\Quote{你们不要彼此抱怨。}仿佛祂说：我知道你们为何不饥饿，也不明白也不寻求这食粮。\Quote{你们不要彼此抱怨：除非派遣我的父吸引他，没有人能到我这里来。}恩宠的高贵卓越！没有人来除非被吸引。有祂所吸引的人，也有祂所不吸引的人；为何祂吸引这一个而不吸引另一个，不要试图判断，如果你不想出错。立刻接受，然后明白；你还没有被吸引？求使你被吸引。我们这里说什么呢，弟兄们？如果我们被\Quote{吸引}到基督那里，意思就是我们违背自己的意愿相信；那么，是强力被施加，而非意志被推动。人可以勉强来到教会，勉强走近祭台，勉强领受圣事：但他不能相信，除非他愿意。如果我们是身体相信，人可能被迫违背意愿相信。但相信不是身体的行为。听宗徒说：\Quote{人心里相信，就可成义。}然后呢？\Quote{口里宣认，就可获得救恩。}\BibleRef{罗 10:10}那宣认发自心的根。有时你听见一个人宣认，却不知道他是否相信。但如果你判断他不相信，就不应称他为宣认者。因为宣认就是说出你心里所存的事：如果你心里有一件事，口里却说另一件事，你是在说话，而不是宣认。既然人心里相信基督，那肯定不是违背他的意愿，而被吸引的人似乎是被强迫违背意愿，我们如何解决这问题：\Quote{除非派遣我的父吸引他，没有人能到我这里来}？\par

3. 有人说：如果他被吸引，他是勉强来的。如果他勉强来，那么他不相信；但如果不相信，他也不来。因为我们不是用脚跑到基督那里，而是借相信；也不是靠身体的运动，而是靠心的倾向接近祂。这就是为什么那位摸祂衣边的女人，比拥挤祂的群众更触及了祂。因此主说：\Quote{谁摸了我？}门徒们惊奇地说：\Quote{群众拥挤你，压迫你，你还说：\InnerQuote{谁摸了我？}}\BibleRef{路 8:45}祂又说：\Quote{有人摸了我。}那女人摸了，群众拥挤了。\Quote{摸了}是什么意思，岂不就是\Quote{相信了}？因此祂也对那复活后要俯伏在祂脚前的女人说：\Quote{不要摸我，因为我还没有升到父那里。}\BibleRef{若 20:17}你以为我只是你所见的那个，\Quote{不要摸我。}这是什么？你假设我只是你眼前所显示的那个：不要这样相信；就是说，\Quote{不要摸我，因为我还没有升到父那里。}对你来说我没有升上去，因为我从未离开那里。当她站在地上时，她并没有摸祂；那么她如何能在祂升到父那里时摸祂呢？然而，祂愿意如此被摸；祂被那些有益地摸祂的人如此摸，祂升到父那里，与父同住，与父同等。\par

4. 因此，祂在此也说，若你们留意，\Quote{除非父所吸引的人，谁也不能到我这里来。}不要以为你们是被勉强地吸引。心灵也被爱所吸引。我们不必害怕，恐怕那些衡量言语却远离事物、尤其是远离神圣事物的人，会在这一点上指责我们和圣经的福音之言；恐怕有人会对我们说：\Quote{我若被吸引，怎能凭意志相信？}我说，仅靠意志被吸引是不够的；你甚至被喜悦所吸引。被喜悦吸引是什么意思？\Quote{你应喜爱上主，祂必赐给你你心中的愿望。}有一种心灵的喜悦，天上的食粮对它是甘甜的。此外，若诗人说得对：\Quote{每个人都被自己的快乐所吸引，}——不是必然，而是快乐；不是义务，而是喜悦——那么，当一个人喜爱真理、喜爱福乐、喜爱正义、喜爱永生（这一切都是基督）时，我们岂不更当大胆地说他被吸引到基督那里？难道肉身感官有其快乐，而心灵却没有自己的快乐吗？若心灵没有自己的快乐，又怎能说：\Quote{世人将投靠在你翅膀的荫下；他们必因你殿里的丰盛而满足；你必使他们喝你快乐河的水。因为在你那里有生命的泉源；在你的光中我们必得见光}？给我一个爱的人，他就体会我所说的话。给我一个渴望的人，一个饥饿的人，一个在这旷野中旅行、渴慕并切望他永恒家园泉源的人；给我这样的人，他就知道我在说什么。若我对冷淡漠然的人说话，他不知道我在说什么。那些彼此窃窃私语的人正是这样。\Quote{凡父所吸引的人，}祂说，\Quote{就到我这里来。}\par

5. 但这是什么意思：\Quote{凡父所吸引的人}，既然基督自己吸引人？祂为何说\Quote{凡父所吸引的人}？若我们必须被吸引，就让我们被那一位所吸引，一位爱祂的人说：\Quote{我们要随你膏油的香气奔跑。}\BibleRef{歌 1:3} 但是，弟兄们，让我们转念，并尽可能领会祂要我们如何理解。父吸引那些相信子的人到子那里，因为他们认为天主是祂的父。因为天主生了与自己相等的子，以致那深思、并在其信仰中感受和默想他所信的那位与父相等的人，正是被父吸引到子那里的人。亚略相信子是受造物：父没有吸引他；因为那不相信子与父相等的人，并不认识父。亚略，你说什么？异端者，你说什么？基督是什么？不是真天主，他说，而是真天主所造的。父没有吸引你，因为你不认识父；你否认祂的子，你所想的不是子本身，而是别的什么。你既不被父吸引，也不被吸引到子那里；因为你所说的子与真相大不相同。佛提乌说：\Quote{基督只是一个人，并非也是天主。}这样相信的人，父没有吸引他。凡被父吸引的人说：\Quote{祢是基督，永生天主子。}不是作为先知，不是作为若翰，不是作为某个伟大正义的人，而是作为独一的、平等的，\Quote{祢是基督，永生天主子。}看，他被吸引了，而且是被父吸引。\Quote{约纳的儿子西满，你是有福的，因为不是血和肉启示了你，而是我在天之父。}\BibleRef{玛 16:16} 这启示本身就是吸引。你向羊展示一根青枝，你就吸引了它。你向孩子展示坚果，他就被吸引；他因他所奔向的东西而被吸引，因爱它而被吸引，身体不受伤害，被内心的绳索所吸引。那么，如果这些在尘世快乐和喜悦中的事物，向喜爱它们的人展示时吸引他们，因为\Quote{每个人都被自己的快乐所吸引}这句话是真的，那么由父启示的基督岂不吸引人吗？因为灵魂还有什么比真理更强烈渴望的呢？灵魂应当贪求什么，以致希望内在有一个健康的味觉来判断真实之事，岂不是为了吃喝智慧、正义、真理、永恒？\par

6. 但这将在何处？在那里更好，在那里更真实，在那里更圆满。因为在这里我们更容易饥饿而非满足，尤其当我们有好希望时：因为祂说：\Quote{饥渴慕义的人是有福的，}这指在此世；\Quote{因为他们必得饱饫，}这指在彼世。因此，当祂说了\Quote{若不是派遣我的父吸引他，谁也不能到我这里来}之后，祂接着说了什么？\Quote{在末日我要叫他复活。}我回报他所爱的、他所希望的：他将看见那尚未看见而相信的；他将吃他所饥饿的；他将因他所渴求的而得饱足。在哪里？在死人复活时；因为\Quote{在末日我要叫他复活。}\par

7. 先知书上记载说：\Quote{他们都要蒙圣父训诲。}我为什么说这话呢，犹太人？圣父并没有训诲你们；你们怎能认识我呢？因为那国度的所有人都要蒙圣父训诲，不是从人学习。虽然他们从人学习，但所领悟的是从内在赐予的，内在光照，内在启示。从外面传报信息的人做什么呢？即使现在我在说话，我是在做什么？我把言语的喧嚷灌入你们耳中。我所说、所讲的，若非那内在者启示，又是什么呢？外面是栽树的人，里面是树的创造者。栽种的和浇灌的，都在外面工作：这就是我们所做的。但\Quote{栽种的不算什么，浇灌的也不算什幺，而是那使生长增长的天主}（\BibleRef{格前 3:7}）。也就是说，\Quote{他们都要蒙圣父训诲。}都是谁？\Quote{凡听了圣父的教导而学习的，都到我这里来。}请看圣父如何吸引：祂以教导来愉悦，而不是强加必需。请看祂如何吸引：\Quote{他们都要蒙圣父训诲。}这就是天主的吸引。\Quote{凡听见了，并从圣父学习的，都到我这里来。}这就是天主的吸引。\par

8. 那么，弟兄们，这是怎么回事？如果凡听见圣父的教导而学习的，都来到基督这里，那么基督在这里没有教导什么吗？对此我们该说什么呢，就是那些没有看见圣父作为老师的人看见了子？子说话，但圣父教导。我，一个人，我教导谁呢？弟兄们，我教导谁呢？岂不是那听见我话的人？如果我这一个人教导那听我话的人，圣父也教导那听祂话的人。如果圣父教导那听祂话的人，那么就请问基督是什么，你们就会找到圣父的圣言。\Quote{在起初已有圣言。}不是在起初天主造了圣言，就像\Quote{在起初天主创造了天地}（\BibleRef{创 1:1}）。看，祂不是受造物。要学习被圣父吸引到子那里：为叫圣父教导你们，你们要听祂的圣言。你说，我听见祂的什么圣言呢？\Quote{在起初已有圣言}（不是\Quote{被造}，而是\Quote{已有}），\Quote{圣言与天主同在，圣言就是天主。}在血肉中的人怎能听见这样的圣言呢？\Quote{圣言成了血肉，居住在我们中间。}\par

9. 祂亲自也解释了这一点，并给我们显示了祂的意思，当祂说：\Quote{凡听见圣父的教导而学习的，都到我这里来。}祂立刻附加了我们能领悟的：\Quote{这不是说有人看见过圣父，只有那从天主来的，祂看见过圣父。}祂所说的是什么意思？我看见过圣父，你们没有看见过圣父；但你们若不受到圣父的吸引，就不来到我这里。被圣父吸引对你们而言是什么呢？岂不是学习圣父？学习圣父是什么呢？岂不是听见圣父？听见圣父是什么呢？岂不是听见圣父的圣言——那就是，听见我？所以，当我向你们说：\Quote{凡听见圣父的教导而学习的}，你们可能在内心说，但我们从未见过圣父，怎能学习圣父呢？你们从我这里听：\Quote{这不是说有人看见过圣父，只有那从天主来的，祂看见过圣父。}我认识圣父，我是从祂来的；但那是按照圣言从祂而来的方式，那里圣言存在，不是那发声并消逝的，而是那与说话者同在并吸引听者的。\par

10. 让接下来的话警戒我们：\Quote{我实实在在告诉你们：信我的，就有永生。}祂愿意启示自己，祂是什么；祂本可以简略地说：信我的，就有我。因为基督本身就是真天主和永生。因此，信我的，祂说，就进入我；进入我的，就有我。但\Quote{拥有我}是什么意思？就是拥有永生。永生接纳了死亡；永生愿意死去；但这是由于你们，而不是由于它自己；从你们那里它接受了那能替你们死去的东西。的确，祂从人取了血肉，但不是像人那样。因为祂在天有圣父，祂在地上选了母亲；在那里无母而生，在这里无父而出生。因此，生命接纳了死亡，为叫生命消灭死亡。因为\Quote{信我的，}祂说，\Quote{就有永生}：不是那显明的，而是那隐藏的。因为永生就是圣言，那\Quote{在起初与天主同在，圣言就是天主，生命就是人的光。}这同一永生也赐予了祂所取的血肉永生。祂来是要死；但第三天祂复活了。在圣言成了血肉与血肉复活之间，那中间的死亡被消耗了。\par

11. \BibleQuote{我是，}祂说，\BibleQuote{生命的食粮。}他们为什么骄傲？祂说：\BibleQuote{你们的祖先，}祂说，\BibleQuote{在旷野里吃过玛纳，却死了。}你们骄傲的是什么呢？\BibleQuote{他们吃过玛纳，还是死了。}为什么他们吃了却死了？因为他们只相信他们看见的；他们所没有看见的，他们不明白。因此，他们是\Quote{你们的}祖先，因为你们像他们一样。因为，我的弟兄们，就这可见的肉身死亡而言，我们这些吃从天上降下来的食粮的人不也死吗？他们死了，正如我们也要死，就我所说的，就这身体可见的和肉身的死亡而言。但至于那种死亡，就是主以恐惧警告我们的，他们的祖先正是在那种死亡中死了：梅瑟吃过玛纳，亚郎吃过玛纳，丕乃哈斯吃过玛纳，许多中悦主的人也吃过玛纳，他们却没有死。为什么？因为他们以属神的方式理解那可见的食物，以属神的方式饥渴，以属神的方式品尝，为能以属神的方式得到饱足。因为就连我们今天也接受可见的食物：但圣事是一回事，圣事的能力是另一回事。有多少人在祭台前领受却死了，而且确实因领受而死？因此宗徒说：\BibleQuote{自己吃、喝，就是吃喝自己的罪案。}\BibleRef{格前 11:29}因为主给犹大的那一口并不是毒药。但他还是领受了；当他领受时，仇敌进入了他里面：不是因为他领受了邪恶的东西，而是因为他这邪恶的人以邪恶的方式领受了善的东西。那么，弟兄们，你们要看到，你们应以属神的方式吃这天上的食粮；要带着纯洁来到祭台前。虽然你们的罪是每日的，但至少不要让它们成为大罪。在你们接近祭台之前，要好好考虑你们将要说的话：\BibleQuote{宽免我们的罪债，如同我们宽免欠我们债的人一样。}\BibleRef{玛 6:12}你们宽免，你们就必被宽免；平安地前来，这是食粮，不是毒药。但也要看你们是否宽免，因为如果你们不宽免，你们就是撒谎，而且是对你们不能欺骗的那一位撒谎。你们可以对天主撒谎，但你们不能欺骗天主。祂知道你们做什么。祂在内心看见你们，在内心省察你们，在内心察看，在内心审判，在内心祂要么定罪，要么加冕。但这些犹太人的祖先，是邪恶儿子的邪恶祖先，是不信者的不信的祖先，是抱怨者的抱怨的祖先。因为那个民族得罪主，没有比抱怨天主更甚的。因此，主为了显示这些人是这样抱怨者的后代，就这样开始对他们讲话：\BibleQuote{你们为什么彼此窃窃私议？}你们这些抱怨者，抱怨者的后代？你们的祖先吃过玛纳，却死了；不是因为玛纳是邪恶的东西，而是因为他们以邪恶的方式吃了它。\par

12. \BibleQuote{这是从天上降下来的食粮。}玛纳预表了这食粮；天主的祭台预表了这食粮。那些是圣事。在标记上，它们不同；在所指代的事物上，它们是相同的。听宗徒说：\BibleQuote{弟兄们，我不愿意你们不知道，}他说，\BibleQuote{我们的祖先都在云彩下，都从海中经过；都曾在云中和海中受洗归于梅瑟；并且都吃过同样的神粮。}当然，是同样的神粮；因为在肉身方面是另一种：因为他们吃玛纳，我们吃别的东西；但神粮与我们吃的是相同的。但\Quote{我们的}祖先，不是那些犹太人的祖先；是我们所像的那些，而不是他们所像的那些。他又加上说：\BibleQuote{并且都喝过同样的神饮。}他们喝一种饮料，我们喝另一种，但只是在可见的形式上，然而，这可见的形式以神性的能力指向同一事物。因为他们怎么可能喝了\Quote{同样的饮料}呢？\BibleQuote{他们喝了，}他说，\BibleQuote{那随着他们的神盘石，那盘石就是基督。}\BibleRef{格前 10:1}由此就有了食粮，由此外有了饮料。盘石是基督的预表；真正的基督是在圣言和肉身内。他们怎么喝的呢？盘石被杖击打了两次；两次击打预表了十字架的两根木梁。\BibleQuote{那么，这就是从天上降下来的食粮，谁若吃了这食粮，他必不死。}但这属于圣事的能力，不属于可见的圣事；是那在内心吃的人，不是那在外表吃的人；是在心中吃的人，不是用牙齿咬的人。\par

13. \BibleQuote{我是生命之粮，自天而降。}正因如此，祂说\BibleQuote{活的}，因为祂自天而降。玛纳也自天而降；但玛纳只是预像，这却是真理。\BibleQuote{谁若吃了这饼，必要生活直到永远；我所要赐给的饼，就是我的肉，为世界的生命。}肉何时理解这祂称为饼的肉？这被称为肉的，肉却不理解。唯其被称为肉，肉更不理解。他们为此惊惶，认为太难，以为不可能。祂说：\BibleQuote{是我的肉，}祂说，\BibleQuote{为世界的生命。}信者认识基督的身体，如果他不忽略成为基督的身体。让他们成为基督的身体，如果他们愿意藉基督的圣神而生活。只有基督的身体才能藉基督的圣神而生活。我的弟兄们，请理解我所说的。你是人，你有灵魂和身体。我称那被称为灵魂的为灵，因它而使你成为人，因为人由灵魂和身体组成。因此你有无形的灵和有形的身体。请告诉我，哪一个依靠另一个生活：是你的灵依靠你的身体生活，还是你的身体依靠你的灵生活？每个活着的人都能回答；而不能回答的人，我不知道他是否活着：每个活着的人怎么回答？我的身体当然依靠我的灵生活。那么你们是否也愿意藉基督的圣神生活？你们要在基督的身体内。因为我的身体不依靠你的灵生活，我的身体依靠我的灵生活，你的身体依靠你的灵生活。基督的身体不能藉基督的圣神生活。正是为此，宗徒保禄解释这饼时说：\BibleQuote{一个饼，}他说，\BibleQuote{我们虽是众多，却是一个身体。}\BibleRef{格前 10:17}啊，虔诚的奥迹！啊，合一的标记！啊，爱德的联系！那愿意生活的人，他有可以生活的地方，他有生活的源头。让他走近，让他相信；让他被组合进去，使他被赋予生命。让他不要退缩于肢体的结合；让他不要成为一个应当被切除的腐朽肢体；让他不要成为一个令其羞耻的畸形肢体；让他成为一个美好、合适、健全的肢体；让他依附于身体，为天主而活，藉天主而活：如今让他在地上劳苦，为的是将来他能在天上为王。\par

14. 犹太人彼此争论说：\BibleQuote{这人怎能把自己的肉给我们吃呢？}他们争论，且是在他们之间争论，因为他们不理解，也不愿领受合一的饼：\Quote{因为吃这种饼的人不会彼此争论；因为我们虽是众多，却是一个饼，一个身体。}藉这饼，\BibleQuote{天主使同心的人住在一屋里。}\par

15. 但他们所问的——他们彼此争论时——即主如何能把自己的肉给人吃，他们并未立即听到回答；反而进一步对他们说：\BibleQuote{我实实在在告诉你们：你们若不吃人子的肉，不喝他的血，在你们内便没有生命。}确实，如何吃，以及吃这饼的方式是什么，你们不知道；然而，\BibleQuote{你们若不吃人子的肉，不喝他的血，在你们内便没有生命。}祂说这些话，当然不是对死人，而是对活人。于是，恐怕他们以为这指的是现世生命，而在这事上也争论，祂接着又说：\BibleQuote{谁吃我的肉，并喝我的血，就有永生。}所以，谁不吃这饼，不喝这血，就没有这生命；因为人能没有那饼而拥有暂时的生命，但绝无可能拥有永生。因此，谁不吃祂的肉，不喝祂的血，在他内就没有生命；谁吃祂的肉，喝祂的血，就有生命。祂所加的\Emph{永生}这形容词，对两方面都适用。对于那为维持现世生命而摄取的食物则不然。因为不吃的人不会活，但吃的人也不一定活。因为许多吃了的人也死了，或死于衰老，或死于疾病，或死于其他意外。但在这种饮食中，即主的体血中，却不然。因为不吃的人没有生命，吃的人有生命，且是永生。因此，祂愿意这种饮食被理解为祂自己的身体和肢体的共融，也就是在祂所预定、蒙召、成义、得荣耀的圣人和信徒中的圣教会。其中，第一项即预定，已经完成；第二和第三项，即蒙召和成义，已经发生，正在发生，并将发生；但第四项，即得荣耀，目前是希望中的，但在实现上是未来的。这事的圣事，即基督体血的合一，在某些地方每日在主的祭台上预备，在某些地方按一定日期间隔预备，从主的祭台领受，一些人得到生命，一些人得到毁灭：但这事本身，即圣事所指向的实体，凡有分的人，对每人都得生命，无人得毁灭。\par

16. 但恐怕他们以为在这种饮食中永生被应许的方式是，领受的人即使在肉身内也永不死亡，祂俯就了这种想法；因为当祂说了\BibleQuote{谁吃我的肉，并喝我的血，就有永生}之后，立即补充道：\BibleQuote{在末日我要使他复活。}这样，同时，按照灵魂，他永生在安息中就有份，圣人们的灵魂被接纳其中；但就身体而言，他不会失去其永生，相反，他将在末日死者复活中拥有永生。\par

17. 祂说：\Quote{因为我的肉，实在是食物；我的血，实在是饮料。}因为人藉着食物和饮料，力求达到既不饥也不渴；但除了这食物和饮料，再也没有什么能真正赐予人这个，它使领受者成为不朽和不腐的；这就是诸圣相通的共融，那里将有圆满而完全的平安与合一。因此，确实，正如天主的人在我们之前所理解的，我们的主耶稣基督藉着那些从许多被归结为同一事物之物，将我们的心神指向祂的身体和血。因为由许多谷粒共同形成的是一个合一；由许多葡萄粒聚在一起形成的是另一个合一。\par

18. 简言之，祂现在解释祂所说的事如何发生，以及吃祂的肉、喝祂的血是什么。\Quote{谁吃我的肉，并喝我的血，便住在我内，我也住在他内。}因此，人吃那食物、喝那饮料，就是住在基督内，并让基督住在他内。所以，那不住在基督内，基督也不住在他内的人，无疑（在灵性上）既不尝祂的肉，也不喝祂的血，【虽然他可能用牙齿肉体地、可见地咬嚼基督身体和血的圣事】；但他反而是吃喝如此伟大之事的圣事，而自取审判，因为他不洁，竟敢前来领受基督的圣事——这圣事除了洁净的人，无人能堪当领受。关于这样的人，经上说：\Quote{心里洁净的人是有福的，因为他们要看见天主。}\BibleRef{玛 5:8}\par

19. 祂说：\Quote{就如生活的圣父派遣了我，我因圣父而生活；照样，那吃我的人，也要因我而生活。}祂并没有说：就如我吃圣父，并因圣父而生活；照样，那吃我的人，也要因我而生活。因为那生而平等的圣子，并不因参与圣父而变得更好；正如我们因参与圣子，藉着祂身体和血的合一而变得更好——这吃与喝所表示的正是这个。因此，我们因吃祂而生活，即因领受祂自己作为永生，这永生我们原本并不从自己拥有。然而祂自己因圣父而生活，因为祂被圣父派遣，因为\Quote{祂贬抑自己，听命至死，且死在十字架上。}\BibleRef{斐 2:8}因为若我们按照祂在另一处所说的\Quote{圣父比我大}来解释这句\Quote{我因圣父而生活}，那么同样，我们也因那比我们大者而生活；这来自祂的被派遣。事实上，派遣就是祂空虚自己，取了奴仆的形体；这应当正确理解，同时也要保持圣子与圣父本性的平等。因为圣父比作为人的圣子大，但祂拥有作为天主的圣子，与祂平等——而这位圣子又是天主又是人，天主子又是人子，同一的基督耶稣。就此而言，若这些话得到正确理解，祂这样说：\Quote{就如生活的圣父派遣了我，我因圣父而生活；照样，那吃我的人，也要因我而生活。}就好像祂说：我的空虚自己（因祂派遣我）导致我因圣父而生活；即把我的生命归于祂作为更大者；但任何人因我而生活，则是藉着那吃我的参与而成就的。因此，我被贬抑，却因圣父而生活；人被高举，则因我而生活。但若\Quote{我因圣父而生活}这句话被理解为祂出自圣父，而非圣父出自祂，则如此说并不损害祂的平等。再者，祂说\Quote{那吃我的人，也要因我而生活}，并非表示祂自己的平等与我们的平等相同，而是藉此显示中保的恩宠。\par

20. \Quote{这是从天上降下来的食粮；}好叫我们吃了它得以生活，因为我们不能从自己得到永生。祂说：\Quote{不像你们的祖先吃过玛纳，却死了；谁吃这食粮，必要生活直到永远。}祂愿意我们理解为，那些祖先死了，是指他们不会永远生活。因为即使那些吃基督的人，也必定会暂时死去；但他们永远生活，因为基督就是永生。\par

\SourceRecord{Translated by John Gibb.；Nicene and Post-Nicene Fathers, First Series；Vol. 7.；Edited by Philip Schaff.；Buffalo, NY: Christian Literature Publishing Co.,；1888.；<http://www.newadvent.org/fathers/1701026.htm>.}

% structure: p0001 source=h1 target=section reason=page-title

\section{讲道集 27（若望福音 6:60-72）}

1. 我们刚才从福音中听到了主继先前讲论之后所说的话。据此，理应对你们的耳朵和心灵作一篇讲道，且今天并非不合时宜；因为这是关于主的身体，祂说祂赐下那身体作为永生之粮。祂解释这赐予和恩赐的方式，祂如何赐下祂的肉作为食物，说：\BibleQuote{谁吃我的肉，并喝我的血，便住在我内，我也在他内。} 证明一个人已经吃了喝了，就在于他是否住在内且被住在内，是否居住且被居住，是否如此紧密结合而不被遗弃。因此，祂教导我们，并用奥秘的话劝诫我们，使我们在祂的身体内，在祂的肢体中，在祂元首之下，吃祂的肉，而不放弃与祂的合一。但当时在场的大多数人，因不理解祂，而起了反感；因为听到这些事，他们只想到肉，就是他们自己所是的肉。但宗徒说，且说得正确：\BibleQuote{随从肉性的切望，是死亡。} \BibleRef{罗 7:6} 主赐给我们祂的肉作为食物，然而按肉性去理解它，便是死亡；而祂却论及祂的肉，说其中有永生。因此，我们不应按肉性去理解这肉。正如随后的话所言：\par

2. \Quote{因此，有许多}，不是祂的敌人，而是\Quote{祂的门徒，听了这话，就说：这话生硬，谁能听它呢？} 如果祂的门徒认为这话生硬，祂的敌人又会怎样想呢？然而，正应该这样说，使众人不能完全理解。天主的奥秘应当使人热切关注，而非敌视。但这些人在主说这样的话时，很快就离开了祂：他们不相信祂是在说某种伟大的事，并用这些话遮盖某种恩宠；他们只是照自己的愿望，并以人的方式去理解，以为耶稣能够或决定这样行，即将圣言所穿上的肉，可以说，一块一块地分给那些相信祂的人。\Quote{这话}，他们说，\Quote{是生硬的，谁能听它呢？}\par

3. \Quote{但耶稣自己知道门徒们私下议论这事，}——因为他们彼此这样说，免得被祂听见；但祂知道他们内心的，内在听着——回答并说：\Quote{这事使你们反感；}因为我曾说：我赐给你们我的肉吃，我的血喝，这着实使你们反感。\BibleQuote{那么，如果你们看到人子升到祂原先所在的地方，又将如何呢？} 这是怎么回事？祂是否借此解答了困扰他们的问题？祂是否借此揭露了他们反感的根源？祂确实揭露了，只要他们理解。因为他们以为祂要将自己的身体分给他们；但祂说祂要升到天上，当然，是完整的：\BibleQuote{当你们看到人子升到祂原先所在的地方时；}那么，至少你们会看到，祂并非如你们所想的那样分施祂的身体；那么，至少你们会明白，祂的恩宠并非被牙齿咀嚼而消耗。\par

4. 祂又说：\BibleQuote{使人生活的是圣神；肉一无所用。} 在主恩准我们讲解这一点之前，不得疏忽地略过祂所说的另一处：\BibleQuote{那么，如果你们看到人子升到祂原先所在的地方，又将如何呢？} 因为基督是人子，由童贞玛利亚所生。因此，人子是在此地、在地上开始存在的，祂从大地取了肉身。为此，先知曾预言说：\BibleQuote{忠信从地下生出。} 那么，祂说\Quote{当你们看到人子升到祂原先所在的地方时，} 是什么意思？因为如果祂这样说：\Quote{如果你们看到天主子升到祂原先所在的地方，} 就不会有问题。但祂既然说：\Quote{人子升到祂原先所在的地方，} 难道人子在开始在地上存在之前，不就已经在天上吗？确实，祂说\Quote{祂原先所在的地方}，仿佛祂说这话时不在那里。但在另一处祂说：\BibleQuote{没有人升过天，除了那自天降下而仍在天上的人子。} \BibleRef{若 3:13} 祂不是说\Quote{曾在}，而是说\Quote{仍在天上的人子}。祂在地上说话，却宣称自己在天上。然而祂并没有这样说：\Quote{没有人升过天，除了那自天降下}的天主子\Quote{，祂在天上。} 这指向何处？难道不是为了让我们明白——我在先前的讲道中也曾向你们，我亲爱的，提示过——基督，既是天主又是人，是一个位格，而非两个位格，免得我们的信德不是天主圣三，而是四位一体？因此，基督是唯一的；圣言、灵魂和肉身，是一个基督；天主子与人子，是一个基督；天主子永远存在，人子在时间中出现，然而在位格的合一中是一个基督。当祂在地上说话时，祂原在天上。祂以那种方式在天上作为人子，正如祂在地上作为天主子一样；祂在地上作为天主子，取了肉身，而祂在天上作为人子，是在位格的合一中。\par

5. 那么，祂又加上了什么呢？\Quote{使人生活的是圣神；血肉一无所益。}让我们对祂说（因为祂允许我们，不是反驳祂，而是渴望知道）：\Person{主}啊，善师！血肉怎么一无所益，而祢却说：\Quote{人若不吃我的肉，不喝我的血，在他内就没有生命}？难道生命也一无所益吗？我们之所以为我们，岂不正是为了获得永生，就是祢藉着祢的血肉所应许的吗？那么，\Quote{血肉一无所益}是什么意思呢？它只有在他们理解的方式下才一无所益。他们确实按着血肉来理解，就像在尸体上被切割成块，或在肉铺里被出售一样；而不是像被\OriginalTerm{圣神}{Spirit}所活化的那样。因此，说\Quote{血肉一无所益}，与说\Quote{知识使人自大}是同样的方式。那么，我们是否应当立即憎恨知识呢？绝不！\Quote{知识使人自大}是什么意思呢？是\OriginalTerm{知识}{knowledge}单独，没有\OriginalTerm{爱德}{charity}。所以祂加上了：\Quote{但爱德却建造人。}\BibleRef{格前 8:1}因此，给知识加上爱德，知识就有益处，不是靠知识本身，而是靠爱德。同样在这里，\Quote{血肉一无所益}也只是单独的时候。让圣神加到血肉上，如同爱德加到知识上，它就大有裨益了。因为如果血肉一无所益，\OriginalTerm{圣言}{Word}就不会成为血肉，居住在我们中间了。如果通过血肉基督大大有益于我们，血肉又怎能一无所益呢？但圣神正是藉着血肉为我们的救恩成就了一些事情。血肉是一个器皿；要思考它所盛装的，而不是它本身。宗徒们被派遣出去；他们的血肉对我们一无所益吗？如果宗徒的血肉有益于我们，难道主的血肉对我们却一无所益吗？因为除了通过血肉的声音，圣言的声音怎能到达我们呢？书写怎能到达我们呢？所有这些都是血肉的活动，但只有当圣神推动它，如同推动它的工具一样。所以，\Quote{使人生活的是圣神；血肉一无所益}，正如他们理解血肉那样，但我给吃我的血肉却不是这样。\par

6. 所以，\Quote{我的话，}祂说，\Quote{我对你们所说的话，就是圣神，就是生命。}因为我们曾说过，弟兄们，这就是主藉着吃祂的肉、喝祂的血所教导我们的：我们应在祂内，祂也应在我们内。但我们是祂的\OriginalTerm{肢体}{members}时，我们就在祂内；我们是祂的圣殿时，祂就在我们内。但为了我们能成为祂的肢体，\OriginalTerm{合一}{unity}将我们连接在一起。除了\OriginalTerm{爱}{love}，还有什么能实现这种合一，使连接我们呢？天主的爱从何而来？请问宗徒：\Quote{天主的爱，他说，藉着所赐与我们的圣神，已倾注在我们心中了。}\BibleRef{罗 5:5}所以，\Quote{使人生活的是圣神}，因为正是\OriginalTerm{圣神}{Spirit}使肢体成为活的。除非圣神在\OriginalTerm{身体}{body}内找到肢体，它不会使任何肢体成为活的，因为就是圣神自己使它们活的。因为在你内的灵魂，人啊，你之所以成为人，就是靠着它，难道它会活化一个与你血肉分离的肢体吗？我把你的灵魂称为你的精神。你的灵魂只活化你血肉内的肢体；如果你取走一个肢体，它就不再被你的灵魂所活化，因为它没有连接到你身体的合一上。说这些话是为了使我们爱慕合一，畏惧分离。因为没有什么比与基督的身体分离更使一个基督徒应当畏惧的了。因为如果他与基督的身体分离，他就不是基督的肢体；如果他不是基督的肢体，他就不会被基督的圣神所活化。\Quote{谁若没有基督的圣神，宗徒说，就不属于基督。}\BibleRef{罗 8:9}所以，\Quote{使人生活的是圣神；血肉一无所益。我对你们所说的话，就是圣神，就是生命。} 那么，\Quote{就是圣神，就是生命}是什么意思呢？它们应被属神地理解。你已属神地理解了吗？\Quote{它们就是圣神，就是生命。}你已属血肉地理解了吗？那么，\Quote{它们就是圣神，就是生命}，但对你却不是这样。\par

7. 祂说：\Quote{但你们中间有些人并不相信。}祂不是说：你们中间有些人不理解；而是指出了他们不理解的原因。\Quote{你们中间有些人不相信，}因此他们不理解，因为不相信。因为先知曾说过：\Quote{如果你们不相信，你们将不会理解。}我们因\OriginalTerm{信德}{faith}而结合，因\OriginalTerm{明悟}{understanding}而获得生命。让我们先藉着信德坚持于祂，好使那因明悟而活化的东西得以存在。因为那不坚持的，就是反抗；反抗的，就不相信。反抗者怎能被活化呢？他是那应穿透他的光线之敌：他不转开他的眼睛，却关闭了他的心智。那么，\Quote{有些人不相信。}让他们相信并开放，让他们开放并被光照。\Quote{因为\Person{耶稣}从起初就知道那些相信的，和谁要出卖祂。}因为\Person{犹达斯}也在那里。有些人确实被冒犯了；但他却留下来等待机会，不是为理解。因为他为此留下来，主并没有对他保持沉默。主没有指名道姓，但也没有对他沉默；这样所有的人都会害怕，即使只有一人丧亡。在祂说话，并区分了相信与不相信的人之后，祂清楚地指出了他们不相信的原因。\Quote{因此，我对你们说过，祂说，除非蒙我\Person{圣父}恩赐，没有人能到我这里来。}因此，连相信也是赐给我们的；因为相信确是一件大事。如果这是一件大事，为你已经相信而喜乐，但不要自高，因为\Quote{你有什么不是领受的呢？}\BibleRef{格前 4:7}\par

8. \BibleQuote{从那时起，他的许多门徒退去了，不再与他同行。}他们退去了，却是随从撒殚，而非随从基督。因为我们的主基督曾称伯多禄为撒殚，其原因是他想走在主的前面，并给出建议叫他不该死亡，而他本是为死亡而来，好叫我们永不死亡；主对他说：\BibleQuote{退到我后面去，撒殚！因为你所体会的，不是天主的事，而是人的事。}\BibleRef{玛 16:23}他并没有驱使他退回去随从撒殚，并因此称他为撒殚，而是使他走到自己后面，好让他因随从主而不成为撒殚。但这些人是退去了，正如宗徒论及某些妇女时所说：\BibleQuote{有些已转去随从了撒殚。}\BibleRef{弟前 5:15}他们不再与他同行。看，他们从身体上被砍下，因为他们也许本不在身体内，他们失去了生命。他们必须被算作不信者，尽管他们被称为门徒。不是少数，而是\BibleQuote{许多人退去了}。这发生，或许是为安慰我们。因为有时会发生这样的事：一个人宣讲真理，而他所说的可能不被理解，于是听到的人就感到冒犯而离去。这时那人后悔说了那真理，对自己说：\Quote{我不该那样说，我不该说这个。}看，这事发生在主身上：他说话，失去了许多人；他仅与少数人留下。但他并不苦恼，因为他从起初就知道谁相信、谁不相信。如果这事发生在我们身上，我们会极其困惑。让我们在主内得到安慰，但仍要以谨慎说话。\par

9. 此时他转向所留下的少数人：耶稣就对那十二人（就是那留下的十二人）说：\BibleQuote{难道你们也愿意走吗？}连犹达斯也没有离去。但主已经清楚知道犹达斯为何留下；后来他才显示给我们。伯多禄代表众人回答，一人代替多人，合一代表整体：\BibleQuote{主！唯你有永生的话，我们去投奔谁呢？我们相信了，也知道你是基督，天主子。}请看伯多禄，因天主的恩赐和圣神的更新，如何理解了他。除了因为他信，还能是什么？\BibleQuote{唯你有永生的话。}因为你在你的身体和血的职务中拥有永生。\BibleQuote{我们相信了，也知道。}不是知道了才相信，而是\BibleQuote{相信了才知道}。因为我们相信是为了知道；因为如果我们想先知道再相信，我们就既不能知道也不能相信。我们相信并知道了什么？\BibleQuote{即你是基督，天主子}；也就是说，你就是那永生，你在你的肉和血中赐下的正是你本身。\par

10. 主\Person{耶稣}便说：\Quote{我不是拣选了你们十二人吗？你们中间有一个是魔鬼。}因此，祂难道应该说：\Quote{我拣选了十一人}？难道一个魔鬼也被拣选，并在蒙拣选者中间？人们通常以赞美的方式称呼某些人为\Quote{蒙拣选者}；或者一个人是因为他做出了某些伟大的善事，且并非出于自己的意志和知识而被拣选？这专属天主；相反的情形则是恶人的特征。因为正如恶人滥用天主的善工；相反，天主也善用恶人的恶行。身体的各个肢体，既然只能由它们的创造者和塑造者——天主——来安排，这是何其美善！然而，淫欲对眼睛作了何等恶用？谎言对舌头作了何等恶用？作假见证的人不是先用舌头杀死自己的灵魂，然后在毁灭自己之后，试图伤害他人吗？他滥用了舌头，但舌头本身并不是邪恶的；舌头是天主的作品，但不义滥用了天主的这美好作品。那些奔向罪恶的人如何用他们的双脚？杀人犯如何用他们的双手？恶人如何滥用那些存在于他们之外的天主的美好受造物？他们用黄金败坏审判，压迫无辜。恶人甚至滥用光明本身；因为通过邪恶的生活，他们甚至将用来观看的光明用于他们的恶行。一个恶人在去做恶事时，希望有光照亮他，以免绊倒；然而他内在已经绊倒并跌倒了；他在身体上所害怕的，已经在心里发生了。因此，为了避免逐一罗列的繁琐，可以说：恶人滥用了天主的一切美好受造物；相反，善人善用了恶人的恶行。有什么像唯一的天主那样善呢？的确，主自己曾说：\Quote{除了唯一的天主外，没有善的。} \BibleRef{谷 10:10} 那么，祂越好，就越善用我们的恶行。有什么比\Person{犹达斯}更坏呢？在所有跟随师傅的人中，在十二人当中，公款的钱囊交给了他；分配施舍给穷人的任务也派给了他。他对如此大的恩惠、如此大的荣耀忘恩负义，拿了钱，失去了义德：虽已死去，却背叛了生命：他先前作为门徒所追随的那一位，如今却以敌人的身份迫害。这一切恶都是\Person{犹达斯}的；但主用了他的恶来成就善。祂忍受被出卖，以救赎我们。看，\Person{犹达斯}的恶变成了善。\Person{撒殚}迫害了多少殉道者！如果\Person{撒殚}停止迫害，我们今天就不会庆祝\Person{圣老楞佐}那极其荣耀的冠冕了。如果天主甚至利用魔鬼本身的恶行来成就善，那么恶人滥用善事所导致的，只是伤害自己，并不能抵触天主的良善。师傅使用了那个人。如果祂不知道如何使用他，这位创造主师傅就不会允许他存在。因此，祂说：\Quote{你们中间有一个是魔鬼，}而我却拣选了你们十二人。这句话：\Quote{我拣选了你们十二人，}可以这样理解：十二是一个神圣的数字。因为那个数字的尊荣并未因失去一个而被剥夺，因为另一个人被选来取代那丧亡者的位置。\BibleRef{宗 1:26} 这个数字仍然是神圣的，一个包含十二的数字：因为他们要向全世界，即世界的四个方向，宣告天主圣三。这就是三乘以四的原因。因此，\Person{犹达斯}只是断绝了自己，并没有亵渎十二这个数字：他遗弃了他的师傅，因为天主指定了一个继承者来接替他的位置。\par

11. 主关于祂的体和血所说的一切——并在那分施的恩宠中，祂应许我们永生，并且祂意指那些吃祂的肉、喝祂的血的人，可从他们住在祂内、祂在他们内的事实来理解；而那些不信的人并不理解；并且他们因按血肉理解属神之事而绊倒；正当这些人绊倒并丧亡时，主临在是为了安慰留下的门徒，为了考验他们，祂问：\Quote{你们也要走吗？}好让他们的坚定答复为我们所知，因为祂知道他们与祂同在——那么，最亲爱的诸位，愿这一切都对我们有益，使我们不仅像许多恶人那样只在圣事中吃基督的体和血，而是为了参与圣神而吃喝，使我们作为肢体留在主的身体内，被祂的圣神所赋予生命，并且我们不至绊倒，即使现在有许多人与我们一同暂时地吃和喝圣事，他们最终将遭受永罚。因为目前基督的身体仿佛混杂在打谷场上：\Quote{但是主认识那些属于祂的人。} \BibleRef{弟后 2:19} 如果你知道你所打的是什么，就知道子粒隐藏在那里，打谷并没有消灭簸箕所扬净的；弟兄们，我们确信，我们所有在主的身体内并住在祂内的人，也愿祂住在我们内，必须在这世上与恶人一同生活直到终结。我不是说那些以口舌亵渎基督的恶人；因为现在很少有人以口舌亵渎，但许多人以生活来亵渎。那么，我们必须在这些人中生活直到终结。\par

12. 但祂所说：\Quote{那住在我内，我也住在他内的}，是什么意思呢？岂不是殉道者们所听到的：\Quote{那坚持到底的，才可得救}？\BibleRef{玛 24:13} 我们今天庆祝其节日的\Person{圣老楞佐}，是如何住在祂内的？他住在祂内，直至受诱惑，直至受暴虐的审问，直至最残酷的威吓，直至毁灭——那是小事一桩，直至野蛮的酷刑。因为他并非被迅速处死，而是在火中受折磨：他被允许活很长时间；不，不是被允许活很长时间，而是被迫缓慢、拖延地死亡。然后，在那拖延的死亡中，在那痛苦中，因为他已好好地吃、好好地喝，如同一位享用了那食粮、陶醉于那杯爵的人，他感觉不到痛苦。因为那说\Quote{使人生活的是圣神}的那一位就在那里。因为肉身确实在燃烧，但圣神却在使灵魂生活。他没有退缩，而进入了天国。但是五天前我们庆祝其节日的\Person{圣西斯笃}曾对他说：\Quote{我儿，不要悲伤}；因为西斯笃是主教，而他（老楞佐）是执事。他说：\Quote{不要悲伤；三天后你将跟随我。}他说三天，意指\Person{圣西斯笃}受难日与今天\Person{圣老楞佐}受难日之间的间隔。三天就是间隔。何等安慰！他不说：\Quote{我儿，不要悲伤；迫害会停止，你将平安无事}；而是说：\Quote{不要悲伤：我先去之处，你将跟随；你的跟随不会延迟：三天之后，你将与我同在。}他接受了神谕，战胜了魔鬼，获得了胜利。\par

\SourceRecord{Translated by John Gibb.；Nicene and Post-Nicene Fathers, First Series；Vol. 7.；Edited by Philip Schaff.；Buffalo, NY: Christian Literature Publishing Co.,；1888.；<http://www.newadvent.org/fathers/1701027.htm>.}

% structure: p0001 source=h1 target=section reason=page-title

\section{讲道集 28（若望福音 7:1-13）}

1. 在这章福音中，弟兄们，我们的主耶稣基督特别将祂自己与祂的人性相关地交托给我们的信德。因为事实上，祂无论在言语还是行动中，始终持守祂应被相信是天主也是人：天主创造了我们，人寻找了我们；与圣父同在，常是天主；与我们同在，在时间内是人。因为祂若不曾成为祂所造的，就不会寻找祂所造的人。但要记住这点，不要让它从你们心中滑落：基督成为人，如此方式并未停止是天主。祂在保持天主的同时，创造了人的那一位取了人性。因此，当祂作为人隐藏自己时，不可认为祂失去了能力，而只是为我们的软弱提供了一个榜样。因为祂在愿意时被拘捕，在愿意时被处死。但由于将要有祂的肢体，即祂的信徒，他们没有我们天主所具有的那种能力；藉着祂的隐藏，藉着祂好像不愿被处死而隐藏自己，祂预示了祂的肢体将这样做，事实上祂自己就在这些肢体中。因为基督不只是在头而不在身体，而是整个基督在头和身体。所以，祂的肢体是什么，祂就是什么；但祂是什么，并不必然意味着祂的肢体也是。因为如果祂的肢体不是祂自己，祂就不会说：\BibleQuote{扫禄，你为什么迫害我？}\BibleRef{宗 9:4} 因为扫禄不是在地上迫害祂自己，而是迫害祂的肢体，即祂的信徒。然而祂不会说：我的圣人、我的仆人，或简言之，我的弟兄（这是更尊贵的称呼）；而是说 \Emph{我}，即我的肢体，我是他们的头。\par

2. 有了这些预备的说明，我想我们不必在这章的含义上过多劳苦；因为那常常在头中预示，将来要在身体中实现。\BibleQuote{此后，} 他说，\BibleQuote{耶稣在加里肋亚行走，因为祂不愿在犹太行走，因为犹太人想要杀死祂。} 这就是我所说的；祂为我们的软弱提供了一个榜样。祂没有失去能力，而是安慰我们的软弱。因为如我所说，会发生这样的事：祂的一些信徒会退隐到隐蔽处，以免被迫害者发现；并且，为了让这种隐藏不被归罪为罪行，这事首先发生在头中，以后要在肢体中得以确认。因为经上说：\BibleQuote{祂不愿在犹太行走，因为犹太人想要杀死祂，} 仿佛基督不能既在犹太人中行走，又不被他们杀死似的。因为祂在愿意时显示了这种能力；当他们想要抓住祂，而祂即将受难时，\BibleQuote{祂对他们说：你们找谁？他们回答：耶稣。于是祂说：我就是，} 没有隐藏，而是显示了自己。然而，这种显示他们无法抗拒，反而\BibleQuote{后退，倒在地上。}\BibleRef{若 18:6} 然而，因为祂来是要受难，他们起来，抓住了祂，把祂带到审判官面前，杀死了祂。但他们做了什么？某段圣经说：\BibleQuote{大地被交在不敬神者的手中。}\BibleRef{约 9:24} 肉身被交在犹太人的权下；这样，那钱袋，好似，可以被撕开，我们的赎价得以流出。\par

3. \BibleQuote{犹太人的帐棚节临近了。} 什么是帐棚节，读圣经的人都知道。他们在圣日搭盖帐棚，模仿他们从埃及被领出后在旷野寄居时所住的帐棚。这是一个圣日，一个隆重的庆节。犹太人正在庆祝这个节日，他们纪念主的恩惠——他们却将要杀害主。在这个圣日上（因为有好几个圣日；但在犹太人看来它被称为一个圣日，虽然它不是一天，而是几天），\BibleQuote{祂的兄弟} 对主基督说话。要理解\BibleQuote{祂的兄弟} 这个短语，如你们知道必须接受的那样，因为这不是你们听到的新鲜事。童贞玛利亚的血亲常被称作主的兄弟。因为圣经的习惯是用\Quote{兄弟}一词称呼血亲和其他所有近亲，这与我们的用法不同，不在我们的言语习惯之内。因为谁会称一位叔父或外甥为\Quote{兄弟}呢？然而圣经称这类亲属为\Quote{兄弟}。因为亚巴郎和罗特被称为兄弟，而亚巴郎是罗特的叔父。\BibleRef{创 11:27} 拉班和雅各伯被称为兄弟，而拉班是雅各伯的叔父。\BibleRef{创 28:2} 因此，当你们听到主的兄弟时，要视他们为玛利亚的血亲，她并没有第二次生育。因为，在安放主身体的坟墓里，之前或之后都没有任何死者躺卧；同样，玛利亚的子宫，之前或之后也没有怀过任何可朽坏的人。\par

4. 我们已经说了兄弟是谁，让我们听听他们说了什么：\BibleQuote{你离开这里，往犹太去，好叫你的门徒也看见你所行的事。} 主的工事并非对门徒隐藏，但对这些人并不明显。他们可能以基督为亲属，但正因那层关系，他们不屑于相信祂。福音告诉了我们；因为我们不敢把这当作仅仅是意见，你们刚才已经听到了。他们继续劝告祂：\BibleQuote{因为没有人暗地里行事，却想公开被人知道；你既然行这些事，就该将自己显示给世界看。} 紧接着说：\BibleQuote{原来连祂的兄弟也不相信祂。} 他们为什么不信祂？因为他们寻求人的光荣。至于祂的兄弟似乎劝告祂的事，他们谋求的是祂的光荣。你行奇妙的事，就让自己被人知道；也就是说，显现给众人，好让众人赞美你。血肉向血肉说话；但无天主的血肉，向有天主的血肉说话。这是血肉的智慧向成为血肉并居住在我们中间的\OriginalTerm{圣言}{Word}说话。\par

5. 主对这些事如何回答呢？于是耶稣对他们说：\Quote{我的时候还没有到，你们的时候却常是现成的。} 这是什么意思呢？基督的时候还没有到吗？那么，如果祂的时候还没有到，基督为什么来了呢？我们不是听说过宗徒说：\BibleQuote{但时期一满，天主就派遣了自己的儿子}？\BibleRef{迦 4:4} 所以，如果祂是在时期一满时被派遣的，祂是在应当被派遣的时候被派遣的，祂是在祂应当来的时候来的。那么，\Quote{我的时候还没有到}是什么意思呢？弟兄们，要懂得他们说话时的意图，当他们似乎是以兄弟的身份劝告祂时。他们是在给祂出主意去追求光荣；以世俗的方式和属地的心态劝告祂，不要让祂埋没于名声，也不要让祂隐藏于暗处。这就是主对那些给祂出主意追求光荣的人的回答：\Quote{我的时候还没有到了——我光荣的时候还没有到。} 请看这是多么深奥：他们是在劝告祂追求光荣；但祂愿意谦卑先于崇高，并愿意通过谦卑准备通向高升的道路。因为那些门徒当然也追求光荣，他们希望坐在祂的右边和左边：他们只想到终点，而看不到必须通过的道路；主把他们召回道路，使他们能按顺序到达祖国。因为祖国在高处，道路却在低处。那土地是基督的生命，道路是基督的死亡；那土地是基督的居所，道路是基督的苦难。那拒绝道路的人，为什么寻求祖国呢？总之，对这些追求高升的人，祂也给了这个回答：\BibleQuote{你们能饮我将要饮的杯吗？}\BibleRef{玛 20:22} 看哪，这就是你们必须通过才能达到你们所渴望的高度的道路。祂所提到的杯的确是祂谦卑和苦难的杯。\par

6. 因此，在此也是如此：\Quote{我的时候还没有到；但你们的时候}，这即是世俗的光荣，\Quote{却常是现成的。} 这就是基督——即基督的身体——在预言中所说的时期：\Quote{当我接受了合宜的时候，我要按正义施行审判。} 因为目前不是审判的时候，而是容忍恶人的时候。所以，让基督的身体目前忍受，并容忍恶人的邪恶。然而，让它现在拥有正义，因为凭着正义它将来到审判。圣经在圣咏中对那些容忍这世界邪恶的肢体说了什么呢？\Quote{主不会抛弃祂的子民。} 事实上，祂的子民在不配的人、不义的人、亵渎者、怨言者、毁谤者、迫害者，以及如果允许的话，毁灭者中间劳苦。是的，它劳苦；但是\Quote{主不会抛弃祂的子民，也不会遗弃祂的产业，直到正义转向审判。} \Quote{直到正义}（现在在祂的圣徒身上的正义）\Quote{转向审判}；那时将应验对他们所说的话：\Quote{你们要坐在十二个宝座上，审判以色列十二支派。}\BibleRef{玛 19:28} 宗徒拥有正义，但还没有那审判，就是他说过的：\Quote{你们不知道我们将审判天使吗？}\BibleRef{格前 6:3} 所以，现在是有德生活的时期；审判恶人的时期将在以后。\Quote{直到正义}，他说，\Quote{被转向审判。} 审判的时期将是主在此所说的：\Quote{我的时候还没有到}的时期。因为将有一个光荣的时期，那时那以谦卑而来的将带着崇高而来；那来受审判的将来到审判；那来被死者杀害的将来到审判活人和死人。圣咏说：\Quote{天主，我们的天主，将要公开来临，且不沉默。} \Quote{将要公开来临}是什么意思？因为祂是隐藏着来的。那时祂将不沉默；因为当祂隐藏着来时，\Quote{祂被牵去宰杀像羊，又像羔羊在剪毛者前不开口。}\BibleRef{依 53:7} 祂将来到，且不保持沉默。\Quote{我沉默了，}祂说，\Quote{难道我要一直沉默吗？}\BibleRef{依 42:14}\par

7. 但那些现在拥有正义的人需要什么呢？就是那篇圣咏所读到的：\Quote{直到正义化为审判，而拥有它的人就是心地正直的。}也许你会问，谁是心地正直的人？我们在\WorkTitle{圣经}中发现，那些忍受世界的邪恶而不指责天主的人，是心地正直的人。请看，弟兄们，我所说的是一件不寻常的事。因为我不知怎么的，当任何灾祸临到一个人时，他跑去指责天主，而他本应指责自己。当你得到任何好处时，你赞美自己；当你遭受任何灾祸时，你指责天主。这就是弯曲的心，不是正直的心。当你从这种扭曲和乖戾中被治愈后，你过去所做的将会转变为相反。你过去做什么？你在天主的恩惠中赞美自己，而在你自己的邪恶中指责天主；当你的心被转变而变得正直时，你将在天主的恩惠中赞美天主，并在你自己的邪恶中指责自己。这些就是心地正直的人。简言之，那个人，当他看到恶人成功、善人受苦时心还不正，后来被纠正时说：\Quote{以色列的天主对心地正直的人多么好啊！但我，}当我心不正时，\Quote{我的脚几乎滑跌，我的脚步险些失足。}为什么？\Quote{因为我嫉妒恶人，看见罪人的平安。}他说，我看见恶人亨通，就对天主不满；因为我希望天主不要让恶人快乐。让人明白：天主从来没有允许过这件事；但一个人被认为是快乐的，是因为人不知道什么是真正的快乐。那么，让我们心正：我们荣耀的时刻还没有到来。告诉这世界的爱慕者，就像主的弟兄那样：\Quote{你们的时刻总是准备好的；}我们的时刻\Quote{还没有到。}因为我们也敢说这话。既然我们是我们的主耶稣基督的身体，既然是祂的肢体，既然我们喜乐地承认我们的头，让我们毫不迟疑地说这话；因为为了我们，祂也屈尊亲自说了这话。当这世界的爱慕者辱骂我们时，让我们对他们说：\Quote{你们的时刻总是准备好的；我们的时刻还没有到。}因为宗徒曾对我们说：\Quote{因为你们已经死了，你们的生命与基督一同藏在天主内。}什么时候我们的时刻来到？\Quote{当基督，}他说，\Quote{你们的生命显现时，那时你们也要与祂一同在光荣中显现。} \BibleRef{哥 3:3}\par

8. 祂又说了什么？\Quote{世界不能恨你们。}这是什么意思？不过是：世界不能恨它的爱慕者、假见证者？因为你们称恶为善，称善为恶。\Quote{但它恨我，因为我指证它，它的工作是邪恶的。你们上这个庆节去吧。}\Quote{这个}是什么意思？你们在那里寻求人的光荣。\Quote{这个}是什么意思？你们愿意延长肉身的喜乐，而不默想永恒的喜乐。\Quote{我不上这个庆节去，因为我的时候还没有完全来到。}在这个庆节上，你们寻求人的光荣；但我的时候，即我光荣的时候，还没有来到。那将是属于我的庆节，不是匆匆赶过这些日子，而是永远存留；那将是节庆，无尽的喜乐，无瑕的永恒，无云的晴朗。\Quote{祂对他们说了这些话以后，仍然留在加里肋亚。但当祂的弟兄们上去以后，祂也上去了，不过不是公开的，而是暗暗的。}因此，\Quote{不是上这个庆节，}因为祂所欲求的不是暂时的光荣，而是教导一些有益处的事，纠正人，劝诫他们想到永恒的庆节，叫他们远离对这世界的爱，而转向天主。但这又是什么意思：\Quote{祂暗暗地上庆节去}？主的这个行动也不是没有意义的。在我看来，即使从祂暗暗地上庆节这个情况，祂已经有意要表示什么；因为后面的事将会显示，祂是在庆节的中期（即那些日子过了一半的时候）上去的，甚至公开施教。但祂说：\Quote{暗暗地，}意思是，不向人们显示自己。基督暗暗地上那个庆节去，不是没有意义的，因为祂自己隐藏在那庆节里。我目前所说的也还隐藏着。那么，让它彰显出来吧，揭开帷幕，让隐藏的事出现吧。\par

9. 在神圣法律的繁复圣经中，对古代选民以色列所说的一切事物——他们在祭献、司祭职务、庆节中所行的一切，总而言之，他们凡用以敬拜天主的任何事物，以及凡作为诫命对他们所说和所赐的一切——都是未来事物的预像。是哪些未来事物？那些在基督内得以实现的事物。因此宗徒说：\Quote{因为天主的许诺在祂都是\InnerQuote{是}的} \BibleRef{格后 1:20}；也就是说，它们在祂内得以实现。又在另一处说：\Quote{这一切发生在他们身上，都是预像；而且是为我们这些活在时代终末的人写的} \BibleRef{格前 10:1}。他在别处又说：\Quote{因为基督是法律的终向} \BibleRef{罗 10:4}；同样在另一处：\Quote{不要让任何人在饮食、节期、月朔或安息日的事上判断你们，这些都是未来事物的影子} \BibleRef{哥 2:16-17}。因此，如果这一切都是未来事物的预像，那么帐棚节也是未来事物的预像。那么，我们来探讨这个节期是未来何事的预像。我已经解释了什么是帐棚节：它是帐棚的庆典，因为百姓从埃及获救后，在穿越旷野前往福地的途中，曾住在帐棚里。让我们观察这是什么，我们就会成为那事物；我说，我们这些基督的肢体，如果我们真是这样的话；但我们是祂使我们配得，不是我们自赚的。那么，弟兄们，让我们反省自己：我们曾被领出埃及，在那里我们曾是魔鬼（如同法郎）的奴隶；我们从事泥工，沉迷于世俗的欲望，在那里极其劳苦。正当我们仿佛在制砖时，基督大声呼喊：\Quote{凡劳苦和负重担的，你们都到我这里来}。此后，我们藉着圣洗，如同穿越红海——红色因基督的血所祝圣——被领出。所有追赶我们的仇敌都死了，即我们所有的罪都被涂抹了，我们被带到了对岸。因此，在如今，在我们到达福地（即永生的王国）之前，我们正在旷野中住在帐棚里。承认这些事的人就在帐棚中；因为必须有人承认这一点。因为那明白自己在此世是旅客的人，就在帐棚中。那看见自己叹息思念家乡的人，明白自己在异乡旅行。但既然基督的身体在帐棚中，基督就在帐棚中；在当时，祂是在帐棚中，但不是明显地，而是隐秘地。因为阴影还遮蔽着光；当光来到时，阴影就被除去。基督是隐秘的：祂在帐棚节中，却是隐藏的。如今，这些事已显明，我们承认自己正在旷野中旅行：因为我们若知道，我们就在旷野中。在旷野中是什么意思？在荒芜的沙漠中。为什么在荒芜的沙漠中？因为在这世上，我们在无水之地干渴。然而，让我们干渴，好使我们得以饱足。因为\Quote{饥渴慕义的人是有福的，因为他们要得饱饫} \BibleRef{玛 5:6}。我们的干渴从旷野的磐石中得到解渴：因为\Quote{那磐石就是基督}，它被杖击打使水流出来。但为使水流出来，磐石被击打了两次：因为十字架有两根木梁。因此，凡以预像方式发生的这一切，如今都向我们显明了。并非毫无意义地论主说：\Quote{祂上了庆节，但不是公开地，而是似乎在隐秘中}。因为祂自己在隐秘中就是那预表之物，因为基督在那同一个节期中被隐藏；因为那个节期本身预表了基督将来要在异地寄居的肢体。\par

10. \Quote{于是犹太人在庆节期间寻找祂}：在祂上去之前。因为祂的兄弟在祂之前上去，而祂没有在他们推测和希望的时候上去：这样也成全了祂所说的，不是这个（即你们希望我去的头一两天）。但后来祂上去了，如福音所述，在节期的中段，也就是说，那节期已过的天数与剩下的天数相等。因为据我们所知，他们连续几天庆祝同一个节期。\par

11. \Quote{他们就问说：\InnerQuote{他在哪里？} 关于他，在群众中起了许多窃窃私议。} 这窃窃私议从何而来？来自纷争。纷争是什么？\Quote{有人说：\InnerQuote{他是好人}；另有人说：\InnerQuote{不，他迷惑群众。}} 我们必须明白这一点适用于所有祂的仆人：现在正是论及他们。因为凡在某项神恩上成为卓越的人，注定会有人这样说：\Quote{他是好人}；另一些人则说：\Quote{不，他迷惑群众。} 这是为何？\BibleQuote{因为我们的生命已与基督一同藏在天主内。} \BibleRef{哥 3:3} 因此，在冬天，人们可能会说：这棵树是死的；例如，一棵无花果树、梨树或某种果树，它如同枯树，只要冬天未过，就无法看出是否如此。但夏天会证明，审判会证明。我们的夏天是基督的显现：\Quote{天主必显然而来，我们的天主，决不缄默；} \Quote{烈火在祂面前开道；} 那火 \Quote{将焚毁祂的仇敌；} 那火将擒住枯树。因为那时枯树将显露出来，当向它们说：\Quote{我饿了，你们没有给我吃；} 但在另一边，即在右边，将看到丰硕的果实和华美的枝叶；绿色将是永恒。因此，就如枯树一般，它们将被告知：\Quote{你们到永火里去。因为看，} 经上说，\BibleQuote{斧子已放在树根上；凡不结好果子的树，必被砍倒投入火中。} \BibleRef{玛 3:10} 因此，如果你们在基督内生长，就让人们论及你们：\Quote{他迷惑群众。} 这话曾论及基督本人；也论及基督的整个身体。要思想基督的身体仍在世上，仍在打谷场上；看它如何被糠秕亵渎。的确，糠秕和麦子一同被打下来；但糠秕被烧毁，麦子被扬净。当时论及主的话，每当它被论及任何基督徒时，都适用于安慰。\par

12. \Quote{但没有人敢公开讲论他，因为害怕犹太人。} 但谁因害怕犹太人而不讲论他？无疑是指那些说\Quote{他是好人}的人，而不是说\Quote{他迷惑群众}的人。至于那些说\Quote{他迷惑群众}的人，他们的喧嚷如同枯叶之声。\Quote{他迷惑群众}，他们越来越大声地叫嚷；\Quote{他是好人}，他们越来越拘谨地低语。但如今，弟兄们，尽管那使我们不朽的基督的荣耀尚未到来，然而如今，我说，祂的教会如此增长，祂已屈尊将其传遍全世界，以致现在\Quote{他迷惑群众}只是低语；而\Quote{他是好人}却越来越响亮地传扬开来。\par

\SourceRecord{Translated by John Gibb.；Nicene and Post-Nicene Fathers, First Series；Vol. 7.；Edited by Philip Schaff.；Buffalo, NY: Christian Literature Publishing Co.,；1888.；<http://www.newadvent.org/fathers/1701028.htm>.}

% structure: p0001 source=h1 target=section reason=page-title

\section{讲道集第29篇（若望福音7:14-18）}

1. 接下来，我们应依次查看福音中今天所读的部分，并依主所赐予的恩宠加以讲解。昨日读到此处：尽管他们在庆节日未在圣殿中见到主耶稣，却仍在谈论祂：\Quote{有人说：\InnerQuote{祂是好人；}但另有人说：\InnerQuote{不然，祂迷惑百姓。}}这话是为安慰日后宣讲天主圣言而将被称为迷惑者、却仍是真诚信实的人而说的（见：\BibleRef{格后 6:8}）。因为若引诱即是欺骗，那么基督不是迷惑者，祂的宗徒也不是，任何基督徒也不应是；但若引诱（即将人引开）是通过劝说使人从一处转向另一处，我们就应当探究其\Emph{从何处}和\Emph{向何处}：若从恶向善，则迷惑者是好人；若从善向恶，则迷惑者是恶人。那么，在那种意义上，即人被从恶引向善，但愿我们所有人既被称为迷惑者，事实上也真是迷惑者！\par

2. 然后主上圣殿去过节，\BibleQuote{约在节期的中间，并施教。}\BibleQuote{犹太人就惊讶说：这人没有学过，怎么通晓经书呢？}祂在暗中施教，祂公开宣讲，且未受阻止。因为那隐藏自己是出于榜样，而这公开显现是祂能力的暗示。但当祂施教时，\BibleQuote{犹太人就惊讶；}我想，确实所有人都惊讶，但并非所有人都皈依。为什么惊讶？因为大家都知道祂在哪里出生、在哪里长大；他们从未见过祂学习文字，却听祂辩论法律，引用法律中的证据，而这些证据除非读过是不能引用的，而除非学过文字是不能阅读的；因此他们惊讶。但他们的惊讶成了导师将真理更深入地渗透到他们思想的契机。实际上，因着他们的惊讶和言语，主说了某些深刻的话，值得更仔细地审视和讨论。为此，我恳劝你们，亲爱的诸位，要热切，不仅亲自聆听，也为我们祈祷。\par

3. 那么，主如何回答那些惊讶于祂没有学过却通晓经书的人呢？祂说：\BibleQuote{我的道理不是我的，而是那派遣我来者的。}这是第一个深奥之处。因为祂似乎用几句话说了相反的话。因为祂不是说：这道理不是我的；而是说：\BibleQuote{我的道理不是我的。}如果不是祢的，如何是祢的？如果是祢的，如何不是祢的？因为祢两者都说：既说\BibleQuote{我的道理，}又说\BibleQuote{不是我的。}因为如果祂说\Quote{这道理不是我的}，就不会有问题了。但现在，弟兄们，首先要好好思考这个问题，然后依次期待解答。因为看不见所提出的问题的人，怎能理解所阐述的内容呢？那么，探究的主题是祂所说的\Quote{我的、不是我的}，这看似矛盾；如何是\Quote{我的}，如何是\Quote{不是我的}？如果我们仔细看圣史自己在福音开头所说的话：\BibleQuote{在起初已有圣言，圣言与天主同在，圣言就是天主；}那么这个问题的解答就系于此。那么，父的道理若不是父的圣言，又是什么呢？因此，基督自己就是父的道理，如果祂是父的圣言。但由于圣言不能来自无，而只能来自某一位，祂便同时说\BibleQuote{自己的道理，}即祂自己，以及\BibleQuote{不是自己的，}因为祂是父的圣言。因为有什么比你自己更\Quote{属于你}的呢？又有什么比你自己更\Quote{不属于你}的呢，如果你的存在来自另一位？\par

4. 圣言就是天主；它也是稳定不变之道的圣言，并非像可由音节发出、转瞬即逝的，而是与父同在的。让我们归向这常在的道理，受声音的短暂声响的提醒。因为那短暂的事物并非如此提醒我们，以至于召叫我们趋向短暂的事物。我们被提醒去爱天主。我刚才所说的一切都是音节；它们穿过空气达到你的听觉，并通过发声而消逝：然而，我所劝告你们的，不应如此消逝，因为我劝你们所爱的那一位并不消逝；当你们被短暂的音节劝化而悔改时，你们将不会消逝，而将和那常在者一同常在。因此，在道理中有这伟大之事，这深邃而永恒、持久不变之物：一切在时间中消逝的事物，当它们用意良善、并非虚假地被提出时，都将我们召向它。因为事实上，我们通过声音产生的一切标记，都意指某种不是声音的东西。因为天主并不是两个短音节\Quote{Deus}，也不是我们所崇拜的两个短音节，也不是我们所钦崇的两个短音节，更不是我们渴望趋赴的两个短音节——这两个音节几乎在开始发声之前就已停止；在发这两个音时，直到第一个音消逝，第二个音才有空间。那么，还剩下某件伟大之物，它被称为\Quote{天主}，虽然当我们说\Quote{天主}这个词时这声音并不存留。因此，让你的思想转向基督的道理，你将到达天主的圣言；当你到达天主的圣言时，要思想这一点：\BibleQuote{圣言就是天主}，你将看到祂曾说\BibleQuote{我的道理}是真实的；也要思想圣言是属于谁的，你将看到祂说\BibleQuote{不是我的}是正确的。\par

5. 因此，简言之，亲爱的诸位，我以为主耶稣基督说：\Quote{我的教训不是我的}，意思等于说：\Quote{我不是从我自己来的。}因为，虽然我们宣称并相信圣子与圣父平等，在他们之间没有任何本性或实体的差异，也没有任何时间间隔隔开生育者与被生育者，然而我们这样说时，仍保持并守护这一点：一位是圣父，另一位是圣子。但若他没有圣子，他就不是圣父；若他没有圣父，他就不是圣子。然而圣子是出自圣父的天主；圣父是天主，却不是出自圣子。圣父是圣子的父，并非出自圣子的天主；而另一位是圣子的子，是出自圣父的天主。因为主基督被称为\OriginalTerm{光从光}{Light from Light}。那么，那不出自光的光，与那不出自光的相等之光，二者同是一光，而非两光。\par

6. 若我们明白了这一点，愿感谢归于天主；若有人未能充分明白，人已尽了其所能；至于其余，让他看看从何处可望明白。作为外面的工人，我们能栽种和浇灌，但使之生长的是天主。他说：\Quote{我的教训不是我的，而是那派遣我者的。}若有人说他尚未明白，请听劝告。因为所说的是伟大而深奥的事，主基督自己确实看到并非所有人都能明白如此深奥的事，于是他在下文给出了劝告。你愿意明白吗？请相信。因为天主借先知说过：\BibleRef{依 7:9}\Quote{除非你们相信，你们必不得明白。}与此相同，主在此处接着又加了这句话——\Quote{谁若愿意承行他的旨意，他就会知道这教训是否出于天主，或者我是凭我自己说话。}\Quote{谁若愿意承行他的旨意}是什么意思？我曾说过：若有人相信；我给了这劝告：若你尚未明白，我说，要相信。因为明白是信德的赏报。所以不要寻求明白为了相信，而要相信为了明白；因为\Quote{除非你们相信，你们必不得明白}。因此，当我为理解的可能性而劝告相信的顺从，并说我们的主耶稣基督在下文中正添加了这一点，我们发现他说：\Quote{谁若愿意承行他的旨意，就会知道这教训。}\Quote{知道}是什么意思？就是\Quote{明白}。但\Quote{谁若愿意承行他的旨意}是什么意思？就是相信。诚然，所有人都看出\Quote{知道}等于\Quote{明白}；但\Quote{谁若愿意承行他的旨意}这句话指相信，并非所有人都看出；为了更准确地看出这一点，我们需要主自己作诠释者，向我们表明承行父的旨意是否确实指相信。然而谁不知道，承行天主的旨意就是做天主的工，即做那中悦他的工？但主自己在另一处公开说：\BibleRef{若 6:29}\Quote{这是天主的工，就是要你们相信他所派遣的那一位。}\Quote{要你们相信他}，不是\Quote{要你们相信他的话}。但若你们相信他，你们就信他的话；然而信他的话的人不一定相信他。因为连魔鬼也信他的话，但他们不相信他。再者，关于他的宗徒，我们可以说：\Quote{我们相信保禄}，但不说\Quote{我们信靠保禄}；\Quote{我们相信伯多禄}，但不说\Quote{我们信靠伯多禄}。因为\BibleRef{罗 4:5}\Quote{对于信靠那使不虔敬者成义的人，他的信德被算为正义。}那么\Quote{信靠他}是什么意思？就是借着相信去爱他，借着相信去珍视他，借着相信进入他并成为他肢体的。因此，信德就是天主向我们要求的；而他若没有先赐予他所要求的，就找不到他所要求的。什么样的信德呢？就是宗徒在另一处最充分地界定的，说：\BibleRef{迦 5:6}\Quote{割损算不了什么，不割损也算不了什么，唯有以爱德行事的信德。}不是任何一种信德，而是\Quote{以爱德行事的信德}：让这信德在你内，你就会明白这教训。你究竟会明白什么呢？就是\Quote{这教训不是我的，而是那派遣我者的}；也就是说，你会明白基督——天主子，圣父的教训——不是从他自己而来，而是圣父之子。\par

7. 这句话推翻了\OriginalTerm{撒伯流异端}{Sabellian heresy}。撒伯流派竟敢断言圣子与圣父是同一位；名字有两个，但实体只有一个。若名字有两个，实体却是一个，就不会说\Quote{我的教训不是我的}。无论如何，主啊，若你的教训不是你的，那是谁的，除非另有其主？撒伯流派不明白你所说的话；因为他们看不到圣三，却追随自己内心的错误。让我们这些敬拜圣父、圣子、圣神之圣三与一体、唯一真主的人，明白基督的教训如何不是他的。他说他不是凭自己说话，原因在于：基督是圣父之子，圣父是基督之父；圣子是出自天主圣父的天主，但天主圣父是天主，并非出自天主圣子的天主。\par

8. \Quote{那凭自己说话的，寻求自己的光荣：\InnerQuote{这就是那被称为\OriginalTerm{假基督}{Antichrist}的,} 高举自己,} 正如宗徒所说，\Quote{超过一切称为天主的和受人崇拜的。}\BibleRef{得后 2:4} \Person{主}指明这同一件事——他将寻求自己的光荣，而非\Person{圣父}的光荣——对犹太人说：\Quote{我因我父的名而来，你们却不接待我；另有一位要因自己的名而来，你们将要接待他。}\BibleRef{若 5:45} 祂暗示他们将接待假基督，他将寻求自己名字的光荣，骄傲自大，并不实在；因此不稳固，却必定是毁灭性的。但我们的主\Person{耶稣基督}给我们显示了谦卑的伟大榜样：因为祂毫无疑问与\Person{圣父}同等，因为\Quote{在起初已有\OriginalTerm{圣言}{Word}，圣言与天主同在，圣言就是天主}；是的，毫无疑问，祂亲自说，并且最真实地说：\Quote{这么长久的时间，我和你们在一起，你还不认识我么，\Person{斐理伯}？看见我就看见了\Person{圣父}。}\BibleRef{若 14:9} 是的，毫无疑问，祂亲自说，并且最真实地说：\Quote{我与\Person{圣父}原是一体。}\BibleRef{若 10:30} 因此，如果祂与\Person{圣父}是一体，与\Person{圣父}同等，出自天主的天主，与天主同在的天主，同永恒，不朽，同样不变，同样无时间，同样创造并安排时间；却因为祂在时间中来临，取了奴仆的形状，与常人相似，\BibleRef{斐 2:7} 祂寻求\Person{圣父}的光荣，而非自己的；那么你，人啊，应当做什么呢？当你行善时，寻求自己的光荣；但当你行恶时，却图谋诬蔑天主？反省你自己：你是受造物，要承认你的造物主；你是仆人，不要轻视你的主；你是被收养的，并非因你自己的功德；你要寻求祂的光荣，是祂赐给了你这恩宠，使你成为被收养的人；你要寻求那一位的光荣，那从祂而出的\OriginalTerm{独生子}{Only-begotten}曾寻求祂的光荣。\Quote{但那寻求派遣祂者的光荣的，是真实的，在祂内毫无不义。} 在假基督内，却有不义，他不真实；因为他将寻求自己的光荣，而非派遣他者的光荣：因为他确实不是被派遣的，只是被准许来的。因此，我们所有属于基督身体的人，不要寻求自己的光荣，免得陷入假基督的圈套。但如果基督寻求了派遣祂者的光荣，我们岂不更应当寻求那造我们的天主的光荣吗？\par

\SourceRecord{Translated by John Gibb.；Nicene and Post-Nicene Fathers, First Series；Vol. 7.；Edited by Philip Schaff.；Buffalo, NY: Christian Literature Publishing Co.,；1888.；<http://www.newadvent.org/fathers/1701029.htm>.}

% structure: p0001 source=h1 target=section reason=page-title

\section{讲道集第三十篇（\BibleRef{若 7:19-24}）}

亲爱的诸位，我们之前向你们讲论的圣福音段落，紧接着就是今天刚才宣读的段落。门徒和犹太人都听见了主在说话；诚实之人和说谎者都听见了真理在说话；朋友和敌人都听见了爱德在说话；好人和坏人都听见了善在说话。他们听见了，但祂辨别；祂看见并预见了谁将因祂的讲论得益，谁将继续得益。在那些当时在场的人中，祂看见；在我们这些将要到来的人中，祂预知。因此，让我们聆听福音，就好比我们正在聆听主亲自临在一般；也不要说：\Quote{啊，那些能够亲眼看见祂的人多么有福！}因为他们中有许多人看见了，并且也杀了祂；而我们中有许多人没有看见祂，却已相信。因为那发自主口中的宝贵真理，既因我们的缘故而被记录下来，又因我们的缘故而被保存下来，因我们的缘故而被诵读，并要因我们的得救而继续诵读，直到世界终结。主在高天；但主，真理，也在这里。因为主的身体——就是祂从死者中复活的那身体——只能在一个地方；但祂的真理却弥漫各处。那么，让我们聆听主，并且也讲论祂关于祂自己的话所赐予我们的。\par

2. \BibleQuote{梅瑟岂不是}祂说，\BibleQuote{给了你们法律，而你们中却没有人遵行法律？你们为什么想要杀我？}因为你们想要杀我的原因正是：你们中没有人遵行法律；因为如果你们遵行法律，你们就会在法律的文字中认出基督，也就不会在祂与你们同在时杀害祂了。他们回答说；群众纷纷回答说，如同骚动的人群，所回答的没有秩序，只有混乱；总而言之，群众被搅乱了。看他们怎么回答：\BibleQuote{你附了魔；谁想要杀你？}仿佛说\Quote{你附了魔}比杀害祂更不恶劣。确实，祂——正在驱魔的那位——被说成附了魔。骚乱无序的群众还能说什么？搅动的污秽除了发臭还能做什么？群众被搅乱了；被什么搅动的？被真理。因为不健康的眼睛无法忍受光明的亮度。\par

3. 但是主，显然没有激动，而是安详地处于祂的真理中，不以恶报恶，不以辱骂还辱骂（\BibleRef{伯前 3:9}）；尽管如果祂对这些人说\Quote{你附了魔}，祂说的确实是真话。因为他们不会对真理说出这样的话，除非是魔鬼的谎言煽动了他们。那么祂如何回答呢？让我们心平气和地聆听，并啜饮这平静的话：\BibleQuote{我作了一件事，你们就都惊奇。}仿佛祂说：如果你们看见我所有的工作呢？因为他们所看见的祂的工作是世界，而他们却未看见那创造万有的祂；祂做了一件事，他们就激动，因为祂在安息日使一个人痊愈。仿佛，在安息日任何病人康复时，那使这人痊愈的并非祂，而他们之所以不快正是因为祂在安息日使一个人痊愈。因为还有谁使其他人痊愈，除了祂就是健康本身——祂甚至赐给走兽那祂赐给这人的健康？因为那是身体的健康。肉体的健康得以修复，而肉体仍然会死；当它被修复时，死亡只是被推迟，并非被除去。然而，即使是同一种健康，弟兄们，也是来自主，无论通过谁赐予；无论通过谁的照顾和事工传授，都是由祂赐予，一切健康都来自祂。在圣咏中向祂祈祷说：\BibleQuote{主啊，祢要拯救人和走兽；天主，祢的仁慈何其丰盛。}因为祢是天主，祢丰盛的仁慈甚至延伸到人类肉体的安全，甚至延伸到哑巴牲畜的安全；但那赐予人与走兽共有的肉体健康的祢，难道没有为人类保留一种健康吗？当然有另一种健康，不仅不是人与走兽共有的，甚至在人之间，也不是好人与坏人共有的。简言之，当祂在那里谈到这种人与牲畜共同接受的健康时，因为那种健康是人——但只有好人——应当盼望的，祂在继续时补充说：\BibleQuote{但人的儿子们将在祢翅膀的庇护下投靠。他们必因祢房屋的丰盛而满足；祢必叫他们喝祢喜悦的急流。因为生命的泉源在祢那里；在祢的光中，我们必得见光。}这就是属于好人的健康，那些祂称为\BibleQuote{人的儿子们}；而祂在前面说：\BibleQuote{主啊，祢要拯救人和走兽。}那么怎样？那些人不是人的儿子吗？为何在说了\Emph{人}之后，祂接着又说\BibleQuote{\Emph{但人的儿子们}}：仿佛\Emph{人}和\Quote{人的儿子们}指不同的东西？然而，我相信圣神说这话并非无意，而是有所区分。\Emph{人}一词指第一个亚当，\Emph{人的儿子们}指基督。也许，的确，\Emph{人}指第一个人；而\Emph{人的儿子们}指人子。\par

4. \Quote{我作了一件事，你们都惊讶。}接着祂立即加上：\Quote{梅瑟所以给你们施行割损。}你们从梅瑟领受割损，这事做得很好。\Quote{不是因为割损是出于梅瑟，而是出于祖先；}因为亚巴郎是从主首先领受了割损。\BibleRef{创 17:10}\Quote{你们也在安息日给人施行割损。}梅瑟已定了你们的罪：你们在律法上领受了要在第八天施行割损；你们在律法上领受了要在第七天停止工作；\BibleRef{出 20:10}如果孩子出生后的第八天正好是七日中的第七天，你们要怎样做？你们是停工守安息日呢，还是施行割损以完成第八天的圣事？但我知道你们所做的是什么，祂说。\Quote{你们给人施行割损。}为什么？因为割损关系到一种救恩的印记，人在安息日不应停止救恩的工作。所以你们不要\Quote{向我发怒，因为我在安息日使一个人完全痊愈了。}\Quote{如果，}祂说，\Quote{人在安息日受割损，为不使律法被破坏}（因为在割损的规定中，梅瑟所命定的是某种救恩的事），那么我在安息日施行医治，你们为什么向我发怒呢？\par

5. 或许，那割损正指向主自己，他们因祂施行医治和痊愈而愤怒。因为割损被命令在第八天施行：割损若不是剥去肉体的皮，又是什么？那么，这割损象征着从心中除去肉欲。因此，它被赐予并命令在那肢体上施行，并非无因；因着那肢体，有死的人类得以繁衍。藉着一人，死亡来到，正如藉着一人，死人复活；\BibleRef{格前 15:21}并且藉着一人，罪恶进入了世界，藉着罪恶，死亡。\BibleRef{罗 5:12}因此，每个人生来都带着包皮，因为每个人都带着繁衍的缺陷出生；而天主洁净我们，无论是生来就有的缺陷，还是我们因恶行所加的缺陷，除非藉着那石刀，主基督。因为基督就是磐石。他们曾用石刀施行割损，藉着磐石之名预表基督；然而当祂在他们中间时，他们却不认识祂，反倒图谋杀害祂。但为何在第八天？除非因为安息日之后，主在主的日子的复活。因此基督的复活，虽然发生于祂受难后的第三天，但在七日中的第八天，正是那复活使我们受割损。听听那些曾被真石所割损的人，宗徒劝诫他们：\Quote{所以，你们既然与基督一同复活了，就该追求天上的事，那里有基督坐在天主的右边；你们要思念天上的事，不要思念地上的事。}\BibleRef{哥 3:1}祂是向受过割损的人说话：基督已复活；祂从你们身上除去了肉欲、邪恶的情欲、你们生来就有的多余之物，以及你们因恶行所加的更糟之物；既然你们被磐石所割损，为何仍思念地上的事？最后，至于\Quote{梅瑟给了你们法律，你们在安息日给人施行割损，}你们要明白，这所象征的善工就是我所作的，即在安息日使一个人完全痊愈；因为他得了医治，身体痊愈；他也信了，灵魂痊愈。\par

6. \Quote{不要按照外貌判断，但要作公义的判断。}这是什么意思？刚才，你们这些按梅瑟法律在安息日施行割损的人，不向梅瑟发怒；而我因在安息日使一个人痊愈，你们却向我发怒。你们是按外貌判断；要留意真理。我并不把自己置于梅瑟之上，主说，祂也是梅瑟的主。所以，请把我们看作如同两个人，都是人；在我们中间判断，但判断真实的判断；不要因尊敬我而定他的罪，而要因认识他而尊敬我。因为祂在另一处对他们说：\Quote{如果你们信梅瑟，也必信我，因为他曾论及我写了书。}\BibleRef{若 5:46}但在此处，祂不愿说这话，而是将自己和梅瑟如同放在他们面前受判断。因着梅瑟的法律，你们在安息日施行割损，即使那天正好是安息日；难道我不该在安息日显示医治的恩惠吗？因为割损的主和安息日的主是同一位，祂也是健康的创造者；你们被禁止在安息日做的是劳作的工作；如果你们真理解什么是劳作的工作，你们就不犯罪。因为犯罪的人是罪恶的奴隶。在安息日医治一个人难道是劳作的工作吗？你们吃、喝，至少根据我们的主耶稣基督的劝诫和祂的话：你们在安息日为什么吃、喝，不就是因为你们所做的事为了健康吗？由此你们表明，健康的工作无论如何不应在安息日被忽略。因此\Quote{不要按外貌判断，但要作公义的判断。}请把我看作同一个人；把梅瑟看作一个人：如果你们按真理判断，你们就不会定梅瑟或我的罪；当你们认识真理时，你们就会认识我，因为我是真理。\par

七、弟兄们，要摆脱主在此所指出的恶习，即不凭外貌判断，而要持守公义的判断，这需要在此世付出巨大的辛劳。主的确曾告诫犹太人，但祂也警告了我们；祂谴责了他们，却教导了我们；祂责备了他们，却鼓励了我们。我们不要因为当时不在场，就以为这话不是对我们说的。这些话被记载下来，被人诵读；当它被宣读时，我们听到了；但我们听到的是对犹太人说的；我们不要置身事外，只看着祂谴责敌人，而自己却做着那真理可能在我们内所谴责的事。犹太人确实凭外貌判断，但正因为如此，他们不属于\WorkTitle{新约}，在基督内没有天国，也没有加入圣天使的行列；他们向主寻求现世的事物：应许之地、战胜仇敌、生育繁盛、子孙增多、果实丰饶——这一切确实由真实而良善的天主应许给了他们，但却是作为属血肉之人应许给他们的——所有这些为他们构成了\WorkTitle{旧约}。什么是\WorkTitle{旧约}？可以说是属于旧人的产业。我们已经被更新，成为了新人，因为那新人的那一位已经来了。有什么比由童贞女所生更新奇呢？因此，因为祂内没有需要教训来更新的事物，因为祂没有罪，所以赐给了祂一个新的出生起源。在祂内是新生的方式，在我们内是新人。什么是新人？是从旧态中更新的人。更新为什么？为渴望天上之事，为渴慕永恒之物，为切慕那至高无上、无惧仇敌的家乡，在那里我们不会失去朋友，也不怕敌人；在那里我们怀着纯良的情感生活，毫无匮乏；在那里不再有进步，因为无人退步；在那里无人出生，因为无人死亡；在那里没有饥渴；不朽是圆满，真理是我们的食粮。我们既拥有这些应许，属于\WorkTitle{新约}，又成为新产业的继承人，并与主自己同为继承人，我们便拥有与他们截然不同的希望：我们不可凭外貌判断，而要持守公义的判断。\par

八、谁是不按外貌判断的？是那平等地爱人的人。平等的爱使人不偏袒。并不是当我们按照人们的不同身份给予不同程度的尊敬时，就应当担心我们是在偏袒人。因为当我们判断两人之间的事，有时是亲属之间，有时会遇到必须在父子之间作出判断；父亲投诉儿子不肖，或儿子投诉父亲苛刻；我们看重儿子对父亲应有的尊敬；我们不使儿子在尊敬上与父亲同等，但如果他有正当的理由，我们也会给他优先：让我们在真理上视儿子与父亲平等，这样我们就能给予应得的尊敬，使公义不破坏功绩。这样，我们因主的话语而获益，而为了我们能获益，我们蒙受祂恩宠的助佑。\par

\SourceRecord{Translated by John Gibb.；Nicene and Post-Nicene Fathers, First Series；Vol. 7.；Edited by Philip Schaff.；Buffalo, NY: Christian Literature Publishing Co.,；1888.；<http://www.newadvent.org/fathers/1701030.htm>.}

% structure: p0001 source=h1 target=section reason=page-title

\section{论讲 31（\BibleBook{若望}{福音} 7:25-36）}

1. 亲爱的诸位，你们记得前几次的讲论——因为福音经文读过了，我们也按自己的能力讨论过了——主耶稣是如何好像秘密地上到庆节去的：不是因为祂害怕被人捉住——祂具有不被捉住的能力——而是为指明，即使在犹太人庆祝的这个庆节上，祂自己也是隐藏的，并且庆节的奥秘就是祂自己。那么，在今日所读的经文中，那原被认为是胆怯的事，如今显为权能；因为祂在庆节上公开说话，以致群众惊讶，并且说了我们在诵读经文时所听到的话：\Quote{这不是他们想要杀的那个人吗？看，祂公开说话，他们什么也不说。难道长官们真的知道这就是基督吗？} 那些知道祂被如何激烈地追捕的人，惊讶于祂被阻止不被捉拿的能力。然后，不完全理解祂的能力，他们猜想这是长官们的知识：认为这些长官知道祂就是基督，因此他们饶恕了他们曾如此急切地想要处死的那一位。\par

2. 然后，那些曾说过\Quote{难道长官们知道这就是基督吗？}的人，在自己中间提出了一个问题，他们因此认为祂不是基督；因为他们又附加说：\Quote{但我们知道这个人是从哪里来的；可是基督来时，没有人知道祂是从哪里来的。} 关于犹太人中间产生的这个意见，即\Quote{基督来时，没有人知道祂是从哪里来的}（因为这意见不是没有理由而产生的），如果我们查考圣经，弟兄们，我们就会发现，圣经已经宣布了基督\BibleQuote{将被称为纳匝肋人}。\BibleRef{玛 2:23} 因此，他们预言了祂是从哪里来的。再者，如果我们寻找祂出生的地方，即祂按出生是从哪里来的，这对犹太人也不是隐藏的，因为圣经预先宣告了这些事。当贤士们因星出现而寻找祂要朝拜祂时，他们来到黑落德那里，告诉他他们寻找什么以及他们的意图；黑落德召集了精通法律的人，问他们基督应该在哪里出生；他们告诉他\BibleQuote{在犹大的白冷}，并且引用了先知的见证。\BibleRef{玛 2:6} 因此，如果先知们既预言了祂肉身的起源地，也预言了祂母亲将生下祂的地方，那么，我们刚听到的犹太人中间这个意见是从哪里来的呢？岂不是因为圣经预先宣告了并且预言了这两点吗？就祂是人而言，圣经预言了祂将从哪里来；就祂是天主而言，这对不虔敬的人是隐藏的，需要虔敬的人去发现。而且，他们说了这话：\Quote{基督来时，没有人知道祂是从哪里来的}，因为依撒意亚所说的话在他们中产生了这个意见，即\BibleQuote{谁能述说祂的世代？} \BibleRef{依 8:8} 总之，主亲自对这两点都作了回答：他们既知道祂是从哪里来的，又不知道；为此，祂要见证那先前关于祂预言的圣预言，既关乎软弱的人性，也关乎威严的神性。\par

3. 因此，请听主的话，弟兄们；请看祂如何向他们确认了他们所说的两点：\Quote{我们知道这个人是从哪里来的}，以及\Quote{基督来时，没有人知道祂是从哪里来的}。\BibleQuote{于是基督在圣殿里大声呼喊说：你们认识我，也知道我从哪里来；但我不是由我自己来的，那派遣我来的那位是真实的，你们却不认识祂。} 这是说：你们既认识我，又不认识我；你们既知道我从哪里来，又不知道我从哪里来。你们知道我从哪里来：纳匝肋的耶稣，祂的父母你们也知道。因为在这种情况下，只有童贞女的诞生是隐藏的，但她的丈夫是见证人；因为那能够忠实地宣告这事的人，也能够像一个丈夫那样嫉妒。因此，除了童贞女的诞生，他们知道关于耶稣属于人的一切：面容是人所知的，故乡是人所知的，家族是人所知的；祂在哪里出生，需要询问才能知道。所以，祂正确地说了：\BibleQuote{你们认识我，也知道我从哪里来}，这是按祂所取的血肉和人的样式而言；但按祂的神性：\BibleQuote{但我不是由我自己来的，那派遣我来的那位是真实的，你们却不认识祂}；但为了你们能认识祂，要相信祂所派遣的那位，你们就会认识祂。因为\BibleQuote{从来没有人见过天主，只有独生子，祂在圣父怀里，将祂彰显出来}：\BibleRef{若 1:8} 以及\BibleQuote{除了圣子，没有人认识圣父，除非圣子愿意启示给谁。} \BibleRef{玛 11:27}\par

4. 最后，当祂说了\BibleQuote{但那派遣我来的那位是真实的，你们却不认识祂}之后，为向他们指明他们从哪里可以认识他们所不认识的那位，祂接着说：\BibleQuote{我认识祂。} 因此，要从我这里认识祂。但为什么我认识祂？\BibleQuote{因为我从祂而来，祂派遣了我。} 祂荣耀地显示了两点。\BibleQuote{我从祂而来}，祂说；因为圣子来自圣父，并且圣子无论是什么，都是来自祂作为圣子的那位。因此，我们说主耶稣是天主出自天主；我们不说圣父是天主出自天主，而只是天主；我们说主耶稣是光出自光；我们不说圣父是光出自光，而只是光。因此，祂所说的\BibleQuote{我从祂而来}属于此。至于你们在肉身中看见我，\BibleQuote{祂派遣了我。} 当你们听到\BibleQuote{祂派遣了我}时，不要理解为本质的不同，而是生者的权威。\par

5. \Quote{他们就想捉拿祂；可是没有人向祂下手，因为祂的时辰还没有到。}这是因为祂不愿意。所谓\Quote{祂的时辰还没有到}是什么意思呢？主并非生于命运。你们不可以这样相信，更何况创造你们的那一位呢？若是你们的时辰就是祂的善意，那么祂的时辰除了祂的善意还能是什么？因此，祂所指的并非祂不得不死的时辰，而是祂愿意被处死的时辰。祂在等待祂应死的时刻，因为祂也在等待祂应生的时刻。宗徒论及这个时刻说：\BibleQuote{但时期一满，天主就派遣了自己的儿子来。}\BibleRef{迦 4:4}为此许多人说：基督为什么不早些来？我们必须回答：因为时期尚未满全，而制定时期的那一位亲自定下了界限；祂知道祂应在何时来临。首先，祂必须经过漫长的年代和岁月被预言；因为将要来临的并不是小事：那永远被拥有的那一位，必须被长期预告。将要来临的审判官越伟大，前驱的队伍就越长。总之，时期一满，那要把我们从时间中拯救出来的那一位也来到了。因为从时间中被拯救出来，我们将会到达那没有时间的永恒：在那里不会说\InnerQuote{时辰何时到来}，因为白昼是永恒的，既无昨日在先，也无明日隔断。但在今世，日子流转不息，有些逝去，有些来临；没有一个停留；我们说话的时刻彼此交替，第一个音节不等第二个音节发出就已消失。自从我们开始说话，我们已经老了一些；毫无疑问，现在比早晨更老了；可见在时间中没有事物固定不变。因此，我们应该爱慕那创造时间的主，以便从时间中被解救出来，固定在永恒中，那里不再有时间的变化。我主耶稣基督的仁慈何其伟大，祂为了我们，在时间中成为那创造时间的主；祂在万物之中成为那创造万物的主；祂成为祂所造的。因为祂成了祂所造的：祂原是造人的天主，却成了人，以免祂所造的遭受灭亡。按照这一救恩计划，祂诞生的时辰已到，祂便诞生了；但祂受难的时辰尚未到，因此祂尚未受难。\par

6. 总之，为使你们知道这话所指的，不是祂必须死的必然性，而是祂的能力——我说这话，是因为有些人在听到\Quote{祂的时辰还没有到}时，就执意相信命运，他们的心就变得愚昧——那么，为使你们知道祂有死的能力，请回想祂的受难，看看被钉在十字架上的祂。祂悬挂在木架上时说：\BibleQuote{我渴。}他们听见后，就用海绵蘸了醋绑在芦苇上送给祂喝。祂接过来，说：\BibleQuote{完成了；}便低下头，交付了灵魂。你们看见了祂死的能力：祂一直等到关于祂所预言的一切都成就了，才在死前完成这一切。因为先知曾说过：\BibleQuote{他们拿苦胆给我作食物，我渴了，他们拿醋给我喝。}祂等到这一切都成就了；之后，祂说：\BibleQuote{完成了；}祂凭能力离去，因为祂不是出于必然而来的。因此，有些人更惊叹祂死的能力，胜过惊叹祂行奇迹的能力。他们来到十字架前，要从木架上取下尸体，因为安息日将近，发现两个强盗还活着。十字架的刑罚之所以更残酷，是因为它使人长时间受折磨，所有被钉十字架的人都是慢慢死去。但为了不让强盗继续挂在木架上，他们打断了他们的腿，使他们死去，以便取下。然而，主被发现已经死了，\BibleRef{若 19:28}人们惊讶不已；那些在祂活着时藐视祂的人，在祂死后竟如此惊讶，以致有人说：\BibleQuote{这人真是天主子。}\BibleRef{玛 27:54}同样，弟兄们，当祂对那些寻找祂的人说\BibleQuote{我就是}时，他们便倒退，都跌倒在地。\BibleRef{若 18:6}由此可见，祂拥有至高无上的权能。祂不是被迫在某个时辰死去；祂等待那个时辰，使祂的意愿得以适当地成就，而非使必然性违背祂的意愿而实现。\par

7. \Quote{但群众中有许多人信了祂。}主医治了谦卑和贫穷的人。首领们却疯狂，不仅不承认这位医师，甚至急于杀害祂。有一群百姓很快看见了自身的疾病，立即认出了祂的良药。请看那群百姓因祂的奇迹而说道：\Quote{基督来的时候，难道会比这人行更多的神迹吗？}当然，除非有两位基督，否则这位就是基督。因此，他们这样说，就信了祂。\par

8. 但是那些首领听见群众的信心和百姓中对基督的赞美声，\Quote{就派遣差役去捉拿祂。}捉拿谁？那位还不愿意被捉拿的。因为那时祂不愿意，他们无法捉拿，就被派去听祂教导。教导什么？\Quote{于是耶稣说：\InnerQuote{我还有不多的时候和你们同在，然后我就往派遣我者那里去。}}你们现在想做的事，你们将要去做，但还不是现在；因为我现在还不愿意。为什么我现在还不愿意？因为\Quote{我还有不多的时候和你们同在，然后我就往派遣我者那里去。}我必须完成我的救恩计划，然后这样去受难。\par

9. \BibleQuote{你们要寻找我，却找不着；我所在的地方，你们不能去。} 此处祂已预言了祂的复活；因为他们当时不肯承认祂，后来看见群众已经相信祂，才寻找祂。因为当主复活并升天之后，行了伟大的奇迹。然后祂的门徒行了大事，但祂藉着他们行事，如同祂亲自行事一样：因为祂曾对他们说：\BibleQuote{没有我，你们什么也不能做。} \BibleRef{若 15:5} 当坐在殿门口的那个瘸子因\Person{伯多禄}的声音起来，用脚行走，使人惊奇时，\Person{伯多禄}对他们这样说：这不是凭他自己的能力做的，而是凭他们所杀害的那位的德能。\BibleRef{宗 3:2} 许多人心感刺痛，说：\Quote{我们该做什么？} 因为他们看见自己被一种极大的不敬之罪所束缚，竟杀害了他们本应敬畏和崇拜的那位；而且他们认为这个罪是不可赦免的。的确，这是一件大恶，想到它可能使人绝望；然而，主在十字架上时曾屈尊为他们祈祷，所以他们不应该绝望。因为祂曾说：\BibleQuote{父啊，宽恕他们吧！因为他们不知道他们做的是什么。} \BibleRef{路 23:34} 祂看见在众多外人中有一些属于自己的人；祂为他们求宽恕，而那时这些人还在伤害祂。祂不顾自己被他们处死，只在乎自己为他们而死。宽恕给他们是一件大事，他们所做并为他们所做也是一件大事，使任何人都不应因自己罪的赦免而绝望，既然杀害基督的人都得到了宽恕。基督为我们死了，但祂确实不是被我们处死的吗？然而那些人确实看见基督死于他们的恶行；并且他们相信了基督赦免他们的恶行。直到他们喝了自己所流的血，他们才对自己的救恩绝望。因此祂说了这话：\BibleQuote{你们要寻找我，却找不着；我所在的地方，你们不能去；} 因为他们在复活后要寻找祂，心中被悔恨刺痛。祂没有说\Quote{我将要去的地方}，而是说\Quote{我所在的地方}。因为基督始终在那将要返回的地方；因为祂的来临是这样的：祂没有离开那个地方。因此祂在别处说：\BibleQuote{没有人升过天，除了那从天上降下来的人子，祂是在天上的。} \BibleRef{若 3:13} 祂没有说\Quote{祂曾在天上}。祂在地上说话，却宣告同时在天上。祂的来临是这样的：祂没有从那里离开；祂的返回也是不抛弃我们。你们惊奇什么？这是天主的工作。因为人，就其身体而言，是在一个地方，并从那个地方离开；当他来到另一个地方时，他就不在原先的地方了；但天主充满万有，处处皆是；祂不被空间所局限。然而主基督，就祂可见的肉身而言，在地上：就祂不可见的尊威而言，在天上也在世上；因此祂说：\BibleQuote{我所在的地方，你们不能去。} 祂没有说\Quote{你们将不能}，而是说\Quote{你们现在不能去}；因为那时他们是不能的。并且你们要知道这话不是要使人绝望，祂也对门徒说了类似的话：\BibleQuote{我去的地方，你们不能去。} \BibleRef{若 13:33} 然而当祂为他们祈祷时，祂说：\BibleQuote{父啊！我愿我在哪里，他们也同我在一起。} \BibleRef{若 17:24} 最后，祂向\Person{伯多禄}解释这话，对他说：\BibleQuote{我去的地方，你现在不能跟随我，但后来却要跟随我。} \BibleRef{若 13:36}\par

10. \BibleQuote{于是犹太人说}不是对祂说，而是\BibleQuote{彼此说：这人要往哪里去，使我们找不着呢？难道祂要往散居在外邦人中的人那里去，教训外邦人吗？} 因为他们不知道自己在说什么；但按照祂的旨意，他们却预言了。主确实要到外邦人那里去，不是以身体临在，但仍以祂的脚去。祂的脚是什么？就是\Person{扫禄}想要通过迫害践踏的，那时元首向他喊道：\BibleQuote{扫禄，扫禄，你为什么迫害我？} \BibleRef{宗 9:4} 祂所说的话是什么意思：\BibleQuote{你们要寻找我，却找不着；我所在的地方，你们不能去？} 为什么主说这话，他们不知道，然而他们却确实在不知情中预言了将要发生的事。因为主所说的是他们不知道那地方，如果那地方可以这样称呼的话，就是圣父的怀抱，基督从未离开过那里；他们也不配想象基督所在的地方，基督从未离开过那里，祂将返回那里，祂一直住在那里。人心如何能想象这事，更不用说用言语解释了呢？这事他们完全不明白；然而藉此机会，他们预言了我们的救恩，即主要到外邦人的散居地去，并要应验他们所读却不明白的经书。\Quote{我所不认识的人民侍奉了我，一听见就服从了我。} 他们亲眼看见祂，却不听祂；那些耳中听见祂声音的人却听从了祂。\par

11. 因为那行将来临的外邦人的教会，患血漏的妇人是一个预象：她触摸了，却没有被看见；她未被认出，却得了痊愈。主问的：\Quote{谁摸了我？}实际上是一个预象，好似不知而医治了她这位陌生者：祂对外邦人也是如此。我们并未在肉身中认识祂，却被堪当吃祂的肉，成为祂肉身的肢体。如何成就的呢？因为祂派遣了。派遣谁？祂的传报者，祂的门徒，祂的仆人，祂所救赎的——由祂所造，却又被祂救赎，也是祂的弟兄。关于他们的一切，我只稍说一句：祂自己的肢体，祂自己；因为祂派遣了祂自己的肢体到我们这里来，并使我们成为祂的肢体。然而，基督并未以犹太人所看见并鄙视的那肉身形态临到我们中间；因为关于祂也说过这样的话，正如宗徒所说：\BibleQuote{我说，基督为了天主的真理，成了割损的仆役，以证实所应许给祖先的恩许。}\BibleRef{罗 15:8}祂必须来到那些曾被许诺的人那里，因为祂是由他们的祖先并为了他们的祖先而被许诺的。为此祂自己也说：\BibleQuote{我被派遣，只是为了以色列家迷失的羊。}\BibleRef{玛 15:24}但宗徒在接下来的话中说了什么？\Quote{而外邦人因天主的仁慈而光荣天主。}再者，主自己说：\BibleQuote{我还有别的羊，不属于这羊栈。}\BibleRef{若 10:16}祂曾说过\BibleQuote{我被派遣，只是为了以色列家迷失的羊}，又怎么会有别的羊——祂未被派遣去的——呢？除非祂暗示说，祂所显示的肉身临在只针对犹太人——那些看见并杀害祂的人。然而他们中许多人，无论是之前还是之后，都相信了。初次收割从十字架上被簸扬，好让种子存在，使另一次收割得以生长。但在现今这时刻，当被福音的名声及其芬芳所唤醒时，祂在各民族中的信者相信了；祂将要成为外邦人的期待，当祂——那已经来过的——再来时；当祂——那当时未被一些人看见、却被另一些人看见的——将被众人看见时；当祂来审判——祂曾来受审判——时；当祂来区分——祂来时是未被区分——时。因为基督未被恶人辨识，却与恶人一同被定罪；因为关于祂曾说：\BibleQuote{祂被列于罪犯之中。}\BibleRef{依 53:12}强盗逃脱了，基督被定罪。被控罪行的人得了赦免，释放所有承认祂之人罪债的祂却被定罪。然而，就连十字架本身，你若善加思量，也是一个审判台；因为审判者被悬在中间，一个相信的强盗得释放，另一个辱骂的被定罪。\BibleRef{路 22:43}祂已经预示了祂将要对活人与死者所行的事：有些祂将置于右边，有些置于左边。那个强盗如同将要站在左边的人，另一个如同站在右边的。祂正在受审，却宣示了审判。\par

\SourceRecord{Translated by John Gibb.；Nicene and Post-Nicene Fathers, First Series；Vol. 7.；Edited by Philip Schaff.；Buffalo, NY: Christian Literature Publishing Co.,；1888.；<http://www.newadvent.org/fathers/1701031.htm>.}

% structure: p0001 source=h1 target=section reason=page-title

\section{讲道录 32（\BibleRef{若 7:37-39}）}

1. 在犹太人对主耶稣基督的纷争与疑惑中，祂所说的话里，有些使一些人困惑，另一些则教训人：\BibleQuote{在那庆节的最后一天}（因为这事就发生在那时），即所谓帐棚节，也就是搭帐篷的庆节，关于这庆节，亲爱的诸位，你们记得我们已经讲论过了。主耶稣基督在呼喊，不是随意地说话，而是高声喊叫，凡口渴的，都可以到祂那里来。如果我们口渴，让我们前去；不是用我们的脚，而是用我们的情感；让我们前去，不是离开我们的地方，而是藉着爱。虽然，按内心的人，那爱的也确实从一处移动。但身体移动是一回事，心移动是另一回事：用身体移动的人，藉身体的运动改变他的地方；用心移动的人，藉心的运动改变他的情感。如果你爱一样东西，而以前爱的是另一样，你现在就不在你过去所在的地方了。\par

2. 因此，主向我们呼喊：因为\BibleQuote{祂站着大声说：谁若口渴，到我这里来喝罢！凡信从我的，就如经上说，从他的心中要流出活水的江河。}我们不必耽搁去探究这是什么意思，因为福音作者已经解释了。因为主说\BibleQuote{谁若口渴，到我这里来喝罢}和\BibleQuote{凡信从我的，从他的心中要流出活水的江河}，福音作者随后解释说：\BibleQuote{祂这话是指那信祂的人将要领受的圣神；因为圣神还没有赐下，因为耶稣还没有受到光荣。}因此，有一种内在的口渴和一个内在的肚腹，因为有一个内心的人。这内心的人固然是看不见的，但外在的人是可见的；然而内心比外在更好。这看不见的反而是更被爱的；因为确实，内心的人比外在的人更被爱。这怎能确定呢？让每个人在自己身上证明它。因为虽然那些生活不善的人可能将他们的心神屈服于身体，但他们仍然愿意活着，而活着只是心神的属性；那些治理的，比那些被治理的更显示自己。是心神治理，身体被治理。每个人都喜乐于快乐，并通过身体接受快乐：但将心神与身体分开，身体就没有什么可快乐的了；如果身体有快乐，那是心神在快乐。如果心神因它的居所而喜乐，它难道不该因它自己而喜乐吗？如果心神有它可以在自身之外享乐的东西，它是否在自身之内没有快乐呢？十分确定，人爱自己的灵魂胜过爱自己的身体。不仅如此，人甚至在另一个人身上爱灵魂胜过爱身体。在朋友中，爱更纯洁、更真诚的，是什么？在朋友中被爱的，是心神，还是身体？如果爱的是忠信，那就是爱心神；如果爱的是善意，心神是善意的所在：如果你在别人身上爱的是他也爱你本身，那么你爱的就是心神，因为不是肉身，而是心神在爱。因此你爱，是因为他爱你：问为什么他爱你，然后看看你爱的是什么。因此，它更被爱，却看不见。\par

3. 我想再说一点，亲爱的诸位，好使你们更清楚地看见，心神多么被爱，如何被置于身体之上。连那些好色之徒，他们喜爱身体之美，被四肢的匀称所吸引，当他们被爱时就爱得更深。因为当一个人爱，却发现自己被憎恨时，他感到愤怒多于喜爱。为什么他感到愤怒多于喜爱？因为他所付出的爱没有得到回报。因此，即使喜爱身体的人也渴望被爱，而且当他们被爱时，这更使他们喜乐，那么，对心神的爱慕者，我们该说什么呢？如果心神的爱慕者是伟大的，那么对使心神美丽的天主的爱慕者，我们又该说什么呢？因为正如心神赋予身体优雅，同样，天主赋予心神优雅。是心神使身体有了可爱之处；当心神离开时，身体就成了你害怕的尸首，无论你多么喜爱它美丽的四肢，你都赶紧将它埋葬。因此，身体的装饰是心神；心神的装饰是天主。\par

4. 因此，主向我们呼喊，要我们若内心口渴，就去喝，并且祂说，当我们喝了，活水的江河要从我们的心中流出。内心之人的肚腹就是心灵的良心。喝了那水，良心被净化，就开始生活；饮入之后，它要拥有一个泉源，它本身要成为一个泉源。什么是泉源？什么是那从内心之人的肚腹流出的江河？就是善意，人藉此顾及邻人的利益。因为如果他认为所喝的只应满足自己，那么活水就不会从他的肚腹流出；但如果他急于顾及邻人的利益，那么他就不干涸，因为有水流。现在我们看看那些信主的人所喝的是什么；因为我们确实是基督徒，如果我们相信，我们就喝。每个人都当知道在自己内他是否喝，是否靠他所喝的而生活；因为泉源不会离弃我们，如果我们不离弃泉源。\par

5. 圣史解释了我所说的主呼喊何事，祂邀请了何种饮品，祂为他们准备了何种饮品，说：\Quote{但祂这话是指着那信祂的人将要领受的圣神说的：因为那时圣神还没有赐下，因为耶稣还没有受到光荣。}祂说到何种神，若不是圣神？因为每个人本身都有自己的心神，我在向你们推荐思想之考虑时曾提及这一点。因为每个人的思想即是其自己的心神：对此，保禄宗徒说：\Quote{除了在人里面的心神，谁能知道人的事呢？}然后他又补充说：\Quote{同样，除了天主的圣神，也没有人知道天主的事。}\BibleRef{格前 2:11} 没有人知道属于我们的事，除了我们自己的心神。我确实不知道你们的想法，你们也不知道我的想法；因为我们内心所思想的事，是属于我们自己的，为我们所独有；每个人的心神是其思想的证人。\Quote{同样，除了天主的圣神，也没有人知道天主的事。}我们借我们的心神，天主借祂的神；然而，天主借祂的圣神也知道我们内心的事；但我们若没有祂的圣神，就不能知道发生在天主内的事。然而，天主知道我们内甚至我们自己不知道的事。因为当伯多禄从主那里听到他将三次否认祂时，他并不知道自己的软弱：病人不知自己的状况；医生知道他生病。因此，有些事是天主在我们内知道的，而我们自己却不知道。然而，就人而言，没有人像自己那样知道另一个人：另一个人不知道他内发生的事，只有他自己的心神知道。但当我们领受天主的圣神时，我们也知道发生在天主内的事：不是全部，因为我们还没有领受全部。我们从抵押中知道许多事；因为我们领受了一个抵押，这抵押的圆满将在以后赐予。同时，让这抵押在此世旅途中安慰我们；因为那屈尊用抵押与我们联系自己的人，准备赐予我们很多。如果抵押如此，那所抵押的又该如何？\par

6. 但祂所说的\Quote{因为那时圣神还没有赐下，因为耶稣还没有受到光荣}是什么意思？祂说这话，显然是显而易见的意义。因为意思不是天主的圣神——祂曾与天主同在——不存在，而是还没有在那些信耶稣的人内。因此，主耶稣决定直到祂复活后才赐给我们所说的圣神；这并非没有原因。或许如果我们探寻，祂会恩赐我们找到；如果我们敲门，祂会为我们开门进去。虔诚敲门，而非手，尽管手若不停于仁慈之工，也在敲门。那么，为什么主耶稣基督决定直到祂受光荣后才赐下圣神？在我们尽可能谈论此事之前，必须先探究，以免有人因以下情况感到困扰：即圣神尚未在圣人们内，而我们在福音中读到，关于新生的主本人，西默盎借圣神认出了祂；寡妇亚纳，一位女先知，也认出了祂\BibleRef{路 2:25}；为祂施洗的若翰认出了祂\BibleRef{若 1:26}；匝加利亚充满圣神说了许多话；玛利亚本人领受圣神以怀孕主\BibleRef{路 1:35}。因此，在主因肉身复活而得光荣之前，我们有许多圣神之前的证据。并且先知们也有另一神，他们预言基督的来临。然而，这种赐予有某种方式，是之前从未出现的。因为在此之前，我们未读到人们聚集在一起，借领受圣神而说各国方言。但在祂复活后，祂首先向门徒显现，对他们说：\Quote{你们领受圣神吧。}关于这次赐予，经上说：\Quote{因为那时圣神还没有赐下，因为耶稣还没有受到光荣。}\Quote{祂向他们面庞嘘气}\BibleRef{若 20:22}，祂曾用气息给第一个人赋予生命，并使他从尘土中起来，以那气息赐给肢体灵魂；这表示祂是那向门徒面庞嘘气的同一位，好使他们从泥土中起来，并弃绝他们泥土般的工作。然后，在祂复活后——圣史称之为祂的光荣——主首先将圣神赐予祂的门徒。然后，如宗徒大事录所载，祂与他们同住了四十天，他们看见祂并与祂同行，祂在他们眼前升天。十天后，在五旬节那天，祂从天上遣发了圣神。他们曾聚集在一处，如我所说，当他们领受了圣神而被充满，便说起各国方言。\par

7. 那么，弟兄们，那在基督内受洗并信祂的人，如今并不说万国的方言，我们就不该相信他领受了圣神吗？愿天主阻止我们的心被这种不信所诱惑。我们确知人人都领受；但只按他带到泉源旁的信德器皿所能容纳的，祂就充满它。因此，既然圣神如今仍被人领受，有人会说：为什么没有一人说万国的方言呢？因为教会本身如今正说万国的方言。从前，教会在一个民族内，在那里它说万国的方言。因此，藉着说万国的方言，它预示了将要发生的事：即它将在各民族中成长，说万国的方言。凡不在这个教会内的，如今就不再领受圣神。因为，若从众肢体的合一中割断和分裂——这合一说万国的方言——那么让他自己判断吧；他并没有圣神。因为若他有，就让他给出那时所赐的标记。我们这样说是什么意思：让他给出那时所赐的标记？让他说万国的方言。他回答我说：那么，你又说万国的方言吗？显然我说；因为每一种语言都是我的，即属于我作为其肢体的身体。教会散居在万民中，说万国的方言；教会是基督的身体，在这个身体中你是一个肢体：因此，既然你是一个说万国方言的身体的肢体，相信你也在说万国的方言。因为肢体的合一是因着爱德而同心一致；那时这个合一就像一个说了话的人那样说话。\par

8. 因此，如果我们爱教会，如果我们因着爱德而彼此结合，如果我们以天主教的名号和信德为乐，我们也同样领受圣神。弟兄们，让我们相信：一个人爱基督的教会有多深，他就有多少圣神。因为圣神是被赐下的，如宗徒所说：\Quote{为显明}。为哪种显明？正如同一位宗徒说：\Quote{这人由圣神领受智慧之言，那人由同一圣神领受知识之言，另一人在同一圣神内领受信德，另一人在同一圣神内领受治病的奇恩，另一人能行奇迹}\BibleRef{格前 12:7}。因为有许多恩赐被赐下为显明，但是你，也许，没有我所说的一切。如果你爱，你所拥有的就不是\Emph{无}：如果你爱合一，凡在合一内拥有任何恩赐的，他也为你拥有。除去嫉妒，则我所有的也是你的。嫉妒的性情使人分离，健全的心智使人合一。在身体内，只有眼睛能看；但眼睛是为自己看吗？它既为手也为脚看，也为所有其他肢体看。如果有一击向脚袭来，眼睛不会避开，以免不加防备。同样，在身体内，只有手能工作，但手是为自己工作吗？它也为了眼睛工作：如果有一击袭来，不是打手，而是打脸，手会说\Quote{我不动，因为它不是冲我来的吗？}同样，脚藉着行走服务于所有肢体；所有其他肢体都沉默，舌头却为所有人说话。因此，如果我们爱教会，我们就拥有圣神；但我们爱教会，如果我们坚定地站在它的合一与爱德中。因为宗徒自己，在他说了不同的恩赐分赐给不同的人，正如几个肢体的职分之后，说：\Quote{我还要指示你们一条更卓越的路}；并开始谈论爱德。他把这放在人言和天使的话语之上，放在信德的奇迹之上，放在知识和预言之上，甚至放在一个人将他所有的一切分施给穷人的伟大慈善工作之上；最后，甚至放在身体的殉道之上：在所有这些伟大的事之上，他放上了爱德。拥有它，你便拥有一切：因为若没有它，你所能拥有的任何东西都毫无益处。但为使你知道我们所谈论的爱德是指圣神（因为福音中当前的问题是关于圣神的），请听宗徒说：\Quote{天主的爱德藉着所赐给我们的圣神，倾注在我们心中}\BibleRef{罗 5:5}。\par

9. 那么，既然圣神在我们内的恩惠是最伟大的，因为藉着祂，天主的爱倾注在我们心中，为什么主的旨意是要在祂复活后才将那圣神赐给我们呢？祂藉此预示什么？为的是在我们的复活中，我们的爱得以炽燃，并从对世俗的爱中脱离，完全奔向天主。因为我们在此出生、死亡：不要爱这世界；让我们藉着爱从此地迁移；藉着爱，让我们居住在至高之处，藉着那使我们爱天主的爱。在此生旅途中，我们不要思念别的，只思念我们不会永远在此，并且藉着善生，我们将在那里为自己预备一个永不再迁移的地方。因为我们的主耶稣基督，在祂复活后，\BibleQuote{不再死；}\BibleQuote{死，}宗徒说，\BibleQuote{再无权柄统治祂。}\BibleRef{罗 6:9} 看，我们应当爱什么。如果我们活着，如果我们相信那复活的主，祂将赐给我们，不是那些不爱天主的人在此所爱的，也不是那些越爱世俗就越不爱天主、或越不爱世俗就越爱天主的东西；而是让我们看看祂应许了我们什么。不是地上的、暂时的财富，不是今世的尊荣和权势；因为你们看见这些皆赐予恶人，使善人不至过分珍视。简而言之，也不是身体的健康本身，尽管那也是祂所赐，但如你们所见，祂甚至赐给走兽。也不是长寿；因为什么是长的呢？那终将结束的算长吗？祂并未应许信徒年岁长久，仿佛那是大事，或应许衰老的晚年——那是人人都盼望着到来，但到来后又人人抱怨的。也不是容貌的美丽，那会被疾病或那被渴望的老年驱走。人希望美丽，又希望活到老年：这两个愿望不能共存；你若老了，便不美丽；当老年到来，美丽逃逸；美貌的活力与老年的呻吟不能同住一躯。因此，这一切都不是祂所说的：\BibleQuote{信我的，让他来喝，从他的腹中要流出活水的江河。} 祂应许了我们永生，在那里我们无所畏惧，不再烦恼，不再迁移，不再死亡；那里既无对已故先人的哀悼，也无对后继者的期待。因此，因为祂向我们这些爱祂的人、并因圣神的爱德而炽热的人所应许的正是如此，所以祂直到受光荣后才赐予我们那同一位圣神，为要在祂的身体中显示那我们现在没有、但在复活中盼望的生命。\par

\SourceRecord{Translated by John Gibb.；Nicene and Post-Nicene Fathers, First Series；Vol. 7.；Edited by Philip Schaff.；Buffalo, NY: Christian Literature Publishing Co.,；1888.；<http://www.newadvent.org/fathers/1701032.htm>.}

% structure: p0001 source=h1 target=section reason=page-title

\section{第33篇 (\BibleRef{若 7:40-8:11})}

1. 亲爱的，你们记得，在上次讲道中，我们曾因当日所读福音的片段，对你们讲论了有关圣神的事。当主邀请那些信祂的人来饮这水的时候，祂是在那些图谋捉拿祂、想要杀祂却不能的人中间说话，因为那并非祂的意愿；好吧，祂说了这些话以后，群众中便起了纷争：有些人认为祂就是基督，另有些人说基督不会从加里肋亚出来。但那些被派去捉拿祂的人回来了，他们清白无罪，且充满钦佩。因为他们甚至为祂神圣的教导作证，当时派他们去的人问：\Quote{你们为什么没有带祂来？}他们回答说，从未听见过人这样说话：\Quote{因为没有任何人这样说话。}但祂这样说话，因为祂既是天主又是人。但法利塞人驳斥他们的见证，对他们说：\Quote{你们也被迷惑了吗？}我们确实看见，你们也被祂的言论迷住了。\Quote{难道有首领或法利塞人相信了祂吗？但这群不懂法律的群众是该受诅咒的。}那些不认识法律的人，相信了那颁布法律的主；而那些教导法律的人却轻视祂，为应验主自己曾说过的话：\Quote{我来是为使看不见的能看见，使看见的反而成为瞎子。}\BibleRef{若 9:39}因为法利塞人，法律的教师，变成了瞎子；而那不认识法律却信了法律之作者的人民，却被光照了。\par

2. 然而，\Quote{尼苛德摩}，\Quote{那个法利塞人，曾在夜间来到主面前的人}——他本人并非不信，而是胆怯；因此他在夜间来到光明之中，因为他希望被光照，却又害怕被人认出——我说，尼苛德摩回答犹太人说：\Quote{难道我们的法律不经审讯，不查明他所做的，就定人的罪吗？}因为他们偏执地希望在审问之前就定祂的罪。尼苛德摩确实知道，或者更确切地说，相信，只要他们愿意耐心听祂，或许会变得像那些被派去捉拿祂的人一样，但那些人却宁愿相信。他们出于内心的偏见，正像他们对那些差役所说的那样回答道：\Quote{你也是加里肋亚人吗？}意思是，被那个加里肋亚人迷惑的人。因为主被称为加里肋亚人，因为祂的父母来自纳匝肋城。我说\Quote{祂的父母}是针对玛利亚说的，而不是关于男人的种子；因为祂在地上只寻求一位母亲，祂在天上早已有一位父亲。因为祂的诞生在这两方面都是奇妙的：神性一面没有母亲，人性一面没有父亲。那么，这些自诩为法律专家的人对尼苛德摩说了什么？\Quote{你去查考圣经，看看从加里肋亚会不会出先知。}然而，先知的主却从那里兴起。福音作者说：\Quote{他们便各自回家去了。}\par

3. \Quote{于是耶稣上了那座山；}即\Quote{橄榄山}——那果实丰盛的山，那膏油之山，那傅油之山。因为除了橄榄山，基督在何处教导更合适呢？因为\Emph{基督}的称号源于\Emph{傅油}；希腊文\OriginalTerm{傅油}{\GreekText{χρισμα}}，拉丁文称为\OriginalTerm{涂油}{unctio}，即\Emph{傅油}。祂曾傅油于我们，正是为了这个缘故，因为祂使我们成为对抗魔鬼的战士。\Quote{清晨，祂又回到圣殿，众百姓都到祂那里来，祂便坐下教训他们。}而祂没有被捉拿，因为祂还尚未甘愿受苦。\par

4. 现在观察主如何被敌人试探祂的温良。\Quote{经师和法利塞人给祂带来一个行淫时被抓的女人，把她放在中间，对祂说：师傅，这女人是行淫时被抓的。法律上梅瑟命令我们，这样的人该用石头砸死；但祢说什么呢？他们说这话，试探祂，为要控告祂。}为什么控告祂？难道他们抓住了祂什么错行？还是那女人与祂有任何关系？那么，\Quote{试探祂，为要控告祂}是什么意思？弟兄们，我们明白，主身上闪耀着一种非凡的温良。他们注意到祂非常温良、非常柔和：因为关于祂早有预言：\BibleQuote{大能者啊，愿祢腰间佩刀，带着荣耀和威严，为真理、温和、正义而策马向前，得胜为王。}因此，作为导师，祂带来了真理；作为解救者，祂带来了温良；作为护卫者，祂带来了正义。先知藉圣神曾预言祂要因这些事而为王。当祂说话时，祂的真理被承认；当祂没有对敌人发怒时，祂的温和受到赞美。因此，关于这两点——即祂的真理和温和——敌人被恶毒和嫉妒所折磨；关于第三点——即正义——他们为祂设下了陷阱。如何设下的呢？因为法律命令奸夫淫妇要用石头打死，而法律当然不会命令不义的事：若有人说与法律命令不同的话，就会被发现是不义的。于是他们彼此说：\Quote{祂被称为真理，祂显得温和；必须针对正义找祂的把柄。我们把一个行淫时被抓的女人带到祂面前，问祂法律对这类人有什么规定：如果祂赞成用石头打死她，祂就不会显出温和；如果祂同意释放她，祂就不会持守正义。但是，他们说，这样祂为了不失去在民众中因温和而获得的爱戴，必定会说必须释放她。这样我们就有了控告祂的机会，就控告祂违反法律：对祂说，你是法律的敌人，你反对梅瑟，不，反对藉梅瑟赐下法律的那一位；你该死，你也必须和这女人一同被石头打死。}用这些话和想法，他们或许能煽动对祂的嫉妒，催促控告，导致人们急切要求定祂的罪。但这是针对谁呢？这是乖谬对抗正直，谎言对抗真理，败坏的心对抗正直的心，愚昧对抗智慧。这样的人何时铺设陷阱而不先把自己的头伸进去呢？看，主在回答他们时，既持守了正义，又没有偏离温和。那被设陷阱捕捉的没有被捉住，反而是设陷阱的人被捉住了，因为他们不信赖那能将他们从网罗中拉出来的主。\par

5. 那么，主耶稣作了什么回答？真理如何回答？智慧如何回答？那正义——针对它的虚假控告已经准备好——如何回答？祂没有说：不要用石头砸她，免得祂似乎反对法律。但绝不能说：用石头砸她，因为祂来不是要失去祂所找到的，而是要寻找失去的。那么祂如何回答呢？你看这回答充满正义、充满温和和真理！\BibleQuote{你们中间谁没有罪，就先拿石头砸她吧。}祂说。啊，智慧的回答！祂是如何将他们的视线转向他们自己的！因为他们站在外面控告、指责，却不审察内心；他们看见了淫妇，却没有看见自己。违犯法律的人，却希望法律被执行，而这并非出于真心去遵行，仿佛通过贞洁来谴责奸淫。你们听见了，犹太人；你们听见了，法利塞人；你们听见了，律法教师——法律的守护者，但尚未理解祂就是立法者。当祂用手指在地上写字时，祂向你们指明什么呢？因为法律是用天主的手指写的，但因人心的顽固而写在石版上。如今主写在地上，因为祂在寻找果实。那么你们已经听见了：让法律成全，让淫妇被砸死。但是，难道法律要由那些该受惩罚的人通过惩罚她来成全吗？你们每个人要省察自己，进入自己的内心，登上自己良心的审判台，强迫自己认罪。因为他知道自己是什么：因为\BibleQuote{除了人内里的心神，没有人知道人的事。}每个人仔细省察自己，就会发现自己是罪人。是的，确实如此。所以，要么释放这女人，要么你们和她一同接受法律的惩罚。如果祂说：不要砸死淫妇，祂会被证明是不义的；如果祂说：砸死她，祂就不会显得温和；那么让祂说那适合祂——既温和又公义——的话吧：\BibleQuote{你们中间谁没有罪，就先拿石头砸她吧。}这是正义的声音：让那罪人受罚，但不是由罪人；让法律成全，但不是由违法者。这确实是正义的声音；那些人如同被箭刺穿，省察自己，发现自己是罪人，\BibleQuote{就从老到少一个一个地都离开了。}只剩下两个人，那可怜的女人和慈悲。但是主用那正义之箭刺透他们之后，不愿理会他们的跌倒，反而转移目光不看他们，\BibleQuote{又用指头在地上写字。}\par

6. 但那妇人独自留下，众人都出去后，祂举目向妇人观看。我们已听见正义的声音，也让我们听见仁慈的声音。因为我想，当那妇人听见主说：\Quote{你们中间谁没有罪的，就先向她投石吧}时，她是更加惊恐的。但他们转而想到自己，因着那退去而承认了自己有罪，便将那妇人连同她的大罪留给了那没有罪的祂。因她听见了这话：\Quote{谁没有罪的，就先向她投石吧}，她指望要受那在她身上找不到罪的祂的惩罚。但祂，以正义的舌头驱退了她的对头，又向祂举起仁慈的眼目，问她说：\Quote{没有人定你的罪吗？}她回答说：\Quote{主，没有人。}祂说：\Quote{我也不定你的罪；}你或许怕被那一位定罪，因为在我身上你找不着罪。\Quote{我也不定你的罪。}这是什么意思，主啊？难道祢因此袒护罪吗？显然不是。注意下文：\Quote{去吧，今后不要再犯罪了。}因此，主也定了罪，但定的是罪，不是人。因为祂若是罪的庇护者，祂就会说：\Quote{我也不定你的罪；去吧，随你怎么活着；在我的救援中安心吧；无论你犯多少罪，我都将从一切地狱的惩罚和阴间的刑吏中救出你。}祂没有这样说。\par

7. 那么，那些在主内喜爱祂的宽仁的人，要留心，也要畏惧祂的真理。因为\Quote{上主是甘饴的、正直的。}你爱祂，因为祂甘饴；畏惧祂，因为祂正直。正如温良者说：\Quote{我缄默不语}；但作为公义者，祂说：\Quote{我岂能永远沉默？}\BibleRef{依 42:14}\Quote{上主是仁慈的、怜悯的。}祂确实如此。再加上\Quote{含忍}；再加上\Quote{极其怜悯}；但要畏惧最后一句：\Quote{且是真实的。}因为祂现今宽容罪人，将来却要审判藐视祂的人。\Quote{难道你藐视祂丰厚的恩慈、宽容与忍耐，不晓得天主的恩慈是领你悔改吗？你竟凭着你刚硬不悔改的心，为自己积蓄忿怒，以致忿怒的日子，就是天主公义审判的日子来到；祂要照各人的行为报应各人。}\BibleRef{罗 2:4}主是温良的，主是含忍的，主是怜悯的；但主也是公义的，主也是真实的。祂给你悔改的时间；但你却喜爱迟来的审判胜过改正你的行为。你昨天是个恶人？今天就做个好人。你今日仍在作恶？无论如何，明天要改变。你总在等待，并从天主的仁慈向自己许下极大的诺言。好像那应许你悔改就得赦免的天主，也应许你更长的寿命。你怎么知道明天会发生什么？你心里说：等我改正了行为，天主必赦免我所有的罪。我们不能否认天主确实应许了赦免那些悔改和归正的人。因为你在哪位先知书中读到天主应许赦免悔改者，你却读不到天主应许你长寿。\par

8. 因此，人处于两种危险之中：既因希望，也因绝望；来自相反的事，来自相反的情感。谁因希望而受骗？那说天主是善的、天主是仁慈的，让我随心所欲、随意而行，放纵我的情欲，满足我灵魂的欲望的人。为什么这样？因为天主是仁慈的，天主是良善的，天主是宽厚的。这些人因希望而处于危险。而那些绝望的人则处于危险，他们陷入重罪，以为悔改后不能再得赦免，相信自己无疑要受罚，便对自己说：我们注定要受罚，为何不随心所欲，像注定赴死的角斗士那样行事？这就是绝望之人危险的原因：因为他们不再有所畏惧，反而极其可怕。绝望杀死这些，希望杀死那些。人的心思在希望与绝望之间摇摆。你必须畏惧：免得希望杀死你；当你过多寄望于仁慈时，免得落入审判；同样，你要畏惧：免得绝望杀死你；当你以为你所犯的重罪不能被赦免时，便不悔改，以致招来智慧的话语：\Quote{我必在你们遭难时发笑。}\BibleRef{箴 1:26}那么，主是如何对待那些因这两种病而处于危险中的人呢？对于因希望而危险的人，祂说：\Quote{归向上主，不可迟缓，不要一天一天地拖延；因为祂的忿怒忽然降来，在报复的时日将你完全消灭。}\BibleRef{德 5:8}对于因绝望而危险的人，祂说什么？\Quote{恶人若悔改，我必不再记念他的一切过犯。}\BibleRef{则 18:21}因此，为了那些因绝望而危险的人，祂给我们提供了赦免的庇护所；并为了那些因希望而危险、被拖延所迷惑的人，祂使死亡的日子不确定。你不知道你的最后一天何时到来。你有今日可以改过，岂不感恩？于是祂对那妇人说：\Quote{我也不定你的罪}；但既然对过去可以安心，且要谨慎将来。\Quote{我也不定你的罪：}我涂抹了你所做的；你要遵守我所吩咐的，好叫你得到我所应许的。\par

\SourceRecord{Translated by John Gibb.；Nicene and Post-Nicene Fathers, First Series；Vol. 7.；Edited by Philip Schaff.；Buffalo, NY: Christian Literature Publishing Co.,；1888.；<http://www.newadvent.org/fathers/1701033.htm>.}

% structure: p0001 source=h1 target=section reason=page-title

\section{讲道辞 34（\BibleRef{若 8:12}）}

刚才我们在恭听神圣福音时，所听到并留心领受的，我相信诸位也都曾努力去理解，而且我们各人按自己的度量，对所宣读的那如此重大之事，各自领会了可能领会的一部分；当圣言的面饼摆出时，没有人能抱怨自己未尝到任何东西。但我也相信，几乎没有人完全理解了全部。然而，即使有谁能充分理解我们主耶稣基督如今从福音中宣读出来的话，也请容忍我们的服务，同时，若有可能，在祂的助佑下，通过我们的讲解，使那少数人因自己理解而喜悦的事，要么全部人，要么多数人也能理解。\par

我认为，主说的\Quote{我是世界的光}，对于拥有眼睛、能分享这光的人而言是清楚的；但那些只在肉身中有眼睛的人，对主耶稣基督所说\Quote{我是世界的光}感到惊奇。或许也有人心中暗自思忖：主基督会不会就是那藉升落而产生昼夜的太阳？因为曾有异端者如此认为。\OriginalTerm{摩尼教徒}{Manichæans}曾猜想主基督就是那肉眼可见、不仅对人、甚至对野兽也都公开显现的太阳。但天主教会的正统信仰拒绝这种虚构，并看出那是魔鬼的教义：它不仅藉信德承认如此，而且在可能的情况下，甚至通过推理来证明。因此，让我们拒绝这类错误，这错误是圣教会从一开始就绝罚的。我们不要以为主耶稣基督就是我们从东方升起、西方落下、随后黑夜来临、其光线被云遮蔽、按固定运动移位的那个太阳：主基督并非这样的事物。主基督不是受造的太阳，而是造太阳的那位。因为\BibleQuote{万物是藉着祂造的；凡受造的，没有一样不是藉着祂造的。}\par

因此，有一个光创造了这个太阳的光：让我们爱这光，让我们渴望理解它，让我们渴慕它；以便以它本身作我们的向导，我们终于能到达它，并使我们生活在它里面，从而永不死亡。这实在就是那光，先祖的预言早已在圣咏中这样歌唱：\BibleQuote{主啊，祢拯救人和兽；天主，祢的仁慈何其丰盛。}这些是圣咏的话：你们要留意古代圣人们关于这光所预先讲述的什么。它说：\BibleQuote{人，}和\BibleQuote{野兽，祢拯救，主啊；天主，祢的仁慈何其丰盛。}因为既然祢是天主，祢有丰盛的仁慈，祢那同样的丰盛仁慈，不仅达到祢按自己肖像所造的人，甚至达到祢给人权下所用的野兽。因为赐予人救恩的，同样也赐予野兽救恩。不要羞耻于这样想你的主天主：不，要相信这事并信赖它，且看你不要有别的想法。那拯救你的，也拯救你的马、你的羊；甚至到最小的，也拯救你的鸡：\BibleQuote{救恩属于主}，天主拯救这一切。你感到不安，你质问，我纳闷你为何怀疑。难道祂不屑于拯救，那曾屈尊创造的一位？主拯救天使、人和野兽：\BibleQuote{救恩属于主}。正如没有人靠自己存在，同样没有人靠自己得救。因此，圣咏说得最真实、最恰当：\BibleQuote{主啊，祢拯救人和兽。}为什么？\BibleQuote{天主，祢的仁慈何其丰盛。}因为祢是天主，祢创造了，祢拯救：祢赐予存在，祢赐予健康的存在。\par

4. 因此，既然天主的仁慈倍增，人畜都因祂而得救，难道人没有天主作为造物主赐予他们、却不赐予牲畜的别的东西吗？按照天主的肖像受造的活物与受造服从天主的肖像的活物之间，难道没有区别吗？显然有区别：除了我们与哑巴牲畜共有的那种救恩之外，还有天主赐予我们、却不赐予它们的东西。这是什么？继续同一圣咏：\Quote{至于人子，他们将在祢翅膀的荫庇下寄望。}现在与他们的牲畜共有救恩，\Quote{人子将在祢翅膀的荫庇下寄望。}他们有一种实际的救恩，另一种是希望的救恩。目前的救恩是人与牲畜共有的；但还有另一种是人寄望的；寄望的人领受，绝望的人不领受。因为经上说：\Quote{人子将在祢翅膀的荫庇下寄望。}那些坚持寄望的人被祢保护，免得他们被魔鬼从希望中推倒：\Quote{在祢翅膀的荫庇下，他们将寄望。}如果他们寄望，他们寄望什么，岂不就是牲畜所没有的吗？\Quote{他们将饱享祢殿中的肥脂，并饮祢欢乐的急流。}那是什么样的酒，令人醉而可赞？那是什么样的酒，不扰乱心智，反而引导心智？那是什么样的酒，使人永葆清醒，饮酒而不致疯狂？\Quote{他们将饱享。}如何？\Quote{祢殿中的肥脂，并祢欢乐的急流，祢将赐他们饮。}何以如此？\Quote{因为祢那里有生命的泉源。}生命的泉源亲自行走在地上，祂说：\Quote{凡口渴的，到我这里来。}看，这泉源！但我们已经开始谈论光，并处理福音中关于光的命题。因为我们读到主说：\Quote{我是世界的光。}由此引发了一个问题，免得有人按肉体理解，以为这光是指太阳：我们来到圣咏，思想之后，暂时发现主是生命的泉源。饮吧，并生活。经上说：\Quote{祢那里有生命的泉源；}因此，\Quote{在祢翅膀的荫庇下，人子寄望，}寻求饱饮这泉源。但我们是在谈论光。那么继续吧；因为先知说了\Quote{祢那里有生命的泉源}后，又接着说：\Quote{在祢的光明中，我们将得见光明——}天主出自天主，光明出自光明。这光明创造了太阳的光芒；而创造太阳的光明，在祂之下祂也创造了我们，为了我们而被造成在太阳之下。那创造太阳的光明，我说，为了我们而被造成在太阳之下。不要轻视肉身的云彩；它被这云彩遮盖，并非为了暗淡，而是为了缓和。\newline{}\par

5. 那永不衰竭的光、智慧的光，藉着肉身的云彩向人说：\Quote{我是世界的光；跟随我的，决不在黑暗中行走，必有生命的光。}祂何等将你们从肉体的眼睛召出，召回心灵的眼睛！因为仅仅说\Quote{跟随我的，决不在黑暗中行走，必有光}还不够；祂还加上了\Quote{生命}，正如那里所说的：\Quote{因为祢那里有生命的泉源。}看，我的弟兄们，主的话如何与圣咏的真理相符：在那里，光与生命的泉源并列；而主说\Quote{生命的光}。但在肉体用途上，\Emph{光}和\Emph{泉源}是不同的东西：我们的口寻求泉源，眼睛寻求光；当我们口渴时，我们寻找泉源；当我们在黑暗中时，我们寻找光；如果我们碰巧在夜间口渴，我们会点灯来到泉源。天主却不是这样：\Emph{光}和\Emph{泉源}是同一回事：祂为你照耀使你看见，祂为你流淌令你饮用。\newline{}\par

6. 那么，我的弟兄们，你们看见，如果你们内在地看见，这是什么样的光，主说：\Quote{跟随我的，决不在黑暗中行走。}跟随太阳，让我们看看你是否不在黑暗中行走。看，它升起向你而来；它沿着轨道向西运行。也许你的旅程是向东：除非你朝着与它行进相反的方向走，否则你必然因跟随它而迷路，不去东方反到西方。如果你在陆地上跟随它，你会走错；如果航海者跟随它，他也会走错。最后，你以为必须跟随太阳，而你自己也向西旅行，也朝着它行的方向；让我们看看它在落山之后，你是否不在黑暗中行走。看，尽管你不愿舍弃它，它却会舍弃你，因它服务的必需而结束白昼。但我们的主耶稣基督，即使当祂并未藉着肉身的云彩向众人显现时，却仍以祂智慧的能力统御万有。你的天主处处全然临在：你若不离弃祂，祂绝不会离弃你。\newline{}\par

7. 因此，祂说：\Quote{跟随我的，}祂说，\Quote{决不在黑暗中行走，却要有生命的光。}祂所应许的，以未来时态的词表达；因为祂不说\Emph{有}，而是说\Quote{却要有生命的光。}然而祂不是说\Quote{\Emph{将要}跟随我的}，而是\Quote{\Emph{实在}跟随我的}。祂将我们应当做的，置于现在时态；但祂应许给实行者的事，则用未来时态的词指明。\Quote{跟随的，必将拥有。}现在跟随，日后拥有：现在因信德跟随，日后因看见而拥有。因为，宗徒说，\BibleQuote{我们住在肉身内，是离主远行的：因为我们凭信德行走，不是凭看见。}\BibleRef{格后 5:6}何时我们将凭看见行走？当我们拥有生命的光时，当我们达到那神视时，当这黑夜过去时。关于那将要升起的日子，经上记载：\BibleQuote{早晨我将侍立在祢面前，并瞻仰祢。}\Quote{早晨}是什么意思？当这世界的黑夜过去，当诱惑的恐惧过去，当那在黑夜中巡行咆哮、寻找可吞噬者的狮子被击败时。\BibleQuote{早晨我将侍立在祢面前，并瞻仰。}如今，弟兄们，我们认为当前的本分是什么？岂不就是那在圣咏中所说的：\BibleQuote{每夜我要洗涤我的床铺，用泪水浸透我的卧榻}？每夜，他说，我要哭泣；我渴慕光明而焦灼。主看见我的渴望：因为另一篇圣咏向祂说：\BibleQuote{我的一切渴望都在祢面前；我的呻吟祢不隐瞒。}你渴望金子？你能被看见；因为寻求金子时，你将显于人前。你渴望粮食？你向拥有它的人求问；你在寻求得到你所渴望之物时，也告知那人。你渴望天主？谁能看见，除了天主？那么，你从谁那里寻求天主，如同你寻求饼、水、金、银、粮食一样？你从哪里寻求天主，除非从天主？祂应许了祂自己，当从祂自己那里寻求。让灵魂伸展她的渴望，以更宽广的胸怀寻求理解那\BibleQuote{眼所未见，耳所未闻，人心未曾想到过的。}\BibleRef{格前 2:9}我们能够渴望它，能够向往它，能够切慕它；但却不能相称地构思它，相称地用言语表达它。\par

8. 因此，我的弟兄们，既然主简短地说：\BibleQuote{我是世界的光：跟随我的，决不在黑暗中行走，却要有生命的光；}在这话中，祂命令了一事，应许了另一事；让我们实行祂所命令的，好让我们不致厚颜无耻地索取祂所应许的；免得祂在审判中对我们说：难道你们实行了我所命令的，就期望我所应许的吗？那么，上主、我们的天主，祢命令了什么？祂对你说：你应跟随我。你曾寻求生命的忠告？什么生命？岂不就是那论及的生命：\BibleQuote{在祢那里有生命的泉源}？有一个人听见对他说：\BibleQuote{你去，变卖你所有的一切，施舍给穷人，你必有宝藏在天上；然后来，跟随我。}他没有跟随，却忧愁地走了；他寻求\Quote{善师}，到祂那里如同到老师那里，却藐视祂的教导；他忧愁地走了，被他的私欲捆绑束缚；他忧愁地走了，肩上有贪婪的重担。他劳苦焦虑；然而他却认为那愿意卸下他重担的祂，不是应当跟随而是应当离弃的。但是自从主藉着福音大声呼喊：\BibleQuote{凡劳苦和负重担的，你们都到我这里来，我要使你们安息；你们背起我的轭，向我学习，因为我是良善心谦的，}有多少人，一听到福音，就做了那富人从主自己口中听到却未做的事？因此，让我们现在去做，让我们跟随主；让我们解开那阻碍我们跟随祂的锁链。谁又能解开这样的束缚呢？除非祂帮助，就是那位祂曾向他说的：\BibleQuote{祢已挣断我的锁链。}关于祂，另一篇圣咏说：\BibleQuote{主释放被囚的；主扶起被压伤和受欺压的。}\par

9. 那些被释放、被举起来的人，所追随的不正是他们所听见的光吗？\BibleQuote{我是世界的光；跟随我的，决不在黑暗中行走。}因为主光照瞎子。所以，弟兄们，我们拥有信德的眼药，如今已被光照。因为祂的唾沫先前与泥土混合，用来涂抹那生来瞎眼者的眼睛。我们也是因亚当而天生盲目，需要祂来光照我们。祂把唾沫与泥土混合：\BibleQuote{圣言成了血肉，寄居在我们中间。}祂把唾沫与泥土混合，因此早有预言：\BibleQuote{真理从地生出。}祂自己说：\BibleQuote{我是道路、真理、生命。}当我们面对面看见时，我们将圆满地享有真理；因为这也曾应许给我们。因为谁敢盼望天主未曾屈尊许诺或给予的事呢？我们要面对面看见。宗徒说：\BibleQuote{我现在所认识的，只是局部的；现在藉着镜子观看，模糊不清；但那时，就要面对面了。}\BibleRef{格前 13:12}宗徒若望在他的书信中说：\BibleQuote{可爱的诸位，现在我们是天主的子女，但我们将来如何，还没有显明；可是我们知道：一显明，我们必要相似祂，因为我们要看见祂实在怎样。}\BibleRef{若一 3:2}这是一个伟大的应许；如果你爱，就跟随。你说：\Quote{我确实爱，但我该走哪条路呢？}如果主你的天主曾对你说：\BibleQuote{我是真理和生命，}你渴望真理、向往生命，你可能会真正地询问通往这些的路，并在心里说：真理是伟大的，生命是伟大的，只要我的灵魂有方法到达！你问方法？最初听祂说：\BibleQuote{我就是道路。}在祂说\InnerQuote{往哪里}之前，祂先声明了方法：\BibleQuote{我是，}祂说，\BibleQuote{道路。}道路往哪里？\BibleQuote{并且是真理和生命。}首先，祂告诉了你来的方法；然后，往哪里去。我是道路，我是真理，我是生命。祂与圣父同在时，是真理和生命；取了肉躯，祂成为道路。没有人对你说：你必须努力寻找通往真理和生命的道路；没有人对你说。懒惰者，起来吧：道路本身已经来到你面前，将你从睡眠中唤醒；如果祂唤醒了你，起来行走。或许你试图行走，却因脚疼而不能。你的脚为何疼？是不是因为贪欲的命令而在崎岖的地方奔跑？但天主的话语也医治了瘸子。看，你说，我的脚健全，但道路本身我看不见。祂也开启了瞎子的眼睛。\par

10. 这一切都凭信德，只要我们远离主，居住在肉身内；但当我们走完了道路，抵达了家乡本身，还有什么比我们更喜乐呢？还有什么比我们更有福呢？因为没有什么比我们更平安；因为再没有人背叛。但现在，弟兄们，我们很难没有争斗。我们确实蒙召走向和睦，被命令彼此和平相处；我们必须为此努力，竭尽全力，好使我们最终达到最完美的和平；但目前我们常在争斗中，常常与那些我们为其谋求益处的人争斗。有人走入歧途，你希望领他走上正路；他抗拒，你与他争斗：外教人抗拒你，你与偶像和魔鬼的错误争辩；异端者抗拒你，你与其他魔鬼的教义争辩；不良的天主教徒不愿正当地生活，你甚至责备你内部的弟兄；他与你同住一屋，却寻求毁灭之路；你热切地渴望纠正他，好使你能向主交出一份关于你和他双方的好账。四面八方的争斗何其多！人常常被疲惫压倒，对自己说：我何必忍受与反对者相处、忍受那些以恶报善的人？我愿意使他们受益，他们却甘愿丧亡；我在争斗中耗尽生命；我没有平安；此外，我把那些本应成为朋友的人变成了敌人，如果他们顾及寻求他们益处者的善意：我何必要忍受这些？让我回到自己里面，我要独处，我要呼求我的天主。你回到自己里面，你在那里也会发现争斗。如果你开始跟随天主，你在那里也会发现争斗。你说，我发现什么争斗？\BibleQuote{肉身与圣神相争，圣神与肉身相争。}\BibleRef{迦 5:17}看，你就是你自己，你独自一人，你与自己同在；看，你没有与任何人相处，但你却看到在你肢体内另有一条法律，与你理智的法律交战，并把你掳去，叫你隶属于那在你肢体内罪的法律。那么，你要大声呼喊，向天主呼喊，求祂赐你从内在争斗中获得平安：\BibleQuote{我这个人真不幸呀！谁能救我脱离这死亡的身体呢？天主的恩宠，藉着我们的主耶稣基督。}\BibleRef{罗 7:23}因为，祂说：\BibleQuote{跟随我的，}祂说，\BibleQuote{决不在黑暗中行走，必有生命的光。}一切争斗结束后，不朽将随之而来；因为最后一个仇敌便是死亡。而这将是怎样的平安？\BibleQuote{这可朽坏的，必须穿上不可朽坏的；这可死的，必须穿上不可死的。}\BibleRef{格前 15:26}为了能到达那里（因为那时将是现实），让我们现在怀着希望跟随祂——祂曾说：\BibleQuote{我是世界的光；跟随我的，决不在黑暗中行走，必有生命的光。}\par

\SourceRecord{Translated by John Gibb.；Nicene and Post-Nicene Fathers, First Series；Vol. 7.；Edited by Philip Schaff.；Buffalo, NY: Christian Literature Publishing Co.,；1888.；<http://www.newadvent.org/fathers/1701034.htm>.}

% structure: p0001 source=h1 target=section reason=page-title

\section{论道35（\BibleBook{若望}{福音}8:13-14）}

1. 你们昨天在场的人，请记住我们曾长时间讨论了我们的主耶稣基督的话，祂说：\BibleQuote{我是世界的光；跟随我的，决不在黑暗中行走，必有生命的光。} 如果我们愿意继续讨论那光，我们还可以讲许久；因为我们不可能简短地解释这事。因此，我的弟兄们，让我们跟随基督，世界的光，免得我们在黑暗中行走。我们必须害怕黑暗——不是眼睛的黑暗，而是品行的黑暗；即使它是眼睛的黑暗，也不是外在眼睛的黑暗，而是内在眼睛的黑暗，就是那些我们用来辨别，不是黑与白，而是是与非的眼睛。\par

2. 当我们的主耶稣基督说了这些话，犹太人回答说：\BibleQuote{你为自己作证；你的证据不真实。} 在我们的主耶稣基督来临之前，祂点燃并派遣了许多先知性的灯在祂前面。其中也有若翰洗者，那大光本身，即主基督，给予了他从未给予任何人的见证；因为祂说：\BibleQuote{在妇女所生者中，没有兴起一位比若翰洗者更大的。}\BibleRef{玛 11:11} 然而，这位在妇女所生者中没有谁比他更大的人，论到主耶稣基督却说：\BibleQuote{我固然用水洗你们，但有一位要来的比我有能力，我就是给祂提鞋也不配。}\BibleRef{若 1:26} 你看这灯如何顺服那白日。主亲自作证说，这位若翰确实是一盏灯：\BibleQuote{他是}，祂说，\BibleQuote{一盏燃烧并照耀的灯，你们情愿暂时在他的光中欢乐。}\BibleRef{若 5:35} 但当犹太人对主说：\BibleQuote{告诉我们，你凭什么权柄做这些事？} 祂知道他们视若翰洗者为大人物，而这他们所视为大人物的人曾对他们为主作证，就回答他们说：\BibleQuote{我也要问你们一句话，告诉我，若翰的洗礼是从哪里来的？是从天上，还是从人间？} 他们困惑了，彼此商议：如果说\BibleQuote{从人间}，恐怕会被百姓用石头砸死，因为百姓相信若翰是先知；如果说\BibleQuote{从天上}，祂会对他们说：\Quote{你们承认他是从天上来的先知，他为我作了见证，你们已从他口中听到我凭什么权柄做这些事。} 他们看到，无论他们选择这两个回答中的哪一个，都会陷入圈套，就说：\BibleQuote{我们不知道。} 主回答他们说：\BibleQuote{我也不告诉你们，我凭什么权柄做这些事。}\BibleRef{玛 21:23} \Quote{我不告诉你们我所知道的，因为你们不愿承认你们所知道的。} 他们当然被最公正地拒绝了，狼狈地离去；于是实现了天主圣父借先知在圣咏中所说的：\BibleQuote{我为我的基督预备了一盏灯}（这灯就是若翰）；\BibleQuote{我要以羞耻给祂的仇敌披上。}\par

3. 所以，主耶稣基督有先祂而来的先知们的见证，也有那在审判者之前的宣告者的见证；祂有若翰的见证；但祂自己为自己作的见证是更大的。然而，那些眼力微弱的人寻求灯，因为他们不能承受那白日；因为那同一位若望宗徒，我们手中的福音就是他的，在他的福音开头论到若翰洗者说：\BibleQuote{曾有一人，由天主派遣，名叫若翰。他来是为作证，为给光作证，为使众人借他而信。他不是那光，只是来给光作证。那普照每人的真光，正在进入世界。} 如果\Quote{每人}，那么也光照若翰。因此，同一位若翰说：\BibleQuote{我们都从祂的圆满中领受了。} 所以你们要辨别这些事，使你们的心智在基督的信德中获益，使你们不再是常求乳吮、回避固体食物的婴孩。你们应当由我们圣洁的母亲基督的教会来滋养和断乳，并借着心智、而非肚腹来到更坚固的食物。那么，你们要辨别这一点：那光照的光是一回事，那被光照的是另一回事。因为我们的眼睛也被称为光；人人都这样指着他的眼睛起誓说：\Quote{愿我的光活着。} 这是一种惯常的起誓。如果这些光真是光，就让他们睁开，并在你们关闭的房间里为你们照亮，即使那里没有光；他们当然不能。因此，正如我们脸上的这些称为光的眼睛，即使又健康又睁开，也需要外在光的帮助——如果那光被移开或不带入，虽然它们健全且睁开，却看不见——同样，我们的心智，即灵魂的眼睛，除非被真理之光照耀，并被那光照人而不被光照的主奇妙地照亮，否则不能达到智慧，也不能达到正义。因为正义地生活对我们来说就是道路本身。但光不照耀的人，怎能不在路上绊倒呢？因此，在这样的路上，我们需要看见；在这样的路上，看见是伟大的事。如今多俾亚脸上的眼睛闭上了，儿子给父亲引路；然而父亲却用他的教导给儿子指出了道路。\BibleRef{多 2:11}\par

4. 犹太人便回答说：\BibleQuote{你为你自己作证；你的见证不真实。} 让我们看看他们听到了什么；让我们也听，却不似他们那样：他们轻视，我们相信；他们想要杀害基督，我们渴望藉基督而生活。愿这区别使我们耳与心与他们有别，并让我们听主如何回答犹太人。\BibleQuote{耶稣回答他们说：我即使为自己作证，我的见证也是真实的，因为我知道我从哪里来，往哪里去。} 光明既显示其他事物，也显示自身。譬如，你点一盏灯找你的外衣，燃着的灯为你提供亮光以找到外衣；你点灯是为了看它燃烧时的自身吗？事实上，一盏燃着的灯既能同时照亮黑暗遮盖的其他事物，也能向你的眼睛展示自身。同样，主基督将祂忠信的人与犹太仇敌分开，如同光与暗分开：将那些祂以信德之光照亮的人，与那些祂将光倾注在其紧闭的眼睛上的人分开。太阳同样照耀着有视力之人的脸和盲人的脸；两人都站着面向太阳，在肉体上受到照耀，但并非两人都在视力中得见光明。一人看见，另一人看不见：太阳对两人都在场，但一人对在场的太阳不在场。同样，天主的智慧、天主的圣言、主耶稣基督，处处都在，因为真理处处都在，智慧处处都在。一人在东方理解正义，另一人在西方理解正义；一人所理解的正义与另一人所理解的正义是不同的吗？在身体上他们相隔遥远，但他们心灵的眼睛都注视着同一对象。我在此处所见到的正义，如果真是正义，与那在肉体上与我相隔千万里行程的义人所见的正义是相同的，并且他在那正义之光中与我合一。因此，光明为自己作证；它开启健康之眼，并成为自己的见证，好让人认出它是光明。但那些不信的人呢？它不也在他们面前吗？它也在他们面前，但他们没有心灵的眼睛去看它。请听取自福音本身关于他们的句子：\BibleQuote{光在黑暗中照耀，黑暗却没有接受它。} 见：\BibleRef{若 1:5} 因此主说，且说得真实：\BibleQuote{我即使为自己作证，我的见证也是真实的，因为我知道我从哪里来，往哪里去。} 祂在这里要我们理解圣父：圣子将光荣归于圣父。那同等的自体荣耀那派遣祂的一位。人该如何荣耀那创造他的一位呢！\par

5. \BibleQuote{我知道我从哪里来，往哪里去。} 那位在你们面前说话者，拥有祂未曾离弃的，然而祂来了；因为祂来并没有从那里离去，祂回去也没有抛弃我们。你们为什么惊奇？这是天主；人做不到；连太阳也做不到。当太阳落向西时，它离开了东方，直到它回到东方、即将升起的时候，它都不在东方；但我们的主耶稣基督既来又在那里，既回去又在这里。请听福音作者本人在另一处说的话，如果可能，就理解；如果不能，就相信：\BibleQuote{天主，从来没有人见过，只有那在圣父怀里的独生子，祂给我们详述了。} 祂没有说\Emph{曾在}圣父怀里，好像祂来就离开了圣父的怀抱一样。祂在这里说话，却宣告祂在那里；当祂要离开这里时，祂说了什么？\BibleQuote{看，我天天与你们在一起，直到今世的终结。} 见：\BibleRef{玛 28:20}\par

6. 因此，光明的见证是真实的，无论它是在显示自身还是显示其他事物；因为没有光，你无法看见光，没有光，你也无法看见任何不是光的其他事物。如果光能够显示不是光的其他事物，它就不能显示自身吗？凡物不藉它便不能显明，它岂不显示自己？先知说了真理；但他从哪里得到它，除非他从真理的泉源汲取？若翰说了真理；但他从哪里说，问他本人：\BibleQuote{我们都从祂的圆满中领受了。} 因此，我们的主耶稣基督配得为自己作证。但无论如何，我的弟兄们，我们这些在现世黑夜中的人，也要热切地听从预言；因为现在主愿意以谦卑来到我们的软弱和我们内心深处深沉的黑夜中：祂来为人所轻视，也为人所尊崇；祂来为人所否认，也为人所承认；为犹太人所轻视和否认，为我们所尊崇和承认；祂来受审判，也要审判；受不公义的审判，作公义的审判。祂如此来临，以致祂需要一盏灯为祂作证。因为如果我们的软弱能直视白日本身，那么若翰作为灯为白日作证还有什么需要呢？但我们不能直视它：祂为软弱者成为软弱；以软弱医治软弱；以必死的肉体除去肉体的死亡；以祂自己的身体为我们的眼睛制作了药膏。因此，既然主已经来了，而我们仍处在现世的黑夜中，我们也应当听从预言。\par

7. 我们正是借着先知书来说服反驳的外邦人。外邦人说：\Quote{基督是谁？} 我们回答：\Quote{就是先知所预言的那一位。} 他问：\Quote{什么先知？} 我们引用 \Person{依撒意亚}、\Person{达尼尔}、\Person{耶肋米亚}及其他圣先知，告诉他们，他们远在基督之前就已来到，距祂来临有多长时间。我们这样回答：\Quote{先知们比祂先来，并预言了祂的来临。} 其中一人回答：\Quote{什么先知？} 我们引用每日向我们宣读的那些先知书。他说：\Quote{这些先知是谁？} 我们回答：\Quote{也是预言了我们所见之事得以实现的那几位。} 他坚持说：\Quote{你们为自己伪造了这些，你们眼见它们应验，就把它们写在你们喜欢的书中，好像它们预先被预言了一样。} 在此，反对外邦仇敌时，其他仇敌的见证自身出现。我们拿出犹太人写的书，回答说：\Quote{无疑你和他们都是我们信德的仇敌。因此他们被分散在万民中，好让我们能用一种仇敌来说服另一种仇敌。} 让犹太人拿出 \Person{依撒意亚}书，让我们看看是否在那里读到：\BibleQuote{祂被牵去如同待宰的羊，又像羔羊在剪毛者前缄默无声，祂总不开口。祂受审判时被夺去了公道；因祂的创伤我们获得了痊愈。我们都像羊一样迷了路，祂为我们的罪被交出了。} \BibleRef{依 53:5} 看，一盏灯！再拿出另一本，打开圣咏，从那里也引用基督预言的苦难：\BibleQuote{他们刺透了我的手和我的脚，他们数算了我的所有骨头；他们注视着我，端详着我，他们分了我的衣服，为我的长衣拈阄。我的赞美归于祢；在伟大的集会中，我将向祢称谢。大地四极将记起并归向主；万国的各族都要在祂面前朝拜；因为王国属于主，祂将统治万民。} 让一位敌人羞愧，因为是另一位敌人给了我这本书。但是，看，从一位敌人提供的书中，我战胜了另一位敌人；也不要放过给我那本书的敌人，让他拿出那本也能使他被战胜的书。我读另一位先知，发现主对犹太人说：\BibleQuote{从日出之地到日落之处，我的名在万民中要受洁净的祭献。} \BibleRef{拉 1:10} 你，犹太人，不前来领受洁净的祭献；我证明你是不洁的。\par

8. 看，甚至灯也为白昼作证，因为我们的软弱，我们无法承受并注视白昼的光辉。的确，与不信者相比，我们基督徒如今就是光；如宗徒所说：\BibleQuote{你们从前原是黑暗，但现在在主内却是光明；生活就像光明之子一样。} \BibleRef{弗 5:8} 他在别处又说：\BibleQuote{黑夜深了，白日已近，所以我们该脱去黑暗的行为，佩带光明的武器；让我们端庄地行事，犹如在白昼中。} \BibleRef{罗 13:12} 然而，即使我们现在所处的白昼，与我们将要达到的光明相比，仍然是黑夜。请听宗徒 \Person{伯多禄}的话：他说，从极其荣耀中有一个声音对主基督说：\BibleQuote{祢是我的爱子，我所喜悦的。这声音从天上传来，我们在圣山上与祂在一起时，曾听见这声音。} 但我们没有在那里，也没有听到这来自天上的声音，同样的 \Person{伯多禄}对我们说：\Quote{我们还有更确实的先知话。} 你没有听到来自天上的声音，但你却有更确实的先知话。因为主耶稣基督预见到，会有某些恶人用巫术来诽谤祂的奇迹，就派先知在祂之前来。因为，假如祂是巫师，用巫术使祂死后受人敬拜，那么祂在出生之前就已经是巫师了吗？听先知们吧，死人啊，繁殖诽谤之虫的人啊，听先知们吧！我读，听那些在主以前来的人。宗徒 \Person{伯多禄}说：\BibleQuote{我们有更确实的先知话，你们要留心听，如同灯照在暗处，直到天亮，晨星在你们心中升起。} \BibleRef{伯后 1:17}\par

9. 所以，当我们的主耶稣基督来临，正如宗徒保禄也说，祂\BibleQuote{要照亮暗中的隐情，并显明心中的意念，好使各人从天主那里获得称赞}时，\BibleRef{格前 4:5} 在那样的日子面前，就不再需要灯盏了：那时没有先知会读给我们听，没有宗徒的书卷会被打开；我们不需要若翰的见证，也不需要福音本身。因此，所有圣经——它们在这世界的黑夜中，如同为我们点燃的灯，使我们不致留在黑暗里——都要被挪开；当这一切都被除去了，好让它们不再发光好似我们还需要它们，而那些为我们传递了这些的天主之人，自己也要和我们一同瞻望那真实而清澈的光。那么，在这些辅助被移除之后，我们将会看见什么？我们的心智将以什么为食？我们的目光将因什么而喜悦？那眼所未见、耳所未闻、人心也未曾想到的喜乐将从何处涌现？我们将看见什么？我恳求你们，与我一同爱慕，以信德与我一同奔跑：让我们渴慕我们的天上家乡，让我们切望我们的天上家乡，让我们感到自己在此是客旅。那时我们将看见什么？让福音现在告诉我们：\BibleQuote{在起初已有圣言，圣言与天主同在，圣言就是天主。} 你将来到那泉源，它的一滴露水已曾洒在你身上：你将看见那真光本身，那曾有一条光线弯弯曲曲地斜射入你黑暗心中的光，现在以其纯净供你观看和承载，而你也正为此被净化。若望自己说——我昨天也引用了：\BibleQuote{亲爱的，我们现在是天主的子女，但我们将来如何，还没有显明；我们知道，当祂显现时，我们必要相似祂，因为我们要看见祂实在怎样。} \BibleRef{若一 3:2} 我感觉你们的情感正与我一同被提升到天上之事；但那可朽坏的身体重压着灵魂，这尘世的寓所抑制着多虑的心。\BibleRef{智 9:15} 我即将放下这本书，你们也要离去，各人回自己的家。我们曾一起在公共之光中，这很好，曾为此欢喜，这很好，曾为此喜乐，这很好；但当我们彼此分离时，愿我们不要离开祂。\par

\SourceRecord{Translated by John Gibb.；Nicene and Post-Nicene Fathers, First Series；Vol. 7.；Edited by Philip Schaff.；Buffalo, NY: Christian Literature Publishing Co.,；1888.；<http://www.newadvent.org/fathers/1701035.htm>.}

% structure: p0001 source=h1 target=section reason=page-title

\section{论道第36篇（\BibleRef{若 8:15-18}）}

1. 在四福音中，或更确切地说，在唯一福音的四卷书中，圣若望宗徒（因其属灵理解而被喻为鹰，并非无端）将他的宣讲提升得比其他三位更高、更超越；在这种提升中，他也愿意我们的心同样被举扬。因为其他三位圣史与主在世同行，如同与人同行；关于他的天主性，他们说得很少；但这位圣史，仿佛不屑于在地上行走，正如在他言论的开端，他就向我们发出雷霆之言，翱翔不仅超越了大地和整个空气与天空的领域，甚至超越了全体天使和无形力量的整个秩序，直达那创造万物者，说：\BibleQuote{在起初已有圣言，圣言与天主同在，圣言就是天主。圣言在起初就与天主同在。万物是藉着他而造成的；凡受造的，没有一样不是由他而造成的。} 与此开端如此崇高的境界相称，他其余的全部宣讲也都吻合；他论及主的天主性，是其他人所不及的。他所汲取的，正是他所流出的。因为在这部福音中记载他在晚餐时斜靠在主的怀里，这并非无因。他从那怀中秘密地汲取；但他秘密汲取的，公开地流出来，为使万民得以知道，不仅天主子降生成人、他的受难与复活，而且他在降生成人之前是什么：圣父的独生子、圣父的圣言，与生他的父同永恒，与派遣他的父同等；然而在这派遣中自谦，好使父更为伟大。\par

2. 无论你听到关于主耶稣基督以卑下方式所说的任何话，都要思考那\OriginalTerm{工程}{economy}，他藉此取了肉身；但无论你在福音中听到或读到关于他所说的高超、超越万物、神圣、与父同等同永恒的话，要确信你所读到的属于天主的形象，而非仆人的形象。因为如果你持守这条准则，你就能理解它（因为并非你们都能理解，但你们全都应当相信）——如果我说，你持守这条准则，作为在光明中行走的人，你将无畏地与异端黑暗的诽谤作战。因为不乏有人，在读福音时，只关注与基督谦卑相关的证据，却对宣告他天主性的证据充耳不闻；他们充耳不闻，好让自己满口恶言。同样也有人，只注意那些论及主崇高的话，即便他们读到主为了我们而降生成人的仁慈，也不相信这些证据，反认为它们是假的，是人捏造的；争辩说我们的主耶稣基督只是天主，并非也是人。一方如此，另一方如彼：两者皆错。但公教信仰，从两者中持守各自持有的真理，并宣讲各自所信的真理，既理解了基督是天主，也相信了他也是人：因为两者都已记载，两者都是真实的。如果你断言基督只是天主，你就否认了使你痊愈的医药；如果你断言基督只是人，你就否认了创造你的德能。因此，两者都持守。噢，忠信的灵魂，公教的心，两者都持守，两者都相信，忠信地宣认两者。基督既是天主，也是人。基督如何是天主？与父同等，与父合一。基督如何是人？由童贞女所生，从人取了可朽性，但未取不义。\par

3. 这些犹太人于是看到了那人；他们既没有察觉，也没有相信他是天主：你已经听到，在其余一切中，他们对他说：\BibleQuote{你为自己作证，你的证据不可信。} 你也听到了他如何回答，正如昨天所读给你们的，并按照我们的能力讨论过了。今天读到了他的这些话：\BibleQuote{你们凭肉身判断。} 因此，他说，这就是为什么你们对我说：\BibleQuote{你为自己作证，你的证据不可信，} 因为你们凭肉身判断，因为你们不察觉天主；你们看到那人，并因迫害那人，而得罪了隐藏在他内的天主。\BibleQuote{你们}，那么，\BibleQuote{凭肉身判断。} 因为我为自己作证，我在你们眼中就显得狂妄。因为每个人，当他想要为自己作证以引人注目时，就显出自大和骄傲。因此经上记载：\BibleQuote{不要让你自己的口赞美你，而要让邻人的口赞美你。}\BibleRef{箴 27:2} 但这是对人说的。因为我们是软弱的，我们向软弱的人发言。我们能讲真理，但也能说谎；虽然我们应讲真理，但我们仍有能力随时说谎。但绝不要认为，在神圣之光的荣耀中能发现谎言的黑暗。他作为光说话，作为真理说话；但光在黑暗中照耀，黑暗却不理解它：因此他们凭肉身判断。\BibleQuote{你们}，他说，\BibleQuote{凭肉身判断。}\par

4. \Quote{我不审判任何人。}难道主耶稣基督不审判任何人吗？祂不就是我们所宣认的那一位——第三天复活，升了天，坐在圣父的右边，日后要从那里降来审判生者死者吗？这不就是我们的信德吗？宗徒说：\Quote{心里相信使人成义，口里宣认使人获得救恩。}（\BibleRef{罗 10:10}）那么，当我们宣认这些事时，我们与主的话有冲突吗？我们说祂要作为审判生者死者的那一位而来，而祂自己却说：\Quote{我不审判任何人。}这个问题可以从两方面解决：要么我们理解这话\Quote{我不审判任何人}——意即\Emph{现在}我不审判任何人，这与祂在另一处所说的一致：\Quote{我来不是为审判世界，乃是为拯救世界；}祂在此并非否认祂的审判，而是将其推迟。要么，当然，当祂说了\Quote{你们是按肉身判断}之后，紧接着说\Quote{我不审判任何人}，意思是应该把\Quote{按肉身}作为省略的部分来理解。因此，不要让任何疑虑的顾忌存留在我们心中，反对我们所持守并宣认的关于基督为审判者的信德。基督已经来了，但首先是来拯救，然后才审判：将那些不愿得救的人判罚，使那些因相信而未拒绝救恩的人获得生命。因此，主耶稣基督的第一项安排是医治性的，而非审判性的；因为如果祂先来审判，祂会发现没有一个人可以赏赐正义的酬报。因为祂看见众人都是罪人，无人免于罪恶的死亡，所以祂的仁慈必须先被祈求，然后祂的审判才得以施行；因为诗篇曾歌唱祂说：\Quote{上主，我要向祢歌颂仁慈和审判。}祂并未说\Quote{审判和仁慈}，因为如果审判在先，就没有仁慈了；但仁慈在先，然后才是审判。什么是先行的仁慈？人的创造者屈尊成为人，成了祂所造的，使祂所造的受造物不至丧亡。这仁慈还能再增加什么？然而祂还增加了。祂成为人还不够，又加上了被人弃绝；被弃绝还不够，又受凌辱；受凌辱还不够，又被处死；但这还不够，是被钉死在十字架上。因为宗徒在称颂祂的服从至死时，说\Quote{祂服从至死}还不够；因为这并非任何一种死亡，他又加上：\Quote{且死在十字架上。}（\BibleRef{斐 2:8}）在所有的死亡方式中，没有比这更惨烈的了。简言之，那种使人受最大痛苦的方式被称为\OriginalTerm{酷刑}{cruciatus}，其名源自\OriginalTerm{十字架}{crux}。因为被钉者悬挂在木头上，被钉在木头上，是缓慢地折磨至死。被钉十字架不仅是处死；受刑者在十字架上存活很久，并非因为选择了更长的生命，而是因为死亡本身被延长，使痛苦不至于太快结束。祂愿意为我们而死，但这样说还不够；祂屈尊被钉十字架，服从至死，且死在十字架上。那将要除去一切死亡的那一位，选择了最低下、最惨烈的死亡方式：祂以最惨烈的死亡消灭了死亡。对于不明白的犹太人来说，这确实是最惨烈的死亡，但主选择了它。因为祂要这十字架本身作为祂的标志；那十字架，可以说，是对被战胜的魔鬼的战利品，祂要把它放在信徒的额头上，所以宗徒说：\Quote{但愿我绝不以别的夸耀，只以我们的主耶稣基督的十字架来夸耀，借着祂，世界于我已被钉在十字架上了，我于世界也被钉在十字架上了。}（\BibleRef{迦 6:14}）那时，在肉体上没有什么比这更难以忍受的了；如今，在额头上没有什么比这更荣耀的了。祂把如此光荣加在祂自己受刑的工具上，祂为忠信的人还保留着什么？现在十字架在罗马人中间已不再用于惩罚罪犯，因为当主的十字架开始被尊敬时，人们认为即使是有罪的人，如果被钉十字架，也会受到尊敬。因此，为此而来的那一位没有审判任何人：祂也遭受了恶人的苦。祂遭受了不义的审判，为要施行公义的审判。但祂以仁慈忍受了不义的审判。总之，祂降得如此卑微，以至于上了十字架；是的，祂放下了自己的权能，却彰显了祂的仁慈。祂在哪里放下了祂的权能？在于祂不愿从十字架上下来，虽然祂有能力从坟墓中复活。祂在哪里彰显了祂的仁慈？在于祂悬在十字架上时说：\Quote{父啊！宽赦他们吧！因为他们不知道他们做的是什么。}（\BibleRef{路 23:34}）因此，无论祂说\Quote{我不审判任何人}是因为祂来不是为审判世界，乃是为拯救世界；还是如同我所提到的，当祂说了\Quote{你们是按肉身判断}之后，加上\Quote{我不审判任何人}，让我们理解基督不是按肉身审判，如同人审判祂那样。\par

5. 但为使你们知道基督如今已是审判者，请听下文：\Quote{若我判断，我的判断是真实的。}看哪，你们有祂作你们的审判者，但要承认祂是你们的救主，免得你们感受到审判者。但祂为什么说祂的判断是真实的？\Quote{因为，}祂说，\Quote{我不是独自一人，而是我与派遣我的父。}弟兄们，我已经对你们说过，这位圣史\Person{若望}翱翔得极高：他很难被理解。但亲爱的，我们需要提醒你们，这翱翔更深的奥秘。在\BibleBook{厄则克耳}{先知书}中，以及在这位\Person{若望}（他的这部\BibleBook{若望}{福音}）的\BibleBook{若望}{默示录}中，都提到一个四面的活物，有四张特征的面孔：人的、牛的、狮子和鹰的。在我们之前处理圣经奥秘的人，大部分都把这活物，或更确切地说，这四个活物，理解为四位福音作者。他们理解狮子代表君王，因为狮子凭其力量和可畏的勇猛，似乎是百兽之王。这种特征归于\Person{玛窦}，因为在主的世代中他遵循王室谱系，表明主如何沿王室谱系，是\Person{达味}的后裔。但\Person{路加}，因为他以\Person{匝加利亚}的司祭职开始，提到\Person{洗者若翰}的父亲，被指明为牛；因为公牛是司祭祭献中的重要祭品。\Person{马尔谷}被恰当地赋予基督的人性，因为他既未说关于王室权威的任何事，也未以司祭职开始，而只是以基督的人性开始。所有这些都很少离开世间的事物，就是说，离开我们的主耶稣基督在地上所行的事；关于祂的天主性，他们说得很少，就像与祂在地上同行的人。剩下的就是鹰；这就是\Person{若望}，崇高真理的宣讲者，以坚定目光凝视内在永恒之光的默观者。据说，幼鹰是这样被双亲测试的：幼鹰被悬在雄鹰的爪下，直接暴露在日光中；若它稳定地看太阳，就被承认为真正的后代；若它的眼睛闪烁，就被允许掉落，视为不纯的后代。因此，现在，想一想，被称为鹰的人应该谈论多么崇高的事物；然而，就连我们这些在地上爬行、软弱、在人中几乎不足道的人，也敢于处理并讲解这些事；并且以为我们或在默想时能领会，或讲论时能被人领会。\par

6. 我为何说这些话？或许在这些话之后，有人可以公义地对我说：那就放下书吧。你为何要从事超过你能力的事？为何要托付你的口舌？对此我回答：异端者众多；天主容许他们如此众多，为的是我们不再常被奶喂养而停留在无知的婴孩状态。因为他们既不明白基督的天主性如何被提供给我们接受，就凭己意下结论；并且因没有正确辨别，就给公教信徒带来了最麻烦的问题；信徒的心开始被扰乱而摇摆。然后立即成为属神之人的必要，他们不仅在福音中读过关于主耶稣基督的天主性的事，而且也已明白，要拿出基督的武器来对抗魔鬼的武器，并竭尽全力为基督的天主性公开争战，反对虚假及欺骗人的教师；免得因他们沉默而别人丧亡。因为凡是认为主耶稣基督是与圣父不同实体的人，或以为只有基督，以致圣父、圣子和圣神是同一位的；凡是选择以为祂只是人，而非天主成为人，或是天主而祂的天主性是可变的，或是天主而非人的人，这些人都在信仰中沉船，并从教会的港口被驱逐出去，免得他们的不安使同行的船失事。这事迫使我们，虽是最小的，就自身而言完全无足轻重，但凭祂的仁慈在祂的管家中算得上有份的，应向你们讲论，使你们或能明白而与我一同喜乐，或若还不能明白，藉着相信它，你们可以安稳地留在港口内。\par

7. 我因此要说话；能理解的，让他理解；不能理解的，让他相信：但我还是要说出主所说的话：\Quote{你们是按肉身判断；我不判断任何人}，无论是现在，还是按肉身。\Quote{但即使我判断，我的判断也是真的。}为什么祢的判断是真的？\Quote{因为我不是独自一人，}祂说，\Quote{而是我与派遣我的圣父在一起。}那么，主耶稣啊？如果祢是独自一人，祢的判断就是假的吗？还是因为祢不是独自一人，而是祢与派遣祢的圣父在一起，祢就判断得真？我该如何回答？让祂亲自回答：祂说：\Quote{我的判断是真的。}为什么？\Quote{因为我不是独自一人，而是我与派遣我的圣父在一起。}如果祂与祢同在，祂又怎能派遣祢？祂派遣了祢，却又与祢同在？是这样的：祢被派遣，却没有离开祂？祢来到我们这里，却仍留在那里？这如何能相信？如何能理解？对这两个问题我回答：你说得对，如何理解；但如何相信，你说得不对。相反，正因如此，相信它是对的，因为它不能立即被理解；因为如果它是能够立即被理解的事，就没有必要相信它，因为它会被看见。正是因为你不能理解，所以你相信，但借着相信，你变得有能力理解。因为如果你不相信，你永远无法理解，因为你将保持更小的能力。让信德洁净你，好使领悟充满你。\Quote{我的判断是真的，}祂说，\Quote{因为我不是独自一人，而是我与派遣我的圣父在一起。}因此，主我们的天主，耶稣基督，祢的派遣就是祢的降生成人。我如此看见，我如此理解；简言之，我如此相信，以免说\InnerQuote{我如此理解}会显得傲慢。无疑，主耶稣基督甚至在这里；更确切地说，祂曾在此按肉身，现在按祂的天主性在此；祂既与圣父同在，又没有离开圣父。因此，祂被说成被派遣并来到我们这里，这是祂的降生成人被启示给我们，因为圣父没有取肉身。\par

8. 因为有某些称为\OriginalTerm{撒贝利派}{Sabellians}的异端者，也称为\OriginalTerm{父受苦派}{Patripassians}，他们断言是圣父自己受了苦。啊，公教徒，你不要如此断言；因为如果你成为父受苦派，你就不清醒。那么，你要理解，圣子的降生成人被称为圣子的派遣；不要相信圣父降生成人，但也不要相信祂离开了降生成人的圣子。圣子取了肉身，圣父与圣子同在。如果圣父在天上，圣子在地上，圣父如何与圣子同在？因为圣父和圣子无所不在：因为天主并不在天上以至于不在世上。听那想逃避天主的审判、却找不到逃避之路的人所说的：\Quote{我往何处去，}他说，\Quote{离开祢的圣神；我往何处逃避祢的面容？我若升到天上，祢在那里。}问题涉及地上；听下文：\Quote{我若下到地狱，祢也在那里。}那么，如果祂甚至在地狱也被说成临在，宇宙中还有什么地方祂不在呢？因为天主借着先知的声音说：\BibleQuote{我充满天地。}\BibleRef{耶 23:24}因此祂无所不在，祂不受任何地方限制。你不要转身离开祂，祂就与你同在。你若愿意到祂那里去，不要迟延去爱；因为你不是用脚，而是用情感奔跑。你虽然停留在一处，却来到了祂那里，如果你相信并爱慕。因此祂无所不在；如果无所不在，祂怎能不与圣子同在？难道祂不与圣子同在，而如果你相信，祂甚至与你同在吗？\par

9. 那么，祂的判断如何是真实的，不正是因为圣子是真实的吗？因为祂曾说：\BibleQuote{如果我来判断，我的判断是真实的；因为我不是独自一人，而是我与那派遣我来的父。} 就如同祂曾说：\BibleQuote{我的判断是真实的}，因为我是天主子。祢如何证明祢是天主子？\BibleQuote{因为我不是独自一人，而是我与那派遣我来的父。} 脸红吧，撒贝里乌派！你听到了 \Emph{子}，你听到了 \Emph{父}。父是父，子是子。祂没有说：我是父，同时我也是子；而是说：\BibleQuote{我不是独自一人。} 为什么祢不是独自一人？因为父与我同在。\BibleQuote{我是，且那派遣我来的父。} 你听到了：\BibleQuote{我是，且那派遣我来的。} 免得你忽视了位格，请区分位格。以理智区分，不要因不信而分隔；免得你如同逃离\OriginalTerm{卡律布狄斯}{Charybdis}，却又撞上\OriginalTerm{斯库拉}{Scylla}。因为撒贝里乌派不敬的漩涡曾吞噬你，说父与子是同一者：刚才你已学到：\BibleQuote{我不是独自一人，而是我与那派遣我来的父。} 你承认父是父，子是子；你承认得对。但不要说父更伟大，子更小；不要说父是金子，子是银子。只有一个本体，一个神性，一个同永世，完全的平等，无任何不像。因为如果你只相信基督是另一位，不是父那同一者，但却以为祂的本性与父有所不同，那么你确实逃脱了卡律布狄斯，却撞上了斯库拉的礁石。要驾驶中道，避免两侧的危险。父是父，子是子。你现在说：父是父，子是子：你幸运地逃脱了吞噬漩涡的危险；为什么你要走向另一边，说父是这个，子是那个？子是与父不同的另一位，这一点你说得对；但如果说祂的本性不同，你就不对了。诚然，子是一位不同的位格，因为祂不是父那同一者；父也是一位不同的位格，因为祂不是子那同一者；然而，在本性上它们并非不同，而是父与子是同一个。什么是\Quote{同一个}？天主是唯一的。你已听到：\BibleQuote{因为我不是独自一人，而是我与那派遣我来的父}：请听你如何相信父与子；听子自己说：\BibleQuote{我与父是同一的} \BibleRef{若 10:30} 祂没有说我是父，或我与父是一个位格；但当祂说\BibleQuote{我与父是同一的}时，你要听到两者：既听到\OriginalTerm{一}{unum}，又听到\OriginalTerm{是}{sumus}，然后你便会从卡律布狄斯和\OriginalTerm{斯库拉}{Scylla}中获救。在这两个词中，因祂说了\OriginalTerm{一}{unum}，祂将你从亚流主义中救出；因祂说了\OriginalTerm{是}{sumus}，祂将你从撒贝里乌主义中救出。因为对一位个格，祂不会说\OriginalTerm{是}{sumus}；另一方面，对不同的实体，祂不会说\OriginalTerm{一}{unum}。因此，祂说\BibleQuote{我的判断是真实的}的理由，就是要你简短地听到：因为我是天主子。但我希望你在相信我是天主子时，能明白父与我同在：我并非如此作为子而离开了祂；我并非如此在这里而不与祂同在；祂也并非如此在那里而不与我同在：我取了奴仆的形象，却没有失去天主的形象；因此祂说：\BibleQuote{我不是独自一人，而是我与那派遣我来的父。}\par

10. 祂已谈论了判断；祂现在要谈论见证。祂说：\BibleQuote{在你们的法律中记载着：两个人的见证是真实的。我是为自己作证的一位，那派遣我来的父也为我作证。} 祂也同样给他们解释了法律，如果他们不是忘恩负义的话。因为这是一个大问题，我的弟兄们，在我看来，这确实是在一个奥秘中预定好的，当时天主说：\BibleQuote{凭两个或三个见证的口，每句话都得成立。} 是否藉两个见证寻求真理？显然是的；这也是人类的习俗：但即使两个见证也可能说谎。贞洁的\Person{苏撒纳}曾被两个假见证所迫：难道因为他们是两个，就不是假的吗？我们说的是两个还是三个？全体人民曾反对基督作证（\BibleRef{路 23:1}）。那么，如果由一大群人组成的全体人民被证明是假见证，那么如何理解\BibleQuote{凭两个或三个见证的口，每句话都得成立}呢？除非在这一方式中，天主圣三被奥秘地显示给我们，在其中真理有永久的稳固？你希望有一个好的理由吗？要有两个或三个见证——圣父、圣子和圣神。简言之，当贞洁的女子和忠信的妻子\Person{苏撒纳}被两个假见证所迫时，天主圣三在她的良心和隐秘中支持了她：那天主圣三从隐秘中兴起一个见证——\Person{达尼尔}，并定罪了那两个。因此，因为你们法律中记载着两个人的见证是真实的，请接受我们的见证，免得你们感受到我们的判断。因为祂说：\BibleQuote{因为我不判断任何人，但为自己作证。} 我推迟判断，我不推迟见证。\par

11. 诸位弟兄，让我们为自己选择天主作我们的审判者，天主作我们的见证者，以对抗人的口舌，对抗人类脆弱的猜疑。因为那审判者并不轻看作见证者，祂在成为审判者时并未受到尊荣的提升；因为那作见证者也将亲自成为审判者。祂如何作见证者？因为祂无需向别人询问以得知你是谁。祂如何作审判者？因为祂有权柄杀死并使人活，定罪并释放，抛入地狱并提升到天堂，与魔鬼结合并用天使加冕。因此，既然祂拥有这权柄，祂就是审判者。现在，因为祂不需要另一位见证者来认识你；而且这位将来审判你的祂，现在正看着你，你没有任何办法能在祂开始审判时欺骗祂。因为你无法提供虚假的见证者来蒙蔽那位审判者，当祂开始审判你时。这就是天主向你所说的：当你轻视时，我看见了；当你不信时，我没有搁置我的判决。我延迟了它，并非撤销了它。你不愿听从我所命令的，你将感受到我所预言的。但如果你听从我所命令的，你就不会感受到我所预言的灾祸，反而会享受我所应许的福乐。\par

12. 任何人都不应感到惊奇，祂说：\Quote{我的判断是真实的，因为我不是独自一个，而是我和那派遣我的圣父；}而在另一处祂却说：\Quote{圣父不判断任何人，但一切判断都已交给了圣子。}关于福音作者的这些相同话语，我们已经讨论过了，现在我们要提醒你们，这话并不是说圣父在圣子来审判时不会与圣子同在，而是说在审判中惟有圣子以祂曾受苦、复活并升入天堂的那一形态显现在善人与恶人面前。因为那时，当门徒们正看着祂升入天堂时，天使的声音在门徒们耳中响起：\Quote{祂怎样你们见祂往天堂去，也要怎样降来；}\BibleRef{宗 1:11}也就是说，祂以曾经受过审判的人形来审判，为使那先知的话也应验：\Quote{他们要瞻望他们所刺透的。}但当义人进入永生时，我们要看见祂实在怎样；那将不是对生者与死者的审判，而仅是生者的赏报。\par

13. 同样地，不要因祂说：\Quote{在你们的法律上记载着：两个人的见证是可凭信的，}而任何人因此以为这不是天主的法律，因为经文并未说\Emph{在天主的法律上}；要知道，当这样说\Emph{在你们的法律上}时，就好像祂说：\Quote{在赐给你们的法律上；}这法律是谁赐给的？除非是天主。正如我们说的：\Quote{我们的日用粮，}而我们又说：\Quote{求你今天赐给我们。}\par

\SourceRecord{Translated by John Gibb.；Nicene and Post-Nicene Fathers, First Series；Vol. 7.；Edited by Philip Schaff.；Buffalo, NY: Christian Literature Publishing Co.,；1888.；<http://www.newadvent.org/fathers/1701036.htm>.}

% structure: p0001 source=h1 target=section reason=page-title

\section{\WorkTitle{讲道集 37（若望福音 8:19-20）}}

1. 圣福音中简略提及的事，不应简略地加以阐释，好使所读的得以理解。主的言语虽然少，却伟大；不应以数量衡量，而应以重量估量；不应因其少而轻视，而应因其伟大而探求。你们昨日在场的人，已经听到我们根据主所说的那番话，尽我们所能讲解的：\BibleQuote{你们凭肉身判断；我不判断任何人。即使我判断，我的判断也是真的，因为我不是独自一个，而是我和那派遣我来的父。你们的法律上记载着：两个人的见证是真的。我是为自己作证的人，那派遣我来的父也为我作证。}正如我所说，昨日已就这些话向你们的耳朵和心灵作了讲道。当主说了这些话后，那些听见\BibleQuote{你们凭肉身判断}的人，显明了他们所听见的是真的。因为他们答覆主——当时主正谈论他的天父——对他说：\BibleQuote{你的父在哪里？}他们对基督的父作了肉身的理解，因为他们凭肉身判断基督的话。然而那说话者外表是肉身，内里却是圣言：可见的人，隐藏的天主。他们看见了覆盖物，却轻视了穿戴者；他们因不认识而轻视；因看不见而不认识；因瞎眼而看不见；因不信而瞎眼。\par

2. 那么，让我们看看主如何回答这一点。他们说：\BibleQuote{你的父在哪里？}因为我们听你说：\BibleQuote{我不是独自一个，而是我和那派遣我来的父。}我们只见你独自一人，不见你的父与你同在；你如何说你不是独自一个，而是与你父同在？否则就向我们显明你的父与你同在。主回答他们说：你们认识我，好让我把父显示给你们吗？这确实是下文；这正是他用他自己的话所作的答复，对此我们已事先作了阐释。请看他说了什么：\BibleQuote{你们既不认识我，也不认识我的父；若是认识我，或许也认识我的父。}那么你们说：\BibleQuote{你的父在哪里？}仿佛你们已经认识我；仿佛你们所看见的就是我的全部。因此，因为你们不认识我，我就不把父显示给你们。事实上，你们以为我只是一个人；因而你们为我寻找一个人类的父亲，因为\BibleQuote{你们凭肉身判断。}然而，按照你们所见的，我是某物；而按照你们所不见的，我是另一物；并且，作为向你们隐藏的那一位，我谈论同样隐藏的父，你们必须先认识我，然后才认识我的父。\par

3. \BibleQuote{若是认识我，或许也认识我的父。}那无所不知的主说\Emph{或许}时并非处于疑惑，而是在责备。请看，这个\Emph{或许}一词，看似表达怀疑，却可以责备地说。是的，当人使用时，它是一个怀疑的词，因为人因不知而怀疑；但当天主说出怀疑的词时——他毫无隐藏——那不信任受到那怀疑的责备，而非神性仅仅表达意见。因为人有时会以责备的方式对确定的事表示怀疑，即使用怀疑的词，心中却不怀疑；正如你对自己的仆人发怒时会说：\Quote{你藐视我；但想想吧，或许我是你的主人。}同样，宗徒对某些藐视他的人说：\BibleQuote{我想我也有天主的神。}\BibleRef{格前 7:40}当他说\Emph{我想}时，他看似怀疑，却是在责备，而非怀疑。在另一处，主耶稣基督自己责备人类将来的不信时说：\BibleQuote{人子来临时，你认为他会找到信德在地上吗？}\BibleRef{路 18:8}\par

4. 你们现在，如我所想，明白这里如何使用了\Emph{或许}一词，以免任何掂量文字、平衡音节的人——仿佛要显摆他的拉丁文知识——挑剔天主圣言所说的词，并因责备天主圣言而不再口若悬河，反成哑口无言。因为谁能像那在起初就与天主同在的圣言那样说话呢？不要按我们使用这些词语的方式来看待它们，并由此去衡量那圣言——即天主。你们确实听见了圣言，却轻视它；听见了天主，就畏惧他：\BibleQuote{在起初已有圣言。}你参考你日常交谈的用法，心里说：什么是词？词是什么了不起的东西？它发出声音，随即消逝；撞击空气后，击中耳朵，便不复存在。再听下去：\BibleQuote{圣言与天主同在；}它常存，并不因发声而消逝。也许你仍然轻视它：\BibleQuote{圣言就是天主。}在你内，人啊，你心里的词与声音是两回事；但你那里的词，为了传递到我这里，需要声音作为载体。它获取声音，乘上这载体，穿过空气，来到我这里，却不离开你。然而声音为了到我这里，离开了你，却未留在我这里。那么你心里的词是否也随着声音的消逝而消逝了呢？你表达了你的思想；为了使你那隐藏在你内的思想传达给我，你发出了音节；音节的声响将你的思想传递到我的耳中；通过我的耳朵，你的思想降入我的心内，中间的声音飞逝了；但那个获取了声音的词，在你发出声音之前已在你内，并且因为你发出了声音，它便在我内，却并未离开你。想想这个吧，无论你是谁，你这斤斤计较声音的人。你轻视天主圣言，因为你连人的言词也不理解。\par

5. 那麼，萬物由祂所造的那位知道一切，但祂卻以疑惑的口吻斥責說：\Quote{你們若認識我，也許就認識我的父了。}祂斥責不信的人。祂對門徒也說過類似的話，但其中沒有一絲疑惑，因為無需斥責不信。因為祂對猶太人所說的這句話：\Quote{你們若認識我，也許就認識我的父了}，祂也對門徒們說過，當斐理伯詢問祂，或者更確切地說，向祂要求說：\Quote{主，請把父顯示給我們，我們就滿足了。}就好像他說：我們自己已經認識祢；祢向我們顯現了；我們看見了祢；祢揀選了我們；我們跟隨了祢，看見了祢的奇蹟，聽到了祢救恩的話語，接受了祢的誡命，我們寄望於祢的許諾；祢藉著祢的親臨，已賜給我們許多恩惠；但我們雖然認識祢，卻還不認識父，我們熱切渴望看見這位我們還不認識的父；因此，因為我們認識祢，但這還不夠，直到我們認識父，請把父顯示給我們，我們就滿足了。主為了使他們明白，他們以為已認識的其實並不認識，便說：\BibleQuote{我與你們同在這麼久，斐理伯，你們還不認識我嗎？誰看見了我，就是看見了父。}\BibleRef{若 14:8} 這句話中有一絲疑惑嗎？祂說：\Quote{誰看見了我，\Emph{也許}就看見了父}嗎？為什麼不？因為聽祂說話的是信者，而非信仰的迫害者；因此主不是斥責，而是教導。\BibleQuote{誰看見了我，就是看見了父}；而在這裡，\Quote{你們若認識我，就必認識我的父}，讓我們移除那表示聽者不信的詞，便成了同一句話。\par

6. 昨天，我們曾將此推薦給各位尊敬的弟兄們默想，並說過，若無可能，福音作者若望敘述他從主所聽到之事的這些話，本無需討論，除非異端者的發明強迫我們這樣做。昨天，我們簡要地向各位尊敬的弟兄們暗示，有稱為\Quote{父受苦說者}或按其創始者稱為\Quote{撒伯流派}的異端者：他們說父與子是同一位；名稱不同，但位格相同。當祂願意時，他們說，祂是父；當祂願意時，祂是子；但仍是同一位。同樣也有其他稱為\Quote{亞流派}的異端者。他們確實承認我們的主耶穌基督是父的獨生子；一位是父，一位是子；那為父的不是子，那為子的不是父；他們承認子是被生的，但否認祂的平等。我們，即大公信仰，來自宗徒的教導，植根於我們內，經由傳承的序列領受，並要完整地傳給後代——我說，大公信仰，在兩派之間，即在兩種謬誤之間，持守了真理。在撒伯流派謬誤中，祂只有一位；父與子是同一位格：在亞流派謬誤中，父與子確實是不同的位格；但子不僅是不同的位格，而且是不同的本性。你處在這兩者之間，你說什麼？你已經排除了撒伯流派，也排除了亞流派。父是父，子是子；另一位格，而非另一本性；因為\BibleQuote{我與父原是一體}，這句話我昨天已盡可能向你們強調了。當他聽到\Quote{我們是}這個詞時，讓撒伯流派羞愧地離去；當他聽到\Quote{一體}這個詞時，讓亞流派羞愧地離去。讓大公信徒在兩者之間駕馭其信仰之舟，因為在兩者中都必須謹防船難。那麼，你要說福音所說的：\BibleQuote{我與父原是一體。}本性不同，因為\Quote{一體}；非一位格，因為\Quote{是}。\par

7. 稍早前祂说：\BibleQuote{我的判断是真实的，因为我不是独自一人，而是我和那派遣我的父。}好像祂说：我的判断之所以真实，是因为我是天主子，因为我讲真理，因为我本身就是真理。那些人，按血肉理解，便说：\BibleQuote{你的父在哪里？}现在请听，亚略异端者：\BibleQuote{你们不认识我，也不认识我的父；}因为\BibleQuote{如果你们认识我，也会认识我的父。}这意味着什么，除非是\BibleQuote{我与父原是一体}？当你们看到某人像另一个人——请注意，亲爱的，这是常见的说法；不要因为你看到的习惯性而觉得困难——当我说的，你们看到某人像另一个人，而你们认识他所像的那个人，你们就惊讶地说：\Quote{这个人多么像那个人！}你们绝不会说\Emph{这样的话}，除非是两个。这里有一个不认识你们所说那个人像谁的人，他说：\Quote{他真像他吗？}你们回答他：什么，你不认识那个人？他说：\Quote{不，我不认识。}你立刻，为了通过他面前所观察的人来让他认识他所不认识的那个人，回答说：\Quote{看见这个人，你就看见了那个人。}你这样说，当然没有断言他们是同一个人，或他们不是两个，而是因为相似而作了这样的回答：\Quote{如果你认识这一个，你就认识那一个；因为他们非常相似，彼此之间毫无差别。}因此，主也说：\BibleQuote{如果你们认识我，也会认识我的父；}不是子就是父，而是子像父。愿亚略异端者脸红。感谢主，连亚略异端者也与撒伯流异端错误分离，并且不是圣父受苦论者：他不承认圣父取了血肉来到人间，圣父受苦、复活并升至祂自己那里；他不承认这一点；他与我一起承认圣父是父，圣子是子。但是，兄弟啊，你逃脱了那次船难，为什么又走向另一个？父是父，子是子；你为什么断言子与父不像，是不同的，是另一个实体？如果祂不像，祂会对祂的门徒说：\BibleQuote{看见我的，就看见了父}吗？祂会对犹太人说：\BibleQuote{如果你们认识我，也会认识我的父}吗？这怎能是真的，除非那另一个也是真的：\BibleQuote{我与父原是一体}？\par

8. \BibleQuote{耶稣在圣殿里，在银库那里说了这些话：}极大的勇气，毫无畏惧。因为祂若不愿，就不能受苦，既然祂若不愿，就不会降生。那么接下来如何？\BibleQuote{没有人拿住祂，因为祂的时候还没有到。}有些人听到这话，又相信主基督是受命运支配的，便说：\Quote{看哪，基督被命运拿住了！}唉，如果你的心不是愚昧的，你就不会相信命运。如果\OriginalTerm{命运}{fatum}，如一些人所理解的，是从\OriginalTerm{源自fari（说）的动名词}{fando}衍生出来的，即从\Quote{说话}而来，那么天主的圣言如何能被命运所束缚，而一切受造之物都在圣言本身之内呢？因为天主没有预定任何祂未曾预知的事；那被造之物原来在祂的圣言内。世界被造了；既被造而又在那里。如何既被造而又在那里？因为建筑师所建造的房屋，先前存在于他的技艺之中；在那里，有一座更好的房屋，没有年岁，没有朽坏：然而，为展示他的技艺，他建造了一座房屋；这样，在某种意义上，房屋出自房屋；如果房屋倒塌，技艺仍在。因此，所有被造之物都是藉着天主的圣言被造的；因为天主以智慧创造了万物，凡祂所造的，祂都知道：因为祂不是因为造而学习，而是因为知道才造。对我们来说，它们被认识，是因为它们被造；对祂来说，如果它们不被认识，就不会被造。因此，圣言先存。那么什么在圣言之前？什么也没有。因为如果有什么在它之前，就不会说\BibleQuote{在起初已有圣言；}而是说\Quote{在起初圣言被造}。简言之，梅瑟论及世界说了什么？\BibleQuote{在起初天主创造了天和地。}创造了原本没有的：嗯，如果祂创造了原本没有的，那么之前有什么？\BibleQuote{在起初已有圣言。}天和地从何而来？\BibleQuote{万物是藉着祂造的。}那么你把基督置于命运之下吗？命运在哪里？你说，在天上，在星辰的次序和变化中。那么命运如何能统治那创造诸天和星辰的主？而你的意志，如果你思想正确，尚且超越星辰？或者，因为你认识基督的血肉是在天下，你就认为基督的能力也被置于诸天之下吗？\par

9. 愚人哪，请听：\Quote{祂的时候还没有到}；不是祂被迫去死的时候，而是祂甘愿被处死的时候。因为祂自己知道祂何时要死：祂思量一切关于祂的预言，等待所有在祂受难之前应预言完成的事都完成，好使当一切应验之后，祂的受难才按预定的次序来临，而非出于命运的必然。简言之，请听，以便你能验证。在关于祂的其他预言中，也写着：\Quote{他们拿苦胆给我当食物，我渴了，他们拿醋给我喝。}这事如何发生，我们从福音书中得知。首先，他们给祂苦胆；祂接过来，尝了，吐了出来。此后，祂悬在十字架上，为使一切预言得以应验，祂说：\Quote{我渴了。}他们拿海绵蘸满了醋，绑在苇子上，送到祂嘴边；祂接了，说：\Quote{完成了。}那是什么意思？我死以前所预言的一切都完成了，那么我还在这里做什么呢？总之，当祂说\Quote{完成了，祂就低下头，断了气。}那些钉在祂旁边的强盗，能按自己的意愿咽气吗？他们被血肉的束缚所捆，因为他们不是血肉的创造者；被钉子固定，他们长时间痛苦，因为他们不能主宰自己的软弱。然而，主却在祂愿意的时候，在童贞女的胎中取了血肉；在祂愿意的时候来到人前；在人间住了祂愿意的时间；在祂愿意的时候离开了血肉。这是权能的一部分，而非必然。所以，祂等待的这一个时辰，不是命运所定的，而是恰当而自愿的时辰，好使在祂去世之前必须先完成的事都先完成。祂在另一处说：\BibleQuote{我有权舍掉我的生命，也有权再取回它来；没有人夺去我的生命，而是我自己舍掉它，并再取回它来。}\BibleRef{若 10:18}祂怎能受命运必然的支配呢？当犹太人寻找祂时，祂显示了这权能。祂说：\Quote{你们找谁？}他们说：\Quote{纳匝肋人耶稣。}祂回答说：\Quote{我就是。}他们一听见这话，\BibleQuote{就倒退跌在地上。}\BibleRef{若 18:6}\par

10. 有人会说：如果祂有这个能力，为什么当犹太人在十字架上侮辱祂说\Quote{如果祂是天主子，让祂从十字架上下来}时，祂没有下来，以显示祂的能力呢？因为祂是在教导我们忍耐，因此祂推迟了显示祂的能力。因为如果祂下来了，好像是因他们的话而动摇，就会被认为是被他们侮辱的刺痛击败了。祂没有下来；祂留在那里，等待祂愿意离开的时刻。因为从十字架上下来对祂算什么呢，既然祂能从坟墓中复活？那么，我们这些受此教导的人应当明白，我们主耶稣基督的能力，当时隐藏的，将在审判中显现，关于这审判经上说：\Quote{天主将显现；我们的天主，祂将不再沉默。}为什么说\Quote{将显现}？因为祂——我们的天主，即基督——曾隐藏而来，将要显现而来。\Quote{将不再沉默}：为什么说\Quote{将不再沉默}？因为起初祂确实保持沉默。什么时候？当祂受审判时；这也是为了应验先知所说的预言：\BibleQuote{祂被牵去宰杀，如同绵羊；又像羔羊在剪毛者前缄默，祂也没有开口。}\BibleRef{依 53:7}如果祂不愿意受苦，祂就不会受苦；如果祂不受苦，那血就不会流；如果那血不流，世界就不会被救赎。因此，让我们感谢祂神性的权能和祂软弱的怜悯；既因为犹太人没有认出的隐藏权能——因此现在对他们说：\Quote{你们不认识我，也不认识我的父}——也因为他们没有认出所取的血肉，却知道祂的族系：因此祂在别处对他们说：\Quote{你们认识我，也知道我是从哪里来的。}让我们在基督身上认识两者，既认识祂与父平等的方面，也认识父比祂大的方面。那便是圣言，这是血肉；那是天主，这是人；但基督却是一位，天主而人。\par

\SourceRecord{Translated by John Gibb.；Nicene and Post-Nicene Fathers, First Series；Vol. 7.；Edited by Philip Schaff.；Buffalo, NY: Christian Literature Publishing Co.,；1888.；<http://www.newadvent.org/fathers/1701037.htm>.}

% structure: p0001 source=h1 target=section reason=page-title

\section{\WorkTitle{Tractate 38 (John 8:21-25)}}

1. 今日先前宣读的圣福音的章节是这样结束的：\BibleQuote{主在银库中教导}，按祂所愿说了你们已听到的话；\BibleQuote{没有人下手拿祂，因为祂的时候还没有到。} 因此，在主的日子，我们以祂自己认为适合赐给我们的话作为我们讲论的主题。我们向你们的爱德指出了为什么经上说\BibleQuote{祂的时候还没有到}，以免任何不敬之人胆敢猜疑基督是受某种致命必然性所束缚。因为祂的时候还没有到，那时按照有关祂的预言，祂不是被迫不情愿地死去，而是预备好被杀。\par

2. 但关于祂自己的受难本身，这并非在于任何祂所受的必然性，而是在于祂自己的权能，祂在对犹太人的讲论中所说的一切就是：\BibleQuote{我去了。} 因为对主基督而言，死亡正是祂前往祂所来自之处，而祂从未离开那地方。\BibleQuote{我去了，}祂说，\BibleQuote{你们要找我，}并非出于对我的渴望，而是出于仇恨。因为当祂从人的视线中消失后，仇恨祂的人和爱祂的人都寻找祂：前者出于迫害的精神，后者出于拥有祂的愿望。在\WorkTitle{圣咏}中，主亲自借先知说：\BibleQuote{避难之处离开了我，没有人寻找我的性命。} 又在\WorkTitle{圣咏}另一处说：\BibleQuote{愿那些寻找我性命的人，蒙羞受辱。} 祂责备前者不寻找，谴责后者因他们寻找。因为不寻找基督的性命是错误的，即门徒寻找的方式；而寻找基督的性命也是错误的，即犹太人寻找的方式：前者寻求占有它，后者寻求毁灭它。因此，因为这些人心术不正，以错误的方式这样寻找，祂接着又说了什么？\BibleQuote{你们要找我，并且}——不是让你们认为你们会善意地找我——\BibleQuote{你们要死在你们的罪中。} 这就是错误地寻找基督的结果：死在罪中；这就是仇恨祂的结果，唯有通过祂才能获得救恩。因为，那些盼望天主的人本不该以恶报恶，这些人却以恶报善。因此，主预先警告他们，并在祂的先知中宣判了他们要死在他们的罪中。然后祂又加上：\BibleQuote{我所去的地方，你们不能去。} 祂在别处也对门徒说了同样的话；然而祂没有对他们说：\BibleQuote{你们要死在你们的罪中。} 但祂说了什么？和这些人一样的话：\BibleQuote{我所去的地方，你们不能去。} 祂没有除去希望，而是预告了延迟。因为当主对门徒说这话时，他们还不能去祂要去的地方，但他们后来能去；而这些人则永远不能，祂在先知中对这些人说：\BibleQuote{你们要死在你们的罪中。}\par

3. 但听到这些话，如同那些属肉体的思想者惯常的那样，他们按肉体判断，按肉体听闻和理解一切，他们说：\BibleQuote{祂要说自己要去的地方我们不能去，难道祂要自杀吗？} 愚蠢的话，充满愚昧！因为假如祂自杀，他们为什么不能去祂将去之处呢？难道他们自己不会死吗？那么，\BibleQuote{祂要说自己要去的地方我们不能去，难道祂要自杀吗？} 这话是什么意思呢？如果祂说人的死亡，哪一个人不死呢？因此，祂说\BibleQuote{我所去的地方}，不是指走向死亡，而是指祂死后自己所去之处。这便是他们的回答，因为他们不明白。\par

4. 主对这些思念属地之事的人说了什么？\BibleQuote{祂对他们说：你们是从下头来的。} 因此你们思念属地之事，因为你们像蛇一样舔土。你们吃土！这是什么意思？你们以属世之事为食，你们以属世之事为乐，你们贪求属世之事，你们没有属天之心。\BibleQuote{你们是从下头来的：我是从上头来的。你们是属这世界的：我不是属这世界的。} 因为世界是藉着祂造的，祂怎能属世界呢？所有属世界的人都是在世界之后来的，因为世界在先；所以人是属世界的。但基督在先，然后有世界；既然基督在世界之前，在基督之前什么都没有：因为\BibleQuote{在起初已有圣言，万物是藉着祂造成的。} 因此，祂是那属上头的。但属什么上头呢？属空气吗？千万不要这样想！鸟在那里飞翔。属我们看见的天空吗？我再说，千万不要这样想！星辰、太阳和月亮在那里运行。属天使吗？也不可这样理解：因为万物是藉着祂造的，天使也是祂造的。那么，基督属什么上头呢？属圣父本身。没有什么在生祂的天主之上，祂生与自己相等的圣言，与自己同永，独生，无时间，以便时间自身的根基因祂而立。那么，你当理解基督是从上头来的，以致在你的思想中超越一切受造之物——整个受造界，一切物质形体，一切受造的精神体，一切以任何方式可变化的：超越这一切，如同\Person{若望}上升，以便达到：\BibleQuote{在起初已有圣言，圣言与天主同在，圣言就是天主。}\par

5. 因此祂说：\Quote{我来自上面。你们属于这世界：我不属于这世界。所以我对你们说过，你们将死在你们的罪中。}弟兄们，祂向我们解释了祂愿意藉\Quote{你们属于这世界}所表明的意思。祂实际上说了：\Quote{你们属于这世界}，因为他们是有罪的，因为他们是不义的，因为他们是不信的，因为他们品味属地之事。你们对圣宗徒们有何看法？犹太人与宗徒之间有何区别？如同黑暗与光明之间，如同信德与不信之间，如同虔敬与不虔敬之间，如同望德与绝望之间，如同爱德与贪吝之间：区别无疑是巨大的。那么，既然有如此区别，宗徒们就不属于这世界吗？如果你们思及他们的出生方式和来源，既然他们都出于亚当，他们就属于这世界。但主亲自对他们说了什么？\Quote{我从世界中拣选了你们。}于是，那些属于世界的，变得不再属于世界，并开始属于那创造世界者。但这些犹太人仍然属于世界，对他们说了：\Quote{你们将死在你们的罪中。}\par

6. 所以，弟兄们，谁也不要自称\Quote{我不属于这世界}。无论你是谁，作为一个人，你就属于这世界；但创造世界的那位来到了你这里，把你从这世界解救出来。如果这世界令你喜悦，你便永远愿意不洁（\OriginalTerm{不洁}{immundus}）；但如果这世界不再令你喜悦，你就已经是洁净的（\OriginalTerm{洁净}{mundus}）。然而，如果由于某种软弱，这世界仍令你喜悦，就让那洁净者（\OriginalTerm{洁净}{mundat}）居住在你内，你也必洁净。但你一旦洁净，就不会继续留在世界中；也不会听到犹太人所听到的：\Quote{你们将死在你们的罪中。}因为我们都是带着罪出生的；我们一生中又加添了出生时已有的罪，从此变得比从父母出生时更属于这世界。若非那完全无罪者来临，以赎尽一切罪，我们将会在哪里？正因为犹太人不信祂，他们理当听到[这句话]：\Quote{你们将死在你们的罪中}；因为你们这些生而有罪的人，绝不可能无罪；然而，祂说，如果你们信我，虽然你们生而有罪这仍然属实，但你们不会死在你们的罪中。因此，犹太人全部的悲惨就在于：不是有罪，而是死在罪中。正是为此，每位基督徒都该力求逃脱；正是为此，我们投奔圣洗圣事；正是为此，那些因疾病或其他原因生命垂危的人急切寻求帮助；也正是为此，婴儿被母亲用虔诚的双手抱着带到教会，以免他未受圣洗就离开这个世界，死在他所出生的罪中。这些人从那些说出真理的嘴唇听到\Quote{你们将死在你们的罪中！}这话，他们的处境是何等可悲，命运又是何等不幸！\par

7. 但祂解释这事何以临到他们：\Quote{因为你们若不信我就是[那一位]，你们将死在你们的罪中。}弟兄们，我相信，在听主讲话的人群中，也有那些后来应当相信的人。然而，那极其严厉的判决仿佛是对所有人发出的：\Quote{你们将死在你们的罪中}；由此，甚至从那些应当相信的人那里，希望也被抽走了：其他人被激怒，他们则恐惧；不但恐惧，他们甚至陷入了绝望。但祂恢复了他们的希望；因为祂接着说：\Quote{你们若不信我就是，你们将死在你们的罪中。}所以，如果你们信我就是，你们就不会死在你们的罪中。希望归还给了绝望的人，沉睡者被唤醒，他们的心重新振作；此后许多人信了，正如福音书下文所证实的。因为那里有基督的肢体，尚未与基督的身体联络；而在那些把祂钉十字架、把祂挂在木头上、在祂被悬挂时嘲笑祂、用长枪刺伤祂、给祂喝苦胆和醋的人群中，有基督的肢体——为了他们，祂说：\Quote{父啊，宽赦他们吧，因为他们不知道他们所做的是什么。}连基督的血被倾流都能被宽恕，还有什么不能宽恕一个回头者呢？连基督都被杀害了，如果那杀害者也能重获希望，哪个杀人犯还该绝望？此后许多人信了；他们领受基督的血作为恩赐，为得救而饮，而非被定流祂血的罪。谁能绝望？如果那个强盗在十字架上得救——片刻前还是杀人犯，之后被指控、定罪、判决、悬挂、交付——不必惊奇。他被定罪的地方也就是他被判刑的地方；而他归化的地方也正是他被释放的地方。\BibleRef{路 18:34} 因此，在主讲这话的人群中，有那些将死在罪中的人；也有那些后来应当信祂并摆脱一切罪的人。\par

8. 且看主基督所说的：\Quote{你们若不信我就是，必要死在你们的罪恶中。}这是什么意思：\Quote{你们若不信我就是？}\Quote{我就是}什么？祂没有附加任何说明；正因为祂没有附加，祂留下了许多有待推论的东西。因为人们期待祂说祂是什么，而祂却没有说。祂应该说什么呢？也许：\Quote{你们若不信我就是}基督；\Quote{你们若不信我就是}天主子；\Quote{你们若不信我就是}圣父的圣言；\Quote{你们若不信我就是}世界的建立者；\Quote{你们若不信我就是}人的塑造者与再造者、创造者与再创造者、制造者与再制造者——\Quote{你们若不信我就是}这个，\Quote{必要死在你们的罪恶中。}祂仅仅说\Quote{我就是}，其中包含了许多含义；因为天主也曾对梅瑟说：\Quote{我是自有者；}谁能充分表达这个\Quote{我是}的意义呢？天主藉着祂的天使派遣祂的仆人梅瑟去解救祂的子民出离埃及（你们读过并知道现在所听的；但我提醒你们）；祂派遣他时，他战战兢兢，推辞自己，却又顺从。当他这样推辞时，他对天主说话，他明白天主是在天使的位格中说话：如果人民对我说：那派遣你的天主是谁？我该对他们说什么？主回答他：\Quote{我是自有者；}又加了一句：\Quote{你要对以色列子民说：那自有者派遣我到你们这里来。}那里祂也没有说：我是天主；或者：我是世界的建造者；或者：我是万物的创造者；或者：我是被解救之民的繁衍者；而只是说：\Quote{我是自有者；}以及\Quote{你要对以色列子民说：那自有者。}祂没有加上：那你们的天主、你们祖先的天主；而只是说：\Quote{那自有者派遣我到你们这里来。}也许连梅瑟自己都难以理解，正如我们也难以理解，甚至更难以理解这样的话的含义：\Quote{我是自有者；}以及\Quote{那自有者派遣我到你们这里来。}而且，就算梅瑟理解了，他受派遣去传达的那些人何时能理解呢？因此，主放下了人所不能理解的，加上了他们所能理解的；因为祂还另外说：\Quote{我是亚巴郎的天主，依撒格的天主，雅各伯的天主。}\BibleRef{出 3:13}这是你们能理解的；因为\Quote{我是自有者}，什么心智能够理解？\par

9. 那么我们呢？我们敢对这样的话说点什么吗？\Quote{我是自有者；}或者，更确切地说，关于你们所听的主说的：\Quote{你们若不信我就是，必要死在你们的罪恶中？}我敢用我这微弱且几乎不存在的能力去讨论主基督所说的\Quote{你们若不信我就是}的含义吗？我要敢于向主自己询问。请听我如同一个询问者而非讨论者，探究者而非假定者，学习者而非教导者，而你们自己也当与我一同或藉着我一同询问。主自己既无所不在，也就在近旁。愿祂听到那催人询问的情感，并赐下领悟的果实。因为即使我有所领悟，又用什么言语传递给你们我所领悟的呢？什么声音足够？什么雄辩足以胜任？什么智能的能力？什么表达的能力？\par

那么，我要向我们的主耶稣基督发言；我要发言，但愿祂乐意垂听我。我相信祂就在眼前，我对此坚信不疑；因为祂亲自说过：\BibleQuote{看，我同你们天天在一起，直到今世的终结。} \BibleRef{玛 28:20} 主，我们的天主，祢所说的\Quote{你们若不信我是}是什么意思？因为有什么不属于祢所造之物呢？难道天不属于祢？地不属于祢？天地间的一切不属于祢？连祢对之说话的人自己，不属于祢？祢所派遣的天使，不属于祢？如果这一切全是祢造的，那么祢所保留为祢独有、未赐给任何其他、使祢独自成为那样的一种存在，究竟是什么？因为我如何听到\BibleQuote{我是自有者}，好像除祢以外没有别的？又如何听到\BibleQuote{你们若不信我是}？难道那些听祂的人不存在吗？对，虽然他们是罪人，他们还是人。那么我能做什么呢？这存在究竟是什么，愿祂告诉我的内心，愿祂告诉我，愿祂在我里面宣示；愿内心的人听见，愿理智领会这真实的存在；因为这样的存在永远是不变的。因为任何事物，无论其卓越如何，如果它是可变的，就不是真实的存在；因为凡是有不存在之处，就没有真实的存在。因为凡是可变的，当其变化时，就不再是原来的样子：如果它不再是原来的样子，一种死亡就在其中发生了；原来存在的东西被消灭了，不再存在。白发老者的黑发消失，愁苦佝偻老人身体中的俊美消失，虚弱者身体中的力量消失，行走者身体中先前站立的姿态消失，站立者身体中的行走消失，躺卧者身体中的行走和站立消失，沉默者舌头中的言语消失——凡是变化而不再是原来样子的，我在那里看到一种存在于现在之中的生命，和一种存在于过去之中的死亡。最后，当我们说某人死了，那人哪里去了？我们回答说：他曾经在了。啊，真理，唯有你真正地存在！因为在我们的一切行动和运动中，是的，在受造物的每一种活动中，我都发现两个时间：过去和未来。我寻找现在，没有东西是静止的：我所说的已不再是现在；我将要说的尚未来到；我所做的已不再是现在；我将要做的尚未来到；我已活过的生命已不再是现在；我还将活的生命尚未来到。我在每一个受造物的运动中都发现过去和未来：在永恒的真理中，我却找不到过去和未来，只有现在，而且是不变的现在，这在受造物中是没有位置的。审察事物的变化，你将会发现\InnerQuote{曾是}和\InnerQuote{将是}；思念天主，你将会发现\InnerQuote{是}，其中\InnerQuote{曾是}和\InnerQuote{将是}不能存在。那么，为使你自己也这样存在，要超越时间的界限。但谁能超越自己存有的能力呢？愿那曾对父说\BibleQuote{我愿他们在那里，同我在一起}的那位提拔我们到那里。因此，在作出这个许诺，叫我们不至于死在我们的罪中时，主耶稣基督，我想，祂用这些话没有说别的，只是说：\BibleQuote{你们若不信我是}；是的，我想祂用这些话无非要表明这一点：\BibleQuote{你们若不信我是}天主，\BibleQuote{你们必要死在你们的罪恶中}。感谢天主，祂说了\BibleQuote{你们若不信}，而没有说\Quote{你们若不明白}。因为谁能明白这事呢？还是说，既然我大胆发言，而你们似乎明白了，你们确实对如此不可言说的话题有所领悟？如果你们不明白，信德就释放你们。因此，主也没有说\Quote{你们若不明白我是}，而是说了他们所能达到的：\BibleQuote{你们若不信我是，你们必要死在你们的罪恶中}。\par

11. 这些总在品味属地之事的人，总是按肉体听取和回答，他们对祂说了什么？\Quote{你是谁？}因为当祢说：\Quote{你们若不相信我就是那一位，}祢并没有告诉我们祢是谁。祢是谁，好让我们相信？祂回答说：\Quote{元始。}这里就是那[永远]存在的存在。元始不能改变：元始是自存而万物之源；也就是说，元始，关于它曾说过：\Quote{但是你自身是同样的，你的年岁不会穷尽。}祂说：\Quote{元始，因为如此我也对你们说话。}你们要相信我[是]元始，好使你们不死在你们的罪恶中。就如当他们藉着说\Quote{你是谁？}时，无非是说：我们该相信祢是什么？祂回答说：\Quote{元始；}也就是说，相信我[是]\Quote{元始}。因为在希腊文表达中，我们分辨出在拉丁文中无法分辨的事。因为在希腊文中，\Quote{元始}（\Emph{principium}, \OriginalTerm{元始}{\GreekText{ἀρχή}}）是阴性，正如在我们语言中\Quote{法律}（\Emph{lex}）是阴性，而在他们的语言中却是阳性（\OriginalTerm{法律}{\GreekText{νόμος}}）；又或者如\Quote{智慧}（\Emph{sapientia}, \OriginalTerm{智慧}{\GreekText{σοφία}}）在两种语言中都是阴性。所以，在不同语言中，词的性别变化是言语的习惯，因为在事物本身没有性别的区分。因为智慧并非真是女性，既然基督是天主的智慧，\BibleRef{格前 1:24}而基督是阳性词，智慧是阴性词。因此，当犹太人说：\Quote{你是谁？}祂知道那里有一些将要相信的人，因此说过：你是谁？好让他们能认识他们该相信祂是什么，祂回答说：\Quote{元始：}不是好像祂说：我是元始；而是好像祂说：相信我[是]元始。正如我所说的，这在希腊文中很明显，在那里元始（\OriginalTerm{元始}{\GreekText{ἀρχή}}）是阴性。就如祂本想说祂是真理，而对他们的问题\Quote{你是谁？}祂回答了\Emph{Veritatem}（真理）；而当祂听到\Quote{你是谁？}时，祂显然应当回答\Emph{Veritas}（真理）；即我是真理。但祂的回答有更深的意义，当祂看见他们提出\Quote{你是谁？}这个问题时，意思是：从祢那里听到\Quote{你们若不相信我就是那一位，}我们该相信祢是什么？对此祂回答说：\Quote{元始：}就好像祂说：相信我[是]元始。祂又说：\Quote{因为[如此]我也对你们说话；}也就是说，为了你们，我谦卑了自己，屈尊到这样的话语。因为如果元始本身与圣父同在，而不接受仆人的形象，并不以人的方式与人交谈，那么他们怎能相信祂呢？因为他们软弱的心智若不藉某种诉诸感官的声音，就不能明智地听到圣言。因此，祂说：相信我[是]元始；为使你们相信，我不但存在，而且也对你们说话。关于这个主题，我还有许多话要告诉你们；愿你们的爱德允许我们保留余下部分，明天藉祂的恩佑再讲。\par

\SourceRecord{Translated by John Gibb.；Nicene and Post-Nicene Fathers, First Series；Vol. 7.；Edited by Philip Schaff.；Buffalo, NY: Christian Literature Publishing Co.,；1888.；<http://www.newadvent.org/fathers/1701038.htm>.}

% structure: p0001 source=h1 target=section reason=page-title

\section{讲道集第39篇（\BibleRef{若 8:26-27}）}

1. 我们的主耶稣基督向犹太人所说的话，祂如此调整祂的言论，以致盲者看不见，而信者的眼睛得以开启，这些就是今天从圣福音中所宣读的：\BibleQuote{那时犹太人说：\InnerQuote{祢是谁？}}因为主先前说过：\BibleQuote{如果你们不信我就是，你们必要死在你们的罪恶中。}对此他们随即回应说：\BibleQuote{祢是谁？}好像想要知道他们应当信靠谁，以免死在他们的罪恶中。祂回答那些问祂\BibleQuote{祢是谁？}的人，说道：\BibleQuote{元始，因为〔我也〕对你们说话。}如果主称自己为元始，那么可以问：圣父是否也是元始？因为如果有一位父的子是元始，那么更容易理解天主圣父是元始，祂确实是子之父，但祂自己并非从谁而来。因为子是圣父的子，圣父当然是圣子的父；但子被称为出自天主的天主——子被称为出自光的光；圣父被称为光，但不是出自光——圣父被称为天主，但不是出自天主。如果出自天主的天主、出自光的光是元始，那么更容易理解那光（光从之而来）和那神（神从之而来）也是元始。因此，亲爱的诸位，称子为元始而不也称圣父为元始，似乎是不合理的。\par

2. 但我们该怎么办？那么，有两个元始吗？让我们谨慎不说这样的话。那么，如果圣父是元始，子也是元始，怎么不是两个元始呢？同样，我们称圣父为天主，称子为天主，却不说是两个天主；然而，那圣父不是圣子，那圣子不是圣父；而圣神、圣父和圣子的神，既不是圣父也不是圣子。因此，如同公教信友在慈母教会的怀抱中所受的教导，那圣父不是圣子，那圣子不是圣父，而圣神、圣父和圣子的神，既不是圣子也不是圣父，然而我们不说有三位天主；虽然，如果我们分别被问及每一位，我们无论被问及哪一位，都必须宣认祂是天主。\par

3. 但这一切在那些将熟悉的事物与不太熟悉的事物、可见的事物与不可见的事物相提并论，并将受造物与造物主比较的人看来，似乎荒谬。因为有时不信者质问我们说：你们称为圣父的那位，你们称祂为天主吗？我们回答：天主。你们称为圣子的那位，你们称祂为天主吗？我们回答：天主。你们称为圣神的那位，你们称祂为天主吗？我们回答：天主。那么，他们说，圣父、圣子和圣神是三位天主吗？我们回答：不是。他们困惑了，因为他们未受光照；他们封闭了心，因为他们缺少信德的钥匙。所以，弟兄们，让我们藉着先行的信德（这信德医治我们心灵的眼睛），毫无晦涩地接受我们所理解的——对于我们所不理解的，则毫无犹豫地相信；让我们不要为达到完善的顶峰而放弃信德的基础。圣父是天主，圣子是天主，圣神是天主：然而那圣父不是圣子，那圣子不是圣父，而圣神、圣父和圣子的神，既不是圣父也不是圣子。天主圣三是一个天主。天主圣三是一个永恒、一个能力、一个威严——三位，〔但不是〕三位天主。愿那诘难者不回答我：\Quote{那么，三位什么？因为，}他补充说，\Quote{如果有三位，你必须说，三位什么？}我回答：圣父、圣子和圣神。\Quote{瞧，}他说，\Quote{你已提到了三位；但请说明这三位是什么？}不，你自己数一数；因为当我提到圣父、圣子和圣神时，我算出了三位。因为圣父就其本身而言是天主，但就圣子而言，祂是父；圣子就其本身而言是天主，但就圣父而言，祂是子。\par

4. 我所说的，你可以从日常的比拟中体会。就拿一个人和另一个人来说，如果一个是父亲，另一个是儿子。就他自身而言，他是人；但就他与儿子的关系而言，他是父亲。儿子就他自身而言是人，但就他与父亲的关系而言是儿子。因为\Quote{父亲}是一个相对的名称，同样，\Quote{儿子}也是相对的；但他们是两个人。当然，天主圣父是相对意义上的父亲，也就是相对于圣子；天主圣子是相对意义上的儿子，也就是相对于圣父；但不像前者是两个人那样，后者是两位神。为什么此处不是这样呢？因为前者属于一个领域，后者属于另一个领域；因为这是神圣的。这里有些不可言喻的东西，无法用言语解释，即应当有数目，又不应当有数目。因为请看，是否显出一种数目：圣父、圣子和圣神——天主圣三？如果说是三位，那么三位什么？这里数目就失效了。因此，天主既不脱离数目，也不被数目所包含。因为有三个，所以有某种数目。如果你问三个什么，数目就停止了。因此经上说：\Quote{我们的主伟大，力量浩大；祂的智慧无法数算。} 当你开始思考时，你就开始数算；当你数算时，你却说不出的你所数算的是什么。圣父是圣父，圣子是圣子，圣神是圣神。这三位，圣父、圣子、圣神，是什么？难道是三位神吗？不。难道是三位全能者吗？不。难道是三位创造世界的吗？不。那么，圣父是全能者吗？显然是全能者。那么圣子不是全能者吗？显然圣子也是全能者。那么圣神不是全能者吗？祂也是全能者。那么，有三位全能者吗？不；只有一位全能者。只有在彼此之间的关系上，它们才显示出数目，而不是在它们的主体存在上。因为虽然天主圣父就祂自身而言，同圣子和圣神一起是天主，但并不是三位神；虽然就祂自身而言，祂是全能者，圣子和圣神也是，但并不是三位全能者；因为事实上，祂不是就自身而言是圣父，而是就圣子而言；圣子也不是就自身而言是圣子，而是就圣父而言；圣神也不是就自身而言是圣神，就祂被称为圣父和圣神的圣神而言。除了圣父、圣子和圣神，一位天主、一位全能者，我无法给出这三位的名称。因此，是一个元始。\par

5. 从圣经取一个比方，好让你们在某种程度上明白我所说的。我们的主耶稣基督复活后，乐意升天，十天后从那里派遣了圣神，使聚集在同一房间的人被充满，开始说万国的语言。杀害主的人被这个奇迹吓住，心中刺痛而忧伤；忧伤而改变；改变而相信。有三千人加入了主的身體，即信徒的数目。借着另一个奇迹，又有五千人加入。形成了一个相当大的团体，在其中，所有领受圣神的人，被属神的爱火点燃，因着他们的爱和神火的热情被熔为一体，并在共融的合一中开始变卖所有的一切，把钱放在宗徒脚前，按照各人的需要分给每个人。圣经这样论到他们说：\Quote{他们一心一意归向天主。} 弟兄们，请留意，并从此认识天主圣三的奥秘：我们如何说，既有圣父、圣子和圣神，却只有一位天主。看！这里有成千上万的人，却只有一颗心；成千上万的人，却只有一个灵魂。但在哪里？在天主内。更何况天主本身呢？当我说两个人是兩個灵魂，或三个人是三个灵魂，或许多人是许多灵魂时，我在措辞上有错吗？我确实说得正确。让他们亲近天主，所有的人都变成一个灵魂。如果借着亲近天主，许多灵魂因爱而成为一个灵魂，许多心成为一颗心，那么，圣父和圣子里那爱的泉源本身，岂不更是如此，以至于圣三是一位天主吗？因为从那里，从那圣神，爱来到我们中间，正如宗徒所说：\BibleQuote{天主的爱借着所赐与我们的圣神，已倾注在我们心中。} \BibleRef{罗 5:5} 如果天主的爱借着所赐与我们的圣神倾注在我们心中，使许多灵魂成为一个灵魂，许多心成为一颗心，那么圣父、圣子和圣神岂不更是唯一的天主，唯一的光明，唯一的元始吗？\par

6. 那么，让我们听听那向我们说话的元始：\Quote{我有}，祂说，\Quote{许多关于你们的事要说，也要审判。} 你们记得他曾说：\Quote{我不审判任何人。}\BibleRef{若 8:15} 看，现在祂说：\Quote{我有许多关于你们的事要说，也要审判。} 但\Quote{我不审判}是一回事，\Quote{我必须审判}是另一回事；因为祂来是为拯救世界，而不是为审判世界。当祂说\Quote{我有许多关于你们的事要说，也要审判}时，祂指的是将来的审判。因为祂升天，正是为了要来审判生者死者。没有人比那位曾被不公义审判过的人更公正地审判。\Quote{许多事}，祂说，\Quote{我有关于你们的事要说，也要审判；但那派遣我的是真实的。} 看看圣子，祂与圣父同等，如何归光荣于圣父。因为祂给我们立了榜样，仿佛在我们心中说：啊，信徒，如果你听我的福音，你的主天主对你说：当我在起初，圣言与天主同在，与圣父同等，与生祂者同等永恒，归光荣于祂，我是祂的儿子，你怎能在那位你作祂僕人的面前骄傲呢？\par

7. \Quote{我有许多事}，祂说，\Quote{要论及你们、判断你们；但那派遣我来的，是真实的；}好像祂说：因此我判断真理，因为，作为真实者的儿子，我就是真理。圣父是真实的，圣子是真理——我们认为哪一个更大？让我们反省，若能，真实者和真理哪一个更大。试举其他例子。一个虔诚的人，还是虔诚之德，哪一个更全面？当然是虔诚之德本身；因为虔诚的人源自虔诚之德，而非虔诚之德源自虔诚的人。因为即使曾经虔诚的人变得不虔，虔诚之德仍然可以存在。他已失去他的虔诚，却未曾从虔诚之德本身取走什么。俊美的人和俊美本身又如何？俊美本身大于俊美的人；因为俊美本身使俊美的人存在，而非俊美的人使俊美本身存在。贞洁的人与贞洁本身也是如此。贞洁本身显然比贞洁的人更大。因为若贞洁本身不存在，人就没有理由贞洁；但即使一个人拒绝贞洁，贞洁本身仍保持不变。既然如此，若虔诚之德一词比虔诚的人含义更多，俊美本身比俊美的人更多，贞洁本身比贞洁的人更多，我们能否说真理比真实者更多？若我们这样说，我们便开始说圣子大于圣父。因为主自己说得最清楚：\BibleQuote{我是道路、真理和生命} \BibleRef{若 14:6}。因此，若圣子是真理，那么，除真理本身所说\Quote{那派遣我来的，是真实的}之外，圣父是什么？圣子是真理，圣父是真实的。我追问哪一个更大，却找到了平等。因为真实的圣父之所以真实，不是因祂包含了真理的一部分，而是因祂生出了整个真理。\par

8. 我看我必须说得更明了。并且，为了不耽误你们太久，今天我只谈论这一点。当我靠天主助佑，说完我想说的话，我的讲论就结束。我说这话，是为了吸引你们的注意。每个灵魂，作为受造物，是可变的；且虽然是伟大的受造物，却仍是受造物；虽优于身体，却仍是受造的。每个灵魂，既是可变的——即有时相信，有时不信；有时愿意，有时拒绝；有时犯奸淫，有时贞洁；此时善，彼时恶——就是可变的。但天主是自有者，因而保留了自己的特殊名字：\BibleQuote{我是自有者} \BibleRef{出 3:14}。圣子亦是如此，当祂说：\Quote{你们若不相信我就是那一位}；与此相关的还有：\BibleQuote{你是谁？——我是元始} \BibleRef{若 8:25}。因此天主是不变的，灵魂是可变的。当灵魂从天主接收其美善的要素时，它因参与而成为美善，正如你的眼睛因参与（光线）而看见。因为光一撤去，它便看不见；而它只要有份于光，就能看见。既然灵魂因参与而成为美善，那么如果\Emph{它}改变而变为邪恶，那使之美善的美善仍保持不变。因为那美善曾在其为善时被它分享；而当它转向邪恶时，那美善仍完整不变。如果灵魂堕落成为邪恶，那美善并未减少；如果它回来成为美善，那美善并未增大。你的眼睛分享这光，你便看见。闭上眼睛？你并未减少光。睁开眼睛？你并未增加光。通过这个例子，弟兄们，要明白如果灵魂是虔诚的，便有天主那里的虔诚之德，灵魂是分享者；如果灵魂是贞洁的，便有天主那里的贞洁本身，灵魂分享它；如果它是美善的，便有天主那里的美善本身，灵魂分享它；如果它是真实的，便有天主那里的真理本身，灵魂是分享者。如果灵魂没有分享这些，那么每个人都是虚假的；如果每个人都可能虚假，那么没有一人是凭自己真实的。但真实的圣父凭自己就是真实的，因为祂生了真理。说那个人是真实的，因为他接受了真理，这是一回事；说天主是真实的，因为祂生了真理，是另一回事。看，天主的真实是这样——不是通过参与，而是通过生出真理。我看你们已明白了，我很高兴。今天你们就以此满足吧。其余部分，照祂所赐，当主愿意时，我们将阐明。\par

\SourceRecord{Translated by John Gibb.；Nicene and Post-Nicene Fathers, First Series；Vol. 7.；Edited by Philip Schaff.；Buffalo, NY: Christian Literature Publishing Co.,；1888.；<http://www.newadvent.org/fathers/1701039.htm>.}

% structure: p0001 source=h1 target=section reason=page-title

\section{论道第40篇（\BibleBook{若望}{福音} 8:28-32）}

1. 关于我们手中所读的圣若望福音，你们的爱德已听过许多，并且赖天主的恩宠，我们已按自己的能力讲解了，提醒你们注意：这位圣史特别选择了讲论主的天主性，在其中祂与圣父同等，是天主的独生子；因此他被比作鹰，因为没有别的鸟被认为飞得更高。因此，按着顺序，当主使我们能够讲解接下来时，请你们全神贯注地聆听。\par

2. 关于前面的段落，我们已经向你们讲过，提示了如何理解圣父是真实的，而圣子是真理。但当主耶稣说：\Quote{派遣我者是真实的，}犹太人却不明白祂是对他们讲论圣父。于是祂对他们说，正如你们刚才在诵读中所听见的：\Quote{当你们高举了人子以后，你们便要知道我是，并且我什么都不由自己作；而是照父所教导我的，我说这些话。}这是什么意思？因为看来祂所说的一切，仿佛是说他们在祂受难之后才会知道祂是谁。因此，毫无疑问，祂看见了那里的一些人，就是祂自己所认识的，并与祂其余的圣人们一起，在创世以前凭祂的先知预知所拣选的，将在祂受难后相信。这些人正是我们时常推荐、并恳切提出供你们效法的。因为在主受难、复活、升天之后，圣神降临时，因着那位曾被逼迫的犹太人视为死人而蔑视的主的名，行了许多奇迹，他们心中刺痛；那些愤怒地杀害祂的人改变了，相信了；那些愤怒地流祂血的人，如今以信德的精神饮了祂的血；也就是那三千和五千犹太人，就是祂现在在那里看见的，当祂说：\Quote{当你们高举了人子以后，你们便要知道我是[那一位]。}如同祂说：我让你们对我的认识延迟，直到我完成我的苦难；你们将在自己的次序上知道我是谁。这并不是说所有听见祂的人都只能在主受难之后才相信；因为稍后经上说：\Quote{祂说这些话时，许多人信了祂；}那时人子还没有被高举。但祂所说的高举是指祂的苦难，而不是祂的光荣；是指十字架，而不是天堂；因为祂在悬于木头上时也被高举了。但那高举是祂的贬抑，因为那时祂服从至死，且死在十字架上。\BibleRef{斐 2:8}这需要由那些后来将要相信的人的手来完成，祂对他们说：\Quote{当你们高举了人子以后，你们便要知道我是[那一位]。}为什么这样呢？无非是叫没有人绝望，不管他的良心如何有罪，当他看见那些杀害基督的人得赦免了杀人之罪时。\par

3. 主因而认出了那人群中的这样的人，就说：\Quote{当你们高举了人子以后，你们便要知道我是。}你们已经知道\Quote{我是}是什么意思；我们不必一再重复，免得如此重大的主题引起厌倦。请回想：\Quote{我是自有者，}以及\Quote{那自有者派遣了我，}\BibleRef{出 3:14}，你们就会认识\Quote{那时你们便要知道我是}这句话的意义。但圣父也是，圣神也是。整个天主圣三都属于同一本体。但因为主是作为圣子说话，为使当祂说\Quote{那时你们便要知道我是}时，没有机会让撒伯里乌派——即圣父受难派——的错误渗入，这个错误我曾告诫你们不要持守，而要提防——我说的是那些说圣父和圣子是同一位的错误；两个名字，一个实体——为防范那个错误，当主说\Quote{那时你们便要知道我是}时，为了不使祂被理解为圣父自己，祂立即加上：\Quote{并且我什么都不由自己作；而是照父所教导我的，我说这些话。}撒伯里乌派已经开始为发现他错误的依据而欣喜；但正当他好像在阴凉处显现时，他立刻被接着的句子之光所击溃。你们以为祂是圣父，因为祂说\Quote{我是}。现在听清楚祂是圣子：\Quote{并且我什么都不由自己作。}这是什么意思：\Quote{我什么都不由自己作}？由我自己我是不存在的。因为圣子是天主，来自圣父；但圣父是天主，却不是来自圣子。圣子是出自天主的天主，圣父是天主，却不是出自天主。圣子是出自光的光，圣父是光，却不是出自光。圣子是存在的，但有祂由之而存在的那一位；圣父是存在的，却没有祂由之而存在的那一位。\par

4. 那么，我的弟兄们，不要让祂进一步的话——\Quote{就如父所教训我的，我就讲论这些事}——成为任何属肉体的思想潜入你们心中的机会。因为人性的软弱不能思想，除非依照它所习惯的作为与听闻。所以，不要将如同两个人（一个是父亲，另一个是儿子，父亲对儿子说话）的景象摆在眼前，如同你们中任何人可能做的那样：当你们对自己的儿子说什么时，告诫并教导他如何说话，将你们所告诉他的存入他的记忆，然后他将它用言语表达出来，清晰地宣述，并将他所领会的东西传递给别人的耳朵。不要这样想，免得你们在内心制造偶像。人的形状、四肢的轮廓、肉身的形态、外在的感官、身体的高矮与动作、舌头的功能、声音的区别——不要认为这些存在于那\OriginalTerm{天主圣三}{Trinity}中，除非它们涉及那独生子所取的\OriginalTerm{仆人之形}{servant-form}，那时圣言成了血肉，居住在我们中间。\newline{}关于那仆人之形，人性的软弱啊，我不禁止你按照你的知识去思想；我反而要求你如此。如果你内所有的信德是真实的，就如此思想基督；但如此思想祂是出于童贞玛利亚，而非出于天主圣父。祂曾是婴孩，祂长大成人，祂行走如人，祂饥饿，祂口渴如人，祂睡眠如人；最后祂受苦如人，悬在木头上，被杀并被埋葬如人。以同样的形体祂复活了；以同样的形体，在门徒眼前，祂升了天；以同样的形体祂还要来审判。因为天使的口已在福音中宣告：\BibleQuote{祂要怎样来，你们就看见祂怎样升了天。} \BibleRef{宗 1:11} 因此，当你们思想基督内的仆人之形时，要想到人的模样，如果你们有信德；但当你们思想\Quote{在起初已有圣言，圣言与天主同在，圣言就是天主}时，要从你们心中除去一切人性的塑造。从你们的思想中驱除所有被形躯界限所束缚、被地方度量所包含、或铺展成一块团块的东西，无论其体积多大。让这样的幻像从你们心中完全消灭。如果你们能够，请思想智慧之美，给自己描绘正义之美。那有形状吗？有大小吗？有颜色吗？它毫无这些，然而它存在；因为如果它不存在，它就不会被爱，也不值得赞美，也不会作为可敬可爱的对象被珍藏在我们心中和生活中。但人在这里成为有智慧的；如果智慧不存在，他们从哪里成为这样的呢？此外，人啊，如果你不能用肉体的眼睛看见你自己的智慧，也不能用你用来思想有形之物的同一心智图象来思想它，你竟敢将一个人身体的形状强加于天主的智慧吗？\par

5. 那么，弟兄们，我们该说什么呢？既然圣子说：\Quote{如同圣父教导了我，我就说这些事}，圣父是如何向圣子说话的？圣父教导圣子时，是否像你教导你的儿子那样，使用了言语？祂怎能对圣言使用言语！众多言语怎能用于那唯一的圣言？圣父的圣言曾靠近祂的耳朵去听圣父的口吗？这些都是属肉体的思虑，从你们心中驱除吧。我说这话，只要你们理解了我的话，我确实说了话，我的言语发出了声音，通过声音传到你们的耳朵，并通过你们的听觉将意义传递到你们的心智，如果你们真的理解了的话。假设一个说拉丁语的人听到了，但只是听到而不理解我所说的。就从我口中发出的声音而言，那不理解的人也和我一样分享了它。他听到了那声音，同样的音节敲击了他的耳朵，却没有在他的心智中产生效果。为什么？因为他不理解。但如果你们理解了，你们的理解从何而来？我的言语在耳中作响：我是否在你们心中点燃了光明？毫无疑问，如果我所说的真实，而这真理你们不仅听到了，而且理解了，那么就有两件事发生了（请区分它们）：听闻和理解。听闻是我做的，但理解是谁做的？我对耳朵说话，使你们听见；谁对你们的心说话，使你们理解？无疑，也有人对你的心说了话，使言语的声音不仅仅敲击耳朵，而且真理的某种东西也降入了你的心。也有人向你的心说了话，但你看不见他。弟兄们，如果你们理解了，那么你们的心也受到了话语。理解是天主的恩赐。如果你们理解了，是谁在你们心中这样说话，不就是那位，诗篇向祂说：\Quote{求赐我明达，好学习祢的诫命}的祂吗？例如，主教说了话。有人问：\Quote{他说了什么？}你重复他所说的，并补充说：\Quote{他讲了真理。}而另一个没有理解的人说：\Quote{他说了什么，或者你在赞扬什么？}两人都听到了我；我对两人都说了话；但天主只对其中一人说了话。如果我们可以拿小事与大事相比（因为我们之于祂算什么？），天主在我们内运作某种非形体且属神的东西，它既不是声音以敲击耳朵，也不是颜色以被眼睛辨别，也不是气味以进入鼻孔，也不是滋味以被口舌评判，也不是任何软硬之物以被触觉感知；然而有某种东西，它容易感受——却难以解释。那么，如果天主，如我所说，在我们心中无声地说话，祂又怎样向祂的圣子说话呢？因此，弟兄们，要这样想，尽你们所能，如果，如我所说，我们可以在某种程度上以小事比大事：要这样想。圣父以非形体的方式向圣子说话，因为圣父以非形体的方式生了圣子。祂并非像生了未受教的圣子那样教导祂；而是教导祂就等同于生了充满知识的祂；而这\Quote{圣父教导了我}等同于\Quote{圣父生了我，使我早已知道}。因为，如果少数人理解的那样，真理的本性是单纯的，那么对圣子而言，存在与知道是同一回事。因此，祂从圣父那里获得了知识，也从圣父那里获得了存在。并非祂先有存在，然后才有知识；而是在生祂时，圣父给了祂存在，也在生祂时，圣父给了祂知识；因为，如前所述，对于真理的单纯本性，存在不是一回事，知道不是另一回事，而是同一回事。\par

6. 于是祂对犹太人说，并且补充道：\Quote{那派遣我的，与我同在。}祂之前也说过这话，但关于这个要点祂不断提醒他们——\Quote{祂派遣了我}和\Quote{祂与我同在}。那么，主啊，如果祂与祢同在，与其说一位被另一位派遣，不如说你们两位都来了。然而，当两位一起时，一位被派遣，另一位是派遣者；因为降生成人就是一种派遣，而这降生成人只属于圣子，不属于圣父。因此，圣父派遣了圣子，但并没有离开圣子。因为并非圣父离开了祂派遣圣子去的地方。因为哪里没有万物的创造者呢？祂说：\Quote{我充满了天地}，哪里没有祂呢？\BibleRef{耶 23:24} 但也许圣父无所不在，而圣子却不是这样？听福音作者的话：\Quote{祂在世界之中，世界也是借着祂造成的。}因此祂说：\Quote{那派遣我的}，即藉着祂作为圣父的权能我成了血肉，\Quote{与我同在——没有离弃我。}为什么祂没有离弃我？祂说：\Quote{祂没有离弃我独自一人，因为我常做祂所喜悦的事。}那平等性\Emph{常}在；不是从某个起始点然后继续；而是无始无终。因为神圣的生成没有时间上的开始，因为时间本身是由独生子创造的。\par

7. \Quote{祂说这些话的时候，有许多人信了祂。}但愿在我说话的时候，也有许多以前持不同意见的人理解并信了祂！因为也许在这个大集会中有一些亚略派。我不敢怀疑有撒伯里乌斯派，他们说圣父与圣子同是一位，因为那个异端太古老了，并且已逐渐被淘汰。但亚略派似乎还有一些活动，像一具腐烂的尸体，或者最多像一个人临终前的气息；因此仍然需要有些人从那里得救，就像许多人从另一个异端中得救一样。的确，这个省份过去没有这些人；但自从许多外国人到来后，其中一些也来到了我们的附近。看，当主说这些话时，许多犹太人信了祂。愿我也看到，在我说话时，亚略派在相信——不是相信我，而是与我一起相信！\par

8. \Quote{那时，主对那些信祂的犹太人说：你们若存留在我的话内。} \Quote{存留，}我说，因为你们已开始入门，开始在那里。\Quote{你们若存留，}就是说，在你们这些信者身上已经开始的那信德中，你们将达到什么？请看这开始的性质，以及它引向何处。你们喜爱了根基，请注意巅峰，并从那低微的境况中寻求另一种崇高。因为信德有谦卑，但认识、不死与永生却没有卑微，而是崇高；也就是说，提升、圆满、永恒的稳固、完全免于仇敌的攻击和对失败的恐惧。那始于信德的伟大，却被人轻视。在建筑中，根基通常被无知者看作微不足道。挖一个大坑，石头被随意抛入各处。那里没有装饰，没有美丽；就像树的根没有美丽的外表。然而，树中一切令你喜悦的都来自那根。你看那根，不感到喜悦；你看树，却赞叹它。愚昧的人啊！你所赞叹的正是从那未令你喜悦的根生长出来的。信者的信德看似微不足道——你没有秤来衡量它。那么，请听它达到什么，看看它的伟大：正如主自己在另一处说：\Quote{如果你们有信德像芥菜子一样。} \BibleRef{玛 17:20} 有什么比这更微不足道？又有什么比这充满更大的能量？什么更微小，却更热切地扩展？同样，祂说：\Quote{你们}也要，\Quote{如果你们存留在我的话内，}即你们所相信的话中，你们将被引向何处？\Quote{你们将确实是我的门徒。} 这为我们有什么益处？\Quote{你们将认识真理。}\par

9. 弟兄们，主向信者应许了什么？\Quote{你们将认识真理。} 为什么这样说？难道他们在主说话时还不曾认识吗？如果他们不认识，怎么相信了？他们相信，不是因为他们认识，而是为了能认识。因为我们相信是为了能认识，我们不是认识后才相信。因为我们还将认识的，是眼所未见，耳所未闻，人心所未想过的。信德是什么，不就是相信你所未见到的吗？信德就是相信你所未见到的；真理，就是见到你所相信的，正如祂自己在某处所说。主在地上行走，首先是为了缔造信德。祂是人，祂处于卑微的境况。祂被众人看见，但并非所有看见祂的人都认识祂。许多人弃绝祂，大众杀害祂，少数人哀悼祂；然而就连那些哀悼祂的人，也仍未认识祂的真体。这一切可说是信德轮廓与未来建造的开始。正如主在某处论及此事说：\Quote{那爱我的人，必遵守我的命令；那爱我的人，必蒙我父所爱，我也要爱他，并将我自己显示给他。} 他们当然已经看见他们正在听的那位；然而对他们，如果他们爱祂，祂应许他们将要看见祂。同样在这里：\Quote{你们将认识真理。} 怎么会这样？难道你所说的不就是真理吗？确实是真理，但此刻只是被相信，而非被看见。如果你们存留在所相信的，你们将达到所见的。因此，若望本人，这位圣史，在他的书信中说：\Quote{亲爱的，我们已是天主的子女；但将来如何，还未显明。} 我们已经是了，而我们还将成为别的。我们将会比现在更多什么？请听：\Quote{将来如何，还未显明；[但]我们知道，当祂显现时，我们将像祂。} 如何？\Quote{因为我们必要看见祂，正如祂所是的。} \BibleRef{若一 3:2} 一个伟大的应许，却是信德的赏报。你寻求赏报，就让工作在先。如果你相信，你就寻求信德的赏报；但如果你不信，你怎能厚颜无耻地寻求信德的赏报？因此，\Quote{如果}你们\Quote{存留在我的话内，你们将确实是我的门徒，}以便你们看见真理本身，不是通过声音的话语，而是在耀眼的光芒中，祂必以这光满足我们：正如我们在圣咏中读到：\Quote{你面容的光辉印在我们身上。} 我们是天主的钱币：我们像钱币从宝库中迷失了。那印在我们身上的肖像因我们的游荡而被磨损了。祂来重塑我们，因为祂起初就是塑造我们的那一位；祂自己来索求祂的钱币，正如凯撒索求他的。因此祂说：\Quote{凯撒的归凯撒，天主的归天主：} \BibleRef{玛 22:21} 给凯撒他的钱币，给天主你们自己。那时，那真理将在我们身上再现。\par

10. 我要对你们的爱德说些什么呢？但愿我们的心多少能向往那不可言喻的光荣！但愿我们在旅途中不断叹息，不爱世界，并以虔诚的心灵不断向着召叫我们的那一位前进！渴望是心灵的内室。如果我们竭尽全力向渴望屈服，我们必将获得。为此，圣经、民众的聚会、圣事的举行、圣洗圣事、赞美天主的歌唱，以及我们此刻的讲解，其目的正是：使这渴望不仅被植入、发芽，而且扩展至足以容纳\Quote{眼所未见，耳所未闻，人心所未想到的}那种容量。但请与我一同爱慕。爱天主的人不会过分爱钱财。我只是略略触及这软弱，不敢说\Quote{他一点也不爱钱财}，而是说\Quote{他不过分爱钱财}；好像钱财是可爱的，但不可过分。啊，若我们相称地爱天主，我们根本就不会爱钱财！那么，钱财将成为你旅途中的工具，而非情欲的刺激；是为了需要而使用，而非作为享乐的原因而喜悦。爱天主吧，如果祂在你内成就了你所听到并称赞的某些事。使用世界，不要让世界俘虏你。你正在经过你已开始的旅程；你已经来了，又要离开，并非停留。你正在经过你的旅程，这生命不过是一个路边的客栈。使用钱财，如同旅客在客栈使用桌子、杯子、水罐和床铺一样，目的不是要留下来，而是要离开它们。如果你们愿意成为这样的人，你们这些能振奋心灵并听我说话的人；如果你们愿意成为这样的人，你们必将获得祂的许诺。这并不超出你们的力量，因为召叫你们的那一位的手是强有力的。祂召叫了你们。你们要呼求祂，对祂说：祢召叫了我们，我们呼求祢；看，我们听到了祢召叫我们，求祢垂听我们呼求祢：带领我们到祢所应许的地方；完成祢所开始的；不要抛弃祢自己的恩赐；不要离开祢自己的田地；让祢的嫩芽仍被收进祢的谷仓。世上的诱惑虽多，但创造世界的那一位更为伟大。诱惑虽多，但那位寄望于祂而祂毫无欠缺的人，必不失败。\par

11. 弟兄们，我用这些话劝勉你们，因为我们的主耶稣基督所说的自由不属于现在这个时期。请看祂加上的话：\Quote{你们如果做我的门徒，便会认识真理，而真理必会使你们获得自由。}\Quote{必会使你们获得自由}是什么意思？它会使你们成为自由人。总之，那些属肉体的、属血气思想的犹太人——不是那些相信了的，而是那些群众中不信的——认为祂对他们说\Quote{真理必会使你们获得自由}是一种侮辱。他们因被指为奴隶而愤慨。他们确实是奴隶；祂向他们解释那是什么奴役，以及祂所应许的将来的自由是什么。但这自由和那奴役，今天说来话太长。\par

\SourceRecord{Translated by John Gibb.；Nicene and Post-Nicene Fathers, First Series；Vol. 7.；Edited by Philip Schaff.；Buffalo, NY: Christian Literature Publishing Co.,；1888.；<http://www.newadvent.org/fathers/1701040.htm>.}

% structure: p0001 source=h1 target=section reason=page-title

\section{讲道词四十一（\BibleRef{若 8:31-36}）}

1. 关于前一课接下来的内容，今日已从圣福音中公开诵读给我们，我当时暂缓讲述，因为我已经说了很多，并且关于救主恩宠所召唤我们进入的自由，有必要以不仓促、不疏忽的态度来处理。今天，靠着主的帮助，我们打算向你们谈论此事。因为主耶稣基督当时所讲话的对象是犹太人，他们中有很大一部分确实是祂的敌人，但也有一部分已经成为、并将要成为祂的朋友；因为祂看到那里有些人，如我们已说过的，将在祂受难后仍会相信。祂看着这些人说：\Quote{当你们举起人子以后，那时你们就会知道我就是那一位。} 那里也有那些人，当祂如此说时，立即相信了祂。祂对他们说了我们今天所听到的话：\Quote{于是耶稣对信祂的犹太人说：你们若继续存留于我的话中，就真是我的门徒了。} 因继续存留而成为如此；因为正如现在你们是相信者，因继续存留你们将成为观看者。因此接着有：\Quote{并且你们将认识真理。} 真理是不变的。真理是食粮，它更新我们的心智，且永不枯竭；它改变食用者，而它本身不被食用者改变。真理本身就是天主的圣言，是偕同天主的，是唯生的圣子。这真理为了我们而穿上肉体，以便由童贞玛利亚诞生，并应验了预言：\Quote{真理从地上生出。} 因此，这真理在向犹太人讲话时，隐藏在肉体中。但祂隐藏并非为了被否认，而是为了延迟（祂的显现）；延迟，是为了在肉体中受苦；在肉体中受苦，是为了使肉体从罪中被救赎。所以，我们的主耶稣基督，在肉体的软弱方面完全显明，但在天主性的威严方面隐藏，对相信祂的人（当祂如此说时）说：\Quote{你们若继续存留于我的话中，就真是我的门徒了。} 因为那坚持到底的，必要得救。\BibleRef{玛 10:22} \Quote{并且你们将认识真理，} 这真理现在向你们隐藏，却对你们讲话。\Quote{并且真理将使你们自由。} 主是从 \OriginalTerm{自由}{libertas} 这个词取用了 \OriginalTerm{将使自由}{liberabit} 这个词。因为 \OriginalTerm{他使自由}{liberat} 严格来说不过就是 \OriginalTerm{使之自由}{liberum facit}。如同 \OriginalTerm{他拯救}{salvat} 不过就是 \OriginalTerm{使他安全}{salvum facit}；如同\Quote{他治愈}不过就是\Quote{使他完整}；\Quote{他使富足}不过就是\Quote{使他富有}；同样，\OriginalTerm{他使自由}{liberat} 不过就是 \OriginalTerm{使之自由}{liberum facit}。这在希腊词中更为清楚。因为在拉丁语惯用法中，我们通常说一个人被解救（\OriginalTerm{解救}{liberari}），并非指自由，而仅指安全，正如一个人被说成从某种软弱中被解救。所以这样说虽属习惯，但并不确切。但主在说\Quote{并且真理将使你们自由}时，以这种方式使用了这个词，以至于在希腊语中无人能怀疑祂是在谈论自由。\par

2. 简言之，犹太人也这样理解，并\Quote{回答祂；}不是那些已经相信的人，而是那群人中尚未相信的人。\Quote{他们回答祂说：我们是亚巴郎的后裔，从来没有做过谁的奴隶：你怎么说：你们将成为自由的？} 但主并未说\Quote{你们将成为自由的}，而是说\Quote{真理将使你们自由。} 然而，因为他们（如我所说，这在希腊语中很清楚）将那词理解为仅指自由，便自夸为亚巴郎的后裔，说：\Quote{我们是亚巴郎的后裔，从来没有做过谁的奴隶：你怎么说：你们将成为自由的？} 哦，膨胀的皮囊！这并非宽宏大量，而是虚浮的膨胀。因为即使是关乎此生的自由，当你们说\Quote{我们从来没有做过谁的奴隶}时，那怎会是真理呢？难道若瑟没有被卖吗？\BibleRef{创 37:28} 难道圣先知们没有被掳去吗？再者，那同一民族在埃及做砖时，难道不也是在服事严酷的统治者，不仅在金银上，而且在泥土上吗？\BibleRef{出 1:14} 如果你们从来没有做过谁的奴隶，忘恩负义的人们，为什么天主不断提醒你们，祂将你们从奴役之家领出来？或者你们或许认为，你们的祖先曾为奴隶，但你们这些讲话者从来没有做过谁的奴隶？那么，你们当时为何向罗马人纳税，并且你们也曾用这税设陷阱给真理本身，仿佛要陷害祂，当你们说\Quote{纳税给凯撒是否合法？} 为的是，如果祂说合法，你们就可以指责祂对亚巴郎后裔的自由心怀敌意；如果祂说不合法，你们就可以在世界的君王面前控告祂，禁止向这样的人纳税。你们被拿出钱来击败，这很公平，迫使你们自己赞同自己的败诉。因为在那里，你们被告知：\Quote{凯撒的归凯撒，天主的归天主}，在你们自己回答说那钱币上有凯撒的肖像之后。\BibleRef{玛 22:15} 因为正如凯撒在钱币上寻找自己的肖像，同样，天主在人类身上寻找自己的肖像。因此，祂这样回答了犹太人。弟兄们，我被人们的空洞骄傲所触动，因为即使是关于他们按肉体所理解的自由，当他们说\Quote{我们从来没有做过谁的奴隶}时，他们也撒了谎。\par

3. 但我们要更好、更用心地留意主自己的回答，免得我们自己也成为奴仆。因为\BibleQuote{耶稣回答他们说：我实实在在告诉你们：凡犯罪的，就是罪的奴仆。}他是奴仆——但愿是人的奴仆，而不是罪的奴仆！谁听了这话能不战栗？愿主我们的天主赐予我们，也就是你们和我，使我能以合适的言辞讲论这当寻求的自由和那当避免的束缚。\BibleQuote{阿们，阿们[实实在在]，我告诉你们。}真理在说话：主我们的天主以什么意义宣称这是祂所说：\BibleQuote{阿们，阿们，我告诉你们}？祂如此宣布，其托付是沉重的。在某种意义上，如果可以这么说，祂发誓的形式是：\BibleQuote{阿们，阿们，我告诉你们。}\Emph{阿们}在某种程度上可以解释为[它是]真的[真实地，确实地]；然而它并没有被翻译，尽管本可以说\Quote{真实地我告诉你们}。希腊文译者和拉丁文译者都不敢这么做；因为这个词\Emph{阿们}既非希腊文也非拉丁文，而是希伯来文。所以它没有被翻译，以保持作为隐藏之物的遮盖的尊荣；不是为了被否认，而是免得它因暴露在眼前而失去信誉。然而主不仅一次，而是两次说：\BibleQuote{阿们，阿们，我告诉你们。}现在就从这重复本身学习，我们面前的托付包含着多大的分量。\par

4. 那么，所给的是什么托付？我实实在在告诉你们，真理说，祂虽然未曾说\Quote{我确实说}，也绝不可能说谎。然而[藉此]祂强调、灌输祂的托付，以某种方式唤醒沉睡者，使他们留心，且不愿被藐视。祂说什么？\BibleQuote{我实实在在告诉你们：凡犯罪的，就是罪的奴仆。}悲惨的奴役！人们常常在邪恶的主人下受苦时，要求将自己卖掉，并非寻求没有主人，而是无论如何要换一个主人。罪的奴仆能做什么？他能向谁提出要求？向谁寻求赔偿？要求谁把自己卖掉？并且有时一个人的奴仆，被无情主人的命令折磨殆尽，在逃跑中找到安息。罪的奴仆能逃往何处？他无论逃到哪里，都带着他自己。邪恶的良心不能逃离自身；它无处可去；它跟随自己。是的，他不能脱离自己，因为他所犯的罪在里面。他犯罪是为了获得某种肉体的享乐。享乐逝去，罪恶存留。令人愉悦的已经消失；刺痛留在了后面。邪恶的束缚！有时人们逃到教会，我们通常允许他们，尽管他们未经教导——这些人希望摆脱主人，却不愿摆脱他们的罪。但有时那些受非法邪恶轭制的人也逃到教会避难；因为，虽然是生来自由的人，却被拘禁在束缚中，并向主教提出申诉。除非他愿意尽一切努力拯救自由出生者免于压迫，否则就被视为不仁慈。让我们都逃向基督，为反对罪恶而向天主申诉，祂是我们的拯救者。让我们寻求将自己出卖，以便能被祂的血赎回来。因为主说：\BibleQuote{你们白白地被卖，也无需金钱被赎}\BibleRef{依 52:3}无需代价，即无需你们自己的，而是因我的。主如此说；因为祂亲自付出了代价，不是金钱，而是祂自己的血。否则我们就仍然是奴仆和贫困者。\par

5. 唯有主才能解救我们脱离这束缚。那没有束缚的主亲自解救我们脱离它；因为只有祂无罪地降生血肉。因为你们所看见被人抱在怀中的婴孩，还不能行走，却已戴上了锁链；因为他们从亚当领受了那要由基督解除的束缚。他们受洗时，也分享了主所应许的恩宠；因为只有那无罪而来、并成为罪的祭献的那一位才能解救脱离罪。因为你们听见宗徒被宣读时说：\Quote{我们作大使，}他说，\Quote{为基督，仿佛天主借我们来劝你们；我们替基督求你们，}——也就是说，仿佛基督在求你们，求什么？——\Quote{你们与天主和好吧。}如果宗徒劝我们并与天主和好，那么我们曾是仇敌。因为除非从敌对状态，无人能和好。我们成为仇敌不是因本性，而是因罪。同样根源我们成为罪的奴仆，成为天主的仇敌。天主没有自由状态的仇敌。他们必须是奴隶；除非被他们因罪而欲与之结仇的那位解救，他们将永远是奴隶。因此，他说：\Quote{我们替基督求你们与天主和好。}但我们如何和好？除非除掉那分离我们与祂之间的障碍？因为祂借先知说：\Quote{祂没有使耳朵沉重以致听不见；但你们的罪孽使你们与你们的天主隔绝。}\BibleRef{依 59:1}因此，除非那中间的障碍被除去，并以别的代替，我们便不和好。因为有分离的障碍，也有和好的中保。分离的障碍是罪，和好的中保是主耶稣基督：\Quote{因为只有一个天主，在天主与人之间的中保也只有一人，就是基督耶稣其人。}\BibleRef{弟前 2:5}那么，为除去那分离的墙壁——也就是罪——那位中保来了，司铎亲自成为了祭品。因为祂成了罪的祭献，在祂受难的十字架上将自己作为全燔祭献上，宗徒在说了\Quote{我们替基督求你们与天主和好}之后——好像我们说：我们怎能和好？——接着又说：\Quote{祂使那无罪的，即基督自己，替我们成为罪，好叫我们在祂内成为天主的义。}\BibleRef{格后 5:20}\Quote{祂，}他说，基督我们自己的天主，\Quote{不知罪。}因为祂降生成人，就是带着罪肉身的形状，\BibleRef{罗 8:3}但不是有罪的肉身，因为祂完全没有罪；因此成为了真实的赎罪祭，因为祂自己无罪。\par

6. 但也许是出于我个人的某种特殊见解，我说过\Emph{罪}是赎罪祭。让那些读过的人自由地承认它；让那些没有读过的人不要迟疑；我说，不要迟疑去读，好使他们在判断时真实。因为当天主命令献赎罪祭时，在这些祭献中没有罪的赦免，而是将来之事的预像，那些同样的祭献，同样的供物，同样的牺牲，同样的动物，被带来为罪而宰杀，在它们的血中预表了那真实的血，它们本身被法律称为\Emph{罪}；甚至到如此地步，在某些段落中如此写道：司铎在献祭时，要按手在罪上，即按手在将要为罪献祭的牺牲的头上。因此，这样的罪，即这样的赎罪祭，我们的主耶稣基督成了，\Quote{不知罪。}\par

7. 那在圣咏中说\Quote{我成了无助的人，在死者中自由的}的一位，以有效的功劳解救我们脱离这罪的束缚。因为只有祂是自由的，因为祂没有罪。因为祂自己在福音中说：\Quote{看，今世的首领来了，}意思是魔鬼即将借着迫害的犹太人来——\Quote{看，}祂说，\Quote{他来了，在我身上找不到什么。}不像他在那些被他杀害的义人身上找到了某些罪；在我身上他找不到什么。正如有人问祂：若他在你身上找不到什么，为何要杀你？祂接着说：\Quote{但为叫众人知道，我承行我父的旨意，起来，我们离开这里吧。}祂说，我不是因我罪性的必然性而付死亡的代价；而是我所死的死亡，是承行我父的旨意。在这事上，我是主动承受而非被动忍受；因为我若不情愿，就不会受这苦难。你们有祂在另一处说：\Quote{我有权柄舍命，也有权柄再取回来。}这里确实有一位\Quote{在死者中自由的。}\par

8. 既然，每个犯罪的人都是罪的奴隶，请听我们得自由的是何种希望。祂说：\BibleQuote{奴隶}，\BibleQuote{不会永远住在家里。}教会就是那家，奴隶就是罪人。许多罪人进入教会。因此祂没有说\BibleQuote{奴隶}不在家中，而是说\BibleQuote{不会永远住在家里。}那么，如果那里没有奴隶，谁会在那里呢？因为正如圣经所说，\BibleQuote{当正义的君王坐在宝座上，谁能自夸有洁净的心呢？谁能自夸纯洁无罪呢？}\BibleRef{箴 20:8}祂说了这句话，大大警醒了我们，我的弟兄们：\BibleQuote{奴隶不会永远住在家里。}但祂又接着说：\BibleQuote{但儿子永远长存。}那么，基督会独自在祂的家中吗？就没有人留在祂身边吗？如果没有身体，祂是谁的头呢？或者，圣子就是这一切，既是头又是身体？祂激起恐惧和希望，并非没有原因：恐惧，是为了叫我们不要喜爱罪恶；希望，是为了叫我们不要不相信罪的赦免。祂说：\BibleQuote{凡犯罪的，就是罪的奴隶。奴隶不会永远住在家里。}那么，我们这些并非无罪的人，还有什么希望呢？听听你的希望：\BibleQuote{圣子永远长存。所以，如果圣子使你们自由，你们就真正自由了。}我们的希望就是，弟兄们：被那自由者所释放；并且，在释放我们时，祂使我们成为祂的仆人。因为我们曾是情欲的奴隶；但获得释放后，我们成为爱德的仆人。宗徒也这样说：\BibleQuote{弟兄们，你们蒙召是为了自由；只是不可将自由当作放纵肉欲的机会，却要凭爱德互相服事。}\BibleRef{迦 5:13}所以，不要让基督徒这样说：\Quote{我是自由的；我蒙召是为了自由；我曾是奴隶，但已被赎回，并且藉着此次赎回，我成了自由的人；我可以为所欲为；如果我是自由的，就没有人能阻挠我的意愿。}但如果你们怀着这样的意愿犯罪，你们就是罪的奴隶。所以不要滥用你们的自由，作为犯罪的自由，而要用来不犯罪。因为只有你们的意愿是虔诚的，它才是自由的。如果你们仍然是仆人——脱离罪，成为正义的仆人，你们就会自由。正如宗徒所说：\BibleQuote{你们从前是罪的奴隶，对于正义是自由的。但现在，你们既然从罪中获得自由，成为天主的仆人，就有了成圣的果实，结局是永生。}让我们追求后者，实行前者。\par

9. 自由的第一阶段是脱离罪行。你们要留意，我的弟兄们，要留意，免得我以任何方式误导你们对目前自由及其未来性质的理解。审视任何在今生具有最高正直的人，无论他可能已多么配得上正直之名，他并非没有罪。请听本福音的作者圣若望本人在其书信中所说的：\BibleQuote{如果我们说我们没有罪，我们就是自欺，真理就不在我们内。}\BibleRef{若一 1:8}只有祂能说出这话，祂是\BibleQuote{在死者中自由}；只有祂能被称为不知罪者。这句话只能用于祂，因为祂也\BibleQuote{在各样事上受过试探，像我们一样，却没有罪。}\BibleRef{希 4:15}只有祂能说：\BibleQuote{看，今世的首领来了，他在我身上一无所获。}审视任何其他被认为正直的人，他并非在所有方面都没有罪；甚至如约伯那样的人，主曾为他作证，以致魔鬼充满嫉妒，要求试探他，并在试探中自己失败，为使约伯得以被证明。\BibleRef{约 1:2}他受证明的原因，并非天主不知道他必然赢得胜利的花冠，而是为了使他成为他人效法的模范。约伯自己说什么？\BibleQuote{谁能洁净呢？连那在世生存只一日的婴儿也不洁净。}但很明显，许多人毫无异议地被称为正直人，因为这个词被理解为没有犯罪行为；因为在人类事务中，那些没有犯罪行为的人没有合理的理由被指责。但罪行是严重的罪，最应受谴责和定罪。然而，这并不是说天主谴责某些罪，而成义和称赞另一些罪。祂不赞成任何罪。祂憎恨所有的罪。正如医生不喜欢病人的疾病，并以治疗手段来除去疾病，使病人康复；同样，天主藉着祂的恩宠在我们内工作，使罪被消灭，人获得自由。但你们会说：什么时候被消灭呢？如果它减少了，为什么不被消灭呢？在那些不断前进的人身上，它不断减少；在那些达到完美的人身上，它被消灭了。\par

10. 所以，自由的第一个阶段就是脱离罪行（罪恶行为）。因此，宗徒保禄在决定有关长老或执事，或任何要被祝圣以监督教会的人的祝圣时，不说\Quote{若有人没有罪}；因为他若这样说，每一个人都会被拒绝，认为不适合，没有人会被祝圣：而是说：\Quote{若有人没有罪行}（英文译本作\Quote{指责}），诸如凶杀、奸淫、淫乱不洁、偷窃、欺诈、亵圣，以及其他类似的事。当一个人开始脱离这些（每个基督徒都应当如此），他开始抬起头来走向自由；但那只是开始的自由，并非完成的自由。有人说，为什么这不是完成的自由？因为\BibleQuote{我见在我的肢体中有另一条法律，反抗我心神的法律；}\BibleQuote{因为我所愿意的，}他说，\BibleQuote{我偏不作；我所憎恨的，我反而去作。}\BibleQuote{肉体，}他说，\BibleQuote{相反圣神，圣神相反肉体，以致你们不能作你们所愿意的事。}\BibleRef{迦 5:17} 部分是自由，部分是束缚：尚未完全，尚未纯洁，尚未圆满的自由，因为尚未达到永恒。因为我们在部分上仍有软弱，在部分上已获得自由。我们所有的罪，早已在圣洗中涤除了。但既然我们的所有不义都被涂抹，难道就没有留下任何软弱吗？若没有留下，我们就该在此世无罪的活着。然而谁敢如此说，除非是骄傲的人，除非是不配得救主仁慈的人，除非是自欺欺人、远离真理的人？因此，由于仍有软弱存在，我敢说：在多少程度上我们事奉天主，我们就是自由的；在多少程度上我们事奉罪的法律，我们就仍在束缚中。因此，宗徒在开始时所说的话：\BibleQuote{按内心来说，我喜爱天主的法律。}\BibleRef{罗 7:22} 正是在这里，我们在其中是自由的，我们在其中喜爱天主的法律；因为自由具有喜乐。只要你是出于畏惧而行善，天主就不能使你喜悦。在祂内找到你的喜悦，你便是自由的。不要畏惧惩罚，而要爱慕正义。你尚不能爱慕正义吗？即使畏惧惩罚，好使你达到爱慕正义。\par

11. 因此，按照以上所说的程度，他感到自己已经是自由的，所以说：\BibleQuote{按内心来说，我喜爱天主的法律。}我喜爱这法律，喜爱它的要求，喜爱正义本身。\BibleQuote{但我见在肢体中有另一条法律}——即这存留的软弱——\BibleQuote{反抗我心神的法律，并把我掳去，隶属那在我肢体中罪的法律。}在这一方面，他感到自己被掳，因为正义尚未完全；因为在他喜爱天主的法律之处，他不是俘虏，而是法律的友人；因此是自由的，因为是友人。那么，对这仍存留的部分该如何处理呢？除了仰望那曾说\Quote{若是子使你们自由，你们就必真正自由}的那一位之外，还能怎样呢？的确，那如此说话的人也仰望祂：\BibleQuote{我真苦啊！}他说：\BibleQuote{谁能救我脱离这该死的身体呢？我感谢天主，藉着我们的主耶稣基督。}\Emph{因此}\Quote{若是子使你们自由，你们就必真正自由。}然后他这样总结：\BibleQuote{这样看来，我本人用内心服事天主的法律，用肉体服事罪的法律。}\BibleRef{罗 7:23} \Emph{我本人}，他说；因为并非有两个我彼此相反，出自不同的根源；而是\BibleQuote{我本人用内心服事天主的法律，用肉体服事罪的法律，}只要软弱仍在对抗救恩。\par

12. 但如果你用肉身顺从罪的法律，就当如宗徒自己所说：\BibleQuote{所以，不要让罪恶在你们必死的身体上为王，而顺从它的情欲；也不要把你们的肢体交给罪恶，作不义的兵器。}\BibleRef{罗 6:12} 他没有说\InnerQuote{不要让罪存在}，而是说\Quote{不要让罪为王}。只要罪仍在你们的肢体中，至少当除去它的王权，不要顺从它的命令。忿怒升起吗？不要将你的舌头交给忿怒去说恶言；不要将你的手或脚交给忿怒去击打。那无理性的忿怒若不因肢体中的罪，就不会升起。但要除去它的统治权，让它没有武器攻击你。那时，当它开始发现缺少武器时，它也将学习不再升起。\BibleQuote{不要把你们的肢体交给罪恶，作不义的兵器，}否则你将被完全掳去，就没有余地说：\Quote{我以心神服事天主的法律。}因为若心神持守兵器，肢体就不会被煽动去服事猖獗的罪。让内在的统治者持守堡垒，因为它在更大的统治者之下，且确知有援助。让它勒住忿怒，让它克制邪恶的贪欲。内在有需要勒住、需要克制、需要管理的东西。那位以心神服事天主法律的义人，他所渴望的，不就是从这需要勒制的事中完全得释放吗？每一个追求成全的人都应当努力争取：让情欲本身，既不再得到肢体的顺从，就在朝圣者的前行中日益减少。他说：\BibleQuote{立志为善由得我，但如何成全善却不行。}\BibleRef{罗 7:18} 他是否说行善由不得我？若他这样说，希望便丧失了。他没有说行由不得我，而是说\BibleQuote{成全善却不行}。因为成全善，不就是消除并终结恶吗？消除恶，不就是法律所说的\BibleQuote{不可贪恋}吗？\BibleRef{出 20:17} 完全没有贪恋就是成全善，因为那就是消除恶。他所说的\BibleQuote{成全善却不行}，是指他的行为尚不能达到从贪恋中得释放的地步。他只能努力勒住贪恋，拒绝同意贪恋，不将肢体交付给它服务。因此他说：\BibleQuote{成全善却不行}，我无法成全\Quote{不可贪恋}这条诫命。那么需要什么呢？要成全这句：\BibleQuote{不要随从你的贪欲。}\BibleRef{德 18:30} 同时，只要还有非法的贪欲在你的肉身中，你就这样做：\Quote{不要随从你的贪欲。}你要持守在事奉天主、在基督的自由中。以心神服事你的天主的法律。不要把自己交给你的贪欲。若跟随它们，你就增加了它们的力量。若给它们力量，当你靠自己的力量养着敌人与你对抗时，你如何能得胜呢？\par

13. 那么，在主耶稣那里所说的圆满而完美的自由是什么呢？祂说：\BibleQuote{如果天主子使你们自由，你们就真正自由了。}这圆满而完美的自由何时能实现呢？当不再有仇敌，当\BibleQuote{最后被毁灭的仇敌是死亡}的时候。\BibleQuote{因为这可朽坏的必须穿上不朽坏的，这可死的必须穿上不死的。当这可死的穿上不死的时候，那时就要应验那经上的话：死亡被胜利吞灭了。死亡啊，你的刺在哪里？}这\BibleQuote{死亡啊，你的刺在哪里？}是什么？\BibleQuote{肉身\Emph{相反}圣神，圣神相反肉身，}但这是在罪的肉身強盛时才如此。\BibleQuote{死亡啊，你的刺如今在哪里？}如今我们要生活，不再死亡，在为我们死而复活的那一位内。他说：\BibleQuote{他们活着，不再为自己活，而是为替他们死而复活的那一位活。}\BibleRef{格后 5:15} 让我们如同受伤者一样向医生祈祷，让我们被抬到客栈去医治。因为正是祂应许了救恩，祂怜悯了那个路上被强盗打得半死的人。祂倒上油和酒，包扎了伤口，将那人扶上自己的牲口，带到客栈，托付给店主照料。这是哪位店主呢？也许是那位说\BibleQuote{我们作基督的使者}的人。祂还给了两个银币，作为医治那受伤者的费用。\BibleRef{路 10:30} 或许这两条诫命就是一切法律和先知所依靠的。\BibleRef{玛 22:37} 因此，弟兄们，教会就是那客栈，那受伤者在此暂时得医治；但在教会之上，乃是主人的产业。\par

\SourceRecord{Translated by John Gibb.；Nicene and Post-Nicene Fathers, First Series；Vol. 7.；Edited by Philip Schaff.；Buffalo, NY: Christian Literature Publishing Co.,；1888.；<http://www.newadvent.org/fathers/1701041.htm>.}

% structure: p0001 source=h1 target=section reason=page-title

\section{论训诲录第四十二（若望福音第8章37-47节）}

1. 我们的主取了仆人的形象，却并非仆人；即使在仆人形象中，他仍是主（那肉身的形象确实似仆人；他虽然\BibleQuote{带着罪恶肉身的形状}\BibleRef{罗 8:3}，却不是罪恶的肉身），他向信他的人应许了自由。但犹太人，仿佛骄傲地夸耀自己的自由，愤怒地拒绝被释放，因为他们自己正是罪恶的奴隶。因此，他们声称自己是自由的，因为他们是亚巴郎的后裔。那么，主对此如何回答，我们今天在经课的诵读中听到了。他说：\BibleQuote{我知道你们是亚巴郎的子孙；但你们却想杀我，因为我的话在你们心中没有立足之地。}我认识你们，他说：\BibleQuote{你们是亚巴郎的子孙，但你们却想杀我。}我承认血统的渊源，却否认信仰的心灵。\BibleQuote{你们是亚巴郎的子孙}，但按肉身而言。因此他说：\BibleQuote{你们想杀我，因为我的话在你们心中没有立足之地。}若我的话被接纳，它就会立足；若你们被接纳，你们就会被网在信德的网中，如同鱼一般。那么，\Quote{没有立足之地}是什么意思？它不占据你们的心，因为不被你们的心接纳。因为天主的话是如此，对信者来说应当像鱼钩一样：它被抓住时就抓住了。被抓住的人不受伤害，因为他们是为得救而被抓住，而非为毁灭。因此主对门徒说：\BibleQuote{来跟随我，我要使你们成为捕人的渔夫。}\BibleRef{玛 4:19} 但这些人并非如此；然而他们是亚巴郎的子孙——一个天主的人的子孙，自己却不义。他们继承了肉身的血统，却因不效法他们祖先的信德而堕落。\par

2. 你们确实听见主说：\BibleQuote{我知道你们是亚巴郎的子孙。}听听他随后说的话：\BibleQuote{我所说的，是我在父那里看见的；而你们所做的，却是在你们父那里看见的。}他已经说过：\BibleQuote{我知道你们是亚巴郎的子孙。}那么，他们做了什么？就是他告诉他们的：\BibleQuote{你们想杀我。}这件事他们从未在亚巴郎身上看见过。但主愿意我们理解为父是天主，当他说：\BibleQuote{我所说的，是我在父那里看见的。}我看见了真理，我说真理，因为我是真理。因为若主说的是他从父那里看见的真理，那么他看见的是他自己——他说的是他自己；因为他本身就是父的真理，这真理是他从父那里看见的。因为他是圣言——那与天主同在的圣言。那么这些人的邪恶，就是主所责斥和谴责的，他们在哪里看见的呢？在他们父那里。当我们随后更清楚地听到谁是他们父时，就会明白他们在这样的父那里看见了什么；因为此时他还没有说出他们的父。稍早他提到亚巴郎，但那是关于他们肉身的渊源，而非生活的相似。他即将论到他们另一个父，那父既没有生他们，也没有造他们成为人。然而他们仍是他的子女，就他们是恶而言，而非就他们是人而言；在于他们效法他，而非受造于他。\par

3. 他们回答他说：\BibleQuote{我们的父是亚巴郎。}仿佛在说：你对亚巴郎有什么可说的？或者，如果你能，你敢责备亚巴郎？并非主不敢责备亚巴郎；但亚巴郎不是一位应受主责备的人，而是蒙主赞许的人。但这些人似乎挑战他说些亚巴郎的坏话，从而为他们的意图提供借口。\BibleQuote{我们的父是亚巴郎。}\par

4. 让我们听听主如何回答他们，在赞扬亚巴郎的同时定他们的罪。耶稣对他们说：\BibleQuote{如果你们是亚巴郎的子孙，就该做亚巴郎所做的事。但现在你们却想杀我——一个向你们讲真理的人，这真理是我从天主那里听来的；亚巴郎没有做过这样的事。}看，亚巴郎被赞扬，他们被定罪。亚巴郎不是杀人者。我不是说（他暗示）我是亚巴郎的主；即使我这样说，也是真理。因为他在另一处说过：\BibleQuote{在亚巴郎出现以前，我就存在}\BibleRef{若 8:58}；那时他们就拿起石头要砸他。但此时他没有这样说。同时，你们看见我，看着我的时候，当你们只把我当作一个人的时候，我是一个人。那么，你们为什么想要杀一个向你们讲论他从天主那里听到的话的人呢？岂不是因为你们不是亚巴郎的子孙吗？然而他上面说过：\BibleQuote{我知道你们是亚巴郎的子孙}。他并不否认他们的血统，却谴责他们的行为。他们的肉身来自他，但他们的生活不是。\par

5. 可是，我们，亲爱的诸位，我们是\Person{亚巴郎}的后裔吗？或者\Person{亚巴郎}按肉体而论可算作我们的父亲吗？犹太人的肉体从他的肉体而生，基督徒的肉体却并非如此。我们来自其他民族，然而，通过效法他，我们成了\Person{亚巴郎}的子女。请听宗徒的话：\Quote{对亚巴郎和他的苗裔，恩许是那样说的。他并不说，}他又说：\Quote{对苗裔们，如同对许多；而是对你的苗裔，那一位就是基督。}因此，我们因天主的恩宠成了\Person{亚巴郎}的苗裔。天主并没有使任何人与\Person{亚巴郎}的肉体同作他的承继人。他剥夺了前者，收养了后者；从那棵根在圣祖中的橄榄树上，他折下了骄傲的本枝，接上了谦卑的野橄榄枝。\BibleRef{罗 11:17}因此，当犹太人来\Person{若翰}那里受洗时，他斥责他们说：\Quote{毒蛇的种类！}他们确实大大夸耀其出身的尊贵，但他却称他们为毒蛇的种类——甚至不是人的种类，而是毒蛇的种类。他看见了人的形体，却察觉了毒汁。然而他们来是要改变，因为他们还是要受洗；他对他们说：\Quote{毒蛇的种类！谁指示你们逃避那即将来临的忿怒？你们要结出与悔改相称的果实。不要心里说：我们有亚巴郎为父。我告诉你们：天主能从这些石头中给亚巴郎兴起子孙来。}\BibleRef{玛 3:7}如果你们不结出与悔改相称的果实，就不要以这样的谱系自夸。天主能够定你们的罪，却不使\Person{亚巴郎}缺少子孙。因为他有办法给\Person{亚巴郎}兴起子孙来。那些效法他信德的人，将被造成为他的子女。\Quote{天主能从这些石头中给亚巴郎兴起子孙来。}我们就是这样的人。在我们的父母中，我们是石头，因为我们曾把石头当作神来敬拜。天主从这样的石头中为\Person{亚巴郎}创造了家庭。\par

6. 那么，这空洞而虚妄的夸耀何以自高自大？让他们停止夸耀自己是\Person{亚巴郎}的子女。他们已经听到了他们应当听到的话：\Quote{如果你们是亚巴郎的子女，}就用你们的行动，而不是言语来证明。\Quote{你们却想要杀我，一个向你们讲真理的人——}我暂且不说天主子；我不说天主；我不说圣言，因为圣言不死。我只说你们所看见的这一个；因为你们所能杀的只是你们所看见的，而你们所不能看见的，你们却可以得罪。\Quote{这，}所以，\Quote{亚巴郎却没有做。}\Quote{你们正做你们父亲的勾当。}而此刻他还没有说出他们的父亲是谁。\par

7. 现在他们给了他什么回答呢？因为他们开始有些意识到主不是在谈论肉身的血统，而是在谈论他们的生活方式。又因为圣经——他们阅读的——有一种习惯，在精神意义上称之为淫乱，即当灵魂仿佛因屈服于许多假神而被玷污时，他们这样回答说：\Quote{于是他们对他说：我们不是从淫乱生的；我们只有一位父亲，就是天主。}\Person{亚巴郎}现在已失去了重要性。因为他们被真理之口驳斥得体无完肤；因为\Person{亚巴郎}是这样的，他们不效法他的行为，却以他的血统为荣。于是他们改变了自己的答复，我猜想，他们在心里说：每当我们提到\Person{亚巴郎}，他就对我们说：你们为什么不像你们所夸耀的血统那样效法他？这样的人，如此圣洁、正义、又无诡诈，我们无法效法。让我们称天主为父亲，看看他会对我们说什么。\par

8. 难道虚谎真的找到了可说的，而真理就不应有适当的回答吗？让我们听听他们说什么：让我们听听他们听到什么。他们说：\Quote{我们只有一位父亲，就是天主。}于是耶稣对他们说：\Quote{如果天主是你们的父亲，你们就〔必然〕爱我；因为我从天主出发而来；我并非由我自己而来，而是他派遣了我。}你们称天主为父亲，那么至少要把我当作兄弟来认识。同时，他刺激了有悟性者的心，因为他触及了他惯常所说的那句话：\Quote{我并非由我自己而来：他派遣了我。我从天主出发而来。}请记住我们惯常所说的：从他而来；而从他而来者，与他同在。因此，基督的受差遣就是他的降生成人。但就圣言从天主出发而言，那是永恒的出发。时间不包含他，时间是由他创造的。谁都不应在心里说：在圣言存在之前，天主如何存在？绝不要说：在天主圣言之前。天主从来没有没有圣言，因为圣言是常存的，不是转瞬即逝的；他是天主，不是声音；天地由他造成，他不与世间所造之物一同消逝。于是，他作为天主、平等者、独生子、圣父的圣言出发而来，并且来到我们中间，因为圣言成了血肉，得以住在我们中间。他的来临彰显了他的人性；他的常住彰显了他的神性。是他的天主性我们朝着它前进，他的人性我们藉以前进。如果他未曾成为我们藉以前进的那一位，我们就永远无法达到那永恒常存的他。\par

9. 他说：\Quote{你们为什么不明白我的话语？就是因为你们不能听我的话。}所以他们不能明白，因为他们不能听。而他们为什么不能听，不就是因为他们拒绝通过相信来改正吗？而为什么这样？\Quote{你们是出于你们的父亲魔鬼。}你们还要说父亲多久？你们还要多少次更换你们的父亲——时而亚巴郎，时而天主？请从天主子那里听你们是谁的子女：\Quote{你们是出于你们的父亲魔鬼。}\par

10. 现在，我们必须警惕摩尼教徒的异端，他们主张存在一个邪恶本原，以及一个拥有其首领的黑暗家族，竟敢与天主争战；但天主为了不让祂的王国被这敌对家族征服，便派出祂自己光明王国的首领，即祂的子孙，去对抗他们；如此便征服了魔鬼起源的那个种族。他们也由此认为我们的肉身起源于该种族，因此认为主说：\Quote{你们是出于你们的父亲魔鬼}，是因为他们按本性是恶的，起源于对立的黑暗家族。他们就这样犯错误，这样他们的眼睛被蒙蔽，这样他们因相信那创造他们者的谎言而自成为黑暗家族。因为每个本性都是善的；但人的本性被邪恶的意志败坏了。天主所造的不能是恶的，除非人自招恶。但当然，造物主是造物主，受造物是受造物。受造物不能与造物主等同。要区分创造者与祂所造的。长凳不能与木匠等同，柱子也不能与建造者等同；然而木匠虽然做了长凳，却并没有创造木材。但主我们的天主，以祂的全能和圣言，创造了祂所造的一切。祂没有材料来造祂所造的一切，却仍造了它们。因为它们因祂的意愿而受造，因祂的言语而受造；但受造物不能与创造者相比。如果你寻找一个合适的比较对象，请将你的心思转向独生子。那么，犹太人如何成为魔鬼的儿女？是因模仿，而非因出生。请听圣经通常的说法。先知对同一群犹太人说：\Quote{你的父亲是阿摩黎人，你的母亲是赫特人。}\BibleRef{则 16:3}阿摩黎人并非生出犹太人的民族。赫特人本身也与犹太人的种族完全不同。但因阿摩黎人和赫特人是不敬神的，而犹太人模仿了他们的不敬，他们便为自己找到了父母——并非由其所生，而是与其同受惩罚，因为他们随从了他们的习俗。但或许你会问：魔鬼本身从何而来？当然与其他天使同源。但其他天使持守了服从；他因不服从而骄傲，如天使般堕落，成了魔鬼。\par

11. 但现在请听主所说的：\Quote{你们，}祂说，\Quote{是出于你们的父亲魔鬼，并且你们父亲的欲望你们愿意去做。}这就是你们如何成为他的儿女：因为你们的欲望如此，并非因你们由他所生。他的欲望是什么？\Quote{他从起初就是杀人者。}这就是对\Quote{你们父亲的欲望你们愿意去做}的解释。\Quote{你们寻找机会要杀我，一个对你们讲真理的人。}他（魔鬼）也对人有恶意，并杀了人。因为魔鬼出于对人的恶意，伪装成蛇，对女人说话，并通过女人将毒注入男人。他们因听从魔鬼而死亡，\BibleRef{创 3:1}若他们听从了主，便不会听从魔鬼；因为人处于创造他者与堕落者之间，应当服从创造者，而非欺骗者。因此\Quote{\Emph{他}从起初就是杀人者。}请看这杀人的种类，弟兄们。魔鬼被称为杀人者，并非因为他佩剑或执刀。他来到人前，播撒恶意的建议，并杀死了他。因此，不要以为当你说服你的兄弟去行恶时，你不是杀人者。如果你说服你的兄弟行恶，你便是杀了他。并且为使你知道你杀了他，请听圣咏：\Quote{人们的牙齿是枪和箭，他们的舌头是锋利的剑。}因此，你们\Quote{愿意去做你们父亲的欲望}；如此你们疯狂地追随肉体，因为你们不能追随圣神。\Quote{他从起初就是杀人者；}至少就人类始祖而言。从杀人可能发生的那一刻起，\Emph{他}就是杀人者。只有人被造之后，杀人才能发生。因为除非人先被造，否则人不能被杀害。因此\Quote{\Emph{他}从起初就是杀人者。}他为何成为杀人者？\Quote{他不站立在真理中。}所以他曾在真理中，但因不站立而堕落。为什么\Quote{他不站立在真理中}？\Quote{因为真理不在他内；}不像在基督内。真理在基督内，以基督本身\Emph{是}真理的方式。因此，如果他曾站立在真理中，他就会站立在基督内；但\Quote{他没有在真理中存留，因为在他内没有真理。}\par

12. \BibleQuote{他说谎时，是出于自己说的，因为他是说谎者，也是谎言之父。}这是什么？你们已听见了福音的话，也留心听了。现在，我重述一遍，好让你们清楚地明白你们所思虑的对象。主说的那些关于魔鬼的事，原是主应当对魔鬼说的。\BibleQuote{他从起初就是杀人者}，这是真的，因为他杀了第一个人；\BibleQuote{他不站在真理中}，因为他从真理中堕落了。\BibleQuote{他说谎时}，即魔鬼自己，\BibleQuote{是出于自己说的}；因为他是说谎者，也是（他的）谎言之父。从这些话中，有些人认为魔鬼有一个父亲，便探究谁是魔鬼的父亲。实在，\OriginalTerm{摩尼教信徒}{Manicheans}这个可恶的错误直到现今仍有办法欺骗简单人。因为他们惯常说：假设魔鬼是一个天使，然后堕落了；如你们所说，罪恶由他开始；但是，谁是他的父亲呢？我们相反答道：我们中谁曾说过魔鬼有父亲呢？他们则反驳说：主说过，福音也宣告了，论及魔鬼：\BibleQuote{他从起初就是杀人者，不站在真理中，因为他里面没有真理。他说谎时，是出于自己说的，因为他是说谎者，也是\Emph{他的}父亲。}\par

13. 你们要听且明白。我不必送你们到远处去寻求意思；只从这些话本身就能明白。主称魔鬼为\OriginalTerm{谎言之父}{father of lies}。这是什么？听这是什么，只需反复思量这些话，便能明白。并非每个说谎的人都是他谎言的父亲。因为如果你从别人那里得到一个谎言并说出来，你固然因说出谎言而说了谎，但你并非那个谎言的父亲，因为你是从别人那里得来的。但魔鬼是出于自己而说谎的。他生出了自己的谎言；他没有从任何人那里听到。正如天主圣父生了圣言作为祂的圣子，同样，魔鬼堕落后，生出了谎言作为他的儿子。听到这些，现在回想并默思主的话。你们公教信友，要思考你们所听到的；注意祂所说的。\BibleQuote{他}——谁？魔鬼——\BibleQuote{从起初就是杀人者。}我们承认——他杀了亚当。\BibleQuote{他不站在真理中。}我们承认，因为他从真理中堕落了。\BibleQuote{因为他里面没有真理。}对的：因他从真理中堕落，便失去了真理的拥有。\BibleQuote{他说谎时，是出于自己说的，因为他是说谎者，也是谎言之父。}他既是说谎者，又是\OriginalTerm{谎言之父}{father of lies}。因为你，可能是说谎者，因为你说了谎，但你并非谎言之父。因为如果你从魔鬼那里得到了你所说的，并相信了魔鬼，你是一个说谎者，但并非谎言之父。但他，因为不是从别处得来那谎言——他曾以蛇的形式像用毒药一样杀害了人——他是谎言之父，正如天主是\OriginalTerm{真理之父}{Father of truth}。所以，离开谎言之父吧，赶快到真理之父那里去，拥抱真理，好使你进入自由。\par

14. 那些犹太人，于是说了他们从他们的父亲那里看见的事。那不是虚假吗？但主从祂的父那里看见祂应当说什么；那不是祂自己吗？不是圣父的圣言吗？不是圣父的真理，它本身是永恒的，并与圣父同永恒吗？那么，他\BibleQuote{从起初就是杀人者，不站在真理中，因为他里面没有真理；他说谎时，是出于自己说的，因为他是说谎者}——不仅是个说谎者，而且是\BibleQuote{那谎言之父}；就是说，他所说的谎言本身，他是其父，因为他自己生出了他的谎言。\BibleQuote{正因为我给你们讲真理，你们却不信我。你们中谁能指证我有罪？}如同我指证你们和你们的父亲一样？\BibleQuote{我若说真理，为什么你们不信我呢？}只因为你们是魔鬼的子女吗？\par

15. \Quote{属于天主的，听从天主的话；你们所以不听，因为你们不是属于天主。}这里，同样，你们不可从他们人性的本质，而要从他们的败坏来看。如此，他们既属于天主，又不属于天主。按本性他们属于天主，按败坏他们不属于天主。请你们留意，我恳求你们。在福音中你们有解毒剂，以对抗异端者的毒害和不敬的谬误。因为关于这些话，摩尼教徒惯常说：看，这里有两种本性——一个善，一个恶；主如此说。主说了什么？\Quote{你们所以不听我，因为你们不是属于天主。}这就是主的话。那么，他反驳道：你对此有何言说？请听我说。他们既属于天主，又不属于天主。按本性他们属于天主；按败坏他们不属于天主；因为那属于天主的美善本性，因听信了魔鬼的诱言而自愿犯罪，于是被败坏；因此它寻求医生，因为它已不再健康。这就是我所说的。但你们认为他们既属于天主又不属于天主是不可能的。请听为什么这不是不可能的。他们既属于天主又不属于天主，正如他们既是亚巴郎的子孙又不是亚巴郎的子孙一样。你们这里有此例。并非如你们所言。请听主亲自说的话；正是他对他们说：\Quote{我知道你们是亚巴郎的子孙。}主岂能说谎？断然不能。那么主所说的是真的吗？是真的。那么他们是亚巴郎的子孙是真的吗？是真的。但请听他否认这一点。那位曾说\Quote{你们是亚巴郎的子孙}的，亲自否认他们是亚巴郎的子孙。\Quote{如果你们是亚巴郎的子孙，就该做亚巴郎的事。但如今你们却想杀我，一个把从天主那里听到的真理告诉你们的人：亚巴郎没有做过这样的事。你们做你们父亲的作为}，即魔鬼的作为。那么，他们如何既是亚巴郎的子孙，又不是他的子孙呢？他在他们身上显明了这两种状态。他们按血统的起源是亚巴郎的子孙，但在随从魔鬼诱惑的罪过中不是他的子孙。同样，应用于我们的主和天主，他们既属于他，又不属于他。他们如何属于他？因为是他创造了他们由之出生的那人。他们如何属于他？因为他是本性的创造者——他自己是肉身和灵魂的创造者。那么，他们如何不属于他？因为他们使自己败坏了。他们不再属于他，因为他们效法魔鬼，成了魔鬼的子孙。\par

16. 因此，主天主来到作为罪人的人那里。你们听到了这两个名称：既是\Emph{人}，又是\Emph{罪人}。作为人，他属于天主；作为罪人，他不属于天主。要把人身上的道德邪恶与他的本性区分开来。要承认那本性，以赞美创造者；要承认那邪恶，以便请医生来医治它。因此，当主说：\Quote{属于天主的，听从天主的话；你们所以不听，因为你们不是属于天主。}他并非在区分不同本性的价值，也不是在人的灵魂和身体之外，发现人有任何未被罪败坏的本性；而是预知那些将要相信的人，称他们为\Quote{属于天主}的，因为他们将藉着再生的义子地位由天主重生。这些人的话是\Quote{属于天主的，听从天主的话。}但接下来的话：\Quote{你们所以不听，因为你们不是属于天主。}是对那些不仅被罪败坏（因为这对所有人都是共同的），而且被预知不会以那种能独自从罪的束缚中释放他们的信心去相信的人说的。因此，他预知那些他如此对他们说话的人，会继续持守他们从魔鬼那里得来的东西，即他们的罪，并会死在与魔鬼相似的不敬之中；他们不会来到再生之中，在那里他们将成为天主的子女，即由那创造他们为人的天主所生。按照这预定的目的，主说了这些话；并非他在他们中间找到了任何人，要么藉着再生已经属于天主，要么按本性已不再属于天主。\par

\SourceRecord{Translated by John Gibb.；Nicene and Post-Nicene Fathers, First Series；Vol. 7.；Edited by Philip Schaff.；Buffalo, NY: Christian Literature Publishing Co.,；1888.；<http://www.newadvent.org/fathers/1701042.htm>.}

% structure: p0001 source=h1 target=section reason=page-title

\section{讲道第43篇（\BibleRef{若 8:48-59}）}

1. 在今日所读的圣福音章节中，我们从能力学习忍耐。因为我们作为仆人对主，作为罪人对正义者，作为受造物对造物主，又算得什么？然而，正如我们在恶的方面是自己如此，同样在善的方面，我们也是由祂、通过祂而成为善的。人所追求的莫过于能力。他在主基督内有莫大的能力；但让他先效法祂的忍耐，以求得能力。我们中谁若听见有人说\Quote{你有魔鬼}而能忍耐倾听呢？正如人对祂所说，祂不仅带领人得救恩，也把魔鬼置于祂的权下。\par

2. 因为当犹太人曾说\Quote{我们说你是撒玛黎雅人，并且有魔鬼，岂不正对吗？}在这两项加于祂的控诉中，祂否认了其中一项，却没有否认另一项。因为祂回答说\Quote{我没有魔鬼}。祂没有说我不是撒玛黎雅人，然而两项指控都已提出。虽然祂不以诅咒还诅咒，不以毁谤还毁谤，但祂否认一项而不否认另一项，是合宜的。这并非没有目的，弟兄们。因为撒玛黎雅人意为守护者。祂知道祂是我们的守护者。因为\Quote{保护以色列的不打盹也不睡觉}；以及\Quote{若不是主看守城池，看守的人就枉然警醒}。祂就是我们的守护者，我们的造物主。因为救赎我们属于祂，保护我们就不属于祂吗？最后，为使你们更完全知晓祂为何不应否认自己是撒玛黎雅人的隐秘原因，请回想那著名的比喻：有一个人从耶路撒冷下耶里哥，落在强盗手中，他们将他重创，把他半死不活地丢在路上。一位司铎经过，没有理会他。一个肋未人上来，也照样过去了。一位撒玛黎雅人——就是我们的守护者——上来了。\Emph{祂}走近那受伤者。\Emph{祂}施行仁慈，对那\Emph{祂}不视为外人的行了好邻舍之份。\BibleRef{路 10:30}所以，祂只是回答说没有魔鬼，而没有说不是撒玛黎雅人。\par

3. 然后在这样的侮辱之后，祂对自己的荣耀仅说了这些：\Quote{但我荣耀我的父，}祂说，\Quote{而你们羞辱我。}就是说，我不荣耀我自己，免得你们以为我傲慢。我有一位可荣耀的；若是你们认识我，就如我荣耀父一样，你们也会荣耀我。我尽我该尽的，你们却不尽你们该尽的。\par

4. \Quote{而我，}祂说，\Quote{不求自己的荣耀：有一位寻求并审判的。}祂希望我们理解为父吗？那么，祂在另一处说\Quote{父不审判任何人，却把一切审判交给了子}，而在这里祂说\Quote{我不求自己的荣耀：有一位寻求并审判的}？那么，如果父审判，祂怎么不审判任何人，却把一切审判交给了子呢？\par

5. 为了解决这一点，留心。可以通过引用一种类似的说话方式来解决。经上写着：\Quote{天主不试探任何人；}\BibleRef{雅 1:13}又经上写着：\Quote{上主你的天主试探你，要知道你是否爱祂。}\BibleRef{申 13:3}这正是争论之点。因为\Emph{天主如何不试探任何人}，而\Emph{上主你的天主如何试探你，要知道你是否爱祂？}又经上写着：\Quote{在爱里没有恐惧，但完美的爱驱除恐惧；}\BibleRef{若一 4:18}而在另一处经上写着：\Quote{敬畏上主是洁净的，存留到永远。}这里也是争论之点。因为如果完美的爱驱除恐惧，那么\Emph{敬畏上主，它是洁净的，如何能存留到永远？}\par

6. 那么，我们当了解试探有两种：一种欺骗，一种证明。就欺骗而言，\Emph{天主不试探任何人}；就证明而言，\Emph{上主你的天主试探你，为要知道你是否爱祂}。但这里又出现另一个问题：祂如何\Emph{试探为要知道}，因为对祂而言，在试探之前没有事物是隐藏的？不是因为天主无知，而是说\Emph{为要知道}，即祂要让你知道。这种说话方式既见于我们的日常谈话，也见于修辞作家。让我对我们的话语方式说几句。我们说\Quote{暗沟}，不是因为沟失去了眼睛，而是因为它隐蔽而使我们看不见它。有人说\Quote{苦羽扇豆}，即\Quote{酸的}；不是它们本身苦，而是因为它们使尝者感到苦。同样，圣经中也有这类表达。那些不辞辛劳去了解这些点的人，解决它们并不困难。所以\Quote{上主你的天主试探你，为要知道}。这是什么，\Quote{为要知道}？即祂要让你知道\Quote{你是否爱祂}。约伯对自己是未知的，但在天主面前并非未知。祂将试探者引到约伯身上，使他认识了自己。\par

7. 那么，这两种恐惧又如何呢？有一种是奴性的恐惧，另一种是纯洁的恐惧：一种是害怕遭受惩罚，另一种是害怕失去义德。那害怕受惩罚的恐惧是奴性的。害怕惩罚有什么了不起呢？最卑贱的奴隶和最残忍的强盗都会如此。害怕惩罚算不上什么大事，但爱慕义德却是大事。那么，爱慕义德的人难道没有恐惧吗？当然有；不是害怕遭受惩罚，而是害怕失去义德。我的弟兄们，你们要确知这一点，并从你们所爱的事物中引出推论。你们中有人爱钱财。你们以为我能找到一个不爱钱财的人吗？然而，从他之所爱，他就能理解我的意思。他怕损失：为什么？因为他爱钱财。他爱钱财到何种程度，就害怕失去它到何种程度。同样，也有人被发现是义德的爱好者，他内心更怕失去义德，比你们怕失去钱财更甚。这就是那清洁的恐惧——这恐惧永远长存。并不是爱将它除去或驱散，而是爱拥抱它，与它同存，并视它为伴侣。因为我们来到主面前，为要面对面看见祂。在那里，正是这纯洁的恐惧保全我们；因为这样的恐惧不会扰乱，反而使人安心。淫妇怕丈夫到来，贞洁的妻子却怕丈夫离去。\par

8. 因此，正如按一种诱惑，\Quote{天主不诱惑任何人}；但按另一种，\Quote{上主你们的天主诱惑你们}；又按一种恐惧，\Quote{爱里没有恐惧；但圆满的爱驱散恐惧}；但按另一种，\Quote{敬畏上主是纯洁的，永远长存}；同样，在此段中，按一种审判，\Quote{父不审判任何人，乃将一切审判交给子}；而按另一种，祂说：\Quote{我不求自己的光荣：有一位为我寻求并审判。}\par

9. 这一点也可以从词本身得到解决。在福音中，你们读到惩罚性的审判：\BibleQuote{不信的人已被定罪}；在另一处：\BibleQuote{时辰要来，那时在坟墓里的都要听见祂的声音，就出来；行善的复活进入生命，作恶的复活进入审判。}你们看见祂如何将审判用作定罪和惩罚。然而，如果审判总是被理解为定罪，我们岂会在圣咏中听到\BibleQuote{天主，求祢审判我}？在前一处，审判意指施加痛苦；在这里，却意指辨别。何以如此？只因那说\BibleQuote{天主，求祢审判我}的人如此解释。因为请读，看看下文。这\BibleQuote{天主，求祢审判我}是什么？岂不是接着说的\BibleQuote{求祢辨明我的案件，脱离不圣洁的国民}？因为当时说\BibleQuote{天主，求祢审判我，并辨明我的案件脱离不圣洁的国民}；同样，现在主基督说：\BibleQuote{我不求自己的光荣：有一位为我寻求并审判。}怎能说\BibleQuote{有一位为我寻求并审判}？那就是圣父，祂辨别并区分我的光荣和你们的光荣。因为你们以今世的精神自夸。但我不这样，我对圣父说：\BibleQuote{父啊，求祢以我在创世之前与祢同有的光荣，荣耀我。}\BibleRef{若 17:5}这\BibleQuote{光荣}是什么？是全然不同于人类的骄矜。这样，圣父审判。所以，\BibleQuote{审判}就是\BibleQuote{辨别}。祂辨别什么？祂儿子的光荣与仅属人的光荣；为此有话说：\BibleQuote{天主，祢的天主，用喜乐油膏了祢，胜过祢的同伴。}因为祂成为人，并非如今要与我们相比。我们作为人，是有罪的，祂是无罪的；我们作为人，从亚当继承了死亡和过犯，祂从童贞女取得了会死的肉体，却没有不义。总之，我们既不是因为愿意而生，也不是如我们所愿地活着，更不是如我们所愿地死去；但祂，在出生之前，就选择了应由谁而生；在出生时，祂使贤士们来朝拜；祂作为婴孩成长，以奇迹显示自己是天主，又以其软弱超越人。最后，祂也选择了祂死亡的方式，即被挂在十字架上，并将十字架本身印在信徒的额上，使基督徒可以说：\BibleQuote{我绝不以别的夸耀，只以我们的主耶稣基督的十字架为夸耀。}\BibleRef{迦 6:14}在十字架上，祂愿意时，就让人取下祂的身体，并离去；在坟墓中，祂愿意多久就躺了多久；当祂愿意时，就如从床榻上复活。所以，弟兄们，就祂的仆人形像而言（因为谁能恰当地谈论那另一形像呢？\BibleQuote{在起初已有圣言，圣言与天主同在，圣言就是天主}），我说，就祂的仆人形像而言，基督的光荣与其他人的光荣之间差异极大。正是论及那光荣，附魔的人听见祂说：\BibleQuote{我不求自己的光荣：有一位为我寻求并审判。}\par

10. 但是，主啊，关于祢自己，祢说了什么呢？\Quote{我实实在在告诉你们：谁如果遵行我的话，永远见不到死亡。}祢说：\Quote{你附了魔鬼。}我却召唤你们到生命那里：遵守我的话，你们就不会死。他们听了\Quote{遵行我的话的人永远见不到死亡}，就发怒，因为他们已经死于那本可逃脱的死亡。\Quote{于是犹太人说：现在我们确知你附了魔鬼。亚巴郎死了，先知们也死了；你却说：谁如果遵行我的话，永远尝不到死味。}请看圣经如何说：\Quote{见不到死亡}，即\Quote{尝到死味}。\Quote{他必见死亡——他必尝死味。}谁看见？谁尝到？人死时用什么眼睛看见？当死亡来临，连那些眼睛也闭上不能见任何东西时，如何说\Quote{他必不见死亡}？此外，死亡能用什么上颚和什么颚骨来品尝，以发现其滋味？当它夺去一切感觉时，上颚还能留下什么？但在本处，\Quote{他必看见}和\Quote{他必尝到}用来指实际情形：他必亲身经验。\par

11. 主如此说（几乎不足以说）如同一个垂死之人向将死之人说话；因为\Quote{主也掌管死亡的出口}，如圣咏所说。那么，既然祂既向注定要死的人说话，又亲自命定要死，祂的话\Quote{遵守我的话的人永远见不到死亡}是什么意思？除非主看到了另一种死亡，祂来正是为了解救我们脱离它——第二次死亡，永恒的死亡，地狱的死亡，与魔鬼及其使者同受咒诅的死亡。这才是真正的死亡；因为另一种只是迁移。那另一种死亡是什么？是离开身体——卸下重担；只要不带走另一个重担，把人拖入地狱。因此，关于那真正的死亡，主说了：\Quote{遵守我的话的人永远见不到死亡。}\par

12. 我们不要为那另一种死亡害怕，而要害怕这个。但非常可悲的是，许多人因对另一种死亡的悖逆恐惧，反而陷入了这个。有人曾对某些人说：朝拜偶像吧；因为若不这样做，你将被处死；或者像拿步高所说：若不然，你们将被投入烈火窑中。许多人害怕而朝拜了。因畏缩死亡而死去。因害怕那无法逃脱的死亡，他们落入了本可幸福逃脱的死亡，假如他们没有不幸地害怕那不可避免的事。作为人，你生来——注定要死。你往哪里去逃避死亡？你将做什么来逃避它？为使你的主在你对死亡的必要顺服中安慰你，祂甘愿屈尊而死。当你看见基督死后躺卧，你还不愿死吗？你必定要死；你无法逃脱。无论是今天还是明天，它终究要来——债必须偿还。那么，人因害怕、逃避、躲藏不让敌人发现，能得到什么呢？能免除死亡吗？不，不过能稍晚一些死。他并没有确保免债，只是请求缓期。你可以尽量拖延，那延迟的事终必来临。让我们害怕那三个人所害怕的死亡，当他们向君王说：\Quote{天主能救我们脱离那火焰；即使不然}等等。\BibleRef{达 3:16}那里有对死亡的恐惧，就是主现在所警告的，当他们说：但如果祂不愿意公开拯救我们，祂可以在暗中用胜利给我们加冕。因此，当主即将任命殉道者并亲自成为元殉道者时，祂说：\Quote{不要怕那些杀害肉身，然后不能再做什么的人。}怎么\Quote{不能再做什么}？如果某人被杀后，他的身体被抛给野兽撕咬，被飞鸟啄碎呢？残忍似乎仍能有所作为。但这是向谁行的？人已离去。身体在那里，但没有感觉。房屋躺在地上，住客已离开。因此\Quote{之后他们不能再做什么}；因为他们不能对无感觉之物做什么。\Quote{但要害怕那有权把灵魂和肉身都毁灭在地狱火里的那位。}这就是祂所说的死亡，当祂说：\Quote{遵守我的话的人永远见不到死亡。}所以，弟兄们，让我们在信德中遵守祂的话，如同那些将要达到看见的人，当我们所领受的自由达到圆满之时。\par

13. 但那些人，愤慨而仍属死亡，且预定要受永死，以侮辱回答，说：\Quote{现在我们确知你附了魔鬼。亚巴郎死了，先知们也死了。}但亚巴郎和先知并非死于主所要人理解的死亡。因为这些人死了，却仍然活着；那些其他人活着，却已经死了。因为主在某处回答撒杜塞人，当他们提出复活问题时，主亲自这样说：\Quote{论到死人复活，你们没有读过天主在荆棘丛中对梅瑟所说的话吗？祂说：我是亚巴郎的天主，依撒格的天主，雅各伯的天主。祂不是死人的天主，而是活人的天主。}那么，如果他们活着，让我们努力生活，以便死后能与他们同活。\Quote{祢把自己当作什么人？}他们又说，以致祢说\Quote{遵守我的话的人永远见不到死亡}，而祢知道亚巴郎和先知都死了？\par

14. \Quote{耶稣回答说：如果我光荣我自己，我的光荣算不了什么；光荣我的是我的父。}祂这样说是因为他们说的\Quote{祢把自己当作什么人？}因为祂将光荣归于父，正是父使祂成为天主。从这句话，\OriginalTerm{亚略派}{Arians}有时也辱骂我们的信德，说：看，父更大；因为无论如何，祂光荣了子。异端者，你没有读到子自己也说祂光荣父吗？如果子光荣父，父也光荣子，放下你的顽固，承认平等，纠正你的乖谬。\par

15. \Quote{是，} 然后祂说，\Quote{我的父荣耀我；你们说祂是你们的天主，你们却不认识祂。} 看，我的弟兄们，祂是如何表明天主自己就是基督之父，祂也曾向犹太人宣告过。我这样说，是因为如今又有某些异端者声称，在旧约中启示的天主不是基督的父，而是某个我不知道的恶天使的首领或其他什么。有摩尼派这样说，有玛尔基翁派这样说。或许还有其他异端者，或者没必要提及，或者我目前无法全部记起；但并非没有说过这话的人。因此，请注意，好让你们也能有反驳他们的话。主基督称那位祂们称为天主却不认识的天主为祂的父；因为如果他们认识那位天主本身，就会接纳祂的圣子。\Quote{但是}，祂说，\Quote{我认识祂。} 对那些按肉身判断的人，祂说这样的话可能显得自以为是，因为祂说\Quote{我认识祂。} 但请看下文：\Quote{如果我说我不认识祂，我就如同你们一样是个撒谎者。} 因此，不要如此防范自以为是，以致放弃真理。\Quote{但我认识祂，并且遵守祂的话。} 祂作为圣子讲说圣父的话；而祂自己就是向人说话的圣父的圣言。\par

16. \Quote{你们的祖先亚巴郎欢欣喜悦地看见我的日子；他看见了，就喜乐。} 亚巴郎的后裔、亚巴郎的造物主给亚巴郎作了一个伟大的见证。祂说，\Quote{亚巴郎欢欣喜悦，} \Quote{要看见我的日子。} 他不是害怕，而是\Quote{欢欣喜悦地看见它。} 因为在他内有那驱除恐惧的爱。\BibleRef{若一 4:18} 祂不说，他\Emph{因}看见而喜乐；而说\Quote{他欢欣喜悦，要看见。} 无论如何，他相信，怀着望德欢欣喜悦地要用理智去看见。\Quote{他看见了。} 主耶稣基督还能说什么，或者还应该说什么呢？\Quote{他看见了，} 祂说，\Quote{而且喜乐。} 我的弟兄们，谁能揭示这喜乐呢？如果那些眼睛被主打开身体眼睛的人喜乐，那么用灵魂的眼睛看见那不可言喻的光、那常存的圣言、那使虔诚者的理智目眩的光辉、那不竭的智慧、那与圣父同在的天主，并且在某个时刻降生成人却没有离开圣父怀抱的那一位——这样的人该有多么大的喜乐呢？这一切亚巴郎都看见了。因为说\Quote{我的日子}，可能不确定祂说的是什么：是主在时间中降生成人的日子，还是那没有黎明也不知黄昏的主的日子。但就我而言，我毫不怀疑亚巴郎祖先知道这一切。我将在哪里找出来呢？我们主耶稣基督的见证不应使我们满足吗？让我们假设我们找不出来，因为或许很难说出在什么意义上亚巴郎\Quote{欢欣喜悦地看见}基督的日子，\Quote{看见了，就喜乐。} 尽管我们找不出，难道真理会说谎吗？让我们相信真理，对亚巴郎应得的赏报毫不怀疑。然而，请听我此时想到的一段经文。当亚巴郎祖先派他的仆人为他的儿子依撒格寻找妻子时，他用这个誓言约束他，要忠实地履行他被命令的事，并且他自己也知道该做什么。因为那是一件大事，当时正在为亚巴郎的后裔寻求婚姻。但为了使仆人领会亚巴郎所知道的——他所渴望的不是肉身的后裔，也不是任何关于他种族的肉身之事——他对所派的仆人说：\Quote{把你的手放在我的大腿下，指着上天之主起誓。} \BibleRef{创 24:2} 上天之主与亚巴郎的大腿有什么关联？你已经了解这奥秘：大腿意指种族。那个起誓岂不是象征着上天之主将从亚巴郎的种族中降生成人吗？愚笨的人指责亚巴郎，因为他说\Quote{把你的手放在我的大腿下}。那些指责基督肉身的人，也指责亚巴郎的行为。但是，弟兄们，如果我们承认基督的肉身值得尊敬，我们就不要轻视那条大腿，而要视之为先知性的话。因为亚巴郎也是一个先知。谁先知？他自己的后裔和他的主。他指着自己的后裔说：\Quote{把你的手放在我的大腿下。} 他指着他的主附加说：\Quote{指着上天之主起誓。}\par

17. 恼怒的犹太人回答说：\Quote{你还不满五十岁，就见过亚巴郎吗？} 主说：\Quote{我实实在在告诉你们：在亚巴郎被造以前，我是。} 请衡量这些话，认识其中的奥秘。\Quote{在亚巴郎被造以前。} 要明白，\Quote{被造}指的是人的形成；而\Quote{我是}指的是天主性的本质。\Quote{他被造，} 因为亚巴郎是一个受造物。祂没有说：在亚巴郎存在以前，我存在；而是说：\Quote{在亚巴郎被造以前，} 那除了我以外没有人被造的那一位，\Quote{我是。} 祂也没有这样说：在亚巴郎被造以前我被造；因为\Quote{起初天主创造了天地；} \BibleRef{创 1:1} 并且\Quote{在起初已有圣言。} \Quote{在亚巴郎被造以前，我是。} 认识造物主——辨别受造物。说话的那一位成了亚巴郎的后裔；并且为了使亚巴郎得以被造，祂自己先在亚巴郎之前。\par

18. 因此，仿佛因对\Person{亚巴郎}最公开的侮辱，他们现在被激起了更大的苦毒。确实，在他们看来，主基督说了亵渎的话，说：\BibleQuote{在\Person{亚巴郎}以前，我是。}\BibleRef{若 8:58}\BibleQuote{于是他们拿起石头要砸他。}\BibleRef{若 8:59}这样刚硬的心还能求助于什么，除了像它自己？\BibleQuote{但耶稣}\BibleRef{若 8:59}作为人行动，作为取奴仆形体的一位，谦卑的，即将受苦，即将死亡，即将以祂的血救赎我们；不是作为祂之所是——不是作为起初的圣言，与天主同在的圣言。因为当他们拿起石头要砸祂时，如果大地裂开将他们吞没，并使地狱的居民代替石头，那将是何等大事？这在天主不算大事；但更好的是赞扬忍耐，而不是施展权能。因此，\BibleQuote{祂隐藏了自己}\BibleRef{若 8:59}离开他们，以免被石头砸死。作为人，祂从石头前逃走；但祸哉，那些天主已从他们石心逃走的人？\par

\SourceRecord{Translated by John Gibb.；Nicene and Post-Nicene Fathers, First Series；Vol. 7.；Edited by Philip Schaff.；Buffalo, NY: Christian Literature Publishing Co.,；1888.；<http://www.newadvent.org/fathers/1701043.htm>.}

% structure: p0001 source=h1 target=section reason=page-title

\section{第44篇讲道（\BibleBook{若望}{福音}第9章）}

1. 我们刚刚读完了关于生来眼瞎者的长篇读经，主耶稣使他重见光明；但若我们试图处理整个读经，并按照我们的能力，以一种与其价值相称的方式来考量每一段，时间将不够。因此，我请求并提醒你们的爱德，不要要求我们对于意义显明的段落发表任何言论；因为停留于每一段会过于冗长。因此，我继续简要阐述这瞎子得光照的奥秘。确实，主耶稣基督所做的一切，无论是工作还是言语，都值得我们惊异和钦佩：祂的工作因为是事实，祂的言语因为是记号。如果我们思考这行为所象征的，那瞎子就是人类；因为这种盲目发生在第一个人身上，因罪而来，我们从祂那里获得起源，不仅是在死亡方面，也是在邪恶方面。因为若不信是盲目，而信德是光照，那么基督来临之时，祂能找到一个有信德的人吗？看，那位属于先知家族的宗徒说：\BibleQuote{我们从前也都在过犯中，生来就是义怒之子，如同别人一样。}\BibleRef{弗 2:3} 若是\BibleQuote{义怒之子}，那么就是报复之子，惩罚之子，地狱之子。为何说是\Quote{生来}，岂不是因为借着第一个人的犯罪，道德邪恶如同本性般扎根在我们里面吗？若邪恶如此扎根在我们里面，那么每个人生来就精神上是瞎眼的。因为若他能看见，他就无需向导。若他需要有人引导和光照他，那么他就是生来瞎眼的。\par

2. 主来了：祂做了什么？祂揭示了一个伟大的奥秘。\BibleQuote{祂向地上吐唾沫}，用唾沫和泥；因为圣言成了血肉。\BibleQuote{祂给瞎子的眼睛傅油}。傅油已经发生，但他仍然看不见。\BibleQuote{祂打发他到那叫作西罗亚的池子}。但福音的作者关心的是引起我们对这池子名字的注意；他接着说：\BibleQuote{那池子被解释为\OriginalTerm{被派遣者}{Sent}}。你们现在明白谁是被派遣的；因为若祂没有被派遣，我们谁也不能从罪恶中获释。因此，他在那个被解释为\Quote{被派遣者}的池子里洗了眼睛——他在基督内受了洗。所以，若祂在某种程度上在祂自己内为他施洗，就光照了他；当祂为他傅油时，也许祂使他成为慕道者。确实，有许多不同的方式可以阐述和处理如此一个圣事性行为的深奥意义；但这一点对于你们的爱德已经足够了。你们听到了一个伟大的奥秘。问一个人：你是基督徒吗？如果他是一个外教人或犹太人，他的回答是：我不是。但若他说：我是；你再问他：你是慕道者还是信徒？如果他回答：慕道者；那么他已被傅油，但尚未洗净。但如何傅油？询问他，他会回答你。问他在谁内相信。就他作为慕道者而言，他说：在基督内。看，我是在同时对信徒和慕道者说话。关于唾沫和泥，我说了什么？那就是圣言成了血肉。这一点连慕道者也听到了；但他们被傅油并不足以满足他们的需要；如果他们寻求光照，就让他们赶快到洗礼池去。\par

3. 现在，由于眼前读经中的某些要点，让我们快速回顾主的话以及整个读经本身，而不是以此作为论述的主题。\BibleQuote{祂经过时，看见一个瞎眼的人}；眼瞎，不是因为任何原因，而是\BibleQuote{从出生就是}。\BibleQuote{祂的门徒问祂说：辣彼，是谁犯了罪，这个人，还是他的父母，使他生来瞎眼？} 你们知道\OriginalTerm{辣彼}{Rabbi}就是老师。他们称祂为老师，因为他们愿意学习。无论如何，他们以老师身份向主提出了这个问题：\Quote{是谁犯了罪，这个人，还是他的父母，使他生来瞎眼？} 耶稣回答说：\BibleQuote{不是这个人犯了罪，也不是他的父母}使他生来瞎眼。这是什么意思？如果没有人是无罪的，难道这瞎子的父母没有罪吗？他自己要么生来没有原罪，要么在生平中没有犯过罪吗？因为他的眼睛闭着，他的情欲就失去警觉了吗？瞎子干多少坏事？一个邪恶的心智，即使眼睛闭着，会停止作恶吗？他看不见，但他知道如何思考，也许还渴望某些因眼瞎而无法获得的东西，因此仍然在心里被洞察人心的那一位审判。那么，如果他的父母和这个人自己都有罪，为什么主说\BibleQuote{不是这个人犯了罪，也不是他的父母}，而只是就他被问的那一点而言，\Quote{使他生来瞎眼}？因为他的父母有罪，但并非因罪本身使他生来瞎眼。那么，如果不是因为父母的罪使他生来瞎眼，为什么他生来瞎眼？听老师教导。祂寻找一个相信的人，赐予他理解。祂自己告诉我们那个人生来瞎眼的原因：祂说：\BibleQuote{不是这个人犯了罪，也不是他的父母，而是为叫天主的工程在他身上彰显出来。}\par

4. 然后，接下来是什么？\Quote{我必须做派遣我来者的工作。}请看，这里就是被派遣的那一位（\OriginalTerm{史罗亚}{Siloam}），盲人在那里洗了脸。请看祂说了什么：\Quote{我必须做派遣我来者的工作，趁着白日。}你们要回想祂如何将普世的光荣归于祂所出自的那一位：因为那一位有祂自己的圣子，祂自己却没有一个人是出自祂的。但是，主啊，祢为什么说\Quote{趁着白日}呢？请听祂为何这样说。\Quote{黑夜一到，就没有人能工作了。}连祢也不能吗，主？难道黑夜有如此大的力量，连祢——黑夜原是祢的作为——在其中也不能工作？因为我想，主耶稣，不，我不是想，而是相信并确信，当天主说\BibleQuote{有光，就有了光}时，祢就在那里。\BibleRef{创 1:3}因为如果祂藉着圣言造了光，祂就是藉着祢造的；因此经上说：\BibleQuote{万物是藉着祂造的；凡受造的，没有一样不是藉着祂造的。}\BibleQuote{天主将光明与黑暗分开：光祂称为昼，黑暗祂称为夜。}\BibleRef{创 1:4}\par

5. 那黑夜是什么？当它来到时，没有人能工作。请听那白日是什么，然后你就会明白黑夜是什么。但我们如何能听到白日是什么呢？让祂亲自告诉我们：\Quote{我在世界上的时候，我是世界的光。}看，祂就是白日。让盲人在白日里洗他的眼，好使他能看见白日。\Quote{我在世界上的时候，}祂说，\Quote{我是世界的光。}然后将是某种我所不知的黑夜，那时基督不再在那里；所以没有人能工作。我还有一个疑问，我的弟兄们；请耐心听我询问。我与你们一同询问：我要与你们一同找到我所询问的那一位。我们一致同意；因为主在此处明确地宣告祂自己是白日，即世界的光。\Quote{我在世界上的时候，}祂说，\Quote{我是世界的光。}因此祂亲自工作。但祂在世界有多久？我们是否以为，弟兄们，祂那时在此，而如今已不在此？如果我们这样想，那么主升天之后，那可怕的黑夜就已经开始了，那时没有人能工作。如果那黑夜在主升天之后开始，宗徒们又怎能做了这么多工作？难道是黑夜，当圣神降临，充满所有在同一地方的人，赐给他们说万国方言的能力？难道是黑夜，当那瘸子藉着伯多禄的话——更确切地说，藉着住在伯多禄内的主的话——而痊愈？\BibleRef{宗 3:6}难道是黑夜，当门徒们经过时，病人被放在床上，好使他们至少能因门徒的影子经过而被触摸？\BibleRef{宗 5:15}然而，当主在此地时，没有人因祂的影子经过而痊愈；但祂亲自对门徒说过：\Quote{你们要做比这些更大的事。}是的，主说过：\Quote{你们要做比这些更大的事；}但不要让血肉之躯自高自大：让他们也听祂说：\Quote{离了我，你们什么也不能做。}\par

6. 那么怎样？我们该如何谈论那黑夜？它何时到来，那时没有人能工作？它将是恶人的黑夜，是那些在结束时将要听到\Quote{进入那为魔鬼及其使者预备的永火里去}的人的黑夜。但这里称之为黑夜，而不是火焰或火。那么，请听为何它也是黑夜。论及某个仆人，祂说：\BibleQuote{捆起他的手脚，把他丢在外面的黑暗中。}\BibleRef{玛 22:13}那么，让人趁活着的时候工作，免得被那无人能工作的黑夜抓住。现在正是以爱德行动的信德；如果我们现在工作，那么这就是白日——基督在此。听祂的许诺，不要以为祂不在。正是祂亲自说过：\Quote{看，我同你们在一起。}多久？我们这些活着的人不要焦虑；若有可能，藉着这句话，我们就能使未来的世世代代都处于完全的安全中。\Quote{看，}祂说，我天天同你们在一起，直到今世的终结。\BibleRef{玛 28:28}那个由太阳的运转所完成的白日只有几个时辰；基督临在的白日却延至世界的终结。但在生者与死者的复活之后，当祂对右边的人说：\BibleQuote{我父所祝福的，你们来吧！承受那国度；}对左边的人说：\BibleQuote{进入那为魔鬼及其使者预备的永火里去；}\BibleRef{玛 25:34}那时就是黑夜，没有人能工作，只能收回他之前所做的事。有时间工作，有时间领受；因为主要照各人的行为予以报应。\BibleRef{玛 16:27}当你活着时，要做，如果你真要做什么；因为那时那可怕的黑夜要来，将恶人包裹在它的黑暗里。但即使现在，每个不信的人，当他死时，就被接进那黑夜：在那里没有工作可做。在那黑夜中，富人在燃烧，用乞丐的手指求一滴水；他哀哭、痛苦、承认，但没有得到缓解。他甚至设法做好事；因为他对亚巴郎说：\BibleQuote{父亲亚巴郎，请打发拉匝禄到我兄弟们那里，告诉他们这里所发生的事，免得他们也来到这痛苦的地方。}\BibleRef{路 16:24}\par

7. \BibleQuote{祂说了这话，就吐唾沫在地上，用唾沫和了泥，把泥抹在他的眼上，对他说：\InnerQuote{你去史罗亚池里洗洗罢。}（史罗亚意即被派遣者。）他就去了，洗了，回来就看见了。}这些话是清楚的，我们可以略过。\par

8. \Quote{因此，邻居们和那些先前见过他的人，因为他是个乞丐，就说：这不是那坐着讨饭的人吗？有的说：是他；有的说：不是，只是像他。}他的眼睛被开启后，容貌改变了。\Quote{他说：我就是。}他的声音发出感激，以免被定为忘恩负义。\Quote{于是他们对他说：你的眼睛是怎样开的呢？他回答说：那名叫耶稣的人做了泥，抹在我的眼睛上，对我说：你到史罗亚池去洗；我去洗了，就看见了。}看，他成了恩宠的传报者；看，他宣讲福音；得到视力，他成了宣认者。那盲人作宣认，而恶人的心就烦乱了；因为他们心里没有他脸上现在所有的。\Quote{他们对他说：那个开了你眼睛的人在哪儿？他说：我不知道。}这些话中，那人的灵魂就像只傅了油却尚未看见的人一样。弟兄们，我们就这样说吧，好像他灵魂里有了那傅油。他宣讲，却不知道他所宣讲的那一位是谁。\par

9. \Quote{他们把从前瞎眼的人带到法利塞人那里。耶稣和泥开他眼睛的那天是安息日。法利塞人也又问他是怎么得看见的。他对他们说：他把泥抹在我眼睛上，我洗了，就看见了。}因此，有些法利塞人说；不是所有的，而是一些；因为有些人已经傅了油。那么，那些既看不见也未傅油的人说了什么呢？\Quote{这个人不是从天主来的，因为他不守安息日。}反倒是他守了安息日，他是没有罪的。因为精神的安息日就是没有罪。事实上，弟兄们，当天主将安息日嘱咐给我们时，正是这样告诫我们：\BibleQuote{你们不可做任何奴役的工作。}\BibleRef{肋 23:8} 这是天主嘱咐安息日时说的话：\BibleQuote{你们不可做任何奴役的工作。}现在请问前面的课程，奴役的工作是什么意思；并听主说：\Quote{凡犯罪的，就是罪的奴仆。}但这些人，如我所说，既看不见也未傅油，他们按肉体守安息日，却按精神亵渎了它。\Quote{另有人说：一个罪人怎能行这样的奇迹？}这些是傅了油的人。\Quote{他们中间就起了纷争。}这一天在白昼与黑暗之间作了划分。\Quote{他们又对瞎子说：他开了你的眼睛，你说他是什么呢？}你对他有什么感觉？你对他有什么意见？你对他有什么判断？他们想辱骂那人，使他被赶出会堂，却被基督寻见。但他坚定地表达了自己的感受。因为他说：\Quote{他是一位先知。}的确，当时他心里只是傅了油，尚且没有宣认天主子，但他说的并非虚言。因为主论到自己说：\BibleQuote{先知除了在自己的本乡外，是没有不受尊重的。}\BibleRef{玛 13:57}\par

10. \Quote{因此，犹太人不信他从前是瞎眼、后来得看见的，直到叫了那得看见者的父母来；}即那曾瞎眼、后来获得视力的人。\Quote{他们就问他们说：这是你们的儿子么？你们说他生来是瞎眼的，如今怎么能看见了呢？他父母回答说：我们知道这是我们的儿子，也知道他是生来瞎眼的；至于他如今怎么能看见，我们却不知道；是谁开了他的眼睛，我们也不知道。你们问他罢，他已经成年，让他自己说。}他确实是我们的儿子，我们本应为他作答，如同婴孩一样，因为那时他还不能为自己说话；他早已能说话，只是如今才看见；我们一向知道他生来是瞎眼的，也认识他早已会说话，只是如今才看见他得了视力；你们问他罢，好受教诲；为何要诬蔑我们呢？\Quote{他父母说这些话，是因为怕犹太人；原来犹太人已经议定，若有人宣认耶稣是基督，就把他赶出会堂。}被赶出会堂已不再是坏事。他们赶出，基督却收纳。\Quote{因此他父母说：他已经成年，你们问他罢。}\par

11. \Quote{于是他们又把那从前瞎眼的人叫来，对他说：你该归光荣于天主。}什么是\Quote{你该归光荣于天主}？就是否认你所领受的。这种做法显然不是归光荣于天主，而是亵渎祂。\Quote{他们说：我们知道这人是个罪人。于是他说：他是不是罪人，我不知道；有一件事我知道：我从前是瞎眼的，如今却看见了。他们又对他说：他向你作了什么？怎么开了你的眼睛？}他如今对犹太人的心硬感到愤慨，并像一个从瞎眼状态被带到视力的人，无法容忍那些瞎眼的人，\Quote{回答说：我已经告诉你们了，你们也听见了；为什么又要听呢？莫非你们也愿意做他的门徒么？}\Quote{你们也愿意}是什么意思？就是说我已经是了。\Quote{你们也愿意这样么？}现在我看得见了，却不斜视。\par

12. \Quote{他们就辱骂他说：你是他的门徒。}但愿这样的咒骂落在我们和我们子孙身上！因为如果你们剖析他们的心，而不是细究言词，这便是一种咒骂。\Quote{我们是梅瑟的门徒。我们知道天主曾对梅瑟说过话；至于这人，我们不知道他是从哪里来的。}但愿你们知道\Quote{天主曾对梅瑟说过话}！你们也会知道天主藉梅瑟所宣讲的。因为你们有主说：\Quote{如果你们信梅瑟，也必信我，因为他曾论及我写了书。}你们这样跟随仆人，却背向主人么？但你们连仆人也没有跟随；因为藉着他，你们本该被引导到主那里。\par

13. \Quote{那人回答说：这事可真奇怪，你们不知道他从哪里来，他却开了我的眼睛。我们知道天主不听罪人；但若有人是敬拜天主的，并遵行他的旨意，天主就听他。}他仍然像一个只被傅了油的人说话。因为天主甚至也听罪人。若天主不听罪人，那么那税吏垂目捶胸说：\Quote{主啊，可怜我这个罪人吧。}岂不是枉然？而那宣认获得了成义，正如这瞎子获得了光明。\Quote{自古未闻有人开了生来瞎子的眼睛。这人若不是从天主来的，他什么也不能做。}他以坦诚、坚定和真诚说了这些话。因为主所做的这些事，不是出于天主，又能是出于谁呢？或者，若主不住在他们内，门徒们何时能做出这样的事呢？\par

14. \Quote{他们回答说：你整个儿生在罪中。}这\Quote{整个儿}是什么意思？甚至到眼瞎的地步。但那开了他眼睛的，也要整个儿拯救他：那赐予他面容光明的主，必使他复活于自己的右边。\Quote{你完全生在罪中，还敢教训我们吗？}于是把他们赶了出去。他们曾把他当作老师；为了自己的教导问了他许多问题，却忘恩负义地把他们的老师赶走了。\par

15. 但是，弟兄们，正如我先前所说，当他们驱逐时，主接纳；他越是遭驱逐，反倒成了基督徒。\Quote{耶稣听说他们把他赶了出去，就找到了他，对他说：你信天主子吗？}现在他要洗净他心灵的面容。\Quote{他回答说：}仍像一个只被傅了油的人一样，\Quote{主，他是谁，好让我信他？}耶稣对他说：\Quote{你已经看见了他，而且现在和你讲话的就是他。}那一位就是被派遣来的；另一位是在史罗亚\OriginalTerm{史罗亚}{Siloam}——这词解释为\Quote{被派遣的}——洗他的脸。现在，终于，他心灵的面容被洗净，良心被净化，不仅承认他是人子——他从前已相信——而且现在承认他是天主子，那取了我们的肉身的，\Quote{他说：主，我信。}说\Quote{我信}还不够；你们也要看看他所信的是什么？\Quote{他就俯伏朝拜了他。}\par

16. \Quote{耶稣对他说。}现在，他就是那日，分别光与暗。\Quote{我为审判来到这世界，叫看不见的可以看见，看得见的反而瞎眼。}主，这是什么意思？你把一个沉重的探究课题放在疲乏的人身上；但求你复兴我们的力量，使我们能理解你所说的话。你来\Quote{叫看不见的可以看见}：很对，因为你是光；很对，因为你是日；很对，因为你从黑暗中拯救：每个灵魂都接受这一点，每个人都理解。后面所说的\Quote{看得见的反而瞎眼}又是什么意思呢？那么，因为你来了，那些看得见的就要瞎眼吗？且听后面的话，或许你就会明白。\par

17. 于是，\Quote{有几个法利塞人}被这些话困扰，\Quote{对他说：难道我们也瞎眼吗？}现在听听是什么触动他们：\Quote{看得见的反而瞎眼。}耶稣对他们说：\Quote{你们若是瞎眼的，就没有罪了；}其实瞎眼本身就是罪。\Quote{你们若是瞎眼的，}也就是说，如果你们自认为是瞎眼的，自称瞎眼，你们也会去找医生；\Quote{如果}这样\Quote{你们是瞎眼的，就没有罪了；}因为我来了是为除去罪。\Quote{但如今你们说：我们看得见，因此你们的罪还在。}为什么呢？因为你们说\Quote{我们看得见}，就不找医生，你们就留在瞎眼中。这就是我们刚才不明白的，当他说：\Quote{我来了，叫看不见的可以看见；}这是什么意思：\Quote{叫看不见的可以看见？}他们承认自己看不见，就去找医生，好得到视力。而\Quote{看得见的反而瞎眼：}这是什么意思：\Quote{看得见的反而瞎眼？}就是那些自以为看得见，却不找医生的人，要留在他们的瞎眼中。因此，这样的分别叫审判，正如他所说：\Quote{我为审判来到这世界，}借此他区分那些相信并宣认的人的原因，与那些骄傲、自以为看得见、因而更严重地瞎眼的人相对；正如那罪人宣认并找医生，对他说：\Quote{天主，求你为我判断，分辨我的冤情，反对不虔诚的民族，}——就是那些说\Quote{我们看得见}，而他们的罪仍在的人。但他带来的并不是那在世界的终末审判活人和死人的审判。因为就这一点他曾说：\Quote{我不审判任何人；}因为他第一次来，\Quote{不是为审判世界，而是为叫世界借着他而得救。}\par

\SourceRecord{Translated by John Gibb.；Nicene and Post-Nicene Fathers, First Series；Vol. 7.；Edited by Philip Schaff.；Buffalo, NY: Christian Literature Publishing Co.,；1888.；<http://www.newadvent.org/fathers/1701044.htm>.}

% structure: p0001 source=h1 target=section reason=page-title

\section{\WorkTitle{论道 45（\BibleRef{若 10:1-10}）}}

1. 吾主向犹太人的讲论，是从那个生来瞎眼而复明的人开始的。故此，你们的爱德应当知道并留意，今日的读经与那篇讲论是交织在一起的。因为当主说：\BibleQuote{我为审判来到这个世界，叫那看不见的可以看见，看得见的反成瞎子}——这经文在宣读时，我们已按自己所能力解说了——有些法利塞人就问：\BibleQuote{难道我们也瞎了吗？}他回答说：\BibleQuote{你们若是瞎了，就没有罪了；但如今你们说\InnerQuote{我们看见}，所以你们的罪还在。}在这些话之后，他又加上了我们今天读经时听到的话。\par

2. \BibleQuote{我实实在在告诉你们：那不从门进入羊栈，却从别处爬进去的，便是贼，是强盗。}因为他们声称自己并不瞎；但他们只能借着做基督的羊才能看见。他们既然是跟光明作对的贼，又怎能拥有光呢？因此，主耶稣因他们虚妄、骄傲、不可救药的狂妄，就接续了这些话，其中也给了我们有益的教训，只要我们铭记在心。因为有许多人，按此世的习惯，被称为好人——好男、好女、无辜，并似乎遵守了法律所吩咐的：孝敬父母，不犯奸淫，不杀人，不偷盗，不作伪证，并遵守法律所要求的一切——但他们不是基督徒；而且多半像这些人一样，自夸地问：\BibleQuote{难道我们也瞎了吗？}但正是因为他们所做的一切，不知道应以什么为目的，所以是徒劳的。主在今天读经中，将自己羊群的比喻和进入羊栈的门摆了出来。外邦人也许会说，我们生活得好。如果他们不从门进入，他们所夸耀的又有什么益处呢？因为行善原是为使每人能得永生：因为谁若没有得着永生，行善又有什么益处呢？那些或出于盲目不知道正生活的目的，或出于骄傲藐视它的人，甚至不能说是生活得好。但除非他认识那生命，就是基督，并从门进入羊栈，没有人能有永远活着的确切希望。\par

3. 这样的人，多半企图劝人生活得好，却又不做基督徒。他们想从别处爬进去，为偷盗、杀害，而不像牧人那样保全和拯救。于是曾有一些哲学家，对德行与恶行作了许多精微的讨论，划分、定义、推演至最缜密的推理过程，著作等身，鼓动舌簧炫耀其智慧；他们甚至敢于对人说：\Quote{跟随我们，信守我们的学派，你就能幸福地生活。}但他们没有从门进入：他们只想毁灭、杀害、谋杀。\par

4. 关于这些人我还能说什么呢？看，法利塞人自己惯于阅读，在他们所读的经文中，他们曾回响着基督，盼望他的来临，但当他来临时却不认识；他们竟自夸是看见的人，即智者之列，却否认基督，不从门进入。因此，这样的人若偶然诱惑了什么人，乃是为使他们被宰杀、被杀害，而不是带他们进入自由。让我们也撇下这些人，来看那些以基督之名自夸的人，看看他们是否也算从门进入。\par

5. 的确有无数的人不仅自夸看见，而且想要显得是被基督光照的；然而他们是异端者。难道他们也从门进入了吗？当然不是。\OriginalTerm{萨贝利乌斯}{Sabellius}说：那为圣子的就是圣父自身；但若是圣子，就没有圣父。那断言圣子就是圣父的人不从门进入。\OriginalTerm{亚略}{Arius}说：圣父是一回事，圣子是另一回事。他若说\Quote{另一位}，就说得对；但却说\Quote{另一物}。因为当他说\Quote{另一物}时，就与他亲耳听见他说\BibleQuote{我与圣父原是一体}的那位相矛盾。因此他也不从门进入；因为他所宣讲的基督是他自己捏造的，而非真理所宣告的基督。你有名，却没有实。基督是某物的名称；你若想从这名得益，就应把握那物本身。另有一个，不知从何处来的，与\OriginalTerm{福提诺}{Photinus}一同说：基督只是人，不是天主。他不从门进入，因为基督既是人又是天主。但我何必多举例子，列举异端者的许多虚妄呢？要抓住这一点：基督的羊栈就是天主教会。凡要进入羊栈的，让他从门进入，让他宣讲真正的基督。不仅要宣讲真正的基督，还要寻求基督的光荣，而非自己的光荣；因为许多人因寻求自己的光荣，反而分散了基督的羊，而不是聚集。因为主基督是低矮的门：凡从这门进入的，必须谦卑自己，才能无伤地进去。但那不谦卑、反而高抬自己的人，想要爬过墙；而爬过墙的人，只会被高举而堕落。\par

6. 然而，主耶稣一直用隐晦的语言说话；至今尚未被理解。祂提到门，提到羊栈，提到羊：祂陈述了这一切，但尚未解释。那么，让我们继续读下去，因为祂即将说到那些话，在那些话中，祂也许认为适合对我们所听的话作一些解释；藉着这解释，祂或许也能使我们理解祂未解释的事。因为祂将清楚的话赐给我们作食粮，将隐晦的话赐给我们作操练。\Quote{那不从门进入羊栈，却从别处爬进去的。}祸哉，那可怜的人，他必定跌倒！那么，让他谦卑吧，让他从门进入；让他走在平地上，他就不致绊倒。\Quote{那个人，}祂说，\Quote{是贼，是强盗。}他觊觎别人的羊，想将其占为己有——即是说，藉着偷窃据为己有，目的不是拯救，而是杀害。因此他是贼，因为把别人的东西称为自己的；是强盗，因为他偷了又杀。\Quote{但那从门进去的，才是羊的牧人；看门的给他开门。}关于这看门的，我们将在亲耳听到主自己说明什么是门、谁是牧人之后再去查考。\Quote{羊听他的声音；他按名字呼唤自己的羊。}因为祂将他们的名字写在生命册上。\Quote{他按名字呼唤自己的羊。}因此，宗徒说：\Quote{主认识属于祂的人。}\BibleRef{弟后 2:19}\Quote{他领他们出来。当他把自己的羊放出来时，他在前面走，羊跟着他，因为认得他的声音。羊却不跟生人，反而逃避他，因为不认得生人的声音。}这些话是隐晦的，充满可探究的主题，孕育着圣事性的标记。那么，让我们跟随，并听导师如何打开这些晦涩之处；也许藉着祂的开启，祂会使我们进入。\par

7. \Quote{耶稣给他们讲了这个比喻，他们却不明白祂所说的。}或许我们也不明白。那么，在我们甚至能理解这些话之前，我们与他们之间有什么区别呢？区别在于：我们这边在敲门，好让门为我们开启；而他们却因否认基督，拒绝进入以得救恩，宁愿留在外面被毁灭。就这一点而言，当我们以虔诚的心听这些话，在我们理解之前就相信它们是真实而神圣的，我们就与这些人相距甚远。因为当两个人听福音的话，一个不虔敬，一个虔敬，而其中有些话两人可能都不理解时，不虔敬者说：它什么也没说；虔敬者说：它说的是真理，它所说的都是好的，只是我们不明白。后者因相信，现在就在敲门，好配得开启，如果他持续敲的话；但前者仍听到这些话：\Quote{你们若不相信，就不得明白。}我为什么提醒你们注意这点呢？正是为了这个原因：当我尽己所能解释这些晦涩的话时，或者因为它们极其深奥，我自己未能达到理解，或者缺乏解释我所理解之事的能力，或者每个人都如此迟钝，即使我在解释，他们也无法跟上我，但人不应对自己绝望；而应持守信德，行走在道路上，并听宗徒说：\Quote{若你们在什么事上存有异见，天主也必指示你们。但我们已经达到什么程度，就该照此行走。}\BibleRef{斐 3:15}\par

8. 那么，让我们从听祂解释我们所听到祂提出的事开始。\Quote{于是耶稣又对他们说：我实实在在告诉你们，我就是羊的门。}看，祂开启了在先前描述中关闭的门。祂自己就是门。我们已经认识它了；让我们进入，或欢喜我们已在其中。\Quote{凡在我以前来的，都是贼，是强盗。}这是什么意思，主？\Quote{凡在我以前来的}？难道您没来吗？但请理解：我说\Quote{凡在我以前来的}，当然是指不包括我自己。那么，让我们回想。在祂来临之前，先知来了：他们是贼和强盗吗？断乎不是。他们不是离开祂而来，而是与祂同来。当祂将要来时，祂派遣了先驱，但保留了祂使者们的心。你想知道他们是与祂同来的吗？祂本身是永远存在的。当然，祂在预定的时间取了人性。但\Quote{以前}是什么意思？\Quote{在起初已有圣言。}\BibleRef{若 1:1}因此，那些带着天主圣言而来的人是与祂同来的。祂说：\Quote{我是道路、真理、生命。}\BibleRef{若 14:6}如果祂是真理，那么那些真诚的人就是与祂同来的。因此，凡离开祂的人，都是\Quote{贼和强盗}，也就是说，是来偷窃和杀害的。\par

9. \Quote{羊群却不听他们。} 这是一个更重要的论点，\Quote{羊群却不听他们。} 在我主耶稣基督来临之前，当祂谦卑地降生成人时，已有义人在祂之前，以同样的方式相信那将要来的那一位，正如我们相信那已经来的一位。时代不同，但信德不变。因为动词本身也随着时态而变化，有不同的词形变化。\Quote{祂将要来}有一个声音，\Quote{祂已经来了}有另一个声音：在\Quote{祂将要来}和\Quote{祂已经来了}之间，声音有变化；然而同样的信德将两者联合起来——既包括那些相信祂将要来的人，也包括那些相信祂已经来了的人。确实，在不同时代，但通过同一道信德之门，即通过基督，我们看到两者都进入了。我们相信主耶稣基督由童贞女所生，祂降生成人，受难，复活，升天：所有这一切，正如你们听到过去时态的动词一样，我们相信已经成就。在那信德中，那些先祖们也与我们一同有份，他们相信祂将由童贞女所生、将受难、将复活、将升天；因为宗徒指着这样的人说：\BibleQuote{但我们既然有同样信德的精神，根据所记载的：\InnerQuote{我信了，因此我说了话；}我们也是信了，因此说话。} \BibleRef{格后 4:13} 先知说：\BibleQuote{我信了，因此我说了话；}宗徒说：\BibleQuote{我们也信了，因此说话。} 但为使你们知道他们的信德是同一个，请听他说的：\BibleQuote{我们有同样信德的精神，我们也信了。} 同样在另一处：\BibleQuote{弟兄们，我不愿意你们不知道，我们的先祖都曾在云柱下，都从海中经过，都曾在云中和海中受洗归于梅瑟，都吃过同样的神粮，都喝过同样的神饮。} 红海象征圣洗；梅瑟，他们经过红海的领袖，象征基督；经过的百姓，象征信徒；埃及人的死亡，象征罪的消灭。在不同的标记下，有同样的信德。不同的标记就像不同的词语（动词）；因为动词通过时态改变声音，而动词无非就是标记。它们之所以是词语，是因为它们所象征的意义：从词语中除去意义，它就变成了无意义的声音。因此，所有都成了标记。难道那些使用这些标记的人，以及那些预言我们所相信之事的人，不也是有同样的信德吗？确实是的：但他们相信那些事将要发生，而我们相信它们已经发生。同样，他也说：\BibleQuote{他们都喝过同样的神饮；} \BibleQuote{同样的神饮，}因为不是同样的物质（饮）。因为他们喝的是什么？\BibleQuote{他们都喝了那随着他们的神石，而那磐石就是基督。} \BibleRef{格前 10:1} 那么请看，当信德保持不变时，标记是如何变化的。在那里磐石是基督；对于我们是那放在天主祭台上的基督。他们，作为同一基督的一个伟大圣事标记，喝了从磐石流出的水：我们所喝的，信徒们是知道的。如果人思想转向可见的形式，那么事物是不同的；如果转向指向理解的意义，那么他们喝了同样的神饮。因此，那时所有相信的人，无论是\Person{亚巴郎}、\Person{依撒格}、\Person{雅各伯}、\Person{梅瑟}，还是其他预言基督的圣祖或先知，都是羊，并听到了基督。他们听到了祂的声音，而不是别人的。审判官以传令官的身份出现。因为即使审判官通过传令官说话，书记也不会说传令官说了，而是说审判官说了。但是还有另外一些人，羊群没有听他们，基督的声音在他们里面没有地位——他们是流浪者，散布谎言，胡言乱语，制造虚妄，误导可怜的人。\par

10. 那么，为什么我说这是一个更重要的论点？其中有什么隐晦难懂之处？请听，我恳求你们。看，主耶稣基督亲自来了并宣讲。更确切地说，那是牧人的声音，是由牧人自己的口发出的。因为如果牧人的声音曾通过先知而来，那么牧人自己的舌头发出牧人的声音，岂不更是如此？然而，并非所有人都听了祂。但我们该怎么想？那些听的人，他们是羊吗？看？犹达斯听了，他却是狼：他跟随，但披着羊皮，他设陷阱陷害牧人。另一方面，一些钉死基督的人没有听，然而他们却是羊；因为祂在群众中看到这样的人时曾说：\BibleQuote{当你们举起人子时，那时你们将知道我就是那一位。} \BibleRef{若 8:28} 现在，这个问题如何解决？那些不是羊的人确实听了，而那些是羊的人却没有听。有些是狼的，跟随牧人的声音；而有些是羊的，却反驳它。最后，羊杀害了牧人。问题解决了；因为有人回答：但当他们不听的时候，他们还不是羊，那时他们是狼：声音，当被听到时，改变了他们，将狼转化为羊；所以，当他们成为羊时，他们听了，找到了牧人，并跟随了祂。他们将希望寄托于牧人的应许，因为他们遵守了祂的诫命。\par

11. 这个问题已有几分解答，或许能使每个人都满意。但我仍然有一个关切，我要把我所关切的事告诉你们，好使我们一同寻求，借着祂的启示，我能与你们一起配得解答。请听，是什么触动了我。主通过先知\Person{厄则克耳}斥责牧者，并在其中论到羊说：\BibleQuote{你们没有召回迷失的羊。}\BibleRef{则 34:4}祂既称它为\Emph{迷途者}，又称它为\Emph{羊}。如果它在迷途中时是一只羊，它听到谁的声音而被引入歧途？因为毫无疑问，如果它听到牧者的声音，它就不会迷失；它迷失正是因为听到了别人的声音；它听到了贼和强盗的声音。诚然，羊不听强盗的声音。祂说：\BibleQuote{那些来的}——我们要理解为\Emph{离我之外}——也就是说，\BibleQuote{那些\Emph{离我之外}来的是贼和强盗，羊没有听到它们。}主，如果羊没有听到它们，羊怎能迷失？如果羊只听到祢，而祢就是真理，那么听到真理的人当然不会陷入错误。但他们错了，却被称为羊。因为如果他们在迷途中时不被叫做羊，厄则克耳就不会说：\BibleQuote{你们没有召回迷失的羊}？它同时是迷途者又是羊，怎么可能？它是否听到了别人的声音？诚然\BibleQuote{羊没有听到它们。}因此，现在有许多人正被聚集到基督的羊栈里，从异端者成为天主教徒。他们从贼手中被救出，被交还给牧者；有时他们抱怨，厌倦那召回他们的人，也不真正认识那要杀害他们的人；然而，经过挣扎之后，那些是羊的人来了，他们认出牧者的声音，欢喜自己的到来，并为自己的迷失感到羞愧。那么，当他们在那错误状态中夸耀如同真理，确实没有听到牧者的声音，而是跟从了别人时，他们是羊，还是不是？如果他们是羊，羊怎能不听陌生人的声音？如果他们不是羊，为何要斥责那些人，对他们说：\BibleQuote{你们没有召回迷失的羊}？对于那些已经成为天主教徒、且有美好盼望的信徒，也时时发生不幸：他们被诱惑陷入错误，并在错误之后被恢复。当他们这样被诱惑，重新受洗，或在主的羊栈中结伴后又转回以前的错误时，他们是羊，还是不是？他们肯定是天主教徒。如果他们是忠信的天主教徒，他们就是羊。如果他们是羊，当主说\BibleQuote{羊没有听到它们}时，他们怎能听陌生人的声音呢？\par

12. 弟兄们，你们听到这个问题的重要性。因此我说：\BibleQuote{主认识那些属于祂的人。}\BibleRef{弟后 2:19}祂认识那些被预知的人，祂认识那些被预定的人；因为论到祂说：\BibleQuote{祂所预知的人，也预定他们与祂儿子的肖像相同，使祂在许多弟兄中作长子。而且祂所预定的人，也召叫了他们；祂所召叫的人，也使他们成义；祂所成义的人，也使他们受光荣。如果天主偕同我们，谁能反对我们呢？}再加上：\BibleQuote{祂没有怜惜自己的儿子，反而为我们众人把祂交出，怎能不也把一切与祂一同白白赐给我们呢？}但\Quote{我们}指的是谁？是那些被预知、被预定、成义、受光荣的人；关于他们接着有：\BibleQuote{谁能控告天主的选民呢？}\BibleRef{罗 7:29}因此，\BibleQuote{主认识那些属于祂的人；}他们就是羊。这些人有时不认识自己，但牧者认识他们，按照这预定，按照天主的预知，按照在世界奠基之前对羊的拣选：因为宗徒也这样说：\BibleQuote{就如祂在世界创立以前，在基督内拣选了我们。}\BibleRef{弗 1:4}那么，按照这神圣的预知和预定，有多少羊在外面，多少狼在里面！又有多少羊在里面，多少狼在外面！有多少人现在过着放荡生活，将来却会贞洁！有多少人在亵渎基督，将来却会相信祂！有多少人沉溺于醉酒，将来却会清醒！有多少人掠夺他人财产，将来却会慷慨施舍自己的财物！然而，此刻他们正在听别人的声音，跟从陌生人。同样，有多少人在内赞美，将来却会亵渎；现在贞洁，将来却行淫；现在清醒，日后却沉溺于酒；现在站立，不久却跌倒！这些人不是羊。（因为我们说的是那些被预定的人——那些主认识属于祂的人。）然而这些人，只要他们持守正道，就听从基督的声音。是的，这些人听见，其他人却不听；但按照预定，这些人不是羊，而那些人却是。\par

13. 仍然有一个问题，我认为目前可以这样解决。有一种声音——我说，牧者有一种特定的声音，羊群听不见生人的声音，而非羊群的人听不见基督的声音。这是什么话语！\BibleQuote{坚持到底的，才可得救。}\BibleRef{玛 10:22}祂自己的人对这样的声音不会无动于衷，生人却听不见：因此祂也向前者宣告，使他能恒心与祂同在到底；但若缺乏这种与祂恒久相联的人，这样的话语就依然未被听见。有人来到基督面前，一次又一次地听到各种话语，都是真的，都是有益的；而在这一切之中，也有这句话：\BibleQuote{坚持到底的，才可得救。}听见这话的人便是羊群中的一只。但也许有人听着，却带着厌恶、冷淡，当作生人的话来听。若是预定得救的人，他暂时迷失了，但并未永远失落：他必回头去听他忽略了的，去做他所听见的。因为若是预定得救的人，那么他的迷失和将来的归正都已为天主所预知：若他迷失了，必会回来听牧者的声音，并跟随那说\BibleQuote{坚持到底的，才可得救}的那一位。弟兄们，这是好的声音；真实而似牧者，正是义人帐幕中的救恩之声。因为听基督容易，赞美福音容易，鼓掌称赞宣讲者容易；但坚持到底，却是听牧者声音的羊群所特有的。试炼临到你，你要坚持到底，因为试炼不会持续到终了。而你所要坚持的终点是什么？直到你抵达你道路的终点。因为只要你还不听基督，祂就是你在道路上的对头，也就是在这必死的生命中。祂说什么？\BibleQuote{趁你还在路上，快与你的对头和好。}\BibleRef{玛 5:25}你听见了，相信了，和好了。若你曾是仇敌，就和好。若你有了和解的机会，就不要继续争吵。因为你不晓得你的道路何时终结，但祂知道。如果你是羊，并坚持到底，你必得救：因此祂自己的人不鄙视那声音，而生人听不见它。按我的能力，如同祂赐予我的力量，我已向你们解释或与你们一同思考了一个非常深奥的主题。若有人未能完全明白，愿他保持虔敬，真理必会显明；而那些已明白的人，不可因更快而向较慢者夸口，免得在夸口时偏离正轨，而较慢者却更易达到目标。但让我们都受祂引导，向祂说：\Quote{主啊，求祢引领我走在祢的路上，我必在祢的真理中行走。}\par

14. 因此，借着主所解释的，即祂自己就是门，让我们进入祂所陈明但未解释的真理。确实，谁是牧者，虽然祂在今天我们诵读的经文中没有告诉我们，但在接下来的经文中祂非常清楚地告诉我们：\BibleQuote{我是善牧。}即使祂没有这样说，我们还能从祂的话中理解为谁呢？祂说：\BibleQuote{从门进去的，才是羊的牧者。看门的就给祂开门；羊也听祂的声音；祂按着名字呼唤自己的羊，并领它们出来。当祂把自己的羊领出来时，祂走在它们前面，羊也跟随祂，因为它们认得祂的声音。}因为除了那知道预定者名字的，还有谁按着名字呼唤自己的羊，并领他们进入永生呢？因此祂对门徒说：\BibleQuote{你们应当欢喜，因为你们的名字已登记在天上。}\BibleRef{路 10:20}，正因如此祂按名字呼唤他们。还有谁领他们出来，除了那除去他们罪孽的，使他们从沉重的束缚中得释放，能够跟随祂？又有谁在他们前面行，到他们将要跟随的地方去，除了那从死者中复活，不再死亡，死亡不再统治祂的，\BibleRef{罗 6:9}以及那在肉身显现时所说的：\BibleQuote{父啊，我愿意祢赐给我的人，也同我在那里}？\BibleRef{若 17:24}因此祂说：\BibleQuote{我就是门；谁若经过我进来，必得救，并且可以出入，找到牧场。}在此祂清楚表明，不仅牧者，羊也经这门进入。\par

15. 但这是什么意思：\BibleQuote{他要出入，并找到牧场}？确实，经由作为门的基督进入教会是极其善的；但离开教会，正如同一圣史\Person{若望}在他的书信中所说：\BibleQuote{他们从我们中间出去，但原不属于我们，}\BibleRef{若一 2:19} 当然是不善的。这样的出去，善牧当时不能赞许，当祂说：\BibleQuote{他要出入，并找到牧场}时。因此，不仅有一种进入，而且也有一种经由作为基督的善门而有的好出去。但那可赞美和蒙福的出去是什么呢？我可以说，确实，当我们从事某种内在的思维锻炼时，我们进入；当我们从事外在的积极工作时，我们出去。并且，正如宗徒所说，基督因着信德住在我们心中，\BibleRef{弗 3:17} 经由基督进入，就是按照那信德将自己交给思维；但经由基督出去，也是按照同一信德从事外在的工作，也就是说，在他人面前。因此，我们也在一篇圣咏中读到：\BibleQuote{人出去工作}；主自己也说：\BibleQuote{让你们的善行闪耀在人前。}\BibleRef{玛 5:16} 但我更满意的是，真理本身如同一位善牧，因此也是一位善师，在某种程度上提醒我们应当如何理解祂的话：\BibleQuote{他要出入，并找到牧场}，当祂在后续加上说：\BibleQuote{贼来，无非是为偷窃、杀害、毁灭；我来了，是为使他们获得生命，且获得更丰盛的生命。}因为在我看来，祂的意思是：他们在进入时获得生命，在离开时获得更丰盛的生命。因为没有人能经过这扇门——即经过基督——到达那将显现于眼前的永生，除非经由同一扇门——即同一基督——他进入了祂的教会，即祂的羊栈，进入那在信德中度过的暂时生命。因此，祂说：\BibleQuote{我来了，是为使他们获得生命}，即信德，那以爱德行事的信德；\BibleRef{迦 5:6} 凭着这信德，他们进入羊栈，为能生活，因为义人因信德而生活：\BibleRef{罗 1:17} \BibleQuote{且获得更丰盛的生命}，那些忍耐到底的人，经由同一扇门，即经由对基督的信德，出去；因为他们作为真信徒死去，并将获得更丰盛的生命，当他们来到牧者已先他们而去的地方，在那里他们不再死去。因此，即使在羊栈内也不缺乏牧场——因为我们可以理解\BibleQuote{并找到牧场}这句话是同时指两者，即既指他们的进入，也指他们的出去——然而只有在那里，他们才会找到真正的牧场，在那里，饥渴慕义的人将得饱饫，\BibleRef{玛 5:6} ——就是那人对他说：\BibleQuote{今天你就要同我在乐园里}的人所找到的牧场。\BibleRef{路 23:43} 至于祂自己如何是门，又如何是牧者，以至于在某种程度上祂也可以被理解为经由祂自己出入，以及谁是看门的，今天探究起来太长，按照祂自己所赐予我们的恩宠，以论文的方式展开也太长了。\par

\SourceRecord{Translated by John Gibb.；Nicene and Post-Nicene Fathers, First Series；Vol. 7.；Edited by Philip Schaff.；Buffalo, NY: Christian Literature Publishing Co.,；1888.；<http://www.newadvent.org/fathers/1701045.htm>.}

% structure: p0001 source=h1 target=section reason=page-title

\section{\WorkTitle{讲道集 46（若望福音 10:11-13）}}

1. 主耶稣正在向祂的羊群说话，向那些已经是的，也向那些将要成为的；因为当时在场的人中，既有已经是祂羊群的，也有后来才成为的。祂同样向当时在场的和将来的，向他们和我们，以及在我们之后凡将成为祂羊群的人，指明谁是那被派遣到他们中间来的。因此，众人都听他们牧人的声音说：\BibleQuote{我是善牧。}祂不会加上\Quote{善}字，除非有恶牧。但恶牧就是那些贼和强盗，或者至多是佣工。因为我们应当查考、辨别并认识祂在这里所描绘的一切角色。主已经揭示了两个要点，祂先前曾以某种隐蔽的形式提出：我们已经知道祂自己就是门，也知道祂自己就是牧人。贼和强盗是谁，在昨天的课文中已经讲明；今天我们又听到了佣工，也听到了狼。昨天还提到了看门的。因此，在善的中间有门、看门的、牧人和羊；在恶的中间有贼和强盗、佣工和狼。\par

2. 我们理解主基督是门，也是牧人；但谁该被理解为看门的呢？因为前两个祂亲自解释了；看门的祂留给我们自己去寻找。关于看门的，祂说了什么？\Quote{给他，}祂说，\BibleQuote{看门的就开门。}祂给谁开门？给牧人。祂给牧人开什么？门。门又是谁？牧人自己。如果主基督没有亲自解释，没有亲自说\BibleQuote{我是牧人}和\BibleQuote{我是门}，我们中有谁敢说基督自己既是牧人又是门呢？因为如果祂说\BibleQuote{我是牧人}，而没有说\BibleQuote{我是门}，我们就会去探求门是什么，也许在错误的见解中，仍然站在门前。祂的恩宠和仁慈向我们启示了牧人，因祂自称如此；也启示了门，因祂宣称如此；但祂却留给我们自己去寻找看门的。那么，我们该称谁为看门的呢？无论我们确定是谁，都必须小心不要认为他比门本身更大；因为在人的家中，看门的比门更大。看门的被放在门前面，而不是门放在看门的前面；因为看门的看守门，而不是门看守看门的。我不敢说任何人比门更大，因为我已经听到了门是什么；这对我来说不再未知，我不再凭自己的猜测，也没有多少纯属人类臆测的余地：天主说了，真理说了，我们不能改变那不变者所宣告的。\par

3. 因此，鉴于这个问题的深奥性质，我将告诉你们我的想法：让各人选择自己所喜欢的，但要虔诚地思考，正如经上所记：\BibleQuote{你们要善意地思念上主，并以纯朴的心寻求祂。}\BibleRef{智 1:1}也许我们应当理解主自己就是看门的：因为牧人和门在人的认知上彼此不同，正如看门和门不同；然而主却称自己既是牧人又是门。那么，为什么我们不能也理解祂是看门的呢？因为如果我们看祂的位格属性，主基督既不是牧人——按我们习惯认识和见到的牧人的方式，也不是门——因为没有工匠制造祂。但是，如果因为某种相似点，祂既是门又是牧人，我敢说，祂也是羊。诚然，羊在牧人之下；然而祂既是牧人又是羊。祂在哪里是牧人？看，这里写着；读福音：\BibleQuote{我是善牧。}祂在哪里是羊？问先知：\BibleQuote{祂被牵去如同待宰的羊。}\BibleRef{依 53:7}问新郎的朋友：\BibleQuote{看，天主的羔羊，除免世罪者。}\BibleRef{若 1:29}此外，我还要说一些更奇妙的事，根据这些相似点。因为羔羊、羊和牧人彼此友善，但羊群由牧人保护免受狮子之害；然而关于基督，祂既是羊又是牧人，经上却说：\BibleQuote{犹大支派的狮子已经得胜了。}\BibleRef{默 5:5}所有这些，弟兄们，都要从相似点来理解，而不是从位格属性。常见牧人坐在磐石上，看守交托给他们的牲畜。当然，牧人比他坐的磐石更好；然而基督既是牧人又是磐石。所有这些都通过比较而言。但如果你问我祂独有的位格属性：\BibleQuote{在起初已有圣言，圣言与天主同在，圣言就是天主。}如果你问我祂独有的位格属性：从永远到永远由圣父所生的独生子，与生者同等，万物的创造者，与圣父不变，因取人性而不变，因降生成人而成为人，人子，以及天主之子。凡我所说的这一切，不是比喻，而是事实。\par

4. 因此，弟兄们，我们不要因理解祂而感到困扰，应与某些比拟一致，即祂自己就是门，也是看门人。因为什么是门？入口的道路。谁是看门人？那开门的人。那么，那开自己的是谁，岂不是那将自己揭示于视野的？看，当主起初说到门时，我们不明白：只要我们不明白，门就是关着的：那打开门的祂自己就是看门人。所以，无需寻求其他意义，不需要；但或许有渴望。若有，就不要离开这条路，不要走出天主圣三之外。若你寻求看门人的其他化身，就想想圣神；因为圣神不会认为当门是不合适的，既然圣子认为祂自己当作门是合适的。将看门人视为或许是圣神：关于祂，主对祂的门徒说：\\BibleQuote{祂要引导你们进入一切真理。}什么是门？基督。什么是基督？真理。那么，谁开门呢，岂不是那位引导进入一切真理的？\par

5. 但我们将对佣工说什么呢？他在此处没有被列入善类。祂说：\\BibleQuote{善牧为羊舍弃自己的生命。但那佣工，不是牧人，羊不是自己的，看见狼来，就撇下羊逃跑了；狼就抓住羊，把羊群驱散。}佣工在此处并不具有好的品性，然而在有些方面是有用的；他也不会被称为佣工，若不是他从雇主那里得到工钱。那么，这既是可责备又是有需要的佣工是谁？在此，弟兄们，愿主亲自光照我们，使我们知道谁是佣工，而自己不要成为佣工。那么谁又是佣工？在教会中有些担任职务的人，关于他们，宗徒\\Person{保禄}说：\\BibleQuote{他们寻求自己的事，而不是耶稣基督的事。}\\Quote{寻求自己的事}是什么意思？就是那些不白白地爱基督，不为了天主本身而寻求天主的人；他们追求暂时的利益，贪图钱财，贪图来自人的荣誉。当一个监督喜爱这类事物，并为这类事物服事天主时，无论他是谁，他就是一个佣工，不能将自己算在子女之中。因为关于这样的人，主也说：\\BibleQuote{我实在告诉你们，他们已经获得了他们的赏报。}\\BibleRef{玛 6:5}请听宗徒\\Person{保禄}关于圣\\Person{弟茂德}所说的：\\BibleQuote{我在主耶稣内希望不久派弟茂德到你们那里去，好使我知道你们的事，而得到安慰；因为我没有一个与我有同样心思的人，会真诚地关心你们。因为众人都寻求自己的事，而不是耶稣基督的事。}\\BibleRef{斐 2:19}牧人在佣工中间哀伤。他寻找一个真诚爱护基督羊群的人，但在当时围绕他的那些人中，他一个也没找到。并不是说在基督的教会中除了宗徒\\Person{保禄}和\\Person{弟茂德}以外就没有人像兄弟一样关心羊群；而是当时在他派遣弟茂德的时候，他身边没有别的儿子，只有佣工与他在一起，\\Quote{寻求自己的事，而不是耶稣基督的事。}然而他本人，怀着对羊群的兄弟般的焦虑，宁愿派遣他的儿子，而自己留在佣工中间。在我们中间也有佣工，但唯独主分辨他们。那洞察人心的，分辨他们；然而有时我们自己也知道他们。因为主自己论到狼时也不是毫无目的地说过：\\BibleQuote{凭他们的果实，你们能认出他们来。}\\BibleRef{玛 7:16}试探使许多人受到审问，然后他们的心思就显露出来；但许多人仍未被发现。主的羊圈必须有监督者，既有子女也有佣工。但那些是儿子的监督者，就是牧人。如果他们是牧人，怎么只有一个牧人呢？除非他们都是那一个牧人的肢体，羊是属于那牧人的。因为他们也是祂本人的肢体，如同那唯一的羊；因为\\BibleQuote{祂像一只羊被牵去宰杀。}\par

6. 但要注意，连佣工也是必要的。因为在教会内，确实有许多人追求地上的利益，却仍宣讲基督，藉着他们可听见基督的声音；羊群跟随的，不是佣工，而是藉佣工说话的牧者的声音。要留心主亲自指出的那些佣工：\Quote{经师}，祂说，\BibleQuote{和法利塞人坐在梅瑟的座位上：凡他们对你们所说的，你们要行要守；但不要照他们的行为去做。}\BibleRef{玛 23:2} 祂还说了什么，岂不就是：要听牧者藉佣工发出的声音？因为他们坐在梅瑟的座位上，教导天主的法律，所以天主藉着他们教导。但如果他们愿意教导自己的事，就不要听他们，也不要效法他们。因为这样的人确实寻求自己的事，而不是耶稣基督的事；但没有一个佣工敢对基督的子民说：寻求你们自己的事，而不是耶稣基督的事。因为他在自己的恶行中不是从基督的座位上宣讲；他的恶行造成伤害，而非他所说的善言。摘葡萄，当心荆棘。我看你们已经明白了；但为了那些较迟钝的人，我要更清楚地把这些话重述一遍。我说\Quote{摘葡萄，当心荆棘}是什么意思呢？主说：\BibleQuote{人岂能从荆棘上摘葡萄，从蒺藜里摘无花果吗？} 这是千真万确的；而我说的话也是真的：摘葡萄，当心荆棘。因为有时葡萄串从葡萄根发出，却攀附在普通的篱笆上；它的枝条生长，被困在荆棘中，而荆棘结的果子并非其自己的。因为荆棘不是从葡萄树生的，而是成了葡萄枝的栖身之所。只管查考根。寻找刺根，你会看见它离葡萄树很远；寻找葡萄的起源，就会发现它从葡萄树根长出。同样，梅瑟的座位是葡萄树，法利塞人的品行是荆棘。纯正的教义藉着恶人传出来，就像篱笆上的葡萄枝，荆棘中的葡萄串。小心采摘，免得在寻找果子时划伤手；当你听到有人讲善言时，不要效法他作恶事。\BibleQuote{凡他们对你们所说的，你们要行}——摘葡萄；\BibleQuote{但不要照他们的行为去做}——当心荆棘。即使藉着佣工，也要听从牧者的声音；但你们自己不要成为佣工，因为你们是牧者的肢体。是的，连圣保禄宗徒——他曾说：\BibleQuote{我没有一个像他那样关切你们的人，因为众人都寻求自己的事，不寻求耶稣基督的事}——在另一处区别了佣工和儿子，且看他说什么：\BibleQuote{有的传基督，是出于嫉妒和纷争，也有的是出于好意：有的出于爱，知道我是为辩护福音而受委派；有的出于私心传基督，并不真诚，以为可以加重我锁链的苦楚。} 这些人就是不喜欢保禄宗徒的佣工。为何如此不喜欢？岂不是因为他们寻求暂时的事物？但请注意他接着说什么：\BibleQuote{这有何妨呢？无论如何，或是假意，或是真诚，基督总被传开了；为此我欢喜，并且还要欢喜。}\BibleRef{斐 1:15} 基督是真理：让真理藉着佣工在假意中被传讲，也藉着儿子在真理中被传讲；儿子耐心等待圣父的永恒产业，佣工却渴望并急于得到雇主的暂时报酬。至于我，让我被剥夺人世的光荣，我看到这是佣工们嫉妒的对象；然而，无论藉着佣工还是儿子的口舌，愿基督的神圣光荣得以广播，因为\BibleQuote{或是假意，或是真诚，基督总被传开了。}\par

7. 我们也看到了谁是佣工。除了魔鬼，谁是狼呢？关于佣工说了什么？\BibleQuote{他一看见狼来，就撇下羊逃走了，因为羊不是他的，他也不关心羊。} 保禄宗徒是这样的人吗？当然不是。伯多禄是这样的人吗？远非如此。其他宗徒，除了丧亡之子犹达斯，有这样的品性吗？当然没有。那么他们就是牧者了？的确，他们是。但怎会只有一个牧者？我已说过，他们是牧者，因为他们是牧者的肢体。他们在那元首内喜乐，在那元首下和谐共处，藉同一圣神生活在同一身体的联系中；因此他们都属于那唯一的牧者。那么，如果他们既是牧者而非佣工，为什么在遭受迫害时逃跑呢？主啊，求你向我们解释。在书信中，我看见保禄逃走了：他被别人用筐子从城墙上缒下，逃脱了迫害者的手。\BibleRef{格后 11:33} 那么，当狼来时，他这样撇下羊群，难道不关心羊吗？显然他是关心的，但他藉祈祷把羊群交托给坐在天堂的牧者；并且为了他们的益处，他藉逃跑保全了自己，正如他在某处所说：\BibleQuote{我留在肉身，对你们更为必要。}\BibleRef{斐 1:24} 因为众人都从牧者亲自听到：\BibleQuote{若有人在这城迫害你们，就逃到另一城去。}\BibleRef{玛 10:23} 愿主乐意向我们解释这一点！主啊，你对那些你确实愿意成为忠实牧者、并塑造成你自己肢体的人说：\BibleQuote{若他们迫害你们，就逃。} 那么，当你责备那看见狼来就逃跑的佣工时，你对他们岂不是不公正吗？我们求你告诉我们，在这问题的深处隐藏着什么意义。我们敲门，守门者——基督——必会显现。\par

8. 谁是那个看见豺狼来就逃跑的佣工？是那寻求自己的事，而非耶稣基督的事的人。他是不敢公然指责犯错者的人。\BibleRef{弟前 5:20}看，某人犯了罪——犯了重罪；本该受责备，被逐出教会：但一旦被逐出，他就会变成仇敌，策划阴谋，尽其所能在一切事上加害。目前，那寻求自己的事，而非耶稣基督的事的人，为了不失去他所追求的人间友谊的好处，并避免招致人间的仇恨，便保持沉默，不施行责备。看，豺狼已咬住了羊的喉咙；魔鬼已引诱一个信徒犯奸淫：你却缄默不语——你不发出任何斥责。哦，佣工，你看见豺狼来就逃跑了！也许他回答说：看，我在这里；我没有逃跑。你逃跑了，因为你沉默了；你沉默了，因为你惧怕了。心灵的逃跑就是恐惧。你身体站着，精神却逃跑了，这与那说\BibleQuote{我身体虽不在你们那里，心却与你们同在}的人的行为不同。\BibleRef{哥 2:5}因为那在身体上虽不在，却在书信中责备行淫的人，他怎能是精神逃跑呢？我们的情感就是我们心灵的动向。喜乐是心灵的扩张；忧伤是心灵的收缩；渴望是心灵的向前运动；恐惧是心灵的逃跑。因为当你喜乐时，心灵扩张；当你忧愁时，心灵收缩；当你热切渴望时，心灵向前；当你惧怕时，心灵逃跑。这就是佣工看见豺狼时被称为逃跑的原因。为什么？\BibleQuote{因为他不关心羊。}为什么\BibleQuote{他不关心羊}？\BibleQuote{因为他是佣工。}\BibleQuote{他是佣工}是什么意思？他寻求暂时的报酬，必不永远住在屋里。这里还有一些事情要与你们查考和讨论，但加重你们的负担是不明智的。因为我们是在向同仆们分施主的食粮；我们是在主的牧场如同羊群一样吃草，也一同被喂养。正如不可扣留必需之物，也不可用丰盛的供应过分重压我们软弱的心。所以，请你们的爱德不要因我今天没有讨论我还认为有待讨论的一切而不悦；在同一位主的名下，这篇同样的道理将在讲道日再次为我们宣读，并在他助佑下更仔细地思考。\par

\SourceRecord{Translated by John Gibb.；Nicene and Post-Nicene Fathers, First Series；Vol. 7.；Edited by Philip Schaff.；Buffalo, NY: Christian Literature Publishing Co.,；1888.；<http://www.newadvent.org/fathers/1701046.htm>.}

% structure: p0001 source=h1 target=section reason=page-title

\section{论道47（\BibleRef{若 10:14-21}）}

你们这些听我们天主圣言的人，不但心甘情愿，而且专心致志，无疑还记得我们的承诺。事实上，今天所读的福音经文与上个主日所读的是相同的；因为我们在某些密切相关的话题上停留太久，未能讨论完所有我们理应向你们的理解力所讲的内容。因此，对于已经说过和论述过的内容，我们今天不再探究，以免因不断重复而妨碍我们触及仍需讲述的内容。你们现在以主的名知道，谁是善牧，善牧如何是祂的肢体，因此牧者只有一位。你们知道谁是我们要忍受的佣工；谁是狼、盗贼和强盗，我们要提防；谁是羊，什么是门，羊和牧者都由此进入；以及我们应如何理解那看门的。你们也知道，凡不从门进来的，就是盗贼和强盗，他来，无非是为偷窃、杀害、毁灭。所有这些言论，我想，都已经充分论述过了。今天，我们应当告诉你们，按主所允许的程度（因为我们的救主\Person{耶稣基督}亲自告诉我们，祂既是牧者，也是门，善牧是从门进入的），祂是如何亲自进入的。因为如果只有从门进入的才是善牧，而祂本人就是至高的善牧，并且祂自己就是门，那么我只能这样理解：祂亲自进入祂的羊群，召唤羊跟随祂，羊出入找到草场，也就是永生。\par

那么，我继续讲下去，不再拖延。当我试图进入你们，即进入你们的心时，我宣讲\Person{基督}；如果我宣讲别的，我就是试图从别的路爬进去。因此，\Person{基督}是我通向你们的门：藉着\Person{基督}，我进入的，不是你们的房屋，而是你们的心。我是藉着\Person{基督}进入的：你们一直心甘情愿聆听的，是那在我内的\Person{基督}。那么，你们为何如此情愿听从在我内的\Person{基督}呢？因为你们是\Person{基督}的羊，是以\Person{基督}的血所赎买的。你们承认自己的身价，这身价不是我付的，而是藉着我宣讲的。祂，唯有祂，是那购买者，祂流了宝贵的血——那无罪的祂的宝血。然而祂也使祂自己人的血成为宝贵的，就是为那些人祂付了血的代价：因为如果祂没有使祂自己人的血成为宝贵，就不会有这话说：\Quote{上主视为宝贵的，是祂圣者的逝世。}同样，当祂说：\Quote{善牧为羊舍命。}祂并不是唯一这样做的人；但是，如果那些这样做的人是祂的肢体，那么只有祂自己才是那行动者。因为祂没有他们也能这样做，但他们离开祂哪里有能力？祂亲自说过：\Quote{没有我，你们什么也不能做。}但从同一源泉，我们也可以显明别人做了什么，因为宗徒\Person{若望}本人，就是宣讲你们所听的这福音的那位，在他的书信中说：\Quote{正如\Person{基督}为我们舍命，我们也应当为弟兄舍命。}\BibleRef{若一 3:16}他说：\Quote{我们应当}：祂先树立了榜样，使我们成了负债者。同一意义，在某处记载说：\Quote{你若坐下赴统治者的筵席，要仔细观察摆在面前的是什么；然后伸手，知道你也当依次准备类似的供应。}你们知道统治者的筵席意指什么：你们在那里找到\Person{基督}的身体和血；让来到这筵席的人准备好类似的供应。这样的类似供应是什么？\Emph{正如祂为我们舍命，我们也应当}，为造就他人、维护信德，\Emph{为弟兄舍命。}同一意义，祂对\Person{伯多禄}说——祂愿意使他成为善牧，不是出于\Person{伯多禄}自己，而是作为祂身体的肢体：\Quote{伯多禄，你爱我吗？喂养我的羊。}祂这样做了一次、两次、三次，令门徒悲伤。当主按祂认为必要次数询问了他，使那三次否认的人能三次承认祂，并第三次将喂养祂羊群的使命交给他时，主对他说：\Quote{你年轻时，自己束上腰带，随意往来；但到年老时，你要伸出手来，别人要给你束上腰，带你到你不愿意去的地方。}而福音作者解释了主的意思：\Quote{祂说这话，是指明他将以怎样的死来光荣天主。}因此，\Quote{喂养我的羊}正是指这一点：你应为我的羊舍命。\par

3. 如今祂说：\\BibleQuote{如同圣父认识我，我也认识圣父一样，}有谁能不明白祂的意思呢？因为祂凭自己认识圣父，而我们凭祂认识圣父。祂凭自己拥有知识，我们已经知道；我们也凭祂获得知识，我们也同样学到了，因为这也是我们从祂学到的。因为祂亲自说过：\\BibleQuote{从来没有人见过天主；只有在圣父怀里的独生圣子，祂已宣告了祂。}因此，我们也藉着祂获得这知识，即祂已向谁宣告了祂。在另一处祂也说：\\BibleQuote{除了圣父之外，没有人认识圣子；除了圣子和圣子所愿意启示的人之外，也没有人认识圣父。}\\BibleRef{玛 11:27}因此，正如祂凭自己认识圣父，我们凭祂认识圣父；同样，祂凭自己进入羊栈，我们也凭祂进入。我们曾说，藉着基督，我们有了进入你们中间的门；为什么呢？因为我们宣讲基督。我们宣讲基督，因此我们藉着门进入。但基督宣讲基督，因为祂宣讲自己；所以牧人凭自己进入。当光显示在光中被看见的其他事物时，它是否还需要其他手段使自己被看见呢？那么，光既展示其他事物，也展示自身。凡我们理解的，我们藉着理智理解：然而，除了藉着理智，我们如何理解理智本身呢？但是，人是否以同样的方式，用肉体的眼睛既看见其他事物，也看见眼睛本身呢？因为人虽然用眼睛看，却看不见自己的眼睛。肉体的眼睛看见其他事物，却不能看见自身；但理智既理解其他事物，也理解自身。正如理智看见自身，同样，基督也宣讲自己。如果祂宣讲自己，并藉着宣讲进入你们内，祂就凭自己进入你们内。祂是通往圣父的门，因为除了藉着祂，没有通往圣父的道路。\\BibleQuote{因为只有一个天主，在天主与人之间只有一个中保，就是那人基督耶稣。}\\BibleRef{弟前 2:5}许多事物藉着言语表达：我刚才所说的一切，自然是藉着言辞说的。如果我想谈论言语本身，除了使用言语，我还能怎么做呢？因此，许多事物藉着言语表达，它们与言语不同，而言语本身也只能藉着言语来表达。藉着主的帮助，我们已做了丰富的说明。那么，请记住，主耶稣基督既是门又是牧人：就呈现自己而言，祂是门；就凭自己进入而言，祂是牧人。确实，弟兄们，因为祂是牧人，祂也赐给祂的肢体成为牧人。因为\\Person{伯多禄}、\\Person{保禄}以及其他宗徒，如同所有善牧一样，都是牧人。但我们中没有人称自己为门。祂保留了——羊羣进入的途径——专属于祂自己。简而言之，\\Person{保禄}在宣讲基督时履行了善牧的职分，因为他进了门。但是当不守纪律的羊开始制造分裂，并在他们面前设立其他的门——不是为了进入共同集会，而是为了堕入分裂——其中有些人说：\\BibleQuote{我是属保禄的；}另一些人说：\\BibleQuote{我是属刻法的；}另一些人说：\\BibleQuote{我是属阿波罗的；}另一些人说：\\BibleQuote{我是属基督的：}保禄为那些说\\BibleQuote{我是属保禄的}的人感到惊恐，仿佛对羊喊道：\\Quote{可怜的人，你们往哪里去？我不是门。}他说：\\BibleQuote{保禄曾为你们被钉十字架吗？你们是受洗归于保禄的名下吗？}\\BibleRef{格前 1:12}但那些说\\BibleQuote{我是属基督的}的人，却找到了门。\par

4. 然而，关于一个羊栈和一位牧人，你们如今确实不断被提醒；因为我们已经大大地赞扬了那一个羊栈，宣讲合一，使所有的羊都藉着基督进入，没有一个跟随\\Person{多纳图}。尽管如此，主为什么说这话，其特殊原因是足够明显的。因为祂是在犹太人中间说话，并且是特别被派遣到犹太人那里，不是为了那些被非人的仇恨捆锁、执意留在黑暗中的那一类人，而是为了那民族中祂称为自己羊的人。关于这些人，祂说：\\BibleQuote{我被派遣，无非是为了以色列家迷失的羊。}\\BibleRef{玛 15:24}祂甚至在狂怒的敌人群众中认识他们，并在信德的平安中预见他们。那么，祂说\\BibleQuote{我被派遣，无非是为了以色列家迷失的羊}是什么意思呢？岂不是说祂只向以色列人民显示了祂肉身的临在？祂自己没有去到外邦人那里，而是派遣了；对于以色列人民，祂既派遣又亲自来临，以便那些证明自己是藐视者的人，因为得享目睹祂实际临在的特恩，而受到更大的审判。主亲自在那里：在那里祂选择了母亲；在那里祂愿意成胎、出生、倾流鲜血；那是祂的足迹，现在成为敬拜的对象，就是祂最后站立并从此升天的地方；但是对外邦人，祂只是派遣了。\par

5. 但或许有人想，既然祂自己并未亲临我们，而是派了使者，我们就没有听见祂自己的声音，只听见祂所派遣者的声音。绝非如此：愿这样的思想从你们心中驱除；因为祂亲自临在于祂所派遣的人中。请听祂所派遣的保禄本人；因保禄是特派为外邦人的宗徒；保禄并非以自己，而是以祂来震慑他们，说：\BibleQuote{你们愿意领受那在我内说话的基督的明证吗？} \BibleRef{格后 13:3} 也请听主自己所说的话。\BibleQuote{我还有别的羊，} 就是在外邦人中的，\BibleQuote{不属于这一羊栈，} 即以色列人民，\BibleQuote{我也必须领回它们。} 因此，即使是藉着祂仆人的服务，仍是祂而非别人领回它们。请再听：\BibleQuote{它们将要听我的声音。} 这里也看到，是祂自己藉着祂的仆人说话，祂的声音在祂所派遣的人中被听见。\BibleQuote{好使成为一牧一栈。} 这两群羊，犹如两堵墙，形成了角石。\BibleRef{弗 2:11} 如此，祂既是门又是角石：这一切都是比喻的说法，没有任何一个是字面的。\par

6. 因为我先前已说过，并曾恳切地向你们强调，那领悟的人是明智的，是的，明智的人必定领悟；但那些还未蒙理智光照的人，仍当以信德持守他们所不能理解的。基督有许多比喻性的名称，严格说并非如此。比喻地说，基督既是盘石，又是门，又是角石，又是牧人，又是羔羊，又是狮子。这样的比喻何其多，数不胜数！但若你按事物通常所见选取其严格意义，那么祂既非盘石，因为祂不是坚硬无知觉的；也非门，因为工匠不曾制造祂；也非角石，因为建造者不曾构造祂；也非牧人，因为祂不是四足牲畜的牧者；也非狮子，因为狮子属于森林野兽；也非羔羊，因为羔羊属于羊群。这一切都是比喻的说法。但祂真正是什么呢？\BibleQuote{在起初已有圣言，圣言与天主同在，圣言就是天主。} 而祂显示于人性中时又如何？\BibleQuote{圣言成了血肉，寄居在我们中间。}\par

7. 也请听以下的话。祂说：\BibleQuote{圣父所以爱我，因为我舍弃我的性命，为能再取回它来。} 祂说的这话是什么意思？\BibleQuote{圣父所以爱我，} 因为我死，为能复活。因为这句话里的\Quote{我}是带有特别强调的：祂说\BibleQuote{因为\Emph{我}舍弃，} \BibleQuote{\Emph{我}舍弃我的性命，} \BibleQuote{\Emph{我}舍弃。} 那\BibleQuote{我舍弃}是什么意思？我舍弃它。犹太人不要再夸口了：他们可以发怒，却不能有任何能力；任凭他们竭力发怒吧；若我不愿舍弃我的性命，他们的一切怒气有何效果？因祂一句话，他们就仆倒在地：当人问\BibleQuote{你们找谁？} 他们回答\BibleQuote{耶稣；} 祂对他们说\BibleQuote{我就是，}他们就仆倒在地。那些在基督临死时因祂一句话就仆倒的人，在祂审判时听见祂的声音将如何？\BibleQuote{我，我，} 我说，\BibleQuote{舍弃我的性命，为能再取回它来。} 不要让犹太人夸口，仿佛他们得胜了；是祂自己舍弃了性命。祂说\BibleQuote{我躺下安睡，} 祂在别处说。你们知道那圣咏：\BibleQuote{我躺下安睡，又醒起，因上主扶持了我。} 那\BibleQuote{我躺下}是什么意思？因为我喜悦这样行。\BibleQuote{我躺下}意指什么？我死了。祂岂不正是安睡，因为祂乐意时就从坟墓起来，如同从床上起来？但祂喜爱将光荣归于圣父，为激励我们光荣我们的造物主。因为祂加上\BibleQuote{我醒起，因上主扶持了我；} 你以为这里显示了祂能力的某种不足，以致祂有权柄死，却没有权柄复活？确实，若不仔细思考，这话似乎如此。\BibleQuote{我躺下安睡；} 就是说，我如此行，因为我喜悦。\BibleQuote{我醒起：} 为何？\BibleQuote{因为上主扶持我。} 那么如何？祢自己没有能力复活吗？若祢没有能力，祢就不会说\BibleQuote{我有权舍弃我的性命，也有权再取回它来。} 但是，为显示不仅圣父使圣子复活，圣子也使自己复活，请听祂在福音另一处所说的话：\BibleQuote{你们拆毁这座圣殿，三天之内我要把它重建起来。} 福音作者补充说：\BibleQuote{但祂所说的，是指祂身体的圣殿。} 因为只有那死去的才恢复生命。圣言不是可死的，祂的灵魂也不是可死的。如果连你们的灵魂都不死，主的灵魂又岂能受制于死亡？\par

8. 你将说：我如何知道自己的灵魂不会死呢？你自己不杀害它，它就不会死。你问：我怎能杀害自己的灵魂呢？暂且不说别的罪，\Quote{说谎的口，杀害灵魂。}\BibleRef{智 1:11} 你说：我如何能确信它不会死呢？请听主亲自向祂的仆人保证的话：\Quote{不要害怕那些杀害身体，此后不能再做什么的人。} 祂用最明白的话说了什么？\Quote{要害怕那有权柄把灵魂和身体都投入地狱里毁灭的。} 这里你看到它既死又不死的事实。什么是它的死？对你的肉身来说，什么是死？肉身之死，是失去它的生命；灵魂之死，是失去它的生命。你肉身的生命是你的灵魂；你灵魂的生命是你的天主。正如肉身因失去作为其生命的灵魂而死亡，同样，灵魂因失去作为其生命的天主而死亡。如此，灵魂确实是不死的。明显是不死的，因为它在死亡时仍活着。因为宗徒论及奢华享乐的寡妇所说的话，也可用在失去天主的灵魂上：\Quote{她虽活着，却是死的。}\BibleRef{弟前 5:6}\par

9. 那么，主如何舍弃祂的性命（灵魂）呢？弟兄们，让我们更仔细地查考此事。时间不像主日那样紧迫：我们有余暇，而今天也聚集来聆听天主圣言的人将获益。祂说：\Quote{我舍弃我的性命。} 谁舍弃？舍弃什么？基督是什么？圣言和人。不是仅仅肉体的人，而是如人由肉体和灵魂组成，同样，在基督内有完整的人性。因为祂不会只取了较低的部分而留下较高的部分，因为人的灵魂显然优于身体。既然基督内有了完整的人性，那么基督是什么？我再重复：圣言和人。圣言和人是什么？圣言、灵魂和肉体。请记住这一点，因为在这方面也不乏异端分子，他们不久前被逐出公教真理，但仍像不从门进入的盗贼和强盗，在羊圈周围设下陷阱。这些异端分子被称为阿波林派，他们竟敢教条式地断言基督只是圣言和肉体，主张祂没有接受人的灵魂。然而他们中有些人不能否认基督内有灵魂。看他们难以容忍的荒谬和疯狂。他们要让祂拥有一个非理性的灵魂，却否认祂有理性的灵魂。他们允许祂有一个纯粹动物性的灵魂，却剥夺了祂的人性灵魂。但他们因失去自己的理性而剥夺了基督的理性。我们则不然，我们是在公教信仰中滋养和建立的。因此，在此场合，我要提醒你们的爱德，如同在以前的讲道中，我们已给予你们充分的教导，反对撒伯流派和亚流派——撒伯流派说圣父与圣子同一，亚流派说圣父是一位，圣子是另一位，好像圣父和圣子不是同一性体——同样，如果你们记得当，也反对佛提诺异端，他们断言基督只是一个人，没有天主性；反对摩尼派，他们主张基督只是天主，没有真正的人性；在此场合，当谈到灵魂时，我们也要给你们一些教导，以反对阿波林派，他们说我们的主耶稣基督没有人的灵魂，即理性的、有智力的灵魂——我是指那灵魂，作为人，我们藉此与禽兽有别。\par

10. 那么，我们的主在此说\Quote{我有权舍弃我的灵魂（生命）}，是什么意思呢？谁舍弃自己的灵魂，又再取回它？是作为圣言的基督这样做吗？还是祂所拥有的人的灵魂自己舍弃并重新取回自身的存在？或是祂的肉体本性舍弃生命又再取回它？让我们筛选我所提出的三个问题，选择符合真理标准的那一个。因为如果我们说天主的圣言舍弃了自己的灵魂，又再取回它，我们就必须担心一种邪恶思想潜入，有人会对我们说：那么，有一个时期那个灵魂是与圣言分离的，并且在祂取了那灵魂之后，有一个时期祂是没有灵魂的。的确，我看见圣言曾一度没有人的灵魂，但那只是在\Quote{在起初已有圣言，圣言与天主同在，圣言就是天主}的时候。但从圣言成了血肉，居住在我们中间，人性被圣言所取，即我们的整个本性——灵魂和肉体——以来，祂的受难和死亡除了使身体与灵魂分离外，还能做什么呢？它并没有使灵魂与圣言分离。因为如果主死了，是的，因为祂死了（祂为我们在十字架上死了），毫无疑问，祂的肉体呼出了那作为其生命的东西：灵魂暂时离开了肉体，虽然注定要因自己的返回而使肉体复活。但我不能说灵魂与圣言分离了。祂对强盗的灵魂说：\Quote{今天你就要同我在乐园里。}\BibleRef{路 23:43} 祂没有离弃那强盗有信德的灵魂，难道祂会离弃自己的灵魂吗？当然不会；但当主将另一个人的灵魂纳入祂的保管时，祂肯定以不可分解的结合保持了自己的灵魂。另一方面，如果我们说灵魂自己舍弃并重新取回自身，我们就陷入最大的荒谬；因为那从未与圣言分离的，也与自身不可分离。\par

11. 那么，让我们转向真实而易于理解的事。以任何一个人为例，他并非由圣言、灵魂和肉身组成，而仅仅由灵魂和肉身；让我们探究这样一个人如何交付自己的灵魂。难道没有普通人能做到吗？你可能会对我说：没有人有权交付自己的灵魂，并再取回它。但是，如果一个人不能交付自己的灵魂，宗徒\Person{若望}就不会说：\BibleQuote{正如基督为我们交付了祂的灵魂，我们也应当为弟兄交付我们的灵魂。} \BibleRef{若一 3:16} 因此，我们也能（只要我们充满祂的勇气，因为离了祂我们什么也不能做）为弟兄交付我们的灵魂。当某位圣殉道者为弟兄交付了他的灵魂时，是谁交付的？他又交付了什么？如果我们明白了这一点，我们就会理解基督所说的 \Quote{我有权交付我的灵魂} 是什么意思。你，人啊，准备为基督而死吗？他回答：我准备好了。让我换种方式重述这个问题。你准备为基督交付你的灵魂吗？他对此给我同样的回答：我准备好了，如同我先前说\Quote{你准备死吗？}时一样。那么，交付灵魂与死是同一回事。但在这种情况下，祭献是为了谁？因为所有人死时都交付了灵魂，但并非所有的人都为基督而交付。并且没有人有权重新取回他已交付的。但是基督既为我们交付了它，并且在祂乐意的时候如此行；当祂乐意的时候，祂又取回了它。那么，交付灵魂就是死。正如宗徒\Person{伯多禄}对主所说：\Quote{我愿为祢交付我的灵魂；} 也就是说，我愿为祢而死。那么，就肉身而言：肉身交付其灵魂，肉身又取回它；当然，不是凭肉身自己的力量，而是凭那居住在其内的那一位的力量。那么，肉身在断气时交付其灵魂。请看主在十字架上：祂说：\Quote{我渴；} 在场的那些人将海绵蘸了醋，绑在芦苇上，送到祂口边；祂接过后说：\Quote{完成了；} 意思是，所有关于我、在我死前尚未来临的预言都已经应验了。并且因为祂有权在祂乐意的时候交付祂的灵魂，在祂说了 \Quote{完成了} 之后，福音记载了什么？\Quote{祂低下头，交付了灵魂。} 这就是交付灵魂。惟愿你们的爱德留意这一点。\Quote{祂低下头，交付了灵魂。} 谁交付了什么？祂交付了灵魂；祂的肉身交付了它。肉身交付它是什么意思？肉身将它发出、呼出。因为当灵魂与肉身分离时，我们就被说成断气。就像离开父辈的土地是被流放，偏离道路是迷失；所以与灵魂分离就是断气；那灵魂就是生命。因此，当灵魂离开肉身，肉身没有灵魂而存留时，这就是人交付了灵魂。基督何时交付了祂的灵魂？当圣言愿意的时候。因为主权权威在于圣言；其中有权决定肉身何时交付其灵魂，何时再取回它。\par

12. 如果肉身交付了其灵魂，那么基督如何交付了祂的灵魂？因为肉身不是基督。当然是这样，基督既是肉身，又是灵魂，又是圣言；然而这三者不是三个基督，而是一个。请问你自身的人性，并从你自己上升到那高于你的，即使还不能理解，至少可以相信。因为正如一个人是灵魂和身体，同样一个基督既是圣言又是人。请思考我所说的，并理解。灵魂和身体是两样东西，但是一个人；圣言和人是两样东西，但是一个基督。那么，将其应用于任何人。宗徒\Person{保禄}现在在哪里？如果一个人回答：与基督一起安息，他说的对。同样，如果有人回答：在罗马的坟墓里，他也同样正确。我得到的第一个回答指向他的灵魂，第二个指向他的肉身。然而我们不说有两个宗徒\Person{保禄}，一个在基督里安息，另一个被葬在坟墓里；虽然我们可以说宗徒\Person{保禄}活在基督里，而同一位宗徒躺在坟墓里死了。有人死了，我们说：他是一个好人，忠诚，与主同享平安；然后立刻说：让我们参加他的葬礼，把他葬在坟墓里。你将要埋葬一个你刚刚宣称与主同享平安的人；因为后一说法指的是永远开放繁荣的灵魂，而前一个说法指的是被交付于腐朽的身体。但是，由于肉身和灵魂的结合得到了\Quote{人}的名称，这个名称现在也被单独用于二者中任何一个。\par

13. 因此，当人听到主说：\BibleQuote{我舍去我的生命，我再取回它来}时，不要困惑。肉躯\OriginalTerm{肉躯}{caro}舍去它，却是靠圣言\OriginalTerm{圣言}{Verbum}的能力；肉躯再取回它，也是靠同一能力。甚至祂自己的名字——主基督，也只用于祂的肉躯。谁能证明这一点？有人问道。我们确实不仅相信天主父，也相信耶稣基督——祂的圣子、我们唯一的主；我刚才所说的包含了整体，即耶稣基督——祂的圣子、我们唯一的主。要明白整体就在这里：圣言、灵魂\OriginalTerm{灵魂}{anima}和肉躯。无论如何，你承认同样信仰所持守的，即你相信那位被钉十字架并被埋葬的\Person{基督}。\Emph{因此}，你不否认\Person{基督}被埋葬；然而被埋葬的只是祂的肉躯。因为如果灵魂在那里，祂就不会死；但如果这是真实的死亡，而且复活是真实的，那么祂先前在坟墓中是没有生命的；然而被埋葬的却是\Person{基督}。因此，与灵魂分离的肉躯也是\Person{基督}，因为被埋葬的只是肉躯。也同样从一位宗徒的话中学到这一点。他说：\BibleQuote{你们该怀有这种心情，那也是在\Person{基督耶稣}内的心情：祂虽具有天主的形像\OriginalTerm{天主的形像}{forma Dei}，却不以自己与天主同等为强夺的。}除了\Person{基督耶稣}，就祂作为圣言的本性而言，谁是与天主同在的天主呢？但请看下文：\BibleQuote{却空虚了自己，取了奴仆的形像\OriginalTerm{奴仆的形像}{forma servi}；成了人的模样，也以人的形态显现。}这人是谁？不就是同一位\Person{基督耶稣}自己吗？但如今我们在这里有了全部部分：既有在天主形像中的圣言——祂取了奴仆的形像，也有在奴仆形像中的灵魂和肉躯——这奴仆形像被天主形像所取。\BibleQuote{祂谦卑了自己，成为服从的，一直到死。} \BibleRef{斐 2:6} 在祂的死亡中，被犹太人杀害的只是祂的肉躯。因为如果祂对门徒说：\BibleQuote{不要害怕那些杀害身体，却不能杀害灵魂的}，\BibleRef{玛 10:28} 那么他们对待祂自己，又怎能超过杀害身体呢？然而在祂肉躯的被杀中，被杀的却是\Person{基督}。因此，当肉躯舍去它的生命时，\Person{基督}舍去了它；当肉躯为了复活而取回它的生命时，\Person{基督}取回了它。然而，这事不是靠肉躯的能力成就的，而是靠那取了灵魂和肉躯的祂的能力，为的是让这些事在它们（灵魂和肉躯）内得以实现。\par

14. 祂说：\BibleQuote{这是我由我父所接受的命令。}圣言不是以言语接受命令，而是在父的独生圣言中，每一命令都存留着。但当圣子被说成从圣父接受祂本质所拥有的东西时，正如经上所说：\BibleQuote{就如父在自己内有生命，同样也赐给子在自己内有生命}，\BibleRef{若 5:26} 而子本身就是生命，祂的权威并未减弱，而是彰显了祂的受生。因为父并非像对一个出生时状态不完全的儿子那样，后来才赐予恩赐；相反，对于祂所生、在完全完美中的那一位，祂在生祂时就赐予了所有恩赐。这样，祂赐予祂与自己平等，并且生祂时并没有处于不平等状态。但当主这样说时，因为光明在黑暗中照耀，黑暗却没有领会它，\BibleQuote{为了这些话，犹太人中间又起了纷争；他们中有许多人说：祂附有魔鬼，且疯了；为什么还听祂？}这是最浓厚的黑暗。另有些人说：\BibleQuote{这些话不是附有魔鬼的人说的；难道魔鬼能开瞎子的眼睛吗？}这样的人的眼睛开始被开启了。\par

\SourceRecord{Translated by John Gibb.；Nicene and Post-Nicene Fathers, First Series；Vol. 7.；Edited by Philip Schaff.；Buffalo, NY: Christian Literature Publishing Co.,；1888.；<http://www.newadvent.org/fathers/1701047.htm>.}

% structure: p0001 source=h1 target=section reason=page-title

\section{论道第48篇（\BibleBook{若望}{福音}10:22-42）}

1. 亲爱的诸位，正如我已屡次嘱咐你们的，你们应当坚定不移地记住：圣史若望不愿我们总是以奶为食，而是要以固体食物喂养。但若有人尚不能承受天主圣言的固体食物，就让他从信德的奶中汲取营养吧；他若不能理解圣言，就毫不迟疑地相信。因为信德是应得的赏报，理解是奖赏。在用心专注的劳作中，我们心灵的眼目奋力摆脱人间迷雾的污秽薄膜，得以看清天主圣言。因此，若有爱德存在，就不会拒绝劳作；因为你们知道，爱自己劳作的人感觉不到劳苦。因为对爱它的人来说，没有什么是痛苦的。如果贪婪者的贪欲能忍受如此大的辛劳，那么在我们身上，爱德又有什么不能忍受的呢？\par

2. 请听福音：\BibleQuote{在耶路撒冷有重建节。}\OriginalTerm{重建节}{Encœnia}是圣殿奉献的庆节。因为在希腊文中，\Emph{kainos}意为\Emph{新}；凡有新的奉献时，就称为\OriginalTerm{重建节}{Encœnia}。如今这个词已进入日常用语：有人穿上新衣，就被说成\OriginalTerm{重新祝圣}{encœniare}（即翻新或举行重建节）。因为犹太人隆重地庆祝圣殿奉献的日子；正是在那个庆节上，主说了刚才所读的话。\par

3. \BibleQuote{那时是冬天。耶稣在圣殿里，在撒罗满游廊下行走。于是犹太人围起祂来，对祂说：祢使我们的心悬疑到几时呢？祢若是基督，就明明地告诉我们。}他们并非渴求真理，而是在准备诽谤。\BibleQuote{那时是冬天，}他们心里冰冷，因为他们迟于接近那神圣的火。因为接近就是相信：相信的人接近，否认的人远离。灵魂不是靠脚移动，而是靠情感。他们对爱祂的甘饴已变得冰冷，却燃烧着加害祂的欲望。他们身在祂旁，心却远离。他们不是以相信而靠近，而是以迫害而紧逼。他们想听主说：\Quote{我是基督}，而他们对基督的想法大概只是人性的。先知们曾预言基督；但连异端者也未能领悟基督的天主性，何况那心中仍蒙着帕子的犹太人呢？\BibleRef{格后 3:15}简言之，在某处，主耶稣知道他们对基督的看法是属人的，而非属神的，便立足于人性的层面，而不是那与人性同时仍为天主的层面，对他们说：\BibleQuote{你们对基督的看法如何？祂是谁的儿子？}他们按自己的意见回答：\BibleQuote{达味的。}因为他们如此读过，也只记得这个；因为当他们读到祂的天主性时，并不明白。但主为了促使他们探求那他们因外表软弱而轻视的天主性，回答他们说：\BibleQuote{那么，达味怎么受圣神感动称祂为主，说：上主对我主说：祢坐在我右边，等我使祢的仇敌作祢的脚凳？那么，如果达味受圣神感动称祂为主，祂怎么又是达味的儿子呢？}\BibleRef{玛 22:42}祂没有否认，而是提问。愿听到这句话的人不要以为主耶稣否认祂是达味的儿子。若主基督曾如此否认，祂就不会光照那些如此称呼祂的瞎子了。因为有一天祂经过时，两个坐在路旁的瞎子喊着说：\BibleQuote{主，达味之子，可怜我们吧！}\BibleRef{玛 20:30}祂听了这话就怜悯了他们。祂站住，治好了他们，开了他们的眼；因为祂承认了这个名字。宗徒保禄也说：\BibleQuote{按肉身，祂是从达味后裔生的；}\BibleRef{罗 1:3}并在致弟茂德书中说：\BibleQuote{你要记着：耶稣基督，按照我所传的福音，是从达味后裔中复活了的。}\BibleRef{弟后 2:8}因为童贞玛利亚出自达味后裔，所以我们的主也出自达味后裔。\par

4. 犹太人向基督提出这个问题，主要是为了：若祂说\Quote{我是基督}，他们就可按他们对这名称的唯一理解（即祂是达味的后裔）来诽谤祂，说祂企图夺取王权。祂对他们的回答还有更多：他们想以祂自称为达味之子来诽谤祂。祂回答说祂是天主子。如何回答呢？请听：\BibleQuote{耶稣回答他们说：我告诉你们，你们却不相信；我因我父之名所行的工程为我作证；但你们不相信，因为你们不属于我的羊。}你们已在上面（第四十五讲）学过谁是羊：你们要成为羊。他们是因相信而成为羊，因跟随牧人而成为羊，因不轻视他们的救赎者而成为羊，因从门进入而成为羊，因出入并找到牧场而成为羊，因享受永生而成为羊。那么，祂对他们说\BibleQuote{你们不属于我的羊}是什么意思呢？那是祂看见他们预定要遭永罚，而非以祂自己的血价所赢得、得享永生。\par

5. \BibleQuote{我的羊听我的声音，我认识他们，他们也跟随我；我赐给他们永生。}这就是牧场。如果你们记得，祂先前曾说过：\BibleQuote{他必出入并找到牧场。}我们因相信而进入——我们因死亡而出离。但正如我们因信德之门进入，同样，作为信徒，我们离开肉身；因为正是从同一道门出去，我们才能找到牧场。好牧场被称为永生；那里没有枯叶——全是青翠茂盛。有一种植物常被称为长生草，只有在那里才能找到它。祂说：\BibleQuote{我赐给他们}，即我的羊，\BibleQuote{永生。}你们寻求诽谤，只是因为你们只思想现世的生命。\par

6. \BibleQuote{他们永不会丧亡：}你们可以听出言外之意，仿佛祂对他们说：你们将永远丧亡，因为你们不是我的羊。\BibleQuote{谁也不能从我手中把他们夺去。}更要留意这一点：\BibleQuote{我父赐给我的，比一切更大。}狼能做什么？盗贼和强盗能做什么？他们只摧毁那些预定遭毁灭的人。但是那些羊，宗徒说：\BibleQuote{主认识那些属于祂的人；}\BibleRef{弟后 2:19} 以及\BibleQuote{祂所预先知道的人，也预定他们；祂所预定的人，也召叫他们；祂所召叫的人，也使他们成义；祂所成义的人，也使他们得光荣；}\BibleRef{罗 8:29}——在这些羊中，没有一只会被狼抓住、被贼偷走、或被强盗杀害。祂知道为他们付出了什么，所以确信他们的数目。正是为此祂说：\BibleQuote{谁也不能从我手中把他们夺去；}论到父也说：\BibleQuote{我父赐给我的，比一切更大。}父赐给圣子什么比一切更大呢？就是作祂自己的独生子。那么，\BibleQuote{赐给}是什么意思？是受赐者先前已经存在，还是父在生育时赐给的？因为如果受赐者承受儿子名分之前已经存在，那么曾有一个时期祂存在却还不是圣子。我们绝不可认为主基督曾存在却还不是圣子。论到我们，可以这样说：曾有一个时期我们是人的儿子，却不是天主的儿子。因为我们因恩宠成为天主的儿子，而祂按其本性，因为祂是这样诞生的。但并非如此，不可说祂在诞生之前不存在；因为祂与圣父同永恒，从未未生。有智慧的，让他领悟；不领悟的，让他相信并得滋养，他终将领悟。天主的圣言永远与圣父同在，永远是圣言；因为是圣言，所以是圣子。所以，祂永远是圣子，永远平等。因为祂平等不是因成长，而是因生育；祂永远是受生的，圣父的圣子，出自天主的真天主，出自永恒者的同永恒者。但圣父不是出自圣子的天主：圣子是出自圣父的天主；因此，圣父在生育圣子时，\BibleQuote{赐给}祂作天主，在生育时赐给祂与己同永恒，在生育时赐给祂与己平等。这就是那比一切更大的。圣子如何是生命，且拥有生命？祂所拥有的，祂就是：至于你，你是一回事，你拥有另一回事。例如，你拥有智慧，但你自己就是智慧吗？简言之，因为你自己不是你拥有的东西，如果你失去你所拥有的，你就回到不再拥有它的状态：有时重新获得，有时失去。如同我们的眼睛本身没有光，它睁开，接受光；闭上，失去光。天主子之为天主并非如此；祂之为圣父的圣言并非如此；祂之为圣言，并非像声音一样消逝，而是存留在其诞生中。祂拥有智慧的方式是祂自己就是智慧，并使人有智慧；祂拥有生命的方式是祂自己就是生命，并使其他人活。这就是那比一切更大的。福音作者若望自己，当要谈论天主子时，他仰望天地；他仰望，并超越这一切。他思想天上成千上万的天使军旅；他思想，如鹰腾云，他的心智超越整个受造界：他超越一切伟大的事物，抵达那比一切更大的，并说：\BibleQuote{在起初已有圣言。}但是因为圣言所出者（即父）并非出于圣言，而圣言出于祂——圣言是祂的圣言——所以祂说：\BibleQuote{父赐给我的，}即作祂的圣言、祂的独生子、祂光明的光辉，\BibleQuote{比一切更大。}因此，祂说：\BibleQuote{谁也不能，}\BibleQuote{把我的羊从我手中夺去。谁也不能从我父手中夺去。}\par

7. \BibleQuote{从我手中，}和\BibleQuote{从我父手中。}这是什么意思？\BibleQuote{谁也不能从我手中把他们夺去，}和\BibleQuote{谁也不能从我父手中把他们夺去？}父和子共有一只手，还是我们可以说，圣子本人就是祂父的手？如果我们把手理解为能力，那么父和子的能力是同一的，因为他们的神性是同一的。但如果我们按照先知所说的意思来理解手：\BibleQuote{主的手臂向谁显示了呢？}\BibleRef{依 53:1}那么父的手就是圣子本身，这不可理解为天主具有人的形状和肢体，而是指万物都是借着祂造的。因为人们也习惯把受他们差遣办事的人称为他们的手。有时人亲手所完成的工作也被称为他的手，就如一个人认出自己写的字时，说认出自己的手。既然论到有手作为身体部分的人，尚有多种说法；那么读到没有身体形状的天主的手时，岂不更有多种理解方式？因此，在这里最好将父和子的手理解为父和子的能力；免得我们把这里父的手理解为圣子时，对圣子本身也产生肉性的想法，以为圣子也可以被类似地看作基督的手。所以，\BibleQuote{谁也不能从我父手中把他们夺去；}也就是说，谁也不能从我手中把他们夺去。\par

8. 但为使人不再犹豫，请听下文：\Quote{我与圣父原是一体。}至此犹太人尚能忍受祂；他们听见\Quote{我与圣父原是一体}，便再也无法忍受；他们顽固于自己的方式，遂诉诸石头。\Quote{他们拿起石头要砸祂。}主，因祂不受祂不愿受的苦，只受祂乐意受的苦，便在他们想要砸祂时仍对他们说话。\Quote{犹太人拿起石头要砸祂。耶稣回答他们说：我从圣父那里显给你们许多善工；你们为了哪一件善工要砸我呢？他们回答说：我们不是为了善工要砸你，而是为了亵渎，并且因为你虽是个人，却把自己当作天主。}这便是他们对祂的话\Quote{我与圣父原是一体}的回答。你们看见，犹太人懂得亚略派所不懂的。因为他们愤怒的原因正是他们感到，只有存在圣父与圣子的平等时，才能说\Quote{我与圣父原是一体}。\par

9. 但请看主如何回答他们迟钝的理解。祂见他们不能承受真理的光辉，便用言语缓和了它。\Quote{你们的法律上不是写着：我说过你们是神吗？}即，赐予你们的法律；主通常把所有圣经都称为法律，尽管别处祂更明确地区分法律与先知，如经上所说：\Quote{法律及先知到若翰为止}（\BibleRef{路 16:16}），以及\Quote{全部法律和先知都系于这两条诫命}（\BibleRef{玛 22:40}）。有时，祂也将同样的圣经分为三部分，如祂说：\Quote{凡在法律、先知和诗篇上关于我所记载的，都必须应验}（\BibleRef{路 24:44}）。但现在祂也将诗篇包含在法律之名下，其中写道：\Quote{我说过你们是神。如果那些承受天主话的人，祂称他们为神——且圣经是不能废去的——那么圣父所祝圣并派遣到世界上来的那位，你们却因为我说了\InnerQuote{我是天主子}而说：你说亵渎的话吗？}如果天主的话临到人，使他们可以被称为神，那么那位与天主同在的天主圣言，难道不是天主吗？如果因天主的话，人成了神；因共融，人成了神；那么他们借以共融的那位难道不是天主吗？如果被点燃的光是神，那么那照亮的光难道不是天主吗？如果借某种方式被救赎之火温暖就成了神，那么赐予他们温暖的那位难道不是天主吗？你接近光就被照亮，被列于天主子民中；你若离开光，就落入黑暗，被归于黑暗；但那光不会靠近，因为它从不离开自身。如果天主的话使你成为神，那么天主圣言怎能不是天主呢？因此，圣父祝圣了祂的子，并派遣祂到世界上来。也许有人会说：如果圣父祝圣了祂，那么是否曾有一段时间祂没有被祝圣？祂祝圣祂的方式如同生了祂。因为在生祂时，祂赐予祂成为圣的能力，因为祂在圣洁中生了祂。因为如果那被祝圣的以前是不洁的，我们怎能向天主圣父说：\Quote{愿你的名被尊为圣}（\BibleRef{玛 6:9}）？\par

10. \Quote{我若不做我圣父的工程，你们不必信我；但我若做了，你们纵然不信我，也要信这些工程，好叫你们知道并相信：圣父在我内，我在圣父内。}子说\Quote{圣父在我内，我在圣父内}，并不像人可以说这话。因为如果我们思想正确，我们就在天主内；如果我们生活良善，天主就在我们内：信徒因分享祂的恩宠，并被祂光照，就在祂内，祂也在我们内。但独生子并非如此：祂在圣父内，圣父在祂内，如同同等者在他同等者内。简言之，我们有时可以说：我们在天主内，天主在我们内；但我们能说：我与天主是一体吗？你在天主内，因为天主包含你；天主在你内，因为你成了天主的殿宇；但因为你是在天主内，天主在你内，你能像独生子那样说：看见我的就是看见父吗？独生子曾说：\Quote{看见我的，就是看见了圣父}；以及\Quote{我与圣父原是一体}。要认清主的特权与仆人的恩典。主的特权是与圣父平等；仆人的恩典是与救主共融。\par

11. 因此他们想要捉拿祂。但愿他们是因信德和理解而捉拿，而非出于愤怒和杀机！因为现在，我的弟兄们，当我这样说时，是软弱者想要捉拿刚强者，渺小者想要捉拿伟大者，脆弱者想要捉拿坚固者；而且正是我们——你们这些与我同质的人，以及我对你们说话的我自己——都想要捉拿基督。捉拿祂是什么意思？如果你理解了，你就捉拿了。但不是像犹太人那样：你们捉拿是为了拥有，他们却想捉拿以除掉祂。因为他们想要的捉拿是这一种，祂向他们做了什么？\Quote{祂从他们手中逃脱了。}他们未能捉拿祂，因为他们缺少信德之手。圣言成了血肉；但对圣言而言，从血肉之手解救祂自己的血肉并非难事。用心智捉拿圣言，才是对基督的正确捉拿。\par

12. \BibleQuote{祂又往约但河外去，到了若翰最初施洗的地方，就在那里住下了。有许多人来到祂那里说：若翰确实没有行过奇蹟。} 你们还记得关于\Person{若翰}所说的话：他是一盏明灯，为光作证。这些人为什么彼此说 \BibleQuote{若翰没有行过奇蹟}？他们说，\Person{若翰}没有以奇蹟著称；他没有驱逐魔鬼，没有退去热症，没有开启瞎子的眼睛，没有复活死者，没有用五饼或七饼使数千人吃饱，没有行走在海面上，没有命令风和浪。这些事\Person{若翰}一件也没有做，而他所说的一切，都是为这人作证。借着灯光，我们可以走向白昼。\BibleQuote{若翰没有行过奇蹟，但若翰关于这人所说的一切，都是真的。} 这些人与犹太人不同，他们抓住了一位留下来的人。犹太人想要抓住一位离开他们的人，而这些人却抓住了那位与他们同住的人。总之，接下来是什么？\BibleQuote{许多人在那里信了祂。}\par

\SourceRecord{Translated by John Gibb.；Nicene and Post-Nicene Fathers, First Series；Vol. 7.；Edited by Philip Schaff.；Buffalo, NY: Christian Literature Publishing Co.,；1888.；<http://www.newadvent.org/fathers/1701048.htm>.}

% structure: p0001 source=h1 target=section reason=page-title

\section{\WorkTitle{讲道集49（\BibleBook{若望}{福音} \BibleRef{若 11:1-54}）}}

1. 在主耶稣基督所行的一切奇迹中，拉匝禄的复活在宣讲中占有首要地位。但如果我们仔细思考是谁行了这奇迹，我们应当欢喜而非惊奇。一个人被那造人者复活了，因为祂是圣父的独生子，正如你们所知，万物都是藉着祂造的。如果万物都是藉着祂造的，那么祂使一个死人复活又有什么稀奇呢？每天有多少人因祂的权能被带到世上来！造人比从死人中复活更伟大。然而祂屈尊俯就，既创造又复活：创造一切，复活一些人。因为主耶稣虽行了许多这样的事，却并未全部记载下来；正如同一圣史若望自己所证实的，基督主所说的和所行的许多事都没有记载；但所选择的记载，似乎是足以使信徒得救。你们刚听到主耶稣使一个死人复活；这足以使你们知道，如果祂愿意，祂可以使所有死人复活。的确，祂将这工作留到世界末日。因为你们听说祂以一个大奇迹使一个已死了四天的人从坟墓中复活，正如祂自己说的：\BibleQuote{时候要到，凡在坟墓里的，都要听见\Emph{祂}的声音，就出来。} 祂复活了一个已经腐烂的人，但在那腐烂的尸体中仍有四肢的形状；但在末日，祂将用一句话将灰烬重新构成人的肉身。但当时只需要行一些这样的奇迹，好让我们接受祂权能的印证，信靠祂，并为那将复活得生命而非受审判的复活作准备。所以，祂确实说：\BibleQuote{时候要到，凡在坟墓里的，都要听见\Emph{祂}的声音，就出来；行过善的，复活得生命；作过恶的，复活受审判。}\par

2. 然而，我们在福音中读到主曾使三个死人复活，但愿这是有某种良好意图的。因为主的作为不单是作为，更是标记。如果它们是标记，除了其奇妙性，它们还具有某种实在的意义：而要找出这些作为的意义，比阅读或听闻它们更为困难。我们怀着惊叹聆听福音的诵读，如同亲眼目睹一件伟大的奇迹，那就是拉匝禄如何被复活。如果我们转而思考基督更奇妙的作为，每个相信的人都复活；如果我们都思考并理解那更可怕的一种死亡——每个犯罪的人都在死亡。但是人人都害怕肉身的死亡，却很少有人害怕灵魂的死亡。关于肉身之死（它迟早必然来临），人人都提防它的临近：这是他们一切劳苦的根源。人注定要死，却努力避免死亡；然而人注定要永生，却不努力停止犯罪。当他努力避免死亡时，他的努力徒劳无功，因为其结果只是暂时延缓死亡，而非逃避死亡；但如果他停止犯罪，他的劳苦将止息，他将永远活着。但愿我们能唤醒人们，也唤醒自己，使我们像人们爱那消逝的生命一样，热爱那恒存的生命！当一个人处于死亡的险境时，他什么不肯做？当刀剑悬在头顶时，人们会放弃他们所有的一切。有谁为了免遭杀戮，不是即刻交出一切？而最终他也许仍被杀了。有谁为了保全生命，不是立刻愿意失去生计，宁可行乞也不愿速死？有谁曾听到这样的吩咐：\Quote{如果你要逃命，就出海去}，而他却迟疑不前？有谁曾听到这样的吩咐：\Quote{如果你要保全生命，就去工作}，而他却继续懒惰？天主要求我们的其实很少，好使我们永远活着，而我们却忽略顺服祂。天主不是对你说：\Quote{放弃你所有的一切，好让你在劳苦中短暂地活着。} 而是说：\Quote{把你所有的给贫穷的人，好让你永远不劳苦地活着。} 贪爱这暂时生命的人（这生命既不是在他们想要的时候，也不是如他们所愿长久地属于他们）控告我们；而我们却不反过来控告自己，我们是如此懒惰，如此冷漠于获得永生（永生如果我们愿意，就会属于我们，并且当我们拥有时不会朽坏），然而我们所害怕的这个死亡，尽管我们百般不愿，却终究会临到我们。\par

3. 那么，如果主以祂恩宠与仁慈的伟大，使我们灵魂复活获得生命，以免我们永远死亡，我们便能明白，祂在肉身所复活的那三个死人，具有某种象征意义，指向那由信德所实现的灵魂复活：祂复活了会堂长还躺在屋里的女儿（\BibleRef{谷 5:41}）；祂复活了寡妇的年轻儿子，当时他正被抬出城门（\BibleRef{路 7:14}）；祂还复活了已在坟墓里四天的拉匝禄。每个人都当留意自己的灵魂：犯罪时，他就死了；罪是灵魂的死亡。但有时罪只发生在思想中。你因邪恶而感到愉悦，你同意了去行恶，你便犯了罪；那同意杀害了你：但死亡是内在的，因为邪恶的思想尚未成熟为行动。主暗示祂要复活这样的灵魂，就像祂复活那女孩一样，她尚未被抬出去埋葬，而是躺在屋里死了，仿佛罪仍隐藏着。但若你不只怀着对邪恶的愉悦，还做了那邪恶之事，那么，你便是已经把死者抬到了门外：你已经在外面，正被送往坟墓。然而，这样的一个人，主也复活了他，并将他归还给他的寡母。你若犯了罪，就悔改吧，主必使你复活，并将你归还给教会你的母亲。第三个死亡的例子是拉匝禄。这是一种严重的死亡，并且被称为恶习。因为陷入罪是一回事，形成犯罪的习惯是另一回事。那陷入罪并立即接受纠正的人，将迅速恢复生命；因为他尚未陷入习惯，尚未被置于坟墓中。但那已习惯于犯罪的人，被埋葬了，并且人们恰当地说他：\Quote{他发臭了；}因为他的品行，如同某种可怕的气味，开始声名狼藉。所有习惯于罪恶、道德沦丧的人都是如此。你对这样的人说：\Quote{不要这样做。}但何时你会被一个被如此厚土覆盖、滋生腐败、被习惯的重量所压垮的人听从呢？然而，基督的能力并非不足以使这样的人复活。我们知道，我们见过，我们每天见到人们改变最恶劣的习惯，并采纳比那些指责他们的人更好的生活方式。你曾厌恶这样的人：请看拉匝禄的妹妹本人（如果确实是那用香液傅抹主的脚，并用头发擦干那用眼泪洗过之处的女人），她获得了比她兄弟更好的复活；她从罪恶品格的沉重负担中被释放。因为她是一个著名的罪人；并且关于她曾说：\Quote{她的许多罪得了赦免，因为她爱得多。}我们见到许多这样的人，我们认识许多：但愿无人绝望，但也不可自恃。两者都是罪人。让你们不愿绝望的心转向选择那唯一可恃靠的主。\par

4. 那么，主也使拉匝禄复活了。你们已经听到他所代表的是什么类型的人；换句话说，拉匝禄的复活意味着什么。因此，现在我们再来读这段经文；既然这一课中许多部分已经清楚，我们就不作详细解说，以便更彻底地处理必要之处。\Quote{有一个病人，名叫拉匝禄，住在伯达尼，就是玛利亚和她姐姐玛尔大所住的村庄。}在上一课中，你们记得主从那些想用石头砸祂的人手中逃脱，去了约但河对岸，就是若翰施洗的地方。因此，当主住在那里时，拉匝禄在伯达尼病了，那是一个靠近耶路撒冷的村庄。\par

5. \Quote{但玛利亚就是那用香液傅抹主，并用头发擦祂脚的人，她的兄弟拉匝禄病了。因此，她的姐妹打发人去见主说：}我们现在明白她们打发人去的地方，就是主所在之处；因为祂离开了，如你们所知，在约但河对岸。她们打发使者去见主，告诉祂她们的兄弟病了。祂延迟医治，以便能够使他复活。但姐妹俩打发人带去的是什么话呢？\Quote{主，请看，祢所爱的人病了。}她们没有说\Quote{来}；因为对于所爱的人来说，只需暗示就够了。她们不敢说\Quote{来医治他}；她们不敢说\Quote{在那里发命令，这里就会成就}。为什么她们不如此说呢？难道百夫长的信德不正是因此而被称赞的吗？因为他说：\Quote{主，我不堪当祢到舍下来；但祢只要说一句话，我的仆人就会痊愈。}\EditorialNote{玛窦福音第八章}这些妇女没有说这样的话，只说：\Quote{主，请看，祢所爱的人病了。}祢知道就够了；因为祢不是那爱了又舍弃的。但有人说：拉匝禄既是一个罪人，怎能被主如此爱呢？让他听主说：\Quote{我来不是召义人，而是召罪人。}\BibleRef{玛 9:13}因为若天主不爱罪人，祂就不会从天降下到世上。\par

6. \Quote{耶稣一听就说：\InnerQuote{这病不至于死，而是为天主的荣耀，为使天主子因此受光荣。}}这样的荣耀自己并没有增加\Emph{祂的}尊威，而是使我们受益。因此祂说\Quote{不至于死}，因为那死亡本身也不至于死，而是为了成就一个奇迹，使人得以信仰基督，从而逃避真正的死亡。请注意主如何间接地称自己为天主，是为了一些否认子是天主的人。因为有异端者否认天主子是天主。让他们在这里聆听：祂说\Quote{这病}，\Quote{不至于死，而是为天主的荣耀，为使天主子因此受光荣。}藉着什么？藉着那病。\par

7. \Quote{那时耶稣爱玛尔大、她的姊妹玛利亚和拉匝禄。}那一位患病，其余的人忧伤，众人都为祂所爱；但爱他们的那一位，既是患者的救主，更是复活死者者和忧伤者的安慰者。\Quote{祂一听说他病了，便仍在原处住了两天。}他们派人告诉祂：祂仍留在那里；时间流逝，直到满了四天。这并非徒然，也许，不，确实，就连这日子的数目也有某种记号性的意义。\Quote{此后，祂又对门徒说，我们往犹太去罢：}在那里祂几乎被石头砸死，而祂显然正是为了逃避被石头砸死才离开的。因为祂以人的身份离去；却仿佛忘记了所有软弱而回来，以彰显祂的能力。\Quote{我们去吧，}祂说，\Quote{往犹太去。}\par

8. 现在请看门徒如何因祂的话而惊恐。门徒对祂说：主，犹太人近来想要用石头砸死祢，祢又要往那里去么？耶稣回答说：白日不是有十二个时辰么？这样的回答是什么意思？他们对祂说：\Quote{犹太人近来想要用石头砸死祢，祢又要往那里去}被砸死么？主却说：\Quote{白日不是有十二个时辰么？人若在白日走路，就不致失足，因为他看见这世界的光；但若在黑夜走路，就会失足，因为他内里没有光。}祂说的确是白日，但按我们的理解，似乎仍是黑夜。让我们呼求白昼驱散黑夜，用光照亮我们的心。主的意思是什么？按我所能判断的，而且随着祂奥义的崇高与深广显露出光明，祂愿意驳斥他们的怀疑和不信。因为他们想以自己的劝告使主免去死亡，而祂正是为死而来，为拯救他们脱离死亡。同样，在另一处经文，圣伯多禄，那爱主却不完全明白祂来临的缘由的门徒，害怕祂死去，因此得罪了生命，即主自己；因为当主向门徒透露祂将在耶路撒冷受犹太人逼迫时，伯多禄在众人中回答说：\Quote{主，请祢远避自己，可怜自己：这事绝不会临到祢身上。}主立刻答道：\Quote{撒殚，退到我后面去：因为你所体会的不是天主的事，而是人的事。}然而，在此之前不久，他因承认天主子而蒙受了赞许：因为他听到了这话：\Quote{西满巴约纳，你是有福的：因为不是血肉启示了你，而是我在天之父。}\BibleRef{玛 16:16}对那曾听到\Quote{你是有福的}的人，祂现在说：\Quote{撒殚，退到我后面去；}因为他之有福不是出于他自己。那是出于什么呢？\Quote{因为不是血肉启示了你，而是我在天之父。}看，你这样是有福的，不是由于你自身任何东西，而是由于属于我的东西。不是由于我是父，而是凡是父所有的都属于我。但如果他的福乐来自主的作为，那么他成为撒殚又来自谁的作为呢？祂在那里告诉我们：因为祂讲明了这种福乐的原因，当祂说：\Quote{不是血肉启示了\Emph{这个}给你，而是我在天之父：}这就是你有福的原因。但我说的\Quote{撒殚，退到我后面去，也请听它的原因：因为你所体会的不是天主的事，而是人的事。}所以，谁也不要自夸：在他本性的东西上，他是撒殚；在来自天主的东西上，他是有福的。因为属于他的一切，从哪里来的呢？岂不来自他的罪？除去那罪，那是属于你自己的。祂说，正义属于我。因为你有什么不是领受的呢？\BibleRef{格前 4:7}因此，当人们想给天主出主意，门徒给主，仆役给主人，病人给医生出主意时，祂斥责他们说：\Quote{白日不是有十二个时辰么？人若在白日走路，就不致失足。}跟随我，如果你们不愿失足：不要给我出主意，你们理应从我这里领受。那么，\Quote{白日不是有十二个时辰么}这话指什么呢？恰就是要指明祂自己是白昼，祂拣选了十二位门徒。祂说，如果我是白昼，你们是时辰，时辰能给白昼出主意么？白昼由时辰跟随，而非时辰由白昼跟随。如果这些门徒是时辰，那么在这样的计算中，犹大是什么？他也列在十二时辰中么？如果他是时辰，就曾有过光；如果有光，那白昼怎么被他出卖至死？但主这样说时，所预见的不是犹大本人，而是他的继任者。因为犹大堕落时，由玛弟亚接替，十二之数得以保持。\BibleRef{宗 1:26}所以，主拣选十二门徒不是没有目的的，而是为了指明祂自己就是那精神的白昼。那么，让时辰侍立于白昼，让他们宣扬白昼，被白昼认识并光照，并且因时辰的宣讲，世界会相信白昼。因此，概括地说，祂所说的正是这一点：跟随我，如果你们不愿失足。\par

9. \\BibleQuote{此后，祂对他们说：\\InnerQuote{我们的朋友拉匝禄睡着了；但我要去唤醒他。}}祂说的是真话。对于他的姐妹，他是死了；对于主，他是在睡眠。他对于不能使他复活的人，是死了；但主唤醒他从坟墓中出来，就如一个人唤醒睡者从床上起来一样容易。因此，祂是就自己的权能将他说成睡觉：因为其他死者，在圣经中也常被说成睡眠；就如宗徒所说的：\\BibleQuote{弟兄们，关于安眠的人，我不愿意你们不知道，免得你们忧伤，像那些没有望德的人一样。}\\BibleRef{得前 4:13} 因此，他也称他们为安眠，因为他预告他们的复活。所以，所有死去的人，无论善恶，都在睡眠。但正如那些天天睡眠又醒来的人，在他们的睡眠中各人所见大不相同：有人做愉快的梦，有人做如此可怕的梦，以致醒来害怕再入睡，怕重蹈覆辙：同样，每个人睡眠和醒来都有各自特殊的情况。而且，对于日后将被带到审判者面前的人，他们所被放置的监禁种类也有区别。因为人被置于何种监禁取决于案件的功过：有些人由扈从看守，这是一种人性温和、符合公民身份的职务；有些人被交给下属；有些人被投入监狱；而在监狱中，并非所有人都被一起关进最底层的地牢，而是按照罪行的轻重和严重程度处理。正如在公务人员中有不同种类的监禁，死者也有不同种类的监禁，复活者的功过也各不相同。那乞丐被拘禁了，那富人也被拘禁了：但一个到了亚巴郎的怀抱；另一个到了他口渴而找不到一滴水的地方。\\BibleRef{路 16:22}\par

10. 因此，为了使这一点能教导你们的爱德：所有灵魂在离开这个世界时，都有各自不同的接纳。善人得享喜乐，恶人忍受痛苦。但到复活时，善人的喜乐将更圆满，恶人的痛苦将更沉重，因为他们将在身体内受折磨。圣祖、先知、宗徒、殉道者、善人和忠信者，都已进入平安；但他们都还要在最终领受天主应许的圆满实现，因为还许诺了他们肉身的复活、死亡的毁灭以及同天使的永生。这一切我们将共同领受；至于人死后立即领受的其余恩惠，凡堪当的人，都在死时领受。圣祖们首先领受——试想他们从何等安息中安息；之后是先知；再近些是宗徒；更近些是圣殉道者；而日复一日，善人和忠信者领受。因此，有些人在那安息中已长久，有些人则不太久；有些人年份更少，而另一些人进入那里则更为近期。但等到他们从这睡眠中醒来时，他们将一起领受那应许的圆满。\par

11. \\BibleQuote{我们的朋友拉匝禄睡着了；但我要去唤醒他。}门徒们便说——按照他们的理解回答——\\BibleQuote{主，如果他睡着了，他会好的。}因为病人的睡眠通常是康复的迹象。\\BibleQuote{然而耶稣说的是他的死，他们却以为祂说的是安眠休息。}于是耶稣就明明地对他们说，——因为祂说得有些隐晦，\\BibleQuote{他睡着了；}——所以祂明明地说：\\BibleQuote{拉匝禄死了。我为你们欢喜，因为我不在那里，好叫你们相信。}我甚至知道他已经死了，而我不在那里：因为人们报告的并非他的死，而是他的病。但那创造了他、死者灵魂已归于祂手中的那一位，有什么事能向祂隐藏呢？因此祂说：\\BibleQuote{我为你们欢喜，因为我不在那里，好叫你们相信}，好让他们现在开始惊奇，主能断言他的死，而祂既未看见也未听说。因为在这里我们特别要记住，门徒们虽然已经信了祂，但他们的信德是通过奇迹被建立的：并非要产生一种以前完全没有的信德，而是要使那已经存在的信德增长；尽管祂用了这样的措辞，好像他们那时才将要相信。因为祂没有说：\\BibleQuote{我为你们欢喜}，好使你们的信德增长或坚定；而是说\\BibleQuote{好叫你们相信}，这应理解为，好使你们的信德更圆满、更坚强。\par

12. \Quote{无论如何，让我们到他那里去。于是那称为迪迪默斯的多默对他的同伴门徒说：我们也去，好与祂同死。因此耶稣来了，发现他已在坟墓里四天了。} 关于这四天，按照圣经中隐晦段落的常例，可以有许多说法，这些段落所承载的意义，与理解它们的人之多样性一样多。也让我们说出对于这\Quote{死了四天}所包含的意义的看法。因为如同在先前的瞎子身上，我们以某种方式理解整个人类，同样，在这死人身上，或许也要理解许多人；因为同一件事可以由不同的预象来表示。当一个人出生时，他已出生在死亡的状态中；因为他从亚当继承了罪。因此宗徒说：\Quote{因一人罪恶进入了世界，死因罪恶而来；于是死亡传给了众人，因为众人都犯了罪。} \BibleRef{罗 5:12} 这里你有了死亡的第一天，因为人从死亡的种子堆继承了它。此后他成长，开始接近理智的年龄，以便可以认识自然律，这自然律每个人都已铭刻在心中：你不愿别人给你做的，不要对他人做。这是从书页中学到的，难道不也在我们的本性中清晰可读吗？你愿意被抢劫吗？当然不。那么，看这在你心中的律法：你不愿遭受的，也不要去做。这律法也被人们违犯；于是，我们有了死亡的第二天。律法也藉天主的仆人梅瑟神圣地赐下；其中说：\Quote{不可杀人；不可奸淫；不可作假见证；孝敬你的父亲和母亲；不可贪恋邻人的财产；不可贪恋邻人的妻子。} \BibleRef{出 20:12} 这里你有了成文的律法，它也受到蔑视：这是死亡的第三天。还有什么呢？福音也来了，天国被宣讲，基督被到处传扬；祂以地狱威胁，祂应许永生；但这也受到蔑视。人们违犯福音；而\Emph{这}是死亡的第四天。现在他理应发臭了。但难道要拒绝怜悯这样的人吗？绝对不可；因为连这样的人也从死者中复活，主认为亲自前来并非不合适。\par

13. \Quote{有许多犹太人来到玛尔大和玛利亚那里，为她们的兄弟安慰她们。玛尔大一听说耶稣来了，就去迎接祂；玛利亚却仍坐在家里。玛尔大对耶稣说：主，祢若在这里，我的兄弟就不会死。但我知道，就是现在，无论祢向天主求什么，天主必会赐给祢。} 她没有说：但即使现在，我求祢使我的兄弟复活。因为她怎知道这样的复活对她的兄弟有益呢？她只说：我知道祢能，并且祢乐意做的，祢就做；因为祢的作为取决于祢自己的判断，而非我的冒昧。\Quote{但即使现在，我知道，无论祢向天主求什么，天主必会赐给祢。}\par

14. \Quote{耶稣对她说：你的兄弟必要复活。}这是含糊的。因为祂没有说：我现在就要使你的兄弟复活；而是说：\Quote{你的兄弟必要复活。玛尔大对祂说：我知道在末日复活时，他必要复活。}对于那复活，我确信，但对这次却不确定。\Quote{耶稣对她说：我就是复活。}你说：我的兄弟在末日必要复活；不错；但藉着祂，就是藉着那时他赖以复活的那一位，他现在也能复活，因为祂说：\Quote{我就是复活和生命。}请听，弟兄们，请听听祂说什么。当然，旁人的普遍期待是，那死了四天的拉匝禄会再活过来；让我们也听，并复活起来。在听众中有多少人被某种罪恶习惯的重荷压垮啊！也许有些人在听我，对他们可以说：\Quote{不要醉酒，醉酒使人放荡；} \BibleRef{弗 5:18} 而他们说：我们不能。另有些人，或许在听我，他们不洁，被情欲和罪恶玷污，对他们说：戒除这样的行为，以免你丧亡；他们回答：我们不能放弃我们的习惯。哦！主啊，使他们复活吧。祂说：\Quote{我就是复活和生命。}复活\Emph{因为}生命。\par

15. \Quote{凡信我的，即使死了，仍要活着；凡活着而信我的，永远不死。}这是什么意思？\Quote{凡信我的，即使死了，}正如\Person{拉匝禄}死了，\Quote{仍要活着；}因为祂不是死人的\OriginalTerm{天主}{God}，而是活人的天主。祂曾就他们很久以前死去的祖先，即\Person{亚巴郎}、\Person{依撒格}和\Person{雅各伯}，如此答复犹太人：我是亚巴郎的天主、依撒格的天主、雅各伯的天主；祂不是死人的天主，而是活人的天主，因为一切人都是为祂而生活。所以，相信吧，即使你死了，你仍要活着；但若你不信，即使你活着，你也是死的。同样，我们来证明这一点：若你不信，即使你活着，你也是死的。有一个人迟疑着不愿跟随祂，说：\Quote{请先让我去埋葬我的父亲。}\OriginalTerm{主}{Lord}便说：\Quote{让死人埋葬他们的死人；你只管来跟随我。}\BibleRef{玛 8:21}那里有一个死人需要埋葬，也有死人去埋葬死人；一个在肉身里死了，其他的在灵魂里死了。灵魂的死亡是怎样发生的？是当\OriginalTerm{信德}{faith}缺失的时候。身体的死亡是怎样来的？是当灵魂缺失的时候。所以，你灵魂的灵魂就是信德。\Quote{凡信我的，}\OriginalTerm{基督}{Christ}说，虽然他在肉身死了，却要在灵魂里活着；直到肉身也复活，永远不再死。这就是\Quote{凡信我的，}即使死了，\Quote{仍要活着。而凡活着}在肉身里，\Quote{而信我的，}虽然他会因肉身之死而暂时死去，\Quote{却永远不死，}因为有了灵魂的生命和\OriginalTerm{复活}{resurrection}的不朽。这就是这句话的意思：\Quote{凡活着而信我的，永远不死。你信这话吗？她对祂说：是的，主，我相信祢是基督、\OriginalTerm{天主子}{Son of God}，来到世界上的那一位。}当我信了这话，我就相信祢是复活，祢是生命；我相信凡信祢的，虽然死了，仍要活着；凡活着而信祢的，永远不死。\par

16. \Quote{她说了这话，就回去，暗暗地叫她的妹妹\Person{玛利亚}说：师傅来了，祂叫你。}值得注意的是，她声音的低语被称作沉默。因为当她说了\Quote{师傅来了，祂叫你}时，怎能是沉默呢？还值得注意的是，福音作者为何没有说明在哪里、何时或如何主叫了玛利亚；也就是说，为了保持叙述的简洁，宁愿从\Person{玛尔大}的话语中理解。\par

17. \Quote{她一听见这话，就急忙起来，到\Person{耶稣}那里去。那时\Person{耶稣}还没有进村镇，仍在\Person{玛尔大}迎接祂的地方。那些同\Person{玛利亚}在家里安慰她的犹太人，见她急忙起来出去，就跟着她，说：她往坟墓那里去哭。}福音作者告诉我们这些是为了什么？为了给我们显示，是什么引起了众人聚集在那里，当\Person{拉匝禄}被复活时。因为犹太人认为她急忙出去是为了以哭泣寻求悲伤的安慰，就跟着她；这样，那使一个死去四天的人复活的伟大奇迹，就能有许多见证人在场。\par

18. \Quote{那时，\Person{玛利亚}到了\Person{耶稣}所在的地方，一看见祂，就俯伏在祂脚前，对祂说：\OriginalTerm{主}{Lord}，祢若早在这里，我的兄弟就不会死了。\Person{耶稣}看见她哭泣，同她一起的犹太人也哭泣，就心神感伤，烦乱起来，说：你们把他放在哪里？}这里，若我们知道的话，祂通过心神感伤和烦乱自己，向我们暗示了什么。因为谁能烦乱祂，除非祂自己？所以，我的弟兄们，首先要注意这里行这事的能力，然后寻求其意义。你们是被迫烦乱；\OriginalTerm{基督}{Christ}是自主地烦乱。\Person{耶稣}确实饿了，但因为是祂愿意的；\Person{耶稣}确实睡了，但因为是祂愿意的；祂确实忧愁，但因为是祂愿意的；祂确实死了，但因为是祂愿意的：成为这样或那样，或者不这样，都在祂自己的权能之内。因为\OriginalTerm{圣言}{Word}取了\OriginalTerm{灵魂}{soul}和肉身，将我们的全部人性结合在祂位格的单一性中。因为\OriginalTerm{宗徒}{apostle}的灵魂被圣言光照；同样，\Person{伯多禄}的灵魂，\Person{保禄}的灵魂，其他宗徒和圣先知们的灵魂——所有灵魂都被圣言光照；但没有一个曾被称为\Quote{圣言成了血肉；}没有一个曾被称为我与\OriginalTerm{圣父}{Father}原是一体。\OriginalTerm{基督}{Christ}的灵魂和肉身是同一圣言的位格，同一基督。并且藉这[圣言]，在其中存有至高权能，软弱被祂意志的示意所使用；这样，\Quote{祂烦乱了自己。}\par

19. 我已论及了能力：现在来看含义。那四天的死亡和埋葬所象征的是一个重大的罪人。那么，基督为什么烦恼自己，岂不是为了向你表明，当你被如此重大的罪孽压垮时，你应当怎样烦恼？因为你已经审视了自己，看到了自己的罪过，自己盘算：我做了这事，天主宽恕了我；我犯了那罪，祂容忍了我；我听了福音，却藐视了它；我受了圣洗，又回到了同样的道路：我在做什么？我往哪里去？我如何逃脱？当你这样说话时，基督已经在呻吟；因为你的信德在呻吟。在这样呻吟者的声音中，他复活的希望就显露出来。如果这样的信德在内，就有基督在呻吟；因为如果我们内有信德，基督就在我们内。因为宗徒还说了什么：\Quote{为使基督因信德住在你们心中。} \BibleRef{弗 3:17} 因此，你对基督的信德就是基督自己在你心中。这就是为什么祂在船上睡觉；为什么当祂的门徒身处危险、濒临船难时，他们来叫醒了祂。基督起来，命令风和浪，就大大平静了。\BibleRef{玛 8:24} 同样，在你身上，风进入你的心，就是你在那里航行，在那里度过此生如风暴危险的海；风进入，波浪翻腾，摇动你的船。这些风是什么？你受了侮辱，就发怒：那侮辱是风，那怒气是波浪。你身处危险，你准备回击，以咒骂还咒骂，你的船已近沉没。唤醒那沉睡的基督。因为你激动不安，准备以恶报恶，因为基督在你的船上沉睡。因为基督在你心中的沉睡就是信德的遗忘。但如果你唤醒基督，就是回忆你的信德，那时基督在你心中醒着，你听到祂对你说什么？祂说，我听说有人对我说：\Quote{你是附了魔的。} 而我已经为他们祈祷了。主听见并忍受；仆人听见却发怒！但你想要报复。为什么？我已经被报复了。当你的信德这样对你说话时，命令仿佛对风和浪发出，就有大大的平静。因此，如同在船上唤醒基督就是唤醒信德，同样，在一个被重大罪孽和习惯压迫的人心中，在甚至违犯神圣福音、藐视永恒惩罚的人心中，让基督呻吟，让他转向自责。再听：基督哭了；让人自悲。因为基督为什么哭，岂不是为了教导人哭？祂为什么呻吟并烦恼自己，岂不是为了表明一个有理不满意自己的人的信德应当因谴责恶行而呻吟，以致罪恶的习惯因补赎的强烈而退让？\par

20. \Quote{祂说：你们把他放在哪里？} 你知道他死了，却不知道埋葬他的地方吗？这里的含义是，这样一个失落的人，在天主眼中仿佛成了未知的。我不敢说，是未知的——因为有什么对祂是未知的呢？但是，仿佛未知。我们如何证明这一点？请听主在审判时还将说的话：\Quote{我不认识你们，你们离开我吧。} \BibleRef{玛 7:23} 那句\Quote{我不认识你们}是什么意思？我在我的光中——在我所认识的正义中——看不见你们。所以这里，祂也仿佛对这样的罪人一无所知，说：\Quote{你们把他放在哪里？} 相似的性质是人在犯罪后天主在乐园中的声音：\Quote{亚当，你在哪里？} \BibleRef{创 3:9} \Quote{他们对祂说：主，请来看。} 这个\Quote{看}是什么意思？就是怜悯。因为主怜悯时就看。因此有人对祂说：\Quote{请看我的卑微（苦难）和我的痛苦，并赦免我所有的罪过。}\par

21. \Quote{耶稣哭了。犹太人就说：看，祂多么爱他！} \Quote{爱他}，是什么意思？\Quote{我来不是召义人，而是召罪人悔改。} \BibleRef{玛 9:13} \Quote{但其中有人说：这个人开了瞎子的眼睛，难道不能也使这个人不死吗？} 但是，祂不愿意阻止他死，是为了更伟大的事：使他从死者中复活。\par

22. \Quote{耶稣又心中悲叹，来到坟墓前。} 愿祂的悲叹也以你为目标，如果你要重新进入生命！每个处于那种可怕道德状况的人，有话说给他：\Quote{祂来到坟墓前。} \Quote{那是一个洞穴，有一块石头盖在上面。} 死在石头下，有罪于法律。因为你知道，赐给犹太人的法律是刻在石头上的。\BibleRef{出 31:18} 所有有罪的人都在法律之下；正直的人与法律和谐。\Quote{法律不是为义人立的。} \BibleRef{弟前 1:9} 那么，\Quote{你们把石头挪开}这句话是什么意思？就是传布恩宠。因为宗徒保禄自称是新约的仆役，不是文字，而是圣神；\Quote{因为文字叫人死，圣神却叫人活。} \BibleRef{格后 3:6} 叫人死的文字就像压碎人的石头。\Quote{祂说：你们把石头挪开。} 挪开法律的重量；传布恩宠。\Quote{因为如果曾有一种法律能赐予生命，那么正义确实就是出于法律了。但圣经把众人都禁锢在罪恶中，为的是使因信耶稣基督而来的恩许赐给相信的人。} \BibleRef{迦 3:21} 所以，\Quote{你们把石头挪开。}\par

23. \Quote{那死者的姊妹玛尔大对祂说：\InnerQuote{主，现在他必定臭了，因为他已经死了四天了。}耶稣对她说：\InnerQuote{我不是对你说过，如果你信，就要看见天主的光荣吗？}}祂说\InnerQuote{你要看见天主的光荣}是什么意思？是叫那已经腐烂、死了四天的人复活。因为\BibleQuote{众人都犯了罪，亏缺了天主的光荣}\BibleRef{罗 3:23} 并且，\BibleQuote{罪恶在那里越多，恩宠在那里也越格外丰富}\BibleRef{罗 5:20}\par

24. \Quote{他们就把石头挪开。耶稣举目向上说：\InnerQuote{父啊，我感谢祢，因为祢俯听了我。我本来知道祢常常俯听我，但我说这话，是为了四周站着的群众，叫他们信我是祢所派遣的。}说了这话，便大声呼喊说。}祂叹息、流泪、大声呼喊。一个人被沉重的罪恶习惯压碎，要起来是何等困难！然而他确能起来：他由内在的隐藏恩宠获得生命，在那大声呼喊之后，他起来了。因为接着是什么？\Quote{祂大声喊说：\InnerQuote{拉匝禄，出来！}那死者便出来了，手脚缠着布条，脸上包着手巾。}你希奇他脚被缠着如何能出来，却不希奇那死了四天的人从死者中复活？两件事都是主的德能在运作，并非死者的力量。他出来了，却仍然被缠着。他仍在殓布中，却已走出坟墓。这是什么意思？当你轻视基督时，你躺在死亡的怀抱中；若你的轻视达到我所说的程度，你也被埋葬了；但当你告解时，你就出来。因为出来，不就是公开承认真实的状况，仿佛离开旧日黑暗的藏身处吗？但你所做的告解是天主促成的，当祂大声呼喊，或用丰盛的恩宠召唤你时。因此，当死者仍然被缚着出来时，他一面告解，一面仍有罪债；为使他的罪也被解除，主对仆人们说：\Quote{解开他，让他去。}祂这话是什么意思？\BibleQuote{凡你们在地上所释放的，在天上也要被释放}\BibleRef{玛 16:19}\par

25. \Quote{那些来到玛利亚那里的犹太人，见了耶稣所行的事，就有许多人信了祂。但其中也有人去到法利塞人那里，把耶稣所行的事告诉他们。}来到玛利亚那里的犹太人并非全都信了，但许多信了。\Quote{但其中也有人}——无论是来的人，或信了的人——\Quote{去到法利塞人那里，把耶稣所行的事告诉他们：}或是为传递消息，好让他们也信；或是出于背叛，激起他们的愤怒。但无论这些人是何许人，无论动机如何，这些事都传到了法利塞人那里。\par

26. \Quote{司祭长和法利塞人聚集了议会，说：\InnerQuote{我们怎么办？}}但他们没有说：让我们信。因为这些丧尽天良的人更思谋如何行恶来毁灭祂，而不思想自己的得救；然而他们害怕，彼此商议。因为\Quote{他们说：\InnerQuote{我们怎么办？这人行了许多奇迹；若让祂这样，众人都会信从祂，罗马人必要来，连我们的圣殿和民族都要除掉。}}他们害怕失去现世财物，而不思想永生；于是两者都失去了。因为在我主受难并进入光荣以后，罗马人夺去了他们的圣殿和民族——他们攻陷了圣殿，迁移了民族；并且现在那在别处所说的也应验在他们身上：\BibleQuote{但本国的子民，却被驱逐到外面的黑暗中}\BibleRef{玛 8:12}但这正是他们所害怕的，即如果人人都信从基督，就没有人留下来保卫天主的城和圣殿以对抗罗马人；因为他们觉得基督的教导是针对圣殿本身和他们祖传法律的。\par

27. \Quote{他们中有一个人，名叫加亚法，就是那年做\OriginalTerm{大司祭}{Pontifex}的，对他们说：\InnerQuote{你们什么都不知道，也不想想，叫一个人替百姓死，免得全民族灭亡，为你们是有利的。}他说这话，不是出于自己，但因他是那年的大司祭，他预言了。} 这里我们受教：预言的神恩甚至利用恶人来预告未来的事；然而，福音作者将这归结于他身为大司祭的神圣圣事性事实。不过，有一个问题：他为何被称为那年的大司祭？因为天主本指定一个人做大司祭，只有他死后别人才继任。但我们必须明白，犹太人中间的野心和争执导致后来不止一人被任命，并且他们每年轮流服务。因为关于匝加利亚也记载说：\Quote{匝加利亚按他的班次，在天主前尽司祭的职务时，照司祭的惯例，他抽中了签，得进主的圣所烧香。}\BibleRef{路 1:8} 由此明显可见，当时不止一位司祭，每人轮流值班；因为只有大司祭才可把香放在祭台上。\BibleRef{出 30:7} 或许在同一年中有几位轮流服务，次年再由另外几位接替，而烧香的任务通过抽签落在其中一人身上。那么，加亚法预言了什么？\Quote{耶稣要替民族死；不但替民族死，也要把天主四散的子女聚集归一。} 这是福音作者补充的；因为加亚法只预言了犹太民族，其中有主亲自说过的羊：\BibleQuote{我奉派遣，只是到以色列家迷失的羊那里去。}\BibleRef{玛 15:24} 但福音作者知道还有别的羊，不属于这羊栈，也必须领来，好使成为一栈一牧。但这是以预定的方式说的；因为那些仍不信的人，既不是祂的羊，也不是天主的子女。\par

28. \Quote{从那一天起，他们就议决要杀害耶稣。因此，耶稣不再公开地在犹太人中间往来，却从那里往近旷野的地方去，到了一座城名叫厄弗辣因，就在那里同门徒住下。} 这并不是由于祂能力的失败，因为只要祂愿意，祂仍可与犹太人往来而不受他们伤害；但祂在人性软弱中，为门徒提供了生活的榜样，借此表明：祂的信徒——祂的肢体——避开迫害者的面，用隐藏来躲避恶人的怒气，而不是因公开露面而激怒它，这并不是罪。\par

\SourceRecord{Translated by John Gibb.；Nicene and Post-Nicene Fathers, First Series；Vol. 7.；Edited by Philip Schaff.；Buffalo, NY: Christian Literature Publishing Co.,；1888.；<http://www.newadvent.org/fathers/1701049.htm>.}

% structure: p0001 source=h1 target=section reason=page-title

\section{讲道辞第50篇（若望福音11:55-12）}

1. 昨日我们在神圣福音中所读的章节，蒙主助佑，我们已作了阐释；今日的章节，我们愿意以同样的依赖之心加以讨论。圣经中有一些经文非常清楚，只需要听众而不需要解释者；对于这样的经文我们不必停留，以便有足够的时间处理那些必然需要更详尽思考的经文。\par

2. \Quote{犹太人的逾越节临近了。}犹太人希望用主的血染红那个庆节。就在那日，那羔羊被宰杀了，祂以自己的血为我们祝圣了那个庆节。犹太人图谋杀害耶稣；而祂，从天上降来受苦的那位，愿意走近祂受苦的地方，因为祂苦难的时刻临近了。因此，\Quote{许多人在逾越节前从乡下上到耶路撒冷，要洁净自己。}犹太人这样做是依照法律中圣梅瑟所传的主的命令，即逾越节庆日时所有人都应从各地聚集，并在举行该日礼仪时洁身自好。但那个庆祝是将来之事的预像。为什么是预像？它是关于将要来临的基督的预言性的预示，是对那位要在那日为我们受苦之人的预言：如此，预像消失，光明来临；预兆过去，真理得以保留。因此犹太人以预像的方式守逾越节，而我们则在光明中守。因为主为什么命令他们在该庆节当天宰杀一只羊，岂不是因为关于祂曾有预言：\BibleQuote{祂被牵去，如同羊被牵到宰杀之地}。\BibleRef{依 53:7}犹太人的门框用被杀牲畜的血涂抹；而我们的额上则用基督的血涂抹。那个印记——因为它有真实的意义——据说能使毁灭者远离被印记的房屋：\BibleRef{出 12:22}如果我们接受救主到心中，基督的印记就为我们驱赶毁灭者。但我为什么说这些？因为许多人的门框已被印记，但内心却没有住客居住：他们轻易在额上印有基督的记号，内心却拒绝接受祂的话。所以，弟兄们，我说过，并且重复说，只要我们将基督作为内心的住客，基督的印记就为我们驱赶毁灭者。我说这些，免得有人心思转向这些犹太人庆节的意义。因此主如同来到牺牲之地，为使我们拥有真正的逾越节，当我们庆祝祂的苦难作为真正的羔羊祭献时。\par

3. \Quote{他们就寻找耶稣。}但怀着恶意。因为那些以善意寻求耶稣的人是有福的。他们寻找祂，目的是使他们和我们都不再拥有祂；但祂离开他们，却被我们接纳。有些寻找祂的人受责备，另一些寻找祂的人受称赞；因为寻找者的精神决定受赞或受责。因此你们也在圣咏集中读到：\Quote{愿那些寻索我灵魂的人，蒙羞受辱。}这样的人就是怀着恶意寻找的人。但在另一处他说：\Quote{我无处可逃，也没有人寻索我的灵魂。}寻找的和不寻找的都同样受责备。因此让我们寻求基督，为使我们拥有祂，持守祂，而不是杀害祂；因为这些人想要抓住祂，只是为了尽快永远除掉祂。\Quote{因此他们寻找祂，彼此谈论：你们想什么？祂不会来过节吗？}\par

4. \Quote{司祭长和法利塞人已经发出命令：若有人知道祂在那里，就要报告，好让他们捉拿祂。}我们呢，我们要向犹太人指明基督在哪里。巴不得所有那些曾下令报告基督所在之人的后裔，都能听见并领会！让他们到教会来听基督在哪里，并抓住祂。他们可以从我们这里听到，也可以从福音中听到。祂被他们的祖先杀害，被埋葬，又复活了，被门徒们认出来，当着他们的面升到天上，现在坐在圣父的右边；那位曾被审判的还要再来审判众人：让他们听见，并牢牢持守。他们也许会回答：我如何抓住不在场的祂？我如何将手伸到天上，抓住坐在那里的那一位？举起你的信德，你就抓住了。你们的祖先以肉体抓住，你当以心抓住；因为不在场的基督也是临在的。若非祂的临在，我们自己也抓不住祂。但既然祂的话是真实的：\BibleQuote{看，我同你们天天在一起，直到今世的终结。}\BibleRef{玛 28:20}祂离开了，又在这里；祂回去了，却不舍弃我们；因为祂将身体带到天上，但祂的尊威从未从世间撤离。\par

5. \Quote{逾越节前六天，耶稣来到伯达尼，那里住着拉匝禄，就是祂从死里复活的那位。他们为祂预备了晚餐；玛尔大服侍，拉匝禄是坐席的人之一。}为了防止人们以为那人成了幻影，因为他从死者中复活了，他竟是坐席者之一；他活着、说话、宴饮：事实彰显，犹太人的不信大为羞愧。因此主与拉匝禄和其他人一同坐席；由拉匝禄的姐妹玛尔大服侍他们。\par

6. 但是，拉匝禄的另一位姐妹玛利亚\Quote{拿了一斤极珍贵的纯哪哒香液，敷抹了耶稣的脚，并用她的头发擦干；屋里便充满了香液的气味。}这就是所发生的事，让我们探究它所隐含的奥义。你们中凡是真心愿意成为信徒的人，都要像玛利亚一样，用珍贵的香液敷抹主的脚。那香液就是正义，因此它恰好是一斤重；但它是纯哪哒（\OriginalTerm{纯哪哒}{nardi pistici}）香液，极其珍贵。由于它被称为\Quote{\OriginalTerm{纯}{pistici}}，我们应当推断有一个地方赋予了它珍贵性：但这并未穷尽其含义，而且它与一件圣事的象征十分吻合。这个词的希腊语词根在我们这里被称为\Quote{信德}。你寻求行正义：义人将因信德而生活。\BibleRef{罗 1:17} 敷抹耶稣的脚：以善生追随主的足迹。用你的头发擦干：把你多余的给穷人，你就擦干了主的脚；因为头发似乎是身体的多余部分。你从你的富余中有些多余：这对你是多余的，但对主的脚却是必要的。也许在这世上，主的脚仍有需要。因为除了祂的肢体，祂在末了还要对谁说：\Quote{你们对我这些最小兄弟中的一个所做的，就是对我做的}？\BibleRef{玛 25:40} 你花费了对你多余的，但你却做了令我的脚喜悦的事。\par

7. \Quote{屋里便充满了气味。}世界充满了善名的声誉；因为善名犹如芬芳的气味。那些生活邪恶却被称为基督徒的人，损害了基督：关于他们，经上说，因着他们\Quote{天主的名被亵渎}。\BibleRef{罗 2:24} 如果因着他们天主的名被亵渎，那么因着善人，主的名被显扬。请听宗徒的话，他说：\Quote{我们在各处都是基督的馨香。}正如《\WorkTitle{雅歌}》中也说：\Quote{你的名字如同倒出的香液。}\BibleRef{歌 1:3} 再听宗徒的话：\Quote{我们在各处都是基督的馨香，在得救的人中，也在丧亡的人中；对后者，我们是致死的馨香；对前者，是生命的馨香：谁堪当这些事呢？}\BibleRef{格后 2:14} 面前的圣福音课程给了我们机会谈论这种馨香，使我们能适当地讲述，你们能殷勤地聆听宗徒自己所说的话：\Quote{谁堪当这些事呢？}但我们能从这些话中推断我们有资格尝试谈论这样一个主题，或你们有资格聆听吗？我们确实不配；但祂是配的，祂乐意藉着我们说出对你们有益的听闻。你看，宗徒，正如他自称，是\Quote{馨香}：但那馨香\Quote{对某些人是生命的馨香，对另一些人是致死的馨香}，而它本身始终是\Quote{馨香}。因为他说过，对某些人是生命的馨香，对另一些人是致死的恶臭吗？他称自己是馨香，不是恶臭；并表明自己是同样的馨香，对某些人至于生命，对另一些人至于死亡。那些在这馨香中找到生命的人有福了！但还有什么比那些对这馨香成为死亡使者的人更悲惨呢？\par

8. 有人说：是谁被这馨香杀死了呢？宗徒正是用这些话暗示：\Quote{谁堪当这些事呢？}天主以何等奇妙的方式使这馨香既为善人带来生命，也为恶人带来死亡；这是如何可能的，就主所乐意激发我的思想而言（因为它可能隐藏着更深的意义，超出我的洞察力）——但我仍就我所洞察到的，你们不应得不到信息。保禄宗徒生活和行为的正直，他在言语和行动中宣讲的正义，他作为导师的奇妙能力和作为管家的忠信，处处被传扬：有些人爱他，有些人嫉妒他。因为他自己在某处说，有些人出于嫉妒宣讲基督，并不真诚，\Quote{想}，他说，\Quote{给我的锁链增加烦恼。}但他又加上什么？\Quote{无论出于假意还是真诚，愿基督被宣讲。}爱我的人宣讲，恨我的人也宣讲；前者在这馨香中生活，后者在其中死亡：然而通过两者的宣讲，基督的名被宣告，世界被这极好的馨香充满。你曾爱过一个其行为证明他良善的人吗？那你就在这馨香中活了。你曾嫉妒过这样的人吗？那你就在这同一馨香中死了。但是，请问，你既选择死亡，是否将这馨香变成了恶臭？从你的嫉妒心中回转，这馨香就不再杀你了。\par

9. 最后，现在请听我们在这段经文中所见，这香液如何对某些人是生命的馨香，对另一些人是致死的馨香。当虔诚的玛利亚向主提供了这项感恩的服务后，立刻，祂的一个门徒，即将出卖祂的依斯加略人犹达斯说：\Quote{为什么不把这香液卖三百块银钱，施舍给穷人？}唉，你这可怜的人！馨香杀了你。因为圣福音作者揭示了他如此说话的原因。但我们也可能认为，要不是福音中揭示了他内心真实的状态，我们会以为是关心穷人促使他如此说话。不是的。那么是什么？请听一个真实的见证：\Quote{他说这话，并不是因为他关心穷人，而是因为他是个贼，掌管钱囊，并偷取所放入的。}他是带着钱囊，还是偷走？为了公共的服务他带着，作为贼他偷走了。\par

10. 现在请注意，并要知道，这个犹大并不是在他向犹太人受贿并出卖了主的那一刻才变得乖戾的。因为不少粗心阅读福音的人以为，犹大只是在他从犹太人那里收受金钱出卖主时才丧亡的。他不是在那时丧亡的，而是当他跟随主的时候，他已经是一个窃贼和丧亡之徒；因为他只是身体上跟随，心里并未跟随。他组成了十二宗徒的数目，但在宗徒的福分中无份：他被造成第十二个只是外表；当他离去、另一个人继任时，宗徒的实质得以完成，数目的完整得以保持。\BibleRef{宗 1:26} 那么，我的弟兄们，我们的主耶稣基督愿意给祂的教会留下什么教训，当祂乐意在十二人中有一个丧亡者时，岂不是要我们容忍恶人，并且不要分裂基督的身体？看，你们有犹大在圣人们中间——就是那个犹大！他是一个窃贼，是的——不要忽略——他不是一般的窃贼，而是窃贼兼亵圣者：他是钱囊的强盗，但那是主的钱囊；是钱囊的强盗，但那是神圣的钱囊。如果在世俗法庭上对普通盗窃和窃取公物（我们所说的窃取公物是指盗窃公共财产）有所区分，并且私人盗窃所受的判决不如公共盗窃严厉，那么对于胆敢盗窃，不是从普通场所，而是从教会盗窃的亵圣窃贼，判决应当多么严厉？从教会盗窃的人，与丧亡的犹大并列。这个犹大就是这样的人，然而他却与十一位圣洁的宗徒一同进出。他甚至与他们一同来到主的筵席前；他被允许与他们交往，但他不能玷污他们。伯多禄和犹大都领受了同一个饼，然而信徒与不信者有什么相通？伯多禄的领受是为了生命，而犹大的领受是为了死亡。因为那好饼就像甘美的香气。正如甘美的香气，好饼也给善人生命，给恶人带来死亡。\Quote{因为那不相称地吃的人，就是吃自己的罪案，喝自己的定案。}\BibleRef{格前 11:29} \Quote{自己的定案，}不是你的定案。那么，如果那是他自己的定案，不是你的，你作为善人就要容忍恶人，使你能得到善人的赏报，而不被投入恶人的惩罚中。\par

11. 要留心我们主与世人共处时的榜样。祂有天使服侍，为什么还要有一个钱囊？无非是预示祂的教会将来要有存放钱款的地方。祂为何容许一个窃贼进来？无非是教导祂的教会耐心容忍窃贼。但这个已经养成从钱囊中窃取钱财习惯的人，为收受金钱而毫不犹豫地出卖了主本身。但让我们看看主对这样的话作了什么回答。请看，弟兄们：祂并没有对他说：你这样说是因为你的盗窃习性。祂知道他是窃贼，却没有揭发他，反而容忍他，为我们展示了在教会中容忍恶人的忍耐榜样。\Quote{于是耶稣说：由她留着罢，到我安葬之日。}祂宣告自己的死亡临近了。\par

12. 但接着是什么？\Quote{因为穷人你们常有，但你们不常有我。}我们当然可以理解\Quote{穷人你们常有；}祂这样说是真实的。教会在什么时候缺少过穷人呢？\Quote{但你们不常有我；}祂这话是什么意思？我们该如何理解\Quote{你们不常有我}？不要惊慌：这是对犹大说的。那么，为什么祂不说\Quote{你将有}，而说\Quote{你们将有}？因为犹大在这里不是一个人。一个恶人代表整个恶人的身体；同样，伯多禄代表整个善人的身体，是的，教会的身体，但这是就善人而言。因为如果在伯多禄身上没有教会的圣事性象征，主就不会对他说：\Quote{我要把天国的钥匙给你：凡你在地上所释放的，在天上也要释放；凡你在地上所束缚的，在天上也要束缚。}\BibleRef{玛 16:19} 如果这话只对伯多禄说，那就没有给教会行动的依据。但如果教会也是如此：凡在地上被束缚的，在天上也被束缚；凡在地上被释放的，在天上也被释放——因为当教会开除教籍时，被开除的人在天上被束缚；当一个人被教会和好时，被和好的人在天上被释放——那么，如果教会也是如此，伯多禄在接收钥匙时代表了圣洁的教会。那么，如果在伯多禄身上代表了教会中的善人，在犹大身上代表了教会中的恶人，那么就是对这些后者说的：\Quote{但你们不常有我。}然而\Quote{不常有}是什么意思，\Quote{常有}又是什么意思？如果你善良，如果你属于伯多禄所代表的身体，你现在和将来都拥有基督：\Emph{现在}借着信德、借着记号、借着圣洗圣事、借着祭台上的饼和酒。你现在拥有基督，但你会永远拥有祂；因为当你离开这里时，你会到祂那里去，祂曾对强盗说：\Quote{今天你要同我在乐园里。}\BibleRef{路 23:43} 但如果你生活邪恶，你似乎现在拥有基督，因为你进入教会，以基督的记号划十字，受基督的洗，与基督的肢体混合，并接近祂的祭台：现在你拥有基督，但因为你生活邪恶，你将不会永远拥有祂。\par

13. 也可以这样理解：\Quote{贫穷的人你们常常有，至于我，你们却不常有。}善人也可以将此话看作是对自己说的，但不必因此而忧虑，因为祂是就祂的肉身临在而言的。因为就祂的尊威、眷顾、不可言喻和不可见的恩宠而言，祂自己的话得以应验：\Quote{看，我同你们天天在一起，直到今世的终结。} \BibleRef{玛 28:20}但就祂作为圣言所取的血肉，就祂作为童贞女之子之所是，就祂被犹太人捉拿、钉在木头上、从十字架上卸下、裹在殓布里、安放在坟墓中、并在复活中显现的那方面而言，\Quote{你们却不常有祂。}为什么呢？因为就祂肉身的临在而言，祂同门徒们一起生活了四十天，然后带他们出去，为叫他们瞻仰祂，而非跟随祂，祂就升了天，不在这里了。祂确实在那里，坐在圣父的右边；祂也在这里，从未撤去祂荣耀的临在。换句话说，就祂神圣的临在而言，我们常拥有基督；就祂在肉身中的临在而言，对门徒们正确地说了：\Quote{我，你们却不常有。}就这方面来说，教会只在几天内享有祂的临在；如今，她是以信德拥有祂，而不以肉眼看见祂。那么，无论\Quote{我，你们却不常有}这句话是如何说的，我想，经过这两种解释之后，它不会再是个疑问了。\par

14. 让我们听听其余的几点：\Quote{因此，许多犹太人知道祂在那里；他们来，不单是为耶稣，也是为看祂从死者中所复活的拉匝禄。}他们为好奇心所驱使，而非因爱德；他们来看。请听人类虚荣的奇异谋划。他们看到拉匝禄从死者中复活——因为主如此奇迹的名声伴随着许多确凿的证据传遍各处，并且这事行得如此公开，以致他们既不能掩饰也不能否认所发生的事——只需想想他们想出的计谋。\Quote{司祭长们商议，也要杀害拉匝禄，因为有许多犹太人为了他而离开，信了耶稣。}愚昧的商议和蒙蔽的愤怒！难道能复活死人的主基督，不能也复活被杀的人吗？当你们为拉匝禄准备暴力死亡时，是否同时也剥夺了主的能力？如果你们认为死人与被杀者是两回事，只需看到这一点：主创造了二者，并复活了死去的拉匝禄，复活了祂自己这被杀者。\par

\SourceRecord{Translated by John Gibb.；Nicene and Post-Nicene Fathers, First Series；Vol. 7.；Edited by Philip Schaff.；Buffalo, NY: Christian Literature Publishing Co.,；1888.；<http://www.newadvent.org/fathers/1701050.htm>.}

% structure: p0001 source=h1 target=section reason=page-title

\section{讲道集第五十一篇（若望福音 12:12–26）}

1. 我们的主使一个死了四天的人复活，令犹太人惊愕不已；其中有些人见了这事便相信了，另一些人却因嫉妒而自取灭亡，这是由于那芬芳的气味——对某些人成为生命的芬芳，对另一些人成为死亡的芬芳（见：\BibleRef{格后 2:15}）。之后，主与拉匝禄——那个死了又复活的人——同席坐席，也斜躺着；有香膏倒在祂脚上，满屋充满香气；犹太人又显示出他们灵性的荒芜，竟想出杀死拉匝禄这般无用的残忍和极其愚蠢、疯狂的罪行。这一切，我们已因主的恩宠，在之前的讲道中尽力论述过了。现在，请您的爱德留意：在主受难之前，祂的宣讲结出了何等丰硕的果实，以色列家迷失的羊群中有多少听见了牧人的声音。\par

2. 因为你们刚才听到的福音经文记载：\Quote{第二天，有许多上来过节的人，听见耶稣来到耶路撒冷，便拿着棕榈枝出去迎接祂，喊说：贺三纳！奉主名而来的以色列王，当受赞美。}棕榈枝是赞美的象征，是胜利的标志，因为主要借着死来战胜死亡，并借着十字架的凯旋战胜死亡的元凶魔鬼。朝圣群众所喊的\Quote{贺三纳}一词，据一些通晓希伯来语的人指出，表达的是一种心境，而非有确切含义；正如在我们的语言中，有一些所谓感叹词，例如悲伤时喊\Quote{哀哉}，喜乐时喊\Quote{哈}，赞叹时喊\Quote{啊，何等美好！}——这句\Quote{啊}只表达赞叹者的感受。同样，我们认为这个词也是如此，因为它无论在希腊文还是拉丁文中都找不到确切的翻译，正如另一处经文：\Quote{谁若对自己的兄弟说拉卡。}这也被认为是一个表达愤怒的感叹词。\par

3. 但当他们说：\Quote{奉主名而来的以色列王，当受赞美}时，\Quote{奉主名}一句更宜理解为\Quote{奉天主圣父的名}，尽管也可以理解为\Emph{奉祂自己的名}，因为祂自己也是主。正如圣经在另一处也说：\Quote{主从天上……降下……火}（见：\BibleRef{创 19:24}）。但祂自己的话更能引导我们理解，因为祂说：\Quote{我因我父的名而来，你们却不接待我；如果有人因自己的名而来，你们反要接待他。}因为谦逊的真正导师是基督，祂谦卑自下，听命至死，且死在十字架上（见：\BibleRef{斐 2:8}）。但祂在教导谦逊时并未失去祂的天主性；一方面，祂与圣父同等；另一方面，祂与我们相似。借祂与圣父同等的位格，祂创造了我们；借祂与我们相似的位格，祂拯救了我们脱离毁灭。\par

4. 这些就是群众向耶稣欢呼的话：\Quote{贺三纳！奉主名而来的以色列王，当受赞美。}当犹太人的首领听见如此众多的群众宣称基督是他们的王时，他们的内心该是多么痛苦啊！但为主来说，作以色列的王算什么光荣呢？永恒之王的尊荣，降格作人类的王，又算得了什么大事呢？因为基督作以色列的王，并非为了征收赋税，不是将刀剑交在士兵手中，也不是以公开战争征服敌人；祂是以王权统治他们的内心，照顾他们永恒的利益，将那些信德、望德、爱德都集中于祂的人带入祂的天国。因此，天主圣子——圣父的同等者、万有借以受造的圣言——乐意作以色列的王，这乃是屈尊俯就，而非高升；是怜悯的表现，而非权力的增加。因为在地上被称为犹太人之王的那一位，在天上是天使的主。\par

5. \Quote{耶稣找到一匹小驴驹，就骑在上面；正如所记载的：\InnerQuote{熙雍女子，不要害怕！看，你的君王骑着驴驹来了。}}那么，在那一民族中确有熙雍女子；因为熙雍就是耶路撒冷。就在那一民族中——尽管他们是被弃绝的、眼瞎的——却有熙雍女子，曾有人对她说：\Quote{熙雍女子，不要害怕！看，你的君王骑着驴驹来了。}这位熙雍女子，蒙受了如此神圣的呼召，就在那些听见牧人声音的羊群之中，也在那支以如此虔诚的热忱庆祝主到来、以如此战斗的队列陪伴主前行的群众之中。对她说：\Quote{不要害怕}；认出你们现在赞美的那一位吧；当祂受苦时，不要害怕；因为借祂的血，你们的罪债要被涂抹，你们的生命要得到恢复。至于那匹从未有人骑过的驴驹（其他福音书对此有记载），我们应理解为尚未领受主法律的各外邦民族；而那匹驴（因为这两只动物都被带到主面前）则代表属于以色列民族、已驯服而认识主人槽头的祂的子民。\par

6. \BibleQuote{这些事祂的门徒起初不明白，但是当耶稣受光荣以后，才想起这些事是记载在祂身上的，并且他们向祂行了这些事。}也就是说，当祂彰显了祂复活的权能时，他们才想起这些事是记载在祂身上的，他们向祂行了这些事。总之，他们心中将主受难之前和受难期间所成就的事与圣经内容相比较，也发现这点，即祂骑在驴驹上符合先知的话。\par

7. \BibleQuote{因此，那些当祂叫拉匝禄出坟墓、使他从死者中复活时与祂同在的群众，都作见证。为此，群众也来迎接祂，因为他们听见祂行了这奇迹。法利塞人彼此说：你们看，我们一事无成？看，整个世界都去跟随祂了。}乌合之众推动乌合之众。\Quote{可是你们这些瞎眼的乌合之众，为什么嫉妒，因为世界去跟随了它的造物主？}\par

8. \BibleQuote{那时，在那些上来过节敬拜的人中，有些外邦人。他们来到加里肋亚的贝特赛达人斐理伯那里，请求他说：先生，我们愿意见耶稣。斐理伯就去告诉安德肋，然后安德肋和斐理伯去告诉耶稣。}我们要听主的回答。注意，犹太人想要杀祂，而外邦人却要见祂；但那些呼喊\BibleQuote{奉主名而来的以色列王，应受赞颂}的人也是犹太人。因此，这里有受割礼的和未受割礼的，就像两堵墙壁从不同方向走来，在基督的一个信德中，以平安之吻相遇。那么，让我们听角石的声音：\BibleQuote{耶稣回答他们说：人子受光荣的时候到了。}也许有人以为祂说祂受光荣是因为外邦人想要见祂。但并非如此。祂是看到了祂受难复活后，外邦人本身将在万民中来信仰祂，因为正如宗徒所说：\BibleQuote{一部分以色列人变得硬心，直到外邦人的全数进入。}\BibleRef{罗 11:25}因此，祂借着那些渴望见祂的外邦人，宣布了外邦民族未来的满全，并许诺祂自己受光荣的时刻临近，当这光荣在天上完成时，外邦民族将被带入信德。这正是那预言所指向的：\BibleQuote{天主，愿祢受尊崇于天，愿祢的荣耀遍及大地。}这就是外邦人的满全，宗徒说：\BibleQuote{一部分以色列人变得硬心，直到外邦人的全数进入。}\par

9. 但祂受光荣的高度必须先有祂受苦的深度。因此，祂接着说：\BibleQuote{我实实在在告诉你们：一粒麦子如果不落在地里死了，仍只是一粒；但如果死了，就结出许多果实。}但祂是指自己说的。祂自己就是那粒必须死了才能增多、因犹太人的不信而受死、又在许多民族的信德中获得增多的麦子。\par

10. 现在，作为鼓励跟随祂自己受苦之路，祂加上：\BibleQuote{爱惜自己性命的，必要丧失性命；}这可以有两种理解：\Quote{爱惜的，必要丧失，}就是说，你若爱惜，就要准备丧失；你若想在基督内拥有生命，就不要怕为基督而死。或者另一种：\BibleQuote{爱惜自己性命的，必要丧失性命。}不要因怕丧失而爱惜；不要在今世爱惜它，免得在永生中丧失它。但我最后所说的似乎更符合福音的意思，因为接着有这些话：\BibleQuote{在现世憎恨自己性命的，必要保存性命直到永生。}因此，当上一句说\BibleQuote{爱惜性命的}，应理解为\Emph{在今世}，他就是那必要丧失的。\BibleQuote{但憎恨的}，即在今世憎恨的，他就是那必保存性命直到永生的。这确实是一个深刻而奇异的宣告，关乎一个人对自己性命的爱导致其丧失，以及对他的恨确保其保全！如果你以有罪的方式爱它，那么你实际上恨它；如果你以符合善的方式恨它，那么你实际上爱它。那些如此恨自己的性命而保存它的人是有福的，他们的爱不会导致他们丧失性命。但要注意，不要怀有这样的观念：你可以因着这种理解你当在今世憎恨生命的义务而自寻死路。因为正是基于这种理由，某些错误而乖僻的人，他们对自己来说是特别残忍和邪恶的杀人犯，投身火海、在水中窒息、撞向悬崖而丧命。这不是基督的教导，相反，祂在回应魔鬼的悬崖试探时说：\BibleQuote{退到我后边去吧，撒殚！因为经上记载：你不可试探上主，你的天主。}\BibleRef{玛 4:7}祂也对伯多禄说，指明他要以怎样的死来光荣天主：\BibleQuote{你年少时，自己束上带，随意往来；但到年老时，别人要给你束上带，带你到你不愿意去的地方。}——祂在此清楚表明，跟随基督脚步的人必须不是由自己而是由他人来杀害。所以，当一个人的情况到了危机时刻，要么他必须违反天主的命令，要么离开此生，而迫害者以死亡威胁他，强迫他选择其中之一；在这样的境况中，让他宁愿选择死于天主的爱，也不愿生活于天主的愤怒；在这样的境况中，让他憎恨自己在今世的生命，以便将其保存到永生。\par

11. \BibleQuote{谁若服事我，就当跟随我。}这是什么意思？\BibleQuote{就当跟随我}，不就是说，当效法我吗？\BibleQuote{因为基督也为你们受了苦，}宗徒伯多禄说，\BibleQuote{给你们留下榜样，叫你们跟随祂的足迹。}\BibleRef{伯前 2:21}这里你便明白了这话的意义：\BibleQuote{谁若服事我，就当跟随我。}但结果如何？什么报酬？什么赏报？\BibleQuote{我在哪里，}祂说，\BibleQuote{我的仆人也必在那里。}让祂被白白地爱，这样服事祂的赏报就是与祂同在。因为离开祂，人能在哪里安好？或者与祂同在，人何时会感到不幸？再听更清楚的话：\BibleQuote{谁若服事我，我父必尊重他。}这尊重不是别的，正是与祂的圣子同在吗？因为祂先前所说\BibleQuote{我在哪里，我的仆人也必在那里}，我们可理解为祂在此\BibleQuote{我父必尊重他}时的解释。因为一个义子所能期望的更大尊重，岂非与独生子同在？固然不是被提升到祂天主性的同等，而是成为祂永生的分享者？\par

12. 但我们更当探究，究竟什么是那伴随如此大赏报的服事基督。因为若我们以为服事基督就是预备身体所需之物，或是烹饪、上菜，或是调酒、在餐桌递杯，那么，那些有幸亲身与祂同在之人确实如此服事过祂，如马尔大和玛利亚，那时拉匝禄也在坐席者中。但连那被弃的犹达斯也曾以这种方式服事基督；因他掌管钱囊，虽极其邪恶地偷取囊中之物，却也曾备办宴席所需。因此，当主对他说：\BibleQuote{你所做的，快去做吧！}时，有人以为祂只是吩咐他去预备节日的必需品，或拿些东西给穷人。所以，主所说\BibleQuote{我在哪里，我的仆人也必在哪里}和\BibleQuote{谁若服事我，我父必尊重他}，绝非指这类仆人；因为我们看见这样服事的犹达斯成了被弃绝的对象，而非被尊重。那么，何必往别处寻找这服事基督的含义，而不从这些话本身看出其启示呢？因为祂说\BibleQuote{谁若服事我，就当跟随我}时，祂愿意这话被理解为：谁若不跟随我，就不服事我。因此，那些不寻求自己的事，只寻求耶稣基督的事的人，才是耶稣基督的仆人。\BibleRef{斐 2:21}因为\BibleQuote{就当跟随我}正是：当行在我的路上，而非自己的路上；正如别处所载：\BibleQuote{那说自己住在祂内的，就应当照祂所行的去行。}\BibleRef{若一 2:6}因为若供应饥饿者食物，就当以怜悯而非夸耀的方式去做，在其中只寻求行善，不让左手知道右手所行；\BibleRef{玛 6:3}换言之，一切自求之心都当从爱德工作中彻底排除。如此服事的人就是服事基督，并要正确地对他说：\BibleQuote{你们对我这些最小兄弟中的一个所做的，就是对我做的。}\BibleRef{玛 25:40}这样，他不仅做那些关乎身体的爱德工作，而且为基督的缘故做一切善工（因为那时一切皆是善的，因为\BibleQuote{基督是法律的终向，使凡信的人都获得正义}\BibleRef{罗 10:4}），他就成为基督的仆人，甚至达到那特别的爱德工作——为弟兄舍命，因为那也是为基督舍命。因为为此祂也将代表祂的肢体说：你们对这些人所做的，就是对我做的。的确，正是论及这样的工作，祂也乐意自称为仆人，当祂说：\BibleQuote{就如人子来不是受服事，而是服事，并交出自己的生命，为众人作赎价。}\BibleRef{玛 20:28}因此，每个人是基督的仆人，正如基督也是仆人。这样服事基督的人，必被祂的父以殊荣尊重，即与祂的圣子同在，永远一无所缺于他的幸福。\par

13. 因此，弟兄们，当你们听见主说：\BibleQuote{我在哪里，我的仆人也必在哪里}时，不要只想到善良的主教和神职人员。你们自己也要以你们的方式，借着善生、施舍、按你们所能宣讲祂的名和道理，来服事基督；每一位家长也当承认这名称中他作为父母对家人的疼爱。为了基督和永生的缘故，让他警告、教导、劝勉、纠正全家的人，显出慈祥，施行管教；这样，在他自己的家中，他就要履行一种教会性的、类似主教的职务，服事基督，好与祂永远同在。因为连那最高贵的受苦服务，你们中许多人也已作出；许多既非主教也非神职人员，而是青年和贞女、年长者和年幼者、已婚男女、许多父母家长，都曾为基督服事，甚至为祂的缘故舍命殉道，并因接受极荣耀的荣冠而被父尊重。\par

\SourceRecord{Translated by John Gibb.；Nicene and Post-Nicene Fathers, First Series；Vol. 7.；Edited by Philip Schaff.；Buffalo, NY: Christian Literature Publishing Co.,；1888.；<http://www.newadvent.org/fathers/1701051.htm>.}

% structure: p0001 source=h1 target=section reason=page-title

\section{讲道集 52（\BibleBook{若望}{福音} 12:27-36）}

1. 主耶稣基督在昨日读经的话语中，曾激励祂的仆人跟随祂，并如此预言了祂自己的受难：除非一粒麦子落在地里死了，它仍旧是一粒；但如果死了，就结出许多子粒来；并曾激励那些愿意跟随祂进入天国的人，在这世界上憎恨自己的生命，如果他们想要保全它直到永生——祂又再次将自己的情感调适至我们的软弱，在我们的今日读经开始处说：\BibleQuote{现在我灵烦乱。}主啊，祢的灵魂从何而烦乱？祂确实在不久之前说过：\BibleQuote{在这世界上憎恨自己生命[灵魂]的，必会保全它直到永生。}难道祢在这世界上爱自己的生命，而当祢离开这世界的时刻临近时，祢的灵魂就烦乱吗？谁敢这样论说主的灵魂[生命]？毋宁说，祂将我们转移到了祂自己身上；祂将我们纳入祂自己的位格，作为我们的元首，并承担了祂肢体的情感；所以，祂并非因任何他人而烦乱，而是如祂在复活拉匝禄时所说的：\BibleQuote{祂在自己内烦乱。}因为天主与人之间的唯一\OriginalTerm{中保}{Mediator}，那人基督耶稣，既然将我们提升到天上的高处，也当与我们一同下降到受苦的最深处。\par

2. 我听见祂不久之前说：\BibleQuote{人子受光荣的时刻到了：一粒麦子若死了，就结出许多子粒来。}我也听见这话：\BibleQuote{在这世界上憎恨自己生命的，必会保全它直到永生。}我不仅被允许赞赏，更被命令效仿，因此，借着随后的话语：\BibleQuote{若有人服事我，就让他跟随我；我在哪里，我的仆人也要在那里，}我如火燃烧般轻视世界，在我眼中，这整个生命，无论多么长久，不过是一口蒸气；与我爱永恒事物相比，一切暂时的都已失去价值。而现在，我的主亲自以这些话奇妙地将我从我的软弱提升到祂的刚强，我却听见祂说：\BibleQuote{现在我灵烦乱。}这是什么意思？祢如何命令我的灵魂跟随祢，如果我在祢身上看见祢自己的烦乱？我怎能忍受被如此大能者所感受到的重负？如果磐石在动摇，我能寻找怎样的根基？但我仿佛在自己思想中听见主回答我说：你将更好地跟随我，因为我这样干预正是为了帮助你承受。你已听见，作为对你自己的训诲，我刚强的声音；现在在我身上聆听你软弱的呼声：我为你奔跑供应力量，我不阻止你的行进，但将你战栗的缘由转移到我自己身上，为你铺设前进的道路。主啊，我们的\OriginalTerm{中保}{Mediator}，在我们之上的天主，为我们而成为人，我承认祢的慈悲！因为祢，如此伟大者，因祢爱情的善意而烦乱，祢以丰富的安慰保全祢身体中许多因不断经验自身软弱而烦乱的人，使他们不至于在自己的绝望中完全丧亡。\par

3. 总之，愿那愿意跟随的人学习他必须行走的道路。也许可怕的考验时刻已经来到，摆在你面前的选择是：要么行不义，要么忍受苦难；软弱的灵魂烦乱了，而那天不能战胜的灵魂[耶稣]曾自愿烦乱；那么，将天主的旨意置于你自己的旨意之前。因为要注意你的创造者和导师随即补充的话语，祂造了你，并为教导你而成为祂所造的；因为那造人的成为了人，但祂仍然是不可改变的天主，并将人性移植到更好的境地。那么，请听祂在\BibleQuote{现在我灵烦乱}之后加上什么：\BibleQuote{我要说什么呢？父啊，救我脱离这时刻：但我正是为这时刻而来的。父啊，荣耀祢的名。}祂在这里教导你该想什么，该说什么，该呼求谁，该寄望于谁，并该将祂那确定而神圣的旨意置于你自己那人性而软弱的旨意之前。所以，不要想祂失去任何崇高地位，祂愿意你从毁灭的深渊中兴起。因为祂认为也应当受\OriginalTerm{魔鬼}{devil}的\OriginalTerm{诱惑}{temptation}，若非愿意，祂本不会受诱惑，正如若非愿意，祂本不会受苦；祂给魔鬼的回答正是你在诱惑之时也当使用的。\BibleRef{玛 4:1}确实，祂受了诱惑，却没有危险，为向你表明，当你在诱惑中处于危险时，如何回答诱惑者，以免被诱惑冲走，而脱离危险。但祂在这里说\BibleQuote{现在我灵烦乱}；以及当祂说\BibleQuote{我的灵魂忧伤至死}；和\BibleQuote{父啊，如果可能，请让这杯离开我}；祂承担了人的软弱，为教导人，当因此而忧伤烦乱时，要接着说：\BibleQuote{然而，父啊，不要照我的意愿，而要照祢的意愿。}\BibleRef{玛 26:38}因为当人将天主的旨意置于自己的旨意之前时，正是这样从人性转向神性。但\Quote{荣耀祢的名}这话所指何事，岂不正是祂自己的受难与复活？因为祂的意思无非是：父应如此荣耀子，子也照样在祂仆人的相似苦难中荣耀自己的名。因此，关于伯多禄，记载说祂对他如此说：\BibleQuote{别人要给你束上带子，带你到你不愿去的地方，}因为祂要指明\BibleQuote{他要以怎样的死来光荣天主。}因此，在他身上，天主也荣耀了祂的名，因为基督也这样在祂的肢体中荣耀了自己。\par

4. \Quote{随后有声音从天上传来，说：我已光荣了它，并要再光荣它。} \Quote{我已光荣了它，} 在创造世界以前，\Quote{并要再光荣它，} 当祂从死者中复活并升天的时候。也可以另作解释。\Quote{我已光荣了它，}——当祂由童贞女所生，当祂行使奇迹时；当贤士们（\OriginalTerm{贤士}{Magi}）在天上星辰的引导下俯伏朝拜祂时；当祂被充满圣神（\OriginalTerm{圣神}{Holy Spirit}）的圣人所认出时；当圣神（\OriginalTerm{圣神}{Spirit}）以鸽子形状降下，并由天上响起的声音公开宣告祂时；当祂在山上显圣容时；当祂行了许多奇迹，医治并洁净多人，用很少的饼喂饱了这么多人，命令风和浪，并复活了死人时——\Quote{并要再光荣它；} 当祂从死者中复活时；当死亡不再管辖祂时；当祂作为天主被高举于诸天之上，祂的光荣充满全地时。\par

5. \Quote{因此站在旁边的众人听见了，就说：打雷了；另有人说：有天使对祂说话。耶稣回答说：这声音不是为我，而是为你们来的。} 祂藉此显示，那声音并未向祂提示祂所已知的事，而是向那些需要得知的人。正如那声音是由天主发出的，不是为祂，而是为他人，同样，祂的灵魂烦乱，不是为自己，而是甘心为他人。\par

6. 请注意下文：祂说：\Quote{现在} \Quote{是这世界的审判。} 那么，我们在时间终结时应当期待什么呢？但最终所期待的审判，将是审判生者死者，给予永恒的赏罚。那么，现在的审判是怎样的呢？亲爱的诸位，我在以前的讲道中，已尽可能提醒你们：有一种审判，不是定罪的审判，而是分辨的审判；正如经上所载：\BibleQuote{天主，求你审判我，为我的案件向不圣洁的民辩诉（即分辨、辨别）。} 天主的审判众多；正如圣咏所说：\BibleQuote{你的判断是伟大的深渊。} 宗徒也说：\BibleQuote{啊，天主的智慧和知识的财富何其深奥！他的判断何其不可测度！} \BibleRef{罗 11:33} 主在此所说的审判也属于这一类：\Quote{现在是这世界的审判；} 而最终的审判则被保留，那时生者死者终将被审判。因此，魔鬼曾拥有对人类的控制权，凭着他们罪的债据把他们当作应受惩罚的罪犯拘留；他在不信者心中作王，欺骗奴役他们，引诱他们离弃造物主而敬拜受造物；但是，藉着对基督的信德——这信德由祂的死亡和复活所确认，并藉着祂为赦罪而倾流的血——成千上万的信徒从魔鬼的权势下获得释放，被联合于基督的身体，在这伟大的元首之下，由祂同一圣神所重生，成为祂忠信的肢体，进入新生命。这就是祂所说的审判，即这正义的分离，这魔鬼从祂所赎者中被驱逐出去。\par

7. 简言之，请注意祂自己的话。因为正如我们曾询问祂说\Quote{现在是这世界的审判}是什么意思，祂接着解释，说：\Quote{现在这世界的元首将被驱逐出去。} 我们如此听见的就是祂所指的审判。因此，不是那最终的审判，那时生者死者将被审判，有些被安置在祂右边，有些在左边；而是那审判，藉此\Quote{这世界的元首将被驱逐出去。} 那么，他在何种意义上是在内，而祂说他将被驱逐到哪里去呢？是否是这样：他是在世界内，并被驱逐到世界之外？因为如果祂说的是那最终将要来临的审判，有人可能会想到那永恒之火，魔鬼将和他的天使以及所有属于他的人被投进去——即不是由于本性，而是由于道德过失；不是因为魔鬼创造了或生了他们，而是因为他说服了他们并抓住了他们：因此，有人可能认为那永恒之火是在世界之外，而这就是\Quote{他将被驱逐出去}这句话的意思。但正如祂所说：\Quote{现在是这世界的审判，} 并且为解释其意，祂补充说：\Quote{现在这世界的元首将被驱逐出去，} 我们就此明白这是现在正在发生的事，而不是那要等到最后一天才发生的事。因此主预言了祂所知道的事：在祂自己受难和受光荣之后，全世界许多民族——魔鬼先前居于他们心中——都将成为信徒，而魔鬼因被信德所弃绝，就被驱逐出去。\par

但有人会说：\newline{}\Quote{那么，祂难道没有被驱逐出圣祖、先知和古时义人的心中吗？当然被驱逐了。那么，为什么说\Quote{现在他要被驱逐出去}呢？我们还能怎样理解，除非是：那时只在极少数人身上所发生的事，如今被预言将迅速在许多强大的民族中发生？正如另一句话：\Quote{因为圣神还没有赐下，因为耶稣还没有受光荣}也引起类似的问题，并得到类似的解答。因为先知们预言未来的事，并非没有圣神；年迈的西默盎和寡居的亚纳也并非没有藉圣神认出婴孩主\BibleRef{路 2:25}；匝加利亚和依撒伯尔也并非没有藉圣神在祂尚未出生、仅被孕育时，就预言了许多关于祂的事。但\Quote{圣神还没有赐下}，意思是说，还没有赐下那样丰富的属神恩宠，使聚集在一起的人能说各种语言\BibleRef{宗 2:4}，从而预先用各种语言宣告未来的教会；并且正是藉着这属神的恩宠，万民才被聚集到会众中，罪过从各处获得赦免，千千万万的人与天主和好。}\par

但有人会说：\newline{}\Quote{既然魔鬼就这样被逐出信徒的心，他现在难道不试探任何信友吗？}不，确实，他并未停止试探。但统治内心是一回事，从外面攻击是另一回事；因为最坚固的城有时也会被敌人攻击而不被攻取。如果他的箭矢有些射中我们，宗徒提醒我们如何使之无效，当他提到信德的甲胄和盾牌时\BibleRef{得前 5:8}。如果有时他伤了我们，我们也有近在咫尺的良药。因为正如战士们被告知：\Quote{我给你们写这些事，是要你们不犯罪}；那些受伤的人也要听接下来的话：\Quote{但若有人犯了罪，我们在父那里有一位中保，就是那义者耶稣基督；祂自己就是赎罪祭，赎我们的罪}\BibleRef{若一 2:1}。当我们说\Quote{宽免我们的罪债}时，我们祈求的岂不是治愈我们的创伤吗？当我们说\Quote{不要让我们陷于诱惑}时\BibleRef{玛 6:12}，我们求的岂不是：那如此窥伺我们或从外面攻击我们的人，从各方面都找不到入口，既不能以诡诈也不能以强力胜过我们吗？然而，无论他筑起什么攻城器械来攻击我们，只要信德所居的心中不再有他的位置，他就是被驱逐出去了。但是\Quote{若不是主看守城邑，守城者就白白醒着。}因此，不要自恃，免得那曾被驱逐的魔鬼又被召回心中。\par

另一方面，我们切不可设想魔鬼被称为\Quote{今世之王}，是认为他有权统治天地。如此说\Quote{世界}，是就散布全地的恶人而言；正如我们称一所房屋为善或恶，是就其居住者而言，并非指责或赞许墙和屋顶的建造，而是指其内住者或善或恶的品行。所以，类似地，称\Quote{今世之王}，即是一切住在此世的恶人之王。世界也被就善人而言，他们也同样散布全地；因此宗徒说：\Quote{天主在基督内，使世界与自己和好}\BibleRef{格后 5:19}。这些人就是今世之王从他们心中被驱逐出去的人。\par

因此，在说了\Quote{现在今世之王要被驱逐出去}之后，祂又接着说：\Quote{至于我，当我从地上被举起来时，我要吸引万有归向我。}这\Quote{万有}是什么，不就是那些从被驱逐者手下出来的人吗？但祂没有说\Quote{万民}，而是说\Quote{万有}；因为并非所有人都有信德\BibleRef{得后 3:2}。所以祂并非指全人类，而是指受造物个人完整的整体，即精神、灵魂和身体；或一切使我们成为有理智、有生命、可见、可触的存有之事物。因为那位说过\Quote{你们连一根头发也不会损失}\BibleRef{路 21:18}的，正是吸引万有归向祂的那一位。或者，如果\Quote{万有}是指人，那么我们可以说是指预定得救的一切人；关于这些人，祂在先前谈论祂的羊时曾宣告，一个也不会失落。诚然，各色人等——无论语言、年龄、地位、才能、合法而有益的艺术职业，以及一切按人类无限差异（除了罪之外）相互区分者，从最高到最低，从君王到乞丐——\Quote{万有，}祂说，\Quote{我要吸引归向我}；使祂成为他们的头，他们成为祂的肢体。但这将实现，祂说，\Quote{若我从地上被举起来}，即当我被举起时；因为祂对所来完成之事必定成就毫无疑虑。祂在此暗示祂先前所说：\Quote{但如果麦粒死了，才结出许多果实。}因为祂被举起除了指祂在十字架上的苦难之外，还意味着什么？福音作者自己并没有忽略这个解释；因为他补充了这句话：\Quote{祂这话是表明祂将要受什么样的死。}\par

12. \BibleQuote{群众回答祂说：我们从法律上得知：基督永远存在。你怎么说：人子必须被举起来？这人子是谁？}他们一直铭记着主常常自称人子。因为，在眼前的这段经文中，祂并没有说如果人子从地上被举起来；但祂在此之前，在昨天宣读和讲解的读经中，当那些渴望见祂的外邦人被宣告时，曾这样自称：\BibleQuote{时候到了，人子应当受光荣} \BibleRef{若 12:23}。因此，他们将这话铭记在心，并理解祂现在所说的\BibleQuote{当我从地上被举起来}是指十字架上的死亡，就询问祂说：\BibleQuote{我们从法律上得知：基督永远存在；你怎么说：人子必须被举起来？这人子是谁？}因为如果祂是基督，那么祂，他们说，永远存在；如果祂永远存在，祂又怎能从地上被举起来，即祂又怎能因十字架的苦难而死呢？因为他们理解祂所说的，正是他们自己图谋要做的事。所以，祂并没有通过传授智慧来为他们消除这些话语的模糊，而是通过激发他们的良心。\par

13. \BibleQuote{于是耶稣对他们说：光在你们中间还有片刻。}正因如此，你们明白基督永远存在。\BibleQuote{你们应当趁着有光的时候行走，免得黑暗笼罩你们。}行走，靠近，达到圆满的理解：基督既会死，也会永远活着；祂会流自己的血来救赎我们，并升到高处，带着被救赎的人同去。但若是你们对基督永恒性的信仰，到了拒绝承认祂曾受屈辱而死的程度，那么黑暗必会笼罩你们。\BibleQuote{那在黑暗中行走的，不知道往哪里去。}于是，他可能绊倒在那块绊脚石和使人跌倒的磐石上，那是主自己向瞎眼的犹太人成了的；正如对那些相信的人，匠人所弃的石头，成了房角的头块石头。\BibleRef{伯前 2:6}因此，他们认为基督不值得相信；因为他们以不敬的态度藐视祂的死亡，嘲笑祂被杀害的想法：然而，正是那粒麦子的死亡，导致了它自身的繁衍，以及那吸引万有跟随祂的一位被高举。\BibleQuote{你们应当趁着有光的时候，}祂接着说，\BibleQuote{相信光，好成为光明之子。}趁着你们还拥有所听见的一些真理，相信真理，好使你们在真理中重生。\par

14. \BibleQuote{耶稣说了这些话，就离开他们，隐藏了。}不是避开了那些开始相信并爱祂的人，也不是避开了那些拿着棕榈枝与赞美诗歌来迎接祂的人；而是避开了那些看见祂却恨祂的人，因为他们并没有看见祂，只是盲目地绊倒在那块石头上。但当耶稣躲避那些想要杀祂的人时（正如你们因健忘而需要经常被提醒一样），祂顾及了我们人性的软弱，却丝毫没有减损自己的权威。\par

\SourceRecord{Translated by John Gibb.；Nicene and Post-Nicene Fathers, First Series；Vol. 7.；Edited by Philip Schaff.；Buffalo, NY: Christian Literature Publishing Co.,；1888.；<http://www.newadvent.org/fathers/1701052.htm>.}

% structure: p0001 source=h1 target=section reason=page-title

\section{讲道集第53篇（若望福音12:37-43）}

当我们的主基督预言自己的受难，以及祂在十字架上被举起之死的丰硕果实，说祂将吸引一切归向祂自己；犹太人明白祂说的是自己的死，就问祂，当他们从法律中听到基督永远长存，祂怎能说死亡等着祂呢？祂劝勉他们，趁他们还有那一点光——那光教导他们基督是永恒的，行走，好使自己熟悉整个主题，免得被黑暗笼罩。祂说了这话，就隐藏起来。你们在以前的主日课程和讲道中已经熟悉了这些要点。\par

此后，福音作者提出构成今日读经简短主题的内容，说：\\Quote{祂在他们面前行了如此多的奇迹，他们却不信祂，为要应验依撒意亚先知所说的话：\\BibleQuote{主，谁相信了我们的报告？主的手臂向谁显露了呢？}}他在这里清楚表明，圣子本身就是主的手臂；并非天主圣父的位格被人的肉体形状限定，而圣子作为祂身体的一部分附着于祂；而是因为万有由祂造成，因此祂被称为主的手臂。因为你用你的手臂工作，同样，天主圣言被称为祂的手臂；因为祂藉圣言创造了世界。因为人为了做某件事伸开手臂，岂不是因为事情不会立即随他的话语成就？倘若他被赋予如此卓越的能力，以至于他说的话不用身体任何动作就能成就，那么他的话语就是他的手臂。但主耶稣，天主圣父的独生子，正如祂不是圣父身体的单纯部分，同样祂也不仅是可用思想理解、可听见、转瞬即逝的话语；因为万有由祂造成，祂就是天主圣言。\par

因此，当我们听到圣子是天主圣父的手臂时，不要让肉性的习俗在我们耳中扬起纷扰的喧嚣；但要按祂的恩宠所能，思考天主的能力和智慧，万有由之所造。确实，这样的手臂既不是伸展出去，也不是收缩回来。因为祂与圣父并非同一，但祂与圣父是一；且与圣父同等，祂在各方面都是圆满的，圣父也是同样：如此便不给那些可恶的错误留有余地，他们断言只有圣父存在，而根据原因的不同，祂自己有时被称为圣子，有时被称为圣神；他们也敢从这些话中推论：你看，你看出只有圣父存在，如果圣子是祂的手臂：因为人和他的手臂不是两个位格，而是一个。他们不理解也不思考，词语是如何因某些相互的相似性而从一个事物转移到另一个事物，即使在我们关于最熟悉可见事物的日常用语中也是如此；更何况为了以言语表达不可言说的事物，它们本身实际存在，根本不能用言语表达，岂不更是如此？因为甚至一个人，如果习惯于通过另一个人办事，就会称那人为自己的手臂：如果失去了他，就悲叹说：我失去了我的手臂；对夺走他的人说：你使我失去了我的手臂。那么，让他们理解圣子被称为圣父的手臂的意义，即圣父藉以完成了一切祂的工程；免得他们因不能理解，继续留在错误黑暗中，像那些犹太人，经上论他们说：\\BibleQuote{主的手臂向谁显露了呢？}\par

在此我们遇到第二个问题，要恰当处理它，探究它所有奥妙的曲折，并以合适的方式揭示出来，我认为这既超出我的能力，也超出有限的时间，以及你们的能力。然而，我们不愿过分使你们的期望落空，而不对此说些什么，就转到其他主题；请接受我们所能提供你们的：在未能满足你们期望的事上，向那指派我们栽种和浇灌的那位求增长；因为宗徒说：\\BibleQuote{栽种的不算什么，浇灌的也不算什么，只有使生长的天主算数。}\\BibleRef{格前 3:7}有些人心存嘀咕，有时在他们能够时公开发言，甚至爆发激烈的争论，说：犹太人做了什么？他们有什么错？如果为要应验依撒意亚先知所说的话：\\BibleQuote{主，谁相信了我们的报告？主的手臂向谁显露了呢？}是必然的。我们的回答是：主以其对未来的预知，藉先知预言了犹太人的不信；祂预见了，但并未造成。因为天主并不因为预先知道人将来的罪，就强迫任何人犯罪。祂预知的是属于他们的罪，而非祂自己的；那些罪不能归咎于任何人，只能归咎于他们自己。因此，如果祂预知为他们所有的，实际上不属他们，那么祂就没有真正的预知；但既然祂的预知绝无错误，无疑就是他们自己，而非别人，才是天主预知其罪恶的罪人。因此，犹太人犯了罪，并不是由于祂的强迫——罪是祂所不喜悦的；但祂预言他们要犯罪，因为没有什么能瞒过祂的知识。并且，倘若他们愿意行善而非作恶，必不会受阻碍；但他们将来要做的，已被那知道每个人将要做什么、并根据各人的行为予以报应的祂所预见。\par

5. 但接着福音的话语更加紧迫，并提出了一个更深奥的问题：因为祂接着说：\Quote{因此他们不能相信，因为依撒意亚又说：祂弄瞎了他们的眼睛，硬了他们的心，免得他们眼睛看见，心里明白，回转过来，我就医治他们。}因为我们被告知：如果他们不能相信，那么人若不能做什么，罪在哪里？如果他们因不信而犯罪，那么他们本有能力相信，却没有使用。那么，如果他们有能力，福音为何说：\Quote{因此他们不能相信，因为依撒意亚又说：祂弄瞎了他们的眼睛，硬了他们的心，}以至于（这是严重的）将不信的原因归于天主自己，因为是祂\Quote{弄瞎了他们的眼睛，硬了他们的心}？因为先知书中如此见证的，至少不是指着魔鬼说的，而是指着天主说的。因为若我们假设这是指着魔鬼说的，说祂\Quote{弄瞎了他们的眼睛，硬了他们的心}，我们就有责任能够指出那些被称为\Quote{他们不能相信}的人有何过错。那么，关于这位先知另一段证言，就是宗徒保禄所引用的，我们该如何回答？他说：\Quote{以色列没有获得他所追求的，但蒙选者获得了，其余的人都瞎了眼，正如经上所载：天主给了他们沉睡的神，眼睛看不见，耳朵听不见，直到今日。}\par

6. 弟兄们，正如你们刚才听到的，这就是我们面前的问题，你们可以看到它有多么深奥；但我们将尽我们所能给予回答。\Quote{他们不能相信，}因为依撒意亚先知预言了；先知预言是因为天主预知事情会如此。但如果有人问我为什么他们不能，我立刻回答：因为他们不愿意；因为天主预见了他们邪恶的意志，并藉着那位没有未来之事能向祂隐藏的先知预言了。但你说，先知指出了另一个原因，而非他们的意志。先知指出了什么原因？即\Quote{天主给了他们沉睡的神，眼睛看不见，耳朵听不见；祂弄瞎了他们的眼睛，硬了他们的心。}对此我也回答：这是他们的意志应得的。因为天主这样瞎了他们的眼、硬了他们的心，仅仅是通过放任不管、收回祂的帮助；天主能以隐藏的、却不是不义的审判这样做。这是一个敬畏天主者的敬虔应当坚定不移、完整无缺地持守的教义；正如宗徒在讨论同样棘手的难题时所说：\Quote{那么，我们说什么呢？难道天主有不义吗？绝对不是。}\BibleRef{罗 9:14}那么，如果我们绝不可认为天主有不义，那么只能是：当祂给予帮助时，祂是仁慈地行事；当祂收回帮助时，祂是正义地行事：因为祂所做的一切，都不是轻率的，而是按照审判。更进一步，如果圣人们的审判是正义的，那么圣化并使人成义的天主的审判岂不更是如此？因此，它们虽然是隐藏的，却是正义的。因此，当这类问题摆在我们面前，为什么这个人这样被对待，那个人那样被对待；为什么这个人因被天主遗弃而瞎眼，那个人因得到天主赐予的帮助而受光照——让我们不要擅自审判这伟大审判者的审判，而要战战兢兢地与宗徒一同呼喊：\Quote{哦，天主的智慧和知识的财富，多么深奥！祂的审判何其难测，祂的道路何其难寻！}\BibleRef{罗 11:33}正如在圣咏中也说：\Quote{祢的审判象深渊一般。}\par

7. 那么，弟兄们，你们爱德的期待不要驱使我试图闯入如此深渊，探测如此深渊，探究那不可探究的。我承认我自己的能力微小，我也认为对你们的能力也同样是微小的。这对我的身量太高，对我的力量太强；对你们，我想也是如此。因此，让我们一同听从圣经的劝告和话语：\Quote{不要探究高于你的事，不要寻求超过你能力的事。}并非这些事被禁止我们，因为神圣的导师说：\Quote{没有隐藏的事不被揭露的：}\BibleRef{玛 10:26}但是，如果我们按照我们目前所得的程度行走，那么正如宗徒告诉我们的，不仅我们所不知道而应当知道的事，而且如果我们还想要知道别的事，天主也会向我们启示这一点。\BibleRef{斐 3:15}但是，如果我们已经走上了信德的道路，让我们坚定不移地持守它：让它引导我们进入君王的居所，在祂里面藏着一切智慧和知识的宝藏。\BibleRef{哥 2:3}因为主耶稣基督自己对待那些伟大而特选的宗徒时，并非出于吝惜，祂说：\Quote{我还有好多事要告诉你们，但你们现在不能承受。}我们必须行走、进步、成长，使我们的心能够承受我们现在无法承受的事物。如果最后一天发现我们已充分长进，那时我们将会学到在此无法知道的事。\par

8. 然而，若有人自认能够并有足够信心，对摆在我们面前的这个问题给出更清楚、更好的阐述，但愿我不至于更乐意学习而非施教。只是，谁也别敢以如此方式辩护自由意志，以至于试图剥夺我们念诵以下祷词的权利：\BibleQuote{不要让我们陷入诱惑}\BibleRef{玛 6:13}；另一方面，谁也不可否认自由意志，从而胆敢为罪寻找借口。但让我们听从主，既在祂的命令中，也在祂的助佑中；祂既告诉我们当尽的本分，又协助我们履行它。因为有些人，祂任凭因过度信赖自己的意志而陷入骄傲，另一些人，祂则任凭因相反的过度不信任而陷入疏忽。前者说：为何我们求天主不让我们被诱惑所胜，既然这完全在我们的能力之内？后者说：为何我们该努力善生，既然行善的能力在天主手中？主啊，父啊，你在天上，不要让我们陷入这些诱惑中的任何一个；但\BibleQuote{救我们脱离凶恶！}\BibleRef{玛 6:13}请听主所说的话：\BibleQuote{我已为你祈求了，伯多禄，为叫你的信德不至丧失；}\BibleRef{路 22:32}这样我们永不会以为我们的信德如此依赖于我们的自由意志，以至于不需要天主的助佑。也让我们听从福音作者，他说：\Quote{祂赐给他们权能，成为天主的子女；}这样我们不会以为相信完全超出我们自己的能力之外；但在两者中，我们都应承认祂的恩惠作为。因为，一方面，我们当感谢祂赐予这权能；另一方面，要祈祷我们自身微薄的力量不至全然失败。正是这以爱德行事的信德，\BibleRef{迦 5:6}按照主赐给各人的尺度，\BibleRef{罗 12:3}好使那夸耀的，不夸耀自己，而夸耀于主。\BibleRef{格前 1:31}\par

9. 因此，他们不能相信，就不足为奇了，因为他们的意志如此骄傲，以至于不认识天主的正义，而想建立自己的正义；正如宗徒论他们说：\BibleQuote{他们没有顺服于天主的正义。}\BibleRef{罗 10:3}因为他们不是因信德，而是仿佛因行为，而自高自大；并因这自满而眼瞎，就绊倒在绊脚石上。因此有话说：\Quote{他们不能，}我们应当明白这是说他们不愿意；正如论到主、我们的天主所说：\BibleQuote{我们若不信，祂仍是信实的，祂不能否认自己。}\BibleRef{弟后 2:13}论到全能者，有话说：\Quote{祂不能。}所以，正如主\Quote{不能否认自己}是对天主圣意的赞美，他们\Quote{不能相信}则是归于人的意志的过错。\par

10. 看哪！我也如此说：那些自视甚高，以至于认为必须将如此多归于自己意志的力量，从而否认他们需要在正义生活中得到天主的助佑的人，不能相信基督。因为基督名号的单纯音节和基督的圣事，在人们对基督的信德本身受到抵抗时，毫无益处。因为对基督的信德，就是相信那使不虔敬的人成义的那一位；\BibleRef{罗 4:5}相信那位中保，没有祂的调停，我们无法与天主和好；相信那位救主，祂来寻找并拯救那失落的；\BibleRef{路 19:10}相信那位曾说：\Quote{没有我，你们什么也不能做。}因此，既然不认识那使不虔敬的人成义的天主的正义，而想建立自己的正义以满足骄傲者的心意，这样的人不能相信基督。所以，那些犹太人\Quote{不能相信：}并非人不能改好；而是只要他们的思想朝着这个方向，他们就不能相信。因此他们被弄瞎并硬了心；因为，他们否认需要天主的助佑，就不被助佑。天主预知了这些关于被弄瞎和硬心的犹太人的事，而先知藉祂的圣神预言了此事。\par

11. 但当祂加上：\Quote{他们应当悔改，我便医治他们，}是否有一个\Quote{不}字应被理解，即他们应当\Emph{不}悔改，与前面所说的\Quote{免得他们眼睛看见，心里明白}相连；因为在此显然也是说\Quote{且免得他们明白}？因为悔改本身同样是祂恩宠的恩赐，正如对祂所说的：\Quote{万军的天主，求祢使我们回转。}或者，我们是否也应理解这正是通过天主治愈方法的慈悲经验而实际发生的：即，他们因骄傲而乖戾的意志，并想建立自己的正义，就被撇下，正是为了被弄瞎；这样被弄瞎，以便他们绊倒在绊脚石上，满面羞惭；如此被谦卑下来后，才寻求主的名，不再寻求那膨胀他们骄傲的自己的正义，而是寻求那使不虔敬的人成义的天主的正义？因为这条道路本身对其中许多人有益，他们后来因恶行而充满悔恨，并相信了基督；祂曾为他们祈祷说：\BibleQuote{父啊！宽恕他们，因为他们不知道他们做的是什么。}\BibleRef{路 23:34}也正是关于他们的那种无知，宗徒说：\Quote{我为他们作证，他们对天主有热忱，但并非按着真知；}因为他接着又说：\BibleQuote{因为他们不认识天主的正义，而想建立自己的正义，便没有顺服于天主的正义。}\BibleRef{罗 10:2}\par

12. \Quote{当依撒意亚看见了他的光荣并谈到了他时，说了这些话。}依撒意亚所看见的，以及它如何指涉主基督，都可在他的书中读到并学会。因为他看见了他，不是照他的本来面目，而是以某种象征的方式，以适合先知所见神视本身必须具有的形态。同样，\Person{梅瑟}也看见了他，然而我们发现他向所看见的那位说：\Quote{如果我在你眼中获得恩宠，求你如今将你自己显示给我，使我清楚看见你；}\BibleRef{出 33:13}因为他看见了他，不是照他的本来面目。然而，当我们将来也要有此经验时，同一位圣史\Person{若望}在他书信中告诉了我们：\Quote{亲爱的诸位，[现在]我们是天主的子女，将来如何还未显明；因为我们知道，当他显现时，我们将相似他，因为我们将看见他与他本来一样。}\BibleRef{若一 3:2}他本可以说\Quote{因为我们将会看见他}，而不加上\Quote{与他本来一样}；但由于他知道，在祖先和先知中有人看见了他，却并非与他本来一样，因此他在说了\Quote{我们将会看见他}之后，又加上了\Quote{与他本来一样。}弟兄们，不要被那些断言圣父不可见、圣子可见的人所欺骗。这种断言是由那些认为圣子是受造物、且其理性不与\Quote{我与父原是一体}这话相协调的人所发出的。因此，就他与父平等的\OriginalTerm{天主的性体}{form of God}而言，圣子也是不可见的；但为了被人看见，他取了\OriginalTerm{奴仆的性体}{form of a servant}，并成了人的模样，\BibleRef{斐 2:7}而成为对人可见的。因此，他甚至在降生成人之前，就按他自己的喜悦，在他所支配的\OriginalTerm{受造物的形态}{creature-form}中，向人眼显示了自己，但并非他本来的样子。让我们藉着信德洁净我们的心，好使我们准备好迎接那不可言传、且可谓不可见的神视。因为\Quote{心里纯洁的人是有福的，因为他们要看见天主。}\BibleRef{玛 5:8}\par

13. \Quote{然而在官长中也有许多人信从了他；但因为法利塞人的缘故，他们不敢承认，以免被赶出会堂：因为他们喜爱人的光荣胜过于天主的光荣。}请看圣史如何指出并责备了一些人，这些人他虽说是信从了他的：他们若曾经从这信德之门前进，也必会克服那被宗徒所克服的对人光荣的贪爱。当宗徒说：\Quote{我绝不夸耀，只夸耀我们的主耶稣基督的十字架，藉着他，世界于我已被钉在十字架上，我于世界也被钉在十字架上。}\BibleRef{迦 6:14}为此，主自己当被人间的骄傲和邪恶的狂态所嗤笑时，也将他的十字架钉在那些信从他的人额头上，那在某种程度上是谦逊的居所，好使信德学会不以他的名为耻，并爱天主的光荣胜过于人的光荣。\par

\SourceRecord{Translated by John Gibb.；Nicene and Post-Nicene Fathers, First Series；Vol. 7.；Edited by Philip Schaff.；Buffalo, NY: Christian Literature Publishing Co.,；1888.；<http://www.newadvent.org/fathers/1701053.htm>.}

% structure: p0001 source=h1 target=section reason=page-title

\section{论讲第54篇（\BibleRef{若 12:44-50}）}

1. 当我们的主耶稣基督在犹太人中间说话，并行了如此多的奇迹时，有些人相信了，他们预定得永生，主也称他们为自己的羊；但有些人不相信，也不能相信，因为藉着天主神秘却并非不义的审判，他们被蒙蔽了心、变得刚硬，因为他们被那抵抗骄傲者、却赐恩宠给谦卑人的天主所舍弃（\BibleRef{雅 4:6}）。但在那些相信的人中，有些人的宣认到了这样的地步：他们拿了棕榈枝，在祂走近时迎接祂，在喜乐中将那宣认化为赞美的礼仪：而另有属于首领的人，没有勇气宣认他们的信德，怕被逐出会堂；福音作者用这样的话谴责他们：他们喜爱人的赞美胜过天主的赞美（\BibleRef{若 12:43}）。在那些不信的人中，有些人后来会相信，祂预见了这一点，当祂说：\Quote{当你们举起人子以后，你们就会知道我就是那一位。}但有些人将留在同样的不信中，并被今日的犹太民族效法；这民族不久以后在战争中被摧毁，按照论及基督的先知性证言，此后几乎分散到全世界。\par

2. 当事情处于这种状态，祂自己的苦难临近时，\Quote{耶稣呼喊说，}正如我们今日的读经开始时所载，\Quote{信我的，不是信我，而是信那派遣我的；看见我的，就是看见那派遣我的。}祂曾在某处说过：\Quote{我的道理不是我的，而是那派遣我的。}我们由此理解，祂称自己的道理就是祂自己，即圣父的圣言；而当祂说\Quote{我的道理不是我的，而是那派遣我的，}时，意思是祂并非出于自己，而是从另一位得到祂的存在。因为祂是天主出于天主，圣父的圣子；但圣父不是天主出于天主，而是天主，圣子的圣父。如今当祂说\Quote{信我的，不是信我，而是信那派遣我的，}时，我们还能如何理解，除非祂以人的形像出现在人前，而作为天主祂仍是不可见的？为使无人以为祂不过是他们所见的那样，祂表明自己愿意被相信为在本质和地位上与圣父相等，当祂说\Quote{信我的，不是信我，}即仅凭我所显现的，\Quote{而是信那派遣我的，}即圣父。但信圣父的人，必须相信祂是圣父；而相信祂是圣父的人，必须相信祂有一位圣子；这样，信圣父的人就必信圣子。但不要让人对独生子所怀的信德，如同对那些因恩宠而非因本性与天性而被称为天主的子女的人所怀的信仰一样，如福音作者所说：\Quote{祂赐给他们权能，成为天主的子女，}并按照主自己在律法中所引述的话：\Quote{我曾说：你们是神，都是至高者的子女。}因为祂说\Quote{信我的，不是信我，}是为了表明我们在基督内的全部信德不应局限于祂的人性。因此，祂说，信我的，不是仅凭所见而信我，而是信那派遣我的；这样，由于信圣父，他相信圣父有一位与自己同等的圣子，从而达致对我的真实信德。因为如果有人以为天主仅有按恩宠而生的子女，他们不过是被造的，而非圣言，而是由圣言所造的，且祂没有一位与祂同等、同永恒、永远生出、同样不变、毫无差异与低下的圣子，那么他就不信那派遣我的圣父，因为那派遣我的圣父并非这样的概念。\par

3. 因此，在说了\BibleQuote{那信我的，不是信我，而是信那派遣我来的}之后，为了使父不会被理解为仅仅是藉恩宠而重生的众子之父，而不是祂自己同等者的独生圣言之父，祂立刻加上说：\BibleQuote{看见我的，就是看见那派遣我来的。}\newline{} 祂这里是否说：看见我的，不是看见我，而是看见那派遣我来的，正如祂所说：\BibleQuote{信我的，不是信我，而是信那派遣我来的}？因为祂说前一句话，是为了使人不是仅仅按祂当时所显现的样式，即作为人子，而信祂；后一句话，是为了使人信祂与父同等。那信我的，不仅信祂所见的我的外表，而是信那派遣我来的。或者，当他信那生我的父（祂自己的同等者）时，让他信我，不是按他所看见的我，而是如同信那派遣我来的；因为真相如此深远，在我与祂之间没有距离，以致看见我的，就是看见那派遣我来的。\newline{} 诚然，主基督亲自派遣了祂的宗徒，正如他们的名称所表明的：因为那些在希腊文中称为\Emph{angeli}的，在拉丁文中称为\Emph{nuntii}（使者），同样，希腊文的\OriginalTerm{宗徒}{apostoli}就成了拉丁文的\Emph{missi}（被派遣者）。但宗徒中绝无一人敢说：\Quote{信我的，不是信我，而是信那派遣我来的}；因为无论如何他不会说\Quote{信我的}。我们信一位宗徒，但我们不信他；因为成义非因宗徒，而是因那使不虔敬者成义的天主，他的信德被算为正义。\BibleRef{罗 4:5} 一位宗徒可以说：\Quote{接纳我的，就是接纳那派遣我来的}；或\Quote{听我的，就是听那派遣我来的}；因为主亲自这样告诉他们：\BibleQuote{接纳你们的，就是接纳我；接纳我的，就是接纳那派遣我来的。}\BibleRef{玛 10:40} 因为主人在仆人身上受尊敬，父亲在儿子身上受尊敬；但那时父亲仿佛在儿子内，主人仿佛在仆人内。然而，独生子能正当地说：\BibleQuote{信天主，也信我}；同样，祂在这里说：\BibleQuote{信我的，不是信我，而是信那派遣我来的。}祂并没有将信者的信心从自己身上转开，而是不愿信者停留在仆人的形象上；因为凡信那派遣祂的父的，立刻也信了子，知道没有子，父就不成其为父，这样他的信心就达到相信祂与父平等，与随后的话相符：\BibleQuote{看见我的，就是看见那派遣我来的。}\par

4. 注意接下来：\BibleQuote{我来到世界上为光，使凡信我的，不留在黑暗中。}祂曾在某处对祂的门徒说：\BibleQuote{你们是世界的光；建在山上的城是不能隐藏的；人点灯，不放在斗底下，而是放在灯台上，就照亮家中所有的人；照样，你们的光也当在人前照耀，好使他们看见你们的善行，光荣你们在天上的父。}\BibleRef{玛 5:14} 但祂并没有对他们说：\Quote{你们来到世界上为光，使凡信你们的，不留在黑暗中。}我断言，这样的说法在任何地方都找不到。因此，诸位圣人都是光，但他们是由祂藉信德而光照；凡与祂分离的，必被黑暗笼罩。但那光照他们的光，不能与自身分离，因为它完全超越变化。我们相信那如此点燃的光，如先知或宗徒；但我们相信他，是为了这个目的：使我们不信那本身被光照的，而是与他一起信那光照了他的光，好叫我们也被光照，不是被他，而是与他一起，被同一光所照。当祂说：\BibleQuote{使凡信我的，不留在黑暗中}，祂充分表明，祂发现众人都在黑暗中；但为了使他们不留在所发现的黑暗中，他们应当信那来到世界的光，因为世界就是藉着祂造成的。\par

5. \BibleQuote{若有人}，祂说，\BibleQuote{听我的话而不遵守，我不审判他。}请记着你们在以前讲座中听过的；若有人忘记了，就回忆起来；那时缺席而如今在场的，请听子如何说\BibleQuote{我不审判他}，而在另一处祂说：\BibleQuote{父不审判任何人，但把一切审判都交给了子}；这就意味着我们应当理解：如今不是审判他的时候。为何不是如今？请听下文：\BibleQuote{因为我来不是为审判世界，而是为拯救世界}；也就是说，为使世界进入得救的状态。如今是仁慈的季节，以后将是审判的时候；因为祂说：\BibleQuote{上主，我要向祢歌颂仁慈和审判。}\par

6. 但也要看祂关于末日审判所说的话：\BibleQuote{谁轻视我，不接受我的话，自有审判他的：我所讲的话，要在末日审判他。}祂并不说：谁轻视我，不接受我的话，我末日不审判他；因为祂若这样说，我看不出怎能不与另一句话矛盾，即祂所说：\BibleQuote{圣父不审判任何人，但把一切审判交给了圣子。}但当祂说：\BibleQuote{谁轻视我，不接受我的话，自有审判他的，}并为了告诉那些等候听谁是这个审判者的人，接着又说：\BibleQuote{我所讲的话，要在末日审判他，}祂便充分表明了自己后来要做审判者。因为祂说的是关于自己，宣布的是自己，并将自己展示为门，祂作为牧人由此进入羊群。因此，那些从未听过这话的人受审判是一种方式，那些听过却轻视的人受审判是另一种方式。\BibleQuote{因为凡没有法律而犯了罪的人，}宗徒说，\BibleQuote{也必没有法律而丧亡；凡在法律之下犯了罪的人，必凭法律受审判。}\BibleRef{罗 2:12}\par

7. \BibleQuote{因为我没有，}祂说，\BibleQuote{凭自己说话。}祂说祂不是凭自己说话，因为祂不是出于自己。关于这一点我们已经多次谈论过；所以现在，无需更多讲解，我们只需提醒你们这个你们熟悉的真理。\BibleQuote{但派遣我的圣父，祂给了我一条诫命，叫我该说什么，该讲什么。}我们本不打算详述这一点，如果我们知道现在是在对以前曾与我们交谈过的人说话，而且不是对所有人，只是对那些还记得所听之事的人；但因为也许现在有一些听众没有听过，也有一些听过但已经忘记的人，为了他们，愿那些记得所听之事的人容忍我们的耽搁。圣父如何给祂的独生子一条诫命？祂以什么言语对圣言说话，既然圣子本身就是独生的圣言？可能是藉着一位天使吗？既然天使是由祂创造的？或者是藉着云彩，当云彩发出声音传给圣子时，这声音并非为祂的缘故，正如祂自己也在别处告诉我们的，而是为了那些需要听见它的别人的缘故\BibleRef{若 12:29}？或者是藉着任何从嘴唇发出的声音，但那里没有形体，也没有把圣子和圣父分开的空间距离，以致有空气居于其间，使其震动生效而成为可听见的声音？让我们抛弃所有关于那无形体且不可言喻之存在的不相称观念。独生子就是圣父的圣言和智慧，其中\Emph{存有}圣父的一切诫命。因为圣子从来没有一个时候不知道圣父的诫命，以致需要祂在时间过程中拥有祂先前所没有的。因为祂从圣父所受的，是在被生时所受的，并在被生育时赐予了祂。因为祂是生命，祂确实在被生时接受了生命，然而并没有一个在先的时候，生命在祂的位格存在中是缺乏的。因为一方面，圣父\Emph{拥有}生命，并且祂\Emph{就是}祂所拥有的：然而祂并未接受生命，因为祂不是出于任何人。但圣子接受了生命作为圣父的恩赐，祂是由圣父而出的：因此祂自己就是祂所拥有的；因为祂拥有生命，并且就是生命。请听祂自己所说的话：\BibleQuote{就如圣父在自己内有生命，同样也赐给圣子在自己内有生命。}难道祂能赐给一个已经存在却迄今仍缺乏生命者吗？相反地，正是在生育本身，由那位生育生命者所赐予，因此生命生了生命。而为了表明祂所生的生命是同等而非低等，这样说：\BibleQuote{就如\Emph{祂}在自己内有生命，同样也赐给圣子在自己内有生命。}祂赐予了生命；因为在生育生命时，祂赐给祂的，除了是生命之外，还能是什么呢？并且因为祂的诞生本身就是永恒的，从来没有一个时候没有那位是生命的圣子，也没有一个时候圣子自己是没有生命的；而且因为祂的诞生是永恒的，所以祂，这样被生的那位，就是永生。因此，圣父并没有赐给圣子一条祂原先没有的诫命；但是，如我所说，在圣父的智慧中，即在圣父的圣言中，存储着圣父的一切诫命。然而，诫命被说成是赐给了祂，因为这样领受的那位不是出于自己：而将圣子从未缺少的东西赐给圣子，其意义等同于生育那位从未缺少存在的圣子。\par

8. 以下是接下来的话：\Quote{我知道祂的诫命就是永生。}因此，如果圣子本身就是永生，而圣父的诫命也是永生，那么这不就是说：我就是圣父的诫命吗？同样地，在祂接着说\Quote{凡我所说的，都是依照圣父对我说的，我就这样说}这句话中，我们不要将\Quote{对我说的}理解为仿佛圣父向唯一的圣言说了话语，或者天主的圣言需要来自天主的言语。圣父对圣子说话的方式，与祂赐予圣子生命的方式相同；并非因为圣子不认识圣父，或没有圣父，而正是因为是\Emph{圣子}。那么，这些话的意思是什么？\Quote{就如祂对我说的，我就这样说}，不就是：我说真理吗？因此，前者作为真实者所说的，后者作为真理也如此说。真实者生了真理。那么，祂现在还能对真理说什么呢？因为真理没有任何缺陷需要用额外的真理来补充。因此，祂对真理说话，是因为祂生了真理。同样地，真理本身也照祂所听到的说话；但只是对那些有理解力、并蒙真理——天主所生的真理——教导的人。但为使人们相信他们尚未有能力理解的事，从祂的人性嘴唇发出了可听见的话语；迅速消逝的声音传入耳中，迅速完成其短暂的持续；但这些声音所标记的真理本身，却仿佛进入了听众的记忆，并且也通过文字符号作为视觉标记传达给我们。但真理说话并非如此；祂内在地对聪明者的灵魂说话；祂无需声音来教导，而是用理解的亮光充满心灵。因此，凡能在那光中看见祂诞生的永恒性的人，他本人也以同样的方式听见真理说话，正如真理听见圣父告诉祂该说什么。祂在我们内唤醒了对祂内在临在的那甘饴体验的巨大渴望；但我们是藉着每日的成长而获得它的；我们是藉着行走而成长的，我们藉着向前努力而行走，以便最终能够达到它。\par

\SourceRecord{Translated by John Gibb.；Nicene and Post-Nicene Fathers, First Series；Vol. 7.；Edited by Philip Schaff.；Buffalo, NY: Christian Literature Publishing Co.,；1888.；<http://www.newadvent.org/fathers/1701054.htm>.}

% structure: p0001 source=h1 target=section reason=page-title

\section{讲道集第五十五（若望福音第十三章第一至五节）}

一、主之晚餐，按若望所载，赖主助佑，当以适当次数的演讲来阐述，并以主乐赐之能力加以解释。\Quote{逾越节前，耶稣知道自己离此世归父的时辰已到，他既然爱了世上属于自己的人，就爱他们到底。} \OriginalTerm{逾越节}{pascha}并非如某些人所想，是希腊文，而是希伯来文；然而在这名词中，两种语言有一种非常恰当的吻合。因为希腊文\OriginalTerm{受苦}{paschein}意为\Quote{受苦}，故\OriginalTerm{逾越节}{pascha}曾被误为意为\Quote{苦难}，似乎这名词源于主的苦难；但在其本语言，即希伯来文中，\OriginalTerm{逾越节}{pascha}意为\Quote{逾越}，因为当主的子民首次过逾越节时，他们从埃及逃跑，\Emph{逾越了}红海。如今这预表的预象在真理中实现了，当基督被牵到宰杀之地，\BibleRef{依 53:7}借他的血涂在我们的门框上，即他的十字架记号印在我们额上，使我们免于这世界将临的毁灭，正如以色列免于埃及人的奴役和毁灭，\BibleRef{出 12:23}而我们做一次极有益的逾越，从魔鬼逾越到基督，从这动荡的世界逾越到他稳固的王国。因此我们确实是逾越到永恒的天主那里，以免与这流逝的世界一同消逝。宗徒在赞美天主赐予我们如此恩宠时说：\Quote{他救了我们脱离黑暗的权势，把我们迁到他爱子的国里。}\BibleRef{哥 1:13}因此，这\OriginalTerm{逾越节}{pascha}之名，正如我所说，拉丁文称为\OriginalTerm{逾越}{transitus}，被蒙福的福音作者为我们解释，当他说：\Quote{逾越节前，耶稣知道自己离此世归父的时辰已到。}这里你看到既有\OriginalTerm{逾越节}{pascha}，又有\Emph{逾越}。他从何处、往何处逾越？即\Quote{离此世归父。}如此，肢体在他们元首中得了希望，他们必将跟随那先行\Quote{逾越}的主。那么，不信者呢？他们完全远离这元首和他的肢体，难道他们不也逾越吗？因为他们也不会永远留在此地。他们显然也逾越：但一是从世界逾越，一是与世界一同消逝；一是逾越到父那里，一是逾越到仇敌那里。因为埃及人也逾越了（海）；但他们并非从海中到王国，而是在海中遭毁灭。\par

二、\Quote{耶稣知道}，\Quote{自己离此世归父的时辰已到，他既然爱了世上属于自己的人，就爱他们到底。}无疑是借着他的爱，使他们也能从他们现今所在的这世界，逾越到已先他们而去的元首那里。那么，\Quote{到底}这话是什么意思？岂不是指基督？因为宗徒说：\Quote{法律的终向就是基督，使所有信者获得正义。}\BibleRef{罗 10:4}终向是成全的，而非消灭的；是我们达到的终向，而非在其中灭亡的终向。我们正应如此理解这话：\Quote{基督，我们的逾越节，已被祭献。}\BibleRef{格前 5:7}他就是我们的终向；我们逾越到他那里。因为我看这些话也可按人的意义来领会，即基督爱自己的人甚至爱到死，这样\Quote{爱他们到底}的意思就是直到死亡。这种意义是属人的，而非属神的：因为他爱我们并非仅仅到此为止，他永恒无尽地爱我们。天主禁止那死亡不能结束他的爱，竟在死亡时结束了他的爱。死后，那骄傲不虔的富人仍爱他的五个兄弟；\BibleRef{路 16:27}难道我们能认为基督只在死亡之前爱我们吗？断乎不可，亲爱的诸位。若他的爱只到死亡为止，那么他以这持续到死的爱来爱我们，便是徒然了。但也许\Quote{爱他们到底}应这样理解：他如此爱他们，以至于为他们死。因为他说：\Quote{人若为自己的朋友舍命，没有比这更大的爱了。}我们当然不反对将\Quote{爱他们到底}这样理解，即正是他的爱把他带到死亡。\par

三、\Quote{吃晚饭的时候，}他说，\Quote{魔鬼已把出卖耶稣的意思放在西满的儿子依斯加略人犹大的心里，耶稣知道父已把万有交在他手中，且知道自己是从天主出来的，又要往天主那里去；就离席站起来，脱了外衣，拿一条手巾束上腰。随后把水倒在盆里，就洗门徒的脚，并用自己所束的手巾擦干。}我们不应将\Quote{吃晚饭的时候}理解为晚餐已结束；因为主离席洗门徒的脚时，晚餐仍在进行。因为他随后又坐下，把那饼块递给出卖他的人，这显然表明晚餐尚未结束，或者说，桌上仍有饼。因此，\Quote{吃晚饭的时候}意思是晚餐已备好，摆在桌上供客人享用。\par

4. 但祂说：\Quote{魔鬼已经把出卖祂的念头投入了西满的儿子犹达斯依斯加略的心里}，如果有人追问：投入犹达斯心里的是什么？无疑就是\Quote{出卖祂。}这种\Quote{投入}是一种精神性的暗示；它非由耳朵进入，而是通过思想；因此不是以有形的方式，而是以精神的方式。因为我们所谓的精神性的，并不总是在褒义上理解。宗徒知道某些邪恶的精神性势力，在天空中的邪恶领域，他证明我们必须与之战斗\BibleRef{弗 6:12}；如果没有邪恶的鬼神，就不会有邪恶的精神性势力。因为精神性事物正是从精神体得名。但这种事如何发生，即魔鬼的暗示如何被引入并如此与人的思想混合，以致人以为那是自己的，这谁能知道？我们也不可怀疑：善的暗示同样由善神以同样不可觉察的精神方式作出；但关键在乎人的心赞同哪一方，或因应受遗弃而缺少天主的助佑，或因蒙恩而得到天主的助佑。因此，在犹达斯心里，因魔鬼的煽动，已经决定了这个门徒要出卖他的师傅，他并未学会认识他是自己的天主。他如今就是怀着这样的心态来赴他们的社交聚餐：他是牧者的间谍，救赎者的阴谋者，救主的出卖者；他如今来，被人看见，被人容忍，自以为未被发现，因为他被他想要欺骗的那位欺骗了。但那位早已洞察那颗心的内在状况，明知地利用着那个不明真相的人。\par

5. \Quote{耶稣知道父已将一切交在祂手中。}因此，连那出卖者也在祂手中：因为如果祂不把他掌握在手中，祂当然不能照祂的意愿使用他。因此，这出卖者已被出卖给那位他想出卖的；他行出卖的恶谋，其结果对那被出卖者却是一种他所不知的善。因为主知道祂为朋友做什么，并耐心地利用祂的敌人：如此，父将一切交在祂手中，包括恶者为现世之用，善者为最终结局。\Quote{也知道祂是从天主来的，又要往天主那里去：}祂来从天主时并未离开天主，回返时也未离开我们。\par

6. 既知道这些事，\Quote{祂就从席间起来，脱了外衣，拿了一条毛巾束在腰间。随后把水倒在盆里，开始洗门徒的脚，并用束在腰间的毛巾擦干。}亲爱的诸位，我们应仔细留意福音作者的意思，因为当他正要谈论主那卓越的谦卑时，他首先愿意彰显祂的尊威。正是为此，他说：\Quote{耶稣知道父已将一切交在祂手中，也知道祂是从天主来的，又要往天主那里去。}因此，正是那父已将一切交在祂手中的那位，如今在洗的，不是门徒的手，而是他们的脚；正是祂在知道自己是从天主而来，并要往天主那里去时，执行了仆人的职分，不是天主主的职分，而是人的职分。他也乐意预先提到祂的出卖者，他如今已这样来了，且不为祂所不知；这样，在知道祂早已预见那双将要浸在恶行中的手之后，仍然不顾尊严地也洗那人的脚，这更彰显祂谦卑的宏大。\par

7. 但我们何必惊奇祂从席间起来，脱了外衣？祂本具有天主的形体，却使自己空虚。又何必惊奇祂用毛巾束腰？祂取了奴仆的形体，成了人的形状\BibleRef{斐 2:6}。何必惊奇祂将水倒在盆里，洗门徒的脚？祂倾流自己的血在地上，以洗去他们罪的污秽。何必惊奇祂用束腰的毛巾擦干所洗的脚？祂以那包裹祂的肉身，为福音使者的脚铺了坚实的道路。的确，为了用毛巾束腰，祂脱下了所穿的外衣；但当祂空虚自己，取了奴仆的形体时，祂并非放下自己所有的，而是取了祂以前所没有的。当祂要被钉十字架时，祂被剥去外衣，死后被用细麻布包裹。祂全部的苦难，就是我们的净化。因此，当祂要受那极端的痛苦时，祂预先以这友善的顺从演示了它；不仅为那些祂将为之忍受死亡的人，也为那决意置祂于死地的人。因为人的谦卑恩惠如此伟大，以致连天主性的尊威也乐意以自己的榜样加以推崇；因为骄傲的人将永远丧亡，若不是被谦卑的天主寻获。因为人子来，是为寻找并拯救那失落的\BibleRef{路 19:10}。他因效法欺骗者的骄傲而失落，如今当他被寻获时，愿他效法救赎者的谦卑。\par

\SourceRecord{Translated by John Gibb.；Nicene and Post-Nicene Fathers, First Series；Vol. 7.；Edited by Philip Schaff.；Buffalo, NY: Christian Literature Publishing Co.,；1888.；<http://www.newadvent.org/fathers/1701055.htm>.}

% structure: p0001 source=h1 target=section reason=page-title

\section{\WorkTitle{讲道集 56（若 13:6-10）}}

1. 当主洗门徒的脚时，\BibleQuote{祂来到西满伯多禄跟前；伯多禄对祂说：主，祢洗我的脚吗？}因为谁不为天主子洗自己的脚而感到战栗呢？因此，虽然仆人违抗主人、受造物违抗天主是极大的胆大妄为，但伯多禄宁愿如此，也不愿自己的脚被主天主洗。我们也不该以为伯多禄是众多表达畏惧和拒绝的人中的一员，因为在他之前，其他人已经欣然并平静地接受了这事。因为更容易这样理解福音的话，在说了\BibleQuote{祂开始洗门徒的脚，并用祂束着的毛巾擦干}之后，接着加上\BibleQuote{随后祂来到西满伯多禄跟前}，仿佛祂已经洗了一些人的脚，然后才来到他们中的首位。因为谁不知道真福伯多禄是宗徒之首呢？但我们不应理解为祂是在其他人之后才来到他跟前，而是祂是从他开始的。因此，当祂开始洗门徒的脚时，祂来到祂开始的那人，即伯多禄跟前；于是伯多禄因任何人都会有的恐惧而害怕，说：\BibleQuote{主，祢洗我的脚吗？}在这个\Quote{祢}字里隐含着什么？在\Quote{我的}里又隐含着什么？这些是思考的问题，而非言语的问题；恐怕灵魂对这些话所形成的任何概念，在言语表达时都会有所欠缺。\par

2. 但\BibleQuote{耶稣回答他说：我所做的，你现在不知道，但以后会明白。}尽管他因主行动的崇高而恐惧，却仍不允许这事发生，因他不明白其目的；他不愿看到，也无法忍受基督如此屈尊至他的脚前。\BibleQuote{祢永远不可}，他说，\BibleQuote{洗我的脚。}这个\Quote{永远}[\OriginalTerm{永远}{in æternum}]是什么意思？我永远不会忍受、永远不会允许、永远不会同意：也就是说，一件事若从未做过，就不是\Quote{\Emph{in æternum}}做的。然后救主为了用他自身得救的危险来恐吓这位不愿的病人，说：\BibleQuote{我若不洗你，你就与我无分。}祂这样说：\Quote{我若不洗你}，当时祂只指他的脚；正如常说的：你在践踏我，实际只是脚被践踏。现在，那另一人在爱与恐惧的激动中，更害怕基督被夺走，甚至比看到祂在自己脚前卑微还要害怕，喊道：\BibleQuote{主，不但是我的脚，连手和头也要洗。}既然这确实是祢的威胁，即我的身体各肢必须由祢来洗，我不仅不再保留最低的部分，而且把最首要的部分也交在祢手中。求祢不要不给我与祢分享的部分，我也不会拒绝祢洗我身体任何部分。\par

3. \BibleQuote{耶稣对他说：洗过澡的人，不需要再洗，只需洗脚，便全然洁净了。}也许有人会因此被激发而说：不，但如果他全然洁净，为何还需要洗脚呢？但主知道祂所说的，即使我们的软弱无法触及祂隐秘的计划。尽管如此，就祂乐意从祂的法律中教导我们而言，按我有限的理解，我也愿在祂的帮助下，对这个问题的深奥之处作出一些回答：首先，我将不难证明在表达方式上并无矛盾。因为谁不能像这里一样，极其恰当地说：他全然洁净，除了他的脚？——尽管更优雅的说法是：他全然洁净，只差他的脚；意思是一样的。因此，主说：\BibleQuote{他不需要再洗，只需洗脚，便全然洁净。}即全然，除了他的脚，仍需洗。\par

4. 但这是什么？它意味着什么？其中有什么需要我们检查的？主说，真理宣布，即使是已经洗过的人，仍需要洗脚。我的弟兄们，你们对此有何看法？岂不是在圣洗圣事中，人全身都被洗了，不是除了脚之外，而是全然洁净；然而，此后在人间生活时，他无法避免用脚踩地。同样，我们人类的情感本身，与我们必死的尘世生活密不可分，就像脚一样，使我们与人事产生感官接触；而且，如果我们说我们没有罪，就是欺骗自己，真理不在我们内。\BibleRef{若一 1:8}因此，每日有那为我们转求者，\BibleRef{罗 8:34}洗我们的脚；而我们自己也需要每日洗脚，即端正我们属灵脚步的道路，这在主祷文中我们也承认，当我们说：\BibleQuote{免我们的债，如同我们也免了欠我们债的人。}\BibleRef{玛 6:12}因为\BibleQuote{如果}，正如经上所写，\BibleQuote{我们承认我们的罪过}，那么，那位洗了门徒脚的主，\BibleQuote{是忠信正义的，必赦免我们的罪过，并洗净我们一切的不义，}\BibleRef{若一 1:9}即甚至洗我们的脚，我们在地上行走的脚。\par

5. 是以教会，基督藉圣言的圣洗所洁净的，没有瑕疵和皱纹（\BibleRef{弗 5:26}），不仅包括那些在重生之圣洗后立刻从今生的传染性影响中被带走、未曾踏足大地而需要洗脚的人，也包括那些从主领受了如此恩慈，以至能带着已洗净的脚离开今生的人。然而，教会对于那些仍居留地上的人也是洁净的，因为他们生活正义；但他们仍需要洗脚，因为他们确实并非无罪。为此，在\BibleBook{}{雅歌}中写道：\BibleQuote{我洗了脚，怎能再玷污呢？}（\BibleRef{歌 5:3}）因为当一个人被迫来到基督面前时，他说这样的话，并且在来的时候必须使他的脚接触地面。但是，又出现了另一个问题：基督不就是在天上吗？祂不是升了天，坐在圣父的右边吗？宗徒不是明确宣告：\BibleQuote{你们既然与基督一同复活，就该思念上面的事，那里有基督坐在天主的右边。你们要思念上面的事，不要思念地上的事。}（\BibleRef{哥 3:1}）那么，我们如何必须踏足大地才能接近基督呢？我们的心岂不是更应转向天上的主，以致能住在祂面前吗？弟兄们，你们看到今天时间短暂，限制了我们对此问题的考虑。如果你们或许在某方面未能做到，但我自己却看出它需要多少澄清。因此，我恳求你们宁可让它延期，也不要现在用过于疏忽和狭窄的方式来处理；你们的期望不会受骗，只会被推迟。因为主这样使我们成为你们的债务人，祂必会临在，使我们也能偿还我们所欠的。\par

\SourceRecord{Translated by John Gibb.；Nicene and Post-Nicene Fathers, First Series；Vol. 7.；Edited by Philip Schaff.；Buffalo, NY: Christian Literature Publishing Co.,；1888.；<http://www.newadvent.org/fathers/1701056.htm>.}

% structure: p0001 source=h1 target=section reason=page-title

\section{论道第57篇（若 13:6-10，歌 5:2-3）}

教会应当如何在走向基督的途中，害怕玷污自己的脚？\par

1. 我并未忘记我的债务，并承认偿还的时刻已到。愿祂赐予我偿还的资财，如同祂使我产生债务的原因一样。因为祂赐予了我爱，关于这爱经上说：\BibleQuote{除了彼此相爱外，你们不可亏欠人什么。} \BibleRef{罗 13:8} 愿祂也赐予我言语，是我感到亏欠于我所爱之人的。我之所以将你们的期待推迟至今，是为了尽我所能解释，我们是如何沿地走向\Person{基督}的，而当我们被命令追求上面的事，而非地上的事时。\BibleRef{哥 3:1} 因为\Person{基督}坐在上面，在\Person{圣父}的右边；但祂确实也在这里；为此祂也对正在地上咆哮的\Person{扫禄}说：\BibleQuote{你为什么迫害我？} \BibleRef{宗 9:4} 但我们当时所谈论的，并引导我们进入这一探讨的主题，是主为祂的门徒洗脚，而门徒们自己已经洗过，只需要洗脚。我们当时看到可以理解为：一个人确实在圣洗中完全被洗净；但此后他生活在这个现世中，用他人性情欲的脚踏在这个地上，即在与他人交往的生活中，他积累了足以唤起祈祷的罪债：\BibleQuote{宽免我们的罪债。} \BibleRef{玛 6:12} 这样，他也从这些罪债中被那位为门徒洗脚、并不停止为我们代求的主所洁净。\BibleRef{罗 8:34} 此时，教会\WorkTitle{雅歌}中的话出现了，她说：\BibleQuote{我洗了我的脚，怎能再污秽呢？} 当她想去开门给那位面貌比世人更美的存在，祂来到她那里敲门，请她开门。这引发了一个问题，我们不愿在有限的时间内压缩它，因此推迟至今，即：教会，在走向\Person{基督}的途中，如何会害怕玷污她已在\Person{基督}的圣洗中洗过的脚？\par

2. 因为教会如此说：\BibleQuote{我睡眠，但我的心醒着；这是我的良人敲门的声音。} 然后祂也说：\BibleQuote{给我开门，我的妹妹，我的近人，我的鸽子，我的完人；因为我的头沾满了露水，我的头发沾满了夜间的滴露。} 她回答说：\BibleQuote{我脱了外衣，怎能再穿上呢？我洗了我的脚，怎能再污秽呢？} \BibleRef{歌 5:2} 啊！奇妙的圣事象征！啊！崇高的奥迹！那么，她害怕在走向那位为门徒洗脚的主时玷污自己的脚吗？她的害怕是真实的；因为她必须沿地走向祂，而祂仍在地上，因为祂不肯离开驻守在此的祂的人。难道不是祂说：\BibleQuote{看！我同你们天天在一起，直到今世的终结。} \BibleRef{玛 28:20} 难道不是祂说：\BibleQuote{你们要看见天开了，天主的使者在人子身上上去下来。} 如果他们上去因为祂在天上，他们又怎能下来，除非祂也在这里？因此教会说：\BibleQuote{我洗了我的脚，怎能再污秽呢？} 她甚至在那些从一切渣滓中洁净的人身上也说：\BibleQuote{我渴望解脱，与\Person{基督}同在一起；但为你们，我留在肉身中更为必要。} \BibleRef{斐 1:23} 她在这班宣讲\Person{基督}、给祂开门、使祂因信德住在人心中的人身上说这话。\BibleRef{弗 3:17} 在这样的时候，当他们斟酌是否承担这样的职务时，他们说这样的话，因为他们认为自己不配毫无指责地履行职务，并免得在给别人宣讲之后，自己反被淘汰。\BibleRef{格前 9:27} 因为听真理比宣讲真理更安全：因为在听的时候，谦卑得以保持；但在宣讲时，几乎无人能阻止某种微小的自夸进入，由此至少玷污了双脚。\par

3. 因此，正如宗徒\Person{雅各伯}所说：\BibleQuote{每个人要快快地听，慢慢地说。} \BibleRef{雅 1:19} 又如另一位天主的人所说：\Quote{你使我听到欢乐和喜乐的声音，你所谦抑的骨头会欢跃。} 这就是我所说的：当听到真理时，谦卑得以保持。另一位又说：\Quote{但新郎的朋友站着听他，因新郎的声音而非常喜乐。} 让我们喜乐地聆听在我们内无声说话的真理。因为虽然当声音在外表发出时，如有人朗读、宣告、宣讲、辩论、命令、安慰、劝勉，甚至有人歌唱或以乐器伴奏时，做这些事的人可能害怕玷污自己的脚，因为他们以隐秘活跃的渴望追求人的掌声。然而，那以乐意和虔诚的心灵听这样的人，在别人的劳苦中没有机会自夸；他不自满，却以谦卑的喜乐因\Person{基督}真理的话语而喜乐。因此，在那些乐意谦卑地听、并在甜美而有益的研读中度过宁静生活的人中，圣教会将感到喜悦，并可以说：\BibleQuote{我睡眠，但我的心醒着。} 而这\BibleQuote{我睡眠，但我的心醒着}是什么呢？岂不是我安静地坐下听？我的闲暇并非用来培养懒惰，而是为了获取智慧。\BibleQuote{我睡眠，但我的心醒着。} 我安静，并看见祢就是主：因为\BibleQuote{文士的智慧来自闲暇的机会；事务少的人将变得智慧。} \BibleRef{德 38:24} \BibleQuote{我睡眠，但我的心醒着：} 我从烦扰的事务中安息，而我的心思转向神圣关注（\Emph{或}启示）。\par

4. 然而，当教会因那些如此温柔谦卑地坐在她脚前的人而得享安息时，有一位前来敲门，说：\BibleQuote{我在暗中告诉你们的，你们要在明处说出来；你们在耳边听到的，要在房顶上宣讲出来。}\BibleRef{玛 10:27} 正是祂的声音在敲门，说：\BibleQuote{给我开门，我的姊妹，我的邻人，我的鸽子，我的成全者；因为我的头满露水，我的发辫满夜露。} 一如祂曾说：你们在闲暇，门却向我关闭；你们关心少数人的闲暇，却因不义的增多，许多人的爱德变冷淡。\BibleRef{玛 24:12} 祂所说的夜就是不义；但祂的露水和夜露是指那些变冷淡而跌倒的人，他们使基督的头变冷淡，即天主的爱德衰败。因为基督的头是天主。\BibleRef{格前 11:3} 但他们被托在祂的发辫上，即在可见的圣事中容忍他们的存在；而他们的感官从未把握内在的实在。因此，祂敲门，要摇动祂那些不活跃的圣徒的安宁，并呼喊：\BibleQuote{给我开门，} 你因我的血成了\BibleQuote{我的姊妹；} 因我的亲近成了\BibleQuote{我的邻人；} 因我的圣神成了\BibleQuote{我的鸽子；} 因我的话语，你在闲暇中已完全学会，成了\BibleQuote{我的成全者：} 给我开门，去向他人宣讲我。因为，若没有人开门，我如何进入那些向我关闭大门的人那里？没有人宣讲，他们如何听见？\BibleRef{罗 10:14}\par

5. 因此，那些喜爱将闲暇用于善学，并回避劳苦工作的困扰，自觉不适合有信用地承担和履行这类服务的人，他们宁愿，若可能的话，希望圣宗徒和古代真理宣讲者再次兴起，以对抗那不义的增多，它已大大减少了基督徒爱德的热度。但关于那些已经离开肉身、脱下血肉之衣的人（因为他们并非完全分离），教会答复说：\BibleQuote{我脱了衣服；怎能再穿上呢？} 那衣服确实将会再得回；而且在那些暂时脱下它的人身上，教会将再次穿上血肉之衣：只是不在此时，当寒冷者需要被温暖时；而是在那时，当死者复活时。因此，意识到她因宣讲者短缺而面临的困难，并记起她那些在言语上纯全、在品格上圣洁的肢体，如今却与身体分离，教会忧伤地说：\BibleQuote{我脱了衣服；怎能再穿上呢？} 我这肢体，曾通过宣讲拥有如此卓越的能力为基督开门，如今怎能回到他们脱下的身体呢？\par

6. 然后，又转向那些宣讲、聚集并治理祂子民会众、尽其所能为基督开门，但在这类工作的困难中害怕陷入罪的人，她说：\BibleQuote{我洗了脚；怎能再玷污呢？} 因为凡在言语上没有过犯的，就是成全的人。然而，谁是成全的呢？在这不义的增多和爱德的冻结中，谁没有过犯？\BibleQuote{我洗了脚；怎能再玷污呢？} 有时我读到和听到：\BibleQuote{我的弟兄们，不要多人作师傅，因为你们要受更重的判断；因为我们许多人都犯有过失。}\BibleRef{雅 3:1} \BibleQuote{我洗了脚；怎能再玷污呢？} 但看，我起身开门。基督，请洗它们。\BibleQuote{赦免我们的罪债，因为我们也赦免我们的欠债者。}\BibleRef{玛 6:12} 当我们聆听祢时，被压伤的骨头与祢一同在天上喜乐。但当我们宣讲祢时，我们不得不踏地前行以便为祢开门：然后，若我们可责备，我们便困扰；若我们受称赞，我们便自高。请洗我们的脚，它们先前已被洗净，但因我们在地上行走为祢开门而再次玷污。亲爱的，今天到此为止。但无论我们在何事上得罪，言非所当言，或因你们的称赞而不当自高，愿祈求我们的脚得洗，并愿你们的祈祷蒙天主悦纳。\par

\SourceRecord{Translated by John Gibb.；Nicene and Post-Nicene Fathers, First Series；Vol. 7.；Edited by Philip Schaff.；Buffalo, NY: Christian Literature Publishing Co.,；1888.；<http://www.newadvent.org/fathers/1701057.htm>.}

% structure: p0001 source=h1 target=section reason=page-title

\section{论道58（\BibleRef{若 13:10-15}）}

亲爱的诸位，我们已经，按主所乐意赐予我们的能力，向你们讲解了福音中的那些话，就是主在洗门徒的脚时所说的：\BibleQuote{一次洗过的人，只需洗脚，就全身洁净了。}现在让我们接着看下文。祂说：\BibleQuote{你们，}祂说，\BibleQuote{是洁净的，但不都是。}为了免去我们探究的必要，福音作者亲自解释了其含义，加上说：\BibleQuote{祂原知道谁要出卖祂，因此祂说：\InnerQuote{你们不都是洁净的。}}还有什么比这更清楚的呢？因此，让我们转到下文。\par

\BibleQuote{然后，祂洗了他们的脚，取了外衣，重新坐下，对他们说：你们知道我给你们做了什么吗？}现在，有福的伯多禄得到那应许的实现：因为他曾受到推延，当时他在战栗和坚持中说：\BibleQuote{你绝不可洗我的脚，}他得到的回答是：\BibleQuote{我所做的，你现在不明白，但日后必会明白}（\BibleRef{若 13:7-8}）。那么，这就是那\Quote{日后}；现在正是讲述不久前所延迟之事的时候。主记起祂先前应许要使他明白祂这一如此出乎意料、如此奇妙、如此令人震惊、且若非祂自己那更令人生畏的回应就绝不可能被允许的行为——即师傅不仅是他们的，也是天使的师傅，主不仅是他们的主，也是万有的主，竟然洗了自己门徒和仆人的脚：祂曾应许要让他知道如此重要行为的意义，当时祂说：\BibleQuote{你日后必会知道，}现在祂开始向他们显示祂所行之事。\par

\BibleQuote{你们称我}，祂说，\BibleQuote{师傅和主；你们说得对，因为我本来就是。}\BibleQuote{你们说得对，}因为你们说的是真理；我确实是你们所说的那样。有一条诫命给人：\BibleQuote{不要让你自己的口赞美你，而要让邻人的口赞美你。}\BibleRef{箴 27:2}因为自满对于一个必须防范陷入骄傲的人来说是危险的事。但那位超越万有的，无论怎样赞美自己，也不能高抬自己超过其实际尊严：天主也不能被正确地称为傲慢。因为认识祂是为了我们的益处，而不是为了祂的益处；也没有人能认识祂，除非那位自知者使自己被认识。因此，如果祂因不作自赞，好似要避免傲慢，祂就会剥夺我们认识祂的能力。确实，没有人会责备祂称自己为师傅，即使相信祂不过是一个人；因为祂只是声称那些人在各种技艺中自称的东西，而不被指控为傲慢，以至于他们被称为教授。但是，祂也称自己为门徒的主——这些人在世俗意义上也是生来自由的——谁能容忍一个人这样呢？但这是天主在说话。在此，如此伟大的崇高不可能有得意，真理不可能有谎言：利益在于我们成为如此崇高的臣民、真理的仆人。祂称自己为主，这对祂而言并非不完美，而是对我们的益处。某位世俗作者的话受到赞赏，他说：\Quote{一切傲慢都是可憎的，尤其令人不悦的是才能和口才的傲慢。}然而，同一个人在谈论自己的口才时说：\Quote{如果我要下判断，我会称它为完美的；而且，事实上，我不会太害怕傲慢的指控。}那么，如果那位最雄辩的人实际上并不害怕被指控傲慢，真理本身怎能有这样的恐惧呢？让那为主者称自己为主，让那为真理者说真理；这样，我便不会因祂对存在之事保持沉默而未能学到有益之事。最有福的保禄——他当然不是天主的独生子，而是那位的仆人和宗徒；不是真理，而是真理的分享者——自由而一贯地宣告：\BibleQuote{纵然我愿意夸耀，我也不算是愚妄，因为我说的是真理。}\BibleRef{格后 12:6}因为他也是在曾下令说：\BibleQuote{夸耀的，应当因主而夸耀}\BibleRef{格前 1:31}的那位。难道爱智慧的人虽然愿意夸耀，却能不担心被指控为愚妄，而智慧本身在夸耀时，会担心这样的指控吗？那说\Quote{我的灵魂要因主而夸耀}的人并不害怕傲慢；主的能力在自我赞许时，怎能害怕呢？而祂的仆人的灵魂正是以此夸耀。\BibleQuote{你们称我}，祂说，\BibleQuote{师傅和主；你们说得对，因为我本来就是。}因此，你们说得对，我就是这样：因为如果我不是你们所说的，你们这么说就是错的，即使是为了赞美我。那么，真理怎能否认真理的弟子所肯定的呢？学习者所说的，怎能被赐予他们学习的真理所否认呢？源头怎能否认饮水者所宣称的呢？光怎能隐藏观看者所宣告的呢？\par

4. \BibleQuote{如果我这主和师傅洗了你们的脚，你们也该彼此洗脚。因为我给你们立了榜样，叫你们照我给你们所做的去做。} 有福的 \Person{伯多禄}，当时你不许这事发生，你便不知道这事的意义。后来当你的师傅和主以威严使你顺服，并洗了你的脚时，祂应许让你明白这事。弟兄们，我们从至高者那里学会了谦卑；让我们这些卑微的人，彼此做那至高者以其谦卑所做的事。我们在此对谦卑有极大的赞赏：弟兄们彼此轮流做这事，甚至在可见的行动中也是如此，就是当他们彼此款待时；因为这种谦卑的实践通常很普遍，并且在使它能被识别的行动本身中表达出来。因此，宗徒在称赞那值得称许的寡妇时说：\BibleQuote{若她好客，若她洗过圣人们的脚。} \BibleRef{弟前 5:10} 而在圣人中，若没有这种实践，他们心里所行的即使手上未行，只要他们是属于那三位有福之人的赞歌中所称呼的：\Quote{你们圣洁而心里谦卑的人，请赞美主。} 但更好且毫无疑问更符合真理的是，也用手去做；基督徒不应认为做基督所做的事是低下的事。因为当身体俯伏在兄弟的脚前时，这种谦卑的情感或是在心中被唤醒，或是若已存在则得以增强。\par

5. 但除了对这经文道德层面的理解外，我们记得，我们曾提请你们注意主这一行为的伟大之处，即：祂洗了那些已洗过且洁净的门徒的脚，是设立了一个标记，为使我们由于占据我们尘世的人性情感，无论我们在理解正义上进步多大，仍知道我们并非免于罪恶；而祂后来透过为我们代祷洗去这罪，当我们向在天之父祈祷：\BibleQuote{宽免我们的罪债，如同我们也宽免得罪我们的人}\BibleRef{玛 6:12}。那么，这样理解经文与祂后来亲自解释自己行为原因的话：\BibleQuote{如果我，你们的主和师傅，洗了你们的脚，你们也该彼此洗脚。因为我给你们立了榜样，叫你们照我给你们所做的去做}之间有何联系呢？难道我们能说一个兄弟能洗净另一个兄弟沾染的罪污吗？是的，诚然，我们知道，在主这行为的深刻意义中，我们也受到劝诫：要彼此告明过错，彼此代祷，就如基督也为我们代祷一样\BibleRef{罗 8:34}。让我们聆听宗徒\Person{雅各伯}，他以最清晰的方式表述这一训诫时说：\BibleQuote{彼此告明你们的过错，彼此祈祷}\BibleRef{雅 5:16}。因为主也给了我们这方面的榜样。既然那位没有、也没有、将来也没有任何罪的主为我们祈求，我们岂不更该彼此代祷呢？既然祂宽恕我们这些对她毫无可亏欠的人，我们这些不能在此世无罪生活的人，岂不更该彼此宽恕呢？因为主在祂圣事标记的深刻意义中，当祂说：\BibleQuote{因为我给你们立了榜样，叫你们照我给你们所做的去做}时，岂不正是明确地暗含了宗徒直截了当所说的话：\BibleQuote{彼此宽恕，若有人对任何人有什么怨恨，如同基督宽恕了你们，你们也要照样做}\BibleRef{哥 3:13}吗？因此，让我们彼此宽恕过错，为彼此的过错祈祷，这样在某种意义上就是在彼此洗脚。借着祂的恩宠，我们的本分是提供爱德和谦卑的服务；祂的本分是垂听我们，并借着基督、在基督内洗净我们一切罪恶的玷污；这样我们在地上所释放的，即我们所宽恕的，在天上也要被释放。\par

\SourceRecord{Translated by John Gibb.；Nicene and Post-Nicene Fathers, First Series；Vol. 7.；Edited by Philip Schaff.；Buffalo, NY: Christian Literature Publishing Co.,；1888.；<http://www.newadvent.org/fathers/1701058.htm>.}

% structure: p0001 source=h1 target=section reason=page-title

\section{论道集59（若 13:16-20）}

1. 我们刚刚在神圣的福音中听到主说话，祂说：\Quote{我实实在在告诉你们：仆人不能大过主人，宗徒（被派遣者）也不能大过派遣他的人。你们如果知道这些事，实行的人便是有福的。}祂说这话，是因为祂洗了门徒的脚，兼以言教和身教成为谦逊的导师。然而，如果我们不停留在已经清楚明白的事情上，我们就能在祂的帮助下，处理那些需要更详尽论述的事。因此，主说了这些话后，又补充说：\Quote{我不是对你们全体说的：我知道我所拣选的是谁；但为应验经上的话：\BibleQuote{同我吃饭的人，要举起脚跟踢我。}}这是什么意思呢？不就是他要践踏我吗？我们知道祂指的是谁：祂指的是出卖祂的犹达斯。所以，祂并没有拣选这个人，因为祂用这些话将这个人从祂所拣选的人中完全分别出来。祂接着说：\Quote{当我说\InnerQuote{你们如果知道这些事，实行的人便是有福的}时，我不是对你们全体说的：}你们中间有一个人将不是有福的，也不会实行这些事。\Quote{我知道我所拣选的是谁。}是哪些人呢？不就是那些实行了祂所命令并指示必须做的事的人吗？祂是唯一能使人有福的那位。祂说，叛徒犹达斯不属于那些被拣选的人。那么，祂在另一个地方所说的话是什么意思：\Quote{我不是拣选了你们十二个人吗？但你们中间有一个是魔鬼。}是否意味着他被拣选是为了某个实际必需的目的，尽管不是为了祂刚才所说的福分：\Quote{你们如果知道这些事，实行的人便是有福的。}？祂并非对全体如此说；因为祂知道祂所拣选，要与祂在福分中结合的是谁。那同我吃饭，为要抬起脚跟踢我的人，不属于这一类。他们吃的面包就是主自己；他怀着敌意吃主的饼：他们吃的是生命，而他吃的是惩罚。\Quote{因为那不相配地吃的人，}宗徒说，\Quote{就是为自己吃审判。}\BibleRef{格前 11:29}\Quote{从今以后，}基督接着说，\Quote{我在事成之前告诉你们，使你们当事成时，相信我就是那一位。}也就是说，我就是先前经上所说：\BibleQuote{同我吃饭的人，要举起脚跟踢我。}的那一位。\par

2. 然后祂接着说：\Quote{我实实在在告诉你们：凡接待我所派遣的，就是接待我；凡接待我的，就是接待那派遣我来的。}祂的意思是要我们了解，被祂派遣的人与祂自己之间的距离，与祂自己与天主圣父之间的距离一样小吗？如果我们这样理解，我不知道我们会采用什么样的距离度量（愿天主禁止！），这简直是亚略派的作风。因为他们听见或读到福音的这些话时，立刻求助于他们教义的度量，借此他们不是升到生命，而是直堕死亡。因为他们马上说：圣子的使者与圣子之间的相对距离，正如圣子与圣父之间的相对距离，正如经上所说：\Quote{凡接待我所派遣的，就是接待我，}以及当祂说\Quote{凡接待我的，就是接待那派遣我来的。}但如果你这样说，异端者，你就忘记了你的度量。因为如果因着主的这些话，你把圣子与圣父之间的距离，等同于使者（宗徒）与圣子之间的距离，那么你打算把圣神放在哪里？你难道忘了，你习惯于把圣神置于圣子之后？那么圣神就要介于使者和圣子之间；如此，圣子与祂的使者之间的距离，将比圣父与圣子之间的距离更大。或者，为了保持圣子与使者之间、圣父与圣子之间的这种距离相等，圣神是否要与圣子相等？但你也不会容许这一点。那么，如果你把圣子置于圣父之下，正如你把使者置于圣子之下，你打算把圣神放在哪里呢？因此，抑制你鲁莽的僭越；不要在这些话中寻求圣子与祂的使者之间，与圣父与圣子之间相同的距离。反而要听圣子自己所说的话：\Quote{我与父原是一体。}因为在那里，真理没有给你留下任何距离的阴影在生者与独生子之间；在那里，基督亲自抹去了你的度量，磐石把你的阶梯砸得粉碎。\par

然而，既然异端的诽谤已被排除，我们应如何理解主的这句话：\Quote{接待我所派遣的，就是接待我；接待我的，就是接待那派遣我的一位}？因为如果我们倾向将这句话，\Quote{接待我的，就是接待那派遣我的一位，}理解为表达圣父与圣子在本性上的合一；那么从另一句相似措辞的语句，\Quote{接待我所派遣的，就是接待我，}推导出的就是圣子与祂的使者在本性上的合一。这样的理解或许并无不妥，因为那强壮的人——他欣然跑完赛程——兼有双重本体；因为圣言成了血肉，也就是说，天主成了人。因此，祂可能曾说过，\Quote{接待我所派遣的，就是接待我，}是就祂的人性而言；\Quote{而接待我}作为天主，\Quote{就是接待那派遣我的一位。}但祂这样说，并非在推崇本性的合一，而是推崇派遣者在被派遣者身上的权威。因此，愿每个人都这样接待那被派遣的，以致在他身上，他能注意那派遣他的。因此，如果你在\Person{伯多禄}身上寻找基督，你将找到门徒的导师；如果你在圣子身上寻找圣父，你将找到独生子的生者；同样，在被派遣者身上，你在接待派遣者时不会犯错。经文中接下来的内容无法在剩余的时间内压缩讲完。因此，亲爱的诸位，如果认为所说的已经足够，愿你们以健康的方式领受，如同神圣羊群的牧草；如果略显不足，愿你们以对更多的炽热渴望再三反刍。\par

\SourceRecord{Translated by John Gibb.；Nicene and Post-Nicene Fathers, First Series；Vol. 7.；Edited by Philip Schaff.；Buffalo, NY: Christian Literature Publishing Co.,；1888.；<http://www.newadvent.org/fathers/1701059.htm>.}

% structure: p0001 source=h1 target=section reason=page-title

\section{讲道集 60（\BibleBook{若望}{福音} \BibleRef{若 13:21}）}

1. 弟兄们，真福\BibleBook{若望}{福音}中遇到了一个并非轻省的问题：圣经记载：\BibleQuote{耶稣说了这些话，心神烦乱，就明明地说：\InnerQuote{我实实在在告诉你们：你们中有一个要出卖我。}}祂心神烦乱，不是肉体，而是精神，难道是因为祂将要说\BibleQuote{你们中有一个要出卖我}吗？这是否那时才初次出现在祂心中，还是那时突然初次启示给祂，因而如此惊人的灾祸的新奇使祂烦乱？不久之前祂不是说过\BibleQuote{同我一起吃饭的，要举脚踢我}吗？而且在这之前，祂不是也说过\BibleQuote{你们是洁净的，但不都是}？福音作者补充说：\BibleQuote{因为祂知道谁要出卖祂}；祂在更早的时候还指着那人说：\BibleQuote{我不是拣选了你们十二个，而你们中有一个是魔鬼吗？}那么，为什么现在祂\Quote{心神烦乱}，当祂\BibleQuote{明明地说：\InnerQuote{我实实在在告诉你们：你们中有一个要出卖我}}？难道是因为现在祂必须如此指出他，使他不再隐藏在其余人中，而从他们中分离出来，因此祂\Quote{心神烦乱}？或者是因为现在叛徒本人即将离去，带那些犹太人，他要将主出卖给他们，祂因祂的苦难临近、危险迫近、以及那已知决意的叛徒阴谋之手而烦乱？毫无疑问，正是因为类似的原因，耶稣\Quote{心神烦乱}，正如祂曾说：\BibleQuote{现在我的灵魂烦乱；我说什么？父啊，救我脱离这个时辰；但我正是为此来到了这个时辰。}因此，正如那时祂的灵魂在祂苦难的时刻临近时感到烦乱，同样现在，当犹大即将离去并回来，叛徒的凶恶罪行接近完成时，祂\Quote{心神烦乱}。\par

2. 那么，祂烦乱了，祂有权舍弃自己的生命，也有权再取回它。那强大的能力烦乱了，那磐石的坚固被搅动了：或者，不如说是我们的软弱在祂内烦乱了？的确如此：愿仆人们不要认为任何不相称于他们主的事，而要承认他们在元首内的肢体。祂为我们而死，也亲自为我们而烦乱。因此，祂以权能而死，在祂的能力中烦乱；祂将我们卑贱的身体改变为相似祂光荣的身体，也将我们软弱的感受转移到自己身上，并以自己灵魂的感受同情我们。所以，当那位伟大、勇敢、确定、不可战胜者烦乱时，我们不要为祂害怕，仿佛祂会失败：祂不是在消逝，而是在寻找我们。我说，是我们；祂正是这样寻找我们，使我们在祂的烦乱中看见自己，这样当烦乱临到我们时，我们不至于陷入绝望而丧亡。藉着祂的烦乱——祂非自愿不能烦乱——祂安慰那些不情愿地烦乱的人。\par

3. 远离那些哲学家的理由，他们断言智者不为心神扰动所影响。天主使这世界的智慧成了愚拙（\BibleRef{格前 1:20}），主知道人的思念是虚妄的。显然，基督徒的心灵可以被扰动，不是出于痛苦，而是出于怜悯：他可能害怕人失去基督，当人正在失去时他可能忧伤，他可能热切渴望赢得人们归向基督，当这样发生时他可能充满喜乐，他可能害怕自己偏离基督，他可能因自己与基督疏远而忧伤，他可能热切期望与基督一同为王，并因希望将来与基督共融而喜乐。这些确实是他们所谓的四种\OriginalTerm{扰动}{perturbations}——恐惧与忧伤、爱与喜乐。基督徒的心灵有充分的理由感受到这些，并表明他们与斯多葛哲学家及其他类似者的错误不同：那些人，正如他们视真理为虚妄，也视麻木为健康；不知道人的心灵，如同身体四肢一样，只有当它连痛苦的感觉都失去时，才更加无可救药地病入膏肓。\par

4. 但有人说：基督徒的心灵甚至对死亡的前景也应当感到烦乱吗？那么宗徒的话：\BibleQuote{我渴求离世，与基督同在}（\BibleRef{斐 1:23}）又成了什么，如果他渴望的对象在来临时竟使他烦乱？我们对此的回答是容易的，尤其是在那些也将喜乐本身称为\OriginalTerm{扰动}{perturbation}的人看来。因为如果他感受到的烦乱完全来自对死亡的喜悦，那又怎样？但那种感受，他们说，应称为快乐，而非喜乐。这不过是换一个名称，而所体验的感受相同。但让我们专注于圣经，并藉着主的帮助，寻求一种符合圣经的解答；然后，既然经上记载\BibleQuote{祂如此说了，就心神烦乱}，我们不会说那是喜乐扰乱了祂；免得祂自己的话使我们相信相反的情况，当祂说：\BibleQuote{我的灵魂忧闷得要死}（\BibleRef{玛 26:38}）。这里也应理解为某种类似的情感，当叛徒正独自离去，并立刻与同伙一同返回时，\BibleQuote{耶稣心神烦乱}。\par

5. 确实，那些基督徒——如果真有这样的人——在面对死亡时全然不受困扰，可谓心志刚毅；但即便如此，他们比基督更刚毅吗？谁敢发狂这样说呢？那么，祂的\Quote{心神烦乱}又意味着什么呢？无非是祂甘愿取了我们软弱的肖像，从而安慰了祂身体——即祂的教会——中的软弱肢体；为的是，如果祂的肢体中有人仍为死亡临近而困扰，他们可以定睛仰望祂，免得因自以为被遗弃而陷入更可怕的死亡——绝望。那么，我们应当在分享祂的神性中期待和盼望的那个善是何等伟大，因为连祂的扰动都使我们得安宁，祂的软弱也使我们得坚固！因此，无论是在这一时刻祂因怜悯犹大自身奔向毁灭而困扰，还是因自己死亡的临近而困扰，都绝无可疑之处：祂的困扰并非出于心志的软弱，而是出于能力的圆满；这样，当我们受困扰时（不是出于能力的拥有，而是处于我们的软弱中），我们心中也不致生出救恩的绝望。祂确实背负了肉身的软弱——这软弱在祂的复活中被吞没了。但祂不仅是人，也是天主，在心志的刚毅上远超整个人类，其距离不可言喻。因此，祂不是被任何外在的人性弱点所困扰，而是祂自己困扰自己；这一点在祂复活拉匝禄时已明确宣示：因为那里记载祂\Quote{心神烦乱}，为使我们在即使经文未如此表达却仍宣称祂困扰之处，也能这样理解。因为祂凭其权力取了我们的完整人性，也凭这权力，每当祂认为适宜时，便在自己内唤醒我们的人性情感。\par

\SourceRecord{Translated by John Gibb.；Nicene and Post-Nicene Fathers, First Series；Vol. 7.；Edited by Philip Schaff.；Buffalo, NY: Christian Literature Publishing Co.,；1888.；<http://www.newadvent.org/fathers/1701060.htm>.}

% structure: p0001 source=h1 target=section reason=page-title

\section{讲道集 61（\BibleBook{若望}{福音} \BibleRef{若 13:21-26}）}

1. 弟兄们，福音这一小节，我们在本篇讲道中提出来加以阐释，认为也应当就主的出卖者说些什么，因为蘸饼递给他这一举动已足够清楚地揭露了他。至于前面的事，即耶稣正要指出他时，心神烦乱，我们已在上一篇讲道中讨论过了；但我在那里或许忽略未提的一点是：主藉着祂自己心神烦乱，也适宜地表明，必须容忍假弟兄以及主田里麦子中的莠子直到收获时节，因为当我们因紧急原因被迫在收获前就分离其中一些时，不能不对教会造成扰动。主仿佛在自己身上预示并预表了祂的圣徒将来因分裂者和异端者而起的扰动，就在那恶人犹大离开、并以此极其公开和决绝的方式结束他长久以来与麦子的混杂之时，祂烦乱了，不是身体，而是心神。因为使祂属灵的肢体在这种丑闻中烦乱的，并非恶意，而是爱德；免得在分离一些莠子时，连一些麦子也一同被拔出来。\par

2. 因此，\BibleQuote{耶稣心神烦乱，就作证说：我实实在在告诉你们，你们中间有一个人要出卖我。} \Quote{你们中间有一个人}，论数目，不是论功德；论外貌，不是论实际；论身体的混合，不是论任何属灵的纽带；因肉体的接近而为同伴，不是因任何心的合一；因此，他不是你们中的一员，而是将要离开你们的人。否则，主如此作证并说的\Quote{你们中间有一个人}如何能成立呢？如果这卷福音的作者在他的书信中所说的\BibleQuote{他们从我们中间出去，却并非属于我们；因为如果他们属于我们，就必会留在我们中间}是真实的？\BibleRef{若一 2:19} 因此，犹大不属于他们；因为如果他属于他们，就会留在他们中间。那么，\Quote{你们中间有一个人要出卖我}这句话是什么意思？无非是说：有一个将从你们中间出去的人要出卖我。正如那说过\BibleQuote{如果他们属于我们，就必会留在我们中间}的人，先前也说\BibleQuote{他们从我们中间出去}。因此，两者都真实：\Quote{属于我们}和\Quote{不属于我们}；在一方面\Quote{属于我们}，在另一方面\Quote{不属于我们}；就圣事共融而言\Quote{属于我们}，但就专属于他们自己的罪行而言\Quote{不属于我们}。\par

3. \BibleQuote{门徒们便彼此对看，猜不透祂所说的是谁。}因为他们虽然对主怀着恭敬的爱，但彼此之间的感情也难免受到人性的软弱影响。各人自己的良心为自己所知，但因不知邻人的心，各人自信的同时也对所有他人不确定，所有他人也对那一个人不确定。\par

4. \BibleQuote{在耶稣怀里侧身靠着祂的，有一个门徒，是耶稣所爱的。}他说\Quote{在祂怀里}是什么意思，稍后他告诉我们，即\Quote{在耶稣的胸膛上}。那正是若望本人，他的福音就在我们面前，他后来明确宣告了。因为那些为我们提供圣经的作者有一个习惯：当他们中任何人叙述神圣历史，说到与己相关的事时，他就像谈论他人一样，并在叙述中给自己安排一个位置，以符合公共事件的记录者身份，而非把自己当作宣讲的主题。圣玛窦也是这样行事的：当他在叙述中说到自己时，他说：\BibleQuote{祂看见一个名叫玛窦的税吏坐在税关上，就对他说：跟随我。}\BibleRef{玛 9:9} 他没有说：祂看见\Emph{我}，并对\Emph{我}说。有福的梅瑟也是如此：他将所有关于自己的历史写得像是关于他人，说：\BibleQuote{主对梅瑟说。}\BibleRef{出 6:1} 宗徒保禄较少这样做，但并非在任何解释公共事件进程的历史中，而是在他自己的书信中。总之，他这样谈论自己：\BibleQuote{我认识一个在基督里的人，十四年前——或在身内、或在身外，我不知道，天主知道——这样的人被提到第三层天上。}\BibleRef{格后 12:2} 因此，当有福的福音作者在这里也说，不是\Quote{我靠着耶稣的怀里}，而是\BibleQuote{有一个门徒靠着}时，让我们认出作者的习惯，而不是对此感到奇怪。因为当事实本身告诉我们时，有什么损于真理呢？而在一定程度上避免了用语的自夸。因为他至少是这样叙说了那最关乎他自己赞美的事。\par

5. 但\BibleQuote{耶稣所爱的}这话是什么意思？难道祂不爱其他门徒吗？同一位若望在前文已经说过：\BibleQuote{祂爱他们到底}\BibleRef{若 13:1}；并且主自己说：\BibleQuote{人若为自己的朋友舍命，没有比这更大的爱。}谁能列举出神圣书页中所有的证据，表明主耶稣不仅爱这一位或当时围绕祂的那几位，也爱将来远处成为祂肢体的那些人，以及祂的普世教会？但这话无疑有些真理，与叙说者所靠的胸膛有关。因为\Emph{胸膛}除了表示某种隐藏的真理，还能表示什么呢？但还有另一段更合适的经文，主或许能使我们对此奥秘说些足够的话。\par

6. \Quote{西满伯多禄遂向他点头，对他说。} 这种表达值得注意，它表明有些话不是通过言语的声音，而仅是通过点头来传达的。\Quote{他点头，说；} 也就是说，他的点头就是他的言语。因为如果有人是在思想中说话，正如圣经所说：\BibleQuote{他们心里说} \BibleRef{智 2:1}，那么他通过点头（以某种外在的记号表达内心先前所构思的）岂不更可以如此？那么，他的点头意味着什么？除了下文所说的还能是什么？\Quote{祂所说的是谁？} 这就是伯多禄点头的语言；因为他不是通过声音，而是通过身体姿势来说话。\Quote{于是他侧身靠近耶稣的胸膛，}——这的确是祂胸怀的怀抱，是智慧的隐秘之所！——\BibleQuote{对祂说：主，是谁？耶稣回答说：我蘸一块饼递给谁，谁就是。祂蘸了饼，递给西满的儿子犹达斯依斯加略。犹达斯接过饼后，撒殚进入了他的心。} 叛徒被揭露了，黑暗的隐匿被显明了。他所领受的是好的，但他却因它而受损，因为他自己是邪恶的，以邪恶的情绪领受了善的。但关于那递给了心怀二意的门徒的蘸饼，以及随之而来的事，我们还有很多要说的；而对此我们需要比现在这次讲道结束时所剩的更多时间。\par

\SourceRecord{Translated by John Gibb.；Nicene and Post-Nicene Fathers, First Series；Vol. 7.；Edited by Philip Schaff.；Buffalo, NY: Christian Literature Publishing Co.,；1888.；<http://www.newadvent.org/fathers/1701061.htm>.}

% structure: p0001 source=h1 target=section reason=page-title

\section{讲道集六二（若 13:26-31）}

1. 亲爱的诸位，我知道有人可能会因这句话而受到触动——无论是以虔诚的心探究其意义，还是以不敬的心加以指责——即：主在将蘸过的饼递给叛徒之后，撒殚就进入了他。因为经上这样记载：\BibleQuote{耶稣蘸了饼，递给西满的儿子依斯加略人犹达斯。他接了那饼以后，撒殚立即进入了他。} 他们说：难道基督的饼，从基督的筵席上赐下的，就有这样的效力，竟使撒殚在之后进入祂的门徒吗？我们给他们的回答是：这反而教导我们，何等需要提防以有罪的精神领受善事。因为至为重要的，不是所领受的东西，而是领受者；也不是所赐之物的性质，而是领受之人的状况。因为按其领受者的品性，善事也可能有害，恶事也可能有益。宗徒说：\BibleQuote{罪过为显示它是罪过，竟藉着善事在我身上产生了死亡。} \BibleRef{罗 7:13} 你看，恶因善而生，是因为善被错误地领受。也正是同一人说道：\BibleQuote{为了使我不因那高超的启示而过于高举自己，就有一根刺加在我身上，就是撒殚的使者来攻击我。关于这事，我三次求主使它离开我；祂对我说：我的恩宠足够你用了，因为能力在软弱中才得以成全。} \BibleRef{格后 12:7} 这里你看到，善因恶而生，是因为恶被以善的精神领受。那么，我们为何惊异于基督的饼被递给犹达斯，从而使他被交给魔鬼呢？我们反而看到，保禄受到了魔鬼使者的攻击，藉此器皿他反而得以在基督内成全。这样，善对恶人有害，恶对善人有益。要记住圣经的含义：\BibleQuote{无论何人，若不相称地吃主的饼，或喝主的杯，就是干犯主的体血。} \BibleRef{格前 11:27} 宗徒说这话时，是针对那些漫不经心、不加分辨地将主的体如同寻常食物领受的人。如果连不分辨——即不将主的体与别的食物区分——的人尚且受到责罚，那么那以朋友面目出现、却成为筵席之敌的人，该受何等刑罚呢！如果食客的疏忽尚且招致谴责，那出卖邀请他的人，又将受到何等惩罚呢？那饼为何赐予叛徒，岂非以此表明他忘恩负义地对待了所领受的恩宠？\par

2. 因此，在这饼之后，撒殚进入了主的叛徒，为要彻底占据这已交付于他权下的人，而在之前，撒殚只进入他引导他犯错。因为我们不应以为，当他去见犹太人、商讨出卖主的价银时，撒殚不在他内；因为福音作者路加很清楚地证实了这一点，他说：\BibleQuote{那时，撒殚进入了那名叫依斯加略的犹达斯，他是十二人之一；他就去同司祭长商议。} \BibleRef{路 22:3} 你看，这里显示撒殚已经进入了犹达斯。因此，他的初次进入，是将出卖基督的念头植入他心中；因着这样的心思，他已经来赴宴。但现在，在饼之后，撒殚进入他，不再只是试探一个属于别人的人，而是将他作为自己的所有物来占有。\par

3. 但并非如某些粗心的读者所想的那样，犹达斯那时领受了基督的身体。因为我们要明白，主早已将祂身体和血的圣事分赐给了所有人，犹达斯也在场，正如圣路加非常清楚地记载的那样；\BibleRef{路 22:19} 此后，才到了按若望记述的时刻，主藉着蘸了饼并递给犹达斯，完全揭露了叛徒，或许也以蘸饼暗示了犹达斯的虚伪。因为蘸东西并不总是表示洗涤；有些东西被蘸是为了染色。但如果在这里蘸饼有一个良好的含义，那么犹达斯对此善的忘恩负义，就理所应当地导致了惩罚。\par

4. 但犹大既非由主，乃由魔鬼所占据，那饼已入其腹，而敌人已入那忘恩者的灵魂；然而，我说，这巨大的邪恶，早已在心中怀胎，只待结出恶果，那该受诅咒的欲望总是先行。因此，当主，这生活的饼，将这饼赐给这死者，并在赐予时揭示了那出卖饼者，祂说：\BibleQuote{你所要做的，快做吧。}祂并非命令那罪行，而是预言了犹大的恶，以及我们的善。因为对犹大而言，有什么比出卖基督更糟？对我们而言，有什么比这更好？那事由他行于自身之毁灭，但撇开他，却为我们而行。\BibleQuote{你所要做的，快做吧。}哦，这话出于一位渴望做好准备而非动怒者！那话！表达的不是对叛徒的惩罚，而是对救赎者的赏报！因为祂说：\BibleQuote{你所要做的，快做吧，}并非愤怒地注视那背信者的毁灭，而是急于完成对信众的救恩；因为祂为了我们的过犯被交付，\BibleRef{罗 4:25}，并且祂爱了教会，为教会舍弃了自己。\BibleRef{弗 5:25}。正如宗徒论及自己时也说：\BibleQuote{祂爱了我，并为我舍弃了自己。}\BibleRef{迦 2:20}。那么，若非基督亲自舍弃自己，无人能交出祂。在犹大的行为中除了罪还有什么？因为他在交出基督时，并未想到我们得救的事——基督正是为此被交付——而只想到自己的钱财，并找到了灵魂的丧失。他得到了他所愿的报酬，但也得到了他不愿的应得报酬。犹大交出了基督，基督交出了自己：前者办理的是出卖自己老师的事，后者办理的是我们救赎的事。\BibleQuote{你所要做的，快做吧，}并非因你自己有权能，而是因祂愿意，祂拥有一切权能。\par

5. \BibleQuote{同席的人没有一个知道耶稣为什么向他说这话。有人因犹大带着钱囊，以为耶稣是对他说：你去买我们过节所应需用的东西，或是叫他拿些东西给穷人。}因此，主也有一个钱箱，祂在那里存放信众的奉献，并分给祂自己和别人的需要。那时，教会金库的惯例首次建立，好使我们明白祂关于不要为明天忧虑的训诫，\BibleRef{玛 6:34}，并非命令圣人们不得存钱，而是不应为任何这类目的事奉天主，并且不应因害怕匮乏而中断行义。因为宗徒也为将来作了如此远见，当他说道：\BibleQuote{若有个信友有寡妇，当供给她们，免得教会受累，好使教会能供给那些真为寡妇的人。}\BibleRef{弟前 5:16}\par

6. \BibleQuote{犹大受了那饼，立刻就出去：那时候是夜间。}而出去的那位，他本身便是黑夜。\BibleQuote{于是当}那黑夜\BibleQuote{出去后，耶稣说：现在人子受到了光荣。}因此，白日向白日发言，即基督向祂忠信的门徒发言，好让他们听见并爱慕祂，追随祂；黑夜向黑夜传知，即犹大向不信的犹太人如此行，好让他们前来作迫害者，并捉拿祂。但现在，在思考主在被恶人逮捕前对虔诚之人所说的这些话时，听者需要特别注意；因此，对于宣讲者而言，现在匆忙思考它们，不如推迟到将来某个时机会更合适。\par

\SourceRecord{Translated by John Gibb.；Nicene and Post-Nicene Fathers, First Series；Vol. 7.；Edited by Philip Schaff.；Buffalo, NY: Christian Literature Publishing Co.,；1888.；<http://www.newadvent.org/fathers/1701062.htm>.}

% structure: p0001 source=h1 target=section reason=page-title

\section{论道第六十三篇（\BibleRef{若 13:31-32}）}

1. 让我们以心智最好的注意力，并赖主的帮助，寻求天主。圣咏说：\Quote{寻求天主，你的灵魂就必生活。}让我们寻找那需要被发现的事物，并探究那已被发现的事物。我们需要发现的那位被隐藏，以便被寻求；而当找到时，祂是无限的，以便仍成为我们寻求的对象。因此，在别处这样说：\Quote{你们要不断寻求祂的面貌。}因为祂按寻求者所能接受的极限满足他；并使寻见者更具能力，使之能根据其接受能力的增长而重新寻求被充满。所以，那话：\Quote{你们要不断寻求祂的面貌，}并不是像某些人的情况那样说的，那些人\Quote{常常学习，却总达不到认识真理的地步；}\BibleRef{弟后 3:7}而是如德训篇所说：\Quote{人做完之后，才刚开始}\BibleRef{德 18:7}，直到我们达到那生命，在那里我们将如此充满，以至于我们的本性将达其极限，因为我们已抵达圆满，不再追求更多。因为那时，凡能满足我们的一切都将显于我们眼前。但在这里，让我们总是寻求，而寻见的赏报不要终止我们的寻求。因为我们并非说，将永远如此，因为只是在这里如此；而是说，在这里我们必须总是寻求，免得我们任何时候以为，在这里我们可以停止寻求。因为对那些被说成\Quote{常常学习，却总达不到认识真理的地步}的人，他们的确在这里常常学习；但当他们离开此生时，他们就不再学习，而是接受他们错误的赏报。因为\Quote{常常学习，却总达不到认识真理的地步}这句话，意思就像是总在行走，却从未走上道路。反之，让我们总在道路上行走，直到抵达它所引向的终点；让我们不在路上任何地方停留，直到我们到达适当的居所：这样，我们既将坚持寻求，也将在寻见中有所得着，并且这样寻求和寻见，我们将继续迈向仍存留的，直到所有寻求的尽头，在那完美境地不容进一步努力的世界里。亲爱的诸位，让这些开场白，使你们的爱德留意我们主在受难前向门徒所说的这篇讲论：因为它本身深奥；尤其是在讲道者计划付出很多劳苦的地方，听者不应懈怠。\par

2. 那么，在犹达斯出去，要快速实行他所计划的，即出卖主之后，主说了什么？当黑夜已离去时，白昼说了什么？当出卖者离去时，救主说了什么？祂说：\Quote{现在人子受光荣了。}为什么是\Emph{现在}？难道仅仅是因为出卖祂的人已经出去，而那些要捉拿并杀害祂的人就在眼前吗？难道就这样祂\Quote{现在受光荣了}，也就是说，祂更深的羞辱临近了，捆锁、审判、定罪、戏弄、十字架和死亡都悬在祂头上？这是光荣还是羞辱？甚至当祂行奇迹时，这位若望不是论祂说：\Quote{圣神还没有赐下，因为耶稣还没有受光荣}吗？因此，那时当祂复活死人时，祂还没有受光荣；而现在当祂亲身临近死亡时，祂受光荣了吗？祂作为天主行事时还没有受光荣，而作为人去受难时却受光荣了？如果这就是伟大的天主、伟大的导师藉这些话所指明和教导的，那会是很奇怪的。我们必须上升到更高处，以揭示至高者的话，祂稍稍启示自己，好让我们能找到祂，又隐藏自己，好让我们寻求祂，从而如同一步一步地从已有的发现推进到尚待我们发现的。我在这里看到某件事，预示着伟大的现实。犹达斯出去，耶稣受光荣了；丧亡之子出去，人子受光荣了。出去的就是那一位，因为他的缘故曾对他们众人说：\Quote{你们是洁净的，但并不都是}（第10节）。因此，当不洁者离去时，所有留下的都是洁净的，并与他们的洁净者同在。类似的情形将发生在世界被基督征服并过去之后，那时在祂的子民中不再留下任何不洁的人；当莠子从麦子中被分离后，义人将在他们父的国里发光如太阳\BibleRef{玛 13:43}。主预见到这样一个未来，并作见证说，现在在莠子的分离中（正如藉犹达斯的离去和留在圣宗徒中的麦子）所预表的正是如此，于是祂说：\Quote{现在人子受光荣了}，好像祂说：看，在我未来光荣的那一天，也将是这样，那时没有一个恶人在场，也没有一个善人缺席。然而，祂的话并不是这样表述：现在人子的受光荣是\Emph{预像}；而是明确地说：\Quote{现在人子受光荣了}；正如并没有说，磐石预表基督；而是说：\Quote{那磐石\Emph{就是}基督}\BibleRef{格前 10:4}。也没有说，好种子预表天国之子，莠子预表恶者之子；而是说：\Quote{好种子就是天国之子；莠子就是恶者之子}\BibleRef{玛 13:38}。因此，按照圣经用语的惯例，即以标记表明所标记的事物，主使用\Quote{现在人子受光荣了}这句话，指明在那大罪人与他们完全分离、圣人们留在祂身边这一已完成的事件中，我们看到了祂受光荣的预像，那时恶人将被最终分离，祂将永远与祂的圣人们同住。\par

3. 但在说了\BibleQuote{现在人子光荣了，}之后，祂又加上\BibleQuote{天主也在祂内光荣了。}因为人子的光荣本身就在于天主在祂内光荣。如果祂不是在自己内光荣，而是天主在祂内光荣，那么祂就是天主在自己内所光荣的那位。正如同要给他们这个解释，祂进一步加上：\BibleQuote{如果天主在祂内光荣了，天主也要在祂内光荣祂。}就是说，\BibleQuote{如果天主在祂内光荣了，}因为祂来不是为承行自己的意愿，而是承行派遣祂者的意愿；\BibleQuote{天主也要在祂内光荣祂，}这样，那\OriginalTerm{永恒圣言}{eternal Word}所取的人性——在其中祂是人子——也要被赋予永恒的不朽。祂说：\BibleQuote{并且，祂要立刻光荣祂；}以此断言预言祂自己的复活即在眼前，并非像我们的复活在世界的终结。因为正是这种光荣，圣史先前曾说过，如我刚才提到的，为此圣神尚未以那种新方式赐给他们——那是在复活后要赐给相信者的——因为耶稣尚未光荣：即必死性尚未披上不朽，暂时的软弱转化为永恒的力量。这种光荣也可能在\BibleQuote{现在人子光荣了；}这句话中表明；如此\BibleQuote{现在}一词可以理解为不是指祂即将来临的苦难，而是指祂紧随而来的复活，仿佛那已经近在眼前的事已经成就了。愿此足以满足你们今天的爱心；当主允许时，我们将继续接下来的话。\par

\SourceRecord{Translated by John Gibb.；Nicene and Post-Nicene Fathers, First Series；Vol. 7.；Edited by Philip Schaff.；Buffalo, NY: Christian Literature Publishing Co.,；1888.；<http://www.newadvent.org/fathers/1701063.htm>.}

% structure: p0001 source=h1 target=section reason=page-title

\section{第64讲（\BibleRef{若 13:33}）}

1. 亲爱的诸位，我们应当注意主的话语的连贯顺序。因为祂先前，在犹大离去并甚至与圣人的外在共融分离之后，曾说：\Quote{现在人子受到了光荣，天主也在祂内受了光荣}——无论祂这样说是指向祂未来的王国，那时恶人将从好人中被分离出来，还是指祂的复活当时就要发生，即不像我们的复活，要延迟到世界末日——然后又加上：\Quote{如果天主在祂内受了光荣，天主也要在祂自己内光荣祂，并且立刻光荣祂}，祂以此毫不含糊地证明了祂自己复活的即刻实现；祂接着又说：\Quote{孩子们，我还与你们同在不久。} 因此，为了不让他们以为天主光荣祂的方式是祂将永远不再与他们有地上的交往，祂说：\Quote{我还与你们同在不久}；好像祂说：我固然要在复活中立刻受到光荣；但我不会立刻升天，而是\Quote{我还与你们同在不久。} 因为，正如我们在\WorkTitle{宗徒大事录}中读到，祂复活后与他们一起度过了四十天，进进出出，也吃喝：\BibleRef{宗 1:3} 当然祂不是有饥渴之感，而是藉着这些证据证实祂肉身的真实性，这肉身虽不再需要，却仍有吃喝的能力。那么，祂说\Quote{我还与你们同在不久}时，是指这四十天呢，还是指别的什么？因为也可以这样理解：\Quote{我还与你们同在不久；} 我仍和你们一样，处于肉身的软弱状态，也就是说，直到祂死而复活为止。因为祂复活后，如前所述，以祂肉身临在的充分显现与他们同在四十天；但祂不再与他们共有人性的软弱。\par

2. 还有另一种祂神圣临在的方式，是凡人感官所不能及的，祂同样说：\BibleQuote{看，我常与你们同在，直到世界的终结。}\BibleRef{玛 28:20} 这至少与\Quote{我还与你们同在不久}不同；因为直到世界终结并不是不久。或者即使这也如此（因为时间飞逝，且千年在天主眼前如同一日，或如同夜间的一更），但我们不能相信祂在此场合有任何这样的意思，尤其当祂接着说：\Quote{你们将要找我，正如我对犹太人说过：\InnerQuote{我去的地方，你们不能去。}} 这就是说，在我与你们同在的这段不久之后，\Quote{你们将要找我，我去的地方，你们不能去。} 难道是在世界终结之后，他们不能去祂要去的地方吗？那么，祂稍后在同一个讲论中所说的：\Quote{父啊，我愿他们也在那里与我同在}，那地方又在哪里呢？因此，当祂说\Quote{我还与你们同在不久}时，祂现在所说的不是祂与祂的子民保持直到世界终结的那种临在，而是指祂在受难之前与他们同在的那种有死软弱的状况，或是祂在升天之前与他们保持的那种肉身的临在。无论谁更赞同哪一种，都不会与信仰相悖。\par

3. 然而，为了不让任何人认为，我们所说主在说\Quote{我还与你们同在不久}时，可能是指祂在受难前与门徒共有的肉身交往这一点与真确的解释不符，也请考虑祂复活后在另一位福音书作者中的话，祂说：\Quote{这是我从前与你们同在时，对你们所说的话。}\BibleRef{路 24:44} 就好像那时祂已不再与他们同在，尽管他们正站在祂身边，看见、触摸并与祂交谈。那么，祂说\Quote{我从前与你们同在时}是什么意思呢？除非是指：我从前与你们同在时，是处于那个你们仍处其中的有死肉身的状态。因为那时，祂确实已在同一肉身中复活；但祂不再与他们共度同样的有死生命。因此，正如在那场合，当祂如今穿着不朽的肉身，真实地说：\Quote{我从前与你们同在时}，我们只能理解为：我从前与你们同在时，是在有死的肉身中；同样，这里我们也毫无荒谬地可以理解祂的话：\Quote{我还与你们同在不久}，好像祂说：我与你们一样，还有不久属于有死之人。那么，让我们接着看后面的话。\par

4. \Quote{你们要寻找我；并且正如我对犹太人说过：我去的地方，你们不能去；现在我也对你们这样说。} 这就是说，你们\Emph{现在}不能去。但当祂对犹太人这样说时，祂没有加上\Quote{现在}。因此，前者在当时不能去祂要去的地方，但后来却能了；因为不久之后，祂以最清楚的话对宗徒伯多禄这样说。因为当伯多禄问：\BibleQuote{主，祢往哪里去？}时，祂回答他说：\BibleQuote{我去的地方，你现在不能跟随我，但后来却要跟随我} \BibleRef{若 13:36}。但这句话的意思不可轻易放过。因为门徒们当时不能跟随主，后来却能跟随，这是往哪里去呢？如果我们说是往死亡里去，那么何时能发现任何人子会不可能死呢？因为在这可朽坏的身体内，人的命运就是这样，在其中生命并不比死亡更容易。所以他们当时不是不能跟随主去死亡，而是不能跟随祂去到那不死的生活。因为主去那里，是为了从死者中复活后，不再死亡，死亡也不再统治祂。\BibleRef{罗 6:9} 因为主将要为义而死，他们那时尚未成熟到能承受殉道的考验，又怎能跟随祂呢？或者，主将要进入肉身的不朽，他们即使愿意死，也要等到世界末日才能复活，此时又怎能跟随祂呢？或者，主将要回到圣父的怀抱，而又不抛弃他们，正如祂来临时从未离开过那怀抱一样，他们此时又怎能跟随祂呢？因为除了在爱德中得以成全的人，没有人能进入那种福乐之境。因此，为了向他们表明如何能获得能力走向祂先去的地方，祂说：\BibleQuote{我给你们一条新诫命：你们要彼此相爱} \BibleRef{若 13:34}。这些就是跟随基督的阶梯；但对此更充分的论述，须留待另一机会。\par

\SourceRecord{Translated by John Gibb.；Nicene and Post-Nicene Fathers, First Series；Vol. 7.；Edited by Philip Schaff.；Buffalo, NY: Christian Literature Publishing Co.,；1888.；<http://www.newadvent.org/fathers/1701064.htm>.}

% structure: p0001 source=h1 target=section reason=page-title

\section{讲道集 65（\BibleBook{若望}{福音}13:34-35）}

1. 主耶稣宣告祂正在将一条新诫命赐给祂的门徒，即他们应当彼此相爱。祂说：\Quote{一条新诫命}，祂说，\Quote{我赐给你们，叫你们彼此相爱。} 但这岂非已在古代天主的法律中受命，其中记载着：\Quote{你应爱人如己}？\BibleRef{肋 19:18} 那么，为何主称它为新的，既然它被证明是如此古老？难道它是新诫命，是因为祂使我们脱去旧人，穿上新人？因为确实不是任何种类的爱都能更新那听它、或更好说是服从它的人，而是那爱，主为将其与一切肉性的情感区分开，附加了：\Quote{如同我爱了你们一样。} 因为夫妻彼此相爱，父母与儿女相爱，以及所有其他将人联系在一起的亲情关系；更不必说那该受谴责和该死的爱，即奸夫淫妇、嫖客娼妓以及所有其他不是因血缘关系而是因人类生活邪恶堕落而彼此纠缠的人所互感的爱。因此，基督赐给了我们一条新诫命，要我们彼此相爱，正如祂也爱了我们。这爱更新我们，使我们成为新人、新约的继承人、新歌的歌唱者。正是这爱，亲爱的弟兄们，也曾更新了古时那些义人，就是族长和先知，正如后来它更新了蒙福的宗徒一样；它也正在更新万民，并从遍布普世的全人类中，聚集一个新民族，即天主独生子新妇的身体，关于她在\WorkTitle{雅歌}中记载着：\Quote{那上升的、被漂白的是谁？} 确实被漂白，因为被更新；如何更新的？岂非是通过新诫命？因此，其肢体彼此互相关心；若一个肢体受苦，所有肢体一同受苦；若一个肢体得荣耀，所有肢体一同欢乐。\BibleRef{格前 12:25} 为此他们听见并遵守：\Quote{我赐给你们一条新诫命，叫你们彼此相爱：} 不是像败坏者那样彼此相爱，也不是像人以属人的方式相爱；而是如同神一样彼此相爱，他们都是至高者的儿子，因此也是祂独生子的兄弟，以那相互的爱，就是祂爱他们的爱，当时祂要引领他们达到那一切丰盈都属于他们的目标，在那里他们的一切渴望都以美善之物得到满足。因为那时，他们所欲求的将一无所缺，当天主成为万物之中的万有时。\BibleRef{格前 15:28} 这样的终结没有终结。没有人死在那样一个地方，因为到达那里的，唯独是向这个世界死去的人，不是那种灵魂离开身体的普遍死亡，而是蒙选者的死亡，通过它，即使仍留在这可朽的肉身内，心却专注于上面的事。关于这样的死亡，宗徒说：\BibleQuote{因为你们已经死了，你们的生命与基督一同藏在天主内。} \BibleRef{哥 3:3} 或许这话也指这点：\BibleQuote{爱情猛烈如死亡。} \BibleRef{歌 8:6} 因为藉着此爱，当我们仍被囚于这腐朽的身体时，我们向世界死去，我们的生命与基督一同藏在天主内；是的，这爱本身就是我们向世界的死，以及与天主的生命。因为如果灵魂离开身体是死亡，那么当我们的爱离开世界时，又怎能不是死亡呢？因此，这样的爱情猛烈如死亡。还有什么比那束缚世界的更强呢？\par

2. 那么，我的弟兄们，不要以为当主说\Quote{我赐给你们一条新诫命，叫你们彼此相爱}时，祂是忽略了那更大的诫命，即要求我们全心、全灵、全意爱上主我们的天主；因为与这种表面上的忽略同时，\Quote{叫你们彼此相爱}这句话似乎也完全没有提及那第二条诫命，即\Quote{你应爱人如己}。因为祂说：\Quote{这两条诫命}，\Quote{是全部法律和先知的基础。} \BibleRef{玛 22:37} 但那些有良好理解的人，可以在其中每一条中找到这两条诫命。因为，一方面，爱天主的人不能轻视祂关于爱近人的诫命；另一方面，以神圣和属灵的方式爱近人的人，他在近人身上所爱的，不是天主吗？这就是那与一切世俗之爱区别开的爱，主特别将其特征化，当祂附加了：\Quote{如同我爱了你们一样。} 因为祂在我们身上所爱的，除了天主，还有什么呢？不是因为我们已经拥有了祂，而是为了使我们能拥有祂；并且为了引领我们，如我刚才所说，到达那天主成为万物之中的万有的地方。同样，医生被恰当地说成爱病人；他在病人身上所爱的不就是健康吗？这健康无论如何是他渴望恢复的，而不是他所来要消除的疾病。那么，让我们也如此彼此相爱，以便尽可能通过我们爱德的关怀，赢得彼此让天主在我们内。这爱是由那说\Quote{如同我爱了你们，你们也要彼此相爱}的那位赐给我们的。因此，正是为了这个目的，祂爱了我们，叫我们也彼此相爱；藉着祂对我们的爱，将这一恩赐赐给我们，使我们以互爱相连结，并因着如此愉快的纽带而结合为肢体，成为如此伟大之首的身体。\par

3. \Quote{凭此，}祂接着说，\Quote{众人因此就可认出你们是我的门徒，如果你们彼此相爱：}好像祂说，我其他的恩赐，那些不属于我的人也同你们一起拥有——不仅有本性、生命、知觉、理性，以及那同样赐予人和走兽的安全；而且还有语言、圣事、预言、知识、信德、将他们的财产施舍给穷人，并将自己的身体焚毁。但因为没有爱德，他们只是像响的钹；他们算不得什么，也一无所益。\BibleRef{格前 13:1}所以，不是凭着我这些恩赐，无论多么好，那些不是我的门徒的人也能同样拥有，而是\Quote{凭此众人就可认出你们是我的门徒，如果你们彼此相爱。}噢，基督的新妇，女性中的佳丽！噢，你洁白地上升，依靠着你的良人！因为藉着祂的光，你被照耀得洁白，藉着祂的帮助，你被保守不致跌倒。你的话语在\WorkTitle{雅歌}中是多么相称，那仿佛是您的婚歌：\Quote{你的爱情在欢愉中！}这就是那不让你的灵魂与不虔敬的人一同灭亡的；这就是那判断你的案件、坚强如死亡、并在你的欢愉中临在的。那死亡的特性是多么奇妙，它几乎被惩罚的苦难吞没，若不是它反而在欢愉中被淹没！但这里，这篇讲论现在必须结束；因为我们必须开始处理接下来的话语。\par

\SourceRecord{Translated by John Gibb.；Nicene and Post-Nicene Fathers, First Series；Vol. 7.；Edited by Philip Schaff.；Buffalo, NY: Christian Literature Publishing Co.,；1888.；<http://www.newadvent.org/fathers/1701065.htm>.}

% structure: p0001 source=h1 target=section reason=page-title

\section{\WorkTitle{讲道集第66篇（\BibleRef{若 13:36-38}）}}

1. 主耶稣正吩咐门徒当以圣爱彼此相爱时，\newline{}\Quote{西满伯多禄对祂说：主，祢往哪里去？}\newline{} 门徒对他的师傅，仆人对他的主，正是这样说的，好像他已预备好要跟随。正因如此，主洞察他心中此问的意图，便这样回答他：\newline{}\Quote{我所去的地方，你现在不能跟我去；}\newline{} 如同在说：关于你所问的事，你现在不能。祂没有说：你不能；而是说：\newline{}\Quote{你现在不能。}\newline{} 祂暗示了延迟，却未剥夺希望；那希望，祂非但没有夺去，反而赐予，随后的言语又加以确认，接着说：\newline{}\Quote{但后来你要跟我去。}\newline{} 伯多禄，为何如此急迫？\newline{}\OriginalTerm{磐石}{petra} 尚未藉祂的圣神坚固你。不要因自负而高傲，\newline{}\Quote{你现在不能；}\newline{} 也不要因绝望而跌倒，\newline{}\Quote{但后来你要跟随。}\newline{} 对此他怎么说呢？\newline{}\Quote{为什么现在我不能跟祢去？我要为祢舍掉我的性命！}\newline{} 他看见了心中愿望的样式，却看不见自己力量的尺度。软弱的人夸口他的意愿，但医师却留意他健康的状况；一人许诺，另一人预知；无知者大胆，那预知一切者屈尊教诲。伯多禄承担了多少，只因他只看自己所愿的，却不知自己所能的！他承担了多少，竟在主将要为朋友舍命、也为他舍命时，敢于保证为主做同样的事；而基督的性命尚未为他舍去，他却许诺要为基督舍掉自己的性命！\newline{}\Quote{耶稣} 于是\newline{}\Quote{回答他说：你要为我舍掉你的性命吗？}\newline{} 你将为我去做我尚未为你做的事吗？\newline{}\Quote{你要为我舍掉你的性命吗？}\newline{} 你不能跟随，却要先行吗？你为何如此推测？你怎样看自己？你把自己想像成什么？请听你是什么：\newline{}\Quote{我实实在在告诉你：鸡未叫以前，你要三次不认我。}\newline{} 看，你如今如此高谈阔论，却不知自己只是个孩童，你很快就将如此显明于自己。你应许我你的死亡，却要以你的生命否认我。你现在以为自己能为我死，当先学习为自己活；因为你害怕肉身的死亡，将导致灵魂的死亡。正如承认基督是生命，否认祂便是死亡。\par

2. 或者，是否宗徒\Person{伯多禄}实际上并未否认基督，如某些人以一种乖僻的好意竭力为他开脱的那样？因为当被使女盘问时，他回答说他不认识那个人，正如其他圣史更明确地证实的。好像否认作为人的基督就是否认基督一样；并且如此否认祂，是针对祂为了我们而成为的那个身份，以免祂所赐予我们的本性失落。因此，凡是承认基督为天主、却否认祂为人的人，基督没有为他而死；因为基督正是作为人而死的。谁否认基督为人，他就不能通过中保与天主和好。因为天主只有一个，在天主与人之间的中保也只有一位，就是那人基督耶稣。\BibleRef{弟前 2:5} 凡否认基督为人的，不能成义；因为正如因一人的悖逆，众人都成了罪人，同样，因一人的服从，众人也都要成为义人。\BibleRef{罗 5:19} 凡否认基督为人的，不能复活进入生命；因为死是藉着人来的，死人的复活也是藉着人来的；因为正如在\Person{亚当}内众人都死了，照样，在基督内众人都要复活。\BibleRef{格前 15:21} 祂凭什么成为教会之首呢？岂不是凭祂的人性吗？因为圣言成了血肉，就是说，天父的独生子天主成了人。那么，谁否认作为人的基督，怎能处于基督的身体内呢？谁否认头，怎能成为肢体呢？但当主亲自推翻了人类推论的种种曲折时，又何必再举诸多理由呢？因为祂并没有说：鸡叫以前，你将否认那个人；也没有像祂在与人更亲密的俯就中常说的：鸡叫以前，你将三次否认人子；但祂说：\Quote{直到你三次否认我。} 这个\Quote{我}是什么呢？不就是祂所是的吗？而祂不是基督吗？因此，祂的任何方面，他否认了，就是否认祂自己，否认了基督，否认了主、他的天主。因为祂的同伴\Person{多默}也曾高呼：\Quote{我主，我的天主！} 他并没有触摸圣言，只触摸了祂的肉身；他没有将好奇的手放在天主无形体的本性上，而是放在祂的人身上。于是他触摸了人，却承认了天主。那么，若后者所触摸的，伯多禄否认了；后者所呼求的，伯多禄冒犯了。\Quote{鸡叫以前，你将三次否认我。} 虽然你说：\Quote{我不认识那个人；} 虽然你说：\Quote{人啊，我不知道你说的是什么；} 虽然你说：\Quote{我不是祂的门徒之一；} 你都是在否认我。如果基督如此说了，并预言了真理（怀疑这一点是罪恶的），那么毫无疑问伯多禄否认了基督。我们不可在为伯多禄辩护时指控基督。让软弱承认自己的罪吧；因为真理中没有谎言。当伯多禄的软弱承认其罪时，他的承认是充分的；他借眼泪表明了他否认基督所犯恶行的重大。他自己谴责他的辩护者们，并为了说服他们，以眼泪作为证人。就我们而言，这样说并非乐于指控首席宗徒；而是当注视他时，我们应当为自己汲取教训，即无人应信赖人的力量。因为我们的导师和救主还有什么其它目的呢？祂不正是以首席宗徒本人为例，向我们表明没有人应当丝毫依靠自己吗？因此，在伯多禄灵魂里确实发生了那件事，他在身体上为它提供了缘由。然而，他并没有如他鲁莽地自以为是的那样，在主面前先行一步，而是做了与他所预料的相反的事。因为在主死亡和复活之前，他当否认时就死了，当哭泣时又活了；他死，是因为他因自己的自恃而骄傲；他活，是因为那另一位以仁慈注视了他。\par

\SourceRecord{Translated by John Gibb.；Nicene and Post-Nicene Fathers, First Series；Vol. 7.；Edited by Philip Schaff.；Buffalo, NY: Christian Literature Publishing Co.,；1888.；<http://www.newadvent.org/fathers/1701066.htm>.}

% structure: p0001 source=h1 target=section reason=page-title

\section{讲道第六十七篇（若望福音第十四章一至三节）}

1. 弟兄们，我们必须将特别的注意力热切地转向天主，以便能够对刚才传入我们耳中的圣福音的话语获得一些明智的理解。因为主耶稣说：\Quote{不要让你们的心烦乱。你们要相信天主，也要相信［\Emph{或}，也相信］我。} 为使他们在人性上不致害怕死亡而烦乱，祂安慰他们，肯定自己也是天主。祂说：\Quote{相信，}祂说，\Quote{在天主内，也要相信我。} 因为由此推得，如果你们相信天主，也应该相信我：如果基督不是天主，就不能这样推论。\Quote{相信天主，也相信}那一位，祂凭本性而非强夺，与天主同等；因为祂空虚了自己，然而并非失去天主的形体，而是取了奴仆的形体。\BibleRef{斐 2:6} 你们对这奴仆形体感到死亡恐惧，\Quote{不要让你们的心烦乱，}天主的形体将复活它。\par

2. 但为什么后面又说：\Quote{在我父的家里有许多住处，} 是因为他们也为自己感到恐惧吗？他们本可听到这句话：\Quote{不要让你们的心烦乱。} 因为，在他们中间，有谁能够免于恐惧呢？当最自信而大胆的伯多禄被告知：\Quote{鸡叫以前，你要三次不认我} 的时候？因此，他们考虑到自己，从伯多禄开始，认为命定要丧亡，他们有理由烦乱；但是当他们现在听到：\Quote{在我父的家里有许多住处；不然，我早就告诉你们了；因为我去是为你们预备地方，}他们从烦乱中复苏，变得确定而自信，在经历了诱惑的一切危险之后，他们将与基督一同住在天主的面前。因为，虽然一人比另一人强壮，一人比另一人有智慧，一人比另一人更义，但在\Quote{父的家里有许多住处}；他们中没有一个会留在那家之外，在那里每个人都按照自己的功德领受一个住处。所有人都得到了那个德纳，即家主吩咐给所有在葡萄园做工的人，不分做多做少；\BibleRef{玛 20:9} 这个德纳当然象征着永生，在那里没有人再活不同的长度，因为永恒中生命的度量没有差异。但许多住处指向那唯一永生中的不同功德等级。因为有一种太阳的光荣，另一种月亮的光荣，另一种星辰的光荣；这星和那星的荣光也有区别；同样，死者的复活也是如此。圣人如同天上的星辰，在王国中获得不同明亮程度的住处；但为了那个德纳，没有人被排除在王国之外；而天主将成为万物中的万有，如此，既然天主是爱，\BibleRef{若一 4:8} 爱将使每人所拥有的成为众人共有的。因为这样，当人爱看见别人身上有自己所没有的东西时，他真正拥有它。因此，在这不同的明亮中不会有嫉妒，因为在众人中掌权的是爱的合一。\par

3. 因此，每一个基督徒的心都必须彻底抛弃那些人的想法，他们想象之所以说到许多住处，是因为在天国之外将有一个地方，作为那些未领受圣洗而离世的蒙福无辜者的居所，因为没有圣洗他们不能进入天国。这样的信仰不是信仰，因为它不是真实而大公的信仰。你们岂不是如此愚昧，被肉体的想象弄瞎了心眼，以至于应受责斥吗？如果你们竟然将住处，我不说伯多禄和保禄或任何宗徒的，甚至任何受洗婴孩的住处，与天国分离；你们岂不认为你们这样做是将这些住处与天主圣父的家分离而应受责斥吗？因为主的原话不是说：\InnerQuote{在全世界}或\InnerQuote{在一切受造物中}或\InnerQuote{在永生和真福中}有许多住处；而是说：\Quote{在我父的家里有许多住处。} 那难道不是我们将要拥有的天主所造的房屋，一座非人手所造、在天上永恒的家吗？\BibleRef{格后 5:1} 那难道不是我们向主歌唱的家吗？\Quote{住在你家里的，真是有福；他们要永远赞美你。} 那么，你们还敢将，不只是每位受洗弟兄的家，而是天主圣父自己的家，即我们所有弟兄向祂说\Quote{我们的天父，你在天上}的家，\BibleRef{玛 6:9} 与天国分离吗？或者以这样的方式分割，使一些住所在天国之内，一些在天国之外？断乎不可！那些渴望居住在天国的人，也会愿意与你们一同居住在这种愚昧中吗？断乎不可！我说，既然统治的儿子的家只能在王国里，王室本身的任何部分就不能在王国之外。\par

4. \Quote{如果我去了，}祂说，\Quote{为你们预备了地方，我必再来，接你们到我那里去，为的是我在哪里，你们也在那里。我去的地方，你们知道；那条路，你们也知道。} 主耶稣啊，如果在你父的家里已经有许多住处，你的子民将与你同住，你怎能去预备地方呢？或者，如果你接他们到你那里去，你怎能再次来临，而你从未离开你的临在呢？亲爱的，这类问题，若我们试图以似乎适合我们今天讲道的适当范围之简短来解释，必定会因篇幅压缩而失去清晰度，而简洁本身也会成为第二种晦涩；因此，我们将延迟这个债务，我们家族之主的慷慨将在更合适的机会使我们能够偿还。\par

\SourceRecord{Translated by John Gibb.；Nicene and Post-Nicene Fathers, First Series；Vol. 7.；Edited by Philip Schaff.；Buffalo, NY: Christian Literature Publishing Co.,；1888.；<http://www.newadvent.org/fathers/1701067.htm>.}

% structure: p0001 source=h1 target=section reason=page-title

\section{讲道集68（\WorkTitle{若望福音14:1-3}）}

1. 亲爱的弟兄们，我们承认亏欠你们，现在应当偿还我们尚待思考的债，即我们如何理解这两句话之间没有真正的矛盾：在说\Quote{在我父的家里有许多住处；不然的话，我就早已告诉你们了，我去为你们预备地方}之后——祂由此清楚表明，祂之所以对他们这样说，正是因为那里已有许多住处，无需再预备任何地方——主又说：\Quote{我若去了，为你们预备了地方，就必再来接你们到我那里去，使我在哪里，你们也在哪里。}既然已有许多住处，祂又怎会去预备地方呢？若没有那些住处，祂就会说\Quote{我去预备}。或者地方仍需预备，祂岂不也应当说\Quote{我去预备}吗？这些住处是已经存在，却仍需要预备吗？因为若不存在，祂就会说\Quote{我去预备}。然而，由于它们现在存在的状态仍然需要预备，祂去预备并不是指它们已经存在；但若祂去预备它们如同将来所是，祂就会再来迎接祂自己的人到自己那里，使祂在哪里，他们也在哪里。那么，在父的家里有许多住处，而且这些不是不同的住处，而是相同的，它们在某种意义上已经存在，可以被认为无需预备，同时又不存在，因为它们仍需预备？我们该如何理解这事呢？惟有像先知那样，他也论及天主说，祂已经造了那将要成为的。因为他不说\Quote{谁将造那将要成为的}，而是说\Quote{谁造了那将要成为的}。因此，祂既已造了它们，也将要造它们。因为若祂不造它们，它们就根本没有被造；若祂自己不造它们，它们也永远不会存在。祂按预定造了它们，祂藉实际工作将要造它们。正如福音在祂拣选门徒时明确暗示的那样，即祂召唤他们的时候（\BibleRef{路 6:13}）；然而宗徒说：\Quote{祂在创世以前就拣选了我们}（\BibleRef{弗 1:4}），即藉预定，而非藉实际召唤。\Quote{祂所预定的，祂也召唤了}（\BibleRef{罗 8:30}）；祂在创世以前藉预定拣选，在世界的结束之前藉召唤拣选。同样，祂也预备了这些住处，并且仍在预备它们；祂已经造了那将要成为的，现在正在预备的，不是不同的住处，而是祂已经预备的那些住处：祂在预定中\Emph{已经}预备的，祂藉实际工作\Emph{正在}预备。因此，从预定的角度看，它们已存在；否则，祂就会说\Quote{我去预备}，即我要预定。但因为它们尚未处于实际预备的状态，祂说：\Quote{我若去了，为你们预备了地方，就必再来接你们到我那里去。}\par

2. 然而，祂在某种意义上是在预备住所，即预备那些住在其中的人。例如，当祂说\Quote{在我父的家里有许多住处}时，我们还能把天主的家理解为什么，不就是天主的殿吗？那殿是什么？请问宗徒，他会回答说：\Quote{天主的殿是圣的，这殿就是你们}（\BibleRef{格前 3:17}）。这也是天主的国，圣子将把它交给父；因此同一宗徒说：\Quote{基督是初果，然后是在基督内属于祂的人；再后，是结局，那时祂将把王国交给天主父}（\BibleRef{格前 15:23}），即祂用祂的血所救赎的人，将交到祂父的面前。这就是那个天国，关于它曾说：\Quote{天国好像一个人把好种子撒在田里；好种子就是天国之子}；虽然现在他们与莠子混杂，但到结局时，君王本人将派遣祂的天使，\Quote{他们将把一切使人跌倒的事从祂国里收集出来；那时义人要在他们父的国里发光如同太阳}。当那国（那些属国的人）到达了国度时，国度将在国度中发出光辉；正如我们现在祈祷时所说：\Quote{愿祢的国来临}（\BibleRef{玛 6:10}）。因此，现在国度已被召叫，但只是正在被聚集。因为若现在不被召叫，就不能说\Quote{他们将把一切使人跌倒的事从祂国里收集出来}。但国度尚未掌权。因此，它已经是国度，以至于当一切绊脚石被收集出来后，它将进入治理，不仅拥有国度的名称，而且拥有治理的权柄。因为到结局时，这个站在右边的国度将听到：\Quote{你们蒙我父降福的，来承受天国吧}（\BibleRef{玛 25:34}），即你们曾是国度，但尚未治理，来治理吧；使你们从前只在希望中所有的，现在能在现实中实现。因此，天主的这房屋，天主的这殿，天主的这国和天国，现在正在建造、施工、预备、聚集的过程中。其中将有住处，正如主现在所预备的；其中已有住处，正如主已经预定的。\par

3. 但是，祂为何离去准备地方呢？既然我们自身是需要准备的客体，祂这样做岂不是因留下我们而受阻？主啊，我尽己所能这样解释：祢准备那些居所，无疑是表明义人应当因信德而生活。\BibleRef{罗 1:17} 因为那旅行在远离主的地方的人，需要因信德而生活，因为藉此我们被准备来瞻仰祂的面容。\BibleRef{格后 5:6} 因为\BibleQuote{心里洁净的人是有福的，因为他们要看见天主}；\BibleRef{玛 5:8} 并且\BibleQuote{祂因信德洁净了他们的心}。\BibleRef{宗 15:9} 前者我们在\WorkTitle{福音}中找到，后者在\WorkTitle{宗徒大事录}中找到。但是那信德，那些还将在远处旅行而尚未看见天主的人，藉此洁净他们的心，相信未见之事；因为若有看见，就不再是信德了。现在功劳为信者积累，然后赏报被付予看者的手中。让主离去，为我们预备地方吧；让祂离去，好使祂不被看见；让祂隐藏，好使信德得以操练。因为那时地方正在预备，若我们因信德而生活。愿对那地方的相信被渴求，好使所渴求的地方本身得以拥有；爱的渴望就是居所的预备。主啊，请这样预备祢所预备的；因为祢预备我们为祢自己，也预备祢自己为我们，正如祢为我们预备一个地方，既在祢内为我们，也在我们内为祢。因为祢曾说：\BibleQuote{你们住在我内，我也住在你们内。} 各人按照对祢的分享程度，有的少些，有的多些，赏报的差别也将按功劳的差别相应；居所的众多也将适应居住者的不平等；但所有居住者都一样，永远活着，无尽享福。祢为何离去？祢为何再来？若我了解祢正确，祢既不离开所去的地方，也不离开所来的地方：祢藉着成为不可见而离去，祢藉着再次显现在我们眼前而再来。但除非祢留下指引我们如何继续在生活的善中前进，地方怎能预备好，使我们能住在圆满的喜乐中？让我们所说的关于从福音读到的话，直到\BibleQuote{我要再来，接你们到我身边}为止，已足够了。但接下来的话，\BibleQuote{我在哪里，你们也在那里；并且我去哪里，你们知道，那条路你们也知道}的意义，我们将在更好的情况下——在那位门徒随后提出的问题之后，我们也藉着他提出这个问题——来聆听，并更适合于作为我们讲论的主题。\par

\SourceRecord{Translated by John Gibb.；Nicene and Post-Nicene Fathers, First Series；Vol. 7.；Edited by Philip Schaff.；Buffalo, NY: Christian Literature Publishing Co.,；1888.；<http://www.newadvent.org/fathers/1701068.htm>.}

% structure: p0001 source=h1 target=section reason=page-title

\section{\WorkTitle{讲道集 69（\BibleRef{若 14:4-6}）}}

1. 亲爱的诸位，现在我们有机会，尽我们所能，从主后来的话理解祂先前的话，并从后续的陈述理解之前的陈述，也就是你们所听到的祂对宗徒多默问题的回答。因为主在前面论及住所时，曾说那些住所已在祂圣父的家里，并说祂要去预备它们；我们在那里理解到，那些住所已存在于\OriginalTerm{预定}{predestination}的计划中，也正通过信德对将来要居住在其中之人内心的净化而预备着，因为他们自己就是天主的住所；而住在天主的家里，除了成为祂子民中的一员，还有什么呢？因为祂的子民同时在天主内，天主也在他们内。为了预备这一切，主离开了，为使我们藉着相信祂，（虽然不再可见，）那外在形式始终隐藏于未来的住所，如今能通过信德预备好。因此，祂曾说：\BibleQuote{我若去了，为你们预备了地方，就必再来接你们到我那里去，为的是我在哪里，你们也在哪里。我去的地方，你们知道，那条路你们也知道。}对此，多默对祂说：\BibleQuote{主！我们不知道祢往哪里去，怎么能知道那条路呢？}主曾说他们两者都知道；而多默却声称他不知道，即主去的地方和所走的路。但他不知道自己在说假话；因此，他们知道，却不知道他们知道。主要使他们确信，他们已经知道他们自认为仍不知道的事。耶稣对他说：\BibleQuote{我就是道路、真理、生命。}弟兄们，这是什么意思？看，我们刚刚听到门徒发问，师傅教导，但即使祂的声音已传入我们耳中，我们仍未领悟祂话语中所隐含的思想。但我们无法领悟什么呢？难道与祂谈话的宗徒们会对祂说：\Quote{我们不认识祢}吗？因此，如果他们认识祂，而祂本身就是道路，那么他们就知道道路；如果他们认识祂（祂本身就是真理），那么他们就知道真理；如果他们认识祂（祂也是生命），那么他们就知道生命。所以，你们看，他们被说服了，知道他们所不知道他们知道的事。\par

2. 那么，在这番话中，我们也没有领悟到什么呢？弟兄们，你们认为还有什么呢？不就是祂说：\BibleQuote{我去的地方，你们知道，那条路你们也知道}吗？在这里我们发现，他们知道那条路，因为他们认识祂（祂就是那条路）：路是我们所走的；但路也是我们去的地方吗？然而祂却说这两者他们都知道，既知道祂去哪里，也知道那条路。因此，祂需要说：\BibleQuote{我就是道路}，为向那些认识祂的人表明，他们知道那条他们自以为不知道的路；但是，祂说：\BibleQuote{我就是道路、真理、生命}还有什么必要呢？因为当知道了祂所走的路之后，他们仍然需要知道祂去哪里，但这正是因为祂要去的是真理和生命。因此，祂是藉着自己到自己那里去。而我们去哪里，除了到祂那里，又走哪条路，除了藉着祂？所以，祂藉着自己到自己那里去，我们藉着祂到祂那里去；是的，同样，祂和我们都是这样到圣父那里去。因为祂在别处也论到自己说：\BibleQuote{我到圣父那里去}；此处为了我们，祂说：\BibleQuote{除非藉着我，没有人能到圣父那里去}。这样，祂藉着自己既到自己那里去，也到圣父那里去；我们藉着祂既到祂那里去，也到圣父那里去。除了具有属灵洞察力的人，谁能领悟这些事呢？即使他如此属灵地分辨，又能领悟多少呢？弟兄们，你们怎能期望我向你们解释这样的事呢？只要想想它们是多么崇高。你们看见我是谁，我看见你们是谁；在我们众人身上，\BibleQuote{会朽坏的肉身使灵魂沉重，这地上的帐棚使多虑的心神负担重重}。\BibleRef{智 9:15}难道我们以为可以说：\Quote{上主，居住在天上的祢，我向祢举起我的灵魂}吗？但我们被如此沉重的负担所压，在这重负之下我们呻吟，我怎能举起我的灵魂，除非那为我放下自己生命的主与我一同举起？那么，我将尽我所能地讲说，你们中谁能领受就领受吧。祂赐予时，我讲说；祂赐予时，领受者领受；祂赐予时，对于尚不能领悟的人，有信德。因为先知说：\Quote{如果你们不信，定然不能理解。}\par

3. 请告诉我，我的主啊，我该对祢的仆人们、我的同仆们说什么。宗徒多默曾面对着祢向祢提问，却仍不能明白祢，除非祢住在他内；我向祢提问，因为我知道祢在我之上；我寻求，尽我所能，让我的灵魂散布在同一个高于我的领域，在那里我可以聆听祢，祢无需用外在的声音来传达祢的教导。请告诉我，我祈求，祢如何走向祢自己。祢从前曾离开自己来到我们这里吗？尤其因为祢不是凭自己来的，而是圣父派遣了祢？诚然，我知道祢空虚了自己；但在取了奴仆的形体时（\BibleRef{斐 2:7}），祢既没有放下天主的性体作为要返回的东西，也没有失去它作为要恢复的东西；然而祢来了，并且不仅将祢自己放在肉体的眼前，甚至放在人的手中。若非藉着祢的血肉，又怎能如此呢？祢藉此来了，却仍停留在原地；祢藉此返回，却没有离开祢所来到的地方。那么，如果祢藉此来去，无疑祢不仅是使我们来到祢那里的道路，也是祢自己来去的道路。因为当祢回到生命（祢就是生命）时，确实祢将那同一个血肉从死亡带到了生命。天主的圣言是一回事，人又是另一回事；但圣言成了血肉，或成了人。因此，圣言的位格与人的位格并无不同，因为基督两者同属一个位格；这样，正如当祂的血肉死亡时，基督死了，当祂的血肉被埋葬时，基督被埋葬了（因为我们心里相信，以致成义；口里宣认，以致救恩——\BibleRef{罗 10:10}），同样，当血肉从死亡来到生命时，基督复活了。因为基督是天主的圣言，祂也就是生命。因此，以一种奇妙而不可言喻的方式，祂从未放下或失去自己，却回到了自己。但如上所说，天主曾藉着血肉来到人那里，真理来到撒谎者那里；因为天主是真实的，人人都是撒谎者（\BibleRef{罗 3:4}）。因此，当祂将祂的血肉从人间撤去，并提升到没有撒谎者的地方时，祂自己——因为圣言成了血肉——也藉着祂自己，即藉着祂的血肉，回到了真理那里，这真理不是别的，正是祂自己。我们毫不怀疑，这真理虽在撒谎者中间，祂甚至在死亡中也保存了它；因为基督曾死过一次，却从未虚假。\par

4. 试举一例，性质很不相同，且完全不合适，但在某种程度上仍有助于理解天主，取材于那些特别亲密地服从天主的事物。请看在我自己的情况中，就我的心灵而言，我与你们一样；如果我保持沉默，我便只对自己存在；但如果我对你们说出一些适合你们理解的话，在某种意义上，我走出到你们那里，却没有离开我自己；同时接近你们，却并未离开我出发的地方。但当我停止说话时，我在某种程度上回到我自己，又在某种程度上仍与你们同在，如果你们保留了我在这篇演说中所听到的。如果连天主所造的肖像都能如此，那么天主本身，天主的肖像本身，不是被造的，而是由天主所生的，又有什么不能呢？祂的身体，祂藉此来到我们中间又离开我们，并没有像我的声音那样消逝，而是停留在那里，在那里不再死亡，死亡也不再统治它（\BibleRef{罗 6:9}）。关于福音的这些话，或许可以也应该说更多；但你们的灵魂不应被神性食物所累，尽管它很甘美，尤其是因为心神固然愿意，肉体却软弱（\BibleRef{玛 26:41}）。\par

\SourceRecord{Translated by John Gibb.；Nicene and Post-Nicene Fathers, First Series；Vol. 7.；Edited by Philip Schaff.；Buffalo, NY: Christian Literature Publishing Co.,；1888.；<http://www.newadvent.org/fathers/1701069.htm>.}

% structure: p0001 source=h1 target=section reason=page-title

\section{论道70（\BibleRef{若 14:7-10}）}

1. 弟兄们，圣福音的话，惟有与前文和谐一致时，才能被正确理解；因为发言者既是真理，前提就当与结论相符。主先前曾说：\BibleQuote{我若去为你们预备了地方，就必再来接你们到我那里去，为的是我在哪里，你们也在那里。}然后又加上说：\BibleQuote{我去哪里，你们知道；那条路，你们也知道。}并表明祂所说的一切，无非是他们认识祂自己。因此，祂独自往祂自己那里去——因为祂也让门徒们看到，他们要靠祂到祂那里去——这一说法的意义，我们在上次讲道中已经尽我们所能告诉了你们。因此，当祂说：\BibleQuote{我在哪里，你们也在哪里}时，他们还能在哪里，除非在祂内？同样，祂也在祂自己内，因此他们正是在祂所在之处，即在祂内。因此，祂自己就是那永生，当祂把我们接到祂自己那里时，这永生必将属于我们；并且，既然祂就是那永生，那么在祂内，我们在哪里也就在那里，就是说，在祂内。因为 \BibleQuote{圣父怎样在自己内有生命}，而祂所拥有的生命与祂作为生命的拥有者自身毫无不同，\BibleQuote{同样，祂也赐给圣子在自己内有生命}，因为圣子本身就是祂在自己内所拥有的生命。然而，当我们开始在那生命，即祂内存在时，我们真的会成为祂所是的（即）那生命吗？当然不是，因为祂凭其作为生命的实际存在而拥有生命，并且祂自身就是祂所拥有的；且作为生命，祂在圣父内，同样，祂也在自己内；但我们不是那生命，而是分享祂的生命，并且将在那里，以至于完全不能在自己内成为祂所是的，而是当我们自己不是生命之时，拥有祂作为我们的生命，而祂自己拥有生命，恰恰因为祂自己就是生命。简言之，祂既不变地存在于自己内，又不可分地存在于圣父内。但我们，当我们想要存在于自己内时，就被抛入有关我们自己的内心困扰中，正如这句话所表达的：\BibleQuote{我的灵魂在我内沮丧}；并且从坏变到更坏，甚至无法保持原来的样子。但是，当我们藉着祂来到圣父面前（按照祂自己的话：\BibleQuote{除非藉着我，没有人能到圣父那里去}），并住在祂内时，就没有人能使我们与圣父分离，也不能与祂分离。\par

2. 因此，祂将前面的话与后面的话连接起来，接着说：\Quote{你们若认识了我，也就必定认识我的父。}这与祂先前的话\Quote{除非经过我，谁也不能到父那里去}相符。然后祂又加上：\Quote{从今以后，你们认识祂，并且已经看见祂。}但宗徒之一的斐理伯不明白他所听到的话，就说：\Quote{主！请将父显示给我们，我们就心满意足了。}主回答他说：\Quote{斐理伯！这么长久的时间，我与你们在一起，而你们还不认识我吗？谁看见了我，就是看见了父。}请看，祂抱怨自己与他们在一起这么久，他们竟不认识祂。然而祂不是亲口说过：\Quote{我去的地方，你们知道那条路}，而当他们说不知道时，祂又加上\Quote{我是道路、真理、生命}的话，使他们确信自己是知道的吗？那么，祂现在怎么说：\Quote{斐理伯！这么长久的时间，我与你们在一起，而你们还不认识我吗？}事实上，他们知道祂去的地方和那条路，无非是因为他们真正认识了祂本人？但这个难题很容易解决：他们当中有些人认识祂，有些人不认识，而斐理伯是那些不认识祂的人之一；因此，当祂说\Quote{我去的地方，你们知道那条路}时，祂被理解为是对那些认识祂的人说的，而不是对斐理伯，因为祂对斐理伯说\Quote{斐理伯！这么长久的时间，我与你们在一起，而你们还不认识我吗？}所以，对那些已经认识圣子的人，现在也提及圣父说：\Quote{从今以后，你们认识祂，并且已经看见祂。}因为这话是由于圣父与圣子之间存在的完全相似；因此，因为他们认识了圣子，就可以说他们从现在起认识圣父。所以，他们已经认识圣子——如果不是所有的人，至少那些被告知\Quote{我去的地方，你们知道那条路}的人认识祂，因为祂自己就是那条路。但他们不认识圣父，因此也得听\Quote{你们若认识了我，也就认识我的父}；也就是说，藉着我，你们也认识了祂。因为我是一个，祂是另一个。但为了不让他们以为祂们不相像，祂又加上：\Quote{从今以后，你们认识祂，并且已经看见祂。}因为他们看见了祂那完全相像的圣子，但需要使他们铭记这个真理：他们所见到的圣子如何，他们所未见的圣父也完全如何。这一点也体现在后来对斐理伯说的话中：\Quote{谁看见了我，就是看见了父。}这并非说祂自己就是父和子——那是撒伯里乌派的观念，他们也称为一位受难论者，被公教信仰所谴责——而是说父和子如此相像，以致认识一位就认识两位。因为我们惯常对那些常常见到两人中之一、并想知道另一人如何的人这样说：看见一个，你就看见了另一个。因此，\Quote{谁看见了我，就是看见了父}是这样说的。当然，不是说那为子的也是父，而是说圣子在各方面都与父的肖像完全一致。因为若非父和子是两位，就不会说：\Quote{你们若认识了我，也就认识我的父。}确实如此，因为祂说：\Quote{除非经过我，谁也不能到父那里去；你们若认识了我，也就认识我的父}；因为我是通往父的唯一道路，我将引你们到祂那里，使祂也为你们所认识。但既然我在各方面是祂的完美肖像，\Quote{从今以后，你们认识祂}，因为认识了我；\Quote{并且已经看见祂}，如果你们用灵魂的属灵眼目看见了我。\par

3. 那么，斐理伯，你为什么说：\Quote{请将父显示给我们，我们就心满意足了？斐理伯！这么长久的时间，我与你们在一起，而你们还不认识我吗？谁看见了我，就是看见了父。}如果你非常渴望看见这个，至少相信你所看不见的。因为祂说：\Quote{你怎么说：请将父显示给我们？}如果你看见了我，我是祂的完美肖像，你就看见了我所肖似的那一位。如果你不能直接看见这个，至少\Quote{你不信我在父内，父在我内吗？}但斐理伯在这里可能会说：\Quote{我确实看见了你，也相信你完全肖似父；但一个人看见一个肖似另一个的人，就希望看见他所肖似的那另一个，这难道该受责备和斥责吗？我确实认识肖像，但至今我只认识一个而不认识另一个；除非我认识那肖像所肖似的另一位，这对我是不够的。所以，请将父显示给我们，我们就心满意足了。}但导师确实责备了这个门徒，因为祂看透了问话者的内心。因为斐理伯是怀着一种想法——仿佛父在某方面比子更优越——而渴望认识父；因此他连子也不认识，因为他相信子比另一位低劣。为了纠正这种观念，才说：\Quote{谁看见了我，就是看见了父。你怎么说：请将父显示给我们？}我明白你话的意思：你不是要看见原本的肖像，而是认为那一位更优越。\Quote{你不信我在父内，父在我内吗？}为什么你渴望在这如此相像的两位之间发现距离？为什么你渴求分别认识那不能被分开的两位？此后祂不仅对斐理伯，也对全体门徒所说的话，现在不可草率放过，以便藉祂的助佑，可以更仔细地予以阐释。\par

\SourceRecord{Translated by John Gibb.；Nicene and Post-Nicene Fathers, First Series；Vol. 7.；Edited by Philip Schaff.；Buffalo, NY: Christian Literature Publishing Co.,；1888.；<http://www.newadvent.org/fathers/1701070.htm>.}

% structure: p0001 source=h1 target=section reason=page-title

\section{讲道集71（\BibleRef{若 14:10-14}）}

1. 请诸位悉心聆听并努力理解，亲爱的诸位；因为虽然是我们说话，但真正教导的却是主自己，祂从不离开我们。主说，你们刚才听到的读经：\BibleQuote{我对你们说的话，不是凭我自己说的；而是住在我内的父，祂做这些工作。}这样说来，祂的话也是工作吗？显然是的。因为一个人用言语建立邻人，就是做善工。但\BibleQuote{我不是凭自己说的}是什么意思呢？乃是：说话的我并非凭自己而存在。因此，祂将所做的一切归于那使祂（行这一切者）存在的一位。因为父不是从任何人（如受生等）而来的天主，而子是与父平等、但从父天主而来的天主。因此，父是天主，却非从天主而来；是光，却非从光而来；而子是从天主而来的天主，从光而来的光。\par

2. 关于这两句话——一句是\BibleQuote{我不是凭自己说的}；另一句是\BibleQuote{但住在我内的父，祂做这些工作}——我们面临两类不同的异端者的反对。他们各自只抓住一句话，朝着相反方向偏离，远离了真理之路。例如，亚略派说：看，子不与父平等，祂不凭自己说话。而撒伯流派（即圣父受苦派）却说：看，父就是子；不然，\BibleQuote{住在我内的父，祂做这些工作}是什么意思呢？难道不是\Quote{我住在我内做这些工作}吗？你们提出相反的主张，不仅因为任何事物都是虚假的（即与真理相反），而且也因为两个虚假的东西彼此矛盾。你们迷失在相反的方向；在两者之间是你们遗弃的道路。你们彼此之间的距离，比你们双方都离弃的正途还要远。你们从这边来，从那边来，不要互相穿越，而是从两边都到我们这里来，使这里成为你们相遇之处。你们撒伯流派，承认你们忽略的位格；亚略派，将你们贬低的那一位置于平等的地位，这样你们就会与我们一同走在真理之路上。因为你们双方都有理由彼此劝诫。撒伯流派，你听：子与父如此不同，以至于亚略派认为祂低于父。亚略派，你听：子与父如此平等，以至于撒伯流派宣称祂就是父。你们要恢复你们所否定的位格，提升你们所贬低的尊严，然后与我们一起站在同一立场：因为你们中的一位（先前是亚略派的），为驳斥撒伯流派，从不错失那与父不同的位格；另一位（先前是撒伯流派的），为驳斥亚略派，谨慎维护与父平等的那一位的尊严。因为祂向你们双方呼喊：\BibleQuote{我与父原是一体。}当祂说\BibleQuote{一体}时，让亚略派倾听；当祂说\BibleQuote{我们}时，让撒伯流派留心，不再继续愚昧地否认祂的平等或祂的位格区分。那么，如果祂说\BibleQuote{我对你们说的话，不是凭我自己说的}被认为表明祂的能力如此低下，以致祂所做的不是祂自己所愿的，就请听祂说的另一句话：\BibleQuote{父怎样复活死人并赐给他们生命，子也照样随意赐给人生命。}同样，如果\BibleQuote{住在我内的父，祂做这些工作}被认为表示祂与父没有位格区别，就请听祂的话：\BibleQuote{凡父所做的，子也照样做。}这样，祂将被理解不是两次指同一位，而是指两位一体。正因为祂们平等而不混淆位格，所以祂不凭自己说话，因为祂不是凭自己存在；而住祂内的父自己做这些工作，因为那与父一同并藉父做这些工作的子，除非从父而来，便不存在。然后祂接着说：\BibleQuote{你们不相信我在父内，父在我内吗？否则，你们也该因那些工作而信我。}此前只有斐理伯受责备，但现在表明当时不止他一人需要责备。祂说：\BibleQuote{因那些工作，你们该信我在父内，父在我内。}因为如果我们彼此分离，我们就无法不可分离地做任何工作。\par

3. 但接着所说的又是什么呢？\BibleQuote{我实实在在告诉你们：凡信我的，我所做的事业，他也要做，并且还要做比这些更大的事业，因为我往父那里去。你们因我的名无论求什么，我必要践行，为叫父在子身上获得光荣。你们若因我的名求任何事物，我必要践行。}祂如此应许，祂自己也要做那些更大的事业。不要使仆人在主以上自高，也不要使门徒在师傅以上自高。祂说他们将要做比祂自己所做的更大的事业；但这一切都是祂在他们内或藉着他们做的，并非像是他们自己做的一样。因此，那向祂歌颂的诗歌说：\BibleQuote{上主，我的力量，我爱慕祢。}那么，那些更大的事业是什么呢？是不是连他们的影子，在他们走过时，也使病人痊愈了？（\BibleRef{宗 5:15}）因为一个影子具有治病的权能，比一件衣裳的繸子更有力量。（\BibleRef{玛 14:36}）前一件事业是由基督自己做的，后一件是由他们做的；然而，仍是祂做了这两件。然而，祂这样说时，是在赞扬祂自己言语的效力；因为正是在这个意义上祂曾说过：\BibleQuote{我对你们所说的话，不是凭我自己说的；而是住在我内的父，祂在做这些事业。}祂所指的事业，不就是祂所说的话语吗？他们正在听并相信，他们的信德正是那些话语的果实；然而，当门徒们宣讲福音时，不仅像他们那样的小数目的人，连万民也相信了；这样的工作，无疑地是更大的事业。然而祂并没有说\Quote{你们将要做比这些更大的事业}，以引导我们以为只有宗徒们会这样做；因为祂接着说：\BibleQuote{凡信我的，我所做的事业，他也要做，并且还要做比这些更大的事业。}那么，事实是，那信基督的人要做基督所做的事业，甚至比祂所做的更大吗？像这样的问题不可草率处理，也不应仓促解答；因此，既然我们现在的讲论必须结束，我们不得不将进一步的思考推迟。\par

\SourceRecord{Translated by John Gibb.；Nicene and Post-Nicene Fathers, First Series；Vol. 7.；Edited by Philip Schaff.；Buffalo, NY: Christian Literature Publishing Co.,；1888.；<http://www.newadvent.org/fathers/1701071.htm>.}

% structure: p0001 source=h1 target=section reason=page-title

\section{第72篇（\BibleRef{若 14:10-14}）}

1. 要理解主的话是什么意思，或该以什么意义来接受，并不容易：\Quote{那信我的人，我所做的作为，他也将要做；}然后，在这个理解上的巨大困难之上，他又加上了另一个更难的：\Quote{并且他要做比这些更大的事。}我们该怎么办？我们还没有找到一个人做了基督所做的作为；难道我们还会找到一个能做更大的吗？但我们在上次讲道中提到，门徒们通过他们的影子经过而治愈病人，比主自己触摸衣边所做的更大；而且相信宗徒的人比相信主自己亲口宣讲的人更多；所以我们可以认为这种作为被理解为更大：不是门徒大于师傅，或仆人大于主人，或义子大于独生子，或人大于天主，而是他通过他们屈尊去做这些更大的作为，同时又在另一处对他们说：\Quote{没有我，你们什么也不能做。}而他自己，暂且不说他无数的其他作为，没有他们的任何帮助就创造了他们，并且没有他们创造了这个世界；因为他自己认为合宜成为人，没有他们也创造了自己。但他们没有他做了什么，除了罪？最后，他立即撤去了所有可能使我们不安的因素；因为说了\Quote{那信我的人，我所做的作为，他也将要做；并且比这些更大的作为，他也将要做}之后，他紧接着加上：\Quote{因为我往父那里去；并且你们因我的名无论求什么，我将要做。}他先说了\Quote{他将要做}，后又说了\Quote{我将要做}；好像说：不要觉得这不可能；因为那信我的人绝不能变得比我大，而是那时我将比现在做更大的事；通过那信我的人而做的事，比脱离他而我自己做的事更大；然而，我自己脱离他时，和我自己通过他时，都是我自己在做工：并且，因为脱离他时，不是他在做；而另一方面，因为通过他，虽然不是通过他自己，也是他在做。此外，通过人比脱离人做更大的事，不是匮乏的标志，而是俯就的标志。因为仆人为一切恩惠能向主回报什么？有时他屈尊将此也算作他对他们的其他恩惠之一，即通过他们做比他脱离他们所做的更大的事。难道那个富人没有因寻求关于永生的忠告而从主面前忧闷地离开吗？他听到了，却拒绝了它；然而在后来的日子里，那落在他耳中的忠告，不是一个人，而是许多人遵行了，当善师通过门徒说话时；他被富人鄙视，当他直接警告时，而被那些人爱，当他通过贫穷的人使那些富人变为贫穷时。所以，你看，他通过信徒宣讲所做的工，比他自己对听众说话所做的更大。\par

2. 但他通过宗徒做这样的更大作为，仍然有值得思考之处；因为他不是说，好像仅仅指着他们：我所做的作为，你们也要做；并且比这些更大的作为，你们也要做：而是愿意被人理解为他所说的包括所有属于他家庭的人，说：\Quote{那信我的人，我所做的作为，他也将要做；并且比这些更大的作为，他也将要做。}那么，若那信的人要做这样的作为，那不做的人当然就不是信的人：正如\Quote{那爱我的，必遵守我的诫命}当然意味着，那遵守的不是爱。在另一处，他也说：\Quote{那听我这些话并实行的，我将把他比作一个聪明人，他把房子盖在磐石上；}（\BibleRef{玛 7:24}）因此，凡不像这聪明人的，无疑要么听了这些话而不实行，要么甚至没有听到。\Quote{那信我的人，即使死了，也要活着；}因此，那不会活着的，现在当然不是信的人。同样，这里也说：\Quote{那信我的人，将要做（这样的作为）：}所以，那不做的人就不是信的人。那么，弟兄们，这是什么意思？难道一个不能做比基督更大的作为的人，就不算基督的信徒吗？这样想是困难的、不合理的、无法容忍的；除非正确理解。那么，让我们听听宗徒的话，当他说：\Quote{对于那信那使不虔敬者成义的天主的人，他的信德就算为正义。}（\BibleRef{罗 4:5}）这就是我们可以做基督作为的作为，因为甚至我们信基督也是基督的作为。这是他成就在我们内的，当然不是没有我们。那么现在，请听和明白：\Quote{那信我的人，我所做的作为，他也将要做：}我先做，然后他也要做；因为我做这样的作为，好让他也做。而这些作为是什么，岂不就是从不虔敬的人中造出义人吗？\par

3. \BibleQuote{他要作比这些更大的事业。}请问，比什么更大？难道我们要说，那个正在战栗畏惧地成就自己救恩的人\BibleRef{斐 2:12}，他所做的比基督所做的一切更大吗？这是一个工作，基督确实在他内工作，但不是没有他；并且我可以毫不犹豫地说，这个工作比天和地以及我们视野内天地间的一切都更大，因为天地都要过去\BibleRef{玛 24:35}，但是那些被预定得救的人的救恩和成义，也就是祂所预知的人的救恩和成义，却要永远存留。在前者中只有天主的工作，但在后者中还有祂的肖像。但天上也有宝座、宰制、率领、掌权、总领天使和天使，这些都是基督的工作；那么，那个与基督在他内工作、在他自己永恒的救恩和成义中成为合作者的人，他所作的竟然比这些还要大吗？我不敢就这一点匆忙下结论：让能明白的人明白，让能判断的人判断：创造正义的存在者与使不虔敬的人成义，哪一个更是更大的工作。因为至少，如果两者运用的是同等的能力，那么在后者中有更大的仁慈。因为\BibleQuote{这就是虔敬的伟大奥迹：祂在肉身内显现，在圣神内成义，被天使看见，被向外邦人宣讲，在世上被相信，被接入荣耀中。}但是当祂说\BibleQuote{他要作比这些更大的事业}时，并没有必要让我们认为基督的一切工作都被包含了。因为祂所说的，或许只是指祂当时正在做的\Emph{这些}工作；而祂当时所做的工作就是说出信仰的话语，并且祂特别提到这些工作，就在之前祂说\BibleQuote{我对你们所说的话，不是凭我自己说的；而是住在我内的父，祂做这些工作。}因此，祂的话就是祂的工作。的确，宣讲正义的话（这是祂不与我们合作而独自做的）比使不虔敬的人成义（这是祂在我们内以这样一种方式做的，以至于我们自己也做这事）要小一些。我们还需要探讨如何理解这句话：\BibleQuote{无论你们因我的名求什么，我必要做。}因为祂的信徒求许多事，却不得应允，这并非一个小问题，需要我们的注意；但由于这篇讲道必须结束，我们至少需要稍微延迟一下来思考和讨论这个问题。\par

\SourceRecord{Translated by John Gibb.；Nicene and Post-Nicene Fathers, First Series；Vol. 7.；Edited by Philip Schaff.；Buffalo, NY: Christian Literature Publishing Co.,；1888.；<http://www.newadvent.org/fathers/1701072.htm>.}

% structure: p0001 source=h1 target=section reason=page-title

\section{第73篇讲道（\BibleRef{若 14:10-14}）}

主藉自己的恩许，使那些把希望寄于祂的人获得一种特殊的确信根基，当祂说：\BibleQuote{因为我往父那里去；凡你们因我名所求的，我必实行。}因此，祂往父那里去，并非有抛弃贫弱者的意图，而是为了垂听并应允他们的祈祷。可是，当我们看见祂的忠实信徒屡次求而不得时，该如何理解\BibleQuote{凡你们所求的}这句话呢？我们可否说，别无他故，只因他们求而不当？因为宗徒雅各伯曾以此作为责备的理由，说：\BibleQuote{你们求而不得，是因为你们求而不当，想要浪费在你们的私欲中。}\BibleRef{雅 4:3}因此，一个人想要得到某物以便用于不当用途，天主反而出于仁慈拒绝赐予。不仅如此，如果一个人所求的若蒙应允只会对自己有害，那恐怕更有理由担忧：唯恐天主会因愤怒而赐予他原本出于仁慈不该拒绝的东西。我们岂不见以色列子民因有罪的贪欲所渴求的，反而给自己带来伤害吗？因为当时天上正给他们降下玛纳，他们却想吃肉。\BibleRef{户 11:32}他们鄙视已有的，无耻地追求所无的：好像他们更好的做法不是求那欠缺的食物来满足不当的欲望，而是求消除自己的厌恶，使自己能妥善接受已提供的食物。因为当邪恶成为我们的喜乐，而善事相反时，我们更应恳求天主使我们重归对善的爱慕，而不是赐给我们邪恶。这并不是说吃肉有错，因为宗徒论及此事时说道：\BibleQuote{凡天主所造的样样都好，若以感恩之心领受，没有一样是可摈弃的。}\BibleRef{弟前 4:4}而是因为，正如他所说的：\BibleQuote{那带着妨碍而吃的人，就是有罪了。}\BibleRef{罗 14:20}既然如此，若对人有妨碍，何况对天主呢？以色列子民拒绝智慧所供给的，而寻求情欲所贪求的，这对天主来说决非小罪。虽然他们实际上并未提出请求，而是因缺乏而抱怨。但是，为让我们知道，错不在天主的任何受造物，而在顽梗的悖逆与无度的欲望：原祖父母并非因猪肉而丧亡，而是因一个苹果；\BibleRef{创 3:6}厄撒乌也并非因一只飞禽，而是因一碗红豆羹而失去了长子名分。\BibleRef{创 25:34}\par

那么，如果有些事信徒祈求，而天主即使有意也替他们留下未做，我们该如何理解\BibleQuote{凡你们所求的，我必实行}呢？还是我们应认为这句话仅是对宗徒们说的？当然不是。因为祂现在所说的，与祂先前所说的完全一致：\BibleQuote{凡信我的，我所做的工程，他也要做，并且还要做比这些更大的工程。}这正是我们上一次讲道的主题。为使没有人将这种能力归给自己，而是显明连这些更大的工程也是由祂完成的，祂接着说：\BibleQuote{因为往父那里去；凡你们因我名所求的，我必实行。}难道只有宗徒们信祂吗？因此，当祂说\BibleQuote{凡信我的}时，祂是对那些人包括我们在内的人说的（我们藉祂的恩宠被包括在内），而我们决非有求必得。如果我们甚至想到最蒙福的宗徒们，我们会发现那位比众人更劳苦的（但不是他，而是与他同在的天主恩宠）\BibleRef{格前 15:10}曾三次求主使那撒殚的使者离开他，却未得所求。\BibleRef{格后 12:8}亲爱的，我们该说什么呢？我们是否认为祂甚至对宗徒们也没有实现这里所作的许诺：\BibleQuote{凡你们因我名所求的，我必实行}？那么，如果祂在这件事上欺骗了自己的宗徒，祂的许诺又将对我们谁实现呢？\par

3. 那么，信徒啊，醒来吧，仔细留意此处所说的：\Quote{\Emph{因我的名}}: 因为在这句话中，祂并没有说\Quote{无论你们求什么}而不加限制，而是说\Quote{因我的名}。如此，那位应许了如此大恩惠的，祂被称作什么呢？当然是基督耶稣：基督意为君王，耶稣意为救主！因为确实，并非任何一位君王都能拯救我们，只有那救主君王才能；因此，凡我们所求的，若有损于救恩的利益，我们便不是因救主的名而求。然而，祂是救主，不仅当祂成就我们所求时，也当祂拒绝成就时；因为祂不成就祂所见有损于我们救恩的事，便更充分地彰显祂是我们的救主。就如医生知道病人的哪些请求有利于其安康，哪些有损，因此当病人求有害之事时，医生不顺从其意愿，以便使他康复。同样，当我们希望祂成就我们一切所求时，我们不要随意去求，而要因祂的名——即因救主的名——呈上我们的请求。那么，我们不可求任何与自身救恩相悖的事；因为若祂成就了那种事，祂便不是以救主的身份成就，而救主正是祂对忠实门徒所怀有的名号。因为那位屈尊成为信徒救主的，也是审判不虔敬者的法官。因此，任何信祂的人，若因那祂对信祂者所怀有的名号而求，祂必成就；因为祂将以救主的身份成就。但若一个信祂的人因无知而求有损于其救恩的事，他便不是因\Emph{救主}之名而求；因为若祂做任何妨碍其救恩的事，祂便不再是他的救主了。因此，在这种情况下，祂不成就所求之事，反而更清晰地保持了祂成就其名号所意味之事的路。为此，祂不仅作为救主，也作为良善的师傅，在祂所教导我们的祈祷中，指示了我们应该求什么，以便无论我们求什么，祂都能成就；同时，我们也因此明白，我们不能因师傅的名求任何与祂自己教训的规则相悖的事。\par

4. 的确，有些事虽真是因祂的名而求——即合乎祂作为救主和师傅的特性——但祂并不在我们求的时候成就，然而祂终必成就。因为我们祈求天主的国来临，这并不意味着祂没有成就我们所求，因为我们并没有立即与祂一同在永远国度中为王：我们所求的是延迟，而非拒绝。无论如何，让我们不要停止祈祷，因为这样做，我们就像撒种的人，到了时节必要收获。\BibleRef{迦 6:9} 即使我们求对了，同时也要求祂不要成就我们所求错的；因为这在主祷文中也有提及，当我们说：\Quote{不要让我们陷入诱惑}。\BibleRef{玛 6:9} 的确，如果你的请求对你有害，那诱惑就不小。但我们不可漠然不顾主的这段话：祂为阻止任何人以为祂应许成就求祂的人的事，是离开圣父而做的，所以在说了\Quote{你们因我的名无论求什么，我必成就}之后，立刻加上：\Quote{为使圣父在圣子内受光荣：你们若因我的名求什么，我必成就}。因此，圣子绝不离开圣父而行事，因为祂如此行事，正是为叫圣父在祂内受光荣。所以，圣父在圣子内行事，为叫圣子在圣父内受光荣；圣子在圣父内行事，为叫圣父在圣子内受光荣；因为圣父与圣子原是一体。\par

\SourceRecord{Translated by John Gibb.；Nicene and Post-Nicene Fathers, First Series；Vol. 7.；Edited by Philip Schaff.；Buffalo, NY: Christian Literature Publishing Co.,；1888.；<http://www.newadvent.org/fathers/1701073.htm>.}

% structure: p0001 source=h1 target=section reason=page-title

\section{论题74（\BibleRef{若 14:15-17}）}

1. 我们听到了，弟兄们，当福音被宣读时，主说：\BibleQuote{如果你们爱我，就要遵守我的诫命；我要求父，祂必赐给你们另一位护慰者［\OriginalTerm{护慰者}{Paraclete}］，使祂永远与你们同在；就是真理之神；世界不能接受祂，因为看不见祂，也不认识祂；你们却认识祂，因为祂与你们同在，并在你们内。}主在这几句简短的话中有许多可以探讨之处；但我们无论要探究其中一切可探究的，或要找出一切我们所探究的，都太难了。然而，按主乐意赐予我们的能力，并按我们和你们相称的容量，请接受我们应当说的和你们应当听的，亲爱的，接受我们所能给予的，并向祂恳求我们所欠缺的。就是圣神，即护慰者，是基督应许给祂宗徒的；但让我们注意祂赐予应许的方式。祂说：\BibleQuote{如果你们爱我，就要遵守我的诫命；我要求父，祂必赐给你们另一位护慰者，使祂永远与你们同在：就是真理之神。}总之，我们在这里看到了在圣三内的圣神，大公信仰承认祂与圣父圣子\OriginalTerm{同体}{consubstantial}、\OriginalTerm{同永}{co-eternal}。就是祂，宗徒论及祂说：\BibleQuote{天主的爱藉着所赐予我们的圣神，已倾注在我们心中。}\BibleRef{罗 5:5}那么，主怎能说\BibleQuote{如果你们爱我，就要遵守我的诫命；我要求父，祂必赐给你们另一位护慰者}，而祂论及圣神这样说？没有圣神，我们既不能爱天主，也不能遵守祂的诫命。我们怎能爱祂以领受祂，而离了祂我们根本不能爱？我们怎能遵守诫命以领受祂，而离了祂我们没有能力遵守？或者，我们爱基督的爱在我们内居于优先地位，以致我们这样爱基督并遵守祂的诫命，就配得领受圣神，好使天父的爱（不是先前已有的基督的爱）藉着所赐予我们的圣神倾注在我们心中？这种想法是完全错误的。因为若有人相信自己爱圣子，却不爱圣父，他当然不是爱圣子，而是爱自己想象中的幻影。此外，宗徒的宣言是：\BibleQuote{除非受圣神感动，没有人能说\InnerQuote{耶稣是主}。}\BibleRef{格前 12:3}那么，谁称祂为主耶稣，除非是爱祂的人，若他以宗徒所意指的方式称呼祂？因为许多人用嘴唇这样称呼，却在心里和行为上否认祂；正如祂论及这样的人所说：\BibleQuote{他们自称认识天主，但在行为上却否认祂。}\BibleRef{铎 1:16}如果因行为否认祂，那么无疑也是因行为真实呼求祂的名。因此，\Quote{没有人}在心里、言语、行为上，用心灵、嘴唇、双手的劳作\Quote{说\InnerQuote{主耶稣}}——除非在圣神内，没有人说\Quote{主耶稣}；除非那爱祂的，没有人如此呼求祂。因此，宗徒们早已称祂为主耶稣：如果他们这样称呼祂，绝非虚假的言词，口里承认，心里和行为上却否认祂；如果他们以完全真诚的心灵这样称呼祂，那么无疑他们是爱祂的。那么，他们是如何爱的呢？难道不是在圣神内吗？然而，他们被命令爱祂并遵守祂的诫命，是在领受圣神之前且为了领受圣神：然而，没有圣神，他们当然不能爱祂并遵守祂的诫命。\par

2. 因此我们应当明白：爱主的人已经拥有圣神，并且因他所拥有的，成为更圆满拥有圣神的配得者，以便他因拥有更多而能爱得更深。因此，门徒们已经拥有主所应许的那位圣神，因为若没有祂，他们就不能称祂为主；但他们尚未如主所应许的方式拥有祂。因此，他们既拥有又未拥有，因为他们尚未像以后将要拥有的那样完全拥有祂。所以，他们以一种更有限的方式拥有祂：祂还将以更丰盛的方式赐予他们。他们以一种隐藏的方式拥有祂，他们将以一种显明的方式领受祂；因为当前的拥有也关系到圣神更圆满的恩赐，以便他们能意识到自己所拥有的。宗徒正是论及这恩赐时说：\Quote{我们领受的，不是这世界的神，而是由天主而来的神，为使我们知道天主白白赐给我们的一切。}\BibleRef{格前 2:12} 因为主在显现圣神时，并非只一次，而是分两次不同的场合。就在祂从死者中复活后不久，祂向他们嘘了一口气，说：\Quote{领受圣神吧。} 祂当时既然赐下了圣神，难道后来就没有按祂的应许派遣祂吗？或者，岂不正是同一位圣神，既被祂亲自嘘气赐予他们，后来又从天上被祂派遣？\BibleRef{宗 2:4} 那么，为什么祂公开赐予的同一恩赐却发生了两次，这是另一个问题：因为祂公开的双重恩赐可能是由于爱的两条诫命，即爱近人和爱天主，以便表明爱属于圣神。如果还要探究其他理由，我们目前不能让我们的论述因这种探问而不合宜地延长；但必须承认，没有圣神，我们既不能爱基督，也不能遵守祂的诫命；我们对祂临在的经验越少，我们所能做的也越少；经验越丰富，能力也越大。因此，这应许对于没有圣神的人或有圣神的人都不是空的。因为它是向没有的人许下的，以便他能拥有；向有的人许下的，以便他能更丰盛地拥有。因为若不是圣神在有些人身上拥有得比在另一些人身上少，\Person{圣厄里叟}就不会对\Person{圣厄里亚}说：\Quote{愿你里面的圣神加倍地在我里面。}\BibleRef{列下 2:9}\par

3. 但当\Person{若翰洗者}说：\Quote{因为天主不按尺度赐予圣神，}他专指天主子，祂不是按尺度领受圣神；因为在祂内住有整个天主性的圆满。\BibleRef{哥 2:9} 同样，天主与人之间的中保、基督耶稣这个人也并非离了圣神的恩宠而独立；\BibleRef{弟前 2:5} 因为祂亲口告诉我们，先知的话已应验在祂身上：\Quote{主的神在我身上，因为祂给我傅了油，并派我去向贫穷人传报福音。}\BibleRef{路 4:18} 因为祂是独生子，与圣父同等，这并非出于恩宠，而是出于本性；但将人性接纳到独生子的位格统一中，并非出于本性，而是出于恩宠，正如福音本身所承认的：\Quote{孩子渐渐长大，强健起来，充满智慧，天主的恩宠在祂身上。}\BibleRef{路 2:40} 但赐予其他人则是按着尺度——这尺度不断增大，直到各人领受了自己达到完美限度的圆满补足。如宗徒也提醒我们：\Quote{不要看自己超过所当看的，但要怀着清醒的心看，照天主所分给各人信德的尺度。}\BibleRef{罗 12:3} 分受的不是圣神本身，而是圣神所赐的恩惠：因为恩惠虽有区别，却是同一圣神。\BibleRef{格前 12:4}\par

4. 但当祂说：\Quote{我要请求圣父，祂必赐给你们另一位护慰者，}祂表明祂自己也是一位护慰者。\OriginalTerm{护慰者}{Paraclete}在拉丁文中称为\Emph{advocatus}（辩护者）；论及基督也说：\Quote{我们在圣父那里有一位辩护者，就是耶稣基督，那义者。}\BibleRef{若一 2:1} 但祂说世界不能接受圣神，正如同样的话也说：\Quote{肉性的思念是与天主为仇，因为它不服从天主的法律，事实上也不能服从；}就好像我们说：不义不能成为义。因为在此段经文中，祂所说的世界，是指那些爱世界的人；这种爱不是出于圣父。\BibleRef{若一 2:16} 因此，这世界的爱——它给我们足够的工作去削弱和摧毁它在内的势力——直接对立于天主的爱，这爱藉着赐予我们的圣神倾注在我们心中。因此，\Quote{世界，}因此\Quote{不能接受祂，因为看不见祂，也不认识祂。}因为世俗的爱没有那些无形的眼睛，除非以无形的方式，圣神是不能被看见的。\par

5. \BibleQuote{可是你们}，祂接着说，\BibleQuote{将要认识祂；因为祂要与你们同住，并在你们内。}祂将在他们内，为能与他们同住；祂不会为要住在他们内而与他们同住：因为存在某处先于居住在那里。但为避免我们以为祂的话\BibleQuote{祂要与你们同住}是如同客人肉眼可见地与常人同住那样的意思，祂在加上\BibleQuote{祂要在你们内}这句话时，解释了\BibleQuote{祂要与你们同住}的含义。因此，祂是以不可见的方式被看见的；除非祂在我们内，我们无法认识祂。因为我们是类似地看见我们内在的良心：我们看见别人的面孔，却看不见自己的；但我们看见的是自己的良心，而非别人的。然而良心除了在我们内，从不在于别处；但圣神也可以与我们分离，因为祂被赐予，为能也在我们内。但除非祂在我们内，我们无法以那唯一可被看见和认识的方式看见和认识祂。\par

\SourceRecord{Translated by John Gibb.；Nicene and Post-Nicene Fathers, First Series；Vol. 7.；Edited by Philip Schaff.；Buffalo, NY: Christian Literature Publishing Co.,；1888.；<http://www.newadvent.org/fathers/1701074.htm>.}

% structure: p0001 source=h1 target=section reason=page-title

\section{Tractate 75 (John 14:18-21)}

1. 在应许圣神之后，以免有人以为主要赐下祂，作为代替自己的某种方式，以至于祂自己也不再与他们同在，祂补充说：\Quote{我不留下你们为孤儿；我要到你们这里来。} \OriginalTerm{孤儿}{Orphani} 是希腊文，\OriginalTerm{孤儿}{pupilli} 是拉丁文，两者指同一事物；因为在圣咏中，我们读到：\Quote{你是孤儿的扶助}（拉丁译本作 \OriginalTerm{孤儿}{pupillo}），希腊文作 \OriginalTerm{孤儿}{orphano}。因此，尽管天主子并没有收养儿子归向圣父，也没有愿意我们因恩宠而拥有那同一位圣父——祂按本性是圣父——，然而祂自己却以某种方式向我们显示出父爱，当祂宣告说：\Quote{我不留下你们为孤儿；我要到你们这里来。} 同样地，祂也称我们为新郎的子女，当祂说：\Quote{日子将到，新郎要从他们那里被取去，那时新郎的子女就要禁食。} \BibleRef{玛 9:15} 新郎是谁，岂不是主基督吗？\par

2. 祂接着说：\Quote{还有一点点时候，世界就再看不见我了。} 这是怎么回事？世界那时看见祂；因为\Quote{世界}这个名字指的是那些祂在上面论及圣神时所说的人：\Quote{世界不能领受祂，因为看不见祂，也不认识祂。} 祂显然被世界的肉眼看见，当祂在肉身中显现时；但世界没有看见隐藏在肉身中的圣言：它看见了人，但没有看见天主；它看见了遮盖物，但没有看见内在者。然而，在复活之后，连祂的肉身——祂曾向自己的门徒显现，让他们看、让他们摸——祂也拒绝向他人显现，我们或许可以这样理解这句话的意思：\Quote{还有一点点时候，世界就再看不见我了；但你们要看见我：因为我活着，你们也要活着。}\par

3. 这句话是什么意思：\Quote{因为我活着，你们也要活着}？为什么祂用现在时态说自己的活着，而用将来时态说他们的活着？岂不是为了许诺：复活身体的生命，先在祂自己身上实现，也必定会在他们身上相随？由于祂自己的复活就在不久的将来，祂用现在时态表示其临近；而他们的复活，因延迟到世界末日，祂不说\Quote{你们活着}，而是说\Quote{你们要活着}。因此，优雅而简洁地，通过两个词，一个现在时，一个将来时，祂应许了两种复活：即祂自己的复活在不久的将来，以及我们的复活尚在未来、在世界末日。祂说：\Quote{因为我活着，你们也要活着}：因为祂活着，所以我们也要活着。因为死因一人而来，死人复活也因一人而来。就如众人都在亚当内死了，照样众人也都要在基督内复活。\BibleRef{格前 15:21} 正如只有通过前者，每个人才难免一死，也只有通过基督，人才能获得生命。因为那时我们没有生命，我们是死的；因为祂活着，我们也要活着。我们曾为祂而死，当我们为自己活着时；但，因为祂为我们死了，祂既为自己、也为我们活着。因为，祂活着，我们也要活着。因为我们自己只能得到死亡，生命不能靠我们自己获得。\par

4. 祂说：\Quote{在那一天，你们将要明白：我在我的圣父内，你们在我内，我也在你们内。} 在哪一天？岂不是祂所说的\Quote{你们也要活着}的那一天？因为那时，我们将能看见我们所相信的。因为现在祂在我们内，我们也在祂内：我们现在相信这一点，但那时也将知道它；尽管我们现在凭信仰所知道的，那时我们将凭实际的看见而知道。因为只要我们在这身体内——如同现在这样，即败坏而重压灵魂——我们就是远离主而生活；因为我们行事为人是凭信德，而非凭目睹。\BibleRef{格后 5:7} 因此，那时将是凭目睹，因为我们将看见祂实在怎样。\BibleRef{若一 3:2} 因为，如果基督如今不在我们内，宗徒就不会说：\Quote{如果基督在你们内，身体固然因罪而死亡，但灵魂因义德而生活。} \BibleRef{罗 8:10} 但即使在那个时刻，我们也在祂内，祂足以清楚地表明，当祂说：\Quote{我是葡萄树，你们是枝条。} 因此，在那一天，当我们活在那生命里——死亡将被吞没——我们将知道，祂在圣父内，我们在祂内，祂在我们内；因为那时将完成那由祂在现在已经开始的状态：祂在我们内，我们在祂内。\par

5. \Quote{那有我的诫命的，}祂接着又说，\Quote{并遵守的，就是爱我的。}那在记忆里拥有并实行在生活中的人；那口头上拥有并道德上遵守的人；那耳闻并行动的人；那在行为上拥有并以恒心遵守的人——\Quote{就是，}祂说，\Quote{爱我的。}爱因行为而显明，绝非徒然挂名。\Quote{那爱我的，}祂说，\Quote{必蒙我圣父爱他，我也要爱他，并且要向他显现。}但什么是\Quote{我\Emph{要}爱}呢？仿佛祂那时才爱，现在却不爱吗？绝对不是。因为圣父怎能离了圣子爱我们，或圣子离了圣父爱我们呢？既然他们不可分离地工作，又怎能分离地爱？但祂说\Quote{我要爱他，}是指后面的话：\Quote{并且我要向他显现。} \Quote{我要爱，并且要显现；}也就是说，我要爱到显现的地步。因为祂的爱目前的目的是要我们相信，并持守信德的诫命；但那时祂的爱将以我们看见、并得见作为我们信德的赏报为目的：因为我们现在也爱，是相信将来要看见的；但那时我们将要在看见我们现在所相信的事物中去爱。\par

\SourceRecord{Translated by John Gibb.；Nicene and Post-Nicene Fathers, First Series；Vol. 7.；Edited by Philip Schaff.；Buffalo, NY: Christian Literature Publishing Co.,；1888.；<http://www.newadvent.org/fathers/1701075.htm>.}

% structure: p0001 source=h1 target=section reason=page-title

\section{讲道集第七十六篇（\BibleBook{若望}{福音}十四章廿二至廿四节）}

1. 当门徒这样发问，耶稣，他们的师父，答复他们时，我们也仿佛与他们一同学习，因为我们正在阅读或聆听神圣福音。因此，因为主曾说：\BibleQuote{再过片刻，世界就看不见我了；但你们要看见我。}犹达——并非那出卖祂的、别号依斯加里犹达斯的那一位，而是那位其书信被列入正典圣经中的——就此事询问祂说：\BibleQuote{主，祢为什么要将自己显示给我们，而不显示给世界呢？}让我们也作为与他们一同发问的门徒，聆听我们共同的师父。因为圣犹达，不是那个不洁的，是追随者，而非逼迫主的人，询问了其原因：为什么耶稣要显示给祂自己的人，而不显示给世界；为什么再过片刻，世界就看不见祂，而他们却要看见祂。\par

2. \BibleQuote{耶稣回答他说：谁爱我，必遵守我的话，我父也必爱他，我们要到他那里去，并要在他那里作我们的住所。那不爱我的，就不遵守我的话。}这里为我们摆明了祂显示给祂自己的人，而不显示给那一类祂称之为世界的其他人的原因；这也是为什么一些人爱祂，而另一些人不爱祂的原因。这正是那在圣咏中宣告的原因：\BibleQuote{天主，求你审断我，为我辩护，对抗不虔诚的民族。}因为那些爱的人是蒙拣选的，因为他们爱；但没有爱的人，即使能说人间的语言和天使的语言，也成了鸣的锣、响的钹；即使有先知之恩，通晓一切奥秘和一切知识，并有全备的信德，甚至能移山，也一无是处；即使把他们所有的财产分施，并把自己的身体交出焚烧，也毫无益处。\BibleRef{格前 13:1}圣人因着那使同心合意的人共居一室的爱，而与世界分别出来。圣父与圣子在此室中建立居所，并将这爱赐给那些，他们最终也必将以这应许的自我显示而尊荣的人；正是这显示，门徒曾询问了师父，为使不仅当时听见的人从祂口中学习，我们也从祂的福音中学习。因为他询问了关于基督的显示，却听到了关于祂的爱与居住。因此有一种内在的天主显示，完全不为人所知晓，他们得不到天主父和圣神的显示：至于子，或许曾有，但只在肉身中；而且，那既非与前者同类的，也无法以任何形式与他们同在，除非短暂片刻；而且，那还是为审判，而非为喜乐；为惩罚，而非为赏报。\par

3. 因此，我们现在必须按祂乐意启示的那样，理解这些话的含义：\BibleQuote{再过片刻，世界就看不见我了；但你们要看见我。}确实，再过片刻，祂甚至要将祂的肉身——不敬虔的人也能看见的——从他们的视线中除去；因为\Emph{他们}之中，没有一个人在祂复活后看见过祂。但是，既然天使曾见证说：\BibleQuote{祂要怎样来，你们就怎样看见祂升天去了，} \BibleRef{宗 1:11}我们的信德也坚持，祂将带着同一身体来审判生者死者；那么，毫无疑问，那时祂将被世界看见，即那些与祂的国度无关的人。因此，我们最好理解为，当祂说\BibleQuote{再过片刻，世界就看不见我了}时，祂立刻指的是那时刻；即，在世界的终结时，祂将从被定罪者的视线中被取走，以便将来只有那些与祂同在的人，即那些爱祂的人，圣父和祂自己在他们心中建立居所的人，才能看见祂。但祂说\Quote{片刻}，因为人在看为长久的，在天主眼中却是极短暂的：关于同一\Quote{片刻}，我们的圣史若望自己说：\BibleQuote{孩子们，这是最后的时刻了。} \BibleRef{若一 2:18}\par

4. 此外，免得有人想象，只有圣父和圣子，而没有圣神，居住在爱他们的人那里，让他回想一下上面关于圣神所说的：\BibleQuote{世界不能领受祂，因为看不见祂，也不认识祂；但你们却要认识祂，因为祂要与你们同在，并在你们内。} \BibleRef{若 14:17}你看，在这里，圣神也与圣父和圣子一同居住在圣人们内；也就是，在他们内，如同天主在祂的圣殿中。圣三的天主，圣父、圣子及圣神，在我们走向他们时，来到我们这里；他们来是要帮助，我们来是要服从；他们来是要光照，我们来是要瞻仰；他们来是要充满，我们来是要容纳：使得我们对他们的看见不是外在的，而是内在的；他们的居住在我们内不是暂时的，而是永恒的。子以这样的方式并不显示给世界：因为在面前的段落中，世界指的是那些人，祂立刻补充说：\BibleQuote{那不爱我的，就不遵守我的话。}这些人永远看不见圣父和圣神；他们看见子也只不过短暂片刻，不是为了获得福乐，而是为了被定罪；而且，他们看见祂，也不是在祂作为天主的形式中——与圣父和圣神同样不可见——而是以人的形式，祂愿意在苦难中成为可轻蔑的，但在审判世界时成为可畏惧的。\par

5. 但当祂加上说：\BibleQuote{你们所听见的话不是我的，而是那派遣我来的圣父的，} 我们不要充满惊奇或恐惧：祂并不亚于圣父，然而祂除了出自圣父之外，什么也不是：祂本身并不不平等，但祂并非出自自己。因为当祂说：\BibleQuote{那不爱我的，不遵守我的话} 时，祂所说的并非虚言。你看，祂称那些话为自己的话；那么，当祂又说：\BibleQuote{你们所听见的话不是我的} 时，祂岂非自相矛盾？大概，是为了一些特意作出的区分：当祂说 \Emph{祂自己的} 时，祂用了复数 \Quote{话}；但当祂说 \Quote{这话}，即 \OriginalTerm{圣言}{Word}，不是祂自己的，而是圣父的，祂希望这话被理解为是指祂自己。因为 \BibleQuote{在起初已有圣言，圣言与天主同在，圣言就是天主。} 因为作为圣言，祂确实不是祂自己的，而是圣父的：正如祂不是自己的肖像，而是圣父的肖像；祂自己也不是自己的儿子，而是圣父的儿子。因此，祂正确地将祂所做的一切，作为平等者，归诸万有的根源，祂从祂那里得到这特权，即在各方面与祂平等。\par

\SourceRecord{Translated by John Gibb.；Nicene and Post-Nicene Fathers, First Series；Vol. 7.；Edited by Philip Schaff.；Buffalo, NY: Christian Literature Publishing Co.,；1888.；<http://www.newadvent.org/fathers/1701076.htm>.}

% structure: p0001 source=h1 target=section reason=page-title

\section{《\BibleBook{若望}{福音}14:25-27》释义 第七十七篇}

1. 在圣福音的上一课之后，紧接着刚才所读的这段。主耶稣曾说过，凡爱祂的人，祂和圣父将要来，并在他们那里建立居所。但祂在此之前也论及圣神说过：\Quote{你们却要认识祂，因为祂要与你们同在，并在你们内。}\BibleRef{若 14:17}由此我们明白，天主圣三是居住在圣人们内，如同住在自己的圣殿中。但现在祂说：\Quote{我还与你们同在的时候，将这些话对你们说了。}因此，祂所应许的未来居住是一种；而祂所宣示的现时居住是另一种。前者是属神的，内在借心灵实现；后者是有形的，外在借眼耳显示。前者为已得救者带来永福，后者在时间中访问那等待救恩者。就前者而言，主从不离开爱祂的人；就后者而言，祂来来去去。祂说：\Quote{我还与你们同在的时候，将这些话对你们说了。}即是在祂肉身临在中，当时祂正有形可见地与他们交谈。\par

2. 祂又加上说：\Quote{但那护慰者，就是圣神，祂将因我的名被圣父派遣，祂要教导你们一切，并使你们记起我所对你们说的一切。}那么，是圣子说话，圣神教导，以致我们仅仅抓住圣子所说的话，然后借着圣神的教导来理解，仿佛圣子能离开圣神说话，或圣神能离开圣子教导？或者，岂不是圣子也教导，圣神也说话，而当天主说话并教导任何事时，正是天主圣三自身在说话和教导？正因为祂是天主圣三，需要分别介绍每一位，使我们能听到祂们个别的位格，并理解其不可分割的本性。请听圣父说话，在你读到\Quote{主对我说：祢是我的儿子}的地方；也听祂教导，在你读到\Quote{凡听了父而学习的，都到我这里来}的地方。另一方面，你刚才听见圣子说话，因祂论自己说：\Quote{我所对你们说的一切}。你若也愿认识祂为师，请回想那位师傅，祂说：\Quote{你们的师傅只有一位，就是基督。}\BibleRef{玛 23:10}此外，论到圣神，你刚才听说祂是师傅，因为说\Quote{祂要教导你们一切}，也请听祂说话，在\BibleBook{}{宗徒大事录}中，你读到圣神对圣伯多禄说：\Quote{你同他们去，因为我派遣了他们。}因此，整个圣三都说话并教导；但若不按个别位格呈现给我们，那必全然超出人类软弱的理解力。因为祂在其自身全然不可分割，若总是不分别地谈论，就永远不能被认识为三位一体。当我们说到父、子和圣神时，我们当然不是同时说出口，但祂们自身却只能是同时的。当祂加上\Quote{祂要使你们记起}时，我们也当明白，我们被命令不要忘记，这些特别有益的训诫，正是圣神使我们记起的那恩宠的一部分。\par

3. 祂说：\Quote{平安我留给你们，我将我的平安赐给你们。}此处我们在先知书中读到：\Quote{平安上加平安。}祂离去时留下平安给我们，末了祂来到时将赐给我们祂自己的平安。祂在这世界上留下平安给我们，祂自己的平安将在来世赐给我们。祂留下祂自己的平安给我们，我们持守其中，战胜仇敌。祂自己的平安将赐给我们，那时再无仇敌可战，我们将为王统治。祂留下平安给我们，好使我们在此也彼此相爱；祂自己的平安将赐给我们，那时我们将超越任何纷争的可能。祂留下平安给我们，好使我们在世上彼此不审判那各人心中隐秘的事；祂自己的平安将赐给我们，那时祂\Quote{要显露人心的计谋；那时各人才从天主那里得到称赞。}\BibleRef{格前 4:5}然而，在我们拥有平安的事上，无论是祂升天时留给我们的，还是祂将我们领到父那里时所赐的，都是因着祂并借着祂的。祂离开我们时留下的，难道不就是祂永不撤回的亲临吗？因为祂自己就是我们的和平，\Quote{使两者合而为一。}\BibleRef{弗 2:14}所以，正是祂成为我们的平安，无论当我们相信祂在，还是当我们\Quote{看见祂那实在的情形。}\BibleRef{若一 3:2}因为，只要我们在这重压灵魂的腐朽身体中，且凭信德而行，并非凭眼见，\Quote{祂并不离弃那些远离祂的旅客，}\BibleRef{格后 5:6}何况当我们达到那一见时，祂将如何更以自己充满我们呢？\par

4. 但为什么祂说：\Quote{我把平安留给你们，}却没有加上\Quote{我的；}而当祂说：\Quote{我赐给你们，}时却用了这个词？是否\Quote{我的}一词即使没有表达出来也应被理解，因为一次表达出来可能同时指涉两者？或者，是不是在这里也隐藏着一些需要我们祈求、寻求，并为叩门者开启的深奥真理？因为，如果祂将自己的平安理解为祂自己所拥有的那种，那又怎样呢？而祂留在这个世界上的平安，更恰当地说，是我们的平安，而非祂的。因为祂完全没有罪，祂内毫无不和谐的因素；而我们所拥有的平安，在其间我们仍需说：\BibleQuote{宽免我们的罪债。}\BibleRef{玛 6:12} 因此，我们的确拥有某种平安，因为按内在我，我们喜爱天主的法律；但这并非圆满的平安，因为我们看到在我们肢体中另有一条法律，与心意中的法律交战。\BibleRef{罗 7:22} 同样，我们彼此之间也有平安，因为彼此相爱，互相信任；但这样的平安也并不完全，因为我们看不见彼此心中的意念；我们各自对他人某些方面有或好或坏的看法，与实际情况不符。因此，那个平安虽然由祂留给我们，却是我们的平安：因为若不是来自祂，我们就不会拥有它，尽管它如此；但祂自己所拥有的平安并非如此。如果我们到底持守所领受的，那么我们将拥有祂所拥有的，那时我们将没有自身的不和谐因素，我们心中也不会彼此隐藏秘密。但我知道，主的话也可以被理解为仅仅是同一思想的重复：\Quote{我把平安留给你们，我将我的平安赐给你们：}这样，在说了\Quote{平安，}之后，祂仅以说\Quote{我的平安；}来重复；而祂在说\Quote{我留给你们，}时所指的，祂在说\Quote{我赐给你们。}时简单地重复了。各人可以按自己的理解去领会；但我的喜乐在于——我相信这也是你们的喜乐，我亲爱的弟兄们——在此世紧紧把握这平安，在此世我们同心对抗仇敌，以便我们永远渴望那没有仇敌搅扰的平安。\par

5. 但当主接着说：\Quote{我所赐给你们的，不像世界所赐的，}祂所意指的不正是：不像那些爱世界的人所赐的，我赐给你们吗？因为他们给自己平安的目的，是免除诉讼和战争的烦扰，以便在世界的友谊中找到快乐，而非在天主内；尽管他们停止迫害义人，给义人平安，但在没有真正和谐的地方，不可能有真正的平安，因为他们的心不合。因为正如一个人被称为伴侣\OriginalTerm{伴侣}{consort}，是将自己的命运\OriginalTerm{命运}{sortem}与他人结合；同样，他的心进入类似结合的人，可称为\OriginalTerm{同心合意}{concordant}的。因此，亲爱的朋友们，基督将平安留给我们，并将祂自己的平安赐给我们，不是按照世界的方式，而是按照那位创造世界者相称的方式，使我们与祂同心，使我们的心合而为一，好使这颗心，专注于天上之事，能逃脱尘世的朽坏。\par

\SourceRecord{Translated by John Gibb.；Nicene and Post-Nicene Fathers, First Series；Vol. 7.；Edited by Philip Schaff.；Buffalo, NY: Christian Literature Publishing Co.,；1888.；<http://www.newadvent.org/fathers/1701077.htm>.}

% structure: p0001 source=h1 target=section reason=page-title

\section{论道78（\BibleRef{若 14:27-28}）}

1. 弟兄们，我们刚才听到主向祂的门徒所说的这些话：\BibleQuote{你们心里不要烦乱，也不要胆怯。你们听见我曾对你们说：我去，并要到你们这里来：如果你们爱我，就必喜欢我往父那里去，因为父比我大。} 他们的心可能会充满烦乱和恐惧，只因为祂离开他们，尽管祂打算回来；恐怕在牧人离开的期间，狼会袭击羊群。但是，作为天主，祂并没有抛弃那些祂以人性身份离开的人：基督自己同时既是人又是天主。所以，就可见的人性而言，祂离开了；就祂的\OriginalTerm{天主性}{Godhead}而言，祂仍留下：就那受限于地方的本性而言，祂离开了；就那无所不在的本性而言，祂仍留下。那么，为什么他们的心要烦乱和害怕呢？祂的离开他们的视线，却并未改变祂在他们心中的临在。尽管即使是没有地方界限临在的天主，也会从那些背弃祂的人心中离开——他们不是用脚，而是用品行背弃；正如祂来到那些归向祂的人那里——他们不是用面孔，而是以信德，以精神而非肉身接近祂。但为使他们明白，祂说\BibleQuote{我去，并要到你们这里来}只是就祂的人性而言，祂接着说：\BibleQuote{如果你们爱我，就必喜欢我往父那里去，因为父比我大。} 所以，正是就圣子与圣父不相等的那一方面，祂要往父那里去，正如祂以后要从父那里来审判活人死人；而就\OriginalTerm{独生子}{Only-begotten}与生祂者相等的那一方面，祂从不离开父；而是在那不知地方限制的天主性中，与父处处完全相等。因为正如宗徒所说：\BibleQuote{祂既然具有\OriginalTerm{天主形象}{form of God}，} \BibleQuote{祂却不以自己与天主同等为僭越。} \BibleQuote{祂却使自己空虚，取了\OriginalTerm{奴仆形象}{servant-form}，} \BibleRef{斐 2:6} 如此，祂没有失去前者，而是取了后者，在自己空虚的那方面，祂向我们显现出比祂仍与父同在时更为谦卑的状态。因为那是奴仆形象的增加，而神圣本性并没有减损：在取其中之一时，并没有消耗另一个。就前者而言，祂说：\BibleQuote{父比我大}；但就后者而言，祂说：\BibleQuote{我与父原是一体。}\par

2. 让亚略人留心这事，并在他的留心得到医治；好使争辩不至于导致虚妄，更坏的是，导致疯狂。因为正是这奴仆形象，天主子在其中更小，不仅小于圣父，也小于圣神；不仅如此，也小于祂自己，因为祂自己，以天主形象，大于祂自己。因为人基督并不停止被称为天主子，这名称被认为甚至适用于祂独自躺在坟墓中的肉身。当我们宣认我们相信天主的独生子，在般雀比拉多执政时被钉十字架并被埋葬时，我们所宣认的岂不正是这个吗？祂有什么被埋葬呢？除了没有灵魂的肉身。所以，当我们相信被埋葬的天主子时，我们确实将天主子这个名称也归于祂的肉身，因为只有肉身躺在坟墓里。因此，基督自己，天主子，因为以天主形象而与圣父相等，由于祂使自己空虚，没有失去天主形象，而是取了奴仆的形象，祂甚至大于祂自己；因为未被失去的天主形象大于所取的奴仆形象。那么，有什么可惊奇的呢？或者有什么不当之处呢？如果天主子就这奴仆形象说\BibleQuote{父比我大}，而就天主形象，同一的天主子宣称\BibleQuote{我与父原是一体}呢？因为祂们是一体，既然\BibleQuote{圣言就是天主}；而父更大，既然\BibleQuote{圣言成了血肉。} 让我补充一点亚略人和欧诺米人不能否认的事：就这奴仆形象而言，基督为孩童时甚至低于祂的父母，因为根据圣经，\BibleQuote{祂服从} \BibleRef{路 2:51} 于祂的长辈。那么，异端者啊，既然基督既是天主又是人，当祂以人性说话时，你为何诽谤天主？祂亲自赞颂我们的人性；你敢在祂身上中伤天主性吗？你这不信且忘恩的人，你竟要贬低那造你的天主，只因为祂正在宣告祂为了你而成了什么？因为祂与圣父平等，天主子，人由祂所造，却成了人，为了小于圣父：若非如此，人会变成什么？\par

愿我们的主和师傅将这句话清晰印在我们心中：\BibleQuote{如果你们爱我，你们必定喜乐，因为我往父那里去；因为父比我大。}让我们与门徒一起聆听师傅的言语，不要与外人一起留意欺骗者的诡计。让我们承认基督的双重本体：即神性的本体，在其中祂与圣父相等；和人性的本体，在这方面父比祂大。然而两者并非两个，因为基督是一；天主不是\OriginalTerm{四位的}{quaternity}，而是天主圣三。因为正如理性灵魂和身体形成一个人，同样地，基督既是天主又是人，却是一；因此基督是天主，拥有理性灵魂和身体。在这一切中，我们宣认祂是基督，我们在每一个方面宣认祂。那么，创造世界的是谁？\Person{基督耶稣}，但按天主的形象。在\Person{般雀比拉多}手下被钉十字架的是谁？\Person{基督耶稣}，但以仆人的形象。同样地，就祂作为人所包含的各部分而言：不被撇在地狱中的是谁？\Person{基督耶稣}，但仅就祂的灵魂而言。被安葬在坟墓后第三天复活的是谁？\Person{基督耶稣}，但仅就祂的肉身而言。因此，按每一个方面，祂同样被称为基督。然而这一切并非两个或三个，而是一个基督。为此，祂说：\BibleQuote{如果你们爱我，你们必定喜乐，因为我往父那里去；}因为人性值得庆贺，它被\OriginalTerm{独生子圣言}{only-begotten Word}如此取纳，以致在天上成为不朽的；并且它本质上是属土的，却被如此升华和高举，以致像不朽的尘土那样，能坐在圣父的右边。正是在这样的意义上，祂说祂要往父那里去。因为事实上，祂往那常与祂同在的父那里去。但祂往父那里去和离开我们，无非就是转变并赋予不朽祂从我们这里取来的那有死的本性，并将那使祂为人在地上生活的本性高举到天上。凡对基督有如此爱的人，有谁不会从这样的源头汲取喜乐呢？他既能庆贺自己的本性在基督内已经不朽，又怀抱希望自己将来也要藉着基督成为不朽的。\par

\SourceRecord{Translated by John Gibb.；Nicene and Post-Nicene Fathers, First Series；Vol. 7.；Edited by Philip Schaff.；Buffalo, NY: Christian Literature Publishing Co.,；1888.；<http://www.newadvent.org/fathers/1701078.htm>.}

% structure: p0001 source=h1 target=section reason=page-title

\section{讲道集79（\BibleRef{若 14:29-31}）}

1. 我们的主、救主\Person{耶稣基督}曾对祂的门徒说：\BibleQuote{如果你们爱我，你们必定欢喜，因为我往父那里去；因为父比我大。}祂这样说，是出于祂的仆人形象，而非出于祂与父同等的天主形象，这是孝爱之心内的信德所熟知的，而非轻慢和愚昧之人的妄加猜测。接着祂又加上：\BibleQuote{现在我在事情发生之前告诉了你们，使你们当事情发生时，可以相信。}祂说这话是什么意思呢？事实上，一个人在事情发生之前就应当相信那要求他相信的事，因为信德的赞美正是在于所信的是看不见的。因为如果看见才信，有什么伟大呢？就如同一主责备一位门徒时所说：\BibleQuote{因为你看见了，你才相信；那些没有看见而相信的，才是有福的。}我几乎不知道是否有人能说他相信他所看见的；因为同样的信德在\BibleBook{希伯来}{书}中如此定义：\BibleQuote{信德是所希望之事的实质，是未见之事的确证。}因此，如果信德是针对所相信且未见之事，那么主的话：\BibleQuote{现在我在事情发生之前告诉了你们，使你们当事情发生时，可以相信}是什么意思？祂难道不应该说：现在我在事情发生之前告诉了你们，使你们可以相信那件事，当它发生时你们将要看见？因为即使那位被告知\BibleQuote{因为你看见了，你才相信}的人，他并非只相信他所看见的；他看见了一件事，相信了另一件事：他看见祂是人，相信祂是天主。他感知并触摸了那活生生的肉体，他曾看见这肉体在死亡中，他相信了隐藏在这肉体中的天主性。这样，他借助身体感官可见之物，在理智中相信了他所未见之物。但是，尽管我们可以说我们相信我们所看见的，就如每个人都说他相信自己的眼睛一样，但这不可与天主在我们灵魂中建立的信德相混淆；而是从可见之物，我们被引导去相信那不可见之物。因此，亲爱的诸位，在目前这段经文中，当主说\BibleQuote{现在我在事情发生之前告诉了你们，使你们当事情发生时，可以相信}时，祂用\BibleQuote{当事情发生时}显然指他们将在祂死后看见祂活着并升到父那里；看见这事后，他们将被迫使相信祂确实是基督、永生天主之子，因为祂能成就这样的事，甚至在预言之后，并能在行这事之前预言。而他们那时相信，并非以一种新的信德，而是以增强的信德；或者至少，这信德曾因祂的死亡而受损，如今因祂的复活而恢复。因为他们并非先前也不相信祂是天主之子，但当祂自己的预言在祂身上应验时，那在祂此刻对他们说话时仍旧软弱、在祂死亡时几乎消失的信德，重新焕发生机，越发强壮。\par

2. 但祂接着说了什么？\Quote{此后我将不再多与你们说话，因为此世界的首领来了；}这首领是谁？不就是魔鬼吗？\Quote{在我内一无所获。}意思就是没有任何罪过。因为祂用这话指明魔鬼是罪人的首领，而不是受造物的首领；祂在此以\Emph{此世界}的名称来指称罪人。每当\OriginalTerm{世界}{world}一词在恶劣意义上被使用时，祂所指的是仅爱此世界的人。关于这些人，经上另处记载：\Quote{凡愿作此世界之友的，就成了天主的仇敌。}\BibleRef{雅 4:4}那么，我们切不可理解为魔鬼是世界的首领，仿佛他掌管整个世界的统治权，即天地及其中的一切；那一类的世界，当我们讲论圣言基督时，曾说过：\Quote{世界是藉着祂造成的。}因此，整个世界，从至高天到最卑地，都隶属于造物主，而非背弃者；隶属于救赎者，而非毁灭者；隶属于解放者，而非奴役者；隶属于导师，而非欺骗者。魔鬼被理解为世界首领的意义，在宗徒保禄那里更清楚地显明了。他在说\Quote{我们不与血肉搏斗，}即不与世人搏斗，接着又说：\Quote{而是与率领者、掌权者、这黑暗世界的统治者搏斗。}因为他在紧接着的词语中解释了他所说的\Quote{世界}是什么意思，当他补充说\Quote{这黑暗的。}这样，没有人会因世界一名而理解为整个受造界，堕落天使绝不是其统治者。他说\Quote{这黑暗的，}即指爱此世界的人。然而其中有一些蒙选的人，不是凭他们自己的功德，而是凭天主的恩宠。他对这些人说：\Quote{你们从前曾是黑暗，但如今在主内是光明。}\BibleRef{弗 5:8}因为所有人都曾属于这黑暗的统治者，即恶人的统治者，或者说黑暗本身，如同隶属于黑暗一般。但是同一位宗徒说：\Quote{感谢天主，祂拯救了我们脱离黑暗的权势，并将我们移置到祂爱子的国度里。}\BibleRef{哥 1:12}在这位爱子内，此世界即此黑暗的首领一无所有；因为祂作为天主，并非带着罪而来，祂的肉身因童贞玛利亚而生于世，也没有任何遗传的罪染。就好像有人对祂说：那么你为何要死，如果你没有罪以应受死的惩罚？祂立即补充说：\Quote{但为使世界知道我爱圣父，并且圣父怎样命令我，我就照样做：起来，我们离开这里吧。}因为祂正与同席的人坐在一起。祂说\Quote{我们离开，}但去哪里？岂不是去那地方，好让那本身无罪不应死的人被交付于死吗？但祂有圣父要他死的命令，正如预言所说：\Quote{那时我偿还了我未曾抢夺的。}如此，祂注定要全额偿还死亡，虽然一无所欠，并救赎我们脱离我们应得的死亡。因为亚当曾攫取罪作为掠物，当他受骗、放肆地伸手摘取树上的果子，并试图侵犯那不被允许的天主性的不可通传之名，而天主子本性上拥有这名，并非靠抢夺所得。\par

\SourceRecord{Translated by John Gibb.；Nicene and Post-Nicene Fathers, First Series；Vol. 7.；Edited by Philip Schaff.；Buffalo, NY: Christian Literature Publishing Co.,；1888.；<http://www.newadvent.org/fathers/1701079.htm>.}

% structure: p0001 source=h1 target=section reason=page-title

\section{讲道集第80篇（若望福音15:1-3）}

1. 弟兄们，这段福音中，主自称为葡萄树，祂的门徒为枝条，这实际是在宣告：\OriginalTerm{天人中保}{Mediator between God and men}、基督耶稣这个人（\BibleRef{弟前 2:5}）是教会的头，我们是祂的肢体。因为葡萄树与枝条同属一性，因此，既然祂的天主性与我们不同，祂便降生成人，使祂的人性成为葡萄树，而我们这些也同样是人的人成为其枝条。那么，祂说\Quote{我是真葡萄树}是什么意思？祂加上\Quote{真}字，难道是要指引他们注意那字面意义上的葡萄树，即这个比喻所源自的对象？因为祂被称为葡萄树是基于相似性，而非任何位格上的专有，正如祂也被称为羊、羔羊、狮子、磐石、基石，以及其它类似的名称，这些名称本身才是\Emph{真}的，而比喻所取用的只是相似性，并非实体。但当祂说\Quote{我是真葡萄树}时，无疑是与那棵被责备的葡萄树区分开来，正如经上的话：\Quote{你怎么竟变成了酸葡萄，成了野葡萄树？}（\BibleRef{耶 2:21}）因为那棵指望结好葡萄却结了野葡萄的（\BibleRef{依 5:4}），怎能是真葡萄树呢？\par

2. 祂说：\Quote{我是真葡萄树，我父是园丁。凡在我身上不结果实的枝条，祂便剪去；凡结果实的，祂就修剪，使它结更多的果实。}那么，园丁和葡萄树是同一吗？基督是葡萄树，是在祂说\Quote{父比我大}的意义上；但在祂说\Quote{我与父原是一体}的意义上，祂也是园丁。然而，祂并非那些只从事外在劳动的人那样的园丁，而是从内在赐予生长的园丁。\Quote{因为栽种的算不得什么，浇灌的也算不得什么，只有那使之生长的天主才算什么。}但基督确实是天主，因为圣言就是天主；所以祂与父是一体；并且即使圣言成了血肉——成了祂原先所不是的——祂依然保持着祂所是。更甚者，在说了父作为园丁剪去不结果实的枝条、修剪结果实的枝条使其结更多果实之后，祂立刻指明自己也是枝条的修剪者，因为祂说：\Quote{你们因我对你们所讲的话，已是清洁的了。}你看，祂也是枝条的修剪者——这属于园丁的工作，而非葡萄树的；不仅如此，祂还使枝条成为祂的同工。因为虽然枝条不赐予生长，但提供了一些帮助；然而并非出于自己：\Quote{因为离了我，你们什么也不能做。}也听听他们自己的承认：\Quote{那么，阿颇罗算什么？保禄算什么？不过是仆役，使你们相信，照主所赐给各人的。我栽种了，阿颇罗浇灌了。}而这也是\Quote{照主所赐给各人的}；所以不是出于他们自己。然而，接下来的\Quote{但天主使之生长}（\BibleRef{格前 3:5}），祂不是藉着他们工作，而是藉着祂自己；因为这样的工作超越人卑微的能力，超越天使崇高的权能，完全全全地掌握在三一园丁手中。\Quote{你们已是清洁的}，即清洁的，但仍需进一步洁净。因为如果他们不清洁，就不能结果实；然而每个结果实的，都被园丁修剪，以便结更多的果实。他结果实是因为他清洁；而要结更多，他需进一步被洁净。因为此世有谁如此清洁，以致不再需要进一步、再进一步的洁净呢？因为\Quote{如果我们说我们没有罪，就是欺骗自己，真理也不在我们内；但如果我们承认自己的罪，祂是忠信公义的，必将赦免我们的罪，并洗净我们的一切不义}；祂的确洁净那些清洁的，即结果实的，使他们更为清洁，从而结出更多的果实。\par

3. 现在你们因我向你们所说的圣言已是洁净的了。祂为什么不说你们因那洗濯你们的圣洗已是洁净的了，却说\BibleQuote{因我向你们所说的圣言}呢？岂不是因为在水里也是圣言洁净么？除去圣言，水就无非是水。圣言加于元素，便形成圣事，仿佛圣言本身也是一种可见的言语。因为祂在洗门徒的脚时，也曾说过类似的话：\BibleQuote{洗过澡的人，除了脚以外，无须再洗，已是完全洁净了。}水怎能有如此大的效力，以致接触身体就能洁净灵魂呢？若非因圣言的运作，而此运作并非因为它被说出，而是因为它被相信？因为即使在圣言本身中，转瞬即逝的声音是一回事，持久有效的效能是另一回事。宗徒说：\BibleQuote{这就是我们所宣讲的信德之言：你若口里承认\Person{耶稣}是主，心里相信天主使祂从死者中复活，就必得救。因为心里相信归于成义，口里承认归于救恩。}\BibleRef{罗 10:10}因此，我们在\BibleBook{}{宗徒大事录}中读到：\BibleQuote{因信德洁净了他们的心。}\BibleRef{宗 15:9}真福\Person{伯多禄}在他的书信中也说：\BibleQuote{这圣洗如今也拯救了我们，并不是除掉肉体的污秽，而是向天主祈求一个无亏的良心。}\BibleQuote{这就是我们所宣讲的信德之言}，圣洗无疑也藉此被祝圣，以获得洁净的能力。因为\Person{基督}——祂与我们一同是葡萄树，与\Person{圣父}一同是园丁——\BibleQuote{爱了教会，并为她舍了自己。}然后读宗徒的书信，看他接着说什么：\BibleQuote{祂为圣化她，用圣言的水洗净她。}\BibleRef{弗 5:25}因此，洁净绝不能归于这易逝且可朽的元素，若不是因那所加的\BibleQuote{藉着圣言}。这信德之言在天主的教会中具有如此大能，以致藉着那凭信德奉献、祝福、洒水之人的媒介，祂甚至洁净那小小的婴儿，虽然婴儿自己还不能心里相信以致成义，口里承认以致得救。这一切都是藉着圣言成就的，关于这圣言，主说：\BibleQuote{现在你们因我向你们所说的圣言已是洁净的了。}\par

\SourceRecord{Translated by John Gibb.；Nicene and Post-Nicene Fathers, First Series；Vol. 7.；Edited by Philip Schaff.；Buffalo, NY: Christian Literature Publishing Co.,；1888.；<http://www.newadvent.org/fathers/1701080.htm>.}

% structure: p0001 source=h1 target=section reason=page-title

\section{讲道集 81（若望福音 15:4-7）}

1. 耶稣自称为葡萄树，称祂的门徒为枝条，称祂的父为园丁；关于这一点，我们已经尽己所能讲解过了。但在当前这段经文中，祂仍以葡萄树和枝条（即门徒）为喻，继续说：\Quote{你们住在我内，我也住在你们内。}他们住在祂内，与祂住在他们内，方式并不相同。然而，这两种方式都有益于他们，而非有益于祂。因为枝条与葡萄树的关系是：枝条对葡萄树毫无贡献，却从葡萄树获得自己生存的养分；而葡萄树对枝条的关系则是：葡萄树供给枝条生命的滋养，却不从枝条收取任何东西。因此，他们既有基督住在他们内，又住在基督内，这两种情形都有益于门徒，而非有益于基督。因为枝条被砍下后，活根上还能再长出另一根枝条；但被砍下的枝条离开根就不能存活。\par

2. 接着祂继续说道：\Quote{正如枝条若不连在葡萄树上，凭自己不能结果实；你们若不住在我内，也照样不能。}我的弟兄们，这是对恩宠莫大的赞颂——它将教导谦卑者的心灵，并堵住骄傲者的口。现在，那些不认识天主的义，而企图建立自己的义，且未顺服于天主的义的人，若敢回答，就让他们回答吧。\BibleRef{罗 10:3}那些自满自足，以为自己行善无需天主的人，让他们回答吧。那些心地败坏、在信仰上被弃绝的人，\BibleRef{弟后 3:8}他们岂不是要反对这样的真理吗？他们只有充满不敬言语的答复，因为他们说：我们之所以成为人，是出于天主；但我们之所以成为义人，却是出于我们自己。你们这些自欺欺人的人，你们在说什么？你们不是建立自由意志，而是将它从自我高举的高处，通过妄自尊大的虚空领域，抛入深海坟墓的深渊！你们主张人靠自己能成就义德，\Emph{那}正是你们自高自大的巅峰。但真理反驳你们，并宣告：\Quote{枝条若不连在葡萄树上，凭自己不能结果实。}现在，你们尽管冲向你们的眩目悬崖，在没有立足之处的地方，以你们的空谈吹嘘。这些就是你们妄自尊大的虚空领域。但请仔细看看紧随你们脚步的是什么，如果你们还有一点理智，就请毛骨悚然吧。因为凡以为自己凭自己结果实的人，就不在葡萄树上；不在葡萄树上的人，就不在基督内；不在基督内的人，就不是基督徒。这就是你们所坠入的深海。\par

3. 请再三思考真理接着所说的：\Quote{我是葡萄树，你们是枝条；那住在我内，我也住在他内的，就结出许多果实；因为离了我，你们什么也不能做。}为了防止有人以为枝条至少能凭自己结出一点点果实，祂在说了\Quote{就结出许多果实}之后，紧接着的话不是\Quote{离了我，你们只能做很少的事}，而是\Quote{你们什么也不能做。}无论是多是少，没有祂都不可能实现；因为离了祂，什么也不能做。因为虽然枝条结的果实很少时，园丁会修剪它，使它结得更多；但如果它不连在葡萄树上，不从根汲取生命，它自己绝不能结出任何果实。虽然基督若不是人，就不会是葡萄树，但祂若不是天主，就不能将这样的恩宠供给枝条。正因为这恩宠对生命如此必要，以至于连死亡本身也不再受自由意志的支配，祂又补充说：\Quote{谁若不住在我内，便像枝条被丢在外面而枯干了；人便把它们拾起来，投入火中焚烧。}因此，葡萄树的枝条若不连在葡萄树上，就相应地变得更加可鄙，正如它连在葡萄上时是荣耀的；总之，正如主在厄则克耳先知书中论到它们被砍下时所说：它们对园丁的任何用途都无用处，也不能用于工匠的任何劳作。\BibleRef{则 15:5}枝条只适合两件事之一：要么是葡萄树，要么是火；若不在葡萄树上，它的去处就是火里；但愿它为了免于后者，而能留在葡萄树上。\par

4. \Quote{如果你们住在我内，}祂说，\Quote{而我的话也住在你们内，你们无论求什么，必给你们成就。}因为这样住在基督内，他们还能求什么不符合基督心意的事呢？这样住在救主内，他们能求什么与救恩不相符的事呢？事实上，有些事我们求是因为我们在基督内，另一些事我们求是因为我们仍在这世界上。因为有时，由于我们现世的居所，我们内心会想求那些我们知道将来得着并不合适的事。但如果我们住在基督内，天主必不会将那样的恩赐给我们；当我们祈求时，祂只做那对我们有益的事。因此，我们自己住在祂内，祂的话也住在我们内时，我们必求我们所愿的，并且必为我们成就。因为我们若求了，却没有成就，我们所求的便与住在祂内无关，也与祂住在我们内的话无关，而是与我们肉体的那欲望和软弱有关，这等事不在祂内，也没有祂的话住在它们内。总之，祂的话包含祂教导我们的祈祷，我们在其中说：\Quote{我们的天父，你在天上。}\BibleRef{玛 6:9}只要我们不离弃这祈祷的话和意义，凡我们所求的，必为我们成就。因为只有当我们将祂所命令的去行，并爱祂所应许的，才可以说祂的话住在我们内。但若祂的话只住在记忆里，在生活中却没有位置，那枝条就算不得在葡萄树上，因为它不从根吸取生命。圣经的话也注意这点：\Quote{并纪念祂的诫命去遵行的人。}因为许多人只将诫命保存在记忆里，为轻慢它们，甚至讥笑并攻击它们。基督的话并不住在那些只有某种接触、却没有连系的人里面；因此，祂的话对他们不是祝福，反而成为反对他们的见证；因为虽然它们在他们里面，却不居留在他们里面，他们抓住它们，正是为了最后要按它们受审判。\par

\SourceRecord{Translated by John Gibb.；Nicene and Post-Nicene Fathers, First Series；Vol. 7.；Edited by Philip Schaff.；Buffalo, NY: Christian Literature Publishing Co.,；1888.；<http://www.newadvent.org/fathers/1701081.htm>.}

% structure: p0001 source=h1 target=section reason=page-title

\section{第八十二篇论道（若望福音第十五章第八至十节）}

1. 救主这样向门徒说话时，越来越推崇我们得救的恩宠，祂说：\BibleQuote{我父受光荣，在于你们多结果实，并成为我的门徒。}我们说\Quote{光荣}或\Quote{显耀}，两者都是对同一个希腊动词的翻译，即\OriginalTerm{光荣}{\GreekText{δοξάζειν}}。因为在希腊文中的\OriginalTerm{光荣}{\GreekText{δόξα}}，在拉丁文中就是\Quote{光荣}。我认为值得提一下这一点，因为宗徒说：\BibleQuote{如果亚巴郎因行为成义，他就有可夸耀的，但并非在天主前。}\BibleRef{罗 4:2}因为这就是\Quote{在天主前}的光荣，光荣的是天主，而非人，当他不是因行为、而是因信德成义时，以至于他行善也是由天主所赐；正如枝条，如我以上所说，不能自己结果实。因为如果天主父受光荣，在于我们多结果实，并成为基督的门徒，我们就不要将此归功于自己的光荣，好像是我们自己得来的。因为这样大的恩宠来自祂，因此光荣不属于我们，而属于祂。因此，在另一处经文，祂说：\BibleQuote{你们的光也当在人前照耀，好使他们看见你们的善行；}为了阻止他们认为这些善行来自他们自己，祂立刻加上：\BibleQuote{并光荣你们在天之父。}\BibleRef{玛 5:16}因为父受光荣，在于我们多结果实，并成为基督的门徒。而我们怎能成为这样的门徒，除非祂的仁慈先预备了我们？因为我们原是祂的化工，在基督耶稣内受造的，为行善工。\BibleRef{弗 2:10}\par

2. \BibleQuote{如同父爱了我，我也爱了你们；你们应存留在我的爱内。}祂说。这里，你们看到，就是我们善行的根源。因为若非信德以爱德行事，我们怎能拥有善行？\BibleRef{迦 5:6}除非我们先被爱，我们又怎能爱？同一福音作者在他的书信中极为清晰地声明了这一点：\BibleQuote{我们爱天主，因为祂先爱了我们。}\BibleRef{若一 4:19}但当祂说：\BibleQuote{如同父爱了我，我也爱了你们，}祂并非表明我们的本性同祂之间有着与祂和父之间那样的平等，而是表明天主与人之间的中保——人基督耶稣的恩宠。\BibleRef{弟前 2:5}因为当祂说\BibleQuote{父——我，我——你们}时，祂被指明为中保。因为父也爱我们，但在祂内；因为父在此受光荣，即我们在葡萄树上结果实，也就是说，在圣子内，从而成为祂的门徒。\par

3. \BibleQuote{你们应存留在我的爱内。}祂说。我们如何存留？请听下文：\BibleQuote{如果你们遵守我的诫命，你们便存留在我的爱内。}爱导致遵守祂的诫命；但遵守祂的诫命导致爱吗？谁能怀疑是爱在先呢？因为缺乏爱的人没有真正遵守诫命的理由。因此，祂说：\BibleQuote{如果你们遵守我的诫命，你们便存留在我的爱内，}祂并非指出爱的来源，而是指出爱如何彰显。就好像祂说：你们不要以为若你们不遵守我的诫命，你们就存留在我的爱内；因为只有在你们遵守了诫命之后，你们才会存留。换句话说，如果你们遵守我的诫命，这样就会明显证明你们存留在我的爱内。因此，没有人可以欺骗自己，说爱祂却不遵守祂的诫命。因为我们在多大程度上遵守祂的诫命，就多大程度上爱祂；遵守得越少，爱得也越少。虽然当祂说\BibleQuote{你们应存留在我的爱内}时，不清楚祂指的是哪一种爱：是指我们爱祂的爱，还是祂爱我们的爱？但从前一句已可看出。因为祂在前面说：\BibleQuote{我也爱了你们}；紧接着这句话，祂立即加上：\BibleQuote{你们应存留在我的爱内}。因此，这指的是祂爱我们的爱。那么，\BibleQuote{你们应存留在我的爱内}是什么意思，不就是\Quote{你们应存留在我的恩宠内}吗？而\BibleQuote{如果你们遵守我的诫命，你们便存留在我的爱内}是什么意思，不就是\Quote{如果你们遵守我的诫命，你们就会知道你们将存留在我对你们的爱内}吗？因此，我们首先遵守祂的诫命，并不是为了唤起祂对我们的爱；而是除非祂爱我们，我们无法遵守祂的诫命。这是一个恩宠，对谦逊的人完全彰显，向骄傲的人隐藏。\par

4. 但我们对随后所说的应该如何理解：\Quote{就如我遵守了我父的诫命，而常在祂的爱内}？在这里，祂也必然愿意我们明白那父爱，即祂为圣父所爱的爱。因为祂刚才说了这话：\Quote{就如父爱了我，我也爱了你们}；接着祂又加上了这话：\Quote{你们应常在我的爱内}；无疑，就是在祂爱了你们的爱内。因此，当祂论及圣父也说：\Quote{我常在祂的爱内}，我们应当理解是指圣父对祂怀有的爱。然而，在这种情况下，圣父对圣子的爱是否也属于同一恩宠，即圣子爱我们的恩宠？因为我们本身成为儿子，不是出于本性，而是出于恩宠；而独生子却是出于本性，而非恩宠。或者，甚至在圣子自己身上，这爱是否也应归于祂作为人的境况？的确如此。因为祂说：\Quote{就如父爱了我，我也爱了你们}，正是指祂作为中保的恩宠。因为基督耶稣是天主与人之间的中保，不是就祂的天主性而言，而是就祂的人性而言。确实，正是就祂的这人性的而言，我们读到：\BibleQuote{耶稣在智慧和身量上，并在天主和人前的喜爱（恩宠）上与日俱增。} \BibleRef{路 2:52}因此，与此一致，我们可以正确地说，尽管人性不属于天主的本性，但这样的人性却因恩宠而属于天主的独生子之位格；并且这恩宠如此巨大，没有什么比它更大，甚至没有近似相等的。因为在获取人性之前并无任何功绩，但祂的一切功绩都始于这一获取本身。因此，圣子常存在于圣父对祂的爱中，并如此遵守了祂的诫命。因为即使作为人，我们应如何看待祂，岂不就是天主是祂的举扬者？因为圣言就是天主，独生子，与生祂者同永恒；但为使祂因那不可言喻的恩宠而赐予我们作中保，圣言成了血肉，居住在我们中间。\par

\SourceRecord{Translated by John Gibb.；Nicene and Post-Nicene Fathers, First Series；Vol. 7.；Edited by Philip Schaff.；Buffalo, NY: Christian Literature Publishing Co.,；1888.；<http://www.newadvent.org/fathers/1701082.htm>.}

% structure: p0001 source=h1 target=section reason=page-title

\section{论道集 83（\BibleRef{若 15:2-12}）}

亲爱的，你们刚才听到主对祂的门徒说：\Quote{我给你们讲了这些事，是要我的喜乐存在你们内，好使你们的喜乐圆满无缺。} 基督在我们内的喜乐，岂不正是祂因我们而欢欣喜乐吗？祂说这喜乐要成为圆满，除了我们与祂共融之外，又是什么呢？为此祂曾对蒙福的伯多禄说：\Quote{我若不洗你，你就与我无分。} 因此，祂在我们内的喜乐，就是祂赐予我们的恩宠：这恩宠也是我们的喜乐。但祂从永远就因这喜乐而欢欣，因为祂在创造世界以前就拣选了我们（\BibleRef{弗 1:4}）。我们也不能正当地说祂的喜乐不是圆满的，因为天主的喜乐从未在任何时刻有所欠缺。但祂那喜乐并不在我们内：因为我们——这喜乐可能在我们内者——当时尚未存在；甚至当我们开始存在时，这喜乐并不开始存在于祂内。但在祂内，这喜乐常是存在的，因为祂在自己的预知无误的真理中，为我们将属于祂而欢欣。因此，当祂因预知和预定我们而欢欣时，祂对我们所有的喜乐已经是圆满的；并且祂那喜乐中几乎没有掺杂任何恐惧，唯恐祂预知自己将完成的事可能失败。当祂开始行祂预知自己将行的事时，祂的喜乐——作为祂真福的表达——并没有增加；否则创造我们必定会增添祂的真福。弟兄们，愿这样的猜测远离我们的思想；因为天主的真福，没有我们时并不更少，因我们而并不更大。因此，祂因我们得救的喜乐——这喜乐永远在祂内，当祂预知并预定我们时——在我们内开始了，当祂召叫我们时；而这喜乐我们正当地称为我们自己的，因为藉着它，我们也将成为有福的；但这喜乐，作为我们的，增长、进步、并恒久地向前推进，直至它自己的完成。因此，它始于重生者的信德，并在他们复活时获得赏报而完成。这就是我对这些话内涵的意见：\Quote{我给你们讲了这些事，是要我的喜乐存在你们内，好使你们的喜乐圆满无缺}；即，\Emph{我的}\InnerQuote{存在你们内}；\Emph{你们的}\InnerQuote{圆满无缺}。因为我的喜乐一直是圆满的，甚至在你们蒙召之前，当你们被预知为我将召叫的人时；但它也在你们内找到位置，当你们被转化为我对你们所预知的样式。并且\Quote{好使你们的喜乐圆满}：因为你们将是有福的，那是你们现在尚未达到的；正如你们现在被创造，而先前并不存在。\par

2. 祂说：\Quote{这是我的诫命，你们要彼此相爱，如同我爱了你们一样。} 无论我们称之为训令或诫命，两者都是同一个希腊词\Emph{entolé}（\OriginalTerm{诫命}{\GreekText{ἐντολή}}）的翻译。但祂已经在先前的一次场合宣布过同样的话，你们应当记得，我曾尽我所能向你们解释过。因为祂在那里说：\Quote{我给你们一条新诫命，你们要彼此相爱；如同我爱了你们，你们也要彼此相爱。} 因此，这诫命的重复就是它的推荐：只不过祂在那里说：\Quote{我给你们一条新诫命}；而这里说：\Quote{这是我的诫命}；\Emph{那里}，似乎从前没有这样的诫命；\Emph{这里}，似乎祂没有别的诫命可以给他们。但那里称之为\Quote{新}，以免我们继续走旧路；这里称之为\Quote{我的}，以免我们轻视它。\par

3. 但当主在此这样说：\Quote{这是我的诫命}，好像再没有别的了，我的弟兄们，我们该怎样想呢？难道那关于彼此相爱之爱德的诫命，是祂唯一的诫命吗？不是还有一条更大的——即我们该爱天主吗？还是天主确实只把关于爱的命令交给了我们，以致我们无需寻找别的？至少宗徒称赞了三样事，他说：\Quote{如今常存的，有信德、望德、爱德，这三样；其中最大的是爱德。}\BibleRef{格前 13:13} 虽然在爱德（即爱）中包含了两条诫命，但这里只称它为最大的，而非唯一的。这样，关于信德，给了我们多少诫命！关于望德，又是多少！谁能将它们汇集起来，或足数清呢？但让我们思考同一位宗徒的话：\Quote{爱是法律的满全。}\BibleRef{罗 13:10} 所以，哪里有了爱，还缺少什么？哪里没有爱，又有什么能有益处？魔鬼相信（\BibleRef{雅 2:19}），却并不爱；没有相信的人，不会爱。确实，一个不爱的人可以希望获得宽恕，但他的希望是徒然的；但爱着的人不会绝望。因此，哪里有了爱，哪里必然有信德和望德；哪里有了爱近人，哪里也必然有爱天主。因为不爱天主的人，怎能爱近人如同自己呢？何况他连自己也不爱！这样的人既不敬虔又不义；爱不义的人，明显是不爱自己的灵魂，反而恨它。因此，让我们持守主的这条命令——彼此相爱；这样，我们就会做一切其他所吩咐的事，因为其他一切都包含在这条命令内。但这种爱有别于人们彼此间的一般之爱；为标明区别，主加上了说：\Quote{如同我爱了你们一样。}基督为什么爱我们？是为了使我们堪当与基督一同为王。因此，让我们也以此为目标彼此相爱，以显明我们的爱与那些没有如此动机的他人之爱的区别，因为他们本身缺乏这爱。但那些以拥有天主本身为彼此相爱目标的人，才真正彼此相爱；因此，甚至为了彼此相爱这目的，他们也爱天主。这种爱并非人人都有；因为很少人有这动机彼此相爱，为使天主成为万物之中的万有。\BibleRef{格前 15:28}\par

\SourceRecord{Translated by John Gibb.；Nicene and Post-Nicene Fathers, First Series；Vol. 7.；Edited by Philip Schaff.；Buffalo, NY: Christian Literature Publishing Co.,；1888.；<http://www.newadvent.org/fathers/1701083.htm>.}

% structure: p0001 source=h1 target=section reason=page-title

\section{讲道集第八十四篇（\BibleBook{若望}{福音}15:13）}

1. 亲爱的弟兄们，主为我们彼此之间应有的圆满之爱定下了界说，祂说：\BibleQuote{人若为朋友舍命，再没有比这更大的爱了。}因此，既然祂以前说过：\BibleQuote{这是我的命令：你们该彼此相爱，如同我爱了你们一样}；并在这些话后又附加上你们刚才听到的：\BibleQuote{人若为朋友舍命，再没有比这更大的爱了}；由此必然得出结论，也就是同一位福音作者\Person{若望}在他的书信中所说的：\BibleQuote{正如基督为我们舍命，我们也当为弟兄舍命}\BibleRef{若一 3:16}；彼此真实相爱，如同祂爱了我们，祂为我们舍命。这也无疑是在\Person{撒罗满}的箴言中所读到的含义：\BibleQuote{你若坐下与一位统治者共进晚餐，要明智地思量摆在你面前的是什么；然后伸手取用，要知道你也必须准备同样的筵席。}因为统治者的筵席，岂不是我们从中领受那为我们舍命者的圣体圣血的那同一筵席吗？在那里就坐，岂不是谦恭地靠近吗？明智地思量摆在你面前的是什么，岂不是恰当地反省恩惠的伟大？至于这样伸手取用，知道你也必须准备同样的筵席，岂不是如我已说过的，正如基督为我们舍命，我们也当为弟兄舍命？正如宗徒\Person{伯多禄}也说：\BibleQuote{基督为我们受苦，给我们留下了榜样，叫我们追随祂的足迹。}\BibleRef{伯前 2:21}这就是准备同样的筵席。这就是真福殉道者们以炽热的爱所做的；如果我们不仅以空洞的形式庆祝他们的纪念，并且在他们自己饱享的筵席中，走近主的筵席，我们就必须像他们一样，也为自己准备同样的筵席。因为正因如此，我们在这同一筵席上纪念他们时，不像纪念那些安息于和平中的其他亡者那样为他们祈祷，反而是请他们为我们祈祷，好使我们能追随他们的足迹；因为他们实际上已达到那圆满的爱，我们的主告诉我们，没有比这更大的爱了。因为他们为弟兄们所表现的爱之标记，正是他们自己在主的筵席上同样领受的。\par

2. 但我们不可被认为是这样说的，好像基于这样的理由，我们就有可能达到与主基督同等的地位，如果我们为祂的缘故而经历了殉道流血的事。\Emph{祂}有能力舍命，也有能力再取回来；但我们却没有能力按自己的意愿活下去；而我们必然要死，即使不情愿；\Emph{祂}借着死亡，立刻在自己内消灭了死亡；我们却借着祂的死亡，从死亡中被解救出来；\Emph{祂的}肉身未见腐朽\BibleRef{宗 2:31}；我们的肉身，在腐朽之后，将在世界穷尽时由祂穿上不朽。\Emph{祂}无需我们就能成就我们的救恩；我们离了祂却一无所能。\Emph{祂}将自己作为葡萄树赐给我们这些枝条；我们离了祂就没有生命。最后，虽然弟兄为弟兄而死，但没有一个殉道者的血是为赦免弟兄的罪而流的，正如祂为我们所做的那样；在这方面，祂赐给我们的不是可效法的事，而是可庆幸的事。因此，就殉道者为弟兄流血而言，他们表现出了他们自己在主的筵席上所领悟的那种爱的标志。（人可以在死上效法祂，却无人能在救赎上效法祂。）那么，在我所说的一切中，虽然我无法一一列举，基督的殉道者远远不如基督自己。但若有人将自己与基督相比，我不是说与祂的能力相比，而是与祂的无罪相比，并且（我不愿说）他以为自己在医治他人的罪，至少以为自己没有罪，那么他的贪求甚至已越过了救恩方法的界限；这对他是相当重要的事，只是他达不到他的愿望。幸好他在\BibleBook{}{箴言}的那段经文中得到了提醒，接着便说：\BibleQuote{但若你的贪求太大，就不要贪图他的美味；因为你什么都不取，比取不恰当的反倒更好。因为这些事}，接着又说，\BibleQuote{有欺骗的生命}，即虚伪的生命。因为当他宣称自己无罪时，他不能证明，只能假装自己是义人。所以经上说：\BibleQuote{因为这些事有欺骗的生命。}唯有那一位能同时具有人性而无罪。我们被命令做随之而来的是合宜的；这样的话和箴言很适合人性的软弱，它说：\BibleQuote{你这贫穷的人，不要与那富足的人对垒。}因为富足的人是基督，祂从未因遗传或个人的债务而应受惩罚，且祂自己是义人，并成义他人。你这个贫穷的人，不要与祂对垒，你在众人眼前每天在祈祷中为罪赦而乞讨。\BibleQuote{但要保守你自己}，他说，\BibleQuote{远离你自己的谋略}（英译：\InnerQuote{远离你自己的智慧}）。远离什么？就是远离这种迷惑的狂妄。因为祂，不仅是人，也是天主，永远不会被指控有恶。\BibleQuote{因为若你转眼注视祂，祂将无处可见。}\BibleQuote{你的眼睛}，即人的眼睛，你用来辨别属人的事物；\BibleQuote{若你转眼注视祂，祂将无处可见}，因为祂不能用你那样视觉器官看见。\BibleQuote{因为祂会为自己预备翅膀如鹰，飞到祂的看守者的家中}，祂正是从那里来到我们这里，发现我们不像祂自己那样。因此，让我们彼此相爱，如同基督爱了我们，并为我们舍弃了自己。\BibleRef{迦 2:20}\BibleQuote{人若为自己的朋友舍命，再没有更大的爱了。}并让我们以敬畏顺服的精神效法祂，以致我们永不敢妄自将祂与我们自己相比。\par

\SourceRecord{Translated by John Gibb.；Nicene and Post-Nicene Fathers, First Series；Vol. 7.；Edited by Philip Schaff.；Buffalo, NY: Christian Literature Publishing Co.,；1888.；<http://www.newadvent.org/fathers/1701084.htm>.}

% structure: p0001 source=h1 target=section reason=page-title

\section{论道 85（\BibleBook{若望}{福音} 15:14-15）}

1. 主耶稣既称赞了祂为我们而死所展现的爱情，并说：\BibleQuote{人若为自己的朋友舍命，再没有比这更大的爱了。}祂又加上：\BibleQuote{你们若遵行我所命令你们的，就是我的朋友。}何其伟大的屈尊！因为一个人甚至连善仆也算不上，除非他遵行主人的命令；而祂愿意那只能证明人是善仆的途径，也成为识别祂朋友的方式。但我所谓的屈尊，是指主屈尊称那些祂知道是祂仆人的人为祂的朋友。因为，为让我们知道仆人应当服从主人命令的本分，祂实际上在另一处责备那些仆人说：\BibleQuote{你们为什么称呼我\InnerQuote{主啊，主啊}，却不照我说的做？}\BibleRef{路 6:46}因此，当你称呼\Quote{主}时，就要以遵行我的命令来证明你的话。难道祂不是将来有一天要对那顺从的仆人说：\BibleQuote{做得好，善良的仆人；因为你既在少许事上忠心，我必要派你管理许多事：进来享受你主人的快乐！}\BibleRef{玛 25:21}所以，一个善仆，既可以是仆人，也可以是朋友。\par

2. 但让我们注意下文。\BibleQuote{以后我不再称你们为仆人，因为仆人不知道他主人所做的事。}那么，当祂说\BibleQuote{以后我不再称你们为仆人，因为仆人不知道他主人所做的事}时，我们该如何理解善仆既是仆人又是朋友呢？祂这样引出\Quote{朋友}这名称，以至于撤回了\Quote{仆人}这名称；并非似乎将两者包含在一个称呼里，而是为了让一个名称取代另一个腾出的位置。这是什么意思？是否意味着即使在遵行主的命令时，我们也不会是仆人？还是说，当我们成了善仆之后，我们就不再是仆人？然而，谁能反驳真理本身的话，当祂说\BibleQuote{以后我不再称你们为仆人}时，并说明祂这样说是因为：\BibleQuote{仆人，}祂接着又说，\BibleQuote{不知道他主人所做的事。}难道一个善良且经过考验的仆人，主人不也托付他秘密吗？那么祂说\BibleQuote{仆人不知道他主人所做的事}是什么意思？就算他\BibleQuote{不知道他主人所做的事}，难道他也不知道主人所命令的事吗？如果是这样，他如何服事？或者一个不做事的仆人又怎么算是仆人？然而主这样说道：\BibleQuote{你们若遵行我所命令你们的，就是我的朋友。以后我不再称你们为仆人。}真是奇妙的言论！既然我们只有遵行祂的命令才能事奉主，那么遵行时又怎么会不再是仆人？如果我遵行祂的命令时不是仆人，然而除非我遵行，又不能事奉祂，那么我正是在事奉中不再是仆人。\par

3. 弟兄们，让我们领悟，愿主使我们领悟，也使我们能实行我们所领悟的。如果我们知道这一点，我们就真知道主所做的事；因为只有主如此使我们能够，也唯有通过这些方法，我们才能获得祂的友谊。正如有两种恐惧，产生两类恐惧者；同样也有两种服役，产生两类仆人。有一种恐惧，是完美的爱所驱逐的；\BibleRef{若一 4:18}另有一种恐惧，是洁净的，且永远存留。那不在爱内的恐惧，宗徒指出说：\BibleQuote{因为你们没有领受再陷入恐惧的服役之灵}。\BibleRef{罗 8:15}但他论及洁净的恐惧时说：\BibleQuote{不可心高气傲，而要恐惧}。\BibleRef{罗 11:20}在被爱驱逐的恐惧中，连同它一起的服役也当被驱逐：因为宗徒将两者联系在一起，即服役和恐惧，当他说：\BibleQuote{因为你们没有领受再陷入恐惧的服役之灵}。主也正是针对与此种服役相连的仆人说：\BibleQuote{以后我不再称你们为仆人，因为仆人不知道他主人所做的事}。当然不是指因洁净的恐惧而特有的仆人，即那被说\BibleQuote{做得好，善良的仆人，进入你主人的喜乐}的仆人；而是指那以被爱驱逐的恐惧为特征的仆人，主在别处说到他：\BibleQuote{仆人不能永远住在家里，但儿子却永远住下}。因此，既然祂赐给我们成为天主子女的权能，让我们不要做仆人，而做儿子：这样，以某种奇妙而不可言喻却真实的方式，我们作为仆人却能有权能不做仆人；确实，以那洁净的恐惧区分那进入主人喜乐的仆人，但不以那必须被驱逐的恐惧，它标记那不能永远住在家里的人。但要记住，是主使我们能服役而不做仆人。这正是那仆人不知道他主人所做的事；他做什么好事时，就自以为是自己做的，而非他的主；因此，他不以主夸耀，而以自己夸耀，从而欺骗自己，因为夸耀，好像没有领受。\BibleRef{格前 4:7}但亲爱的，为了我们能成为主的朋友，让我们知道我们的主所做的事。因为是祂不仅使我们成为人，也使我们成为义人，不是我们自己。除了祂，还有谁能是引导我们达到这种认识的行事者？因为\BibleQuote{我们领受的，不是世界的精神，而是那从天主来的圣神，为使我们知道天主白白赐予我们的事}。\BibleRef{格前 2:12}所有的善，都是由祂白白赐予的。因为这善也是由那恩赐一切善的祂所赐予的这种认识之恩；为使那夸耀的，在主内夸耀。\BibleRef{格前 1:31}但接下来的话：\BibleQuote{我却称你们为朋友，因为凡我从我父听来的一切，我都显示给你们了}。是如此深奥，我们绝不能压缩在本讲演的范围内，而要留到下一次。\par

\SourceRecord{Translated by John Gibb.；Nicene and Post-Nicene Fathers, First Series；Vol. 7.；Edited by Philip Schaff.；Buffalo, NY: Christian Literature Publishing Co.,；1888.；<http://www.newadvent.org/fathers/1701085.htm>.}

% structure: p0001 source=h1 target=section reason=page-title

\section{讲道集 86（若 15:15-16）}

1. 探讨主这些话应如何理解，确实是值得研究的问题：\BibleQuote{但我称你们为朋友，因为凡我从父那里听到的，我都显示给你们了。}因为谁敢断言或相信，有人知道独生子从父那里听到的一切呢？既然没有人能理解祂如何听到父的任何话，因为祂自己就是父的独一圣言。不仅如此，稍后在同一篇讲论中，即祂在晚餐与受难之间向门徒所说的，祂说：\BibleQuote{我还有许多事要告诉你们，但你们现在不能担负。}那么，既然有许多事祂没有说，只是因为祂知道门徒现在不能担负，我们又怎能理解祂已将自己从父那里听到的一切都告诉了门徒呢？无疑地，祂说自己已做了祂将来要做的事，正如那使尚未存在之事存在的同一存在者。\BibleRef{依 45:11}因为正如祂通过先知所说：\BibleQuote{他们刺穿了我的手和我的脚，}而不是说\InnerQuote{他们将刺穿}；祂似乎是说过去的事，却预言着尚未到来之事：同样，在我们面前的这段经文中，祂说祂已把一切告诉了门徒，而祂知道祂将在知识的圆满中才使门徒知道一切，关于此宗徒说：\BibleQuote{但当那圆满的来到时，那局部的就要消逝。}因为在同一处他补充说：\BibleQuote{现在我知道的只是局部，但那时我将知道，如同我全被知道一样；现在是通过镜子在谜语中，但那时是面对面。}因为同一位宗徒也说，我们因重生的洗涤而得救，\BibleRef{铎 3:5}他又在另一处宣称：\BibleQuote{我们得救在于希望：但看见的希望就不是希望；因为人看见了，为什么还希望呢？但如果我们希望那看不见的，我们就耐心等待它。}\BibleRef{罗 8:24}与此相似，他的同僚宗徒伯多禄也说：\BibleQuote{虽然你们现在看不见祂，却相信祂；当你们看见祂时，你们要因那不可言喻、充满光荣的喜乐而欢欣：获得信德的酬报，就是灵魂的救恩。}\BibleRef{伯前 1:8}因此，如果现在正是信德的时节，信德的酬报就是灵魂的救恩；那么，在那以爱德工作的信德中，\BibleRef{迦 5:6}谁能怀疑那日子必将结束，并在结束时获得酬报呢？不仅是身体的救赎，宗徒保禄所说的，\BibleRef{罗 8:23}而且是灵魂的救恩，如宗徒伯多禄告诉我们的。因为从两者而来的幸福，在现今这个必死的状态中，与其说是实际拥有，不如说是希望的事。但我们需要记住这一点：我们外在的人，即身体，正在衰败；而内在的人，即灵魂，却日日更新。\BibleRef{格后 4:16}因此，当我们等候将来的肉身不死和灵魂救恩时，凭着已领受的抵押，我们可以说已经得救了；所以，对于独生子从父那里听到的一切知识，我们应视为仍属于未来的希望之事，尽管基督已宣告祂已传给了我们。\par

2. \BibleQuote{不是你们拣选了我，}祂说，\BibleQuote{而是我拣选了你们。}这种恩宠是无法言喻的。因为当基督尚未拣选我们时，我们是什么？我们因此仍缺乏爱。因为那拣选祂的人，怎能爱祂？你们想，我们是否处于圣咏中所唱的状态：\Quote{我宁愿在主的殿中做小卒，也不愿住在恶人的帐幕里？}当然不是。那么我们是什么？不过是罪人、丧亡者。我们尚未信祂，以至于祂拣选我们；因为如果祂拣选的是已经相信的人，那么祂在拣选之前自己先被拣选了。但祂怎能说\Quote{不是你们拣选了我}，除非祂的仁慈先于我们？那些维护天主的预知以反对祂恩宠的人，他们的虚妄推理在此失误，他们宣称我们是在世界创立之前被拣选的，\BibleRef{弗 1:4}因为天主预知我们会成为善的，而不是祂自己要使我们成为善的。说\Quote{不是你们拣选了我}的那位，并不这样说。因为如果祂基于预知我们会成为善而拣选我们，那么祂也会预知我们不会首先拣选祂。因为我们不可能以其他方式成为善：除非有人从未行善却被称为善。那么，在那些不善的人中祂拣选了什么呢？他们不是因善而被拣选，因为他们若不被拣选就不能成为善。否则，如果我们坚持功劳的优先性，恩宠就不再是恩宠。确实，这就是恩宠的拣选，宗徒说：\BibleQuote{同样，现在也按照恩宠的拣选，留下了残余。}他又补充：\BibleQuote{既然是出于恩宠，就不是出于作为，否则恩宠就不再是恩宠。}\BibleRef{罗 11:5}你这忘恩负义的人，听吧：\BibleQuote{不是你们拣选了我，而是我拣选了你们。}并非让你说：我已被拣选，因为我已经相信。因为如果你信祂，那么你已经拣选了祂。但是听吧：\BibleQuote{不是你们拣选了我。}并非让你说：在我相信之前，我已经做了善工，因此我被拣选。因为在信德之前能有什么善工呢？宗徒说：\BibleQuote{凡不出于信德的，都是罪。}\BibleRef{罗 14:23}那么，听到\Quote{不是你们拣选了我}这些话，我们该说什么呢？只能是我们是邪恶的，而祂拣选我们，是为了使我们因拣选我们的祂的恩宠而成为善。因为如果功劳在先，就不是恩宠；但这是恩宠，因此恩宠不是找到功劳，而是生成功劳。\par

3. 那么，亲爱的，请看：祂不是拣选好人，而是使祂所拣选的人成为好人。祂说：\BibleQuote{我拣选了你们，}祂说，\BibleQuote{并派你们该去结果实，并使你们的果实存留。}这果实不就是祂曾说过的\BibleQuote{没有我，你们什么也不能做}中的果实吗？因此，祂拣选了并派我们该去结果实；而我们原本没有任何果实能促使祂拣选我们。祂说：\BibleQuote{你们该去，}祂说，\BibleQuote{并结果实。}我们去结果实，而祂自己就是我们所行的道路，也是祂派我们行的道路。祂的慈悲在一切事上已先于我们。祂说：\BibleQuote{并使你们的果实存留，好使你们因我的名无论向圣父求什么，祂都会赐给你们。}因此，让爱德存留，因为祂自己就是我们的果实。这爱德如今在于渴慕之情，尚未臻于圆满的享乐；凡我们怀着此渴慕之情因独生子的名而求的，圣父必赐予我们。但若所求无益于我们的救恩，我们不要以为我们是在以救主的名求；我们只有在救主的名内求那真正属于救恩之途的事物。\par

\SourceRecord{Translated by John Gibb.；Nicene and Post-Nicene Fathers, First Series；Vol. 7.；Edited by Philip Schaff.；Buffalo, NY: Christian Literature Publishing Co.,；1888.；<http://www.newadvent.org/fathers/1701086.htm>.}

% structure: p0001 source=h1 target=section reason=page-title

\section{第87篇（\BibleRef{若 15:17-19}）}

1. 在之前的福音课中，主曾说过：\BibleQuote{你们不是拣选了我，而是我拣选了你们，并委派你们去结果实，并且你们的果实应当存留；使你们因我的名无论向父求什么，祂必赐给你们。}对于这些话，你们记得我们已经按主所允许的讲过了。但在这里，也就是在你们刚才听到的后续课文中，祂说：\BibleQuote{我命令你们这些事，就是要你们彼此相爱。}由此我们应明白这就是我们的果实，就是祂所说的：\BibleQuote{我拣选了你们，要你们去结果实，并且你们的果实应当存留。}祂接着说的：\BibleQuote{使你们因我的名无论向父求什么，祂必赐给你们，}祂一定会赐给我们，只要我们彼此相爱；因为正是这件事祂也赐给了我们，当我们还没有果实时就拣选了你们，因为我们没有拣选祂；并委派我们去结果实——也就是要彼此相爱——这果实没有祂我们不可能有，就像枝条离了葡萄树不能做什么。因此，我们的果实就是爱德，宗徒解释为\BibleQuote{出于纯洁的心、良好的良心和真诚的信德。}\BibleRef{弟前 1:5}所以我们要彼此相爱，也要爱天主。因为如果我们不爱天主，我们就不会以真爱彼此相爱。因为如果一个人爱天主，他就爱近人如同自己；如果他不爱天主，他就不爱自己。因为全部法律和先知都系于这两条爱的诫命：\BibleRef{玛 22:40}这就是我们的果实。因此，关于这种果实，祂命令我们说：\BibleQuote{我命令你们这些事，就是要你们彼此相爱。}同样，宗徒保禄在希望将圣神的果实与肉性的行为对比时，也将此作为起点，说：\BibleQuote{圣神的果实就是爱；}然后，似乎从这点出发并与之相连，他编织了其他果实，即\BibleQuote{喜乐、平安、忍耐、良善、仁慈、信德、柔和、节制。}\BibleRef{迦 5:22}因为谁能真正喜乐，如果他爱的不是作为其喜乐源泉的善？谁能有真正的平安，如果他没有与他真正所爱的那一位平安？谁能因坚持不懈地行善而长久忍耐，如果不是出于炽热的爱？谁能仁慈，如果他不爱他所帮助的人？谁能良善，如果他不因爱而成为良善？谁能有健康的信德，如果没有那藉着爱德而运作的信德？谁的柔和能有益于品格，如果不是由爱所调节？谁又能远离那使人卑贱的事，如果他不爱那使人高贵的事？因此，善师如此频繁地称赞爱，作为唯一需要称赞的事，因为没有爱，其他一切善事都无功效，而拥有爱不可能不带来那些使人真正成为善人的其他善事。\par

2. 但是，在这爱之外，我们也应当忍耐地承受世界的仇恨。因为世界必然仇恨那些它觉察到在回避它所爱之事的人。但是主以自己的事例给我们特别的安慰，因为祂在说了\BibleQuote{我命令你们这些事，就是要你们彼此相爱}之后，加上说：\BibleQuote{如果世界恨你们，你们要知道，它恨你们以前，已经恨了我。}那么，肢体为何要自高于头呢？如果你不愿与头一同忍受世界的仇恨，你就是拒绝在身体内。\BibleQuote{如果你们属于世界，}祂说，\BibleQuote{世界必爱自己的。}当然，祂是说整个教会，祂也常用世界的名称来称呼教会本身；正如经上说：\BibleQuote{天主在基督内使世界与自己和好。}\BibleRef{格后 5:19}又说：\BibleQuote{人子来不是为审判世界，而是为叫世界藉着祂而得救。}\BibleRef{若 3:17}若望在他的书信中说：\BibleQuote{我们在父那里有一位中保，就是正义的耶稣基督；祂为我们的罪作了赎罪祭，不仅为我们的罪，也为全世界的罪。}\BibleRef{若一 2:1}那么，全世界就是教会，然而全世界恨教会。所以，世界恨世界：敌对的恨那和好的，被定罪的恨那得救的，被玷污的恨那被洁净的。\par

3. 但是，那个天主在基督内使之与自己和好的世界，那个被基督拯救、所有罪都被基督白白赦免的世界，是从敌对、被定罪、被玷污的世界中被拣选出来的。因为从那个全都在亚当内灭亡的团块中，形成了仁慈的器皿；那和好的世界就是由这些器皿组成的，它被属于忿怒的器皿的世界所恨，那些忿怒的器皿是从同一团块中形成并注定要毁灭的。最后，在说了\BibleQuote{如果你们属于世界，世界必爱自己的}之后，祂立刻加上：\BibleQuote{但因为你们不属于世界，而是我从世界中拣选了你们，所以世界恨你们。}所以这些人本身也曾属于那个世界，而为了不再属于它，他们被从中拣选出来，不是由于他们自己的任何功德，没有任何善行先于；也不是由于本性，因为本性藉着自由意志已经从根源上完全败坏了；而是白白地，即出于实际的恩宠。因为那位从世界中拣选世界的主，为自己创造了而非找到了祂所应当拣选的；因为\BibleQuote{按照恩宠的拣选，有剩余的遗民得救了。既是出于恩宠，}祂接着说，\BibleQuote{就不是出于作为；不然，恩宠就不再是恩宠了。}\BibleRef{罗 11:5}\par

4. 但若有人问我们，那仇恨救赎世界的灭亡世界所抱持的自爱是什么，我们回答：它爱自己，当然是一种虚假的爱，而非真实的。因此，它虚假地爱自己，真实地恨自己。因为那爱邪恶的，就是恨自己的灵魂。然而，它被称为爱自己，是因它爱那使它邪恶的邪恶；另一方面，它被称为恨自己，是因它爱那伤害它的。因此，它恨自己内真实的本性，而爱恶习：它恨自己所是的——天主美善所造之物，而爱那由自由意志在自己内造成的。所以，若我们正确理解，我们既被禁止去爱它，又被命令去爱它：当我们被告知\BibleQuote{不可爱世界；}时，我们被禁止；\BibleRef{若一 2:15}当我们被告知\BibleQuote{爱你们的仇敌。}时，我们被命令。\BibleRef{路 6:27}这些就是那恨我们的世界。因此，我们被禁止爱它在自身所爱的，而被命令爱它在自身所恨的，即天主的化工，以及祂美善的各种安慰。因为我们被禁止爱其中的恶习，而被命令爱其本性，而它爱自身内的恶习，恨那本性：如此我们能够正当地爱它和恨它，而它却悖逆地爱自己和恨自己。\par

\SourceRecord{Translated by John Gibb.；Nicene and Post-Nicene Fathers, First Series；Vol. 7.；Edited by Philip Schaff.；Buffalo, NY: Christian Literature Publishing Co.,；1888.；<http://www.newadvent.org/fathers/1701087.htm>.}

% structure: p0001 source=h1 target=section reason=page-title

\section{\WorkTitle{讲道集 88（若望福音 15:20-21）}}

1. 主在劝勉祂的仆人耐心忍受世人的仇恨时，给他们提出了不比祂自己的榜样更伟大或更好的榜样；因为正如宗徒伯多禄所说：\BibleQuote{基督为我们受苦，给我们留下了榜样，叫我们跟随祂的脚踪。}\BibleRef{伯前 2:21} 如果我们果真如此行，乃是藉着祂的助佑而行的，祂曾说：\Quote{离了我，你们什么也不能做。} 但更进一步，对祂已经说过\Quote{世界若恨你们，你们该知道，它已先恨了我}的那些人，祂现在又在你们刚才听到的福音经文中说：\Quote{你们要记念我从前对你们所说的话：仆人不能大于主人；人们若迫害了我，也要迫害你们；若遵守了我的话，也要遵守你们的话。} 现在，当祂说\Quote{仆人不能大于主人}时，难道不是清楚地指明祂要我们如何理解祂在上文所说的\Quote{我不再称你们为仆人}吗？因为，你们看，祂称他们为仆人。因为\Quote{仆人不能大于主人：人们若迫害了我，也要迫害你们}这句话还能有什么别的含义呢？因此，很明显，当经上说\Quote{我不再称你们为仆人}时，祂是指那位不能永远住在屋里、且具有被爱所逐出的恐惧的仆人而言的；\BibleRef{若一 4:18} 然而，当此处说\Quote{仆人不能大于主人：人们若迫害了我，也要迫害你们}时，所指的仆人是那位具有常存的洁净敬畏的仆人。因为正是这个仆人将听到：\BibleQuote{好！善良的仆人，进入你主人的福乐吧！}\BibleRef{玛 25:21}\par

2. \Quote{但这一切事，}祂说，\Quote{他们都要因我的名向你们施行，因为他们不认识那派遣我来的。} 那\Quote{这一切事}到底是指什么，他们\Quote{将要行}的呢？不就是祂刚才所说的话，即他们会恨你们、迫害你们，并轻视你们的话吗？因为他们如果不遵守他们的话，却又既不恨也不迫害他们；或者即使他们恨，却不迫害他们：那就不算是\Emph{这一切事}了。但\Quote{他们因我的名向你们施行这一切事}——这不等于说，他们会在你们身上恨我，在你们身上迫害我；而你们的话，正是因为属于我，他们就不遵守吗？因为\Quote{他们因我的名向你们施行这一切事}——不是为了你们的名，而是为了\Emph{我的}名。因此，因那名而行这些事的人是多么可怜，而因那名而受苦的人是多么有福：正如祂自己在别处所说：\BibleQuote{为义德而受迫害的人是有福的。}\BibleRef{玛 5:10} 因为那正是为了我的缘故，或是\Quote{因我的名}。因为正如宗徒教导我们的：\BibleQuote{基督成了我们的智慧、义德、圣化和救赎，正如经上所载：凡要夸耀的，应因主而夸耀。}\BibleRef{格前 1:30} 因为恶人对恶人做这些事，却不是为义德之故；因此两者同样可怜，行事的和受苦的。善人也对恶人做这些事：在那里，虽然前者是为义德而行，后者却并非因同一缘故而受苦。\par

3. 有人会说：如果当恶人为基督的名迫害善人时，善人是为义德而受苦，那么显然恶人对他们这样做也是为义德；如果是这样，那么当善人为义德而迫害恶人时，恶人同样也是为义德而受苦。因为如果恶人能因基督的名而迫害善人，为什么恶人不能因同一缘故而遭受善人的迫害？那岂不是为义德？因为如果善人行动的原因不同于恶人受苦的原因（善人是为义德而行，恶人却是为不义而受苦），那么同样，恶人行动的原因也不同于善人受苦的原因（恶人是由于不义而行，善人却为义德而受苦）。那么，\Quote{他们因我的名向你们施行这一切事}这句话怎么可能是真的呢？因为前者不是因基督的名（即义德）而行事，而是因他们自己的邪恶。这样的问题只有当我们理解\Quote{他们因我的名向你们施行这一切事}这句话完全是指义人而言时才能解决，如同说：你们将因我的名而受他们的苦，所以\Quote{他们向你们施行}等同于\Quote{你们受他们的苦}。但如果\Quote{因我的名}应理解为祂曾说过的\Quote{因我的名，就是他们在你们身上所恨的}，那么另一句话\Quote{为义德之故}也可以理解为\Quote{因他们在你们身上所恨的义德}；这样，当善人发动对恶人的迫害时，可以正确地说是既为义德（因爱义德而迫害恶人）也是为恶（因恨恶人自身的邪恶）；同样，恶人受苦可以说是既因惩罚于他们身上的不义，也因施行于惩罚中的义德。\par

4. 或可追问：若恶人也迫害恶人，正如不敬天主者的王侯与审判者，他们既是虔敬者的迫害者，当然也惩罚杀人犯、奸淫者以及所有被查明违反公共法律的各类作恶者，那么我们如何理解主的话：\BibleQuote{你们若是属于世界，世界必爱其所属}？\BibleRef{若 15:19} 因为世界所惩罚的，不能为世界所爱；但我们看到，世界通常惩罚上述各类罪行，而世界既存在于惩罚这些罪行的人中，也存在于爱这些罪行的人中。因此，那被理解为存在于恶人与不敬天主者之中的世界，一方面恨其所属（当它使罪犯受损害时），另一方面爱其所属（当它向自己的犯罪同伙施恩时）。因此，\BibleQuote{他们因我的名要向你们行这一切}这句话，或是指你们受苦的缘故，或是指他们如此对待你们的缘故，因为他们既恨又迫害你们内的那一位。祂又说：\BibleQuote{因为他们不认识那派遣我来的。}这句话应理解为指那种认识，关于这种认识，别处也记载着：\BibleQuote{认识祢，就是完美的智慧。}\BibleRef{智 6:16} 因为凡以这种认识认识圣父（基督被祂所派遣）的人，绝不可能迫害那些基督正在聚集的人；因为他们自己也与其他人一同被基督聚集。\par

\SourceRecord{Translated by John Gibb.；Nicene and Post-Nicene Fathers, First Series；Vol. 7.；Edited by Philip Schaff.；Buffalo, NY: Christian Literature Publishing Co.,；1888.；<http://www.newadvent.org/fathers/1701088.htm>.}

% structure: p0001 source=h1 target=section reason=page-title

\section{讲道集第八十九篇（若望福音第十五章22-23节）}

1. 主上面对祂的门徒说过：\BibleQuote{如果他们迫害了我，也要迫害你们；如果他们遵守了我的话，也要遵守你们的话。但这一切他们都要为了我的名而加给你们，因为他们不认识那派遣我来的。}如果我们问祂这是对谁说的，我们就会发现祂是从之前的话引到这些来的：\BibleQuote{如果世界憎恨你们，你们要知道，它在我以先已经憎恨了你们；}如今又加上：\BibleQuote{如果我没有来，也没有对他们说话，他们就没有罪。}祂更明确地是指犹太人说的。因此，前面的话也是指着他们说的，因为上下文本身就暗示了这一点。因为祂说：\BibleQuote{如果我没有来，也没有对他们说话，他们就没有罪；}所指的正是同一批人；祂也对他们说了：\BibleQuote{如果他们迫害了我，也要迫害你们；如果他们遵守了我的话，也要遵守你们的话；但这一切他们都要为了我的名而加给你们，因为他们不认识那派遣我来的；}因为祂紧接着也说了以下的话：\BibleQuote{如果我没有来，也没有对他们说话，他们就没有罪。}因此，犹太人迫害了基督，正如福音非常清楚地指出的，而基督是对犹太人说话，而不是对其他民族；所以他们正是祂所指的、那憎恨基督和祂门徒的\Quote{世界}；而且，确实，不仅他们，就连门徒也被祂显示为属于同一个世界。那么，祂说\BibleQuote{如果我没有来，也没有对他们说话，他们就没有罪}是什么意思？难道是说，在基督以血肉之身来到他们中间之前，犹太人没有罪吗？即使是最愚蠢的人，谁会这样说呢？但祂要我们理解的，乃是一种大罪，而非所有罪，仿佛是用一个总称来涵盖。因为这一罪包含了所有其他的罪；谁脱离了它，他的一切罪都获得了赦免：这就是他们不信基督——基督来正是为了赢得他们的信德。如果祂没有来，他们当然不会犯这一罪。祂的来临对于不信者而言成为灭亡，正如对于信者成为救恩；因为祂，宗徒的元首和首领，自己也仿佛成了他们自己所宣称的：\BibleQuote{对一些人，却是生命的芬芳，使人活；对另一些人，却是死亡的芬芳，使人死。}\BibleRef{格后 2:16}\par

2. 但祂接着说：\BibleQuote{但现在他们对自己的罪再没有借口了。}有人可能会被触动而追问：那些基督既没有来也没有对他们说话的人，对自己的罪是否有借口？因为如果没有，那么这里说这些人没有借口，恰恰是因为祂确实来并对他们说话了，那又是为什么？如果他们有借口，他们能否因此免于惩罚，或者只是减轻惩罚？对于这些问题，在主帮助下，尽我所能，我回答：这样的人是有借口的，但并非对他们所有的罪，而是对不信基督这一罪，因为祂没有来也没有对他们说话。但那些祂借着祂的门徒而来并对他们说话的人，不应当被列在其中，正如祂现今所做的；因为祂借着祂的教会而来，并借着祂的教会对外邦人说话。因为祂所说的话应该引用到这里：\BibleQuote{谁接纳你们，就是接纳我；}\BibleRef{玛 10:40}和\BibleQuote{谁拒绝你们，就是拒绝我。}\BibleRef{路 10:16} 宗徒保禄说：\BibleQuote{或者，你们愿意得到一个证明，证明在我内说话的乃是基督吗？}\BibleRef{格后 13:3}\par

3. 还剩下一点需要我们探究：那些在基督借着祂的教会来到外邦人中间以及他们听到祂的福音之前，已经或正在被今生的终结所接去的人，能否有这样的借口？显然可以，但他们并不能因此逃脱惩罚。\BibleQuote{因为凡没有法律而犯了罪的，也必不按法律而丧亡；凡在法律之下犯了罪的，就必按法律受审判。}\BibleRef{罗 2:12} 宗徒的这些话，由于他说\BibleQuote{他们必丧亡}比他说\BibleQuote{他们必受审判}听起来更可怕，似乎表明这样的借口不仅对他们毫无益处，反而成了进一步的加重。因为那些以没有听见为借口的人，\BibleQuote{必不按法律而丧亡}。\par

4. 但这也是一个值得探究的问题：那些以轻蔑甚至反对来回应所听言辞的人——不仅通过驳斥，更通过仇恨而逼迫那些使他们听到这些言辞的人——是否应被算在那些\Quote{他们将受法律审判，}这句话含有较温和意味的人当中？但如果没有法律而丧亡与受法律审判是两回事，前者较严厉，后者较轻微；那么，这类人无疑不会被置于那较轻的惩罚尺度中；因为他们不仅不在法律内犯罪，反而完全拒绝接受基督的法律，并尽其所能地企图将其彻底消灭。而那些在法律内犯罪的人，则是那些在法律之中的人，也就是说，他们接受法律，承认法律是圣洁的，诫命是圣洁、公义和良善的（见：\BibleRef{罗 7:12}），但因软弱而未能履行他们毫不怀疑其中所命令的极其正义之事。这些人的命运或许与那些没有法律之人的丧亡有所不同；然而，如果宗徒的话\Quote{他们将受法律审判，}可以被理解为意味着他们不会丧亡，那该是多么令人惊奇！因为他的论述并非关于不信者和信者而说出这样的话，而是关于外邦人和犹太人；诚然，如果他们在那位来寻找迷失之物的救主（见：\BibleRef{路 19:10}）中找不到救恩，他们无疑将成为丧亡的猎物；尽管可以说，有些人的丧亡更可怕，有些人的丧亡更缓和；换言之，有些人将在他们的丧亡中承受更重的惩罚，有些人则承受更轻的惩罚。因为凡因受罚而与天主赐予圣者的真福分离的人，都可说是在天主面前丧亡；惩罚的多样性如同罪恶的多样性一样多；但其方式在天主智慧中过于深奥，非人力猜测所能探究或表达。无论如何，基督所来到并向他们说话的那些人，对于他们不信的大罪，没有任何借口可以说：他们没有看见，他们没有听见；无论是这种借口不会被他那不可测度的审判所支持，还是会被支持，即使不是为了他们完全脱离惩罚，至少是为了部分减轻惩罚。\par

5. \Quote{恨我的，}祂说，\Quote{也恨我的父。}这里可能会有人对我们说：谁能够恨一个他所不认识的人呢？而确实，在说\Quote{我若没有来向他们说话，他们就没有罪，}之前，祂曾对祂的门徒说：\Quote{这些人要向你们行这些事，因为他们不认识那派遣我来的。}那么，他们又如何既\Quote{不认识}又\Quote{恨}呢？因为如果他们对于祂所形成的概念不是祂本身，而是某种他们自己的未知臆测，那么显然他们并非恨祂本身，而是恨他们因错误而臆造或更确切说是猜疑的那个虚构之物。然而，如果人不能恨他们所不认识的东西，那么真理就不会断言这两点，即他们既\Quote{不认识}又\Quote{恨}祂的父。但是这种可能性，即使我们赖主的助佑能够指明它，现时也无法证明，因为这篇讲论现在必须结束。\par

\SourceRecord{Translated by John Gibb.；Nicene and Post-Nicene Fathers, First Series；Vol. 7.；Edited by Philip Schaff.；Buffalo, NY: Christian Literature Publishing Co.,；1888.；<http://www.newadvent.org/fathers/1701089.htm>.}

% structure: p0001 source=h1 target=section reason=page-title

\section{讲道集第90篇（\BibleRef{若 15:23}）}

1. 主说，正如你们刚才所听到的：\BibleQuote{那恨我的，也恨我的圣父。}但祂先前又说：\BibleQuote{他们向你们做这些事，是因为不认识那派遣我来的。}因此产生一个不可忽略的问题：他们如何能恨一个不认识的人呢？因为，如果他们不是按天主真实的样子，而是按他们所猜疑或相信的、我不知为何物的某个东西来恨祂，那么他们恨的确实不是天主本身，而是他们从自己错误的猜疑或无凭的信心里所虚构的东西；而如果他们想祂正如祂所是的，又怎能说他们不认识祂呢？固然，对于人来说，我们常常爱那些从未见过的人；另一方面，我们恨那些从未见过的人，也并非不可能。例如，关于某位宣讲者的好或坏的消息，自然地引导我们去爱或恨那位未知者。但如果消息是真实的，那么怎能说一个我们对其有如此真实叙述的人是不认识的呢？是因为我们没有见过他的脸吗？然而，虽然他自己看不到自己的脸，但没有人能比他自己更认识他。因此，对任何人的认识，并不是通过他的面容传达给我们，而只有在他的生活和品格被揭示时，才向我们敞开。否则，没有人能够认识自己，因为无法看见自己的脸。但他肯定比他人对他的认识更确定地认识自己，因为通过内在的审视，他更能确定地看到自己所意识到的、所渴望的、为活着的目的；而正是当这些也同样向我们揭示时，他才真正为我们所认识。相应地，这些关于缺席者甚至死者的情况，通常通过传闻或通信传达给我们，因此，我们经常爱或恨那些从未谋面的人（尽管我们对他们并非完全无知）。\par

2. 但在这种情况下，我们的轻信常常出错；因为有时甚至历史，尤其是普通的传闻，结果是虚假的。然而，为避免被有害的意见误导，既然我们无法洞察人的良心，我们应当对事物本身有一种真实而确定的感受。我的意思是，关于这个或那个人，如果我们不知道他是无节制的还是有节制的，我们无论如何都应当恨恶无节制而喜爱节制；如果关于某个人，我们不知道他是不义的还是正义的，我们无论如何都应当喜爱正义而憎恶不义；不是我们错误地向自己幻想的那种，而是我们信从地按天主的真理所感知的，一种值得渴求，另一种应当回避；这样，当对事物本身我们确实渴求应当渴求的，彻底避免应当避免的时，我们对于他人隐藏于我们视线之外的实际状态有时（甚至常常）所怀的错误感受，可以得到宽恕。因为我认为这涉及到人的试探，没有它我们无法度过此世，因此宗徒说：\BibleQuote{不要让你们遇到超过人所能承受的试探。}（\BibleRef{格前 10:13}）因为有什么比不能洞察人的心更常见于人的呢？因此，不是去仔细审视其最深处，而是在大多数情况下猜疑其中很不同的东西？尽管在这些人类现实的黑暗区域（即他人内心的思想）中，我们无法澄清我们的猜疑，因为我们只是人，但我们应当约束我们的判断，即一切确定而固定的意见，不要在任何时候以前判断，直到主来到，照亮黑暗中的隐秘之事，揭示心中的计谋；那时每个人都要从天主那里得到称赞。（\BibleRef{格前 4:5}）因此，当我们对事物本身没有犯错误，以致在谴责恶和赞同善上与正义一致时，那么，如果在个人身上犯了错误，那么这种典型的人的试探就在可宽恕的范围内。\par

但在人心的这一切黑暗之中，发生了一件极其令人惊异和哀叹的事：我们视为不义的人，其实却是义人，我们不知不觉中爱着正义，却有时躲避他、转身离开他、阻止他接近我们、拒绝与他共同生活和居住；如果基于需要而施行管教，为拯救他人免受伤害或引导他本人归正，我们甚至以有益的严厉对待他；这样，我们如同对待恶人一样折磨一个善人，而我们在不知不觉中爱着他。例如，假设我们相信一个谦逊的人是不谦逊的，就会发生这种情况。因为毫无疑问，如果我爱一个谦逊的人，他本人就是我所爱的对象；因此我爱这个人本身，却不知道。如果我恨一个不谦逊的人，那我所恨的并不是他本人，因为他不是我所恨的对象；然而，对于我所爱的对象——我的心在爱慕谦逊中不断居住于他——我却在无知中加害于他，因为我在分辨德行与恶习上没有犯错，但在人心的幽暗中犯了错。因此，可能发生这样的情况：一个善人可能在不知不觉中恨一个善人，或者更确切地说，不知不觉地爱着他（因为他爱那个人，是在爱那善的方面；因为对方所是的，正是他所爱的）；而不知不觉中，他所恨的并非那人本身，而是他所假想的那个人。同样，一个不义的人也可能恨一个义人，而当他以为自己爱一个像他一样不义的人时，却不知不觉地爱着这个义人；然而，只要他相信他是不义的，他所爱的就不是那人本身，而是他所想象的那个人。对人如此，对天主也是如此。总之，如果问犹太人是否爱天主，他们会怎样回答？只能回答他们爱天主，而且不是故意说谎，而是因为他们错误地以为自己在爱天主。那些充满对真理本身的仇恨的人，怎能爱真理的圣父呢？因为他们不愿自己的行为受谴责，而谴责这种行为的正是真理的职责；所以他们恨真理，如同恨自己的惩罚，真理将惩罚判给这样的人。但他们不知道那谴责他们这类人的就是真理；因此他们恨他们所不知道的；既然恨它，当然也不能不恨那生它的那一位。这样，因为他们不认识那审判他们的真理，即由天主圣父所生的那一位；确信他们既不认识圣父本人，也恨他。可怜的人！因为他们愿意作恶，而否定那谴责恶人的真理。他们拒绝承认真理之所是，而他们本应拒绝自己之所是；这样，真理保持不变，他们自己得到改变，免得因真理的审判而堕入定罪。\par

\SourceRecord{Translated by John Gibb.；Nicene and Post-Nicene Fathers, First Series；Vol. 7.；Edited by Philip Schaff.；Buffalo, NY: Christian Literature Publishing Co.,；1888.；<http://www.newadvent.org/fathers/1701090.htm>.}

% structure: p0001 source=h1 target=section reason=page-title

\section{《讲道集》第九十一篇（若望福音 15:24-25）}

1. 主曾说：\BibleQuote{恨我的，也恨我的父。}因为恨真理的人，必定也恨那生于真理者；关于这一点，我们已按所赐的能力讲论过了。接着祂又加上我们现在要讨论的话：\BibleQuote{我若没有在他们中行过别人未曾行的事，他们就没有罪。}这指的是那大罪，就是祂先前所说的：\BibleQuote{我若没有来对他们说话，他们就没有罪。}他们的罪就是不信那说话并行事的那位。因为在祂这样对他们说话并在他们中行这些事之前，他们并非没有罪；但他们的这个罪，即不信祂，被特别提及，因为它实际上包含了所有其他的罪。因为如果他们免除了这个罪，信了祂，那么一切其他的罪也都得到了赦免。\par

2. 但是，当祂说了\BibleQuote{我若没有在他们中行这些事}后，立刻加上\BibleQuote{别人未曾行过}是什么意思呢？确实，在基督的一切事工中，没有比复活死人更大的了；然而我们知道，古代的先知也行过同样的事。例如，厄里亚曾这样做（\BibleRef{列上 17:21}）；厄里叟也是如此，生前（\BibleRef{列下 4:35}）以及埋葬在坟墓中时（因为有人在躲避敌人袭击时，抬着一个死人逃到那里，将他放在墓中，死人立刻复活了（\BibleRef{列下 13:21}））。然而，基督行过一些别人未曾行的事：例如，用五个饼喂饱五千人，用七个饼喂饱四千人；在水面上行走，并赐给伯多禄同样的能力（\BibleRef{玛 14:25}）；把水变成酒（\BibleRef{若 2:9}）；开了一个生来瞎眼的人的眼（\BibleRef{若 9:7}）；还有许多别的事，一时说不完。但是有人回答说，别人也行过祂没有行过的事，并且是别人未曾行过的。因为除了梅瑟，还有谁用那么多、那么大的灾祸打击埃及人（\EditorialNote{出谷纪七至十二章}），当祂带领百姓走过分开的海水时（\BibleRef{出 14:21}），当祂在饥饿中从天上为他们取得玛纳（\EditorialNote{出谷纪第十六章}），并在干渴中从磐石中取得水呢？（\BibleRef{出 17:6}）除了农的儿子若苏厄，还有谁分开约但河的水让百姓过去（\EditorialNote{若苏厄书第三章}），并藉着向天主的祈祷使运转的太阳停止呢？（\BibleRef{苏 10:12}）除了三松，还有谁从死驴的腮骨流出的水中解了渴？（\BibleRef{民 15:19}）除了厄里亚，还有谁乘着火马车升天呢？（\BibleRef{列下 2:11}）除了厄里叟，正如我刚才提到的，还有谁在自己埋葬后使另一个死人复活呢？除了达尼尔，还有谁在饥饿狮子口中安然无恙地活着，与它们一同被关？（\BibleRef{达 6:22}）除了阿纳尼雅、阿匝黎雅和米沙耳三人，还有谁在炽烈而不燃烧的火焰中安然行走呢？（\BibleRef{达 3:23}）\par

3. 我略过其他例子，认为这些足以表明有些圣人行过别人未曾行过的奇妙事工。但我们读到，古人中没有一个人能以如此的大能治好如此多的身体缺陷、疾病和人的痛苦。因为，暂且不提祂用命令治好那些个别病例，福音史家马尔谷在某处说：\BibleQuote{到了晚上，日落之后，人把所有患病的和附魔的都带到祂跟前，全城的人都聚集在门前。祂治好了许多患各种病症的人，并驱逐了许多魔鬼。}（\BibleRef{谷 1:32}）玛窦在记载同一件事时，加上了先知的话，说：\BibleQuote{这应验了依撒意亚先知所说的话：\InnerQuote{祂承担了我们的懦弱，背负了我们的疾病。}}（\BibleRef{玛 8:17}）在另一处，马尔谷又说：\BibleQuote{凡祂所到的地方，或村庄，或城市，或乡间，人都把患病的人放在街市上，求祂容许他们至少摸摸祂的衣边；凡摸到祂的，都痊愈了。}（\BibleRef{谷 6:56}）别人未曾在他们中行过这样的事。因为\Quote{在他们中}一词应当理解为不仅是在他们中间或在他们面前，而是\Quote{在他们身上}，因为祂治愈了他们。因为祂愿意他们明白这些事工不仅引起惊叹，而且赋予明显的治愈，是恩惠，他们理应以爱而非以恨来回报。确实，祂的事工超越所有其他人的，因为祂由童贞女所生，并在怀孕和降生中独有能力保持祂母亲的无玷完整；但那件事既没有在他们眼前发生，也没有在他们身上发生。因为对这一奥迹之真理的认识，是宗徒们通过他们单独的门徒训练获得的，而不是通过与他人共同的观察。此外，祂在第三天将自己从坟墓中复活，带着曾被杀害的肉身，此后不再死亡，并升天，这超越祂所做的一切；但这件事同样也没有在犹太人身上或在他们眼前发生；当祂说\BibleQuote{我若没有在他们中行过别人未曾行过的事}时，这事尚未发生。\par

4. 那么，这些工作无疑就是祂为治疗他们身体上的各种疾病所行的奇蹟，其数量之多，以前没有人在他们中间行过。因为他们看见了这些，而主正是在以此责备他们时继续说：\Quote{如今他们却看见了，也恨了我和我的圣父；但这正应验了他们法律上所记载的话：\BibleQuote{他们无缘无故地恨了我。}} 祂称之为\Quote{他们的法律}，不是指由他们自创，而是赐给他们的；正如我们说的\Quote{我们的日用粮}，然而我们却是在祈求天主时加上\Quote{求祢赐给我们}这句话 \BibleRef{玛 6:11}。一个人\OriginalTerm{无缘无故地}{gratuitously}恨人，是指他既不想从仇恨中得到好处，也不想避免不便。恶人就是这样恨主的；同样，义人也是\OriginalTerm{无缘无故地}{gratis, freely,}爱主，因为他们除了主自身之外不期待别的恩赐，因为祂将成为万有之中的万有。但若有人愿意在基督的话中寻找更深奥的意义：\Quote{我若没有在他们中间行过别人没有行过的工作}（因为虽然这些工作是由圣父或圣神行的，但别人没有行过，因为整个天主圣三在本质上是同一且相同的），他会发现，即使是某位天主的人行了类似的事，仍然是祂在行动。因为祂自己能凭自己行一切事；但离了祂，没有人能行任何事。因为基督与圣父及圣神不是三位神，而是一位天主，正如经上所载：\BibleQuote{愿以色列的上主天主受赞颂，唯有祂行伟大奇事。} 因此，没有别人真正凭自己行了祂在他们中间所行的工作；因为任何其他行了此类工作的人，都只是借着祂的行动而行。而祂自己却是在他们毫无行动的情况下行了这些工作。\par

\SourceRecord{Translated by John Gibb.；Nicene and Post-Nicene Fathers, First Series；Vol. 7.；Edited by Philip Schaff.；Buffalo, NY: Christian Literature Publishing Co.,；1888.；<http://www.newadvent.org/fathers/1701091.htm>.}

% structure: p0001 source=h1 target=section reason=page-title

\section{《讲道集第九十二篇（若 15:26-27）》}

1. 主耶稣在晚餐后向门徒们讲论时，自己已临近受难，仿佛即将离去，虽要使他们失去祂肉体的同在与陪伴，却仍以精神的方式临在于所有门徒中，直到世界终结。祂劝勉他们忍受恶人的迫害，祂以\Quote{世界}之名区分了恶人，并告诉他们，祂正是从这世界中拣选了门徒们，好让他们明白，他们之所以成为现在这样，是出于天主的恩宠，而他们过去之所以是那样，则是由于自身的恶习。然后，祂清楚指明迫害祂和他们的就是犹太人，好使人们完全明白，那些迫害圣徒的可咒骂的世界也包括他们。当祂论及他们不认识那派遣祂者，并且既恨圣子也恨圣父——即既恨那被派遣者，也恨那派遣祂者——关于这一切我们在先前几篇讲道中已经讨论过了——便来到了这段话：\Quote{这是为应验他们法律上所记载的话：\InnerQuote{他们无故地恨了我。}} 然后，祂仿佛顺理成章地加上了我们现下要讨论的话：\Quote{但当护慰者，就是我从父那里要给你们派遣的，那发自父的真理之神来到时，祂必要为我作证；你们也要作证，因为你们从起初就与我同在。} 但这与祂刚说的话有何关联？祂说：\Quote{但现在他们连我和我的父都看见了，也都恨了；但这是为应验他们法律上所记载的话：\InnerQuote{他们无故地恨了我。}} 难道是为了使护慰者，即真理之神，在来临时，以更清晰的见证定他们的罪吗？是的，确实，甚至那些看见了却仍然恨祂的人中，也有一些因着祂这一显现而皈依了那藉爱德工作的信德。\BibleRef{迦 5:6} 为使你们理解这段经文的意思，我们提醒你们，事实确是如此。因为在五旬节那天，圣神降臨在一百二十人的集会上，其中包含所有宗徒；他们充满圣神，用各种语言说话时，许多原本恨祂的人，因着这伟大奇迹而惊愕（尤其是当他们从伯多禄的讲辞中，听到为基督所作的如此伟大而神圣的见证，即那位被他们杀害、被视为死者的人，已被证明复活且活着），便感到心中刺痛而悔改；于是他们认识到那被他们不敬而残忍地流出的宝贵血是有益的，因为他们自己也被自己所流出的血赎回了。\BibleRef{宗 2:2} 因为基督的血为了赦免一切罪而有效地流出，以至于它甚至能抹去流这血本身的罪。因此，主心中念及此事，便说：\Quote{他们无故地恨了我；但当护慰者来到时，祂必要为我作证；} 这就是说：他们恨我，当我在他们眼前可见地站立时，他们杀了我；但护慰者为我所作的见证将是如此有力，以致当我从他们眼前消失时，祂会引领他们信从我。\par

2. \Quote{你们也要，}祂说，\Quote{作证，因为你们从起初就与我同在。}圣神将作证，你们也将如此。因为，正因为你们从起初就与我同在，你们能宣讲你们所知道的；但你们目前还不能这样做，因为那圣神的圆满尚未临于你们内。\Quote{祂因此要为我作证，你们也要作证：}因为天主的爱，借着将要赐予你们的圣神倾注在我们心中，\BibleRef{罗 5:5}将赐给你们为这样的作证所需的勇气。那勇气确实尚缺乏于伯多禄身上，当他因一位侍女的问题而惊惶失措时，他无法做出真实的见证；反而违背自己的承诺，被极大的恐惧所驱迫，三次否认了祂。\BibleRef{玛 26:69}但在爱中没有这样的恐惧，因为圆满的爱驱逐恐惧。\BibleRef{若一 4:18}总之，在主受难之前，他的奴性恐惧被一个女仆质问；但在主复活之后，他的自由之爱却被自由的主亲自质问：\BibleRef{若 21:15}因此在前一场合，他受困扰；在另一场合，他得安息；在那里他否认了他所爱的，在这里他爱了他所否认的。但即使那时，那爱仍然微弱而狭窄，直到被圣神加强并扩展。那时那圣神，以更丰富恩宠的圆满充满他，点燃了他之前冰冷的心，使他为基督作证，并解开了那些因之前颤抖而压抑真理的嘴唇，以至于当所有领受圣神的人都在说万国的方言，向周围聚集的犹太人宣讲时，他独自在众人之前，因急于为基督作证而发言，以祂的复活之事叫那些杀害祂的人惊惶失措。若有人愿享受凝视如此神圣魅力的景象，就让他读\WorkTitle{宗徒大事录}：\EditorialNote{宗二至五}那里，他将充满惊讶地看着蒙福的伯多禄的宣讲，他刚才还在哀悼伯多禄对主的否认；那里让他看到那舌头，本身从胆怯变为勇敢，从束缚变为自由，使如此多祂的敌人的舌头改而宣认基督，而这些敌人中没有一个是他能忍受的，当他陷入否认时。我还要说什么呢？在他身上闪耀出如此恩宠的光辉，如此圣神的圆满，如此宝贵的真理从宣讲者口中流出，以至于他将那大群曾是基督仇敌和谋杀者的犹太人转变为甘愿为祂的名而死的人，而他自己从前却害怕与祂的主同死。这一切都是圣神被派遣后所作的，而祂先前只是被预许的。主预见了自己的这些伟大奇妙的恩赐，因此祂说：\Quote{他们既看见了我，也看见了我的父，却都恨了我和我的父：这为应验他们法律上所记载的话：\InnerQuote{他们无故地恨了我。}但当护慰者，就是我从父那里要给你们派遣的，那发于父的真理之神来到时，祂要为我作证；并且你们也要作证。}因为祂，亲自作证，并激励这样的见证人具有不可战胜的勇气，使基督的朋友们卸下恐惧，并将祂仇敌的恨转变为爱。\par

\SourceRecord{Translated by John Gibb.；Nicene and Post-Nicene Fathers, First Series；Vol. 7.；Edited by Philip Schaff.；Buffalo, NY: Christian Literature Publishing Co.,；1888.；<http://www.newadvent.org/fathers/1701092.htm>.}

% structure: p0001 source=h1 target=section reason=page-title

\section{讲道集 93（\BibleRef{若 16:1-4}）}

1. 主在前面的话中，坚固了祂的门徒，使他们能忍受敌人的仇恨，并用自己的榜样预备他们更勇敢地效法祂；并加上了应许：圣神将要来为祂作证，而且他们自己也因圣神在他们心中有效地工作，能成为祂的证人。因为当祂说：\BibleQuote{祂要为我作证，你们也要作证。}意思就是这样：因为祂要作证，你们也要作证；祂在你们心中，你们在你们言语上；祂藉感动，你们藉宣讲：好使这话得以应验：\BibleQuote{他们的声音传遍了普世。}因为若不用自己的榜样鼓励他们，而不以圣神充满他们，那是徒劳的。正如我们看见宗徒伯多禄，在听到祂的话，当祂说：\BibleQuote{仆人不能大于主人：他们若迫害了我，也要迫害你们；}并看见那话已在祂身上应验（如果榜样足够，他本该效法主的忍耐），却屈从并陷入否认，无法忍受他看见主所忍受的。但当真正领受了圣神的恩赐，他便宣讲了他曾否认的那位；他以前害怕承认的，现在却不再害怕公开宣告了。确实，他早已被榜样充分教训，知道该做什么；但还没有被赋予力量去做他所知道的：他得到了站立的教导，却没有得到不跌倒的力量。但圣神补充了这力量后，他宣讲基督甚至至死，而先前因怕死他否认了祂。因此，主在这接着的一章中，我们现在要向你们讲的，祂说：\BibleQuote{我给你们说了这些，是要你们不致跌倒。}正如圣咏所唱的：\BibleQuote{爱祢法律的人享有大和平，他们不会跌倒。}因此，在应许圣神（藉祂在心中的运行，他们将成为祂的证人）之后，祂恰当地加上：\BibleQuote{我给你们说了这些，是要你们不致跌倒。}因为天主的爱借着赐给我们的圣神，已倾注在我们心中，\BibleRef{罗 5:5} 爱天主法律的人享有大和平，以致什么也不能使他们跌倒。\par

2. 然后祂明确宣告他们要受什么苦：\BibleQuote{他们要把你们逐出会堂。}但宗徒们被逐出犹太人的会堂有什么害处呢？好像即使没有人驱逐他们，他们自己就不会离开会堂一样？无疑祂要谴责地宣告，犹太人将拒绝接受基督，而他们也必定不肯离开祂；于是那些不能没有祂而存在的人，也将与祂一同被那些不愿以祂为住所的人驱逐出去。因为当然，除了亚巴郎的苗裔，没有别的天主之民；如果他们承认并接受基督，他们就会作为自然的枝条留在橄榄树上；\BibleRef{罗 11:17} 基督的教会也不会与犹太人的会堂不同，它们本是同一的，如果他们也愿意住在祂内的话。但他们拒绝了，除了他们自己留在基督之外，把那些不肯放弃基督的人逐出会堂，还剩下什么呢？因为他们领受了圣神，成为祂的证人，当然不属于这类人，关于这类人经上说：\BibleQuote{犹太人的许多首领信了祂，但为了怕犹太人，不敢承认祂，免得被逐出会堂：因为他们爱人的光荣胜过爱天主的光荣。}所以他们信了祂，但不是以祂愿意他们相信的方式，当祂说：\BibleQuote{你们既然彼此求光荣，而不寻求唯独来自天主的光荣，怎么能信呢？}因此，那些门徒这样信祂，充满圣神，或者说充满天主恩宠的恩赐，就不再属于那些\BibleQuote{不知道天主的正义，企图建立自己的正义，而不服从天主的正义}的人；\BibleRef{罗 10:3} 也不属于那些经上所说\BibleQuote{他们爱人的光荣胜过爱天主的光荣}的人：这样，预言就与他们在自己身上所应验的和谐一致：\BibleQuote{上主，他们要在祢容颜的光中行走；他们因祢的名终日欢跃；因祢的正义被高举：因为祢是他们力量的光荣。}对这样的人恰当地说：\BibleQuote{他们要把你们逐出会堂；}就是说，那些\BibleQuote{对天主有热心，但不按着真知识}的人；因为\BibleQuote{不知道天主的正义，企图建立自己的正义}，\BibleRef{罗 10:2} 他们驱逐那些被高举的人，不是因自己的正义，而是因天主的正义，且不因被人驱逐而羞愧，因为祂是他们力量的光荣。\par

3. 最后，祂对所告诉的这一切又补充说：\BibleQuote{然而时候要到，凡杀害你们的，还以为是服事天主；他们做这些事，是因为不认识圣父，也不认识我。}这就是说，他们不认识圣父，也不认识祂的圣子；他们以为杀害你们是在服事祂。主加上这句话，作为对祂自己人的安慰，因为他们将被赶出犹太会堂。因为正是这样预先宣布他们为祂作见证将要忍受的苦难时，祂说：\BibleQuote{他们要把你们赶出会堂。}祂并没有说：然而时候要到，凡杀害你们的，还以为是服事天主。那么是什么呢？\BibleQuote{然而时候要到：}仿佛祂正在预言某种好事会跟随这样的苦难而来。那么，祂说\BibleQuote{他们要把你们赶出会堂：然而时候要到}是什么意思？就好像祂本来会继续说：他们确实会驱散你们，\Emph{但是}我将聚集你们；或，他们确实会驱散你们，\Emph{但是}你们喜乐的时候来到。那么，祂所用的词\BibleQuote{\Emph{然而}时候要到}在这里做什么，好像祂正要应许他们在苦难之后的安慰，而显然祂本应以连续叙述的形式说\Emph{并且}时候要到呢？但祂没有说并且时候要到，尽管预言了一种苦难接另一种苦难，而不是苦难之后的安慰。难道是因为被赶出会堂会如此扰乱他们，以至于他们宁愿死也不愿离开犹太会堂继续活在世上？那些寻求的并不是人的称赞，而是天主的称赞的人，当然不会这样沮丧。那么，\BibleQuote{他们要把你们赶出会堂：然而时候要到}这些词是什么意思？显然祂本应说\Emph{并且}时候要到，\BibleQuote{凡杀害你们的，还以为是服事天主}？因为甚至没有说然而时候要到，他们必杀害你们，仿佛他们因分离而得的安慰在于他们将要遭受的死亡；而是说\BibleQuote{时候要到}，祂说，\BibleQuote{凡杀害你们的，还以为是服事天主。}总的来说，我认为祂不想传达任何更深的意思，只是让他们明白并欢喜：他们自己被赶出犹太会堂，将赢得许多人归向基督，以至于驱逐他们变得不够，犹太人还不会让他们活着，恐怕他们所有的听众都会被他们的宣讲归向基督的名，从而离开犹太教义，仿佛那是天主的真理一样。因为我们应该这样理解祂的话在指犹太人时所说：\BibleQuote{他们要把你们赶出会堂。}因为基督的见证者，换言之，殉道者，也同样被外邦人杀害：但外邦人以为他们这样做不是在服事真天主，而是在服事他们自己的假神。但每一个杀害基督宣讲者的犹太人都以为自己是在服事天主；因为他相信凡归向基督的人都是在背弃以色列的天主。正是出于同样的理由，他们也被煽动去杀害基督自己：因为关于这件事他们自己的话也被记载了下来。看哪，全世界都跟祂去了：\BibleQuote{如果让祂活着，罗马人必来，夺去我们的土地和人民。}以及\Person{盖法}的话：\BibleQuote{叫一个人替百姓死，免得整个民族灭亡，为我们是有利的。}因此，在这篇讲论中，祂以身作则来激励门徒，祂刚才对他们说：\BibleQuote{如果他们迫害了我，也会迫害你们；}这样，正如杀害祂时他们以为服事了天主，对你们也会如此。\par

4. 这就是这些话语的意义：\BibleQuote{他们会把你们逐出会堂；}但不要惧怕孤独：因为当你们与他们的聚会分离时，你们将以我的名字聚集如此众多的人，以至于他们极其担心那与他们同在的圣殿以及旧律法的一切圣事会被遗弃，便会杀死你们：实际上，他们如此流你们的血，满以为是在服事天主。这无疑正是宗徒话语的例证：\BibleQuote{他们对天主有热心，但不是按着真知；}\BibleRef{罗 10:2}因为他们以为在杀害天主的仆人时是在服事天主。可怕的错误！难道你们竟要藉击打那中悦天主的人来中悦天主吗？难道你们要将天主的活殿宇击倒在地，好使天主的石殿不致被遗弃吗？可咒的盲目！但这只是部分地临到以色列，好使外邦人的丰盈得以进来：我说是部分地，而非整体地临到。因为并非所有的枝子都被折下，只有一些，好让野橄榄得以接上。因为正当基督的门徒充满圣神，说着万国的方言，行了许多神圣的奇迹，并将神圣的话语四散传播时，基督——即使被杀害——仍如此被爱，以至于祂的门徒在被逐出犹太人的会堂后，却为自己聚集了一大群犹太人，并且毫不惧怕被遗弃在孤独中。\EditorialNote{宗二至四章}于是，其余那些被弃绝而盲目的人，被愤怒所点燃，对天主有热心，但不是按着真知，并相信自己是在服事天主，便杀害了他们。但那为他们而被杀的一位，却聚集了他们；正如祂在被杀之前，也曾教诲他们将要发生的事，免得他们的心思因无知和毫无准备，被那些虽短暂却出乎意料且未蒙预备的凶恶投入困苦；反而因预先知道并耐心承受，得以被引领至永远的福乐。因为这就是祂预先告诉他们这些事的原因，这由祂随后的话语所表明：\BibleQuote{但你们已经听我告诉了你们这事，好使时候到了，你们可以想起我告诉过你们了。}他们的时辰是黑暗的时辰，是午夜时分。但主在白天命令了祂的慈爱，使他们在夜间歌颂它：当犹太人的黑夜并未将任何混乱的黑暗投入基督徒的白天之中——这白天已与他们分离；而那能杀害肉身的，并没有力量使他们的信德变为黑暗。\par

\SourceRecord{Translated by John Gibb.；Nicene and Post-Nicene Fathers, First Series；Vol. 7.；Edited by Philip Schaff.；Buffalo, NY: Christian Literature Publishing Co.,；1888.；<http://www.newadvent.org/fathers/1701093.htm>.}

% structure: p0001 source=h1 target=section reason=page-title

\section{讲道集94（\BibleRef{若 16:44-47}）}

1. 主耶稣预告了祂的门徒，在祂离去后将要遭受的迫害，接着说道：\BibleQuote{这些事，我起初没有告诉你们，因为我与你们同在；但现在我往派遣我者那里去。}这里我们首先要看的是，祂是否以前没有预告过他们将来的苦难。其他三位福音作者已经足够清楚地表明，祂在逾越节晚餐前就说过这样的预言；按照\BibleBook{若望}{福音}的记载，晚餐结束时祂说了这些话，并补充说：\BibleQuote{这些事，我起初没有告诉你们，因为我与你们同在。}那么，我们是否应该这样解决这个问题：他们也告诉我们，祂说这些话时已临近祂的苦难？那么，祂起初与他们同在时并没有这样说，因为祂正即将离去，到圣父那里去；所以，即使根据这些福音作者的记载，这里所说的话也是严格真实的：\BibleQuote{这些事，我起初没有告诉你们。}但我们对\BibleBook{玛窦}{福音}的可信性又该如何处理呢？他记载了主不仅即将与门徒同坐逾越节晚餐时，而且在起初，当十二位宗徒第一次被提名并受派遣去从事天主的工作时，就向他们宣布了这些事情。\BibleRef{玛 10:17}那么，祂在这里所说的：\BibleQuote{这些事，我起初没有告诉你们，因为我与你们同在}究竟是什么意思？岂不是指祂在这里所说的关于圣神将要降临到他们身上，当他们要忍受这样的苦难时为他们作见证，这件事祂起初没有告诉他们，因为祂自己与他们同在？\par

2. 那时，护慰者，或辩护者（因为两者都是对希腊词\OriginalTerm{护慰者}{\GreekText{Παράκλητος}}的翻译），在基督离去后成为必要：因此，祂起初与他们同在时，并没有提到祂，因为祂自己的同在就是他们的安慰；但在祂即将离去之际，祂必须提到祂的到来，藉着祂，爱被倾注在他们心中，使他们能以极大的胆量宣讲天主的圣言；并且，祂内在地与他们一同为基督作见证，他们也要作见证，并且当他们的犹太仇敌将他们赶出会堂并杀害他们，以为是在事奉天主时，他们不会感到绊倒；因为爱德含忍一切，\BibleRef{格前 13:7}这爱德将藉着圣神的恩赐倾注在他们心中。\BibleRef{罗 5:5}因此，全部意义就在于此：祂要使他们成为祂的殉道者，即通过圣神成为祂的见证人；这样，藉着祂在他们内有效的运作，他们要忍受各种迫害的困苦，并因那神圣之火而炽热，在宣讲的爱中不失去任何温暖。因此祂说：\BibleQuote{这些事我已经告诉你们了，为的是当时候来到时，你们可以想起我告诉了你们。}\BibleRef{若 16:4}这些事，我说，我告诉你们，不仅是因为你们将要忍受这样的事，而且因为当护慰者来到时，祂要为我作见证，使你们不致因恐惧而退缩，并且藉着祂，你们自己也能作见证。\BibleQuote{这些事，我起初没有告诉你们，因为我与你们同在，}我自己就是你们的安慰，通过我的肉身临在向你们的人性感官显示，而你们作为婴孩，能够理解。\par

3. 祂说：\BibleQuote{但现在我往派遣我者那里去；你们中没有人问我：你往哪里去？}祂的意思是，祂的离去将是这样的，没有人会问祂关于他们将会亲眼看见光天化日之下发生的事：因为在此之前，他们曾问祂往哪里去，并被告知祂要往他们那时不能去的地方。然而，现在祂许诺，祂将以这样一种方式离去，以致他们中没有人会问祂往哪里去。因为当祂从他们身边升天时，有一朵云彩接住了祂；对于祂升天的事，他们并没有口头询问，而是亲眼看见了。\BibleRef{宗 1:9}\par

4. \Quote{但我把这事告诉了你们，}祂接着说，\Quote{忧愁充满了你们的心。} 祂确实看到祂的话在他们心中产生的效果；因为他们还没有内在的属神的安慰，就是后来藉着圣神将要获得的，他们仍然在基督身上客观地看到的东西，他们害怕失去；而且因为他们毫不怀疑将要失去祂（祂的宣告总是真实的），他们人性的情感感到悲伤，因为他们对祂属血肉的看法将变得空白。但祂知道什么对他们最有益，因为那内在的眼界，就是圣神将要用来安慰他们的，无疑更为优越；不是将一个属人的身体带入观看者的身体，而是将自己注入相信者的心中。然后祂又说：\Quote{然而我告诉你们真理，我去为你们是有益的。因为如果我不去，\OriginalTerm{护慰者}{Comforter}就不会到你们这里来；但如果我去了，我就派遣祂到你们这里来。} 好像祂说：这个仆人的形象从你们那里被带走，对你们是有益的；作为确实成为血肉的圣言，我居住在你们中间；但我不愿意你们继续以属血肉的方式爱我，并且满足于这样的奶，总是渴望保持婴儿状态。\Quote{我去为你们是有益的：因为如果我不去，护慰者就不会到你们这里来。} 如果我不撤回我用以养育你们的温柔营养品，你们就不会获得对固体食物的敏锐品味；如果你们以属血肉的方式依附于肉体，你们就没有容纳圣神的地方。这是什么意思呢？\Quote{如果我不去，护慰者就不会到你们这里来；但如果我去了，我就派遣祂到你们这里来。} 难道是祂自己在这里就不能派遣祂吗？谁敢这样说呢？也不是说，祂所在的地方，另一位就离开了，或者说祂从圣父那里来，以至于祂没有与圣父同住。更进一步，即使在祂在地上居住的时候，祂怎能不能派遣祂呢？我们知道祂在受洗时临于祂并停留在祂身上；是的，更多，我们知道祂从未与祂分离？那么，\Quote{如果我不去，护慰者就不会到你们这里来}是什么意思呢？岂不是说，只要你们继续按属血肉的方式认识基督，你们就不能领受圣神吗？因此，一个已经分享圣神的人说：\Quote{固然我们曾按血肉认识基督，但现在不再这样认识祂了。} \BibleRef{格后 5:16} 因为现在，当被带到对成为血肉的圣言的属神认识中时，他连基督的本身肉体也不再以属血肉的方式认识。这无疑就是良善的师傅想要暗示的，当祂说：\Quote{如果我不去，护慰者就不会到你们这里来；但如果我去了，我就派遣祂到你们这里来。}\par

5. 但随着基督身体的离去，圣父、圣子和圣神都属神地临在他们中间。因为如果基督离开他们到这样的程度，以至于圣神要代替祂的位置，而不是与祂一起临在他们内，那么祂的应许又成了什么呢？祂说：\Quote{看，我常与你们同在，直到世界的终结；}\BibleRef{玛 28:20} 以及，我和圣父\Quote{要到他那里去，并要同他住在那里；}可见祂也应许要派遣圣神，使祂永远与他们同在。另一方面，就是这样，看到他们尚且处于现在的属血肉或属魂的状态，要成为属神的，他们无疑也还要以更全面的方式拥有圣父、圣子和圣神。但我们绝不可相信圣父没有圣子和圣神而临在，或圣父和圣子没有圣神而临在，或圣子没有圣父和圣神而临在，或圣神没有圣父和圣子而临在，或圣父和圣神没有圣子而临在；而是无论其中一位在哪里，那里也有天主圣三，唯一的天主。但在这里，必须以这样的方式提示天主圣三：虽然本质没有差异，但每一位各自的位格区别应当被注意；而那些有正确理解的人绝不会想象本性的分离。\par

6. 但随后的话：\Quote{当祂来了，就要指证世界关于罪、关于义、关于审判：关于罪，是因为他们不信我；关于义，是因为我往父那里去，你们将不再看见我；关于审判，是因为这世界的首脑已被审判。}\BibleRef{若 16:8} 好像罪就是单纯不信基督；义就是不再看见基督；审判就是这世界的首脑，即魔鬼，已被审判：这一切都非常晦涩，不能包括在当前的讲道中，以免简短反而增加晦涩；而必须推迟到另一个场合，愿主使我们能给予解释。\par

\SourceRecord{Translated by John Gibb.；Nicene and Post-Nicene Fathers, First Series；Vol. 7.；Edited by Philip Schaff.；Buffalo, NY: Christian Literature Publishing Co.,；1888.；<http://www.newadvent.org/fathers/1701094.htm>.}

% structure: p0001 source=h1 target=section reason=page-title

\section{第95篇讲道（若望福音16:8-11）}

1. 主在应许要派遣圣神时说：\Quote{当祂来到时，祂要指证世界关于罪、关于正义、关于审判。}这是什么意思？难道主耶稣基督没有指证世界关于罪吗？祂曾说：\Quote{我若没有来向他们说话，他们就没有罪；但现在他们对自己的罪没有推诿之辞了。}为使没有人敢说这话只适用于犹太人，不适用于世界，祂不是在另一处说过：\Quote{你们若是属于世界，世界必爱属自己的}吗？祂不是也曾指证世界关于正义吗？祂说：\Quote{正义的父啊！世界没有认识祢。}祂不是也曾指证世界关于审判吗？祂宣布要对左边的人说：\BibleQuote{你们到那为魔鬼和他的使者预备的永火里去吧！}\BibleRef{玛 25:41}在神圣福音中还有许多其他段落，记载基督指证世界这些事。那么，为何祂将此归于圣神，似乎这是圣神特有的特权呢？是不是因为基督只在犹太民族中讲话，所以祂似乎没有指证世界，因为只有实际听见指证者的人才算被指证；而圣神当时住在分散全世界的门徒内，应该被理解为指证了不是一个民族，而是世界？请看祂在升天前对他们说的话：\BibleQuote{父以自己的权柄所定的时候和日期，不是你们应当知道的。但当圣神降临于你们身上时，你们将充满圣神的能力；并在耶路撒冷、全犹太和撒玛黎雅，直到地极，作我的证人。}\BibleRef{宗 1:7}这无疑就是指证世界。但谁敢说圣神通过基督的门徒指证世界，而基督自己却不指证呢？因为宗徒喊道：\BibleQuote{你们既然要寻求基督在我内说话的凭证。}\BibleRef{格后 13:3}因此，凡圣神所指证的人，基督也同样指证。但我认为，因为圣神要将那爱倾注在我们心中\BibleRef{罗 5:5}，这爱驱散恐惧\BibleRef{若一 4:18}，而恐惧可能阻止他们敢于指证那个充满逼迫的世界，所以祂说：\Quote{祂要指证世界。}仿佛祂要说：祂要将爱倾注在你们心中，并由此驱逐恐惧，你们就有了指证的自由。然而，我们常说过：天主圣三的作为是不可分的；但圣三的位格需要逐一展示，使我们不仅不分开它们，也不将它们混淆，从而正确理解它们的合一与圣三。\par

2. 祂接着解释祂所说的\Quote{论到罪、论义德、论审判}。祂说：\Quote{论到罪，是因为他们不信我}。祂将这一罪放在其他罪之前，好像它是唯一的一种；因为只要这罪存留，所有其他的罪都被保留；而一旦这罪被除去，其他的罪都被赦免。祂又加上：\Quote{论义德，是因为我往圣父那里去，而你们将不再看见我}。我们首先必须考虑：如果有人因罪而受到合理的指责，他是否也可能因义德而受到合理的指责？因为如果罪人因其是罪人而应受指责，那么有谁会认为义人因其是义人而应受指责呢？当然不会。因为即使有时义人也受到指责，他之所以受到合理指责，是因为按照圣经，\Quote{世上没有不犯罪而常行善的义人}。因此，当义人受指责时，他是因罪受指责，而非因义德。因为在那段神圣的话语中，我们读到\Quote{不要过分自显为义}，所指责的不是智者的义德，而是妄自尊大者的骄傲。因此，那\Quote{过分自显为义}的人，正因这过分而成为不义。因为他宣称自己无罪，或者想象自己成为义人是出于自己意志的充足，而非天主的恩宠，这样他就成了过分自显为义；他并非因正义地生活而成为义人，而是因自以为是的想象而自我膨胀。那么，世界如何因义德受指责，如果不是因信徒的义德？因此，它因不信基督而被定罪；它因那些相信的人的义德而被定罪。因为与信徒的比较本身就是在指责不信者。这解释本身已充分表明。因为在解释祂所说的话时，祂加上：\Quote{论义德，是因为我往圣父那里去，而你们将不再看见我}。祂不是说\Quote{他们将不再看见我}，即那些祂曾论到罪时说\Quote{因为他们不信我}的人。在解释祂所谓的罪时，祂说到他们：\Quote{因为他们不信我}；但当解释祂所谓的义德（世界因此而受指控）时，祂转向那些祂正在对他们说话的人，说：\Quote{因为我往圣父那里去，而你们将不再看见我}。因此，世界是因自己的罪，但因他人的义德而受指控，正如黑暗被光揭露：宗徒说：\Quote{凡事受了责备，就被光显明出来}。\BibleRef{弗 5:13} 因为不信者的罪恶之严重性，不仅凭其本身显明，也通过相信者的良善显明。因为不信者通常叫嚷：\Quote{我们怎么能相信我们看不见的事？} 所以不信者的义德正需要这个定义：\Quote{因为我往圣父那里去，而你们将不再看见我}。因为那些没有看见却相信的人是有福的。因为那些看见基督的人中，受到称赞的对祂的信德不是因为他们相信他们所看见的，即人子，而是因为他们相信他们所没有看见的，即天主之子。但是，当祂的仆役形象也从他们眼前被撤去之后，那时这句话才真正得以实现：\Quote{义人因信德而生活}。因为根据\WorkTitle{希伯来书}的定义，\Quote{信德是所希望之事的担保，是未见之事的确证}。\par

3. 但\Quote{你们将不再看见我}应该如何理解？因为祂不是说\Quote{我往圣父那里去，而你们将看不见我}，好像是指一段看不见的时间，无论长短，但终有结束；而是说\Quote{你们将不再看见我}，仿佛预言他们将永远不再看见基督。这就是我们所谓的义德吗？即永远不再看见基督，却仍信祂？因为义人赖以生活的信德正是基于相信那暂时看不见的基督某一天将会看见。再者，关于这义德，我们是否可以说宗徒\Person{保禄}不义，因为他承认自己见到基督升天之后\BibleRef{格前 15:8}（这无疑是在主说\Quote{你们将不再看见我}之后），保禄说\Quote{末了也显现给我}？\Person{斯德望}，这位卓越非凡的英雄，难道在这义德的精神下不是义的吗？当他被石头砸时，他喊道：\BibleQuote{看，我看见天开了，并看见人子站在天主的右边}\BibleRef{宗 7:56}。那么，\Quote{我往圣父那里去，而你们将不再看见我}的意思，无非就是：就我现在与你们同在的样式而言？因为那时祂仍是有罪肉身的相似形体的必朽之人\BibleRef{罗 8:3}。祂能忍受饥渴、疲倦、睡眠；这样的基督，即处于这种状况的基督，在祂从这世界前往圣父那里去之后，他们将不再看见；而这样的，也是信德的义德，宗徒说：\Quote{纵然我们曾按肉体认识基督，但如今不再这样认识祂了}\BibleRef{格后 5:16}。祂说：\Quote{这就是你们的义德，世界将因此受责备，因为我往圣父那里去，而你们将不再看见我：因为你们将相信一位你们看不见的我；而当你们看见那时我将成为的我时，你们将不会看见我现在与你们同在的样式；你们不会看见我的谦卑，而是我的高举；不会看见我的可朽性，而是我的永恒；不会看见审判台，而是审判宝座。借着你们的这信德，也就是你们的义德，圣神将责备不信的世界}。\par

4. 祂也要叫世界\BibleQuote{为审判受责备，因为今世之王受到了审判。}这今世之王是谁？不就是祂在另一处所说的那位吗？\BibleQuote{看，今世之王来到，在我内一无所获；}也就是说，在他的权下、属于他的什么也没有；事实上，丝毫罪也没有。因为魔鬼正是因此成了今世之王。因为魔鬼为王，并非指由天和地以及其中万有所构成的、当经上所说\BibleQuote{世界是藉着祂造的}时所意指的那个世界；而是指那同一个段落中祂紧接着添上的话\BibleQuote{世界却不认识祂}所意指的世界，即不信的人，整个世界极广之处都充满了他们；信者的世界就在他们中间叹息，这信者的世界是造世界的那位从世界中拣选出来的；对于这信者的世界，祂亲自说：\BibleQuote{人子来不是为审判世界，而是为叫世界藉着祂而得救。}祂是审判者，世界藉祂而被定罪；祂是救助者，世界藉祂而得救：因为正如一棵树长满叶和果，或一块田长满糠秕和麦子，同样，世界也充满了信者和不信者。因此，今世之王，即这世界的黑暗之王，或不信者之王，那世界（即信者的世界）从他手中被救出；对这世界曾说过：\BibleQuote{你们从前原是黑暗，但现在你们在主内却是光明；}\BibleRef{弗 5:8}今世之王，即祂在另一处所说\BibleQuote{现在这今世之王要被赶出去}的那一位，确实受到了审判，因为他已无可挽回地被判给永火之罚。因此，对于这审判——即今世之王所受的审判——世界也要被圣神责备；因为世界与它的王一同受审判，它效法了王的骄傲和不敬。\BibleQuote{因为如果天主，没有宽免犯罪的天使，反而把他们投入地狱的黑暗深渊，将他们交出来，拘留着等候审判的惩罚；}\BibleRef{伯后 2:4}那么，世界如何不被圣神为这审判而责备呢？因为宗徒如此说话，正是在圣神内。因此，但愿人们相信基督，免得他们因自己不信的罪而被定罪，这不信的罪使一切罪得以存留；但愿他们加入信者的行列，免得他们因那些成义者的义德而被定罪，因为他们未能效法这些成义者；但愿他们提防那将来的审判，免得他们与今世之王一同受审判，他们继续效法那已经受了审判的王。因为凡人的倔强骄傲，决不会想到自己可获宽免，当它这样被召唤去恐惧地思想那降临在天使骄傲身上的惩罚。\par

\SourceRecord{Translated by John Gibb.；Nicene and Post-Nicene Fathers, First Series；Vol. 7.；Edited by Philip Schaff.；Buffalo, NY: Christian Literature Publishing Co.,；1888.；<http://www.newadvent.org/fathers/1701095.htm>.}

% structure: p0001 source=h1 target=section reason=page-title

\section{讲道第九十六篇（\BibleRef{若 16:12-13}）}

1. 在圣福音的这一段中，主对祂的门徒说：\Quote{我还有许多事要告诉你们，然而你们现在还不能承担。}我们首先遇到一个必须探究的问题：祂不久前曾说：\Quote{凡我从圣父所听见的，我都已使你们知道了。}然而此处却说：\Quote{我还有许多事要告诉你们，但你们现在还不能承担。}但祂如何将尚未成就的事说得如同已经成就，正如先知见证天主使那尚未发生的事成为事实，当祂说：\Quote{祂使那尚未发生的事成为事实}；我们在处理那几句话时，已尽力作了解释。然而现在，你们或许想知道那些宗徒当时不能承担的事究竟是什么。但谁又敢断言自己现在有能力接受他们当时无力领受的事呢？因此，你们既不要期望我告诉你们那些连我自己也可能无法理解的事（即使别人告诉我），也不要期望我有足够的天赋让你们听到超乎你们理解力的事。的确，你们中可能有人已经足够领悟别人尚不能把握的事；即使不是全部，至少也可能是那些神圣的师傅所说\Quote{我还有许多事要告诉你们}中的一部分。但祂自己省略而未告诉他们的那些事，即使有愿望去猜测也是鲁莽的。因为那时宗徒们甚至还不适合为基督而死，当祂对他们说：\Quote{你们现在不能跟随我}时，他们中最杰出的伯多禄曾傲慢地宣称自己已经能够，却遭遇了出乎意料的结果。然而后来，无数男女、孩童、青年、老人，都获得了殉道的荣冠；羊群竟然能承担那些当主说这些话时牧人们尚且不能承担的事。那么，在那种极端的考验中，当要求他们为真理战斗至死，并为基督的名或教义流血时，难道应该问那些羊群：\Quote{你们中谁敢把自己算作准备好殉道的人？而伯多禄即使由主亲自面授时，尚且不适合殉道。}同样，可以说基督徒，即使渴望听讲，也不应被告知主当时所说\Quote{我还有许多事要告诉你们，但你们现在不能承担}的那些事是什么。如果宗徒们尚且不能承担，你们更甚；尽管可能现在许多人能承担伯多禄那时不能承担的事，正如许多人能够获得殉道的荣冠，而那时伯多禄尚不能；尤其现在圣神已被派遣，而那时尚未派遣，紧接着祂又加上这些话：\Quote{然而当祂，真理之神来到时，祂将教导你们一切真理。}这明确显示他们不能承担祂还要说的话，因为圣神尚未降临到他们身上。\par

2. 那么好，我们姑且承认如此：现在圣神已被派遣，许多人能承担那些圣神降临前门徒们不能承担的事。但我们因此就知道祂不愿说的是什么，如同我们读到或听到祂亲口说出时那样知道吗？因为知道我们或你们是否能承担是一回事，知道那件事是什么（不论是否能承担）则是另一回事。但当祂自己对此保持沉默时，我们中谁能说：\Quote{是这事或那事}呢？即使有人胆敢如此说，他又如何证明呢？因为谁能如此虚荣或鲁莽，在随意向人述说自己的意见时，即使说的是真理，却毫无神圣权威地断言这正是主当时拒绝说出的事呢？我们中谁能这样做而不犯下最严重的轻率之罪？这种事不可能从先知或宗徒的权威得到支持。因为如果我们曾在教会权威认可的书籍中读到任何此类事，而这些书是在主升天后写成的，那么仅仅读到这样的陈述是不够的，除非我们在同一处还读到这确实是主当时不愿告诉门徒、因为他们不能承担的那些事之一。例如，假如我说我们在本福音开篇读到的话：\BibleQuote{在起初已有圣言，圣言与天主同在，圣言就是天主；这圣言在起初就与天主同在。}以及后面的内容，因为这些内容后来写成，且未提及是主耶稣在肉身时所说，而是由祂的一位宗徒受圣神启示写成，这些正是主当时不愿说出、因为门徒不能承担的话之一，那么有谁会听我如此鲁莽的断言呢？但如果我们读到其中一处时也读到了另一处，谁又会不相信这样的宗徒呢？\par

3. 但我认为还有一种说法也十分荒谬，即：门徒们当时无法承担我们后来在宗徒书信中所读到的那一切关于不可见及最深邃事物的事迹，这些书信是后来写成的，并且从未提及主在他们可见地临在时曾宣讲过这些事。因为，如果现在每个人都读他们的书，甚至不理解也能承担，那么他们当时为什么不能承担呢？确实，在圣经中有一些事物，不信的人当阅读或听到时不仅不能理解，而且也无法承担：例如外邦人无法承担世界是由被钉十字架的那一位所创造的；犹太人无法承担祂能是天主子，并打破了他们守安息日的方式；撒伯里乌派无法承担父、子和圣神是一个天主圣三；亚略派无法承担圣子与圣父平等，圣神与圣父和圣子平等；佛提诺派无法承担基督不仅是一个像我们一样的人，而且也是天主，与天主父平等；摩尼派无法承担基督耶稣（我们必须靠祂得救）屈尊降生为人并由人的血肉所生；以及其他各种乖谬异端的人，他们无论如何也无法承担圣经及天主教信仰中与他们的谬误相悖的一切，正如我们也无法承担他们亵渎的胡言乱语和谎言的疯狂。因为，不能承担，不是别的，正是不能以平静和沉着的心态保持于心中吗？但是，自我们的主升天以来，所有以正典真理和权威写下的东西，不是每一个信徒和慕道者（甚至在他领受圣洗前）都以平静的心情阅读和聆听吗？即使他们尚未理解得当。那么，即使圣神尚未赐予他们，门徒们怎么就不能承担那些在主升天后所写的任何东西呢？而如今，慕道者在领受圣神之前，都能承担这一切？因为，虽然信徒的圣事特权没有向他们展示，但这并非因此他们就不能承担；而是为了让他们更加热切地渴求，这些特权被荣耀地隐藏起来。\par

4. 所以，亲爱的诸位，你们不必期望从我们口中听到主当时没有告诉门徒的那些事，因为他们尚且不能承担；反而应当努力在因所赐予你们的圣神而倾注在你们心中的爱德中成长；\BibleRef{罗 5:5}这样，你们由于心神火热，热爱属神的事物，就能——不是藉着任何你们肉眼可见的迹象，或任何冲击你们肉耳的声音，而是藉着内心的眼睛和耳朵——认识那属神的光和属神的话，那是属血肉的人不能承担的。因为完全不知道的事物是无法被爱的。但是，当那已知的（哪怕只有一点点）被爱的时候，因着这同样的爱，人就被引向更好、更圆满的知识。如果你们在圣神倾注在你们心中的爱德中成长，祂\BibleQuote{将教导你们一切真理}；或者，如同其他抄本所记，\BibleQuote{祂将指引你们进入一切真理}：正如经上所说：\BibleQuote{主啊，求祢引领我走在祢的道路上，我必行在祢的真理中。}其结果将是，你们不是从外在的老师那里学到主当时没有说出的事，而是都由天主所教导；这样，那些你们通过外在的教训和宣讲所学习并相信的关于天主本性的事——即祂是无形的，不受限制，却又不是像物质团块那样延展至无限空间，而是处处完整、完美、无限，没有颜色的闪烁，没有形体的轮廓，没有文字记号或音节连续——你们的心灵本身可能就有能力察觉。好了，我刚才说了一些也许也属于同类性质的事，而你们却接受了；你们不仅能够承担，而且还愉快地聆听了。但是，如果那位内在的老师（祂曾外在性地对门徒说：\BibleQuote{我还有许许多多的事要告诉你们，但你们现在不能承担}）愿意以祂对天使说话的方式（天使常常见到天父的面，\BibleRef{玛 18:10}），内在地对我们说出我刚才所说的天主无形本性的事，我们仍然会不能承担。因此，当祂说：\BibleQuote{祂将教导你们一切真理}，或\BibleQuote{将指引你们进入一切真理}时，我认为在今生任何人的心灵中都不可能完全实现（因为谁活在这样可朽坏而压抑灵魂的身体里，\BibleRef{智 9:15}能知道一切真理呢？即使宗徒也说：\BibleQuote{我们现在所知道的，只是局部}），而是因为藉着圣神——我们现已领受了祂的抵押，\BibleRef{格后 1:22}——我们将实际达到知识的圆满；同样一位宗徒论及此事说：\BibleQuote{但那时就要面对面了}；以及\BibleQuote{我现在所知道的只是局部，但将来我要清清楚楚地知道，如同我被知道一样}；\BibleRef{格前 13:9}这并不是他在此生完全知道的事，而是作为达到那圆满境界时仍为未来的事，主藉着圣神的爱许诺给我们，当祂说：\BibleQuote{祂将教导你们一切真理}，或\BibleQuote{将指引你们进入一切真理}时。\par

5. 事既如此，亲爱的，我在基督的爱内劝告你们，要提防不洁的诱惑者和污秽可憎的派别，宗徒论及这些说：\BibleQuote{他们暗中所行的，就是连提起，也是可耻的。}\BibleRef{弗 5:12} 免得当他们开始教导他们可怕的污秽——任何人的耳朵都无法忍受这些时——他们竟宣称这些正是主所说的\BibleQuote{我还有许多事要告诉你们，然而你们现在不能担负}，并断言是圣神的作为使这些不洁可憎之事变得可以忍受。邪恶之事是没有人性的谦逊所能忍受的，而另一类则是人的微薄悟性无法承受的美好事物：前者在不清贞的身体中施行，后者超越一切身体；一个在肉体的污秽中犯下，另一个却连纯洁的心灵也几乎无法觉察。因此，\BibleQuote{你们应在心思念虑上更新}\BibleRef{弗 4:23}，并\BibleQuote{领会什么是天主的旨意，什么是良善、可悦、成全的}\BibleRef{罗 12:2}，好使你们\BibleQuote{在爱德上根植并奠立基础，能与众圣徒一同理解何为长、宽、高、深，并认识基督那超越知识的爱，为使你们充满天主的一切富裕}\BibleRef{弗 3:17}。因为圣神将如此教导你们一切真理，当祂将那爱日益广施在你们心中时。\par

\SourceRecord{Translated by John Gibb.；Nicene and Post-Nicene Fathers, First Series；Vol. 7.；Edited by Philip Schaff.；Buffalo, NY: Christian Literature Publishing Co.,；1888.；<http://www.newadvent.org/fathers/1701096.htm>.}

% structure: p0001 source=h1 target=section reason=page-title

\section{《论若望福音》第97讲（若 16:12-13）}

1. 圣神，即主应许派遣给祂门徒的那一位，为教导他们当时祂向他们说话时他们尚不能承受的全部真理。关于这圣神，正如宗徒所说，我们现在已领受了\Quote{抵押品，}\BibleRef{格后 1:22}，藉此表达我们明白祂的圆满是为我们保留到另一生命。因此，圣神也在现世教导信徒，按各人所能领悟属灵之事；并在他们心中燃起不断增长的渴望，各人依照在那爱德中的进展，这爱德将引领他既爱他所已知的，又渴慕那尚待认知的。这样，他现在略有认识的那些事，他也知道自己仍一无所知，因为它们的认知有待于那\Quote{眼所未见，耳所未闻，人心所未曾想到}的生命\BibleRef{格前 2:9}。但若内在的导师现在愿意以那种认知的方式说出那些事，即展开并使之显明于我们的心智，我们人性的软弱将无法承受。关于这一点，亲爱的诸位，你们记得我已讲过，当我们研讨圣福音的话语时，主说：\Quote{我还有许多事要告诉你们，然而你们现在不能承担。}并非在这些话语中我们应当怀疑主过分谨慎地隐藏了无人知晓的秘密，这些秘密可说得出却无法被学习者承受；而是那些与宗教教义相关、我们阅读书写、听闻谈论的事，若按某些人的知识，倘若基督愿意以祂向圣天使们说话的方式，即作为圣父的独生圣言、与祂同等的永恒者，亲自向我们说出，又有谁能承受呢？即使他们已是属神的人，而当初主这样对他们说话时宗徒们还不是，后来当圣神降临时他们才成为那样。因为，凡受造物所能知道的，总小于造物主自己，祂是无上、真实、不变的天主。然而有谁缄默不言祂呢？祂的名号哪里不在读者、辩论者、探询者、回答者、朝拜者、歌唱者、各类演讲者，甚至亵渎者的口中？虽然无人对祂缄默，但谁能理解祂之所是呢？尽管祂从未离开人们的口和耳？有谁的心智之敏锐能接近祂？若祂自己不愿意被认识，有谁会认识祂是天主圣三呢？如今有谁对那天主圣三闭口不言？然而又有谁对那天主圣三具有天使那样的概念呢？因此，那些关于天主的永恒、真理、圣德而不断被随口公开谈论的事，有人理解得好，有人理解得差；更确切地说，有人理解，有人全然不理解。因为理解得差的人根本不算理解。就连那些正确理解的人，有人领悟得较模糊，有人较清晰，而在世上无人能像天使那样领悟。因此，在心智中，即在内在的人中，有一种增长，不仅在于从奶到固体食物的过渡，也在于食物本身越来越大的摄取。但这种增长并非物质占据空间体积的方式，而是理解光照的方式；因为那食物本身就是理解的光。为使你们增长并领悟天主，为使你们的领悟与你们不断进步的增长同步，你们应当向祂祈祷，并寄望于祂，而非寄望于那声音只达于你们耳中的导师，即仅藉外在工作栽种和浇灌的，而是寄望于那赐予成长的\BibleRef{格前 3:6}。\par

2. 因此，正如我在上次讲道中劝诫你们的，你们（尤其是那些仍是婴孩、需要奶食的人）要当心，不要好奇地去听那些在主的话中找到自欺欺人把柄的人；主说：\BibleQuote{我还有许多事要告诉你们，但你们现在不能担负。}他们用这话来发现未知之事，而你们的心智尚不能辨别真假。尤其要当心因撒殚藉天主的允许，灌输给不稳定和属血肉的灵魂的淫秽污秽之事；为的是使天主的审判在各地成为可畏的对象，使纯洁的纪律在与邪恶的污秽对比中最充分地显示其甘甜；并使荣耀归于祂，每个人都怀着敬畏和谦逊的态度，无论是那些被祂的王权保护免于陷入如此邪恶的人，还是那些被祂提拔的手从其中提升起来的人。要怀着敬畏和祈祷，谨防陷入所罗门的奥秘，那里\BibleQuote{那愚蠢、厚颜无耻、断粮的女人}用这些话邀请过路的人：\BibleQuote{来，吃隐藏的饼，偷来的水是甘甜的。}因为这所说的女人，就是不敬虔者的虚荣；他们毫无感觉，自以为知道些什么，正如那女人所说的，她\BibleQuote{断粮了}；她虽然连一块饼也没有，却应许饼；换句话说，她虽不知道真理，却应许真理的知识。但她应许的是隐藏的饼，并声称它是被愉快地分享的，还有偷来的水的甘甜；为的是那些在教会中公开禁止说出或相信的事，可以被乐意和津津有味地听从和实行。因为通过这种秘密，亵渎的教师给好奇者一种调味品，使他们以为学到了什么了不起的东西，因为被当作守秘密的人，并更甘甜地吸收他们所视为智慧的愚昧，而他们被视为偷听了禁止听的事。\par

3. 因此，魔法技艺的系统通过亵渎的好奇心，向那些被欺骗或即将被欺骗的人推荐其邪恶的仪式。同样，那些通过检查宰杀动物的内脏、鸟的叫声和飞行、或各种魔鬼标志的非法占卜，通过与堕落者的交谈注入到濒临毁灭之人的耳中。正是由于这些非法和可罚的秘密，上面提到的女人才不仅被称为\Quote{愚蠢}，而且被称为\Quote{厚颜无耻}。但这类事情不仅与真实无关，而且与我们宗教的名称也格格不入。我们该怎么说这个愚蠢厚颜的女人呢？她竟用基督教的表达方式调配如此多的邪恶异端，端出如此多可憎的寓言！但愿它们只是剧院里发现的那种，作为歌曲或舞蹈的主题，或由模仿的滑稽戏弄而成笑料；而不至于其中有些让我们为愚蠢而悲伤，同时又惊奇于竟有人敢如此反对天主而编造它们！然而所有这些完全无感觉的异端者，他们希望被称为基督徒，试图借主所说的福音那句话来粉饰他们完全违反人类情感的厚颜无耻的发明，那句话是：\BibleQuote{我还有许多事要告诉你们，但你们现在不能担负。}仿佛正是这些事是宗徒们当时不能担负的，仿佛圣神教导了他们，而那不洁之神尽管可以尽其厚颜无耻之所能，却羞于在光天化日之下教导和宣讲这些事。\par

4. 宗徒藉圣神预先看见的正是这种人，当他说：\Quote{时日将到，人必不容忍健全的道理，反而放纵自己的私欲，为自己聚拢许多师傅，耳朵发痒，且掩耳不听真理，而偏向无稽之谈。}\BibleRef{弟后 4:3} 因那提及隐密与偷窃的话，所谓\Quote{偷来的水是甜的，暗中的饼是美味的}，使那些以追逐灵性奸淫之耳听者发痒，正如肉欲的瘙痒败坏贞洁的健全。因此，请听宗徒如何预见这些事，并给予有益的劝诫以避之，他说：\Quote{要远避亵渎的言论新说，因为它们使人日趋不敬神，而他们的话如同毒疮蔓延。}他并非只说\Quote{言论新说}，而加上了\Quote{亵渎的}。因为也有与宗教教义完全和谐的言论新说，正如圣经所载基督徒之名初用时所述。原来是在主升天之后，安提约基雅的门徒首先被称为基督徒，如我们在宗徒大事录中所读到的：\BibleRef{宗 11:26} 此后，某些房屋被称为收容所和隐修院的新名；但事物本身先于名称，且由宗教真理所确认，这真理也是抵抗恶者的保障。为反对亚略异端者的不敬，他们造了新的术语 \OriginalTerm{父的同性同体}{Patris Homousios}；但此名所指并非新事；因为所谓\Quote{同性同体}正是指\Quote{我与父是一体}，即同一实体。因为若每一新奇皆是亵渎，则主说\Quote{我给你们一条新命令}亦不应存留，新约也不应称为新，新歌也不应遍唱大地。但言论新说中有亵渎，正如\Quote{愚昧狂妄的妇人}所说：\Quote{来，享受暗中的饼和偷来的水的甘饴。}宗徒从这种虚假学问的诱语中也给予禁止的警告，在那段话中他说：\Quote{啊，弟茂德，要保管所受的寄托，远避亵渎的言论新说及虚假学问的反对；有些人自命不凡，竟然背离了信德。}\BibleRef{弟前 6:20} 因为这些人的爱好莫过于自称有学问，并嘲笑年轻人被命相信的那些真理为纯然愚蠢。\par

5. 但有人会说：属神的人在教义方面难道没有一些事，是决不可对属血肉的人说，而只可对属神的人明言吗？我若回答\Quote{没有}，便会立即面对宗徒保禄在致格林多人书中所言：\Quote{我从前对你们说话，不能把你们当作属神的，只能当作属血肉的人，当作在基督内的婴孩。我给你们喝的是奶，并非饭食，因为那时你们还不能吃，就是如今还是不能，因为你们仍是属血肉的}；\BibleRef{格前 3:1} 以及\Quote{我们在成全的人中讲智慧}，又\Quote{以属神的事解释属神的事：然而属血气的人不能领受天主圣神的事，因为为他乃是愚妄。}凡此种种，为使宗徒的话不再因亵渎的言论新说而引人追求隐密，且为使那本该为贞洁者心神身体所常避讳者，不被称为仅为血肉者所不能承受，我们蒙主恩准，将在另一篇讲道中对此加以论述，以便此刻先结束本篇。\par

\SourceRecord{Translated by John Gibb.；Nicene and Post-Nicene Fathers, First Series；Vol. 7.；Edited by Philip Schaff.；Buffalo, NY: Christian Literature Publishing Co.,；1888.；<http://www.newadvent.org/fathers/1701097.htm>.}

% structure: p0001 source=h1 target=section reason=page-title

\section{论文第九十八篇（若十六12－33）}

1. 主说：\Quote{我还有许多事要告诉你们，然而你们现在不能担负。}从这些话引起了一个难解的问题，我记得曾在上一篇讲道中搁置，因为它已达到适当的长度，需要结束，所以留待以后更有余暇时处理。如今，我们既有时间履行诺言，就让我们在主的恩赐下讨论这个问题，是祂把这个提议放在我们心中的。问题是：属神的人是否在教义中有所保留，不告诉属血肉的人，却只向属神的人揭示？因为如果我们说\Quote{没有}，就会遇到宗徒在致格林多人书信中的话的对答：\Quote{我从前不能向你们说话，如同对属神的人，只能如同对属血肉的人，如同对基督内的婴孩。我给你们喝的是奶，不是食物，因为那时你们还不能吃，就是如今还是不能，因为你们仍是属血肉的。}（\BibleRef{格前 3:1}）但如果我们说\Quote{有}，我们就有理由恐惧和谨慎，免得在这种借口下，秘密教导可憎的教义，并借着属神的名号，似乎不仅可以用貌似合理的借口粉饰，还值得在讲道中颂扬。\par

2. 首先，你们的爱德应当知道，宗徒说他是用奶喂养那些作为婴孩的，而作为奶的正是被钉在十字架上的基督自己；但祂的肉身本身，在其中见证了祂真实的死亡，即祂被刺穿时的真实伤口和祂被刺透时所流的血，向属血肉之人的心灵所呈现的方式，与向属神之人的心灵不同：因此对前者是奶，对后者是干粮；因为他们即使不比别人听得更多，他们也理解得更好。因为对于那些同样凭信德领受的事物，心灵的领悟能力并不相等。因此，宗徒所宣讲的被钉十字架的基督，对犹太人来说是绊脚石，对外邦人是愚妄，但对那蒙召的，不论犹太人或希腊人，却是天主的德能和天主的智慧（\BibleRef{格前 1:23}）。然而对于属血肉的人，如同那些只凭信德持守的婴孩，而对于属神的人，如同那些更有能力的人，却当作可理解的事物去领悟：因此，对前者是如奶一般的饮品，对后者是固体食物。这并不是说前者在公开场合以一种方式认识，后者在密室中以另一种方式认识；而是说，当两者同样听到公开讲论时，各自按照自己的程度去领会。因为基督被钉十字架，正是为了流出祂的血以赦免罪过，并借着祂独生子的苦难彰显天主的恩宠，使人不以任何人夸耀，那么那些仍在说\Quote{我是属保禄的}的人（\BibleRef{格前 1:12}）对基督被钉十字架有何理解呢？难道他们具有保禄本人那样的理解吗？保禄能说：\Quote{至于我，我决不以别的夸耀，只以我们的主耶稣基督的十字架来夸耀}（\BibleRef{迦 6:14}）。因此，即使是关于基督被钉十字架，他本人也按自己的能力找到食物，并按照他们的软弱用奶喂养他们。更进一步，他知道他写给格林多人的话无疑会被那些尚是婴孩的人以一种方式理解，而被那些更有能力的人以另一种方式理解，于是他说：\Quote{若有人自以为是先知，或属神的人，就该承认我写给你们的，是主的诫命；但不承认的，就由他不承认吧。}（\BibleRef{格前 14:37}）他当然希望属神的人的知识是实在的，只要不仅信德找到了适当的居所，而且还拥有某种理解的能力；借着这种能力，他们相信那些作为属神的人同样承认的事。至于\Quote{不承认的，就由他不承认}，他说，因为他还未被启示去认识他所相信的。当这事发生在人的心灵中时，他便被称为被天主所认识；因为是天主赐予他这种理解的能力，正如别处所说：\Quote{如今你们认识了天主，更好说已被天主所认识了。}（\BibleRef{迦 4:9}）因为并不是在那时天主才初次认识那些在世界创立以前就被预知和拣选的人（\BibleRef{弗 1:4}）；而是在那时祂使他们认识祂自己。\par

3. 因此，一开始就确定了这一点：那些同等地被属神的人和属血肉的人听到的事物，是各人按自己狭窄的容量接受的——对一些人如同婴儿，对另一些人如同较成熟的人——对一人如同奶的滋养，对另一人如同固体食物——似乎没有必要将任何教义问题作为秘密保留沉默，隐瞒于婴孩般的信徒，而只对较年长或拥有更成熟理解力的人私下谈论；我们应视为必须这样做，只因宗徒的话：\Quote{我未能对你们讲话如同对属神的人，而如同对属血肉的人。} 因为就连他这句话本身——他在他们中间不知道别的，只知道耶稣基督，且是被钉在十字架上的基督\BibleRef{格前 2:2}——他也未能对他们讲话如同对属神的人，而如同对属血肉的人，因为就连那一件事他们也不能作为属神的人而接受。但他们中间所有属神的人，都以属神的理解接受了那些真理，而其他人只以属血肉的方式听见；这样，我们可以理解这句话：\Quote{我未能对你们讲话如同对属神的人，而如同对属血肉的人}，好像他说：我所说的话，你们不能作为属神的人接受，而只能作为属血肉的人。因为\Quote{属血气的人}——即那人的智慧纯属人的种类，从灵魂称为\OriginalTerm{属魂的}{soulish}，从肉体称为\OriginalTerm{属血肉的}{carnal}，因为完整的人由灵魂和肉体构成——\Quote{不能领受天主圣神的事}；\BibleRef{格前 2:14}反而认为那个十字架所成就的一切，只是提供给我们一个榜样，让我们为真理而争战甚至致死。因为如果这类人——他们不愿成为别的，只愿成为人——知道了被钉在十字架上的基督如何成为\Quote{由天主给我们成了智慧、正义、圣化、救赎，为应验经上所记载的：夸耀的，应因主而夸耀}，\BibleRef{格前 1:30}他们无疑不会再夸耀于人，也不会以属血肉的精神说：\Quote{我是属保禄的，我是属阿颇罗的，我是属刻法的}，而是以属神的方式说：\Quote{我是属基督的}。\BibleRef{格前 1:12}\par

4. 但由我们在\WorkTitle{希伯来书}中读到的，问题进一步被提出：\Quote{按时间你们本应做导师，却仍需有人教导你们天主圣言的开端；你们成了需要奶，而不是坚硬食物的人。因为凡用奶的，对正义之言没有经验，因为他是婴孩。但坚硬食物属于成全的人，就是那些因习惯而感官锻炼了，能辨别善恶的人。}\BibleRef{希 5:12} 因为我们在这里看到，似乎清楚地定义了，他所说的成全之人的坚硬食物是什么；这与他在\WorkTitle{格林多书}中所写的是相同的：\Quote{我们在成全的人中讲智慧。}\BibleRef{格前 2:6} 但他希望在这段经文中被理解为成全的人是谁，他接着用这些话指出：\Quote{就是那些因习惯而感官锻炼了，能辨别善恶的人。} 因此，那些因软弱和未受训练的心智而缺乏这能力的人，除非被那可称为信德的奶所扶持，以相信他们未见的不显之事，以及他们尚未领悟的可领悟之事，就必定会被知识的许诺轻易引诱到虚妄而亵圣的神话中：以致只以形体样式思考善恶，对天主本身也没有任何概念，只当作某种形体，并且只能视恶为一种实体；而实际上，对于一切可变的实体——它们是由不变且至高的实体本身（即天主）从无中创造的——是一种从不变本体的坠落。而且，凡是不但相信，而且通过内心锻炼的感官，理解、领悟并知道这点的，就无需再担心会被那些认为恶是天主未创造的实体、从而使天主本身也成为可变实体的人所诱惑，正如摩尼教徒所做的，或者任何其他类似愚妄的害虫——若有的话。\par

5. 但对于那些在心智上仍是婴孩，属于血肉的人，宗徒说他们需要用乳哺养，所有关于此类主题的论述——其中我们不仅涉及相信，也涉及对所言之事的理解和认识——对他们来说必定是沉重的，因为他们尚不能领会这样的事，而这样的论述更适合压迫而非滋养他们。由此便发生这样的事：属神的人虽然并不完全对属于血肉的人缄口不言这类主题，因为公教信仰应当向所有人宣讲，但他们处理这些主题时，并不因为希望简化给尚缺乏能力的理解力而损及他们关于真理的言论，反而损及他们言论中的真理在听众理解中的感知。因此，他在\WorkTitle{致哥罗森人书}中说：\Quote{我身体虽不在你们那里，心灵却与你们同在，欣喜欢乐，看见你们的秩序，以及你们对基督的信德所缺少的。} \BibleRef{哥 2:5} 又在\WorkTitle{致得撒洛尼人书}中说：\Quote{黑夜白日，}他说，\Quote{更恳切地祈求，使我们能见你们的面，并补全你们信德所缺的。} \BibleRef{得前 3:10} 当然，我们在此是指那些处于初级教理训导之下的人，其所接受的是乳的滋养而非干粮；关于前者，在\WorkTitle{致希伯来人书}中提到丰足的供应，给那些他本应此刻已以固体食物喂食的人。因此他说：\Quote{因此，让我们离开基督的开端之言，进到完成；不再重立悔改死亡之行的根基，以及对天主的信德、圣洗池的教义、按手、死人复活、和永恒的审判。} \BibleRef{希 6:1} 这就是丰足的乳，没有它，甚至那些已经充分运用理性相信、却还未能辨别善恶、以致不仅需要信德也需要理解（这属于干粮的范围）的人也无法存活。但当他也在对乳的描述中包含了教义时，那就是传授给我们的信经和天主经中的教义。\par

6. 但我们切不可认为，这乳与必须由健全理智接受的神性事物之食粮之间有任何矛盾，这食粮正是哥罗森人和得撒洛尼人所缺少、尚待补足的。因为补足缺乏并不表示对已有之物的不认可。因为即使在我们所取用的食物中，乳与干粮之间也绝无矛盾，反而后者本身变为乳，以便适合婴孩，通过母亲或乳母的身体到达他们；同样，母亲智慧本身——她在天使的高天是干粮——也屈尊俯就，仿佛成为婴孩的乳，那时圣言成了血肉，居住在我们中间。但人基督自己——在其真肉身、真十字架、真死亡、真复活中被称为婴孩的纯洁之乳——当被属神的人正确理解时，便被发现是天使的主。因此，婴孩不应总以乳喂养，以致始终不能理解基督的天主性；也不应如此断乳，以致背弃祂的人性。同样的事也可以用另一种方式这样表述：他们既不应以乳喂养到从不理解基督是创造者的地步，也不应断乳到永远背弃基督是中保的地步。在这方面，事实上，母乳与干粮的类比很难与如此陈述的现实相协调，但根基的类比更为合适：因为当孩子被断乳，脱离婴儿的滋养时，他再也不会在干粮中寻找他吮吸过的乳房；但被钉十字架的基督对婴孩是乳，对更成熟者是干粮。因此，根基的类比更为合适，因为建筑是在建造过程中加上的，而根基并未被抽走。\par

7. 既然是这样，无论你是谁，你们中无疑有许多人仍然是在基督内的婴孩，你们要向着心灵的固体食物长进，不是腹中的。在辨别善恶的能力上增长，并越来越依附于那从邪恶中拯救你们的中保；这种拯救不是将你们与邪恶在空间上分开，而是在你们内治愈邪恶。但无论谁对你们说：不要相信基督是真的人，或任何人的身体或动物的身体是由真天主创造的，或旧约是由真天主赐下的，以及诸如此类的事，因为当你们以奶为养料时，这些事先前没有告诉你们，因为你们的心还不能接受真理：这样的人不是给你们干粮，而是毒药。因此，蒙福的宗徒在向那些在他看来已经成熟的人讲话时，甚至在称自己不完全之后，说：\BibleQuote{所以，我们凡是成熟的人，都应怀有这种心情；如果你们在什么事上另有看法，天主也会将此启示给你们。} 而且，为了使他们不致落入欺骗者的手中，这些欺骗者想藉着应许他们认识真理，使他们背弃信德，并以为宗徒的话\BibleQuote{天主也会启示给你们}就是此意，他立刻加上说：\BibleQuote{但是，我们已达到什么程度，就应该照着什么程度进行。} \BibleRef{斐 3:15} 那么，如果你们对某些不与公教信仰的准则相悖的事有所领悟，你们已达到这准则，作为引导你们走向天乡的道路；并且你们如此领悟，以致感到有责任消除对这主题的一切疑虑：那么，就在建筑物上加添，但不要放弃根基。而且，年长者给婴孩的任何教导，其性质必定是不得包含这样的论断：即万有之主基督、以及比他们年长许多的先知和宗徒，在任何方面说了假话。你们不仅应当避免那些喋喋不休的心灵欺骗者，他们喋喋不休地讲述自己的寓言和谎言，并在这些虚妄中，装模作样地应许深奥的学问，违反我们视为公教的信仰准则；但也要避免那些人，他们比前者更具阴险的祸害，因为他们虽然真实地讨论了神性不变性、无形受造物或造物主，并用最有说服力的文献和推理充分证明了他们肯定的事，但却试图使你们离开天主与人之间唯一的中保。因为宗徒论及这样的人说：\BibleQuote{因为他们虽然认识了天主，却没有以祂为天主而予以光荣。} \BibleRef{罗 1:21} 一个人若没有抓住那使人从邪恶中得救的一位，他即使有对不变之善的真实理解，又有什么益处呢？最蒙福的宗徒的劝诫，绝不可在你们心中失去其地位：\BibleQuote{若有人给你们宣讲与你们所领受不同的福音，当受诅咒。} \BibleRef{迦 1:9} 他并没有说\BibleQuote{多于你们所领受的}，而是说\BibleQuote{不同于你们所领受的}。因为若他说了前者，他就会判断自己，因为他渴望来到得撒洛尼人那里，以补充他们信德中缺乏的。但补充的人，是在欠缺上加添，而不除去已有的；而违反信仰准则的人，不是在道路上进步，而是偏离了道路。\par

8. 因此，当主说：\BibleQuote{我还有许多事要告诉你们，然而你们现在不能担负。} 祂的意思是，他们仍无知的事，日后必须补充给他们，而不是要推翻他们已经学得的。事实上，正如我在前一篇讲道中已指出的，祂之所以能这样说，是因为祂所教导他们的事，如果祂愿意像天使对祂所理解的那样向他们启示，他们依然脆弱的人性将无法承受。但任何属神的人可以向另一个人传授他所知道的，只要圣神赐给他更大的获益能力，在这过程中教师自己也可能得到进一步的增长，为使两人同受天主的教导。虽然即使在属神的人中，无疑也有些人的能力比其他人更强、状态更好，以致他们中有人达到了人不可言传的境界。有些人利用这一点，凭着暴露极端愚昧的狂妄，伪造了\WorkTitle{保禄启示录}，充满了各种无稽之谈，被正统教会所拒绝；他们声称这就是保禄所说的他被提到第三层天，在那里听到不可言传的话\BibleQuote{是人不可说出的}。然而，如果他说的他是听到了一些现在还不容许人说出的的话，这种人的胆大妄为或许还可容忍；但当他已说了\BibleQuote{是人不可说出的}，谁又敢如此肆无忌惮、毫无效果地说出来呢？但我现在要以这些话结束这篇讲道；我愿你们在善事上确实明智，在恶事上却不受玷污。\par

\SourceRecord{Translated by John Gibb.；Nicene and Post-Nicene Fathers, First Series；Vol. 7.；Edited by Philip Schaff.；Buffalo, NY: Christian Literature Publishing Co.,；1888.；<http://www.newadvent.org/fathers/1701098.htm>.}

% structure: p0001 source=h1 target=section reason=page-title

\section{\WorkTitle{讲道集 99（若 16:13）}}

1. 主论及圣神，当祂应许将要来到并教导祂的门徒一切真理，或引导他们进入一切真理时，说了什么？\BibleQuote{因为祂不凭自己说话，而是祂所听见的，祂才讲论}？因为这与祂论及自己所说的相似：\BibleQuote{我凭自己不能做什么：我怎样听见，就怎样审判}。但当我们解释那句话时，我们曾说，那可以理解为指着祂的人性而言；如此，祂作为圣子似乎预先宣告祂自己的服从（祂服从至死，且死在十字架上\BibleRef{斐 2:8}）也将出现在审判中，那时祂将审判活人死人；因为祂之所以如此行，正是由于祂是人子。因此祂说：\BibleQuote{父不审判任何人，但将一切审判交给了子}；因为在审判时，祂将显现，不是以天主的形像——即祂与父平等，恶人不能看见——而是以人的形像，在祂甚至暂时比天使微小一点的形式中；虽然那时祂将在光荣中来临，而非起初的谦卑，但却是善人和恶人都能看见的。因此祂进一步说：\BibleQuote{并且祂给了祂执行审判的权柄，因为祂是人子}。在祂自己的这些话中，清楚表明那不是那位在祂不以自己与天主同等为僭越时所拥有的形像将出现在审判中，而是祂在空虚自己时所取的那形像。因为祂空虚自己，取了奴仆的形像\BibleRef{斐 2:6}；在那形像中，为了执行审判，祂似乎也引证了祂的服从，当祂说：\BibleQuote{我凭自己不能做什么：我怎样听见，就怎样审判}。因为亚当，因他的不顺服，如同一个人，使许多人成为罪人，他并没有按他所听见的审判；因为他歪曲了他所听见的，并凭自己做了他所做的恶；因为他没有遵行天主的旨意，而是自己的；而这位（基督），因祂的服从，如同一个人，使许多人成为义人\BibleRef{罗 5:19}，祂不仅服从至死，且死在十字架上——就这一点，祂被视为从死者中活着受审判——而且许诺在审判本身中也显示服从，那时祂将作为活人死人的审判者，当祂说：\BibleQuote{我凭自己不能做什么：我怎样听见，就怎样审判}。但是，当论及圣神时，说：\BibleQuote{因为祂不凭自己说话，而是祂所听见的，祂才讲论}，我们敢不敢认为这句话是针对祂的任何人性或任何受造形像的假设？因为只有圣子在圣三内取了奴仆的形像，那形像在祂身上被纳入祂位格的合一，或者说，为了我们不至于宣讲一个\OriginalTerm{四体}{quaternity}而非圣三——这是天主禁止的。正是由于这同一个位格由两种性体（天主性和人性）组成，祂有时按照祂作为天主的性体说话，如当祂说：\BibleQuote{我与父原是一体}；有时按照祂的人性说话，如：\BibleQuote{父比我大}；据此，我们也理解了祂目前正在讨论的那些话：\BibleQuote{我凭自己不能做什么：我怎样听见，就怎样审判}。但是，关于圣神的位格，有一个相当大的难题，即我们如何理解\BibleQuote{因为祂不凭自己说话，而是祂所听见的，祂才讲论}这句话；因为在那位格内，并没有一个天主性体加上另一个人性或任何其他受造物的性体。\par

2. 因为圣神以形体出现，像鸽子一样，\BibleRef{玛 3:16}，是一件在当时开始并结束的景象：同样地，当祂降临在门徒们身上时，他们看到如同火的舌头分开来，停留在他们每人身上。\BibleRef{宗 2:3} 因此，任何人若说鸽子与圣神在祂位格的合一中相连，如同它和天主性（因为圣神是天主）共同构成圣神的一个位格，也不得不断言关于那火的同样事情；这样他就明白他两者都不该断言。对于关于天主实体的事，在任何时候需要用某种外在方式表现，并如此向人的身体感官显示，然后消逝，都是由天主的能力从可受役使的受造物中暂时形成，而非来自主导本性本身；这本性永远同一，激发它所愿意的任何行动；它自身不变，却随意改变万物。同样，那来自云中的声音也确实击打在身体耳朵上，即称为听觉的身体感官上；\BibleRef{路 9:35} 然而，我们决不可相信天主的圣言，即独生子，因为祂被称为圣言，是由音节和声音界定的：因为当一篇讲道在进行时，所有声音不能同时发出；而是各个声音依次诞生，取代消逝的声音，以致我们所说的一切只由最后一个音节完成。确实，圣父对圣子说话的方式，即天主对天主，祂的圣言，与此截然不同。但这，就人所能理解的，是一个适合那些接受固体食物而非奶的人的理解。因此，既然圣神并非借着取人性而成为人，也并非借着取天使性而成为天使，也并未借着取任何受造物形式而进入受造物状态，那么关于祂，我们如何理解我们主的话：\Quote{因为祂不凭自己说话，但祂所听见的，祂就说话}？一个困难的问题；是的，太过困难了。愿圣神亲自临在，使我们至少按我们思考此类事物的能力，能表达我们的思想，并使这些思想，按我微小的能力，进入你们的理解。\par

3. 那么，首先你们应当被告知，并且你们中能理解的，去理解，而其余尚不能理解的人，去相信：在那实质体（即天主）内，感官并非像借着身体物质构造那样分布在适当部位；如同在所有动物的可朽肉体中，一处有视觉，另一处有听觉，另一处有味觉，另一处有嗅觉，遍布全身的是触觉。我们断不可对那无形不变的本性作如此相信。因此，在其中，听觉和视觉是同一回事。同样，嗅觉也被说存在于天主内；如宗徒所说：\Quote{如同基督也爱了我们，并将祂自己献上，作为奉献和祭品，呈于天主，成为馨香的祭品。} \BibleRef{弗 5:2} 味觉也可包括在内，按此天主憎恨心态苦毒，并将那些不冷不热、既非冷也非热的人从祂口中吐出：\BibleRef{默 3:16} 而我们的天主基督说：\Quote{我的食物就是承行派遣我者的旨意。} 还有那神圣的触觉，按此新娘论新郎说：\Quote{他的左手在我头下，他的右手将拥抱我。} \BibleRef{歌 2:6} 但这些在天主内并非在身体的不同部分。因为当祂被称为知道时，一切都包括在内：视觉、听觉、嗅觉、味觉、触觉；没有任何实体的改变，也没有任何在一个地方较大、在另一个地方较小的物质要素的存在：并且即使在年长者心中，若有任何关于天主的此类思想，那也只是孩童般的思想。\par

4. 你无需对那不可言说的天主知识感到惊奇，祂凭借这知识洞悉万事，因人类言语的不同方式，用所有那些身体感官的名称来指称；因为我们自己的心灵，亦即内在的人——它以一种统一的方式运用其认知能力时，其知识的各个不同对象由身体的五位使者传达给它，当它理解、选择并爱慕那不变的真理时——据说既能看见光，关于这光曾说过：\BibleQuote{那是真光；}也能听到圣言，关于这圣言曾说过：\BibleQuote{在起初已有圣言；}并有嗅觉，关于这曾说过：\BibleQuote{我们愿追随你香膏的芬芳；}并饮用这泉源，关于这曾说过：\BibleQuote{在祢那里有生命之源；}并享受触觉，当时曾说过：\BibleQuote{亲近于天主是我的福乐；}在这一切中，并不是不同的事物，而是同一个理智，藉着如此众多感官的名称来表达。因此，当论及圣神说：\BibleQuote{因为祂不凭自己讲话，而是祂所听到的，祂就讲出来，}一个在真正意义上单纯的（不复合的）本性，就更应当被理解或相信；这本性在其广延与崇高上远超我们心灵的本性。因为在我们心灵中存在着可变性：它通过学问获得对先前无知之事的感知，又因遗忘而失去原先所知；它被与真理相似的事物所欺骗，以致将虚假当作真理而认同，并因自身的昏暗如同被一种黑暗所阻碍，无法达至真理。所以，那存在与认知不等同的实体，并非在真正意义上单纯的；因为它可以存在而不拥有知识。但神圣实体却非如此，因为祂就是祂所拥有的。因此，祂的认知并非那种使其认知的知识是一回事，使其存在的本质是另一回事；而是两者为一。那单纯为一的，不应被称为两者。\BibleQuote{正如圣父在自己内有生命，}而祂自己与祂内的生命并非不同；\BibleQuote{同样也赐给圣子在祂自己内有生命，}即生了圣子，使祂自己也是生命。因此，我们应当如此接受论及圣神的话：\BibleQuote{因为祂不凭自己说话，而是祂所听到的，祂就说出来，}由此理解祂不是出于自己。因为只有圣父不是出于另一个。因为圣子由圣父所生，圣神由圣父所发；但圣父不受生，也不由另者所发。然而，至高的天主圣三中却不应因此在人的思想中产生任何不平等的观念；因为圣子与生祂的那一位平等，圣神与祂所由发出的那一位平等。但发出与受生在这种情况下有什么区别，若成为探究与论述的主题就太冗长了，并且即使讨论之后，我们的定义也易受鲁莽的指责；因为这样的事极其困难，既令心灵难以充分理解，即使心灵达到了某种理解，也难以用舌头解释，无论那担任教师的人或在场听讲的人多么有才能。因此，\BibleQuote{祂不凭自己说话；}因为祂不是出于自己。\BibleQuote{但祂所听到的，祂就说出来：}祂将从祂所由发出的那一位听到。对祂而言，听就是知，而知就是存在，如上所述。所以，祂既不是出于自己，而是出于祂所由发出的那一位，并从祂获得本质，就从祂获得知识；因此，祂从祂获得听，而这听无非就是知识。\par

5. 你不可因动词用的是将来时而困扰。因为经上所说的不是\Quote{祂所听到的}或\Quote{祂所听见的}，而是\BibleQuote{凡祂将听到的，祂将说出。}因为这样的听是永恒的，因为知是永恒的。但对于永恒者，没有开始也没有终结，无论动词使用何种时态，过去、现在或将来，都不因此暗示谬误。因为虽然对于那不变且不可言说的本性，没有\Quote{曾是}和\Quote{将是}的恰当应用，而只有\Quote{是}；因为唯有那本性真正地\Emph{是}，因为它不能改变；因此，它特别适合于说：\BibleQuote{我是自有者}，以及\BibleQuote{你要对以色列子民说：那\Quote{我是}的派遣我到你们这里来：}\BibleRef{出 3:14}然而，由于我们所度过的有死且易变生命的时代之变化，当我们说祂曾是、将是并同时是时，并无虚假。祂在过去曾是，在现在是，在未来将是。祂曾是，因为祂从未缺席；祂将是，因为祂永不会缺席；祂是，因为祂永远是。因为祂不像不再存活的人那样随过去死去；也不像不能常住者那样随现在消逝；也不像先前不存在者那样随未来升起。因此，正如我们人类的说话方式随时间的流转而变动，那位在一切时间中都不曾、现在不会、将来也必不缺席的祂，可以正确地用任何动词时态来谈论。因此，圣神总是在听，因为祂总是知道：\Emph{所以}，祂既曾知道，也知道，并将知道；同样，祂既曾听见，也听见，并将听见；因为，如我们已说过的，对祂而言听与知是一回事，知与祂的存在是一回事。因此，祂从祂所隶属的那一位，听见过、听见并将听见；而祂也是从祂所由发出的那一位而存在。\par

6. 有人或许会问：圣神是否也由圣子所发？因为圣子唯独是圣父的圣子，圣父唯独是圣子的圣父；但圣神并非二者中某一位的神，而是二者的神。你们有主亲自说：\Quote{因为说话的并不是你们，而是你们父的圣神在你们内说话。}\BibleRef{玛 10:20}又有宗徒说：\Quote{天主派遣了自己圣子的圣神，到你们心中。}\BibleRef{迦 4:6}那么，有两个吗？一个是圣父的，另一个是圣子的？当然不是。因为他说，当论及教会时：\Quote{一个身体}，随即加上\Quote{一个圣神}。请注意他如何在那里构成天主圣三。他说：\Quote{正如你们蒙召，}\Quote{同有一个希望。}\Quote{一个主，}他当然是指基督；但还需提及圣父，于是接着说：\Quote{一个信德，一个圣洗，一个天主，众人之父，祂超越众人，贯通众人，且在众人之内。}\BibleRef{弗 4:4}因此，既然只有一个父，一个主即圣子，也只有一个圣神；祂无疑是属于二者的，尤其是基督耶稣亲自说：\Quote{你们父的圣神住在你们内；}宗徒亦宣称：\Quote{天主派遣了自己圣子的圣神到你们心中。}同一位宗徒在另一处说：\Quote{但如果那使耶稣从死者中复活者的圣神住在你们内，}他显然是指圣父的圣神；然而他在另一处又说：\Quote{谁若没有基督的圣神，就不属于基督。}还有许多其他证据，清楚表明那在天主圣三中被称为圣神的，是圣父和圣子二者的神。\par

7. 我认为，祂之所以被称为圣神，并无其他原因；因为尽管依次询问每一位，我们也不得不承认圣父和圣子各自都是神，因为天主是神，即天主不是属血肉的，而是属神的。因此，借着他们各自共有的名称，有必要明确称呼那一位，祂不是二者中的任何一个，而是在祂内显明了二者的共通之处。那么，我们为何不信圣神也由圣子所发，既然祂同样是圣子的神？因为若祂不是这样发，祂在复活后显现给门徒时，就不能向他们嘘气，说：\Quote{你们领受圣神吧！}这样的嘘气除了表明圣神也从祂发出，还能意味着什么？同样，关于患血漏的妇人，祂所说的话也是同一性质：\Quote{有人摸了我，因为我感觉有能力从我身上出去了。}\BibleRef{路 8:46}因为圣神也被称为能力，这既清楚见于天使回应玛利亚的疑问\Quote{这事如何成就？因为我不认识男人}时所说的\Quote{圣神要临于你，至高者的能力要荫庇你；}\BibleRef{路 1:34}也见于我们的主在应许门徒圣神时说：\Quote{你们要留在城中，直到你们领受从上而来的能力；}\BibleRef{路 24:49}又在另一场合说：\Quote{你们要领受圣神降临于你们的能力，并要为我作证人。}那能力，我们相信就是福音作者所说的：\Quote{有能力从祂身上出来，治好了众人。}\BibleRef{路 6:19}\par

8. 那么，如果圣神既由圣父所发，也由圣子所发，为什么圣子说：\Quote{祂由圣父所发}？你以为是什么原因？岂不是因为祂惯于将属于自己的一切都归于祂（圣父），而圣子本身也是从祂而来？因此我们有祂说：\Quote{我的教训不是我的，而是那派遣我者的。}如果在这段经文中我们应理解那是祂的教训，但祂却声称不是祂的，而是圣父的，那么在另一段经文中，我们岂不更应理解圣神是由祂自己所发，因为祂的话\Quote{祂由圣父所发}并非暗示祂不是由我而发。但圣子从祂（圣父）那里得以是天主（因为祂是出自天主的天主），当然也从祂那里得以使圣神也由祂所发：这样，圣神从圣父本身获得也由圣子所发，正如祂由圣父所发一样。\par

关于这事，我们也多少领悟了进一步的一点，即——就我们这类存在所能领悟的程度而言——为什么圣神不被称为\Quote{生}，而被称为\Quote{发}：因为，如果祂也被称为子的名字，祂就免不了被称为双方的儿子，那是极其荒谬的。因为没有人是两个人的儿子，除非是父亲和母亲。但若怀疑天主圣父与天主圣子之间有此类干预，那是极其可憎的。因为即使是人的儿子，也不是同时从父亲和母亲发出的：当他从父亲发出进入母亲时，那时他并非从母亲发出；当他从母亲出来进入光天化日时，那时他并非从父亲发出。但圣神并不是从圣父发出进入圣子，然后从圣子发出为受造物的圣化；祂同时从两者发出：虽然圣父已赐给圣子，使祂也从圣子发出，如同从祂自己发出一样。我们也同样不能说圣神不是生命，既然圣父是生命，圣子也是生命。正如圣父自己内有生命，也赐给圣子自己内有生命一样，祂也赐给了生命从祂发出，如同从祂自己发出一样。现在我们来到主随后的话，祂说：\Quote{祂将要告诉你们未来的事。祂要光荣我，因为祂要从我领受，并告诉你们。凡父所有的，都是我的；所以我曾说，祂要从我领受，并告诉你们。}但由于当前的论述已经拉得相当长，这些话必须留待另一次。\par

\SourceRecord{Translated by John Gibb.；Nicene and Post-Nicene Fathers, First Series；Vol. 7.；Edited by Philip Schaff.；Buffalo, NY: Christian Literature Publishing Co.,；1888.；<http://www.newadvent.org/fathers/1701099.htm>.}

% structure: p0001 source=h1 target=section reason=page-title

\section{讲道100（\BibleRef{若 16:13-15}）}

1. 当我们的主应许圣神将要来临时，祂说：\Quote{祂要教导你们一切真理}，或者，如我们在一些抄本中读到的：\Quote{祂要引导你们进入一切真理。因为祂不是凭自己说话，而是凡祂所听见的，祂就讲说。}关于这些福音的话，我们已按主所赐的能力讲论过了；现在请留意接下来的一段。祂又说：\Quote{祂要显示给你们将要来的事。}对于这一节，意思非常清楚，无需多留；因为它不包含任何需要正式讲解的问题。但祂接着说的话：\Quote{祂要显扬我，因为祂要从我领受，并显示给你们。}不可轻易放过。因为\Quote{祂要显扬我}这句话，我们可以理解为：圣神把天主的爱倾注在信徒心中，使他们成为属神的人，从而向他们显明圣子与圣父平等——他们先前只按肉体认识祂，且自认为祂仅仅是凡人。或者至少，他们因这爱而充满勇气，摆脱一切恐惧，向人宣讲基督；这样，祂的名声便传遍普世。因此，祂说\Quote{祂要显扬我}，意思是：祂要使你们脱离恐惧，赐给你们一种爱，这爱将点燃你们宣讲我的热情，以致你们将我的荣耀的芬芳传遍世界，令人称颂。因为圣神自己也要做门徒们藉圣神所做的事，正如这话所暗示的：\Quote{因为不是你们说话，而是你们父的圣神在你们内说话。}\BibleRef{玛 10:20} 希腊文动词\OriginalTerm{显扬}{\GreekText{δοξάσει}}，拉丁文译者分别译作\OriginalTerm{显扬}{clarificabit}和\OriginalTerm{光荣化}{glorificabit}：因为希腊文\OriginalTerm{光荣}{\GreekText{δόξα}}（由此衍生出动词\OriginalTerm{显扬}{\GreekText{δοξάσει}}）既可解作\OriginalTerm{光明}{claritas}，也可解作\OriginalTerm{光荣}{gloria}。因为光荣使人光明，光明使人光荣；因此，这两个词所指的是同一件事。正如古时最著名的拉丁作家所定义的：光荣是伴随着赞美的、普遍流传且被接受的关于某人的名声。当此事在世上关于基督发生时，我们不应以为这是赐予基督什么大事，而是赐予世界。因为赞美善事，并非使受赞者受益，而是使赞美者受益。\par

2. 但也存在虚假的光荣，即赞美出于错误，无论是对事物、对人物，或对两者皆错。对事物的错误，是指人以为恶是善；对人物的错误，是指人以为恶人是善人；对两者的错误，是指事实上是恶习却被当作美德，并且被赞美者并不具备人们所认为的优点——无论他是善是恶。将虚荣者所宣称的归于他们，这当然是大恶而非美德；然而你知道这类人的赞美之声多么常见，正如圣经所说：\BibleQuote{罪人在他灵魂的愿望中受赞美，行不义的人受祝福。}这里赞美者并非错认了人，而是错认了事；因为他们相信恶是善。那些道德上被浪费之恶败坏的人，无疑就是赞美者所并非怀疑、而是亲眼所见的人。再者，若有人假装为义人，其实不然，并且他一切在人前可称赞的行为并非为天主而做（即为真义），而只是寻求和贪图人的光荣；而那些传播他美名的人，以为他是为天主而活——他们不是错认了事，而是被骗了人。因为他们相信是善的事，确实是善；但他们相信是善的人，却是相反。然而，若将巫术视为善，而某人因被认为用这些他完全不懂的巫术拯救了国家，在不敬神的人中得到那种普遍流传的名声（即所谓光荣），那么那些赞美的人就在两方面都错了：既错在事（以为恶为善），也错在人（他根本不是他们所想的那样）。但若一个人因天主的恩宠并为天主而成为义人（即真正的义人），因这义而得到普遍的赞美名声，那么这光荣确实是真实的；然而我们不应以为这义人因此而蒙福，反而应祝贺那些赞美者，因为他们判断正确，并爱义人。那么，主基督因自己的光荣，岂不更是造福了别人，而非自己？连祂的死也造福了他人。\par

3. 但祂在异端者中所享有的并非真光荣，尽管在那里祂似乎拥有一种伴随着赞美的普遍声望。这种光荣不是真的，因为他们在两方面都错了：他们既把不善的当作善，又以为基督是基督所不是的。因为说独生子与生祂者不相等，是不善；说天主的独生子仅是凡人而非天主，是不善；说真理的肉不是真肉，是不善。在我所陈述的三种学说中，第一种为亚略派所持守，第二种为福提诺派，第三种为摩尼派。但由于在它们任何一派中都没有善，且基督与它们毫无关系，因此他们在两方面都是错的；他们并未将真光荣归于基督，尽管在他们中间似乎有一种对基督的赞美性的普遍声望。因此，所有对基督没有正确见解的异端者（若要一一列举就太冗长了），之所以在这件事上犯错，是因为他们对善与恶的看法都不真实。此外，外邦人中也有许多赞美基督者，他们自身也在两方面都错了，因为他们不是按照天主的真理，而是按照自己的臆测，说基督是个术士。他们指责基督徒缺乏技能，却赞美基督为术士，从而暴露了他们所爱的是什么：他们爱的确实不是基督，因为他们所爱的是基督从未成为的东西。因此，他们在两方面都错了，因为做术士是邪恶的，而基督是善的，祂并非术士。所以，既然此处我们不必谈论那些辱骂和亵渎基督的人——因为我们所说的是祂在世上所受的光荣——只有在圣而公教会中，圣神才以真光荣光荣了祂。因为在别处，即或在异端者或某些外邦人中，祂在世上所享的光荣不可能是真的，即使那里有伴随着赞美的普遍声望。因此，祂在教会中的真光荣由先知以下的话所颂扬：\BibleQuote{天主，愿祢被尊高于天，愿祢的光荣超越全地。}所以，在祂被举扬之后，圣神将来到并光荣祂，这是圣咏和独生子自己所应许的未来事件，我们见其已成就。\par

4. 但当祂说：\Quote{祂要从我所领受的，并传告给你们}时，你们要怀着公教会的耳朵聆听，并以公教会的心智领受。因为肯定不是如某些异端者所想象的那样，圣神低于圣子，仿佛圣子从父领受，圣神从圣子领受，是按照某种本性的等级。我们决不信、不说，也不从基督徒心中想这样的事。总之，祂自己立刻解决了这个问题，并解释祂为何这样说：\Quote{凡父所有的都是我的；为此我说，祂要从我所领受的，并传告给你们。}你们还要什么呢？因此，圣神从父领受，圣子也从父领受；因为在 天主圣三中，圣子由父所生，圣神由父所发。唯独父不是由谁所生，也不是由谁所发。但是，独生子所说\Quote{凡父所有的都是我的}（这当然不同于对那个不仅为独生、且是兄弟中较年长之子所说的\Quote{你常同我在一起，凡我所有的都是你的}），\BibleRef{路 15:31} 将得到我们仔细的考虑，如果主愿意，我们将在联系独生子对父说\Quote{凡我所有的都是祢的，凡祢所有的都是我的}那段经文处进行，以便我们现在的讨论在此结束，因为接下来的话需要另一个开端来讨论。\par

\SourceRecord{Translated by John Gibb.；Nicene and Post-Nicene Fathers, First Series；Vol. 7.；Edited by Philip Schaff.；Buffalo, NY: Christian Literature Publishing Co.,；1888.；<http://www.newadvent.org/fathers/1701100.htm>.}

% structure: p0001 source=h1 target=section reason=page-title

\section{讲道辞101（\BibleBook{若望}{福音} 16:16-23）}

主说这些话时——\Quote{片刻之后，你们就再也看不见我了；再过片刻，你们又要看见我；因为我往父那里去}——在祂所说的这句话应验之前，对门徒们来说是如此隐晦，以至于他们彼此询问祂所说的是什么，并不得不承认自己完全不明白。因为福音接着说：\Quote{于是祂的几个门徒彼此说：\InnerQuote{祂给我们所说的\InnerQuote{片刻之后，你们就再也看不见我了；再过片刻，你们又要看见我；并且，因为我往父那里去}，这是什么意思？}他们便说：\InnerQuote{祂所说的\InnerQuote{片刻之后}是什么意思？我们不晓得祂在说什么。}} 使他们困惑的就是祂所说的：\Quote{片刻之后，你们就再也看不见我了；再过片刻，你们又要看见我。} 因为在前面，祂没有说\Quote{片刻之后}，而只说\Quote{我往父那里去，你们就再也看不见我了}，这在祂们看来似乎说得相当明白，所以他们对此没有彼此询问。但现在，那时对他们隐晦、不久后便启示给我们的，如今已经向我们完全显明了：因为再过片刻，祂受难，他们便不见祂；再过片刻，祂复活，他们便见祂。但祂所说的话\Quote{你们就再也看不见我了}，通过\Quote{再}这个词，祂希望让人明白他们此后将不再见祂，这一点我们在祂所说之处已经解释：圣神\BibleQuote{要指证关于正义，因为我往父那里去，你们就再也看不见我了}，以此指明，他们此后将永不再见基督处于祂当时那受死亡束缚的状态。\par

\Quote{现在耶稣知道（正如福音作者接着说的）他们渴望问祂，便对他们说：你们彼此议论我所说的\InnerQuote{片刻之后，你们就看不见我了；再过片刻，你们又要看见我}。我实实在在告诉你们：你们将要痛哭哀号，世界却要喜乐；你们将要忧愁，但你们的忧愁要变为喜乐。} 这可以这样理解：门徒们因主的死而陷入忧愁，随即因祂的复活而充满喜乐；但世界——即指那些杀害基督的敌人——当然在基督被杀时欣喜若狂，而门徒们却正充满忧愁。因为藉着世界之名，可以理解为这个世界的邪恶，换言之，就是那些属于这个世界的人。正如宗徒雅各伯在他的书信中所说：\Quote{谁若愿意作这个世界的朋友，就成了天主的仇敌；} \BibleRef{雅 4:4} 那对天主的仇恨的结果是，连祂的独生子也没有被放过。\par

祂接着说：\Quote{妇人生产的时候有忧愁，因为她的时辰来到了；但一旦生了孩子，她就不再记得那痛苦，因为喜乐一个男人生到世上。你们现在因此也有忧愁；但我要再见你们，你们的心将喜乐，并且你们的喜乐无人能夺去。} 这里所用的比喻似乎不难理解；因为祂自己所给的解释提供了理解的关键。因为生产的阵痛被比作忧愁，而生产本身被比作喜乐；通常当生的是男孩而不是女孩时，喜乐更大。但祂说\Quote{你们的喜乐无人能夺去}，因为他们的喜乐就是耶稣自己，这隐含了宗徒所说的：\BibleQuote{基督既从死者中复活，就不再死；死亡不再统治祂。} \BibleRef{罗 6:9}\par

到目前为止，我们今天所讨论的这段福音中，一切内容，我可以说，都是容易理解的；对于接下来的话，则需要更加密切的注意。因为祂所说的\Quote{在那一天，你们什么也不必问我}是什么意思？这里使用的动词\Emph{问}不仅指\Emph{求}，也指\Emph{询问}；其所译自的希腊文福音中有一个词也可作两种理解，因此歧义并未消除；即使如此，也并非所有困难都会消失。因为我们读到，主基督复活之后，既被询问也被祈求。在祂升天前夕，门徒们问祂何时显明自己，以色列国何时来临；\BibleRef{宗 1:6} 甚至在祂已经在天上时，\Person{圣斯德望}祈求祂接收自己的灵魂。\BibleRef{宗 7:59} 谁敢想或敢说，基督坐在天上时不应被祈求，而祂在地上居住时却被祈求？或说，在祂不朽的状态下不应被祈求，尽管祂仍处于死亡束缚时，人们有义务祈求祂？不，亲爱的，让我们祈求祂亲手解开我们当前疑问的结，藉着光照我们的内心，使我们能领悟祂所说的话。\par

5. 因为我认为，祂的这句话：\BibleQuote{但我将再见到你们，你们的心要喜乐，并且你们的喜乐无人能从你们夺去}，不是指祂复活的时候，也不是指祂将祂的肉身显示给他们看并触摸的时候；而是更指向祂已经说过的那句话：\BibleQuote{那爱我的，必蒙我父所爱；我也要爱他，并要将我自己显示给他}。因为祂已经复活了，已经以肉身向他们显现了，并且已经坐在圣父的右边，而此时那同一位宗徒若望，即这部\WorkTitle{福音}的作者，在他的\WorkTitle{书信}中说：\BibleQuote{亲爱的，现在我们已是天主的子女；将来如何，还未显明；但我们知道，一旦祂显现出来，我们必要相似祂，因为我们要看见祂实在怎样。}\BibleRef{若一 3:2}那景象不属于此生，而属于来世；它不是暂时的，而是永恒的。这就是永生，那生命本身说：\BibleQuote{认识祢，唯一的真天主，和祢所派遣的耶稣基督。}关于这景象和知识，宗徒说：\BibleQuote{现在我们借着镜子，\OriginalTerm{在谜语中}{in a riddle}观看，但那时就要面对面了；现在我所认识的是局部的，但那时我将认识，如同我全然被认识一样。}\BibleRef{格前 13:12}目前教会正怀着热望，期盼她一切劳苦的果实，但那时她将在实际默观中生产；现在她呻吟着生产，那时她将在喜乐中生产；现在她通过祈祷生产，那时她将在赞美中生产。因此，它也是男婴；因为当前行为的全部义务都指向默观中的这种果实。因为唯独祂是自由的；因为祂是因自身而被渴望，而非为了其他任何事物。这样的行为是在服事祂；因为凡出于善意的行为都与祂有关，因为是为祂而做的；而另一方面，祂是因自身而被获得和拥有，而非为了其他任何事物。因此，我们在那里找到唯一能满足我们自身的终向。祂将因此是永恒的；因为除非在无终者身上，否则没有终能满足我们。\Person{斐理伯}被这所感动，说：\BibleQuote{请把圣父显示给我们，我们就满足了。}而在那显示中，圣子也承诺了祂自己的临在，说：\BibleQuote{你不信我在圣父内，圣父在我内吗？}因此，关于那唯独满足我们的事，我们被很恰当地告知：\BibleQuote{你们的喜乐无人能从你们夺去。}\par

6. 关于这一点，参照上述的话，我想我们也能更好地理解这句话：\BibleQuote{片刻，你们就不再看见我；再片刻，你们又要看见我。}因为现世这一恩典时期所涵盖的全部时间，不过是一片刻；因此，这位同一位福音作者在他的书信中说：\BibleQuote{是\OriginalTerm{最后的时辰}{the last hour}}\BibleRef{若一 2:18}。因为在这种意义上，祂也加上了：\BibleQuote{因为我到圣父那里去}，这应指向前面的语句，即祂说：\BibleQuote{片刻，你们就不再看见我}；而不是后面的语句，即祂说：\BibleQuote{再片刻，你们又要看见我}。因为祂到圣父那里去，将导致他们不再看见祂。因此，祂的话并非指祂即将死亡，被从他们的视线中隐去直到复活；而是指祂即将到圣父那里去，这是祂在复活之后，与他们交往四十天，然后升天时所行的。所以祂对当时在肉身中看见祂的人说了这句话：\BibleQuote{片刻，你们就不再看见我}；因为祂即将到圣父那里去，此后永远不再以他们当时目睹祂的那种可朽状态显现。但是祂加上的一句话：\BibleQuote{再片刻，你们又要看见我}，则是作为应许给予普世教会的：正如祂也赐给教会另一个应许：\BibleQuote{看，我同你们天天在一起，直到今世的终结。}\BibleRef{玛 28:20}主并不迟延祂的应许：片刻之后，我们将看见祂，在那里我们将不再有任何请求，不再提出任何问题；因为没有什么可渴望的，没有什么隐藏的可以询问。这片刻在我们看来很长，因为它仍在延续；当它结束，我们将感受到它是多么短暂。因此，不要让我们的喜乐像世界的喜乐，关于这世界曾说过：\BibleQuote{但是世界要喜乐}；也不要在怀着这种渴望生产时，让我们的忧愁毫无喜乐；但正如宗徒所说，要\BibleQuote{在希望中喜乐，在磨难中忍耐}\BibleRef{罗 12:12}；因为就连我们比作的那位临产妇女，她对于即将诞生的婴儿的喜乐，也大于对当前痛苦的忧愁。但让我们在此结束当前的讨论，因为接下来的话包含一个非常棘手的问题，不应不恰当地缩短，以便，主若愿意，它们能获得更合适的解释。\par

\SourceRecord{Translated by John Gibb.；Nicene and Post-Nicene Fathers, First Series；Vol. 7.；Edited by Philip Schaff.；Buffalo, NY: Christian Literature Publishing Co.,；1888.；<http://www.newadvent.org/fathers/1701101.htm>.}

% structure: p0001 source=h1 target=section reason=page-title

\section{论道 102（\BibleBook{若望}{福音} 16:23-28）}

1. 我们现在要思考主的这些话：\Quote{我实实在在告诉你们：你们若因我的名向父求什么，他必赐给你们。} 在吾主此篇讲道的前面部分，已针对那些因基督的名向父求某些事物却未得着的人说过：凡是违反救恩之道的祈求，就不是因救主之名而求。因为他说的\Quote{因我的名}，所宣告的不是字母和音节的声音，而是这声音本身所传达的，以及由这声音应正确而真实领悟的。因此，凡是对基督持有不应当对天主的独生子持有的想法的，即使他嘴上一再提到基督，也并非因他的名祈求；因为只有当他实际所想的是某一位时，他才是在因那一位的名祈求。但若他对基督持有应当持有的想法，他就是因他的名祈求，并且他所求的，只要不违背他永久的救恩，就必得着。而且他会在应当得着的时候得着。因为有些事物并非被拒绝，而是被延缓，等到合宜的时候才赐予。我们当然应当这样理解他的话：\Quote{他必赐给你们}，好让我们明白，所指的是那些特别适用于祈求者的恩惠。因为所有圣人在关乎自身的事上，其祈求都蒙垂听，但并非在所有关于他人（无论是朋友、仇敌或其他任何人）的事上都如此蒙垂听；因为他的话不是笼统地说\Quote{他必赐给}，而是说\Quote{他必赐给你们}。\par

2. 他说：\Quote{直到现在，你们没有因我的名求什么；求吧，你们将得到，以使你们的喜乐得以圆满。} 他所说的圆满喜乐当然不是肉身的喜乐，而是属神的喜乐；当这喜乐大到不能再增加时，就无疑圆满无缺了。因此，凡是属于达至这喜乐所祈求的，就应当因基督的名祈求，只要我们理解天主的恩宠，并真正追求幸福的生命。若祈求与此不同的事物，那就等于没有祈求：不是那事物本身毫无价值，而是与如此伟大的事物相比，任何其他渴望的事物都几乎等于虚无。因为宗徒说\BibleQuote{人若自以为有什么，其实什么也没有}，\BibleRef{迦 6:3} 那个人当然并非绝对虚无。但诚然，与那靠天主的恩宠认识自己真实身份的属神之人相比，怀着虚妄自大的人算不得什么。这样，我们也可以正确理解这话：\Quote{我实实在在告诉你们：你们若因我的名向父求什么，他必赐给你们；}即，通过\Quote{若求什么}这句话，不应理解为无论什么，而是与幸福生命并非真正无关的事物。接着的话\Quote{直到现在，你们没有因我的名求什么}可以有双重理解：要么，你们没有因我的名祈求，因为这名你们还不认识，尚未如将来所要认识的那样；要么，你们什么也没有求，因为与你们应当祈求的事物相比，你们所祈求的算不得什么。因此，为了使他们因他的名祈求，不是求那虚无，而是求圆满的喜乐（因为任何与此不同的事物他们求的实际上都是虚无），他劝勉他们说：\Quote{求吧，你们将得到，以使你们的喜乐得以圆满}；即，因我的名求这个，使你们的喜乐圆满，你们就必得着。因为他的圣人，凡恒心祈求如此美善之事的，绝不会被天主的仁慈所亏负。\par

3. 他说：\Quote{这些事我以比喻对你们说了；但时候将到，我不再用比喻对你们说话，而要明明地告诉你们关于父的事。} 我或许会说，他所说的这个\Quote{时候}必须理解为将来我们得以明明看见的那段时期，正如蒙福的宗徒所说：\BibleQuote{面对面}，\BibleRef{格前 13:12}；而他说的\Quote{这些事我以比喻对你们说了}，与同一宗徒所说的\BibleQuote{我们现在藉着镜子观看，如同猜谜}，\BibleRef{格前 13:12} 是一致的；而\Quote{我要告诉你们}，因为父将藉着子被看见，类似于他别处所说的：\BibleQuote{除了子，没有人认识父，除非子愿意向他启示的人。} \BibleRef{玛 11:27} 但这种理解似乎被后面的话所妨碍：\BibleQuote{在那一天，你们要因我的名祈求。} 因为在那未来的世界，当我们到达那国度，我们将相似他，因为我们要看见他实在怎样，\BibleRef{若一 3:2} 那时我们还有什么可祈求的呢？因为我们的欲望已饱飨美物？正如在另一篇圣咏中也说：\BibleQuote{当我得见你的光荣时，我便得到满足。} 因为祈求涉及某种需要，而在那丰盛掌权的境界，是没有需要之处的。\par

4. 因此，按我的理解能力所及，我们当理解耶稣应许要使门徒从属血肉和属本性的变为属神的；虽然还不是将来我们拥有属神身体时的那样，而是像那说\Quote{我们在成全的人中讲智慧}\BibleRef{格前 2:6}以及\Quote{我不能把你们当作属神的人，只能当作属血肉的人}\BibleRef{格前 3:1}以及\Quote{我们所领受的，不是这世界的神，而是出于天主的圣神，为使我们知道天主白白赐给我们的事。我们也讲论这些事，不是用人的智慧所教的言语，而是用圣神所教的言语，以属神的事解释属神的事。然而属本性的人不能领受天主圣神的事。}的人。\newline{}如此，属本性的人既不能领受天主圣神的事，无论怎样听到论天主本性的事，他也只能设想某种形体，无论多么广袤或巨大，多么光辉灿烂，总归是形体；因此，他将凡论及智慧的非物质不变本体的言论都当作箴言；并非他认为它们是箴言，而是他的思想追随那些惯于听箴言却不理解的人同一方向。但当属神的人开始辨别一切，而他自身却不受任何人辨别时，他——即使在此生仍是通过镜子观看，模糊不清——不是凭任何肉体的感官，也不是凭任何抓住或捏造各种形体相似物的想象概念，而是凭心灵最清晰的理解，认识到天主不是物质的，而是属神的；圣子如此将圣父显明给我们，以致那显明者自己也显为同一本体。那时，凡祈求的，是因祂的名祈求；因为在那名号的声音中，他们只理解那名所称呼的实在，且不因心灵的虚妄或软弱，心存幻想，以为圣父在一处，圣子在另一处，站在圣父面前为我们代求，以为二者各有物质本体占据各自位置，圣言以言语向祂所是的圣父为我们代求，而说话者口与听者耳之间有一定的距离相隔；以及其他诸如此类的荒谬，系属本性且同时属血肉的人在自己心中所捏造的。因为属神的人想到天主时，任何由肉体习惯经验所提示的这类事物，他们都会否认、拒绝，像驱赶恼人的昆虫一样从心灵的眼目驱走；而将自己委身于那光的纯洁，藉其见证和判断，他们证明那些侵入他们内心视觉的肉体形象完全是虚假的。他们能在一定程度上思想我们的主耶稣基督，就祂的人性而论，是代表我们向圣父发言；但就祂的天主性而论，是与圣父一同聆听并回应我们。我认为祂说\Quote{我不说我要为你们祈求圣父}时，所指的就是这点。但对这点的直觉领悟，即圣子如何不祈求圣父，而圣父和圣子同样垂听祈求者，是只有心灵的神目才能达到的高度。\par

5. 祂说：\Quote{因为圣父自己爱你们，因为你们爱了我。}那么，是祂爱，因为我们爱；还是我们爱，因为祂爱？让同一位福音作者从他自己的书信中回答我们吧：\Quote{我们爱祂，因为祂先爱了我们。}\BibleRef{若一 4:19}因此，我们爱的动力原因，是我们被爱。爱天主确实是天主的恩赐。那爱的，祂是那在我们尚不被爱时赐予爱祂恩宠的一位。当我们不中悦祂时，祂已经爱了我们，为使我们里面能有变成中悦祂的东西。因为我们不能爱圣子，除非也爱圣父。圣父爱我们，因为我们爱圣子；因为我们从圣父和圣子领受了爱圣父和圣子的能力：因为爱藉着两者的圣神倾注在我们心中（\BibleRef{罗 5:5}），藉着那圣神我们爱圣父和圣子，并和圣父圣子一起爱那圣神。因此，是天主在我们内成就了这宗教性的爱，藉以朝拜天主；祂看见这爱是好的，因此祂爱了祂所造的。但若祂在造作之前不爱我们，祂就不会在我们内造作祂所能爱的。\par

6. 祂又说：\Quote{你们已经相信我是从天主出来的。我从圣父出来，来到世界；我又离开世界，往圣父那里去。}我们确实已经相信了。因为，因祂这样来到世界，从圣父出来而不离开圣父；又因祂这样离开世界，往圣父那里去而不实际遗弃世界，这不应视为不可信之事。因为祂从圣父出来，是因祂出自圣父；祂来到世界，是向世界显示祂从童贞玛利亚所取的形体。祂藉身体的撤离离开了世界，祂作为人藉升天往圣父那里去，但祂并未在祂临在的治权活动中遗弃世界。\par

\SourceRecord{Translated by John Gibb.；Nicene and Post-Nicene Fathers, First Series；Vol. 7.；Edited by Philip Schaff.；Buffalo, NY: Christian Literature Publishing Co.,；1888.；<http://www.newadvent.org/fathers/1701102.htm>.}

% structure: p0001 source=h1 target=section reason=page-title

\section{论道103（\BibleRef{若 16:29-33}）}

1. 基督门徒的内心状态，在主受难前与祂交谈，如同与孩童谈论大事，但以适合向孩童讲述大事的方式时，因他们尚未领受圣神（如同祂复活后，或亲自向他们嘘气，或由天上降下时所领受的），他们的心智更倾向于人性而非神性，——这在福音中到处有众多证言显明；眼前的经文所说的，也与此一致。因为，福音记载说：\BibleQuote{门徒对祂说：看，如今你明明白白地讲论，不说任何比喻了。现在我们确信你晓得一切，也不需要任何人问你：由此我们相信你是从天主而来的。}主自己不久之前曾说过：\BibleQuote{这些事我以比喻对你们讲了；时候来到，我不再以比喻对你们讲论了。}那么，他们怎么又说：\BibleQuote{看，如今你明明白白地讲论，不说任何比喻了}呢？难道那祂应许不再以比喻讲话的时候果真已经来了吗？当然还没有，这由祂随后的话可以证明，祂这样说：\BibleQuote{这些事我以比喻对你们讲了；时候来到，我不再以比喻对你们讲论了，但要明明白白地把圣父的事告诉你们。在那一天，你们要因我的名祈求；我不对你们说，我要为你们求圣父，因为圣父自己爱你们，因为你们爱了我，并且相信我是由天主而来的。我由圣父出来，来到了世界；我又离开世界，往圣父那里去。}（\BibleRef{若 16:25}）既然在所有这些话中，祂仍应许那个不再以比喻讲论，而要明明白白地告诉他们关于圣父的事的时候；祂说那时他们要因祂的名祈求，且祂不为他们求圣父，因为圣父自己爱他们，而他们也爱了基督，并相信祂是由圣父而来，来到世界，又将离开世界往圣父那里去：既然那个祂应许明明白白讲论的时候仍是未来的事，为什么他们说\BibleQuote{看，如今你明明白白地讲论，不说任何比喻了}呢？只因为那些祂知道对没有领悟力的人是比喻的事，他们仍远未能理解，甚至不懂他们自己不懂！因为他们是婴孩，还没有属灵的辨识力去领悟所听的关于非属身体、而是属神的事。\par

2. 当他们仍因内心的人而年幼体弱时，\Person{耶稣}进一步告诫他们：\BibleQuote{\Person{耶稣}回答他们说：现在你们相信吗？看，时辰要来，而且已经到了，你们要各自散到自己的地方，留下我独自一人。但我不是独自一人，因为\Person{圣父}与我同在。}他刚刚说过：\BibleQuote{我离开世界，到\Person{圣父}那里去；}现在他说：\BibleQuote{\Person{圣父}与我同在。}谁到与他同在的那一位那里去呢？这对理解的人是言语，对不理解的人是箴言；然而，即使目前婴孩不能理解，也在某种程度上吸收了它；即使它不能给他们提供他们尚不能接受的固体食物，至少也不拒绝给他们乳类的饮食。正是从这种饮食中，他们得知他知道一切，不需要任何人问他；而且，他们为什么这样说，确实是一个值得探讨的话题。因为人们会认为他们更应该说：你不需要问任何人；而不是任何人问你。他们刚才说：我们确信你知道一切；当然，知道一切的人通常是被那些不知道的人提问，以便他们通过提问从知道一切的那一位那里听到他们想知道的事；而不是他自己提问，好像想知道什么，而他明明知道一切。那么，我们应该如何理解这一点：当他们显然应该对这位他们知道全知的\Person{主}说你不需要问任何人时，他们却认为更合适说\BibleQuote{你不需要任何人问你}呢？是的，难道我们不是读到两者都发生过吗？即\Person{主}既问过问题，也被问过问题？但这很快就能回答：因为这对他不是必要的，而是对他所提问的那些人，或对他提问的那些人。因为他从不向任何人提问以图从他们那里学到什么，而是为了教导他们。而对那些向他请教的人，可以肯定他们需要从他那里了解一些事。同样，毫无疑问，也正是因为这个原因，他不需要任何人问他。就像我们，当被那些想从我们这里获得信息的人提问时，正是通过他们的提问才发现他们想知道什么，因此我们需要被那些我们想教导的人提问，以便了解他们需要回答的问题；但他，知道一切的那一位，甚至连这个也不需要，他同样不需要通过他们的问题来发现有人想从他那里知道什么，因为在问题提出之前，他已经知道提问者的意图。但他允许自己被提问，是为了向当时在场的人或将要听到或读到这些话的人表明，向他提问的那些人的特征；这样我们就可以知道那些无法欺骗他的欺诈，以及那些在他眼中对我们有益的接近方式。但是，预见人的思想，因而不需要任何人问他，对\Person{天主}来说不算什么大事，但对那些对他说\BibleQuote{由此我们相信你从天主而来}的婴孩来说，却是够大的了。有一件更大的事，为了理解它，他希望他们的心智得以扩展和扩大：当他们说，并且说得对：\BibleQuote{你从天主而来，}他回答说：\BibleQuote{\Person{圣父}与我同在；}以便他们不会认为\Person{圣子}以某种方式从\Person{圣父}出来，以致他们以为他也离开了\Person{圣父}的同在。\par

3. 然後，在結束這篇沉重而冗長的談話時，祂說：\Quote{我對你們說了這些事，是要使你們在我內有平安。在世界上你們要受苦難，但你們放心，我已戰勝了世界。}這種患難的開端，見於祂先前所說的話，是為了顯示他們如同嬰孩，他們仍然缺乏智慧，易將一事誤認作另一事，以至於祂所說的一切偉大而神聖的事，對他們而言與比喻相差無幾。祂先前說：\Quote{現在你們信嗎？看，時辰要來，而且已經到了，你們要四散，各人歸自己的地方去。}我説，這就是患難的開端，但並非以他們同樣的恆心為尺度。因為祂加上說：\Quote{你們也要撇下我獨自一人，}祂並非意指他們在以後的患難中（即祂升天後他們要在世上忍受的）會有那樣的本性，以至於離棄祂；而是說他們仍住在祂內，在祂內應有平安。但在祂被捕時，他們不僅外在拋棄了祂的肉身臨在，而且也內心拋棄了他們的信德。而祂的話正是指此而說：\Quote{現在你們信嗎？看，時辰要來，你們要四散回自己的地方去，撇下我：}好像祂說：那時你們會如此驚慌失措，以至於連你們現在所信的都要拋棄。因為他們陷入了如此的絕望，可以説，他們舊日的信德已經死了，這在\Person{克羅帕}身上很清楚；他在耶穌復活後，不知道自己在和耶穌說話，敘述祂所遭遇的事，說：\Quote{我們原指望祂就是要拯救以色列的那一位。}\BibleRef{路 24:21}他們當時就是這樣離棄了祂，甚至拋棄了他們從前信祂的那種信德。但在那患難中——在祂受光榮之後、他們自己領受了聖神之後所遭遇的——他們沒有離棄祂；雖然他們從一個城市逃到另一個城市，卻沒有逃離祂自己；而是為了使他們在世界上受苦難時，能在祂內有平安，他們非但不逃離祂，反而以祂自己為避難所。因為在領受聖神時，他們內裡產生了現在用這句話對他們描述的那種狀態：\Quote{放心，我已戰勝了世界。}他們放心了，並且得勝了。但靠誰呢？除非靠祂。因為如果祂的肢體還要被世界戰勝，那麼祂就沒有戰勝世界。所以宗徒說：\Quote{感謝天主，使我們獲得了勝利；}並立刻加上：\Quote{藉著我們的主耶穌基督：}\BibleRef{格前 15:57}正是藉著那對祂自己的人說\Quote{放心，我已戰勝了世界}的那一位。\par

\SourceRecord{Translated by John Gibb.；Nicene and Post-Nicene Fathers, First Series；Vol. 7.；Edited by Philip Schaff.；Buffalo, NY: Christian Literature Publishing Co.,；1888.；<http://www.newadvent.org/fathers/1701103.htm>.}

% structure: p0001 source=h1 target=section reason=page-title

\section{讲道集104（\BibleRef{若 17:1}）}

在这些话之前——我们现在靠着主的帮助即将以这些话为讲论的主题——耶稣曾说过：\Quote{我给你们讲了这些事，为使你们在我内获得平安。}我们应当认为，这并非指祂紧接着说出的那些话，而是指祂从起初认他们为门徒时起，或者至少是从晚餐后祂开始这个卓越而长篇的讲论时起，对祂们所说的一切话。祂确实给了他们这样一个说话的理由，以至于要么祂曾经对他们说的一切话都可以恰当地归于这个目的，要么特别是这些作为祂最后的话，当祂即将为他们而死时，在出卖祂的人离开他们之后所说的。因为祂给出了祂讲论的原因，即在他们内祂可能有平安，正如我们成为基督徒完全是为了这个原因。这个平安没有时间上的终结，而是本身将成为我们目前所有虔诚意图和行为的终结。为了它，我们领受了祂的圣事；为了它，我们受教于祂的言行；为了它，我们领受了圣神的抵押；为了它，我们相信并盼望祂，并按照祂恩赐的恩宠被祂的爱点燃；我们的一切苦难都因这平安得安慰，也因它脱离一切苦难；为了它，我们勇敢地忍受一切磨难，以便在其中幸福地统治，没有任何磨难。祂恰当地以此结束了祂的话，这些话对门徒来说是比喻，他们当时还不太明白，但后来会明白，当祂赐予他们所应许的圣神时，关于这圣神祂先前说过：\Quote{我还与你们同在的时候，给你们讲了这些事。但那护慰者，就是父因我的名所要派遣来的圣神，祂要教导你们一切，也要使你们想起我所对你们所说的一切。}毫无疑问，那就是祂所应许的时刻，那时祂不再用比喻对他们说话，而是将父明明地显示给他们。因为当圣神启示时，祂的那些话语对那些明白的人不再是比喻。因为当圣神在他们心中说话时，独生子并没有沉默，祂曾说在那时刻祂要明明地向他们显示父，这当然不再是对现在有了理解的他们成为比喻。但这也是这样的话：天主子和圣神如何同时在属神的人心中说话，甚至天主圣三本身永远不可分离地工作，对拥有领悟力的人来说是话语，对没有领悟力的人来说却是比喻。\par

因此，当祂告诉他们祂所说一切的理由，即为使他们在祂内获得平安，同时在世界中却有苦难，并劝勉他们要鼓起勇气，因为祂已战胜了世界；这样结束了祂对他们的讲论后，祂便向父说话，开始祈祷。因为福音作者接着说：\Quote{耶稣说了这些话，便举目向天说：\InnerQuote{父啊，时候到了，请光荣祢的子。}}主、独生子和圣父的同永生者，能以奴仆的形体并从奴仆的形体中，在必要时默默祈祷；但祂以另一种方式表明祂是在向父祈祷，以便祂记得祂仍然是我们的导师。因此，祂为我们献上的祈祷，也使我们知晓；因为如此伟大的导师不仅向门徒发表讲论，而且为他们向父祈祷，这也是教导门徒的一种方式。如果这对当时在场听这些话的人如此，那么对我们这些将要读这些文字的人也是如此。因此，当祂说\Quote{父啊，时候到了，请光荣祢的子}时，祂表明了一切时间和祂行事或受难的每一个时机，都是由那位不受时间约束者所安排；因为那些在时间上各自未来的事物，其有效原因存在于天主的智慧中，其中没有时间的区别。因此，不要以为这个时刻是因任何命运的迫切性而来，而是出于天主的安排。没有天体运行的必然法则将基督的苦难束缚于其时间；因为我们不敢想象星辰能迫使它们的造物主死亡。因此，不是时间驱迫使基督去死，而是基督选择了死亡的时间：祂与父一起也确定了祂由童贞女诞生的时间，父是祂超越时间而生的。根据这真实而有益的教义，宗徒保禄也说：\BibleQuote{但是时期一满，天主就派遣了祂的子来，} \BibleRef{迦 4:4} 天主藉先知说：\BibleQuote{在悦纳的时候我俯允了你，在救恩的日子我帮助了你。} \BibleRef{依 49:8} 而宗徒又说：\BibleQuote{看，如今正是悦纳的时候；看，如今正是救恩的日子。} \BibleRef{格后 6:2} 因此，那位与父一起安排每一时刻的祂可以说：\Quote{父啊，时候到了，} 如同说：\Quote{父啊，我们所共同为人类并为我得以在他们中受光荣而固定的时刻到了，请光荣祢的子，好使祢的子也光荣祢。}\par

3. 圣子受圣父的光荣化，有些人以其为父不怜惜祂，却为众人把祂交出了。\BibleRef{罗 8:32} 但若我们说祂因苦难而受光荣，那么祂因复活所受的光荣岂不更大！因为在祂的苦难中，我们的目光更专注于祂的谦卑而非祂的光荣，正如宗徒所见证的，他说：\Quote{祂谦抑自己，听命至死，且死在十字架上。}接着他论及祂的光荣说：\Quote{为此，天主极其举扬祂，赐给祂一个名字，超越一切名字，为使天上、地上和地下的一切，因耶稣之名，屈膝叩拜，并且一切唇舌承认主耶稣基督是在天主圣父的光荣中。}这就是我们的主耶稣基督的光荣化，从祂的复活开始。因此，祂的谦卑在宗徒的话语中始于他说：\Quote{祂使自己空虚，取了奴仆的形体；}并达到\Quote{死于十字架。}但祂的光荣始于他说：\Quote{为此，天主也举扬了祂；}并延伸至\Quote{是在天主圣父的光荣中。}因为就连名词本身，若查考希腊文抄本的用语（宗徒书信是从希腊文译成拉丁文的），在拉丁文中读作光荣，在希腊文中读作\OriginalTerm{光荣}{\GreekText{δόξα}}；由此我们得到希腊文的动词，即在此处说\OriginalTerm{光荣化}{\GreekText{δόξασον}}（光荣化），拉丁文译作\Quote{\Emph{clarifica}}（使显耀），虽然他也可以说\Quote{\Emph{glorifica}}（光荣化），含义相同。同理，在宗徒的书信中，我们发现\Quote{\Emph{gloria}}的地方，本来也可改用\Quote{\Emph{claritas}}，因为这样做同样能保留含义。但为了不偏离字音，正如\Quote{\Emph{clarificatio}}（使光辉）源于\Quote{\Emph{claritas}}（光辉），同样\Quote{\Emph{glorificatio}}（使光荣）源于\Quote{\Emph{gloria}}（光荣）。那么，为使天主与人之间的中保、人耶稣基督因复活而成为光辉或光荣的，祂先因苦难而谦卑；因为若祂不死，就不会从死者中复活。谦卑是赢得光荣，光荣是谦卑的报酬。然而，这些都是发生在奴仆的形体中；但祂始终具有天主的形体，祂的光荣永远延续：是的，祂的光荣不是过去如此而现在不再，也不是将来才有而目前未存，而是无始无终，永远常存。因此，当祂说：\Quote{父啊，时候到了，光荣祢的圣子，}这应理解为祂说：谦卑种子的时辰到了，不要迟延我光荣的果实。但接着的话是什么意思：\Quote{为使圣子光荣祢？}难道天主圣父也遭受了身体或苦难的屈辱，必须由此被提升至光荣吗？若非如此，圣子又如何光荣祂呢？圣父的永恒光荣既不会因取了人性而显得减少，也不会因神性而增大。但我不会把这样的问题局限在这篇讲道中，也不会因这样的讨论而使讲道变得冗长。\par

\SourceRecord{Translated by John Gibb.；Nicene and Post-Nicene Fathers, First Series；Vol. 7.；Edited by Philip Schaff.；Buffalo, NY: Christian Literature Publishing Co.,；1888.；<http://www.newadvent.org/fathers/1701104.htm>.}

% structure: p0001 source=h1 target=section reason=page-title

\section{第105篇（\BibleBook{若望}{福音} \BibleRef{若 17:1-5}）}

1. 圣子在祂仆人的形体中受圣父光荣，圣父又使祂从死者中复活，并安置在自己右边，此事本身已昭然若揭，基督徒中无人置疑。但祂不仅说：\Quote{圣父，光荣您的圣子}，还加上说：\Quote{使圣子也光荣您}，因此值得探究：圣子如何光荣了圣父？因为圣父永恒的光荣，在任何人形中都未曾减损，就其自身的神性完美而言，也不能增加。的确，圣父的光荣本身既不能减损，也不能增加；但在世人中间无疑更为微小，当时天主仅在犹太被认识；从日出之地到日落之处，孩童们尚未颂扬上主的名。然而，这事藉基督的福音得以成就，即圣父通过圣子为外邦人所认识，所以无疑圣子也光荣了圣父。然而，倘若圣子只是死了，而没有复活，那么祂既不会受圣父光荣，也不会光荣圣父；但现在，祂通过复活受圣父光荣，祂便藉着宣讲祂的复活来光荣圣父。因为这话的次序本身就揭示了这一点：\Quote{光荣}，祂说，\Quote{您的圣子，使圣子也光荣您}；这如同说：再次把我举起，好藉着我，使您被普世认识。\par

2. 然后祂进一步阐明了圣父如何受圣子光荣，说：\Quote{您怎样赐给祂权柄掌管一切血肉，使祂将永生赐给一切您所赐给祂的人。}祂用\Quote{一切血肉}指代每一个人，以部分代表整体；同样，宗徒说：\Quote{所有灵魂都要服从更高的权势}\BibleRef{罗 13:1}时，也是用优越部分指代整个人。因为祂用\Quote{所有灵魂}还指别的什么吗？不就是指每一个人吗？因此，基督由圣父领受掌管一切血肉的权柄，应就祂的人性来理解；因为就祂的神性而言，万有都是藉祂造的，并且天上地下、有形无形的一切都是在祂内受造的。\BibleRef{哥 1:16} 所以，祂说：\Quote{既然您赐给祂权柄掌管一切血肉}，愿您的圣子这样光荣您，也就是说，使您为一切您所赐给祂的血肉所认识。因为您如此赐予，\Quote{使祂将永生赐给一切您所赐给祂的人。}\par

3. 祂接着说：\Quote{永生就是：他们认识您，唯一的真天主，和您所派遣的耶稣基督。}经文的正确顺序是：\Quote{他们认识您和您所派遣的耶稣基督是唯一的真天主。}因此，圣神也被包含在内，因为祂是圣父和圣子的圣神，是二者本质的、同体的爱。因为圣父和圣子不是两位天主，圣父、圣子和圣神也不是三位天主；而是这三位一体的天主圣三本身是唯一真天主。然而，圣父不是圣子，圣子不是圣父，圣神也不是圣父和圣子；因为圣父、圣子和圣神是三个位格，而圣三本身却是一位天主。那么，如果圣子这样光荣您，\Quote{如同您赐给祂权柄掌管一切血肉}，并如此赐予，\Quote{使祂将永生赐给一切您所赐给祂的人}，而\Quote{永生就是认识您}；那么，圣子便是这样光荣您：使您为一切您所赐给祂的人所认识。因此，既然认识天主是永生，我们便是在迈向生命的途中越加前进，越是在这种认识中增长。我们不会在永生中死亡；因为那时没有死亡，对天主的认识将达至圆满。那时将实现天主的圆满光辉，因为那圆满的光荣，希腊文称之为\OriginalTerm{光荣}{\GreekText{δόξα}}。由此而来的是这里所用的\OriginalTerm{光荣化（命令式）}{\GreekText{δόξασον}}，有些拉丁译者将其译为\OriginalTerm{使之光辉}{clarifica}，有些译为\OriginalTerm{使之光荣}{glorifica}。但古人将\Quote{光荣}（人由此被称为光荣者）这样定义：光荣是广为流传的、伴随赞美的名声。如果一个人在被相信的传言中得到赞美，那么当天主本身被看见时，祂将如何被赞美呢？因此经上写道：\Quote{住在祢殿中的人是有福的，他们要永远赞美祢。}那里将有无尽的天主赞美，那时将有对天主的圆满认识；因为圆满认识，所以也有圆满的光辉或光荣化。\par

4. 然而天主首先在此受到光荣，就是当祂藉口传被人认识，并因信徒的信德而被宣讲时。因此，祂说：\BibleQuote{我在地上光荣了祢：我完成了祢交给我做的工。}祂没有说\BibleQuote{祢命令了}，而是说\BibleQuote{祢交给了}；其中明显的恩宠被推荐注意。因为人性即使在独生子身上，岂非领受了什么？难道不是领受了这个，即当它被圣言——万物藉祂而造——假定到祂位格的统一体时，它就不行恶，只行一切善事？但祂如何完成了托付给祂的工，当受难的考验仍存，在考验中祂特别为祂的殉道者提供了他们要效法的榜样？关于此事，宗徒\Person{伯多禄}说：\BibleQuote{基督为我们受苦，给我们留下了榜样，叫我们跟随祂的足迹；}\BibleRef{伯前 2:21} 但祂说已经完成了，岂不是因为祂确确实实知道祂将完成？正如很久以前，在预言中，祂使用过去时态的词语，当祂所说的事情要在许多年后发生时：\BibleQuote{他们刺穿了我的手、我的脚，他们数算了我的所有骨头；}祂没有说\BibleQuote{他们将刺穿}和\BibleQuote{他们将数算}。就在这部福音中祂说：\BibleQuote{凡我从我父所听见的，我都已告诉你们；}后来祂又对这些人宣称：\BibleQuote{我还有许多事要告诉你们，但你们现在不能担当。}因为祂，那藉着确定且不变的缘由预定了所有将要发生之事者，已经做了祂将要做的一切：正如先知论祂所宣告的：\BibleQuote{祂创造了将要存在的事物。}\par

5. 与此类似，祂接着说：\BibleQuote{父啊，现在求祢以祢自身的光荣光荣我，就是我在世界未有之前，同祢所有的光荣。}因为祂前面说过：\BibleQuote{父啊，时候到了，光荣祢的子，使子也光荣祢；}在这词句的安排中，祂表示了父要先受子的光荣，为使子能光荣父。但现在祂说：\BibleQuote{我在地上光荣了祢：我完成了祢交给我做的工；现在求祢光荣我；}仿佛祂自己先光荣了父，然后要求被父光荣。因此我们要明白，祂前面使用了那句话，是按将来之事以及将来的次序说的：\BibleQuote{光荣祢的子，使子也光荣祢；}但现在祂用过去时态说那尚属将来的事，当祂说\BibleQuote{我在地上光荣了祢：我完成了祢交给我做的工}时。然后，当祂说\BibleQuote{父啊，现在求祢以祢自身的光荣光荣我}时，仿佛祂之后要受父的光荣，而祂自己先光荣了父；祂暗示了什么？岂不是暗示，当祂上面说\BibleQuote{我在地上光荣了祢}时，祂那样说是如同祂做了那尚待做的事；但在这里祂要求父去做那使子仍要这样做的事，换言之，父应光荣子，藉此子的光荣，子还要光荣父？总之，如果我们把动词也放在将来时，关联那尚属将来的事，而祂用了过去时代替将来时，句子将不再含糊：仿佛祂曾说：\Quote{我将在世上光荣祢：我将完成祢交给我做的工；父啊，现在求祢以祢自身的光荣光荣我。}这样，这就像祂说\BibleQuote{光荣祢的子，使子也光荣祢}一样明白；而这里也告诉我们同一光荣的方式，那在前面被忽略了；仿佛后者解释了前者，为那些心被触动的人，使他们明白父如何光荣子，尤其是子如何光荣父。因为祂说父在地上由祂自己光荣，而祂自己由父以父自身光荣，祂确实向他们展示了两种光荣的方式。因为祂自己在世上光荣了父，藉着向万民宣讲父；但父以父自身光荣了祂，使祂坐在自己的右边。但正因此，当祂后来论及光荣父时说\BibleQuote{我光荣了祢}，祂宁可把动词放在过去时，为表示在预定中这事已实现，那确定将要发生的事应视为已经发生；即子受父以父光荣之后，也要在地上光荣父。\par

6. 但关于祂自己的光荣，即祂受父光荣的方式，祂更清楚地揭示了这预定，当祂补充说：\BibleQuote{就是我在世界未有之前，同祢所有的光荣。}这些话的正常次序是\BibleQuote{我在世界未有之前，同祢所有的}。祂的话\BibleQuote{现在求祢光荣我}即指此意；也就是说，如同那时，也如现在：那时藉预定，现在藉完成：求祢在世上成就那在世界之前已同祢完成了的事：在它的时机中成就祢在万世之前所决定的事。有人曾想象，这事应如此理解，仿佛那被圣言所摄取的人性变成了圣言，那人被改变成天主；是的，若我们稍加留意他们所提出的见解，仿佛人性在神性中消失了。因为没有人会走到如此地步，说从人性这样的变化中，天主的圣言或者加倍或者增加，以致那原本是一的现在成为二，或者那较小的现在变为较大。因此，如果祂的人性改变并转变为圣言，天主的圣言仍将如祂先前一样伟大，和祂先前一样，那么人性在哪里，如果它没有消失？\par

7. 但对于这种我确实认为不符合真理的看法，我们没有任何理由坚持，如果当圣子说：\BibleQuote{父啊，现在求你以你自己荣耀我，就是在创世以前，我同你所有的荣耀}，我们理解为祂人性的荣耀的\OriginalTerm{预定}{predestination}，即此后从可朽变成与父一同不朽；而且这在创世以前已藉预定完成，正如在适当的时候也在世界内完成。因为如果宗徒论到我们说：\BibleQuote{照祂在创世以前在祂内拣选了我们} \BibleRef{弗 1:4}，那么，如果父在祂内拣选我们作祂的肢体时，也荣耀了我们的元首，有什么理由认为这与真理不符呢？因为我们被拣选的方式与祂被荣耀的方式相同；因为在世界之先，无论是我们，还是天主与人之间的\OriginalTerm{中保}{Mediator}，那\Emph{人}基督耶稣，\BibleRef{弟前 2:5}，都尚未存在。但祂，就其是圣言而言，亲自\BibleQuote{创造了那些尚未存在的事物}，并\BibleQuote{叫那不存在的如同存在} \BibleRef{罗 4:17}，当然，就祂作为天主与人之间中保的人性而言，在创世以前，已被天主父为我们荣耀了，如果我们在祂内被拣选的话。因为宗徒说什么？\BibleQuote{我们知道：天主使一切协助那些爱祂的人，就是那些按祂的旨意蒙召的人，获得益处。因为祂所预知的人，也预定他们与祂儿子的肖像相同，为使祂在许多弟兄中作长子。祂所预定的人，也召叫了他们} \BibleRef{罗 8:28}。\par

8. 但是，我们或许会有些害怕说祂被预定，因为宗徒似乎只在我们被塑造成祂肖像的意义上说了这话。事实上，如果有人忠信地考虑信德的标准，他却否认天主子被预定，而与此同时他又不能否认祂是人。因为正确地说，祂作为天主的圣言、天主与天主同在，在这方面祂不被预定。因为祂怎能被预定呢？祂已经是祂所是的，无始无终，永恒。但是，那尚不存在的东西必须被预定，以便在它适当的时候成为现实，正如它被预定在万世之前如此来临一样。因此，凡否认天主子预定的，也否认祂本人也是人子。但是，为了那些好争论的人，让我们在这件事上也听听宗徒在他的书信开端所说的。因为无论是在他的第一封书信，即致罗马人的书信中，还是在书信的开端，我们都读到：\BibleQuote{保禄，耶稣基督的仆人，蒙召为宗徒，被选拔去传天主的福音，这福音是天主先前藉着祂的先知在圣经上所预许的，论及祂的儿子，祂按肉身是达味的后裔，按圣德的神，藉着从死者中的复活，被预定为具有权能的天主子。} 那么，论到这种预定，祂在世界之前已被荣耀，为要使祂的荣耀藉着从死者中的复活与父同在，祂坐在父的右边。因此，当祂看到这预定的荣耀的时刻已到，为要使已在预定中成就的事也在实际完成中成就，祂在祈祷中说：\BibleQuote{父啊，现在求你以你自己荣耀我，就是我在创世以前同你所有的荣耀。} 好像祂说：我同你所有的荣耀，即我在你的预定中同你所有的荣耀，现在是我也该同你拥有它，即在坐在你右边的时候。但是，因为对这个问题的讨论已经使我们很长久，其余的内容必须在另一篇讲论中加以考虑。\par

\SourceRecord{Translated by John Gibb.；Nicene and Post-Nicene Fathers, First Series；Vol. 7.；Edited by Philip Schaff.；Buffalo, NY: Christian Literature Publishing Co.,；1888.；<http://www.newadvent.org/fathers/1701105.htm>.}

% structure: p0001 source=h1 target=section reason=page-title

\section{讲道集106（若望福音17:6-8）}

1. 在这次讲论中，我们打算依照祂所赐的恩宠，讲论主所说的这些话：\Quote{我已将祢的名显示给祢从世界中所赐给我的人。} 如果祂这话仅指那些与祂共进晚餐的门徒，并在开始祈祷前对他们说了许多话，那么这与祂先前所说的\Quote{澄清}（或如他人所译的\Quote{光荣}）无关，即圣子藉此澄清或光荣圣父。因为让十二个，或更确切地说是十一个必朽的人认识祂，算是什么伟大的光荣，或类似的光荣呢？但是，如果祂说\Quote{我已将祢的名显示给祢从世界中所赐给我的人}时，愿意所有人都被理解，甚至那些将要相信祂的人，都属于祂伟大的教会（这教会尚未由万民组成，且在圣咏中记载：\Quote{我要在伟大的教会中向祢谢恩}），那么这显然就是圣子光荣圣父的那种光荣，即当祂将祂的名显示给万民和如此多世代的人时。祂在此所说\Quote{我已将祢的名显示给祢从世界中所赐给我的人}与祂稍早所说\Quote{我在地上已光荣了祢}\BibleRef{若 17:4}相似，皆用过去时表示未来，因为祂知道这事已预定要完成，因此祂说祂已经做了祂尚需做的事，尽管对未来毫无不确定性。\par

2. 但接下来的话使祂所说的\Quote{我已将祢的名显示给祢从世界中所赐给我的人}更可信是祂为那些已是祂门徒的人说的，而非为所有将要相信祂的人。因为紧接这些话，祂又加上了：\Quote{他们本是祢的，祢将他们赐给了我，他们也遵守了祢的话。现在他们知道，凡祢所赐给我的，都是从祢来的，因为祢所赐给我的话，我已赐给了他们，他们也领受了，也确实知道我是从祢出来的，并且信了祢派遣了我。} 虽然这些话也可以指所有将要来临的信徒，当那时只是盼望的事变成了事实，因为这些话仍指向未来；但我们更倾向于理解祂只是对那些当时是祂门徒的人说了这些话，因为祂随后不久又说了：\Quote{我与他们同在的时候，我因祢的名保全了祢所赐给我的人；我护卫了他们，其中除了丧亡之子以外，没有一个丧亡，为应验经上的话}\BibleRef{若 17:12}，这里指的是出卖祂的犹达斯，因为他是十二宗徒中唯一丧亡的。然后祂又说：\Quote{现在我到祢那里去}，由此明显看出祂的肉身临在，祂说\Quote{我与他们同在的时候，我保全了他们}，仿佛那临在已不再与他们同在。因为祂这样暗示自己的升天即将到来，当祂说\Quote{现在我到祢那里去}，即往父的右边去；祂将在同一肉身临在中审判活人死人，按信德准则和纯正道理；因为从属灵临在来说，祂升天后当然仍与他们同在，并与祂整个在世上的教会同在，直到世界终结。\BibleRef{玛 28:20} 因此，我们不能正确理解祂说\Quote{我与他们同在的时候，我保全了他们}是指谁，除非是指那些相信祂的人，祂已开始用肉身临在保全他们，但现在要离开这种临在，以便藉着属灵临在与父一起保全他们。此后，祂确实也将其他门徒与他们联合，当祂说：\Quote{我不但为他们祈求，也为那些因他们的话而信从我的人祈求。} 这里祂更清楚地表明，祂先前说\Quote{我已将祢的名显示给祢所赐给我的人}时，不是指所有属于祂的人，而仅指当时听祂说话的人。\par

3. 因此，从祂祈祷一开始，当祂\Quote{举目望天说：父啊，时候到了，请光荣祢的子，使子也光荣祢}，一直到随后祂所说的\Quote{父啊，现在求祢在祢自己面前光荣我，使我在世界未有之前同祢所有的光荣}，祂愿意所有门徒都被理解，祂使他们认识父，从而光荣父。因为祂说了\Quote{使子光荣祢}之后，立刻显示了如何成就这事，加上了\Quote{正如祢赐给祂权柄掌管一切有血肉的，使祂将永生赐给祢所赐给祂的人：永生就是认识祢，唯一的真天主，和祢所派遣的耶稣基督。} 因为父不能藉着人所获得的任何知识而被光荣，除非那光荣祂的人也被认识，即被世上的万民所认识。父的光荣并非仅限于宗徒们所显示的那一种，而是藉着所有人而显示，基督作为元首，他们是祂的肢体。因为\Quote{祢赐给祂权柄掌管一切有血肉的，使祂将永生赐给祢所赐给祂的人}这句话不能仅理解为适用于宗徒，而肯定适用于所有因相信祂而获得永生的人。\par

4. 因此，现在我们来看祂关于当时正在听祂说话的那些门徒所说的话。祂说：\BibleQuote{我已将祢的名显示给祢赐给我的人。}难道他们作犹太人时还不认识天主的名字吗？我们所读的\BibleQuote{天主在犹大为人所知；祂的名在以色列为至大}又是什么意思呢？因此，\BibleQuote{我已将祢的名显示给祢从世界中所赐给我、现在正听我话语的这些人：}不是显示祢被称为天主的那个名，而是显示祢被称为我圣父的那个名；若不显示圣子本身，就无法显示这个名。因为天主这名（祂以此被称呼）在未信基督之前，不可能不为整个受造界、从而为每个民族以某种方式所知。真天主的能量是如此的，以致任何理性的受造物，只要运用其理性，就不能完全且彻底地隐藏它。因为除了少数本性极度堕落的人外，全人类都承认天主是这世界的创造者。因此，就祂是这周围天地间可见世界的创造者而言，天主在万民尚未受教于基督信仰之前就已为他们所知。但就祂不可在严重亏负祂自身的情况下与假神同受敬拜这一点而言，天主仅在犹太为人所知。而就祂是那位除去世界之罪的基督的圣父而言，祂的这个名（以前对所有人隐藏）如今显示给圣父亲自从世界中所赐给祂的人。但祂怎能已经显示了呢？如果那时辰尚未到来——祂曾说过那时辰将到，\BibleQuote{到那时我不再以比喻对你们说话，而要明白地把我的父的事告诉你们}——怎么办呢？难道比喻本身包含了这样明白的宣告吗？那么为何说\BibleQuote{我要明明地告诉你们}？正因为\BibleQuote{以比喻}不是\BibleQuote{明明地}。但当它不再隐藏在比喻中，而是以明白的话语说出时，那就无疑是在明明地说了。那么，祂如何已将祂尚未明明宣告的事显示了呢？所以必须这样理解：过去时态被用来代替将来时态，正如另一些话：\BibleQuote{凡从我的父所听见的，我都已经告诉了你们}；这是一件祂尚未做的事，却当作已经做了来说，因为祂知道这事必定成就。\newline{}\par

5. 但\BibleQuote{祢从世界中所赐给我的人}这些话该如何理解呢？因为经上说他们不属于世界。但他们是通过重生而非出生达到这一点的。此外，接下来的\BibleQuote{他们原是祢的，祢将他们赐给了我}又是什么意思呢？难道有一段时间他们只属于圣父，而不属于祂的独生子？难道圣父曾有过任何不与圣子同在的事物吗？当然不是。然而，曾有一段时间天主圣子拥有某种东西，而作为人的同一圣子却不拥有；因为当祂与圣父共同拥有一切时，祂尚未由地上的母亲成为人。因此，当祂说\BibleQuote{他们原是祢的}时，并非天主圣子自我分裂——离开祂，圣父从未拥有任何事物；而是祂习惯将祂拥有的一切权能归于那一位，祂自身由祂而来，那位拥有权能。因为祂由祂而有存在，也由祂而有能力；二者祂一直同有，因为祂从未有存在而无能力。因此，凡是圣父能做的，圣子总能与祂同时做；因为那位从未有存在而无能力的，从未没有圣父，正如圣父从未没有祂。因此，正如圣父是永远全能，圣子也是同永远全能；如果全能，那么必定全有。确切地说，如果我们精确地表达，就是希腊人所谓\OriginalTerm{全能者（持有万有者）}{\GreekText{παντοκράτωρ}}的那个词，我们的作者本不会将其译为\Quote{全能者}，因为\OriginalTerm{全能者（持有万有者）}{\GreekText{παντοκράτωρ}}意为\Quote{持有万有者}，若非他们觉得意思等同的话。那么，永远持有万有者曾有什么是那同永远持有万有者所没有的呢？因此，当祂说\BibleQuote{祢将他们赐给了我}时，祂暗示作为人，祂领受了拥有他们的权能；因为那永远全能者并非永远是人。因此，虽然祂似乎将此归因于圣父——祂从圣父领受他们，因为万有都由祂而来，祂由祂而有存在——但祂也将他们赐给了自己，即作为与圣父同在的天主的基督，将人赐给了基督的人性，而这人性并非与圣父同在。最终，在此说\BibleQuote{他们原是祢的，祢将他们赐给了我}的祂，先前曾对同一批门徒说：\BibleQuote{我从世界中拣选了你们。}因此，要让一切属肉体的思想被粉碎和消灭。圣子说那些人被圣父从世界赐给了祂，而祂在别处对他们说：\BibleQuote{我从世界中拣选了你们。}那些天主圣子与圣父一同从世界中所拣选的人，正是同一圣子作为人从圣父那里从世界中所领受的；因为圣父若不是拣选了他们，就不会将他们赐给圣子。如此，当圣子说\BibleQuote{我从世界中拣选了你们}时，祂并未因此排除圣父，因为他们也同时被圣父拣选；同样，当祂说\BibleQuote{他们原是祢的}时，祂也未因此排除自己，因为他们同样也是圣子的产业。但现在，作为人的同一圣子领受了那些原本不属于祂的人，因为作为天主的祂也领受了原本不属于祂的奴仆形象。\newline{}\par

6. 祂接着说：\Quote{他们已遵守了祢的话：现在他们知道，凡祢所赐给我的一切，都是出于祢；}这就是说，他们知道我是出于祢的。因为圣父在生育那将要拥有一切的那一位时，就赐下了一切。祂说：\Quote{我已将祢赐给我的话，赐给了他们；他们已接受了；}这就是说，他们已理解和持守了它们。因为话被心灵领悟时，就被接受了。祂又加了一句：\Quote{他们确实知道我是出于祢，并且相信是祢派遣了我。}在这最后一句中，我们也必须补上\Quote{确实}；因为当祂说\Quote{他们确实知道}时，祂意图通过加上\Quote{并且他们相信}来加以解释。因此，\Quote{他们确实相信}就是\Quote{他们确实知道}；正如\Quote{我出于祢}等同于\Quote{祢派遣了我}。因此，当祂说\Quote{他们确实知道}时，为避免有人以为这种知识已是通过目睹而非通过信德获得的，祂附加了解释：\Quote{并且他们相信}，这样我们就应补上\Quote{确实}，而将\Quote{他们确实知道}理解为等同于\Quote{他们确实相信}：不是像祂不久前所示的方式，当时祂说：\Quote{现在你们信吗？看，时辰要来，且已经到了，你们要各自散开，留下我独自一人。}\Quote{他们确实相信}，即按理应相信的方式，不受强制，坚定、恒常、刚毅：不再各自散去，留下基督独自一人。的确，那时门徒们尚未成为祂在这里用过去时态所描述的那样子，好像他们已经是那样了；而是以此预先宣告他们将来会成为怎样的人，即当他们领受了圣神之后——圣神按应许将教导他们一切。因为在领受圣神之前，他们如何遵守了祂关于他们所说的话，好像他们已经遵守了，而当他们中为首的那位三次否认了祂 \BibleRef{玛 26:69}，而祂亲口说过否认祂于人前之人的未来命运 \BibleRef{玛 10:33}？因此，正如祂所说，祂已把圣父赐给祂的话赐给了他们；但直到他们以属神的方式接受它们，不是外在用耳朵，而是内里在心中，那时他们才真正接受了它们，因为那时他们真正知道了它们；而他们真正知道了它们，是因为他们真正相信了。\par

7. 然而，何种人的言语足以解释圣父如何将那些话赐给圣子？当然，如果我们设想祂是作为人子接受这些话的，问题似乎容易些。然而，虽然祂如此由童贞玛利亚诞生，谁敢讲述祂在何时、如何学会了它们？因为祂由童贞玛利亚所生的那个诞生本身，谁又敢宣称呢？但若我们以为祂是作为圣父所生、与圣父同永恒的那一位，而接受圣父的这些话，那么让我们排除一切时间性的想法，仿佛祂在拥有它们之前已经存在，然后接受了祂以前未曾拥有的事物；因为天主圣父所赐给天主圣子的一切，祂都是在生育时赐予的。因为圣父赐给圣子那些没有它们祂就不能成为圣子的事物，正如祂赐予祂存在本身一样。否则，祂如何将话语赐给圣言，而圣言正是祂以不可言喻的方式谈论万有的？但现在，关于下文，你们必须将期待推迟到另一篇讲道。\par

\SourceRecord{Translated by John Gibb.；Nicene and Post-Nicene Fathers, First Series；Vol. 7.；Edited by Philip Schaff.；Buffalo, NY: Christian Literature Publishing Co.,；1888.；<http://www.newadvent.org/fathers/1701106.htm>.}

% structure: p0001 source=h1 target=section reason=page-title

\section{\WorkTitle{论若望福音}第一〇七讲（\BibleRef{若 17:9-13}）}

1. 主向圣父谈到那些祂已收为门徒的人时，在其它话中也说了这句话：\BibleQuote{我为他们祈求。我不为世界祈求，只为祢所赐给我的人祈求。}这里，祂愿意藉\Quote{世界}指那些顺从世俗私欲、不蒙恩立于祂从世界拣选之列的人。因此，祂表示自己不是为世界，而是为圣父赐给祂的人祈求；因为他们既然由圣父赐给了祂，就已不再属于祂不为他们祈求的那个世界了。\par

2. 接着祂又说：\BibleQuote{因为他们属于祢。}圣父把那些人赐给圣子，并未失去他们；因为圣子接着说：\BibleQuote{我的一切都是祢的，祢的一切也是我的。}由此足以显明圣父的一切如何也属于圣子：即祂自己也是天主，由圣父所生，与圣父平等；而不像对两个儿子中的长子所说的：\BibleQuote{你常同我在一起，凡我所有的，都是你的。}\BibleRef{路 15:31}因为那是就所有低于神圣理性受造物、当然从属于教会的受造物而言的；其中包含长子与幼子两个儿子以及所有圣天使，我们在基督和天主的国里将与他们平等：\BibleRef{路 20:36}但这里说：\BibleQuote{我的一切都是祢的，祢的一切也是我的}，其含义是连理性受造物本身也包含在内，它只服从于天主，所以它以下的一切也都服从于祂。既然它属于天主圣父，倘若圣子不与圣父平等，就不会同时也属于圣子；因为祂说：\BibleQuote{我不为世界祈求，只为祢所赐给我的人祈求：因为他们属于祢。我的一切都是祢的，祢的一切也是我的。}从道德上说，祂这样谈论的圣徒不可能属于任何别的主，只能属于创造并圣化他们的那一位；同样，他们的一切也必然属于他们所归属的那一位。因此，既然他们同时属于圣父和圣子，就证明了他们同等归属之对象的平等。但当祂论及圣神说：\BibleQuote{凡圣父所有的，都是我的；为此我说，祂要从我取，并告诉你们}时，祂指的是关乎圣父实际天主性的事，在那些事上祂与圣父平等，拥有祂所有的一切。而圣神所要领受的，并非祂所说\BibleQuote{祂要从我取}的那个，是属于服从圣父和圣子的受造物；而绝对是来自圣父，圣神由祂发出，圣子也由祂所生。\par

3. 祂接着说：\BibleQuote{我在他们内已得了光荣。}现在祂说到自己的光荣时，好像已经完成，尽管尚未来临；而稍前祂还请求圣父完成此事。但这一光荣是否与祂所说的\BibleQuote{父啊，现在求祢在祢自己内光荣我，赐给我在世界未有以前，我在祢面前所有的光荣}相同，确实值得探讨。因为若说\Quote{在祢面前}，又怎能\Quote{在他们内}？岂不是当这一认识传达给他们，并藉他们作为见证人传达给所有信从他们的人时吗？这样我们就能清楚地理解为基督论及宗徒时说祂在他们内得了光荣；祂说这事已经完成，表明它已预定，只不过是希望未来之事被视为确定。\par

4. 祂又说：\BibleQuote{现在我不在世界上，他们却在世界上。}如果你想到祂说话的那个时辰，祂自己和祂所谈论的那些人都仍在世上；因为我们不能或不应从心灵和生活的倾向来理解，以致说他们仍在世上是因为他们还贪恋世俗，而祂不在世上是因为祂的心志是神圣的。因为这里用了一个词，使得任何这样的理解绝不可接受；祂不是说\Quote{我不在世上}，而是说\BibleQuote{我不再在世上}；由此表明祂曾在世上，但不再在了。那么我们能否相信祂一度贪恋世俗，最终摆脱了这种错误，不再保留旧习？谁敢于把自己封闭在如此亵渎的理解中呢？因此，剩下的理解是：就祂先前也在世界上的同样意义而言，祂宣告自己不再在世上，即祂的肉身临在；换言之，由此表明祂自己离开世界即将到来，而他们的离开则在以后，尽管祂和他们仍都在眼前。因为祂讲这话，是照常人的说法，与人一致。我们不是每天都这样讲一个即将离去的人\Quote{他不再在这里}吗？特别是我们惯常这样讲到那些临近死亡的人。此外，主自己仿佛预见到日后读这些话的人可能产生的思想，又加了一句：\BibleQuote{我往祢那里去}，从而在某种程度上解释祂为何说\BibleQuote{我不再在世界上}。\par

5. 因此，祂将因身体缺席而即将离开的那些人托付给圣父照顾，说：\BibleQuote{圣父啊，求祢因祢的名，保全祢所赐给我的人。}也就是说，作为人，祂为祂从天主那里领受的门徒向天主祈祷。但请注意下文：\BibleQuote{使他们合而为一，如同我们一样。}祂不是说：使他们与我们合而为一，或使他们与我们合而为一，如同我们一样；而是说：\BibleQuote{使他们合而为一，如同我们一样。}这当然是指，在他们的本性上，他们可以合而为一，如同我们在我们的本性上合而为一；这若不是指祂作为天主，与圣父具有同一本性，则这句话就不可能是真实的，正如祂在其他地方所说的：\BibleQuote{我与圣父原是一体。}而并非指祂作为人时的身份，因为在这方面祂说：\BibleQuote{圣父比我大。}但由于同一位既是天主又是人，我们应将祂的祈求理解为出于人性；而祂与祂所祈求的那位是一体，则出于天主性。但在下文还有一处，我们需要更仔细地讨论这个问题。\par

6. 但在这里祂继续说道：\BibleQuote{我与他们同在时，我因祢的名保全了他们。}祂说：我正到祢那里去，求祢因祢的名保全他们，就是我自己与他们同在时所因以保全他们的那名。圣父的名下，圣子作为人，在与人同处时保守了祂的门徒；但圣父也因圣子的名，保全了那些因圣子的名祈求而祂垂听并应允的人。因为圣子自己也对他们说过：\BibleQuote{我实实在在告诉你们：你们因我的名无论向圣父求什么，祂必赐给你们。}但我们不可用如此属肉体的方式来理解，以为圣父和圣子轮流看顾我们，在守护我们的交替更迭中，好像一位离开了另一位才接替；因为我们是同时被圣父、圣子和圣神守护的，祂们是唯一真实有福的天主。但圣经除非降卑到我们的水平，就无法提升我们：正如圣言成为血肉，降下来提拔我们，而祂自己并没有跌落到深层。如果我们认识这位降下的那一位，就让我们与那提拔我们的祂一同上升；并让我们明白，祂这样说，是在位格上加以区别，却没有分开本性。因此，当圣子以身体同在的方式保守祂的门徒时，圣父并非等待着圣子离开以接替守护，而是两者藉着属神的能力共同保守他们；当圣子收回祂的身体同在时，祂与圣父一同保留了属神的守护。因为当圣子作为人也担任守护者的角色时，祂并没有将门徒从圣父的守护中撤出；而当圣父将门徒交托给圣子守护时，在交托中祂并非与接受者分离，而是将他们交托给作为人的圣子，却并非与作为天主的同一圣子分离。\par

7. 因此圣子继续说道：\BibleQuote{祢所赐给我的人，我保全了他们，除了那丧亡之子，没有失落一个，为应验圣经。}基督的背叛者被称为丧亡之子\OriginalTerm{丧亡之子}{\GreekText{υἱὸς} \GreekText{τῆς} \GreekText{ἀπωλείας}}，预定要丧亡，正如圣经所预言的，在第一百零九篇圣咏中特别预言了他。\par

8. \BibleQuote{现在，我往祢那里去；我仍在这世上说这些话，为叫他们心里充满我的喜乐。}请看！祂说祂在这世上说话，而祂刚才还说：\BibleQuote{我不再在这世上了。}其理由我们已在前面解释过，或者不如说，已经表明祂自己解释了。因此，一方面，祂尚未离开，所以还在这里；另一方面，因为即将离开，所以某种程度上已经不再在这里。但祂说：\BibleQuote{叫他们的心里充满我的喜乐，}这喜乐是什么，上面在祂说\BibleQuote{使他们合而为一，如同我们一样}时已经阐明。祂说，祂要将自己的喜乐赐给他们，使之在他们内圆满；为此祂宣称祂在这世上说了这些话。这就是来世和平与福乐，为达到这境界，我们在今世必须生活节制、正义和虔敬。\par

\SourceRecord{Translated by John Gibb.；Nicene and Post-Nicene Fathers, First Series；Vol. 7.；Edited by Philip Schaff.；Buffalo, NY: Christian Literature Publishing Co.,；1888.；<http://www.newadvent.org/fathers/1701107.htm>.}

% structure: p0001 source=h1 target=section reason=page-title

\section{\WorkTitle{论道集108}（\BibleBook{若望}{福音} 17:14-19）}

1. 当主仍在向圣父说话，并为祂的门徒祈祷时，祂说：\Quote{我已经将祢的话赐给了他们，世界却恨了他们。}他们尚未在自己后来的那些苦难中经历这种恨意，但祂以惯常的方式说话，用过去时态预言未来。然后，祂补充了世界恨他们的原因，说：\Quote{因为他们不属于世界，如同我不属于世界一样。}这是藉着重生赋予他们的；因为按出生，他们属于世界，正如祂曾对他们说：\Quote{我从世界中拣选了你们。}因此，这是一种恩宠的特权赋予他们，使他们像祂自己一样，不属于世界，藉着祂将赐予他们的从世界中的解救。然而，祂从不属于世界；因为即使在祂仆人形态的方面，祂也是由圣神所生，而他们是藉着圣神重生的。因为如果他们因由圣神重生而不再属于世界，同样地，祂因由圣神所生也从不属于世界。\par

2. \Quote{我不求}祂接着说，\Quote{祢将他们从世界提去，但求祢保护他们免于凶恶。}因为他们仍然认为有必要留在世界中，尽管他们已不再属于世界。然后祂重申同一句话：\Quote{他们不属于世界，如同我不属于世界一样。请以真理圣化他们。}这样，他们便如祂先前所求的，得以免于凶恶。但有人可能会问，如果他们尚未在真理中被圣化，他们如何不再属于世界；或者，如果他们已经圣化，为何祂还要如此祈求？岂不是因为那些已经圣化的人仍在这同一圣化中不断进步，在圣德中成长，并且除非有天主的恩宠相助，靠祂圣化他们的进步，如同祂圣化他们的开端，否则便不能如此吗？因此，宗徒也同样说：\BibleQuote{在你们中间开始了美好工作的那位，必会完成这事，直到耶稣基督的日子。}\BibleRef{斐 1:6}新约的继承者便是在那旧约的洁净礼所预表的真理中被圣化；当他们被圣化在真理中时，换句话说，他们是在基督内被圣化，祂曾真实地说：\Quote{我就是道路、真理和生命。}同样，当祂说：\Quote{真理将使你们自由，}在解释祂的话时，祂随即补充说：\Quote{如果圣子使你们自由，你们必会真正自由，}以表明祂先前所称呼的真理，片刻后就称为圣子。那么，祂在我们面前的话中意谓了什么：\Quote{请以真理圣化他们，}不就是说：在祂内圣化他们？\par

3. 最后，祂继续讲论，并不忘记以越来越清楚的同义反复提示：\Quote{祢的言语就是真理。}祂还能意指别的什么吗，除了\Quote{我是真理}？因为希腊福音中用的是\OriginalTerm{Logos}{\GreekText{λόγος}}，这也是在\Quote{在起初已有圣言，圣言与天主同在，圣言就是天主}这段经文中所出现的那个词。而我们至少知道那\OriginalTerm{圣言}{Word}是天主的独生子，祂\Quote{成了血肉，居住在我们中间}。因此，这里也可能像某些抄本那样写作\Quote{祢的圣言就是真理}，正如在某些抄本中另一经文写作\Quote{在起初已有言语}。但在希腊文中，两者没有差异，都是\GreekText{λόγος}。因此，圣父是在真理中圣化，即在祂自己的圣言、祂的独生子内圣化祂的继承者和（圣子的）共同继承者。\par

4. 但此刻祂继续谈论宗徒，因为祂随后补充说：\Quote{祢怎样派遣我到世界上来，我也这样派遣他们到世界上去。}祂派遣的若不是祂的宗徒，还能是谁？因为宗徒这个名称，原是希腊词，在拉丁文中无非意为\Quote{被派遣者}。因此，天主派遣了祂的圣子，不是在罪恶的血肉中，而是在罪恶血肉的模样下；\BibleRef{罗 8:3}而圣子派遣了那些在罪恶血肉中出生、却被祂从罪恶的玷污中圣化了的人。\par

5. 然而，由于天主与人之间的中保——基督耶稣的人性——已成了教会的元首，因此他们便是祂的肢体；故此，祂在随后的话中说：\Quote{我为他们的缘故，圣化我自己。} 祂说\Quote{我为他们的缘故，圣化我自己} 是什么意思呢？岂不是我在他们内圣化他们，因为他们在某种意义上也是我自己？因为祂所论及的那些人，如我所说，是祂的肢体，而头和身体同是一个基督，正如宗徒论及亚巴郎的后裔时所说：\Quote{如果你们属于基督，那么你们就是亚巴郎的后裔。} 在此之前他曾说：\Quote{祂并没有说\InnerQuote{与后裔们}，好像为许多人的，而是如同为一个的：\InnerQuote{与你的后裔}，这指的是基督。} 因此，如果亚巴郎的后裔是基督，那么那对听他话的人所宣告的\Quote{你们就是亚巴郎的后裔}，岂不意味着\Quote{你们就是基督}？同一位宗徒在另一处也说过类似的话：\Quote{现在我因着为你们所受的苦而喜乐，并在我的肉身上补满基督苦难所欠缺的。} \BibleRef{哥 1:24} 他没有说\Quote{我的苦难}，而是说\Quote{基督的}，因为他自己是基督的肢体，在他的迫害中（这些苦难正是基督该在祂整个身体上所承受的），他也在补满基督苦难中属于他自己的那一份。要确知本章中这段话的确切含义，请留意下文。因为在说完\Quote{我为他们的缘故，圣化我自己}之后，为了让我们明白祂的意思是祂要在自己内圣化他们，祂立刻加上：\Quote{好叫他们也在真理中被圣化。} 这岂不就是在\Quote{我}内，因为真理就是那在起初就与天主同在的圣言？人子自身也从受造之初就在圣言内被圣化了，当时圣言成了血肉，因为圣言与人成为一位。因此，祂是在自己内圣化自己，即祂的人性在圣言内圣化自己；因为圣言与人同是一位基督，祂在人身上以圣言圣化人性。但为了祂的肢体，祂说：\Quote{我为他们的缘故，我}——即为了使这恩惠也临到他们身上，因为他们在\Quote{我}中也占有一份，正如这在我自己身上有益一样，因为若没有他们，我只是一个单独的人——\Quote{圣化我自己}，即我圣化他们，如同圣化我自己的自我一样，因为在我内，他们也是我。\Quote{好叫他们也在真理中被圣化。} 因为\Quote{他们也在} 一词岂不意味着\Quote{他们} \Emph{同样地} 像我一样被圣化？而\Quote{在真理中}，这\Quote{真理}就是我。此后，祂开始不仅论及宗徒，也论及祂的其余肢体，我们将就着所赐予的恩宠在另一篇讲道中加以讨论。\par

\SourceRecord{Translated by John Gibb.；Nicene and Post-Nicene Fathers, First Series；Vol. 7.；Edited by Philip Schaff.；Buffalo, NY: Christian Literature Publishing Co.,；1888.；<http://www.newadvent.org/fathers/1701108.htm>.}

% structure: p0001 source=h1 target=section reason=page-title

\section{\WorkTitle{讲道集 109（若 17:20）}}

1. 主耶稣，在祂受难已近之时，为祂的门徒（祂也称之为宗徒）祈祷后，与这些门徒共进了最后的晚餐，出卖者因面包蘸汁被揭露而离开；此后，在为他们祈祷之前，祂已与他们说了许多话，现在又将所有将来要信祂的人联结起来，对圣父说：\Quote{我不但为他们祈求，}即为当时与祂同在的门徒，\Quote{也为那些因他们的话而信从我的人}。借此，祂愿意所有属于祂的人都被包括：不仅当时在世的人，也包括将来的人。因为所有自此以后信祂的人，无疑都是因宗徒的话而信，且将继续信直到祂来临；因为祂曾对他们说：\Quote{你们也要作证，因为你们从起初就与我同在}；福音正是通过他们得以传扬，甚至在写成文字之前；确实，凡信基督的人都是信福音。因此，祂说因他们的话而信祂的人，不仅指那些亲耳听见宗徒们现世宣讲的人，也包括他们去世后的其他人；我们这些很久以后出生的人，也是因他们的话而信了基督。因为当时与祂同在的人将他们从祂那里听来的话传给了别人；如此，他们的话传到我们这里，使我们也能相信；无论教会在何处存在，这话还将传给后代——无论将来谁在何处信祂。\par

2. 在这篇祈祷中，耶稣似乎忽略了为祂的一些子民祈祷，除非我们仔细审视祂祈祷中的话语。因为如果祂先为那些当时与祂同在的人祈祷（如我们已指出的），然后也为那些因他们的话而信祂的人祈祷，那么有人可能会说，祂没有为那些既不在祂说话时与祂同在、此后也不是因他们的话而信、而是在更早时候以某种方式（或凭自己，或通过其他可以设想的方式）信了的人祈祷。例如，\Person{纳塔乃耳}那时与祂同在吗？\Person{阿黎玛特雅人若瑟}（他从\Person{比拉多}那里求得耶稣的身体，同一位福音作者\Person{若望}见证他已是耶稣的门徒）呢？祂的母亲\Person{玛利亚}，以及我们从福音中知道早已是祂门徒的其他妇女呢？那些\Person{若望}福音中多次提到\Quote{许多人信了祂}的人那时与祂同在吗？那些手持树枝，有的在前、有的在后，随着祂骑在驴上高呼\Quote{奉主名而来的，当受赞美}的群众，以及那些孩童——祂亲自宣称预言应验于他们——\Quote{从婴孩和吃奶者的口中，祢预备了赞美}——从何而来？那五百弟兄（祂复活后曾同时显现给他们——\BibleRef{格前 15:6}）又是从哪里来的？那一百零九人加上十一人共一百二十人（祂升天后聚集等候并领受了圣神的应许）又是从哪里来的？所有这些人不都是出自那些\Quote{许多人信了祂}的人吗？因此，救主此时并没有为他们祈祷，因为祂是为那些当时与祂同在的人祈祷，并为那些尚未信、但将因他们的话而信的人祈祷。然而，那些人当时并不与祂同在，而且已经在更早的时候信了祂。我不提年迈的\Person{西默盎}（他在耶稣婴孩时就信了祂）、女先知\Person{亚纳}（\BibleRef{路 2:25}）、\Person{匝加利亚}和\Person{依撒伯尔}（他们在耶稣由童贞女诞生前就预言了祂）、他们的儿子\Person{若翰}——新郎的朋友、祂的前驱——他因圣神认出祂，并在祂不在时宣讲祂，在祂显现时向他人指明；我不提这些人，因为有人可能回答说：祂不应为已死的人祈祷，他们带着大功绩离世，已安息于接纳中；对于古时的义人也有类似的回答。因为若不是藉圣神的启示，相信那将来在肉身中降世的天主与人之间的唯一中保，他们中有谁能从那由一人招致的整个毁灭之罚中得救呢？但祂有必要为宗徒祈祷，却不为此众多仍然活着、但当时不与祂同在、且已在更早时候被引入信仰的人祈祷吗？谁会这样说呢？\par

3. 因此我们应当理解，他们对祂的信心尚未达到祂所期望的程度，因为就连\Person{伯多禄}本人——祂曾因\Person{伯多禄}宣认\Quote{祢是默西亚，永生天主之子}而给予他极卓越的见证——也宁愿阻止祂死去，而不相信祂死后的复活，因此随即被祂称为\Person{撒殚}。这样看来，那些早已去世、却因圣神的启示毫不怀疑基督会复活的人，在信德上比那些在相信祂要救赎\Place{以色列}后、见到祂的死便失去先前所有希望的人更大。因此，我们最好相信，在祂复活后，圣神被赐下，宗徒们受教导、被坚固，并从起初就成为教会中的教师，其他人因他们的话而获得了对基督应有的信德，换言之，那时他们才牢牢把握了祂复活的信德。如此，所有那些似乎早已信了祂的人，实际上也属于祂祈祷时所说的\Quote{我不但为他们祈求，也为那些因他们的话而信从我的人}之列。\par

4. 为更进一步解决这问题，我们尚保有一位荣福的宗徒，以及那个在邪恶中为恶徒、在十字架上却成信者的右盗。因为宗徒保禄告诉我们，他成为宗徒，不是由于人，也不是借于人，而是由于耶稣基督；他在谈及自己的福音时说：\BibleQuote{我既不是从人领受，也不是借于人，而是借着耶稣基督的启示而来的。}那么，他怎能列在那些经上所说\BibleQuote{他们将因他们的话而信从我}的人中呢？另一方面，那个右盗在教师们本人先前所有的信德全然崩溃之时相信了。因此，连他也不是因他们的话而信从基督；然而他的信德达到如此地步，以至于他承认自己所看见被钉在十字架上的那位不但要复活，而且还要为王，他说：\BibleQuote{当祢进入祢的王国时，请记念我。}\BibleRef{路 23:42}\par

5. 因此，倘若我们要相信主耶稣在此祈祷中是为所有属于祂的人，即那些当时或在世上（此为地上考验之境）将生活的人祈祷，我们就必须如此理解\BibleQuote{因他们的话}这一表达：相信它在这里是指在世上所宣讲的信德之言本身，它之所以被称为他们的话，是因为它主要是由他们首先宣讲的。因为当保禄借着耶稣基督的启示领受同样的话时，这话已经在世上由他们宣讲了。由此也发生了这样的事：他将福音与他们加以比较，免得他先前或将来是空跑；他们向他伸出右手，因为在保禄身上，他们虽然并未亲自给他，却找到了他们自己已在宣讲、且在其中已得坚固的话。关于基督复活的话，同一位宗徒说：\BibleQuote{无论是我，或是他们，我们都这样宣讲，你们也这样相信了。}\BibleRef{格前 15:11} 他又说：\BibleQuote{这是信德之言，我们宣讲的，就是：你若口里承认耶稣是主，心里相信天主使祂从死者中复活了，你便可得救。}\BibleRef{罗 10:8} 在\BibleBook{宗徒}{大事录}中我们也读到：天主在基督内已经向所有人标明了信德的基础，因为祂使祂从死者中复活了。\BibleRef{宗 17:31} 因此，这信德之言，由于主要是由那些依附于祂的宗徒们首先宣讲的，故被称为他们的话。然而，它并不因被称为他们的话而不再是天主圣言；因为同一位宗徒说，得撒洛尼人从他领受了这话，\BibleQuote{不是人的话，而是，正如它实实在在是天主的话一样。}\BibleRef{得前 2:13} \Quote{天主的话}，正是因为它是天主所白白赐予的；但被称为他们的话，则是因为天主首先主要把它交给了他们去宣讲。同样，上面提到的那个右盗，在他自己的信德方面，也拥有他们的话；这话恰被称为他们的，是因为其宣讲首先主要属于他们所担负的职务。还有，当希腊寡妇因膳食服务的事发怨言时，是在保禄被带到基督信仰之前；那时，早已依附于主的宗徒们回答道：\BibleQuote{我们放弃天主圣言而去服侍膳桌，并不相宜。}\BibleRef{宗 6:1} 当时他们便设立了执事职，使自身不致被从宣讲圣言的职务中拉走。因此，那信德之言，无论是谁从何处听到而相信了基督，或是尚待听到而将要相信，都恰如其分地被称为他们的话。因此，在这次祈祷中，我们的救主在为当时与祂同在的宗徒们祈祷时，也将那些将因他们的话而信从祂的人联合起来，从而为所有祂所救赎的、不论是当时有肉身生命的还是以后要在肉身生活的人都作了祈祷。至于在这样结合之后，祂接着所说的话，必须留待另一篇讲道来讨论。\par

\SourceRecord{Translated by John Gibb.；Nicene and Post-Nicene Fathers, First Series；Vol. 7.；Edited by Philip Schaff.；Buffalo, NY: Christian Literature Publishing Co.,；1888.；<http://www.newadvent.org/fathers/1701109.htm>.}

% structure: p0001 source=h1 target=section reason=page-title

\section{论训110（若望福音 17:21-23）}

1. 主耶稣为当时与他同在的门徒祈祷后，又将其他属于他的人也联合在他们中，说道：\Quote{我不但为他们祈祷，也为那些因他们的话而信从我的人祈祷，}就好像我们追问他为那些人祈祷什么或为什么祈祷，他立刻接下去说：\Quote{使他们众人合而为一；如同你，圣父，在我内，我在你内，使他们也在我们内合而为一。}稍前，他仍仅为当时与他同在的门徒祈祷时，曾说：\Quote{圣父啊，请因你的名保全你所赐给我的人，使他们合而为一，如同我们一样}\BibleRef{若 17:11}。因此，他如今也为我们祈祷同一件事，正如他当时为他们祈祷的一样，就是使所有人——即我们和他们——合而为一。在这里我们必须特别注意：主并没有说我们众人合而为一，而是说\Quote{使他们众人合而为一；如同你，圣父，在我内，我在你内}（其中应理解为\Emph{合而为一}，正如后来更清楚表明的）；因为他先前也曾论及与他同在的门徒说：\Quote{使他们合而为一，如同我们一样。}因此，圣父在圣子内，圣子在圣父内，如此成为一体，因为他们属于同一实体；然而我们虽然确实能在他们内，却不能与他们成为一体；因为他们和我们不属于\OriginalTerm{同一实体}{unius substantiae}，就圣子与圣父同为天主而言。但就他作为人而言，他与我们属于同一实体。但此刻他更愿唤起我们对他在另一处所说的另一句话的注意：\Quote{我与圣父是一体，}在那里他表明自己的本性同圣父的本性一样。因此，虽然圣父和圣子，甚至圣神，在我们内，我们绝不能认为他们与我们具有同一本性。所以，他们在我们内，或我们在他们内，其意义是：他们在自己的本性内是一体，我们也在我们的本性内是一体。因为他们在我们内，如同天主在他的圣殿内；但我们在他们内，如同受造物在其造物主内。\par

2. 但接着在说了\Quote{使他们也在我们内合而为一}之后，他加上了\Quote{为使世界相信是你派遣了我}。他这是什么意思？难道当我们都在圣父和圣子内合而为一时，世界就会因此被带入信仰吗？这种状态不就是永恒的平安和信德的赏报，而非信德本身吗？因为我们合而为一，不是为了让我们相信，而是因为我们已经相信了。但虽然在此生，因着共同的信仰本身，所有相信的人都合而为一，正如宗徒的话说：\Quote{因为你们众人在基督耶稣内都是一体；}\BibleRef{迦 3:28}。即使这样，我们合而为一，不是为了让我们相信，而是因为我们确实相信。那么，这话是什么意思：\Quote{使他们都合而为一，为使世界相信}？这无疑就是说，\Quote{全体}本身就是那相信的世界。因为将要合而为一的人并非属于一类，而之后基于这些人将要合而为一这一理由而相信的世界是另一类；因为十分确定的是，他说\Quote{使他们都合而为一}是指那些他先前说过\Quote{我不但为他们祈祷，也为那些因他们的话而信从我的人祈祷}的人，并且紧接着就说\Quote{使他们都合而为一}。这\Quote{全体}若不是世界，又是什么？当然不是那敌对的世界，而是那相信的世界。因为你们在这里看到，那位曾说\Quote{我不为世界祈祷}的，现在却为世界祈祷，为使它相信。因为有一个世界，关于它经上记载说：\Quote{使我们不致与这世界一同受罚。}\BibleRef{格前 11:32}。他不为那个世界祈祷，因为他完全知道它被预定成什么。还有另一个世界，关于它经上记载说：\Quote{因为人子来到世上，不是为审判世界，而是为叫世界藉着他而得救。}因此宗徒也说：\Quote{天主在基督内使世界与自己和好。}\BibleRef{格后 5:19}。正是为这个世界他祈祷，说：\Quote{为使世界相信是你派遣了我。}因为藉着这个信仰，世界在相信天主所派遣的基督时，就与天主和好了。那么，我们该如何理解他所说的\Quote{使他们也在我们内合而为一，为使世界相信是你派遣了我}呢？就只能这样理解：他并没有把世界相信的原因归于那些他人合而为一的事实，好像世界是因看见他们合而为一才相信的；因为这里的\Quote{世界}本身包括所有因自己的相信而合而为一的人；但在他的祈祷中，他说\Quote{为使世界相信}，正如他在祈祷中也说\Quote{使他们都合而为一}；并且在同一祈祷中还说了\Quote{使他们也在我们内合而为一}。因为\Quote{使他们都合而为一}这句话等同于\Quote{使世界相信}，因为正是藉着相信，他们才合而为一，完美地合而为一；也就是说，那些虽然按本性是一体，却因彼此纷争而不再合一的人。最后，如果他所用的动词\Quote{我祈祷}在第三分句中被理解为，或者更确切地说，为了使整个表达更完整，处处补充进去，那么这句话的解释就会更清楚：我祈祷\Quote{使他们众人合而为一；如同你，圣父，在我内，我在你内；}我祈祷\Quote{使他们也在我们内合而为一；}我祈祷\Quote{为使世界相信是你派遣了我。}并且请注意，他加上了\Quote{在我们内}这句话，为使我们知道，我们在那不变信实的爱中合而为一应当归功于天主的恩宠，而非我们自己：正如宗徒在说了\Quote{因为你们从前曾是黑暗，但现在你们却是光明}之后，为使无人将此事归于自己，便加上了\Quote{在主内}。\BibleRef{弗 5:8}\par

3. 此外，我们的救主这样向圣父祈祷，显示出祂是人；同时祂现在也表明，祂自己作为与圣父同等的天主，成就了祂所祈求的事，当祂说：\BibleQuote{父赐给我的光荣，我赐给了他们}。那光荣若不朽，还能是什么呢？这原是人性此后要在祂内接受的。因为祂自己那时尚未接受，但按祂惯常的方式，出于预定的绝对确定性，祂用过去时态的动词暗示未来的事，因为祂正被光荣，也就是说，被圣父复活，而祂自己最终也要将我们复活到同一光荣中。这与祂在别处所说的相似：\BibleQuote{就如父复活死人，赐予他们生命，同样，子也随自己的意愿赐予生命}。以及\BibleQuote{谁}，岂非与父同样？\BibleQuote{因为父所做的事}，不是别的事，而是\BibleQuote{子也照样做}，不是用不同的方式，而是\BibleQuote{同样}。祂也以这种方式复活了祂自己。为此祂曾说过：\BibleQuote{拆毁这座圣殿，三天内我要把它重建起来}。因此，祂说圣父所赐给祂的光荣，即不朽的光荣，也必须理解为祂自己赐给自己的，尽管祂没有明说。正是由于这个原因，祂更频繁地说唯独圣父做的事，其实祂自己也与圣父一同做，这样祂便将一切所能归功于祂所从出的那一位。但有时祂也不提圣父，而说祂自己做那些祂只与圣父一同做的事：好让我们明白，当祂不提自己而将某件事工归于圣父时，圣子与圣父的事工不可分离；同样，当提到圣子在做某件圣父也平等参与的事工，却丝毫未提及圣父时，圣父与圣子的事工也不可分离。因此，在圣父的某件工作中，圣子若绝口不提自己的工作，祂便是彰显谦卑，好成为我们更健全健康的泉源；但反过来，在自己某件工作中，圣子若绝口不提圣父的工作，祂便是彰显自己的平等，好叫我们不以为祂是低等的。就这样，在本段经文中，祂并没有使自己与圣父的事工分离，尽管祂说\BibleQuote{圣父赐给我的光荣}；因为祂也把这光荣赐给了自己。祂也没有使圣父与自己的事工分离，尽管祂说\BibleQuote{我赐给了他们}；因为圣父也赐给了他们。原来不仅圣父和圣子，而且圣神的事工，都是不可分离的。但正如祂出于自己的意愿，为所有祂的子民向圣父祈祷，好使\BibleQuote{他们都合而为一}；同样，出于祂自己的恩惠，正如祂所说：\BibleQuote{圣父赐给我的光荣，我赐给了他们}，成就这事也同样出于祂的意愿；因为祂紧接着说：\BibleQuote{使他们合而为一，如同我们原为一体}。\par

4. 然后祂接着说：\Quote{我在他们内，祢在我内，使他们完全合而为一。} 在此祂简略地表明自己是天主与人之间的中保。这话并不是说圣父不在我们内，或我们不在圣父内；因为祂在另一处也曾说：\Quote{我们要到祂那里去，并要在祂那里作我们的住所；} 而且在这同一段稍前，祂并没有像现在这样说：\Quote{我在他们内，祢在我内；} 或：他们在祢内，我在祢内；而是：\Quote{祢在我内，我在祢内，他们在我们内。} 因此，当祂现在说：\Quote{我在他们内，祢在我内，} 这些话是就中保位格而言的，如同宗徒所用的另一句话：\Quote{你们属于基督，基督属于天主。} \BibleRef{格前 3:23} 但祂补充说：\Quote{使他们完全合而为一，} 表明中保所完成的修和，带领我们达到圆满真福之境，此后不能再有增加。因此，紧随其后的：\Quote{为使世界知道祢派遣了我，} 我想不应理解为祂再次说了\Quote{为使世界相信；} 因为有时\Quote{知道}也与\Quote{相信}同义，正如祂早些时候所说的话：\Quote{他们真正知道了我从祢出来，并且相信了祢派遣了我。} 祂后来用\Quote{他们相信了}表达与先前\Quote{他们知道了}相同的意思。但既然这里祂说的是圆满，那么\Quote{知道}必须理解为那时的眼见，而非现在的信德。因为似乎保存了一种次序，针对祂稍前所说的\Quote{为使世界相信，} 而这里则是\Quote{为使世界知道。} 因为虽然祂在那里说\Quote{使他们众人合而为一}和\Quote{在他们内合而为一，} 但祂并没有说\Quote{他们完全合而为一，} 然后再加上\Quote{为使世界相信祢派遣了我；} 但这里祂说\Quote{使他们完全合而为一，} 然后添加的不是\Quote{为使世界相信，} 而是\Quote{为使世界知道祢派遣了我。} 因为只要我们相信所不见的，我们就还没有完全，如同我们将要配得见我们所信的那样。因此，祂在前一处正确地说\Quote{为使世界相信，} 而这里说\Quote{为使世界知道；} 但两处都说\Quote{祢派遣了我；} 使我们知道，就圣父与圣子不可分离的爱而言，我们目前只相信那藉着相信而正在认识的事。如果祂说\Quote{使他们知道祢派遣了我}，这与祂实际所说的\Quote{为使世界知道}具有同等效力。因为他们是那不存留于敌意中的世界，不像那预定受罚的世界；而是从仇敌转化为朋友的世界，因这世界的缘故\Quote{天主在基督内使世界与自己和好。} 因此祂说：\Quote{我在他们内，祢在我内；} 如同祂说：我在祢所派遣给我的人内；祢在我内，藉着我使世界与祢自己和好。\par

5. 与这些话密切相关的还有祂接下来所说的：\Quote{祢爱了他们，如同祢爱了我。} 这就是说，圣父在圣子内爱了我们，因为祂在创世以前就在祂内拣选了我们。\BibleRef{弗 1:4} 因为那爱独生子的，当然也爱祂的肢体——祂藉着祂的中介而用义子收养所嫁接于祂的肢体。但我们并不因此等于独生子——我们藉着他被创造和再造——以致于说：\Quote{祢爱了他们，如同祢也爱了我。} 因为人们说\Quote{如同这样，那样也如此}时，并不总表示平等，有时只是说\Quote{因为这样，所以那样也如此}或\Quote{为了某物如此，使另一物也如此。} 因为谁能说宗徒们被基督派遣到世界上，完全如同圣父派遣了祂一样？因为，姑且不提其他差异——这些差异提起来冗长——他们总之是在已经为人时被派遣的；但祂受派遣却是为了成为人；然而祂在上文说：\Quote{如同祢派遣我到世界上，我也派遣了他们到世界上；} 如同祂说：因为祢派遣了我，我便派遣了他们。同样，在当前的经文中祂说：\Quote{祢爱了他们，如同祢爱了我；} 这无非就是：祢爱了他们，因为祢也爱了我。因为祂不能不爱圣子的肢体，既然祂爱圣子本身；并且爱祂的肢体没有别的理由，只因祂爱圣子本身。但祂爱圣子，就祂的天主性而言，因为祂生了祂与自己同等；祂也爱祂的人性，因为独生子圣言亲自成了血肉，并因圣言的缘故，圣言的血肉为祂所爱；但我们，祂爱我们，因为我们是祂所爱者的肢体；并且为了使我们成为肢体，祂在未有我们之前就爱了我们。\par

6. 因此，天主所怀有的爱是不可理解、永不改变的。因为祂并非从我们藉祂圣子的血与祂修和之时才开始爱我们；而是在世界奠基之前就已爱我们，为使我们能与祂的独生子一同作祂的子女，尽管当时我们尚未有任何存在。因此，我们藉天主圣子的死亡与天主修和这件事，不应被听作或被理解成：圣子使我们与祂修和，是在于祂从此开始爱那些先前祂所恨恶的人，如同仇敌与仇敌修和，之后他们成为朋友，彼此之爱取代了相互之恨。但我们是与那位早已爱我们的天主修和，只是我们因自己的罪而与祂为敌。关于这一点，我所说的是否真实，让宗徒来作证，他说：\BibleQuote{天主却在我们还是罪人的时候，就为我们死，以此证明祂对我们的爱。}\BibleRef{罗 5:8}因此，即使当我们向祂行敌意、作不义之时，祂仍对我们怀有爱；然而，对祂又确确实实说过：\Quote{上主，祢憎恨一切作恶的人。}因此，以一种奇妙而神圣的方式，即使祂恨我们时，祂也爱我们；因为祂恨我们，是就我们不是祂自己所造的那个样子而言；而我们自己的不义并未完全消灭祂的工程，祂同时知道如何在每个人身上，恨我们所做的，而爱祂所做的。这一点，的确可以从所有受造物身上理解，因为那一位是真正被称为：\BibleQuote{祢爱一切所有，不恨祢所造的。}\BibleRef{智 11:25}因为祂绝不会愿意自己憎恨的任何事物存在；若非在祂所憎恨的事物中也有祂所爱之处，全能者所不愿的事物根本不会存在。因为祂正义地憎恨并摈弃恶习，视之为与祂的治理原则完全相悖；然而，祂甚至在堕落的人身上，也爱那能经由医治而接受祂的恩惠，或经由审判而接受祂的判决的部分。如此，天主丝毫也不恨祂所造的任何事物；因为祂是本性的创造者，而非恶习的创造者；祂所恨的恶并非祂所造；而对于这些恶本身，祂真正所做的一切——无论是仁慈地医治，还是公正地整顿——都是好的。既然祂不恨自己所造的任何事物，谁能相称地描述祂多么爱祂独生子的肢体，尤其是多么爱独生子本身？在祂内隐藏着一切可见与不可见的事物，它们按不同的类别被安排，祂也以与这种安排完全和谐的爱来爱它们。因为祂正以其恩宠的慷慨，引领祂独生子的肢体达到与圣天使平等的地位；而独生子本身，作为万有之主，无疑是天使的主，祂按本性作为天主，并非与天使平等，而是与圣父同等；而按恩宠——就祂是人而言——祂如何能不超越一切天使的卓越呢？因为祂内人的肉身与圣言构成了同一个位格。\newline{}\par

7. 然而，不乏有人也将我们置于天使之前；因为他们说，基督为我们而死，并非为天使而死。但这样的想法除了是夸耀我们自身的不敬之外，还能是什么？因为正如宗徒所说：\BibleQuote{基督就在我们尚不敬虔时，按所定的日期为我们死了。}\BibleRef{罗 5:6}这里所称赞的不是我们的功德，而是天主的仁慈。因为一个人因自己的邪恶而病入膏肓，唯有藉医生的死亡才能得医治，却希望自己因此受夸赞，这会是怎样的人？这当然不是我们功德的荣耀，而是我们疾病的医药。或者我们因此就自认为比天使更优越？因为虽有天使也犯了罪，但并没有如此大的劳苦用于他们的医治。仿佛为他们曾尽了至少一点力，而为我们却尽了更大之力。但即使如此，仍可探究：是否因为我们曾一度处于更优越的地位，还是因为我们现在处于更绝望的境地？然而，既然我们知道一切美善的创造者并未赐予任何恩宠来修复天使的恶，我们为何不推断：他们的过犯被判定为更当受罚，是因为犯罪者的本性更为崇高？因为他们在本性上比我们优越，相应地，他们比我们更不该陷入罪恶。然而，他们冒犯创造者时，由于被造时能行更大的善，却对祂的恩惠忘恩负义，就显得更加可憎；而且他们不仅自己离弃祂，还成了我们的欺骗者。因此，那位爱我们如同爱基督一样的主，将使我们成为这伟大美善的承受者；为祂的缘故——祂愿意我们作祂的肢体——我们能等同于圣天使\BibleRef{路 20:36}，尽管我们在本性上被造得比天使低微，又因自己的罪落入如此更大的不配之中，以致我们理应成为他们某种形式的伴侣。\newline{}\par

\SourceRecord{Translated by John Gibb.；Nicene and Post-Nicene Fathers, First Series；Vol. 7.；Edited by Philip Schaff.；Buffalo, NY: Christian Literature Publishing Co.,；1888.；<http://www.newadvent.org/fathers/1701110.htm>.}

% structure: p0001 source=h1 target=section reason=page-title

\section{讲道集111（\BibleRef{若 17:24-26}）}

1. 主耶稣将祂的子民提升到极大的望德中，再也没有比这更大的望德了。请听，并要在望德中喜乐，因为现世的生活不是应被爱慕的，而是应被容忍的，好使你在它的一切患难中拥有忍耐的力量。\BibleRef{罗 12:12}我说，请听，并仔细思想我们的望德被提升到何等境地。基督耶稣说——天主之子、独生子，与圣父同等永恒且相等的那一位说：祂为了我们而成为人，却不似其他所有人那样成为说谎者，祂说——祂就是道路、生命、真理——那战胜了世界的一位，论及祂为之战胜世界的人们说：请听，相信，盼望，切慕祂所说的：\Quote{父啊，}祂说，\Quote{我愿意我在哪里，祢赐给我的人也与我在哪里。}祂说这些人是父赐给祂的，他们是哪些人呢？岂不就是祂在另一处所说的人吗：\Quote{除非派遣我的父吸引他，没有人能到我这里来。}如果我们在这部福音中已经取得了有益的进步，我们就已经知道，祂说父所做的事，祂自己也同父一起同样去做。因此，他们就是祂从父那里领受的人，也是祂亲自从世界中拣选出来的人，并且拣选他们，使他们不再属于世界，如同祂也不属于世界一样；并且使他们也成为相信并知道基督是由天主父所派遣的世界，好使这世界从世界中被拯救出来，这样，作为一个将要与天主和好的世界，它就不会与那处于敌意中的世界一同被定罪。因为祂在这篇祈祷的开头就是这样说的：\Quote{祢赐给了祂权柄管辖凡有血肉的人，}即管辖每一个人，\Quote{使祂将永生赐给祢所赐给祂的一切人。}在此祂清楚地表明，祂确实接受了管辖所有人的权柄，作为将来审判生者死者的那一位，祂可以任意释放祂所愿意的人，任意定罪祂所愿意的人；但这些人被赐给祂，是为了使祂将永生赐给他们全体。因为祂这样说：\Quote{使祂将永生赐给祢所赐给祂的一切人。}因此，这些人被赐给祂，并不是为了使祂扣留永生不给他们；虽然祂也拥有管辖他们的权柄，因为祂已接受了管辖凡有血肉的人的权柄，换言之，管辖每一个人。这样，那已经和好的世界将从敌对的世界中被释放出来，当祂运用祂管辖它的权柄，将它遣送至永死中时；但另一个世界，祂使之成为自己的，好赐给它永生。因此，善牧向祂的每一只羊，伟大的元首向祂的每一个肢体，都确定地应许了这项赏报：我在哪里，你们也将与我同在哪里；全能圣子向全能圣父所宣告的祂的旨意，不可能不成就。因为在那里也有圣神，同样永恒，同样为天主，是二者的唯一圣神，是二者意志的实体。因为我们读到祂在受难前夕所说的话：\Quote{然而不要照我的意愿，而要照祢的意愿，} \BibleRef{玛 26:39}仿佛父有一个意愿，圣子有另一个意愿，这是我们软弱无能的回声，但这软弱已被信德充满，我们的元首在自己的身上改变了它，那时祂也同样背负了我们的罪过。但圣父和圣子的旨意是同一的，圣父和圣子也只有一位圣神，当我们包含圣神时，我们就得以认识天主圣三，愿虔敬相信这一点，尽管同时我们的软弱不允许我们理解。\par

2. 然而，由于我们已经以与我们简短讲述相称的方式，谈论了许诺的对象及其本身的稳定性；现在让我们尽我们所能来看这一点：祂所乐意许诺的是什么，当祂说：\Quote{我愿祢赐给我的人，也与我同在，在我所在的地方。}就祂按肉身成为\Person{达味}后裔的受造物性而言（\BibleRef{罗 1:3}），连祂自己那时也尚未在祂以后将要去的地方；但祂之所以能这样说：\Quote{我所在的地方}，是为了让我们明白祂即将升天，因此祂说到自己时，仿佛已经在那里，正如祂即将在那里一样。祂也能以同样的方式，如同祂先前对\Person{尼苛德摩}讲话时所说：\Quote{没有人升过天，除了那从天上降下来、仍在天上的人子。}因为在那里祂也没有说\Quote{将要}，而说\Quote{是}，这是由于位格的合一，在天主内同时是人，而在人内同时是天主。因此，祂许诺我们将在天堂；因为那里，祂从童贞女所取的仆人之形已被高举，并坐在圣父的右边。由于同一有福的望德，宗徒也说：\Quote{然而富于慈悲的天主，因祂疼爱我们的大爱，竟在我们因过犯死了的时候，使我们同基督一起生活起来——你们得救是由于恩宠——且使我们同祂一起复活，在基督耶稣内使我们一同坐在天上。}\BibleRef{弗 2:4}因此我们可以理解主所说的话：\Quote{我在哪里，他们也在哪里。}祂，确实说到自己已经在那里；但关于我们，祂仅仅宣告祂愿意我们与祂同在那里，没有任何迹象表明我们已经在那里。但是主说祂愿意成就的事，宗徒却说已经完成。因为祂没有说祂将要复活我们，并使我们坐在天上；而是说\Quote{祂已复活了我们，并使我们坐在天上}；因为并非没有充分的理由，而是出于信德的确信，他视那必然将成的事为已成就。但是，如果我们倾向于按照天主性——在此祂与圣父平等——来理解祂的话：\Quote{我愿他们与我同在，在我所在的地方}，那么让我们的心思摆脱一切物质影像的念头：无论灵魂曾呈现给它什么，具有长度、宽度或厚度，被光线染上任何形体的色彩，或弥漫于任何有限或无限的地方空间，都让它尽可能从这些观念中移开其默观的视线，转向内心思想的倾向。我们也不要追问那与圣父同等的圣子在哪里，因为还没有人发现祂不在哪里。如果有人要追问，就让他追问如何能与他同在；不是像祂那样处处都在，而是在祂可能所在的地方。因为当祂对那在树上赎罪并宣认以得救的人说：\Quote{今天你就要同我在乐园里。}\BibleRef{路 23:43}就祂的人性而言，祂自己的灵魂当天要下到地狱，祂的肉身在坟墓里；但就祂的天主性而言，祂当然也在乐园里。因此，那贼的灵魂，解脱了以往的罪行，已幸福地享有祂的恩宠，虽然不能像祂那样处处都在，但当天就能与祂同在乐园里，祂总是处处都在，并未从乐园撤离。无疑地，为此祂说\Quote{我愿他们也在我在的地方}是不够的；祂还加上了\Quote{与我同在}。因为与祂同在是至善。因为甚至不幸的人也能在祂所在的地方，因为无论何人在何处，祂也在那里；但只有有福的人才与祂同在，因为他们只能因祂而有福。难道不是确实对天主说过：\Quote{我若升到天上，祢在那里；我若下到地狱，祢也在那里。}或者基督不就是那\Quote{因纯洁而渗透万有}的天主智慧吗？\BibleRef{智 7:24}然而光明在黑暗中照耀，黑暗却没有领会它。同样地，取一个可见事物的类比，虽然远非相似，就如盲人，即使他在光明所在的地方，他自己却不与光明同在，而是实际上缺席于那临在的事物；同样，不信者和亵渎者，甚至信徒和虔诚者，因为尚不能注视智慧之光，尽管他不能在任何基督不在的地方，他自己却不与基督同在——我指的是实际看见。因为我们不能怀疑，真信徒因信德与基督同在；因为关于此事，祂说：\Quote{不与我在一起的，就是反对我的。}\BibleRef{玛 12:30}但当祂对天主圣父说：\Quote{我愿祢赐给我的人，也与我同在，在我所在的地方}时，祂专指我们将要看见祂实体的那景象。\BibleRef{若一 3:2}\par

3. 不要让人用模糊的矛盾来扰乱意思的清晰；而要让下文为前面的言语提供证据。因为祂在说了\BibleQuote{我愿意他们也在那里，同我在一起}之后，立即接着加上\BibleQuote{为叫他们看见祢赐给我的光荣：因为祢在创世之前就爱了我}。\BibleQuote{为叫他们看见}，祂说；而不是叫他们相信。这是信德的报酬，而非信德本身。因为如果信德在\WorkTitle{希伯来书}中被正确定义为\BibleQuote{所望之事的保证，未见之事的确证}，\BibleRef{希 11:1}那么为什么不可以将信德的报酬定义为看见那些在信德中所希望的事物呢？因为当我们看见圣父所赐给圣子的光荣时——尽管我们会理解此处所说的光荣不是圣父在生祂的同等的圣子时所赐的，而是当祂成为人子、在十字架死亡之后所赐的；——我说，当我们看见圣子的那光荣时，生者与死者的审判就必然会实行，那时恶人将被带走，使他不得见主的荣光\BibleRef{依 26:10}；那光荣，除了祂的天主性，还能是什么呢？因为心地纯洁的人是有福的，因为他们将看见天主\BibleRef{玛 5:8}；而因为恶人心地不纯洁，所以他们不得看见。那时他们将进入永罚；因为恶人将这样被带走，使他不得见主的荣光；但义人将进入永生\BibleRef{玛 25:46}。而永生是什么？\BibleQuote{认识祢，唯一的真天主，和祢所派遣来的耶稣基督}\BibleRef{若 17:3}；不是像那些心地不纯洁的人那样认识祂——他们虽然能看见祂坐在审判中、穿着光荣的仆人形状；而是像心地纯洁的人将要认识的那样，作为唯一真天主，圣子偕同圣父及圣神，因为天主圣三本身就是唯一真天主。那么，如果我们是就祂作为天主圣子、与圣父同等且同永恒的天主性来理解这话\BibleQuote{我愿意他们也在那里，同我在一起}，我们就将与基督同在圣父内；但祂以祂自己的方式，我们以我们的方式，无论我们的身体在哪里。因为如果要理解地方，以及容纳无形体的存有的地方，而每样事物都有其所处之地，那么基督永恒的地方、祂永远所在之处，就是圣父自己，而圣父的地方就是圣子；因为祂说：\BibleQuote{我在父内，父在我内}；并在这祈祷中说：\BibleQuote{父啊，祢在我内，我在祢内}；而他们是我们的地方，因为紧接着：\BibleQuote{使他们也合而为一，在我们内}；而我们是天主的地方，因为我们是祂的圣殿；正如祂为我们死而为我们活着的那样，也为我们祈祷，使我们能合而为一在他们内；因为\BibleQuote{祂的居所设在平安中，祂的住处设在熙雍}，我们就是那熙雍。但是，谁能有资格思考这样的地方或其中的事物，而不带有空间界定的容量和物质团块的概念呢？然而，只要稍有一点进步，那就是当任何这样的观念出现在心目之前时，它被否认、拒绝、排斥；并且尽可能思量一种光，在其中这些事物被感知为只配被否认、拒绝、排斥；而那光的确定性被认识、被爱慕，以便由此在我们内开始一个向上的运动，并努力到达更深处的地方；而当心智因其自身的软弱和仍较低下的纯洁未能穿透它们时，被驱回，不是没有爱情的叹息和热切渴望的眼泪，并且继续耐心忍受，直到它被信德净化，被内在生活的圣德预备好，能够在那里居住。\par

4. 那么，当我们与祂同在圣父内时，我们怎么会不与基督同在祂所在之处呢？在这方面，宗徒也不是无话对我们说，尽管我们尚未拥有现实，只怀着望德。因为他说：\BibleQuote{你们既然与基督一同复活了，就该追求天上的事，在那里有基督坐在天主的右边；你们该思念天上的事，不该思念地上的事。因为你们已经死了，}他加上说，\BibleQuote{你们的生命与基督一同藏在天主内。}看，我们的生命暂时借着信德和望德与基督同在祂所在之处；因为它是与基督一同在天主内。你看，这仿佛已经成就了祂祈祷所求的，当祂说\BibleQuote{我愿意他们也在那里，同我在一起}时；但如今只是借着信德。那么，何时会借着实际看见而成就呢？\BibleQuote{当基督——你们的生命——显现时，}他说，\BibleQuote{那时你们也要与祂一同在光荣中出现。}\BibleRef{哥 3:1}那时我们将显现为我们将是的那个样子；因为那时将明显，我们相信并希望我们将成为那样，并非没有充分的理由，在它实际发生之前。祂将这样做，圣子在说了\BibleQuote{为叫他们看见祢赐给我的光荣}之后，立即加上\BibleQuote{因为祢在创世以前就爱了我}。因为祂在祂内也在创世以前就爱了我们，并预定了祂在世界的终结时将要做的事。\par

5. \Quote{正义的圣父啊}——祂说——\Quote{世界没有认识祢。}正因祢是正义的，世界便没有认识祢。那被预定为受罚的世界理应如此，即没有认识祂；而祂藉着基督与自己和好的世界，并非因功绩，而是因恩宠认识了祂。因为认识祂，不就是永生吗？祂无疑将永生扣留于受罚的世界，而赐予了和好的世界。因此，正因如此，世界没有认识祢，因为祢是正义的，并按其所应得而报应它，使牠不能认识祢；同样，和好的世界认识了祢，因为祢是仁慈的，并非因其任何功绩，而是因恩宠，供应了认识祢所需的帮助。随后接着：\Quote{但我已认识祢。}祂是恩宠的泉源，祂在本性上是天主，又藉着不可言喻的恩宠，由圣神和童贞玛利亚成为人；然后以祂自己的名义，因为天主的恩宠是藉着我们的主耶稣基督，祂又加上：\Quote{这些人也知道是祢派遣了我。}这就是和好的世界。但因祢派遣了我，他们才知道：因此，他们是因恩宠而知道的。\par

6. \Quote{我已将祢的名显示给他们}——祂说——\Quote{并且还要显示。}我已藉信德显示，我将藉亲眼看见显示；我已显示给那些现今在异乡寄居、期限已定的人，我将显示给那些将作王、统治永无止境的人。\Quote{使爱}——祂又加上——\Quote{祢所爱我的爱［字面意为：那］，在他们内，我也在他们内。}（这种说法不寻常：\Quote{爱，\Emph{祢所爱}我的，在他们内，我也在他们内；}因为通常的说法是\Quote{祢用以爱我的爱}。在这里，当然是从希腊文翻译过来的；但在拉丁文中也有类似的表达，如同我们说：祂尽了一份忠信的职守，祂打了一场激烈的仗；而实际上我们理当说：祂以忠信的职守尽了职，祂打了一场激烈的仗。但正如\Quote{祢所爱我的爱}这种表达方式，宗徒也用了类似的表达：\Quote{我打了美好的仗；}\BibleRef{弟后 4:7}他并没有说\Emph{在}美好的仗中，这或许是更常用且更正确的表达方式。）但父爱子的爱，除了因为我们是祂的肢体、在祂内被爱之外，还能如何在我们内呢？因为祂以完整的位格被爱，既是元首又是肢体。因此祂加了：\Quote{我也在他们内；}仿佛说：既然我也在他们内。因为在一种意义上，祂在我们内如同在祂的圣殿里；但在另一种意义上，因为我们自己也是祂，因为祂使自己成为人，为作我们的元首，我们是祂的身体。救主的祈祷结束了，祂的苦难开始了；因此，我们也将结束本讲道，以便在祂恩准之下，在接下来的讲道中论述祂的苦难。\par

\SourceRecord{Translated by John Gibb.；Nicene and Post-Nicene Fathers, First Series；Vol. 7.；Edited by Philip Schaff.；Buffalo, NY: Christian Literature Publishing Co.,；1888.；<http://www.newadvent.org/fathers/1701111.htm>.}

% structure: p0001 source=h1 target=section reason=page-title

\section{论道第112篇（\BibleRef{若 18:1-12}）}

当主在晚餐后、即将为我们倾流祂的血前夕，向当时与祂同在的门徒们发表了那宏大而冗长的讲论，并加上向天父的祈祷之后，圣史\Person{若望}便开始以这些话叙述祂的受难：\BibleQuote{耶稣说了这话，就和门徒出去，过了克德龙溪，在那里有一个园子，祂和门徒都进去了。出卖祂的犹达斯也知道那地方，因为耶稣屡次同祂的门徒在那里聚集。}他在这里所叙述的主与门徒进入园子的事，并非发生在祂结束祈祷之后立即发生，他曾说\BibleQuote{耶稣说了这些话}，但中间插入了其他事件，这些事件被这位圣史略过，而在其他圣史那里可以找到；正如在这位圣史中，也发现许多其他圣史在自己的叙述中同样沉默不提的事。但任何人若愿意知道它们如何彼此一致，以及由一位圣史所提出的真理绝不被另一位所反驳，可以在这些现前的讲论中寻找他所要的，不过不是在现前的讲论里，而是在其他详尽论著中；但若要掌握这门学问，不是站着听，而是坐下来阅读，或对如此行的人给予最密切的关注和思考。然而，让他先相信而后知道，无论他是否也能在此生获得这种知识，或因某些现存的纠葛觉得不可能，即就那些已被教会接纳为正典权威的圣史而言，任何一位圣史所写的，都不可能与他自己的或其他同样真实的叙述相矛盾。因此，目前让我们专注于有福的\Person{若望}的叙述，这是我们承担来阐释的，不作与他人比较，也不在已经足够清楚的事情上逗留；这样，在需要如此行时，我们可以更好地回应要求。因此，让我们不要误解祂的话\BibleQuote{耶稣说了这些话，就和门徒出去，过了克德龙溪，在那里有一个园子，祂和门徒都进去了}，仿佛祂是在说了这些话之后立即进入园子；而是让\BibleQuote{耶稣说了这些话}这个从句表达这样的意思：我们不应假定祂在结束这些话之前就进入了园子。\par

2. \BibleQuote{出卖祂的犹达斯}他说，\BibleQuote{也知道那地方，因为耶稣屡次同祂的门徒在那里聚集。}因此，那披着羊皮的狼，因家主的深谋远虑而被容忍在羊群中，学到了在何处可以适时分散这脆弱的羊群，并为牧羊人设下他垂涎的陷阱。\BibleQuote{犹达斯于是}他接着说，\BibleQuote{领了一队兵，从司祭长和法利塞人那里来的差役，带着灯笼、火把和武器，来到那里。}那是一队兵，不是犹太人，而是士兵。因此我们应当理解，这是从总督那里领来的，仿佛是为了逮捕一个罪犯，通过保持法律权力的形式，来威慑任何人敢于抵抗逮捕者：虽然同时聚集了如此大的队伍，而且如此武装前来，要么是为了恐吓，要么甚至攻击任何敢于挺身保卫\Person{基督}的人。因为祂的能力只是被隐藏，而祂的软弱被突出，以至于祂的敌人认为必须采取这些措施来对付祂，而实际上，除了祂自己所喜悦的，没有任何东西能够伤害祂；祂以自身的良善善用恶人，并在对待恶人时行善，以便将恶转化为善，并区别善与恶。\par

3. \BibleQuote{耶稣因此}圣史接着说道，\BibleQuote{知道一切要临到祂身上的事，就上前问他们说：你们找谁？他们回答说：纳匝肋的耶稣。耶稣对他们说：我就是。出卖祂的犹达斯也同他们站在一起。耶稣一对他们说：我就是，他们就倒退，跌倒在地。}那队兵和司祭长及法利塞人的差役，如今在哪里？武器带来的恐怖和保护在哪里？祂独自的嗓音说出\BibleQuote{我就是}这句话，没有任何武器，就击打、击退、打倒了那一大群人，连同他们所有的仇恨凶残和武器的恐怖。因为天主隐藏在那人性肉身中；永恒的白昼如此隐藏在那些人的肢体中，以至于在灯笼和火把下，祂被黑暗所寻求而杀害。祂说\BibleQuote{我就是}，就将恶人打倒在地。当祂作为审判者来临时，祂会做什么？祂在交付自己受审时就做了这事。当祂来为王时，祂的能力会是怎样？祂在赴死时就拥有了这能力。现在，\Person{基督}透过福音各处仍在说\BibleQuote{我就是}；犹太人却在寻找假\Person{基督}，以便他们可以倒退跌倒在地，如同那些离弃天上之事、贪恋地上之事的人。迫害者跟随叛徒，目的正是要捉拿\Person{耶稣}：他们找到了所寻求的那一位，因为他们听到了\BibleQuote{我就是}。那么，他们为什么没有捉住祂，反而倒退跌倒了？只是因为如此祂乐意，祂能随心所欲行事。但如果祂从未允许他们捉拿祂，他们肯定做不成他们来要做的事，但祂也做不成祂来要做的事。他们实在出于疯狂怒气，寻求置祂于死地；而祂，在交付自己于死亡时，也是在寻求我们。因此，祂向那些有意但无力捉拿祂的人显示了祂的能力；现在让他们捉拿祂，以便祂可以对那些不认识祂旨意的人施行祂的旨意。\par

4. \Quote{然後祂又問他們：\InnerQuote{你們找誰？}他們說：\InnerQuote{納匝肋人耶穌。}耶穌回答說：\InnerQuote{我已經告訴你們，我就是。你們若找我，就讓這些人走吧。}這就應驗了祂所說的話：\Quote{你所賜給我的人，我一個也沒有失落。}} \Quote{你們若找我，}祂說，\Quote{就讓這些人走吧。} 祂看見祂的仇敵，他們就聽從祂的命令：他們讓那些人走了，那些祂不願喪亡的人。但他們後來不是要死嗎？那麼，如果他們現在死了，祂怎會失去他們呢？豈不是因為他們那時還沒有相信祂，就如所有不喪亡的人都相信祂一樣？\par

5. \Quote{西滿伯多祿帶着一把劍，就拔出來，擊中大司祭的僕人，削掉了他的右耳；那僕人名叫瑪耳曷。}這是唯一一位福音作者告訴我們這僕人的名字，正如路加是唯一一位告訴我們主摸了他的耳朵並醫治了他。\BibleRef{路 22:51} 瑪耳曷的解釋是注定要統治的人。那麼，在主的冒犯下被割掉、又被主醫治的耳朵，除了象徵那被修剪了舊人的更新的聽覺，好從此在聖神的新穎中，而不是在文字的陳舊中，還意味着什麼呢？\BibleRef{羅 7:6} 誰能懷疑，那被基督如此對待的人，注定要與基督一同為王呢？而他被發現是僕人，也屬於那生出為奴之身的舊人，即哈加爾。\BibleRef{迦 4:24} 但當醫治來臨時，自由也被預示了。然而，伯多祿的行為被主所不悅，主用這些話阻止他繼續下去：\Quote{把你的劍收回鞘裏；我父給我的杯，我豈能不喝嗎？}因為在那樣的行為中，那門徒只是想保護他的師傅，絲毫沒有考慮到所要象徵的意義。因此，他必須被勸勉操練忍耐，而事件本身也要被記錄下來，作為操練理解的功課。但當祂說苦難之杯是父賜給祂時，我們所得到的真理正是宗徒所說的：\Quote{如果天主為我們，誰能反對我們呢？祂沒有憐惜自己的兒子，反而為我們眾人把祂交出來。}\BibleRef{羅 8:31} 而這杯的創始者也與那飲杯者是同一的；因此同一宗徒也說：\Quote{基督愛了我們，並為我們把自己交出，作為馨香的供物和祭獻，獻於天主。}\BibleRef{弗 5:2}\par

6. \Quote{於是那營隊、千夫長和猶太人的差役捉住了耶穌，把祂捆綁起來。}他們捉住了那位他們從未接近過的：因為祂保持為白日，而他們卻停留在黑暗中；他們也沒有留意那句話：\Quote{你們來接近祂，並被光照。}因為他們若這樣接近祂，就會捉住祂，不是用手去殺害，而是用心去歡迎接納。然而，當他們這樣捉住祂時，他們與祂的距離反而大大增加了：他們捆綁了那本應釋放他們的主。或許在他們中間，有人那時給基督繫上了鎖鏈，後來卻被祂釋放，並能說：\Quote{你解開了我的捆綁。}今天就講到這裏；如果天主願意，我們將在另一篇講道中繼續討論下文。\par

\SourceRecord{Translated by John Gibb.；Nicene and Post-Nicene Fathers, First Series；Vol. 7.；Edited by Philip Schaff.；Buffalo, NY: Christian Literature Publishing Co.,；1888.；<http://www.newadvent.org/fathers/1701112.htm>.}

% structure: p0001 source=h1 target=section reason=page-title

\section{讲道集113（\BibleRef{若 18:13-27}）}

1. 在祂的迫害者通过犹达斯的背叛逮捕并捆绑了爱我们、为我们舍己的主（\BibleRef{弗 5:2}），并且圣父没有顾惜祂，却为我们众人把祂交出（\BibleRef{罗 8:32}），为使我们明白：犹达斯不应因其背叛的用处而受赞美，却应因其恶意的执迷而受惩罚；\Quote{他们把祂带到了}，若望福音作者告诉我们，\Quote{首先去见亚纳斯}。他也没有隐瞒这样做的理由：\Quote{因为他是盖法的岳父，盖法就是那一年的大司祭。盖法就是}，他说，\Quote{向犹太人出过主意，叫一个人替百姓死是有利的}。因此，玛窦相当恰当地，在希望用更少的话说同样的事时，告诉我们祂被带到了盖法面前（\BibleRef{玛 26:57}）；因为祂起初也被带到亚纳斯那里，只因为他是他的岳父；而我们只需明白，这正是盖法所希望的事。\par

2. \Quote{但是耶稣被}，他说，\Quote{西满伯多禄和另一个门徒跟着}。那另一个门徒是谁，我们无法肯定，因为此处没有指明。但若望通常这样提到自己，加上\Quote{耶稣所爱的那个}。因此，也许在此也是他；但不管是谁，让我们看看下文。\Quote{那个门徒}，他说，\Quote{是被大司祭认识的，就同耶稣进了大司祭的庭院；伯多禄却站在门外。那另一个门徒，就是被大司祭认识的，出来对看门的使女说了一声，就领伯多禄进去。于是那看门的使女对伯多禄说：你也是这个人的门徒吗？他说：我不是}。看，那最坚固的支柱，被一口气吹得根基动摇。那许愿者的一切勇敢，和他先前过分的自信，如今在哪里呢？他当时所说的那些话：\Quote{为什么现在我不能跟随祢？我要为祢舍命}，如今又怎样呢？这就是跟随师傅的方式吗？否认自己是门徒？当一个人被一个使女的声音吓到，唯恐迫使我们牺牲时，难道就是这样为主舍命吗？但是，如果天主预言了真相，而人妄自想象了虚妄，又有什么奇怪呢？无疑地，在宗徒伯多禄的这次否认中——现在已进入其第一阶段——我们应当注意：基督不只被那说祂不是基督的人否认，也被那虽然是真正的基督徒，却否认自己是基督徒的人否认。因为主没有对伯多禄说：你将否认你是我的门徒；而是说：\Quote{你将否认我}（\BibleRef{玛 26:34}）。所以，当他否认自己是祂的门徒时，他就否认了祂。这种否认的形式，除了否认他自己的基督徒身份之外，还暗示了什么别的呢？因为虽然基督的门徒那时尚未被称为这个名字——因为是在祂升天以后，在安提约基雅，门徒才初次开始被称为基督徒（\BibleRef{宗 11:26}）——然而那事物本身，后来取得这个名字的，已经存在，那些后来被称为基督徒的人已经是门徒；他们把这个共同的名称，如同共同的信仰，传给了后代。因此，那否认自己是基督门徒的人，就否认了那事物的现实，而被称为基督徒只是其名称。此后有多少人——且不说那些对现世生活感到厌倦的老翁老妇，他们可能更容易为宣认基督而勇敢赴死；也不只说男女青年，在似乎相当需要施展勇气的年龄；就连男孩女孩都能做到——正如无数圣殉道者以勇敢的心和暴烈的死亡进入了天国——而在那一刻，那接受了天国钥匙的人却未能做到（\BibleRef{玛 16:19}）。在此我们看到为什么说：\Quote{让这些人走吧}，当祂以自己的血救赎了我们，为我们舍己时，使祂所说的话得以应验，\Quote{祢赐给我的人，我一个也没有失落}。\par

3. \Quote{仆人和差役站在炭火旁边，因为天冷，就烤火取暖。} 虽然并非冬季，天气却寒冷：这在春分时节也常有发生。\Quote{伯多禄也同他们站着烤火。大司祭就询问耶稣有关祂的门徒和祂的教训。耶稣回答说：\InnerQuote{我向来公开对世界讲话；我常在会堂和圣殿里，就是犹太人聚集的地方教导人，在暗地里我没有说过什么。你为什么问我呢？可以问那些听过我讲话的人，我对他们说了什么；看，他们知道我所说的话。}} 有一个问题不可忽视，就是主耶稣如何说\Quote{我公开对世界讲话}，尤其祂后来加上\Quote{在暗地里我没有说过什么}。难道祂在晚餐后给门徒的最后训话中，不是对他们说：\Quote{这些事我用比喻对你们说了；但时候将到，我不再用比喻对你们说，而要明明地告诉你们关于圣父的事}？那么，如果祂甚至对更亲近的门徒团体也不是公开讲话，而是应许一个将来公开讲话的时候，祂又怎能公开对世界讲话呢？再者，正如其他圣史所证实的，对那些真正属于祂的人（相对于那些不是祂门徒的人），祂独自与他们远离群众的时候，确实讲得更加明白；因为那时祂为他们解释祂用隐晦言语对别人所说的比喻。那么，\Quote{在暗地里我没有说过什么}这句话是什么意思？我们要这样理解祂所说的\Quote{我公开对世界讲话}，仿佛祂说：有许多人听过我。而\Quote{公开}这个词在某种意义上是公开的，在另一种意义上又不是公开的。说它公开，是因为许多人听见祂；说它不公开，是因为他们不明白祂。甚至祂私下对门徒所说的话，也肯定不是暗地里说的。因为谁在许多人面前说话，还能说是暗地里说的呢？正如经上记载：\Quote{凭两三个人的口供，每句话都可成立：}\BibleRef{申 19:15} 尤其是如果对少数人说的话，是希望他们传开让多人知道；正如主本人对祂那时的少数门徒说：\Quote{我在暗中告诉你们的，你们要在光天化日之下说出来；你们在耳边听见的，要在房顶上宣扬出来}？\BibleRef{玛 10:27} 所以，那看似祂在暗地里说的话，在某种意义上也不是暗地里说的；因为祂这样说，并不是要那些听见的人闭口不言，而是要他们尽力宣讲。因此，一件事可以同时是公开的而又不是公开的，或者是暗地里说的而又不是暗地里说的，正如经上所说：\Quote{他们看了又看，却看不见。}\BibleRef{谷 4:12} 因为\Quote{他们看了}，岂不是因为公开、不是暗地里？而同一批人\Quote{看不见}，岂不是因为不公开、而是暗地里？至于那些他们听见却不明白的话，则不能公正或真实地构成对祂的刑事指控。每当他们企图通过问题寻找指控祂的把柄时，祂的回答都完全挫败了他们的阴谋，使他们捏造的诽谤毫无根据。因此祂说：\Quote{你为什么问我呢？可以问那些听过我讲话的人，我对他们说了什么；看，他们知道我所说的话。}\par

4. \Quote{耶稣说了这话，旁边站着一个差役，用手掌打了耶稣，说：\InnerQuote{你就这样回答大司祭吗？}耶稣回答说：\InnerQuote{我若说得不对，你可以指证不对；若说得对，你为什么打我？}} 有什么比这样的回答更真实、更温良、更正义的呢？因为这是祂的回答，从祂口中早已发出先知的声音：\Quote{为了真理、温柔和正义，你当预备你的目标（字面意思：瞄准），顺利前进，为王统治。} 如果我们思量是谁挨了打，难道我们不希望那打祂的人被天火烧灭、或被裂开的地吞没、或被魔鬼抓住带走、或遭受其他甚至更重的这类惩罚吗？创造世界的那一位，若祂不愿意教导我们那战胜世界的忍耐，难道祂不能凭祂的大能命令这一切发生吗？有人会说：为什么祂不遵行祂自己所命令的呢？\BibleRef{玛 5:39} 因为对于打祂的人，祂不应该这样回答，而应该把另一边脸也转过去。不仅如此，难道祂不是既真实、温良、正义地回答了，同时又不但把另一边脸预备给那将要再打的人，而且把祂的整个身体预备被钉在十字架上吗？祂借此反而表明了需要表明的，即那些伟大的诫命不是靠身体上的表现，而是靠内心的准备来完成的。因为即使是一个愤怒的人，也可能在外表上伸出另一边脸。那么，一个内心平静的人真实地回答，并以平静的心预备承受将来更重的苦难，岂不是好得多吗？那为正义而遭受不义的人，若能真诚地说：\Quote{我的心已准备妥当，天主，我的心已准备妥当；} 这一句话正是引出下文的原因：\Quote{我要歌唱，我要赞美；} 保禄和巴尔纳伯即使在最残酷的捆锁中也能做到这一点。\par

5. 但让我们回到福音记述的后续。\Quote{亚纳斯遂把祂捆绑着，解送到大司祭盖法那里。} 按玛窦的记述，祂起初被带到盖法那里，因为他是那一年的大司祭。因为这两位大司祭习惯逐年轮流执行职务，也就是说，作为大司祭；那时正是亚纳斯和盖法，如福音作者路加所记载的，当他讲述主的先驱若翰开始宣讲天国并召集门徒的时候。因为他这样说：\Quote{在大司祭亚纳斯和盖法治下，天主的话临到匝加利亚的儿子若翰在旷野里，}\BibleRef{路 3:2}等等。因此这两位大司祭轮流履行了他们的年份：基督受难时正是盖法的年份。所以，按玛窦的记述，当祂被捕时，祂被带到盖法那里；但首先，按若望的记述，他们带着祂来到亚纳斯那里；不是因为他（亚纳斯）是他的同僚，而是因为他是他的岳父。我们必须设想，这样做是出于盖法的意愿；或者他们的房屋位置如此，以至于他们在经过的路上不能不顾及亚纳斯。\par

6. 但福音作者在说了亚纳斯把祂捆绑着解送到盖法那里之后，又回到他叙事中离开伯多禄的地方，以便解释在亚纳斯家里关于他三次否认所发生的事。\Quote{伯多禄却站着，}他说，\Quote{正在烤火。}他这样重复了他先前已经说过的话；然后接着叙述下文。\Quote{于是他们对他说：你不也是祂的门徒之一吗？他否认说：我不是。}他已经否认过一次；这是第二次。然后，为使第三次否认也应验，\Quote{大司祭的一个仆人，是被伯多禄砍掉耳朵那人的亲戚，说：我不是在园子里看见你同祂在一起吗？伯多禄于是又否认，立刻鸡就叫了。}看，医师的预言应验了，病人的自负被揭露了。因为后者所宣称的\Quote{我愿为祢舍命}并没有实现，但前者所预言的\Quote{你要三次否认我}却实现了。但随着伯多禄三次否认的完成，让这篇讲道现在也结束吧，以便以后我们重新开始思考关于主在总督般雀比拉多面前所发生的事。\par

\SourceRecord{Translated by John Gibb.；Nicene and Post-Nicene Fathers, First Series；Vol. 7.；Edited by Philip Schaff.；Buffalo, NY: Christian Literature Publishing Co.,；1888.；<http://www.newadvent.org/fathers/1701113.htm>.}

% structure: p0001 source=h1 target=section reason=page-title

\section{论道 114（若 18:28–32）}

1. 现在让我们依照若望福音的指示，思考我们的主耶稣基督被带到总督般雀比拉多面前时，祂遭遇了什么，或与祂相关的事。因为福音作者回到他叙述中曾搁置的地方，以解释伯多禄的否认。你们知道，他已经说过：\Quote{亚纳斯就把耶稣捆绑着，解送到大司祭盖法那里。}然后，从他离开伯多禄在庭院中烤火的那段叙述回来后，在完成伯多禄三次否认的全部经过后，他说：\Quote{于是他们把耶稣从盖法那里带到总督府（\OriginalTerm{总督府}{pretorium}）。}因为他曾说过，耶稣是由亚纳斯——他的同僚和岳父——送到盖法那里的。但如果是送到盖法那里，为何又进入总督府呢？这只能理解为总督比拉多居住的地方。因此，要么是因为某种紧急原因，盖法从亚纳斯家里（他们两人曾在那里审问耶稣）出发去了总督的官邸，而将审问耶稣的事留给了他的岳父；要么是比拉多将他的总督府设在盖法的房子里，那房子大到足以容纳主人自己的房间和法官的房间。\par

2. \Quote{早晨到了；他们自己}——即那些带耶稣来的人——\Quote{没有进入总督府}，也就是说，没有进入房子中比拉多占据的那部分，他们以为那是盖法的房子。然后，在解释他们为什么不进入总督府的原因时，他说：\Quote{免得受玷污，但好能吃逾越节晚餐。}因为那是无酵节日的开始；他们认为进入外邦人的住所是一种玷污。不虔敬的盲目！他们真的会因为陌生人的住所而受玷污，却不会因自己的邪恶而受玷污吗？他们害怕因外国法官的总督府而受玷污，却不怕因无辜兄弟的血而受玷污；暂且不说更多，这一点已足以让恶人的良心定罪。此外，主因他们的不虔敬而被带去受死，生命的赐予者即将被杀害，这事实可以归咎于他们的无知，而非他们的良心。\par

3. \Quote{于是比拉多出去到他们那里，说：\InnerQuote{你们带什么控告来告这个人？}他们回答他说：\InnerQuote{这人若不是作恶的（\OriginalTerm{作恶的}{malefactor}），我们就不把他交在你手里。}}让那些从邪灵中被释放的人、被治愈的病人、被洁净的癞病人、听见的聋子、说话的哑巴、看见的瞎子、从死里复活的人，尤其是变成智者的愚拙者，来回答这个问题：耶稣是否作恶？但这些话是由那些主早已借先知预言过的人说的：\Quote{他们以恶报善。}\par

4. \Quote{于是比拉多对他们说：\InnerQuote{你们自己带他去，按照你们的法律审判他罢。}犹太人对他说：\InnerQuote{我们不可以杀人。}}他们疯狂的残酷说的是什么？难道他们没有杀死祂吗？他们正是为此目的把祂带来受审的。难道十字架不能杀人吗？这就是那些不追求智慧、反而迫害智慧的人的愚昧。那么，\Quote{我们不可以杀人}这句话是什么意思呢？如果祂是作恶的（\OriginalTerm{作恶的}{malefactor}），为什么不合法？难道法律没有命令他们不要宽恕作恶者，尤其是那些（如他们所认为的）引诱他们离开天主的人吗？然而，我们必须明白，他们说他们不可以杀人，是因为节日的圣洁——他们刚刚开始庆祝这个节日，并且因为担心进入总督府（\OriginalTerm{总督府}{pretorium}）而受玷污。假以色列人啊，你们变得如此刚硬了吗？你们因过度的恶意而如此丧失理智，以致认为你们没有因无辜者的血而受玷污，因为你们把祂交给别人去流吗？难道比拉多会亲手杀死祂，而你们把祂交给他正是为此目的？如果你们不希望祂被杀，如果你们没有设谋陷害祂，如果你们没有用钱把祂出卖给你们，如果你们没有动手拿祂、捆绑祂、带祂到这里来，如果你们没有亲手把祂交出来、用声音要求杀死祂——那么你们就可以夸口说不是你们杀害了祂。但是，如果除了你们先前所做的一切之外，你们还喊叫：\Quote{钉死，钉死祂}，那么听听先知对你们宣告的是什么：\Quote{人的儿子，他们的牙齿是枪和箭，他们的舌头是利剑。}看哪，这些就是枪、箭、剑，你们用它们杀了那义人，而你们却说你们不可以杀人。因此，当大司祭们为捉拿耶稣而本人没有来、只是派人来时，路加福音在同一段叙述中说：\Quote{耶稣对那些来到祂跟前的人，即大司祭、圣殿警官和长老们说：\InnerQuote{你们带着刀剑棍棒出来，如同对付强盗吗？}}\BibleRef{路 22:52} 因此，正如大司祭们不是亲身前往，而是通过他们所派的人去捉拿耶稣，这不就等于他们本人以其权位来参与吗？同样，所有以不虔敬的声音呼喊钉死基督的人，确实不是直接亲手杀了祂，而是通过那因他们的喊叫而被推动去犯此罪的人，亲自杀了祂。\par

5. 然而，当\Person{圣史若望}补充说：\Quote{为应验耶稣所说的话，祂所指明自己要怎样死}；若我们要把这段话理解为是指十字架的死，似乎犹太人说过：\Quote{我们没有权柄处死任何人}，是因为处死与钉十字架是两回事：我不明白这如何能成为后果，因为这是他们对\Person{比拉多}刚才对他们所说的话的回答：\Quote{你们自己把他带去，按照你们的法律审判他吧。}若是如此，他们难道不可以把祂抓去，自己钉死祂吗？如果他们想用这种刑罚方式来避免处死祂？但谁会看不出，那些不被允许处死任何人的人，却允许有人钉死任何人，这有多么荒谬？不仅如此，主自己岂不也称祂的这种死亡，就是十字架的死亡，为处死，正如我们在\BibleBook{马尔谷}{福音}中读到的，祂说：\BibleQuote{看，我们上耶路撒冷去，人子要被交给司祭长和经师；他们要定祂死罪，要把祂交给外邦人；他们要戏弄祂，唾污祂，鞭打祂，杀害祂，但第三天祂要复活。}\BibleRef{谷 10:33}因此，毫无疑问，主这样说是指明祂自己要怎样死：祂在这里并不是要人理解十字架的死，而是指犹太人要把祂交给外邦人，也就是交给罗马人。因为比拉多是罗马人，由罗马人派到犹太作总督。所以，为使耶稣的话得以应验，即祂被交出去后，被外邦人处死，正如耶稣所预言的那样；因此，当比拉多这位罗马法官想把他交还给犹太人，让他们按照自己的法律审判祂时，他们拒绝接受祂，说：\Quote{我们没有权柄处死任何人。}如此，耶稣关于自己的死所说的话便应验了，即祂被犹太人交出去，被外邦人处死。后者的罪比犹太人轻，犹太人用这种办法试图表明自己不愿祂死，为的是显明他们的疯狂，而非无辜。\par

\SourceRecord{Translated by John Gibb.；Nicene and Post-Nicene Fathers, First Series；Vol. 7.；Edited by Philip Schaff.；Buffalo, NY: Christian Literature Publishing Co.,；1888.；<http://www.newadvent.org/fathers/1701114.htm>.}

% structure: p0001 source=h1 target=section reason=page-title

\section{\WorkTitle{讲道集 115（若望福音 18:33-40）}}

1. 本讲论需思量和讨论比拉多对基督所说的话，以及基督对比拉多的回答。因为向犹太人说了\Quote{拿他去，按你们的法律审判他}这些话之后，而犹太人回答说\Quote{我们处死任何人是不合法的}，比拉多又进了审判厅，叫来耶稣，对他说：\Quote{你是犹太人的国王吗？}耶稣回答说：\Quote{你这话是凭你自己说的，还是别人告诉你的？}主确实知道祂自己所问的，以及对方将如何回答；但祂愿意这些话被说出来，并非为使自己得知，而是为使我们得知的话被记录在圣经中。比拉多回答说：\Quote{难道我是犹太人吗？你的民族和司祭长把你交付给了我：你做了什么？}耶稣回答说：\Quote{我的国不属于这世界；若是我的国属于这世界，我的臣仆必要争斗，使我不致被交付给犹太人；但如今我的国不是从这里来的。}这就是善师愿意我们知道的；但首先必须向我们显明人们关于祂的国所怀的虚妄念头，无论外邦人或犹太人，比拉多正是从他们那里听说的；仿佛祂应因觊觎不合法的王国而被处以死刑；或是如同那些拥有王权的人通常对尚未获得王权者显示敌意，例如要提防，以免祂的国对罗马人或犹太人造成不利。但是主能够回答总督的第一个问题，当他问祂：\Quote{你是犹太人的国王吗？}时，用这话：\Quote{我的国不属于这世界}，等等；但祂转而反问祂，是否这话是凭自己说的，还是从别人听到的，祂愿意通过他的回答表明，犹太人曾以此作为罪行在祂面前控告祂：向我们揭示人心里的思念，这一切祂都知道，不过是虚幻；并且如今，在比拉多回答之后，给他们，无论犹太人还是外邦人，一个更为合理和相称的回答：\Quote{我的国不属于这世界。}但若祂立即回答比拉多的问题，祂的回答似乎只针对外邦人，而不包括犹太人也持有如此看法。但当比拉多回答说：\Quote{难道我是犹太人吗？你的民族和司祭长把你交付给了我}时，他排除了自己可能被视为凭己意说话、称耶稣是犹太人之王的嫌疑，表明这话是由犹太人转告他的。然后通过说：\Quote{你做了什么？}他充分表明这被控告为罪行：仿佛在说，你若否认这样的王权，你做了什么以致被交付给我？因为若有人自称是王，被交付给审判官受罚，不足为奇；但若没有这样的声称，就需要盘问祂，是否做了什么事，以致应当被交付给审判官。\par

2. 那么，你们犹太人、外邦人，听吧；受割损的，听吧；未受割损的，听吧；地上万国，听吧：我不干涉你们在这世界的统治，\Quote{我的国不属于这世界。}不要怀有那极其虚妄的恐惧，当基督的诞生被宣告时，这恐惧曾使大黑落德陷入惊慌，导致他杀害许多婴儿，希望将基督也包括在那罪恶的数量中，他的恐惧比愤怒更残忍：\Quote{我的国}祂说，\Quote{不属于这世界。}你们还要怎样？来归于那不属于这世界的国吧；来相信吧，不要因恐惧而陷入愤怒的疯狂。祂确实预言性地论及天主圣父说：\Quote{然而我被祂立于祂的圣熙雍山为君王}；但熙雍山不属于这世界。因为祂的国是什么，岂不就是那些相信祂的人，祂对他们说：\Quote{你们不属于这世界，如同我不属于这世界}？但祂愿意他们在这世界内；正因此，祂对圣父论及他们说：\Quote{我不求祢将他们从世界接去，但求祢保护他们免于邪恶。}所以祂在这里也不是说：\Quote{我的国不在这世界}，而是\Quote{不属于这世界。}当祂以这话证明这一点：\Quote{若是我的国属于这世界，我的臣仆必要争斗，使我不致被交付给犹太人}，祂不是说：\Quote{但如今我的国不在这里}，而是\Quote{不是从这里来的。}因为祂的国在这里，直到世界的终结，其中有莠子混杂直到收割；收割就是世界的终结，那时收割的人，即天使，要来从祂的国中收集一切绊脚石；\BibleRef{玛 13:38} 若祂的国不在这里，这当然不会发生。但它仍然不是从这里来的；它只是在世界内作客旅；因为祂对祂的国说：\Quote{你们不属于这世界，而是我从世界中拣选了你们。}因此，他们原是属世界的，当他们还不是祂的国时，而是属于今世的首领。所以，所有人类都是属世界的，固然由真天主创造，但从亚当生为败坏和受诅咒的根源；凡在基督内重生的，都成为不再属于这世界的国。因为天主这样从黑暗的权势中拯救了我们，将我们迁移到祂爱子的国里：\BibleRef{哥 1:13} 正是关于这国，祂说：\Quote{我的国不属于这世界}；或\Quote{我的国不是从这里来的。}\par

3. \BibleQuote{因此比拉多对他说：\InnerQuote{那么，你是君王吗？}耶稣回答说：\InnerQuote{你说我是君王。}} 祂并非害怕承认自己是君王，而是\Quote{你说}一语如此平衡，以致祂既不否认自己是君王（因为祂是君王，祂的国不属于这世界），也不承认自己是那样的君王，以致让人以为祂的国属于这世界。因为当比拉多问道：\Quote{那么，你是君王吗？}时，这正是他心中的想法；所以他得到的回答是：\Quote{你说我是君王。}因为\Quote{你说}这句话，就好像是说：\Quote{你既是属血肉的，你便按血肉的方式来说。}\par

4. 此后祂又加上说：\BibleQuote{我为此而生，也为此来到世界，为给真理作证。}由此明显可见，祂在这里指的是祂自己暂时的诞生，即祂降生成人来到世界，而非那无始的诞生，凭此祂是天主，圣父藉祂创造了世界。因此祂声明，祂为此（即为此缘故）而生，为此来到世界，就是由童贞女所生，为给真理作证。但因并非所有人都有信德，\BibleRef{得后 3:2}祂进一步说：\BibleQuote{凡属真理的人，必听我的声音。}祂听，就是说，以内心之耳聆听，换言之，祂服从我的声音，这等于说，祂相信我。因此，当基督为真理作证时，祂当然是为自己作证，因为出自祂自己口中的话是：\BibleQuote{我是真理；}又如祂在另一处所说：\BibleQuote{我为自己作证。}但当祂说：\BibleQuote{凡属真理的人，必听我的声音}时，祂是在称赞那恩宠，祂凭此恩宠按自己的旨意召唤人。关于这个旨意，宗徒说：\BibleQuote{我们知道，凡爱天主的人，天主使一切协助他们获得益处，就是那些按天主的旨意蒙召的人}，\BibleRef{罗 8:28}即召唤者的旨意，而非被召者的旨意；这在另一处说得更清楚：\BibleQuote{你要依赖天主的大能，为福音共同劳苦，祂拯救了我们，以圣召召叫了我们，并非按照我们的行为，而是按照祂自己的旨意和恩宠。}\BibleRef{弟后 1:8}因为如果我们的思想转向我们所被造的本性，既然我们都是被真理所造，谁不属于真理呢？但并非所有人都得了真理的恩赐去听，即服从真理、相信真理；当然，绝没有任何先行功劳，免得恩宠不再是恩宠。因为如果祂说：\Quote{凡听我声音的，就是属于真理的}，那么就会有人认为他之所以被宣告属于真理，是因为他符合真理；然而祂所说的并非如此，而是：\BibleQuote{凡属真理的人，必听我的声音。}这样，他并非仅仅因为他听祂的声音而属于真理，而是他之所以听，是因为他属于真理，即这是真理赐予他的恩赐。这若不是藉基督的恩赐他信基督，又是什么呢？\par

5. \BibleQuote{比拉多对他说：\InnerQuote{什么是真理？}} 他没有等待回答：但\BibleQuote{说了这话，又出去到犹太人那里，对他们说：\InnerQuote{我在祂身上查不出什么罪来。但你们有个惯例，在逾越节我该给你们释放一个人；那么，你们愿意我给你们释放犹太人的君王吗？}} 我相信，当比拉多说\Quote{什么是真理？}时，他立刻想起了犹太人的惯例，即他在逾越节时常给他们释放一个人；因此他没有等待听耶稣对他那问题\Quote{什么是真理？}的回答，以免耽搁，因为他想起了那个惯例，他显然十分希望藉此在逾越节释放祂。但是，耶稣是犹太人的君王这一事实，却无法从他心中抹去，而是如同在标题上那样，被真理本身牢牢地固定在那里，而这真理正是他刚才问\Quote{什么是真理？}的真理。\BibleQuote{但他们听了这话，再次喊叫说：\InnerQuote{不要这人，而要巴拉巴！}巴拉巴原是个强盗。} 犹太人啊，我们并不责怪你们在逾越节释放有罪的人，而是责怪你们杀害无辜的人；然而，除非这事成就，真正的逾越节就不会发生。但误入歧途的犹太人保留了真理的影子，而藉着天主上智奇妙的安排，这同一影子的真理却被受骗的人实现了；因为为了举行真正的逾越节，基督被牵到祭献的屠宰场，如同羊一般。因此接下来是基督在比拉多及其兵士手下所受凌辱的记述；但这将在另一篇讲道中继续。\par

\SourceRecord{Translated by John Gibb.；Nicene and Post-Nicene Fathers, First Series；Vol. 7.；Edited by Philip Schaff.；Buffalo, NY: Christian Literature Publishing Co.,；1888.；<http://www.newadvent.org/fathers/1701115.htm>.}

% structure: p0001 source=h1 target=section reason=page-title

\section{Tractate 116 (John 19:1-16)}

1. 当犹太人在逾越节上喊叫说，他们不愿耶稣被释放给他们，而要强盗巴辣巴，不是救主，而是杀人凶手；不是生命的赐予者，而是毁灭者时，\BibleQuote{那时比拉多命人把耶稣带去鞭打了。}我们必须相信，比拉多这样做，没有别的原因，只是因为犹太人饱足了加给祂的凌辱，可能就满意了，不再疯狂地迫害祂，甚至要置祂于死地。怀着类似的目的，作为总督，他也允许（甚或命令，尽管福音作者对此保持沉默）他的营队做了随后的事。因为福音作者告诉我们士兵随后做了什么，却没有说比拉多命令了这些事。他说：\BibleQuote{士兵们用荆棘编了一个冠冕，放在祂头上，给祂披上一件紫红袍。他们来到祂跟前说：\InnerQuote{犹太人的君王，万岁！}并用手掌打祂。}这样，基督自己所预言的关于祂的事都实现了；这样，殉道者们被塑造，得以忍受迫害者乐意施加给他们的一切；这样，祂暂时隐藏了祂大能的威严，却把祂的忍耐的优先效法托付给了我们；这样，那不属于这世界的王国，不是藉着战斗的凶暴，而是藉着受苦的谦卑，战胜了那骄傲的世界；这样，那即将倍增的谷粒，在耻辱的恐怖中被播种，以便在荣耀的奇观中结出果实。\par

2. \BibleQuote{比拉多又出去，对他们说：\InnerQuote{看，我带祂出来，为使你们知道我查不出祂有什么罪。}于是耶稣戴着荆棘冠冕、披着紫红袍出来了。比拉多对他们说：\InnerQuote{看，这个人！}}由此可见，这些事是士兵们并非不知比拉多而做的，无论是他命令他们，还是仅仅允许他们，其原因正如我们上面所说的：为使祂的仇敌更愿意看到这种嘲弄的待遇，并止息他们对祂的血的进一步渴求。耶稣戴着荆棘冠冕、披着紫红袍出来，不是在君王的权力中闪耀，而是满负凌辱；对他们说的话是：\BibleQuote{看，这个人！}如果你们憎恨你们的君王，现在看见祂如此卑微，就饶恕祂吧；祂被鞭打，戴荆棘冠，被穿上嘲弄的衣服，受到最痛苦的辱骂，被手掌击打；祂的耻辱已达到顶点，但愿你们的恶意降至零点。但后者并没有冷却，反而更加炽热和猛烈。\par

3. \BibleQuote{因此，司祭长和差役看见祂，就喊叫说：\InnerQuote{钉死祂，钉死祂！}比拉多对他们说：\InnerQuote{你们自己把祂带去钉死吧，因为我查不出祂有什么罪。}犹太人回答他说：\InnerQuote{我们有法律，按那法律祂应该死，因为祂自命为天主子。}}看，还有另一个更大的仇恨理由。前者似乎是一件小事，即针对藉着非法的大胆行为篡夺王权；然而，这两样耶稣都没有虚假地声称拥有，但每一样都真实地属于祂，因为祂既是天主的独生子，又是被祂立为君王在祂的圣山熙雍上；现在祂本可以证明这两样都属祂，只是若与祂的大能相称，祂宁愿显明相应的大忍耐。\par

4. \BibleQuote{比拉多听到这话，越发害怕，又进了总督府，对耶稣说：\InnerQuote{你从那里来的？}耶稣却没有回答他。}比较所有福音作者的叙述，我们发现我们的主耶稣基督的这种沉默不止一次发生，既在司祭长面前，也在黑落德面前——路加暗示比拉多曾把耶稣送给他审问——还在比拉多面前；所以，关于祂的预言并非徒然预先说过：\BibleQuote{祂像在剪毛者前缄默的羔羊，祂不开口。}\BibleRef{依 53:7}，特别是在祂没有回答那些问话者的时候。因为虽然祂常常回答向祂提出的问题，但由于有些问题祂拒绝回答，羔羊的比喻就被提供了，为使在祂的沉默中，祂被看为无罪，而非有罪。因此，当祂经过审判的过程时，凡祂不开口之处，祂是作为羔羊这样做的；也就是说，不是作为一个良心有罪、被自己的罪证定罪的人，而是作为一个以祂的温良为别人的罪被祭献的人。\par

5. \BibleQuote{比拉多于是对他说：\InnerQuote{你不向我说话吗？你不知道我有权钉死你，也有权释放你吗？}耶稣回答说：\InnerQuote{除非从上头赐给你，你对我毫无权柄；因此，把我交付给你的人，罪过更大。}} 你看，祂回答了；然而在祂不回答的地方，并不是像一个罪犯或狡诈之人，而是如同一只羔羊；就是说，祂在纯朴与无辜中不开口。因此，在未作回答之处，祂如羊般缄默；在回答之处，祂如牧者般教导。所以，让我们努力去明白祂所说的，以及祂也借着宗徒所教导的：\BibleQuote{没有权柄不是出于天主的};\BibleRef{罗 13:1}，而且那个恶意将无辜者交付于有权力者去杀害的人，其罪过比那权力本身（若它因惧怕另一个更大的权力而杀害他）更大。确实，天主赐给比拉多的权力是这样的，以至于他也处于凯撒的权下。因此祂说：\BibleQuote{你对我毫无权柄}，就是说，连你真正拥有的那一点微小的权力，\BibleQuote{除非}这权力，不论有多大，\BibleQuote{从上头赐给你}。但我清楚其大小，因为它并不大得足以让你完全独立自主，\BibleQuote{因此，把我交付给你的人，罪过更大。}的确，他是出于嫉妒将我交付给你的权力，而你却要出于恐惧来行使你的权力对付我。然而，即使出于恐惧，人也不该杀害另一个人，尤其是无辜者；但出于多管闲事的热忱去这样做，比起受恐惧的约束，是更大的恶。因此，这位真理的导师不是说：\BibleQuote{把我交付给你的人}只有他有罪，好像另一个人无罪；而是说：\BibleQuote{罪过更大}，让他明白自己并非没有责任。因为后者的罪并没有因为前者更大而化为乌有。\par

6. \BibleQuote{从此，比拉多设法要释放祂。} 这里所用的词\Quote{从此}应如何理解，好像他先前并没有设法释放祂？读一读前面的经文，你就会发现他已经有一段时间在设法释放耶稣了。因此，我们应该把原文的词理解为\Emph{基于此}，即\Emph{为此缘故}——为了不让自己因杀害一个被交付到他手中的无辜者而犯罪，尽管他的罪要比那些把他交付给他去处死的犹太人的罪小。\BibleQuote{从此}，因此，即为此缘故，为不犯这样的罪，他\BibleQuote{设法}不是现在第一次，而是从一开始就\BibleQuote{要释放祂}。\par

7. \BibleQuote{但犹太人喊叫说：\InnerQuote{你若释放这人，就不是凯撒的朋友；因为凡自命为王的，就是反对凯撒。}}他们以为，用关于凯撒的恐吓对比拉多施加更大的恐惧，好让他处死基督，比之前他们说的\BibleQuote{我们有法律，按法律他应该死，因为他自充为天主子}更有力。确实，不是他们的法律因恐惧驱使他去杀人，而是他对天主之子的恐惧阻止他犯罪。但现在他不能像漠视别国的法律那样，漠视凯撒——他自己权力的来源。\par

8. 然而，福音作者继续叙述：\BibleQuote{比拉多听了这些话，就带耶稣出来，到了一个名叫石铺地——希伯来话叫\OriginalTerm{加巴达}{Gabbatha}——的地方，坐在审判座上。那天正是逾越节的预备日，约莫第六时辰。}关于主被钉死在哪一个时辰的问题，由于另一位福音作者提供证据，他说：\BibleQuote{他们把祂钉在十字架上时，是第三时辰}，\BibleRef{谷 15:25}，我们将尽可能在走到记载祂受难的段落时，若主许可，加以考虑。因此，当比拉多坐在审判座上后，\BibleQuote{他对犹太人说：\InnerQuote{看，你们的君王！}但他们喊叫说：\InnerQuote{除掉，除掉，钉他在十字架上！}比拉多对他们说：\InnerQuote{要我把你们的君王钉在十字架上吗？}}他仍试图克服他们关于凯撒带给他的恐惧，想藉着说明他们这样做的无耻来打破他们的企图，说了\BibleQuote{要我把你们的君王钉在十字架上吗？}当他未能以对基督的羞辱来软化他们时，他最终被恐惧所压倒。\par

9. 因为\BibleQuote{司祭长回答说：\Quote{除了凯撒，我们没有君王。}于是比拉多把祂交给他们去钉十字架。}因为如果他们宣称除了凯撒没有别的君王，而他却想通过释放那个因这些图谋而被交给处死的人，从而将另一个君王强加给他们，那就会显得他是在对抗凯撒。\BibleQuote{于是他把祂交给他们去钉十字架。}然而，这与祂先前所说的\BibleQuote{你们把祂带去钉十字架；}或更早所说的\BibleQuote{你们把祂带去，按照你们的法律审判祂？}有何不同呢？为什么当他们说\BibleQuote{我们不可处死任何人，}时表现如此大的抵触，千方百计地要祂不被他们自己杀死，而是由总督处死，因此拒绝接受祂以便处死祂，而现在他们却实际接受祂来达到同样的目的？如果不是这样，为什么说\BibleQuote{于是他把祂交给他们去钉十字架？}这又有什么重要呢？显然重要。因为所说的不是\BibleQuote{于是他把祂交给他们}让他们去钉祂十字架，而是\BibleQuote{祂要被钉十字架，}即祂要借总督的判决和权力被钉十字架。但福音作者之所以说祂被交给他们，正是要表明他们参与了自己试图置身事外的罪行；因为比拉多绝不会这样做，除非他察觉到这是他们不可动摇的意愿。然而接下来的话：\BibleQuote{他们就把耶稣带去了，}现在可能指的是总督的随从士兵。因为后面更明确地说：\BibleQuote{当兵士们把祂钉在十字架上以后，}尽管福音作者合理地这样做，即使他把一切归咎于犹太人，因为犹太人接受了自己贪婪所求的，他们做了自己所迫使的一切。但接下来的事件必须在另一篇讲道中加以考虑。\par

\SourceRecord{Translated by John Gibb.；Nicene and Post-Nicene Fathers, First Series；Vol. 7.；Edited by Philip Schaff.；Buffalo, NY: Christian Literature Publishing Co.,；1888.；<http://www.newadvent.org/fathers/1701116.htm>.}

% structure: p0001 source=h1 target=section reason=page-title

\section{讲道集 117（若望福音 19:17-22）}

1. 在比拉多于审判席前作出判决和定罪之后，他们带走了主耶稣基督，大约在第六时辰，把祂带走了。\BibleQuote{祂背着自己的十字架，出来到了一个地方，名叫髑髅，希伯来语叫哥耳哥达；他们就在那里把祂钉在十字架上。} 那么，马尔谷福音所说的\BibleQuote{到了第三时辰，他们就把他钉在十字架上}\BibleRef{谷 15:25}，又作何解释呢？岂非是说：主在第三时辰被犹太人的舌头钉死，在第六时辰被士兵的手钉死？这样我们可以理解为：当比拉多登上审判席时，第五时辰已经结束，第六时辰刚开始，若望所说的\BibleQuote{大约第六时辰}即指此时；然后当祂被带出去，与两个强盗一起被钉在木架上，以及记载的其它事件在祂的十字架旁发生时，第六时辰完全结束；而从这一时辰到第九时辰，太阳昏暗，黑暗降临，我们有三位福音书作者——玛窦、马尔谷和路加的共同见证。但因犹太人企图把杀害基督的罪责从自己身上转嫁到罗马人，即比拉多和他的士兵身上，所以马尔谷省略了基督被士兵钉在十字架上的时辰（那正是第六时辰的开始），反而特意记下了第三时辰，因为犹太人在那时曾在比拉多面前喊叫\BibleQuote{钉死他，钉死他}\BibleRef{若 19:6}，这样不仅可以看出是前者（即第六时辰把祂挂在木架上的士兵）把耶稣钉死，也可以看出犹太人（在第三时辰呼喊把祂钉死）也有份。\par

2. 这个难题还有另一种解释：这里不应理解为白天的第六时辰，因为若望说的不是\BibleQuote{那天大约是第六时辰}或\BibleQuote{大约第六时辰}，而是说\BibleQuote{正是逾越节的\Emph{预备日}，大约第六时辰}\BibleRef{若 19:14}。而\OriginalTerm{预备日}{parasceve} 在拉丁文里是 \OriginalTerm{准备}{præparatio}；但犹太人，即使是那些说拉丁语而非希腊语的，也喜欢在这类礼仪中使用希腊词。因此这是逾越节的预备日。但\BibleQuote{我们的逾越节——基督，已被祭献了}\BibleRef{格前 5:7}，正如宗徒所说；如果我们将这个逾越节的预备日从夜晚第九时辰算起（因为那时司祭长们似乎已经判决了主的祭献，当时他们说\BibleQuote{他该死}\BibleRef{玛 26:66}，并且案件仍在司祭长家中审理；因此有一种和谐的理解：真逾越节的预备日，也就是基督被祭献的预备，从此开始——那是犹太人逾越节的预像——当司祭们宣判祂该被祭献时），那么，从那个据说大约是夜晚第九时辰的时刻，到白天第三时辰（马尔谷证实基督被钉在十字架上），共有六小时：夜间的三小时和白天的三小时。因此，对于这个逾越节的预备日，即基督祭献的预备（从夜晚第九时辰开始），大约是第六时辰；也就是说，第五时辰已结束，第六时辰已开始，此时比拉多登上了审判席：因为这同一个预备，从夜晚第九时辰开始，一直持续到基督的祭献（即在预备中的事件）完成，而这一完成发生在第三时辰——按照马尔谷，不是预备日的第三时辰，而是白天的第三时辰；但同时也是第六时辰，不是白天的第六时辰，而是预备日的第六时辰，当然是从夜晚第九时辰到白天第三时辰的六个小时。对于这个难题的这两种解释，各人可以自行选择喜欢的一种。但谁若阅读了非常详尽的关于\WorkTitle{福音书的和谐}的讨论，就能更好地判断该选哪一个。如果还能找到其他解释，那么福音真理的稳固性将赢得累积的辩护，以对抗不信和渎神虚妄的诽谤。现在，在简略讨论之后，让我们回到若望福音的叙述。\par

3. \BibleQuote{他们就把耶稣带了去，} 他说，\BibleQuote{带走了；祂背着自己的十字架，出来到了一个地方，名叫髑髅，希伯来语叫哥耳哥达；他们就在那里把祂钉在十字架上。} 因此，耶稣背着自己的十字架，走向祂要被钉死的地方。一个壮观的景象！但若观看者是亵渎者，那就是绝大的笑柄；若是虔诚者，那就是极大的奥迹：若观看者是亵渎者，那就是极为可耻的显示；若是虔诚者，那就是坚固的信德堡垒：若观看者是亵渎者，他就嘲笑那位王，祂手中拿的不是王杖，而是祂受刑的木架；若是虔诚者在观看，他就看见这位王背着祂自己的十字架，而这十字架祂还要放在众王的额上，虽然遭受亵渎者的轻蔑一瞥，但圣人们的心日后却要以此夸耀。因为保罗日后要说：\BibleQuote{至于我，我只以我们的主耶稣基督的十字架来夸耀}\BibleRef{迦 6:14}，祂正是以自己的肩膀扛起这同一十字架，背着即将燃烧的光明之灯台，且不会放在斗底下\BibleRef{玛 5:15}。因此，\BibleQuote{祂背着自己的十字架，出来到了一个地方，名叫髑髅，希伯来语叫哥耳哥达；他们就在那里把祂钉在十字架上，并同祂一起钉了两个人，一边一个，耶稣在中间。} 这两个人，正如我们从其他福音书作者的叙述中得知的，是和祂一同被钉的强盗，祂被钉在他们中间；先前的预言早已宣告：\BibleQuote{祂被列于叛逆者之中}\BibleRef{依 53:12}。\par

4. \BibleQuote{\Person{比拉多}又写了一个牌子，放在十字架上，写的是：\InnerQuote{\Place{纳匝肋}人\Person{耶稣}，犹太人的君王}。这个牌子有许多犹太人读了，因为\Person{耶稣}被钉十字架的地方离城很近，它是用希伯来文、希腊文、拉丁文写的：\InnerQuote{犹太人的君王}。}因为在这地方，这三种语言超越所有其他语言显得突出：希伯来文是为了犹太人，他们因天主的法律而夸耀；希腊文是为了外邦人中的智者；拉丁文是为了罗马人，他们当时正统治着许多、几乎所有的国家。\par

5. \BibleQuote{那时，犹太人的司祭长对\Person{比拉多}说：\InnerQuote{不要写\InnerQuote{犹太人的君王}，但要写他说：\InnerQuote{我是犹太人的君王}}。\Person{比拉多}回答说：\InnerQuote{我所写的，我已经写了。}}哦，天主作为的不可言喻的大能，甚至在无知者的心中！难道没有一种隐藏的声音，以一种（如果可以这样说）响亮的沉默，透过\Person{比拉多}的内心，说出了在圣咏中早已预言的话语：\BibleQuote{不要破坏牌子的题字}？看，他没有破坏；他所写的，他已经写了。但是司祭长们希望它被破坏，他们说了什么？\BibleQuote{\InnerQuote{不要写\InnerQuote{犹太人的君王}，但要写他说：\InnerQuote{我是犹太人的君王}}}。什么事，疯狂的人们，你们在说什么？为什么你们反对做那件你们完全无法改变的事？难道因此\Person{耶稣}所说的\BibleQuote{我是犹太人的君王}就变得不那么真实了吗？如果\Person{比拉多}所写的不能篡改，那么真理所发出的言语就能篡改吗？但\Person{基督}只是犹太人的君王，还是也是外邦人的君王？是的，也是外邦人的。因为当他在预言中说：\BibleQuote{我被祂立为君王，在祂的圣山熙雍上，宣告上主的法令}，这样没有人会因为\Place{熙雍}山而说他只是被立为犹太人的君王，他立刻接着说：\BibleQuote{上主对我说：你是我的儿子，我今日生了你。你向我请求，我就将外邦人赐你为产业，地极赐你为基业。}所以他自己，现在在犹太人中间用自己的口说：\BibleQuote{我还有别的羊，不是这圈里的；我也必须把他们带来，他们要听我的声音，并且要成为一群，一个牧人。}那么，为什么我们要在这个写上\Quote{犹太人的君王}的题字中理解某种伟大的奥秘呢，如果\Person{基督}也是外邦人的君王？因为这个缘故，因为野橄榄枝接上了橄榄树的肥脂，而不是橄榄树接上了野橄榄的苦汁。\BibleRef{罗 11:17}因此，既然\Quote{犹太人的君王}这个头衔真实地写在\Person{基督}身上，那么应该被理解为\Quote{犹太人}的是谁呢？岂不就是\Person{亚巴郎}的后裔，应许的儿女，他们也是天主的儿女？因为宗徒说：\BibleQuote{他们，那些按肉身作儿女的，不是天主的儿女；但应许的儿女才算是后裔。}\BibleRef{罗 9:7}而那些外邦人，他对他们说：\BibleQuote{如果你们属于基督，那么你们就是\Person{亚巴郎}的后裔，按照应许是承继人。}\BibleRef{迦 3:29}\Person{基督}因此是犹太人的君王，但那些在心里受割损、按精神而非按文字、他们的赞誉不是来自人而是来自天主的犹太人；\BibleRef{罗 2:29}他们属于那自由的\Place{耶路撒冷}，我们在天上的永恒母亲，精神的\Person{撒辣}，她把婢女和她的孩子从自由的家中赶出去。\BibleRef{迦 4:22}因此，\Person{比拉多}所写的，他写了，因为主所说的，他说了。\par

\SourceRecord{Translated by John Gibb.；Nicene and Post-Nicene Fathers, First Series；Vol. 7.；Edited by Philip Schaff.；Buffalo, NY: Christian Literature Publishing Co.,；1888.；<http://www.newadvent.org/fathers/1701117.htm>.}

% structure: p0001 source=h1 target=section reason=page-title

\section{讲道集118（\BibleRef{若 19:23-24}）}

1. 我们现在靠主的助佑，在本讲道中要讨论主被钉十字架后所发生的事。\BibleQuote{那时，士兵们把耶稣钉在十字架上后，拿了他的衣服，分成四份，每个士兵一份；又拿了他的内衣；那内衣是无缝的，从上到下浑然织成。他们便彼此说：我们不要撕破它，我们抽签，看归谁。这应验了经上的话：\InnerQuote{他们分了我的外衣，为我的内衣拈阄。}\BibleRef{咏 22:19}}\BibleRef{若 19:23-24} 这事的发生正如犹太人所愿；其实不是他们自己，而是服从比拉多（他亲自担任审判官）的士兵钉死了耶稣。但若我们反省他们的意愿、他们的阴谋、他们的努力、他们的交付，以及他们强逼的喊声，那么钉死耶稣的，确实是犹太人，而非其他人。\par

2. 但我们不可轻率地谈论分祂衣服和拈阄的事。因为虽然四位福音作者都提到了这事，但其他三位比若望更简略；他们的记述隐晦，而若望的记述最为清晰。因为玛窦说：\BibleQuote{他们把祂钉死后，分了祂的衣服，拈阄。} \BibleRef{玛 27:35} 马尔谷说：\BibleQuote{他们把祂钉在十字架上，分了祂的衣服，拈阄看谁得什么。} \BibleRef{谷 15:24} 路加说：\BibleQuote{他们分了祂的衣服，拈阄。} \BibleRef{路 23:34} 但若望也告诉我们，他们把祂的衣服分成了几份，即四份，每人一份。由此可知，有四个士兵执行总督的命令钉死祂。因为他清楚地说：\BibleQuote{那时，士兵们把耶稣钉在十字架上后，拿了他的衣服，分成四份，每个士兵一份；也拿了内衣。} 这里省略了\Quote{他们拿了}，意思是他们拿了祂的衣服，分成四份，每个士兵一份；他们也拿了祂的内衣。他这样说，让我们看到，在祂的其他衣服上没有拈阄；但祂的内衣，他们与其他衣服一起拿了，却没有同样分开。因为关于这件内衣，他接着解释：\BibleQuote{那内衣是无缝的，从上到下浑然织成。} 然后告诉我们为什么他们拈阄，他说：\BibleQuote{他们彼此说：我们不要撕破它，我们抽签，看归谁。} 由此清楚，在其他衣服上，他们各得一份，不需要拈阄；但关于这件，若不撕破，就不能每人一份，从而只得到无用的碎片；为避免此，他们宁愿让它通过拈阄归给其中一人。这位福音作者的叙述也与预言的见证相符，他立刻加上说：\BibleQuote{这应验了经上的话：\InnerQuote{他们分了我的外衣，为我的内衣拈阄。}\BibleRef{咏 22:19}} 因为祂不是说\Quote{他们拈阄}，而是\Quote{他们分了}；祂也不是说\Quote{拈阄分了}；而是，对于其他衣服，祂完全没有提拈阄，之后祂说\Quote{为我的内衣拈阄}，只指那件剩余的内衣。关于这事，待我先驳斥那可能出现的诽谤，即认为福音作者们彼此矛盾，通过指出其他任何一人的话与若望的叙述并不冲突后，按主所赐的能力再讲。\par

3. 因为玛窦说：\BibleQuote{分了祂的衣服，拈阄}，他愿意人理解，在整个分衣服的事上，那件内衣也包括在内，他们拈阄就是为它；因为在分所有衣服的过程中，它也是其中一件，而唯有它他们拈了阄。路加的话也有同样的目的：\BibleQuote{分了祂的衣服，拈阄}；因为在分的过程中，他们遇到了要拈阄的那件内衣，以便完成祂的衣服在他们之间的完全分配。那么，照路加说\Quote{分的时候拈阄}，或照玛窦说\Quote{分了，拈阄}，有什么区别呢？除非路加说\Quote{lots}（阄）时用了复数代替单数——这种说法在圣经中并不罕见，虽然有些抄本有\Quote{lot}（阄）而不是\Quote{lots}。所以，只有马尔谷似乎带来了某种困难；因为他说：\BibleQuote{拈阄，看谁得什么}，他的话似乎暗示阄是拈在所有衣服上，而不仅是那件内衣。但这里也是简略造成隐晦；因为\Quote{拈阄}这个词，就好像是说，在他们进行分配的过程中拈阄；事实也是如此。因为所有衣服的分割不会完整，除非通过拈阄决定他们中谁也应得到那件内衣，从而结束或更确切地防止分者之间任何争执的发生。因此，当他说\Quote{看谁得什么}时，就拈阄而言，我们不可将它理解为指所有被分的衣服；因为拈阄是为了决定谁得那件内衣：由于他省略了描述具体形式，以及如何在平均分配各部分后，那件内衣单独留下，为免撕破而由拈阄决定归属，所以他使用了\Quote{看谁得什么}这个说法，换句话说，就是谁该得到它：就好像整个是这样表达的：他们分了祂的衣服，拈阄，看谁得那件内衣，它是在他们平等分得其余部分后剩下的。\par

4. 也许有人会问，将祂的衣服分成这么多部分，以及为长袍拈阄，这有什么含义呢。主耶稣基督的衣服分成四份，象征祂的四分教会，即遍布于由四部分组成的世界，并且均匀地，也就是说，和谐地分布于这四方。为此，祂在别处说，祂将派遣天使从四方聚集祂的选民：\BibleRef{玛 24:31}——这不是指从世界的四方：东、西、北、南，又是什么？但拈阄所得的长袍，则象征所有部分的合一，这合一包含在爱德的联系中。当宗徒要谈论爱德时，他说：\Quote{我指给你们一条更优越的道路；} \BibleRef{格前 12:31} 又在另一处说：\Quote{并认识基督那远超知识的爱；} \BibleRef{弗 3:19} 并且还在别处说：\Quote{在这一切之上，还要有爱德，它是完善的联系。} \BibleRef{哥 3:14} 既然爱德既有更优越的道路，又远超知识，且在一切之上被命令，那么那以之象征的衣服被描述为从上到下编织而成，就非常恰当了。而且它没有缝，免得它的缝线被分开；它归于一人，因为祂将众人聚为一体。正如宗徒的情形，他们恰好是十二人，换言之，可分四个三人组，当向所有人提问时，只有伯多禄回答：\Quote{祢是基督，永生天主之子；} 并向他说：\Quote{我要把天国的钥匙交给你，} 仿佛他独自接受了束缚和释放的权力：这样，既然一人代表众人发言，并与众人一同接受后者，如同代表合一本身；因此一人代表众人，因为合一在众人之中。因此，在说过\Quote{从上到下编织}之后，他又加上\Quote{遍及全衣。} 这若按其含义而言，也意味着凡被发现属于整体的人，无人被排除在分享之外：从这整体，正如希腊文所示，教会得名为公教会。而拈阄所赞扬的，不是天主的恩宠，又是什么？因为这样，藉着一个人的位格，它达到了众人，因为阄满足了众人，同样，天主的恩宠在其合一中达到众人；当拈阄时，赏赐的决定不是根据各人的功绩，而是根据天主隐秘的判决。\par

5. 但仍不要说这些事是因为恶人所作的，即不是因跟随基督的人，而是因迫害祂的人，就没有好的意义。因为对于十字架本身，我们都知道它同样是由仇敌和罪人制造并钉在基督身上的，我们又能说什么呢？然而，正是十字架，我们可以正确地理解宗徒的话适用其上：\Quote{那宽、长、高、深是什么。} \BibleRef{弗 3:18} 因为其宽度在于横木，被钉者的双手伸展其上，象征善工在爱的全部广度中；其长度从横木延伸到地面，是背部和双脚所固定的地方，象征从时间之初到末了的恒心；其高度在于顶端，向上高出横木，象征超越的目标，所有善工都指向它，因为一切在宽度和长度上做得好且持之以恒的事，也要相应顾及神圣赏报的崇高性质；其深度在于插入地面的部分，因为那里既隐藏又看不见，然而所有突出且可感知的部分都由此生出；正如我们内的一切善都出自天主恩宠的深奥，这深奥非人的理智和判断所能及。但即使基督的十字架所象征的，并不比宗徒所说的更多：\Quote{那些属于耶稣基督的人，已经把肉身连同私欲偏情钉在十字架上了，} \BibleRef{迦 5:24} 这是多么大的善事！然而，除非善灵与肉身相争，否则十字架并不成就这事，因为正是敌对的，换言之，恶灵建造了基督的十字架。最后，如众所知，基督的标记若不是基督的十字架，又是什么？因为除非那标记被应用，无论印在信徒的额上，还是在他们重生所用的水上，或在他们接受傅油所用的油上，或在滋养他们的祭献上，否则没有一件能恰当地施行。那么，由恶人所作的事怎能不象征任何善事呢？因为藉着恶人所造的基督的十字架，在我们举行祂的圣事时，每一件善事为我们封印了。但我们在此停住；接下来的事，我们将在下次讨论中，如天主赐予我们助佑，再作考虑。\par

\SourceRecord{Translated by John Gibb.；Nicene and Post-Nicene Fathers, First Series；Vol. 7.；Edited by Philip Schaff.；Buffalo, NY: Christian Literature Publishing Co.,；1888.；<http://www.newadvent.org/fathers/1701118.htm>.}

% structure: p0001 source=h1 target=section reason=page-title

\section{论著119（\BibleRef{若 19:24-30}）}

1. 主既已被钉在十字架上，祂的衣裳也由拈阄分配完毕，现在我们来看福音圣史若望随后所记载的事。他说：\Quote{兵士就这样做了。在耶稣的十字架旁，站着祂的母亲、祂母亲的姐妹克罗帕的妻子玛利亚、以及玛利亚玛达肋纳。耶稣看见母亲和祂所爱的门徒站在旁边，就对母亲说：\InnerQuote{女人，看，你的儿子！}然后又对那门徒说：\InnerQuote{看，你的母亲！}就从那时起，那门徒把她接到自己家里去。}无疑地，这正是耶稣在将要变水为酒时对祂母亲所说的那个时辰：\Quote{女人，这于我和你有什么关系？我的时辰尚未来到。}因此，祂预言的这个时辰，那时尚未来到，就是祂要在死亡时认她的时辰，而祂也正是为此作为有死之人而生的。那时，祂将要行天主性的事迹，便如同不认识一般地推开她——她不是祂天主性的母亲，而是祂人性软弱的母亲；但如今，在人性痛苦之中，祂以属人的情感托付了那使祂成为人的母亲。因为那时，创造玛利亚的祂以大能显明了；而此刻，玛利亚所生的祂正悬在十字架上。\par

2. 因此，这里插入了一段道德的教训。善师教导我们当做之事，并以祂自己的榜样教训门徒：孝爱儿女应当关心父母，正如那悬着将死之人的木头，正是导师施教的座位。宗徒保禄正是从这有益的教训学得他所教导的，他说：\Quote{若有人不照顾自己的亲属，尤其不照顾自己的家人，就是背弃了信仰，比不信的人更坏。}（\BibleRef{弟前 5:8}）对任何人来说，有什么比父母之于子女、子女之于父母更属于自己家的事呢？圣人的导师亲身以身作则，立下了这最有益的诫命：祂不是作为天主为自己所创造并治理的婢女，而是作为人为那生祂的、如今要留下的母亲，在某种意义上另立了一个儿子替代自己。祂为何这样做，在接下来的话中表明了：福音圣史说：\Quote{就从那时起，那门徒把她接到自己家里。}这是指着他自己说的。他通常这样称自己为耶稣所爱的门徒：主确实爱他们众人，但爱他超越其余，且更为亲密，甚至使他在晚餐时靠在自己的怀里；我相信，这是为了更显明地推崇这部福音的天主性卓越，祂以后要藉着他宣讲出来。\par

3. 但若望接去主的母亲，这\Quote{自己的家}是什么呢？他并不在那些对主说\Quote{看，我们舍弃了一切，跟随了你}的人之外。相反，他当时也听见这话：\Quote{凡为我的缘故舍弃这些的，在今世必得百倍。}因此，那门徒所获得的，比起他所舍弃的，已多了一百倍，可以接待那位赐予这一切恩惠之主的母亲。原来，真福若望是在那个社会中获得百倍的，那里没有人称什么为自己的，却一切公用，正如\WorkTitle{宗徒大事录}所载。宗徒们好似一无所有，却拥有万有（\BibleRef{格后 6:10}）。那么，门徒和仆人如何能接他的主和导师的母亲到\Quote{自己的家}里，而那里没有人称什么为自己的呢？或者，我们在这同一部书稍后读到：\Quote{凡拥有田地或房屋的，都卖了，把价银拿来放在宗徒脚前，照每人的需要分配}（\BibleRef{宗 4:32}），难道我们不应理解为，分给这门徒所需之物时，真福玛利亚的份也加给了他，好像她是他的母亲；而\Quote{从那时起，那门徒把她接到自己家里}这句话，我们更应理解为：一切所需都托付给他照料吗？因此，他不是接到自己的田地（因为他没有自己的田地），而是接到自己孝爱的服务，这服务的执行，藉着特殊的安排，托付给了他。\par

4. 然后他接着说：\BibleQuote{此后，耶稣知道一切事都已成就，为使经上的话应验，就说：我渴了。那里放着一个盛满了醋的器皿；他们便将海绵蘸满了醋，绑在牛膝草上，送到祂口边。耶稣一尝了那醋，便说：完成了。就低下头，交付了灵魂。}谁能像这人那样，有能力如此安排自己所行的一切，正如祂安排自己所受的一切？但这个人就是天主与人之间的中保；这人在预言中被记载说：祂也是人，谁能认识祂呢？因为做这些事的人并未承认这人是天主。因为那显现为人的，是隐藏的天主；那显现的祂承受了这一切，而隐藏的祂亲自安排了这一切。因此，祂看见所有在祂接受醋、交付灵魂之前必须完成的事都已成就；并且为使经上预言的话也应验，祂说：\BibleQuote{他们在我的渴中给我醋喝}，祂说：\Quote{我渴了}，好像是在说：还有一件事你们没有做到，给我你们所是的。因为犹太人自己就是那醋，他们从先祖和先知们的酒中堕落；像充满的器皿一样充满此世的邪恶，心像海绵，在其洞穴和曲折的隐蔽处有诡诈。但牛膝草，他们放置蘸满醋的海绵在上面，是一种卑微的草药，能洁净心脏，我们适当地将其视为基督自己的谦卑；他们这样包围了它，以为完全困住了它。因此诗篇上说：\BibleQuote{求你用牛膝草洁净我，我就洁净}。因为我们是因基督的谦卑而得洁净；因为如果祂没有谦抑自己，服从至死，且死在十字架上\BibleRef{斐 2:8}，祂的血定然不会为赦免罪过而倾流，也就是说，为我们的洁净。\par

5. 我们也不必因这个问题而困扰：当祂被举离地面钉在十字架上时，海绵怎能送到祂口边？因为正如我们在其他福音书中所读到的（这本福音书省略了这一点），海绵是绑在芦苇上，以便将海绵中所盛的饮品送到十字架的顶端。然而，芦苇所象征的是圣经，这行动正应验了经上的话。因为正如舌头被称为希腊语、拉丁语或其他语言，指的是舌头所发出的声音；同样，芦苇也可以以其命名由芦苇写成的文字。然而，我们通常称那些表达人声的为语言；而称圣经为芦苇，这事的稀罕性更增强了它所象征的奥秘性质。一个邪恶的百姓做了这些事，一位慈悲的基督承受了它们。那些做这些事的人不知道自己在做什么；但那位承受苦难的祂，不仅知道所做的事以及为何如此，还藉着那些行恶的人，成就了美善。\par

6. \BibleQuote{耶稣一尝了那醋，便说：完成了。}除此以外，还有什么呢？不就是预言许久以前所预言的一切吗？这时，因为祂死前再没有什么需要完成的了，仿佛祂——那有权舍弃自己生命并重新取回它——终于完成了祂所等待的一切，\BibleQuote{就低下头，交付了灵魂。}有谁能像耶稣那样，想睡时就睡？有谁能像祂那样，想脱去衣服时就脱去？有谁能像祂那样，想离开时就离开？如果祂在死时展现了这样的权能，那么作为审判者，祂的权能该有多么伟大——无论我们盼望它还是畏惧它！\par

\SourceRecord{Translated by John Gibb.；Nicene and Post-Nicene Fathers, First Series；Vol. 7.；Edited by Philip Schaff.；Buffalo, NY: Christian Literature Publishing Co.,；1888.；<http://www.newadvent.org/fathers/1701119.htm>.}

% structure: p0001 source=h1 target=section reason=page-title

\section{《若望福音释义》第一二〇篇（若 19:31-20:9）}

1. 主耶稣既已完成了祂在受死以前预知必须完成的一切，并在祂自己愿意的时刻交付了灵魂，此后所发生的事，正如福音所记载的，现在我们当加以思考。他说：\Quote{犹太人因那日是预备日（\OriginalTerm{预备日}{parasceve}），免得安息日（那安息日是个大日）尸首留在十字架上，就求比拉多打断他们的腿，将他们拿去。}这不是说将腿拿去，而是指那些被打断腿以便致死其人的本人被拿去，使他们从木架上被取下，免得他们继续悬在十字架上，以其终日的痛苦景象污秽这大节日。\par

2. \Quote{于是兵士来了，将第一个人的腿，以及与祂同钉的另一个人的腿，都打断了。但当他们来到耶稣面前，看见祂已经死了，就没有打断祂的腿；但有一个兵士用长枪刺开了祂的肋旁，立刻流出血和水来。}福音记述者运用了一个意味深长的词，他没有说\OriginalTerm{刺穿}{pierced}或\OriginalTerm{刺伤}{wounded}祂的肋旁，或别的，而是说\Quote{刺开}，好使，在某种意义上，生命之门得以敞开，从那里涌流出教会的圣事，没有这些圣事，便无法进入那真实的生命。那血是为赦罪而流；那水构成了赐健康的杯，同时提供了圣洗池和饮用的水。这在诺厄受命在方舟的侧面开一道门时便已预表，\BibleRef{创 6:16} 使那些不注定在洪水中灭亡的动物得以进入，而方舟预表了教会。正因为此，第一个女人是从睡着的人的肋骨形成的，\BibleRef{创 2:22} 并被称为生命和众生的母亲。\BibleRef{创 3:20} 这确实指向一件伟大的善事，先于那伟大的悖逆之恶（以如此睡着者的形象）。这第二位亚当垂下头，在十字架上睡着，为要使一位配偶从睡着者的肋旁流出的东西形成。啊，死亡！你使死者重新获得生命！有什么比这血更纯洁？有什么比这创伤更给予健康？\par

3. \Quote{看见这事的人，}他说，\Quote{就作证，他的见证是真实的；并且他知道自己所说的真实，为使你们也相信。}他没有说：为使你们也知道，而是说\Quote{为使你们相信}；因为看见的人知道，那没有看见的人可以相信他的见证。而相信更属于信德的本性，而非看见。因为相信岂不就是给信德以适当的接纳吗？\Quote{这些事发生了，}他接着又说，\Quote{是为应验经上的话：你们不可折断祂的一根骨头。另一处经上又说：他们要仰望他们所刺透的那位。}他为所记载的每件事都提供了两个圣经的见证。因为对于\Quote{但当他们来到耶稣面前，看见祂已经死了，就没有打断祂的腿}这句话，对应的见证是\Quote{你们不可折断祂的一根骨头}：这是古时律法中那命令以羔羊为祭献庆祝逾越节的人所领受的诫命，那羔羊预表了基督的受难。因此，\Quote{我们的逾越节，基督，已被祭献，}\BibleRef{格前 5:7} 关于祂，依撒意亚先知也曾预言：\Quote{祂如同羔羊被牵去宰杀。}\BibleRef{依 53:7} 同样，对于他随后加上的话\Quote{但有一个兵士用长枪刺开了祂的肋旁}，对应的另一个见证是\Quote{他们要仰望他们所刺透的那位}；此处基督是以那祂后来要被钉死的同一肉身被预许的。\par

4. \Quote{此后，阿黎玛特雅人若瑟（他是耶稣的门徒，但因怕犹太人而暗暗地）求比拉多，要取耶稣的身体；比拉多准许了他。于是他来了，取走了耶稣的身体。又来了尼苛德摩，他起初是在夜间来到耶稣那里的，带着没药和沉香调和的香料，大约一百斤。}我们不应以解释为\Quote{起初带着没药}来理解，而应将\Quote{起初}一词归属于前面的分句。因为尼苛德摩起初是在夜间来到耶稣那里的，正如同一若望在他福音的前面部分所记载的。因此，根据此处给我们的叙述，我们应当理解为：尼苛德摩来到耶稣那里，并非仅在那一次，而是那时是第一次；之后他经常来，为借聆听而成为门徒；这为如今，至少对几乎所有民族而言，在蒙福的斯德望身体的启示中，得到了证实。\Quote{他们就取了耶稣的身体，用细麻布和香料裹好，如犹太人的殡葬习俗。}我想，福音记述者并非毫无用意地如此措辞：\Quote{如犹太人的殡葬习俗}；因为照我看来，他以此提醒我们，在对待死者所行的此类义务中，各民族的习俗皆当遵守。\par

5. \Quote{在祂被钉的地方有一个园子，园子里有一座新坟墓，里面从未有人安放过。}正如在童贞玛利亚的子宫中，祂以前无人受孕，祂以后也无人受孕，同样，在这坟墓中，祂以前无人埋葬，祂以后也无人埋葬。\Quote{因为那日是犹太人的预备日，又因那坟墓近了，他们就把耶稣安放在那里。}他要我们明白，葬埋是匆忙的，唯恐暮色降临；那时因预备日，不再许可做任何这类的事，这预备日我们之中的犹太人更惯于以拉丁文称之为\OriginalTerm{纯餐}{cœna pura}。\par

6. \BibleQuote{一周的第一天，清晨，天还黑的时候，\Person{玛利亚玛达肋纳}来到坟墓那里，看见石头已从坟墓挪开了。}一周的第一天，因主的复活，基督信徒现今称之为主日。\BibleQuote{于是她跑去，来到\Person{西满伯多禄}和\Person{耶稣}所爱的那个门徒那里，对他们说：他们把主从坟墓里搬走了，我们不知道他们把祂放在哪里了。}有些希腊文抄本作：\InnerQuote{他们把我的主搬走了，}这可能是出于爱情与婢女之情的超常情感；但我们在所查阅的多种抄本中并未见到此种读法。\par

7. \BibleQuote{于是\Person{伯多禄}和那个门徒出来，往坟墓那里去。两人一起跑，那个门徒比\Person{伯多禄}跑得快，先到了坟墓。}此处重复值得注意，值得称道，因为他先省略了某事，后又回头补述，且仿佛按顺序接续。因为先已说了\BibleQuote{他们到了坟墓，}然后回头叙述他们怎样到达，并说\BibleQuote{两人一起跑}等等。此处他显示，那个门徒比同伴跑得快，先到了坟墓，这个门徒即指他自己，却以第三人称叙述。\par

8. \BibleQuote{他弯腰往里看，}他说，\BibleQuote{看见殓布放在那里，却没有进去。随后\Person{西满伯多禄}也来了，进了坟墓，看见殓布放在那里，也看见那裹在祂头上的头巾，没有和殓布放在一起，而是卷着放在另一个地方。}难道我们以为这些事没有意义吗？我绝不这么认为。但我们急于转向其他问题，这些问题因需要探究或其他某种晦暗而迫使我们停留。因为对于自明之事，探求每处细节的意义固然是神圣的喜悦，却只适合有闲暇的人，而我不在此列。\par

9. \BibleQuote{于是先到坟墓的那个门徒也进去了。}他先到，却最后进去。这无疑也必有其深意，但我无暇详加解释。\BibleQuote{他看见了，就相信了。}此处有人因未加注意，便以为\Person{若望}相信\Person{耶稣}已经复活了；但接下来的话并不支持这一理解。因为紧接着他便说：\BibleQuote{因为他们还不明白圣经，即祂必须从死者中复活。}他当时既不知道祂必须复活，就不可能相信祂已经复活了。那么他看见了什么？他相信了什么？无非是看见坟墓空了，并相信了妇人的话，即祂被人从坟墓中移走了。\BibleQuote{因为他们还不明白圣经，即祂必须从死者中复活。}同样，当他们从主亲耳听到时，虽然话语极其明白，但因他们惯常听祂用比喻说话，便不明白，以为祂另有所指。\par

\SourceRecord{Translated by John Gibb.；Nicene and Post-Nicene Fathers, First Series；Vol. 7.；Edited by Philip Schaff.；Buffalo, NY: Christian Literature Publishing Co.,；1888.；<http://www.newadvent.org/fathers/1701120.htm>.}

% structure: p0001 source=h1 target=section reason=page-title

\section{121讲（若望福音20:10-29）}

1. 玛利亚玛达肋纳已将主从坟墓中被带走的消息告诉了祂的门徒伯多禄和若望；他们来到那里，只找到了包裹身体的殓布；除了相信她所说的话以及她自己所相信的之外，他们还能相信什么呢？\BibleQuote{门徒又回到自己的地方去}（即他们居住之处，他们就是从那里跑到坟墓的）。\BibleQuote{玛利亚却站在坟墓外面哭泣。} 因为当男人们返回时，那较弱的性别却被更强烈的爱牢牢系于那地方。那双曾寻找主而未寻见的眼睛，现在除了哭泣别无他事，她们因主被从坟墓中带走而忧伤，比祂被钉在树上时更深；因为即使是面对这样一位师傅，当祂活生生的临在从他们眼前消失时，连对祂的记忆也消失了。因此，这样的悲伤使那妇人留在坟墓旁。\BibleQuote{她一面哭，一面俯身向坟墓里窥看。} 我不知道她为何这样做。因为她并非不知道她所寻找的那位已不在那里，因为正是她亲自告诉门徒说祂已被从那里搬走了；而且他们也来过坟墓，不仅用眼查看，还进去寻找了主的身体，却没有找到。那么，她一面哭一面又俯身向坟墓里窥看，这意味着什么？是因为她的悲伤如此之大，几乎使她不敢相信他们或她自己的眼睛？还是出于某种神圣的冲动，她的心灵引导她向里面观看？因为她确实看了，\BibleQuote{并看见两个穿着白衣的天使，坐在安放过耶稣身体的地方，一个在头，一个在脚。} 为什么一个坐在头，另一个坐在脚？是否因为那些在希腊语中称为\OriginalTerm{天使}{angels (Greek: \GreekText{ἄγγελοι})}，在拉丁语中称为\OriginalTerm{信使}{nuntii (Latin)}的，以此表明基督的福音要从头到脚、从起初直到终结被宣讲？\BibleQuote{他们对她说：\InnerQuote{女人，你为什么哭？}她对他们说：\InnerQuote{因为他们把我的主搬走了，我不知道他们把祂放在哪里。}} 天使禁止她哭泣：因为这样的位置除了预示某种未来的喜乐之外，还能宣布什么呢？因为他们问\BibleQuote{你为什么哭？}，仿佛在说\Quote{不要哭}。但她以为他们是出于无知而问，便说明了她哭泣的原因。她说：\BibleQuote{因为他们把我的主搬走了。}她称她主的无生命的身体为她的主，这是以部分代整体；正如我们所有人都承认耶稣基督——天主的独生子、我们的主，祂当然同时是圣言、灵魂和肉体——曾被钉十字架并被埋葬，而只有祂的肉体被放在坟墓里。她接着说：\BibleQuote{我不知道他们把祂放在哪里。}这是更大的悲伤，因为她不知道该去哪里缓解她的痛苦。但此时，那喜乐——某种程度上由禁止她哭泣的天使所宣告——即将取代哭泣的时刻已经到来。\par

2. 最后，\BibleQuote{她说了这话，就转过身来，看见耶稣站在那里，却不知道那就是耶稣。耶稣对她说：\InnerQuote{女人，你为什么哭？你找谁？}她以为祂是园丁，就对祂说：\InnerQuote{先生，若是你把祂搬走了，请告诉我你把祂放在哪里，我好去取回祂。}耶稣对她说：\InnerQuote{玛利亚！}她转过身来，对祂说：\InnerQuote{辣步尼！}——就是说\InnerQuote{师傅。}} 谁也不可指责这妇人，因为她称园丁为先生（\OriginalTerm{主}{domine (Latin)}），而称耶稣为师傅。因为在那里她是在请求，在这里她是在认主；在那里她是在向一个她请求帮助的人表示尊敬，在这里她是在回忆那教导她分辨人性和神性之事的老师。她称那一位为主（先生），而她并非其婢女，为的是藉着他找到那属于她的主。因此，在一种意义上，她用\Quote{主}这个词时说了\BibleQuote{他们把祂搬走了}；在另一种意义上，当她说\BibleQuote{先生，若是你把祂搬走了}时，也是如此。因为先知也称那些仅仅是人的为\Quote{主}，但不同于那一位，关于祂经上写着：\BibleQuote{主是祂的名}。但这妇人，她已转过身来看见耶稣，并且正与祂交谈，却以为祂是园丁，怎么又说她再次转身去对祂说\Quote{辣步尼}呢？这岂不是因为，当她那时在身体上转身时，她以为祂并非祂所是的；而现在，当她在心中转身时，她认出了祂的真正所是。\par

3. \BibleQuote{耶稣对她说：\InnerQuote{你不要摸我；因为我还没有升到父那里去；但你去到我的弟兄那里，对他们说：我升到我的父和你们的父那里去，升到我的天主和你们的天主那里去。}} 这几句话中的某些点，我们必须加以考察；虽要简短，却需比平时更加留意。因为耶稣正在教导这妇人信德，她已认出祂是自己的老师，并在回答中这样称呼祂；这位园丁正在她心中，如同在自己的园子里，播下芥菜种子。那么\BibleQuote{不要摸我}是什么意思？似乎是为了寻求这种禁止的理由，祂接着说：\BibleQuote{因为我还没有升到父那里去。}这是什么意思？如果说，当祂还站在地上时，不可触摸祂，那么当祂坐在天上时，人又怎能触摸祂呢？因为当然，在祂升天之前，祂曾让门徒们触摸自己，如福音作者路加所证实的，祂说：\BibleQuote{你们摸我看看！因为鬼神没有骨肉，如同你们看我所有的。}\BibleRef{路 24:39} 或者当祂对门徒多默说：\BibleQuote{伸过你的指头来，看我的手；伸出你的手来，探入我的肋旁。}\BibleRef{若 20:27} 有谁会荒谬到宣称，祂在升到父那里之前，确实愿意让门徒触摸，而对妇女却在升天之后才允许？但无人，即使有人愿意，也不许如此愚妄。因为我们也读到，在复活之后、升天之前，妇女也触摸了耶稣，其中就有玛利亚玛达肋纳本人；因为玛窦记载，耶稣迎接她们说：\BibleQuote{愿你们平安！她们便上前抱住祂的脚，朝拜了祂。}\BibleRef{玛 28:9} 这事若望略过不提，但玛窦宣告为真理。因此，这些话中必然隐藏着某种神圣的奥秘；无论我们能否发现，都不应怀疑它的存在。故此，或者\BibleQuote{不要摸我，因为我还没有升到父那里去}这句话有这样的含义：这妇人象征外邦人的教会，他们直到耶稣实际升到父那里才相信基督；或者，基督这样愿意自己被信仰，换言之，被属灵地触摸，即祂与父原是一体。因为凡是在认识基督上进步到承认祂与父平等的人，祂在某种意义上是升到父那里去了；否则，祂就不被正确地触摸，也就是说，不被正确地信仰。但玛利亚可能仍然相信祂与父不相等，这话的确禁止她：\BibleQuote{不要摸我}；即，不要照你现在的想法这样信我；不要让你的思想向外延伸到那为你而成为的（人），而超越不到那使你成为的（天主）。否则，她仍然以血肉的方式相信祂，因她所哀哭的是一位人。祂说：\BibleQuote{因为我还没有升到父那里去}：当你相信我与父同等、绝不不平等时，你就触摸我了。\BibleQuote{但你去到我的弟兄那里，对他们说：我升到我的父和你们的父那里去。}祂不说\BibleQuote{我们的父}：因此，在一种意义上，祂是我的，在另一种意义上，是你们的；按本性是我的，按恩宠是你们的。\BibleQuote{升到我的天主和你们的天主那里去。}祂在这里也不说\BibleQuote{我们的天主}：因此，在这里，在一种意义上祂是我的，在另一种意义上祂是你们的：我的天主，我作为人也在祂之下；你们的天主，我是你们与祂之间的中保。\par

4. \BibleQuote{玛利亚玛达肋纳就去告诉门徒们说：\InnerQuote{我看见了主，祂对我说了这些话。}}\BibleRef{若 20:18} \BibleQuote{就在那一周的第一天晚上，门徒所在的地方，因为怕犹太人，门都关着，耶稣来了，站在中间，对他们说：\InnerQuote{愿你们平安！}}\BibleRef{若 20:19} \BibleQuote{说了这话，就把手和肋旁指给他们看。}\BibleRef{若 20:20} \BibleQuote{门徒们一看见主，就欢喜起来。}\BibleRef{若 20:20} \BibleQuote{耶稣又对他们说：\InnerQuote{愿你们平安！}}\BibleRef{若 20:21} \BibleQuote{就如父派遣了我，} 祂接着说，\BibleQuote{我也同样派遣你们。}\BibleRef{若 20:21} \BibleQuote{说了这话，就向他们嘘了一口气，说：\InnerQuote{你们领受圣神吧！}}\BibleRef{若 20:22} \BibleQuote{你们赦免谁的罪，就给谁赦免；你们保留谁的罪，就给谁保留。}\BibleRef{若 20:23} 教会的爱，借着圣神倾注在我们心中，赦免所有与她共融者的罪，却保留那些不共融者的罪。因此，在说了\BibleQuote{你们领受圣神吧}之后，祂立刻接着说了关于赦罪和保留罪的话。\par

5. \BibleQuote{十二人中的一个，号称狄狄摩的多默，当耶稣来时，却没有和他们在一起。别的门徒向他说：\InnerQuote{我们看见了主。}但他对他们说：\InnerQuote{我除非看见他手上的钉孔，用我的指头探入钉孔，用我的手探入他的肋旁，我决不信。}八天以后，耶稣的门徒又在屋里，多默也和他们在一起。门户关着，耶稣来了，站在中间说：\InnerQuote{愿你们平安！}然后对多默说：\InnerQuote{把你的指头伸到这里来，看看我的手罢！并伸过你的手来，探入我的肋旁，不要作无信的人，但要作个有信德的人。}多默回答说：\InnerQuote{我主！我天主！}}他看见了并触摸了那人，也认出了他既看不见也触摸不到的天主；但藉着他所看见和触摸的，他现在消除了一切疑虑，并相信了其他。\BibleQuote{耶稣对他说：因为你看见了我，你才相信。}祂没有说：你触摸了我，而是说：\BibleQuote{你看见了我，}因为视觉是一种普遍的官能。视觉也习惯性地与其他四种感觉联系在一起：正如我们常说，听一听，看它听起来多好；闻一闻，看它闻起来多好；尝一尝，看它尝起来多好；摸一摸，看它有多热。到处都有\OriginalTerm{看}{see}这个词出现，尽管严格说来，视觉只属于眼睛。因此这里主自己也说：\BibleQuote{把你的指头伸到这里来，看看我的手罢！}祂的意思除了\OriginalTerm{触摸并看见}{Touch and see}之外还有什么？然而他的指头上并没有眼睛。所以无论是通过观看，还是通过触摸，祂说：\BibleQuote{因为你看见了我，}祂说：\BibleQuote{你才相信。}虽然可以断言，当主主动提供（身体）时，那位门徒不敢那样去触摸；因为并未记载多默触摸了祂。但无论是仅凭注视，还是也通过触摸，他看见并相信了，随后的话更宣扬并称赞了外邦人的信德：\BibleQuote{那些没有看见而相信的，才是有福的！}祂使用了过去时态，如同一位在祂预定旨意中知道未来之事的人，仿佛已经发生了。但当前的讲道必须避免冗长之嫌：主将赐给我们机会在另一时间讨论剩余的主题。\par

\SourceRecord{Translated by John Gibb.；Nicene and Post-Nicene Fathers, First Series；Vol. 7.；Edited by Philip Schaff.；Buffalo, NY: Christian Literature Publishing Co.,；1888.；<http://www.newadvent.org/fathers/1701121.htm>.}

% structure: p0001 source=h1 target=section reason=page-title

\section{论道 122（\BibleRef{若 20:30-21:11}）}

在讲述了与门徒多默触摸基督身上伤处的地方相关的事——他看见了他所不愿相信的，就相信了——之后，福音作者若望插入了这些话，说：\BibleQuote{耶稣在门徒面前还行了许多别的神迹，没有记在这书上；但记这些，是要你们信耶稣是基督，天主子；并使你们信的人，因祂的名获得生命。}这一段似乎标示着本书的结束；但后来却记述了主如何在提庇黎雅海显现自己，并在捕鱼奇事中特别提及教会的奥迹，关乎其未来在死者最后复活中的特质。因此，我认为，这样安排是合适的：在中间插入一个如同本书的结尾，又为随后要叙述的内容加上一种序言，以便在一定程度上赋予它更显著的位置。叙述本身是这样开始的：\BibleQuote{这些事后，耶稣在提庇黎雅海再次向门徒显现；祂显现的方式是这样的：西满伯多禄，号称狄狄摩的多默，加里肋亚加纳的纳塔乃耳，载伯德的儿子们，以及另外两位门徒，都在一起。西满伯多禄对他们说：\InnerQuote{我去打鱼。}他们对他说：\InnerQuote{我们也跟你一起去。}}\par

关于门徒们的这次捕鱼，通常有人会问：为什么伯多禄和载伯德的儿子们又回到了他们在被主召唤之前的状态；因为就在主对他们说\BibleQuote{来跟随我，我要使你们成为渔人的渔夫。}\BibleRef{玛 4:19}时，他们还是渔夫。那时他们如此真诚地跟随祂，以致舍弃了一切，只为依附祂作为他们的师傅；以至于当那个富有的人因祂对他说的\BibleQuote{你去变卖你所有的，施舍给穷人，你必有宝藏在天上，然后来跟随我，}而悲伤地离开时，伯多禄对祂说：\BibleQuote{看，我们舍弃了一切，跟随了祢。}那么，为什么现在，当他们好像放弃了宗徒职务，就重新变成他们原来的样子，再次寻求他们曾舍弃的东西，仿佛忘记了他们曾听到的那句话：\BibleQuote{手扶着犁向后看的人，不适合于天主的国}\BibleRef{路 9:62}？假如他们是在耶稣躺在坟墓里、尚未从死者中复活的时候这样做——当然，他们不可能这样做，因为祂被钉十字架的那一天，他们都全神贯注直到祂在傍晚前被埋葬；第二天是安息日，对于遵守祖传习俗的人来说，完全不得工作；而在第三天，主复活了，并重新召叫他们归向那他们尚未开始对祂怀有的希望——然而，假如他们那时就做了，我们可以认为那是由于占据了他们心灵的绝望所致。但现在，在祂从坟墓中活生生地回到他们身边之后，在祂复活的肉身最明显的事实呈现在他们的眼前和手中——不仅可见，还可触摸和抚摸——之后，在查看了伤处的痕迹，甚至到了宗徒多默的宣认（他曾事先声明否则绝不相信）之后，在祂向他们嘘气、领受圣神之后，并从祂口中流入他们耳中的话：\BibleQuote{就如父派遣了我，我也同样派遣你们。你们赦免谁的罪，就给谁赦免；你们保留谁的罪，就给谁保留。}之后，他们突然又变成了他们原来的样子，渔夫，不是渔人的渔夫，而是鱼的渔夫。\par

3. 因此，我们必须给那些因此事而不安的人回答：他们并非被禁止以手工技艺寻求必要的生计，只要那技艺本身合法，且只要他们保持宗徒的职务不变，若在任何时候他们别无谋生之道。除非有人敢于想象或断言，宗徒\Person{保禄}未能达到那舍弃一切跟随基督之人的成全，因为为了不成为他所传福音之人的负担，他亲手做工以维持生计：\BibleRef{得后 3:8} 其中我们反而看到他自己话语的应验：\Quote{我比他们众人更劳碌；}并且他加上：\Quote{不是我，而是同我一起的天主的恩宠。}\BibleRef{格前 15:10} 为要显明，这也应归功于天主的恩宠：他能在心灵和身体上比他们众人更丰盛地劳碌，以致他既未停止宣讲福音，也未像他们那样从福音中获取现世的供养，同时却将福音更广泛、更丰硕地播撒在无数先前从未听闻基督之名的民族中。由此他表明：靠福音生活，即从福音获得生计，并非加于宗徒身上的一项强制义务，而是赐予他们的一种权能。这位同一位宗徒在论及这权能时说：\Quote{如果我们给你们种下了属神的事物，从你们收获属血肉的事物，难道是大事吗？如果别人在你们身上分享这权能，我们岂不更该当呢？但是，}他接着说，\Quote{我们没有使用这权能。}不久之后他又说：\Quote{那些供职祭坛的，分享祭坛；主也是这样命令的：传福音的人应靠福音生活；但我没有用过这些权利。}所以很清楚，这并非命令宗徒必须只能靠福音生活，而是赋权给他们，使他们可以从那些因宣讲福音而播种属神事物的人那里收获属血肉的事物；即获得身体的供养，并且作为基督的士兵，从属于基督的省份的居民那里领取应得的薪饷。因此，那位杰出的士兵在不久之前论及此事时说：\Quote{有谁自备粮饷当兵呢？}然而他本人却这样做了，因为他比众人更劳碌。所以，如果蒙福的\Person{保禄}——为不向众人使用他确实与其他福音宣讲者同有的权能，而是自备粮饷当兵，以免那些对基督之名极为抗拒的外邦人，认为他的教导是某种索取钱财的东西而反感——以一种与他所受教育非常不同的方式，学习了一门全新的技艺，以便教师亲手养活自己，不让任何听者受累；那么，蒙福的\Person{伯多禄}，他先前曾是渔夫，若当时找不到其他谋生之道，岂不更应做他所熟悉的事吗？\par

4. 但有人会回答说：为什么他找不到生计，当时主曾应许说：\Quote{你们先寻求天主的国和它的义德，这一切自会加给你们}？\BibleRef{玛 6:33} 主正是以这种方式应验了他的应许。因为若不是他，谁把那些要捕的鱼放在那里呢？我们深信，他使他们陷入迫使他们去捕鱼的贫穷，唯一的原因是他要向他们展示他预备的奇迹，以便他既能喂养他福音的宣讲者，同时又借着那将要印在他们心中的鱼数所蕴含的伟大奥秘，来提升那福音本身？关于这事，我们现在也应当向你们讲论他自己摆在我们面前的事。\par

5. \BibleQuote{\Person{西满伯多禄}，因此，说：\InnerQuote{我去打鱼。}与他一起的人对他说：\InnerQuote{我们也要同你去。}他们就出去，上了船；但那整夜什么也没有捕到。到了早晨，\Person{耶稣}站在岸上；门徒们却没有认出他是\Person{耶稣}。\Person{耶稣}就对他们说：\InnerQuote{孩子们，你们有吃的吗？}他们回答说：\InnerQuote{没有。}他对他们说：\InnerQuote{向船右边撒网，你们就会找到。}他们便撒下网，竟拉不上来了，因为鱼很多。于是\Person{耶稣}所爱的那个门徒对\Person{伯多禄}说：\InnerQuote{是主。}\Person{西满伯多禄}一听说是主，便束上外衣（因为他赤着身子），跳进海里。其他门徒坐着小船（因为他们离岸不远，大约两百肘）来了，拖着一网鱼。他们一上岸，就看见那里有一堆炭火，上面放着鱼和饼。\Person{耶稣}对他们说：\InnerQuote{把你们刚才捕的鱼，拿一些来。}\Person{西满伯多禄}便上去，把网拉到岸上，里面装满了大鱼，共一百五十三条；虽然这么多，网却没有破。}\par

6. 这是伟大的\BibleBook{若望}{福音}中的一个伟大奥迹；为了使它更有力地引起我们的注意，最后一章成了它的纪录所在。因此，既然有七位门徒参与了那次捕鱼：伯多禄、多默、纳塔乃耳、载伯德的两个儿子，以及另外两位没有提及名字的，他们藉着七数指向时间的终结。因为所有时间都在七天中循环。与此相关的还有陈述：到了早晨的时候，耶稣站在岸上；因为岸也是海的边界，所以象征世界的终结。伯多禄将网拖到陆地（即岸上）的行动，也显示同样的世界终结。主自己曾解释过这一点，当祂在另一处用一个下到海里的网作比喻时：祂说：\Quote{他们将网拖到岸上。} 祂又解释道：\Quote{世界终结时也将如此。} \BibleRef{玛 13:48}\par

7. 然而，那是一个言语中的比喻，而不是体现在外在行动中的比喻；正如在我们面前的这段经文中，主藉着一个外在行动指出了教会将在世界终结时所拥有的特性，同样地，藉着另一次捕鱼，祂指出了教会现在的特性。祂在传道初期做了前者，在复活后做了后者，从而表明：在前者中，捕鱼象征现今存在于教会中的善人与恶人；但在后者中，只象征善人，即当死者的复活在世界终结时完成之后，教会将永远拥有的善人。此外，在先前那次，当耶稣命令捕鱼时，祂并没有像这里一样站在岸上，而是\Quote{进了西满的一只船，请他稍微划离岸边；祂便坐在船上，教训群众。讲完了，祂对西满说：划到深处去，撒网捕鱼。} 在那次，他们也把捕到的鱼装在船上，而不是像这里一样将网拖到岸上。藉着这些标记以及可能找到的其他标记，前一次预表了教会如同存在于现世中，后一次预表了教会如同在世界终结时那样；因此，前者发生在主复活之前，后者发生在主复活之后；因为在前者中，我们被基督象征为被召叫的人，在后者中，象征为从死者中复活的人。在那次，网没有撒在右边，以免只象征善人，也没有撒在左边，以免仅限于恶人；而是不偏左右，祂说：\Quote{撒网捕鱼}，好让我们明白善恶混杂；而在这次，祂说：\Quote{把网撒在船的右边}，以象征站在右边的人，即善人。在那里，网破裂了，因为要象征分裂；但在这里，因为在那至高圣人的和平中不再有分裂，所以福音作者有理由说：\Quote{虽然鱼这样多}（即如此之大），\Quote{网却没有破}；这似乎是与之前网破的时候相比，并称赞这里的善胜过了之前的恶。在那里，捕到的鱼数量极多，两只船都装满了，并开始下沉\BibleRef{路 5:3}，即被压到快要沉的程度；因为它们并没有真的沉下去，而是处于极度的危险中。因为除了无法抵挡那庞大的人群（他们几乎完全颠覆纪律，其德行与圣人的道路全然相悖）涌入教会之外，我们所叹息的巨大邪恶还能从哪里来呢？然而在这里，他们把网撒在右边，\Quote{竟不能把网拉上来，因为鱼太多。} \Quote{竟不能把网拉上来}这句话是什么意思呢？岂不就是说，那些属于生命复活、即属于右边的人，在基督徒名号的网中离开此生，只有在岸上——换句话说，当他们在世界终结时从死者中复活时——才会显现出来吗？因此，他们不能把网拉上来，以便将所捕的鱼倒入船中，像那些网破、船沉的情景中那样。但教会在这生命结束后，在平安的睡眠中拥有那些右边的人，他们仿佛隐藏在深处，直到网被拖到岸上，即约两百肘。而且，正如第一次捕鱼是用两只船，分别象征割礼与未割礼的人；同样在这里，藉着两百肘，我认为是象征着这两类选民——割礼与未割礼的人——如同两个分开的百；因为往右边去的人数被总结性地用百来代表。最后，在先前那次捕鱼中，鱼的数目没有说明，仿佛先知的话在那里应验了：\Quote{我发言说了；他们多得不可胜数}；而在这里，鱼是可以计算的，确定数目是一百五十三；关于这个数目的理由，我们现在必须依靠主的帮助来解释一下。\par

8. 因为如果我们指定一个数字来表示法律，那除了十还能是什么呢？因为我们确知，法律的十诫，即那十条众所周知的诫命，原是由天主的指头写在两块石版上的。\BibleRef{申 9:10}但是，当法律没有恩宠辅助时，它就使人成为犯法者，且仅仅是文字而已。为此，宗徒特别宣称：\Quote{文字叫人死，但圣神却叫人活。}\BibleRef{格后 3:6}因此，让圣神加在文字上，免得文字杀死那圣神所不使其活着的人，并且让我们成就法律的诫命，不是靠我们自己的力量，而是靠救主的恩宠。但当恩宠加在法律上，即圣神加在文字上时，在某种意义上，七这个数字就加在十上了。因为圣书交给我们的文件证明，七这个数字表示圣神。例如，圣洁或圣化原本属于圣神，因此，正如父是神，子是神，因为天主是神，所以父是圣的，子也是圣的，但两者的圣神却特别被称为圣神。那么，在法律中，第一次明确提到圣化是在哪里呢？不就是在第七天吗？因为天主没有圣化第一天，当他造光的时候；也没有圣化第二天，当他造穹苍的时候；也没有圣化第三天，当他分开海洋和陆地，土地生出草木的时候；也没有圣化第四天，当星辰被造的时候；也没有圣化第五天，当水中和空中的动物被造的时候；也没有圣化第六天，当地上的活物和人本身被造的时候；但他圣化了第七天，当他停止一切工作安息的时候。因此，圣神正好由\OriginalTerm{七}{septenary}这个数字来代表。同样，\Person{依撒意亚}先知说：\Quote{天主的圣神将安息在祂身上；}然后他借着说：\Quote{智慧和聪敏的神，超见和刚毅的神，明达和孝爱之神；他将充满敬畏上主的神。}\BibleRef{依 11:2}而\WorkTitle{默示录}又如何呢？那里不是称它们为\OriginalTerm{天主的七神}{seven Spirits of God}吗？\BibleRef{默 3:1}然而，只有同一圣神，随他的心意分给每人一份。\BibleRef{格前 12:11}但那同一圣神的\OriginalTerm{七重化工}{sevenfold operation}之所以这样称呼，是由于圣神自己，他的临在使作者写出它们被称为七神。因此，当我们将代表法律的十加上代表圣神的七时，就得到十七；而当这个数字用来相加它所包含的每个数字，从1加到它本身时，总和是一百五十三。因为如果你将2加于1，当然得到3；如果再将这些加上3和4，总和是10；然后如果你加上所有直到17的数字，总和就是上述数字；也就是说，如果你将1到4加起来得到10，再加上5，得到15；加上6，结果是21；再加上7，得到28；再加上8、9、10，得到55；再加上11、12、13，得到91；再加上14、15、16，得到136；然后再加上我们一直说的剩余数字，即17，就得到了鱼的总数。但这并不意味着只有一百五十三个圣人将来从死者中复活获得永生，而是成千上万的圣人，他们都分享了圣神的恩宠；借着这恩宠，人得以与天主的法律，如同与敌人，和好；因此，借着赋予生命的圣神，文字不再杀死人，而是文字所命令的借着圣神的帮助得以成就，如果有任何亏缺，也被赦免。因此，所有分享这种恩宠的人，都由这个数字象征，即被象征性地代表。此外，这个数字还有三个五十，再加三，这涉及天主圣三的奥迹；而五十这个数字又是通过7乘7再加1得到的，因为7乘7是49。加1是为了表明有一位用七来表示的，因为他的七重化工；并且我们知道，在主升天后的第五十天，圣神被派遣了，门徒们曾遵照应许等候他。\par

9. 那么，描述这些鱼的数量如此之多、尺寸如此之大，即一百五十三条且都是大鱼，绝非偶然。因为经上记载说：\Quote{祂便将网拖上岸，满网大鱼，共一百五十三条。}当主说：\Quote{我来不是为废除法律，而是为成全}时，因为祂即将赐下圣神，使法律得以成全，并借此增加，仿佛将七加于十；在插入几句别的话后，祂接着说：\Quote{所以，谁若废除这些诫命中最小的一条，也这样教训人，在天国里他将称为最小的；但谁若实行并教训人，在天国里他将称为大的。}因此，后者可能属于大鱼之列。而那最小的，即言行不一者，可能存在于像首次捕鱼所象征的那种教会中，那教会善恶皆有，因为它也被称为天国，如祂所说：\Quote{天国好像撒在海里的网，网罗各种鱼。}（\BibleRef{玛 13:47}）在此，祂意指善者与恶者并存，并宣称他们将在岸上被分开，即在世界终结时。最后，为了表明那些最小的、口传善道却行恶败坏的人是可弃绝的，他们甚至在永生中也不会被视为最小的，而将毫无地位；在说了\Quote{他在天国里将称为最小的}之后，祂立刻补充说：\Quote{我告诉你们：除非你们的义德超过经师和法利塞人的义德，你们绝不能进天国。}（\BibleRef{玛 5:17}）无疑，这些经师和法利塞人就是那些坐在梅瑟座位上的人，关于他们祂说：\Quote{凡他们对你们所说的，你们要遵行，但不要效法他们的行为；因为他们能说不能行。}（\BibleRef{玛 23:2}）他们在讲道中教导的，却被自己的品行破坏。因此，在天国现今的形态中为最小的人，将不能进入未来形态的天国；因为他教导别人遵守自己却惯于违背的，在那些言行一致者的团体中便无份，因此他不属于大鱼之列，因为\Quote{谁若实行并教训人，在天国里他将称为大的}。正因为他在此地为大，所以将在那里——那里最小的不存在。是的，他们在那里必然如此伟大，以至于那里最小的比此地最大的还大。（\BibleRef{玛 11:11}）然而，那些在此地为大的人，即在网罗善恶的天国中行善教导的人，将在永恒的天国状态中更为伟大——我指的是那些以此处的鱼为代表、属于右方和复活得生命的人。我们还要按天主所赐的能力，论述主与那七位门徒的用餐、餐后祂所说的话，以及福音本身的结尾；但这些题目无法纳入本次讲道之中。\par

\SourceRecord{Translated by John Gibb.；Nicene and Post-Nicene Fathers, First Series；Vol. 7.；Edited by Philip Schaff.；Buffalo, NY: Christian Literature Publishing Co.,；1888.；<http://www.newadvent.org/fathers/1701122.htm>.}

% structure: p0001 source=h1 target=section reason=page-title

\section{讲道录 123（\BibleRef{若 21:12-19}）}

主在复活后第三次显现给祂的门徒，真福宗徒若望的福音就此结束。我们之前已尽其所能讲解了其较早部分，直到记载主显现给门徒，他们捕获了一百五十三条鱼，尽管鱼如此之大，网却没有破。接下来我们要在主所赐的恩惠下，按照各点所需，进行思考和讨论。捕鱼结束后，\BibleQuote{耶稣对他们说：\InnerQuote{来吃吧。}坐席的人中，没有一个敢问祂：\InnerQuote{祢是谁？}因为知道是主。}如果他们已经知道，何必还要问？如果没有必要，为何说\BibleQuote{他们不敢}，好像有必要，却因某种恐惧而不敢？因此，其中的意思是：耶稣亲自显现给这些门徒的真理证据如此确凿，以致他们中无人敢否认，甚至无人敢怀疑；因为若是有人怀疑，他必定会问。因此，经上说\BibleQuote{没有人敢问祂：祢是谁？}，意思就是：没有人敢怀疑那就是祂本人。\par

\BibleQuote{耶稣就来，拿起饼，递给他们，鱼也同样。}你看，这里也告诉我们他们吃了什么；关于这餐饭，我们也要说些甜美有益的话，只要我们也得以由祂分享这食物。前面记载，这些门徒来到岸边时，\BibleQuote{看见炭火已经生好，上面放着鱼，还有饼。}这里我们不应理解为饼也放在了炭火上，而只是要补充\Quote{他们看见}。如果我们把这个动词补在适当的位置，整句话就可以这样读：他们看见炭火生好，上面放着鱼，又看见饼。或者这样说：他们看见炭火生好，上面放着鱼；他们也看见了饼。遵照主的命令，他们也把自己捕的鱼拿来；虽然史家并未明确记述他们这样做了，但对主的命令并未缄默。因为祂说：\BibleQuote{把你们刚才捕的鱼拿一些来。}既然我们如此确定祂下了命令，谁会认为他们没有服从呢？因此，主为这七位门徒准备了晚餐，就是用他们看见放在炭火上的鱼，再加上他们捕来的鱼，以及同样清楚记载他们看见的饼。烤鱼是受苦的基督；祂自己也是自天而降的饼。教会与祂结合，以分享永福。正因如此，经上说\BibleQuote{把你们刚才捕的鱼拿一些来}，好使我们所有怀有这希望的人知道，我们自己也借着这七位门徒（在此段落中我们的普世团体可以理解为由其象征）参与此伟大圣事，并分享同一福乐。这就是主与祂门徒的晚餐；若望虽有许多关于基督的事可以说，却在此以深刻的思考和重要的教训结束了他的福音。因为在此，教会（将来唯独由善人组成）由一百五十三条鱼的捕获所预示；而对于那些如此相信、盼望和爱慕的人，这晚餐表明了他们分享如此超绝的福乐。\par

\BibleQuote{这已经是}，他说，\BibleQuote{耶稣从死人中复活后第三次显现给门徒。}这我们不应理解为显现的次数本身，而应理解为日子（即：以祂复活的那一天为第一天，八天后的第二次（当门徒多默看见并相信时），以及今天的第三次（祂与鱼如此行动的日子），尽管我们不知道这发生在多少天之后）；因为在那第一天，祂显现了不止一次，正如所有福音书作者相互印证所表明的那样；但正如我们所说，应按照日子计算祂的显现，使这一次成为第三次；因为那第一次显现（包括在第一天内的一切显现）应算作第一次，无论祂在复活那天显现了多少次、向多少人显现；第二次是八天之后，这一次是第三次，此后直到第四十日升天时，祂随心所欲地显现了多次，但并未全部记载在圣经中。\par

4. \BibleQuote{他们吃完了早饭以后，\Person{耶稣}对\Person{西满}\Person{伯多禄}说：\InnerQuote{西满，\Person{若望}的儿子，你爱我比这些更深吗？}他回答说：\InnerQuote{是的，主，你知道我爱你。}\Person{耶稣}对他说：\InnerQuote{你喂养我的羔羊。}他第二次又问他说：\InnerQuote{西满，\Person{若望}的儿子，你爱我吗？}他回答说：\InnerQuote{是的，主，你知道我爱你。}\Person{耶稣}对他说：\InnerQuote{你喂养我的羔羊。}他第三次问他说：\InnerQuote{西满，\Person{若望}的儿子，你爱我吗？}\Person{伯多禄}因\Person{耶稣}第三次问他说\InnerQuote{你爱我吗}便忧愁起来，他向他说：\InnerQuote{主，你知道一切，你知道我爱你。}\Person{耶稣}对他说：\InnerQuote{你牧放我的羊群。}\InnerQuote{我实实在在告诉你：你年轻时，自己束上腰，随意往来；但到了老年，你要伸出手来，别人要给你束上腰，带你到你不愿意去的地方。}他说这话，是指明\Person{伯多禄}要以怎样的死来光荣天主。} 那否认又爱的人，结果便是如此：他曾因自负而心高气傲，因否认而俯伏在地，因痛哭而得以洁净，因宣认而获认可，因受苦而戴上冠冕；这便是他达致的结局：为那位的名字，以成全的爱而死；他曾因颠倒的卤莽，许诺要与那位同死。他在自己的软弱中过早许诺的事，要等到因主的复活而坚强时才去实行。因为必要的次序是：\Person{基督}先为\Person{伯多禄}的救恩而死，然后\Person{伯多禄}再为宣讲\Person{基督}而死。这种由人的鲁莽开始的胆大妄为，完全颠倒了真理所建立的秩序。\Person{伯多禄}想为\Person{基督}舍命，即被救赎者代替救赎者舍命，然而\Person{基督}原是为所有属于祂的人、也包括\Person{伯多禄}，舍掉自己的生命；你看，如今这事已经成就了。从今以后，我们可以怀着真实的心力——因为这是恩赐的——为主的名而赴死，而不是凭对自己错误的估计而僭取虚假的心力。如今我们无须再害怕离开现世的生命，因为主的复活已经预先展示了将来的生命。如今，\Person{伯多禄}，你有理由不再害怕死亡，因为你曾哀悼祂的死亡，并曾以血肉之爱试图阻止祂为我们而死的那位，现在活着。\BibleRef{玛 16:21}你曾胆敢抢先走在导师前面，却在迫害者面前战栗：如今你的赎价已经付清，你应当跟随那位赎主，跟随祂直至十字架上的死亡。你已经亲耳听见那位你已证实是可信者的话；祂曾预言你的背弃，如今又预言你的苦难。\par

5. 主首先问祂所知道的事，且不只一次，而是第二次和第三次，问圣\Person{伯多禄}是否爱祂；同样多次，祂得到相同的回答，即祂是被爱的；同样多次，祂给圣\Person{伯多禄}同样的命令：喂养祂的羊。如今，在三次否认之后，接上了三次宣认，好使他的舌头为爱所奉献的服务不弱于为恐惧所献的，并且即将来临的死亡不显得比现世的生命从口中引出更多。如果恐惧的标志是否认牧者，那么喂养主的羊群就当是爱的职责。那些喂养基督羊群的目的，是使其成为自己的羊群，而非基督的羊群的人，被证明是爱自己，而非爱基督，因为他们出于自夸、掌权或得利的欲望，而非出于顺服、服务和悦乐天主的爱。因此，基督的三次叮嘱成了警醒的哨兵来防范这样的人，宗徒曾抱怨他们寻求自己的事，而非耶稣基督的事。\BibleRef{斐 2:21}因为\BibleQuote{你爱我吗？喂养我的羊，}这句话，若不是等于说：你若爱我，就不要想喂养你自己，而要喂养我的羊，当作我的羊而不是你的羊；在它们身上寻求我的光荣而非你的光荣；我的主权而非你的主权；我的利益而非你的利益；否则你就会被发现与那些属于危险时期的人为伍，他们爱自己，并且所有其他与这万恶之源相连的恶。因为宗徒在说了\BibleQuote{因为人将成为爱自己者}之后，接着加上\BibleQuote{爱财者、自夸者、骄傲者、亵渎者、忤逆父母者、忘恩者、邪恶者、不虔敬者、无情者、诬告者、不能自制者、残暴者、无恩慈者、出卖者、卤莽者、盲目者；爱逸乐胜于爱天主者；有虔敬的外表，却否认其能力。}\BibleRef{弟后 3:1}所有这些邪恶都从他首先指出的源头——即\BibleQuote{爱自己者}——涌流而出。因此，极恰当地，圣\Person{伯多禄}被问\BibleQuote{你爱我吗？}，并回答\BibleQuote{我爱你}，于是命令他\BibleQuote{喂养我的羔羊}，并且这是第二次和第三次。我们在此也得到证明，爱和喜欢是同一回事；因为主在最后一个问题中说的不是\OriginalTerm{你爱我吗？}{Diligis me?}，而是\OriginalTerm{你爱我吗？}{Amas me?}。那么，让我们不爱我们自己，而爱祂；并在喂养祂的羊群时，寻求属于祂的事，而非属于我们自己的事。因为以某种无法解释的方式，我不知其所以然，凡是爱自己而不爱天主的人，并不爱自己；凡爱天主而不爱自己的人，他才是爱自己。因为那不能靠自己生活的人，必定因爱自己而死亡；因此，那爱自己以至于丧失生命的人，并不爱自己。但当那位保存生命者被爱时，一个人因不爱自己而更爱，正是为了这个原因他不爱自己，即他可以爱那使他生活的那一位。那么，那些喂养基督羊群的人，不要成为\BibleQuote{爱自己者}，以免他们喂养羊群如同喂养自己的羊，而非祂的羊，并且想从羊群中获利，如同\BibleQuote{爱财者}；或管辖他们，如同\BibleQuote{自夸者}；或因从他们得到的荣誉而自夸，如同\BibleQuote{骄傲者}；或甚至走到创立异端的地步，如同\BibleQuote{亵渎者}；而不给圣教父让位，如同\BibleQuote{忤逆父母者}；对那些想要纠正他们的人以恶报善，因为他们不愿让人毁灭，如同\BibleQuote{忘恩者}；杀害自己的灵魂和他人的灵魂，如同\BibleQuote{邪恶者}；侮辱教会的慈母心怀，如同\BibleQuote{不虔敬者}；对软弱者没有同情心，如同\BibleQuote{无情者}；企图诽谤圣人的品格，如同\BibleQuote{诬告者}；放纵最卑鄙的私欲，如同\BibleQuote{不能自制者}；以诉讼为常事，如同\BibleQuote{残暴者}；不知爱德的服务，如同\BibleQuote{无恩慈者}；向虔敬者的敌人泄露他们明知应当保密的事，如同\BibleQuote{出卖者}；以无耻的讨论扰乱人的廉耻，如同\BibleQuote{卤莽者}；不明白他们所说的话和所坚持的事，\BibleRef{弟前 1:7}如同\BibleQuote{盲目者}；并喜爱肉体的享乐过于属灵的喜乐，如同\BibleQuote{爱逸乐胜于爱天主者}。因为这些以及类似的恶习，无论是全部集中在一个人身上，还是某些恶习在某个人身上占主导，都从这一根源以某种形式萌发出来，就是当人是\BibleQuote{爱自己者}时。这是一种特别需要警惕的恶习，尤其是喂养基督羊群的人，免得他们寻求自己的事，而非耶稣基督的事，并且将基督的宝血所赎买的人转向满足自己的私欲。在喂养祂羊群的人身上，这份爱应当增长到如此大的属灵热忱，以至能克服对死亡的本能恐惧，这恐惧让我们即使渴望与基督同活时也不愿死去。因为宗徒\Person{保禄}也说他渴望被解脱，与基督同在，\BibleRef{斐 1:23}但他也叹息，深感重负，不愿被脱去，而愿穿上，好使有死的被生命吞没。\BibleRef{格后 5:4}因此，主对当时爱祂的人说：\BibleQuote{当你年老时，你要伸出手来，别人要给你束上带子，带你到你不愿意去的地方。}祂说这话，是指明他要以怎样的死来光荣天主。祂说：\BibleQuote{你要伸出手来}；换句话说，你将被钉在十字架上。但为了达到这一步，\BibleQuote{别人要给你束上带子，带你}，不是到你所愿意的地方，而是\BibleQuote{到你所不愿意的地方}。祂先告诉他将发生的事，然后这事如何成就。因为不是在钉十字架之后，而是在正要被钉十字架时，他被带到他不愿意去的地方；因为钉十字架之后，他走的路不是他不愿意的，而是他愿意的。虽然当他脱离肉体时，他渴望与基督同在，但是，若有可能，他仍渴望永生而无死亡的痛苦；他被不情愿地携入这痛苦的经历，但从其中（当一切结束后）他情愿地被带走了；他不情愿地来到它面前，但情愿地征服了它，并把这软弱的感受留在身后，这感受使每个人都不愿死——这是一种如此永恒自然的感受，甚至连年老本身也不能使圣\Person{伯多禄}摆脱它的影响，正如对他所说的：\BibleQuote{当你年老时}，你将被领\BibleQuote{到你所不愿意的地方}。为安慰我们，救主自己也在祂身上转化了这同样的感受，当祂说：\BibleQuote{父啊！如果可能，请免去我这杯吧！}\BibleRef{玛 26:39}而祂当然是来受死，不是出于任何必然性，而是出于自愿去死，有权舍弃祂的生命，也有权再取回来。但无论死亡的痛苦有多大，它都应当被对祂的爱所产生的力量所克服，祂是我们的生命，却甘愿为我们忍受死亡。因为如果死亡没有任何痛苦，即使是最小的痛苦，殉道者的光荣就不会如此巨大。但如果善牧为自己的羊舍命，并从这些羊中为自己兴起了如此众多的殉道者，那么，那些受托喂养、即教导和管理这些羊的人，岂不更应当为真理奋争至死，甚至流血抵抗罪恶吗？因此，结合祂自己受难的先前榜样，谁能看不出牧者们应当更加紧密地模仿牧者，如果连许多羊都这样效法了祂，而在这唯一的牧者和唯一的羊群中，牧者们自己同样也是羊？因为祂使所有祂为之而死的人都成为祂的羊，因为祂自己也成了羊，好为众人受苦。\par

\SourceRecord{Translated by John Gibb.；Nicene and Post-Nicene Fathers, First Series；Vol. 7.；Edited by Philip Schaff.；Buffalo, NY: Christian Literature Publishing Co.,；1888.；<http://www.newadvent.org/fathers/1701123.htm>.}

% structure: p0001 source=h1 target=section reason=page-title

\section{论道一二四（\BibleRef{若 21:19-25}）}

1. 主第三次显现给门徒时，对宗徒伯多禄说：\Quote{跟随我吧。}但对宗徒若望却说：\Quote{我这样愿意他留待我来，这关你什么事？}这是一个不小的问题。我们按照主的赐予，在这部著作的最后讲论中，致力于讨论或解决这个问题。主既已预先告诉伯多禄，他将以何种死来光荣天主，\Quote{就对他说：跟随我吧。伯多禄转过身来，看见主所爱的那个门徒跟着，就是那晚餐时靠在他胸前，说：主，出卖你的是谁？伯多禄看见他，就对耶稣说：主，这个人怎样？耶稣对他说：我这样愿意他留待我来，这关你什么事？你跟随我吧。于是这话在弟兄中间传开了，说那个门徒不死；其实耶稣并没有说他不死，而是说：我这样愿意他留待我来，这关你什么事？}可见在这部福音中，这个问题所涉及的广度，其深度必然使探究者倍加用心。因为为什么对伯多禄说\Quote{跟随我吧}，却不对同样在场的其他人说呢？当然，门徒们也跟随他们的师傅。但如果这仅指他的受苦而言，难道只有伯多禄为基督信仰的真理受苦吗？在那一同在场的七人中，不是还有另一个载伯德的儿子、若望的兄弟雅各伯吗？据记载，他在主升天后显然被黑落德所杀。\BibleRef{宗 12:2}但有人可能说，雅各伯不是被钉死的，因此对伯多禄说\Quote{跟随我吧}是恰当的，因为他不仅经历了死亡，而且像基督一样，经历了十字架上的死亡。假定如此，如果没有更令人满意的解释。那么，为什么论到若望时说\Quote{我这样愿意他留待我来，这关你什么事？}并重复说\Quote{你跟随我吧}，仿佛那另一个因此不会跟随，因为主要他留待主来？谁能轻易相信，这所指的是别的意思，而非当时的弟兄们所相信的，即那个门徒不会死，而要活在世上直到耶稣来？但若望自己消除了这种想法，他明确否认主曾那样说。因为他为什么加上\Quote{耶稣并没有说他不死}，岂不是为了防止虚假观念占据人心？\par

2. 但无论谁，若他乐意，仍可拒绝同意，并宣称若望所说的是够真实的，即主并没有说那个门徒不死，然而这正是此处所记主的话语的含义；并进一步宣称若望宗徒仍然活着，坚持说他是在厄弗所的坟墓中沉睡而非死卧其中。他可以引用流行的传闻作为论据，说那里的土地有明显震动，呈现一种起伏的外观，并断言这无论是坚定地还是固执地，是由于他的呼吸所致。因为我们不会缺少一些这样相信的人，如果也不缺少那些声称梅瑟仍然活着的人；因为经上记载他的坟墓找不到了，\BibleRef{申 34:6}并且他与主一起在山上同厄里亚一同显现，\BibleRef{玛 17:3}而厄里亚我们读到并没有死，而是被提升天了。\BibleRef{列下 2:11}仿佛梅瑟的身体不可能藏在某处，以致人完全找不到它的位置，并在厄里亚和他与基督一同显现的时候，由天主的能力从那里被举起，正如在基督受难时，许多圣人的身体复活了，并在祂复活后，按照圣经，在圣城中显现给许多人。\BibleRef{玛 27:52}但是，正如我开始说的，如果有些人否认梅瑟的死——经上在那段记载他坟墓无处可寻的文字中，明确说他死了——那么，从主说\Quote{我这样愿意他留待我来}这句话，更有理由相信若望在地底下沉睡，但仍然活着？关于他，我们也有传统（见于某些伪经中），说他身体健康时，下令为自己造一座坟墓；当坟墓被挖好并精心预备后，他像躺在床上一样躺下，立即就断了气；然而那些这样理解主话的人认为，他并非真的断气，而只是像睡着了一样；当被认为是死的时候，他实际上是睡着的被埋葬了，并且他将这样直到基督再来，同时通过灰尘的冒泡表明他活着的事实，这被认为是由沉睡者的呼吸从坟墓深处推到表面。我认为与这种意见争论是多余的。因为认识那地方的人可以自己看看那里的地面是否像所说的那样或经历那样，因为我们确实也从并非完全不可靠的见证人那里听说了这回事。\par

3. 同时，让我们姑且接受这一看法，因为我们无法以任何确凿的证据加以反驳，以免又引发另一个问题：为什么土地本身似乎在某种程度上对掩埋的尸首具有活力和呼吸？但如果借着伟大的奇迹，全能的天主能使活着的身体长眠于地下，直到世界终结，这样的问题还能解决如此重大的疑问吗？不，反倒产生了更广泛、更困难的问题：为什么耶稣将长久安眠于身体的恩赐赐予祂所爱的、超过其他门徒、以至于他配得靠在祂胸膛上的那位门徒，却以殉道的崇高荣耀使蒙福的\Person{伯多禄}脱离身体的重担，并赐予他宗徒\Person{保禄}所渴望并写下的话，即\Quote{解开束缚，与基督同在}？\BibleRef{斐 1:23}但如果——这更值得相信——圣\Person{若望}声称主并未说\Quote{他不死}，正是为了不至于有人将祂的话理解成这样；并且他的身体像其他死者一样毫无生气地躺在坟墓里；那么，如果传闻中关于土地的事确实发生——土地不断被取走却又重新长出——那么要么是为了彰显他死亡的珍贵，因为他的死亡缺乏殉道的荣耀（他并非因基督信仰死于迫害者之手），要么出于其他我们不得而知的原因。然而，问题依然存在：主为何对一个注定要死的人说：\Quote{我愿意他这样存留直到我来。}\par

4. 此外，有谁不会倾向于在这两位宗徒——\Person{伯多禄}和\Person{若望}——身上提出进一步的疑问：为什么主更爱\Person{若望}，而主自己却更被\Person{伯多禄}所爱？因为每当\Person{若望}谈到自己时，为了使人明白所指而不提自己的名字，他总是加上这句话：耶稣爱他，仿佛他是唯一被爱的人，以此将他与其他所有当然也被基督所爱的门徒区分开来；他这样说时，除了表明自己更被爱之外，还想让人明白什么呢？他绝不会说谎。耶稣还能给出什么比这更大的证据表明祂对他更深的爱呢：这位与其余门徒同享伟大救恩的人，竟是唯一靠在救主本人胸膛上的那位？此外，从许多文字证据都可以证明，宗徒\Person{伯多禄}爱基督胜于其他人；但不必远求，这一点在紧接着前面那段经文中就足够明显，就是主第三次显现时，祂问他说：\Quote{你爱我比这些更深吗？}主当然知道，却仍然提问，为使我们这些读福音的人也能从主的问话和\Person{伯多禄}的回答中，知道\Person{伯多禄}对基督的爱。但是当\Person{伯多禄}只回答说\Quote{我爱祢}，而没有加上\Quote{比这些人更深}时，他的回答包含了他自己所知道的一切。因为他无法知道别人被爱的程度，无法洞察他人的内心。但在他第一次回答中说\Quote{主啊，是的，祢知道}，这就足够清楚地表明，主所问的一切祂都完全知道。因此，主不仅知道\Person{伯多禄}爱祂，而且知道\Person{伯多禄}爱祂胜于其他人。然而，如果我们自问并探究：两个人中谁更好？是更爱基督的人还是爱得较少的人？谁会犹豫回答：更爱基督的人更好？另一方面，如果我们提出这个问题：两个人中谁更好？是更少被基督所爱的人，还是更多被基督所爱的人？毫无疑问，我们会回答：更被基督所爱的人更好。因此，在我首先提出的比较中，\Person{伯多禄}优于\Person{若望}；但在后一种比较中，\Person{若望}优于\Person{伯多禄}。于是，我们提出第三种形式的问题：哪一位门徒更好？是爱基督少于同伴、但比同伴更被基督所爱的人？还是比同伴更被基督所爱较少、却比同伴更爱基督的人？在这里，回答显然停滞不前，问题变得更加重大。然而，就我自己的智慧而言，我可以轻易回答：更爱基督的人更好，但更被基督所爱的人更幸福；只要我能完全明白如何为我们的救主的正义辩护：为什么祂更少爱那更爱祂的人，却更多爱那更少爱祂的人。\par

5. 因此，我将在祂那公义隐藏而仁慈彰显的眷顾下，按照祂可能恩赐的力量，着手讨论，以解决如此重要的问题；因为迄今为止，它只是被提出，尚未被阐释。那么，就让这成为它阐释的开端，即我们要记住：在这个败坏的、沉重地压迫灵魂的身体中（\BibleRef{智 9:15}），我们过着悲惨的生活。但我们这些如今被中保救赎、并已领受圣神作为凭据的人，虽有福乐的生命在望，却尚未实际拥有它。因为所见的希望不是希望；一个人所见的事物，为什么还希望呢？但我们若希望那未见的，就要耐心等待。（\BibleRef{罗 8:24}）人需要在所受的恶中，而非所享的善中，才有忍耐的必要。因此，现今的生命，正如经上所记：\BibleQuote{人的生命岂不是一个考验的时期吗？}（\BibleRef{约 7:1}）我们每天向主呼喊：\BibleQuote{救我们免于凶恶}（\BibleRef{玛 6:13}），即使在人的罪过被赦免之后，仍被迫忍受这生命，尽管是原罪导致他堕入如此悲惨的境地。因为惩罚比过犯更持久；免得如果惩罚与过犯同时结束，过犯就被视为微小。正因如此，要么是为了显示我们所欠的悲苦之债，要么是为了改正我们短暂的生命，要么是为了操练必要的忍耐，人即使不再因他的罪而受永罚的束缚，仍被暂时留在惩罚中。这就是我们在必死状态中度过的现今邪恶日子的真实可悲却无可指责的境况，即使在这期间我们以爱慕的眼光注视着善的日子。因为这来自天主的义怒，关于这义怒，圣经说：\BibleQuote{妇女所生的人，寿命不长，且充满怒气。}（\BibleRef{约 14:1}）因为天主的愤怒不像人的愤怒，不是激动之人的骚动，而是平静地施行公义的惩罚。在祂的这种愤怒中，天主并不抑制祂的慈悲，如经上所记；但祂不断赐予人类的其它对悲苦之人的安慰外，在时期圆满时，当祂知道必须如此行时，祂派遣了祂的独生子（\BibleRef{迦 4:4}），万物是藉着祂造的，使祂成为人而仍是天主，成为天主与人之间的中保，即降生成人的基督耶稣（\BibleRef{弟前 2:5}），好使那些信祂的人，藉着再生之洗除去一切罪恶的罪责——即是他们因出生而继承的原罪，为应对此罪特别设立了再生之洗，以及所有因恶行而招致的其它罪过——得以脱离永罚，并在寄居世界期间生活于信德、望德和爱德中，走向祂的可见临在，一方面经历劳苦和危险的诱惑，另一方面经历天主肉身和精神的安慰，始终持守基督为他们所成的道路。并且因为即使在祂内行走，他们也免不了因这生命的软弱而潜入的罪过，祂赐给他们施舍的救药，使他们藉此得到祈祷的帮助，因为祂教导他们说：\BibleQuote{宽免我们的罪债，如同我们也宽免得罪我们的人。}（\BibleRef{玛 6:12}）教会就这样在蒙福的希望中度过这患难的生命；而教会在其普遍性中所象征的，被具体化在宗徒\Person{伯多禄}身上，因着他宗徒职的首位。因为，就他自己的个人而言，他按本性是一个人，按恩宠是一个基督徒，按更丰盛的恩宠是第一位宗徒；但当主对他说：\BibleQuote{我要将天国的钥匙交给你；凡你在地上所捆绑的，在天上也要被捆绑；凡你在地上所释放的，在天上也要被释放。}时，他代表了普世教会，这教会在此世受各种诱惑的摇撼，如同暴雨、洪水和风暴袭来，却不倒塌，因为它建立在磐石（\OriginalTerm{磐石}{petra}）上，伯多禄从这磐石得名。因为磐石（\OriginalTerm{磐石}{petra}）不是从伯多禄而来的，而是伯多禄从磐石而来的；正如基督不是因基督徒而得名，而是基督徒因基督而得名。正是为此，主说：\BibleQuote{我要在这磐石上建立我的教会}，因为伯多禄曾说：\BibleQuote{你是默西亚，永生天主之子。}（\BibleRef{玛 16:16}）因此，祂说，在你所承认的这磐石上，我要建立我的教会。因为那磐石（\OriginalTerm{磐石}{Petra}）就是基督（\BibleRef{格前 10:4}）；伯多禄自己也建立在这根基上。因为除了已经奠定的根基，即耶稣基督，谁也不能奠立别的根基。（\BibleRef{格前 3:11}）因此，建立在基督内的教会从主通过伯多禄接受了天国的钥匙，即捆绑和释放罪过的权柄。因为教会在基督内本质上是什么，伯多禄在磐石上代表的就是什么；在这种代表中，基督应被理解为磐石，伯多禄应被理解为教会。因此，伯多禄所代表的这个教会，只要生活在邪恶中，藉着爱慕和跟随基督，就从邪恶中得到解救。但其跟随对于那些为真理奋斗至死的人更为密切。但这是对教会的普遍性说的：\BibleQuote{跟随我}，正如基督为了这同一普遍性而受苦：关于这，伯多禄自己说：\BibleQuote{基督为我们受苦，给我们留下了榜样，使我们追随祂的足迹。}（\BibleRef{伯前 2:21}）因此，你看，这就是为什么对他说：\BibleQuote{跟随我}的原因。但还有另一种生命，即不朽的生命，它不在邪恶之中：在那里我们将面对面地看见那在此处透过镜子观看、如同谜语所见（\BibleRef{格前 13:12}）的，即使在真理的默观中已有很大进步。因此，有两种生命状态，从天上被宣讲并推荐给教会，为教会所认识：一种在信德中，另一种在直观中；一种在暂时的异乡客居，另一种在永恒的天乡居所；一种在劳苦中，另一种在安息中；一种在道路上，另一种在父家中；一种在行动中，另一种在默观的报酬中；一种远离邪恶并趋向良善，另一种无恶可离，有至善可享；一种与仇敌作战，另一种无仇敌作王；一种在逆境中勇敢，另一种无逆境经历；一种克制肉性情欲，另一种尽享精神喜乐；一种为制胜而忧虑，另一种在胜利的平安中安稳；一种在诱惑中得到帮助，另一种无诱惑，因帮助者本身而喜乐；一种致力于救济贫困，另一种在无贫困之处；一种宽恕他人的罪过，使自己得到赦免，另一种无所可宽恕，也无须求赦免；一种被邪恶鞭打，以免因善自夸，另一种因恩宠的满盈而脱离一切邪恶，毫无骄傲诱惑，能与至善紧密相连；一种辨别善恶，另一种只看见善。因此，一种是善的，却仍悲苦；另一种更善且有福。前者由宗徒\Person{伯多禄}所象征，后者由\Person{若望}所象征。前者全部在此世度过直到世界终结，并在那里找到其终点；后者延迟到世界终结后才完成，但在来世没有终结。因此，对后者说：\BibleQuote{跟随我}；但对前者却说：\BibleQuote{我要他存留直到我来，与你何干？你跟随我。}这最后一句是什么意思？按我的智慧所及，按我所理解的，难道不就是：你跟随我，藉着忍受暂时的邪恶而效法我；让他存留直到我来恢复永恒的善？这可以更清楚地表述为：让由我苦难榜样所塑造的完美的行动跟随我；但让仅仅是开始的默观存留直到我来，在我来时得以圆满完成。因为那直到死亡的虔诚的忍耐完满跟随基督；但知识的丰满则等待直到基督来临时才显明。因为在此世，邪恶在死者之地被忍受，而在那里，将在生者之地看见主的善事。因为当祂说：\BibleQuote{我要他存留直到我来}时，我们不应当理解祂的意思是永久停留或常住，而是等待；因为由他象征的之事肯定不会在现在实现，而是在基督来时实现。但由那被吩咐\BibleQuote{你跟随我}的人所象征的，除非现在去做，否则永不能达到所期望的终点。在这行动的生命中，我们越爱基督，就越容易从邪恶中得到解救。但祂爱我们如我们现在的样子较少，因此解救我们，使我们不至于总是如此。然而在那里，祂更爱我们；因为我们身上将没有任何使祂不悦的事，也没有任何祂必须将我们与之分离的事；祂在此世爱我们，也无非是为了医治我们，并将我们从祂所不爱的一切中迁移出去。因此，在此世，祂爱我们较少，因为祂不愿我们停留；在那里，祂爱我们更多，因为祂愿我们进入，并从祂不愿我们灭亡的境地出来。因此，让\Person{伯多禄}爱祂，以便我们获得从现世必死性中的解救；让\Person{若望}被祂所爱，以便我们得以保存到来世的不朽。\par

6. 但通过这一论证思路，我们已表明\Person{基督}为何爱\Person{若望}胜过\Person{伯多禄}，而非\Person{伯多禄}为何爱\Person{基督}胜过\Person{若望}。因为如果\Person{基督}在来世——我们将与他永久同住——比在此世——我们正从中被解救，以便永远在另一世界中——更爱我们，那并不意味着当我们变得更好时会少爱他；因为我们不可能变得更好，除非更爱他。那么，\Person{若望}为何比\Person{伯多禄}爱\Person{基督}少呢？如果他所象征的生命中，\Person{基督}必须被更丰盛地爱，难道不正是因为如此才说：\Quote{我愿意他存留}，也就是等待，\Quote{直到我来}？因为我们尚未拥有那爱本身——它将来要远大于现在——而是期待那未来，以便在他来临时获得它。正如同一宗徒在其书信中所宣告的：\BibleQuote{我们将来如何，还未显明；但我们知道，他若显现，我们必要相似他，因为我们要看见他实在怎样。}\BibleRef{若一 3:2} 那么，我们将会更爱我们所看见的。但主自己，凭其预定之知，更爱我们那未来尚未来到的生命，如他所知它将来在我们内会如何，以便藉此爱我们，吸引我们走向对其的拥有。因此，正如主的一切道路是仁慈和真理，我们知道我们现在的苦难，因为我们感受到它；所以我们更爱主的仁慈——我们希望在从苦难中被解救时彰显出来——我们每日祈求并体验它，尤其是在罪的赦免中：这就是\Person{伯多禄}所象征的，他爱得更多，却被爱得更少；因为\Person{基督}在我们苦难中爱我们少于在我们真福中。但对真理的默观——如它将来要成为的那样——我们爱得较少，因为至今我们既不知也不拥有它：这就是\Person{若望}所象征的，他爱得较少，因此既等待那状态本身，也等待在我们内那份对所应得的他的爱得以成全，直到主来临；但他被爱得更多，因为正是他所象征的那一位使他蒙福。\par

7. 然而，不可将这些杰出的宗徒分开。在\Person{伯多禄}所象征的事上，他们两者相似；在\Person{若望}所象征的事上，他们将来也要相似。在其代表身份中，一个跟随，另一个存留；但在他们个人的信德中，两者都忍受此世的苦难之恶，两者都期待来世真福的美善。情况不仅限于他们，而是涉及整个圣而公教会——\Person{基督}的新娘，她仍需从现今的考验中被解救，并在未来的幸福中得以保存。这两个生命状态由\Person{伯多禄}和\Person{若望}所象征，一个由一人，另一个由另一人；但在此世，他们两者都一度凭信德行走，而另一状态他们将永远凭目睹享受。因此，为整个诸圣团体——不可分割地属于\Person{基督}的身体——以及为他们在现今风暴般的生命中获得安全航行，\Person{伯多禄}——宗徒之长——领受天国的钥匙，以束缚和释放罪过；而为同一诸圣团体，面对那神秘来世怀抱中的完美安息，福音作者\Person{若望}倚靠在\Person{基督}的胸怀。因为不仅仅前者，而是整个教会，束缚和释放罪过；也不仅仅是后者，独自从主胸怀的泉源畅饮，在宣讲中再次发出\BibleQuote{在起初已有圣言，与天主同在的天主}，以及其他关于\Person{基督}天主性和整个天主圣三合一性的崇高真理——这些在将来王国中要面对面地看见，但与此同时，直到主来临，只能在镜子和谜语中看见；但主自己将这部福音传遍全世界，以便每个属于他的人都能按各自的能力畅饮。有些人持有一种看法——这些人在圣善雄辩方面并非无足轻重——认为宗徒\Person{若望}更被\Person{基督}所爱，因为他从未娶妻，且从少年时期起就生活在完美的贞洁中。诚然，正典圣经中并无确凿证据；然而，这看法对上述观点的适切性不无贡献，即他所象征的生命是那没有婚姻的生命。\par

8. \BibleQuote{这就是为这些事作证，并写下这些事的门徒；我们知道他的见证是真的。还有，} 他接着说，\BibleQuote{许多别的事，耶稣也行了，如果每一件都写下来，我想，即使世界本身也容不下所要写的书。} 我们不应认为，从地方空间来看，世界无法容纳它们；因为如果世界在它们写成之后不能承载，它们又怎能被写在里面呢？但也许是因为读者的理解力无法领会它们：尽管我们对某些事物本身的信仰仍然完好无损，但我们对它们所使用的言辞，可能常常显得超乎相信。这种情况不会发生在这样的情形：当任何晦暗或可疑的事物，通过陈述其根据和理由而得到解释时；而只发生在以下情形：当本身清楚的事物被夸大或缩小，却没有实质上偏离所要暗示的真理路径；因为言辞可能超越所指的事物本身，但仅限于这样的方式，即说话者的意愿（没有欺骗的意图）变得明显，因此，他知道自己会被相信到什么程度，便通过口述，将主题缩小或夸大，超出了人们会给予相信的限度。这种说话方式，在希腊文学和拉丁文学的大家们口中，被称为希腊语名称\OriginalTerm{夸张法}{hyperbole}。这种方式不仅在此处出现，在神圣著作的其他地方也多次出现：例如，\BibleQuote{他们设口顶天}；以及\BibleQuote{那些在过犯中前行之人的发顶}；还有许多其他同类的例子，它们在圣经中并不比其他修辞格或说话方式更少见。关于这些，我本可以进行更详尽的讨论，但既然福音作者在此结束了他的福音，我也被迫结束我的论述。\par

\SourceRecord{Translated by John Gibb.；Nicene and Post-Nicene Fathers, First Series；Vol. 7.；Edited by Philip Schaff.；Buffalo, NY: Christian Literature Publishing Co.,；1888.；<http://www.newadvent.org/fathers/1701124.htm>.}
